\documentclass[11pt]{article}
\usepackage[margin=1in]{geometry}
\usepackage{amsmath,amssymb,amsthm,mathtools}
\usepackage{enumitem}
\usepackage[hidelinks]{hyperref}

\newtheorem{theorem}{Theorem}[section]
\newtheorem{lemma}[theorem]{Lemma}
\newtheorem{proposition}[theorem]{Proposition}
\newtheorem{corollary}[theorem]{Corollary}
\newtheorem{definition}[theorem]{Definition}
\newtheorem{assessment}[theorem]{Assessment}
\newtheorem{firewall}[theorem]{Firewall}
\newtheorem{programme}[theorem]{Programme}
\newtheorem{audit}[theorem]{Audit}

\newcommand{\Tr}{\operatorname{Tr}}
\newcommand{\spec}{\operatorname{spec}}
\newcommand{\sgn}{\operatorname{sgn}}

\title{Renewal Standard Model--Einstein Continuum Research Notes\\
Seventh Series, V1}
\author{}
\date{\today}

\begin{document}
\maketitle
\tableofcontents

\section{Context and purpose of this restart}
\label{sec:s7v1-context}

These notes continue a constructive mathematical programme whose intended
endpoint is a cutoff-independent Standard Model--Einstein continuum with a
well-defined chiral fermion measure, controlled zero modes, gauge and
diffeomorphism covariance, clustering, reflection positivity, and
Osterwalder--Schrader reconstruction.  The endpoint has not been reached.

The preceding active file reached V100 and was frozen without alteration as
\begin{center}
\texttt{Renewal\_SM\_Einstein\_Continuum\_Research\_Notes\_V100\_f6.tex}.
\end{center}
Its SHA--256 digest is
\begin{center}
\texttt{B5A5540CBDA10D6371506F77948DD85B99DFF3C09A96A3FB30A2E65F77D564DE}.
\end{center}
It has \(35\,459\) logical source lines and \(1\,362\,633\) bytes.

This compact V1 starts a new active series because the append-only source
had again become difficult to navigate.  It first records the major proved
interfaces and the open gates, then continues the attack selected by the
V100 companion audit.  Subsequent active versions will append to this file
and replace their predecessors.  The newest Yang--Mills and Navier--Stokes
notes will be audited at V5 and every fifth active version.  At V100 this
series will again be frozen with the next \(\texttt{\_fN}\) suffix and
restarted from a summarized V1.

A theorem described as a card below is a proved implication from named
hypotheses.  It is not a claim that those hypotheses already hold for the
complete continuum measure.  Exact finite-cutoff identities, conditional
analytic estimates, and continuum existence statements remain logically
distinct.

\section{Major mathematical result map}
\label{sec:s7v1-result-map}

\subsection{Fermion topology, anomalies, and zero modes}

The archived programme treats a chiral determinant as a section of its
determinant line.  It computes the one-generation Standard Model local
anomaly arithmetic in the all-left-handed convention and keeps the residual
global phase as a mapping-torus condition.  Exact Schur--Berezin and
Riesz-projector constructions retain every zero or crossing mode as a
finite Grassmann sector before the complementary block is inverted.
Changing a spectral chart moves a crossing mode by an exact finite-rank
transfer identity rather than dividing by its vanishing determinant.

At finite cutoff, zeta/eta bookkeeping identifies the phase of a Hermitian
Dirac determinant, spectral flow fixes its sign changes, and the
Dai--Freed-type endpoint is isolated as a determinant-line holonomy gate.
Local anomaly cancellation does not by itself prove that global gate.

\subsection{Overlap operator and admissible gauge fields}

For the four-dimensional Wilson kernel at \(\rho=1\), the exact
forward--backward commutator expansion gives, whenever every represented
plaquette satisfies
\(\|I-U_p\|\le\delta<1/30\),
\begin{equation}
 H_W^2\succeq{1-30\delta\over a^2}I,
 \qquad
 \|H_W\|\le{7\over a}.                                    \label{eq:s7v1-Wilson-window}
\end{equation}
Consequently the overlap sign, Ginsparg--Wilson operator, moving chiral
projectors, and finite index are analytic, gauge covariant, and
exponentially local with constants independent of lattice volume.
Admissible index sectors are separated.

A smooth central autocorrelation on the compact gauge group produces a
nonnegative plaquette weight supported strictly inside the common Standard
Model admissibility ball.  Its character coefficients are nonnegative.
The resulting pure-gauge finite-cutoff measure has exact admissible support
and the standard link-reflection-positive factorization.  Universality of
this hard-support regulator and reflection positivity after chiral matter
and gravity are still open.

\subsection{Higgs reduction and massive overlap blocks}

Any protected set containing the free periodic gauge field and zero Higgs
field has chiral least singular value zero.  Antiperiodic spin boundary
conditions remove the exact finite-volume zero but do not produce a
physical-volume-uniform floor.

For
\[
 D_m=\left(1-{am\over2}\right)D_{\rm ov}+mI,
 \qquad 0<m\le2/a,
\]
the overlap spectral circle proves
\begin{equation}
 s_{\min}(D_m)\ge m,
 \qquad
 \|D_m^{-1}\|\le m^{-1},                                  \label{eq:s7v1-mass-floor}
\end{equation}
uniformly in the admissible gauge field and lattice volume.  Its vectorlike
determinant is positive.

On a Higgs chart separated from zero, the normalized one-doublet field
reduces the electroweak group to \(U(1)_{\rm em}\).  The charge ledger proves
left--right representation matching for the charged quarks and electron,
while the minimal neutrino row remains unmatched.  An equivariant tubular
projection gives a lift-independent local reduced electromagnetic
connection and an explicit broken-link remainder.

For a Ginsparg--Wilson operator \(D\), the improved field
\(A_D=I-aD/2\) obeys
\[
 D\widehat P_\pm=P_\mp D,
 \qquad
 A_D\widehat P_\pm=P_\pm A_D.
\]
The four chiral kinetic and Yukawa corners therefore assemble exactly into
\(D+mA_D=D_m\).  Singular-value decomposition of a constant Yukawa matrix
gives one massive-overlap block per nonzero Yukawa singular value.  Radial
Higgs variation is bounded by the largest Yukawa singular value times the
supremum radial fluctuation.  The nonzero tubular chart, its probability,
the unmatched neutrino sector, and the full chiral phase remain separate
gates.

\subsection{Metric, constraints, and long-range sectors}

The earlier volumes insert an elliptic order-zero metric normalizer before
forming a relative Dirac determinant; the residual perturbation has the
critical weak-Schatten order in four dimensions.  Heat and symbol
expansions assign cosmological, Einstein--Hilbert, gauge, Higgs,
nonminimal, curvature-squared, and boundary divergences to local owners,
while finite nonlocal Born terms are retained.

A nonlinear York--Lichnerowicz chart solves the local mean-zero constraint
near flat constant-mean-curvature data, leaving the global scalar
constraint.  Elliptic constraint substitution has Coulomb tails and belongs
to a polynomial rather than exponentially local grade.  Monopole and dipole
cancellations produce the first absolutely summable three-dimensional
multipole row.  The photon requires the analogous Ward endpoint moments.
The real Euclidean conformal direction, nonlinear global Einstein measure,
and Lorentzian reconstruction remain open.

Explicit high-frequency and high-amplitude tests show that a global
linearized Gaussian metric completion over the unbounded Higgs field is not
normalizable.  A nonlinear gravitational constraint construction or a
declared effective-field-theory horizon is required.

\subsection{Transport, aggregation, and ultraviolet trajectory}

Finite-depth conditional-density identities, positive-tilt cumulant
identities, Gaussian BKAR formulas, rooted tree estimates, and
Koteck\'y--Preiss cards reduce interacting transport to explicit local
activities and no-escape bounds.  These are conditional implications; the
full Standard Model--Einstein action has not been shown to satisfy all
uniform constants.

The programme records three structural obstructions to the present route:
the perturbative Gaussian-boundary hypercharge trajectory reaches its
ultraviolet boundary, a finite gravitational curvature list is not closed
under perturbative renormalization, and the Euclidean conformal kinetic
direction is unbounded below.  These facts require changed ultraviolet or
gravitational input; they are not a theorem excluding every ultraviolet
completion with the same infrared physics.

\subsection{The V100 assembled-transport result}

For an invertible chiral block transported to fixed source and row frames,
V100 proves
\begin{equation}
 {\det\widetilde B_1\over\det\widetilde B_0}
 =
 \exp\left\{\int_0^1
 \Tr(\widetilde B_t^{-1}\dot{\widetilde B}_t)\,dt\right\}. \label{eq:s7v1-Jacobi-summary}
\end{equation}
A singular floor controls the integrand only through the trace norm:
\[
 |\Tr(\widetilde B^{-1}\dot{\widetilde B})|
 \le s_{\min}(\widetilde B)^{-1}\|\dot{\widetilde B}\|_1.
\]
Replacing the trace norm by operator norm introduces the full matrix
dimension, sharply.  Thus the overlap floor is necessary but does not by
itself control the continuum determinant.  One must subtract local
extensive pieces before taking absolute bounds, or prove an exact holonomy
formula.

\section{Current open-gate ledger}
\label{sec:s7v1-open}

\begin{theorem}[Claim boundary at the seventh-series restart]
\label{thm:s7v1-open}
The archived results do not construct the full Standard Model--Einstein
continuum.  At least the following gates remain:
\begin{enumerate}[label=\textup{(G\arabic*)},leftmargin=4.5em]
\item a global regulator-level chiral determinant section with trivial
      combined mapping-torus holonomy;
\item cutoff- and volume-uniform control of the assembled relative
      determinant connection after local subtraction;
\item a probability and uniform-integrability theorem for the admissible
      nonzero-Higgs tubular chart;
\item a complete treatment of unmatched massless fermion rows and spectral
      crossings;
\item interacting photon Ward moments and gravitational multipole/contact
      identities sufficient for every long-range edge;
\item a nonperturbative ultraviolet replacement for the hypercharge
      Gaussian-boundary trajectory;
\item a nonlinear, positive, reflection-compatible gravitational
      construction which handles the conformal direction and the generated
      curvature tower;
\item clustering, reflection positivity of the complete mixed theory,
      cutoff removal, and Osterwalder--Schrader reconstruction.
\end{enumerate}
\end{theorem}

\begin{proof}
Each item is an explicit firewall in the frozen archives.  None follows from
the finite-cutoff overlap gap, massive singular floor, pure-gauge
positive-type regulator, or conditional cluster cards.  Therefore a
continuum conclusion cannot yet be drawn.
\end{proof}

\section{New attack: subtract the massive base before taking traces}
\label{sec:s7v1-relative}

Let \(E\) and \(F\) be equal-dimensional Hermitian spaces.  Let
\(B_0:E\to F\) be the fibrewise massive-overlap reference assembled in the
Higgs-reduced frames, and let \(B:E\to F\) be the complete massive chiral
block on the same chart.  Define
\begin{equation}
 \mathcal K:=B_0^{-1}(B-B_0),
 \qquad
 B=B_0(I+\mathcal K).                                     \label{eq:s7v1-K}
\end{equation}

\begin{lemma}[The relative determinant is a frame-independent scalar]
\label{lem:s7v1-relative-scalar}
Although \(\det B\) and \(\det B_0\) are individually elements of
\(\det F\otimes(\det E)^*\), their quotient
\begin{equation}
 \mathcal R(B,B_0):={\det B\over\det B_0}
 =\det(I+\mathcal K)                                      \label{eq:s7v1-relative-det}
\end{equation}
is a canonical scalar.  Under arbitrary unitary changes of source and row
frames, \(\mathcal K\) changes by unitary conjugation and
\(\mathcal R(B,B_0)\) is unchanged.
\end{lemma}

\begin{proof}
In frames, let
\(B'=U_F^*BU_E\) and \(B_0'=U_F^*B_0U_E\).  The determinant factors
\(\overline{\det U_F}\det U_E\) cancel in the quotient.  Moreover,
\[
 (B_0')^{-1}(B'-B_0')
 =U_E^*B_0^{-1}(B-B_0)U_E
 =U_E^*\mathcal KU_E.
\]
The determinant of \(I+\mathcal K\) is invariant under this conjugation.
\end{proof}

\begin{theorem}[Exact relative determinant series]
\label{thm:s7v1-relative-series}
Suppose
\begin{equation}
 \|\mathcal K\|\le\theta<1.                               \label{eq:s7v1-theta}
\end{equation}
Then \(B\) is invertible, the path
\(B_z=B_0(I+z\mathcal K)\), \(0\le z\le1\), has a canonical logarithm
starting at zero, and
\begin{equation}
 \log\mathcal R(B,B_0)
 =\sum_{n=1}^{\infty}{(-1)^{n+1}\over n}\Tr(\mathcal K^n).
                                                               \label{eq:s7v1-log-series}
\end{equation}
The series is absolutely convergent and obeys
\begin{equation}
 |\log\mathcal R(B,B_0)|
 \le{\|\mathcal K\|_1\over1-\theta}.                       \label{eq:s7v1-log-bound}
\end{equation}
The same bound controls the logarithmic modulus and the continuous phase
lift separately.
\end{theorem}

\begin{proof}
For \(|z|\le1\), the Neumann series makes \(I+z\mathcal K\) invertible.
Finite-dimensional holomorphic functional calculus gives
\[
 \log(I+\mathcal K)
 =\sum_{n\ge1}{(-1)^{n+1}\over n}\mathcal K^n.
\]
Taking the trace and using
\(\det e^X=e^{\Tr X}\) proves
\eqref{eq:s7v1-log-series} with the logarithm selected by the path.
Furthermore,
\[
 |\Tr(\mathcal K^n)|
 \le\|\mathcal K^n\|_1
 \le\|\mathcal K\|^{n-1}\|\mathcal K\|_1.
\]
Summing and using \(1/n\le1\) gives
\[
 \sum_{n\ge1}{\theta^{n-1}\over n}\|\mathcal K\|_1
 \le{\|\mathcal K\|_1\over1-\theta}.
\]
Real and imaginary parts are bounded by the modulus of the logarithm.
\end{proof}

\begin{corollary}[HYFC supplies the convergence radius]
\label{cor:s7v1-HYFC-radius}
If the fibrewise base obeys
\(\|B_0^{-1}\|\le m_*^{-1}\) and the complete assembled remainder obeys
\[
 \|B-B_0\|\le\theta m_*,
 \qquad\theta<1,
\]
then \eqref{eq:s7v1-theta} holds.  Hence the relative determinant has the
canonical series \eqref{eq:s7v1-log-series}.
\end{corollary}

\begin{proof}
Equation \eqref{eq:s7v1-K} gives
\[
 \|\mathcal K\|
 \le\|B_0^{-1}\|\|B-B_0\|
 \le\theta.
\]
Apply \Cref{thm:s7v1-relative-series}.
\end{proof}

\section{Relative transport and the trace-class target}
\label{sec:s7v1-relative-transport}

Let \(B(t)\) and \(B_0(t)\) be \(C^1\) families on the same moving source and
row bundles, transported by the same unitary frames.  Define
\[
 I+\mathcal K(t):=B_0(t)^{-1}B(t).
\]

\begin{theorem}[Subtracted Jacobi identity]
\label{thm:s7v1-subtracted-Jacobi}
Whenever both blocks are invertible,
\begin{align}
 {d\over dt}\log{\det B(t)\over\det B_0(t)}
 &=
 \Tr\{B(t)^{-1}\dot B(t)-B_0(t)^{-1}\dot B_0(t)\}          \label{eq:s7v1-subtracted-trace}\\
 &=
 \Tr\{(I+\mathcal K(t))^{-1}\dot{\mathcal K}(t)\}.         \label{eq:s7v1-K-connection}
\end{align}
If \(\|\mathcal K(t)\|\le\theta<1\), then
\begin{equation}
 \left|
 \Tr\{(I+\mathcal K)^{-1}\dot{\mathcal K}\}
 \right|
 \le {1\over1-\theta}\|\dot{\mathcal K}\|_1.               \label{eq:s7v1-relative-length}
\end{equation}
Thus the continuum target is the trace norm of the derivative of the
assembled relative operator, after subtraction, rather than the two
extensive traces in \eqref{eq:s7v1-subtracted-trace} estimated separately.
\end{theorem}

\begin{proof}
Apply Jacobi's formula to the determinant quotient.  Alternatively,
\[
 {\det B\over\det B_0}=\det(I+\mathcal K)
\]
and Jacobi's formula gives \eqref{eq:s7v1-K-connection}.  Equality of the
two expressions follows.  The Neumann bound
\(\|(I+\mathcal K)^{-1}\|\le(1-\theta)^{-1}\) and trace duality yield
\eqref{eq:s7v1-relative-length}.
\end{proof}

\begin{assessment}[Cancellation must precede the trace norm]
\label{ass:s7v1-cancellation-order}
Each trace on the right side of \eqref{eq:s7v1-subtracted-trace} may scale
with the number of lattice degrees of freedom even when their difference
has a local or trace-class limit.  Bounding them separately loses the
subtraction.  Equation \eqref{eq:s7v1-K-connection} performs the exact
operator assembly before the absolute trace-norm estimate, following the
V100 audit.
\end{assessment}

\begin{corollary}[Two sufficient trace-class mechanisms]
\label{cor:s7v1-trace-mechanisms}
Under \(\|\mathcal K\|\le\theta<1\), either of the following controls the
relative logarithm and connection:
\begin{enumerate}[label=\textup{(\roman*)},leftmargin=3.5em]
\item a direct uniform bound on
      \(\|\mathcal K\|_1+\|\dot{\mathcal K}\|_1\);
\item factorizations
      \(\mathcal K=X Y\) and
      \(\dot{\mathcal K}=\dot X Y+X\dot Y\) with uniform
      Hilbert--Schmidt bounds on all displayed factors.
\end{enumerate}
In the second case,
\begin{align}
 \|\mathcal K\|_1&\le\|X\|_2\|Y\|_2,\\
 \|\dot{\mathcal K}\|_1
 &\le\|\dot X\|_2\|Y\|_2+\|X\|_2\|\dot Y\|_2.              \label{eq:s7v1-HS-factor}
\end{align}
\end{corollary}

\begin{proof}
The first statement is
\Cref{thm:s7v1-relative-series,thm:s7v1-subtracted-Jacobi}.
The Schatten ideal product inequality
\(\|XY\|_1\le\|X\|_2\|Y\|_2\) gives
\eqref{eq:s7v1-HS-factor}.
\end{proof}

\begin{proposition}[Finite-support rank bound and its limit]
\label{prop:s7v1-rank}
If \(\mathcal K\) has rank at most \(r\), then
\begin{equation}
 \|\mathcal K\|_1\le r\|\mathcal K\|,
 \qquad
 |\log\mathcal R(B,B_0)|
 \le {r\theta\over1-\theta}.                              \label{eq:s7v1-rank-bound}
\end{equation}
A perturbation supported on \(q\) lattice fibres of internal dimension
\(d_{\rm int}\) has rank at most \(d_{\rm int}q\) when its range is
supported there.  This is volume independent only when \(q\) is.  It does
not control a fixed physical region containing \(O(a^{-4})\) sites without
additional scaling or local subtraction.
\end{proposition}

\begin{proof}
The trace norm is the sum of at most \(r\) singular values, each bounded by
the operator norm.  Combine this with
\Cref{thm:s7v1-relative-series}.  The range dimension gives the support
rank.  In four dimensions, a fixed nonzero physical volume contains a
number of sites proportional to \(a^{-4}\).
\end{proof}

\section{Relative determinant acquisition card}
\label{sec:s7v1-card}

\begin{definition}[Relative determinant subtraction card, RDSC]
\label{def:s7v1-RDSC}
RDSC consists of:
\begin{enumerate}[label=\textup{(\roman*)},leftmargin=3.5em]
\item one AOP and Higgs-reduced chart with a fibrewise massive base \(B_0\);
\item an exact complete chiral block \(B=B_0(I+\mathcal K)\), including all
      projector transports and discretization terms;
\item a uniform contraction \(\|\mathcal K\|\le\theta<1\);
\item a local counterterm or factorization for which the renormalized
      \(\mathcal K\) and its field derivatives obey the required trace-class
      estimates;
\item compatible formulas on chart overlaps and trivial complete
      mapping-torus holonomy; and
\item uniform integrability on the complement of the protected chart.
\end{enumerate}
\end{definition}

\begin{theorem}[What RDSC rows \textup{(i)}--\textup{(iv)} prove]
\label{thm:s7v1-RDSC-local}
At fixed cutoff, RDSC rows \textup{(i)}--\textup{(iv)} define a
frame-independent nonzero relative determinant on each protected chart,
with logarithm and path variation bounded by
\eqref{eq:s7v1-log-bound} and
\eqref{eq:s7v1-relative-length}.  These rows do not by themselves glue the
phase around a noncontractible gauge loop or control the measure outside the
chart; those are exactly rows \textup{(v)} and \textup{(vi)}.
\end{theorem}

\begin{proof}
Frame independence is \Cref{lem:s7v1-relative-scalar}; nonvanishing and the
logarithmic bound are \Cref{thm:s7v1-relative-series}; path control is
\Cref{thm:s7v1-subtracted-Jacobi}.  The proofs are local to one chart and
make no estimate on its complement, so the final two rows remain logically
independent.
\end{proof}

\begin{programme}[V2 acquisition]
\label{prog:s7v1-V2}
The next version will derive a kernel formula for
\(\dot{\mathcal K}\) using the overlap-sign resolvent identity and the
improved-field Yukawa block.  It will test the Hilbert--Schmidt
factorization in \eqref{eq:s7v1-HS-factor} for a compactly supported
bosonic variation, identify the precise powers of \(a\), and separate local
heat-kernel counterterms from the trace-class remainder.  If the raw
factorization is still extensive, the attack will move upstream to the
relative heat-semigroup representation before taking any absolute norm.
The next companion audit is V5.
\end{programme}

\begin{firewall}[Claim boundary after seventh-series V1]
\label{fire:s7v1-boundary}
This V1 summarizes the archived programme, proves that the determinant
quotient against the fibrewise massive base is a canonical scalar, proves
its exact convergent logarithmic series under the HYFC contraction, derives
the subtracted Jacobi identity, and gives trace-class and
Hilbert--Schmidt sufficient bounds.

It does not prove the required cutoff-uniform trace-class estimate, local
counterterm subtraction, chart-complement uniform integrability,
mapping-torus trivialization, full reflection positivity, cutoff removal, or
the Standard Model--Einstein continuum.
\end{firewall}

\begin{theorem}[Seventh-series V1 result]
\label{thm:s7v1-result}
Under the local protected-chart hypotheses, subtracting the fibrewise
massive overlap block before taking a trace converts two
determinant-line-valued objects into the scalar
\(\det(I+\mathcal K)\).  HYFC supplies the operator-norm convergence radius.
A cutoff-uniform continuum bound is reduced to trace-class control of the
assembled relative operator and its derivative after local subtraction,
together with the independent global holonomy and no-escape rows of RDSC.
\end{theorem}

\begin{proof}
Use \Cref{lem:s7v1-relative-scalar,
thm:s7v1-relative-series,
cor:s7v1-HYFC-radius,
thm:s7v1-subtracted-Jacobi,
cor:s7v1-trace-mechanisms,
thm:s7v1-RDSC-local}.
\end{proof}

\paragraph{Archive-parent record.}
This V1 follows the immutable archive
\texttt{Renewal\_SM\_Einstein\_Continuum\_Research\_Notes\_V100\_f6.tex},
with \(35\,459\) logical source lines, \(1\,362\,633\) bytes, and SHA--256
\[
 \texttt{B5A5540CBDA10D6371506F77948DD85B99DFF3C09A96A3FB30A2E65F77D564DE}.
\]


% =============================================================================
% BEGIN APPENDED SEVENTH-SERIES V2 MATERIAL
% =============================================================================

\clearpage
\section{Seventh-series V2: sign-resolvent trace ideals and the local
variation ledger}
\label{sec:v7v2}

\subsection{Differentiating the overlap sign before estimating it}
\label{sec:newS7V2-sign}

Let \(H(s)\) be a \(C^1\) family of finite-dimensional self-adjoint
operators with
\begin{equation}
 H(s)^2\succeq g^2I                                       \label{eq:newS7V2-gap}
\end{equation}
for a common \(g>0\).  Put
\[
 R_s(t):=(H(s)-it)^{-1}.
\]

\begin{theorem}[Exact sign derivative and trace-norm bound]
\label{thm:newS7V2-sign-derivative}
The sign derivative is
\begin{equation}
 {d\over ds}\sgn H(s)
 =-{1\over\pi}\int_{-\infty}^{\infty}
       R_s(t)\dot H(s)R_s(t)\,dt,                          \label{eq:newS7V2-sign-derivative}
\end{equation}
where the integral converges in every finite-dimensional Schatten norm.
Moreover,
\begin{align}
 \left\|{d\over ds}\sgn H(s)\right\|
 &\le {1\over g}\|\dot H(s)\|,                             \label{eq:newS7V2-sign-op}\\
 \left\|{d\over ds}\sgn H(s)\right\|_1
 &\le {1\over g}\|\dot H(s)\|_1.                           \label{eq:newS7V2-sign-trace}
\end{align}
\end{theorem}

\begin{proof}
The resolvent sign formula gives
\[
 \sgn H(s)-\sgn H(s_0)
 ={1\over\pi}\int_{\mathbb R}
 \{R_s(t)-R_{s_0}(t)\}\,dt.
\]
The difference is absolutely norm integrable because the resolvent identity
contains two resolvents.  Differentiate that identity:
\[
 {d\over ds}R_s(t)=-R_s(t)\dot H(s)R_s(t).
\]
The gap gives
\[
 \|R_s(t)\|\le(g^2+t^2)^{-1/2}.
\]
Dominated convergence proves \eqref{eq:newS7V2-sign-derivative} and
\[
 \|R_s(t)\dot H(s)R_s(t)\|
 \le{\|\dot H(s)\|\over g^2+t^2}.
\]
For the trace norm, write the polar decomposition
\(\dot H=U|\dot H|\) and factor
\[
 R_s\dot H R_s
 =\{R_sU|\dot H|^{1/2}\}
  \{|\dot H|^{1/2}R_s\}.
\]
The Schatten product inequality gives
\[
 \|R_s\dot H R_s\|_1
 \le\|R_s\|^2\|\dot H\|_1.
\]
Finally,
\(\pi^{-1}\int_{\mathbb R}(g^2+t^2)^{-1}dt=g^{-1}\).
\end{proof}

\begin{corollary}[Kernel formula]
\label{cor:newS7V2-kernel}
In a lattice site decomposition,
\begin{equation}
 \dot\varepsilon(x,y)
 =-{1\over\pi}\int_{\mathbb R}
   \sum_{u,v}R(t;x,u)\dot H(u,v)R(t;v,y)\,dt.              \label{eq:newS7V2-kernel}
\end{equation}
Thus a localized Wilson-kernel variation is inserted once between two
resolvents before any site or Schatten sum is taken.
\end{corollary}

\begin{proof}
Take the \((x,y)\) block of
\eqref{eq:newS7V2-sign-derivative}.  Finite dimensionality and the
integrable majorant in its proof justify the block sum and integral.
\end{proof}

\begin{assessment}[Why this is the required assembly]
\label{ass:newS7V2-assembly}
Bounding two separate sign operators would retain their volume-sized
identity parts.  Formula \eqref{eq:newS7V2-sign-derivative} differentiates
the sign first and contains the field variation exactly once.  This is the
resolvent analogue of forming the relative operator
\(\mathcal K\) before taking the trace norm.
\end{assessment}

\subsection{Wilson and overlap trace-norm rows}
\label{sec:newS7V2-Wilson}

Let \(d_{\rm fib}\) be the full spin--internal dimension of one lattice site.
Let a differentiable link variation be supported on a set
\(\mathcal L_0\) of oriented positive-direction links.  Write
\[
 \|\dot U\|_{\ell^1(\mathcal L_0)}
 :=\sum_{\ell\in\mathcal L_0}\|\dot U_\ell\|.
\]

\begin{lemma}[Finite-rank Wilson variation]
\label{lem:newS7V2-Wilson-rank}
At \(\rho=1\),
\begin{align}
 \operatorname{rank}\dot H_W
 &\le2d_{\rm fib}|\mathcal L_0|,                           \label{eq:newS7V2-H-rank}\\
 \|\dot H_W\|_1
 &\le {2d_{\rm fib}\over a}
       \|\dot U\|_{\ell^1(\mathcal L_0)}.                  \label{eq:newS7V2-H-trace}
\end{align}
\end{lemma}

\begin{proof}
Vary the shift form
\[
 A=-3I+\sum_{\mu=1}^4(P_\mu^-T_\mu+P_\mu^+T_\mu^*).
\]
One link variation changes only the forward block between its two endpoint
fibres and the adjoint backward block.  Their combined rank is at most
\(2d_{\rm fib}\), and their operator norm is at most
\(\|\dot U_\ell\|\).  Since
\(\dot H_W=-a^{-1}\gamma_5\dot A\), the trace norm of that link
contribution is at most
\(2d_{\rm fib}a^{-1}\|\dot U_\ell\|\).  Sum the contributions before using
the triangle inequality.
\end{proof}

\begin{theorem}[Localized overlap trace derivative]
\label{thm:newS7V2-overlap-trace}
On the V94 admissible class with
\(c_\delta=\sqrt{1-30\delta}\), the overlap derivative satisfies
\begin{align}
 \|\dot\varepsilon\|_1
 &\le {a\over c_\delta}\|\dot H_W\|_1,                     \label{eq:newS7V2-eps-trace}\\
 \|\dot D_{\rm ov}\|_1
 &\le {1\over c_\delta}\|\dot H_W\|_1                     \label{eq:newS7V2-D-trace}\\
 &\le {2d_{\rm fib}\over ac_\delta}
       \|\dot U\|_{\ell^1(\mathcal L_0)}.                  \label{eq:newS7V2-link-trace}
\end{align}
\end{theorem}

\begin{proof}
The Wilson gap is \(g=c_\delta/a\).  Apply
\eqref{eq:newS7V2-sign-trace} with \(H=H_W\) to obtain the first line.
Since
\(\dot D_{\rm ov}=a^{-1}\gamma_5\dot\varepsilon\), the second line follows.
Insert \eqref{eq:newS7V2-H-trace} for the last.
\end{proof}

\begin{corollary}[Smooth-link scaling and the site-count obstruction]
\label{cor:newS7V2-scaling}
If a continuum-sized variation is parametrized by
\begin{equation}
 \|\dot U_\ell\|\le a\,C_A                               \label{eq:newS7V2-smooth-link}
\end{equation}
on \(L_a\) links, then
\begin{equation}
 \|\dot D_{\rm ov}\|_1
 \le {2d_{\rm fib}C_A\over c_\delta}L_a.                  \label{eq:newS7V2-count}
\end{equation}
For a fixed nonzero physical four-volume,
\(L_a=O(a^{-4})\).  Hence localized trace class at each cutoff does not give
a cutoff-uniform trace norm on a fixed physical region.
\end{corollary}

\begin{proof}
Substitute \eqref{eq:newS7V2-smooth-link} into
\eqref{eq:newS7V2-link-trace}.  A four-dimensional region of fixed physical
volume contains \(O(a^{-4})\) lattice sites and links.
\end{proof}

\begin{assessment}[Sharp logical content of the divergence]
\label{ass:newS7V2-divergence-scope}
Equation \eqref{eq:newS7V2-count} is an upper-bound ledger.  It shows that
the raw trace-norm proof retains the full link count; it does not prove that
the renormalized relative trace actually diverges.  Cancellation by local
symbol coefficients has not yet been used and must occur before the final
absolute norm.
\end{assessment}

\subsection{The improved-field Yukawa derivative}
\label{sec:newS7V2-Yukawa}

Let
\[
 A_D=I-{a\over2}D_{\rm ov},
 \qquad
 B_H=D_{\rm ov}+M_HA_D
\]
in the Higgs-reduced matched representation.

\begin{lemma}[Exact assembled Yukawa derivative]
\label{lem:newS7V2-Yukawa-derivative}
For any differentiable field path,
\begin{equation}
 \dot B_H
 =\left(I-{a\over2}M_H\right)\dot D_{\rm ov}
   +\dot M_HA_D.                                           \label{eq:newS7V2-Yukawa-derivative}
\end{equation}
If \(\dot M_H\) is supported on \(q\) site fibres and
\(\sup_x\|\dot M_H(x)\|\le C_M\), then
\begin{equation}
 \|\dot M_HA_D\|_1
 \le d_{\rm fib}qC_M.                                      \label{eq:newS7V2-Yukawa-trace}
\end{equation}
\end{lemma}

\begin{proof}
Differentiate \(B_H=D_{\rm ov}+M_H(I-aD_{\rm ov}/2)\) without commuting
\(M_H\) through \(D_{\rm ov}\).  This gives
\eqref{eq:newS7V2-Yukawa-derivative}.  The multiplication operator
\(\dot M_H\) has range dimension at most \(d_{\rm fib}q\), operator norm at
most \(C_M\), and
\(\|A_D\|\le1\).  Hence
\(\operatorname{rank}(\dot M_HA_D)\le d_{\rm fib}q\) and its trace norm is
bounded by rank times operator norm.
\end{proof}

\begin{assessment}[No ordering substitution]
\label{ass:newS7V2-order}
The term \(\dot M_HA_D\) in
\eqref{eq:newS7V2-Yukawa-derivative} has the ordering fixed in V99.
Replacing it by \(A_D\dot M_H\) would change the chiral operator by a
commutator and is not used in the trace estimate.
\end{assessment}

\subsection{Derivative of the relative operator}
\label{sec:newS7V2-relative}

Let \(\Delta B=B-B_0\) and
\(\mathcal K=B_0^{-1}\Delta B\), as in seventh-series V1.

\begin{theorem}[Exact base-covariant derivative of the relative operator]
\label{thm:newS7V2-relative-derivative}
One has
\begin{equation}
 \dot{\mathcal K}
 =B_0^{-1}\{\dot{\Delta B}-\dot B_0\mathcal K\}.            \label{eq:newS7V2-relative-derivative}
\end{equation}
If \(\|B_0^{-1}\|\le m_*^{-1}\), then
\begin{equation}
 \|\dot{\mathcal K}\|_1
 \le {1\over m_*}
 \|\dot{\Delta B}-\dot B_0\mathcal K\|_1.                  \label{eq:newS7V2-relative-trace}
\end{equation}
The subtraction inside braces must be performed before its trace norm is
estimated.
\end{theorem}

\begin{proof}
Differentiate
\(\mathcal K=B_0^{-1}\Delta B\) and use
\[
 {d\over dt}B_0^{-1}=-B_0^{-1}\dot B_0B_0^{-1}.
\]
Since \(B_0^{-1}\Delta B=\mathcal K\), the two terms assemble into
\eqref{eq:newS7V2-relative-derivative}.  The ideal property
\(\|XY\|_1\le\|X\|\|Y\|_1\) proves
\eqref{eq:newS7V2-relative-trace}.
\end{proof}

\begin{corollary}[Relative determinant connection bound]
\label{cor:newS7V2-relative-connection}
If \(\|\mathcal K\|\le\theta<1\), then
\begin{equation}
 \left|
 \Tr\{(I+\mathcal K)^{-1}\dot{\mathcal K}\}
 \right|
 \le {1\over(1-\theta)m_*}
 \|\dot{\Delta B}-\dot B_0\mathcal K\|_1.                 \label{eq:newS7V2-relative-connection}
\end{equation}
\end{corollary}

\begin{proof}
Combine the seventh-series V1 subtracted Jacobi bound with
\eqref{eq:newS7V2-relative-trace}.
\end{proof}

\begin{assessment}[What can cancel]
\label{ass:newS7V2-cancel}
The two terms in braces in
\eqref{eq:newS7V2-relative-derivative} contain the response of the complete
block and the response induced by moving the massive base.  They can share
the same local ultraviolet symbol.  Estimating
\(\|\dot{\Delta B}\|_1\) and
\(\|\dot B_0\mathcal K\|_1\) separately would discard precisely that
relative cancellation.
\end{assessment}

\subsection{Finite-cutoff trace card and next upstream step}
\label{sec:newS7V2-next}

\begin{definition}[Localized resolvent trace card, LRTC]
\label{def:newS7V2-LRTC}
LRTC consists of:
\begin{enumerate}[label=\textup{(\roman*)},leftmargin=3.5em]
\item the V94 Wilson gap on a common admissible chart;
\item the exact sign kernel \eqref{eq:newS7V2-kernel};
\item the improved-field derivative
      \eqref{eq:newS7V2-Yukawa-derivative};
\item the base-covariant subtraction
      \eqref{eq:newS7V2-relative-derivative}; and
\item a local symbol or heat-kernel subtraction after which the expression
      in braces has a cutoff-uniform trace norm on fixed physical supports.
\end{enumerate}
\end{definition}

\begin{theorem}[What LRTC rows \textup{(i)}--\textup{(iv)} establish]
\label{thm:newS7V2-LRTC-finite}
At every finite cutoff, LRTC rows \textup{(i)}--\textup{(iv)} give an exact
trace-class derivative for every finitely supported bosonic variation and
the determinant-connection bound
\eqref{eq:newS7V2-relative-connection}.  The raw bound grows with the number
of varied links or sites.  A cutoff-uniform conclusion requires row
\textup{(v)}.
\end{theorem}

\begin{proof}
Finite support gives finite trace norm by
\Cref{thm:newS7V2-overlap-trace,
lem:newS7V2-Yukawa-derivative,
thm:newS7V2-relative-derivative}.  Apply
\Cref{cor:newS7V2-relative-connection}.  The remaining site-count dependence
is \Cref{cor:newS7V2-scaling}.
\end{proof}

\begin{programme}[V3 acquisition]
\label{prog:newS7V2-V3}
V3 will go upstream to a relative heat-semigroup formula for
\(\sgn H-\sgn H_0\) and its derivative.  The local small-time coefficients
will be kept inside the difference until they are identified with the
already owned gauge, Higgs, and gravitational counterterms.  Only the
post-subtraction remainder will be tested in trace norm.  No termwise
absolute heat-trace estimate will be used before this cancellation.
\end{programme}

\begin{firewall}[Claim boundary after seventh-series V2]
\label{fire:newS7V2-boundary}
V2 proves the exact sign-resolvent derivative, operator- and trace-norm
bounds, the finite-rank Wilson variation ledger, the localized overlap and
improved-Yukawa trace rows, and the exact base-covariant derivative of the
relative operator.  It proves finite-cutoff trace class for finitely
supported variations.

V2 does not remove the \(O(a^{-4})\) fixed-physical-volume count, identify
the local heat coefficients of the relative chiral block, prove the
cutoff-uniform trace-class row, global determinant holonomy, chart
concentration, reflection positivity, cutoff removal, or the full
Standard Model--Einstein continuum.
\end{firewall}

\begin{theorem}[Seventh-series V2 result]
\label{thm:newS7V2-result}
For a gapped Wilson kernel, a field variation enters the overlap sign once
between two resolvents and its trace norm costs exactly one inverse gap.
A finite link or Higgs support therefore gives a finite-cutoff trace-class
relative determinant connection.  On a fixed physical four-volume the raw
bound retains the \(a^{-4}\) site count.  The next necessary operation is
local heat/symbol subtraction inside the base-covariant difference
\(\dot{\Delta B}-\dot B_0\mathcal K\).
\end{theorem}

\begin{proof}
Use \Cref{thm:newS7V2-sign-derivative,
lem:newS7V2-Wilson-rank,
thm:newS7V2-overlap-trace,
cor:newS7V2-scaling,
lem:newS7V2-Yukawa-derivative,
thm:newS7V2-relative-derivative,
cor:newS7V2-relative-connection,
thm:newS7V2-LRTC-finite}.
\end{proof}

\paragraph{Parent-version record.}
V2 was formed from
\texttt{Renewal\_SM\_Einstein\_Continuum\_Research\_Notes\_V1.tex}.
The V1 parent had \(551\) logical source lines, \(22\,231\) bytes, and
SHA--256
\[
 \texttt{FAF1E713CEC5543260468F6634347C6E242B14522BED869AFE0EDFF8E7C0FAFD}.
\]
The next compulsory companion audit is V5.

% =============================================================================
% END APPENDED SEVENTH-SERIES V2 MATERIAL
% =============================================================================

% =============================================================================
% BEGIN APPENDED SEVENTH-SERIES V3 MATERIAL
% =============================================================================

\clearpage
\section{Seventh-series V3: relative sign heat flow and counterterm
ownership}
\label{sec:v7v3}

\subsection{Scope relative to the archived proper-time theorem}
\label{sec:newS7V3-scope}

The frozen sixth-series V90 theorem treated a positive scalar Laplace-type
operator on a three-dimensional compact slice.  It formed a relative heat
trace, subtracted its exact linear term, and identified a
Hilbert--Schmidt modified determinant.  The present note does not repeat
that result.  It treats the four-dimensional overlap sign
\(\varepsilon=\sgn H_W\), whose square is the identity and whose dependence
on the gauge field is carried by the Wilson kernel.  The goal is to expose
the relative heat object whose local coefficients must enter the chiral
determinant subtraction card.

Put
\begin{equation}
 A:=aH_W,\qquad A_0:=aH_{W,0}.                             \label{eq:newS7V3-dimensionless}
\end{equation}
On the common V94 admissible class,
\begin{equation}
 A^2,A_0^2\succeq c_\delta^2I,\qquad
 \|A\|,\|A_0\|\le7,\qquad
 c_\delta=\sqrt{1-30\delta}.                              \label{eq:newS7V3-window}
\end{equation}
The dimensionless heat parameter below is
\(\tau=t/a^2\), where \(t\) is physical proper time.

\subsection{Exact heat representation of the sign}
\label{sec:newS7V3-sign-heat}

\begin{lemma}[Proper-time sign formula]
\label{lem:newS7V3-sign-heat}
For every finite-dimensional invertible self-adjoint \(A\),
\begin{equation}
 \sgn A
 ={1\over\sqrt\pi}\int_0^\infty
   \tau^{-1/2}A e^{-\tau A^2}\,d\tau.                     \label{eq:newS7V3-sign-heat}
\end{equation}
The integral converges in operator norm.  If
\(A^2\succeq c^2I\), its tail after \(T>0\) obeys
\begin{equation}
 \left\|
 {1\over\sqrt\pi}\int_T^\infty
 \tau^{-1/2}Ae^{-\tau A^2}\,d\tau
 \right\|
 \le {1\over\sqrt\pi}
 \int_T^\infty\tau^{-1/2}
 \sup_{\lambda\in\spec A}|\lambda|e^{-\tau\lambda^2}\,d\tau.
                                                               \label{eq:newS7V3-tail}
\end{equation}
In particular the tail tends to zero exponentially in \(T\).
\end{lemma}

\begin{proof}
For a nonzero scalar \(\lambda\),
\[
 {1\over\sqrt\pi}\int_0^\infty
 \tau^{-1/2}\lambda e^{-\tau\lambda^2}\,d\tau
 ={\lambda\over|\lambda|}
\]
by the change of variables \(u=\tau\lambda^2\) and
\(\Gamma(1/2)=\sqrt\pi\).  Apply the finite-dimensional spectral theorem.
The same diagonalization gives \eqref{eq:newS7V3-tail}; the positive spectral
gap supplies exponential large-\(T\) decay.
\end{proof}

\begin{theorem}[Heat derivative of the sign]
\label{thm:newS7V3-sign-heat-derivative}
Let \(A(s)\) be \(C^1\), self-adjoint, and obey
\(A(s)^2\succeq c^2I\) and \(\|A(s)\|\le M\) on a parameter interval.
Then
\begin{align}
 {d\over ds}\sgn A
 ={1\over\sqrt\pi}\int_0^\infty\tau^{-1/2}
 \Bigg[
 &\dot A e^{-\tau A^2}                                    \notag\\
 &-A\int_0^\tau e^{-(\tau-u)A^2}
       (A\dot A+\dot A A)e^{-uA^2}\,du
 \Bigg]\,d\tau.                                            \label{eq:newS7V3-sign-heat-derivative}
\end{align}
The integral converges in every finite-dimensional Schatten norm.
\end{theorem}

\begin{proof}
Duhamel's derivative formula is
\begin{equation}
 {d\over ds}e^{-\tau A^2}
 =-\int_0^\tau e^{-(\tau-u)A^2}
   (A\dot A+\dot A A)e^{-uA^2}\,du.                       \label{eq:newS7V3-Duhamel-derivative}
\end{equation}
Differentiate the integrand in
\eqref{eq:newS7V3-sign-heat}.  Near \(\tau=0\), the first term is
\(O(\tau^{-1/2})\) and the Duhamel term is
\(O(\tau^{1/2})\), so both are integrable.  At infinity the common gap
gives \(e^{-c^2\tau}\); finite dimensionality converts the same majorant to
every Schatten norm.  Dominated convergence proves the formula.
\end{proof}

\begin{proposition}[The termwise heat bound loses a cancellable factor]
\label{prop:newS7V3-termwise-loss}
Under the hypotheses of \Cref{thm:newS7V3-sign-heat-derivative},
termwise trace-norm estimation gives
\begin{equation}
 \left\|{d\over ds}\sgn A\right\|_1
 \le
 \left({1\over c}+{M^2\over c^3}\right)\|\dot A\|_1.       \label{eq:newS7V3-termwise-loss}
\end{equation}
The resolvent theorem of V2 gives the sharper bound
\(c^{-1}\|\dot A\|_1\).  Thus the \(M^2c^{-3}\) term is an artefact of
separating the two heat contributions before their cancellation.
\end{proposition}

\begin{proof}
The first term in
\eqref{eq:newS7V3-sign-heat-derivative} is bounded in trace norm by
\(\|\dot A\|_1e^{-c^2\tau}\).  The Duhamel integrand is bounded by
\(2M^2\|\dot A\|_1e^{-c^2\tau}\), and integration over \(u\) contributes
\(\tau\).  Therefore the right side is at most
\[
 {1\over\sqrt\pi}
 \left\{\int_0^\infty\tau^{-1/2}e^{-c^2\tau}\,d\tau
 +2M^2\int_0^\infty\tau^{1/2}e^{-c^2\tau}\,d\tau\right\}
 \|\dot A\|_1.
\]
Use
\(\Gamma(1/2)=\sqrt\pi\) and
\(\Gamma(3/2)=\sqrt\pi/2\).
The comparison with V2 is its exact resolvent trace-norm theorem, applied to
the dimensionless kernel.
\end{proof}

\subsection{The relative sign is one assembled heat object}
\label{sec:newS7V3-relative-sign}

Set
\begin{equation}
 V:=A-A_0,\qquad
 Q:=A^2-A_0^2=A_0V+VA_0+V^2.                             \label{eq:newS7V3-VQ}
\end{equation}

\begin{theorem}[Exact relative sign Duhamel formula]
\label{thm:newS7V3-relative-sign}
If \(A\) and \(A_0\) are invertible self-adjoint matrices, then
\begin{align}
 \sgn A-\sgn A_0
 ={1\over\sqrt\pi}\int_0^\infty\tau^{-1/2}
 \Bigg[
 &V e^{-\tau A^2}                                         \notag\\
 &-A_0\int_0^\tau e^{-(\tau-u)A^2}
       Qe^{-uA_0^2}\,du
 \Bigg]\,d\tau.                                            \label{eq:newS7V3-relative-sign}
\end{align}
No factor in \(Q=A_0V+VA_0+V^2\) is estimated before that sum is formed.
\end{theorem}

\begin{proof}
Start from
\[
 Ae^{-\tau A^2}-A_0e^{-\tau A_0^2}
 =Ve^{-\tau A^2}
  +A_0\{e^{-\tau A^2}-e^{-\tau A_0^2}\}.
\]
Duhamel's formula gives
\[
 e^{-\tau A^2}-e^{-\tau A_0^2}
 =-\int_0^\tau e^{-(\tau-u)A^2}
       (A^2-A_0^2)e^{-uA_0^2}\,du.
\]
Insert these identities into
\eqref{eq:newS7V3-sign-heat} for \(A\) and \(A_0\), then subtract.  The
near-zero integrability follows as in
\Cref{thm:newS7V3-sign-heat-derivative}; the spectral gaps give large-time
integrability.
\end{proof}

\begin{corollary}[Finite-cutoff relative heat trace bound]
\label{cor:newS7V3-relative-bound}
If
\[
 A^2,A_0^2\succeq c^2I,\qquad \|A_0\|\le M_0,
\]
and \(V,Q\) are trace class, then
\begin{equation}
 \|\sgn A-\sgn A_0\|_1
 \le {1\over c}\|V\|_1
      +{M_0\over2c^3}\|Q\|_1.                             \label{eq:newS7V3-relative-bound}
\end{equation}
\end{corollary}

\begin{proof}
Estimate the first term of
\eqref{eq:newS7V3-relative-sign} by
\(\|V\|_1e^{-c^2\tau}\).  The inner integral in the second term is at most
\(\tau e^{-c^2\tau}\|Q\|_1\), followed by the factor \(M_0\).  Integrate
against \(\tau^{-1/2}/\sqrt\pi\).
\end{proof}

\begin{assessment}[Role of the heat formula]
\label{ass:newS7V3-relative-role}
The bound \eqref{eq:newS7V3-relative-bound} is not intended to replace the
sharper V2 resolvent bound.  Formula
\eqref{eq:newS7V3-relative-sign} exposes the relative local heat kernel
before the regulator is removed.  Its purpose is coefficient ownership and
counterterm matching.
\end{assessment}

\subsection{An exact cancellation visible only after integration}
\label{sec:newS7V3-index-cancellation}

\begin{theorem}[Trace of the sign is constant on a gapped chart]
\label{thm:newS7V3-sign-trace-zero}
For the gapped family of
\Cref{thm:newS7V3-sign-heat-derivative},
\begin{equation}
 {d\over ds}\Tr(\sgn A(s))=0.                              \label{eq:newS7V3-sign-trace-zero}
\end{equation}
Equivalently, the trace of the complete heat integrand in
\eqref{eq:newS7V3-sign-heat-derivative}, integrated over \(\tau\), is zero.
\end{theorem}

\begin{proof}
The gap prevents any eigenvalue from crossing zero.  The numbers of positive
and negative eigenvalues are therefore constant, and their difference is
\(\Tr(\sgn A)\).  Differentiation gives
\eqref{eq:newS7V3-sign-trace-zero}.  Equality with the integrated heat
formula follows from
\Cref{thm:newS7V3-sign-heat-derivative}.
\end{proof}

\begin{lemma}[Cyclic collapse of the traced heat derivative]
\label{lem:newS7V3-cyclic-collapse}
At each fixed \(\tau>0\),
\begin{align}
 &\Tr\Bigg[
 \dot A e^{-\tau A^2}
 -A\int_0^\tau e^{-(\tau-u)A^2}
       (A\dot A+\dot A A)e^{-uA^2}\,du
 \Bigg]                                                     \notag\\
 &\hspace{7em}
 =\Tr\{\dot A(I-2\tau A^2)e^{-\tau A^2}\}.                 \label{eq:newS7V3-cyclic-collapse}
\end{align}
\end{lemma}

\begin{proof}
The operator \(A\) commutes with every function of \(A^2\).  Cyclicity of
the finite trace moves the two semigroup factors together.  Each of the two
terms containing \(A\dot A\) or \(\dot A A\) becomes
\(\Tr(\dot A A^2e^{-\tau A^2})\), independently of \(u\).
Integration over \(u\in[0,\tau]\) supplies the factor \(2\tau\).
\end{proof}

\begin{corollary}[Local heat-jet traces]
\label{cor:newS7V3-heat-jets}
The traced integrand in
\eqref{eq:newS7V3-cyclic-collapse} has the convergent expansion
\begin{equation}
 \Tr\{\dot A(I-2\tau A^2)e^{-\tau A^2}\}
 =
 \sum_{n=0}^\infty
 {(-1)^n(2n+1)\over n!}\tau^n\Tr(\dot A A^{2n}).           \label{eq:newS7V3-heat-jets}
\end{equation}
For the nearest-neighbour Wilson kernel, the coefficient
\(\Tr(\dot A A^{2n})\) is a finite sum of closed lattice walks of length at
most \(2n+1\) which meet the varied links.  It is gauge invariant only after
the complete closed trace has been formed.
\end{corollary}

\begin{proof}
Multiply the power series for \(e^{-\tau A^2}\) by
\(I-2\tau A^2\); the coefficient of
\(\tau^nA^{2n}\) is
\((-1)^n(2n+1)/n!\).  A product of \(2n+1\) nearest-neighbour-or-diagonal
kernel factors has that propagation length.  The diagonal trace closes its
endpoint, and a Wilson link variation forces the walk to meet the varied
link.  Gauge transformations conjugate every site block; conjugations
cancel around the closed trace.
\end{proof}

\begin{assessment}[Index cancellation is not a determinant cancellation]
\label{ass:newS7V3-index-vs-det}
Equation \eqref{eq:newS7V3-sign-trace-zero} concerns the unweighted trace of
the sign.  The relative determinant connection contains inverse massive
blocks, chiral projectors, Yukawa factors, and frame transports, so its heat
coefficients are different inserted traces.  The index identity proves the
need to preserve complete heat combinations; it does not cancel the
determinant counterterms.
\end{assessment}

\subsection{A regulator-uniform subtraction theorem}
\label{sec:newS7V3-subtraction}

The next statement isolates exactly what must be computed for the assembled
relative determinant connection.  Let \(a\downarrow0\) and suppose a
regulated scalar connection integrand has already been formed as one
relative object:
\begin{equation}
 \Omega_a=\int_{a^2}^{\infty}j_a(t)\,dt.                   \label{eq:newS7V3-Omega}
\end{equation}
The proper-time weight appropriate to the operator under study is included
in \(j_a\).

\begin{theorem}[Local-coefficient ownership and convergent remainder]
\label{thm:newS7V3-subtraction}
Fix \(t_0>0\).  Suppose that on \(a^2\le t\le t_0\),
\begin{equation}
 j_a(t)=
 \sum_{\ell=1}^N c_{\ell,a}t^{\alpha_\ell-1}
 +r_a(t),
 \qquad
 \alpha_\ell\le0,                                         \label{eq:newS7V3-expansion}
\end{equation}
where:
\begin{enumerate}[label=\textup{(\roman*)},leftmargin=3.5em]
\item each \(c_{\ell,a}\) is the derivative along the field path of one
      declared local gauge-covariant functional \(C_{\ell,a}\);
\item for some \(\eta>0\),
      \(|r_a(t)|\le Ct^{-1+\eta}\), uniformly in \(a\);
\item \(\int_{t_0}^{\infty}|j_a(t)|dt\le C_\infty\), uniformly in \(a\);
      and
\item \(r_a(t)\) and the large-time integrand converge pointwise along the
      cutoff sequence.
\end{enumerate}
Define
\begin{equation}
 I_\alpha(a,t_0):=
 \begin{cases}
 (t_0^\alpha-a^{2\alpha})/\alpha,&\alpha<0,\\
 \log(t_0/a^2),&\alpha=0,
 \end{cases}                                               \label{eq:newS7V3-Ialpha}
\end{equation}
and the subtracted connection
\begin{equation}
 \Omega_a^{\rm ren}
 :=\Omega_a-\sum_{\ell=1}^Nc_{\ell,a}
                         I_{\alpha_\ell}(a,t_0).            \label{eq:newS7V3-ren-Omega}
\end{equation}
Then the small-time remainder is uniformly integrable.  If the pointwise
limits in \textup{(iv)} have an integrable common large-time majorant, then
\(\Omega_a^{\rm ren}\) converges to
\begin{equation}
 \int_0^{t_0}r_0(t)\,dt+\int_{t_0}^{\infty}j_0(t)\,dt.      \label{eq:newS7V3-ren-limit}
\end{equation}
Moreover, the subtraction is the derivative of the local counterterm
\(\sum_\ell C_{\ell,a}I_{\alpha_\ell}(a,t_0)\).
\end{theorem}

\begin{proof}
Integrating each monomial in
\eqref{eq:newS7V3-expansion} from \(a^2\) to \(t_0\) gives
\eqref{eq:newS7V3-Ialpha}.  After subtraction,
\[
 \Omega_a^{\rm ren}
 =\int_{a^2}^{t_0}r_a(t)\,dt
  +\int_{t_0}^{\infty}j_a(t)\,dt.
\]
The small-time majorant is integrable because
\(\int_0^{t_0}t^{-1+\eta}dt=t_0^\eta/\eta\).
Dominated convergence on the two intervals proves
\eqref{eq:newS7V3-ren-limit}; the large-time conclusion uses the additional
majorant stated in the theorem.  Item \textup{(i)} identifies the subtracted
one-form as the field derivative of the displayed local functional.
\end{proof}

\begin{corollary}[Closed-loop compatibility of owned counterterms]
\label{cor:newS7V3-loop-counterterm}
For a closed field loop on one chart, the integral of every owned
counterterm derivative in
\Cref{thm:newS7V3-subtraction} is zero.  Hence such local subtractions do
not change the loop holonomy on that chart.  A coefficient which is not an
exact derivative of a single-valued local functional cannot be removed by
this argument and belongs to the anomaly or chart-transition ledger.
\end{corollary}

\begin{proof}
The integral of \(dC_{\ell,a}\) around a closed loop is zero when
\(C_{\ell,a}\) is single valued.  Sum over \(\ell\).  The final statement is
the contrapositive of the ownership hypothesis used in the proof.
\end{proof}

\begin{assessment}[What remains to be computed]
\label{ass:newS7V3-coefficients-open}
The theorem proves the subtraction implication but does not supply the
coefficients \(c_{\ell,a}\) for the full overlap--Yukawa determinant
connection.  They must be computed from the complete inserted heat trace,
including Standard Model representations, improved Yukawa factors, ghosts,
regulators, and metric normalization.  The archived scalar V90 coefficients
cannot be substituted for them.
\end{assessment}

\subsection{Overlap heat-subtraction card and next acquisition}
\label{sec:newS7V3-next}

\begin{definition}[Overlap relative heat-subtraction card, ORHSC]
\label{def:newS7V3-ORHSC}
ORHSC consists of:
\begin{enumerate}[label=\textup{(\roman*)},leftmargin=3.5em]
\item the common AOP spectral window
      \eqref{eq:newS7V3-window};
\item the exact relative sign formula
      \eqref{eq:newS7V3-relative-sign};
\item insertion of that formula into the complete base-covariant
      determinant connection of V2 before any trace norm;
\item computation of every nonintegrable four-dimensional small-time
      coefficient in one gauge-covariant regulator;
\item ownership of each nonanomalous coefficient by the existing local
      counterterm basis; and
\item a uniform integrable bound and cutoff limit for the post-subtraction
      remainder.
\end{enumerate}
\end{definition}

\begin{theorem}[ORHSC implication]
\label{thm:newS7V3-ORHSC}
If ORHSC holds and the HYFC contraction satisfies
\(\|\mathcal K\|\le\theta<1\), then the renormalized relative determinant
connection has a finite cutoff limit on the protected chart.  Its owned
local subtractions do not alter contractible-loop holonomy.  Global
mapping-torus trivialization and chart-complement uniform integrability
remain separate requirements.
\end{theorem}

\begin{proof}
Rows \textup{(i)}--\textup{(iii)} define the exact finite-cutoff relative
connection by V1--V2.  Rows \textup{(iv)}--\textup{(vi)} meet the hypotheses
of \Cref{thm:newS7V3-subtraction}.  The loop statement is
\Cref{cor:newS7V3-loop-counterterm}.  Neither proof addresses a
noncontractible gauge loop or fields outside the chart.
\end{proof}

\begin{programme}[V4 acquisition]
\label{prog:newS7V3-V4}
V4 will compute the first variation coefficients that are fixed without a
full pseudodifferential expansion: the gauge-covariant closed-walk terms
generated by \eqref{eq:newS7V3-heat-jets}, their charge traces over one
Standard Model generation, and the distinction between anomaly-odd and
modulus-even insertions.  It will test which low-order coefficients cancel
by the already verified hypercharge sums and which instead require gauge
kinetic or Higgs counterterms.  V5 will perform the compulsory companion
audit.
\end{programme}

\begin{firewall}[Claim boundary after seventh-series V3]
\label{fire:newS7V3-boundary}
V3 proves the dimensionless proper-time sign formula, its exact derivative,
the relative sign Duhamel formula, the sign-trace cancellation, the
closed-walk heat-jet expansion, and a regulator-uniform local-coefficient
subtraction theorem with loop-compatible ownership.

V3 does not compute the full four-dimensional Standard Model coefficients,
prove their BRST-compatible ownership, bound the post-subtraction remainder,
trivialize the global determinant line, establish chart concentration or
full reflection positivity, remove the cutoff, or construct the
Standard Model--Einstein continuum.
\end{firewall}

\begin{theorem}[Seventh-series V3 result]
\label{thm:newS7V3-result}
The overlap sign and its gauge variation admit exact relative heat
representations on the V94 admissible window.  Local heat jets are closed
gauge loops only after the complete trace is assembled, and an unweighted
sign trace has an exact cancellation which termwise absolute bounds lose.
A cutoff-uniform relative determinant connection is reduced to computing
and owning the nonintegrable coefficients of the complete inserted chiral
heat trace and proving an integrable remainder.
\end{theorem}

\begin{proof}
Use \Cref{lem:newS7V3-sign-heat,
thm:newS7V3-sign-heat-derivative,
prop:newS7V3-termwise-loss,
thm:newS7V3-relative-sign,
thm:newS7V3-sign-trace-zero,
lem:newS7V3-cyclic-collapse,
cor:newS7V3-heat-jets,
thm:newS7V3-subtraction,
cor:newS7V3-loop-counterterm,
thm:newS7V3-ORHSC}.
\end{proof}

\paragraph{Parent-version record.}
V3 was formed from
\texttt{Renewal\_SM\_Einstein\_Continuum\_Research\_Notes\_V2.tex}.
The V2 parent had \(953\) logical source lines, \(36\,277\) bytes, and
SHA--256
\[
 \texttt{24497DB4ECD645E16D2B1E686C183AF4597DF94823BE73D0D07F6838A176A5AF}.
\]
The next compulsory companion audit is V5.

% =============================================================================
% END APPENDED SEVENTH-SERIES V3 MATERIAL
% =============================================================================

% =============================================================================
% BEGIN APPENDED SEVENTH-SERIES V4 MATERIAL
% =============================================================================

\clearpage
\section{Seventh-series V4: charge-moment projection of the relative
heat response}
\label{sec:v7v4}

\subsection{Purpose and limits of the charge template}
\label{sec:newS7V4-purpose}

V3 reduced the next step to coefficients of one complete inserted
overlap--Yukawa heat trace.  This version performs the part of that
calculation fixed solely by one-generation representation arithmetic.  It
does not assume that all Standard Model multiplets have the same mass or
kinetic kernel.  Instead, it first proves an exact common-kernel template,
then states precisely which conclusions survive without that degeneracy.

Use the all-left-handed hypercharges and multiplicities
\begin{equation}
 \begin{array}{c|ccccc}
 \text{multiplet}&Q_L&u^c&d^c&L&e^c\\ \hline
 n_f&6&3&3&2&1\\
 Y_f&1/6&-2/3&1/3&-1/2&1 .
 \end{array}                                                \label{eq:newS7V4-hypercharges}
\end{equation}
These are inherited from the archived anomaly ledger.

\subsection{Exact charge moments}
\label{sec:newS7V4-moments}

Define
\begin{equation}
 S_n:=\sum_fn_fY_f^n.                                      \label{eq:newS7V4-Sn}
\end{equation}

\begin{lemma}[The first five hypercharge moments]
\label{lem:newS7V4-moments}
For \eqref{eq:newS7V4-hypercharges},
\begin{align}
 S_1&=0,                                                    \label{eq:newS7V4-S1}\\
 S_2&={10\over3},                                          \label{eq:newS7V4-S2}\\
 S_3&=0,                                                    \label{eq:newS7V4-S3}\\
 S_4&={95\over54},                                         \label{eq:newS7V4-S4}\\
 S_5&={5\over9}.                                           \label{eq:newS7V4-S5}
\end{align}
\end{lemma}

\begin{proof}
The first three values are the archived gravitational, quadratic, and cubic
hypercharge sums.  Direct common-denominator calculations for the new rows
give
\begin{align*}
 S_4
 &= {1\over216}+{16\over27}+{1\over27}+{1\over8}+1
  ={95\over54},\\
 S_5
 &= {1\over1296}-{32\over81}+{1\over81}-{1\over16}+1
  ={5\over9}.
\end{align*}
\end{proof}

\begin{assessment}[What anomaly cancellation does not say]
\label{ass:newS7V4-not-all-odd}
The equalities \(S_1=S_3=0\) are the Abelian gravitational and cubic local
anomaly arithmetic.  Equation \eqref{eq:newS7V4-S5} proves that the Standard
Model hypercharge list does not cancel every odd charge moment.  Higher
gauge-invariant phase responses are organized by Ward identities and the
four-dimensional anomaly cohomology, not by the assertion
\(S_{2k+1}=0\) for all \(k\).
\end{assessment}

\subsection{All-order common-kernel determinant identity}
\label{sec:newS7V4-template}

Let \(K\) be a finite matrix with \(\|zY_fK\|<1\) for every \(f\).  Define
the degenerate Abelian template
\begin{equation}
 \mathscr W(z):=\sum_fn_f\Tr\log(I+zY_fK).                 \label{eq:newS7V4-W}
\end{equation}

\begin{theorem}[Charge moments project every Taylor coefficient]
\label{thm:newS7V4-template}
The function \(\mathscr W\) is analytic near zero and
\begin{equation}
 \mathscr W(z)
 =\sum_{n=1}^{\infty}
 {(-1)^{n+1}\over n}S_nz^n\Tr(K^n).                       \label{eq:newS7V4-template}
\end{equation}
Equivalently,
\begin{equation}
 \mathscr W^{(n)}(0)
 =(-1)^{n+1}(n-1)!S_n\Tr(K^n).                            \label{eq:newS7V4-derivative}
\end{equation}
\end{theorem}

\begin{proof}
For each \(f\), the norm-convergent logarithm is
\[
 \Tr\log(I+zY_fK)
 =\sum_{n\ge1}{(-1)^{n+1}\over n}
  z^nY_f^n\Tr(K^n).
\]
The flavour set is finite, so it may be summed term by term.  The charge
coefficient is exactly \(S_n\).  Differentiate at zero.
\end{proof}

\begin{corollary}[First nonzero terms]
\label{cor:newS7V4-first-terms}
For one generation,
\begin{align}
 \mathscr W(z)
 ={}&
 -{5\over3}z^2\Tr(K^2)
 -{95\over216}z^4\Tr(K^4)
 +{1\over9}z^5\Tr(K^5)
 +O(z^6),                                                   \label{eq:newS7V4-first-terms}
\end{align}
with no linear or cubic term.
\end{corollary}

\begin{proof}
Insert \Cref{lem:newS7V4-moments} into
\eqref{eq:newS7V4-template}.  The fifth coefficient is
\(S_5/5=1/9\).
\end{proof}

\begin{assessment}[Degeneracy firewall]
\label{ass:newS7V4-degeneracy}
Equations \eqref{eq:newS7V4-template}--\eqref{eq:newS7V4-first-terms}
factor out one common matrix \(K\).  Actual Standard Model blocks have
different non-Abelian representations, Yukawa masses, and Higgs couplings.
Their modulus coefficients cannot be cancelled by inserting charge moments
unless a common operator coefficient has first been proved.  The anomaly
coefficients are representation-theoretic and survive mass deformations,
but generic finite determinant coefficients do not share that protection.
\end{assessment}

\subsection{Odd phase and even modulus in an anti-Hermitian direction}
\label{sec:newS7V4-parity}

\begin{lemma}[Trace parity]
\label{lem:newS7V4-trace-parity}
If \(K^*=-K\), then
\begin{equation}
 \overline{\Tr(K^n)}=(-1)^n\Tr(K^n).                       \label{eq:newS7V4-trace-parity}
\end{equation}
Hence \(\Tr(K^{2j})\) is real and
\(\Tr(K^{2j+1})\) is purely imaginary.
\end{lemma}

\begin{proof}
Use
\[
 \overline{\Tr(K^n)}
 =\Tr((K^*)^n)
 =(-1)^n\Tr(K^n).
\]
\end{proof}

\begin{theorem}[Template separation of modulus and phase]
\label{thm:newS7V4-modulus-phase}
Under the common-kernel hypotheses of
\Cref{thm:newS7V4-template}, if \(z\in\mathbb R\) and \(K^*=-K\), then the
even-\(n\) terms of \(\mathscr W(z)\) are real and the odd-\(n\) terms are
purely imaginary.  Therefore:
\begin{enumerate}[label=\textup{(\roman*)},leftmargin=3.5em]
\item \(S_1=S_3=0\) removes the first and third template phase
      coefficients;
\item \(S_2=10/3\) leaves a nonzero quadratic modulus response;
\item \(S_4=95/54\) leaves a nonzero fourth-order modulus response; and
\item \(S_5=5/9\) permits a fifth-order finite phase response compatible
      with cancellation of the four-dimensional local anomaly.
\end{enumerate}
\end{theorem}

\begin{proof}
Apply \Cref{lem:newS7V4-trace-parity} term by term in
\eqref{eq:newS7V4-template}, then use
\Cref{lem:newS7V4-moments}.  The final item is not an anomaly claim: it says
only that the charge arithmetic does not force the fifth template
coefficient to vanish.
\end{proof}

\begin{assessment}[Power counting is separate from charge arithmetic]
\label{ass:newS7V4-power-counting}
A nonzero charge moment does not prove that the associated coefficient is
ultraviolet divergent.  In four dimensions, power counting and the
small-time heat degree decide which local operators require counterterms.
The charge ledger supplies the representation trace only after the common
local operator tensor has been identified.
\end{assessment}

\subsection{Vectorlike pairing cancels every odd template term}
\label{sec:newS7V4-vectorlike}

\begin{lemma}[Exact conjugate-charge pairing]
\label{lem:newS7V4-conjugate-pair}
For a pair of charges \(y,-y\),
\begin{equation}
 \Tr\log(I+zyK)+\Tr\log(I-zyK)
 =\Tr\log(I-z^2y^2K^2)                                    \label{eq:newS7V4-pair}
\end{equation}
on the logarithm branch through \(z=0\).  All odd Taylor coefficients
cancel exactly.
\end{lemma}

\begin{proof}
The two factors commute because both are polynomials in \(K\), and
\[
 (I+zyK)(I-zyK)=I-z^2y^2K^2.
\]
Take determinants and the analytic logarithm through the identity.
\end{proof}

\begin{corollary}[Why the fibrewise massive base has no phase series]
\label{cor:newS7V4-base-phase}
After the V98 left--right matching, each charged massive species is
vectorlike under the unbroken colour--electromagnetic group.  In
all-left-handed notation its two Weyl representations are dual, so the
common unbroken-gauge template has conjugate charges and every odd term
cancels as in \eqref{eq:newS7V4-pair}.  This is the Taylor form of the exact
positive massive-overlap determinant proved in the frozen V95 theorem.
\end{corollary}

\begin{proof}
The V98 charge ledger pairs the physical left and right fields with equal
electric charge, equivalently the two all-left-handed fields with dual
charges.  Apply \Cref{lem:newS7V4-conjugate-pair}.  Exact positivity of the
full finite-cutoff base is stronger and was proved from
gamma-five Hermiticity and the massive overlap spectral circle.
\end{proof}

\begin{assessment}[Scope of the vectorlike subtraction]
\label{ass:newS7V4-vectorlike-scope}
The subtraction by \(B_0\) removes the unbroken vectorlike base response
from the relative determinant.  It does not make the broken electroweak
chiral block vectorlike, remove its local anomaly by itself, or trivialize
the determinant line across Higgs zeros and noncontractible gauge loops.
\end{assessment}

\subsection{The anomaly rows that are representation exact}
\label{sec:newS7V4-anomaly-rows}

\begin{proposition}[One-generation local gauge-anomaly coefficients]
\label{prop:newS7V4-anomaly-rows}
For the one-generation all-left-handed Standard Model representation:
\begin{align}
 [U(1)_Y]^3 &: \quad \sum_fn_fY_f^3=0,                     \label{eq:newS7V4-U1cube}\\
 \mathrm{grav}^2U(1)_Y &: \quad \sum_fn_fY_f=0,             \label{eq:newS7V4-gravU1}\\
 [SU(2)]^2U(1)_Y &: \quad 3(1/6)+(-1/2)=0,                 \label{eq:newS7V4-SU2U1}\\
 [SU(3)]^2U(1)_Y &: \quad 2(1/6)-2/3+1/3=0.                \label{eq:newS7V4-SU3U1}
\end{align}
The pure \(SU(3)^3\) anomaly cancels between the two triplet components of
\(Q_L\) and the two anti-triplets \(u^c,d^c\).  The perturbative
\(SU(2)^3\) anomaly vanishes because the doublet representation is
pseudoreal, and the mod-two \(SU(2)\) global count is even: three coloured
quark doublets plus one lepton doublet.
\end{proposition}

\begin{proof}
The first two equations are
\eqref{eq:newS7V4-S3} and \eqref{eq:newS7V4-S1}.  For the mixed anomalies,
the common Dynkin indices factor out, leaving the displayed hypercharge
sums.  Fundamental and anti-fundamental cubic \(SU(3)\) indices have
opposite signs and equal magnitude.  The \(SU(2)\) fundamental is
pseudoreal, and the number of left doublets is \(3+1=4\).
\end{proof}

\begin{assessment}[Local arithmetic versus regulator theorem]
\label{ass:newS7V4-regulator}
The proposition supplies the representation coefficients of the local
anomaly polynomial.  A regulator-level theorem must still show that the
complete lattice phase curvature converges to that polynomial plus an exact
local variation, with uniform control of the remainder.  Only then can the
zero coefficients be used in ORHSC.  The global mapping-torus holonomy is
not determined by these local sums.
\end{assessment}

\subsection{Coefficient ownership ledger}
\label{sec:newS7V4-ownership}

\begin{definition}[Charge-arithmetic projection card, CAPC]
\label{def:newS7V4-CAPC}
CAPC consists of:
\begin{enumerate}[label=\textup{(\roman*)},leftmargin=3.5em]
\item an assembled inserted heat coefficient decomposed into common local
      operator tensors and finite representation traces;
\item use of \Cref{prop:newS7V4-anomaly-rows} only on the anomaly tensors;
\item retention of nonzero even traces such as \(S_2\) as gauge, Higgs, or
      metric counterterm coefficients;
\item proof that every removed modulus coefficient is the variation of one
      gauge-covariant local functional;
\item a regulator-level bound on the coefficient remainder; and
\item separate transport of the finite gauge-invariant phase and the global
      determinant-line holonomy.
\end{enumerate}
\end{definition}

\begin{theorem}[Consequences of CAPC rows \textup{(i)}--\textup{(iv)}]
\label{thm:newS7V4-CAPC-local}
If CAPC rows \textup{(i)}--\textup{(iv)} hold for a nonintegrable
small-time coefficient, then its locally anomalous representation part
vanishes for one Standard Model generation, while its nonzero even modulus
part is assigned to a declared local counterterm.  This subtraction does not
alter contractible-loop holonomy when the counterterm is single valued.
It does not bound the heat remainder or the finite/global phase.
\end{theorem}

\begin{proof}
Apply the zero anomaly coefficients of
\Cref{prop:newS7V4-anomaly-rows} to the anomaly tensors and retain the
nonzero even coefficients from
\Cref{lem:newS7V4-moments}.  Counterterm ownership and contractible-loop
invariance are the V3 local-coefficient subtraction theorem.  Rows
\textup{(v)}--\textup{(vi)} are not among the hypotheses and therefore
remain open.
\end{proof}

\begin{programme}[V5 acquisition and compulsory audit]
\label{prog:newS7V4-V5}
V5 will first take immutable snapshots of the newest settled Yang--Mills and
Navier--Stokes notes.  It will then choose between two upstream coefficient
calculations: a covariant lattice Ward identity for the inserted
overlap-sign heat trace, or a relative heat remainder estimate after
subtracting the common closed-walk tensors.  The choice will be made from
the companion bottlenecks.  No cancellation will be inferred from charge
moments before the common operator tensor is proved.
\end{programme}

\begin{firewall}[Claim boundary after seventh-series V4]
\label{fire:newS7V4-boundary}
V4 computes the first five one-generation hypercharge moments, proves their
all-order common-kernel determinant projection, separates odd phase and
even modulus terms in an anti-Hermitian direction, proves exact
vectorlike odd-term cancellation, and records the complete
one-generation local anomaly arithmetic.

V4 does not derive the common-kernel factorization for the actual
overlap--Yukawa heat coefficients, compute their four-dimensional
small-time tensors, prove regulator convergence to the anomaly polynomial,
bound the post-subtraction remainder, establish global holonomy or chart
concentration, prove full reflection positivity, remove the cutoff, or
construct the Standard Model--Einstein continuum.
\end{firewall}

\begin{theorem}[Seventh-series V4 result]
\label{thm:newS7V4-result}
The Standard Model charge ledger cancels the linear and cubic Abelian
anomaly traces and all mixed local gauge anomalies, but leaves the quadratic
and fourth hypercharge moments nonzero and does not cancel every odd finite
phase coefficient.  The unbroken fibrewise massive base cancels all odd
terms by exact vectorlike pairing.  Applying these facts to ORHSC requires
first factoring each complete inserted heat coefficient into its common
operator tensor and representation trace.
\end{theorem}

\begin{proof}
Use \Cref{lem:newS7V4-moments,
thm:newS7V4-template,
thm:newS7V4-modulus-phase,
lem:newS7V4-conjugate-pair,
cor:newS7V4-base-phase,
prop:newS7V4-anomaly-rows,
thm:newS7V4-CAPC-local}.
\end{proof}

\paragraph{Parent-version record.}
V4 was formed from
\texttt{Renewal\_SM\_Einstein\_Continuum\_Research\_Notes\_V3.tex}.
The V3 parent had \(1\,454\) logical source lines, \(55\,429\) bytes, and
SHA--256
\[
 \texttt{E1B5395DFA65C8A2C56DFC07704E66EFFDAAD2D1949AFA92094FC70E7FA19EB7}.
\]
The next compulsory companion audit is V5.

% =============================================================================
% END APPENDED SEVENTH-SERIES V4 MATERIAL
% =============================================================================

% =============================================================================
% BEGIN APPENDED SEVENTH-SERIES V5 MATERIAL
% =============================================================================

\clearpage
\section{Seventh-series V5: companion audit, the exact relative Ward
identity, and the heat tail}
\label{sec:v7v5}

\subsection{Compulsory companion audit}
\label{sec:newS7V5-audit}

The companion files were read as immutable byte snapshots.  The settled
Yang--Mills body was
\begin{center}
\begin{tabular}{ll}
file &
\texttt{Renewal\_Yang\_Mills\_Mass\_Gap\_Research\_Notes\_V77.tex},\\
bytes & \(1\,144\,602\),\\
logical source lines & \(35\,887\),\\
SHA--256 &
\texttt{D9416F63F155CC90A99BBC6C684301368998AB361683BD9CA71B29782A83D224}.
\end{tabular}
\end{center}
A simultaneously visible file named
\texttt{Renewal\_Yang\_Mills\_Mass\_Gap\_Research\_Notes\_V78.tex}
had exactly the same bytes and digest and ended with the V77 marker and
parent record.  It contained no V78 body and was discarded as a
transitional duplicate.

The settled Navier--Stokes snapshot was
\begin{center}
\begin{tabular}{ll}
file &
\texttt{Renewal\_NS\_Stability\_Research\_Notes\_V55.tex},\\
bytes & \(658\,662\),\\
logical source lines & \(12\,925\),\\
SHA--256 &
\texttt{E46B3CEAAA668CEA75543823746A283A38F0B5F39AD523452CF000FF7EC065F2}.
\end{tabular}
\end{center}

\begin{audit}[Transfer at V5]
\label{audit:newS7V5-transfer}
Yang--Mills V77 expands every local child-star power into weighted Cayley
trees, glues them to one global tree, and proves that the gluing fibre has
the sharp exponential size \(2^{p-2}\).  This closes the full heterogeneous
incidence-hypertree core.  The decisive order is exact global gluing before
coefficientwise domination; cyclic incidence cores remain open.

Navier--Stokes V55 consolidates the result that a spread type-I window is
uniformly weak and cannot grow or carry a travelling concentration while
its \(H^2\) budget is available.  The unresolved spike can begin only after
that window, when the controlling budget lapses.  Thus a local window
estimate cannot be extrapolated without a separate continuation or tail
bound.

For the present programme, the exact gauge-covariant relative determinant
must be assembled before its heat coefficients are bounded.  The
large-proper-time part must also be controlled independently rather than
being inferred from a small-time expansion.  The next sections carry out
both operations.  No numerical estimate is imported from either companion
programme.
\end{audit}

\subsection{Gauge covariance of a relative determinant}
\label{sec:newS7V5-relative-Ward}

Let a gauge transformation \(g\) act unitarily on the source and row fibres
by \(G_E(g)\) and \(G_F(g)\).  Suppose two invertible blocks
\(B,B_0:E\to F\) have the same covariance law
\begin{equation}
 B^g=G_F(g)B\,G_E(g)^{-1},
 \qquad
 B_0^g=G_F(g)B_0\,G_E(g)^{-1}.                            \label{eq:newS7V5-common-covariance}
\end{equation}
Define
\[
 \mathcal K=B_0^{-1}(B-B_0),
 \qquad
 \mathcal R(B,B_0)=\det(I+\mathcal K).
\]

\begin{theorem}[Exact finite-cutoff relative Ward identity]
\label{thm:newS7V5-relative-Ward}
Under \eqref{eq:newS7V5-common-covariance},
\begin{align}
 \mathcal K^g
 &=G_E(g)\mathcal K G_E(g)^{-1},                           \label{eq:newS7V5-K-covariance}\\
 \mathcal R(B^g,B_0^g)
 &=\mathcal R(B,B_0).                                      \label{eq:newS7V5-R-invariant}
\end{align}
For an infinitesimal gauge path with source generator \(X_E\),
\begin{equation}
 \delta_\xi\mathcal K=[X_E,\mathcal K],
 \qquad
 \delta_\xi\log\mathcal R
 =\Tr\{(I+\mathcal K)^{-1}[X_E,\mathcal K]\}=0.            \label{eq:newS7V5-infinitesimal-Ward}
\end{equation}
The conclusion is exact at each finite cutoff and does not use an
operator-norm or heat-kernel estimate.
\end{theorem}

\begin{proof}
Invert the second covariance law and multiply:
\begin{align*}
 (B_0^g)^{-1}(B^g-B_0^g)
 &=
 G_EB_0^{-1}G_F^{-1}
 G_F(B-B_0)G_E^{-1}\\
 &=G_E\mathcal KG_E^{-1}.
\end{align*}
Taking the determinant of \(I+\mathcal K\) proves
\eqref{eq:newS7V5-R-invariant}.  Differentiation gives the commutator in
\eqref{eq:newS7V5-infinitesimal-Ward}.  Since
\((I+\mathcal K)^{-1}\) is a function of \(\mathcal K\), it commutes with
\(\mathcal K\), and cyclicity of the finite trace makes the last expression
zero.
\end{proof}

\begin{corollary}[The two extensive Ward traces cancel before estimation]
\label{cor:newS7V5-two-traces}
Under the same hypotheses,
\begin{equation}
 \Tr(B^{-1}\delta_\xi B)
 -\Tr(B_0^{-1}\delta_\xi B_0)=0.                           \label{eq:newS7V5-two-traces}
\end{equation}
Neither trace needs to have a cutoff-uniform absolute bound.
\end{corollary}

\begin{proof}
Differentiate
\(\log(\det B/\det B_0)=\log\mathcal R\) and use
\eqref{eq:newS7V5-infinitesimal-Ward}.
\end{proof}

\begin{assessment}[Where a chiral anomaly can still enter]
\label{ass:newS7V5-anomaly-scope}
The theorem applies when \(B_0\) is a nonvanishing covariant reference in
the same determinant line as \(B\).  Constructing such a base globally is
part of the problem: the Higgs reduction fails at Higgs zeros, index sectors
have different ranks, and a gauge loop may return the chosen frames with
nontrivial determinant-line holonomy.  A regulator or fermion measure whose
Jacobian is not contained in the common covariance law also contributes an
anomaly term.  Thus the local relative Ward identity does not prove global
trivialization.
\end{assessment}

\subsection{Gauge covariance inside the sign heat flow}
\label{sec:newS7V5-sign-Ward}

Let \(A^g=G(g)AG(g)^{-1}\) be a gauge-covariant dimensionless Wilson
kernel.  Along \(g(s)=e^{sX}\),
\begin{equation}
 \dot A=[X,A].                                             \label{eq:newS7V5-pure-gauge-A}
\end{equation}

\begin{theorem}[The complete sign heat flow is a commutator]
\label{thm:newS7V5-sign-Ward}
For the pure-gauge variation
\eqref{eq:newS7V5-pure-gauge-A},
\begin{equation}
 {d\over ds}\sgn A(s)\bigg|_{s=0}
 =[X,\sgn A].                                              \label{eq:newS7V5-sign-commutator}
\end{equation}
The complete proper-time expression in V3 equals this commutator.  Hence
\begin{equation}
 \Tr\left({d\over ds}\sgn A(s)\bigg|_{s=0}\right)=0.       \label{eq:newS7V5-sign-Ward-trace}
\end{equation}
\end{theorem}

\begin{proof}
Functional calculus commutes with unitary conjugation:
\[
 \sgn(GAG^{-1})=G(\sgn A)G^{-1}.
\]
Differentiate at the identity.  V3 proved that its proper-time integral is
the derivative of the same sign operator, so it equals the commutator.
The trace of a commutator vanishes.
\end{proof}

\begin{lemma}[Every closed-walk heat jet obeys the pure-gauge Ward
identity]
\label{lem:newS7V5-jet-Ward}
For every \(n\ge0\),
\begin{equation}
 \Tr\{[X,A]A^{2n}\}=0.                                    \label{eq:newS7V5-jet-Ward}
\end{equation}
Thus the closed-walk coefficients in V3 vanish separately on a pure-gauge
tangent after the complete trace is formed.
\end{lemma}

\begin{proof}
Cyclicity gives
\[
 \Tr\{[X,A]A^{2n}\}
 =\Tr(XA^{2n+1})-\Tr(AXA^{2n})
 =0.
\]
\end{proof}

\begin{assessment}[Local Ward form]
\label{ass:newS7V5-local-Ward}
When \(X\) is the diagonal multiplication operator determined by arbitrary
site parameters \(\xi(x)\), expanding
\eqref{eq:newS7V5-jet-Ward} in those parameters gives the lattice
covariant-divergence identity for the corresponding link current, including
all endpoints of the closed walks.  Dropping one endpoint before the trace
is closed would destroy that identity.
\end{assessment}

\subsection{Coefficientwise inheritance of an assembled Ward identity}
\label{sec:newS7V5-coefficient-Ward}

\begin{lemma}[Uniqueness of a finite asymptotic scale]
\label{lem:newS7V5-asymptotic-unique}
Let
\(\beta_1<\cdots<\beta_N<\beta_{N+1}\), and suppose
\begin{equation}
 \sum_{j=1}^Nc_jt^{\beta_j}
 +O(t^{\beta_{N+1}})=0
 \quad(t\downarrow0).                                     \label{eq:newS7V5-asymptotic-zero}
\end{equation}
Then \(c_1=\cdots=c_N=0\).
\end{lemma}

\begin{proof}
Divide by \(t^{\beta_1}\) and let \(t\downarrow0\) to obtain \(c_1=0\).
Remove that term and repeat.
\end{proof}

\begin{theorem}[Assembled heat coefficients satisfy Ward identities]
\label{thm:newS7V5-coefficient-Ward}
Suppose a gauge-covariant regulator gives the pure-gauge variation of the
complete relative connection an asymptotic expansion
\begin{equation}
 \delta_\xi j_a(t)
 =\sum_{j=1}^Nc_{j,a}(\xi)t^{\beta_j}
  +O(t^{\beta_{N+1}}),                                    \label{eq:newS7V5-Ward-expansion}
\end{equation}
with distinct increasing exponents.  If the exact relative Ward identity
makes the left side zero for every \(t>0\), then
\begin{equation}
 c_{j,a}(\xi)=0
 \quad\hbox{for every }j,\ a,\ \xi.                        \label{eq:newS7V5-coefficient-zero}
\end{equation}
Thus every nonintegrable coefficient is gauge invariant after the full
relative operator, representations, projectors, and regulator factors have
been assembled.
\end{theorem}

\begin{proof}
For fixed \(a\) and \(\xi\), apply
\Cref{lem:newS7V5-asymptotic-unique} to
\eqref{eq:newS7V5-Ward-expansion}.  The arbitrary gauge parameter yields
the Ward identity for each coefficient functional.
\end{proof}

\begin{assessment}[Anomaly interpretation]
\label{ass:newS7V5-coefficient-anomaly}
If the regulated chiral measure has an anomaly, the exact left side of
\eqref{eq:newS7V5-Ward-expansion} contains its Jacobian and is not zero
until all Standard Model multiplets and counterterms are included.  V4's
anomaly arithmetic predicts cancellation of the local cohomology class;
V5 proves that, once an exact common-covariance regulator has been
constructed, coefficientwise Ward identities follow automatically.  It
does not construct that global regulator.
\end{assessment}

\subsection{The independent large-time heat tail}
\label{sec:newS7V5-tail}

Let
\[
 \Gamma(\alpha,x):=\int_x^\infty u^{\alpha-1}e^{-u}\,du
\]
be the upper incomplete gamma function.

\begin{theorem}[Uniform tail of the sign derivative]
\label{thm:newS7V5-heat-tail}
Under the V3 hypotheses
\[
 A^2\succeq c^2I,\qquad \|A\|\le M,
\]
the trace norm of the part of the proper-time sign derivative with
\(\tau\ge T\) is at most
\begin{equation}
 \left\{
 {\Gamma(1/2,c^2T)\over\sqrt\pi\,c}
 +{2M^2\Gamma(3/2,c^2T)\over\sqrt\pi\,c^3}
 \right\}\|\dot A\|_1.                                    \label{eq:newS7V5-heat-tail}
\end{equation}
The coefficient tends to zero exponentially as \(T\to\infty\), uniformly
in lattice volume and throughout a common spectral window.
\end{theorem}

\begin{proof}
The termwise majorant established in V3 is
\[
 {1\over\sqrt\pi}
 \{\tau^{-1/2}+2M^2\tau^{1/2}\}e^{-c^2\tau}\|\dot A\|_1.
\]
Integrate from \(T\) to infinity and substitute \(u=c^2\tau\).  This gives
\eqref{eq:newS7V5-heat-tail}.  Standard integration by parts, or the
defining integral, shows exponential decay of
\(\Gamma(\alpha,c^2T)\) for fixed \(\alpha\).
\end{proof}

\begin{corollary}[Only the short-time relative coefficient problem remains]
\label{cor:newS7V5-short-time}
On an AOP chart, the large-\(\tau\) part of the inserted overlap-sign heat
flow is uniformly controlled by
\eqref{eq:newS7V5-heat-tail}.  Therefore the unresolved cutoff dependence
in ORHSC lies in:
\begin{enumerate}[label=\textup{(\roman*)},leftmargin=3.5em]
\item the conversion between dimensionless \(\tau\) and physical
      \(t=a^2\tau\);
\item the local short-time coefficient tensors of the complete relative
      determinant connection; and
\item the uniform remainder after their owned counterterms are removed.
\end{enumerate}
\end{corollary}

\begin{proof}
Use \(c=c_\delta\) and \(M=7\) from the V94 window.  The tail coefficient is
independent of lattice volume and cutoff in dimensionless variables.  The
three listed rows are not part of that tail estimate.
\end{proof}

\subsection{Audit decision and next acquisition}
\label{sec:newS7V5-next}

\begin{definition}[Relative Ward--heat card, RWHC]
\label{def:newS7V5-RWHC}
RWHC consists of:
\begin{enumerate}[label=\textup{(\roman*)},leftmargin=3.5em]
\item a covariant base and full chiral block in the same determinant line;
\item the exact finite-cutoff relative Ward identity
      \eqref{eq:newS7V5-infinitesimal-Ward};
\item a common regulator whose inserted relative heat trace has a
      coefficient expansion;
\item coefficientwise Ward identities from
      \Cref{thm:newS7V5-coefficient-Ward};
\item the uniform large-time bound
      \eqref{eq:newS7V5-heat-tail}; and
\item owned local subtractions and a uniform short-time remainder.
\end{enumerate}
\end{definition}

\begin{theorem}[RWHC reduction]
\label{thm:newS7V5-RWHC}
Rows \textup{(i)}--\textup{(v)} close the exact gauge covariance and
large-time part of the relative heat problem at finite cutoff.  If row
\textup{(vi)} also holds, the renormalized relative determinant connection
has the protected-chart limit supplied by ORHSC.  These rows do not imply a
global base across Higgs zeros or trivial mapping-torus holonomy.
\end{theorem}

\begin{proof}
Use \Cref{thm:newS7V5-relative-Ward,
thm:newS7V5-coefficient-Ward,
thm:newS7V5-heat-tail} and then the V3 ORHSC implication.
The scope limitation is \Cref{ass:newS7V5-anomaly-scope}.
\end{proof}

\begin{programme}[V6 acquisition]
\label{prog:newS7V5-V6}
The audit selects the coefficient calculation only after the exact global
Ward assembly.  V6 will construct the first local lattice tensor by varying
a closed Wilson loop, prove its covariant-divergence identity with both
endpoints retained, and identify the continuum \(F_{\mu\nu}F^{\mu\nu}\)
owner under a smooth-link scaling.  The calculation will remain
representationwise until the common tensor is explicit; only then will the
V4 charge traces be inserted.  The next companion audit is V10.
\end{programme}

\begin{firewall}[Claim boundary after seventh-series V5]
\label{fire:newS7V5-boundary}
V5 performs the compulsory companion audit, proves exact gauge invariance of
the finite-cutoff relative determinant under a common covariance law,
derives its infinitesimal Ward identity, proves the corresponding sign and
closed-walk heat-jet identities, transfers an exact Ward identity to every
distinct asymptotic coefficient, and proves a uniform incomplete-gamma
large-time tail.

V5 does not construct a global covariant massive base, derive the full
chiral regulator Jacobian, compute the short-time coefficient tensors,
prove the local anomaly limit or the post-subtraction remainder, establish
mapping-torus triviality or chart concentration, prove full reflection
positivity, remove the cutoff, or construct the Standard Model--Einstein
continuum.
\end{firewall}

\begin{theorem}[Seventh-series V5 result]
\label{thm:newS7V5-result}
The relative determinant on a protected chart has an exact finite-cutoff
Ward identity whenever its full and base blocks share the same covariance
law.  The identity is inherited coefficientwise by any unique heat
asymptotic expansion after complete assembly.  The AOP Wilson gap separately
controls the entire large-time sign-heat tail.  The remaining ORHSC work is
therefore the gauge-covariant short-time tensor and its renormalized
remainder, together with the global determinant-line and no-escape gates.
\end{theorem}

\begin{proof}
Use \Cref{audit:newS7V5-transfer,
thm:newS7V5-relative-Ward,
cor:newS7V5-two-traces,
thm:newS7V5-sign-Ward,
lem:newS7V5-jet-Ward,
thm:newS7V5-coefficient-Ward,
thm:newS7V5-heat-tail,
cor:newS7V5-short-time,
thm:newS7V5-RWHC}.
\end{proof}

\paragraph{Parent-version record.}
V5 was formed from
\texttt{Renewal\_SM\_Einstein\_Continuum\_Research\_Notes\_V4.tex}.
The V4 parent had \(1\,836\) logical source lines, \(70\,715\) bytes, and
SHA--256
\[
 \texttt{14D63CD7DF0FA081BE02394D630A7090B94D133C545B4A118A6070B20D6C1CCA}.
\]
The next compulsory companion audit is V10.

% =============================================================================
% END APPENDED SEVENTH-SERIES V5 MATERIAL
% =============================================================================

% =============================================================================
% BEGIN APPENDED SEVENTH-SERIES V6 MATERIAL
% =============================================================================

\clearpage
\section{Seventh-series V6: the closed-loop current, its exact lattice
Ward identity, and the quadratic-curvature owner}
\label{sec:v7v6}

V5 proved the Ward identity at the level of a complete relative
determinant and showed that a unique short-time expansion inherits that
identity coefficient by coefficient.  The first unresolved local object is
therefore not an isolated matrix entry but the current obtained after a
closed walk has been traced.  This version constructs that current, retains
both endpoints of every varied link, and then takes its smooth-link limit.

\subsection{Oriented links and left-trivialized currents}
\label{sec:newS7V6-current-setup}

Let \(K\) be a compact Lie group with a faithful unitary representation on
\(V_R\).  Write
\[
 \mathfrak k_R\subseteq\mathfrak u(V_R),
 \qquad
 \langle X,Y\rangle_R=-\Re\tr_R(XY).
\]
This is a positive definite, adjoint-invariant inner product.  On a finite
lattice choose one orientation of each geometric link and denote the
resulting set by \(E^+\).  A link \(e=(s(e),t(e))\in E^+\) carries
\(U_e\in K\), represented on \(V_R\), and \(U_{\bar e}=U_e^{-1}\).

For a differentiable real functional \(S(U)\), define its
left-trivialized link current \(J_e(S)\in\mathfrak k_R\) by
\begin{equation}
 \delta S
 =\sum_{e\in E^+}\langle J_e(S),\eta_e\rangle_R,
 \qquad
 \delta U_e=\eta_eU_e,
 \quad \eta_e\in\mathfrak k_R.                            \label{eq:newS7V6-current-definition}
\end{equation}
Faithfulness makes this current unique.  For a nonfaithful representation
the same definition is made in the quotient of the Lie algebra by the
kernel of the representation.

Let \(C=(e_1,\ldots,e_n)\) be a closed directed walk, with inverse links
inserted as \(U_{\bar e}=U_e^{-1}\), and let
\[
 H_C=U_{e_1}\cdots U_{e_n},
 \qquad
 S_C(U)=w_C\Re\tr_R(I-H_C),
 \qquad w_C\ge0.
\]
When the occurrence \(e_k\) is traversed in its displayed orientation, put
\begin{equation}
 L_{C,k}=U_{e_k}\cdots U_{e_n}U_{e_1}\cdots U_{e_{k-1}}.  \label{eq:newS7V6-cyclic-holonomy}
\end{equation}
It is the holonomy of the same closed walk based at \(s(e_k)\).  Let
\(\operatorname{skew}M=(M-M^*)/2\), and let
\(\operatorname{pr}_{\mathfrak k_R}\) be orthogonal projection with respect
to the Hilbert--Schmidt inner product.

\begin{lemma}[One-occurrence Wilson insertion]
\label{lem:newS7V6-one-insertion}
If only the displayed occurrence \(e_k\) is varied by
\(\delta U_{e_k}=\eta U_{e_k}\), then
\begin{equation}
 \delta S_C
 =\left\langle
 w_C\operatorname{pr}_{\mathfrak k_R}
       \operatorname{skew}L_{C,k},\eta
 \right\rangle_R.                                       \label{eq:newS7V6-one-insertion}
\end{equation}
For an occurrence with inverse orientation, its contribution to the
current of the chosen representative in \(E^+\) is obtained from
\begin{equation}
 \delta U_{\bar e}
 =-U_e^{-1}\eta_eU_e\,U_{\bar e}
 =-\operatorname{Ad}_{U_e^{-1}}(\eta_e)U_{\bar e}.        \label{eq:newS7V6-inverse-variation}
\end{equation}
Thus repeated and oppositely oriented occurrences are summed without
dropping an endpoint.
\end{lemma}

\begin{proof}
Write \(P=U_{e_1}\cdots U_{e_{k-1}}\) and
\(Q=U_{e_k}\cdots U_{e_n}\).  Cyclicity gives
\[
 \delta S_C=-w_C\Re\tr_R(P\eta Q)
            =-w_C\Re\tr_R(\eta QP).
\]
For anti-Hermitian \(\eta\), the Hermitian part of \(QP=L_{C,k}\) has
purely imaginary trace against \(\eta\), so
\[
 \Re\tr_R(\eta L_{C,k})
 =\Re\tr_R\{\eta\operatorname{skew}L_{C,k}\}.
\]
Projection onto \(\mathfrak k_R\) does not change the pairing with
\(\eta\in\mathfrak k_R\), and the sign in the definition of
\(\langle\cdot,\cdot\rangle_R\) proves
\eqref{eq:newS7V6-one-insertion}.  Differentiating
\(U_{\bar e}=U_e^{-1}\) proves
\eqref{eq:newS7V6-inverse-variation}.
\end{proof}

\subsection{The endpoint-complete lattice Ward identity}
\label{sec:newS7V6-lattice-Ward}

For a site field \(\xi(x)\in\mathfrak k_R\), the infinitesimal gauge
variation of a positively oriented link is
\begin{equation}
 \delta_\xi U_e=\xi(s(e))U_e-U_e\xi(t(e))
 =\{\xi(s(e))-\operatorname{Ad}_{U_e}\xi(t(e))\}U_e.     \label{eq:newS7V6-link-gauge-variation}
\end{equation}
Define the covariant lattice divergence of a link current by
\begin{equation}
 (\operatorname{div}_U J)(x)
 :=\sum_{\substack{e\in E^+\\s(e)=x}}J_e
 -\sum_{\substack{e\in E^+\\t(e)=x}}
       \operatorname{Ad}_{U_e^{-1}}J_e.                  \label{eq:newS7V6-lattice-divergence}
\end{equation}

\begin{theorem}[Exact closed-loop divergence identity]
\label{thm:newS7V6-loop-divergence}
For every finite real linear combination of closed Wilson actions
\(S=\sum_C S_C\), its current from
\eqref{eq:newS7V6-current-definition} satisfies
\begin{equation}
 (\operatorname{div}_U J(S))(x)=0
 \qquad\hbox{at every site }x.                            \label{eq:newS7V6-divergence-zero}
\end{equation}
The identity is exact for every lattice spacing, volume, link
configuration, and loop length.
\end{theorem}

\begin{proof}
A finite gauge transformation acts by
\(U_e\mapsto g(s(e))U_eg(t(e))^{-1}\).  The interior factors of a closed
walk cancel pairwise, leaving
\(H_C\mapsto g(x_0)H_Cg(x_0)^{-1}\).  Hence every \(S_C\), and therefore
\(S\), is gauge invariant.

Insert \eqref{eq:newS7V6-link-gauge-variation} into
\eqref{eq:newS7V6-current-definition}.  Adjoint invariance of the pairing
gives
\begin{align*}
 0=\delta_\xi S
 &=\sum_{e\in E^+}
   \langle J_e,\xi(s(e))-\operatorname{Ad}_{U_e}\xi(t(e))\rangle_R\\
 &=\sum_x\langle(\operatorname{div}_U J)(x),\xi(x)\rangle_R.
\end{align*}
The site parameters are arbitrary and the pairing is nondegenerate, which
proves \eqref{eq:newS7V6-divergence-zero}.
\end{proof}

\begin{corollary}[The two endpoint terms are individually necessary]
\label{cor:newS7V6-two-endpoints}
For a nonconstant site parameter, neither sum in
\eqref{eq:newS7V6-lattice-divergence} is gauge invariant by itself.  The
incoming term must be transported to the site \(x\) before it is combined
with the outgoing term.  Omitting either endpoint generally leaves a
nonzero open-walk insertion.
\end{corollary}

\begin{proof}
Equation \eqref{eq:newS7V6-link-gauge-variation} contains two independent
endpoint values.  The coefficient of \(\xi(t(e))\) is transported by
\(\operatorname{Ad}_{U_e^{-1}}\) when the pairing is moved to the target.
Only their sitewise sum is the coefficient of an arbitrary gauge
parameter.  A single term has one uncancelled endpoint and transforms as
an open transporter.
\end{proof}

\begin{assessment}[Relation to the V3 heat jets]
\label{ass:newS7V6-heat-jet-current}
Every traced monomial in the V3 sign-heat expansion is a finite weighted
sum of closed walks.  Applying \Cref{thm:newS7V6-loop-divergence} term by
term gives the local form of V5's identity
\(\Tr\{[X,A]A^{2n}\}=0\).  The present theorem adds the missing information:
it identifies the link insertion and the parallel transport required to
form its sitewise divergence.  Absolute convergence or a controlled
remainder is still required before an infinite heat series may be
differentiated term by term.
\end{assessment}

\subsection{Smooth plaquettes and the unique quadratic-curvature tensor}
\label{sec:newS7V6-plaquette-limit}

Work first on a compact flat four-torus, or in a compactly supported
coordinate chart.  Let \(A=A_\mu\,dx^\mu\) be a \(C^3\)
anti-Hermitian connection and let \(U_{x,\mu}(a)\) be its exact parallel
transport along a coordinate edge of length \(a\).  Denote the positively
oriented \(\mu\nu\)-plaquette holonomy by \(U_{\mu\nu}(x,a)\), and set
\[
 F_{\mu\nu}=\partial_\mu A_\nu-\partial_\nu A_\mu
             +[A_\mu,A_\nu].
\]

\begin{lemma}[Uniform small-plaquette expansion]
\label{lem:newS7V6-small-plaquette}
On every set on which the \(C^3\)-norm of \(A\) is bounded,
\begin{equation}
 U_{\mu\nu}(x,a)
 =I+\epsilon_{\mu\nu}a^2\rho_*(F_{\mu\nu}(x))+R_{\mu\nu}(x,a),
 \qquad
 \|R_{\mu\nu}(x,a)\|\le C_Aa^3,                          \label{eq:newS7V6-plaquette-expansion}
\end{equation}
where \(\epsilon_{\mu\nu}=\pm1\) depends only on the convention for
parallel transport and orientation.  Consequently
\begin{equation}
 \Re\tr_R(I-U_{\mu\nu}(x,a))
 ={a^4\over2}
   \|\rho_*(F_{\mu\nu}(x))\|_{\rm HS}^2+O(a^5),          \label{eq:newS7V6-plaquette-quadratic}
\end{equation}
uniformly on that set.
\end{lemma}

\begin{proof}
Taylor expansion of the edge transport ODE, multiplied in plaquette order,
cancels the four order-\(a\) edge terms.  The order-\(a^2\) terms are the two
derivatives and the commutator in \(F_{\mu\nu}\); the integral remainder in
the edge ODE is bounded by the stated \(C^3\) norm.  This proves
\eqref{eq:newS7V6-plaquette-expansion}.

For every unitary matrix \(H\), there is the exact identity
\begin{equation}
 \Re\tr_R(I-H)={1\over2}\|I-H\|_{\rm HS}^2.              \label{eq:newS7V6-unitary-square}
\end{equation}
Apply it to the plaquette.  Squaring
\(I-U_{\mu\nu}=-\epsilon_{\mu\nu}a^2\rho_*(F_{\mu\nu})+O(a^3)\)
gives \eqref{eq:newS7V6-plaquette-quadratic}.  In particular the orientation
sign disappears.
\end{proof}

Fix an invariant inner product
\(\langle\cdot,\cdot\rangle_0\) on each gauge-algebra factor.  Define the
quadratic representation index \(I_2(R)\) by
\begin{equation}
 -\Re\tr_R\{\rho_*(X)\rho_*(Y)\}
 =I_2(R)\langle X,Y\rangle_0.                             \label{eq:newS7V6-index-definition}
\end{equation}
On a product group this equation is applied factor by factor.  For an
Abelian factor, \(I_2(R)\) is the multiplicity-weighted sum of squared
charges.

\begin{theorem}[Continuum owner of a plaquette coefficient]
\label{thm:newS7V6-F2-owner}
Let
\begin{equation}
 S_a^{(R)}(A)=\kappa_R
 \sum_x\sum_{1\le\mu<\nu\le4}
 \Re\tr_R\{I-U_{\mu\nu}(x,a)\}.                          \label{eq:newS7V6-plaquette-action}
\end{equation}
For a periodic \(C^3\) connection on a fixed four-torus,
\begin{equation}
 S_a^{(R)}(A)
 \longrightarrow
 {\kappa_R I_2(R)\over2}
 \int\sum_{\mu<\nu}
 \langle F_{\mu\nu},F_{\mu\nu}\rangle_0\,d^4x.           \label{eq:newS7V6-F2-limit}
\end{equation}
The same statement holds locally on a bounded-geometry Riemannian chart
after inserting the cell volume and the inverse-metric contraction.  Thus
the only parity-even, dimension-four tensor supplied by this closed
plaquette coefficient is the Yang--Mills density
\(\sqrt g\,g^{\mu\rho}g^{\nu\sigma}
 \langle F_{\mu\nu},F_{\rho\sigma}\rangle_0\); the
representation dependence is entirely in \(I_2(R)\).
\end{theorem}

\begin{proof}
By \eqref{eq:newS7V6-plaquette-quadratic} and
\eqref{eq:newS7V6-index-definition}, each summand is
\[
 {a^4I_2(R)\over2}
 \langle F_{\mu\nu}(x),F_{\mu\nu}(x)\rangle_0+O(a^5).
\]
There are \(O(a^{-4})\) plaquettes in a fixed physical volume, so the sum of
the uniform remainders is \(O(a)\).  The leading sum is a Riemann sum and
converges to \eqref{eq:newS7V6-F2-limit}.  In a metric chart, use an
orthonormal frame at the base point and the smooth expansions of cell
volume and inverse metric.  Coordinate invariance then gives the displayed
metric contraction.
\end{proof}

\subsection{The continuum shadow of the lattice divergence}
\label{sec:newS7V6-continuum-current}

\begin{theorem}[Weak current limit and Noether identity]
\label{thm:newS7V6-current-limit}
Let \(A(t)=A+t\alpha\), where \(\alpha\) is a \(C^3\) compactly supported
Lie-algebra-valued one-form.  Under the hypotheses of
\Cref{thm:newS7V6-F2-owner}, the first variations converge:
\begin{align}
 {d\over dt}S_a^{(R)}(A(t))\bigg|_{t=0}
 &\longrightarrow
 \kappa_RI_2(R)
 \int\sum_{\mu<\nu}
 \langle F_{\mu\nu},D_\mu\alpha_\nu-D_\nu\alpha_\mu\rangle_0\,d^4x
                                                                    \label{eq:newS7V6-first-variation}\\
 &=-\kappa_RI_2(R)
 \int\sum_\nu\langle D_\mu F_{\mu\nu},\alpha_\nu\rangle_0\,d^4x. \label{eq:newS7V6-continuum-current}
\end{align}
Thus the rescaled lattice link currents converge weakly to
\begin{equation}
 j_\nu=-\kappa_RI_2(R)D_\mu F_{\mu\nu}.                  \label{eq:newS7V6-j-limit}
\end{equation}
This continuum current has the off-shell identity
\begin{equation}
 D_\nu j_\nu=0.                                           \label{eq:newS7V6-continuum-divergence}
\end{equation}
\end{theorem}

\begin{proof}
The plaquette expansion and its \(t\)-derivative are uniform for \(t\) in a
small compact interval because \(A(t)\) has a common \(C^3\) bound.  The
Riemann-sum proof of \Cref{thm:newS7V6-F2-owner} may therefore be
differentiated at \(t=0\).  Since
\(\delta F_{\mu\nu}=D_\mu\alpha_\nu-D_\nu\alpha_\mu\), this gives
\eqref{eq:newS7V6-first-variation}.  Covariant integration by parts, with no
boundary term, gives \eqref{eq:newS7V6-continuum-current}.  This proves weak
convergence of the currents by their pairing with arbitrary test
\(\alpha\).

Finally, antisymmetry of \(F\) gives
\begin{align*}
 D_\nu D_\mu F_{\mu\nu}
 &=\frac12(D_\nu D_\mu-D_\mu D_\nu)F_{\mu\nu}\\
 &=\frac12[F_{\nu\mu},F_{\mu\nu}]=0,
\end{align*}
where the last equality is summed over \(\mu,\nu\).  Equation
\eqref{eq:newS7V6-continuum-divergence} follows and is the continuum shadow
of \eqref{eq:newS7V6-divergence-zero}.
\end{proof}

\subsection{Standard Model representation ledger at the quadratic owner}
\label{sec:newS7V6-SM-index}

Use the all-left-handed one-generation multiplets from V4 and normalize the
fundamental quadratic index of \(SU(N)\) to \(1/2\).  Multiplicity from the
other gauge factors gives
\begin{equation}
 \begin{array}{c|c}
 \text{factor}&\displaystyle\sum_{\text{one generation}} I_2(R)\\ \hline
 SU(3)_c&2\cdot\frac12+\frac12+\frac12=2\\
 SU(2)_L&3\cdot\frac12+\frac12=2\\
 U(1)_Y&\displaystyle\sum_R d_3(R)d_2(R)Y_R^2={10\over3}.
 \end{array}                                               \label{eq:newS7V6-SM-quadratic-ledger}
\end{equation}
The first row receives weak multiplicity two from the quark doublet; the
second receives colour multiplicity three.  The Abelian row is exactly the
second hypercharge moment proved in V4.  A sterile left-handed
antineutrino, if added, has \(Y=0\) and changes none of the rows.

\begin{corollary}[Charge assembly occurs after the common tensor]
\label{cor:newS7V6-charge-after-tensor}
For any assembled short-time coefficient whose closed-walk quadratic term
has the plaquette form \eqref{eq:newS7V6-plaquette-action}, the geometric
owner is the common \(F^2\) tensor in
\eqref{eq:newS7V6-F2-limit}.  The one-generation fermion representation
weights are \(2,2,10/3\) for \(SU(3)_c,SU(2)_L,U(1)_Y\), respectively,
before generation number, spin, chirality, statistics, and regulator
constants are inserted.
\end{corollary}

\begin{proof}
Apply \Cref{thm:newS7V6-F2-owner} representationwise and sum the quadratic
indices using \eqref{eq:newS7V6-SM-quadratic-ledger}.  The factors listed at
the end are common heat-kernel or field-content coefficients and are not
part of the representation trace computed here.
\end{proof}

\begin{definition}[Closed-loop curvature card, CLCC]
\label{def:newS7V6-CLCC}
CLCC consists of:
\begin{enumerate}[label=\textup{(\roman*)},leftmargin=3.5em]
\item the cyclic insertion \eqref{eq:newS7V6-one-insertion}, including the
      inverse-orientation rule \eqref{eq:newS7V6-inverse-variation};
\item the exact site divergence \eqref{eq:newS7V6-divergence-zero};
\item the uniform plaquette expansion
      \eqref{eq:newS7V6-plaquette-quadratic};
\item the continuum \(F^2\) owner \eqref{eq:newS7V6-F2-limit};
\item weak convergence of the current and
      \(D_\nu j_\nu=0\); and
\item the Standard Model quadratic-index ledger
      \eqref{eq:newS7V6-SM-quadratic-ledger}.
\end{enumerate}
\end{definition}

\begin{programme}[V7 acquisition]
\label{prog:newS7V6-V7}
V6 identifies the lattice tensor and its representation weight but not the
coefficient with which the overlap--Yukawa determinant produces it.  V7
will move one step upstream to a continuum Dirac-square model: derive the
gauge part of the four-dimensional \(a_4\) heat coefficient from the
Laplace-type recursion, keep the spin and chirality traces explicit, and
then state exactly which comparison theorem would be required to transfer
that coefficient to the admissible overlap lattice with a uniform
remainder.  This separates the universal local target from the unresolved
lattice-to-continuum estimate.
\end{programme}

\begin{firewall}[Claim boundary after seventh-series V6]
\label{fire:newS7V6-boundary}
V6 proves an exact finite-lattice covariant-divergence identity for every
finite sum of closed Wilson loops, including repeated links, inverse
orientations, and both link endpoints.  It proves the uniform smooth
plaquette expansion, identifies its continuum Yang--Mills \(F^2\) owner,
derives the weak continuum current and its off-shell Noether identity, and
computes the one-generation Standard Model quadratic representation
weights.

V6 does not prove that the overlap--Yukawa heat coefficient equals a Wilson
plaquette sum with a cutoff-uniform remainder, compute its spin or chiral
constant, derive its parity-odd \(F\wedge F\) partner, control Higgs and
Riemann-curvature coefficients, construct a global determinant-line base,
remove the lattice cutoff, or construct the Standard Model--Einstein
continuum.
\end{firewall}

\begin{theorem}[Seventh-series V6 result]
\label{thm:newS7V6-result}
The first local closed-walk tensor required by RWHC has an exact
finite-cutoff current whose covariant divergence vanishes site by site.  In
a smooth-link scaling, the plaquette part converges to the unique
parity-even quadratic gauge-curvature density, and its current converges
weakly to \(-\kappa_RI_2(R)D_\mu F_{\mu\nu}\).  Standard Model charge
assembly enters only through the quadratic indices \(2,2,10/3\).  The
remaining local bottleneck is therefore the universal spin--chiral heat
coefficient and the uniform comparison from the overlap lattice to that
continuum target.
\end{theorem}

\begin{proof}
Combine \Cref{lem:newS7V6-one-insertion,
thm:newS7V6-loop-divergence,
lem:newS7V6-small-plaquette,
thm:newS7V6-F2-owner,
thm:newS7V6-current-limit,
cor:newS7V6-charge-after-tensor}.
\end{proof}

\paragraph{Parent-version record.}
V6 was formed from
\texttt{Renewal\_SM\_Einstein\_Continuum\_Research\_Notes\_V5.tex}.
The V5 parent had \(2\,266\) logical source lines, \(86\,825\) bytes, and
SHA--256
\[
 \texttt{57F0163FE00429D6E7D37DEDF60EF1188CEFF4D3BAF73F3CE36A728EEE2F91DA}.
\]
The next compulsory companion audit is V10.

% =============================================================================
% END APPENDED SEVENTH-SERIES V6 MATERIAL
% =============================================================================

% =============================================================================
% BEGIN APPENDED SEVENTH-SERIES V7 MATERIAL
% =============================================================================

\clearpage
\section{Seventh-series V7: the four-dimensional Dirac heat coefficient
and exact relative gauge subtraction}
\label{sec:v7v7}

V6 identified the closed-plaquette \(F^2\) tensor and its representation
index without assigning the spin, chirality, or proper-time coefficient
produced by a fermion determinant.  This version computes that universal
continuum coefficient.  The computation is deliberately made before any
claim that the overlap lattice has reached it.

\subsection{Laplace convention and the quadratic heat coefficient}
\label{sec:newS7V7-Laplace}

Let \(M\) be a closed oriented Riemannian four-manifold and let
\[
 L=\nabla^*\nabla+E
\]
be a Laplace-type operator on a Hermitian vector bundle.  Write
\(\Omega_{\mu\nu}=[\nabla_\mu,\nabla_\nu]\).  Our heat convention is
\begin{equation}
 \Tr(e^{-tL})
 \sim(4\pi t)^{-2}
 \sum_{j\ge0}t^jA_{2j}(L)
 \qquad(t\downarrow0).                                    \label{eq:newS7V7-heat-convention}
\end{equation}
Only the terms quadratic in a gauge curvature are needed below.

\begin{lemma}[Universal quadratic part of \(A_4\)]
\label{lem:newS7V7-universal-A4}
For the convention in \eqref{eq:newS7V7-heat-convention}, the part of the
local \(A_4\) density quadratic in \(E\) and in the connection curvature is
\begin{equation}
 \tr\left\{{1\over2}E^2
             +{1\over12}\Omega_{\mu\nu}\Omega_{\mu\nu}\right\}.       \label{eq:newS7V7-universal-A4}
\end{equation}
Terms omitted from \eqref{eq:newS7V7-universal-A4} contain Riemann
curvature, covariant derivatives, or the identity endomorphism.
\end{lemma}

\begin{proof}
The diagonal heat recursion in normal coordinates is local and covariant.
At weight four, its derivative-free quadratic endomorphisms are
\(c_EE^2\) and \(c_\Omega\Omega_{\mu\nu}\Omega_{\mu\nu}\), up to terms
involving Riemann curvature.  A constant commuting endomorphism factors
from the scalar heat kernel as \(e^{-tE}\).  Its second Taylor coefficient
is \(E^2/2\), so \(c_E=1/2\).

To obtain \(c_\Omega\), take a flat bundle with a constant Abelian unitary
curvature supported in one two-plane, \(\Omega_{12}=ib\) and
\(\Omega_{21}=-ib\).  The exact diagonal magnetic heat kernel differs from
the free one by
\[
 {bt\over\sinh(bt)}
 =1-{b^2t^2\over6}+O(t^4).
\]
Here
\(\Omega_{\mu\nu}\Omega_{\mu\nu}=-2b^2I\), so the displayed coefficient is
\((1/12)\Omega_{\mu\nu}\Omega_{\mu\nu}\).  Local covariance and
polarization fix the same coefficient for a general unitary curvature
after the fibre trace.  This proves \eqref{eq:newS7V7-universal-A4}.
\end{proof}

\begin{remark}[Compatibility with the earlier three-dimensional
calculation]
\label{rem:newS7V7-convention}
The archived V89 formula used the same operator convention:
\(A_2(L)=\int\tr(RI/6-E)\).  Thus no sign conversion is needed when the V7
coefficient is inserted into the existing proper-time ledger.
\end{remark}

\subsection{Lichnerowicz reduction and spin traces}
\label{sec:newS7V7-Dirac}

Let \(S\) be the rank-four complex spin bundle.  Choose Hermitian Euclidean
gamma matrices and set
\[
 \gamma^{\mu\nu}={1\over2}[\gamma^\mu,\gamma^\nu],
 \qquad
 \gamma_5=\gamma^1\gamma^2\gamma^3\gamma^4,
 \qquad
 P_\chi={1\over2}(I+\chi\gamma_5),\quad\chi=\pm1.
\]
Let \(R\) be a unitary gauge representation with anti-Hermitian curvature
\(F_{\mu\nu}\).  For the twisted Dirac square
\(L_R=\slashed D_R^{\,*}\slashed D_R\), the gauge parts of the
Lichnerowicz endomorphism and bundle curvature are
\begin{equation}
 E_F={1\over2}\gamma^{\mu\nu}\rho_*(F_{\mu\nu}),
 \qquad
 (\Omega_{\mu\nu})_F=I_S\otimes\rho_*(F_{\mu\nu}).         \label{eq:newS7V7-Lichnerowicz-gauge}
\end{equation}
Changing the overall sign convention for \(\slashed D\) changes the sign of
\(E_F\) but not the parity-even result below.

\begin{lemma}[Four-dimensional gamma traces]
\label{lem:newS7V7-gamma-traces}
With \(\varepsilon^{1234}=1\),
\begin{align}
 \tr_S I_S&=4,                                             \label{eq:newS7V7-spin-rank}\\
 \tr_S(\gamma^{\mu\nu}\gamma^{\rho\sigma})
 &=4(g^{\mu\sigma}g^{\nu\rho}
       -g^{\mu\rho}g^{\nu\sigma}),                        \label{eq:newS7V7-two-gamma}\\
 \tr_S(\gamma_5\gamma^{\mu\nu}\gamma^{\rho\sigma})
 &=4\varepsilon^{\mu\nu\rho\sigma}.                       \label{eq:newS7V7-gamma5-trace}
\end{align}
Consequently
\begin{align}
 \tr_{S\otimes R}(E_F^2)
 &=-2\,\tr_R\{\rho_*(F_{\mu\nu})\rho_*(F_{\mu\nu})\},      \label{eq:newS7V7-E-square}\\
 \tr_{S\otimes R}\{(\Omega_{\mu\nu})_F
                    (\Omega_{\mu\nu})_F\}
 &=4\,\tr_R\{\rho_*(F_{\mu\nu})\rho_*(F_{\mu\nu})\},       \label{eq:newS7V7-Omega-square}\\
 \tr_{S\otimes R}(\gamma_5E_F^2)
 &=\varepsilon^{\mu\nu\rho\sigma}
   \tr_R\{\rho_*(F_{\mu\nu})\rho_*(F_{\rho\sigma})\}.      \label{eq:newS7V7-gamma5-E-square}
\end{align}
\end{lemma}

\begin{proof}
The Clifford relations and cyclicity of the trace give
\eqref{eq:newS7V7-spin-rank}--\eqref{eq:newS7V7-two-gamma}.
The only completely antisymmetric rank-four tensor is a scalar multiple of
\(\varepsilon\); evaluation at \(1234\) gives the factor \(4\) in
\eqref{eq:newS7V7-gamma5-trace}.  Insert
\eqref{eq:newS7V7-Lichnerowicz-gauge}.  Antisymmetry of \(F_{\mu\nu}\)
contracts \eqref{eq:newS7V7-two-gamma} to
\(-8F_{\mu\nu}F_{\mu\nu}\); the prefactor \(1/4\) in \(E_F^2\) yields
\eqref{eq:newS7V7-E-square}.  The spin identity in the connection curvature
gives \eqref{eq:newS7V7-Omega-square}, and
\eqref{eq:newS7V7-gamma5-trace} gives
\eqref{eq:newS7V7-gamma5-E-square}.
\end{proof}

\begin{theorem}[Dirac and Weyl \(A_4\) gauge coefficients]
\label{thm:newS7V7-Dirac-A4}
The pure-gauge part of the integrated \(A_4\) coefficient of one complex
four-component Dirac square is
\begin{equation}
 A_4(L_R)\big|_{F^2}
 =-{2\over3}\int_M
 \tr_R\{\rho_*(F_{\mu\nu})\rho_*(F_{\mu\nu})\}\,d\mu_g.   \label{eq:newS7V7-Dirac-A4}
\end{equation}
If the massless square is restricted by \(P_\chi\), its gauge part is
\begin{align}
 A_4(P_\chi L_R)\big|_{\rm gauge}
 &=-{1\over3}\int_M
 \tr_R\{\rho_*(F_{\mu\nu})\rho_*(F_{\mu\nu})\}\,d\mu_g    \notag\\
 &\quad+{\chi\over4}\int_M
 \varepsilon^{\mu\nu\rho\sigma}
 \tr_R\{\rho_*(F_{\mu\nu})\rho_*(F_{\rho\sigma})\}\,d\mu_g.
                                                                    \label{eq:newS7V7-Weyl-A4}
\end{align}
The first line is parity even.  The second is the local index-density
component in these gamma and orientation conventions.
\end{theorem}

\begin{proof}
Insert \eqref{eq:newS7V7-E-square} and
\eqref{eq:newS7V7-Omega-square} into
\eqref{eq:newS7V7-universal-A4}:
\[
 {1\over2}(-2)+{1\over12}(4)=-{2\over3}.
\]
Cross terms between spin curvature and gauge curvature vanish because
\(\tr_S\gamma^{\mu\nu}=0\).  This proves
\eqref{eq:newS7V7-Dirac-A4}.  For \(P_\chi\), the parity-even spin trace is
one half of the Dirac trace.  The \(\gamma_5\) trace of the
\(\Omega^2\) term vanishes, while the \(E^2/2\) term contributes
\[
 {1\over2}\tr(P_\chi E_F^2)\big|_{\rm odd}
 ={\chi\over4}\tr(\gamma_5E_F^2).
\]
Use \eqref{eq:newS7V7-gamma5-E-square}.
\end{proof}

\begin{corollary}[Positive-norm form and the V6 owner]
\label{cor:newS7V7-positive-F2}
With the quadratic index convention of V6,
\eqref{eq:newS7V7-Dirac-A4} becomes
\begin{equation}
 A_4(L_R)\big|_{F^2}
 ={2I_2(R)\over3}
 \int_M\langle F_{\mu\nu},F_{\mu\nu}\rangle_0\,d\mu_g.    \label{eq:newS7V7-positive-F2}
\end{equation}
The parity-even Weyl coefficient is half of this.  It is exactly the
quadratic-curvature tensor identified by the Wilson-plaquette limit in V6.
\end{corollary}

\begin{proof}
Use
\(-\Re\tr_R\{\rho_*(X)\rho_*(Y)\}
=I_2(R)\langle X,Y\rangle_0\).
\end{proof}

\subsection{Yukawa zero-order terms do not change the pure \(F^2\)
coefficient}
\label{sec:newS7V7-Yukawa}

\begin{proposition}[Stability of the pure gauge owner]
\label{prop:newS7V7-Yukawa-stability}
Let \(\mathcal D_\Phi=\slashed D+\gamma_5\Phi\), where \(\Phi\) is a smooth
gauge-covariant internal endomorphism commuting with Clifford multiplication.
In the square, the new endomorphism terms are linear combinations of
\(\Phi^2\), \(\gamma^\mu\nabla_\mu\Phi\), and their chirality-decorated
variants.  None changes the coefficient multiplying the monomial \(F^2\)
in \eqref{eq:newS7V7-Dirac-A4}.  They instead generate the separate local
owners
\[
 (\nabla\Phi)^2,\qquad \Phi^4,\qquad R\Phi^2,
\]
and total derivatives, with representation-compatible contractions.
\end{proposition}

\begin{proof}
The connection-curvature term in \eqref{eq:newS7V7-universal-A4} is
unchanged by a zero-order perturbation.  In \(E^2\), a cross term between
\(E_F=(1/2)\gamma^{\mu\nu}F_{\mu\nu}\) and \(\Phi^2\) has spin trace
proportional to \(\tr_S\gamma^{\mu\nu}=0\).  A cross term with
\(\gamma^\rho\nabla_\rho\Phi\) contains three gamma matrices and also has
zero ordinary spin trace.  Thus no \(\Phi\)-dependent cross term has the
pure \(F^2\) tensor.  Squaring the remaining pieces and including the
standard curvature and derivative terms yields only the displayed owner
types at weight four.
\end{proof}

\begin{assessment}[Scope of the parity-odd term]
\label{ass:newS7V7-odd-scope}
The second line of \eqref{eq:newS7V7-Weyl-A4} is an index density for a
chiral heat trace.  It is not by itself the four-dimensional gauge anomaly.
An infinitesimal gauge variation inserts an additional gauge generator;
the resulting descent coefficient contains the cubic and mixed charge
traces checked in V4.  The phase of a chiral determinant also requires the
eta and determinant-line data that the present modulus calculation does
not supply.
\end{assessment}

\subsection{Exact consequence for a relative determinant}
\label{sec:newS7V7-relative}

Let \(L=L_R(A,\Phi)\) and \(L_0=L_R(A_0,\Phi_0)\) have the same principal
symbol, spin convention, and gauge representation.  Their relative heat
trace is
\[
 H_{\rm rel}(t)=\Tr(e^{-tL}-e^{-tL_0}).
\]

\begin{theorem}[Relative \(F^2\) coefficient]
\label{thm:newS7V7-relative-F2}
The parity-even pure-gauge part of the relative Dirac coefficient is
\begin{equation}
 A_4(L)-A_4(L_0)\big|_{F^2}
 =-{2\over3}\int_M\tr_R
 \{\rho_*(F_{\mu\nu})\rho_*(F_{\mu\nu})
  -\rho_*(F^0_{\mu\nu})\rho_*(F^0_{\mu\nu})\}\,d\mu_g.    \label{eq:newS7V7-relative-F2}
\end{equation}
In particular, if the full and base blocks use the same gauge connection
\(A_0=A\), their universal pure \(F^2\) logarithmic coefficients cancel
exactly, even when \(\Phi_0\ne\Phi\).
\end{theorem}

\begin{proof}
Heat coefficients are additive under the difference of traces.  Apply
\eqref{eq:newS7V7-Dirac-A4} to \(L\) and \(L_0\).
\Cref{prop:newS7V7-Yukawa-stability} shows that changing only the zero-order
Yukawa--Higgs field cannot alter the pure \(F^2\) coefficient.
\end{proof}

\begin{corollary}[Gauge normalization must have an explicit owner]
\label{cor:newS7V7-gauge-normalization}
For a same-\(A\) covariant base, a relative fermion determinant cannot by
itself fix the absolute coefficient of the physical Yang--Mills action.
That coefficient must be supplied by an owned bare term or local
counterterm and fixed by a renormalization condition.  Changing the base
changes this local normalization but does not change the finite relative
Ward identity of V5.
\end{corollary}

\begin{proof}
The logarithmic coefficient from the relative heat trace is zero by
\Cref{thm:newS7V7-relative-F2}.  Adding a gauge-invariant local multiple of
\(\int\langle F,F\rangle\) preserves the exact Ward identity and changes
the absolute gauge normalization.  A normalization condition is therefore
needed to choose that multiple.
\end{proof}

\subsection{One-generation Standard Model coefficient ledger}
\label{sec:newS7V7-SM-ledger}

\begin{theorem}[Parity-even Weyl heat ledger]
\label{thm:newS7V7-SM-ledger}
For one all-left-handed Standard Model generation, the positive-norm
parity-even \(A_4\) gauge coefficients, before the common factor
\((4\pi)^{-2}\) in the heat trace, are
\begin{equation}
 {2\over3}\quad\hbox{for }SU(3)_c,\qquad
 {2\over3}\quad\hbox{for }SU(2)_L,\qquad
 {10\over9}\quad\hbox{for }U(1)_Y.                        \label{eq:newS7V7-SM-ledger}
\end{equation}
Three identical generations multiply these numbers by three.  A sterile
zero-hypercharge antineutrino changes none of them.
\end{theorem}

\begin{proof}
One Weyl coefficient is \(I_2(R)/3\) by
\Cref{cor:newS7V7-positive-F2}.  Multiply the V6 one-generation quadratic
indices \(2,2,10/3\) by \(1/3\).
\end{proof}

\begin{assessment}[Determinant exponents and signs]
\label{ass:newS7V7-determinant-sign}
Equation \eqref{eq:newS7V7-SM-ledger} is a heat coefficient.  Its
contribution to an effective action also carries the determinant exponent,
the fermionic sign, the factor \(1/2\) used when a modulus is written as a
Dirac square, and the chosen proper-time subtraction convention.  These
factors must be applied once at the assembled determinant level; they are
not silently included in \eqref{eq:newS7V7-SM-ledger}.
\end{assessment}

\subsection{The exact lattice-to-continuum estimate still required}
\label{sec:newS7V7-transfer}

Let \(H_a^{\rm rel}(t)\) denote the physical relative overlap--Yukawa heat
trace on a fixed four-volume, and let \(C_{0,a},C_{2,a},C_{4,a}\) be its
owned local coefficient functionals.  The continuum target just computed
is the gauge part of \(C_4\).

\begin{definition}[Mesoscopic \(A_4\) transfer card, MATC]
\label{def:newS7V7-MATC}
MATC consists of a scale \(t_0>0\) and functions
\(\varepsilon(a)\downarrow0\), with
\(a^2/\varepsilon(a)\to0\), such that on each bounded smooth AOP family:
\begin{enumerate}[label=\textup{(\roman*)},leftmargin=3.5em]
\item for \(\varepsilon(a)\le t\le t_0\), the relative heat trace has the
      four-dimensional expansion
      \[
       H_a^{\rm rel}(t)
       =(4\pi t)^{-2}
       \{C_{0,a}+tC_{2,a}+t^2C_{4,a}+r_a(t)\};
      \]
\item \(C_{4,a}\) converges to the continuum coefficient of
      \Cref{thm:newS7V7-relative-F2}, together with its Higgs and metric
      partners;
\item the subtracted remainder is logarithmically integrable uniformly:
      \[
       \lim_{\tau\downarrow0}\limsup_{a\downarrow0}
       \int_{\varepsilon(a)}^\tau {|r_a(t)|\over t^3}\,dt=0;
      \]
\item the microscopic interval \(0<t<\varepsilon(a)\) is exhausted by the
      declared local counterterms; and
\item the large-time interval is controlled by V5's incomplete-gamma
      estimate.
\end{enumerate}
\end{definition}

\begin{theorem}[What MATC would close]
\label{thm:newS7V7-MATC}
If MATC holds for the complete relative overlap--Yukawa block, then its
renormalized proper-time logarithm has the continuum logarithmic gauge
coefficient \eqref{eq:newS7V7-relative-F2}.  For a same-\(A\) base that
coefficient vanishes, and the chosen owned Yang--Mills normalization is the
only surviving absolute \(F^2\) coefficient.
\end{theorem}

\begin{proof}
Subtract the \(C_0\), \(C_2\), and declared \(C_4\) local terms in the
proper-time integral.  Row (iii) permits the limits \(a\downarrow0\) and
\(t\downarrow0\) on the mesoscopic remainder.  Row (iv) accounts for the
remaining microscopic contribution by local owners, and row (v) controls
the opposite tail.  Row (ii) identifies the logarithmic coefficient.
The same-\(A\) statement is \Cref{thm:newS7V7-relative-F2}.
\end{proof}

\begin{programme}[V8 acquisition]
\label{prog:newS7V7-V8}
V8 will attack MATC upstream at the free translation-invariant overlap
kernel.  It will derive the Brillouin-zone symbol, isolate the physical
small-momentum branch from the Wilson doublers, and prove a quantitative
mesoscopic comparison for the free heat trace.  Only after that baseline is
closed will gauge insertions be added by Duhamel expansion.  This avoids
assuming the uniform interacting remainder whose absence is now the main
local bottleneck.
\end{programme}

\begin{firewall}[Claim boundary after seventh-series V7]
\label{fire:newS7V7-boundary}
V7 derives the universal four-dimensional Dirac and Weyl \(A_4\) gauge
coefficients, including their spin, chirality, and representation traces.
It proves that zero-order Yukawa--Higgs terms do not change the pure
\(F^2\) coefficient, computes the one-generation Standard Model
parity-even heat ledger, and proves exact cancellation of that universal
coefficient in a relative determinant whose base uses the same gauge
connection.

V7 does not prove MATC for the overlap lattice, compute the complete Higgs,
gravity, gauge, ghost, and fermion supertrace, fix a physical running
coupling, construct the chiral determinant phase, prove global
determinant-line triviality, remove the cutoff, or construct the Standard
Model--Einstein continuum.
\end{firewall}

\begin{theorem}[Seventh-series V7 result]
\label{thm:newS7V7-result}
The continuum target of the V6 plaquette owner is fixed: one Dirac square
has parity-even gauge coefficient
\(-{2\over3}\tr_R(F_{\mu\nu}F_{\mu\nu})\), and one Weyl block has half that
coefficient, with the separate local index density shown in
\eqref{eq:newS7V7-Weyl-A4}.  For one Standard Model generation the
positive-norm Weyl weights are \(2/3,2/3,10/9\).  A same-connection
relative base cancels the universal \(F^2\) logarithm exactly.  The next
unresolved step is the mesoscopic overlap-to-Dirac comparison encoded by
MATC.
\end{theorem}

\begin{proof}
Use \Cref{lem:newS7V7-universal-A4,
lem:newS7V7-gamma-traces,
thm:newS7V7-Dirac-A4,
prop:newS7V7-Yukawa-stability,
thm:newS7V7-relative-F2,
thm:newS7V7-SM-ledger,
thm:newS7V7-MATC}.
\end{proof}

\paragraph{Parent-version record.}
V7 was formed from
\texttt{Renewal\_SM\_Einstein\_Continuum\_Research\_Notes\_V6.tex}.
The V6 parent had \(2\,721\) logical source lines, \(105\,637\) bytes, and
SHA--256
\[
 \texttt{449E5741AB52974204FADEA3036A9B306E41C5112830C1E02A7F2C704B55FE1E}.
\]
The next compulsory companion audit is V10.

% =============================================================================
% END APPENDED SEVENTH-SERIES V7 MATERIAL
% =============================================================================

% =============================================================================
% BEGIN APPENDED SEVENTH-SERIES V8 MATERIAL
% =============================================================================

\clearpage
\section{Seventh-series V8: free overlap dispersion and a quantitative
mesoscopic heat comparison}
\label{sec:v7v8}

V7 computed the continuum coefficient that an interacting overlap heat
trace must approach.  The first analytic obstruction in MATC is already
present at zero gauge and Higgs field: the lattice trace contains the full
Brillouin zone, while the continuum heat trace contains an unbounded
momentum lattice.  This version proves a uniform comparison between those
two traces and shows exactly how slowly the lower proper-time endpoint may
shrink.

\subsection{Exact free overlap symbol}
\label{sec:newS7V8-symbol}

Take a periodic flat four-torus of side length \(\ell\), lattice spacing
\(a=\ell/N\), and Wilson parameter and overlap mass parameter
\(r=\rho=1\).  For a lattice momentum
\[
 p={2\pi n\over\ell},\qquad
 \theta=ap\in[-\pi,\pi)^4,
\]
put
\begin{equation}
 s_\mu(\theta)=\sin\theta_\mu,\qquad
 W(\theta)=\sum_{\mu=1}^4(1-\cos\theta_\mu),\qquad
 z(\theta)=W(\theta)-1,                                   \label{eq:newS7V8-sWz}
\end{equation}
and
\begin{equation}
 \omega(\theta)=\{\,|s(\theta)|^2+z(\theta)^2\,\}^{1/2}.  \label{eq:newS7V8-omega}
\end{equation}
The free Wilson block and its overlap operator are
\begin{align}
 A_W(\theta)
 &={1\over a}\{i\gamma^\mu s_\mu(\theta)+z(\theta)\},      \label{eq:newS7V8-AW}\\
 D_{\rm ov}(\theta)
 &={1\over a}\left\{I+{i\gamma^\mu s_\mu(\theta)+z(\theta)
                         \over\omega(\theta)}\right\}.     \label{eq:newS7V8-Dov}
\end{align}
At \(\theta=0\), the quotient in \eqref{eq:newS7V8-Dov} is understood by
continuity as \(-I\).

\begin{lemma}[Exact scalar overlap square]
\label{lem:newS7V8-square}
For every Brillouin momentum,
\begin{equation}
 D_{\rm ov}(\theta)^*D_{\rm ov}(\theta)
 =\lambda_a(\theta)I_S,\qquad
 \lambda_a(\theta)
 ={2\over a^2}\left(1+{z(\theta)\over\omega(\theta)}\right).           \label{eq:newS7V8-lambda}
\end{equation}
The only zero of \(\lambda_a\) is \(\theta=0\).
\end{lemma}

\begin{proof}
Hermiticity of the Euclidean gamma matrices gives
\[
 \{i\gamma\cdot s+z\}^*\{i\gamma\cdot s+z\}
 =(|s|^2+z^2)I_S=\omega^2I_S.
\]
Multiplication of \eqref{eq:newS7V8-Dov} by its adjoint then yields
\[
 a^{-2}\{1+2z/\omega+(|s|^2+z^2)/\omega^2\}I_S,
\]
which is \eqref{eq:newS7V8-lambda}.  It can vanish only if
\(s=0\) and \(z<0\).  The zeros of \(s\) are the sixteen corners with
\(\theta_\mu\in\{0,\pi\}\).  If \(k\) components equal \(\pi\), then
\(z=2k-1\), which is negative only for \(k=0\).  Hence the sole zero is the
physical corner.
\end{proof}

\begin{corollary}[Exact doubler mass]
\label{cor:newS7V8-doublers}
At every nonzero Brillouin corner,
\begin{equation}
 D_{\rm ov}(\theta)={2\over a}I_S,\qquad
 \lambda_a(\theta)={4\over a^2}.                          \label{eq:newS7V8-doubler-mass}
\end{equation}
Thus no doubler supplies a second low-energy heat branch.
\end{corollary}

\begin{proof}
At a corner with \(k\ge1\), \(s=0\), \(z=2k-1>0\), and
\(\omega=z\).  Substitute in \eqref{eq:newS7V8-Dov} and
\eqref{eq:newS7V8-lambda}.
\end{proof}

\subsection{Global coercivity and the physical Taylor branch}
\label{sec:newS7V8-bounds}

Let \(|\theta|\) denote Euclidean distance from the origin in the principal
Brillouin cube.

\begin{lemma}[Uniform global dispersion comparison]
\label{lem:newS7V8-global-comparison}
There are numerical constants \(0<c_0\le C_0<\infty\), independent of
\(a,N,\ell\), such that
\begin{equation}
 c_0{|\,\theta\,|^2\over a^2}
 \le\lambda_a(\theta)
 \le C_0{|\,\theta\,|^2\over a^2}
 \qquad(\theta\in[-\pi,\pi)^4).                           \label{eq:newS7V8-global-comparison}
\end{equation}
\end{lemma}

\begin{proof}
The function
\[
 q(\theta)
 ={2(1+z(\theta)/\omega(\theta))\over|\theta|^2}
\]
is continuous away from zero.  By
\Cref{lem:newS7V8-square}, it is strictly positive there.  The Taylor
calculation below shows that it extends continuously to \(q(0)=1\).
Compactness of the closed Brillouin cube then gives positive finite minimum
and maximum values, which are \(c_0,C_0\).
\end{proof}

\begin{lemma}[Fourth-order physical-branch error]
\label{lem:newS7V8-low-expansion}
There are \(\delta>0\) and \(C_1<\infty\) such that, whenever
\(|\theta|\le\delta\),
\begin{equation}
 \left|a^2\lambda_a(\theta)-|\theta|^2\right|
 \le C_1|\theta|^4.                                      \label{eq:newS7V8-low-expansion}
\end{equation}
Equivalently, for \(|ap|\le\delta\),
\begin{equation}
 |\lambda_a(ap)-|p|^2|\le C_1a^2|p|^4.                   \label{eq:newS7V8-low-p}
\end{equation}
\end{lemma}

\begin{proof}
Taylor expansion gives
\begin{align*}
 |s(\theta)|^2
 &=|\theta|^2-{1\over3}\sum_\mu\theta_\mu^4+O(|\theta|^6),\\
 z(\theta)
 &=-1+{1\over2}|\theta|^2
      -{1\over24}\sum_\mu\theta_\mu^4+O(|\theta|^6).
\end{align*}
Hence
\[
 \omega(\theta)^2
 =1+{1\over4}\left\{|\theta|^4-\sum_\mu\theta_\mu^4\right\}
   +O(|\theta|^6),
\]
so \(\omega=1+O(|\theta|^4)\).  Therefore
\[
 2\{1+z/\omega\}=|\theta|^2+O(|\theta|^4).
\]
Taylor's theorem on a sufficiently small closed ball makes the remainder
uniform and proves the two statements.
\end{proof}

\subsection{Finite-volume Gaussian sums}
\label{sec:newS7V8-Gaussian}

Let
\[
 \mathcal P={2\pi\over\ell}\mathbb Z^4,\qquad
 \mathcal P_a=\mathcal P\cap[-\pi/a,\pi/a)^4.
\]
The scalar continuum torus heat trace is
\[
 H_{\rm c}(t)=\sum_{p\in\mathcal P}e^{-t|p|^2},
\]
and the scalar overlap heat trace is
\[
 H_a(t)=\sum_{p\in\mathcal P_a}e^{-t\lambda_a(ap)}.
\]

\begin{lemma}[Gaussian moment bounds]
\label{lem:newS7V8-Gaussian}
For every integer \(m\ge0\) and \(0<t\le\ell^2\), there is a numerical
constant \(C_m\) such that
\begin{equation}
 \sum_{p\in\mathcal P}|p|^{2m}e^{-c_0t|p|^2}
 \le C_m\ell^4t^{-2-m}.                                  \label{eq:newS7V8-Gaussian}
\end{equation}
\end{lemma}

\begin{proof}
Partition momentum space into cubes of side \(2\pi/\ell\) centered at
points of \(\mathcal P\).  On each cube, a Gaussian with a slightly smaller
exponent controls the value at the center up to a constant depending only
on \(m,c_0\).  Summing and comparing with the Euclidean integral gives
\[
 C\left({\ell\over2\pi}\right)^4
 \int_{\mathbb R^4}|p|^{2m}e^{-(c_0/2)t|p|^2}\,d^4p
 =C_m\ell^4t^{-2-m}.
\]
The finitely many cubes near the origin are absorbed into the same bound
because \(t\le\ell^2\).
\end{proof}

\begin{theorem}[Quantitative free overlap heat comparison]
\label{thm:newS7V8-heat-comparison}
There are constants \(C,c>0\) such that for
\(0<a\le\ell/4\) and \(a^2\le t\le\ell^2\),
\begin{equation}
 |H_a(t)-H_{\rm c}(t)|
 \le C\ell^4t^{-2}
 \left\{
 {a^2\over t}
 +\left({t\over a^2}\right)^2e^{-ct/a^2}
 \right\}.                                                \label{eq:newS7V8-heat-comparison}
\end{equation}
For spin rank four and internal rank \(r_R\), both sides are multiplied by
\(4r_R\).
\end{theorem}

\begin{proof}
Split both momentum sums at \(|ap|=\delta\), with \(\delta\) from
\Cref{lem:newS7V8-low-expansion}.  On the low region, the mean-value
theorem, \eqref{eq:newS7V8-global-comparison}, and
\eqref{eq:newS7V8-low-p} give
\[
 |e^{-t\lambda_a(ap)}-e^{-t|p|^2}|
 \le Cta^2|p|^4e^{-c_0t|p|^2}.
\]
Summation and \Cref{lem:newS7V8-Gaussian} with \(m=2\) bound this by
\(C\ell^4a^2t^{-3}\).

On the complementary lattice region,
\eqref{eq:newS7V8-global-comparison} gives
\(\lambda_a\ge c/a^2\).  There are at most \(C\ell^4a^{-4}\) lattice
momenta, so this tail is bounded by
\[
 C\ell^4a^{-4}e^{-ct/a^2}
 =C\ell^4t^{-2}(t/a^2)^2e^{-ct/a^2}.
\]
The missing continuum momenta outside \(\mathcal P_a\), and the continuum
part with \(|ap|\ge\delta\), obey the same bound after comparison with a
Gaussian integral.  Adding the low and high estimates proves
\eqref{eq:newS7V8-heat-comparison}.
\end{proof}

\begin{corollary}[No hidden Brillouin multiplicity]
\label{cor:newS7V8-no-hidden-rank}
For every fixed \(t>0\),
\begin{equation}
 \lim_{a\downarrow0}H_a(t)=H_{\rm c}(t).                  \label{eq:newS7V8-fixed-t}
\end{equation}
As \(t\downarrow0\),
\begin{equation}
 4r_RH_{\rm c}(t)
 ={4r_R\ell^4\over(4\pi t)^2}
 \{1+O(e^{-\ell^2/(Ct)})\}.                              \label{eq:newS7V8-continuum-leading}
\end{equation}
Thus the free overlap branch has exactly the continuum spin--internal rank;
the fifteen nonphysical corners do not multiply \(A_0\).
\end{corollary}

\begin{proof}
The first claim follows from
\eqref{eq:newS7V8-heat-comparison}.  The second is the Poisson summation
formula for the flat torus heat kernel.  The doubler statement also follows
directly from \eqref{eq:newS7V8-doubler-mass}.
\end{proof}

\subsection{An integrable shrinking proper-time window}
\label{sec:newS7V8-window}

Fix \(q\) with \(0<q<2/3\) and define the dimensionally consistent
mesoscopic lower endpoint
\begin{equation}
 \varepsilon_q(a)
 :=\ell^2\left({a\over\ell}\right)^q.                     \label{eq:newS7V8-epsilon}
\end{equation}
Then \(\varepsilon_q(a)\downarrow0\) and
\(a^2/\varepsilon_q(a)\to0\).

\begin{theorem}[Free proper-time transfer]
\label{thm:newS7V8-proper-time}
For every fixed \(t_0\le\ell^2\),
\begin{equation}
 \lim_{a\downarrow0}
 \int_{\varepsilon_q(a)}^{t_0}
 {\,|H_a(t)-H_{\rm c}(t)|\over t}\,dt=0.                 \label{eq:newS7V8-proper-time}
\end{equation}
The same is true after multiplication by the spin--internal rank.
\end{theorem}

\begin{proof}
The first term in \eqref{eq:newS7V8-heat-comparison} contributes at most
\[
 C\ell^4a^2
 \int_{\varepsilon_q(a)}^{t_0}t^{-4}\,dt
 \le C{\ell^4a^2\over\varepsilon_q(a)^3}
 =C\left({a\over\ell}\right)^{2-3q},
\]
which tends to zero precisely for \(q<2/3\).  For the second term, put
\(u=t/a^2\).  Its integral is bounded by a polynomial in
\(\varepsilon_q(a)/a^2\) times
\(\exp\{-c\varepsilon_q(a)/a^2\}\), and hence also tends to zero.
\end{proof}

\begin{corollary}[Free rows of MATC]
\label{cor:newS7V8-free-MATC}
At zero gauge and Higgs field, the free overlap square satisfies the
mesoscopic comparison and logarithmic integrability rows of MATC with
\(\varepsilon=\varepsilon_q\), \(0<q<2/3\).  Its only local coefficient on
the flat torus is the continuum rank coefficient in
\eqref{eq:newS7V8-continuum-leading}; all doubler contributions are
exponentially confined to the microscopic interval.
\end{corollary}

\begin{proof}
Use \Cref{thm:newS7V8-heat-comparison,
thm:newS7V8-proper-time,cor:newS7V8-no-hidden-rank}.  The high-momentum
term in \eqref{eq:newS7V8-heat-comparison} is exponentially small whenever
\(t/a^2\to\infty\), which holds throughout the chosen mesoscopic interval.
\end{proof}

\subsection{Why the interacting coefficient is still open}
\label{sec:newS7V8-interacting}

\begin{assessment}[The free theorem does not produce \(F^2\)]
\label{ass:newS7V8-no-F2}
At \(A=0\), the gauge curvature vanishes, so
\Cref{thm:newS7V8-proper-time} tests the principal symbol, the physical
rank, and doubler decoupling but cannot determine the quadratic gauge
coefficient.  To recover V7's \(A_4\) term one needs two gauge insertions,
or one curvature-level closed-loop insertion, with estimates uniform in
the mesoscopic window.  Pointwise convergence of
\(D_{\rm ov}(p)\) at fixed \(p\) is insufficient because the relevant
momenta grow like \(t^{-1/2}\).
\end{assessment}

\begin{definition}[Free overlap baseline card, FOBC]
\label{def:newS7V8-FOBC}
FOBC consists of:
\begin{enumerate}[label=\textup{(\roman*)},leftmargin=3.5em]
\item the exact scalar dispersion \eqref{eq:newS7V8-lambda};
\item one and only one zero momentum branch;
\item the doubler mass \eqref{eq:newS7V8-doubler-mass};
\item the global coercive comparison
      \eqref{eq:newS7V8-global-comparison};
\item the fourth-order physical-branch error
      \eqref{eq:newS7V8-low-p}; and
\item the integrated mesoscopic comparison
      \eqref{eq:newS7V8-proper-time}.
\end{enumerate}
\end{definition}

\begin{programme}[V9 acquisition]
\label{prog:newS7V8-V9}
V9 will take the first and second Fr\'echet derivatives of the finite
overlap heat trace about the free connection.  It will use the exact
Duhamel simplex, prove vanishing of the one-point term by the lattice Ward
identity, and isolate the transverse two-point polarization tensor.  The
target is a uniform momentum-space bound on the mesoscopic window, stated
separately from the extraction of the V7 coefficient.  The next compulsory
YM/NS audit remains V10.
\end{programme}

\begin{firewall}[Claim boundary after seventh-series V8]
\label{fire:newS7V8-boundary}
V8 derives the exact free overlap dispersion, proves that the physical
corner is its only zero, proves an exact \(2/a\) doubler mass, establishes
global and fourth-order dispersion bounds, and obtains a quantitative
finite-volume heat comparison.  For every \(0<q<2/3\), the free lattice and
continuum proper-time traces agree after integration over the shrinking
mesoscopic interval \([\varepsilon_q(a),t_0]\).

V8 does not treat a nonzero gauge or Higgs field, derive the \(F^2\)
polarization coefficient, control a chiral determinant phase, prove the
microscopic counterterm classification, establish the full MATC card,
remove the cutoff in the interacting theory, or construct the Standard
Model--Einstein continuum.
\end{firewall}

\begin{theorem}[Seventh-series V8 result]
\label{thm:newS7V8-result}
The free overlap regulator has precisely one continuum Dirac branch and no
hidden doubler multiplicity.  Its heat trace differs from the continuum
torus trace by a relative \(O(a^2/t)\) physical-branch error plus an
exponentially small Brillouin tail.  Choosing
\(\varepsilon_q(a)=\ell^2(a/\ell)^q\) with \(0<q<2/3\) makes that error
integrable in the proper-time logarithm.  The remaining local task is the
uniform two-insertion gauge polarization about this free baseline.
\end{theorem}

\begin{proof}
Combine \Cref{lem:newS7V8-square,
cor:newS7V8-doublers,
lem:newS7V8-global-comparison,
lem:newS7V8-low-expansion,
thm:newS7V8-heat-comparison,
thm:newS7V8-proper-time}.
\end{proof}

\paragraph{Parent-version record.}
V8 was formed from
\texttt{Renewal\_SM\_Einstein\_Continuum\_Research\_Notes\_V7.tex}.
The V7 parent had \(3\,163\) logical source lines, \(124\,266\) bytes, and
SHA--256
\[
 \texttt{21B1E395DEB9FE565EC19D24A8FC54C6B8CB787CC3D8D4CC8DEDB9C40F720179}.
\]
The next compulsory companion audit is V10.

% =============================================================================
% END APPENDED SEVENTH-SERIES V8 MATERIAL
% =============================================================================

% =============================================================================
% BEGIN APPENDED SEVENTH-SERIES V9 MATERIAL
% =============================================================================

\clearpage
\section{Seventh-series V9: exact overlap polarization, lattice
transversality, and the failed termwise bound}
\label{sec:v7v9}

V8 closes the free principal-symbol and doubler part of MATC.  The first
coefficient capable of seeing gauge curvature is the Hessian of the heat
trace with respect to the link connection.  This version derives that
Hessian exactly and imposes the Ward identity before any norm estimate is
taken.

\subsection{Analytic gauge variations on the AOP chart}
\label{sec:newS7V9-variations}

Let \(D_a(A)\) be the overlap operator on a finite periodic lattice, written
in exponential link coordinates near the trivial connection, and put
\[
 L_a(A)=D_a(A)^*D_a(A),\qquad
 \mathcal H_{a,t}(A)=\Tr(e^{-tL_a(A)}).
\]
The AOP Wilson gap makes these maps analytic on a common small link chart.
For gauge one-forms \(\alpha,\beta\), define at \(A=0\)
\begin{align}
 \dot D_\alpha&=D D_a(0)[\alpha],
 &\ddot D_{\alpha\beta}&=D^2D_a(0)[\alpha,\beta],          \label{eq:newS7V9-D-derivatives}\\
 V_\alpha&=D L_a(0)[\alpha],
 &W_{\alpha\beta}&=D^2L_a(0)[\alpha,\beta].               \label{eq:newS7V9-L-derivatives}
\end{align}

\begin{lemma}[Derivative ledger for the overlap square]
\label{lem:newS7V9-square-derivatives}
Writing \(D_0=D_a(0)\),
\begin{align}
 V_\alpha
 &=D_0^*\dot D_\alpha+\dot D_\alpha^*D_0,                 \label{eq:newS7V9-V}\\
 W_{\alpha\beta}
 &=\dot D_\alpha^*\dot D_\beta
   +\dot D_\beta^*\dot D_\alpha
   +D_0^*\ddot D_{\alpha\beta}
   +\ddot D_{\alpha\beta}^*D_0.                           \label{eq:newS7V9-W}
\end{align}
Both are self-adjoint when \(\alpha=\beta\), and
\(W_{\alpha\beta}=W_{\beta\alpha}\).
\end{lemma}

\begin{proof}
Differentiate the product \(D_a(A)^*D_a(A)\) once and twice.  Adjoint and
Fr\'echet differentiation commute in finite dimension.
\end{proof}

\subsection{The exact Duhamel Hessian}
\label{sec:newS7V9-Duhamel}

\begin{theorem}[Finite-cutoff heat polarization formula]
\label{thm:newS7V9-polarization}
Let \(L_0=L_a(0)\).  The first and mixed second variations of the heat trace
are
\begin{align}
 D\mathcal H_{a,t}(0)[\alpha]
 &=-t\Tr(V_\alpha e^{-tL_0}),                             \label{eq:newS7V9-first-heat}\\
 \Pi_{a,t}(\alpha,\beta)
 :=D^2\mathcal H_{a,t}(0)[\alpha,\beta]
 &=t\int_0^t
   \Tr\{V_\alpha e^{-uL_0}
             V_\beta e^{-(t-u)L_0}\}\,du
   -t\Tr(W_{\alpha\beta}e^{-tL_0}).                       \label{eq:newS7V9-polarization}
\end{align}
The right side of \eqref{eq:newS7V9-polarization} is symmetric in
\(\alpha,\beta\), even though the symmetry is not termwise visible at a
fixed \(u\).
\end{theorem}

\begin{proof}
For \(L(s)=L_0+sV+O(s^2)\), Duhamel's formula gives
\[
 {d\over ds}e^{-tL(s)}\bigg|_{s=0}
 =-\int_0^t e^{-(t-u)L_0}Ve^{-uL_0}\,du.
\]
Taking the trace and using cyclicity proves
\eqref{eq:newS7V9-first-heat}.  For
\[
 L(s,r)=L_0+sV_\alpha+rV_\beta+srW_{\alpha\beta}
        +O(s^2+r^2),
\]
the mixed coefficient contains the two time orderings of
\(V_\alpha,V_\beta\) and one insertion of \(W_{\alpha\beta}\).  After the
trace, one ordering is
\[
 \int_0^t(t-u)
 \Tr\{V_\alpha e^{-uL_0}V_\beta e^{-(t-u)L_0}\}\,du.
\]
Cyclicity identifies the opposite ordering at \(u\) with the displayed
integrand at \(t-u\).  The two simplex weights therefore add to \(t\).
The single \(W\) insertion gives
\(-t\Tr(W_{\alpha\beta}e^{-tL_0})\), proving
\eqref{eq:newS7V9-polarization} and its symmetry.
\end{proof}

\begin{assessment}[The contact term is part of the Ward identity]
\label{ass:newS7V9-contact}
The last term in \eqref{eq:newS7V9-polarization} is the second derivative of
the lattice operator itself.  Discarding it leaves only the two separated
insertions and generally creates a spurious gauge-boson mass.  It is the
finite-cutoff analogue of the contact or seagull term in continuum
polarization.
\end{assessment}

\subsection{Vanishing of the free one-point functional}
\label{sec:newS7V9-tadpole}

\begin{lemma}[Translation reduction of a linear functional]
\label{lem:newS7V9-translation}
Any translation-invariant linear functional of a lattice one-form has the
form
\begin{equation}
 \mathcal T(\alpha)
 =\sum_{\mu=1}^4 c_\mu\sum_x\alpha_\mu(x),                \label{eq:newS7V9-linear-form}
\end{equation}
with the gauge indices paired against an invariant linear functional on
the represented Lie algebra.
\end{lemma}

\begin{proof}
Write a general linear functional as
\(\sum_{x,\mu}\ell_{x,\mu}(\alpha_\mu(x))\).
Translation invariance makes \(\ell_{x,\mu}\) independent of \(x\).
Invariance under constant gauge conjugation makes its gauge-fibre component
an invariant linear functional.  Summing gives
\eqref{eq:newS7V9-linear-form}.
\end{proof}

\begin{theorem}[The free overlap heat tadpole vanishes]
\label{thm:newS7V9-tadpole-zero}
For every finite lattice, \(t>0\), and link perturbation \(\alpha\),
\begin{equation}
 D\mathcal H_{a,t}(0)[\alpha]=0,
 \qquad
 \Tr(V_\alpha e^{-tL_0})=0.                              \label{eq:newS7V9-tadpole-zero}
\end{equation}
\end{theorem}

\begin{proof}
The free operator and trace are translation invariant, so
\Cref{lem:newS7V9-translation} applies.  Reflect the lattice in the
\(\mu\)-coordinate and reverse the reflected \(\mu\)-links.  The free
Wilson--overlap square is unitarily conjugate to itself, whereas the zero
momentum component \(\sum_x\alpha_\mu(x)\) changes sign.  Hence
\(c_\mu=-c_\mu\), so \(c_\mu=0\) for every \(\mu\).  This proves the first
identity.  The second follows from \eqref{eq:newS7V9-first-heat}.
\end{proof}

\begin{remark}[Abelian factors]
\label{rem:newS7V9-Abelian}
The reflection proof does not require traceless gauge generators and
therefore includes \(U(1)_Y\).  A proof based only on
\(\tr_R X=0\) would miss this case.
\end{remark}

\subsection{Exact Hessian Ward identity}
\label{sec:newS7V9-Ward}

At the trivial link field, define the lattice gauge gradient
\begin{equation}
 (d_a\xi)_\mu(x)
 ={1\over a}\{\xi(x)-\xi(x+a\hat\mu)\}.                  \label{eq:newS7V9-gradient}
\end{equation}
The factor \(a^{-1}\) matches continuum gauge-potential coordinates.

\begin{theorem}[Finite-cutoff transverse polarization]
\label{thm:newS7V9-transverse}
For every site field \(\xi\), lattice one-form \(\beta\), and \(t>0\),
\begin{equation}
 \Pi_{a,t}(d_a\xi,\beta)=0
 =\Pi_{a,t}(\beta,d_a\xi).                                \label{eq:newS7V9-transverse}
\end{equation}
This identity uses the sum of the separated-insertion and contact terms in
\eqref{eq:newS7V9-polarization}.
\end{theorem}

\begin{proof}
Gauge covariance of \(D_a\) makes
\(\mathcal H_{a,t}(A)\) gauge invariant.  If
\(\mathcal G_A\xi\) is the infinitesimal gauge tangent at \(A\), then
\[
 D\mathcal H_{a,t}(A)[\mathcal G_A\xi]=0.
\]
Differentiate this identity at \(A=0\) in the direction \(\beta\):
\[
 D^2\mathcal H_{a,t}(0)[\beta,\mathcal G_0\xi]
 +D\mathcal H_{a,t}(0)[D\mathcal G_0[\beta]\xi]=0.
\]
Here \(\mathcal G_0\xi=d_a\xi\), while the second term vanishes by
\Cref{thm:newS7V9-tadpole-zero}.  This proves the first equality; symmetry
of the Hessian proves the second.
\end{proof}

Use link-midpoint Fourier transforms and put
\begin{equation}
 \widehat q_\mu={2\over a}\sin{aq_\mu\over2}.              \label{eq:newS7V9-lattice-momentum}
\end{equation}
Translation invariance represents the Hessian as
\begin{equation}
 \Pi_{a,t}(\alpha,\beta)
 =\sum_q
 \overline{\widehat\alpha_\mu^A(q)}
 \,\widehat\Pi_{\mu\nu}^{AB}(q;a,t)\,
 \widehat\beta_\nu^B(q).                                  \label{eq:newS7V9-Fourier-polarization}
\end{equation}

\begin{corollary}[Momentum-space Ward identity]
\label{cor:newS7V9-momentum-Ward}
For every lattice momentum,
\begin{equation}
 \widehat q_\mu\widehat\Pi_{\mu\nu}^{AB}(q;a,t)=0,
 \qquad
 \widehat\Pi_{\mu\nu}^{AB}(q;a,t)\widehat q_\nu=0.         \label{eq:newS7V9-momentum-Ward}
\end{equation}
\end{corollary}

\begin{proof}
In midpoint convention the Fourier symbol of \(d_a\) is a harmless unit
phase times \(\widehat q_\mu\).  Insert arbitrary Fourier modes in
\eqref{eq:newS7V9-transverse}; the phases may be absorbed into the link
Fourier transforms.
\end{proof}

\subsection{The unique leading local tensor}
\label{sec:newS7V9-leading-tensor}

\begin{theorem}[Quadratic transverse form]
\label{thm:newS7V9-leading-tensor}
Suppose the free polarization kernel has a finite second spatial moment, so
that \(\widehat\Pi(q;a,t)\) is twice differentiable at \(q=0\).
Hypercubic symmetry, link reflection, constant gauge covariance, and
\eqref{eq:newS7V9-momentum-Ward} then imply
\begin{equation}
 \widehat\Pi_{\mu\nu}^{AB}(q;a,t)
 =c_{a,t}\,\delta^{AB}
   \{\widehat q^{\,2}\delta_{\mu\nu}
      -\widehat q_\mu\widehat q_\nu\}
   +O(a^2|q|^4)+o(|q|^2).                                 \label{eq:newS7V9-leading-tensor}
\end{equation}
For a product gauge group, \(c_{a,t}\delta^{AB}\) is replaced by an
invariant quadratic form on each factor.  There is no constant mass tensor
and no term linear in \(q\).
\end{theorem}

\begin{proof}
Analyticity and reflections remove odd powers.  A constant hypercubic
rank-two tensor is \(m^2\delta_{\mu\nu}\).  Transversality for all small
nonzero \(q\) forces \(m^2=0\).  At quadratic order the most general
hypercubic symmetric tensor is a linear combination of
\(q^2\delta_{\mu\nu}\), \(q_\mu q_\nu\), and
\(\delta_{\mu\nu}q_\mu^2\) with no sum in the last monomial.  Contracting
with \(q_\mu\) forces the last coefficient to vanish and makes the first
two coefficients opposite.  Constant gauge covariance makes the gauge
indices an invariant quadratic form.  Replacing \(q\) by
\(\widehat q=q+O(a^2q^3)\) gives
\eqref{eq:newS7V9-leading-tensor}.
\end{proof}

\begin{corollary}[Identification with the V6--V7 owner]
\label{cor:newS7V9-owner}
The position-space local functional represented by the leading tensor in
\eqref{eq:newS7V9-leading-tensor} is a multiple of
\(\int\langle F_{\mu\nu},F_{\mu\nu}\rangle\).  Thus the exact lattice Ward
identity selects the same quadratic owner as the Wilson-loop and continuum
Dirac computations of V6 and V7.
\end{corollary}

\begin{proof}
The Hessian at \(A=0\) of
\((1/4)\int\langle F_{\mu\nu},F_{\mu\nu}\rangle\) is the transverse Fourier
form \(q^2\delta_{\mu\nu}-q_\mu q_\nu\).  Gauge invariance supplies its
nonlinear completion.
\end{proof}

\subsection{Why separate norm estimates are insufficient}
\label{sec:newS7V9-failed-bound}

\begin{lemma}[Natural AOP derivative sizes]
\label{lem:newS7V9-derivative-sizes}
On a fixed small exponential-link chart about the free field, the sign
resolvent bounds imply
\begin{align}
 \|D_0\|&\le {2\over a},&
 \|\dot D_\alpha\|&\le C\|\alpha\|_\infty,&
 \|\ddot D_{\alpha\beta}\|
 &\le Ca\|\alpha\|_\infty\|\beta\|_\infty,                \label{eq:newS7V9-D-sizes}\\
 \|V_\alpha\|&\le {C\over a}\|\alpha\|_\infty,&
 \|W_{\alpha\beta}\|
 &\le C\|\alpha\|_\infty\|\beta\|_\infty.                 \label{eq:newS7V9-L-sizes}
\end{align}
The constants are independent of lattice volume.
\end{lemma}

\begin{proof}
An exponential-link variation satisfies
\(\|\delta U\|\le a\|\alpha\|_\infty\) and
\(\|\delta^2U\|\le a^2\|\alpha\|_\infty\|\beta\|_\infty\).
The dimensionless Wilson kernel is finite range, so its first and second
variations have the same orders.  The V3 resolvent formula for the sign,
with the free dimensionless gap equal to one, preserves these orders.
The factor \(a^{-1}\) in \(D_{\rm ov}\) gives
\(\dot D=O(1)\) and \(\ddot D=O(a)\).  Insert these estimates into
\eqref{eq:newS7V9-V}--\eqref{eq:newS7V9-W}.
\end{proof}

\begin{proposition}[The naive termwise estimate loses transversality]
\label{prop:newS7V9-naive-fails}
If the two terms in \eqref{eq:newS7V9-polarization} are bounded separately
only by operator norm and matrix dimension, one obtains
\begin{equation}
 |\Pi_{a,t}(\alpha,\beta)|
 \le C\dim(\mathcal H_a)
 \left({t^2\over a^2}+t\right)
 \|\alpha\|_\infty\|\beta\|_\infty.                       \label{eq:newS7V9-naive-bound}
\end{equation}
Since \(\dim(\mathcal H_a)=O(\ell^4a^{-4})\), this estimate neither forbids
a mass term nor yields a mesoscopic continuum bound.
\end{proposition}

\begin{proof}
Use \(\|e^{-uL_0}\|\le1\),
\(|\Tr X|\le\dim(\mathcal H_a)\|X\|\), and
\eqref{eq:newS7V9-L-sizes} in the two terms of
\eqref{eq:newS7V9-polarization}.  The \(u\)-integral has length \(t\) and
the outside factor is \(t\), giving \eqref{eq:newS7V9-naive-bound}.
The estimate has discarded the exact cancellation in
\eqref{eq:newS7V9-transverse}.
\end{proof}

\begin{definition}[Transverse polarization card, TPC]
\label{def:newS7V9-TPC}
TPC consists of:
\begin{enumerate}[label=\textup{(\roman*)},leftmargin=3.5em]
\item the complete Duhamel Hessian
      \eqref{eq:newS7V9-polarization}, including \(W_{\alpha\beta}\);
\item the tadpole identity \eqref{eq:newS7V9-tadpole-zero};
\item exact transversality \eqref{eq:newS7V9-momentum-Ward};
\item a volume-uniform second spatial moment of the assembled polarization
      kernel;
\item a mesoscopic bound for its transverse scalar coefficient
      \(c_{a,t}\); and
\item convergence of the logarithmic part of \(c_{a,t}\) to the V7
      coefficient.
\end{enumerate}
\end{definition}

\begin{programme}[V10 audit and acquisition]
\label{prog:newS7V9-V10}
V10 is a compulsory companion-audit version.  It will first snapshot the
newest settled Yang--Mills and Navier--Stokes notes.  It will then attack
TPC row (iv): rewrite the sum of the separated and contact kernels before
estimation, use overlap locality and the heat semigroup to seek a uniform
second spatial moment, and compare the resulting scale with the free
\(O(a^2/t)\) error from V8.  If that route becomes circular, the programme
will move upstream to a divided-difference formula for the complete
gauge-covariant function \(e^{-tD_{\rm ov}^*D_{\rm ov}}\).
\end{programme}

\begin{firewall}[Claim boundary after seventh-series V9]
\label{fire:newS7V9-boundary}
V9 derives the exact first and second gauge variations of the finite
overlap heat trace, proves the free one-point functional vanishes, proves
exact lattice transversality of the complete polarization, and identifies
its unique leading local transverse tensor.  It also proves quantitatively
that separate norm estimates of the two-insertion and contact terms lose
the cancellation needed for MATC.

V9 does not prove a uniform second kernel moment, compute the scalar
polarization coefficient, establish its logarithmic limit, treat the
nonperturbative interacting AOP family, construct the determinant phase,
remove the cutoff, or construct the Standard Model--Einstein continuum.
\end{firewall}

\begin{theorem}[Seventh-series V9 result]
\label{thm:newS7V9-result}
At every finite cutoff the complete overlap heat polarization is the
Duhamel two-insertion term minus its contact term.  Its tadpole vanishes and
its Hessian annihilates every lattice gauge gradient.  Consequently its
leading analytic local tensor is exactly transverse and has the
\(F_{\mu\nu}F_{\mu\nu}\) form selected independently in V6 and V7.  A
termwise operator-norm estimate destroys this information; the next
analytic target is the second spatial moment of the already assembled
polarization kernel.
\end{theorem}

\begin{proof}
Use \Cref{thm:newS7V9-polarization,
thm:newS7V9-tadpole-zero,
thm:newS7V9-transverse,
cor:newS7V9-momentum-Ward,
thm:newS7V9-leading-tensor,
prop:newS7V9-naive-fails}.
\end{proof}

\paragraph{Parent-version record.}
V9 was formed from
\texttt{Renewal\_SM\_Einstein\_Continuum\_Research\_Notes\_V8.tex}.
The V8 parent had \(3\,576\) logical source lines, \(138\,973\) bytes, and
SHA--256
\[
 \texttt{CD10969468820AB7EB41D557C0BA399D0E9F08375A37D9438FA60B3D0C230532}.
\]
The next compulsory companion audit is V10.

% =============================================================================
% END APPENDED SEVENTH-SERIES V9 MATERIAL
% =============================================================================

% =============================================================================
% BEGIN APPENDED SEVENTH-SERIES V10 MATERIAL
% =============================================================================

\clearpage
\section{Seventh-series V10: companion audit, divided-difference
polarization, and the assembled second-moment target}
\label{sec:v7v10}

\subsection{Compulsory companion audit}
\label{sec:newS7V10-audit}

The companion files were read as immutable byte snapshots.  The newest
settled Yang--Mills note was
\begin{center}
\begin{tabular}{ll}
file &
\texttt{Renewal\_Yang\_Mills\_Mass\_Gap\_Research\_Notes\_V82.tex},\\
bytes & \(1\,208\,484\),\\
logical source lines & \(37\,745\),\\
SHA--256 &
\texttt{CEE404577EEC967DAC992EA9B5F606A61F4029C1112F8BF892CF66BE0EB5CA7F}.
\end{tabular}
\end{center}
Its V82 body is a genuine append to V81.  It expands every percolative
label root by exact connectivity gates, proves the simultaneous Russo
connected-quotient formula, and completely evaluates the acyclic quotient:
a root survives precisely when its block support spans the quotient tree.
It also gives a three-block shortcut counterexample to the tempting claim
that every differentiated forest edge is individually pivotal.  The
remaining cyclic quotient must be combined with the squarefree background
roots that own its Steiner blocks before norms are taken.  This is the
SCOC bottleneck; the Yang--Mills mass gap remains unproved.

The newest settled Navier--Stokes note was
\begin{center}
\begin{tabular}{ll}
file &
\texttt{Renewal\_NS\_Stability\_Research\_Notes\_V59.tex},\\
bytes & \(733\,865\),\\
logical source lines & \(14\,270\),\\
SHA--256 &
\texttt{284F666998334C6581AC1B40950CB64701409A585C05EABF3A158749769BE32E}.
\end{tabular}
\end{center}
V59 proves an explicit Morrey-energy floor at every time before a putative
singularity:
\[
 \mathfrak M(t)\ge
 c_0\varepsilon_*\{\nu,\max(\mathfrak Y(t),1)\},
\]
with
\(\varepsilon_*(\nu,\Lambda)\ge\exp\{-C(\nu)\Lambda^3\}\).
It follows that vanishing parabolic-scale energy forces logarithmically
diverging scaled enstrophy.  It also decomposes the energetic cells of an
intermediate window into finitely many separating clusters and constructs a
nonzero strong limit for each cluster.  The limit exists only on the
controlled window and does not close a global trap; full regularity remains
unproved.

\begin{audit}[V10 transfer decision]
\label{audit:newS7V10-transfer}
The two audits impose the same order on the V9 polarization problem:
\begin{enumerate}[label=\textup{(\roman*)},leftmargin=3.5em]
\item combine the separated Duhamel insertion and its contact owner before
      applying absolute values, just as a cyclic quotient must be combined
      with the background roots that own its Steiner blocks; and
\item do not infer a scale-weighted second spatial moment from unweighted
      locality alone, just as small local energy without the scaled
      enstrophy budget does not close the Navier--Stokes trap.
\end{enumerate}
Accordingly V10 first assembles the exact spectral divided difference, then
converts the Ward identity into a second-difference formula.  The remaining
estimate is explicitly the weighted second moment of that assembled
kernel.
\end{audit}

\subsection{A smooth divided difference for the heat Hessian}
\label{sec:newS7V10-divided-difference}

For \(x,y\ge0\), define
\begin{equation}
 k_t(x,y)
 :=t\int_0^t e^{-ux-(t-u)y}\,du
 =t^2\int_0^1e^{-t\{sx+(1-s)y\}}\,ds.                    \label{eq:newS7V10-kt}
\end{equation}

\begin{lemma}[Divided-difference identities and bounds]
\label{lem:newS7V10-kt}
The kernel \(k_t\) is symmetric, nonnegative, and smooth on
\([0,\infty)^2\), including the diagonal.  It obeys
\begin{align}
 k_t(x,y)
 &=t\,{e^{-ty}-e^{-tx}\over x-y}
 \quad(x\ne y),                                           \label{eq:newS7V10-kt-quotient}\\
 k_t(x,x)&=t^2e^{-tx},                                    \label{eq:newS7V10-kt-diagonal}\\
 |\partial_x^r\partial_y^s k_t(x,y)|
 &\le t^{2+r+s}e^{-t\min(x,y)}
 \qquad(r,s\ge0).                                         \label{eq:newS7V10-kt-derivatives}
\end{align}
\end{lemma}

\begin{proof}
The second integral in \eqref{eq:newS7V10-kt} gives symmetry after
\(s\mapsto1-s\), as well as positivity and smoothness.  Direct integration
gives \eqref{eq:newS7V10-kt-quotient}, whose continuous diagonal value is
\eqref{eq:newS7V10-kt-diagonal}.  Differentiation under the integral
introduces factors \((ts)^r\{t(1-s)\}^s\), each bounded by the corresponding
power of \(t\), and the exponential is at most
\(e^{-t\min(x,y)}\).
\end{proof}

Let \(\{\psi_m\}\) be any orthonormal eigenbasis of the free overlap square
\(L_0\), with \(L_0\psi_m=\lambda_m\psi_m\).

\begin{theorem}[Assembled spectral polarization]
\label{thm:newS7V10-spectral-polarization}
The V9 heat Hessian has the exact representation
\begin{align}
 \Pi_{a,t}(\alpha,\beta)
 &=
 \sum_{m,n}k_t(\lambda_m,\lambda_n)
 (V_\alpha)_{mn}(V_\beta)_{nm}
 -t\sum_m e^{-t\lambda_m}(W_{\alpha\beta})_{mm}.          \label{eq:newS7V10-spectral-polarization}
\end{align}
The first summand is the smooth divided difference of the complete heat
function; it has no singularity when eigenvalues coincide.  The second
summand is its contact owner and must remain in the same assembled
functional.
\end{theorem}

\begin{proof}
Insert the spectral resolution of \(L_0\) into
\eqref{eq:newS7V9-polarization}.  The \(u\)-integral is precisely
\(k_t(\lambda_m,\lambda_n)\).  The contact term is already diagonal after
the trace.
\end{proof}

\begin{corollary}[Degeneracies cost no inverse spectral gap]
\label{cor:newS7V10-no-denominator}
Coincident free eigenvalues and the physical zero mode create no
denominator in \eqref{eq:newS7V10-spectral-polarization}.  In particular,
\[
 0\le k_t(\lambda_m,\lambda_n)
 \le t^2e^{-t\min(\lambda_m,\lambda_n)}.
\]
Any inverse difference \((\lambda_m-\lambda_n)^{-1}\) introduced by using
\eqref{eq:newS7V10-kt-quotient} must be recombined through
\eqref{eq:newS7V10-kt} before estimates.
\end{corollary}

\begin{proof}
Use \Cref{lem:newS7V10-kt}.  The integral definition is regular at every
pair of nonnegative eigenvalues.
\end{proof}

\subsection{Pure-gauge cancellation inside the spectral formula}
\label{sec:newS7V10-gauge-cancellation}

\begin{proposition}[Exact orbit cancellation]
\label{prop:newS7V10-orbit-cancellation}
Let \(X^*=-X\) and
\(L(s)=e^{sX}L_0e^{-sX}\).  Then
\[
 V=[X,L_0],\qquad W=[X,[X,L_0]],
\]
and substitution into \eqref{eq:newS7V10-spectral-polarization} gives zero.
The cancellation occurs pairwise in the unordered eigenvalue pairs.
\end{proposition}

\begin{proof}
In the eigenbasis,
\[
 V_{mn}=(\lambda_n-\lambda_m)X_{mn},
\]
while
\[
 W_{mm}
 =2\sum_n(\lambda_n-\lambda_m)|X_{mn}|^2.
\]
For an unordered pair \(\{m,n\}\), the two ordered separated terms equal
\[
 2k_t(\lambda_m,\lambda_n)
 (\lambda_n-\lambda_m)^2|X_{mn}|^2
 =2t(\lambda_n-\lambda_m)
   (e^{-t\lambda_m}-e^{-t\lambda_n})|X_{mn}|^2,
\]
where \eqref{eq:newS7V10-kt-quotient} is used only as an exact identity.
The two diagonal contact contributions are the negative of this
expression.  Diagonal pairs vanish separately.
\end{proof}

\begin{assessment}[The Yang--Mills shortcut analogue]
\label{ass:newS7V10-shortcut-analogue}
The separated insertion is nonzero on a pure-gauge tangent, and the contact
term is also nonzero.  Their sum vanishes.  Assigning a bound to either one
as though it owned the gauge variation by itself is the exact analogue of
assigning two independent endpoint marks to each Russo edge in the V82
shortcut cycle.  The ownership is visible only after the full connected
object has been assembled.
\end{assessment}

\subsection{Plane-wave vertices and the complete momentum kernel}
\label{sec:newS7V10-vertices}

For a link plane wave of momentum \(q\), let
\[
 \mathcal V_\mu^A(p,q)
 :=\langle p+q|V_\mu^A(q)|p\rangle,
\]
and let
\[
 \mathcal W_{\mu\nu}^{AB}(p;q,-q)
 :=\langle p|W_{\mu\nu}^{AB}(q,-q)|p\rangle.
\]
The free eigenvalue is the scalar \(\lambda_a(ap)\) of V8.

\begin{theorem}[Exact Brillouin polarization formula]
\label{thm:newS7V10-Brillouin}
With spin and internal traces understood,
\begin{align}
 \widehat\Pi_{\mu\nu}^{AB}(q;a,t)
 &=
 \sum_{p\in\mathcal P_a}
 k_t\{\lambda_a(a(p+q)),\lambda_a(ap)\}
 \tr\{\mathcal V_\mu^A(p,q)
          \mathcal V_\nu^B(p+q,-q)\}                     \notag\\
 &\quad
 -t\sum_{p\in\mathcal P_a}e^{-t\lambda_a(ap)}
 \tr\mathcal W_{\mu\nu}^{AB}(p;q,-q).                    \label{eq:newS7V10-Brillouin}
\end{align}
Momenta are read periodically in the Brillouin zone.  Formula
\eqref{eq:newS7V10-Brillouin} contains the physical branch, every doubler
region, and the contact term exactly once.
\end{theorem}

\begin{proof}
Translation invariance makes a first vertex shift momentum by \(q\), the
second shift it back by \(-q\), and the mixed contact vertex diagonal in
the total zero-momentum sector.  Insert these matrix elements into
\eqref{eq:newS7V10-spectral-polarization}.
\end{proof}

\begin{lemma}[First vertex Ward--Takahashi identity]
\label{lem:newS7V10-vertex-Ward}
After absorbing the harmless midpoint phases into the link Fourier
transform,
\begin{equation}
 \widehat q_\mu\mathcal V_\mu^A(p,q)
 =\rho_*(T^A)\lambda_a(ap)
  -\lambda_a(a(p+q))\rho_*(T^A).                          \label{eq:newS7V10-vertex-Ward}
\end{equation}
The differentiated gauge-covariance identity also relates the contraction
of \(\mathcal W_{\mu\nu}\) to the gauge variation of
\(\mathcal V_\nu\).  Together these two identities give the complete
polarization Ward identity of V9 directly in
\eqref{eq:newS7V10-Brillouin}.
\end{lemma}

\begin{proof}
At the trivial connection, a gauge plane wave with generator \(T^A\)
acts on \(L_0\) by the commutator
\([\rho_*(T^A)e^{iqx},L_0]\).  Its link tangent is the lattice gradient, whose
midpoint symbol is \(\widehat q_\mu\).  Taking the matrix element from
\(p\) to \(p+q\) gives \eqref{eq:newS7V10-vertex-Ward}.  Differentiating the
same covariance law once more gives the contact identity; insertion into
\Cref{prop:newS7V10-orbit-cancellation} proves the contracted two-point
formula.
\end{proof}

\subsection{Second differences before absolute values}
\label{sec:newS7V10-second-difference}

Let \(K_{\mu\nu}^{AB}(z;a,t)\) be the translation-invariant position kernel
whose Fourier transform is
\(\widehat\Pi_{\mu\nu}^{AB}(q;a,t)\).  Distances below use the shortest
periodic representative.

\begin{proposition}[Finite-torus holonomy correction]
\label{prop:newS7V10-holonomy-correction}
On a finite periodic torus, the Ward identity at nonzero lattice momentum
does not force
\(\widehat\Pi_{\mu\nu}^{AB}(0;a,t)=0\).  A constant connection can have
nontrivial torus holonomy and need not be a periodic pure gauge.  Therefore
the no-mass conclusion of V9 applies directly to the infinite-lattice
symbol or to the local, zero-winding part of the finite-volume heat kernel.
At finite volume the zero-momentum holonomy term must be retained and
estimated separately.
\end{proposition}

\begin{proof}
A periodic infinitesimal gauge parameter has lattice gradient with zero
average, so it cannot generate an arbitrary constant link potential.  Flat
constant connections are distinguished by their holonomy around the
periodic cycles.  A gauge-invariant heat trace may depend on those
holonomies.  Such dependence is global rather than a local mass owner, but
it can give a nonzero Hessian at the isolated momentum \(q=0\).
\end{proof}

\begin{theorem}[Assembled second-difference representation]
\label{thm:newS7V10-second-difference}
Assume
\begin{equation}
 M_2(a,t):=
 \sum_z|z|^2
 \max_{\mu,\nu,A,B}|K_{\mu\nu}^{AB}(z;a,t)|<\infty.       \label{eq:newS7V10-M2}
\end{equation}
Link reflection implies that the first spatial moment vanishes.  Hence
\begin{align}
 &\widehat\Pi_{\mu\nu}^{AB}(q;a,t)
  -\widehat\Pi_{\mu\nu}^{AB}(0;a,t)                       \notag\\
 &\qquad
 =\sum_z
 \{e^{iq\cdot z}-1-iq\cdot z\}
 K_{\mu\nu}^{AB}(z;a,t).                                  \label{eq:newS7V10-second-difference}
\end{align}
Consequently
\begin{equation}
 |\widehat\Pi_{\mu\nu}^{AB}(q;a,t)
  -\widehat\Pi_{\mu\nu}^{AB}(0;a,t)|
 \le {1\over2}|q|^2M_2(a,t).                              \label{eq:newS7V10-M2-bound}
\end{equation}
Moreover,
\begin{equation}
 \partial_{q_\rho}\partial_{q_\sigma}
 \widehat\Pi_{\mu\nu}^{AB}(0;a,t)
 =-\sum_z z_\rho z_\sigma K_{\mu\nu}^{AB}(z;a,t).         \label{eq:newS7V10-M2-derivative}
\end{equation}
\end{theorem}

\begin{proof}
Reflection of the separation variable gives
\(\sum_z z_\rho K_{\mu\nu}^{AB}(z)=0\).  Subtract the exact constant
Fourier value
\(\widehat\Pi(0)=\sum_zK(z)\) and this zero first moment from the Fourier
series to obtain \eqref{eq:newS7V10-second-difference}.  The elementary
inequality \(|e^{iu}-1-iu|\le u^2/2\) gives
\eqref{eq:newS7V10-M2-bound}.  The absolute second moment justifies twice
differentiating the Fourier series and proves
\eqref{eq:newS7V10-M2-derivative}.
\end{proof}

Define the quadratic moment tensor
\begin{equation}
 \mathcal Q_{\mu\nu;\rho\sigma}^{AB}(a,t)
 :=-\sum_z z_\rho z_\sigma K_{\mu\nu}^{AB}(z;a,t).        \label{eq:newS7V10-Q}
\end{equation}

\begin{corollary}[Precise sufficient local and winding bounds]
\label{cor:newS7V10-sufficient-moment}
Suppose throughout the V8 mesoscopic window that
\begin{align}
 M_2(a,t)&\le C_R\{1+a^2/t\},                              \label{eq:newS7V10-uniform-M2}\\
 \|\mathcal Q(a,t)-\mathcal Q_{\rm Dirac}(t)\|
 &\le C_Ra^2/t,                                           \label{eq:newS7V10-Q-convergence}\\
 \|\widehat\Pi(0;a,t)\|
 &\le C_R\exp\{-c\ell^2/t\}
      +C_R\exp\{-ct/a^2\}.                                \label{eq:newS7V10-holonomy-bound}
\end{align}
Then the local transverse coefficient is uniformly bounded, has
\(O(a^2/t)\) lattice error, and the finite-volume holonomy Hessian is
invisible to every local short-time coefficient.
\end{corollary}

\begin{proof}
Equation \eqref{eq:newS7V10-M2-derivative} identifies the quadratic
coefficient with \(\mathcal Q\); use
\eqref{eq:newS7V10-Q-convergence}.  Equation
\eqref{eq:newS7V10-M2-bound} supplies uniform quadratic control away from
zero momentum.  Both terms in \eqref{eq:newS7V10-holonomy-bound} vanish
faster than every power in the mesoscopic short-time limit, so they cannot
contribute to a local heat coefficient.
\end{proof}

\begin{assessment}[What is proved and what remains in the moment row]
\label{ass:newS7V10-moment-status}
At each finite lattice, \(M_2(a,t)\) is finite simply because the lattice is
finite.  V10 proves that this is exactly the weighted quantity controlling
the local transverse coefficient after all Ward cancellations.  It does
not prove the cutoff-uniform estimate
\eqref{eq:newS7V10-uniform-M2}, the tensor convergence
\eqref{eq:newS7V10-Q-convergence}, or the winding estimate
\eqref{eq:newS7V10-holonomy-bound}.  Exponential locality of each overlap
vertex controls separation in units of \(a\), while the heat factors spread
over the larger scale \(\sqrt t\).  Those scales must be convolved before
the second moment is estimated.
\end{assessment}

\subsection{The next upstream kernel estimate}
\label{sec:newS7V10-next}

\begin{definition}[Assembled polarization moment card, APMC]
\label{def:newS7V10-APMC}
APMC consists of:
\begin{enumerate}[label=\textup{(\roman*)},leftmargin=3.5em]
\item the smooth spectral formula
      \eqref{eq:newS7V10-spectral-polarization};
\item the pairwise pure-gauge cancellation of
      \Cref{prop:newS7V10-orbit-cancellation};
\item the exact Brillouin formula \eqref{eq:newS7V10-Brillouin};
\item the vertex Ward identity \eqref{eq:newS7V10-vertex-Ward};
\item the finite-volume-correct second-difference representation
      \eqref{eq:newS7V10-second-difference};
\item the uniform assembled moment estimate
      \eqref{eq:newS7V10-uniform-M2};
\item the quadratic tensor convergence
      \eqref{eq:newS7V10-Q-convergence}; and
\item the finite-volume holonomy bound
      \eqref{eq:newS7V10-holonomy-bound}.
\end{enumerate}
\end{definition}

\begin{theorem}[APMC reduction]
\label{thm:newS7V10-APMC}
Rows \textup{(i)}--\textup{(v)} of APMC hold at finite cutoff.  If rows
\textup{(vi)}--\textup{(viii)} hold, then TPC rows
\textup{(iv)}--\textup{(vi)} reduce to identifying
\(\mathcal Q_{\rm Dirac}\) with the V7 coefficient.  No local gauge-boson
mass subtraction is needed; the separate finite-volume holonomy Hessian
vanishes beyond all orders in the local short-time expansion.
\end{theorem}

\begin{proof}
Rows \textup{(i)}--\textup{(v)} are
\Cref{thm:newS7V10-spectral-polarization,
prop:newS7V10-orbit-cancellation,
thm:newS7V10-Brillouin,
lem:newS7V10-vertex-Ward,
thm:newS7V10-second-difference}.  Under the remaining rows,
\eqref{eq:newS7V10-M2-derivative} controls and identifies the quadratic
local tensor, while \eqref{eq:newS7V10-holonomy-bound} removes the global
zero-momentum term from every local coefficient.  V7 supplies the
continuum target.
\end{proof}

\begin{programme}[V11 acquisition]
\label{prog:newS7V10-V11}
V11 will prove separate weighted bounds that are strong enough to be
recombined at the APMC level:
\begin{enumerate}[leftmargin=3em]
\item derive a Gaussian-plus-lattice-tail estimate for the free overlap heat
      kernel from the V8 symbol;
\item prove first and second vertex kernels have exponential moments on the
      scale \(a\) using the Wilson gap and the sign resolvent;
\item convolve the heat scale \(\sqrt t\) with the vertex scale \(a\);
\item insert the Ward second difference before summing the resulting
      second moment; and
\item determine whether the contact term cancels the leading
      \(t/a^2\) contribution uniformly.
\end{enumerate}
If the last cancellation cannot be seen in position space, the calculation
will return upstream to the Brillouin divided difference
\eqref{eq:newS7V10-Brillouin}.
\end{programme}

\begin{firewall}[Claim boundary after seventh-series V10]
\label{fire:newS7V10-boundary}
V10 performs the compulsory YM/NS audit, rewrites the complete finite
overlap heat polarization as a smooth spectral divided difference plus its
contact owner, proves the pure-gauge cancellation pairwise, derives the
exact Brillouin vertex formula and Ward--Takahashi identity, and converts
the local momentum dependence into a second-difference representation
controlled by one assembled second spatial moment.  It also corrects the
finite-torus zero-momentum claim by retaining the global holonomy Hessian.

V10 does not prove the cutoff-uniform second-moment estimate, the quadratic
tensor convergence, or the finite-volume winding bound; extract the
numerical overlap polarization coefficient; establish interacting MATC;
construct the chiral determinant phase; remove the cutoff; prove either
companion programme; or construct the Standard Model--Einstein continuum.
\end{firewall}

\begin{theorem}[Seventh-series V10 result]
\label{thm:newS7V10-result}
The V9 polarization bottleneck has an exact nonsingular spectral
representation.  Degenerate eigenvalues cost no inverse gap, and a
pure-gauge separated insertion cancels its contact owner pairwise.  After
that assembly, the local momentum difference
\(\widehat\Pi(q)-\widehat\Pi(0)\) is a second difference whose quadratic
coefficient equals the second spatial moment
\eqref{eq:newS7V10-M2-derivative}.  The isolated finite-torus value
\(\widehat\Pi(0)\) is a global holonomy term and requires the separate
winding estimate \eqref{eq:newS7V10-holonomy-bound}.  The remaining targets
are therefore the uniform moment, its V7 tensor limit, and winding
suppression.
\end{theorem}

\begin{proof}
Use \Cref{audit:newS7V10-transfer,
lem:newS7V10-kt,
thm:newS7V10-spectral-polarization,
prop:newS7V10-orbit-cancellation,
thm:newS7V10-Brillouin,
lem:newS7V10-vertex-Ward,
thm:newS7V10-second-difference,
thm:newS7V10-APMC}.
\end{proof}

\paragraph{Parent-version record.}
V10 was formed from
\texttt{Renewal\_SM\_Einstein\_Continuum\_Research\_Notes\_V9.tex}.
The V9 parent had \(3\,991\) logical source lines, \(155\,101\) bytes, and
SHA--256
\[
 \texttt{7FBE1D2664EB4BC6995C96554B1F852625C0B54D4FA30F3E1DE7B6A0C66F4F41}.
\]
The next compulsory companion audit is V15.

% =============================================================================
% END APPENDED SEVENTH-SERIES V10 MATERIAL
% =============================================================================

% =============================================================================
% BEGIN APPENDED SEVENTH-SERIES V11 MATERIAL
% =============================================================================

\clearpage
\section{Seventh-series V11: free overlap heat-kernel moments and
finite-torus winding suppression}
\label{sec:v7v11}

V10 reduced the local polarization coefficient to a second spatial moment
and separated the isolated torus-holonomy Hessian.  This version proves the
free heat-kernel estimate needed for both tasks.  The result is stated with
rapid off-diagonal decay rather than an unproved pointwise Gaussian bound;
rapid decay is already sufficient for every polynomial moment and for
removing winding terms from local heat coefficients.

\subsection{The infinite-lattice physical heat kernel}
\label{sec:newS7V11-kernel}

Let \(\lambda_a(ap)\) be the exact free overlap dispersion from V8.  On the
infinite lattice \(a\mathbb Z^4\), define the physical scalar heat kernel
\begin{equation}
 \mathsf K_a(t;x)
 :=\int_{[-\pi/a,\pi/a]^4}{d^4p\over(2\pi)^4}
 e^{ip\cdot x}e^{-t\lambda_a(ap)}.                        \label{eq:newS7V11-kernel}
\end{equation}
It has dimension length\(^{-4}\) and is normalized for the lattice measure
\(a^4\sum_x\).

\begin{lemma}[Low and high symbol derivative bounds]
\label{lem:newS7V11-symbol-derivatives}
Choose a smooth cutoff \(\chi(\theta)\) supported in
\(|\theta|<2\delta\) and equal to one on \(|\theta|\le\delta\), where
\(\delta\) is the V8 physical-branch radius.  For every multi-index
\(\alpha\) and \(t\ge a^2\), the low symbol
\[
 \chi(ap)e^{-t\lambda_a(ap)}
\]
has, after the rescaling \(k=\sqrt t\,p\), uniformly integrable
\(k\)-derivatives through order \(|\alpha|\).  The complementary symbol
obeys
\begin{equation}
 \int_{[-\pi,\pi]^4}
 \left|\partial_\theta^\alpha
 \{(1-\chi(\theta))e^{-(t/a^2)\,a^2\lambda_a(\theta)}\}\right|\,d^4\theta
 \le C_\alpha e^{-c_\alpha t/a^2}.                        \label{eq:newS7V11-high-derivatives}
\end{equation}
\end{lemma}

\begin{proof}
On \(|\theta|<2\delta\), V8 gives
\(c|\theta|^2\le a^2\lambda_a(\theta)\le C|\theta|^2\) and
\[
 a^2\lambda_a(\theta)=|\theta|^2+O(|\theta|^4).
\]
The exact formula in V8 is smooth there because
\(\omega(\theta)\) is bounded away from zero.  Repeated chain rules after
\(\theta=(a/\sqrt t)k\) give a polynomial in \(k\), with coefficients
bounded by powers of \(a/\sqrt t\le1\), times \(e^{-c|k|^2}\).  This proves
the low assertion.

On the support of \(1-\chi\), global V8 coercivity gives
\(a^2\lambda_a\ge c>0\).  Every \(\theta\)-derivative produces only a
polynomial in \(t/a^2\), since the symbol and all its derivatives are
bounded on that compact set.  Absorb that polynomial into a slightly
weaker exponential to obtain
\eqref{eq:newS7V11-high-derivatives}.
\end{proof}

\begin{theorem}[Rapid heat-scale decay plus lattice tail]
\label{thm:newS7V11-kernel-bound}
For every integer \(M\ge0\), there are constants \(C_M,c_M>0\), independent
of \(a,t,x\), such that for \(t\ge a^2\),
\begin{align}
 |\mathsf K_a(t;x)|
 &\le C_Mt^{-2}
 \left(1+{|x|\over\sqrt t}\right)^{-M}                    \notag\\
 &\quad
 +C_Ma^{-4}e^{-c_Mt/a^2}
 \left(1+{|x|\over a}\right)^{-M}.                        \label{eq:newS7V11-kernel-bound}
\end{align}
\end{theorem}

\begin{proof}
Split \eqref{eq:newS7V11-kernel} with \(\chi(ap)\).
In the low term set \(k=\sqrt t\,p\).  Its prefactor is \(t^{-2}\).
For \(|x|\le\sqrt t\), the \(L^1\) bound from
\Cref{lem:newS7V11-symbol-derivatives} gives the first term in
\eqref{eq:newS7V11-kernel-bound}.  For
\(|x|>\sqrt t\), integrate by parts \(M\) times in a coordinate for which
\(|x_\mu|\ge|x|/2\).  The same derivative bounds give the factor
\((\sqrt t/|x|)^M\).

For the high term set \(\theta=ap\), producing the prefactor \(a^{-4}\).
The zero-derivative estimate gives the bound for \(|x|\le a\).
For \(|x|>a\), integrate by parts in \(\theta\) and use
\eqref{eq:newS7V11-high-derivatives}.  The cutoff makes the integrand
periodic and smooth at the Brillouin boundary, so no boundary term remains.
\end{proof}

\subsection{Uniform zeroth and second moments}
\label{sec:newS7V11-moments}

For a scalar lattice kernel \(f\), put
\begin{equation}
 m_j(f):=a^4\sum_{x\in a\mathbb Z^4}|x|^j|f(x)|,
 \qquad j=0,1,2.                                          \label{eq:newS7V11-moments}
\end{equation}

\begin{lemma}[Lattice radial summation]
\label{lem:newS7V11-radial-sum}
If \(M>4+j\) and \(\sigma\ge a\), then
\begin{equation}
 a^4\sum_{x\in a\mathbb Z^4}
 |x|^j\sigma^{-4}
 \left(1+{|x|\over\sigma}\right)^{-M}
 \le C_{M,j}\sigma^j.                                    \label{eq:newS7V11-radial-sum}
\end{equation}
\end{lemma}

\begin{proof}
Partition the lattice into shells
\(n\sigma\le|x|<(n+1)\sigma\).  Since \(\sigma\ge a\), the number of lattice
points in the \(n\)-th shell is at most
\(C(\sigma/a)^4(1+n)^3\).  The summand is bounded by
\(C\sigma^{j-4}(1+n)^{j-M}\).  Multiplication by \(a^4\) and summation in
\(n\) proves the result when \(M>4+j\).
\end{proof}

\begin{theorem}[Free overlap heat moments]
\label{thm:newS7V11-heat-moments}
For \(t\ge a^2\),
\begin{align}
 m_0\{\mathsf K_a(t;\cdot)\}&\le C,                       \label{eq:newS7V11-m0}\\
 m_1\{\mathsf K_a(t;\cdot)\}&\le C\sqrt t,                \label{eq:newS7V11-m1}\\
 m_2\{\mathsf K_a(t;\cdot)\}&\le Ct.                      \label{eq:newS7V11-m2}
\end{align}
The constants are independent of the cutoff.
\end{theorem}

\begin{proof}
Apply \Cref{lem:newS7V11-radial-sum} to the first term of
\eqref{eq:newS7V11-kernel-bound} with \(\sigma=\sqrt t\).
For the lattice tail use \(\sigma=a\).  Its \(j\)-th moment is bounded by
\(Ca^je^{-ct/a^2}\), which is at most \(Ct^{j/2}\) for \(t\ge a^2\).
Choose \(M>6\).
\end{proof}

\begin{lemma}[Weighted convolution algebra]
\label{lem:newS7V11-convolution}
For lattice convolution
\[
 (f*g)(x)=a^4\sum_y f(y)g(x-y),
\]
one has
\begin{align}
 m_0(f*g)&\le m_0(f)m_0(g),                               \label{eq:newS7V11-conv0}\\
 m_1(f*g)&\le m_1(f)m_0(g)+m_0(f)m_1(g),                  \label{eq:newS7V11-conv1}\\
 m_2(f*g)&\le
 2m_2(f)m_0(g)+2m_0(f)m_2(g).                            \label{eq:newS7V11-conv2}
\end{align}
\end{lemma}

\begin{proof}
Take absolute values, use \(x=y+(x-y)\), and apply
\(|u+v|\le|u|+|v|\) and
\(|u+v|^2\le2|u|^2+2|v|^2\).  Tonelli's theorem separates the two sums.
\end{proof}

\begin{corollary}[Heat and vertex scales add at the level of moments]
\label{cor:newS7V11-scale-addition}
If a vertex kernel \(v_a\) obeys
\[
 m_0(v_a)\le C_v,\qquad m_2(v_a)\le C_va^2,
\]
then
\begin{equation}
 m_2(v_a*\mathsf K_a(t;\cdot)*v_a)
 \le C'_v(t+a^2).                                        \label{eq:newS7V11-scale-addition}
\end{equation}
\end{corollary}

\begin{proof}
Apply \eqref{eq:newS7V11-conv2} twice, together with
\Cref{thm:newS7V11-heat-moments}.
\end{proof}

\begin{assessment}[Why this does not yet prove the polarization moment]
\label{ass:newS7V11-not-yet-polarization}
The separated term in the polarization contains two vertices and two heat
segments and therefore has the expected combined scales \(a\) and
\(\sqrt t\).  The contact term has a different kernel.  Applying
\Cref{cor:newS7V11-scale-addition} to their absolute values separately
would again lose the pairwise cancellation proved in V10.  V11 supplies
the heat factor of APMC; a vertex estimate must be inserted into the
already assembled second difference.
\end{assessment}

\subsection{Exact periodization and the zero-momentum holonomy term}
\label{sec:newS7V11-winding}

Let \(\mathsf K_{a,\ell}\) be the free heat kernel on the lattice torus of
side \(\ell=Na\).

\begin{lemma}[Exact periodization]
\label{lem:newS7V11-periodization}
For lattice points \(x,y\) on the torus,
\begin{equation}
 \mathsf K_{a,\ell}(t;x,y)
 =\sum_{n\in\mathbb Z^4}
 \mathsf K_a(t;x-y+n\ell).                                \label{eq:newS7V11-periodization}
\end{equation}
The sum converges absolutely for every \(t\ge a^2\).
\end{lemma}

\begin{proof}
Apply the finite Fourier sampling identity to the Brillouin-periodic
function \(e^{-t\lambda_a(ap)}\).  Momentum sampling at
\((2\pi/\ell)\mathbb Z^4\) selects precisely the position Fourier modes
whose displacement differs by \(n\ell\).  Absolute convergence follows
from \eqref{eq:newS7V11-kernel-bound} with \(M>4\).
\end{proof}

For a constant commuting gauge potential \(h\), the free Abelianized trace
is
\[
 \mathcal H_{a,t}(h)
 =r\sum_{p\in\mathcal P_a}
 e^{-t\lambda_a\{a(p+h)\}},
\]
where \(r\) is the spin--internal multiplicity.  This is sufficient for the
quadratic holonomy Hessian of each gauge-algebra direction at the trivial
connection.

\begin{theorem}[Zero winding is exactly constant-shift invariant]
\label{thm:newS7V11-zero-winding}
Poisson periodization gives
\begin{equation}
 \mathcal H_{a,t}(h)
 =r\ell^4\sum_{n\in\mathbb Z^4}
 e^{-i\ell n\cdot h}\mathsf K_a(t;n\ell).                 \label{eq:newS7V11-holonomy-series}
\end{equation}
The zero-winding term is independent of \(h\).  Therefore the entire
zero-momentum Hessian is a sum over \(n\ne0\):
\begin{equation}
 \partial_{h_\mu}\partial_{h_\nu}\mathcal H_{a,t}(0)
 =-r\ell^4\sum_{n\ne0}
 (\ell n)_\mu(\ell n)_\nu\mathsf K_a(t;n\ell).            \label{eq:newS7V11-holonomy-Hessian}
\end{equation}
\end{theorem}

\begin{proof}
Expand the Brillouin-periodic function in its position Fourier series and
sum it over the torus momentum grid.  This gives
\eqref{eq:newS7V11-holonomy-series}.  The \(n=0\) phase is one for every
\(h\).  Twice differentiate the absolutely convergent nonzero-winding
sum, using \Cref{thm:newS7V11-kernel-bound} with \(M>6\).
\end{proof}

\begin{theorem}[Holonomy is beyond every local heat order]
\label{thm:newS7V11-holonomy-bound}
For every integer \(N_0\ge1\), there are constants \(C_{N_0},c>0\) such
that, whenever \(a^2\le t\le\ell^2/4\),
\begin{align}
 \left\|
 \partial_h^2\mathcal H_{a,t}(0)
 \right\|
 &\le C_{N_0}r\ell^2
 \left\{
 \left({t\over\ell^2}\right)^{N_0}
 +e^{-ct/a^2}
 \right\}.                                                \label{eq:newS7V11-holonomy-bound}
\end{align}
Thus the finite-torus value \(\widehat\Pi(0;a,t)\) vanishes beyond every
power relevant to a local heat coefficient in the joint mesoscopic limit.
\end{theorem}

\begin{proof}
Insert \eqref{eq:newS7V11-kernel-bound} into
\eqref{eq:newS7V11-holonomy-Hessian}.  In its heat-scale term choose
\(M=2N_0+4\).  Since \(n\ne0\) and \(\sqrt t\le\ell/2\),
\[
 \ell^4\sum_{n\ne0}|\ell n|^2t^{-2}
 \left(1+{|\ell n|\over\sqrt t}\right)^{-M}
 \le C_{N_0}\ell^2(t/\ell^2)^{N_0}.
\]
For the lattice tail choose any \(M>6\).  The factor
\((1+|\ell n|/a)^{-M}\), summed with \(|n|^2\), cancels the prefactor
\(\ell^6a^{-4}\) up to a bounded power of \(a/\ell\), leaving
\(C\ell^2e^{-ct/a^2}\).
\end{proof}

\begin{corollary}[Correction of the V10 sufficient winding row]
\label{cor:newS7V11-V10-correction}
The exponential finite-volume hypothesis
\(\exp\{-c\ell^2/t\}\) in V10 was sufficient but stronger than needed.
For APMC it may be replaced by the proved bound
\eqref{eq:newS7V11-holonomy-bound}: decay faster than every power of
\(t/\ell^2\) removes the holonomy Hessian from all local heat
coefficients.
\end{corollary}

\begin{proof}
A coefficient at any fixed heat order is extracted after multiplication
by a fixed power of \(t\).  Choose \(N_0\) larger than that power and use
\eqref{eq:newS7V11-holonomy-bound}; the lattice-tail term vanishes because
\(t/a^2\to\infty\) in the mesoscopic window.
\end{proof}

\begin{definition}[Free weighted heat card, FWHC]
\label{def:newS7V11-FWHC}
FWHC consists of:
\begin{enumerate}[label=\textup{(\roman*)},leftmargin=3.5em]
\item the low/high symbol derivative bounds of
      \Cref{lem:newS7V11-symbol-derivatives};
\item the rapid heat-scale and lattice-tail estimate
      \eqref{eq:newS7V11-kernel-bound};
\item the cutoff-uniform moments
      \eqref{eq:newS7V11-m0}--\eqref{eq:newS7V11-m2};
\item the weighted convolution algebra
      \eqref{eq:newS7V11-conv0}--\eqref{eq:newS7V11-conv2};
\item exact torus periodization
      \eqref{eq:newS7V11-periodization}; and
\item the all-orders winding estimate
      \eqref{eq:newS7V11-holonomy-bound}.
\end{enumerate}
\end{definition}

\begin{programme}[V12 acquisition]
\label{prog:newS7V11-V12}
V12 will attack the vertex half of APMC.  It will write the first and
second overlap vertices through the resolvent derivative of the sign,
insert exponential weights before the proper-time integral, and prove
zeroth and second spatial moments on the scale \(a\).  The final estimate
will be made only after the vertex Ward identity combines the separated
and contact owners.  If a second sign derivative loses two inverse gaps,
the calculation will return upstream to a joint divided difference of the
complete overlap square rather than estimate those derivatives
separately.
\end{programme}

\begin{firewall}[Claim boundary after seventh-series V11]
\label{fire:newS7V11-boundary}
V11 proves rapid off-diagonal decay for the free overlap heat kernel,
uniform zeroth through second spatial moments, a weighted convolution
algebra, exact finite-torus periodization, exact independence of the
zero-winding trace under a constant momentum shift, and suppression of the
remaining holonomy Hessian beyond every local heat order.

V11 does not prove the required weighted bounds for overlap gauge vertices,
assemble the interacting polarization moment, extract the V7 coefficient,
establish interacting MATC, construct the chiral determinant phase, remove
the cutoff, or construct the Standard Model--Einstein continuum.
\end{firewall}

\begin{theorem}[Seventh-series V11 result]
\label{thm:newS7V11-result}
On \(t\ge a^2\), the free overlap heat kernel has spatial scale
\(\sqrt t\), a cutoff-uniform \(L^1\) bound, and second absolute moment
\(O(t)\), plus an exponentially small lattice-scale tail.  Its finite-torus
zero-momentum Hessian consists entirely of nonzero windings and is smaller
than every power of \(t/\ell^2\) in the mesoscopic limit.  The free heat
and winding rows of APMC are therefore closed; the remaining local
bottleneck is the assembled first/second overlap vertex moment.
\end{theorem}

\begin{proof}
Use \Cref{thm:newS7V11-kernel-bound,
thm:newS7V11-heat-moments,
lem:newS7V11-convolution,
lem:newS7V11-periodization,
thm:newS7V11-zero-winding,
thm:newS7V11-holonomy-bound,
cor:newS7V11-V10-correction}.
\end{proof}

\paragraph{Parent-version record.}
V11 was formed from
\texttt{Renewal\_SM\_Einstein\_Continuum\_Research\_Notes\_V10.tex}.
The V10 parent had \(4\,503\) logical source lines, \(176\,061\) bytes, and
SHA--256
\[
 \texttt{A2D59C982EB0B4A459FA2A68F705D5AC1873190DC040D526087AA50D1B4B13D8}.
\]
The next compulsory companion audit is V15.

% =============================================================================
% END APPENDED SEVENTH-SERIES V11 MATERIAL
% =============================================================================

% =============================================================================
% BEGIN APPENDED SEVENTH-SERIES V12 MATERIAL
% =============================================================================

\clearpage
\section{Seventh-series V12: weighted overlap vertices from the Wilson
gap}
\label{sec:v7v12}

V11 proves the weighted heat-kernel half of APMC.  The remaining input to
the polarization kernel is the response of the overlap operator to one and
two localized gauge links.  This version proves their exponential
locality, physical scaling, and spatial moments directly from the AOP
Wilson gap.

\subsection{Weighted resolvents for a finite-range gapped kernel}
\label{sec:newS7V12-resolvent}

Use lattice distance in units of one lattice edge and denote it by
\(d(x,y)\).  Let \(\mathsf A(U)=aH_W(U)\) be the dimensionless Hermitian
Wilson kernel.  On the AOP class,
\begin{equation}
 \mathsf A(U)^2\succeq c_\delta^2I,\qquad
 \|\mathsf A(U)\|\le7,\qquad
 c_\delta=\sqrt{1-30\delta}>0.                            \label{eq:newS7V12-window}
\end{equation}
The hopping range of \(\mathsf A\) is one.

\begin{lemma}[Combes--Thomas conjugation]
\label{lem:newS7V12-CT}
Let \(\varphi\) be a real lattice function satisfying
\(|\varphi(x)-\varphi(y)|\le d(x,y)\).  There are
\(\mu_0>0\) and \(C_{\rm CT}<\infty\), depending only on the window
\eqref{eq:newS7V12-window}, such that for
\(0\le\mu\le\mu_0\) and every \(z\) at distance at least
\(c_\delta/4\) from \(\operatorname{spec}\mathsf A\),
\begin{equation}
 \left\|
 e^{\mu\varphi}(z-\mathsf A)^{-1}e^{-\mu\varphi}
 \right\|\le C_{\rm CT}.                                  \label{eq:newS7V12-CT}
\end{equation}
The constants are independent of lattice volume and spacing.
\end{lemma}

\begin{proof}
Finite hopping range and the Lipschitz condition give
\[
 \|e^{\mu\varphi}\mathsf A e^{-\mu\varphi}-\mathsf A\|
 \le C_H(e^\mu-1),
\]
where \(C_H\) is the uniform sum of the finitely many Wilson hopping
coefficients.  Choose \(\mu_0\) so that the right side is at most
\(c_\delta/8\).  The unweighted resolvent has norm at most
\(4/c_\delta\).  The resolvent identity and a Neumann series then give
\eqref{eq:newS7V12-CT}, uniformly in volume.
\end{proof}

\begin{corollary}[Uniform exponential resolvent kernel]
\label{cor:newS7V12-resolvent-kernel}
Under the same hypotheses,
\begin{equation}
 \|(z-\mathsf A)^{-1}(x,y)\|
 \le C_{\rm CT}e^{-\mu_0d(x,y)}.                          \label{eq:newS7V12-resolvent-kernel}
\end{equation}
\end{corollary}

\begin{proof}
In \Cref{lem:newS7V12-CT} choose
\(\varphi(w)=d(w,y)\), take the matrix element from \(y\) to \(x\), and
use \(\varphi(y)=0\).
\end{proof}

\subsection{Holomorphic sign derivatives}
\label{sec:newS7V12-sign}

Let \(\Gamma_+\) and \(\Gamma_-\) be fixed contours surrounding the positive
and negative spectral intervals of \(\mathsf A\), respectively, remaining
at distance at least \(c_\delta/4\) from the spectrum.  On their disjoint
interiors the sign is the holomorphic function equal to \(+1\) and \(-1\).

\begin{lemma}[First and second sign derivatives]
\label{lem:newS7V12-sign-derivatives}
For bounded perturbations \(B,C\),
\begin{align}
 D\sgn(\mathsf A)[B]
 &={1\over2\pi i}\int_{\Gamma_+\cup\Gamma_-}
 \sgn(z)R_zBR_z\,dz,                                     \label{eq:newS7V12-first-sign}\\
 D^2\sgn(\mathsf A)[B,C]
 &={1\over2\pi i}\int_{\Gamma_+\cup\Gamma_-}\sgn(z)
 \{R_zBR_zCR_z+R_zCR_zBR_z\}\,dz,                        \label{eq:newS7V12-second-sign}
\end{align}
where \(R_z=(z-\mathsf A)^{-1}\).  In operator norm,
\begin{align}
 \|D\sgn(\mathsf A)[B]\|&\le C_\delta\|B\|,               \label{eq:newS7V12-first-sign-norm}\\
 \|D^2\sgn(\mathsf A)[B,C]\|
 &\le C_\delta\|B\|\,\|C\|.                              \label{eq:newS7V12-second-sign-norm}
\end{align}
\end{lemma}

\begin{proof}
The holomorphic functional calculus gives
\[
 \sgn(\mathsf A)
 ={1\over2\pi i}\int_{\Gamma_+\cup\Gamma_-}
 \sgn(z)(z-\mathsf A)^{-1}\,dz.
\]
Differentiate the resolvent once and twice.  The two possible orders of
\(B,C\) give \eqref{eq:newS7V12-second-sign}.  The contours have fixed
length and the resolvent is uniformly bounded, proving the norm estimates.
\end{proof}

\begin{remark}[No artificial second inverse-gap integral]
\label{rem:newS7V12-gap-cost}
The constants in \eqref{eq:newS7V12-first-sign-norm} and
\eqref{eq:newS7V12-second-sign-norm} depend on the fixed AOP gap, but the
contour representation has no large proper-time tail and no singular
eigenvalue divided difference.  It is therefore the appropriate upstream
formula for localized vertex moments.
\end{remark}

\subsection{Localized link insertions and tree decay}
\label{sec:newS7V12-tree}

For a finite set of sites and links \(S\), let
\(\tau(S)\) be the number of edges in a shortest lattice tree meeting every
element of \(S\).  A link is met when either endpoint belongs to the tree.
For a positively oriented link \(e\), vary
\begin{equation}
 U_e(s)=e^{saX_e}U_e,\qquad \|X_e\|\le1,                  \label{eq:newS7V12-link-path}
\end{equation}
and leave all other links fixed.  The parameter \(X_e\) has the physical
dimension of a gauge potential.

\begin{lemma}[Wilson insertion sizes]
\label{lem:newS7V12-Wilson-insertions}
For the dimensionless Wilson kernel,
\begin{align}
 B_e:={d\over ds}\mathsf A(U(s))\bigg|_{0}
 &\quad\hbox{is supported within one edge of }e,\qquad
 \|B_e\|\le C_Wa,                                         \label{eq:newS7V12-Be}\\
 C_{ef}:={\partial^2\over\partial s\partial r}
 \mathsf A(U_{e}(s),U_f(r))\bigg|_{0}
 & =0\quad(e\ne f),\qquad
 \|C_{ee}\|\le C_Wa^2.                                   \label{eq:newS7V12-Cef}
\end{align}
\end{lemma}

\begin{proof}
Each link occurs linearly in one forward shift and its adjoint occurs
linearly in the reverse shift.  Differentiating
\(e^{saX_e}\) gives \(aX_e\), and differentiating twice gives \(a^2X_e^2\).
All remaining Wilson coefficients are dimensionless and bounded.  Two
distinct links never occur in the same matrix element of the nearest-
neighbor Wilson kernel, proving the cross derivative is zero.
\end{proof}

\begin{theorem}[Tree-localized sign vertices]
\label{thm:newS7V12-sign-tree}
There are \(C,\mu>0\), uniform on the AOP chart, such that
\begin{align}
 \|D\sgn(\mathsf A)[B_e](x,y)\|
 &\le Ca\,e^{-\mu\tau(\{x,y,e\})},                        \label{eq:newS7V12-sign-tree-one}\\
 \|D^2\sgn(\mathsf A)[B_e,B_f](x,y)\|
 &\le Ca^2e^{-\mu\tau(\{x,y,e,f\})}.                      \label{eq:newS7V12-sign-tree-two}
\end{align}
The full second derivative along link coordinates, including
\(D\sgn(\mathsf A)[C_{ee}]\), obeys the same \(Ca^2\) bound with the
appropriate repeated-link tree.
\end{theorem}

\begin{proof}
Insert the localized kernel \(B_e\) between the two resolvents in
\eqref{eq:newS7V12-first-sign}.  By
\eqref{eq:newS7V12-resolvent-kernel}, a summand through sites adjacent to
\(e\) costs
\[
 Ca\exp\{-\mu(d(x,e)+d(e,y))\},
\]
which dominates the first tree length after decreasing \(\mu\).

For \eqref{eq:newS7V12-second-sign}, each ordered product contains three
resolvents and the two localized insertions.  Summing their adjacent
endpoints gives the exponential of the length of an ordered connecting
walk.  The shorter of the two insertion orders dominates a spanning tree
through \(x,y,e,f\).  Polynomially many choices at fixed excess length are
absorbed by decreasing \(\mu\).  The factor is \(a^2\) by
\eqref{eq:newS7V12-Be}.  The \(C_{ee}\) contribution follows from
\eqref{eq:newS7V12-first-sign} and \eqref{eq:newS7V12-Cef}.
\end{proof}

\subsection{Overlap and overlap-square vertex scales}
\label{sec:newS7V12-overlap}

Recall
\[
 D_{\rm ov}(U)={1\over a}\{I+\gamma_5\sgn\mathsf A(U)\}.
\]
Let \(\dot D_e\) and \(\ddot D_{ef}\) denote its physical link derivatives
along \eqref{eq:newS7V12-link-path}.

\begin{corollary}[First and second overlap vertices]
\label{cor:newS7V12-overlap-vertices}
Uniformly on the AOP chart,
\begin{align}
 \|\dot D_e(x,y)\|
 &\le C e^{-\mu\tau(\{x,y,e\})},                          \label{eq:newS7V12-D-one}\\
 \|\ddot D_{ef}(x,y)\|
 &\le Ca e^{-\mu\tau(\{x,y,e,f\})}.                       \label{eq:newS7V12-D-two}
\end{align}
Consequently, after summing the fermion endpoints and, in the second line,
the second link,
\begin{align}
 \sup_e\sum_{x,y}
 \{1+a^2d(x,e)^2+a^2d(y,e)^2\}\|\dot D_e(x,y)\|
 &\le C,                                                  \label{eq:newS7V12-D-one-moment}\\
 \sup_e\sum_f\sum_{x,y}
 \{1+a^2\tau(\{x,y,e,f\})^2\}\|\ddot D_{ef}(x,y)\|
 &\le Ca.                                                 \label{eq:newS7V12-D-two-moment}
\end{align}
\end{corollary}

\begin{proof}
Divide \eqref{eq:newS7V12-sign-tree-one} and
\eqref{eq:newS7V12-sign-tree-two} by the overlap factor \(a\).  Exponential
tree sums on a fixed-dimensional lattice are finite uniformly in the
chosen root link.  Multiplication by a squared physical distance contributes
\(a^2n^2\), whose supremum against \(e^{-\mu n}\) is bounded.  This proves
the moment statements.
\end{proof}

\begin{assessment}[Counting convention]
\label{ass:newS7V12-counting}
The sums in \eqref{eq:newS7V12-D-one-moment} and
\eqref{eq:newS7V12-D-two-moment} use the dimensionless orthonormal lattice
basis.  When the vertices are paired with smooth link fields and converted
to physical integral kernels, the standard factors \(a^4\) from site sums
must be inserted once.  The estimates state the invariant content: the
first physical overlap vertex is \(O(1)\), the second is \(O(a)\), and both
are localized on the physical scale \(a\).
\end{assessment}

For \(L=D_{\rm ov}^*D_{\rm ov}\), let \(V_e,W_{ef}\) be the V9 first and
second vertices.

\begin{theorem}[Overlap-square vertex tree bounds]
\label{thm:newS7V12-square-vertices}
There are \(C,\mu>0\) such that
\begin{align}
 \|V_e(x,y)\|
 &\le {C\over a}e^{-\mu\tau(\{x,y,e\})},                 \label{eq:newS7V12-V-tree}\\
 \|W_{ef}(x,y)\|
 &\le C e^{-\mu\tau(\{x,y,e,f\})}.                        \label{eq:newS7V12-W-tree}
\end{align}
Their physical second spatial moments carry an additional factor \(a^2\)
relative to their zeroth moments.
\end{theorem}

\begin{proof}
The unvaried overlap kernel obeys
\(\|D_{\rm ov}(x,y)\|\le Ca^{-1}e^{-\mu d(x,y)}\), by the same contour
argument without an insertion.  Insert this and
\eqref{eq:newS7V12-D-one}--\eqref{eq:newS7V12-D-two} into the exact product
rules \eqref{eq:newS7V9-V}--\eqref{eq:newS7V9-W}.  Convolution of
exponential trees preserves exponential tree decay after reducing
\(\mu\).  The amplitudes are \(a^{-1}\) for \(V_e\) and one for both
\(\dot D_e^*\dot D_f\) and \(D_{\rm ov}^*\ddot D_{ef}\).  The moment
statement follows as in
\Cref{cor:newS7V12-overlap-vertices}.
\end{proof}

\subsection{Combination with the free heat scale}
\label{sec:newS7V12-combination}

\begin{corollary}[Every unassembled polarization term has a finite second
moment]
\label{cor:newS7V12-term-moments}
At every \(t\ge a^2\), the separated and contact kernels in the V9
polarization each have a volume-uniform finite second spatial moment.  The
only physical spatial scales entering those moments are \(a\) from
\Cref{thm:newS7V12-square-vertices} and \(\sqrt t\) from V11.
\end{corollary}

\begin{proof}
Use the tree bounds
\eqref{eq:newS7V12-V-tree}--\eqref{eq:newS7V12-W-tree}, the free heat
moments \eqref{eq:newS7V11-m0}--\eqref{eq:newS7V11-m2}, and the convolution
algebra \eqref{eq:newS7V11-conv0}--\eqref{eq:newS7V11-conv2} in each
Duhamel term.
\end{proof}

\begin{proposition}[The amplitude cancellation remains essential]
\label{prop:newS7V12-amplitude}
The vertex moment bounds alone retain the amplitudes
\[
 V_e=O(a^{-1}),\qquad W_{ef}=O(1).
\]
In the separated Duhamel term the outside simplex factor \(t^2\) therefore
permits a crude contribution of order \(t^2/a^2\), while the contact term
has order \(t\).  These estimates do not yield the \(O(1+a^2/t)\) APMC
bound and cannot determine the V7 coefficient.
\end{proposition}

\begin{proof}
The amplitudes are
\eqref{eq:newS7V12-V-tree}--\eqref{eq:newS7V12-W-tree}; the time factors
are explicit in \eqref{eq:newS7V9-polarization}.  Bounding the two
structures separately discards the exact spectral cancellation of V10.
\end{proof}

\begin{definition}[Overlap vertex locality card, OVLC]
\label{def:newS7V12-OVLC}
OVLC consists of:
\begin{enumerate}[label=\textup{(\roman*)},leftmargin=3.5em]
\item the weighted resolvent bound \eqref{eq:newS7V12-CT};
\item the first and second sign derivatives
      \eqref{eq:newS7V12-first-sign}--\eqref{eq:newS7V12-second-sign};
\item the localized sign trees
      \eqref{eq:newS7V12-sign-tree-one}--\eqref{eq:newS7V12-sign-tree-two};
\item the physical overlap vertices
      \eqref{eq:newS7V12-D-one}--\eqref{eq:newS7V12-D-two};
\item their weighted moments
      \eqref{eq:newS7V12-D-one-moment}--\eqref{eq:newS7V12-D-two-moment};
      and
\item the overlap-square trees
      \eqref{eq:newS7V12-V-tree}--\eqref{eq:newS7V12-W-tree}.
\end{enumerate}
\end{definition}

\begin{programme}[V13 acquisition]
\label{prog:newS7V12-V13}
V13 will return to the assembled Brillouin formula
\eqref{eq:newS7V10-Brillouin}.  It will use the Ward--Takahashi identity to
subtract the \(q=0\) integrand before taking two momentum derivatives,
combine the derivative of \(k_t\), the two first vertices, and the contact
vertex algebraically, and only then apply the V8 Gaussian sums.  The target
is an explicit uniform bound for
\(\mathcal Q(a,t)\), with the coefficient limit left separate.  If
individual momentum derivatives create powers \(t/a^2\), they will be
reassembled through the pure-gauge identity of V10 before estimation.
\end{programme}

\begin{firewall}[Claim boundary after seventh-series V12]
\label{fire:newS7V12-boundary}
V12 proves a uniform weighted resolvent bound for the dimensionless Wilson
kernel, first and second holomorphic sign-derivative formulae, exponential
tree localization for localized link insertions, their exact physical
scales after overlap normalization, and corresponding first and second
overlap-square vertex bounds.  Combined with V11, every individual
finite-cutoff polarization term has a finite volume-uniform spatial
moment.

V12 does not prove cancellation of the leading separated/contact
amplitudes, the uniform assembled APMC moment, the V7 coefficient limit,
interacting MATC, the chiral determinant phase, cutoff removal, or the
Standard Model--Einstein continuum.
\end{firewall}

\begin{theorem}[Seventh-series V12 result]
\label{thm:newS7V12-result}
The AOP Wilson gap controls the full localized vertex calculus: one
physical gauge derivative of \(D_{\rm ov}\) is \(O(1)\), two derivatives
are \(O(a)\), and their kernels decay exponentially with the shortest tree
joining all fermion and link insertions.  The induced overlap-square
vertices have sizes \(O(a^{-1})\) and \(O(1)\).  Their moments combine with
the V11 heat scale, but separate absolute estimates still lose the exact
contact cancellation.  The remaining APMC step is therefore an assembled
momentum-space subtraction, not another locality estimate.
\end{theorem}

\begin{proof}
Use \Cref{lem:newS7V12-CT,
lem:newS7V12-sign-derivatives,
thm:newS7V12-sign-tree,
cor:newS7V12-overlap-vertices,
thm:newS7V12-square-vertices,
cor:newS7V12-term-moments,
prop:newS7V12-amplitude}.
\end{proof}

\paragraph{Parent-version record.}
V12 was formed from
\texttt{Renewal\_SM\_Einstein\_Continuum\_Research\_Notes\_V11.tex}.
The V11 parent had \(4\,893\) logical source lines, \(191\,083\) bytes, and
SHA--256
\[
 \texttt{B6140B9259BC0E156E408E3F57FD6A597497DF1185B87D4BA9E6192743A5502B}.
\]
The next compulsory companion audit is V15.

% =============================================================================
% END APPENDED SEVENTH-SERIES V12 MATERIAL
% =============================================================================

% =============================================================================
% BEGIN APPENDED SEVENTH-SERIES V13 MATERIAL
% =============================================================================

\clearpage
\section{Seventh-series V13: analytic Ward homotopy and exact curvature
factorization of the local polarization}
\label{sec:v7v13}

V10--V12 isolate the complete local polarization, its zero-winding
subtraction, and the weighted heat and vertex calculus.  A direct absolute
estimate still obscures the separated/contact cancellation.  This version
moves upstream: it proves that any analytic symmetric tensor satisfying the
two Ward identities factors through the linearized curvature twice.  The
factorization is exact near zero momentum and requires no division by
\(q^2\).

\subsection{A radial homotopy for transverse one-forms}
\label{sec:newS7V13-one-homotopy}

Let \(B_r\subset\mathbb R^4\) be a ball centered at the origin.

\begin{lemma}[Transverse one-form homotopy]
\label{lem:newS7V13-one-homotopy}
Let \(f_\mu\in C^1(B_r)\) satisfy
\begin{equation}
 k_\mu f_\mu(k)=0,\qquad f_\mu(0)=0.                      \label{eq:newS7V13-one-transverse}
\end{equation}
Define
\begin{equation}
 g_{\rho\mu}(k)
 :=\int_0^1 t\{
 \partial_\rho f_\mu(tk)-\partial_\mu f_\rho(tk)\}\,dt.   \label{eq:newS7V13-one-homotopy}
\end{equation}
Then \(g_{\rho\mu}=-g_{\mu\rho}\) and
\begin{equation}
 f_\mu(k)=k_\rho g_{\rho\mu}(k).                          \label{eq:newS7V13-one-factor}
\end{equation}
If \(f\) is analytic, so is \(g\).
\end{lemma}

\begin{proof}
Differentiate the identity
\(u_\rho f_\rho(u)=0\) with respect to \(u_\mu\):
\[
 f_\mu(u)+u_\rho\partial_\mu f_\rho(u)=0.
\]
At \(u=tk\), this gives
\[
 k_\rho\partial_\mu f_\rho(tk)=-t^{-1}f_\mu(tk).
\]
Also
\[
 k_\rho\partial_\rho f_\mu(tk)
 ={d\over dt}f_\mu(tk).
\]
Contract \eqref{eq:newS7V13-one-homotopy} with \(k_\rho\).  The integrand
becomes
\[
 t{d\over dt}f_\mu(tk)+f_\mu(tk)
 ={d\over dt}\{t f_\mu(tk)\}.
\]
Integration from zero to one and \(f(0)=0\) prove
\eqref{eq:newS7V13-one-factor}.  Antisymmetry is immediate.  Analyticity
follows by integrating the locally uniformly convergent Taylor series.
\end{proof}

\begin{corollary}[No inverse momentum is introduced]
\label{cor:newS7V13-no-inverse}
The factor \(g\) in \eqref{eq:newS7V13-one-factor} is bounded by first
derivatives of \(f\):
\begin{equation}
 \sup_{B_r}|g|
 \le \sup_{B_r}|\nabla f|.                               \label{eq:newS7V13-one-bound}
\end{equation}
In particular the construction is regular at \(k=0\).
\end{corollary}

\begin{proof}
Use \eqref{eq:newS7V13-one-homotopy} and
\(\int_0^1 2t\,dt=1\), with the component norm adjusted by a fixed
dimension-dependent constant.
\end{proof}

\subsection{Double homotopy for a Ward tensor}
\label{sec:newS7V13-double}

\begin{theorem}[Analytic double-curvature factorization]
\label{thm:newS7V13-double-factor}
Let \(P_{\mu\nu}(k)\) be a \(C^2\) matrix-valued tensor on \(B_r\) such
that
\begin{align}
 P_{\mu\nu}(k)&=P_{\nu\mu}(k),                            \label{eq:newS7V13-symmetric}\\
 k_\mu P_{\mu\nu}(k)&=0,
 &P_{\mu\nu}(k)k_\nu&=0,                                 \label{eq:newS7V13-two-Ward}\\
 P_{\mu\nu}(0)&=0,
 &\partial_\rho P_{\mu\nu}(0)&=0.                        \label{eq:newS7V13-two-zero}
\end{align}
Then there is a \(C^0\) tensor
\(R_{\rho\mu,\sigma\nu}(k)\), antisymmetric in
\((\rho,\mu)\) and in \((\sigma,\nu)\), such that
\begin{equation}
 P_{\mu\nu}(k)
 =k_\rho k_\sigma
 R_{\rho\mu,\sigma\nu}(k).                               \label{eq:newS7V13-double-factor}
\end{equation}
It may be chosen symmetric under exchange of the two antisymmetric index
pairs.  If \(P\) is analytic, then \(R\) is analytic.  Moreover
\begin{equation}
 \sup_{B_r}|R|
 \le C\sup_{B_r}|\nabla^2P|.                             \label{eq:newS7V13-R-bound}
\end{equation}
\end{theorem}

\begin{proof}
Apply \Cref{lem:newS7V13-one-homotopy} to the second index of \(P\):
\begin{equation}
 H_{\mu,\sigma\nu}(k)
 :=\int_0^1t\{
 \partial_\sigma P_{\mu\nu}(tk)
 -\partial_\nu P_{\mu\sigma}(tk)\}\,dt.                  \label{eq:newS7V13-H}
\end{equation}
Then \(H_{\mu,\sigma\nu}=-H_{\mu,\nu\sigma}\) and
\[
 P_{\mu\nu}=k_\sigma H_{\mu,\sigma\nu}.
\]
The first-index Ward identity and symmetry imply that \(H\) is transverse
in its first index.  Indeed, differentiating
\(u_\mu P_{\mu\nu}(u)=0\) gives
\[
 k_\mu\partial_\sigma P_{\mu\nu}(tk)
 =-t^{-1}P_{\sigma\nu}(tk),
\]
and the two terms in \(k_\mu H_{\mu,\sigma\nu}\) cancel because
\(P_{\sigma\nu}=P_{\nu\sigma}\).
Condition \eqref{eq:newS7V13-two-zero} also gives \(H(0)=0\).

Apply the one-form homotopy once more:
\begin{equation}
 R_{\rho\mu,\sigma\nu}(k)
 :=\int_0^1s\{
 \partial_\rho H_{\mu,\sigma\nu}(sk)
 -\partial_\mu H_{\rho,\sigma\nu}(sk)\}\,ds.             \label{eq:newS7V13-R}
\end{equation}
Then \(H_{\mu,\sigma\nu}=k_\rho R_{\rho\mu,\sigma\nu}\), proving
\eqref{eq:newS7V13-double-factor}.  Antisymmetry in both pairs follows from
\eqref{eq:newS7V13-H} and \eqref{eq:newS7V13-R}.  Average \(R\) with its
pair-exchanged transpose; this leaves \eqref{eq:newS7V13-double-factor}
unchanged because \(P\) is symmetric.

The two integrals contain only second derivatives of \(P\), with bounded
weights \(s,t\); this proves \eqref{eq:newS7V13-R-bound}.  Analyticity is
preserved at each homotopy step.
\end{proof}

\begin{remark}[Real gauge Hessians]
\label{rem:newS7V13-real}
The theorem is stated on the real gauge-component space.  For the
quadratic heat functional one has
\(P_{\mu\nu}^{AB}=P_{\nu\mu}^{BA}\).  Complex Fourier coefficients are
obtained afterward by complexification and the usual conjugate pairing.
\end{remark}

\subsection{Factorization of the zero-winding overlap polarization}
\label{sec:newS7V13-overlap}

Work on the infinite lattice, or equivalently with the zero-winding part
of the finite-torus kernel from V11.  Use
\begin{equation}
 k_\mu=\widehat q_\mu={2\over a}\sin{aq_\mu\over2}.        \label{eq:newS7V13-khat}
\end{equation}
For \(|aq_\mu|<\pi\), the map \(q\mapsto k\) is an analytic local
diffeomorphism at the origin.

\begin{lemma}[Analyticity of the local polarization]
\label{lem:newS7V13-polarization-analytic}
For fixed \(t\ge a^2\), the zero-winding polarization
\(\widehat\Pi^{\,0}_{\mu\nu}(q;a,t)\) is analytic in \(k\) near zero.
It is even under simultaneous momentum reflection, obeys both Ward
identities, and satisfies
\begin{equation}
 \widehat\Pi^{\,0}_{\mu\nu}(0;a,t)=0,\qquad
 \partial_{k_\rho}\widehat\Pi^{\,0}_{\mu\nu}(0;a,t)=0.    \label{eq:newS7V13-local-zero}
\end{equation}
\end{lemma}

\begin{proof}
The free quantity \(\omega(\theta)\) of V8 is bounded away from zero on the
real Brillouin torus.  By compactness it remains nonzero in a fixed complex
neighborhood.  The free overlap symbol is therefore analytic there.  The
V12 contour formula makes its first and second link vertices analytic in
the external momentum, and the divided difference \(k_t\) is entire in its
two spectral arguments.  The zero-winding Brillouin integral consequently
preserves analyticity near \(q=0\).

The V9 Ward identity supplies transversality.  Reflection gives evenness
and the first-derivative identity.  For a constant commuting potential, the
infinite-lattice Brillouin integral is invariant under the momentum
translation \(p\mapsto p+h\).  V11 identified this integral as exactly the
zero-winding term.  Its constant-field Hessian is therefore zero, proving
the first identity in \eqref{eq:newS7V13-local-zero}.
\end{proof}

\begin{theorem}[Exact local lattice-curvature owner]
\label{thm:newS7V13-overlap-factor}
There is an analytic tensor
\(\mathcal R_{\rho\mu,\sigma\nu}^{AB}(k;a,t)\), antisymmetric in each
spacetime index pair and symmetric under exchange
\((\rho,\mu,A)\leftrightarrow(\sigma,\nu,B)\), such that
\begin{equation}
 \widehat\Pi_{\mu\nu}^{0,AB}(q;a,t)
 =\widehat q_\rho\widehat q_\sigma
 \mathcal R_{\rho\mu,\sigma\nu}^{AB}(\widehat q;a,t).     \label{eq:newS7V13-overlap-factor}
\end{equation}
For a gauge potential \(\alpha\), define its linearized lattice curvature
\begin{equation}
 \widehat{\mathcal F}_{\rho\mu}^A(q)
 :=\widehat q_\rho\widehat\alpha_\mu^A(q)
   -\widehat q_\mu\widehat\alpha_\rho^A(q).               \label{eq:newS7V13-linear-curvature}
\end{equation}
Then the zero-winding quadratic heat response is exactly
\begin{equation}
 \overline{\widehat\alpha_\mu^A}
 \widehat\Pi_{\mu\nu}^{0,AB}
 \widehat\alpha_\nu^B
 ={1\over4}
 \overline{\widehat{\mathcal F}_{\rho\mu}^A}
 \mathcal R_{\rho\mu,\sigma\nu}^{AB}
 \widehat{\mathcal F}_{\sigma\nu}^B.                     \label{eq:newS7V13-curvature-quadratic}
\end{equation}
\end{theorem}

\begin{proof}
Use \(k=\widehat q\) and apply
\Cref{thm:newS7V13-double-factor} to
\Cref{lem:newS7V13-polarization-analytic}.  This gives
\eqref{eq:newS7V13-overlap-factor}.  Antisymmetry of \(\mathcal R\) gives
\[
 \widehat q_\rho\widehat\alpha_\mu
 \mathcal R_{\rho\mu,\sigma\nu}
 \widehat q_\sigma\widehat\alpha_\nu
 ={1\over4}
 \widehat{\mathcal F}_{\rho\mu}
 \mathcal R_{\rho\mu,\sigma\nu}
 \widehat{\mathcal F}_{\sigma\nu},
\]
which is \eqref{eq:newS7V13-curvature-quadratic}, with conjugation inserted
in the real Fourier pairing.
\end{proof}

\begin{corollary}[The local mass channel is absent exactly]
\label{cor:newS7V13-no-mass}
The zero-winding overlap polarization contains two explicit lattice
derivatives and therefore has no local gauge-boson mass tensor at any
finite cutoff.  The finite-torus constant-field Hessian is the separate
nonzero-winding term bounded in V11.
\end{corollary}

\begin{proof}
Set \(\widehat q=0\) in
\eqref{eq:newS7V13-overlap-factor}.  The finite-volume qualification is
\Cref{thm:newS7V11-holonomy-bound}.
\end{proof}

\subsection{Hypercubic reduction at zero momentum}
\label{sec:newS7V13-hypercubic}

\begin{theorem}[Unique parity-even curvature tensor]
\label{thm:newS7V13-unique-tensor}
For the reflection-even overlap modulus and one simple gauge factor,
hypercubic symmetry and constant gauge covariance give
\begin{align}
 \mathcal R_{\rho\mu,\sigma\nu}^{AB}(0;a,t)
 &=c_{a,t}^{(R)}\delta^{AB}
 \{\delta_{\rho\sigma}\delta_{\mu\nu}
   -\delta_{\rho\nu}\delta_{\mu\sigma}\}.                 \label{eq:newS7V13-unique-tensor}
\end{align}
For \(U(1)\), \(\delta^{AB}\) is omitted.  A parity-odd chiral trace may
also contain an \(\varepsilon_{\rho\mu\sigma\nu}\) term, but it is absent
from the modulus.
\end{theorem}

\begin{proof}
An invariant gauge bilinear is proportional to the chosen invariant
quadratic form on a simple factor.  A hypercubic-invariant rank-four tensor
antisymmetric in each pair and symmetric under pair exchange is a linear
combination of the displayed metric wedge and the orientation tensor.
Link reflection removes the orientation tensor from a parity-even
functional.
\end{proof}

\begin{corollary}[Scalar coefficient and the \(F^2\) normalization]
\label{cor:newS7V13-scalar}
At quadratic order near zero momentum,
\begin{align}
 \widehat\Pi_{\mu\nu}^{0,AB}(q;a,t)
 &=c_{a,t}^{(R)}\delta^{AB}
 \{\widehat q^{\,2}\delta_{\mu\nu}
 -\widehat q_\mu\widehat q_\nu\}
 +O(|\widehat q|^4),                                     \label{eq:newS7V13-scalar}\\
 {1\over4}\widehat{\mathcal F}\mathcal R(0)
 \widehat{\mathcal F}
 &={c_{a,t}^{(R)}\over2}
 \widehat{\mathcal F}_{\mu\nu}^A
 \widehat{\mathcal F}_{\mu\nu}^A.                        \label{eq:newS7V13-F2-normalization}
\end{align}
Thus convergence of the single scalar \(c_{a,t}^{(R)}\) is equivalent to
convergence of the parity-even local curvature owner.
\end{corollary}

\begin{proof}
Contract \eqref{eq:newS7V13-unique-tensor} with two momenta for the first
identity and with two antisymmetric curvatures for the second.
\end{proof}

\subsection{A bounded homotopy criterion}
\label{sec:newS7V13-criterion}

\begin{proposition}[Signed second derivatives suffice]
\label{prop:newS7V13-signed-criterion}
Let
\[
 B_{a,t}(r)
 :=\sup_{|\widehat q|\le r}
 \left|\nabla_{\widehat q}^{\,2}
 \widehat\Pi^{\,0}(\widehat q;a,t)\right|.
\]
Then
\begin{equation}
 \sup_{|\widehat q|\le r}
 |\mathcal R(\widehat q;a,t)|
 \le C B_{a,t}(r).                                       \label{eq:newS7V13-signed-criterion}
\end{equation}
This criterion uses the signed, already assembled polarization and is
strictly weaker than an absolute position-kernel moment bound.
\end{proposition}

\begin{proof}
Apply \eqref{eq:newS7V13-R-bound} to the local polarization.
An absolute second moment implies the displayed derivative bound by
termwise Fourier differentiation, but cancellations may make
\(B_{a,t}\) finite even when that stronger absolute estimate is wasteful.
\end{proof}

\begin{definition}[Curvature factorization card, CFC]
\label{def:newS7V13-CFC}
CFC consists of:
\begin{enumerate}[label=\textup{(\roman*)},leftmargin=3.5em]
\item zero-winding subtraction from V11;
\item analytic local polarization and the exact Ward identities;
\item the double radial homotopy
      \eqref{eq:newS7V13-H}--\eqref{eq:newS7V13-R};
\item exact curvature factorization
      \eqref{eq:newS7V13-curvature-quadratic};
\item hypercubic reduction to \(c_{a,t}^{(R)}\); and
\item a uniform signed derivative bound and convergence of
      \(c_{a,t}^{(R)}\) to the V7 coefficient.
\end{enumerate}
\end{definition}

\begin{programme}[V14 acquisition]
\label{prog:newS7V13-V14}
V14 will target the last CFC row directly in the Brillouin formula.  It
will differentiate the complete divided-difference integrand twice in
\(\widehat q\), retain the separated and contact pieces in one expression,
and use the V8 low/high momentum split.  The high branch should be
exponentially small for \(t/a^2\to\infty\); the low branch will be rescaled
by \(p=k/\sqrt t\) and compared with the continuum Dirac polarization.
The immediate goal is a uniform bound for \(B_{a,t}(r)\); numerical
coefficient convergence will be asserted only if the rescaled integrand
has an integrable common majorant.
\end{programme}

\begin{firewall}[Claim boundary after seventh-series V13]
\label{fire:newS7V13-boundary}
V13 proves a two-stage radial homotopy for symmetric transverse tensors,
applies it to the analytic zero-winding overlap polarization, factors that
polarization exactly through two linearized lattice curvatures, removes the
local mass channel at finite cutoff, and reduces the parity-even
hypercubic tensor to one scalar coefficient.

V13 does not bound or compute that scalar uniformly, prove its V7 limit,
control nonzero background gauge fields, establish interacting MATC,
construct the chiral determinant phase, remove the cutoff, or construct the
Standard Model--Einstein continuum.
\end{firewall}

\begin{theorem}[Seventh-series V13 result]
\label{thm:newS7V13-result}
After the V11 winding subtraction, the complete finite-cutoff overlap heat
polarization admits the exact factorization
\(\widehat\Pi^0=d_1^*\mathcal R\,d_1\), obtained without division by
\(\widehat q^2\).  Its parity-even zero-momentum value is one scalar
\(c_{a,t}^{(R)}\) multiplying the lattice \(F^2\) tensor.  The strong
absolute moment target of V10 may therefore be replaced, for coefficient
extraction, by the signed assembled derivative bound
\eqref{eq:newS7V13-signed-criterion}.  The remaining task is the
low/high Brillouin comparison of that scalar with V7.
\end{theorem}

\begin{proof}
Use \Cref{lem:newS7V13-one-homotopy,
thm:newS7V13-double-factor,
lem:newS7V13-polarization-analytic,
thm:newS7V13-overlap-factor,
cor:newS7V13-no-mass,
thm:newS7V13-unique-tensor,
prop:newS7V13-signed-criterion}.
\end{proof}

\paragraph{Parent-version record.}
V13 was formed from
\texttt{Renewal\_SM\_Einstein\_Continuum\_Research\_Notes\_V12.tex}.
The V12 parent had \(5\,286\) logical source lines, \(206\,976\) bytes, and
SHA--256
\[
 \texttt{4DCEACD3B8B3DBAE5E86EE01751DC63964535BC10BBB5463CB4650BDBB8BE72A}.
\]
The next compulsory companion audit is V15.

% =============================================================================
% END APPENDED SEVENTH-SERIES V13 MATERIAL
% =============================================================================

% =============================================================================
% BEGIN APPENDED SEVENTH-SERIES V14 MATERIAL
% =============================================================================

\clearpage
\section{Seventh-series V14: mesoscopic convergence of the free-background
overlap curvature coefficient}
\label{sec:v7v14}

V13 factors the local two-point heat response exactly through the
linearized lattice curvature and reduces its parity-even part to one
scalar \(c_{a,t}^{(R)}\).  This version compares that scalar with the
continuum Dirac-square coefficient of V7.  The result concerns the
zero-winding, free-background, vectorlike overlap modulus; extension to a
general AOP background and to a chiral determinant remains separate.

\subsection{Continuum and lattice two-point densities}
\label{sec:newS7V14-densities}

Let \(\pi_{a,t}^{AB}{}_{\mu\nu}(q)\) be the zero-winding polarization per
unit physical volume, obtained from the Brillouin integral version of
\eqref{eq:newS7V10-Brillouin}.  Thus
\begin{equation}
 \pi_{a,t}^{AB}{}_{\mu\nu}(q)
 =c_{a,t}^{(R)}\delta^{AB}
 \{\widehat q^{\,2}\delta_{\mu\nu}
  -\widehat q_\mu\widehat q_\nu\}
 +O(|\widehat q|^4).                                     \label{eq:newS7V14-lattice-c}
\end{equation}
Let \(\pi_{D,t}\) denote the corresponding heat polarization density for
the flat continuum Dirac square in representation \(R\).

\begin{lemma}[Continuum normalization of the scalar]
\label{lem:newS7V14-continuum-c}
With the V7 heat and quadratic-index conventions,
\begin{equation}
 c_D^{(R)}
 ={8I_2(R)\over3(4\pi)^2}                                \label{eq:newS7V14-continuum-c}
\end{equation}
for one complex four-component Dirac square.  The parity-even Weyl value is
\begin{equation}
 c_{D,\mathrm W}^{(R)}
 ={4I_2(R)\over3(4\pi)^2}.                               \label{eq:newS7V14-continuum-Weyl}
\end{equation}
\end{lemma}

\begin{proof}
V7 gives the Dirac heat coefficient
\[
 (4\pi)^{-2}{2I_2(R)\over3}
 \int F_{\mu\nu}^AF_{\mu\nu}^A.
\]
For \(A=s\alpha\), its second \(s\)-derivative is twice this quadratic
functional.  In Fourier variables,
\[
 F_{\mu\nu}^AF_{\mu\nu}^A
 =2\overline{\alpha_\mu^A}
 (q^2\delta_{\mu\nu}-q_\mu q_\nu)\alpha_\nu^A.
\]
The Hessian coefficient is therefore four times the coefficient of
\(\int F^2\), proving \eqref{eq:newS7V14-continuum-c}.  A Weyl heat trace
has one half of the parity-even spin trace.
\end{proof}

\subsection{Physical-branch consistency of overlap vertices}
\label{sec:newS7V14-consistency}

Let \(\mathcal D_a(p)\) be the free overlap symbol and let
\(\mathcal V_{a,\mu}^A(p,q)\) and
\(\mathcal W_{a,\mu\nu}^{AB}(p;q,-q)\) be the first and contact vertices of
the overlap square from V10.  Add a subscript \(D\) for the corresponding
continuum Dirac-square vertices.

\begin{lemma}[Low-momentum vertex consistency]
\label{lem:newS7V14-vertex-consistency}
There are \(\delta,C>0\) and an integer \(M\) such that, for
\(|ap|+|aq|\le\delta\), the following hold through two derivatives in the
external momentum \(q\):
\begin{align}
 \mathcal D_a(p)
 &=i\gamma^\mu p_\mu+O(a|p|^2),                           \label{eq:newS7V14-D-consistency}\\
 \partial_q^\beta
 \{\mathcal V_{a,\mu}^A(p,q)-\mathcal V_{D,\mu}^A(p,q)\}
 &\le Ca(1+|p|+|q|)^M,                                   \label{eq:newS7V14-V-consistency}\\
 \partial_q^\beta
 \{\mathcal W_{a,\mu\nu}^{AB}(p;q,-q)
       -\mathcal W_{D,\mu\nu}^{AB}(p;q,-q)\}
 &\le Ca(1+|p|+|q|)^M,                                   \label{eq:newS7V14-W-consistency}
\end{align}
for \(|\beta|\le2\), after inserting the fixed powers of a reference length
needed to make the displayed component norms dimensionally homogeneous.
More precisely, at the heat scale \(p=k/\sqrt t\), each relative error in
the dimensionless assembled integrand is bounded by
\begin{equation}
 C{a\over\sqrt t}(1+|k|)^M.                              \label{eq:newS7V14-scaled-error}
\end{equation}
\end{lemma}

\begin{proof}
The free Wilson symbol is
\[
 aA_W(ap)=i\gamma^\mu\sin(ap_\mu)
 +\sum_\mu\{1-\cos(ap_\mu)\}-1.
\]
Its Taylor expansion is
\[
 -I+ia\gamma^\mu p_\mu+{a^2\over2}|p|^2+O(a^3|p|^3).
\]
The positive square root in the overlap formula is analytic near this
physical branch and equals \(I+O(a^2|p|^2)\).  Substitution gives
\eqref{eq:newS7V14-D-consistency}.

A link insertion replaces one shift factor by its derivative
\(aT^A\), and a second insertion replaces it by \(a^2T^AT^B\).  Differentiate
the holomorphic sign formula of V12 before taking the Taylor expansion.
The fixed Wilson gap makes the contour resolvents and their first two
momentum derivatives uniformly bounded.  The leading first insertion is
the continuum covariant derivative vertex; the leading second insertion
is the continuum contact vertex.  Every omitted monomial contains at least
one additional factor \(a(|p|+|q|)\).  Taylor's theorem on the compact ball
\(|ap|+|aq|\le\delta\) gives
\eqref{eq:newS7V14-V-consistency}--\eqref{eq:newS7V14-W-consistency}.
Putting \(p=k/\sqrt t\), with \(a/\sqrt t\le1\), gives
\eqref{eq:newS7V14-scaled-error}.
\end{proof}

\begin{remark}[Why a polynomial bound is enough]
\label{rem:newS7V14-polynomial}
The exponent \(M\) is immaterial.  The heat divided difference supplies a
Gaussian \(e^{-c|k|^2}\) after the rescaling \(p=k/\sqrt t\), so every fixed
polynomial is integrable.  No uniform Taylor expansion is asserted at
momenta of order \(a^{-1}\); those belong to the high branch below.
\end{remark}

\subsection{The twice-differentiated assembled integrand}
\label{sec:newS7V14-integrand}

Write the zero-winding Brillouin formula as
\begin{equation}
 \pi_{a,t}(q)
 =\int_{[-\pi/a,\pi/a]^4}{d^4p\over(2\pi)^4}
 \mathfrak I_a(p,q;t),                                   \label{eq:newS7V14-integrand}
\end{equation}
where \(\mathfrak I_a\) is the sum of the divided-difference two-vertex
term and the negative contact term.  Let \(\mathfrak I_D\) be the analogous
continuum integrand.

\begin{lemma}[Low-branch common majorant]
\label{lem:newS7V14-low-majorant}
Choose a smooth cutoff supported where \(|ap|<2\delta\) and equal to one
where \(|ap|\le\delta\).  For \(t\ge a^2\), after \(p=k/\sqrt t\),
\begin{align}
 t^{-2}
 \left|
 \partial_{q_\rho}\partial_{q_\sigma}
 \{\chi(ap)\mathfrak I_a(p,q;t)\}_{q=0}
 \right|
 &\le C(1+|k|)^{M_1}e^{-c|k|^2},                         \label{eq:newS7V14-low-majorant}\\
 t^{-2}
 \left|
 \partial_q^2
 \{\chi(ap)\mathfrak I_a-\mathfrak I_D\}_{q=0}
 \right|
 &\le C{a\over\sqrt t}
 (1+|k|)^{M_1}e^{-c|k|^2}.                               \label{eq:newS7V14-low-difference}
\end{align}
The constants are independent of \(a,t\) on a fixed mesoscopic window.
\end{lemma}

\begin{proof}
Use the integral definition
\[
 k_t(x,y)=t^2\int_0^1e^{-t\{sx+(1-s)y\}}\,ds.
\]
On the low branch, V8 coercivity gives
\(\lambda_a(ap)\ge c|p|^2\), also after a sufficiently small external
momentum shift.  Each \(q\)-derivative of a spectral argument produces at
most a polynomial in \(p\); each derivative of \(k_t\) produces one
additional power of \(t\).  The first overlap-square vertex is at most
linear in \(p\) on this branch and the contact vertex is bounded.

At \(p=k/\sqrt t\), the terms with two external derivatives all have the
common net factor \(t^2\): for example,
\(\partial_x^2k_t\) contributes \(t^4\), its two spectral derivatives
contribute \(t^{-1}\), and the two first vertices contribute another
\(t^{-1}\).  Other derivative placements have the same or better power.
What remains is a polynomial in \(k\) times \(e^{-c|k|^2}\), proving
\eqref{eq:newS7V14-low-majorant}.  Replace one lattice factor at a time by
its continuum counterpart and use
\eqref{eq:newS7V14-scaled-error}; a finite telescoping sum proves
\eqref{eq:newS7V14-low-difference}.
\end{proof}

\begin{lemma}[High-branch suppression after two external derivatives]
\label{lem:newS7V14-high}
There are \(C,c>0\) such that
\begin{equation}
 \left|
 \int{d^4p\over(2\pi)^4}
 \partial_q^2\{(1-\chi(ap))\mathfrak I_a(p,q;t)\}_{q=0}
 \right|
 \le C e^{-ct/a^2}                                       \label{eq:newS7V14-high}
\end{equation}
for \(t\ge a^2\), after a polynomial in \(t/a^2\) has been absorbed by
decreasing \(c\).
\end{lemma}

\begin{proof}
On the support of \(1-\chi\), either momentum entering a heat factor lies a
fixed distance from the physical corner.  V8 gives a spectral floor
\(c/a^2\).  V12 bounds the first two vertex derivatives by fixed powers of
\(a^{-1}\), while
\eqref{eq:newS7V10-kt-derivatives} bounds every divided-difference
derivative by a polynomial in \(t\) times \(e^{-ct/a^2}\).  The Brillouin
volume is \(O(a^{-4})\).  All powers of \(a^{-1}\) and \(t\) combine into a
fixed polynomial in \(t/a^2\), which is absorbed into
\(e^{-c't/a^2}\).
\end{proof}

\subsection{Convergence of the curvature scalar}
\label{sec:newS7V14-convergence}

\begin{theorem}[Free-background mesoscopic coefficient convergence]
\label{thm:newS7V14-convergence}
For one complex vectorlike overlap fermion in representation \(R\), there
are \(C_R,c>0\) such that, whenever \(a^2/t\) is sufficiently small,
\begin{equation}
 \left|c_{a,t}^{(R)}-c_D^{(R)}\right|
 \le C_R\left\{
 {a\over\sqrt t}+e^{-ct/a^2}
 \right\},                                                \label{eq:newS7V14-convergence}
\end{equation}
where \(c_D^{(R)}\) is
\eqref{eq:newS7V14-continuum-c}.  In particular,
\begin{equation}
 \lim_{\substack{a\downarrow0,\ t\downarrow0\\a^2/t\to0}}
 c_{a,t}^{(R)}
 ={8I_2(R)\over3(4\pi)^2}.                               \label{eq:newS7V14-limit}
\end{equation}
\end{theorem}

\begin{proof}
By V13, \(c_{a,t}^{(R)}\) is a fixed contraction of
\(\partial_q^2\pi_{a,t}(0)\).  On the low branch, substitute
\(p=k/\sqrt t\) in \eqref{eq:newS7V14-integrand}.  The measure contributes
\(t^{-2}\), which cancels the \(t^2\) in
\eqref{eq:newS7V14-low-majorant}.  Dominated convergence applies, and
\eqref{eq:newS7V14-low-difference} bounds the error by
\(C_Ra/\sqrt t\).  The continuum integral is the flat Dirac-square heat
polarization, whose scalar is \eqref{eq:newS7V14-continuum-c}.
\Cref{lem:newS7V14-high} supplies the remaining exponential error.
\end{proof}

\begin{corollary}[One-generation vectorlike heat weights]
\label{cor:newS7V14-SM-weights}
If the Standard Model representations are temporarily paired into
four-component vectorlike test blocks, the one-generation limiting scalar
weights are
\begin{equation}
 {16\over3(4\pi)^2}\quad\hbox{for }SU(3)_c,\qquad
 {16\over3(4\pi)^2}\quad\hbox{for }SU(2)_L,\qquad
 {80\over9(4\pi)^2}\quad\hbox{for }U(1)_Y.               \label{eq:newS7V14-vectorlike-weights}
\end{equation}
These are the V6 indices \(2,2,10/3\) multiplied by the Dirac scalar
\(8/\{3(4\pi)^2\}\).  They are test-block values, not yet the chiral
Standard Model determinant.
\end{corollary}

\begin{proof}
Insert the V6 quadratic-index ledger into
\eqref{eq:newS7V14-continuum-c}.
\end{proof}

\subsection{Integrated mesoscopic error}
\label{sec:newS7V14-integrated}

Use the V8 endpoint
\[
 \varepsilon_q(a)=\ell^2(a/\ell)^q.
\]

\begin{theorem}[Logarithmic integrability of the coefficient error]
\label{thm:newS7V14-integrated}
For every \(0<q<2\) and fixed \(t_0>0\),
\begin{equation}
 \lim_{a\downarrow0}
 \int_{\varepsilon_q(a)}^{t_0}
 {\,|c_{a,t}^{(R)}-c_D^{(R)}|\over t}\,dt=0.             \label{eq:newS7V14-integrated}
\end{equation}
In particular the V8 choice \(0<q<2/3\), already required by the free trace
comparison, satisfies this condition.
\end{theorem}

\begin{proof}
The algebraic error in \eqref{eq:newS7V14-convergence} gives
\[
 Ca\int_{\varepsilon_q(a)}^{t_0}t^{-3/2}\,dt
 \le C{a\over\sqrt{\varepsilon_q(a)}}
 =C(a/\ell)^{1-q/2}\longrightarrow0.
\]
For the high-branch error set \(u=t/a^2\).  Its lower limit
\(\varepsilon_q(a)/a^2=(\ell/a)^{2-q}\) tends to infinity, and the
exponential integral tends to zero.
\end{proof}

\begin{corollary}[Perturbative gauge row of MATC]
\label{cor:newS7V14-perturbative-MATC}
At quadratic order about the free connection, the zero-winding vectorlike
overlap heat trace has the continuum Dirac \(F^2\) coefficient and an
integrable mesoscopic lattice error.  The finite-volume zero-momentum term
is removed by V11's all-orders winding estimate.
\end{corollary}

\begin{proof}
Combine \Cref{thm:newS7V14-convergence,
thm:newS7V14-integrated,thm:newS7V11-holonomy-bound}.
\end{proof}

\begin{assessment}[Relative subtraction remains exact]
\label{ass:newS7V14-relative}
If the full and covariant base blocks use the same gauge connection and
representation, V7's relative \(F^2\) cancellation occurs before the
limit.  V14 then compares two equal universal coefficients and their
difference is zero.  For bases with different connections,
\eqref{eq:newS7V14-limit} identifies the coefficient of
\(F^2-(F^0)^2\).  The absolute physical gauge coupling still requires the
owned normalization condition isolated in V7.
\end{assessment}

\begin{definition}[Free two-point transfer card, FTTC]
\label{def:newS7V14-FTTC}
FTTC consists of:
\begin{enumerate}[label=\textup{(\roman*)},leftmargin=3.5em]
\item exact zero-winding curvature factorization from V13;
\item low-momentum overlap and vertex consistency
      \eqref{eq:newS7V14-D-consistency}--\eqref{eq:newS7V14-scaled-error};
\item the common low-branch majorant
      \eqref{eq:newS7V14-low-majorant};
\item high-branch suppression \eqref{eq:newS7V14-high};
\item scalar convergence \eqref{eq:newS7V14-convergence}; and
\item logarithmic integrability \eqref{eq:newS7V14-integrated}.
\end{enumerate}
\end{definition}

\begin{programme}[V15 audit and acquisition]
\label{prog:newS7V14-V15}
V15 is a compulsory companion-audit version.  After snapshotting the
newest settled Yang--Mills and Navier--Stokes notes, it will extend FTTC
from the free background to a bounded smooth AOP family.  The first target
is a covariant Duhamel expansion in the background curvature with one
remainder, controlled in the weighted calculus of V11--V12.  The expansion
will be made at the assembled heat operator, not as independent sign
series.  If uniformity fails because the background varies on the lattice
scale, the admissible family will be separated into a smooth-link core and
a rough admissible remainder rather than silently equating the two.
\end{programme}

\begin{firewall}[Claim boundary after seventh-series V14]
\label{fire:newS7V14-boundary}
V14 proves low-momentum consistency of the free overlap operator and its
first two gauge vertices, constructs an integrable common majorant for the
twice-differentiated assembled heat integrand, suppresses the high
Brillouin branch, and proves convergence of the zero-winding vectorlike
curvature scalar to the continuum Dirac value with an integrable
mesoscopic error.

V14 does not prove the corresponding chiral moving-projector coefficient,
control a nonzero or rough AOP background, establish the full nonlinear
heat expansion, fix the absolute gauge coupling, construct the determinant
phase, remove the interacting cutoff, or construct the Standard
Model--Einstein continuum.
\end{firewall}

\begin{theorem}[Seventh-series V14 result]
\label{thm:newS7V14-result}
The free-background vectorlike overlap heat polarization reaches the V7
Dirac target:
\[
 c_{a,t}^{(R)}
 \longrightarrow {8I_2(R)\over3(4\pi)^2}
 \quad\hbox{when }a^2/t\to0.
\]
Its error is \(O(a/\sqrt t)\) plus an exponentially small doubler tail and
is integrable on the V8 shrinking proper-time window.  Together with V11
and V13, this closes the free zero-winding two-point \(F^2\) transfer.
The next step is background-uniform covariant expansion, preceded by the
V15 companion audit.
\end{theorem}

\begin{proof}
Use \Cref{lem:newS7V14-continuum-c,
lem:newS7V14-vertex-consistency,
lem:newS7V14-low-majorant,
lem:newS7V14-high,
thm:newS7V14-convergence,
thm:newS7V14-integrated,
cor:newS7V14-perturbative-MATC}.
\end{proof}

\paragraph{Parent-version record.}
V14 was formed from
\texttt{Renewal\_SM\_Einstein\_Continuum\_Research\_Notes\_V13.tex}.
The V13 parent had \(5\,710\) logical source lines, \(223\,523\) bytes, and
SHA--256
\[
 \texttt{96F8D3DCAAFFE9EE0225BACB23CA6254F6D25C635A93A1DD6E71F27674BF4FD5}.
\]
The next compulsory companion audit is V15.

% =============================================================================
% END APPENDED SEVENTH-SERIES V14 MATERIAL
% =============================================================================

% =============================================================================
% BEGIN APPENDED SEVENTH-SERIES V15 MATERIAL
% =============================================================================

\clearpage
\section{Seventh-series V15: companion audit, factorial-free background
insertions, and radial heat-scale gauges}
\label{sec:v7v15}

\subsection{Compulsory companion audit}
\label{sec:newS7V15-audit}

The companion files were read as immutable byte snapshots.  The newest
settled Yang--Mills note was
\begin{center}
\begin{tabular}{ll}
file &
\texttt{Renewal\_Yang\_Mills\_Mass\_Gap\_Research\_Notes\_V87.tex},\\
bytes & \(1\,277\,244\),\\
logical source lines & \(39\,681\),\\
SHA--256 &
\texttt{AC1148FA6B3FAB623FCFB66641B2C67CCB1D9A78FDF9D2CCDCB90D5CB71A8DE5}.
\end{tabular}
\end{center}
V87 proves a doubly polarized Duhamel majorant, an exact weighted-Cayley
extension formula for a prescribed forest, and domination of every forest
gate by the product of its one-edge endpoint probabilities.  Its
small-activity Touchard estimate collapses all repeated occurrences of one
high-arity support to a single activity with a factor two, independently
of multiplicity; differentiated pair supports cannot repeat.  The
remaining NACC step must convert normalized endpoint choices into raw
local anchors without spending the saturated tree reserve.  The
Yang--Mills mass gap remains unproved.

The newest settled Navier--Stokes note was
\begin{center}
\begin{tabular}{ll}
file &
\texttt{Renewal\_NS\_Stability\_Research\_Notes\_V67.tex},\\
bytes & \(861\,468\),\\
logical source lines & \(16\,308\),\\
SHA--256 &
\texttt{5AE2E187107B27048CA2AECC00851CD147CCDF5FC33B1F1554E7D65D087EAD59}.
\end{tabular}
\end{center}
V67 closes negatively the proposed common-active-scale argument: the
available minimal log-period is a lower bound, not an upper bound on gaps
between active scales, so distinct singular points need not share a
comparable active scale.  It proves instead a quiet-scale dichotomy.  At
each scale a singular point either lies near an energetic seed or its
nearby cells are trapped on that window with explicitly small
dissipation; the missing activity may be carried by later satellite
windows.  Full regularity remains unproved.

\begin{audit}[V15 transfer decision]
\label{audit:newS7V15-transfer}
The audit changes the background-field attack in two ways.
\begin{enumerate}[label=\textup{(\roman*)},leftmargin=3.5em]
\item Repeated gauge-background insertions must be normalized at the
      holomorphic expansion level.  Expanding an \(n\)-th Fr\'echet
      derivative into \(n!\) ordered resolvent words and bounding afterward
      would recreate the multiplicity loss removed in YM V87.
\item A smooth radial gauge on one heat-scale ball does not control every
      scale or every rough plaquette-admissible field.  The smooth-link
      family and the rough AOP remainder must remain separate, just as the
      NS quiet window cannot be promoted to aligned activity at all later
      scales.
\end{enumerate}
\end{audit}

\subsection{Why a direct overlap-square Duhamel expansion is not uniform}
\label{sec:newS7V15-failed-L}

Let \(A\) be a bounded smooth gauge potential in a fixed local
trivialization and let \(U(A)\) be its edge parallel transports.

\begin{lemma}[Sizes of a bounded smooth background perturbation]
\label{lem:newS7V15-sizes}
For \(a\|A\|_\infty\) sufficiently small,
\begin{align}
 \|\mathsf A\{U(A)\}-\mathsf A(1)\|
 &\le Ca\|A\|_\infty,                                    \label{eq:newS7V15-Wilson-size}\\
 \|D_{\rm ov}\{U(A)\}-D_{\rm ov}(1)\|
 &\le C\|A\|_\infty,                                     \label{eq:newS7V15-D-size}\\
 \|L_a(A)-L_a(0)\|
 &\le C\left\{{\|A\|_\infty\over a}
                 +\|A\|_\infty^2\right\}.                \label{eq:newS7V15-L-size}
\end{align}
\end{lemma}

\begin{proof}
Parallel transport along an edge solves an ODE of length \(a\), so
\(\|U_e(A)-I\|\le Ca\|A\|_\infty\).  The dimensionless Wilson kernel has a
uniform finite row sum, giving \eqref{eq:newS7V15-Wilson-size}.  The V3 or
V12 sign derivative bound and the factor \(a^{-1}\) in \(D_{\rm ov}\) give
\eqref{eq:newS7V15-D-size}.  Finally write
\[
 L_a(A)-L_a(0)
 =D_0^*\Delta D+\Delta D^*D_0+\Delta D^*\Delta D
\]
and use \(\|D_0\|\le2/a\).
\end{proof}

\begin{proposition}[The naive heat Duhamel parameter diverges]
\label{prop:newS7V15-naive-Duhamel}
A direct operator-norm Duhamel series for
\(e^{-tL_a(A)}\) about \(L_a(0)\) is controlled only by
\begin{equation}
 t\|L_a(A)-L_a(0)\|
 \le C\left\{{t\|A\|_\infty\over a}
              +t\|A\|_\infty^2\right\}.                 \label{eq:newS7V15-bad-parameter}
\end{equation}
On the V8 lower scale
\(t=\ell^2(a/\ell)^q\), \(0<q<2/3\), the first term grows like
\((a/\ell)^{q-1}\).  Hence this route cannot prove a uniform mesoscopic
background expansion.
\end{proposition}

\begin{proof}
Use \eqref{eq:newS7V15-L-size}.  Since \(q-1<0\), the stated factor
diverges as \(a\downarrow0\).  This is an estimate failure, not a physical
divergence: V10--V14 show that Ward and curvature assembly cancel the
corresponding free-background power.
\end{proof}

\subsection{Factorial-free holomorphic background expansion}
\label{sec:newS7V15-holomorphic}

Put
\[
 \mathsf A_0=\mathsf A(1),\qquad
 \Delta\mathsf A=\mathsf A\{U(A)\}-\mathsf A_0.
\]
Choose the fixed sign contours of V12 and write
\(R_0(z)=(z-\mathsf A_0)^{-1}\).

\begin{theorem}[Normalized resolvent-word expansion]
\label{thm:newS7V15-resolvent-series}
If
\begin{equation}
 \sup_{z\in\Gamma_+\cup\Gamma_-}
 \|\Delta\mathsf A\,R_0(z)\|\le\eta<1,                   \label{eq:newS7V15-eta}
\end{equation}
then
\begin{align}
 \sgn(\mathsf A_0+\Delta\mathsf A)
 -\sgn(\mathsf A_0)
 &=
 {1\over2\pi i}
 \sum_{n=1}^\infty
 \int_{\Gamma_+\cup\Gamma_-}
 \sgn(z)R_0(z)\{\Delta\mathsf A R_0(z)\}^n\,dz.          \label{eq:newS7V15-resolvent-series}
\end{align}
The series converges absolutely in operator norm and has remainder
\begin{equation}
 \|\mathcal E_{N+1}\|
 \le C_\delta{\eta^{N+1}\over1-\eta}.                    \label{eq:newS7V15-resolvent-remainder}
\end{equation}
There is one ordered word at each power \(n\); no factorial occurs.
\end{theorem}

\begin{proof}
On each contour,
\[
 (z-\mathsf A_0-\Delta\mathsf A)^{-1}
 =R_0(z)\sum_{n=0}^\infty
 \{\Delta\mathsf A R_0(z)\}^n
\]
by the Neumann series.  Insert it into the holomorphic sign calculus and
subtract the \(n=0\) term.  Uniform convergence permits termwise contour
integration.  Sum the geometric tail and use the fixed contour length and
resolvent bound to obtain
\eqref{eq:newS7V15-resolvent-remainder}.
\end{proof}

\begin{corollary}[Cancellation of the Fr\'echet factorial]
\label{cor:newS7V15-factorial}
For one repeated perturbation \(B\),
\begin{equation}
 {1\over n!}D^n\sgn(\mathsf A_0)[B,\ldots,B]
 ={1\over2\pi i}
 \int_{\Gamma_+\cup\Gamma_-}
 \sgn(z)R_0(z)\{BR_0(z)\}^n\,dz.                         \label{eq:newS7V15-factorial}
\end{equation}
Thus the Taylor normalization cancels the \(n!\) orderings before norms,
exactly as the repeated-support normalization in the YM V87 audit.
\end{corollary}

\begin{proof}
Insert \(\Delta\mathsf A=sB\) in
\eqref{eq:newS7V15-resolvent-series} and compare the coefficient of \(s^n\)
with the Fr\'echet Taylor formula.
\end{proof}

\begin{theorem}[Background expansion of the overlap operator]
\label{thm:newS7V15-D-expansion}
For every bounded smooth background and all sufficiently small \(a\),
\begin{equation}
 D_{\rm ov}\{U(A)\}
 =D_{\rm ov}(1)+\sum_{n=1}^N\mathcal D_{a,n}(A)
 +\mathcal R_{a,N+1}(A),                                 \label{eq:newS7V15-D-expansion}
\end{equation}
where
\begin{align}
 \|\mathcal D_{a,n}(A)\|
 &\le C_\delta^n a^{n-1}\|A\|_\infty^n,                  \label{eq:newS7V15-Dn}\\
 \|\mathcal R_{a,N+1}(A)\|
 &\le C_{\delta,A}a^N\|A\|_\infty^{N+1}.                 \label{eq:newS7V15-D-remainder}
\end{align}
The kernels obey the corresponding exponential tree bounds with all
background link insertions included before their sums are estimated.
\end{theorem}

\begin{proof}
\Cref{lem:newS7V15-sizes} gives
\(\|\Delta\mathsf A\|\le Ca\|A\|_\infty\), so
\eqref{eq:newS7V15-eta} holds for small \(a\).  Multiply the \(n\)-th word
of \eqref{eq:newS7V15-resolvent-series} by the overlap factor \(a^{-1}\);
this proves \eqref{eq:newS7V15-Dn}.  The geometric tail gives
\eqref{eq:newS7V15-D-remainder}.  For kernels, use the weighted resolvent
bound of V12 between successive finite-range insertions.  Summing the
complete word produces the shortest-tree exponential exactly as in
\Cref{thm:newS7V12-sign-tree}; the Taylor coefficient is normalized before
that sum.
\end{proof}

\begin{assessment}[What the series does and does not control]
\label{ass:newS7V15-series-scope}
The expansion \eqref{eq:newS7V15-D-expansion} is uniform for the first-order
Dirac operator because its small parameter is \(a\|A\|_\infty\).  Squaring
it creates the \(O(a^{-1})\) cross term in
\eqref{eq:newS7V15-L-size}; inserting those square terms independently
into a heat Duhamel series would lose covariance again.  The series must be
used inside a local covariant parametrix or inside the already assembled
curvature response.
\end{assessment}

\subsection{Radial gauge on a heat-scale patch}
\label{sec:newS7V15-radial}

Let \(A\) be a \(C^1\) connection on a star-shaped Euclidean coordinate
ball \(B_R(x_0)\), and let \(U_{xy}\) be exact edge parallel transport.
For every lattice vertex \(x\) in the ball, let \(\gamma_x\) be the
straight continuum radial segment from \(x_0\) to \(x\), and gauge transform
by its exact parallel transporter.  Denote the transformed links by
\(U_e^{\rm rad}\).

\begin{lemma}[Curvature control of radial links]
\label{lem:newS7V15-radial-links}
There is a geometric constant \(C\) such that every edge
\(e=(x,x+a\hat\mu)\) whose radial comparison loop lies in \(B_R(x_0)\)
obeys
\begin{equation}
 \|I-U_e^{\rm rad}\|
 \le Ca\{|x-x_0|+a\}
 \sup_{B_R(x_0)}\|F_A\|
 \exp\{CR\|A\|_\infty\}.                                 \label{eq:newS7V15-radial-links}
\end{equation}
\end{lemma}

\begin{proof}
After the gauge transformation, \(U_e^{\rm rad}\) is the holonomy around
the loop formed by \(\gamma_x\), the edge \(e\), and the reverse of
\(\gamma_{x+a\hat\mu}\).  Fill this loop by the radial strip between the two
radial paths.  Its area is at most
\(Ca(|x-x_0|+a)\).

Differentiate the parallel transport as the strip is swept.  The
non-Abelian variation formula inserts the curvature once, conjugated by
partial path transports.  Unitarity gives norm one for the conjugations in
a unitary gauge; in an arbitrary fixed matrix norm Gronwall contributes at
most the displayed exponential.  Integrating the curvature norm over the
strip proves \eqref{eq:newS7V15-radial-links}.
\end{proof}

\begin{corollary}[Smallness on the heat scale]
\label{cor:newS7V15-heat-scale}
For edges in \(B_{\Lambda\sqrt t}(x_0)\) and \(a^2\le t\),
\begin{equation}
 \|I-U_e^{\rm rad}\|
 \le C_{A,\Lambda}a\sqrt t\,\|F_A\|_{L^\infty}
 +C_{A,\Lambda}a^2\|F_A\|_{L^\infty}.                    \label{eq:newS7V15-heat-scale}
\end{equation}
Thus the gauge potential seen after radial trivialization is of order
\(\sqrt t\,F_A\), rather than the gauge-dependent global size of \(A\).
\end{corollary}

\begin{proof}
Insert \(|x-x_0|\le\Lambda\sqrt t\) in
\eqref{eq:newS7V15-radial-links}.
\end{proof}

\begin{lemma}[First curvature jet in radial gauge]
\label{lem:newS7V15-radial-jet}
If \(A\in C^2\), then in the continuum radial gauge centered at \(x_0\),
\begin{equation}
 A_\mu^{\rm rad}(x_0+y)
 =\int_0^1s\,y^\nu F_{\nu\mu}(x_0+sy)\,ds
 ={1\over2}y^\nu F_{\nu\mu}(x_0)
 +O(|y|^2\|\nabla F\|_\infty).                           \label{eq:newS7V15-radial-jet}
\end{equation}
\end{lemma}

\begin{proof}
Radial gauge is \(y^\mu A_\mu^{\rm rad}(x_0+y)=0\).
Differentiate \(sA_\mu^{\rm rad}(x_0+sy)\) in \(s\) and use the radial
condition to replace the antisymmetric derivative by
\(s y^\nu F_{\nu\mu}(x_0+sy)\).  Integrate from zero to one.  Taylor's
theorem for \(F\) gives the second equality.
\end{proof}

\subsection{Smooth links versus rough admissibility}
\label{sec:newS7V15-smooth-rough}

\begin{proposition}[Plaquette admissibility alone does not give the radial
heat-scale estimate]
\label{prop:newS7V15-rough-AOP}
The AOP condition
\(\|I-U_p\|\le\delta\) with fixed \(\delta>0\) does not imply
\eqref{eq:newS7V15-heat-scale} uniformly as \(a\downarrow0\).  A radial
strip of physical area \(O(a\sqrt t)\) may contain
\(O(\sqrt t/a)\) plaquettes, so a product estimate based only on
\(\delta\) permits an error of order
\begin{equation}
 {\sqrt t\over a}\,\delta,                                \label{eq:newS7V15-rough-loss}
\end{equation}
which is not small in the mesoscopic regime.
\end{proposition}

\begin{proof}
Tile the radial comparison strip by elementary plaquettes.  The only
general product inequality available from fixed admissibility is
\[
 \|I-U_{p_1}\cdots U_{p_m}\|
 \le\sum_{j=1}^m\|I-U_{p_j}\|
 \le m\delta.
\]
Here \(m=O(\sqrt t/a)\).  This proves
\eqref{eq:newS7V15-rough-loss} and shows that a continuum curvature bound,
whose plaquette defect is \(O(a^2)\), is additional information.
\end{proof}

\begin{assessment}[No scale-alignment shortcut]
\label{ass:newS7V15-no-alignment}
A field may satisfy the uniform AOP gap while failing every fixed smooth
curvature bound as the cutoff changes.  V14 therefore extends first to
bounded smooth-link families.  A later probabilistic or action-coercive
argument must show that the rough AOP remainder is negligible.  The NS V67
audit warns against assuming that control on one heat-scale patch aligns
with all smaller scales.
\end{assessment}

\subsection{The covariant smooth-background gate}
\label{sec:newS7V15-gate}

\begin{definition}[Covariant smooth-background card, CSBC]
\label{def:newS7V15-CSBC}
For every bounded \(C^2\) smooth-link family, CSBC consists of:
\begin{enumerate}[label=\textup{(\roman*)},leftmargin=3.5em]
\item radial link control \eqref{eq:newS7V15-heat-scale};
\item the curvature jet \eqref{eq:newS7V15-radial-jet};
\item the factorial-free overlap expansion
      \eqref{eq:newS7V15-D-expansion};
\item a patchwise heat parametrix on radius \(\Lambda\sqrt t\), with the
      outside error controlled by V11;
\item identification of its two-curvature term with the V14 coefficient;
\item a uniform error
      \[
       C_A\{a/\sqrt t+\sqrt t\,\|\nabla F\|_\infty\}
      \]
      after the local lower-order owners are subtracted; and
\item a separate estimate for the rough AOP complement.
\end{enumerate}
\end{definition}

\begin{theorem}[Closed and open CSBC rows]
\label{thm:newS7V15-CSBC}
Rows \textup{(i)}--\textup{(iii)} of CSBC hold.  V11 supplies the
off-patch decay required in row \textup{(iv)} for the free heat factors.
Rows \textup{(iv)}--\textup{(vi)} still require a variable-coefficient
covariant lattice parametrix.  Row \textup{(vii)} is a distinct
concentration or universality problem and does not follow from AOP.
\end{theorem}

\begin{proof}
Use \Cref{cor:newS7V15-heat-scale,
lem:newS7V15-radial-jet,
thm:newS7V15-D-expansion,
thm:newS7V11-kernel-bound,
prop:newS7V15-rough-AOP}.
\end{proof}

\begin{programme}[V16 acquisition]
\label{prog:newS7V15-V16}
V16 will construct the first patchwise parametrix row of CSBC.  Freeze the
radial-gauge curvature jet at \(x_0\), use the V14 free two-point kernel as
the model, and write one Duhamel remainder containing
\(F(x)-F(x_0)\).  The V11 weighted heat moments should give the
\(O(\sqrt t\,\|\nabla F\|_\infty)\) error, while V12 gives
\(O(a/\sqrt t)\) for the lattice vertex mismatch.  All repeated frozen
insertions will remain in the normalized resolvent series of V15.  No
statement about the rough AOP complement will be made without a separate
measure estimate.
\end{programme}

\begin{firewall}[Claim boundary after seventh-series V15]
\label{fire:newS7V15-boundary}
V15 performs the compulsory companion audit, proves that direct
operator-norm Duhamel expansion of the overlap square is nonuniform,
constructs a factorial-free holomorphic background expansion of the
first-order overlap operator with a geometric remainder, proves curvature
control and the first curvature jet in radial gauge, and proves that fixed
plaquette admissibility alone cannot supply the smooth heat-scale bound.

V15 does not construct the variable-coefficient lattice parametrix, prove
the full smooth-background coefficient transfer, control the rough AOP
complement, establish a chiral phase, remove the interacting cutoff, prove
either companion programme, or construct the Standard Model--Einstein
continuum.
\end{firewall}

\begin{theorem}[Seventh-series V15 result]
\label{thm:newS7V15-result}
Bounded smooth backgrounds admit a factorial-free overlap expansion with
\(n\)-th term \(O(a^{n-1}\|A\|^n)\).  On a heat-scale ball, radial gauge
replaces the gauge-dependent potential by the curvature jet
\((1/2)y^\nu F_{\nu\mu}(x_0)\) plus
\(O(t\|\nabla F\|)\).  A direct expansion of the overlap square remains
nonuniform, and fixed AOP admissibility does not imply this smooth radial
bound.  The next rigorous gate is therefore a patchwise covariant
parametrix for the smooth-link family, with rough admissible fields kept as
a separate no-escape problem.
\end{theorem}

\begin{proof}
Use \Cref{audit:newS7V15-transfer,
prop:newS7V15-naive-Duhamel,
thm:newS7V15-resolvent-series,
cor:newS7V15-factorial,
thm:newS7V15-D-expansion,
lem:newS7V15-radial-links,
lem:newS7V15-radial-jet,
prop:newS7V15-rough-AOP,
thm:newS7V15-CSBC}.
\end{proof}

\paragraph{Parent-version record.}
V15 was formed from
\texttt{Renewal\_SM\_Einstein\_Continuum\_Research\_Notes\_V14.tex}.
The V14 parent had \(6\,123\) logical source lines, \(239\,952\) bytes, and
SHA--256
\[
 \texttt{291948294930B1EFB7BD6A154B7037A05145ABDA0E4920A858C4BF80B912F12D}.
\]
The next compulsory companion audit is V20.

% =============================================================================
% END APPENDED SEVENTH-SERIES V15 MATERIAL
% =============================================================================

% =============================================================================
% BEGIN APPENDED SEVENTH-SERIES V16 MATERIAL
% =============================================================================

\clearpage
\section{Seventh-series V16: a radial-gauge lattice parametrix for the
smooth-background \(F^2\) coefficient}
\label{sec:v7v16}

V15 supplies the factorial-free overlap expansion and the radial curvature
jet but leaves their heat-scale assembly open.  This version performs that
assembly for a bounded smooth gauge background on a flat chart.  The
result is local and zero-winding; rough AOP fields, Higgs terms, and metric
variation remain outside its scope.

\subsection{All polynomial moments of the free heat kernel}
\label{sec:newS7V16-moments}

\begin{lemma}[All free heat moments]
\label{lem:newS7V16-all-moments}
For every integer \(j\ge0\), the free overlap heat kernel of V11 satisfies
\begin{equation}
 a^4\sum_{x\in a\mathbb Z^4}
 |x|^j|\mathsf K_a(t;x)|
 \le C_jt^{j/2},
 \qquad t\ge a^2.                                        \label{eq:newS7V16-all-moments}
\end{equation}
\end{lemma}

\begin{proof}
In \eqref{eq:newS7V11-kernel-bound} choose the rapid-decay exponent
\(M>4+j\) and apply the radial summation lemma of V11.  The heat-scale term
has moment \(C_jt^{j/2}\).  The lattice tail has moment
\(C_ja^je^{-ct/a^2}\le C_jt^{j/2}\).
\end{proof}

\begin{corollary}[Heat-scale Taylor remainder]
\label{cor:newS7V16-Taylor-moment}
If \(G\in C^{m+1}\) near \(x_0\), then convolution against any finite
product of free heat kernels with total time \(t\) permits Taylor expansion
of \(G(x_0+y)\) through order \(m\), with remainder bounded by
\begin{equation}
 C_mt^{(m+1)/2}
 \sup_{|y|\le C\sqrt t}|\nabla^{m+1}G(x_0+y)|,             \label{eq:newS7V16-Taylor-moment}
\end{equation}
plus a term smaller than any prescribed power from the complement of the
heat-scale ball.
\end{corollary}

\begin{proof}
Taylor's theorem bounds the local remainder by
\(C|y|^{m+1}\sup|\nabla^{m+1}G|\); use
\eqref{eq:newS7V16-all-moments} and the weighted convolution algebra of
V11.  Outside a ball of radius \(\Lambda\sqrt t\), choose the exponent
\(M\) in V11 larger than the requested power and then let \(\Lambda\)
increase.
\end{proof}

\subsection{The lattice radial curvature jet}
\label{sec:newS7V16-link-jet}

Let \(A\in C^2\) be in continuum radial gauge about \(x_0\), let
\(y=x-x_0\), and let \(U_{x,\mu}\) be exact parallel transport along the
edge from \(x\) to \(x+a\hat\mu\).

\begin{lemma}[Frozen-curvature link expansion]
\label{lem:newS7V16-link-jet}
For \(|y|\le\Lambda\sqrt t\),
\begin{align}
 U_{x,\mu}
 &=I+{a\over2}y^\nu F_{\nu\mu}(x_0)
   +\mathcal E_{x,\mu},                                  \label{eq:newS7V16-link-jet}\\
 \|\mathcal E_{x,\mu}\|
 &\le C_{A,\Lambda}\{
 a|y|^2\|\nabla F\|_\infty
 +a^2\|F\|_\infty
 +a^2|y|^2\|F\|_\infty^2\}.                              \label{eq:newS7V16-link-error}
\end{align}
Consequently,
\begin{equation}
 \|\mathcal E_{x,\mu}\|
 \le C_{A,\Lambda}\{
 at\|\nabla F\|_\infty
 +a^2\|F\|_\infty
 +a^2t\|F\|_\infty^2\}.                                  \label{eq:newS7V16-link-error-heat}
\end{equation}
\end{lemma}

\begin{proof}
V15 gives
\[
 A_\mu(x_0+y)
 ={1\over2}y^\nu F_{\nu\mu}(x_0)
 +O(|y|^2\|\nabla F\|_\infty).
\]
The edge transport ODE gives
\[
 U_{x,\mu}
 =I+\int_0^a A_\mu(x+s\hat\mu)\,ds
 +O\left\{
 \left(\int_0^a\|A_\mu(x+s\hat\mu)\|\,ds\right)^2
 \right\}.
\]
Changing \(A_\mu(x+s\hat\mu)\) to \(A_\mu(x)\) costs
\(Ca^2(\|F\|_\infty+|y|\|\nabla F\|_\infty)\).  Insert the radial jet and
absorb the mixed derivative term into the first two displayed errors.
The quadratic ODE remainder is
\(Ca^2|y|^2\|F\|_\infty^2\).  This proves
\eqref{eq:newS7V16-link-error}; use
\(|y|\le\Lambda\sqrt t\) for
\eqref{eq:newS7V16-link-error-heat}.
\end{proof}

Define the frozen radial potential
\begin{equation}
 A_\mu^{F_0}(y):={1\over2}y^\nu F_{\nu\mu}(x_0),
 \qquad F_0=F(x_0).                                      \label{eq:newS7V16-frozen-potential}
\end{equation}

\subsection{The rescaled insertion budget}
\label{sec:newS7V16-budget}

On the heat scale it is natural to use
\[
 \mathscr D_a(A):=\sqrt t\,D_{\rm ov}\{U(A)\}.
\]

\begin{lemma}[Frozen and error insertion sizes]
\label{lem:newS7V16-insertion-budget}
Inside the heat-scale patch, the factorial-free V15 series can be organized
as
\begin{equation}
 \mathscr D_a(A)
 =\mathscr D_a(0)+\mathscr V_a(F_0)
 +\mathscr E_a(A;x_0),                                   \label{eq:newS7V16-rescaled-D}
\end{equation}
where, after convolution with the heat kernels and vertex trees relevant
to a diagonal heat coefficient,
\begin{align}
 \|\mathscr V_a(F_0)\|_{\rm loc}
 &\le C_A t\|F_0\|,                                      \label{eq:newS7V16-frozen-size}\\
 \|\mathscr E_a(A;x_0)\|_{\rm loc}
 &\le C_A\{
 t^{3/2}\|\nabla F\|_\infty
 +a\sqrt t\,\|F\|_\infty
 +at^{3/2}\|F\|_\infty^2\}.                              \label{eq:newS7V16-error-size}
\end{align}
The \(n\)-fold repeated frozen insertion is bounded by
\begin{equation}
 C_A^n\{t\|F_0\|\}^n                                     \label{eq:newS7V16-n-frozen}
\end{equation}
with no factorial.
\end{lemma}

\begin{proof}
The leading link perturbation in
\eqref{eq:newS7V16-link-jet} is \(aA^{F_0}(y)\).  V12 shows that division by
the overlap normalization \(a\) turns a physical gauge-potential
insertion into an \(O(A^{F_0})\) Dirac vertex.  On the heat scale,
\(|A^{F_0}(y)|\le C\sqrt t\,|F_0|\); multiplication by the outside
\(\sqrt t\) gives \eqref{eq:newS7V16-frozen-size}.

Divide \eqref{eq:newS7V16-link-error-heat} by \(a\), multiply by
\(\sqrt t\), and use the V12 weighted vertex trees.  This gives
\eqref{eq:newS7V16-error-size}.  Finally, insert the frozen link
perturbation into the normalized resolvent-word series of V15.  Each
additional word supplies another heat-scale factor \(tF_0\); the Taylor
factorial has already been cancelled by
\eqref{eq:newS7V15-factorial}.  This proves
\eqref{eq:newS7V16-n-frozen}.
\end{proof}

\begin{corollary}[Heat-order separation]
\label{cor:newS7V16-order-separation}
In the local diagonal heat density:
\begin{enumerate}[label=\textup{(\roman*)},leftmargin=3.5em]
\item the term linear in \(\mathscr V_a(F_0)\) vanishes by the V9 tadpole
      and spin/reflection symmetry;
\item two frozen insertions have order \(t^2F_0^2\) inside the
      \((4\pi t)^{-2}\) heat prefactor and therefore supply the
      dimension-four \(F^2\) coefficient; and
\item three or more frozen insertions are bounded after the heat prefactor
      by \(C_At\|F_0\|^3\) for \(t\|F_0\|\le1\).
\end{enumerate}
\end{corollary}

\begin{proof}
The first item is \Cref{thm:newS7V9-tadpole-zero}, applied to the frozen
linear gauge perturbation.  The other two follow from
\eqref{eq:newS7V16-n-frozen} and the four-dimensional prefactor \(t^{-2}\).
The geometric tail beginning at \(n=3\) is bounded by its first term when
\(C_At\|F_0\|<1/2\).
\end{proof}

\subsection{Local smooth-background coefficient transfer}
\label{sec:newS7V16-transfer}

Let
\[
 h_a(t;x_0,A)
 :=\tr\{e^{-tD_{\rm ov}(A)^*D_{\rm ov}(A)}\}(x_0,x_0)
\]
denote the physical zero-winding diagonal heat density, including the
lattice normalization that makes its leading term of order \(t^{-2}\).

\begin{theorem}[Smooth-background local \(F^2\) parametrix]
\label{thm:newS7V16-local-parametrix}
Let \(A\) range over a bounded \(C^2\) family on a flat coordinate chart,
with \(t\|F\|_\infty\le c_0\) and \(a^2/t\le c_0\).  After subtracting the
free rank density, the parity-even gauge part satisfies
\begin{align}
 &\left|
 h_a(t;x_0,A)-h_a(t;x_0,0)
 -{2I_2(R)\over3(4\pi)^2}
  \langle F_{\mu\nu}(x_0),F_{\mu\nu}(x_0)\rangle_0
 \right|                                                  \notag\\
 &\qquad\le C_A\left[
 {a\over\sqrt t}\|F\|_\infty^2
 +\sqrt t\,\|F\|_\infty\|\nabla F\|_\infty
 +t\{\|F\|_\infty^3+\|\nabla F\|_\infty^2\}
 \right]
 +\mathcal T_{a,t},                                      \label{eq:newS7V16-local-parametrix}
\end{align}
where for every \(N\)
\begin{equation}
 |\mathcal T_{a,t}|
 \le C_{A,N}\{(t/R^2)^N+e^{-ct/a^2}\}                    \label{eq:newS7V16-parametrix-tail}
\end{equation}
when the chart contains a ball of fixed radius \(R\) about \(x_0\).
\end{theorem}

\begin{proof}
Use the continuum radial gauge at \(x_0\), which does not change the
diagonal trace.  By V11 and
\Cref{lem:newS7V16-all-moments}, paths leaving the fixed chart contribute
\eqref{eq:newS7V16-parametrix-tail}.  Within the heat-scale ball insert
\eqref{eq:newS7V16-rescaled-D} into the normalized V15 resolvent words and
then into the complete heat divided difference, retaining the contact
owners as in V10.

The linear frozen term vanishes and the quadratic frozen term is exactly
the free-background two-point coefficient of V14 evaluated on the radial
jet.  Equations \eqref{eq:newS7V14-convergence} and
\eqref{eq:newS7V14-continuum-c}, together with the factor \(1/4\) in the
curvature quadratic form \eqref{eq:newS7V13-curvature-quadratic}, give
\[
 {2I_2(R)\over3(4\pi)^2}|F_0|^2
\]
and the first error on the right of
\eqref{eq:newS7V16-local-parametrix}.

A cross term between one frozen insertion of size \(tF_0\) and the first
jet error \(t^{3/2}\nabla F\), divided by the heat factor \(t^2\), gives
\(C\sqrt t\,|F_0||\nabla F|\).  Two jet errors and the higher Taylor
remainders are bounded by the last bracket using
\Cref{cor:newS7V16-Taylor-moment}.  Three or more frozen insertions give
\(Ct|F_0|^3\) by
\Cref{cor:newS7V16-order-separation}.  The lattice part of
\eqref{eq:newS7V16-error-size} combines with one frozen insertion and,
after division by \(t^2\), is bounded by
\(C(a/\sqrt t)\|F\|^2\) plus the displayed higher-order term.  The
factorial-free geometric bounds justify summing all omitted words.
\end{proof}

\begin{corollary}[Pointwise continuum gauge coefficient on smooth links]
\label{cor:newS7V16-pointwise}
For every fixed smooth background and every joint limit
\[
 a\downarrow0,\qquad t\downarrow0,\qquad a^2/t\to0,
\]
the zero-winding parity-even gauge part obeys
\begin{equation}
 h_a(t;x_0,A)-h_a(t;x_0,0)
 \longrightarrow
 {2I_2(R)\over3(4\pi)^2}
 \langle F_{\mu\nu}(x_0),F_{\mu\nu}(x_0)\rangle_0.       \label{eq:newS7V16-pointwise}
\end{equation}
\end{corollary}

\begin{proof}
Every term on the right of
\eqref{eq:newS7V16-local-parametrix} tends to zero.  For the chart tail,
first choose \(N\) and use \(t/R^2\to0\); the lattice tail vanishes because
\(t/a^2\to\infty\).
\end{proof}

\subsection{Integrated smooth-link error}
\label{sec:newS7V16-integrated}

\begin{theorem}[Mesoscopic integrability after local subtraction]
\label{thm:newS7V16-integrated}
Let \(0<q<2/3\) and
\(\varepsilon_q(a)=\ell^2(a/\ell)^q\).  On a bounded \(C^2\) smooth-link
family, subtract the local \(F^2\) owner in
\eqref{eq:newS7V16-local-parametrix}.  The terms proportional to
\(a/\sqrt t\) are logarithmically integrable with vanishing norm:
\begin{equation}
 \int_{\varepsilon_q(a)}^{t_0}
 {a\,dt\over t^{3/2}}
 \le {2a\over\sqrt{\varepsilon_q(a)}}\longrightarrow0.   \label{eq:newS7V16-integrated-lattice}
\end{equation}
The \(\sqrt t\) and \(t\) remainders are integrable at zero.  Hence the
smooth-link local gauge row of CSBC has a renormalized mesoscopic limit.
\end{theorem}

\begin{proof}
Equation \eqref{eq:newS7V16-integrated-lattice} is the same integral as in
V14 and vanishes because \(q<2\).  Dividing the other errors by the
proper-time measure \(t\) gives \(t^{-1/2}\) and a constant, both integrable
at zero.  The rapid chart and lattice tails are treated as in V11 and V14.
\end{proof}

\begin{assessment}[What smooth-background closure means]
\label{ass:newS7V16-scope}
The theorem is a local coefficient and remainder statement for a
deterministic bounded smooth family.  It does not show that the interacting
lattice measure concentrates on such a family.  Fixed AOP admissibility
allows plaquette defects of order one rather than \(a^2\), so the rough
complement remains a separate no-escape problem exactly as isolated in
V15.
\end{assessment}

\begin{definition}[Smooth-link gauge coefficient card, SLGCC]
\label{def:newS7V16-SLGCC}
SLGCC consists of:
\begin{enumerate}[label=\textup{(\roman*)},leftmargin=3.5em]
\item all free heat moments \eqref{eq:newS7V16-all-moments};
\item the frozen radial link jet
      \eqref{eq:newS7V16-link-jet}--\eqref{eq:newS7V16-link-error-heat};
\item the rescaled insertion budget
      \eqref{eq:newS7V16-frozen-size}--\eqref{eq:newS7V16-n-frozen};
\item the local parametrix
      \eqref{eq:newS7V16-local-parametrix};
\item pointwise coefficient convergence
      \eqref{eq:newS7V16-pointwise}; and
\item logarithmic integrability
      \eqref{eq:newS7V16-integrated-lattice}.
\end{enumerate}
\end{definition}

\begin{programme}[V17 acquisition]
\label{prog:newS7V16-V17}
V17 will add the zero-order Higgs--Yukawa field to the same radial
parametrix.  It will compute the \((\nabla\Phi)^2\), \(\Phi^4\), and
\(R\Phi^2\) owners in a flat first pass, keep gauge--Higgs cross terms
through the complete internal trace, and prove the analogues of
\eqref{eq:newS7V16-local-parametrix}.  The moving chiral projectors will
remain separate until the vectorlike square and its relative subtraction
are fully assembled.
\end{programme}

\begin{firewall}[Claim boundary after seventh-series V16]
\label{fire:newS7V16-boundary}
V16 proves all polynomial moments of the free overlap heat kernel, a
frozen-curvature expansion of radial-gauge links, a factorial-free
heat-scale insertion budget, a local smooth-background lattice parametrix,
pointwise convergence to the V7 vectorlike \(F^2\) coefficient, and
integrability of its mesoscopic lattice error.

V16 does not control the rough AOP complement, add Higgs--Yukawa or metric
coefficients, construct moving chiral projectors or the determinant phase,
prove measure concentration, remove the full interacting cutoff, or
construct the Standard Model--Einstein continuum.
\end{firewall}

\begin{theorem}[Seventh-series V16 result]
\label{thm:newS7V16-result}
For bounded \(C^2\) smooth gauge backgrounds, radial gauge and the
factorial-free overlap expansion yield the local zero-winding limit
\[
 h_a(t;x,A)-h_a(t;x,0)
 \longrightarrow
 {2I_2(R)\over3(4\pi)^2}|F_{\mu\nu}(x)|^2
\]
whenever \(a^2/t\to0\).  The lattice error is
\(O(a/\sqrt t)\), the curvature-variation error is
\(O(\sqrt t)\), and both are integrable after the local \(F^2\)
subtraction.  This closes the deterministic smooth-link gauge coefficient
row; the rough AOP no-escape problem and the remaining Standard Model and
Einstein owners are still open.
\end{theorem}

\begin{proof}
Use \Cref{lem:newS7V16-all-moments,
lem:newS7V16-link-jet,
lem:newS7V16-insertion-budget,
cor:newS7V16-order-separation,
thm:newS7V16-local-parametrix,
cor:newS7V16-pointwise,
thm:newS7V16-integrated}.
\end{proof}

\paragraph{Parent-version record.}
V16 was formed from
\texttt{Renewal\_SM\_Einstein\_Continuum\_Research\_Notes\_V15.tex}.
The V15 parent had \(6\,583\) logical source lines, \(258\,564\) bytes, and
SHA--256
\[
 \texttt{350CEA56A30A9B0E8705FCAFCA10E65C8F384034152DAAA2E72991F39DF7F3D3}.
\]
The next compulsory companion audit is V20.

% =============================================================================
% END APPENDED SEVENTH-SERIES V16 MATERIAL
% =============================================================================

% =============================================================================
% BEGIN APPENDED SEVENTH-SERIES V17 MATERIAL
% =============================================================================

\clearpage
\section{Seventh-series V17: the Dirac--Yukawa \(A_4\) owners and their
relative subtraction}
\label{sec:v7v17}

V16 closes the deterministic smooth-link transfer of the vectorlike gauge
coefficient.  The next local Standard Model tensors arise from the
zero-order Higgs--Yukawa field.  This version computes their continuum
coefficients in the V7 heat convention and separates them from the still
open lattice and chiral-projector transfer.

\subsection{The Dirac--Yukawa square}
\label{sec:newS7V17-square}

Let \(S\) be the rank-four Euclidean spin bundle, \(E_R\) a unitary
internal gauge bundle, and \(\Phi=\Phi^*\) a smooth gauge-covariant internal
endomorphism commuting with Clifford multiplication.  Use the self-adjoint
Dirac--Yukawa operator
\begin{equation}
 \mathcal D_\Phi=i\gamma^\mu\nabla_\mu+\gamma_5\Phi.       \label{eq:newS7V17-DPhi}
\end{equation}

\begin{lemma}[Lichnerowicz endomorphism with Yukawa field]
\label{lem:newS7V17-Lichnerowicz}
The square is of Laplace type,
\begin{equation}
 \mathcal D_\Phi^2=\nabla^*\nabla+E_\Phi,                 \label{eq:newS7V17-square}
\end{equation}
with
\begin{equation}
 E_\Phi={R\over4}I
 -{1\over2}\gamma^{\mu\nu}\rho_*(F_{\mu\nu})
 +\Phi^2+i\gamma^\mu\gamma_5D_\mu\Phi.                   \label{eq:newS7V17-E}
\end{equation}
Changing the sign convention for the gauge term changes neither its
parity-even square nor the Higgs coefficients below.
\end{lemma}

\begin{proof}
The ordinary twisted Lichnerowicz formula gives the first two terms.
Because \(\gamma_5\gamma^\mu=-\gamma^\mu\gamma_5\) and \(\Phi\) commutes
with the gamma matrices, the first-order cross terms cancel:
\[
 i\gamma^\mu\nabla_\mu(\gamma_5\Phi)
 +\gamma_5\Phi\,i\gamma^\mu\nabla_\mu
 =i\gamma^\mu\gamma_5D_\mu\Phi.
\]
The square of the zero-order term is \(\Phi^2\).
\end{proof}

\subsection{The integrated \(A_4\) formula in the fixed convention}
\label{sec:newS7V17-A4}

\begin{lemma}[Relevant integrated Laplace coefficient]
\label{lem:newS7V17-general-A4}
For \(L=\nabla^*\nabla+E\) on a closed four-manifold, the terms in the
integrated \(A_4\) coefficient that can contain gauge or Yukawa fields are
\begin{equation}
 A_4(L)\big|_{\rm matter}
 =\int_M\tr\left\{
 {1\over2}E^2-{1\over6}RE
 +{1\over12}\Omega_{\mu\nu}\Omega_{\mu\nu}
 \right\}\,d\mu_g,                                       \label{eq:newS7V17-general-A4}
\end{equation}
modulo integrated total derivatives and terms containing only Riemann
curvature and the identity endomorphism.
\end{lemma}

\begin{proof}
Write the standard Laplace operator as
\(-(\nabla^2+\mathcal E)\).  In the present convention
\(\mathcal E=-E\).  The weight-four recursion gives the matter-dependent
terms
\[
 {1\over360}\tr\{
 60R\mathcal E+180\mathcal E^2
 +30\Omega_{\mu\nu}\Omega_{\mu\nu}
 +60\nabla^2\mathcal E\}.
\]
Substitute \(\mathcal E=-E\).  The last term integrates to zero on a closed
manifold, yielding \eqref{eq:newS7V17-general-A4}.
\end{proof}

\begin{lemma}[Spin traces of the Yukawa pieces]
\label{lem:newS7V17-spin-traces}
Put
\[
 E_F=-{1\over2}\gamma^{\mu\nu}\rho_*(F_{\mu\nu}),\qquad
 E_Y=\Phi^2+i\gamma^\mu\gamma_5D_\mu\Phi.
\]
Then
\begin{align}
 \tr_S(E_Y^2)
 &=4\Phi^4+4D_\mu\Phi\,D_\mu\Phi,                         \label{eq:newS7V17-EY-square}\\
 \tr_S(E_FE_Y)&=0,                                        \label{eq:newS7V17-cross-zero}\\
 \tr_S(E_Y)&=4\Phi^2.                                     \label{eq:newS7V17-EY-trace}
\end{align}
\end{lemma}

\begin{proof}
The cross term between \(\Phi^2\) and
\(i\gamma^\mu\gamma_5D_\mu\Phi\) has one gamma matrix and zero spin trace.
Furthermore
\[
 (i\gamma^\mu\gamma_5)(i\gamma^\nu\gamma_5)
 =\gamma^\mu\gamma^\nu,
\]
so its spin trace is \(4g^{\mu\nu}\), proving
\eqref{eq:newS7V17-EY-square}.  The \(E_F\Phi^2\) trace contains
\(\tr_S\gamma^{\mu\nu}=0\); the other cross term contains three gamma
matrices with \(\gamma_5\) and also has zero trace.  This proves
\eqref{eq:newS7V17-cross-zero}.  The derivative part of \(E_Y\) has zero
trace, giving \eqref{eq:newS7V17-EY-trace}.
\end{proof}

\begin{theorem}[Dirac--Yukawa matter owners]
\label{thm:newS7V17-owners}
For one complex four-component Dirac--Yukawa square, the parity-even gauge
and Higgs part of \(A_4\), modulo total derivatives and pure-gravity terms,
is
\begin{align}
 A_4(\mathcal D_\Phi^2)\big|_{\rm gauge+Higgs}
 =\int_M\Big[
 &-{2\over3}\tr_R\{\rho_*(F_{\mu\nu})\rho_*(F_{\mu\nu})\}
 +2\tr_R(D_\mu\Phi\,D_\mu\Phi)                            \notag\\
 &+2\tr_R(\Phi^4)
 +{1\over3}R\,\tr_R(\Phi^2)
 \Big]\,d\mu_g.                                           \label{eq:newS7V17-owners}
\end{align}
There is no parity-even monomial \(F_{\mu\nu}\Phi^2\) and no
\(F_{\mu\nu}D_\rho\Phi\) cross owner.
\end{theorem}

\begin{proof}
The gauge term is V7.  For the Yukawa terms, insert
\eqref{eq:newS7V17-EY-square} into \(E^2/2\), producing
\(2\tr(D\Phi)^2+2\tr\Phi^4\).

The curvature--Higgs cross term in \(E^2/2\) comes from
\((R/4)I\) and \(\Phi^2\).  Its spin trace is
\(R\tr_R\Phi^2\).  The term \(-RE/6\) contributes
\(-(2/3)R\tr_R\Phi^2\), leaving
\((1/3)R\tr_R\Phi^2\).  All gauge--Higgs cross terms vanish by
\eqref{eq:newS7V17-cross-zero}; the bundle-curvature term contains no
\(\Phi\).
\end{proof}

\begin{corollary}[Constant-mass check]
\label{cor:newS7V17-mass-check}
If \(\Phi=mI_R\) is constant and commutes with the connection, then the
\(\Phi\)-dependent terms in \eqref{eq:newS7V17-owners} are
\begin{equation}
 2r_Rm^4+{r_R\over3}Rm^2.                                \label{eq:newS7V17-mass-check}
\end{equation}
These agree with multiplying the massless heat series by \(e^{-tm^2}\).
\end{corollary}

\begin{proof}
The rank-four massless coefficients are
\(A_0=4r_R\) and
\(A_2=-r_RR/3\).  In
\[
 e^{-tm^2}\{A_0+tA_2+t^2A_4+\cdots\},
\]
the order-\(t^2\) mass terms are
\((m^4/2)A_0-m^2A_2\), which give
\eqref{eq:newS7V17-mass-check}.
\end{proof}

\subsection{Equivariant Higgs--Yukawa invariants}
\label{sec:newS7V17-invariants}

Let \(\mathcal H\) be the real Higgs representation and suppose
\(H\mapsto\Phi(H)\) is real-linear and gauge equivariant, including the
conjugate-doublet entries used for up-type Yukawa couplings.  Use the
invariant real norm \(|H|\).

\begin{lemma}[Quadratic and quartic Yukawa traces]
\label{lem:newS7V17-Yukawa-invariants}
On each irreducible Higgs orbit sector there are nonnegative constants
\(a_Y,b_Y\) such that
\begin{align}
 \tr_R\{\Phi(H)^2\}&=a_Y|H|^2,                            \label{eq:newS7V17-aY}\\
 \tr_R\{\Phi(H)^4\}&=b_Y|H|^4,                            \label{eq:newS7V17-bY}\\
 \tr_R\{D_\mu\Phi(H)D_\mu\Phi(H)\}
 &=a_Y|D_\mu H|^2.                                       \label{eq:newS7V17-kineticY}
\end{align}
The constants are sums of squared and fourth-power Yukawa singular values,
with colour and generation multiplicities determined by the internal
representation.
\end{lemma}

\begin{proof}
The first trace is a gauge-invariant positive quadratic form on the real
Higgs representation.  Irreducibility and invariance make it a scalar
multiple \(a_Y\) of the chosen norm.  Polarization gives the same scalar
for the derivative quadratic form, proving
\eqref{eq:newS7V17-kineticY}.  The fourth trace is a gauge-invariant
homogeneous quartic.  For one Standard Model Higgs doublet, the invariant
polynomial algebra at degree four is generated by \(|H|^4\), giving
\eqref{eq:newS7V17-bY}.  Diagonalizing the generation-space Yukawa matrices
identifies the constants with the stated singular-value sums.
\end{proof}

\begin{corollary}[Higgs owner form]
\label{cor:newS7V17-Higgs-owner}
For one Higgs doublet, the Yukawa part of
\eqref{eq:newS7V17-owners} is
\begin{equation}
 2a_Y|D_\mu H|^2
 +2b_Y|H|^4
 +{a_Y\over3}R|H|^2.                                    \label{eq:newS7V17-Higgs-owner}
\end{equation}
\end{corollary}

\begin{proof}
Use \Cref{lem:newS7V17-Yukawa-invariants}.
\end{proof}

\begin{assessment}[Normalization scope]
\label{ass:newS7V17-normalization}
The constants \(a_Y,b_Y\) include exactly the internal states present in
the chosen Dirac test block.  Particle--antiparticle doubling, a missing
right-handed neutrino, generation number, and the passage from a
four-component Dirac square to a chiral determinant change these traces.
No numerical Standard Model value is asserted until that finite internal
block is fixed.
\end{assessment}

\subsection{Relative Yukawa subtraction}
\label{sec:newS7V17-relative}

Let \(\Phi,\Phi_0\) be two covariant Yukawa endomorphisms on the same gauge
and spin bundle.

\begin{theorem}[Exact relative Higgs owners]
\label{thm:newS7V17-relative}
When the full and base Dirac squares use the same gauge connection, their
relative parity-even \(A_4\) coefficient has no pure \(F^2\) term and its
Higgs part is
\begin{align}
 \Delta A_4^{Y}
 =\int_M\Big[
 &2\tr_R\{(D\Phi)^2-(D\Phi_0)^2\}
 +2\tr_R(\Phi^4-\Phi_0^4)                                \notag\\
 &+{R\over3}\tr_R(\Phi^2-\Phi_0^2)
 \Big]\,d\mu_g.                                           \label{eq:newS7V17-relative}
\end{align}
\end{theorem}

\begin{proof}
Subtract \eqref{eq:newS7V17-owners} for the two blocks.  The pure gauge
terms are identical and cancel, as already proved in V7.
\end{proof}

\begin{corollary}[Owned Higgs normalizations remain necessary]
\label{cor:newS7V17-owned}
A same-\((A,\Phi)\) covariant base cancels all universal matter terms in
\eqref{eq:newS7V17-owners}.  Therefore the absolute Higgs kinetic,
quartic, and nonminimal-curvature coefficients must be retained as owned
bare or counterterm normalizations and fixed by physical conditions.
\end{corollary}

\begin{proof}
Set \(\Phi_0=\Phi\) in
\eqref{eq:newS7V17-relative}.  Gauge-invariant local additions preserve the
relative Ward identity but change the absolute normalizations.
\end{proof}

\subsection{The overlap--Yukawa transfer gate}
\label{sec:newS7V17-transfer}

In the rescaled V16 heat operator, a smooth zero-order Yukawa insertion has
size \(\sqrt t\,\Phi\), while its square enters with size \(t\Phi^2\).

\begin{definition}[Smooth overlap--Yukawa coefficient card, SOYCC]
\label{def:newS7V17-SOYCC}
SOYCC consists of:
\begin{enumerate}[label=\textup{(\roman*)},leftmargin=3.5em]
\item a lattice Yukawa block converging to \(\gamma_5\Phi(H)\) on every
      smooth-link Higgs family;
\item exact gauge covariance and the Ginsparg--Wilson chiral
      intertwiners;
\item a factorial-free joint gauge--Yukawa expansion at the first-order
      Dirac level;
\item heat-scale moments for \(H\), \(D H\), and their radial Taylor
      remainders;
\item convergence to the three owners in
      \eqref{eq:newS7V17-Higgs-owner}; and
\item a logarithmically integrable post-subtraction remainder, followed
      separately by the moving-projector and determinant-line phase.
\end{enumerate}
\end{definition}

\begin{proposition}[Continuum target fixed, lattice transfer open]
\label{prop:newS7V17-SOYCC}
Rows \textup{(i)}--\textup{(iv)} of SOYCC are sufficient to run the V16
radial parametrix with zero-order insertions.  If its joint remainder is
uniform, row \textup{(v)} must equal
\eqref{eq:newS7V17-Higgs-owner}; row \textup{(vi)} is still needed for the
renormalized chiral determinant.
\end{proposition}

\begin{proof}
The heat-scale sizes \(\sqrt t\,\Phi\) and \(t\Phi^2\), together with V16
moments, make the quadratic derivative and quartic potential terms occur at
weight four.  Gauge covariance combines ordinary derivatives into
\(D_\mu H\).  The continuum coefficient is unique and was computed in
\Cref{thm:newS7V17-owners,cor:newS7V17-Higgs-owner}.  None of these local
steps supplies the moving-projector phase.
\end{proof}

\begin{programme}[V18 acquisition]
\label{prog:newS7V17-V18}
V18 will instantiate SOYCC on the archived overlap--Yukawa block.  It will
write the exact finite internal matrix \(\Phi(H)\), compute \(a_Y,b_Y\)
with the existing generation and colour multiplicities, and verify the
Ginsparg--Wilson chiral intertwining before any heat estimate.  If the
archived block omits a right-handed neutrino or uses a reduced Higgs chart,
both branches will be recorded separately rather than completed by an
unstated field.
\end{programme}

\begin{firewall}[Claim boundary after seventh-series V17]
\label{fire:newS7V17-boundary}
V17 derives the full parity-even gauge and Higgs part of the continuum
Dirac--Yukawa \(A_4\) coefficient in the established convention, proves
all gauge--Higgs cross traces vanish, identifies the Higgs kinetic,
quartic, and \(R|H|^2\) owners, and computes their exact relative
subtraction.

V17 does not fix the finite Standard Model Yukawa matrix or its numerical
traces, transfer these coefficients through the lattice overlap--Yukawa
block, construct moving chiral projectors or the determinant phase,
control rough AOP fields, remove the interacting cutoff, or construct the
Standard Model--Einstein continuum.
\end{firewall}

\begin{theorem}[Seventh-series V17 result]
\label{thm:newS7V17-result}
For a four-component Dirac--Yukawa square, the local matter coefficient is
\[
 -{2\over3}\tr_R(F_{\mu\nu}F_{\mu\nu})
 +2\tr_R(D_\mu\Phi D_\mu\Phi)
 +2\tr_R(\Phi^4)
 +{R\over3}\tr_R(\Phi^2).
\]
For one equivariant Higgs doublet the last three owners are
\(2a_Y|DH|^2+2b_Y|H|^4+(a_Y/3)R|H|^2\).
A same-gauge relative base cancels the pure \(F^2\) term, and a
same-\((A,\Phi)\) base cancels all of these universal local owners.  The
next task is to instantiate and transfer the actual overlap--Yukawa
internal block without assuming its chiral phase.
\end{theorem}

\begin{proof}
Use \Cref{lem:newS7V17-Lichnerowicz,
lem:newS7V17-general-A4,
lem:newS7V17-spin-traces,
thm:newS7V17-owners,
lem:newS7V17-Yukawa-invariants,
thm:newS7V17-relative,
prop:newS7V17-SOYCC}.
\end{proof}

\paragraph{Parent-version record.}
V17 was formed from
\texttt{Renewal\_SM\_Einstein\_Continuum\_Research\_Notes\_V16.tex}.
The V16 parent had \(6\,989\) logical source lines, \(274\,001\) bytes, and
SHA--256
\[
 \texttt{20A43AB5FEECB2DF6A27575FC95B7B6AB0C6BC765AD6FF3C77AC2BA8B3BBB4B7}.
\]
The next compulsory companion audit is V20.

% =============================================================================
% END APPENDED SEVENTH-SERIES V17 MATERIAL
% =============================================================================

% =============================================================================
% BEGIN APPENDED SEVENTH-SERIES V18 MATERIAL
% =============================================================================

\clearpage
\section{Seventh-series V18: the finite Standard Model Yukawa block,
exact trace constants, and its improved-overlap realization}
\label{sec:v7v18}

V17 fixed the continuum Dirac--Yukawa owners but deliberately left the
finite internal trace unspecified.  The frozen second-series archive fixes
the minimal one-generation representation and the improved-field overlap
ordering, but does not write the Higgs-dependent finite matrix.  This
version supplies that missing matrix for an arbitrary number of generations,
proves its gauge equivariance, computes its quadratic and quartic traces,
and distinguishes the doubled heat trace from the physical chiral weight.

\subsection{The physical left and right representation spaces}
\label{sec:newS7V18-spaces}

Let \(F=\mathbb C^{N_g}\) be generation space and
\(C=\mathbb C^3\) colour space.  In the physical right-handed convention,
put
\begin{align}
 K_L&=(\mathbb C^2\otimes C\otimes F)_q
       \oplus(\mathbb C^2\otimes F)_\ell,                 \label{eq:newS7V18-KL}\\
 K_R^{\min}&=(C\otimes F)_u\oplus(C\otimes F)_d
       \oplus F_e.                                        \label{eq:newS7V18-KRmin}
\end{align}
The corresponding hypercharges are
\[
 Y(q_L)={1\over6},\quad Y(\ell_L)=-{1\over2},\quad
 Y(u_R)={2\over3},\quad Y(d_R)=-{1\over3},\quad Y(e_R)=-1.
\]
The optional Dirac-neutrino extension is
\begin{equation}
 K_R^\nu=K_R^{\min}\oplus F_\nu,
 \qquad Y(\nu_R)=0.                                       \label{eq:newS7V18-KRnu}
\end{equation}
No Majorana mass is included in this version.

For \(H=(h_1,h_2)^T\in\mathbb C^2\), define
\begin{equation}
 \widetilde H:=i\sigma_2\overline H
   =\binom{\overline h_2}{-\overline h_1}.                 \label{eq:newS7V18-Htilde}
\end{equation}
Then
\begin{equation}
 |\widetilde H|=|H|,
 \qquad \widetilde H^*H=0.                                \label{eq:newS7V18-Horthogonal}
\end{equation}

Let \(Y_u,Y_d,Y_e\in\End(F)\), and in the optional branch let
\(Y_\nu\in\End(F)\).  These matrices include the flavour mixing; no
diagonal flavour basis is chosen.

\begin{definition}[Finite Standard Model Yukawa map]
\label{def:newS7V18-MH}
The minimal map \(M_H^{\min}:K_R^{\min}\to K_L\) is
\begin{equation}
 M_H^{\min}(u,d,e)
 =\left(
   \widetilde H\otimes Y_u u+H\otimes Y_d d,
   H\otimes Y_e e
  \right),                                                \label{eq:newS7V18-MHmin}
\end{equation}
where the identity acts on colour in both quark terms.  In the
Dirac-neutrino branch,
\begin{equation}
 M_H^\nu(u,d,e,n)
 =\left(
   \widetilde H\otimes Y_u u+H\otimes Y_d d,
   H\otimes Y_e e+\widetilde H\otimes Y_\nu n
  \right).                                                \label{eq:newS7V18-MHnu}
\end{equation}
On \(K_L\oplus K_R^\bullet\), with
\(\bullet\in\{\min,\nu\}\), define the self-adjoint completion
\begin{equation}
 \Phi_\bullet(H):=
 \begin{pmatrix}0&M_H^\bullet\\ (M_H^\bullet)^*&0\end{pmatrix}.
                                                               \label{eq:newS7V18-Phi}
\end{equation}
\end{definition}

\begin{proposition}[Exact Standard Model gauge equivariance]
\label{prop:newS7V18-equivariance}
For \(g\in SU(3)_c\times SU(2)_L\times U(1)_Y\),
\begin{equation}
 M_{gH}^\bullet
 =\rho_L(g)M_H^\bullet\rho_R(g)^{-1}.                    \label{eq:newS7V18-equivariance}
\end{equation}
Consequently \(\Phi_\bullet(H)\) is a gauge-covariant endomorphism of
the doubled internal bundle.
\end{proposition}

\begin{proof}
Colour covariance is immediate because the quark maps contain the identity
on \(C\), and the lepton maps are colour singlets.  Under \(SU(2)_L\),
both \(H\) and \(\widetilde H\) transform as doublets; the identity
\(i\sigma_2\overline U=U i\sigma_2\) holds for \(U\in SU(2)\).

It remains to check hypercharge.  An entry from a right representation of
charge \(y_R\) to a left representation of charge \(y_L\) must carry
charge \(y_L-y_R\).  The four differences are
\begin{align*}
 {1\over6}-{2\over3}&=-{1\over2} &&\text{for }u,\\
 {1\over6}-\left(-{1\over3}\right)&={1\over2} &&\text{for }d,\\
 -{1\over2}-(-1)&={1\over2} &&\text{for }e,\\
 -{1\over2}-0&=-{1\over2} &&\text{for }\nu.
\end{align*}
These are precisely the hypercharges of
\(\widetilde H,H,H,\widetilde H\), respectively.  This proves
\eqref{eq:newS7V18-equivariance}; taking the adjoint proves covariance of
both off-diagonal corners in \eqref{eq:newS7V18-Phi}.
\end{proof}

\subsection{Exact quadratic and quartic trace constants}
\label{sec:newS7V18-traces}

Define the two flavour invariants
\begin{align}
 \mathfrak a_{\min}
 &:=\Tr_F\!\left[
  3Y_u^*Y_u+3Y_d^*Y_d+Y_e^*Y_e\right],                   \label{eq:newS7V18-amin}\\
 \mathfrak b_{\min}
 &:=\Tr_F\!\left[
  3(Y_u^*Y_u)^2+3(Y_d^*Y_d)^2+(Y_e^*Y_e)^2\right].        \label{eq:newS7V18-bmin}
\end{align}
For the Dirac-neutrino branch put
\begin{align}
 \mathfrak a_\nu&:=\mathfrak a_{\min}
   +\Tr_F(Y_\nu^*Y_\nu),                                 \label{eq:newS7V18-anu}\\
 \mathfrak b_\nu&:=\mathfrak b_{\min}
   +\Tr_F\{(Y_\nu^*Y_\nu)^2\}.                          \label{eq:newS7V18-bnu}
\end{align}

\begin{lemma}[Orthogonal-column identity]
\label{lem:newS7V18-column}
For either branch,
\begin{align}
 \Tr_{K_R^\bullet}\{(M_H^\bullet)^*M_H^\bullet\}
 &=\mathfrak a_\bullet |H|^2,                             \label{eq:newS7V18-M2}\\
 \Tr_{K_R^\bullet}\{((M_H^\bullet)^*M_H^\bullet)^2\}
 &=\mathfrak b_\bullet |H|^4.                             \label{eq:newS7V18-M4}
\end{align}
In particular, no mixed \(Y_uY_d\) or \(Y_eY_\nu\) quartic occurs.
\end{lemma}

\begin{proof}
By \eqref{eq:newS7V18-Horthogonal}, the two columns entering the same weak
doublet are orthogonal.  Thus the quark part of
\((M_H^\bullet)^*M_H^\bullet\) is block diagonal with diagonal blocks
\[
 |H|^2(I_C\otimes Y_u^*Y_u),
 \qquad |H|^2(I_C\otimes Y_d^*Y_d),
\]
and the lepton part has diagonal blocks
\(|H|^2Y_e^*Y_e\) and, when present,
\(|H|^2Y_\nu^*Y_\nu\).  Taking the trace gives
\eqref{eq:newS7V18-M2}; colour supplies the factor three.  Squaring this
block-diagonal matrix before taking the trace gives
\eqref{eq:newS7V18-M4} and proves the absence of mixed quartics.
\end{proof}

\begin{theorem}[The constants \(a_Y,b_Y\) fixed]
\label{thm:newS7V18-ab}
For the self-adjoint completion \eqref{eq:newS7V18-Phi},
\begin{align}
 \Tr_{K_L\oplus K_R^\bullet}\Phi_\bullet(H)^2
 &=2\mathfrak a_\bullet |H|^2,                            \label{eq:newS7V18-Phi2}\\
 \Tr_{K_L\oplus K_R^\bullet}\Phi_\bullet(H)^4
 &=2\mathfrak b_\bullet |H|^4.                            \label{eq:newS7V18-Phi4}
\end{align}
For the induced covariant derivatives,
\begin{equation}
 \Tr_{K_L\oplus K_R^\bullet}
   \{D_\mu\Phi_\bullet(H)D_\mu\Phi_\bullet(H)\}
 =2\mathfrak a_\bullet |D_\mu H|^2.                      \label{eq:newS7V18-DPhi2}
\end{equation}
Hence the V17 constants are exactly
\begin{equation}
 a_Y=2\mathfrak a_\bullet,
 \qquad b_Y=2\mathfrak b_\bullet.                         \label{eq:newS7V18-ab-values}
\end{equation}
\end{theorem}

\begin{proof}
The even powers of \(\Phi_\bullet\) are block diagonal:
\[
 \Phi_\bullet^2=
 \begin{pmatrix}M_H^\bullet(M_H^\bullet)^*&0\\
 0&(M_H^\bullet)^*M_H^\bullet\end{pmatrix}.
\]
The nonzero eigenvalues of \(MM^*\) and \(M^*M\), including
multiplicity, agree.  Therefore each trace is twice the corresponding
right-space trace in
\Cref{lem:newS7V18-column}, proving
\eqref{eq:newS7V18-Phi2,eq:newS7V18-Phi4}.

Gauge equivariance gives
\(D_\mu M_H=M_{D_\mu H}\).  Moreover
\(D_\mu\widetilde H=\widetilde{D_\mu H}\), because the conjugate
doublet uses the conjugate hypercharge connection.  Repeating the quadratic
calculation with \(H\) replaced by \(D_\mu H\), and again counting the two
self-adjoint corners, proves \eqref{eq:newS7V18-DPhi2}.  Comparison with
the definitions in V17 gives \eqref{eq:newS7V18-ab-values}.
\end{proof}

\begin{corollary}[Explicit doubled Dirac--Yukawa owners]
\label{cor:newS7V18-owners}
For either branch, the Higgs part of the V17 integrated \(A_4\) coefficient
is
\begin{equation}
 4\mathfrak a_\bullet |D_\mu H|^2
 +4\mathfrak b_\bullet |H|^4
 +{2\mathfrak a_\bullet\over3}R|H|^2.                    \label{eq:newS7V18-explicit-owners}
\end{equation}
The minimal branch is obtained from
\eqref{eq:newS7V18-amin,eq:newS7V18-bmin}; the optional Dirac-neutrino
branch adds precisely the terms in
\eqref{eq:newS7V18-anu,eq:newS7V18-bnu}.
\end{corollary}

\begin{proof}
Substitute \eqref{eq:newS7V18-ab-values} into the V17 owner formula
\(2a_Y|DH|^2+2b_Y|H|^4+(a_Y/3)R|H|^2\).
\end{proof}

\begin{assessment}[Doubled trace versus chiral determinant]
\label{ass:newS7V18-doubling}
The factor two in \eqref{eq:newS7V18-ab-values} is an exact algebraic
consequence of the self-adjoint \(K_L\oplus K_R\) completion.  It is not by
itself the multiplicity of the physical chiral functional integral.  If a
specific determinant modulus is represented as
\(\frac12\log\det(\mathcal B^*\mathcal B)\), its heat coefficient carries
the corresponding factor \(1/2\); any further charge-conjugate or mirror
doubling must be counted from the declared finite Hilbert space.  The
determinant-line phase is not recovered from this positive square.
\end{assessment}

\subsection{Exact improved-overlap chiral block}
\label{sec:newS7V18-overlap}

Let \(D_L(U)\) and \(D_R(U)\) be the vectorlike overlap extensions in the
left and right Standard Model representations.  For either one set
\begin{equation}
 \widehat\gamma_{5,r}=\gamma_5(I-aD_r),\quad
 \widehat P_{\pm,r}={I\pm\widehat\gamma_{5,r}\over2},\quad
 A_r=I-{a\over2}D_r,
 \qquad r\in\{L,R\}.                                      \label{eq:newS7V18-GW-data}
\end{equation}
The Ginsparg--Wilson identity gives
\begin{equation}
 D_r\widehat P_{\pm,r}=P_\mp D_r,
 \qquad A_r\widehat P_{\pm,r}=P_\pm A_r.                 \label{eq:newS7V18-GW-intertwine}
\end{equation}

\begin{definition}[Gauge-covariant improved Yukawa block]
\label{def:newS7V18-block}
With source columns \((\psi_L,\psi_R)\) and row components
\((\overline\psi_L,\overline\psi_R)\), define
\begin{equation}
 \mathcal B_a(U,H):=
 \begin{pmatrix}
 P_+D_L\widehat P_{-,L}&
 P_+M_H^\bullet A_R\widehat P_{+,R}\\
 P_-(M_H^\bullet)^* A_L\widehat P_{-,L}&
 P_-D_R\widehat P_{+,R}
 \end{pmatrix}.                                           \label{eq:newS7V18-block}
\end{equation}
\end{definition}

\begin{proposition}[Exact covariance and chiral typing]
\label{prop:newS7V18-block-covariance}
Every entry in \eqref{eq:newS7V18-block} has the displayed source and row
chirality, and the full block is gauge covariant.  In particular, the
upper-right Yukawa operator obeys
\begin{equation}
 M_{gH}^\bullet A_R(gU)
 =\rho_L(g)M_H^\bullet A_R(U)\rho_R(g)^{-1},               \label{eq:newS7V18-improved-covariance}
\end{equation}
with the analogous adjoint identity in the other corner.
\end{proposition}

\begin{proof}
The chirality statements follow from
\eqref{eq:newS7V18-GW-intertwine}; \(M_H^\bullet\) acts only on internal
indices and therefore commutes with the fixed spin projectors.  Overlap
covariance implies
\(A_R(gU)=\rho_R(g)A_R(U)\rho_R(g)^{-1}\).  Multiply this identity by
\eqref{eq:newS7V18-equivariance} to obtain
\eqref{eq:newS7V18-improved-covariance}.  The kinetic and adjoint corners
are covariant by the same calculation.
\end{proof}

\begin{theorem}[Constant broken-phase reduction to massive-overlap blocks]
\label{thm:newS7V18-massive-blocks}
Let \(H=v n_0\) with \(|n_0|=1\), and restrict the gauge field to the
unbroken \(SU(3)_c\times U(1)_{\rm em}\) subgroup.  After the exact
left--right representation intertwiners and singular-value decompositions
of \(Y_u,Y_d,Y_e\) are applied, the charged part of
\eqref{eq:newS7V18-block} is an orthogonal sum of standard massive-overlap
operators
\begin{equation}
 D_m=\left(1-{am\over2}\right)D_{\rm ov}+mI,
 \qquad m=v\sigma,                                        \label{eq:newS7V18-Dm}
\end{equation}
one for each Yukawa singular value \(\sigma\).  The minimal neutrino source
has no right partner and remains in the unmatched massless chiral block.
In the optional branch, its nonzero singular values give the same
construction.
\end{theorem}

\begin{proof}
At \(H=vn_0\), Higgs reduction identifies the left and right colour and
electromagnetic representations for each charged species.  The overlap
kinetic operator is then the same vectorlike operator \(D\) in the two
identified spaces.  Constant flavour unitaries commute with \(D\), with
\(A_D=I-aD/2\), and with all spin projectors, so singular-value
decomposition reduces each Yukawa matrix to scalar entries
\(m=v\sigma\).

For one such entry, summing the two kinetic corners gives \(D\) by the
first identity in \eqref{eq:newS7V18-GW-intertwine}, while summing the two
Yukawa corners gives \(mA_D\) by the second.  Hence the full-space operator
is
\[
 D+mA_D=\left(1-{am\over2}\right)D+mI,
\]
which is \eqref{eq:newS7V18-Dm}.  The representation ledger contains no
\(\nu_R\) in the minimal branch, so no unitary left--right identification
or mass corner exists there.  Adding \(F_\nu\) and applying the same
singular-value argument proves the optional statement.
\end{proof}

\begin{assessment}[Scope before symmetry reduction]
\label{ass:newS7V18-prereduction}
Before Higgs reduction, \(D_L\) and \(D_R\) act in inequivalent
electroweak representations.  Gauge covariance of
\eqref{eq:newS7V18-block} is exact, but it does not identify that block with
a single vectorlike \(D+mA_D\).  The latter identity is asserted only after
the unbroken-representation intertwiners in
\Cref{thm:newS7V18-massive-blocks} have been supplied.
\end{assessment}

\subsection{Closed SOYCC rows and next estimate}
\label{sec:newS7V18-SOYCC}

\begin{proposition}[Algebraic SOYCC rows]
\label{prop:newS7V18-SOYCC}
For the discretization \eqref{eq:newS7V18-block}:
\begin{enumerate}[label=\textup{(\roman*)},leftmargin=3.5em]
\item the finite Higgs map is exactly
      \eqref{eq:newS7V18-MHmin} or \eqref{eq:newS7V18-MHnu};
\item its Standard Model gauge covariance is
      \eqref{eq:newS7V18-equivariance};
\item its Ginsparg--Wilson source/row typing is exact by
      \eqref{eq:newS7V18-GW-intertwine};
\item its doubled continuum trace constants are
      \eqref{eq:newS7V18-ab-values}; and
\item its constant broken-phase normalization is the standard massive
      overlap operator by \Cref{thm:newS7V18-massive-blocks}.
\end{enumerate}
Thus the algebraic parts of SOYCC are closed.  The joint smooth-background
heat remainder and the physical chiral determinant weight remain analytic
tasks.
\end{proposition}

\begin{proof}
The five statements are respectively
\Cref{def:newS7V18-MH,
prop:newS7V18-equivariance,
prop:newS7V18-block-covariance,
thm:newS7V18-ab,
thm:newS7V18-massive-blocks}.
\end{proof}

\begin{programme}[V19 acquisition]
\label{prog:newS7V18-V19}
V19 will derive the factorial-free joint gauge--Yukawa word expansion for
\eqref{eq:newS7V18-block} on a smooth Higgs family.  It will keep the
zero-order heat-scale budget \(\sqrt t\,M_H\) separate from the first-order
gauge budget, prove which word lengths can contribute to weight four, and
seek an integrable post-subtraction remainder.  If the chiral square creates
an ordering commutator \([A_D,M_H]\), that commutator will be estimated
before cyclicity is used.
\end{programme}

\begin{firewall}[Claim boundary after seventh-series V18]
\label{fire:newS7V18-boundary}
V18 writes the exact minimal and optional-Dirac-neutrino Standard Model
Yukawa maps, proves their gauge equivariance, computes the exact quadratic
and quartic trace constants with colour and generation multiplicities,
constructs the improved Ginsparg--Wilson chiral block, and proves its
constant broken-phase reduction to massive-overlap summands.

V18 does not add a Majorana sector, identify the doubled heat trace with an
undeclared physical determinant multiplicity, prove the joint
gauge--Yukawa heat remainder, construct the determinant-line phase, control
rough AOP fields, prove measure concentration, remove the interacting
cutoff, or construct the Standard Model--Einstein continuum.
\end{firewall}

\begin{theorem}[Seventh-series V18 result]
\label{thm:newS7V18-result}
Let
\[
 \mathfrak a_{\min}=\Tr_F(3Y_u^*Y_u+3Y_d^*Y_d+Y_e^*Y_e),
\]
and define \(\mathfrak b_{\min}\) by replacing each quadratic matrix with
its square.  For the self-adjoint finite Standard Model Yukawa completion,
\[
 a_Y=2\mathfrak a_{\min},\qquad
 b_Y=2\mathfrak b_{\min},
\]
so its doubled Higgs owners are
\[
 4\mathfrak a_{\min}|DH|^2
 +4\mathfrak b_{\min}|H|^4
 +{2\mathfrak a_{\min}\over3}R|H|^2.
\]
The optional Dirac-neutrino branch adds
\(\Tr(Y_\nu^*Y_\nu)\) and
\(\Tr\{(Y_\nu^*Y_\nu)^2\}\) to the two flavour invariants.  The improved
Yukawa ordering is exactly gauge covariant and becomes the standard
massive-overlap operator after constant Higgs reduction and flavour SVD.
\end{theorem}

\begin{proof}
Use \Cref{prop:newS7V18-equivariance,
lem:newS7V18-column,
thm:newS7V18-ab,
cor:newS7V18-owners,
prop:newS7V18-block-covariance,
thm:newS7V18-massive-blocks}.
\end{proof}

\paragraph{Parent-version record.}
V18 was formed from
\texttt{Renewal\_SM\_Einstein\_Continuum\_Research\_Notes\_V17.tex}.
The V17 parent had \(7\,372\) logical source lines, \(288\,490\) bytes, and
SHA--256
\[
 \texttt{7775CE7AA7FE9273BFAC0762EB2C4A467C396AF635BCCCCC430B8C256FF02C51}.
\]
The next compulsory companion audit is V20.

% =============================================================================
% END APPENDED SEVENTH-SERIES V18 MATERIAL
% =============================================================================


% =============================================================================
% BEGIN APPENDED SEVENTH-SERIES V19 MATERIAL
% =============================================================================

\clearpage
\section{Seventh-series V19: factorial-free joint gauge--Yukawa words and
the finite weight-four gate}
\label{sec:v7v19}

V18 closes the finite Standard Model Yukawa algebra.  Directly expanding
the square of the chiral block would recreate the V15 loss: large
first-order terms would be separated before Ward, Clifford, and
off-diagonal cancellations act.  This version instead expands the assembled
improved-Yukawa operator, proves an exact commutator estimate, and identifies
the finite set of heat words which can contribute to the flat-space
weight-four coefficient.

\subsection{An exact joint first-order expansion}
\label{sec:newS7V19-joint}

Work first in one Higgs-reduced matched representation.  Write
\begin{equation}
 B(A,H):=D(A)+M(H)A_D(A),
 \qquad A_D(A):=I-{a\over2}D(A).                           \label{eq:newS7V19-B}
\end{equation}
Let \(D_0=D(0)\), \(A_0=I-aD_0/2\), fix a constant Higgs value
\(H_0\), and put
\[
 M_0=M(H_0),\qquad K=M(H)-M_0.
\]
Use the V15 normalized expansion
\begin{equation}
 D(A)=D_0+\sum_{n\ge1}\mathcal D_{a,n}(A),                \label{eq:newS7V19-Dseries}
\end{equation}
where
\(\|\mathcal D_{a,n}(A)\|
 \le C_\delta^n a^{n-1}\|A\|_\infty^n\).

\begin{theorem}[Exact factorial-free joint word expansion]
\label{thm:newS7V19-joint-series}
Whenever the V15 resolvent smallness condition holds,
\begin{align}
 B(A,H)-B(0,H_0)
 &=KA_0+\sum_{n\ge1}\mathcal W_{a,n}(A,K),                \label{eq:newS7V19-joint-series}\\
 \mathcal W_{a,n}(A,K)
 &:=(I-aM_0/2)\mathcal D_{a,n}(A)
   -{a\over2}K\mathcal D_{a,n}(A).                        \label{eq:newS7V19-Wn}
\end{align}
There is one pure-gauge and one mixed Higgs--gauge word for each normalized
gauge word.  In particular, no factorial and no independent permutation
sum occurs.  The bounds are
\begin{align}
 \|KA_0\|&\le\|K\|,                                      \label{eq:newS7V19-KA0}\\
 \|\mathcal W_{a,n}(A,K)\|
 &\le C_\delta^n a^{n-1}\|A\|_\infty^n
 \left(1+{a\over2}\|M_0\|+{a\over2}\|K\|\right).       \label{eq:newS7V19-Wn-bound}
\end{align}
If the series is truncated after \(N\), the remainder obeys
\begin{equation}
 \|\mathcal R^{B}_{a,N+1}\|
 \le C_{\delta,A}a^N\|A\|_\infty^{N+1}
 \left(1+{a\over2}\|M_0\|+{a\over2}\|K\|\right).       \label{eq:newS7V19-B-remainder}
\end{equation}
\end{theorem}

\begin{proof}
The identity
\[
 B(A,H)=(I-aM(H)/2)D(A)+M(H)
\]
is exact and preserves the declared ordering.  Substitute
\(M(H)=M_0+K\) and \eqref{eq:newS7V19-Dseries}, then subtract
\(B(0,H_0)=D_0+M_0A_0\).  The term independent of the gauge perturbation is
\(KA_0\); at gauge order \(n\) the two remaining terms are exactly
\eqref{eq:newS7V19-Wn}.  The contraction
\(\|A_0\|\le1\) gives \eqref{eq:newS7V19-KA0}; the V15 estimate gives
\eqref{eq:newS7V19-Wn-bound}.  Multiplying its geometric remainder by the
same bounded prefactor proves \eqref{eq:newS7V19-B-remainder}.
\end{proof}

\begin{corollary}[Only one Higgs insertion occurs in a first-order word]
\label{cor:newS7V19-one-Higgs}
At the assembled first-order level, every term in
\eqref{eq:newS7V19-joint-series} contains zero or one factor of the Higgs
increment \(K\).  Multiple Higgs insertions arise only when complete heat
divided differences of \(B^*B\), or equivalently complete products of the
first-order words, are assembled.
\end{corollary}

\begin{proof}
The map \(M\mapsto D+M(I-aD/2)\) is affine.  Formula
\eqref{eq:newS7V19-Wn} displays its only mixed term.
\end{proof}

\subsection{The improved-field ordering defect}
\label{sec:newS7V19-commutator}

Use the physical lattice distance \(d(x,y)\).  The overlap locality already
proved in the active series has the Schur-summable form
\begin{equation}
 \|D(A;x,y)\|\le {C_\delta\over a}
 e^{-c_\delta d(x,y)/a},                                  \label{eq:newS7V19-locality}
\end{equation}
with the harmless lattice-measure normalization absorbed in the kernel.

\begin{lemma}[First moment of an overlap commutator]
\label{lem:newS7V19-commutator}
Let \(M\) be a multiplication operator satisfying
\begin{equation}
 \|M(x)-M(y)\|\le L_M d(x,y).                             \label{eq:newS7V19-M-Lipschitz}
\end{equation}
Then, uniformly in the lattice volume,
\begin{align}
 \|[D(A),M]\|&\le C_\delta L_M,                          \label{eq:newS7V19-DM}\\
 \|[A_D(A),M]\|&\le C_\delta aL_M.                      \label{eq:newS7V19-ADM}
\end{align}
For the V18 Standard Model map one may take
\begin{equation}
 L_M\le C y_*\|DH\|_\infty,
 \qquad
 y_*:=\max\{\|Y_u\|,\|Y_d\|,\|Y_e\|,\|Y_\nu\|\},     \label{eq:newS7V19-LM}
\end{equation}
with the absent neutrino term omitted.
\end{lemma}

\begin{proof}
The commutator kernel is
\[
 [D,M](x,y)=D(x,y)\{M(y)-M(x)\}.
\]
By \eqref{eq:newS7V19-locality,eq:newS7V19-M-Lipschitz}, its row Schur sum
is bounded by
\[
 {C_\delta L_M\over a}
 \sum_y e^{-c_\delta d(x,y)/a}d(x,y)\le C_\delta L_M.
\]
The same estimate holds for column sums.  The Schur test proves
\eqref{eq:newS7V19-DM}.  Since
\([A_D,M]=-(a/2)[D,M]\), \eqref{eq:newS7V19-ADM} follows.

For the Standard Model map, parallel-transport the source and target fibres
before comparing two sites.  Each entry is linear in \(H\) or
\(\widetilde H\); their covariant Lipschitz constants agree, and their
operator coefficients are bounded by \(y_*\).  The finite direct-sum
multiplicity is absorbed in \(C\), proving \eqref{eq:newS7V19-LM}.
\end{proof}

\begin{corollary}[Ordering may be changed only with an explicit lattice
remainder]
\label{cor:newS7V19-ordering}
On a bounded smooth Higgs family,
\begin{equation}
 \|M A_D-A_D M\|
 \le C_\delta a y_*\|DH\|_\infty.                         \label{eq:newS7V19-ordering}
\end{equation}
Thus the ordering defect tends to zero at fixed physical regularity, but it
cannot be discarded at a fixed cutoff.
\end{corollary}

\begin{proof}
Use \eqref{eq:newS7V19-ADM,eq:newS7V19-LM}.
\end{proof}

\subsection{Covariant Higgs jets on the heat patch}
\label{sec:newS7V19-Higgs-jet}

Trivialize the Higgs bundle by exact radial parallel transport to the base
point \(x_0\), using the same radial gauge as V15--V16.  Write
\(H_0=H(x_0)\).

\begin{lemma}[Radial covariant Higgs jet]
\label{lem:newS7V19-Higgs-jet}
For \(|y|\le\Lambda\sqrt t\),
\begin{equation}
 H(x_0+y)=H_0+y^\mu D_\mu H(x_0)+\mathcal E_H(y),          \label{eq:newS7V19-Higgs-jet}
\end{equation}
where
\begin{equation}
 \|\mathcal E_H(y)\|
 \le C|y|^2\|D^2H\|_{L^\infty(B_{\Lambda\sqrt t})}.      \label{eq:newS7V19-Higgs-error}
\end{equation}
Consequently, after multiplication by the outside heat scaling
\(\sqrt t\), the constant, linear, and remainder Yukawa insertions have
local sizes
\begin{align}
 \mathsf h_0&:=C y_*\sqrt t\,|H_0|,                       \label{eq:newS7V19-h0}\\
 \mathsf h_1&:=C y_*t\,|DH(x_0)|,                         \label{eq:newS7V19-h1}\\
 \mathsf h_2&:=C_{\Lambda}y_*t^{3/2}\|D^2H\|_\infty.    \label{eq:newS7V19-h2}
\end{align}
\end{lemma}

\begin{proof}
Along the radial geodesic, exact parallel transport turns the covariant
derivative into the ordinary derivative of the transported vector.  Taylor's
formula with integral remainder proves
\eqref{eq:newS7V19-Higgs-jet,eq:newS7V19-Higgs-error}.  The V18 map is
linear and has operator norm at most \(Cy_*|H|\).  On the heat patch,
\(|y|\le\Lambda\sqrt t\); multiply the three Taylor terms by
\(\sqrt t\) to obtain
\eqref{eq:newS7V19-h0,eq:newS7V19-h1,eq:newS7V19-h2}.
\end{proof}

Together with the V16 frozen gauge size, put
\begin{equation}
 \mathsf g:=C t|F(x_0)|.                                  \label{eq:newS7V19-g}
\end{equation}
The improved-ordering remainder has rescaled size
\begin{equation}
 \mathsf c_a:=C a\sqrt t\,y_*\|DH\|_\infty
 ={a\over\sqrt t}\,O(\mathsf h_1).                       \label{eq:newS7V19-ca}
\end{equation}

\subsection{The finite weight-four word list}
\label{sec:newS7V19-weight}

Assign heat weights
\begin{equation}
 w(\mathsf g)=1,qquad
 w(\mathsf h_0)={1\over2},qquad
 w(\mathsf h_1)=1,qquad
 w(\mathsf h_2)={3\over2}.                               \label{eq:newS7V19-weights}
\end{equation}
A product of total weight two sits at order \(t^2\) before the
four-dimensional diagonal heat prefactor and can contribute to \(A_4\).

\begin{lemma}[Complete flat-space weight-two enumeration]
\label{lem:newS7V19-enumeration}
After the vacuum and weight-one Higgs mass terms are separated, every
flat-space word of total heat weight two belongs to one of the following
classes:
\begin{equation}
 \mathsf g^2,qquad
 \mathsf g\mathsf h_0^2,qquad
 \mathsf g\mathsf h_1,qquad
 \mathsf h_0^4,qquad
 \mathsf h_0^2\mathsf h_1,qquad
 \mathsf h_1^2,qquad
 \mathsf h_0\mathsf h_2.                                 \label{eq:newS7V19-word-list}
\end{equation}
The classes \(\mathsf g\mathsf h_1\) and
\(\mathsf h_0^2\mathsf h_1\) have an odd number of Yukawa off-diagonal
insertions and vanish under the doubled internal trace.  The class
\(\mathsf g\mathsf h_0^2\) vanishes under the parity-even spin trace.
The surviving classes are therefore
\begin{equation}
 \mathsf g^2,qquad \mathsf h_0^4,qquad
 \mathsf h_1^2,qquad \mathsf h_0\mathsf h_2.             \label{eq:newS7V19-survivors}
\end{equation}
\end{lemma}

\begin{proof}
Multiplying the weights by two reduces the enumeration to the nonnegative
integer equation
\[
 2n_g+n_0+2n_1+3n_2=4.
\]
Listing its solutions gives \eqref{eq:newS7V19-word-list}.  Each Yukawa
insertion exchanges \(K_L\) and \(K_R\), so a product with an odd number of
such insertions has zero diagonal block and zero internal trace.  This
removes the two stated odd classes.  In the remaining mixed gauge class,
the continuum frozen-curvature insertion contains
\(\gamma^{\mu\nu}F_{\mu\nu}\), while the two constant Higgs insertions are
Clifford scalars.  Since \(\tr_S\gamma^{\mu\nu}=0\), its parity-even spin
trace vanishes.  This proves \eqref{eq:newS7V19-survivors}.
\end{proof}

\begin{proposition}[Identification of the four surviving owners]
\label{prop:newS7V19-owner-identification}
The four word classes in \eqref{eq:newS7V19-survivors} have the following
continuum owners:
\begin{enumerate}[label=\textup{(\roman*)},leftmargin=3.5em]
\item \(\mathsf g^2\) gives the gauge curvature square fixed in V14--V16;
\item \(\mathsf h_0^4\) gives
      \(4\mathfrak b_\bullet|H|^4\) in the doubled convention;
\item the assembled \(\mathsf h_1^2\) and
      \(\mathsf h_0\mathsf h_2\) words give
      \(4\mathfrak a_\bullet|DH|^2\), modulo the local total derivative
      generated by the heat recursion.
\end{enumerate}
On a curved background the connection coefficient additionally assembles
the already computed
\((2\mathfrak a_\bullet/3)R|H|^2\) owner.
\end{proposition}

\begin{proof}
The gauge statement is the V16 local parametrix.  The constant Yukawa word
is fixed by
\(\Tr\Phi^4=2\mathfrak b_\bullet|H|^4\) and the factor
\(E^2/2\) in the V17 heat coefficient, giving
\(4\mathfrak b_\bullet|H|^4\).  Taylor expansion of a multiplication
operator produces both two first jets and the constant--second-jet pair.
The local heat recursion combines them into the covariant kinetic invariant
and a Laplacian total derivative.  Its integrated coefficient is fixed by
the V17 spin trace and the V18 identity
\(\Tr(D\Phi)^2=2\mathfrak a_\bullet|DH|^2\), giving the stated factor four.
The curved owner is the V17 curvature cross term with
\(a_Y=2\mathfrak a_\bullet\).
\end{proof}

\subsection{A finite-symbol transfer gate}
\label{sec:newS7V19-transfer}

\begin{definition}[Finite gauge--Yukawa symbol card, FGYSC]
\label{def:newS7V19-FGYSC}
FGYSC consists of uniform local divided-difference bounds for complete
products of the first-order words in
\eqref{eq:newS7V19-joint-series}, together with explicit convergence of the
four classes \eqref{eq:newS7V19-survivors} to the coefficients in
\Cref{prop:newS7V19-owner-identification}.  The bounds must be applied to
each complete word before vertex positions or orderings are summed.
\end{definition}

\begin{theorem}[All-order transfer reduces to FGYSC]
\label{thm:newS7V19-finite-gate}
Assume FGYSC on a bounded \(C^2\) smooth gauge--Higgs family and the V16
heat-kernel moment bounds.  After subtracting the vacuum, the weight-one
Higgs mass term, and the weight-four owners in
\Cref{prop:newS7V19-owner-identification}, all continuum Taylor words have
total heat weight at least \(5/2\).  Their diagonal contribution after the
\(t^{-2}\) prefactor is bounded by
\begin{equation}
 C\sqrt t\,\mathcal P(F,H,DH,D^2H),                        \label{eq:newS7V19-continuum-tail}
\end{equation}
for a bounded polynomial \(\mathcal P\).  Lattice ordering replacements
carry the additional relative factor \(a/\sqrt t\) from
\eqref{eq:newS7V19-ca}.  On the mesoscopic window
\(t\ge\varepsilon_q(a)=\ell^2(a/\ell)^q\), \(0<q<2/3\), the continuum tail is integrable at zero and vanishes when its upper
endpoint tends to zero, while the lattice error vanishes as \(a\downarrow0\):
\begin{align}
 \int_0^{t_0}\sqrt t\,{dt\over t}&<\infty,                \label{eq:newS7V19-sqrt-integrable}\\
 \int_{\varepsilon_q(a)}^{t_0}{a\over\sqrt t}{dt\over t}
 &\le {2a\over\sqrt{\varepsilon_q(a)}}
 =O(a^{1-q/2}).                                            \label{eq:newS7V19-order-integrable}
\end{align}
\end{theorem}

\begin{proof}
The normalized first-order series is geometric by
\Cref{thm:newS7V19-joint-series}.  FGYSC permits complete heat words to be
bounded after their cancellations, so the sum of words above weight two is
dominated by a convergent multivariate geometric series in
\(\mathsf g,\mathsf h_0,\mathsf h_1,\mathsf h_2\).  The smallest retained
power after subtraction is \(t^{5/2}\); division by the diagonal heat
factor \(t^2\) proves \eqref{eq:newS7V19-continuum-tail}.  Equation
\eqref{eq:newS7V19-ca} gives the ordering factor.  Direct integration proves
\eqref{eq:newS7V19-sqrt-integrable,eq:newS7V19-order-integrable}; since
\(q<2\), the last power vanishes, and the established window
\(q<2/3\) is more than sufficient.
\end{proof}

\begin{assessment}[What is conditional in the finite gate]
\label{ass:newS7V19-conditional}
The word expansion, commutator estimate, heat weights, and finite
enumeration are proved in V19.  The hypothesis FGYSC is not proved here: it
contains the remaining lattice free-symbol limits for the Higgs-gradient
and quartic classes and their complete divided-difference domination.
Therefore \Cref{thm:newS7V19-finite-gate} is a reduction theorem, not a
claim that the smooth lattice Yukawa coefficient has already converged.
\end{assessment}

\begin{programme}[V20 audit and acquisition]
\label{prog:newS7V19-V20}
V20 is the compulsory companion audit.  It will inspect the newest settled
Yang--Mills and Navier--Stokes notes, test whether their normalization or
localization mechanisms sharpen FGYSC, and then compute the simplest
unclosed survivor: the constant-Higgs quartic free-symbol word.  The
Higgs-gradient pair will remain separate until its ordering and contact
terms have been assembled.
\end{programme}

\begin{firewall}[Claim boundary after seventh-series V19]
\label{fire:newS7V19-boundary}
V19 proves the exact factorial-free joint first-order gauge--Yukawa
expansion, the volume-uniform improved-ordering commutator bound, the
covariant Higgs heat-scale jet, and the complete finite enumeration of
flat-space weight-four words.  It reduces the all-order smooth transfer to
FGYSC and proves the resulting tail would be logarithmically integrable.

V19 does not prove FGYSC, transfer the Higgs quartic or gradient coefficients
through the lattice symbol, identify the doubled trace with the physical
chiral determinant multiplicity, construct the determinant-line phase,
control rough AOP fields, prove measure concentration, remove the
interacting cutoff, or construct the Standard Model--Einstein continuum.
\end{firewall}

\begin{theorem}[Seventh-series V19 result]
\label{thm:newS7V19-result}
The improved overlap--Yukawa operator has an exact joint background series
with one normalized gauge word and at most one Higgs increment per
first-order term, and its ordering defect obeys
\[
 \|[A_D,M_H]\|\le C_\delta a y_*\|DH\|_\infty.
\]
At heat weight two, all flat gauge--Higgs words either vanish by internal or
spin trace or reduce to
\[
 F^2,\qquad |H|^4,\qquad |DH|^2
\]
with the constant--second-jet term included in the kinetic assembly.  Once
the four surviving complete lattice words are verified, all higher words
are bounded by an integrable \(O(\sqrt t)+O(a/\sqrt t)\) remainder on the
established mesoscopic window.
\end{theorem}

\begin{proof}
Use \Cref{thm:newS7V19-joint-series,
lem:newS7V19-commutator,
lem:newS7V19-Higgs-jet,
lem:newS7V19-enumeration,
prop:newS7V19-owner-identification,
thm:newS7V19-finite-gate}.
\end{proof}

\paragraph{Parent-version record.}
V19 was formed from
\texttt{Renewal\_SM\_Einstein\_Continuum\_Research\_Notes\_V18.tex}.
The V18 parent had \(7\,824\) logical source lines, \(306\,311\) bytes, and
SHA--256
\[
 \texttt{7BFD396F174142755956647463FA89F246316C50A76B280CBE848D4429AC07C5}.
\]
The next compulsory companion audit is V20.

% =============================================================================
% END APPENDED SEVENTH-SERIES V19 MATERIAL
% =============================================================================


% =============================================================================
% BEGIN APPENDED SEVENTH-SERIES V20 MATERIAL
% =============================================================================

\clearpage
\section{Seventh-series V20: companion audit and exact constant-Higgs
quartic transfer}
\label{sec:v7v20}

V20 performs the compulsory companion audit and then closes the simplest
surviving FGYSC word.  The computation is made on the constant broken-phase
free symbol, where V18 gives an exact orthogonal sum of massive-overlap
operators.  The lower-order subtraction is kept at the lattice level until
the quartic coefficient is isolated.

\subsection{Companion audit}
\label{sec:newS7V20-audit}

The newest settled companion snapshots read for this audit were:
\begin{center}
\small
\begin{tabular}{p{0.26\textwidth}p{0.66\textwidth}}
Yang--Mills &
\texttt{Renewal\_Yang\_Mills\_Mass\_Gap\_Research\_Notes\_V93.tex},
\(1\,365\,160\) bytes, \(42\,180\) logical source lines, SHA--256
\texttt{ECB8A8145A4022DC25662956FAC775CCF320CB8EDD1C0A04DBCECECEC6CCA47F}.
\\ [0.4em]
Navier--Stokes &
\texttt{Renewal\_NS\_Stability\_Research\_Notes\_V83.tex},
\(1\,029\,002\) bytes, \(19\,657\) logical source lines, SHA--256
\texttt{7F60723920E0FC6B8D8658808347F1B48AEA5FDE64BFECFF15DD1436C481616B}.
\end{tabular}
\end{center}

\begin{audit}[Yang--Mills V93]
\label{audit:newS7V20-YM}
YM V93 proves the intermediate-wedge Duhamel contraction before insertion
norms are multiplied, retains the unrelaxed primitive-root reserve, and
closes every acyclic local--transversal incidence hypertree at arbitrary
total mark.  Its remaining sector has positive transversal incidence excess.
The warning is decisive: replacing the exact reserve by a global
\(B^{2g}\) estimate before quotient-cycle choices are summed loses the
needed exponential ledger.  No mass gap is claimed.
\end{audit}

\begin{audit}[Navier--Stokes V83]
\label{audit:newS7V20-NS}
NS V83 re-derives and validates the rate-free V26--V30 identities and the
Rayleigh V31--V35 identities, including the distinction between two
normalized flux conventions.  It confirms the exact assembled bridge and
Rayleigh formulas but makes no global-regularity claim.  Its conclusions
remain audits of exact identities; they do not turn invariant averages into
pointwise estimates.
\end{audit}

\begin{audit}[Transfer to the SM programme]
\label{audit:newS7V20-transfer}
The shared usable rule is to retain the exact normalizing factor through the
complete object whose cancellation it controls.  Here the analogue is the
factor \(1-a^2m^2/4\) in the massive-overlap square.  Expanding that factor
and subtracting only the continuum mass owner would leak a cutoff-dependent
quadratic term into the quartic scale.  V20 instead subtracts the exact
lattice derivative in the mass-squared parameter.  The NS normalization
audit also motivates an explicit three-level ledger: doubled internal heat
trace, one physical massive-block heat trace, and determinant modulus.  No
numerical inequality is imported from either companion programme.
\end{audit}

\subsection{The exact massive-overlap square}
\label{sec:newS7V20-massive-square}

Let \(D\) be the free overlap operator for one vectorlike internal state,
and let
\begin{equation}
 D_m=\left(1-{am\over2}\right)D+mI,
 \qquad 0\le am<2.                                        \label{eq:newS7V20-Dm}
\end{equation}
Set \(L_0=D^*D\).

\begin{lemma}[Exact scalar reduction of the massive square]
\label{lem:newS7V20-massive-square}
The free overlap operator is normal and satisfies
\begin{equation}
 D+D^*=aD^*D.                                              \label{eq:newS7V20-circle}
\end{equation}
Consequently
\begin{equation}
 D_m^*D_m=left(1-{a^2m^2\over4}\right)L_0+m^2I.          \label{eq:newS7V20-massive-square}
\end{equation}
\end{lemma}

\begin{proof}
Gamma-five Hermiticity and the Ginsparg--Wilson identity give
\eqref{eq:newS7V20-circle}; its adjoint shows that \(D^*D=DD^*\).
Put \(c_m=1-am/2\).  Then
\begin{align*}
 D_m^*D_m
 &=c_m^2D^*D+mc_m(D+D^*)+m^2I\\
 &=\{c_m^2+amc_m\}L_0+m^2I\\
 &=\left(1-{a^2m^2\over4}\right)L_0+m^2I.
\end{align*}
\end{proof}

Let
\begin{equation}
 k_a(t):=\tr_S(e^{-tL_0})(x,x)                            \label{eq:newS7V20-k0}
\end{equation}
be the translation-invariant free diagonal heat density for one internal
state, with spin rank four, and let \(k_{a,m}(t)\) be the corresponding
density for \(D_m^*D_m\).

\begin{corollary}[Exact massive heat identity]
\label{cor:newS7V20-massive-heat}
For \(0\le am<2\),
\begin{equation}
 k_{a,m}(t)=e^{-tm^2}
 k_a\!\left(t\left(1-{a^2m^2\over4}\right)\right).       \label{eq:newS7V20-massive-heat}
\end{equation}
\end{corollary}

\begin{proof}
The two terms on the right of
\eqref{eq:newS7V20-massive-square} commute.  Apply the functional calculus
and take the diagonal spin trace.
\end{proof}

\subsection{Free heat moments with time derivatives}
\label{sec:newS7V20-moments}

\begin{lemma}[Differentiated free heat limit]
\label{lem:newS7V20-heat-derivatives}
For \(j=0,1,2,3\), on every regime
\(a^2/t\to0\),
\begin{equation}
 t^{2+j}k_a^{(j)}(t)
 \longrightarrow
 {4(-1)^j(j+1)!\over(4\pi)^2}.                            \label{eq:newS7V20-heat-derivatives}
\end{equation}
On \(a/\sqrt t\le c_0\), the convergence has the common majorant
\begin{equation}
 \left|t^{2+j}k_a^{(j)}(t)
 -{4(-1)^j(j+1)!\over(4\pi)^2}\right|
 \le C_j\left({a\over\sqrt t}+e^{-ct/a^2}\right).        \label{eq:newS7V20-heat-derivative-error}
\end{equation}
\end{lemma}

\begin{proof}
In the free Fourier representation,
\(k_a^{(j)}(t)\) inserts \((-\lambda_a(p))^j\) into
\(4\int_{\mathrm{BZ}_a}e^{-t\lambda_a(p)}\,dp/(2\pi)^4\).
Rescale \(p=q/\sqrt t\).  The physical overlap branch converges with all
fixed polynomial moments to \(|q|^2\), while the complementary
Brillouin zone has the V14 Gaussian majorant and contributes
\(O(e^{-ct/a^2})\).  The V16 all-moment estimate dominates the factors
\(|q|^{2j}\).  Dominated convergence gives
\[
 4{d^j\over dt^j}\{(4\pi t)^{-2}\}
 ={4(-1)^j(j+1)!\over(4\pi)^2}t^{-2-j}.
\]
The quantitative low-momentum symbol comparison used in V14 is unchanged
by a fixed polynomial insertion and gives \(O(a/\sqrt t)\), proving
\eqref{eq:newS7V20-heat-derivative-error}.
\end{proof}

\subsection{Exact lattice subtraction and the quartic coefficient}
\label{sec:newS7V20-quartic}

Use the mass-squared variable \(s=m^2\) and define
\begin{equation}
 F_{a,t}(s):=e^{-ts}k_a\!\left(t\left(1-{a^2s\over4}\right)\right).
                                                               \label{eq:newS7V20-F}
\end{equation}
Thus \(F_{a,t}(m^2)=k_{a,m}(t)\).

\begin{lemma}[Exact quadratic and quartic lattice responses]
\label{lem:newS7V20-responses}
At \(s=0\),
\begin{align}
 C^{(2)}_{a,t}:=F_{a,t}'(0)
 &=-tk_a(t)-{a^2t\over4}k_a'(t),                          \label{eq:newS7V20-C2}\\
 C^{(4)}_{a,t}:={1\over2}F_{a,t}''(0)
 &={t^2\over2}\left{
 k_a(t)+{a^2\over2}k_a'(t)+{a^4\over16}k_a''(t)
 \right}.                                                \label{eq:newS7V20-C4}
\end{align}
Moreover,
\begin{equation}
 C^{(4)}_{a,t}\longrightarrow {2\over(4\pi)^2}
 \qquad(a^2/t\to0),                                      \label{eq:newS7V20-C4-limit}
\end{equation}
and, for \(a/\sqrt t\le c_0\),
\begin{equation}
 \left|C^{(4)}_{a,t}-{2\over(4\pi)^2}\right|
 \le C\left({a\over\sqrt t}+e^{-ct/a^2}\right).        \label{eq:newS7V20-C4-error}
\end{equation}
\end{lemma}

\begin{proof}
Differentiate \eqref{eq:newS7V20-F}.  If
\(\beta=a^2t/4\), then
\[
 F'(0)=-tk_a(t)-\beta k_a'(t),
 \qquad
 F''(0)=t^2k_a(t)+2t\beta k_a'(t)+\beta^2k_a''(t),
\]
which gives \eqref{eq:newS7V20-C2,eq:newS7V20-C4}.  By
\Cref{lem:newS7V20-heat-derivatives}, the first term in
\eqref{eq:newS7V20-C4} tends to \(2/(4\pi)^2\); the second and third are
bounded respectively by \(C a^2/t\) and \(C(a^2/t)^2\).  Both are bounded
by \(Ca/\sqrt t\) when \(a/\sqrt t\le1\).  This proves
\eqref{eq:newS7V20-C4-limit,eq:newS7V20-C4-error}.
\end{proof}

For \(m\ne0\), define the exact quadratically subtracted quotient
\begin{equation}
 \mathcal Q_{a,t}(m):=
 {k_{a,m}(t)-k_a(t)-m^2C^{(2)}_{a,t}\over m^4},           \label{eq:newS7V20-Q}
\end{equation}
and set \(\mathcal Q_{a,t}(0)=C^{(4)}_{a,t}\).

\begin{theorem}[Constant-mass quartic lattice transfer]
\label{thm:newS7V20-quartic-transfer}
For \(|m|\le M\), \(am\le1\), and \(a/\sqrt t\le c_0<1\),
\begin{equation}
 \left|\mathcal Q_{a,t}(m)-{2\over(4\pi)^2}\right|
 \le C_M\left{
 {a\over\sqrt t}+tm^2+e^{-ct/a^2}
 \right}.                                                \label{eq:newS7V20-quartic-transfer}
\end{equation}
In particular, the constant-Higgs quartic free-symbol word of one physical
massive-overlap block converges to \(2m^4/(4\pi)^2\) after its exact lattice
quadratic response has been subtracted.
\end{theorem}

\begin{proof}
Taylor's theorem applied to \eqref{eq:newS7V20-F} gives
\[
 F(m^2)=F(0)+m^2F'(0)+m^4{F''(0)\over2}
 +{m^6\over6}F'''(\theta m^2)
\]
for some \(\theta\in(0,1)\).  Three differentiations produce a sum of
terms
\[
 t^3\left({a^2\over4}\right)^r
 k_a^{(r)}\!\left(t\left(1-{a^2s\over4}\right)\right),
 \qquad 0\le r\le3,
\]
times bounded exponentials and numerical coefficients.  The time argument
is comparable to \(t\) because \(am\le1\).  The moment bounds in
\Cref{lem:newS7V20-heat-derivatives} show the complete third derivative is
bounded by \(C_Mt\): the \(r\)-th term has size
\(Ct(a^2/t)^r\).  Divide the Taylor remainder by \(m^4\), then use
\eqref{eq:newS7V20-C4-error}.  The definition at \(m=0\) follows by
continuity.
\end{proof}

\begin{corollary}[Constant Standard Model Higgs quartic]
\label{cor:newS7V20-SM-quartic}
Let \(H\) be constant on the broken chart.  Sum
\Cref{thm:newS7V20-quartic-transfer} over the massive-overlap blocks of
V18.  In the minimal or Dirac-neutrino branch,
\begin{equation}
 \sum_jm_j^4=\mathfrak b_\bullet|H|^4,                    \label{eq:newS7V20-mass-four-sum}
\end{equation}
including colour multiplicity.  Hence the physical block heat trace has
quartic coefficient
\begin{equation}
 {2\mathfrak b_\bullet\over(4\pi)^2}|H|^4.                \label{eq:newS7V20-physical-quartic}
\end{equation}
The error is bounded by
\begin{equation}
 C_Y|H|^4\left({a\over\sqrt t}+e^{-ct/a^2}
ight)
 +C_Yt|H|^6.                                               \label{eq:newS7V20-SM-error}
\end{equation}
\end{corollary}

\begin{proof}
The singular masses are \(m_j=|H|\sigma_j\).  Summing their fourth powers
gives precisely the trace invariant
\(\mathfrak b_\bullet\) computed in V18, including three identical colour
copies of each quark singular value.  Sum
\eqref{eq:newS7V20-quartic-transfer}; the finite generation and species
sum is absorbed in \(C_Y\).
\end{proof}

\begin{corollary}[The three-level multiplicity ledger]
\label{cor:newS7V20-multiplicity}
For the quartic Higgs owner, before the common factor
\((4\pi)^{-2}\) and any overall effective-action sign, the coefficients are
\begin{equation}
 \begin{array}{c|c}
 \text{object}&|H|^4\text{ coefficient}\\ \hline
 \text{self-adjoint doubled }K_L\oplus K_R\text{ heat square}
   &4\mathfrak b_\bullet\\
 \text{one physical massive-block heat trace}
   &2\mathfrak b_\bullet\\
 \log|\det\mathcal B|={1\over2}\log\det(\mathcal B^*\mathcal B)
   &\mathfrak b_\bullet.
 \end{array}                                               \label{eq:newS7V20-multiplicity}
\end{equation}
\end{corollary}

\begin{proof}
The first row is V18.  The second is
\Cref{cor:newS7V20-SM-quartic}.  The last follows from the factor
\(1/2\) in the determinant-modulus identity.  The phase remains a separate
determinant-line object.
\end{proof}

\begin{corollary}[Mesoscopic integrability]
\label{cor:newS7V20-integrability}
On the established window
\(t\ge\varepsilon_q(a)=\ell^2(a/\ell)^q\), \(0<q<2/3\), the first error in
\eqref{eq:newS7V20-SM-error} is logarithmically integrable with vanishing
norm:
\begin{equation}
 \int_{\varepsilon_q(a)}^{t_0}
 {a\over\sqrt t}{dt\over t}
 =O(a^{1-q/2}).                                            \label{eq:newS7V20-integrability}
\end{equation}
The \(t|H|^6\) remainder is integrable at zero.
\end{corollary}

\begin{proof}
The first statement is the direct integral already evaluated in V19.
The second follows from \(\int_0^{t_0}t\,dt/t=t_0\).
\end{proof}

\begin{assessment}[FGYSC progress]
\label{ass:newS7V20-FGYSC}
The constant-Higgs quartic survivor is now closed on the free broken-phase
overlap symbol, including its exact lattice lower-order subtraction and
physical multiplicity.  This does not close the Higgs-gradient survivor:
there the two first jets, the constant--second-jet contact, and the
improved-ordering commutator must be combined before a continuum kinetic
coefficient is asserted.
\end{assessment}

\begin{programme}[V21 acquisition]
\label{prog:newS7V20-V21}
V21 will compute the Higgs-gradient free-symbol response.  It will introduce
a plane-wave Higgs perturbation, keep both Yukawa corners and the
\(A_D\) ordering, differentiate the complete heat trace twice in its
amplitude, and only then take the small external-momentum Hessian.  This
upstream two-point calculation will decide the coefficient of
\(|DH|^2\) and expose every contact term explicitly.
\end{programme}

\begin{firewall}[Claim boundary after seventh-series V20]
\label{fire:newS7V20-boundary}
V20 audits YM V93 and NS V83, proves the exact massive-overlap square and
heat identity, proves differentiated free heat convergence, closes the
constant-mass quartic lattice coefficient with an exact quadratic
subtraction, sums it over the Standard Model Yukawa singular values, and
fixes the three-level multiplicity ledger.

V20 does not prove the Higgs-gradient free-symbol response, the full FGYSC,
the pre-broken electroweak heat transfer, the determinant-line phase, rough
AOP control, measure concentration, interacting cutoff removal, or the
Standard Model--Einstein continuum.
\end{firewall}

\begin{theorem}[Seventh-series V20 result]
\label{thm:newS7V20-result}
For one massive-overlap block,
\[
 D_m^*D_m=\left(1-{a^2m^2\over4}\right)D^*D+m^2I.
\]
After subtracting the exact lattice response linear in \(m^2\), its
quartic heat coefficient converges with error
\[
 O(a/\sqrt t)+O(tm^2)+O(e^{-ct/a^2})
\]
to \(2m^4/(4\pi)^2\).  Summing Standard Model blocks gives
\(2\mathfrak b_\bullet|H|^4/(4\pi)^2\) in the physical block heat trace,
half the doubled V18 heat coefficient; the determinant modulus carries one
further factor \(1/2\).
\end{theorem}

\begin{proof}
Use \Cref{lem:newS7V20-massive-square,
cor:newS7V20-massive-heat,
lem:newS7V20-heat-derivatives,
lem:newS7V20-responses,
thm:newS7V20-quartic-transfer,
cor:newS7V20-SM-quartic,
cor:newS7V20-multiplicity}.
\end{proof}

\paragraph{Parent-version record.}
V20 was formed from
\texttt{Renewal\_SM\_Einstein\_Continuum\_Research\_Notes\_V19.tex}.
The V19 parent had \(8\,252\) logical source lines, \(323\,807\) bytes, and
SHA--256
\[
 \texttt{78D16412440AEB8DAA6D127F1364CE9DE67D6A05AF9AA13E2E509B79F2A683A7}.
\]
The next compulsory companion audit is V25.

% =============================================================================
% END APPENDED SEVENTH-SERIES V20 MATERIAL
% =============================================================================


% =============================================================================
% BEGIN APPENDED SEVENTH-SERIES V21 MATERIAL
% =============================================================================

\clearpage
\section{Seventh-series V21: the Higgs-gradient two-point symbol and its
physical coefficient}
\label{sec:v7v21}

V20 closes the constant-Higgs quartic word.  The remaining pure-Higgs
weight-four word is the covariant gradient pair.  A Taylor estimate of two
separated first jets would miss the constant--second-jet contact identified
in V19.  This version therefore uses a plane-wave Higgs perturbation and
differentiates the complete heat trace in its amplitude before taking the
external-momentum Hessian.

\subsection{The complete scalar-amplitude Hessian}
\label{sec:newS7V21-Hessian}

Work in one matched left--right representation at zero gauge and Higgs
background.  Let \(D=D_{\rm ov}(1)\),
\(A_D=I-aD/2\), and let \(J_+\) be a normalized chiral Yukawa corner whose adjoint pair
describes one physical Dirac mass; put \(J_-=J_+^*\).  For an allowed torus momentum \(q\),
define the two oriented multiplication operators
\begin{equation}
 M_+(q;x):=e^{iq\cdot x}J_+,
 \qquad M_-(q;x):=e^{-iq\cdot x}J_- .                     \label{eq:newS7V21-oriented-M}
\end{equation}
The zero-winding Brillouin formula extends the response analytically
to continuous \(q\) near zero.  For formal complex amplitudes \(z,\bar z\), set
\begin{equation}
 M_{z,q}:=zM_+(q)+\bar zM_-(q),\qquad
 B_{z,\bar z}(q):=D+M_{z,q}A_D,qquad
 L_{z,\bar z}(q):=B_{z,\bar z}(q)^*B_{z,\bar z}(q).       \label{eq:newS7V21-Bz}
\end{equation}
On the physical slice \(\bar z=\overline z\), the two Fourier components
are adjoints and the mass perturbation is self-adjoint in the doubled
left--right fibre.  Wirtinger derivatives at the origin give the analytic
oriented two-point response, including at \(q=0\).

\begin{lemma}[First and mixed contact vertices]
\label{lem:newS7V21-vertices}
The first amplitude derivatives and the mixed contact derivative are
\begin{align}
 V_+(q)&:=\left.\partial_zL_{z,\bar z}(q)\right|_0
 =D^*M_+(q)A_D+A_D^*M_+(q)D,                              \label{eq:newS7V21-Vplus}\\
 V_-(q)&:=\left.\partial_{\bar z}L_{z,\bar z}(q)\right|_0
 =D^*M_-(q)A_D+A_D^*M_-(q)D,                              \label{eq:newS7V21-Vminus}\\
 W(q)&:=\left.\partial_z\partial_{\bar z}
 L_{z,\bar z}(q)\right|_0                                \notag\\
 &=A_D^*\{M_+(q)M_-(q)+M_-(q)M_+(q)\}A_D.                \label{eq:newS7V21-W}
\end{align}
No commutation of a mass multiplication and \(A_D\) is used.
\end{lemma}

\begin{proof}
Expand \(B_{z,\bar z}^*B_{z,\bar z}\) as a polynomial in
\(z,\bar z\).  The coefficients linear in the two amplitudes are the first
two displayed expressions.  The coefficient of \(z\bar z\) is the sum of
the two possible mass orders, giving \eqref{eq:newS7V21-W}.
\end{proof}

Define the oriented mixed Hessian per unit physical volume
\begin{equation}
 \Pi_{a,t}(q):={1\over|M|}
 \left.\partial_z\partial_{\bar z}
 \Tr e^{-tL_{z,\bar z}(q)}\right|_{z=0}.                  \label{eq:newS7V21-Pi}
\end{equation}

\begin{theorem}[Assembled heat Hessian]
\label{thm:newS7V21-assembled}
At every finite cutoff,
\begin{align}
 \Pi_{a,t}(q)={1\over|M|}\Bigg[
 &\int_0^t(t-u)\Tr\big\{
 V_+(q)e^{-uL_0}V_-(q)e^{-(t-u)L_0}                       \notag\\
 &\hspace{7.4em}+
 V_-(q)e^{-uL_0}V_+(q)e^{-(t-u)L_0}\big\}\,du            \notag\\
 &-t\Tr\{W(q)e^{-tL_0}\}\Bigg],                         \label{eq:newS7V21-assembled}
\end{align}
where \(L_0=D^*D\).  The formula is real and even in \(q\).
\end{theorem}

\begin{proof}
Differentiate Duhamel's formula first in one amplitude and then in the
other.  The two orders of the first vertices give the two integrands with
the simplex weight \(t-u\); differentiating the operator itself gives the
negative contact term.  This proves \eqref{eq:newS7V21-assembled} without
splitting a complete word.  Adjunction exchanges \(V_+(q)\) and
\(V_-(q)\), while \(q\mapsto-q\) exchanges the two orientations, proving
reality and evenness.
\end{proof}

\begin{assessment}[The contact is retained]
\label{ass:newS7V21-contact}
For the declared ordering,
\(M_+(q)M_-(q)+M_-(q)M_+(q)\) is independent of \(q\), so its direct
contribution cancels in \(\Pi_{a,t}(q)-\Pi_{a,t}(0)\).  It is nevertheless
retained in \eqref{eq:newS7V21-assembled}: it fixes the weight-one mass
response and is part of the complete amplitude Hessian before that
subtraction.
\end{assessment}

\subsection{Brillouin representation and low-branch consistency}
\label{sec:newS7V21-Brillouin}

Let \(\mathcal D_a(p)\) be the free overlap symbol and
\begin{equation}
 \mathcal A_a(p):=I-{a\over2}\mathcal D_a(p).              \label{eq:newS7V21-A-symbol}
\end{equation}
The oriented vertices shift momentum by \(q\) and \(-q\).  Thus
\eqref{eq:newS7V21-assembled} is a Brillouin integral of a complete
two-vertex divided difference plus its contact term.  Denote its assembled
integrand by \(\mathfrak J_a(p,q;t)\).

\begin{lemma}[Yukawa-vertex consistency through two external derivatives]
\label{lem:newS7V21-consistency}
There are \(\delta,C>0\) and an integer \(N\) such that, whenever
\(|ap|+|aq|\le\delta\), the complete first and contact Yukawa vertices
obtained from
\eqref{eq:newS7V21-Vplus,eq:newS7V21-Vminus,eq:newS7V21-W} differ from their continuum
counterparts by
\begin{equation}
 C a(1+|p|+|q|)^N                                         \label{eq:newS7V21-vertex-error}
\end{equation}
through two derivatives in \(q\).  After
\(p=k/\sqrt t\), the relative dimensionless error is bounded by
\begin{equation}
 C{a\over\sqrt t}(1+|k|)^N.                               \label{eq:newS7V21-scaled-error}
\end{equation}
\end{lemma}

\begin{proof}
The V14 physical-branch expansion gives
\[
 \mathcal D_a(p)=i\gamma\cdot p+O(a|p|^2),
 \qquad
 \mathcal A_a(p)=I+O(a|p|).
\]
The Yukawa multiplication vertex has no lattice derivative: its only
symbol factors are \(J_\pm\), \(\mathcal A_a(p)\), and one adjacent
\(\mathcal D_a(p\pm q)\).  Differentiate the analytic physical-branch
symbols in \(q\) at most twice and replace one lattice factor at a time by
its continuum limit.  Taylor's theorem gives
\eqref{eq:newS7V21-vertex-error}; the contact vertex is simpler because it
contains only two \(\mathcal A_a\) factors.  Heat scaling gives
\eqref{eq:newS7V21-scaled-error}.
\end{proof}

\begin{lemma}[Common Gaussian majorant for the assembled Higgs Hessian]
\label{lem:newS7V21-majorant}
Choose the V14 physical-branch cutoff \(\chi(ap)\).  For
\(t\ge a^2\), after \(p=k/\sqrt t\),
\begin{align}
 t^{-2}\left|
 \partial_{q_\mu}\partial_{q_\nu}
 \{\chi(ap)\mathfrak J_a(p,q;t)\}_{q=0}
 \right|
 &\le C(1+|k|)^{N_1}e^{-c|k|^2},                         \label{eq:newS7V21-majorant}\\
 t^{-2}\left|
 \partial_q^2\{\chi(ap)\mathfrak J_a-\mathfrak J_D\}_{q=0}
 \right|
 &\le C{a\over\sqrt t}(1+|k|)^{N_1}e^{-c|k|^2}.          \label{eq:newS7V21-majorant-error}
\end{align}
On the complementary branch,
\begin{equation}
 \left|\int {d^4p\over(2\pi)^4}
 \partial_q^2\{(1-\chi(ap))\mathfrak J_a(p,q;t)\}_{q=0}
 \right|
 \le Ce^{-ct/a^2}.                                       \label{eq:newS7V21-high}
\end{equation}
\end{lemma}

\begin{proof}
Use the smooth heat divided difference in the assembled formula
\eqref{eq:newS7V21-assembled}.  On the low branch, the free dispersion is
bounded below by \(c|p|^2\).  Each differentiated vertex is polynomial in
\(p\), while every derivative of a heat spectral argument supplies the
matching power of \(t\).  The two external derivatives have net factor
\(t^2\), as in the V14 gauge-polarization proof, leaving a polynomial in
\(k\) times \(e^{-c|k|^2}\).  This proves
\eqref{eq:newS7V21-majorant}.  The telescoping replacement based on
\Cref{lem:newS7V21-consistency} proves
\eqref{eq:newS7V21-majorant-error}.

Outside the physical branch, every heat factor has spectral floor
\(c/a^2\).  The overlap and improved-field vertices and their first two
momentum derivatives have only fixed powers of \(a^{-1}\).  The Brillouin
volume and these powers form a polynomial in \(t/a^2\), which is absorbed
by decreasing the exponent in \(e^{-ct/a^2}\).  This proves
\eqref{eq:newS7V21-high}.
\end{proof}

\subsection{The kinetic coefficient}
\label{sec:newS7V21-coefficient}

Hypercubic symmetry and analyticity give a scalar \(\kappa_{a,t}\) such
that
\begin{equation}
 \Pi_{a,t}(q)-\Pi_{a,t}(0)
 =\kappa_{a,t}|q|^2+O(|q|^4).                             \label{eq:newS7V21-kappa-def}
\end{equation}

\begin{lemma}[Continuum normalization]
\label{lem:newS7V21-continuum}
For one physical four-component Dirac block with a real scalar mass field,
the continuum value is
\begin{equation}
 \kappa_D={2\over(4\pi)^2}.                               \label{eq:newS7V21-continuum}
\end{equation}
\end{lemma}

\begin{proof}
For the continuum square, the mass-dependent Laplace endomorphism is
\(E_m=m^2-\gamma^\mu\partial_\mu m\), up to the harmless sign convention
for the derivative term.  In the integrated \(A_4\) coefficient,
\(\frac12\tr_S(E_m^2)\) contributes
\[
 {1\over2}\tr_S
 \{(-\gamma^\mu\partial_\mu m)^2\}
 =2|\partial m|^2
\]
in the positive Euclidean square convention established in V17.  For the oriented perturbation, the
mixed amplitude Hessian of
this quadratic functional is exactly
\(2|q|^2/(4\pi)^2\), proving
\eqref{eq:newS7V21-continuum}.
\end{proof}

\begin{theorem}[Higgs-gradient free-symbol transfer]
\label{thm:newS7V21-gradient-transfer}
There are \(C,c>0\) such that, whenever \(a^2/t\) is sufficiently small,
\begin{equation}
 \left|\kappa_{a,t}-{2\over(4\pi)^2}\right|
 \le C\left\{{a\over\sqrt t}+e^{-ct/a^2}\right\}.       \label{eq:newS7V21-gradient-transfer}
\end{equation}
Thus the physical massive-block heat trace produces
\(2|Dm|^2/(4\pi)^2\) for a slowly varying mass field.
\end{theorem}

\begin{proof}
By \eqref{eq:newS7V21-kappa-def}, \(\kappa_{a,t}\) is a fixed contraction
of \(\partial_q^2\Pi_{a,t}(0)\).  On the physical branch substitute
\(p=k/\sqrt t\).  The measure factor \(t^{-2}\) cancels the factor in
\eqref{eq:newS7V21-majorant}; dominated convergence applies, and
\eqref{eq:newS7V21-majorant-error} gives the error
\(Ca/\sqrt t\).  The limiting integral is the continuum assembled
amplitude Hessian, whose value is
\eqref{eq:newS7V21-continuum}.  Equation
\eqref{eq:newS7V21-high} supplies the complementary-branch error.
\end{proof}

\begin{corollary}[Standard Model Higgs kinetic coefficient]
\label{cor:newS7V21-SM-kinetic}
Summing the V18 Yukawa intertwiners, the physical Standard Model block heat
trace has the smooth Higgs kinetic coefficient
\begin{equation}
 {2\mathfrak a_\bullet\over(4\pi)^2}|D_\mu H|^2,          \label{eq:newS7V21-SM-kinetic}
\end{equation}
with error
\begin{equation}
 C_Y\left\{{a\over\sqrt t}+e^{-ct/a^2}\right\}|DH|^2.   \label{eq:newS7V21-SM-error}
\end{equation}
The optional Dirac-neutrino contribution is included by
\(\mathfrak a_\nu=\mathfrak a_{\min}+\Tr(Y_\nu^*Y_\nu)\).
\end{corollary}

\begin{proof}
Polarization of the quadratic response replaces the unit intertwiner by
each Yukawa matrix.  The flavour trace is
\(\Tr(Y_s^*Y_s)\); colour gives three copies for each quark.  Their sum is
exactly \(\mathfrak a_\bullet\) from V18.  Gauge covariance replaces the
ordinary derivative by \(D_\mu H\), and V19's finite word enumeration shows
that no additional background-Higgs insertion can enter at heat weight two.
Sum \eqref{eq:newS7V21-gradient-transfer} over the finite species list.
\end{proof}

\begin{corollary}[Kinetic multiplicity ledger]
\label{cor:newS7V21-multiplicity}
Before the common factor \((4\pi)^{-2}\) and any overall effective-action
sign, the coefficients of \(|DH|^2\) are
\begin{equation}
 \begin{array}{c|c}
 \text{object}&|DH|^2\text{ coefficient}\\ \hline
 \text{self-adjoint doubled }K_L\oplus K_R\text{ heat square}
   &4\mathfrak a_\bullet\\
 \text{physical chiral-block heat trace}&2\mathfrak a_\bullet\\
 \log|\det\mathcal B|={1\over2}\log\det(\mathcal B^*\mathcal B)
   &\mathfrak a_\bullet.
 \end{array}                                               \label{eq:newS7V21-multiplicity}
\end{equation}
\end{corollary}

\begin{proof}
The doubled value is V18, the physical heat value is
\Cref{cor:newS7V21-SM-kinetic}, and the determinant modulus supplies the
factor one half.
\end{proof}

\begin{corollary}[Mesoscopic kinetic integrability]
\label{cor:newS7V21-integrability}
On \(t\ge\varepsilon_q(a)=\ell^2(a/\ell)^q\),
\(0<q<2/3\), the coefficient error in
\eqref{eq:newS7V21-SM-error} is logarithmically integrable and tends to
zero in norm, with contribution \(O(a^{1-q/2})\) plus an exponentially
small term.
\end{corollary}

\begin{proof}
Integrate \(a/\sqrt t\) against \(dt/t\) as in V19--V20.  Substitute
\(u=t/a^2\) for the exponential term.
\end{proof}

\begin{assessment}[FGYSC status]
\label{ass:newS7V21-FGYSC}
The free broken-phase \(|H|^4\) and \(|DH|^2\) symbol rows are now closed,
as is the earlier \(F^2\) row.  The remaining flat weight-four check is the
lattice suppression of the mixed classes which vanish in the continuum.
After that finite check, the smooth local covariant parametrix must still be
run for the full chiral block; rough AOP fields remain separate.
\end{assessment}

\begin{programme}[V22 acquisition]
\label{prog:newS7V21-V22}
V22 will treat the mixed gauge--Higgs words.  It will differentiate the
complete assembled heat trace jointly in one gauge amplitude and two Higgs
amplitudes, retain every contact derivative, and prove that the parity-even
\(F|H|^2\) coefficient tends to zero.  The odd-Yukawa classes will be
removed at finite cutoff by the exact \(K_L\oplus K_R\) grading before any
momentum integral is estimated.
\end{programme}

\begin{firewall}[Claim boundary after seventh-series V21]
\label{fire:newS7V21-boundary}
V21 constructs the complete improved-overlap Higgs amplitude Hessian,
retains its contact term, proves low-branch vertex consistency and a common
Gaussian majorant, suppresses the high branch, transfers the physical
Higgs-gradient coefficient, and fixes its Standard Model and determinant
multiplicities.

V21 does not close the mixed gauge--Higgs lattice words, the full smooth
FGYSC parametrix, the pre-broken electroweak determinant phase, rough AOP
control, measure concentration, cutoff removal, or the Standard
Model--Einstein continuum.
\end{firewall}

\begin{theorem}[Seventh-series V21 result]
\label{thm:newS7V21-result}
For the improved overlap--Yukawa ordering, the complete plane-wave Higgs
two-point heat response has a mesoscopic external-momentum coefficient
\[
 \kappa_{a,t}={2\over(4\pi)^2}
 +O(a/\sqrt t)+O(e^{-ct/a^2}).
\]
Consequently the physical Standard Model block heat trace produces
\[
 {2\mathfrak a_\bullet\over(4\pi)^2}|D_\mu H|^2,
\]
half the doubled V18 coefficient; the determinant modulus carries one
further factor \(1/2\).  The error is logarithmically integrable on the
established mesoscopic window.
\end{theorem}

\begin{proof}
Use \Cref{thm:newS7V21-assembled,
lem:newS7V21-consistency,
lem:newS7V21-majorant,
lem:newS7V21-continuum,
thm:newS7V21-gradient-transfer,
cor:newS7V21-SM-kinetic,
cor:newS7V21-multiplicity}.
\end{proof}

\paragraph{Parent-version record.}
V21 was formed from
\texttt{Renewal\_SM\_Einstein\_Continuum\_Research\_Notes\_V20.tex}.
The V20 parent had \(8\,657\) logical source lines, \(339\,111\) bytes, and
SHA--256
\[
 \texttt{4737F516E6586F92C8651784BAC453085027C134F04A70D5F4D53A02A489CF71}.
\]
The next compulsory companion audit is V25.

% =============================================================================
% END APPENDED SEVENTH-SERIES V21 MATERIAL
% =============================================================================


% =============================================================================
% BEGIN APPENDED SEVENTH-SERIES V22 MATERIAL
% =============================================================================

\clearpage
\section{Seventh-series V22: exact Yukawa parity and suppression of the
curvature--two-Higgs channel}
\label{sec:v7v22}

V21 closes the pure Higgs-gradient two-point symbol.  V19 left three mixed
word classes: two have odd Yukawa parity, and the third contains one frozen
curvature insertion and two constant Higgs insertions.  This version removes
the odd classes at finite cutoff and proves mesoscopic convergence of the
curvature channel to its zero continuum Clifford trace.  The ordinary gauge
completion of \(|DH|^2\) is retained and is not part of that curvature
projection.

\subsection{The exact finite-cutoff internal grading}
\label{sec:newS7V22-grading}

On the doubled finite internal space \(K_L\oplus K_R^\bullet\), define
\begin{equation}
 \tau:=\begin{pmatrix}I_{K_L}&0\\0&-I_{K_R^\bullet}\end{pmatrix}.
                                                               \label{eq:newS7V22-tau}
\end{equation}
Every gauge kinetic overlap operator and improved-field factor is block
diagonal in this decomposition, while the V18 Yukawa map is off diagonal.
Thus
\begin{equation}
 [\tau,D_L\oplus D_R]=[\tau,A_L\oplus A_R]=0,
 \qquad \{\tau,\Phi(H)\}=0.                               \label{eq:newS7V22-parities}
\end{equation}

\begin{lemma}[Odd Yukawa trace vanishes before estimates]
\label{lem:newS7V22-odd-trace}
Let \(X_0,\ldots,X_n\) commute with \(\tau\), and let
\(Y_1,\ldots,Y_{2r+1}\) anticommute with \(\tau\).  Whenever the product is
trace class,
\begin{equation}
 \Tr\{X_0Y_1X_1\cdots Y_{2r+1}X_n\}=0.                   \label{eq:newS7V22-odd-trace}
\end{equation}
The statement remains true when the \(X_j\) contain heat factors,
resolvents, gauge vertices, fixed spin projectors, or moving
Ginsparg--Wilson projectors.
\end{lemma}

\begin{proof}
Conjugation by \(\tau\) preserves every \(X_j\) and changes the sign of
every \(Y_j\).  Cyclicity and \(\tau^2=I\) therefore give
\[
 \Tr(X_0Y_1\cdots Y_{2r+1}X_n)
 =-\Tr(X_0Y_1\cdots Y_{2r+1}X_n).
\]
The trace is zero.  The listed operators commute with \(\tau\) because
they preserve the declared left and right internal representations; their
spin chirality does not alter the internal grading.
\end{proof}

\begin{corollary}[Two V19 mixed classes vanish exactly]
\label{cor:newS7V22-odd-classes}
Every complete heat word in the classes
\(\mathsf g\mathsf h_1\) and
\(\mathsf h_0^2\mathsf h_1\) has zero finite-cutoff trace.  This remains
true after all contact derivatives are included.
\end{corollary}

\begin{proof}
The first class has one Yukawa insertion and the second has three.  Gauge
and contact factors preserve \(K_L\oplus K_R\).  Apply
\Cref{lem:newS7V22-odd-trace} to each complete term of the assembled
divided difference.
\end{proof}

\subsection{Separating curvature from the kinetic gauge completion}
\label{sec:newS7V22-projection}

The third derivative with one gauge and two Higgs amplitudes also contains
the cubic gauge completion of the already determined invariant
\(|DH|^2\).  It must not be set to zero.  In the V15 radial gauge at
\(x_0\), that completion is carried by Higgs first jets.  The V19 class
\(\mathsf g\mathsf h_0^2\), by contrast, keeps both Higgs insertions
constant and replaces the gauge link jet by the frozen antisymmetric
curvature
\begin{equation}
 A_\mu^{F_0}(y)={1\over2}y^\nu F_{\nu\mu}(x_0).            \label{eq:newS7V22-radial-curvature}
\end{equation}

\begin{definition}[Curvature--two-Higgs projection]
\label{def:newS7V22-CHP}
CHP is the complete third amplitude derivative of the diagonal heat density
with one gauge and two Yukawa derivatives, after:
\begin{enumerate}[label=\textup{(\roman*)},leftmargin=3.5em]
\item the two Higgs vertices are evaluated on the constant jet \(H_0\);
\item every derivative of the overlap operator and of
      \(A_D=I-aD/2\) is retained;
\item all first, second, and third derivatives of
      \(L=B^*B\) are included in the complete heat divided difference; and
\item the gauge jet is factored through the V13 linearized lattice
      curvature and contracted with the antisymmetric radial factor in
      \eqref{eq:newS7V22-radial-curvature}.
\end{enumerate}
Denote the resulting scalar per unit volume by
\(\mu_{a,t}(F_0;H_0)\).
\end{definition}

\begin{assessment}[What CHP excludes]
\label{ass:newS7V22-CHP-scope}
CHP excludes terms in which an external derivative lands on a Higgs field.
Those terms assemble the gauge completion of
\(|DH|^2\) and are fixed by V21 plus exact gauge covariance.  CHP isolates
only a possible local owner linear in curvature and quadratic in the
constant Higgs value.
\end{assessment}

\subsection{The complete third divided difference}
\label{sec:newS7V22-third}

Let \(\alpha,z,\bar z\) be respectively a gauge amplitude and the two
oriented Higgs amplitudes of V21.  Put
\begin{equation}
 L(\alpha,z,\bar z)
 :=B(\alpha,z,\bar z)^*B(\alpha,z,\bar z).                 \label{eq:newS7V22-Lthree}
\end{equation}

\begin{lemma}[Finite complete contact ledger]
\label{lem:newS7V22-contact-ledger}
The derivative
\begin{equation}
 \left.\partial_\alpha\partial_z\partial_{\bar z}
 \Tr e^{-tL(\alpha,z,\bar z)}\right|_0                    \label{eq:newS7V22-third-derivative}
\end{equation}
is the finite sum indexed by set partitions of
\(\{\alpha,z,\bar z\}\):
\begin{equation}
 \sum_{\pi}
 (-1)^{|\pi|}
 \int_{\Delta_{|\pi|}(t)}
 \Tr\left\{
 e^{-s_0L_0}
 \prod_{C\in\pi}^{\longrightarrow}
 L_C e^{-s_CL_0}
 \right\}d\mathbf s,                                    \label{eq:newS7V22-partition-formula}
\end{equation}
summed over the distinct time orders of the blocks, where
\(L_C=\partial_C L|_0\) and
\(\Delta_k(t)=\{s_j\ge0:\sum_{j=0}^ks_j=t\}\).
Thus all one-, two-, and three-parameter contacts occur before any norm is
taken.
\end{lemma}

\begin{proof}
Apply Duhamel's formula successively in the three parameters.  Whenever a
new derivative lands on a heat factor it creates a new time-ordered block;
when it lands on an existing vertex it joins that derivative to the same
block.  These alternatives are exactly the set partitions and their time
orders.  The sign is the number of heat insertions.  Since there are three
parameters, the ledger is finite and contains every derivative of
\(B^*B\) forced by the product rule.
\end{proof}

\begin{remark}[No factorial loss]
\label{rem:newS7V22-no-factorial}
Formula \eqref{eq:newS7V22-partition-formula} is differentiated only to the
fixed order three.  The arbitrary gauge background dependence inside each
\(L_C\) is still represented by the normalized V19 resolvent words, so the
all-order gauge sum remains geometric rather than factorial.
\end{remark}

\subsection{Low-branch convergence of CHP}
\label{sec:newS7V22-convergence}

After the V13 curvature factorization, write the zero-winding CHP as a
Brillouin integral of its complete integrand
\(\mathfrak K_a(p;t;F_0,H_0)\).

\begin{lemma}[Physical-branch consistency of the complete mixed word]
\label{lem:newS7V22-low}
On \(|ap|\le\delta\), after \(p=k/\sqrt t\),
\begin{align}
 t^{-2}|\chi(ap)\mathfrak K_a(p;t;F_0,H_0)|
 &\le C|F_0|\,|H_0|^2(1+|k|)^Ne^{-c|k|^2},               \label{eq:newS7V22-low-majorant}\\
 t^{-2}|\chi(ap)(\mathfrak K_a-\mathfrak K_D)|
 &\le C{a\over\sqrt t}|F_0|\,|H_0|^2
 (1+|k|)^Ne^{-c|k|^2}.                                   \label{eq:newS7V22-low-error}
\end{align}
\end{lemma}

\begin{proof}
Differentiate the physical-branch overlap and improved-field symbols in the
three amplitudes before expanding them.  V14 supplies the gauge and contact
vertex comparison through the curvature derivative; V21 supplies the two
oriented Yukawa vertices and their contacts.  Replacing one lattice factor
at a time by its continuum counterpart costs
\(Ca/\sqrt t\).  The smooth heat divided differences in
\Cref{lem:newS7V22-contact-ledger} supply the matching powers of \(t\) for
every derivative placement.  After the heat rescaling, a fixed polynomial
in \(k\) remains and is dominated by the free Gaussian.  This proves both
estimates.
\end{proof}

\begin{lemma}[Complementary-branch suppression]
\label{lem:newS7V22-high}
For \(t\ge a^2\),
\begin{equation}
 \left|\int {d^4p\over(2\pi)^4}
 (1-\chi(ap))\mathfrak K_a(p;t;F_0,H_0)\right|
 \le C|F_0|\,|H_0|^2e^{-ct/a^2}.                         \label{eq:newS7V22-high}
\end{equation}
\end{lemma}

\begin{proof}
The complementary branch has the V8 spectral floor \(c/a^2\).  The finite
contact ledger and the first two amplitude derivatives of the overlap and
improved-field vertices contribute only a fixed polynomial in
\(a^{-1}\) and \(t\).  Absorb that polynomial and the Brillouin volume into
the exponential.
\end{proof}

\begin{lemma}[The continuum CHP is zero]
\label{lem:newS7V22-continuum-zero}
For every \(F_0,H_0\),
\begin{equation}
 \mu_{D}(F_0;H_0)=0.                                      \label{eq:newS7V22-continuum-zero}
\end{equation}
\end{lemma}

\begin{proof}
In the continuum Dirac--Yukawa square, the curvature insertion in the
Laplace endomorphism is
\(-\frac12\gamma^{\mu\nu}\rho_*(F_{\mu\nu})\), while two constant Higgs
insertions assemble \(\Phi(H_0)^2\), a Clifford scalar.  Their only
weight-four cross trace is proportional to
\[
 \tr_S(\gamma^{\mu\nu})
 \tr_{\rm int}\{\rho_*(F_{\mu\nu})\Phi(H_0)^2\}=0.
\]
The connection-curvature square contains no Higgs field.  Hence the complete
continuum coefficient in the CHP channel vanishes, as independently found
in V17.
\end{proof}

\begin{theorem}[Mesoscopic suppression of curvature times two Higgs fields]
\label{thm:newS7V22-mixed-zero}
Whenever \(a^2/t\) is sufficiently small,
\begin{equation}
 |\mu_{a,t}(F_0;H_0)|
 \le C|F_0|\,|H_0|^2
 \left\{{a\over\sqrt t}+e^{-ct/a^2}\right\}.             \label{eq:newS7V22-mixed-zero}
\end{equation}
Therefore the lattice overlap--Yukawa heat trace has no continuum
\(F|H|^2\) owner.
\end{theorem}

\begin{proof}
On the physical branch, rescale \(p=k/\sqrt t\).  The measure cancels the
factor \(t^2\) in \eqref{eq:newS7V22-low-majorant}.  Dominated convergence
and \eqref{eq:newS7V22-low-error} compare CHP with its continuum value,
which is zero by \Cref{lem:newS7V22-continuum-zero}.  Add the high-branch
bound \eqref{eq:newS7V22-high}.
\end{proof}

\begin{corollary}[Integrated mixed error]
\label{cor:newS7V22-integrability}
On the V8 window
\(t\ge\varepsilon_q(a)=\ell^2(a/\ell)^q\), \(0<q<2/3\), the right side of
\eqref{eq:newS7V22-mixed-zero} is logarithmically integrable with norm
tending to zero as \(a\downarrow0\).
\end{corollary}

\begin{proof}
Integrate the two terms against \(dt/t\) as in V20--V21.
\end{proof}

\subsection{The completed flat coefficient ledger}
\label{sec:newS7V22-ledger}

\begin{theorem}[All flat weight-four coefficient rows]
\label{thm:newS7V22-flat-ledger}
For the improved overlap--Yukawa block on the free physical branch:
\begin{enumerate}[label=\textup{(\roman*)},leftmargin=3.5em]
\item \(F^2\) has the V14--V16 coefficient;
\item \(|H|^4\) has the V20 physical coefficient
      \(2\mathfrak b_\bullet/(4\pi)^2\);
\item \(|DH|^2\) has the V21 physical coefficient
      \(2\mathfrak a_\bullet/(4\pi)^2\);
\item odd-Yukawa words vanish exactly by
      \Cref{cor:newS7V22-odd-classes}; and
\item the curvature--two-Higgs channel vanishes in the continuum with the
      error \eqref{eq:newS7V22-mixed-zero}.
\end{enumerate}
Thus the finite flat coefficient part of FGYSC is closed.  Its smooth local
covariant parametrix and rough-field extension remain separate.
\end{theorem}

\begin{proof}
Combine the cited results and
\Cref{thm:newS7V22-mixed-zero} with the V19 complete enumeration.
\end{proof}

\begin{programme}[V23 acquisition]
\label{prog:newS7V22-V23}
V23 will insert the completed flat coefficient ledger into a joint radial
gauge--Higgs parametrix.  It will keep constant Higgs, Higgs jets, frozen
curvature, and their lattice errors as separate heat-scale variables, then
prove a local smooth-background remainder after subtracting the physical
\(F^2\), \(|DH|^2\), and \(|H|^4\) owners.  Metric curvature and the
\(R|H|^2\) owner will remain outside that flat-chart theorem.
\end{programme}

\begin{firewall}[Claim boundary after seventh-series V22]
\label{fire:newS7V22-boundary}
V22 proves exact finite-cutoff vanishing of odd-Yukawa heat words, defines a
complete contact-preserving curvature--two-Higgs projection, proves its
low-branch convergence and high-branch suppression, and closes its zero
continuum coefficient.  Together with V14--V16 and V20--V21, this completes
the flat free-symbol weight-four coefficient ledger.

V22 does not prove the joint smooth local parametrix, the curved
\(R|H|^2\) lattice transfer, the determinant-line phase, rough AOP control,
measure concentration, interacting cutoff removal, or the Standard
Model--Einstein continuum.
\end{firewall}

\begin{theorem}[Seventh-series V22 result]
\label{thm:newS7V22-result}
Odd Yukawa word classes vanish exactly at finite cutoff.  After the complete
third heat derivative is projected onto one antisymmetric curvature and two
constant Higgs insertions,
\[
 |\mu_{a,t}(F;H)|
 \le C|F|\,|H|^2
 \{a/\sqrt t+e^{-ct/a^2}\},
\]
so no \(F|H|^2\) owner survives.  The flat physical overlap--Yukawa
coefficient ledger now consists precisely of the previously established
\(F^2\), \(2\mathfrak a_\bullet|DH|^2/(4\pi)^2\), and
\(2\mathfrak b_\bullet|H|^4/(4\pi)^2\) rows.
\end{theorem}

\begin{proof}
Use \Cref{lem:newS7V22-odd-trace,
cor:newS7V22-odd-classes,
lem:newS7V22-contact-ledger,
lem:newS7V22-low,
lem:newS7V22-high,
lem:newS7V22-continuum-zero,
thm:newS7V22-mixed-zero,
thm:newS7V22-flat-ledger}.
\end{proof}

\paragraph{Parent-version record.}
V22 was formed from
\texttt{Renewal\_SM\_Einstein\_Continuum\_Research\_Notes\_V21.tex}.
The V21 parent had \(9\,045\) logical source lines, \(354\,894\) bytes, and
SHA--256
\[
 \texttt{385E336333AFAEDC86B5005D7602FDE51E9A6B7ED1324141DDD82D4AB0D89F64}.
\]
The next compulsory companion audit is V25.

% =============================================================================
% END APPENDED SEVENTH-SERIES V22 MATERIAL
% =============================================================================


% =============================================================================
% BEGIN APPENDED SEVENTH-SERIES V23 MATERIAL
% =============================================================================

\clearpage
\section{Seventh-series V23: the joint smooth gauge--Higgs local
parametrix}
\label{sec:v7v23}

V22 completes the flat free-symbol coefficient ledger.  This version inserts
that ledger into one covariant radial parametrix for the full improved
chiral block.  The weight-one Higgs mass term is subtracted with its exact
lattice coefficient from V20; only after that subtraction are the physical
weight-four gauge, Higgs-gradient, and Higgs-quartic owners compared with
their continuum values.

\subsection{Physical Standard Model heat conventions}
\label{sec:newS7V23-conventions}

Let \(\mathscr W_{\rm SM}\) be the declared list of physical Weyl
multiplets, including generation and colour multiplicities but no undeclared
right-handed neutrino.  The optional Dirac-neutrino branch appends that
singlet explicitly.  For a block-diagonal gauge curvature define
\begin{equation}
 \mathfrak G_{\rm SM}(F)
 :=\sum_{f\in\mathscr W_{\rm SM}}
 \tr_{R_f}\{\rho_{f,*}(F_{\mu\nu})^*
ho_{f,*}(F_{\mu\nu})\}.             \label{eq:newS7V23-G}
\end{equation}
Thus one physical Weyl heat block contributes one half of the V16
four-component vectorlike coefficient, and the complete physical gauge
owner is
\begin{equation}
 \mathcal O_F(F):={1\over3(4\pi)^2}\mathfrak G_{\rm SM}(F).              \label{eq:newS7V23-OF}
\end{equation}

Let \(h_a^{\rm SM}(t;x;A,H)\) denote the zero-winding diagonal heat density
of \(\mathcal B_a(A,H)^*\mathcal B_a(A,H)\), summed over the physical
chiral source blocks.  Its determinant modulus has an additional overall
factor \(1/2\), which is not included in \(h_a^{\rm SM}\).

Use the exact free lattice response
\begin{equation}
 C^{(2)}_{a,t}=-tk_a(t)-{a^2t\over4}k_a'(t)                \label{eq:newS7V23-C2}
\end{equation}
from V20, and define the local weight-one subtraction
\begin{equation}
 \mathcal O^{(2)}_{a,t}(H)
 :=C^{(2)}_{a,t}\,\mathfrak a_\bullet|H|^2.               \label{eq:newS7V23-O2}
\end{equation}
The physical weight-four Higgs owner is
\begin{equation}
 \mathcal O_H(H)
 :={1\over(4\pi)^2}
 \{2\mathfrak a_\bullet|D_\mu H|^2
   +2\mathfrak b_\bullet|H|^4\}.                         \label{eq:newS7V23-OH}
\end{equation}

\begin{lemma}[The exact quadratic trace behind the mass subtraction]
\label{lem:newS7V23-quadratic-trace}
At every point,
\begin{equation}
 \sum_jm_j(H)^2=\mathfrak a_\bullet|H|^2,                 \label{eq:newS7V23-mass-two}
\end{equation}
where the sum runs over the physical massive blocks, with colour and
generation multiplicities.  Hence
\eqref{eq:newS7V23-O2} is precisely the sum of their exact free lattice
responses linear in mass squared.
\end{lemma}

\begin{proof}
The V18 orthogonal-column identity gives
\(\Tr(M_H^*M_H)=\mathfrak a_\bullet|H|^2\).  Its eigenvalues are the
squares of the physical singular masses, including multiplicities.  Apply
the V20 response \eqref{eq:newS7V23-C2} to each singular block and sum.
\end{proof}

\subsection{Joint radial jets and their budgets}
\label{sec:newS7V23-jets}

Fix a flat coordinate ball \(B_R(x_0)\).  Put the continuum gauge field in
exact radial gauge and radially parallel-transport the Higgs bundle to
\(x_0\).  Write
\begin{align}
 F_0&:=F(x_0),& H_0&:=H(x_0),& H_{1,\mu}&:=D_\mu H(x_0).                  \label{eq:newS7V23-jets}
\end{align}
For \(|y|\le\Lambda\sqrt t\), V16 and V19 give the heat-scale variables
\begin{align}
 \mathsf g&=O(t|F_0|),                                    \label{eq:newS7V23-g}\\
 \mathsf h_0&=O(\sqrt t\,y_*|H_0|),                      \label{eq:newS7V23-h0}\\
 \mathsf h_1&=O(t y_*|H_1|),                              \label{eq:newS7V23-h1}
\end{align}
and the joint jet error
\begin{align}
 \mathsf e_{a,t}
 =O\big(&t^{3/2}\|\nabla F\|
       +a\sqrt t\|F\|
       +at^{3/2}\|F\|^2                                      \notag\\
      &+t^{3/2}y_*\|D^2H\|
       +a\sqrt t\,y_*\|DH\|
 \big).                                                       \label{eq:newS7V23-error-budget}
\end{align}
Here every unlabelled norm is the supremum on \(B_R(x_0)\); the line breaks
in \eqref{eq:newS7V23-error-budget} use a common omitted subscript
\(L^\infty(B_R)\).  The final term is the improved-ordering commutator of
V19.

\begin{lemma}[Factorial-free joint local series]
\label{lem:newS7V23-local-series}
On a bounded smooth family and for sufficiently small \(a,t\), every
complete local heat word is dominated, after all of its vertex positions
and contacts have been assembled, by a coefficient of the convergent
multivariate series
\begin{equation}
 \sum_{n_g,n_0,n_1,n_e\ge0}
 C^{n_g+n_0+n_1+n_e}
 \mathsf g^{n_g}\mathsf h_0^{n_0}
 \mathsf h_1^{n_1}\mathsf e_{a,t}^{n_e}.                  \label{eq:newS7V23-geometric}
\end{equation}
No factorial depending on the all-order gauge word occurs.
\end{lemma}

\begin{proof}
V19 expands the assembled first-order block into one normalized resolvent
word at each gauge order and at most one Higgs multiplication per
first-order term.  The V12 weighted vertex trees sum the positions inside a
complete word.  The heat Duhamel simplex has volume \(t^r/r!\) at fixed
number \(r\) of assembled first-order insertions; retaining it can only
improve the geometric bound.  Insert the heat-scale estimates
\eqref{eq:newS7V23-g}--\eqref{eq:newS7V23-error-budget} and sum the resulting
products.  Smallness of the four variables makes
\eqref{eq:newS7V23-geometric} convergent.
\end{proof}

\subsection{The local smooth matter theorem}
\label{sec:newS7V23-local}

For compact notation, let \(\mathcal P_j(x_0)\) be a positive polynomial
in the bounded covariant jets of \(F,H\), the fixed Yukawa matrices, and a
reference inverse length, homogeneous of engineering dimension \(j\) in
its nonconstant terms.  It may be chosen uniformly on a bounded
\(C^2\)-gauge and \(C^3\)-Higgs family.  In particular, take
\begin{align}
 \mathcal P_4&\ge
 \mathfrak G_{\rm SM}(F_0)
 +\mathfrak a_\bullet|DH_0|^2
 +\mathfrak b_\bullet|H_0|^4,                             \label{eq:newS7V23-P4}\\
 \mathcal P_5&\ge
 |F_0|\,|\nabla F_0|+|DH_0|\,|D^2H_0|
 +|H_0|\,|D^3H_0|+y_*^4|H_0|^3|DH_0|.                   \label{eq:newS7V23-P5}
\end{align}

\begin{theorem}[Joint smooth gauge--Higgs local parametrix]
\label{thm:newS7V23-local-parametrix}
Let \((A,H)\) range over such a bounded smooth family on a flat chart, with
\(a^2/t\le c_0\).  Then
\begin{align}
 \Big|&h_a^{\rm SM}(t;x_0;A,H)-h_a^{\rm SM}(t;x_0;0,0)
 -\mathcal O^{(2)}_{a,t}(H_0)                             \notag\\
 &\hspace{4em}-\mathcal O_F(F_0)-\mathcal O_H(H_0)\Big|
 \le C\left\{
 {a\over\sqrt t}\mathcal P_4(x_0)
 +\sqrt t\,\mathcal P_5(x_0)
 +t\mathcal P_6(x_0)
 \right\}
 +\mathcal T_{a,t}(x_0),                                  \label{eq:newS7V23-local-parametrix}
\end{align}
where, for every \(N\),
\begin{equation}
 |\mathcal T_{a,t}(x_0)|
 \le C_N\{(t/R^2)^N+e^{-ct/a^2}\}.                       \label{eq:newS7V23-tail}
\end{equation}
\end{theorem}

\begin{proof}
Use exact continuum radial gauge and radial Higgs transport.  Paths leaving
the chart give \eqref{eq:newS7V23-tail} by the V11 localization and the
V16 all-moment bound.  Inside the chart use
\Cref{lem:newS7V23-local-series}.

Heat weights below two consist of the free rank term and the quadratic
constant-Higgs response.  The former cancels against
\(h_a^{\rm SM}(t;x_0;0,0)\); the latter is exactly
\(\mathcal O^{(2)}_{a,t}(H_0)\) by
\Cref{lem:newS7V23-quadratic-trace}.  At weight two, the complete V19
enumeration applies.  Its gauge square is V16, with the physical Weyl
factor recorded in \eqref{eq:newS7V23-OF}.  Its constant Higgs quartic is
V20, its two Higgs jets including the constant--second-jet contact are V21,
and its odd and curvature--two-Higgs classes are V22.  These are precisely
\(\mathcal O_F+\mathcal O_H\).

Replacing one retained lattice symbol by its continuum counterpart costs
\(Ca/\sqrt t\) and yields the first term on the right.  A retained frozen
jet paired with one next Taylor jet yields
\(\sqrt t\mathcal P_5\); two errors or a word of engineering dimension at
least six yield \(t\mathcal P_6\).  The ordering commutator in
\eqref{eq:newS7V23-error-budget} has the same relative factor
\(a/\sqrt t\).  Sum all higher complete words by
\eqref{eq:newS7V23-geometric}.
\end{proof}

\begin{corollary}[Pointwise smooth matter coefficient]
\label{cor:newS7V23-pointwise}
For every fixed smooth gauge--Higgs background and every joint limit
\[
 a\downarrow0,
 \qquad t\downarrow0,
 \qquad a^2/t\to0,
\]
the exactly mass-subtracted diagonal density converges to
\begin{equation}
 \mathcal O_F(F(x))
 +{1\over(4\pi)^2}
 \{2\mathfrak a_\bullet|D_\mu H(x)|^2
   +2\mathfrak b_\bullet|H(x)|^4\}.                       \label{eq:newS7V23-pointwise}
\end{equation}
\end{corollary}

\begin{proof}
Every term on the right of
\eqref{eq:newS7V23-local-parametrix} tends to zero.  Use a fixed chart
radius in \eqref{eq:newS7V23-tail}.
\end{proof}

\begin{theorem}[Logarithmic integrability of the joint smooth remainder]
\label{thm:newS7V23-integrated}
On the mesoscopic window
\(t\ge\varepsilon_q(a)=\ell^2(a/\ell)^q\), \(0<q<2/3\), the local
remainder after the exact weight-one and physical weight-four subtractions
is integrable against \(dt/t\).  Uniformly on the bounded smooth family,
\begin{align}
 \int_{\varepsilon_q(a)}^{t_0}
 {a\over\sqrt t}{dt\over t}&=O(a^{1-q/2}),                \label{eq:newS7V23-a-integral}\\
 \int_0^{t_0}\sqrt t\,{dt\over t}&=2\sqrt{t_0},
 &\int_0^{t_0}t\,{dt\over t}&=t_0.                       \label{eq:newS7V23-t-integrals}
\end{align}
Thus the cutoff error vanishes as \(a\downarrow0\), and the continuum
Taylor tail vanishes as the upper local proper-time scale
\(t_0\downarrow0\).
\end{theorem}

\begin{proof}
Integrate \eqref{eq:newS7V23-local-parametrix}.  The exponential tail is
handled by \(u=t/a^2\); choose the chart-moment power \(N\) large enough for
the polynomial tail.  The displayed elementary integrals give the remaining
terms.
\end{proof}

\begin{assessment}[What the exact mass subtraction accomplishes]
\label{ass:newS7V23-mass-subtraction}
Replacing \(\mathcal O^{(2)}_{a,t}\) by its continuum leading term before
taking the joint limit would leave derivatives of the lattice time
rescaling in V20.  The exact subtraction removes that lower-order leakage
before the weight-four coefficient is read.  It is a local regulator
counterterm; its continuum leading part is the ordinary quadratic Higgs
mass divergence.
\end{assessment}

\begin{assessment}[Scope of the smooth theorem]
\label{ass:newS7V23-scope}
The theorem is deterministic and applies to bounded smooth backgrounds in a
flat chart.  It does not estimate the complement of the smooth chart inside
the rough AOP field ensemble.  It also omits spin-connection curvature and
therefore does not yet transfer the \(R|H|^2\) owner.
\end{assessment}

\begin{programme}[V24 acquisition]
\label{prog:newS7V23-V24}
V24 will add a slowly varying orthonormal frame and spin connection to the
same local construction.  It will first derive the overlap spin-connection
vertex and its exact local Lorentz covariance, then isolate the mixed
scalar-curvature--two-Higgs coefficient.  If a direct lattice vierbein
symbol is absent from the archived construction, V24 will stop at a typed
metric transfer card rather than assume one.
\end{programme}

\begin{firewall}[Claim boundary after seventh-series V23]
\label{fire:newS7V23-boundary}
V23 defines the physical Weyl gauge sum and exact lattice Higgs mass
subtraction, proves a factorial-free joint local series, constructs the
smooth radial gauge--Higgs parametrix, transfers the full flat matter
weight-four coefficient, and proves logarithmic integrability of its
remainder.

V23 does not transfer metric or spin-connection coefficients, prove the
\(R|H|^2\) lattice owner, control rough AOP configurations, prove measure
concentration, construct the determinant-line phase, remove the interacting
cutoff, or construct the Standard Model--Einstein continuum.
\end{firewall}

\begin{theorem}[Seventh-series V23 result]
\label{thm:newS7V23-result}
After subtraction of the free density and the exact lattice quadratic Higgs
response, the smooth flat-chart physical Standard Model heat density has
local continuum coefficient
\[
 {1\over3(4\pi)^2}\mathfrak G_{\rm SM}(F)
 +{1\over(4\pi)^2}
 \{2\mathfrak a_\bullet|DH|^2+2\mathfrak b_\bullet|H|^4\}.
\]
Its deterministic remainder is
\(O(a/\sqrt t)+O(\sqrt t)+O(t)+O(e^{-ct/a^2})\) with the
corresponding bounded jet polynomials and is logarithmically integrable on
the established mesoscopic window.
\end{theorem}

\begin{proof}
Use \Cref{lem:newS7V23-quadratic-trace,
lem:newS7V23-local-series,
thm:newS7V23-local-parametrix,
cor:newS7V23-pointwise,
thm:newS7V23-integrated}.
\end{proof}

\paragraph{Parent-version record.}
V23 was formed from
\texttt{Renewal\_SM\_Einstein\_Continuum\_Research\_Notes\_V22.tex}.
The V22 parent had \(9\,410\) logical source lines, \(369\,218\) bytes, and
SHA--256
\[
 \texttt{6958752E3B5525C205187C20E29FCB67CAC316BB3C3509B4E2B10CF1F65ECB9A}.
\]
The next compulsory companion audit is V25.

% =============================================================================
% END APPENDED SEVENTH-SERIES V23 MATERIAL
% =============================================================================


% =============================================================================
% BEGIN APPENDED SEVENTH-SERIES V24 MATERIAL
% =============================================================================

\clearpage
\section{Seventh-series V24: the scalar-curvature--Higgs owner and the
missing lattice metric operator}
\label{sec:v7v24}

V23 closes the deterministic smooth matter coefficient on flat charts.  A
metric extension requires more than a continuum coframe: the overlap kernel,
its Hilbert-space adjoint, Wilson term, sign gap, volume weights, and stress
contacts must all be defined for the same discrete metric.  The attached
metric continuation supplies a controlled continuum coframe and spin
connection, but it does not supply that lattice operator.  V24 fixes the
physical continuum target, proves that flat data cannot determine it, and
states a sufficient metric overlap card with no hidden row.

\subsection{Acquisition record}
\label{sec:newS7V24-acquisition}

The current source inspected for the metric acquisition was
\begin{center}
\small
\begin{tabular}{ll}
file &
\texttt{Renewal\_Einstein\_Standard\_Model\_Physical\_Source\_Metric\_Duhamel\_Ancestry\_and\_Quantum\_Continuum\_Programme\_V35.tex},\\
bytes & \(1\,410\,598\),\\
logical source lines & \(36\,310\),\\
SHA--256 &
\texttt{AE1FBC0A8125C37465E38E7649E2992725F4579A86F0BB17483B9AE829CFB17C}.
\end{tabular}
\end{center}

\begin{audit}[What the metric source actually supplies]
\label{audit:newS7V24-source}
The V35 source contains an analytic scheduler map to a continuum metric and
coframe, proves that a summable \(C^{2,\alpha}\) scheduler tail transports
the Levi--Civita and torsion-free spin connections, their curvatures, and a
densitized Einstein writer, and gives an \(H^1\) incidence test for an
independently occurring connection.  These results provide regularity and
cofinal comparison of continuum geometric fields.

The source contains no formula for a metric-dependent lattice Wilson--Dirac
operator, no weighted lattice spinor inner product, no spin lift on lattice
edges, no metric-uniform Wilson sign gap, and no overlap heat vertex obtained
by differentiating such an operator.  Its translation-Ward firewall also
states that the breaking functional depends on the chosen lattice metric
coupling and may not be set to zero before that realization is constructed.
The Grand Tensor manuscript contains no replacement discretization.
\end{audit}

\subsection{The physical continuum target}
\label{sec:newS7V24-target}

\begin{theorem}[Scalar-curvature--Higgs multiplicity ledger]
\label{thm:newS7V24-RH-ledger}
For the V18 Standard Model Yukawa map, the coefficients of
\(R|H|^2\), before the common factor \((4\pi)^{-2}\) and any overall
effective-action sign, are
\begin{equation}
 \begin{array}{c|c}
 \text{object}&R|H|^2\text{ coefficient}\\ \hline
 \text{self-adjoint doubled }K_L\oplus K_R\text{ heat square}
    &{2\over3}\mathfrak a_\bullet\\
 \text{physical chiral-block heat trace}
    &{1\over3}\mathfrak a_\bullet\\
 \log|\det\mathcal B|={1\over2}\log\det(\mathcal B^*\mathcal B)
    &{1\over6}\mathfrak a_\bullet.
 \end{array}                                               \label{eq:newS7V24-RH-ledger}
\end{equation}
\end{theorem}

\begin{proof}
V17 gives \((a_Y/3)R|H|^2\) for the self-adjoint Dirac--Yukawa square,
and V18 gives \(a_Y=2\mathfrak a_\bullet\).  This proves the first row.
One physical massive chiral block is one half of the doubled internal heat
square, exactly as verified independently for the quartic and kinetic rows
in V20--V21.  The determinant modulus supplies the final factor one half.
\end{proof}

\begin{corollary}[Constant-mass curvature check]
\label{cor:newS7V24-mass-check}
For one physical Dirac block with constant mass \(m\), the heat coefficient
contains
\begin{equation}
 {1\over3(4\pi)^2}Rm^2.                                  \label{eq:newS7V24-mass-check}
\end{equation}
\end{corollary}

\begin{proof}
The massless spin-rank-four coefficient in the V17 convention is
\(A_2=-R/3\).  Multiplication by \(e^{-tm^2}\) contributes
\(-m^2A_2=Rm^2/3\) at weight four.  Restore \((4\pi)^{-2}\).
\end{proof}

\subsection{Why the flat overlap symbol is insufficient}
\label{sec:newS7V24-underdetermined}

\begin{proposition}[Flat data do not determine the mixed curvature owner]
\label{prop:newS7V24-underdetermined}
Knowing the complete flat metric heat operator, all of its gauge--Higgs
vertices, and its principal continuum symbol does not determine the
coefficient of \(R|H|^2\).  More precisely, if
\(L_g=\nabla_g^*\nabla_g+E_g\) is one covariant metric extension, then
\begin{equation}
 L_g^{(\lambda)}:=L_g+\lambda R_g I                       \label{eq:newS7V24-Llambda}
\end{equation}
has the same flat restriction and the same principal symbol for every
\(\lambda\in\mathbb R\), but its integrated \(A_4\) coefficient changes by
\begin{equation}
 \lambda\int_M R_g\,\tr(\Phi(H)^2)\,d\mu_g               \label{eq:newS7V24-shift}
\end{equation}
in the mixed curvature--Higgs channel.
\end{proposition}

\begin{proof}
When \(g\) is flat, \(R_g=0\), so all operators
\eqref{eq:newS7V24-Llambda} have the same flat restriction.  Adding a
zero-order endomorphism does not change the principal symbol.  In the
matter part of the integrated heat coefficient,
\[
 {1\over2}\tr(E^2)-{1\over6}R\tr(E),
\]
replace \(E\) by \(E+\lambda RI\).  The cross term between
\(\lambda RI\) and \(\Phi^2\) in \(E^2/2\) is exactly
\(\lambda R\tr\Phi^2\); the change in \(-RE/6\) contains only pure
curvature.  This proves \eqref{eq:newS7V24-shift}.
\end{proof}

\begin{corollary}[A metric discretization is a necessary datum]
\label{cor:newS7V24-necessary}
The coefficient in \eqref{eq:newS7V24-RH-ledger} cannot be transferred from
V23 by covariance or flat-symbol universality alone.  A declared lattice
metric Dirac operator, or an equivalent condition fixing its complete
subprincipal curvature endomorphism, is necessary.
\end{corollary}

\begin{proof}
The family \eqref{eq:newS7V24-Llambda} obeys all stated flat conditions but
has arbitrary mixed coefficient by
\Cref{prop:newS7V24-underdetermined}.
\end{proof}

\subsection{The lattice metric overlap datum}
\label{sec:newS7V24-card}

\begin{definition}[Lattice metric overlap card, LMOC]
\label{def:newS7V24-LMOC}
LMOC consists of the following typed rows on one uniformly nondegenerate
Euclidean coframe chart:
\begin{enumerate}[label=\textup{(\roman*)},leftmargin=3.5em]
\item positive site volume weights \(w_a(x;e)\) and a spinor inner product
      whose adjoint is used in every square and determinant modulus;
\item oriented coframes and \(\operatorname{Spin}(4)\) edge transports which transform
      exactly under sitewise frame rotations and converge to the
      torsion-free spin connection;
\item one explicit metric Wilson--Dirac operator
      \(D_{W,a}(e,\omega,A)\), including its Wilson term and all volume and
      coframe placements;
\item weighted gamma-five Hermiticity and a sign-kernel gap uniform on the
      declared smooth metric--gauge chart, so that the overlap sign and
      moving projectors exist;
\item exact gauge and local-frame covariance of the improved Yukawa block;
\item physical-branch consistency, through two metric derivatives and all
      mixed Higgs contacts, with the continuum Dirac principal and
      subprincipal symbols;
\item a Gaussian majorant and high-branch suppression for those complete
      metric heat divided differences;
\item convergence of the scalar-curvature--two-Higgs projection to the
      physical middle row of \eqref{eq:newS7V24-RH-ledger}; and
\item a stress definition obtained by differentiating the same weighted
      action, including measure, adjoint, Wilson, sign, projector, and
      Yukawa contacts.
\end{enumerate}
\end{definition}

\begin{lemma}[The order of LMOC is forced]
\label{lem:newS7V24-order}
Rows \textup{(i)}--\textup{(iv)} must be fixed before a metric overlap
operator or its heat square is defined.  Rows \textup{(v)}--\textup{(viii)}
then determine its local coefficient.  Row \textup{(ix)} must use the same
choices if the heat effective action is to source the Einstein equation.
\end{lemma}

\begin{proof}
The adjoint in \(D^*D\) depends on the inner product and site weights; the
sign depends on the Wilson kernel and its gap.  Hence the first four rows
are definitional prerequisites.  Heat vertices are derivatives of that
operator, so changing any of those choices changes the subprincipal and
contact terms by \Cref{prop:newS7V24-underdetermined}.  Finally, the stress
is the metric derivative of the action; differentiating a different metric
coupling would violate the common-action identity even if its flat value
agreed.
\end{proof}

\begin{theorem}[LMOC is sufficient for the smooth mixed transfer]
\label{thm:newS7V24-sufficiency}
Assume LMOC and bounded \(C^{2,\alpha}\) metric/coframe data, bounded
\(C^2\) gauge data, and bounded \(C^3\) Higgs data.  Then the V23 radial
parametrix extends to the scalar-curvature--two-Higgs channel and, after the
same exact lattice mass subtraction, has physical local owner
\begin{equation}
 {\mathfrak a_\bullet\over3(4\pi)^2}R_g|H|^2.             \label{eq:newS7V24-physical-RH}
\end{equation}
Its coefficient error is
\begin{equation}
 O(a/\sqrt t)+O(\sqrt t)+O(e^{-ct/a^2})                  \label{eq:newS7V24-conditional-error}
\end{equation}
times bounded covariant jet polynomials, and is logarithmically integrable
on the established mesoscopic window.
\end{theorem}

\begin{proof}
Rows (i)--(v) define one covariant overlap--Yukawa square.  In normal
coframe gauge at \(x_0\), the first nonflat metric jet is curvature; rows
(vi)--(vii) permit the V14 and V21 low/high-branch argument for the complete
mixed metric--Higgs divided difference.  Row (viii) identifies its limit as
\eqref{eq:newS7V24-physical-RH}.  The next metric and Higgs Taylor jets cost
\(O(\sqrt t)\), while symbol replacement costs \(O(a/\sqrt t)\); the high
branch is exponential.  Integrate as in V23.
\end{proof}

\begin{assessment}[Conditional status]
\label{ass:newS7V24-conditional}
Theorem \ref{thm:newS7V24-sufficiency} is a sufficiency theorem.  The
searched manuscripts do not establish LMOC rows (i)--(viii), so V24 does
not claim the lattice transfer \eqref{eq:newS7V24-physical-RH}.  The
continuum coefficient and the necessity of the missing datum are proved
unconditionally.
\end{assessment}

\subsection{Use of the existing continuum scheduler}
\label{sec:newS7V24-scheduler}

\begin{lemma}[Cofinal stability of the target owner]
\label{lem:newS7V24-scheduler}
Let \((g_n,e_n,H_n)\to(g,e,H)\) with
\begin{equation}
 \|g_n-g\|_{C^{2,\alpha}}
 +\|e_n-e\|_{C^{2,\alpha}}
 +\|H_n-H\|_{C^0}\le\rho_n,qquad \rho_n\to0,            \label{eq:newS7V24-scheduler-tail}
\end{equation}
inside a uniformly nondegenerate chart.  Then
\begin{equation}
 \|R_{g_n}|H_n|^2-R_g|H|^2\|_{C^0}
 \le C\rho_n.                                             \label{eq:newS7V24-RH-tail}
\end{equation}
If \(\sum_n\rho_n<\infty\), the corresponding adjacent owner differences
are summable.
\end{lemma}

\begin{proof}
Scalar curvature is locally Lipschitz from a uniformly nondegenerate
bounded \(C^{2,\alpha}\) metric set to \(C^{0,\alpha}\).  The squared Higgs
norm is locally Lipschitz in the metric/coframe chart and in \(H\).  Write
\[
 R_{g_n}|H_n|^2-R_g|H|^2
 =(R_{g_n}-R_g)|H_n|^2+R_g(|H_n|^2-|H|^2)
\]
and use the uniform bounds.  Apply the same estimate to adjacent pairs and
sum.
\end{proof}

\begin{corollary}[What V35 can supply after LMOC]
\label{cor:newS7V24-V35-use}
The V35 \(C^{2,\alpha}\) scheduler tail is sufficient to transport the
continuum target \eqref{eq:newS7V24-physical-RH} and its cofinal test
functionals once LMOC supplies the lattice-to-continuum coefficient.  It
does not supply LMOC itself.
\end{corollary}

\begin{proof}
Apply \Cref{lem:newS7V24-scheduler} to the metric/coframe convergence proved
in the inspected source.  The final sentence is the acquisition audit.
\end{proof}

\begin{programme}[V25 audit and upstream acquisition]
\label{prog:newS7V24-V25}
V25 is the compulsory companion audit.  It will then choose one concrete
Euclidean metric lattice datum for LMOC rows (i)--(iii), beginning with site
volume weights and a \(\operatorname{Spin}(4)\) edge lift.  The sign-gap row will be tested
before any metric heat coefficient is differentiated.  If uniform
admissibility fails, the construction will return to the Wilson kernel
rather than assume an overlap sign.
\end{programme}

\begin{firewall}[Claim boundary after seventh-series V24]
\label{fire:newS7V24-boundary}
V24 records the absence of a lattice metric overlap operator in the searched
sources, fixes the physical \(R|H|^2\) continuum coefficient and its three
multiplicities, proves flat data cannot determine that coefficient, defines
LMOC, proves its sufficiency, and proves cofinal stability of the target
under the existing continuum scheduler.

V24 does not construct LMOC, transfer the mixed curvature coefficient on
the lattice, define the common lattice stress, control rough metric or AOP
fields, prove measure concentration, construct the determinant-line phase,
remove the interacting cutoff, or construct the Standard Model--Einstein
continuum.
\end{firewall}

\begin{theorem}[Seventh-series V24 result]
\label{thm:newS7V24-result}
The physical chiral-block heat target for the mixed geometric Higgs owner is
\[
 {\mathfrak a_\bullet\over3(4\pi)^2}R_g|H|^2,
\]
and the determinant modulus carries one half of this value.  The target is
stable under the existing \(C^{2,\alpha}\) scheduler convergence.  Its
lattice transfer is underdetermined by all flat overlap data: a concrete
weighted metric Wilson--Dirac kernel, spin lift, gap, and complete metric
contacts are necessary.  LMOC lists those rows and is sufficient for the
smooth transfer when proved.
\end{theorem}

\begin{proof}
Use \Cref{thm:newS7V24-RH-ledger,
cor:newS7V24-mass-check,
prop:newS7V24-underdetermined,
cor:newS7V24-necessary,
lem:newS7V24-order,
thm:newS7V24-sufficiency,
lem:newS7V24-scheduler,
cor:newS7V24-V35-use}.
\end{proof}

\paragraph{Parent-version record.}
V24 was formed from
\texttt{Renewal\_SM\_Einstein\_Continuum\_Research\_Notes\_V23.tex}.
The V23 parent had \(9\,739\) logical source lines, \(383\,046\) bytes, and
SHA--256
\[
 \texttt{DB6F97A4D944E3D2DDCCEC6411A5986E181128FD0EA155CB879304D11B49126B}.
\]
The next compulsory companion audit is V25.

% =============================================================================
% END APPENDED SEVENTH-SERIES V24 MATERIAL
% =============================================================================


% =============================================================================
% BEGIN APPENDED SEVENTH-SERIES V25 MATERIAL
% =============================================================================

\clearpage
\section{Seventh-series V25: companion audit and a weighted metric
Wilson kernel with a proved local sign gap}
\label{sec:v7v25}

V24 showed that the scalar-curvature--Higgs coefficient cannot be recovered
from the flat overlap symbol.  The missing datum begins one level earlier
than a heat calculation: one must first put the spinors, transports, Clifford
hops, Wilson term, and adjoint in a single weighted finite-dimensional
Hilbert space.  This note performs the compulsory companion audit and then
constructs that algebraic datum.  The construction closes LMOC rows
\textup{(i)}--\textup{(iii)} on a Euclidean coordinate lattice and proves
row \textup{(iv)} on a declared perturbative normal chart.  No curvature heat
coefficient is inferred from these algebraic facts.

\subsection{Compulsory companion audit}
\label{sec:newS7V25-audit}

Because the companion tasks were replacing their active files during the
audit, each file was read as one coherent byte snapshot.  The inspected
snapshots were
\begin{center}
\small
\begin{tabular}{lll}
programme & source & snapshot data \\ \hline
Yang--Mills &
\texttt{Renewal\_Yang\_Mills\_Mass\_Gap\_Research\_Notes\_V98.tex}
& (1\,441\,088) bytes, (44\,404) newlines,\\
&&\texttt{D6E4A8A89A38E4BF7CC687FF7DADC58E982762B53D68B832922FD8E1B31BC945}
\\ [0.4em]
Navier--Stokes &
\texttt{Renewal\_NS\_Stability\_Research\_Notes\_V15.tex}
& (147\,056) bytes, (2\,878) newlines,\\
&&\texttt{D0433B54C1409B3EE28D870A43A33A43DFF76B7426C182B788290B53D82BEC5F}.
\end{tabular}
\end{center}

\begin{audit}[Decision from the V25 companion snapshots]
\label{audit:newS7V25-companions}
YM V98 proves the Bernoulli lift of differentiated interpolation gates,
complete optional-support collapse on the uniform interpolation line,
support localization at partition endpoints, and the nested-partition
representation of a BKAR interpolation matrix.  Its general ultrametric
layer-collapse card remains open.  NS V15 records that its type-I cross term
has been localized to a bounded configuration neighbourhood but that the
localized push still has no sign; its next step concerns the geometry of the
configuration at a homogeneous point.

Neither note contains a metric Wilson--Dirac kernel, a spin-lift gap, or a
curvature--Higgs heat coefficient, so no mathematical estimate or constant
is imported.  Their common structural lesson is applicable: complete the
underlying object and its sign property before expanding its derived
vertices.  Accordingly V25 proves the metric kernel and its sign gap before
the metric Duhamel expansion planned for V26.
\end{audit}

\subsection{Half-edge data and the weighted Hilbert space}
\label{sec:newS7V25-halfedge}

Let \(X_a\) be a finite four-dimensional periodic coordinate lattice; every
site has eight nearest neighbours.  The periodic choice is only an infrared
closure for the local construction.  A normal-coordinate ball of interest
will be separated from the closing collar.

\begin{definition}[Weighted half-edge Clifford datum]
\label{def:newS7V25-halfedge}
For every site \(x\in X_a\), let \(S_x\) be a finite-dimensional Hermitian
spin--internal fibre with a self-adjoint involution \(\Gamma_x\).  For every
unoriented nearest-neighbour edge \(e=\{x,y\}\), let \(S_e\) have chirality
\(\Gamma_e\), and choose:
\begin{enumerate}[label=\textup{(\roman*)},leftmargin=3.5em]
\item a positive site weight \(w_x\);
\item unitary half-edge maps \(P_{ex}:S_x\to S_e\) satisfying
      \(P_{ex}\Gamma_x=\Gamma_eP_{ex}\);
\item a self-adjoint edge Clifford map \(c_e\) satisfying
      \(c_e\Gamma_e=-\Gamma_ec_e\); and
\item an orientation sign \(\sigma_{xy}=-\sigma_{yx}\in\{ -1,1\}\).
\end{enumerate}
Set
\begin{equation}
 U_{xy}:=P_{ex}^*P_{ey},\qquad
 J_{xy}:=\sigma_{xy}P_{ex}^*c_eP_{ey}.                  \label{eq:newS7V25-UJ}
\end{equation}
The spinor Hilbert space is
\begin{equation}
 \mathscr H_{a,w}:=\bigoplus_{x\in X_a}S_x,qquad
 \langle\psi,\phi\rangle_w
 :=\sum_xw_x\langle\psi_x,\phi_x\rangle_{S_x}.        \label{eq:newS7V25-weighted-H}
\end{equation}
\end{definition}

\begin{lemma}[Exact reversal and chirality identities]
\label{lem:newS7V25-edge-identities}
The half-edge datum obeys
\begin{align}
 U_{yx}&=U_{xy}^*,&
 J_{yx}&=-J_{xy}^*,                                     \label{eq:newS7V25-reversal}\
 \Gamma_xU_{xy}&=U_{xy}\Gamma_y,&
 \Gamma_xJ_{xy}&=-J_{xy}\Gamma_y.                     \label{eq:newS7V25-edge-chiral}
\end{align}
\end{lemma}

\begin{proof}
Take adjoints in \eqref{eq:newS7V25-UJ}, use
\(c_e^*=c_e\) and \(\sigma_{yx}=-\sigma_{xy}\).  Move chirality through
the two half-edge maps; it commutes with their product and changes sign
once when it crosses \(c_e\).
\end{proof}

\subsection{The metric Wilson--Dirac operator}
\label{sec:newS7V25-Wilson}

For \(r>0\), define
\begin{align}
 (K_a\psi)_x
 &:= {1\over2a}\sum_{y\sim x}
       \sqrt{w_y\over w_x}\,J_{xy}\psi_y,              \label{eq:newS7V25-K}\
 (L_a\psi)_x
 &:= {1\over a^2}\sum_{y\sim x}
       \left(\psi_x-\sqrt{w_y\over w_x}\,U_{xy}\psi_y\right),
                                                               \label{eq:newS7V25-L}\
 W_a&:={ra\over2}L_a,qquad D_{W,a}:=K_a+W_a.            \label{eq:newS7V25-DW}
\end{align}
Let \(\Gamma=\bigoplus_x\Gamma_x\).

\begin{theorem}[Weighted adjoints and positivity]
\label{thm:newS7V25-adjoints}
On \(\mathscr H_{a,w}\),
\begin{equation}
 K_a^*=-K_a,\qquad L_a^*=L_a\ge0,qquad
 K_a\Gamma=-\Gamma K_a,qquad L_a\Gamma=\Gamma L_a,    \label{eq:newS7V25-adjoints}
\end{equation}
and hence
\begin{equation}
 D_{W,a}^*=\Gamma D_{W,a}\Gamma.                        \label{eq:newS7V25-g5H}
\end{equation}
More precisely,
\begin{equation}
 \langle\psi,L_a\psi\rangle_w
 ={1\over a^2}\sum_{e=\{x,y\}}
 \left\|\sqrt{w_x}P_{ex}\psi_x-sqrt{w_y}P_{ey}\psi_y\right\|_{S_e}^2.
                                                               \label{eq:newS7V25-L-positive}
\end{equation}
\end{theorem}

\begin{proof}
For a block operator \(A_{xy}:S_y\to S_x\), its weighted adjoint has
block \((A^{*w})_{yx}=(w_x/w_y)A_{xy}^*\).  Applied to the off-diagonal
block in \eqref{eq:newS7V25-K}, the weight ratio cancels the square roots;
\eqref{eq:newS7V25-reversal} gives \(K_a^{*w}=-K_a\).  Expanding the
right side of \eqref{eq:newS7V25-L-positive}, once per unoriented edge,
gives the diagonal and off-diagonal terms in
\(\langle\psi,L_a\psi\rangle_w\).  This proves self-adjointness and
positivity.  The two chirality relations follow sitewise from
\eqref{eq:newS7V25-edge-chiral}.  Therefore
\(D_{W,a}^*=-K_a+W_a=\Gamma(K_a+W_a)\Gamma\).
\end{proof}

\begin{lemma}[Exact sitewise covariance]
\label{lem:newS7V25-covariance}
Let \(G_x:S_x\to S_x\) be sitewise unitary frame--gauge rotations which
commute with chirality, and transform
\begin{equation}
 U_{xy}'=G_xU_{xy}G_y^*,\qquad
 J_{xy}'=G_xJ_{xy}G_y^*.                                \label{eq:newS7V25-transform}
\end{equation}
If \((\mathcal G\psi)_x=G_x\psi_x\), then
\begin{equation}
 K_a'=\mathcal GK_a\mathcal G^*,\quad
 L_a'=\mathcal GL_a\mathcal G^*,\quad
 D_{W,a}'=\mathcal GD_{W,a}\mathcal G^*.               \label{eq:newS7V25-covariance}
\end{equation}
\end{lemma}

\begin{proof}
Insert \eqref{eq:newS7V25-transform} in
\eqref{eq:newS7V25-K}--\eqref{eq:newS7V25-L}.  The site weights are frame
and gauge scalars, so every summand is conjugated by \(G_x\) at its output
site and by \(G_y^*\) at its input site.
\end{proof}

\subsection{Realization from a smooth coframe}
\label{sec:newS7V25-realization}

\begin{definition}[Midpoint realization]
\label{def:newS7V25-realization}
On a uniformly nondegenerate Euclidean coframe chart set
\begin{equation}
 w_x:=a^4\sqrt{\det g(x)}.                               \label{eq:newS7V25-volume}
\end{equation}
For the coordinate edge in direction \(\mu\), let \(m_e\) be its midpoint,
take \(P_{ex}\) to be parallel transport from \(x\) to \(m_e\) for the
torsion-free spin connection tensored with the gauge connection, and put
\begin{equation}
 c_e:=c_{g(m_e)}(dx^\mu)\otimes I_{\rm int}.             \label{eq:newS7V25-mid-Clifford}
\end{equation}
Choose \(\sigma_{x,x+a\widehat\mu}=1\).  On a closing collar the data may
be extended to the flat trivial datum; all local statements below are made
at positive distance from that collar.
\end{definition}

\begin{proposition}[The midpoint datum supplies LMOC rows
\textup{(i)}--\textup{(iii)}]
\label{prop:newS7V25-LMOC123}
The realization \eqref{eq:newS7V25-volume}--
\eqref{eq:newS7V25-mid-Clifford} has positive geometric volume weights,
exact spin--gauge covariance, and edge transports converging to the
torsion-free spin and gauge connections.  Equations
\eqref{eq:newS7V25-K}--\eqref{eq:newS7V25-DW} therefore define one explicit
metric Wilson--Dirac operator with a declared adjoint and Wilson placement.

In the convention in which the continuum geometric Dirac operator
\(\mathcal D_g=c_g(dx^\mu)\nabla_\mu\) is skew-adjoint, for every compactly
supported \(C^2\) spinor \(\psi\) in the interior chart,
\begin{equation}
 \|K_aI_a\psi-I_a\mathcal D_g\psi\|_{\mathscr H_{a,w}}
 \le Ca\|\psi\|_{C^2},qquad
 \|W_aI_a\psi\|_{\mathscr H_{a,w}}
 \le Ca\|\psi\|_{C^2}.                                 \label{eq:newS7V25-consistency}
\end{equation}
Here \(I_a\) is site sampling and \(C\) is uniform on bounded
\(C^2\) coframe--connection families with a fixed support.
\end{proposition}

\begin{proof}
Positivity of \(w_x\) follows from uniform nondegeneracy.  Parallel
transport is unitary, preserves chirality, and transforms at its endpoints;
midpoint Clifford multiplication is self-adjoint in the present convention,
anticommutes with chirality, and transforms in the spin representation.
Thus \Cref{def:newS7V25-halfedge,lem:newS7V25-covariance} apply.

For consistency, Taylor-expand the two oriented contributions of each
coordinate direction about the midpoint.  Their odd part is the covariant
central difference and their even part cancels.  The factors
\(\sqrt{w_y/w_x}\) give the volume half-density correction.  Equivalently,
the resulting first-order expression is the skew symmetrization of
\(c_g(dx^\mu)\nabla_\mu\); metric compatibility and the continuum
skew-adjoint identity identify that symmetrization with \(\mathcal D_g\).
The Taylor remainder is \(O(a)\) under the stated \(C^2\) bounds.  The
covariant graph Laplacian applied to a sampled smooth spinor is uniformly
\(O(1)\); multiplication by \(ra/2\) proves the second estimate.
Riemann-sum equivalence between \(\mathscr H_{a,w}\) and \(L^2(d\mu_g)\)
makes both estimates uniform.
\end{proof}

\begin{remark}[What consistency does not determine]
\label{rem:newS7V25-subprincipal}
The first-order consistency in \eqref{eq:newS7V25-consistency} does not by
itself determine a curvature heat coefficient.  V24's family
\(L_g^{(\lambda)}\) has the same principal consistency and different
\(R|H|^2\) owners.  The complete second metric derivative of this particular
operator must still be computed.
\end{remark}

\subsection{The free and perturbative sign gaps}
\label{sec:newS7V25-gap}

Fix \(0<\rho<2r\) and define
\begin{equation}
 X_a:=D_{W,a}-{\rho\over a}I,qquad H_a:=\Gamma X_a.     \label{eq:newS7V25-H}
\end{equation}

\begin{theorem}[Exact flat Wilson gap]
\label{thm:newS7V25-flat-gap}
For the flat trivial datum on the periodic four-lattice,
\begin{equation}
 \|H_{a,0}\psi\|\ge{\delta_{\rho,r}\over a}\|\psi\|,
                                                               \label{eq:newS7V25-flat-gap}
\end{equation}
where
\begin{equation}
 \delta_{\rho,r}^2
 :=\min_{q\in[-\pi,\pi]^4}
 \left\{\sum_{\mu=1}^4\sin^2q_\mu+
 \left(r\sum_{\mu=1}^4(1-\cos q_\mu)-\rho\right)^2\right\}>0.
                                                               \label{eq:newS7V25-delta}
\end{equation}
\end{theorem}

\begin{proof}
Fourier transformation gives
\begin{equation}
 K_{a,0}(q)={i\over a}\sum_\mu\gamma^\mu\sin q_\mu,qquad
 W_{a,0}(q)={r\over a}\sum_\mu(1-\cos q_\mu)I.          \label{eq:newS7V25-flat-symbol}
\end{equation}
The Clifford relations make \(H_{a,0}(q)^2\) equal \(a^{-2}\) times the
quantity in braces in \eqref{eq:newS7V25-delta}.  A zero would require
every \(q_\mu\in\{0,\pi\}\) and
\(2rk=\rho\) for an integer \(k\in\{0,\ldots,4\}\).  This is impossible
when \(0<\rho<2r\).  Compactness gives a strictly positive minimum, and
the finite periodic momentum set inherits the same lower bound.
\end{proof}

Let \(\mathcal Q_w:\mathscr H_{a,w}\to\bigoplus_xS_x\) be the unitary
half-density map \((\mathcal Q_w\psi)_x=\sqrt{w_x}\psi_x\).  In this
representation the square-root weight ratios in
\eqref{eq:newS7V25-K}--\eqref{eq:newS7V25-L} cancel exactly.  Relative to
the flat oriented direction \(\mu(x,y)\), put
\begin{equation}
 \eta_a:=\max_{x\sim y}
 \left\{\|J_{xy}-\sigma_{xy}\gamma^{\mu(x,y)}\|
       +r\|U_{xy}-I\|\right\}.                           \label{eq:newS7V25-eta}
\end{equation}

\begin{theorem}[Quantitative perturbative sign gap]
\label{thm:newS7V25-perturb-gap}
For the weighted half-edge datum,
\begin{equation}
 \|\mathcal Q_w(H_a-H_{a,0})\mathcal Q_w^{-1}\|
 \le {4\eta_a\over a}.                                  \label{eq:newS7V25-H-perturb}
\end{equation}
Consequently, if \(4\eta_a<\delta_{\rho,r}\), then \(H_a\) is
self-adjoint and
\begin{equation}
 \|H_a\psi\|_w
 \ge {\delta_{\rho,r}-4\eta_a\over a}\|\psi\|_w.     \label{eq:newS7V25-curved-gap}
\end{equation}
In particular its sign is well-defined.
\end{theorem}

\begin{proof}
Self-adjointness of \(H_a\) follows from
\eqref{eq:newS7V25-g5H}.  After conjugation by \(\mathcal Q_w\), each of
the eight hopping blocks differs from its flat value by at most
\((2a)^{-1}\) times the corresponding summand in
\eqref{eq:newS7V25-eta}; the Wilson diagonal is unchanged.  The block Schur
bound therefore gives \(8/(2a)=4/a\), proving
\eqref{eq:newS7V25-H-perturb}.  The reverse triangle inequality and
\Cref{thm:newS7V25-flat-gap} give \eqref{eq:newS7V25-curved-gap}.
\end{proof}

\begin{corollary}[A nonempty geometric admissibility chart]
\label{cor:newS7V25-geometric-chart}
In coframe normal gauge and gauge radial gauge on a ball of radius \(R\),
the midpoint realization satisfies, for bounded curvature and first jets,
\begin{equation}
 \eta_a\le C\{R^2\|{\rm Rm}\|_{C^0}
        +aR(\|{\rm Rm}\|_{C^1}+\|F\|_{C^0})+a^2\mathcal P_2\},
                                                               \label{eq:newS7V25-geometric-eta}
\end{equation}
where \(\mathcal P_2\) is a bounded polynomial in the declared second
jets.  Hence one may first choose \(R\) and then \(a\) so that
\(4\eta_a<\delta_{\rho,r}/2\), uniformly on a bounded smooth family.
\end{corollary}

\begin{proof}
Normal-frame expansion gives
\(c_g(dx^\mu)=\gamma^\mu+O(R^2\|{\rm Rm}\|_{C^0})\).
The torsion-free connection is \(O(R\|{\rm Rm}\|_{C^0})\), while radial
gauge gives \(A=O(R\|F\|_{C^0})\).  Integrating either connection over a
half-edge of length \(a/2\), and using the second-order remainder of the
path-ordered exponential, gives \eqref{eq:newS7V25-geometric-eta}.
Choose \(R\) to control the first term and then \(a\) to control the
remaining terms.  Apply \Cref{thm:newS7V25-perturb-gap}.
\end{proof}

\subsection{The resulting overlap operator}
\label{sec:newS7V25-overlap}

On the admissibility chart of \Cref{thm:newS7V25-perturb-gap}, define
\begin{equation}
 \epsilon_a:=\operatorname{sign}(H_a),\qquad
 D_{\mathrm{ov},a}:={\rho\over a}(I+\Gamma\epsilon_a).   \label{eq:newS7V25-overlap}
\end{equation}

\begin{proposition}[Exact overlap algebra]
\label{prop:newS7V25-overlap-algebra}
The operator \eqref{eq:newS7V25-overlap} is sitewise frame--gauge
covariant, satisfies weighted gamma-five Hermiticity, and obeys
\begin{equation}
 \Gamma D_{\mathrm{ov},a}+D_{\mathrm{ov},a}\Gamma
 ={a\over\rho}D_{\mathrm{ov},a}\Gamma D_{\mathrm{ov},a}. \label{eq:newS7V25-GW}
\end{equation}
\end{proposition}

\begin{proof}
Functional calculus preserves unitary covariance, and the gap makes the
sign continuous on the declared chart.  Since \(\epsilon_a^*=\epsilon_a\)
and \(\epsilon_a^2=I\), put \(V=\Gamma\epsilon_a\).  Then
\(V^*=\Gamma V\Gamma\), \(V^*V=I\), and
\(V\Gamma V=\Gamma\).  Substitution of
\(D_{\mathrm{ov},a}=(\rho/a)(I+V)\) gives both weighted gamma-five
Hermiticity and \eqref{eq:newS7V25-GW}.
\end{proof}

\begin{assessment}[Rows closed and rows still open]
\label{ass:newS7V25-status}
V25 unconditionally constructs the weighted half-edge Wilson datum and
proves its adjoint, positivity, chirality, and covariance identities.  The
midpoint realization supplies LMOC rows \textup{(i)}--\textup{(iii)} on a
smooth local chart.  The free calculation and norm-stability theorem prove
LMOC row \textup{(iv)} on the explicit perturbative chart
\(4\eta_a<\delta_{\rho,r}\).

This is not a metric heat-coefficient theorem.  In particular it does not
yet prove LMOC rows \textup{(vi)}--\textup{(viii)}, extend the gap to an
arbitrary global metric chart, or show that the midpoint realization has
the target \(R|H|^2\) owner.  Those questions require differentiating this
same operator, including its weights, half-edge maps, Clifford maps,
Wilson term, sign, and moving overlap projectors.
\end{assessment}

\begin{programme}[V26: complete metric first and second variations]
\label{prog:newS7V25-V26}
V26 will remain upstream of the heat coefficient.  On the proved gap chart
it will derive the first and second Fr\'echet derivatives of
\(\epsilon_a=\operatorname{sign}(H_a)\) from one resolvent integral, with
explicit gap bounds.  It will then list the derivatives of the volume
weights, half-edge transports, midpoint Clifford maps, and Wilson term and
assemble the complete metric vertex and contact.  Only after that contact
ledger is closed will the scalar-curvature--two-Higgs Duhamel projection be
taken.
\end{programme}

\begin{firewall}[Claim boundary after seventh-series V25]
\label{fire:newS7V25-boundary}
V25 performs the compulsory YM/NS audit, constructs a finite weighted
metric Wilson--Dirac operator from half-edge spin--gauge data, proves its
weighted adjoint and positivity identities and exact covariance, realizes
the data from a smooth coframe, proves the exact flat Wilson gap, proves a
quantitative perturbative curved-chart gap, and defines the resulting
covariant overlap operator with its exact Ginsparg--Wilson identity.

V25 does not compute or transfer the lattice \(R|H|^2\) coefficient, prove
the complete metric contact ledger, define the common lattice stress,
globalize the perturbative gap, control rough metric or AOP fields, prove
measure concentration, construct the determinant-line phase, remove the
interacting cutoff, or construct the Standard Model--Einstein continuum.
\end{firewall}

\begin{theorem}[Seventh-series V25 result]
\label{thm:newS7V25-result}
There is an explicit weighted midpoint Wilson--Dirac datum on every
uniformly nondegenerate smooth coordinate chart.  It is exactly
frame--gauge covariant and gamma-five Hermitian.  For
\(0<\rho<2r\), its flat sign kernel has gap
\(\delta_{\rho,r}/a\); on the geometric chart
\(4\eta_a<\delta_{\rho,r}\), the curved kernel has gap
\((\delta_{\rho,r}-4\eta_a)/a\).  Thus the overlap sign and moving
projectors exist on a nonempty, quantitatively declared metric--gauge
neighbourhood.  This closes the operator and local sign prerequisites for
the metric Duhamel calculation, but not its curvature coefficient.
\end{theorem}

\begin{proof}
Use \Cref{lem:newS7V25-edge-identities,
thm:newS7V25-adjoints,
lem:newS7V25-covariance,
prop:newS7V25-LMOC123,
thm:newS7V25-flat-gap,
thm:newS7V25-perturb-gap,
cor:newS7V25-geometric-chart,
prop:newS7V25-overlap-algebra}.
\end{proof}

\paragraph{Parent-version record.}
V25 was formed from
\texttt{Renewal\_SM\_Einstein\_Continuum\_Research\_Notes\_V24.tex}.
The V24 parent had (10\,081) logical source lines, (397\,734) bytes,
and SHA--256
\[
 \texttt{70117850268C2BCD4C17E720384D2E0252FADA71CC348FB2E5B833A84A688432}.
\]
The next compulsory companion audit is V30.

% =============================================================================
% END APPENDED SEVENTH-SERIES V25 MATERIAL
% =============================================================================


% =============================================================================
% BEGIN APPENDED SEVENTH-SERIES V26 MATERIAL
% =============================================================================

\clearpage
\section{Seventh-series V26: fixed-half-density metric contacts and the
complete second variation of the overlap sign}
\label{sec:v7v26}

V25 constructs the metric Wilson kernel on a family of weighted Hilbert
spaces.  A derivative between different Hilbert spaces is not defined until
an identification is chosen.  V26 uses the canonical half-density
identification.  It proves that every derivative of the explicit
square-root volume ratios cancels there, derives the first and second
variations of the half-edge transports and Clifford hops, and obtains the
complete first and second Fr\'echet derivatives of the overlap sign from a
single fixed contour.  This closes the finite-cutoff metric contact ledger
needed before a curvature projection can be attempted.

\subsection{One fixed Hilbert space}
\label{sec:newS7V26-fixed-space}

Let \(s=(s_1,\ldots,s_N)\) range in a smooth finite-dimensional metric--
coframe parameter chart contained in the V25 gap region.  Write
\(w_x(s),U_{xy}(s),J_{xy}(s)\) for the V25 midpoint data and
\(\mathscr H_{a,w(s)}\) for its weighted Hilbert space.  Fix the unweighted
space \(\mathscr H_a^0=\bigoplus_xS_x\), after a smooth local spin-frame
trivialization, and define
\begin{equation}
 \mathcal Q_s:\mathscr H_{a,w(s)}\longrightarrow\mathscr H_a^0,
 \qquad (\mathcal Q_s\psi)_x=\sqrt{w_x(s)}\psi_x.         \label{eq:newS7V26-Q}
\end{equation}

\begin{lemma}[Exact half-density cancellation]
\label{lem:newS7V26-half-density}
For the V25 operators,
\begin{align}
 (\widehat K_a(s)\psi)_x
 &:=(\mathcal Q_sK_a(s)\mathcal Q_s^{-1}\psi)_x
 ={1\over2a}\sum_{y\sim x}J_{xy}(s)\psi_y,             \label{eq:newS7V26-Khat}\
 (\widehat L_a(s)\psi)_x
 &:=(\mathcal Q_sL_a(s)\mathcal Q_s^{-1}\psi)_x
 ={1\over a^2}\sum_{y\sim x}\{\psi_x-U_{xy}(s)\psi_y\}.
                                                               \label{eq:newS7V26-Lhat}
\end{align}
Thus
\begin{equation}
 \widehat D_{W,a}=\widehat K_a+{ra\over2}\widehat L_a  \label{eq:newS7V26-Dhat}
\end{equation}
contains no explicit site-weight ratio.  This is an identity at every
lattice spacing, not an asymptotic cancellation.
\end{lemma}

\begin{proof}
At an off-diagonal block \(x\leftarrow y\), conjugation multiplies the
V25 factor \(\sqrt{w_y/w_x}\) by
\(\sqrt{w_x}\) on the left and \(w_y^{-1/2}\) on the right.  The product is
one.  The diagonal of \(L_a\) is unchanged.
\end{proof}

For comparison with a calculation in the moving coordinate fibres, put
\begin{equation}
 q_h:=\mathcal Q_s^{-1}\partial_h\mathcal Q_s
 ={1\over2}\partial_h\log w,
                                                               \label{eq:newS7V26-qh}
\end{equation}
where the last expression is diagonal in the sites.

\begin{proposition}[Volume contacts are identification contacts]
\label{prop:newS7V26-volume-cancel}
If \(T(s)\) is an operator on \(\mathscr H_{a,w(s)}\) and
\(\widehat T(s)=\mathcal Q_sT(s)\mathcal Q_s^{-1}\), then
\begin{equation}
 \partial_h\widehat T
 =\mathcal Q_s\{\partial_hT+[q_h,T]\}\mathcal Q_s^{-1}. \label{eq:newS7V26-moving-derivative}
\end{equation}
For \(T=K_a,L_a,D_{W,a}\), the derivatives of
\(\sqrt{w_y/w_x}\) in \(\partial_hT\) cancel the commutator in
\eqref{eq:newS7V26-moving-derivative} edge by edge, leaving precisely the
derivatives of \eqref{eq:newS7V26-Khat}--\eqref{eq:newS7V26-Lhat}.
Consequently cyclic spectral traces of the conjugated kernel contain no
independent ``volume-ratio vertex.''
\end{proposition}

\begin{proof}
Differentiate \(\mathcal Q_sT\mathcal Q_s^{-1}\) and use
\(\partial_h\mathcal Q_s=\mathcal Q_sq_h\) and
\(\partial_h\mathcal Q_s^{-1}=-q_h\mathcal Q_s^{-1}\).  For an
\(x\leftarrow y\) block,
\begin{equation}
 \partial_h\log\sqrt{w_y/w_x}=q_h(y)-q_h(x),             \label{eq:newS7V26-ratio-derivative}
\end{equation}
whereas \([q_h,T]_{xy}=(q_h(x)-q_h(y))T_{xy}\).  Their sum is zero.
The trace statement follows because conjugation preserves finite-dimensional
functional calculus and its cyclic trace.
\end{proof}

\begin{remark}[The geometric volume is not discarded]
\label{rem:newS7V26-volume-status}
The weights determine the Hilbert-space adjoint and the half-density
identification, and their variation is exactly what produces the
commutator cancellation above.  Bosonic volume factors and a chosen
fermionic integration measure still have their own metric variations.  The
proposition concerns the spectral trace of the declared overlap kernel; it
does not set those other variations to zero.
\end{remark}

For reference, if \(g(s,t)=g+sh+tk\), then at a site
\begin{align}
 {\partial_hw\over w}&={1\over2}\tr(g^{-1}h),            \label{eq:newS7V26-w1}\
 {\partial_{hk}^2w\over w}
 &= {1\over4}\tr(g^{-1}h)\tr(g^{-1}k)
   -{1\over2}\tr(g^{-1}hg^{-1}k).                      \label{eq:newS7V26-w2}
\end{align}
These identities verify directly every first and mixed second ratio contact
which cancels in \Cref{prop:newS7V26-volume-cancel}.

\subsection{Complete half-edge variations}
\label{sec:newS7V26-halfedge-variation}

Let \(P_{ex}(s)\) solve parallel transport along the oriented half-edge
\(\gamma_{ex}:[0,1]\to M\):
\begin{equation}
 {dP(\tau,0;s)\over d\tau}=-\mathcal A(\tau;s)P(\tau,0;s),
 \qquad P(0,0;s)=I,                                    \label{eq:newS7V26-transport-ode}
\end{equation}
where \(\mathcal A\) is the pulled-back torsion-free spin connection
tensored with the gauge connection.  Write
\(P(\tau_2,\tau_1;s)\) for the corresponding partial transport.

\begin{lemma}[First and second transport contacts]
\label{lem:newS7V26-P-contacts}
For parameter directions \(h,k\),
\begin{align}
 \partial_hP(1,0)
 &=-\int_0^1P(1,\tau)(\partial_h\mathcal A)(\tau)
                  P(\tau,0)\,d\tau,                    \label{eq:newS7V26-P1}\
 \partial_{hk}^2P(1,0)
 &=-\int_0^1P(1,\tau)(\partial_{hk}^2\mathcal A)(\tau)
                  P(\tau,0)\,d\tau                    \notag\\
 &\quad+\int_{0<\tau_2<\tau_1<1}
 P(1,\tau_1)(\partial_h\mathcal A)(\tau_1)
 P(\tau_1,\tau_2)(\partial_k\mathcal A)(\tau_2)
 P(\tau_2,0)\,d\tau_2d\tau_1                         \notag\\
 &\quad+\int_{0<\tau_2<\tau_1<1}
 P(1,\tau_1)(\partial_k\mathcal A)(\tau_1)
 P(\tau_1,\tau_2)(\partial_h\mathcal A)(\tau_2)
 P(\tau_2,0)\,d\tau_2d\tau_1.                        \label{eq:newS7V26-P2}
\end{align}
\end{lemma}

\begin{proof}
Differentiate the Volterra equation associated with
\eqref{eq:newS7V26-transport-ode}.  One differentiation inserts
\(-\partial_h\mathcal A\) at every possible parameter time.  A second
differentiation either inserts \(-\partial_{hk}^2\mathcal A\) once or
inserts the two first variations in either order.  Time ordering gives the
two simplices in \eqref{eq:newS7V26-P2}.
\end{proof}

Let a dot with subscript denote a parameter variation.  Since
\(U_{xy}=P_{ex}^*P_{ey}\) and
\(J_{xy}=\sigma_{xy}P_{ex}^*c_eP_{ey}\), their complete contacts are
\begin{align}
 (U_{xy})_h
 &=(P_{ex,h})^*P_{ey}+P_{ex}^*P_{ey,h},                 \label{eq:newS7V26-U1}\
 (U_{xy})_{hk}
 &=(P_{ex,hk})^*P_{ey}+(P_{ex,h})^*P_{ey,k}
  +(P_{ex,k})^*P_{ey,h}+P_{ex}^*P_{ey,hk},              \label{eq:newS7V26-U2}\
 (J_{xy})_h
 &=\sigma_{xy}\{(P_{ex,h})^*c_eP_{ey}
                 +P_{ex}^*(c_e)_hP_{ey}
                 +P_{ex}^*c_eP_{ey,h}\},               \label{eq:newS7V26-J1}\
 (J_{xy})_{hk}
 &=\sigma_{xy}\bigl\{(P_{ex,hk})^*c_eP_{ey}
 +(P_{ex,h})^*(c_e)_kP_{ey}
 +(P_{ex,h})^*c_eP_{ey,k}                               \notag\\
 &\qquad +(P_{ex,k})^*(c_e)_hP_{ey}
 +(P_{ex})^*(c_e)_{hk}P_{ey}
 +P_{ex}^*(c_e)_hP_{ey,k}                               \notag\\
 &\qquad +(P_{ex,k})^*c_eP_{ey,h}
 +P_{ex}^*(c_e)_kP_{ey,h}
 +P_{ex}^*c_eP_{ey,hk}\bigr\}.                         \label{eq:newS7V26-J2}
\end{align}
If \(E_a{}^\mu(s)\) is the inverse coframe at the midpoint, then
\begin{equation}
 c_e(s)=\gamma^aE_a{}^\mu(s),\quad
 (c_e)_h=\gamma^a(E_a{}^\mu)_h,quad
 (c_e)_{hk}=\gamma^a(E_a{}^\mu)_{hk}.                  \label{eq:newS7V26-c-contacts}
\end{equation}

\begin{theorem}[Complete Wilson metric vertex and contact]
\label{thm:newS7V26-Wilson-contacts}
On the fixed half-density space,
\begin{align}
 (\widehat D_{W,a})_h\psi_x
 &= {1\over2a}\sum_{y\sim x}(J_{xy})_h\psi_y
   -{r\over2a}\sum_{y\sim x}(U_{xy})_h\psi_y,          \label{eq:newS7V26-DW1}\
 (\widehat D_{W,a})_{hk}\psi_x
 &= {1\over2a}\sum_{y\sim x}(J_{xy})_{hk}\psi_y
   -{r\over2a}\sum_{y\sim x}(U_{xy})_{hk}\psi_y.      \label{eq:newS7V26-DW2}
\end{align}
Equations \eqref{eq:newS7V26-P1}--\eqref{eq:newS7V26-c-contacts}
contain every transport, connection, Clifford, and Wilson contribution.
They preserve the differentiated reversal, chirality, and covariance
identities inherited from V25.
\end{theorem}

\begin{proof}
Differentiate \eqref{eq:newS7V26-Khat}--
\eqref{eq:newS7V26-Dhat}; the Wilson diagonal is parameter-independent.
The product rules \eqref{eq:newS7V26-U1}--\eqref{eq:newS7V26-J2} and the
transport identities supply all terms.  Differentiating the exact V25
algebraic identities proves the final assertion.
\end{proof}

\subsection{A fixed contour for the sign}
\label{sec:newS7V26-sign}

Set
\begin{equation}
 \widehat H_a:=\Gamma\left(\widehat D_{W,a}-{\rho\over a}I\right),
 \qquad \widetilde H_a:=a\widehat H_a.                  \label{eq:newS7V26-Htilde}
\end{equation}
On a closed subchart of the V25 region suppose
\begin{equation}
 \sigma(\widetilde H_a)
 \subset[-\Lambda,-g]\cup[g,\Lambda],qquad g>0,       \label{eq:newS7V26-uniform-spectrum}
\end{equation}
uniformly in \(a\) and in the metric parameters.  V25 permits
\(g=\delta_{\rho,r}-4\sup\eta_a\).

Choose disjoint positively oriented contours \(\mathcal C_+\) and
\(\mathcal C_-\) enclosing respectively the positive and negative spectral
intervals in \eqref{eq:newS7V26-uniform-spectrum}.  Let
\(\mathcal C=\mathcal C_+\cup\mathcal C_-\), let
\(d=\operatorname{dist}(\mathcal C,
[-\Lambda,-g]\cup[g,\Lambda])>0\), let \(L_{\mathcal C}\) be its total
length, and set \(f=1\) inside \(\mathcal C_+\), \(f=-1\) inside
\(\mathcal C_-\).

\begin{theorem}[Complete first and second sign variations]
\label{thm:newS7V26-sign-contacts}
With \(R_z=(z-\widetilde H_a)^{-1}\),
\begin{align}
 \epsilon_a:=\operatorname{sign}(\widehat H_a)
 &= {1\over2\pi i}\int_{\mathcal C}f(z)R_z\,dz,         \label{eq:newS7V26-sign-contour}\
 (\epsilon_a)_h
 &= {1\over2\pi i}\int_{\mathcal C}
 f(z)R_z(\widetilde H_a)_hR_z\,dz,                      \label{eq:newS7V26-sign1}\
 (\epsilon_a)_{hk}
 &= {1\over2\pi i}\int_{\mathcal C}f(z)
 \bigl\{R_z(\widetilde H_a)_{hk}R_z
 +R_z(\widetilde H_a)_hR_z(\widetilde H_a)_kR_z       \notag\\
 &\hspace{42mm}
 +R_z(\widetilde H_a)_kR_z(\widetilde H_a)_hR_z\bigr\}\,dz.
                                                               \label{eq:newS7V26-sign2}
\end{align}
Moreover,
\begin{align}
 \|(\epsilon_a)_h\|
 &\le {L_{\mathcal C}\over2\pi d^2}
       \|(\widetilde H_a)_h\|,                         \label{eq:newS7V26-sign1-bound}\
 \|(\epsilon_a)_{hk}\|
 &\le {L_{\mathcal C}\over2\pi}
 \left\{d^{-2}\|(\widetilde H_a)_{hk}\|
 +2d^{-3}\|(\widetilde H_a)_h\|
          \|(\widetilde H_a)_k\|\right\}.            \label{eq:newS7V26-sign2-bound}
\end{align}
\end{theorem}

\begin{proof}
The locally constant function \(f\) is holomorphic on the disconnected
neighbourhood bounded by \(\mathcal C\), so the finite-dimensional
holomorphic functional calculus gives \eqref{eq:newS7V26-sign-contour}.
The resolvent identity gives
\begin{equation}
 (R_z)_h=R_z(\widetilde H_a)_hR_z.                       \label{eq:newS7V26-resolvent1}
\end{equation}
Differentiate once and twice under the fixed finite contour to obtain
\eqref{eq:newS7V26-sign1}--\eqref{eq:newS7V26-sign2}.
On \(\mathcal C\), self-adjointness gives \(\|R_z\|\le d^{-1}\); taking
norms under the integral proves the two bounds.
\end{proof}

\begin{corollary}[Uniform finite-cutoff metric bounds]
\label{cor:newS7V26-uniform}
Suppose the coframe, connection, and their first two parameter derivatives
are uniformly bounded on the V25 chart.  Then
\begin{equation}
 \|(\widetilde H_a)_h\|\le C\|h\|_{C^1},\qquad
 \|(\widetilde H_a)_{hk}\|
 \le C\{\|h\|_{C^1}\|k\|_{C^1}+\|(h,k)\|_{C^1,2}\}, \label{eq:newS7V26-H-bounds}
\end{equation}
with no negative power of \(a\).  Hence the right sides of
\eqref{eq:newS7V26-sign1-bound}--\eqref{eq:newS7V26-sign2-bound} are uniform
in the cutoff.
\end{corollary}

\begin{proof}
Each row of \eqref{eq:newS7V26-DW1}--\eqref{eq:newS7V26-DW2} has eight
blocks of size \(O(a^{-1})\), while multiplication by \(a\) in
\(\widetilde H_a\) cancels that scale.  The transport formulas and unitarity
bound their first and second variations by the corresponding connection
and coframe variations.  Apply the block Schur bound and then
\Cref{thm:newS7V26-sign-contacts}.
\end{proof}

\subsection{Overlap, Yukawa, and heat-square contacts}
\label{sec:newS7V26-full-contact}

On the fixed space define
\begin{equation}
 \widehat D_{\mathrm{ov},a}={\rho\over a}(I+\Gamma\epsilon_a),
 \qquad
 \widehat{\mathcal B}_a
 =\widehat D_{\mathrm{ov},a}
  +M\left(I-{a\over2\rho}\widehat D_{\mathrm{ov},a}\right),
                                                               \label{eq:newS7V26-B}
\end{equation}
where \(M=\Phi(H)\) is held fixed under a pure metric variation.  Put
\(\widehat{\mathcal L}_a=widehat{\mathcal B}_a^*widehat{\mathcal B}_a\).

\begin{theorem}[Complete pure-metric vertex and contact ledger]
\label{thm:newS7V26-complete-ledger}
For \(j=h,hk\),
\begin{equation}
 (\widehat D_{\mathrm{ov},a})_j
 ={\rho\over a}\Gamma(\epsilon_a)_j,qquad
 (\widehat{\mathcal B}_a)_j
 =\left(I-{a\over2\rho}M\right)
   (\widehat D_{\mathrm{ov},a})_j.                      \label{eq:newS7V26-B-contacts}
\end{equation}
The heat-square contacts are
\begin{align}
 (\widehat{\mathcal L}_a)_h
 &=(\widehat{\mathcal B}_a)_h^*\widehat{\mathcal B}_a
   +\widehat{\mathcal B}_a^*(\widehat{\mathcal B}_a)_h,
                                                               \label{eq:newS7V26-L1}\
 (\widehat{\mathcal L}_a)_{hk}
 &=(\widehat{\mathcal B}_a)_{hk}^*\widehat{\mathcal B}_a
   +\widehat{\mathcal B}_a^*(\widehat{\mathcal B}_a)_{hk}
   +(\widehat{\mathcal B}_a)_h^*(\widehat{\mathcal B}_a)_k
   +(\widehat{\mathcal B}_a)_k^*(\widehat{\mathcal B}_a)_h. \label{eq:newS7V26-L2}
\end{align}
Together with \eqref{eq:newS7V26-P1}--\eqref{eq:newS7V26-sign2}, these
form the complete finite-cutoff pure-metric vertex and contact of the V25
operator.  No sign, projector, Wilson, transport, Clifford, adjoint, or
site-weight identification term is omitted.
\end{theorem}

\begin{proof}
Differentiate \eqref{eq:newS7V26-B}.  Since \(M\) is fixed and stands to
the left of \(\widehat D_{\mathrm{ov},a}\), both derivatives factor as
shown.  Apply the product rule once and twice to
\(\widehat{\mathcal B}_a^*\widehat{\mathcal B}_a\).  Completeness follows
from the exact half-density cancellation, the exhaustive product rules for
\(U,J\), and the holomorphic functional-calculus derivative of the sign.
\end{proof}

\begin{corollary}[Gauge--frame Ward covariance of the contacts]
\label{cor:newS7V26-contact-covariance}
If the metric parameter directions are transformed together with the
background, every operator in
\eqref{eq:newS7V26-DW1}--\eqref{eq:newS7V26-L2} transforms by sitewise
conjugation.  Therefore the corresponding differentiated heat traces are
gauge and local-frame invariant at finite cutoff.
\end{corollary}

\begin{proof}
Differentiate the exact covariance identity of V25.  Holomorphic functional
calculus commutes with conjugation, adjoints commute with unitary
conjugation, and the finite-dimensional trace is cyclic.
\end{proof}

\begin{assessment}[What V26 closes]
\label{ass:newS7V26-status}
V26 closes the algebraic part of LMOC row \textup{(vi)} through two metric
derivatives on the V25 gap chart: the vertex and contact are now explicitly
defined and uniformly bounded at finite cutoff after rescaling the sign
kernel.  It also removes a possible ambiguity in LMOC rows \textup{(i)} and
\textup{(ix)}: site-volume ratios cancel from spectral vertices under the
canonical half-density identification, while bosonic volumes and the
chosen fermionic integration measure remain separate stress owners.

V26 does not prove the low-branch continuum symbol of the complete contact,
its Gaussian majorant, high-branch suppression, or its
scalar-curvature--two-Higgs projection.  Those analytic rows remain
necessary before the coefficient in V24 can be transferred.
\end{assessment}

\begin{programme}[V27: curvature projection of the complete contact]
\label{prog:newS7V26-V27}
V27 will use normal coframe gauge, one compactly supported metric plane-wave
direction, and a constant Higgs background.  It will expand the first metric
heat response \eqref{eq:newS7V26-L1} to quadratic Higgs order and project its
order-two external-momentum part onto the linearized scalar-curvature
combination \(q_\mu q_\nu h^{\mu\nu}-q^2\operatorname{tr}h\).  Both Yukawa
insertions and every induced overlap contact will be retained.  The second
metric contact proved here is reserved for stress differentiability and
later pure-gravity owners.  The low physical branch will be compared with
the continuum constant-mass check in V24; the remaining task will be an
integrable majorant uniform in the lattice momentum and in the V25 gap
chart.
\end{programme}

\begin{firewall}[Claim boundary after seventh-series V26]
\label{fire:newS7V26-boundary}
V26 fixes the varying-Hilbert-space problem by the canonical half-density
map, proves exact cancellation of site-volume-ratio contacts in the spectral
kernel, derives all first and mixed second half-edge, Clifford, Wilson, sign,
overlap, Yukawa, and heat-square contacts, proves explicit contour bounds
uniform in the cutoff on a closed gap subchart, and proves finite-cutoff
contact covariance.

V26 does not compute the lattice \(R|H|^2\) projection, prove its continuum
limit or majorant, define the complete bosonic and fermionic-measure stress,
globalize the gap, control rough metric or AOP fields, prove measure
concentration, construct the determinant-line phase, remove the interacting
cutoff, or construct the Standard Model--Einstein continuum.
\end{firewall}

\begin{theorem}[Seventh-series V26 result]
\label{thm:newS7V26-result}
On every closed subchart of the V25 metric gap region, the V25 overlap--
Yukawa heat square has a complete, gauge--frame covariant first and mixed
second metric contact on one fixed half-density Hilbert space.  The sign
contacts are given by \eqref{eq:newS7V26-sign1}--
\eqref{eq:newS7V26-sign2} and obey cutoff-uniform bounds after the natural
rescaling \(\widetilde H_a=a\widehat H_a\).  Explicit site-volume-ratio
vertices cancel exactly under this identification.  Thus the next
curvature projection can be performed without an undefined adjoint,
moving-space derivative, or omitted sign contact.
\end{theorem}

\begin{proof}
Use \Cref{lem:newS7V26-half-density,
prop:newS7V26-volume-cancel,
lem:newS7V26-P-contacts,
thm:newS7V26-Wilson-contacts,
thm:newS7V26-sign-contacts,
cor:newS7V26-uniform,
thm:newS7V26-complete-ledger,
cor:newS7V26-contact-covariance}.
\end{proof}

\paragraph{Parent-version record.}
V26 was formed from
\texttt{Renewal\_SM\_Einstein\_Continuum\_Research\_Notes\_V25.tex}.
The V25 parent had (10\,555) logical source lines, (418\,111) bytes,
and SHA--256
\[
 \texttt{EB5046D313B70739E7A3E77ACA92D23861FB81EF3757F97012B4F67D5BE7D1C3}.
\]
The next compulsory companion audit is V30.

% =============================================================================
% END APPENDED SEVENTH-SERIES V26 MATERIAL
% =============================================================================


% =============================================================================
% BEGIN APPENDED SEVENTH-SERIES V27 MATERIAL
% =============================================================================

\clearpage
\section{Seventh-series V27: exact massive metric factorization and reduction
of the mixed owner to the massless (A_2) card}
\label{sec:v7v27}

The direct curvature--two-Higgs expansion announced in V26 is algebraically
unnecessary.  On a constant Higgs background, each Yukawa singular mode is
a constant overlap mass, and the Ginsparg--Wilson relation factorizes its
square exactly for every metric in the V25 gap chart.  V27 proves that
factorization, separates the power-divergent time-rescaling response from
the logarithmic mixed owner, and reduces the desired (R|H|^2) transfer to
one massless scalar-curvature (A_2) statement.  This is an upstream
reduction: it does not assume that remaining massless statement.

\subsection{Exact constant-mass algebra on the metric chart}
\label{sec:newS7V27-mass-algebra}

Let (D_a=widehat D_{\mathrm{ov},a}(g,A)) be the V25--V26 overlap
operator on the fixed half-density Hilbert space.  It satisfies
\begin{equation}
 D_a^*=\Gamma D_a\Gamma,qquad
 \Gamma D_a+D_a\Gamma={a\over\rho}D_a\Gamma D_a.        \label{eq:newS7V27-GW}
\end{equation}

\begin{lemma}[Scalar form of the overlap identity]
\label{lem:newS7V27-scalar-GW}
Equation \eqref{eq:newS7V27-GW} implies
\begin{equation}
 D_a+D_a^*={a\over\rho}D_a^*D_a.                        \label{eq:newS7V27-Dplus}
\end{equation}
\end{lemma}

\begin{proof}
Multiply the Ginsparg--Wilson identity on the left by \(\Gamma\), use
\(\Gamma^2=I\), and replace \(\Gamma D_a\Gamma\) by \(D_a^*\).
\end{proof}

For a real constant mass \(m\), define the improved massive block
\begin{equation}
 D_{a,m}:=D_a+m\left(I-{a\over2\rho}D_a\right)
 =\left(1-{am\over2\rho}\right)D_a+mI.                 \label{eq:newS7V27-Dm}
\end{equation}

\begin{theorem}[Exact massive overlap square on every metric background]
\label{thm:newS7V27-mass-square}
For every metric--gauge datum in the V25 gap chart,
\begin{equation}
 D_{a,m}^*D_{a,m}
 =c_a(m)D_a^*D_a+m^2I,qquad
 c_a(m):=1-{a^2m^2\over4\rho^2}.                       \label{eq:newS7V27-mass-square}
\end{equation}
The identity uses no translation invariance and remains true after any
metric variation for which the overlap sign exists.
\end{theorem}

\begin{proof}
Put \(\alpha=1-am/(2\rho)\).  Expanding the left side of
\eqref{eq:newS7V27-mass-square} gives
\[
 \alpha^2D_a^*D_a+\alpha m(D_a+D_a^*)+m^2I.
\]
Use \eqref{eq:newS7V27-Dplus}.  The coefficient of \(D_a^*D_a\) is
\[
 \alpha^2+{am\over\rho}\alpha
 =\left(1-{am\over2\rho}\right)
  \left(1+{am\over2\rho}\right)=c_a(m).
\]
No step refers to a flat symbol.
\end{proof}

\begin{corollary}[Exact heat-kernel factorization]
\label{cor:newS7V27-heat-factor}
If \(|am|<2\rho\), then for every \(t>0\),
\begin{equation}
 e^{-tD_{a,m}^*D_{a,m}}
 =e^{-tm^2}e^{-tc_a(m)D_a^*D_a}.                        \label{eq:newS7V27-heat-factor}
\end{equation}
The identity holds for the full trace, localized diagonal trace, and every
metric derivative of either side.
\end{corollary}

\begin{proof}
The two operators on the right of \eqref{eq:newS7V27-mass-square} commute,
one being scalar.  Functional calculus proves the factorization.  Equality
of smooth finite-dimensional operator families may be differentiated and
then traced or localized.
\end{proof}

\subsection{The exact quadratic-mass response}
\label{sec:newS7V27-mass-response}

Let \(u=m^2\) and, for any fixed localization operator \(\chi\), put
\begin{equation}
 \Theta_a(t,u;g):=Tr\!\left(\chi
 e^{-tD_{a,\sqrt u}^*D_{a,\sqrt u}}\right),qquad
 \Theta_{a,0}(t;g):=\Tr\!\left(\chi e^{-tD_a^*D_a}\right). \label{eq:newS7V27-Theta}
\end{equation}
Although \(D_{a,\sqrt u}\) contains \(\sqrt u\), its square and heat trace
are smooth in \(u\) by \eqref{eq:newS7V27-mass-square}.

\begin{proposition}[Exact derivative at quadratic mass order]
\label{prop:newS7V27-u-response}
At \(u=0\),
\begin{equation}
 \partial_u\Theta_a(t,0;g)
 =-t\Theta_{a,0}(t;g)
  -{a^2t\over4\rho^2}\partial_t\Theta_{a,0}(t;g).      \label{eq:newS7V27-u-response}
\end{equation}
The same identity holds after one or two metric derivatives.
\end{proposition}

\begin{proof}
By \Cref{cor:newS7V27-heat-factor},
\[
 \Theta_a(t,u;g)=e^{-tu}
 \Theta_{a,0}\!\left(t\left(1-{a^2u\over4\rho^2}\right);g\right).
\]
Differentiate at \(u=0\).  Metric differentiation commutes with
\(\partial_u\), \(\partial_t\), localization, and the finite trace.
\end{proof}

\begin{lemma}[The time-rescaling term carries no continuum logarithmic
\(R m^2\) owner]
\label{lem:newS7V27-rescale-no-log}
Suppose the massless local heat density has the continuum expansion
\begin{equation}
 \Theta_0(t;g)\sim{1\over(4\pi t)^2}
 \{A_0(g)+tA_2(g)+t^2A_4(g)+\cdots\}.                  \label{eq:newS7V27-heat-expansion}
\end{equation}
In \eqref{eq:newS7V27-u-response}, the finite coefficient linear in
\(A_2\) is \(-A_2/(4\pi)^2\), arising from
\(-t\Theta_0\).  The term
\(-a^2t\partial_t\Theta_0/(4\rho^2)\) contributes only
\begin{equation}
 {a^2\over4\rho^2(4\pi)^2}
 \{2A_0t^{-2}+A_2t^{-1}+O(t)\},                         \label{eq:newS7V27-rescale-orders}
\end{equation}
and hence has no (t^0A_2) term.
\end{lemma}

\begin{proof}
Insert \eqref{eq:newS7V27-heat-expansion} into the two terms of
\eqref{eq:newS7V27-u-response} and differentiate powers of \(t\).
The derivative of the constant \(A_4\) term is zero.  The displayed powers
are exact at the orders relevant to the assertion.
\end{proof}

\begin{assessment}[Power subtraction versus the mixed logarithmic owner]
\label{ass:newS7V27-subtraction}
At finite cutoff the second term of \eqref{eq:newS7V27-u-response} changes
the exact quadratic-mass power subtraction, as already encountered in V20.
It must be retained in that subtraction.  It cannot change the continuum
logarithmic (Rm^2) coefficient because
\Cref{lem:newS7V27-rescale-no-log} places its (A_2) contribution at power
(t^{-1}) with an explicit (a^2) factor.
\end{assessment}

\subsection{Reduction to the massless metric (A_2) card}
\label{sec:newS7V27-A2-card}

For a flat-background metric direction
\(h_{\mu\nu}(x)=\widehat h_{\mu\nu}e^{iq\cdot x}\), define the linearized
scalar-curvature symbol
\begin{equation}
 \mathscr S_q(\widehat h)
 :=q^2\tr\widehat h-q_\mu q_\nu\widehat h^{\mu\nu}.      \label{eq:newS7V27-Sq}
\end{equation}
This is the Fourier symbol of
\(\partial_\mu\partial_\nu h^{\mu\nu}-\partial^2\tr h\) in the stated
sign convention.

\begin{definition}[Massless metric (A_2) card, MMAC]
\label{def:newS7V27-MMAC}
MMAC is the following statement for one physical massless Dirac block of
the V25 midpoint overlap operator.  After subtracting the zero-momentum
volume response, the linear metric response of its localized heat density,
denoted \(\mathscr R_{a,t}(q;\widehat h)\), satisfies for
\(a^2\ll t\ll |q|^{-2}\),
\begin{equation}
 (4\pi)^2t\,
 {\mathscr R_{a,t}(q;\widehat h)
       -\mathscr R_{a,t}(0;\widehat h)
  \over \mathscr S_q(\widehat h)}
 =-{1\over3}+\mathcal E_{a,t,q},                        \label{eq:newS7V27-MMAC}
\end{equation}
whenever \(\mathscr S_q(\widehat h)\ne0\), with
\begin{equation}
 |\mathcal E_{a,t,q}|
 \le C\left({a\over\sqrt t}+\sqrt t\,|q|+a^2|q|^2\right) \label{eq:newS7V27-MMAC-error}
\end{equation}
on bounded smooth normal-chart families, together with a Gaussian
majorant and high-branch bound which make this estimate integrable against
\(dt/t\) on the V23 mesoscopic window.
\end{definition}

\begin{remark}[Normalization of MMAC]
\label{rem:newS7V27-MMAC-normalization}
The right side \(-1/3\) is exactly the V24 massless spin-rank-four value
\(A_2=-R/3\).  The subtraction at (q=0) removes the volume response of
\(A_0\); it does not remove scalar curvature because
\(\mathscr S_0=0\).
\end{remark}

\begin{theorem}[MMAC is sufficient for the physical
scalar-curvature--Higgs transfer]
\label{thm:newS7V27-MMAC-sufficient}
Assume MMAC for the V25 midpoint kernel.  Then, after the exact V20
quadratic Higgs subtraction, the physical Standard Model chiral-block heat
density has mixed continuum owner
\begin{equation}
 {\mathfrak a_\bullet\over3(4\pi)^2}R_g|H|^2,            \label{eq:newS7V27-RH-result}
\end{equation}
and the determinant modulus has one half of this coefficient.  The cutoff
error is logarithmically integrable on the V23 mesoscopic window.
\end{theorem}

\begin{proof}
To isolate the \(R|H|^2\) invariant, first set the gauge background to zero.
Freeze \(H\) at the observation point and use the V18 singular-value
decomposition of the constant Standard Model Yukawa map.  The resulting spin
Dirac operator acts identically on each constant internal singular-value
projector, so \Cref{thm:newS7V27-mass-square} applies blockwise with
\(m_j^2=y_j^2|H|^2\).  By
\Cref{prop:newS7V27-u-response,lem:newS7V27-rescale-no-log}, the mixed
logarithmic owner of each physical block is (-m_j^2A_2/(4\pi)^2).
MMAC gives (A_2=-R/3), hence (Rm_j^2/[3(4\pi)^2]).

Summing quark colour and family multiplicities gives exactly
\(\mathfrak a_\bullet|H|^2\) by V18.  This proves
\eqref{eq:newS7V27-RH-result}.  The determinant modulus supplies the usual
factor one half, as in the V24 ledger.  The MMAC majorant and
\eqref{eq:newS7V27-MMAC-error} integrate by the V23 estimates.  Replacing
the frozen Higgs by its Taylor expansion changes only derivative-Higgs or
higher-weight owners and is controlled by the established smooth
parametrix.
\end{proof}

\begin{corollary}[The three-insertion curvature calculation is redundant]
\label{cor:newS7V27-redundant}
For the mixed owner at quadratic Higgs order, no independent evaluation of
one metric vertex with two Yukawa vertices is required.  Such an evaluation
must reproduce \eqref{eq:newS7V27-u-response}; its only undetermined
geometric input is MMAC.
\end{corollary}

\begin{proof}
The exact operator identity \eqref{eq:newS7V27-heat-factor} holds before
any Duhamel expansion.  Taylor coefficients of two identical smooth
finite-dimensional operator functions are identical.
\end{proof}

\subsection{What remains in MMAC}
\label{sec:newS7V27-MMAC-remaining}

\begin{proposition}[Finite list of the remaining analytic rows]
\label{prop:newS7V27-MMAC-rows}
Given V25--V26, MMAC reduces to the following four assertions:
\begin{enumerate}[label=\textup{(\roman*)},leftmargin=3.5em]
\item the first metric vertex from \eqref{eq:newS7V26-L1}, at zero Higgs,
      has the continuum Dirac-square low-branch symbol through order
      \(q^2\);
\item its order-(q^0) part equals the differentiated (A_0) volume term
      and is removed by the subtraction in \eqref{eq:newS7V27-MMAC};
\item the complement of the physical momentum ball obeys an integrable
      divided-difference bound, including low--high mixed propagators; and
\item hypercubic non-(O(4)) tensors at quadratic external momentum are
      (O(a^2q^2)) after the normalized projection.
\end{enumerate}
No Yukawa, mass-rescaling, multiplicity, sign-contact, or moving-Hilbert
row remains.
\end{proposition}

\begin{proof}
The mass and multiplicity rows are closed by
\Cref{thm:newS7V27-mass-square,thm:newS7V27-MMAC-sufficient}.  V26 closes
the sign and Hilbert-space contacts.  A local heat coefficient is then
determined by its low symbol after removal of the volume owner, provided a
dominating high-branch estimate justifies the joint limit.  Hypercubic
symmetry permits lattice artifacts, so their suppression is the final
tensor-consistency row.
\end{proof}

\begin{programme}[V28: prove the low branch of MMAC]
\label{prog:newS7V27-V28}
V28 will expand the midpoint half-edge data in coframe normal gauge through
second external momentum order.  It will insert those expansions into the
V26 first metric contact, use the flat overlap symbol and its sign-resolvent
derivative, subtract the (q=0) volume response, and identify the surviving
continuum tensor with \(-\mathscr S_q/(3(4\pi)^2t)\).  The high and mixed
branches will be kept as an explicit remainder card if their uniform
majorant is not yet proved.
\end{programme}

\begin{firewall}[Claim boundary after seventh-series V27]
\label{fire:newS7V27-boundary}
V27 proves the exact massive overlap square and heat factorization on every
V25 metric background, derives the exact quadratic-mass response, proves
that the lattice time-rescaling term carries no continuum logarithmic
(Rm^2) owner, defines MMAC, and proves that MMAC alone transfers the full
physical Standard Model (R|H|^2) coefficient with the V24 multiplicities.

V27 does not prove MMAC, the lattice (R|H|^2) transfer, the massless metric
high-branch majorant, the complete bosonic and fermionic-measure stress,
the global overlap gap, rough metric or AOP control, measure concentration,
the determinant-line phase, interacting cutoff removal, or the Standard
Model--Einstein continuum.
\end{firewall}

\begin{theorem}[Seventh-series V27 result]
\label{thm:newS7V27-result}
The only new lattice coefficient needed for the mixed geometric Higgs owner
is the massless physical overlap value (A_2=-R/3).  At constant Higgs
background the improved massive square factorizes exactly as
\[
 D_{a,m}^*D_{a,m}
 =\left(1-{a^2m^2\over4\rho^2}\right)D_a^*D_a+m^2I,
\]
and its exact quadratic-mass response is
\eqref{eq:newS7V27-u-response}.  Therefore MMAC implies the physical owner
\(\mathfrak a_\bullet R|H|^2/[3(4\pi)^2]\) and the determinant-modulus
owner one half as large.  The remaining bottleneck is precisely the four
massless analytic rows in \Cref{prop:newS7V27-MMAC-rows}.
\end{theorem}

\begin{proof}
Use \Cref{lem:newS7V27-scalar-GW,
thm:newS7V27-mass-square,
cor:newS7V27-heat-factor,
prop:newS7V27-u-response,
lem:newS7V27-rescale-no-log,
thm:newS7V27-MMAC-sufficient,
prop:newS7V27-MMAC-rows}.
\end{proof}

\paragraph{Parent-version record.}
V27 was formed from
\texttt{Renewal\_SM\_Einstein\_Continuum\_Research\_Notes\_V26.tex}.
The V26 parent had (11\,016) logical source lines, (437\,982) bytes,
and SHA--256
\[
 \texttt{40A401256F5ED77F7D5C170F7F7D9AD0F1C4399AA7ECB17CB499EBD1A932A45E}.
\]
The next compulsory companion audit is V30.

% =============================================================================
% END APPENDED SEVENTH-SERIES V27 MATERIAL
% =============================================================================


% =============================================================================
% BEGIN APPENDED SEVENTH-SERIES V28 MATERIAL
% =============================================================================

\clearpage
\section{Seventh-series V28: the low physical branch of the massless metric
\(A_2\) card}
\label{sec:v7v28}

V27 reduces the mixed (R|H|^2) problem to MMAC, a massless one-metric
response.  V28 proves the low physical branch of that card.  The proof uses
the exact V26 sign contact rather than differentiating a formal continuum
sign, and it keeps a Fourier cutoff which separates the result from the
unproved low--high and high--high contributions.  The remaining gate is
therefore a single divided-difference estimate, stated at the end.

\subsection{Flat Fourier localization and the response}
\label{sec:newS7V28-response}

Work on the flat periodic lattice at zero gauge field and perturb the
metric by
\begin{equation}
 g_{\mu\nu}(s,x)=\delta_{\mu\nu}
 +s\widehat h_{\mu\nu}e^{iq\cdot x}.                    \label{eq:newS7V28-metric-wave}
\end{equation}
The real perturbation is obtained by adding the conjugate mode; the complex
mode is used only to read one Fourier coefficient.  Let
\(\widehat{\mathcal L}_a(s)=D_a(s)^*D_a(s)\) be the massless physical
overlap square.  If \(\chi_{-q}\) denotes multiplication by
\(e^{-iq\cdot x}\), define the normalized response
\begin{equation}
 \mathscr R_{a,t}(q;\widehat h)
 :=-{t\over V_a}\int_0^1
 \Tr\!\left[
  \chi_{-q}e^{-(1-\theta)t\widehat{\mathcal L}_a(0)}
  (\widehat{\mathcal L}_a)_h
  e^{-\theta t\widehat{\mathcal L}_a(0)}\right]d\theta, \label{eq:newS7V28-response}
\end{equation}
where \(V_a\) is the coordinate volume.  Duhamel's formula identifies this
with the first variation of the corresponding local heat-density Fourier
coefficient.

Choose an even \(\zeta\in C_c^\infty((-\pi,\pi)^4)\) which is one on
\(|k|\le\kappa/2\) and zero on \(|k|\ge\kappa\), for a fixed sufficiently
small \(\kappa\).  Let \(P_{\rm lo,a}\) be the flat Fourier multiplier
\(\zeta(ap)\), and insert it immediately on both sides of
\((\widehat{\mathcal L}_a)_h\) in \eqref{eq:newS7V28-response}.  The
resulting quantity is \(\mathscr R^{\rm lo}_{a,t}\); the difference is
\(\mathscr R^{\rm hi}_{a,t}\).  The smooth overlap between the two cutoffs
is included in the high remainder, so no sharp momentum boundary is used.

\subsection{Low symbol of the midpoint Wilson kernel}
\label{sec:newS7V28-Wilson-symbol}

Let \(p'=p+q\) and put
\(P=1+|p|+|p'|\).  All estimates below are for
\(|ap|,|ap'|\le\kappa\), with constants uniform for bounded
\(\widehat h\).

\begin{lemma}[Flat and metric-vertex Wilson expansions]
\label{lem:newS7V28-Wilson-expansion}
The V25 midpoint Wilson kernel and its first metric contact obey
\begin{align}
 \widehat D_{W,a}(p)
 &=i\gamma^\mu p_\mu+{ra\over2}|p|^2I
   +a^2\mathcal R_{W,0}(a,p),                            \label{eq:newS7V28-DW-flat}\
 (\widehat D_{W,a})_h(p',p)
 &=\mathcal V_W(h;p',p)+a\mathcal R_{W,1}(a;h,p',p),    \label{eq:newS7V28-DW-vertex}
\end{align}
where \(\mathcal V_W\) is the first variation, in the fixed half-density
identification, of the continuum skew Dirac operator at the flat metric,
and
\begin{equation}
 \|\mathcal R_{W,0}\|\le CP^3,qquad
 \|\mathcal R_{W,1}\|\le C\|\widehat h\|P^2.           \label{eq:newS7V28-Wilson-remainders}
\end{equation}
The same bounds hold after up to two derivatives in the external momentum
\(q\), with the corresponding powers of \(P\) lowered.
\end{lemma}

\begin{proof}
For the flat datum, V25 gives
\[
 {i\over a}\sum_\mu\gamma^\mu\sin(ap_\mu)
 +{r\over a}\sum_\mu\{1-\cos(ap_\mu)\}I.
\]
Taylor's theorem proves \eqref{eq:newS7V28-DW-flat} and the first bound.

For \eqref{eq:newS7V28-DW-vertex}, differentiate the midpoint formula
\eqref{eq:newS7V26-DW1}.  The coframe contact is evaluated at the edge
midpoint, producing the symmetric phase
\(e^{iaq_\mu/2}\); the two half-edge transport contacts are the ordered
integrals \eqref{eq:newS7V26-P1}.  Taylor expansion of these phases and of
the spin connection of \eqref{eq:newS7V28-metric-wave} gives exactly the
fixed-half-density continuum Dirac variation \(\mathcal V_W\).  Every
discarded trigonometric term contains one additional factor
\(aP\).  The Wilson metric vertex already carries its explicit factor
\(a\), so it belongs to the displayed remainder.  Differentiating the
finite trigonometric formula in \(q\) proves the last assertion.
\end{proof}

\begin{remark}[Continuum vertex identification]
\label{rem:newS7V28-continuum-vertex}
The phrase ``continuum Dirac variation'' in
\Cref{lem:newS7V28-Wilson-expansion} is not an assumption: it is the sum of
the differentiated inverse-coframe principal term and the differentiated
torsion-free spin transport, in the same half-density identification as
V26.  Metric compatibility is the identity which combines their
zero-order pieces into the variation of the geometric Dirac operator.
\end{remark}

\subsection{Passing the vertex through the overlap sign}
\label{sec:newS7V28-overlap-symbol}

For the flat kernel set
\begin{equation}
 \widetilde H_{a,0}(p)
 =\Gamma\{a\widehat D_{W,a}(p)-\rho I\}.                \label{eq:newS7V28-Hflat}
\end{equation}
Its spectrum stays in the V25 gap region.  Apply the V26 fixed-contour
formula pointwise in \((p',p)\) to the metric vertex.

\begin{theorem}[Low overlap vertex consistency]
\label{thm:newS7V28-overlap-vertex}
On the low support,
\begin{align}
 D_a(p)&=i\gamma^\mu p_\mu+a\mathcal R_{D,0}(a,p),       \label{eq:newS7V28-D-low}\
 (D_a)_h(p',p)
 &=\mathcal V_{\mathcal D}(h;p',p)
   +a\mathcal R_{D,1}(a;h,p',p),                        \label{eq:newS7V28-D-vertex}
\end{align}
where \(\mathcal V_{\mathcal D}=\mathcal V_W\) is the continuum metric
Dirac vertex and
\begin{equation}
 \|\mathcal R_{D,0}\|\le CP^2,qquad
 \|\mathcal R_{D,1}\|\le C\|\widehat h\|P^2.          \label{eq:newS7V28-D-remainders}
\end{equation}
The estimates remain valid after two \(q\)-derivatives.
\end{theorem}

\begin{proof}
By \eqref{eq:newS7V28-DW-flat},
\(\widetilde H_{a,0}=-\rho\Gamma+O(aP)\).  On the fixed contour of V26,
the resolvent admits a uniformly convergent Neumann expansion for
\(aP\le\kappa\) after \(\kappa\) is chosen small.  Insert this expansion in
\eqref{eq:newS7V26-sign-contour} and
\eqref{eq:newS7V26-sign1}.  The zeroth sign is \(-\Gamma\), so it cancels
the identity in
\(D_a=(\rho/a)(I+\Gamma\epsilon_a)\).  The first off-chiral divided
difference equals \(1/\rho\); after multiplication by \(\rho/a\), the
order-(a) Wilson and metric-vertex terms are respectively
\(i\gamma\cdot p\) and \(\mathcal V_{\mathcal D}\).

Every remaining resolvent word contains either two factors (aP) or the
explicit \(a^2P^3\) remainder of
\Cref{lem:newS7V28-Wilson-expansion}.  Multiplication by \(\rho/a\) leaves
the bounds in \eqref{eq:newS7V28-D-remainders}.  The resolvent gap and the
differentiated bounds in
\Cref{lem:newS7V28-Wilson-expansion} justify two \(q\)-derivatives term by
term.
\end{proof}

\begin{corollary}[Low heat-square vertex]
\label{cor:newS7V28-square-vertex}
Let \(L_0(p)=|p|^2I\) be the continuum flat Dirac square and
\(\mathcal V_L(h;p',p)\) its first metric variation.  Then
\begin{align}
 D_a(p)^*D_a(p)&=|p|^2I+a\mathcal R_{L,0}(a,p),          \label{eq:newS7V28-L-low}\
 (D_a^*D_a)_h(p',p)
 &=\mathcal V_L(h;p',p)+a\mathcal R_{L,1}(a;h,p',p),    \label{eq:newS7V28-L-vertex}
\end{align}
with
\begin{equation}
 \|\mathcal R_{L,0}\|\le CP^3,qquad
 \|\mathcal R_{L,1}\|\le C\|\widehat h\|P^3,         \label{eq:newS7V28-L-remainders}
\end{equation}
also after two external-momentum derivatives.
\end{corollary}

\begin{proof}
Apply the product rule to \Cref{thm:newS7V28-overlap-vertex}.  The leading
terms are the continuum square and its continuum metric variation.  Every
other term contains one factor (a) and at most three powers of the two
momenta.
\end{proof}

\subsection{The low heat response}
\label{sec:newS7V28-low-heat}

\begin{lemma}[Gaussian low-branch majorant]
\label{lem:newS7V28-low-majorant}
There are \(c,C>0\), independent of \(a,t,p,q\) on the low support, such
that the integrand of \(\mathscr R^{\rm lo}_{a,t}\), together with up to two
\(q\)-derivatives after the zero-momentum subtraction, is bounded by
\begin{equation}
 C\|\widehat h\|,tP^2(1+aP)^C
 e^{-ct\{|p|^2+|p+q|^2\}}.                              \label{eq:newS7V28-Gaussian}
\end{equation}
The same expression with one additional factor (aP) bounds the
lattice-to-continuum difference.
\end{lemma}

\begin{proof}
For sufficiently small \(\kappa\),
\eqref{eq:newS7V28-L-low} implies
\(D_a(p)^*D_a(p)\ge c|p|^2\) on the physical branch.  The heat factors in
\eqref{eq:newS7V28-response} therefore give the Gaussian.  The vertex bound
comes from \eqref{eq:newS7V28-L-vertex}.  Duhamel differentiation of either
heat factor introduces another vertex and a simplex integral; the V28
symbol bounds absorb the resulting polynomial.  The difference estimate
uses the explicit factor (aP) in
\eqref{eq:newS7V28-L-remainders}.
\end{proof}

\begin{theorem}[Low physical branch of MMAC]
\label{thm:newS7V28-low-MMAC}
For \(a^2\ll t\ll |q|^{-2}\),
\begin{equation}
 (4\pi)^2t\,
 {\{\mathscr R^{\rm lo}_{a,t}(q;\widehat h)
       -\mathscr R^{\rm lo}_{a,t}(0;\widehat h)\}
  \over\mathscr S_q(\widehat h)}
 =-{1\over3}+\mathcal E^{\rm lo}_{a,t,q},               \label{eq:newS7V28-low-MMAC}
\end{equation}
and, uniformly away from the null set
\(\mathscr S_q(\widehat h)=0\),
\begin{equation}
 |\mathcal E^{\rm lo}_{a,t,q}|
 \le C\left({a\over\sqrt t}+\sqrt t\,|q|+a^2|q|^2\right).
                                                               \label{eq:newS7V28-low-error}
\end{equation}
\end{theorem}

\begin{proof}
Insert \Cref{cor:newS7V28-square-vertex} in the Duhamel response and use
\Cref{lem:newS7V28-low-majorant}.  After \(p=k/\sqrt t\), the additional
factor \(aP\) in the lattice error integrates to \(O(a/\sqrt t)\).
Taylor expansion in the external momentum under the Gaussian gives the
next continuum heat coefficient with error \(O(\sqrt t|q|)\).

The continuum integral is the first metric response of the flat geometric
Dirac heat kernel.  Its \(A_0\) response is independent of \(q\) and cancels
in the displayed subtraction.  V17's Lichnerowicz calculation gives the
remaining spin-rank-four coefficient \(A_2=-R/3\); under
\eqref{eq:newS7V28-metric-wave}, \(\delta R=\mathscr S_q(\widehat h)\).
This proves the leading term.  Finally, Taylor expansion of sine and cosine
shows that the first hypercubic tensor absent from the continuum integral
has two additional powers of \(aq\), giving \(O(a^2|q|^2)\) after division
by \(\mathscr S_q\).
\end{proof}

\begin{corollary}[Suppression of low-branch hypercubic competitors]
\label{cor:newS7V28-hypercubic}
At quadratic external momentum, every non-\(O(4)\) tensor in the normalized
low-branch response is contained in the vanishing
\(O(a/\sqrt t)\) remainder of \eqref{eq:newS7V28-low-error}; the leading
even trigonometric correction is more sharply \(O(a^2/t)\).  An anisotropy
formed only from the external momentum begins at relative order
\(a^2|q|^2\).
\end{corollary}

\begin{proof}
Hypercubic symmetry alone would allow, for example,
\(\sum_\mu\widehat h_{\mu\mu}q_\mu^2\), so symmetry by itself is
insufficient.  The explicit expansions
\eqref{eq:newS7V28-DW-flat}--\eqref{eq:newS7V28-L-remainders} show that the
leading symbol and its order-\(a\) Wilson correction are continuum
covariant.  Direction-dependent sine and cosine terms first carry two
additional powers of the dimensionless lattice momenta.  Two loop-momentum
powers integrate to \(O(a^2/t)\); two purely external powers give
\(O(a^2|q|^2)\).  Both are bounded by
\eqref{eq:newS7V28-low-error} on the mesoscopic window.
\end{proof}

\subsection{The one remaining MMAC gate}
\label{sec:newS7V28-high-gate}

\begin{definition}[Metric high-branch card, MHBC]
\label{def:newS7V28-MHBC}
MHBC is the existence of \(\sigma>0\) and \(c,C,N>0\) such that, after the
same zero-momentum subtraction as in MMAC,
\begin{equation}
 (4\pi)^2t\left|
 {\mathscr R^{\rm hi}_{a,t}(q;\widehat h)
       -\mathscr R^{\rm hi}_{a,t}(0;\widehat h)
  \over\mathscr S_q(\widehat h)}\right|
 \le C\left\{left({a^2\over t}\right)^\sigma
 +(1+t/a^2)^Ne^{-ct/a^2}\right\}.                      \label{eq:newS7V28-MHBC}
\end{equation}
The high response includes low--high and high--high divided differences of
the complete V26 metric vertex; it is not merely the trace of the
high--high heat block.
\end{definition}

\begin{proposition}[Low theorem plus MHBC proves MMAC]
\label{prop:newS7V28-MHBC-sufficient}
\Cref{thm:newS7V28-low-MMAC} and MHBC imply MMAC and therefore, by V27, the
physical Standard Model (R|H|^2) transfer.
\end{proposition}

\begin{proof}
Add the low and high responses.  On the V23 mesoscopic window
\(a^2/t\to0\); both terms in \eqref{eq:newS7V28-MHBC} vanish and are
integrable against (dt/t) after the established lower cutoff is inserted.
Combine with \eqref{eq:newS7V28-low-error} and apply
\Cref{thm:newS7V27-MMAC-sufficient}.
\end{proof}

\begin{assessment}[Status after the low-branch proof]
\label{ass:newS7V28-status}
V28 proves MMAC rows \textup{(i)}, \textup{(ii)}, and the low-branch part of
row \textup{(iv)} from V27.  The coefficient (-1/3) is obtained from the
actual V25 midpoint kernel through the V26 sign contact.  The only remaining
MMAC assertion is MHBC.  It cannot be replaced by pointwise doubler
suppression because the metric vertex couples distinct momenta and creates
low--high divided differences.
\end{assessment}

\begin{programme}[V29: low--high resolvent separation]
\label{prog:newS7V28-V29}
V29 will write the first heat divided difference in the spectral projectors
of the flat overlap square.  It will use the Wilson separation
\(\lambda_{\rm hi}\ge c/a^2\), integrate the interpolation parameter before
taking norms, and keep the denominator
\((\lambda_{\rm hi}-\lambda_{\rm lo})^{-1}\).  The target is the power
\(a^2/t\) for one low--high bridge and exponential suppression for a
high--high bridge, followed by two external-momentum differences required
by the scalar-curvature projection.
\end{programme}

\begin{firewall}[Claim boundary after seventh-series V28]
\label{fire:newS7V28-boundary}
V28 defines the exact localized metric response and its smooth low/high
split, proves the low Wilson and overlap metric-vertex expansions through
two external-momentum derivatives, proves a Gaussian low-branch majorant,
and proves the low physical branch of MMAC with coefficient (-1/3) and an
integrable mesoscopic error.

V28 does not prove MHBC, full MMAC, the lattice (R|H|^2) transfer, the
complete bosonic and fermionic-measure stress, a global overlap gap, rough
metric or AOP control, measure concentration, the determinant-line phase,
interacting cutoff removal, or the Standard Model--Einstein continuum.
\end{firewall}

\begin{theorem}[Seventh-series V28 result]
\label{thm:newS7V28-result}
The low physical momentum branch of the V25 midpoint overlap kernel has the
massless scalar-curvature response
\[
 \mathscr R^{\rm lo}_{a,t}(q;\widehat h)
 -\mathscr R^{\rm lo}_{a,t}(0;\widehat h)
 =-{\mathscr S_q(\widehat h)\over3(4\pi)^2t}
 \left\{1+O\!\left({a\over\sqrt t}+\sqrt t|q|+a^2|q|^2\right)\right\}.
\]
All Wilson, spin-transport, sign, adjoint, and half-density contacts are
included.  The complete MMAC and hence the mixed (R|H|^2) transfer now
reduce to the single low--high/high--high estimate MHBC.
\end{theorem}

\begin{proof}
Use \Cref{lem:newS7V28-Wilson-expansion,
thm:newS7V28-overlap-vertex,
cor:newS7V28-square-vertex,
lem:newS7V28-low-majorant,
thm:newS7V28-low-MMAC,
cor:newS7V28-hypercubic,
prop:newS7V28-MHBC-sufficient}.
\end{proof}

\paragraph{Parent-version record.}
V28 was formed from
\texttt{Renewal\_SM\_Einstein\_Continuum\_Research\_Notes\_V27.tex}.
The V27 parent had (11\,369) logical source lines, (452\,369) bytes,
and SHA--256
\[
 \texttt{2860A332030636C9B54C08F5969C3E87D1F92DF77620D286509B8989CB186412}.
\]
The next compulsory companion audit is V30.

% =============================================================================
% END APPENDED SEVENTH-SERIES V28 MATERIAL
% =============================================================================


% =============================================================================
% BEGIN APPENDED SEVENTH-SERIES V29 MATERIAL
% =============================================================================

\clearpage
\section{Seventh-series V29: kinematic high-branch separation and closure
of the smooth (R|H|^2) transfer}
\label{sec:v7v29}

V28 left MHBC as the only missing row of the massless metric \(A_2\) card.
The anticipated low--high denominator is not needed for the one-mode
curvature projection.  The metric vertex conserves momentum up to the
single external momentum \(q\); if \(|aq|\) is smaller than the buffer in the
smooth low cutoff, a term outside that cutoff has both adjacent momenta a
fixed distance from the unique overlap zero.  Both heat factors are then
uniformly high and the complete remainder is exponentially small.  V29
proves this statement, closes MMAC, and combines it with the exact V27 mass
factorization to transfer the smooth mixed geometric Higgs owner.

\subsection{The unique flat overlap zero}
\label{sec:newS7V29-zero}

Let (k=ap\in[-\pi,\pi]^4).  For the flat V25 Wilson kernel write
\begin{align}
 s^2(k)&:=\sum_\mu\sin^2k_\mu,                           \label{eq:newS7V29-s}\
 b(k)&:=r\sum_\mu(1-\cos k_\mu)-\rho,                  \label{eq:newS7V29-b}\
 \omega(k)&:=\sqrt{s^2(k)+b(k)^2}.                      \label{eq:newS7V29-omega}
\end{align}
The flat overlap symbol is
\begin{equation}
 aD_a(k)=\rho\left[I+{b(k)+i\gamma^\mu\sin k_\mu\over\omega(k)}\right].
                                                               \label{eq:newS7V29-overlap-symbol}
\end{equation}

\begin{lemma}[Only the physical origin is massless]
\label{lem:newS7V29-unique-zero}
For (0<\rho<2r), (D_a(k)=0) if and only if (k=0) modulo the periodic
Brillouin lattice.  Consequently, for every (0<\kappa_0<\pi), there is
\(c_{\kappa_0}>0\), independent of \(a\), such that
\begin{equation}
 D_a(k)^*D_a(k)\ge {c_{\kappa_0}\over a^2}I
 \quad\hbox{whenever }|k|\ge\kappa_0.                  \label{eq:newS7V29-high-gap}
\end{equation}
\end{lemma}

\begin{proof}
Equation \eqref{eq:newS7V29-overlap-symbol} vanishes only if every
\(\sin k_\mu=0\) and \(b(k)/|b(k)|=-1\).  Thus each component is zero or
\(\pi\), and if \(n\) components equal \(\pi\), then
\(b=2rn-\rho\).  This is negative only for \(n=0\) because
\(0<\rho<2r\).  Hence all components vanish.  The smallest eigenvalue of
\(a^2D_a(k)^*D_a(k)\) is a continuous positive function on the compact set
\(|k|\ge\kappa_0\); its minimum is positive.
\end{proof}

\subsection{The cutoff buffer eliminates a low--high bridge}
\label{sec:newS7V29-buffer}

Use the V28 cutoff \(\zeta\), equal to one on
\(|k|\le\kappa/2\) and zero on \(|k|\ge\kappa\).  The high response is the
full response minus the term containing
\(\zeta(ap)\zeta(a(p+q))\).

\begin{lemma}[Buffered momentum separation]
\label{lem:newS7V29-buffer}
Assume \(|aq|\le\kappa/8\).  If
\begin{equation}
 1-\zeta(ap)\zeta(a(p+q))\ne0,                           \label{eq:newS7V29-high-support}
\end{equation}
then, for the original smooth cutoff,
\begin{equation}
 |ap|\ge{3\kappa\over8},\qquad
 |a(p+q)|\ge{3\kappa\over8}.                            \label{eq:newS7V29-both-high}
\end{equation}
\end{lemma}

\begin{proof}
Because \(\zeta=1\) on \(|k|\le\kappa/2\), a nonzero left side in
\eqref{eq:newS7V29-high-support} forces at least one of \(|ap|\) and
\(|a(p+q)|\) to exceed \(\kappa/2\).  Their distance on the Brillouin
torus is at most \(|aq|\le\kappa/8\), so the triangle inequality bounds
the other one below by \(3\kappa/8\).
\end{proof}

\begin{corollary}[Both heat energies are high]
\label{cor:newS7V29-both-high}
On every term of the high response and for \(|aq|\le\kappa/8\),
\begin{equation}
 D_a(p)^*D_a(p)\ge{c_\kappa\over a^2}I,qquad
 D_a(p+q)^*D_a(p+q)\ge{c_\kappa\over a^2}I.             \label{eq:newS7V29-two-gaps}
\end{equation}
\end{corollary}

\begin{proof}
Apply \Cref{lem:newS7V29-unique-zero} to
\eqref{eq:newS7V29-both-high}.
\end{proof}

\subsection{Complete high-response bound}
\label{sec:newS7V29-high-bound}

Let
\begin{equation}
 \mathcal V_{a,h}(p+q,p):=(D_a^*D_a)_h(p+q,p)           \label{eq:newS7V29-vertex}
\end{equation}
be the complete V26 first metric heat-square vertex.  It includes the
coframe, spin transport, Wilson, overlap sign, adjoint, and half-density
identification contacts.

\begin{lemma}[Global vertex and momentum-derivative bounds]
\label{lem:newS7V29-vertex-bounds}
For (j=0,1,2),
\begin{equation}
 \|\partial_q^j\mathcal V_{a,h}(p+q,p)\|
 \le C_j\|\widehat h\|\,a^{j-2}                        \label{eq:newS7V29-vertex-bounds}
\end{equation}
on the Brillouin zone and on the closed V25 gap subchart.  Derivatives of a
flat heat factor obey
\begin{equation}
 \|\partial_q^j e^{-\theta tD_a(p+q)^*D_a(p+q)}\|
 \le C_j a^j(1+t/a^2)^j
 e^{-c\theta tD_a(p+q)^*D_a(p+q)}                      \label{eq:newS7V29-heat-derivatives}
\end{equation}
with a possibly smaller positive constant in the exponential.
\end{lemma}

\begin{proof}
The dimensionless overlap symbol \(aD_a(k)\) and its metric sign contact
are smooth on the compact Brillouin zone because the V25 Wilson sign gap is
uniform.  Each physical \(q\)-derivative is \(a\) times a derivative in
(k'=a(p+q)).  Since the square vertex contains two factors of scale
\(a^{-1}\), this proves \eqref{eq:newS7V29-vertex-bounds}.

For the heat factor, apply Duhamel's formula to a \(q\)-derivative.  Each
derivative of (D_a^*D_a) is bounded by (Ca^{-1}) for the first derivative
and by (C) for the second, while the simplex integral supplies the factor
\(t\).  Iteration gives a polynomial in \(t/a^2\).  The elementary inequality
(x^me^{-x}\le C_{m,c}e^{-cx}) absorbs the portions of this polynomial
needed to retain the smaller exponential.
\end{proof}

\begin{theorem}[MHBC with exponential remainder]
\label{thm:newS7V29-MHBC}
There are \(c,C>0\) and an integer \(N\), independent of \(a,t,q\), such
that for \(|aq|\le\kappa/8\) and
\(\mathscr S_q(\widehat h)\ne0\),
\begin{equation}
 (4\pi)^2t\left|
 {\mathscr R^{\rm hi}_{a,t}(q;\widehat h)
       -\mathscr R^{\rm hi}_{a,t}(0;\widehat h)
  \over\mathscr S_q(\widehat h)}\right|
 \le C(1+t/a^2)^Ne^{-ct/a^2}.                           \label{eq:newS7V29-MHBC}
\end{equation}
Thus MHBC holds, with its power term omitted.
\end{theorem}

\begin{proof}
Write the high response using \eqref{eq:newS7V28-response} and the smooth
partition in \Cref{lem:newS7V29-buffer}.  By
\Cref{cor:newS7V29-both-high}, for every \(0\le\theta\le1\),
\[
 e^{-(1-\theta)t\lambda_{p+q}}e^{-\theta t\lambda_p}
 \le e^{-c_\kappa t/a^2};
\]
there is no endpoint loss because both eigenvalues are high.  Insert
\Cref{lem:newS7V29-vertex-bounds}.  The normalized momentum sum has at most
(Ca^{-4}) modes per unit coordinate volume, and all remaining negative
powers of \(a\) can be absorbed into a fixed polynomial in \(t/a^2\) times
the exponential when \(t\ge\varepsilon_q(a)\gg a^2\).

It remains to justify division by the quadratic scalar symbol.  The even
cutoff, the pairing \(p\leftrightarrow-p\), and reflection covariance make
the first external-momentum derivative of the zero-momentum-subtracted
response vanish.  Taylor's formula with integral remainder expresses the
numerator as two \(q\)-derivatives evaluated on the segment from \(0\) to
\(q\).  Apply \eqref{eq:newS7V29-vertex-bounds}--
\eqref{eq:newS7V29-heat-derivatives}; the assumed exclusion of
\(\mathscr S_q(\widehat h)=0\) controls the fixed tensor projection.  The
result is \eqref{eq:newS7V29-MHBC} after increasing \(N\) and decreasing
\(c\).
\end{proof}

\begin{remark}[Why the feared low--high denominator is absent]
\label{rem:newS7V29-no-bridge}
A general localized metric perturbation contains many Fourier momenta and
can connect the physical origin to a high branch.  MMAC reads one Taylor
coefficient using one external \(q\) with \(|aq|\to0\).  In that acquisition
geometry, \Cref{lem:newS7V29-buffer} forces both endpoints of every cutoff
remainder away from the origin.  This special kinematic fact is what permits
the exponential estimate; it is not asserted for an arbitrary rough metric
field.
\end{remark}

\subsection{Closure of MMAC and the mixed smooth owner}
\label{sec:newS7V29-closure}

\begin{theorem}[MMAC for the V25 midpoint overlap kernel]
\label{thm:newS7V29-MMAC}
On the perturbative smooth metric--gauge chart of V25, at zero gauge
background for the isolated scalar-curvature channel, the massless physical
overlap heat response satisfies
\begin{equation}
 (4\pi)^2t\,
 {\mathscr R_{a,t}(q;\widehat h)
       -\mathscr R_{a,t}(0;\widehat h)
  \over\mathscr S_q(\widehat h)}
 =-{1\over3}+O\!\left({a\over\sqrt t}+\sqrt t|q|+a^2|q|^2
ight)
 +O\!\left((1+t/a^2)^Ne^{-ct/a^2}\right).              \label{eq:newS7V29-MMAC}
\end{equation}
The error is logarithmically integrable on the V23 mesoscopic window.
\end{theorem}

\begin{proof}
Add the low theorem \eqref{eq:newS7V28-low-MMAC} and the high estimate
\eqref{eq:newS7V29-MHBC}.  The first three errors integrate exactly as in
V23.  For the last, set \(u=t/a^2\); the lower endpoint tends to infinity,
and every polynomial times \(e^{-cu}\) is integrable against \(du/u\).
\end{proof}

\begin{theorem}[Smooth lattice transfer of the mixed geometric Higgs owner]
\label{thm:newS7V29-RH-transfer}
For the V25 midpoint overlap kernel, bounded smooth Standard Model
backgrounds in the perturbative normal-chart region, and the exact V20
quadratic Higgs subtraction, the physical chiral-block heat density has
mixed continuum owner
\begin{equation}
 {\mathfrak a_\bullet\over3(4\pi)^2}R_g|H|^2,            \label{eq:newS7V29-RH-physical}
\end{equation}
while the determinant modulus has owner
\begin{equation}
 {\mathfrak a_\bullet\over6(4\pi)^2}R_g|H|^2.            \label{eq:newS7V29-RH-det}
\end{equation}
The joint cutoff error is logarithmically integrable and vanishes in the
order \(a\downarrow0\), then local proper-time scale \(t_0\downarrow0\).
\end{theorem}

\begin{proof}
\Cref{thm:newS7V29-MMAC} proves MMAC.  Apply the exact blockwise mass
factorization and multiplicity theorem
\Cref{thm:newS7V27-MMAC-sufficient}.  V27 shows that the finite-\(a\) time
rescaling affects only the exact power subtraction, not the logarithmic
mixed owner.  The physical/determinant factor two is the V24 ledger.
\end{proof}

\begin{corollary}[The smooth local matter coefficient including curvature]
\label{cor:newS7V29-full-smooth}
Combining V23 and \Cref{thm:newS7V29-RH-transfer}, the physical smooth local
Standard Model heat coefficient through the gauge, Higgs, and mixed scalar
curvature owners is
\begin{equation}
 {1\over3(4\pi)^2}\mathfrak G_{\rm SM}(F)
 +{1\over(4\pi)^2}
 \left\{2\mathfrak a_\bullet|DH|^2
       +2\mathfrak b_\bullet|H|^4
       +{\mathfrak a_\bullet\over3}R_g|H|^2\right\}.    \label{eq:newS7V29-smooth-total}
\end{equation}
This formula lists the physical heat coefficient.  Passing to the
determinant modulus applies its separate one-half ledger to every physical
block.
\end{corollary}

\begin{proof}
Add the independent invariant owners proved in V23 and V29.  Their mixed
cross channels vanish or were subtracted by the V22--V23 ledger.  The final
sentence is the physical/determinant multiplicity distinction maintained
since V20.
\end{proof}

\begin{assessment}[Scope of the completed transfer]
\label{ass:newS7V29-scope}
The result is a deterministic local theorem for bounded smooth backgrounds
in the explicit perturbative normal-chart gap region.  Locality then permits
a finite smooth atlas whose smaller chart interiors carry the coefficient,
provided the global lattice spin structure and the overlap gap are supplied
on overlaps.  V29 has not supplied that global gluing.  It also has not
identified the pure gravitational \(R^2\), \(|{\rm Ric}|^2\), and
\(|{\rm Rm}|^2\) owners or the complete common stress.
\end{assessment}

\begin{programme}[V30: compulsory audit and the global metric interface]
\label{prog:newS7V29-V30}
V30 will perform the compulsory YM/NS audit.  It will then compare the V35
continuum scheduler with the local V25 metric lattice datum and formulate
the exact global spin-atlas compatibility, collar-independence, and
uniform-gap requirements.  If these cannot be derived from the current
scheduler, the programme will move upstream to a global triangulated or
good-cover spin lift before attempting the Einstein stress writer.
\end{programme}

\begin{firewall}[Claim boundary after seventh-series V29]
\label{fire:newS7V29-boundary}
V29 proves the unique flat overlap zero, buffered high-momentum separation,
global metric-vertex derivative bounds on the Brillouin zone, exponential
MHBC, full MMAC, and the smooth local lattice transfer of the physical and
determinant-modulus (R|H|^2) owners.  Together with V23 it gives the smooth
local matter heat coefficient \eqref{eq:newS7V29-smooth-total}.

V29 does not globalize the lattice spin structure or overlap gap, construct
the complete common metric stress, transfer pure gravitational heat owners,
control rough metric or AOP fields, prove measure concentration, construct
the determinant-line phase, remove the interacting cutoff, or construct the
Standard Model--Einstein continuum.
\end{firewall}

\begin{theorem}[Seventh-series V29 result]
\label{thm:newS7V29-result}
The high complement of the V28 curvature response has both heat energies at
least (c/a^2) and is exponentially small on the mesoscopic window.  Hence
MMAC holds with massless coefficient (-1/3).  The exact V27 massive
factorization then proves, for the explicit V25 midpoint overlap kernel on
bounded smooth perturbative normal charts, the physical mixed owner
\[
 {\mathfrak a_\bullet\over3(4\pi)^2}R_g|H|^2
\]
and determinant-modulus owner one half as large, with a logarithmically
integrable joint-cutoff error.
\end{theorem}

\begin{proof}
Use \Cref{lem:newS7V29-unique-zero,
lem:newS7V29-buffer,
cor:newS7V29-both-high,
lem:newS7V29-vertex-bounds,
thm:newS7V29-MHBC,
thm:newS7V29-MMAC,
thm:newS7V29-RH-transfer,
cor:newS7V29-full-smooth}.
\end{proof}

\paragraph{Parent-version record.}
V29 was formed from
\texttt{Renewal\_SM\_Einstein\_Continuum\_Research\_Notes\_V28.tex}.
The V28 parent had (11\,753) logical source lines, (468\,670) bytes,
and SHA--256
\[
 \texttt{3AADB805D94D79F93060531FA54E9D74FF0C95EA6A7DC7F7CBF95C761E2591A0}.
\]
The next compulsory companion audit is V30.

% =============================================================================
% END APPENDED SEVENTH-SERIES V29 MATERIAL
% =============================================================================


% =============================================================================
% BEGIN APPENDED SEVENTH-SERIES V30 MATERIAL
% =============================================================================

\clearpage
\section{Seventh-series V30: compulsory audit and the global spin-mesh
interface}
\label{sec:v7v30}

V29 completes the deterministic smooth mixed owner on a normal chart.  A
continuum Einstein source cannot be assembled by declaring those local
charts independent: the spin fibres, edge transports, Wilson stencil,
volume quadrature, and sign gap must belong to one global finite-cutoff
operator.  V30 performs the compulsory companion audit, constructs the
intrinsic spin-cocycle part of such an operator on any embedded graph, and
states the remaining global mesh card.  It then isolates a quantitative
local-to-global gap mechanism for V31.

\subsection{Compulsory companion audit}
\label{sec:newS7V30-audit}

The companion tasks were active during the audit, so the newest coherent
files were read as immutable byte snapshots:
\begin{center}
\small
\begin{tabular}{lll}
programme & source & snapshot data \\ \hline
Yang--Mills &
\texttt{Renewal\_Yang\_Mills\_Mass\_Gap\_Research\_Notes\_V5.tex}
& (85\,750) bytes, (2\,464) newlines,\\
&&\texttt{6F6CBD7532FA75C668C3B9D1E2A96A47018542B1EA59A7DD09D11415F609E88A}
\\ [0.4em]
Navier--Stokes &
\texttt{Renewal\_NS\_Stability\_Research\_Notes\_V28.tex}
& (276\,340) bytes, (5\,433) newlines,\\
&&\texttt{BA05D9D0AD6B7C0F22C8243E448C7D252CA24916C82651255E58E2B589044C32}.
\end{tabular}
\end{center}

\begin{audit}[V30 companion decision]
\label{audit:newS7V30-companions}
YM sixth-series V5 proves an exact conditional quotient-tree law and a
measure-preserving one-node orientation factorization.  Its local topology
counts have disappeared; the remaining QTC gate is canonical highest-node
ownership when quotient skeletons are glued across nested nodes.  NS V28
proves that a small variance locks one spectral level, obtains an explicit
window floor and an invariant-mean floor, and records that these local floors
remain far below the size needed for its global gate.

Neither note contains a spin-mesh, metric overlap, or Einstein-stress
estimate, so no constant is imported.  Their common applicable point is that
normalized local objects and local lower bounds do not determine a global
gluing.  For the SM programme, the analogue is canonical ownership of every
edge by one global spin mesh and a gap estimate for the resulting whole
operator.  V30 adopts that direction.
\end{audit}

\subsection{Acquisition from the metric source}
\label{sec:newS7V30-source}

\begin{audit}[No global lattice spin mesh in the attached metric source]
\label{audit:newS7V30-metric-source}
The V35 metric source inspected in V24 conditions several sector statements
on fixed spin structure and boundary data and separately develops a
determinant/Pfaffian line atlas.  It also proves continuum coframe and spin-
connection scheduler convergence.  A targeted search finds no good-cover
spinor lattice, triangulation or cubulation, global lattice Wilson stencil,
or uniform global overlap gap.  A partition-of-unity construction in that
source concerns a positive determinant-line control density and explicitly
does not supply the missing phase or cocycle.  It is not a spin-mesh gluing
of the V25 operator.
\end{audit}

\subsection{Intrinsic half-edge data on a spin manifold}
\label{sec:newS7V30-intrinsic}

Let \((M,g)\) be a compact oriented Riemannian spin four-manifold with a
fixed spin structure, spinor bundle \(S\to M\), chirality \(\Gamma\), and
the torsion-free spin connection tensored with a unitary Standard Model
connection.  Let \(G_a=(X_a,E_a)\) be a finite graph embedded in \(M\), each
edge supplied with a path and a distinguished midpoint \(m_e\).

\begin{definition}[Global weighted spin-graph datum]
\label{def:newS7V30-spin-graph}
Choose positive site weights \(w_x\), symmetric positive edge numbers
\(\kappa_e\) and \(\lambda_e\), and for every edge \(e=\{x,y\}\):
\begin{enumerate}[label=\textup{(\roman*)},leftmargin=3.5em]
\item intrinsic parallel transports \(P_{ex}:S_x\to S_{m_e}\) and
      \(P_{ey}:S_y\to S_{m_e}\);
\item a self-adjoint Clifford map \(c_e=c_{g(m_e)}(\vartheta_e)\) for a
      declared midpoint covector \(\vartheta_e\); and
\item an orientation sign \(\sigma_{xy}=-\sigma_{yx}\).
\end{enumerate}
Define \(U_{xy},J_{xy}\) by V25 and set
\begin{align}
 (K_{G,a}\psi)_x
 &:=\sum_{y\sim x}\kappa_e\sqrt{w_y\over w_x}\,J_{xy}\psi_y,
                                                               \label{eq:newS7V30-KG}\
 (L_{G,a}\psi)_x
 &:=\sum_{y\sim x}\lambda_e
 \left(\psi_x-\sqrt{w_y\over w_x}\,U_{xy}\psi_y\right),
                                                               \label{eq:newS7V30-LG}\
 D_{W,G,a}&:=K_{G,a}+{ra\over2}L_{G,a}.                 \label{eq:newS7V30-DWG}
\end{align}
\end{definition}

\begin{theorem}[The spin-cocycle row glues intrinsically]
\label{thm:newS7V30-spin-gluing}
For any good cover with local spin frames and transition lifts
\(g_{ij}:U_i\cap U_j\to\operatorname{Spin}(4)\), the local matrices of
\eqref{eq:newS7V30-KG}--\eqref{eq:newS7V30-DWG} are related on overlaps by
sitewise unitary conjugation.  On triple overlaps the relations satisfy the
spin cocycle exactly.  Thus the operators are independent of the chosen
local frames and define global operators on
\(\bigoplus_{x\in X_a}S_x\).
\end{theorem}

\begin{proof}
Parallel transport and Clifford multiplication are bundle maps.  If the
frames at \(x,y,m_e\) are changed, the midpoint change cancels between
\(P_{ex}^*\), \(c_e\), and \(P_{ey}\), leaving
\[
 U_{xy}^{(j)}=g_{ij}(x)U_{xy}^{(i)}g_{ij}(y)^*,\qquad
 J_{xy}^{(j)}=g_{ij}(x)J_{xy}^{(i)}g_{ij}(y)^*.
\]
This is precisely the sitewise covariance proved in V25.  The transition
lifts satisfy \(g_{ij}g_{jk}=g_{ik}\) for the fixed spin structure, so the
conjugations satisfy the same identity on triple overlaps.  Scalar edge and
site weights do not change.
\end{proof}

\begin{theorem}[Global weighted adjoint algebra]
\label{thm:newS7V30-global-algebra}
On the global weighted space with inner product
\(\sum_xw_x\langle\cdot,\cdot\rangle_{S_x}\),
\begin{equation}
 K_{G,a}^*=-K_{G,a},\qquad L_{G,a}^*=L_{G,a}\ge0,qquad
 D_{W,G,a}^*=\Gamma D_{W,G,a}\Gamma.                   \label{eq:newS7V30-global-algebra}
\end{equation}
Moreover,
\begin{equation}
 \langle\psi,L_{G,a}\psi\rangle_w
 =\sum_{e=\{x,y\}}\lambda_e
 \|\sqrt{w_x}P_{ex}\psi_x-\sqrt{w_y}P_{ey}\psi_y\|^2. \label{eq:newS7V30-global-positive}
\end{equation}
\end{theorem}

\begin{proof}
The V25 proof uses only reversal of \(U,J\), symmetry of the scalar edge
coefficient, and positivity of the conductance.  Those properties hold
edgewise here.  Sum the identical two-site calculation over the global edge
set.
\end{proof}

\begin{assessment}[What the spin cocycle does and does not solve]
\label{ass:newS7V30-spin-status}
A chosen continuum spin structure is enough to glue fibre maps and
parallel transports; no independent chart-boundary spin phase remains in
the finite graph operator.  It does not choose the graph, midpoint
covectors, Dirac weights, Wilson conductances, or volume quadrature.  Those
data determine consistency and doubler removal and cannot be obtained by
conjugating overlapping copies of the coordinate lattice, whose stencils
need not agree.
\end{assessment}

\subsection{The global spin-mesh card}
\label{sec:newS7V30-card}

\begin{definition}[Global spin-mesh card, GSMC]
\label{def:newS7V30-GSMC}
GSMC consists of a cofinal sequence of data in
\Cref{def:newS7V30-spin-graph} satisfying:
\begin{enumerate}[label=\textup{(\roman*)},leftmargin=3.5em]
\item quasi-uniform embedded meshes of size \(a\), uniformly bounded degree,
      and collars subordinate to one finite good cover;
\item positive quadrature \(\sum_xw_xf(x)\to\int_Mf\,d\mu_g\), uniformly
      through two metric derivatives on the declared smooth family;
\item local first-moment and Wilson second-moment identities for
      \(\kappa_e,\lambda_e,\vartheta_e\), giving
      \(K_{G,a}\to\mathcal D_g\) and \(L_{G,a}\) a uniformly elliptic
      covariant Laplace principal symbol, through the contacts used in V26;
\item exact spin--gauge covariance and the global algebra
      \eqref{eq:newS7V30-global-algebra};
\item a global Wilson sign gap \(\|H_{G,a}\psi\|\ge c_*a^{-1}\|\psi\|\)
      for \(H_{G,a}=\Gamma(D_{W,G,a}-\rho/a)\);
\item Gaussian heat locality, including collar replacement errors
      \(O(e^{-d^2/(Ct)})\), and a global high-branch majorant; and
\item conditioning on the fixed spin structure, orientation, boundary
      record, and determinant/Pfaffian line sector used by the physical
      measure.
\end{enumerate}
\end{definition}

\begin{theorem}[GSMC globalizes the V29 mixed owner]
\label{thm:newS7V30-GSMC-sufficient}
Assume GSMC on a compact smooth background.  Then the local V29 coefficient
patches to the global physical heat owner
\begin{equation}
 {\mathfrak a_\bullet\over3(4\pi)^2}
 \int_MR_g|H|^2\,d\mu_g,                                \label{eq:newS7V30-global-RH}
\end{equation}
with determinant-modulus coefficient one half as large.  The result is
independent of the good cover and of the closing collars.
\end{theorem}

\begin{proof}
Rows (i)--(v) put every local calculation on restrictions of one global
overlap operator.  On the interior of a normal chart, rows (ii)--(iii) reduce
the complete metric contact to the V25--V29 symbol with the same low owner;
row (vi) makes replacement by its periodic closing collar exponentially
small.  On overlaps, \Cref{thm:newS7V30-spin-gluing} gives unitary
conjugacy, while \(R|H|^2\) is a scalar, so the local coefficients agree.
Multiply by a subordinate partition of unity and sum; its sum is one, giving
\eqref{eq:newS7V30-global-RH}.  A second cover has a common refinement and
therefore gives the same integral.  Row (vii) keeps the discrete global
sector fixed.  Apply the V24 physical/determinant ledger.
\end{proof}

\subsection{Why a separate global gap proof is necessary}
\label{sec:newS7V30-gap-need}

\begin{proposition}[Gapped chart blocks need not give a gapped gluing]
\label{prop:newS7V30-gap-no-go}
Lower bounds for isolated chart operators do not, without a common global
operator and commutator control, imply a global gap.  Indeed, the two
one-dimensional chart blocks \(H_1=H_2=1\) have gap one, whereas the coupled
self-adjoint operator
\begin{equation}
 H=\begin{pmatrix}1&1\\1&1\end{pmatrix}                \label{eq:newS7V30-gap-counterexample}
\end{equation}
has eigenvalue zero.
\end{proposition}

\begin{proof}
The two diagonal restrictions are the scalar one.  The vector
\((1,-1)\) lies in the kernel of the displayed matrix.  The off-diagonal
gluing is invisible in the isolated restrictions.
\end{proof}

\begin{definition}[IMS gap-gluing row]
\label{def:newS7V30-IMS-row}
For a quadratic partition of unity \(\sum_i\chi_i^2=1\) subordinate to
enlarged mesh collars, the IMS row is:
\begin{align}
 \|H_{G,a}\chi_i\psi\|&\ge{g_0\over a}\|\chi_i\psi\|,
                                                               \label{eq:newS7V30-local-gap}\
 \left(\sum_i\|[H_{G,a},\chi_i]\psi\|^2\right)^{1/2}
 &\le C_\chi\|\psi\|.                                  \label{eq:newS7V30-commutator}
\end{align}
Here \(g_0>0\) and \(C_\chi<\infty\) are independent of \(a\).
\end{definition}

\begin{proposition}[The IMS row is sufficient for the global sign gap]
\label{prop:newS7V30-IMS-sufficient}
If \eqref{eq:newS7V30-local-gap}--
\eqref{eq:newS7V30-commutator} hold, then
\begin{equation}
 \|H_{G,a}\psi\|\ge\left({g_0\over a}-C_\chi\right)\|\psi\|. \label{eq:newS7V30-IMS-target}
\end{equation}
In particular the global sign kernel is gapped for \(a<g_0/C_\chi\).
\end{proposition}

\begin{proof}
By the local gaps and \(\sum_i\|\chi_i\psi\|^2=\|\psi\|^2\),
\[
 {g_0\over a}\|\psi\|
 \le\left(\sum_i\|H_{G,a}\chi_i\psi\|^2\right)^{1/2}.
\]
Use \(H\chi_i=\chi_iH+[H,\chi_i]\), the triangle inequality in the direct
sum over \(i\),
\(\sum_i\|\chi_iH\psi\|^2=\|H\psi\|^2\), and
\eqref{eq:newS7V30-commutator}.  Rearrangement gives
\eqref{eq:newS7V30-IMS-target}.
\end{proof}

\begin{programme}[V31: prove the IMS row for a subordinate global mesh]
\label{prog:newS7V30-V31}
V31 will estimate the commutator of the global nearest-neighbour Wilson
kernel with lattice cutoffs sampled from a fixed smooth quadratic partition
of unity.  Since an edge difference of \(\chi_i\) is \(O(a)\), the
\(a^{-1}\) hopping scale should leave an \(O(1)\) commutator.  It will then
show that a collar-supported vector sees the V25 local gap up to an
\(O(1)\) boundary replacement error and apply
\Cref{prop:newS7V30-IMS-sufficient}.  If the graph moment row cannot supply
that collar comparison, it will remain as the next upstream mesh gate.
\end{programme}

\begin{firewall}[Claim boundary after seventh-series V30]
\label{fire:newS7V30-boundary}
V30 performs the compulsory YM/NS audit, records the absence of a global
lattice spin mesh in the attached metric source, constructs an intrinsic
weighted spin-graph Wilson operator on any fixed spin manifold, proves its
spin-cocycle independence and global weighted adjoint algebra, defines GSMC,
proves GSMC sufficient to globalize the V29 mixed owner, proves that
separate local gaps do not imply a global gap, and proves the IMS row
sufficient for one.

V30 does not construct a cofinal mesh satisfying the GSMC moment rows,
prove the IMS hypotheses, globalize the overlap sign, construct the common
Einstein stress, transfer pure gravitational owners, control rough metric or
AOP fields, prove measure concentration, construct the determinant-line
phase, remove the interacting cutoff, or construct the Standard Model--
Einstein continuum.
\end{firewall}

\begin{theorem}[Seventh-series V30 result]
\label{thm:newS7V30-result}
On a fixed spin manifold, the half-edge parallel transports and Clifford
maps glue exactly across a good cover and give the global weighted Wilson
algebra \eqref{eq:newS7V30-global-algebra}.  The unresolved global data are
the cofinal mesh moments, heat locality, and whole-operator sign gap listed
in GSMC.  The latter gap follows quantitatively from the local collar gap
and commutator rows by
\(|H_{G,a}\psi\|\ge(g_0/a-C_\chi)\|\psi\|\).
\end{theorem}

\begin{proof}
Use \Cref{audit:newS7V30-companions,
audit:newS7V30-metric-source,
thm:newS7V30-spin-gluing,
thm:newS7V30-global-algebra,
thm:newS7V30-GSMC-sufficient,
prop:newS7V30-gap-no-go,
prop:newS7V30-IMS-sufficient}.
\end{proof}

\paragraph{Parent-version record.}
V30 was formed from
\texttt{Renewal\_SM\_Einstein\_Continuum\_Research\_Notes\_V29.tex}.
The V29 parent had (12\,098) logical source lines, (483\,296) bytes,
and SHA--256
\[
 \texttt{67EB65610F9AD52BB799DF830646A3AD9F76A4D2800BBAD8EB9264C9E85797D1}.
\]
The next compulsory companion audit is V35.

% =============================================================================
% END APPENDED SEVENTH-SERIES V30 MATERIAL
% =============================================================================


% =============================================================================
% BEGIN APPENDED SEVENTH-SERIES V31 MATERIAL
% =============================================================================

\clearpage
\section{Seventh-series V31: the global Wilson commutator and IMS gap
gluing}
\label{sec:v7v31}

V30 proves that intrinsic half-edge spin data glue, and reduces the global
sign problem to a local collar gap plus a cutoff commutator bound.  V31
proves that commutator bound for the complete graph Wilson kernel, proves
exact collar replacement for a nearest-neighbour operator, and obtains the
whole-operator sign gap.  The result is conditional only on the existence
of compatible local closing models at every mesh star; that remaining seam
condition is stated separately and is not hidden in the IMS estimate.

\subsection{Uniform spin-graph geometry}
\label{sec:newS7V31-geometry}

Work in the fixed half-density space, where V26 removes all square-root
site-weight ratios.  Assume the global spin graph of V30 has degree at most
\(d_*\) and
\begin{equation}
 \kappa_e\le{K_*\over a},\qquad
 \lambda_e\le{L_*\over a^2},\qquad
 \|c_e\|\le C_c.                                       \label{eq:newS7V31-edge-bounds}
\end{equation}
The edge transports are unitary.  For a scalar site function \(\chi\), let
\begin{equation}
 \operatorname{Lip}_a(\chi)
 :=\max_{e=\{x,y\}}{|\chi_x-\chi_y|\over a}.            \label{eq:newS7V31-Lip}
\end{equation}

\begin{lemma}[Exact Wilson commutator]
\label{lem:newS7V31-commutator-formula}
Let
\(\widehat H_{G,a}=\Gamma(\widehat D_{W,G,a}-\rho/a)\) be the global
half-density Wilson sign kernel.  Then
\begin{align}
 ([\widehat K_{G,a},\chi]\psi)_x
 &=\sum_{y\sim x}\kappa_eJ_{xy}(\chi_y-\chi_x)\psi_y,
                                                               \label{eq:newS7V31-Kcomm}\
 ([\widehat L_{G,a},\chi]\psi)_x
 &=-\sum_{y\sim x}\lambda_eU_{xy}(\chi_y-\chi_x)\psi_y,
                                                               \label{eq:newS7V31-Lcomm}\
 [\widehat H_{G,a},\chi]
 &=\Gamma\left([\widehat K_{G,a},\chi]
       +{ra\over2}[\widehat L_{G,a},\chi]\right).      \label{eq:newS7V31-Hcomm}
\end{align}
\end{lemma}

\begin{proof}
The diagonal Wilson degree term and the mass shift \(\rho/a\) commute with
sitewise scalar multiplication.  Subtracting \(\chi\widehat K\) and
\(\chi\widehat L\) from the two hopping formulas in V30 gives the first two
identities; multiplication by chirality gives the third.
\end{proof}

\begin{theorem}[Cutoff-uniform single-commutator bound]
\label{thm:newS7V31-one-commutator}
Under \eqref{eq:newS7V31-edge-bounds},
\begin{equation}
 \|[\widehat H_{G,a},\chi]\|
 \le d_*\left(K_*C_c+{rL_*\over2}\right)
 \operatorname{Lip}_a(\chi).                            \label{eq:newS7V31-one-bound}
\end{equation}
In particular the bound contains no negative power of \(a\).
\end{theorem}

\begin{proof}
For an edge, \eqref{eq:newS7V31-edge-bounds} and
\(|\chi_y-\chi_x|\le a\operatorname{Lip}_a(\chi)\) bound the \(K\) block by
\(K_*C_c\operatorname{Lip}_a(\chi)\).  The Wilson block, including its
factor \(ra/2\), is bounded by
\(rL_*\operatorname{Lip}_a(\chi)/2\).  Every block row and column has at
most \(d_*\) entries.  The block Schur test proves
\eqref{eq:newS7V31-one-bound}.
\end{proof}

\subsection{A quadratic lattice partition}
\label{sec:newS7V31-partition}

Let \(\{U_i\}_{i=1}^{N_*}\) be a finite good cover and choose smooth
functions \(\chi_i\) with
\begin{equation}
 \sum_{i=1}^{N_*}\chi_i^2=1,qquad
 \operatorname{supp}\chi_i\Subset U_i.                 \label{eq:newS7V31-quadratic-POU}
\end{equation}
Sample them at the mesh sites.  Renormalizing the sampled vector by
\((\sum_i\chi_i(x)^2)^{-1/2}\), if necessary, preserves
\eqref{eq:newS7V31-quadratic-POU} exactly and changes the edge Lipschitz
constant by a cover-dependent factor.

Define
\begin{equation}
 \Lambda_\chi^2
 :=\max_{e=\{x,y\}}{1\over a^2}
 \sum_{i=1}^{N_*}|\chi_i(y)-\chi_i(x)|^2.               \label{eq:newS7V31-Lambda}
\end{equation}

\begin{lemma}[Uniform quadratic-partition Lipschitz bound]
\label{lem:newS7V31-partition-bound}
For all sufficiently small \(a\), \(\Lambda_\chi\le C_{\mathcal U}\), where
\(C_{\mathcal U}\) depends only on the fixed good cover, its subordinate
partition, and the quasi-uniform edge-length constant.
\end{lemma}

\begin{proof}
The mean-value theorem bounds each edge difference by the edge length times
\(\|d\chi_i\|_\infty\).  Quasi-uniformity bounds the edge length by \(Ca\).
Only finitely many functions occur.  The sampled normalization has a
denominator bounded away from zero and bounded first derivative, so the
same estimate holds after normalization.
\end{proof}

\begin{theorem}[Direct-sum commutator row]
\label{thm:newS7V31-direct-sum}
There is a constant
\begin{equation}
 C_\chi\le
 d_*\left(K_*C_c+{rL_*\over2}\right)C_{\mathcal U}\sqrt{N_*}             \label{eq:newS7V31-Cchi}
\end{equation}
such that
\begin{equation}
 \left(\sum_i\|[\widehat H_{G,a},\chi_i]\psi\|^2\right)^{1/2}
 \le C_\chi\|\psi\|.                                  \label{eq:newS7V31-direct-sum}
\end{equation}
The coarser factor \(\sqrt{N_*}\) may be removed by an edgewise Schur
estimate using \eqref{eq:newS7V31-Lambda}, but is harmless for the gap.
\end{theorem}

\begin{proof}
Apply \Cref{thm:newS7V31-one-commutator} to every \(\chi_i\), square, and
sum.  Finiteness of the cover and
\Cref{lem:newS7V31-partition-bound} give the displayed constant.
\end{proof}

\subsection{Exact collar replacement}
\label{sec:newS7V31-collar}

For each \(i\), let \(\widetilde H_{i,a}\) be a self-adjoint local closing
model on an enlarged chart \(\widetilde U_i\Subset U_i\).  It may use a
periodic closing collar as in V25.  Say that it is \emph{one-star compatible}
with the global kernel on \(\operatorname{supp}\chi_i\) if every site and
edge incident to that support has the same fibre map, Clifford hop, edge
coefficient, Wilson conductance, and mass shift in the two operators.

\begin{lemma}[Nearest-neighbour collar identity]
\label{lem:newS7V31-collar-identity}
If the local model is one-star compatible, then for every global spinor
\(\psi\), after extending the local output by zero in the common fibres,
\begin{equation}
 \widehat H_{G,a}(\chi_i\psi)
 =\widetilde H_{i,a}(\chi_i\psi).                        \label{eq:newS7V31-collar-identity}
\end{equation}
\end{lemma}

\begin{proof}
The vector \(\chi_i\psi\) vanishes outside its support.  A nearest-neighbour
operator can obtain a nonzero term from this vector only at a site in the
support or at a site joined to it by one edge.  One-star compatibility makes
every such block identical.  All remaining blocks multiply zero.
\end{proof}

\begin{definition}[Seam-compatible local gap row, SLGR]
\label{def:newS7V31-SLGR}
SLGR asserts that every \(\widetilde H_{i,a}\) is one-star compatible and
obeys, with one \(g_0>0\) independent of \(i,a\),
\begin{equation}
 \|\widetilde H_{i,a}\phi\|\ge{g_0\over a}\|\phi\|    \label{eq:newS7V31-local-gap}
\end{equation}
for every local spinor \(\phi\).  For chart interiors that are the V25
midpoint coordinate stencil, the perturbative gap theorem supplies this
row.  SLGR additionally requires it at mesh seams.
\end{definition}

\begin{corollary}[SLGR gives the localized global gap]
\label{cor:newS7V31-localized-gap}
Under SLGR,
\begin{equation}
 \|\widehat H_{G,a}\chi_i\psi\|
 \ge{g_0\over a}\|\chi_i\psi\|.                        \label{eq:newS7V31-localized-gap}
\end{equation}
\end{corollary}

\begin{proof}
Use \Cref{lem:newS7V31-collar-identity} and then
\eqref{eq:newS7V31-local-gap}.
\end{proof}

\subsection{The whole-operator gap}
\label{sec:newS7V31-global-gap}

\begin{theorem}[IMS globalization of the Wilson sign gap]
\label{thm:newS7V31-global-gap}
Assume the uniform graph bounds \eqref{eq:newS7V31-edge-bounds}, the fixed
quadratic partition \eqref{eq:newS7V31-quadratic-POU}, and SLGR.  Then
\begin{equation}
 \|\widehat H_{G,a}\psi\|
 \ge\left({g_0\over a}-C_\chi\right)\|\psi\|.          \label{eq:newS7V31-global-gap}
\end{equation}
In particular, if \(a\le g_0/(2C_\chi)\),
\begin{equation}
 \|\widehat H_{G,a}\psi\|\ge{g_0\over2a}\|\psi\|.    \label{eq:newS7V31-global-half-gap}
\end{equation}
\end{theorem}

\begin{proof}
By \Cref{cor:newS7V31-localized-gap} and the quadratic partition,
\[
 {g_0\over a}\|\psi\|
 \le\left(\sum_i\|\widehat H_{G,a}\chi_i\psi\|^2\right)^{1/2}.
\]
Use \(\widehat H\chi_i=\chi_i\widehat H+[\widehat H,\chi_i]\), the triangle
inequality in the direct sum, and
\[
 \sum_i\|\chi_i\widehat H\psi\|^2=\|\widehat H\psi\|^2.
\]
The second direct-sum term is at most \(C_\chi\|\psi\|\) by
\Cref{thm:newS7V31-direct-sum}.  Rearrangement gives the first estimate;
the stated restriction on \(a\) gives the second.
\end{proof}

\begin{corollary}[Global overlap algebra and smooth sign contacts]
\label{cor:newS7V31-global-overlap}
Under the hypotheses of \Cref{thm:newS7V31-global-gap}, the global overlap
operator
\begin{equation}
 D_{\mathrm{ov},G,a}={\rho\over a}
 \{I+\Gamma\operatorname{sign}(\widehat H_{G,a})\}      \label{eq:newS7V31-global-overlap}
\end{equation}
exists, is globally spin--gauge covariant, and obeys the weighted
Ginsparg--Wilson identity.  On a closed parameter subfamily preserving
SLGR and the graph bounds, its first and second metric contacts are the V26
fixed-contour formulas with a contour uniform in \(a\).
\end{corollary}

\begin{proof}
The global gap permits sign functional calculus and supplies a uniform gap
for the rescaled kernel \(a\widehat H_{G,a}\).  Apply the V25 overlap algebra
and V26 contour differentiation to the global operator.  The intrinsic
spin-cocycle theorem of V30 supplies covariance.
\end{proof}

\begin{assessment}[The exact remaining seam gate]
\label{ass:newS7V31-seam}
The commutator part of the V30 IMS row is now unconditional under standard
bounded-degree mesh scaling, and one-star compatibility makes collar
replacement exact.  The remaining gap input is SLGR at irregular mesh
seams.  Local coordinate stencils satisfy it by V25; arbitrary
triangulation stars do not inherit it from spin covariance or from the
moment conditions alone.  A seam construction must either be a small
operator-norm perturbation of one of the gapped coordinate models or carry
an independent Wilson-symbol gap proof.
\end{assessment}

\begin{programme}[V32: seam stencil by norm-controlled interpolation]
\label{prog:newS7V31-V32}
V32 will test a collar interpolation between two chart-coordinate Wilson
stencils after transporting both to the same intrinsic spin fibres.  It
will preserve the skew/self-adjoint edge pairing exactly, estimate the
interpolated sign kernel against one endpoint in operator norm, and determine
whether a partition transition of fixed width keeps the perturbation below
\(g_0/(2a)\).  If the two stencils are not norm-close at scale \(a^{-1}\),
the programme will return upstream to a single global cubulation rather than
average incompatible local operators.
\end{programme}

\begin{firewall}[Claim boundary after seventh-series V31]
\label{fire:newS7V31-boundary}
V31 derives the exact global Wilson cutoff commutator, proves its cutoff-
uniform norm bound, constructs a sampled quadratic partition with a uniform
edge Lipschitz constant, proves exact nearest-neighbour collar replacement,
formulates SLGR, and proves that SLGR gives the whole-operator gap
\(g_0/a-C_\chi\) and hence a global covariant overlap sign with uniform
metric contacts.

V31 does not prove SLGR at irregular seams, construct the GSMC cofinal mesh
moments, complete global heat locality, construct the common Einstein
stress, transfer pure gravitational owners, control rough metric or AOP
fields, prove measure concentration, construct the determinant-line phase,
remove the interacting cutoff, or construct the Standard Model--Einstein
continuum.
\end{firewall}

\begin{theorem}[Seventh-series V31 result]
\label{thm:newS7V31-result}
For every bounded-degree global spin graph with the natural
\(a^{-1}\) Dirac and \(a^{-2}\) Wilson edge scaling, a fixed smooth quadratic
partition satisfies
\[
 \left(\sum_i\|[\widehat H_{G,a},\chi_i]\psi\|^2\right)^{1/2}
 \le C_\chi\|\psi\|.
\]
If every supported one-star has a compatible local closing model with gap
\(g_0/a\), then the global kernel has gap at least \(g_0/a-C_\chi\), and at
least \(g_0/(2a)\) for all sufficiently small \(a\).  The sole remaining
sign obstruction is construction of those gapped seam stencils.
\end{theorem}

\begin{proof}
Use \Cref{lem:newS7V31-commutator-formula,
thm:newS7V31-one-commutator,
lem:newS7V31-partition-bound,
thm:newS7V31-direct-sum,
lem:newS7V31-collar-identity,
cor:newS7V31-localized-gap,
thm:newS7V31-global-gap,
cor:newS7V31-global-overlap}.
\end{proof}

\paragraph{Parent-version record.}
V31 was formed from
\texttt{Renewal\_SM\_Einstein\_Continuum\_Research\_Notes\_V30.tex}.
The V30 parent had (12\,439) logical source lines, (498\,420) bytes,
and SHA--256
\[
 \texttt{0936C37A0195B57D9CD53CE381D4D0471584C40DA3B9897728CDB4D27B96C2E4}.
\]
The next compulsory companion audit is V35.

% =============================================================================
% END APPENDED SEVENTH-SERIES V31 MATERIAL
% =============================================================================


% =============================================================================
% BEGIN APPENDED SEVENTH-SERIES V32 MATERIAL
% =============================================================================

\clearpage
\section{Seventh-series V32: failure of partition interpolation and reduction
to a finite global star catalogue}
\label{sec:v7v32}

V31 proposes testing an interpolation between chart-coordinate Wilson
stencils.  V32 proves the exact stability criterion for such an
interpolation and shows why increasing the physical transition width does
not enforce it: two stencils may differ by order \(a^{-1}\) even when the
mixing function varies by only \(O(a)\) per edge.  A two-dimensional
counterexample shows that convex interpolation can close a gap between
unitarily equivalent endpoints.  The programme therefore moves upstream to
one global combinatorial mesh.  For cofinal subdivisions of a fixed mesh,
the seam problem is reduced to checking a finite catalogue of frozen star
types.

\subsection{What an algebra-preserving interpolation can be}
\label{sec:newS7V32-interpolation}

Let \(H_a^{(0)}\) and \(H_a^{(1)}\) be two self-adjoint Wilson sign kernels
represented on the same intrinsic spin fibres and on the union of their edge
sets, with a missing edge represented by a zero hopping block.  For a
constant \(0\le\theta\le1\), the affine interpolation
\begin{equation}
 H_a^{(\theta)}=(1-\theta)H_a^{(0)}+\theta H_a^{(1)}    \label{eq:newS7V32-affine}
\end{equation}
is self-adjoint.  At the Dirac/Wilson level, positivity of the Wilson part is
preserved more transparently by interpolating its quadratic edge forms:
\begin{equation}
 q_\theta(\psi)
 =(1-\theta)\langle\psi,L_a^{(0)}\psi\rangle
 +\theta\langle\psi,L_a^{(1)}\psi\rangle\ge0.          \label{eq:newS7V32-form-interpolation}
\end{equation}
Skew-adjoint Dirac edge pairs may be interpolated affinely.  Thus algebraic
Hermiticity and Wilson positivity do not themselves prevent interpolation.

\begin{theorem}[Sharp norm condition for a safe affine path]
\label{thm:newS7V32-safe-path}
Suppose
\begin{equation}
 \|H_a^{(0)}\psi\|\ge{g_0\over a}\|\psi\|,qquad
 \|H_a^{(1)}-H_a^{(0)}\|\le{\eta\over a}.              \label{eq:newS7V32-closeness}
\end{equation}
Then for every \(\theta\in[0,1]\),
\begin{equation}
 \|H_a^{(\theta)}\psi\|
 \ge{g_0-\theta\eta\over a}\|\psi\|.                 \label{eq:newS7V32-path-gap}
\end{equation}
In particular the whole path is gapped if \(\eta<g_0\).
\end{theorem}

\begin{proof}
Write
\(H_a^{(\theta)}=H_a^{(0)}+\theta(H_a^{(1)}-H_a^{(0)})\) and apply the
reverse triangle inequality and \eqref{eq:newS7V32-closeness}.
\end{proof}

\begin{proposition}[A wide collar does not provide dimensionless stencil
closeness]
\label{prop:newS7V32-width-no-help}
Let a sitewise mixing function vary on a fixed physical width \(R\), so that
\(|\theta_y-\theta_x|\le Ca/R\) on every edge.  For two bounded-degree
Wilson stencils whose dimensionless hopping coefficients differ by
\(\Delta_*\), the interpolated kernel obeys an estimate of the form
\begin{equation}
 \|H_{a,\theta}-H_a^{(0)}\|
 \le {C_1\sup_x|\theta_x|\Delta_*\over a}
      +{C_2\over R}.                                    \label{eq:newS7V32-width-bound}
\end{equation}
The second term comes from variation of the mixing function, but the first
does not decrease when \(R\) increases.  Therefore collar width alone cannot
verify \eqref{eq:newS7V32-closeness}.
\end{proposition}

\begin{proof}
After half-density conjugation, every Dirac and Wilson hopping block has
scale \(a^{-1}\).  Split the difference into the value of \(\theta\) times
the difference of the two edge blocks and a commutator in which an edge
difference \(\theta_y-\theta_x\) occurs.  The block Schur estimate gives the
first term.  The second contains \(a^{-1}|\theta_y-\theta_x|\le C/R\).
Bounded degree absorbs the number of incident edges.
\end{proof}

\begin{theorem}[Unitary-equivalent endpoints can have a gapless affine
midpoint]
\label{thm:newS7V32-affine-no-go}
There is no lower-gap theorem for \eqref{eq:newS7V32-affine} based only on
equal endpoint gaps, equal inertia, or unitary equivalence.  Indeed, set
\begin{equation}
 H_a^{(0)}={1\over a}\begin{pmatrix}1&0\\0&-1\end{pmatrix},qquad
 H_a^{(1)}={1\over a}\begin{pmatrix}-1&0\\0&1\end{pmatrix}.              \label{eq:newS7V32-two-by-two}
\end{equation}
Both endpoints have gap \(a^{-1}\), have one positive and one negative
eigenvalue, and are unitarily conjugate, but
\(H_a^{(1/2)}=0\).
\end{theorem}

\begin{proof}
Conjugation by the matrix exchanging the two coordinate vectors maps the
first endpoint to the second.  Their arithmetic mean is the zero matrix.
\end{proof}

\begin{corollary}[Decision on chart-stencil averaging]
\label{cor:newS7V32-no-average}
Two chart stencils may be averaged only after the dimensionless estimate
\(a\|H_a^{(1)}-H_a^{(0)}\|<g_0\) is proved.  Smooth chart transitions and a
wide partition collar do not imply this estimate.  In the absence of that
proof, averaging is not an admissible construction of SLGR.
\end{corollary}

\begin{proof}
Sufficiency is \Cref{thm:newS7V32-safe-path}; insufficiency of the proposed
surrogates follows from
\Cref{prop:newS7V32-width-no-help,thm:newS7V32-affine-no-go}.
\end{proof}

\subsection{The frozen anisotropic coordinate gap}
\label{sec:newS7V32-anisotropic}

The no-go concerns different combinatorial stencils.  A change of smooth
coframe on one fixed coordinate stencil is harmless at the frozen-symbol
level and need not be close to the identity.

\begin{theorem}[Uniform anisotropic Wilson gap]
\label{thm:newS7V32-anisotropic-gap}
Let \(E=(E_a{}^\mu)\in\operatorname{GL}(4,\mathbb R)\) range in a compact
uniformly elliptic set
\begin{equation}
 \sigma_-|v|\le|Ev|\le\sigma_+|v|,qquad \sigma_->0.     \label{eq:newS7V32-elliptic-E}
\end{equation}
On the fixed hypercubic stencil define the frozen symbol
\begin{equation}
 aH_E(k)=\Gamma\left{
 i\gamma^aE_a{}^\mu\sin k_\mu
 +r\sum_\mu(1-\cos k_\mu)-\rho\right}.                \label{eq:newS7V32-anisotropic-symbol}
\end{equation}
If \(0<\rho<2r\), there is a number
\(\delta_{\rm an}=\delta_{\rm an}(\rho,r,\sigma_-,\sigma_+)>0\) such that
\begin{equation}
 \|H_E(k)v\|\ge{\delta_{\rm an}\over a}\|v\|          \label{eq:newS7V32-anisotropic-gap}
\end{equation}
for every \(k\in[-\pi,\pi]^4\) and every allowed \(E\).
\end{theorem}

\begin{proof}
Clifford multiplication gives
\begin{equation}
 a^2H_E(k)^2
 =|E\sin k|^2+
 \left\{r\sum_\mu(1-\cos k_\mu)-\rho\right\}^2.       \label{eq:newS7V32-anisotropic-square}
\end{equation}
A zero requires \(E\sin k=0\).  Invertibility of \(E\) gives
\(\sin k_\mu=0\) for all \(\mu\), so each component is zero or \(\pi\).
If \(n\) components are \(\pi\), the second term vanishes only if
\(2rn=\rho\), impossible for \(0<\rho<2r\).  The square is therefore
positive.  Its minimum over the compact product of the Brillouin zone and
the declared compact coframe set is strictly positive; take its square root.
\end{proof}

\begin{corollary}[Large frame changes are not seam defects]
\label{cor:newS7V32-frame-not-seam}
On one combinatorial coordinate stencil, a uniformly elliptic frozen
coframe can be handled directly by \Cref{thm:newS7V32-anisotropic-gap}; it
need not be interpolated with the Euclidean coframe.  Spin-frame changes are
already exact conjugations by V30.  A genuine seam defect is a change in the
edge incidence or normalized difference stencil.
\end{corollary}

\begin{proof}
The anisotropic theorem treats the coframe matrix without a smallness
assumption.  V30 treats spin-frame representatives by conjugation.  Neither
operation changes the underlying edge set.
\end{proof}

\subsection{One mesh and a finite catalogue of star types}
\label{sec:newS7V32-star-catalogue}

Let \(\mathcal K\) be one finite global cellular decomposition of \(M\), and
let \(G_a\) be obtained by a declared uniform subdivision rule.  A
\emph{normalized star type} records the incidence pattern and the
dimensionless Dirac and Wilson edge coefficients after rescaling lengths by
\(a^{-1}\) and freezing the coframe and connection.

\begin{definition}[Finite star gap card, FSGC]
\label{def:newS7V32-FSGC}
FSGC consists of:
\begin{enumerate}[label=\textup{(\roman*)},leftmargin=3.5em]
\item a fixed global cell complex and a cofinal subdivision rule whose
      one-star neighbourhoods belong to a finite set
      \(\mathfrak T=\{\tau_1,\ldots,\tau_N\}\) independent of \(a\);
\item for every \(\tau\in\mathfrak T\), a self-adjoint finite or periodic
      frozen closing kernel \(H_{\tau,a}^{\rm fr}\) which agrees with that
      star and has
      \begin{equation}
       \|H_{\tau,a}^{\rm fr}\phi\|
       \ge{\delta_\tau\over a}\|\phi\|,qquad
       \delta_\tau>0;                                  \label{eq:newS7V32-star-gap}
      \end{equation}
\item one-star geometric coefficients which, on a physical collar of radius
      \(R\), differ from the appropriate frozen type by
      \begin{equation}
       a\|H_{i,a}-H_{\tau(i),a}^{\rm fr}\|
       \le\omega(R)+C{a\over R},qquad \omega(R)\to0;   \label{eq:newS7V32-star-modulus}
      \end{equation}
      and
\item exact intrinsic spin-fibre identification with the global operator on
      the supported one-star, as in V31.
\end{enumerate}
Put \(\delta_*:=\min_{\tau\in\mathfrak T}\delta_\tau>0\).
\end{definition}

\begin{theorem}[FSGC proves the seam-compatible local gap row]
\label{thm:newS7V32-FSGC-SLGR}
Assume FSGC.  Choose \(R\) so that
\(\omega(R)<\delta_*/4\), and then \(a\) so that
\(Ca/R<\delta_*/4\).  Every local closing model then satisfies
\begin{equation}
 \|H_{i,a}\phi\|\ge{\delta_*\over2a}\|\phi\|,          \label{eq:newS7V32-SLGR-gap}
\end{equation}
uniformly through regular chart interiors and mesh seams.  Thus SLGR holds
with \(g_0=\delta_*/2\).
\end{theorem}

\begin{proof}
For the type \(\tau(i)\), use the reverse triangle inequality,
\eqref{eq:newS7V32-star-gap}, and
\eqref{eq:newS7V32-star-modulus}:
\[
 \|H_{i,a}\phi\|
 \ge {\delta_*-\omega(R)-Ca/R\over a}\|\phi\|
 \ge{\delta_*\over2a}\|\phi\|.
\]
The intrinsic identification in FSGC(iv) is precisely the one-star
compatibility required by V31.
\end{proof}

\begin{corollary}[FSGC gives the whole-operator overlap gap]
\label{cor:newS7V32-FSGC-global}
Under the V31 bounded-degree hypotheses and FSGC, for all sufficiently small
\(a\),
\begin{equation}
 \|H_{G,a}\psi\|
 \ge\left({\delta_*\over2a}-C_\chi\right)\|\psi\|
 \ge{\delta_*\over4a}\|\psi\|.                         \label{eq:newS7V32-global-gap}
\end{equation}
Hence the global overlap sign and its V26 metric contacts exist.
\end{corollary}

\begin{proof}
Apply \Cref{thm:newS7V32-FSGC-SLGR} and then the V31 IMS theorem, decreasing
\(a\) once more so that \(C_\chi\le\delta_*/(4a)\).
\end{proof}

\begin{assessment}[What is now finite and what is not yet proved]
\label{ass:newS7V32-status}
The seam problem no longer calls for an uncontrolled partition interpolation.
For subdivisions of one fixed global mesh it has two finite inputs: enumerate
the normalized star catalogue and prove the frozen gap of each member.  The
smooth geometric perturbation and global IMS gluing are then proved by V31--
V32.  The current manuscripts do not provide a subdivision rule with the
FSGC moment consistency and frozen gaps, so FSGC itself remains open.
\end{assessment}

\begin{programme}[V33: acquire a concrete subdivision and its star symbols]
\label{prog:newS7V32-V33}
V33 will begin from a smooth finite triangulation or cubulation of the fixed
spin manifold and select one explicit uniform subdivision.  It will derive
the dimensionless Dirac first moments and positive Wilson conductances for
each normalized vertex-star type.  For every translation-periodic regular
type it will apply the anisotropic theorem above; irregular inherited
vertices will be listed separately with their finite frozen matrices.  If
the chosen subdivision generates infinitely many shapes or a nonpositive
Wilson form, it will be rejected before any heat calculation.
\end{programme}

\begin{firewall}[Claim boundary after seventh-series V32]
\label{fire:newS7V32-boundary}
V32 proves the exact norm criterion for safe stencil interpolation, proves
that transition width cannot supply the required dimensionless closeness,
gives a gap-closing counterexample with unitarily equivalent endpoints,
proves a uniform frozen Wilson gap for arbitrary uniformly elliptic coframes
on one coordinate stencil, defines FSGC, and proves FSGC sufficient for SLGR
and the global overlap gap.

V32 does not construct or verify FSGC, choose a valid global subdivision,
prove the GSMC moment rows or heat locality, construct the common Einstein
stress, transfer pure gravitational owners, control rough metric or AOP
fields, prove measure concentration, construct the determinant-line phase,
remove the interacting cutoff, or construct the Standard Model--Einstein
continuum.
\end{firewall}

\begin{theorem}[Seventh-series V32 result]
\label{thm:newS7V32-result}
A partition interpolation between unrelated \(a^{-1}\)-scale chart stencils
is not gap-safe; even unitarily equivalent gapped endpoints can have a zero
affine midpoint.  On a fixed stencil, however, every uniformly elliptic
frozen coframe has a uniform Wilson gap.  For cofinal subdivisions of one
global mesh, the complete seam and whole-operator gap problem reduces to the
finite star card FSGC.  Once its finitely many frozen gaps and geometric
modulus are established, the global gap is at least
\(\delta_*/(4a)\) for all sufficiently small \(a\).
\end{theorem}

\begin{proof}
Use \Cref{thm:newS7V32-safe-path,
prop:newS7V32-width-no-help,
thm:newS7V32-affine-no-go,
cor:newS7V32-no-average,
thm:newS7V32-anisotropic-gap,
cor:newS7V32-frame-not-seam,
thm:newS7V32-FSGC-SLGR,
cor:newS7V32-FSGC-global}.
\end{proof}

\paragraph{Parent-version record.}
V32 was formed from
\texttt{Renewal\_SM\_Einstein\_Continuum\_Research\_Notes\_V31.tex}.
The V31 parent had (12\,766) logical source lines, (511\,737) bytes,
and SHA--256
\[
 \texttt{B7815FC841BB4768D6D84D5C01645249D15BB0C00BD851FE296B5E290FC0A775}.
\]
The next compulsory companion audit is V35.

% =============================================================================
% END APPENDED SEVENTH-SERIES V32 MATERIAL
% =============================================================================


% =============================================================================
% BEGIN APPENDED SEVENTH-SERIES V33 MATERIAL
% =============================================================================

\clearpage
\section{Seventh-series V33: a concrete global simplicial spin mesh and its
Wilson forms}
\label{sec:v7v33}

V32 reduces seam control to a finite catalogue only after one global
subdivision rule is fixed.  V33 fixes such a rule and defines the associated
spin operator without choosing incompatible chart lattices.  The Dirac part
is the mass-lumped Galerkin form of the intrinsic geometric Dirac operator;
the Wilson part is the positive connection-stiffness form.  Their locality,
spin--gauge covariance, adjoint algebra, cutoff scaling, quadrature, and
first consistency moments are proved.  What remains is a finite family of
frozen star-gap problems, rather than an unspecified seam interpolation.

\subsection{Face-compatible edgewise subdivision}
\label{sec:newS7V33-subdivision}

Fix a finite smooth ordered triangulation \(\mathcal K_0\) of the compact
spin four-manifold \(M\), fine enough that every simplex lies in a convex
normal neighbourhood.  In the barycentric coordinates
\((t_0,\ldots,t_4)\) of each ordered four-simplex, subdivide at level \(n\)
using the lattice points
\begin{equation}
 t_j={\alpha_j\over n},\qquad
 \alpha_j\in\mathbb N_0,\qquad \sum_{j=0}^4\alpha_j=n,  \label{eq:newS7V33-barycentric-points}
\end{equation}
and the standard edgewise/Freudenthal ordering of the resulting small
simplices.  Use the induced rule on every face.  Put \(a=n^{-1}\) after a
fixed reference-length normalization, and denote the resulting complex by
\(\mathcal K_a\) and its vertices by \(X_a\).

\begin{lemma}[Global compatibility and shape regularity]
\label{lem:newS7V33-mesh}
The induced subdivisions of two original simplices agree on every common
face.  The family \(\{\mathcal K_a\}\) is cofinal, has uniformly bounded
vertex degree, and is uniformly shape regular after the fixed finite set of
original simplex charts is included in the constants.
\end{lemma}

\begin{proof}
The barycentric lattice and its ordering restrict to the same construction
when one barycentric coordinate is zero; hence the two face subdivisions
coincide.  In a reference simplex, every small simplex is an affine image,
up to a permutation from a finite list, of one of finitely many standard
edgewise cells scaled by \(n^{-1}\).  Its aspect ratio is therefore bounded
independently of \(n\).  The original triangulation is finite and its smooth
simplex maps have uniformly bounded derivatives and inverses, which
preserves shape regularity for sufficiently fine subdivisions.  Each
vertex meets a bounded number of reference small cells, so its degree is
uniformly bounded.  Their diameters tend to zero with \(a\).
\end{proof}

\begin{lemma}[Finite normalized combinatorial star catalogue]
\label{lem:newS7V33-finite-stars}
There is a finite list \(\mathfrak T_{\rm ed}\), independent of \(a\), such
that every closed one-star of \(\mathcal K_a\), after rescaling barycentric
integer coordinates and recording the incident original faces, is
isomorphic to one member of \(\mathfrak T_{\rm ed}\).
\end{lemma}

\begin{proof}
Inside an original simplex the edgewise pattern is translation-periodic in
the barycentric integer lattice, up to one of finitely many coordinate
orderings.  A radius-one star can distinguish an original face only by
which barycentric coordinates vanish at its centre or at an adjacent
vertex; there are finitely many such subsets.  Along an original face, the
gluing is determined by the finite incidence data and vertex orderings of
\(\mathcal K_0\).  Combining the finite interior patterns, zero-coordinate
patterns, and original face incidences gives a finite list.  No distance in
numbers of cells beyond one enters a one-star.
\end{proof}

\subsection{Intrinsic nodal spin interpolation}
\label{sec:newS7V33-interpolation}

Let \(\phi_x\) be the continuous scalar piecewise-affine hat function of
the vertex \(x\).  For a nodal spinor
\(\psi=(\psi_x)_{x\in X_a}\), define on a small simplex \(\sigma\)
\begin{equation}
 (I_a\psi)(z)
 :=\sum_{x\in\operatorname{Vert}(\sigma)}
 \phi_x(z)P_{z\leftarrow x}\psi_x,                     \label{eq:newS7V33-spin-interpolant}
\end{equation}
where \(P_{z\leftarrow x}\) is spin--gauge parallel transport along the
unique short geodesic in the containing convex normal neighbourhood.

\begin{lemma}[Well-defined spin interpolation]
\label{lem:newS7V33-interpolant}
The traces of \eqref{eq:newS7V33-spin-interpolant} from adjacent simplices
agree on their common face.  Thus \(I_a\psi\) is a global continuous
piecewise-\(H^1\) spinor.  It commutes with chirality and is exactly covariant
under changes of spin frame and Standard Model gauge frame.
\end{lemma}

\begin{proof}
On a common face, hat functions belonging to vertices outside that face
vanish and the remaining scalar hat functions agree.  The short geodesic
from a face vertex to \(z\) and its intrinsic parallel transport do not
depend on which adjacent simplex is used.  Hence the two traces coincide.
Parallel transport preserves chirality.  Under a frame change, its endpoint
matrices conjugate, so every summand transforms as a spinor at \(z\).
\end{proof}

Define the positive lumped weights
\begin{equation}
 w_x(g):=\int_M\phi_x\,d\mu_g>0,qquad
 \langle\psi,\eta\rangle_{a,w}
 :=\sum_xw_x\langle\psi_x,\eta_x\rangle_{S_x}.          \label{eq:newS7V33-lumped}
\end{equation}

\begin{lemma}[Quadrature and metric contacts]
\label{lem:newS7V33-quadrature}
For \(f\in C^2(M)\),
\begin{equation}
 \left|\sum_xw_xf(x)-\int_Mf\,d\mu_g\right|
 \le Ca^2\|f\|_{C^2}.                                  \label{eq:newS7V33-quadrature}
\end{equation}
On a bounded \(C^2\) metric family the estimate is uniform after one or two
metric parameter derivatives.  The weight derivatives are
\begin{align}
 \partial_hw_x
 &=\frac12\int_M\phi_x\tr_g(h)\,d\mu_g,                \label{eq:newS7V33-weight1}\
 \partial_{hk}^2w_x
 &=\int_M\phi_x\left\{
 {1\over4}\tr_g(h)\tr_g(k)-{1\over2}\tr(g^{-1}hg^{-1}k)
 \right\}d\mu_g.                                      \label{eq:newS7V33-weight2}
\end{align}
\end{lemma}

\begin{proof}
Because \(\sum_x\phi_x=1\), the lumped sum is
\(\int_M I_af\,d\mu_g\).  The standard affine interpolation remainder on
a shape-regular simplex is bounded by \(Ca^2\|f\|_{C^2}\); sum over the
mesh.  Metric differentiation of the finite sum commutes with integration,
and the bounded-family estimates are unchanged.  The two displayed
formulas are the first and mixed second variations of the Riemannian volume
density.
\end{proof}

\subsection{The global Galerkin Dirac and Wilson forms}
\label{sec:newS7V33-forms}

Use the skew-adjoint Euclidean convention for the geometric Dirac operator
\(\mathcal D_{g,A}\).  Define sesquilinear forms on nodal spinors by
\begin{align}
 \mathfrak d_a(\eta,\psi)
 &:=\sum_{\sigma\in\mathcal K_a}
 \int_\sigma\langle I_a\eta,\mathcal D_{g,A}I_a\psi\rangle\,d\mu_g,
                                                               \label{eq:newS7V33-Dirac-form}\
 \mathfrak q_a(\eta,\psi)
 &:=\sum_{\sigma\in\mathcal K_a}
 \int_\sigma\langle\nabla I_a\eta,\nabla I_a\psi\rangle\,d\mu_g. \label{eq:newS7V33-stiffness-form}
\end{align}
Let \(K_a^{\rm FE}\) and \(L_a^{\rm FE}\) be the unique operators satisfying
\begin{equation}
 \langle\eta,K_a^{\rm FE}\psi\rangle_{a,w}=\mathfrak d_a(\eta,\psi),
 \qquad
 \langle\eta,L_a^{\rm FE}\psi\rangle_{a,w}=\mathfrak q_a(\eta,\psi),     \label{eq:newS7V33-operators}
\end{equation}
and put
\begin{equation}
 D_{W,a}^{\rm FE}:=K_a^{\rm FE}+{ra\over2}L_a^{\rm FE}.                 \label{eq:newS7V33-FE-Wilson}
\end{equation}

\begin{theorem}[Exact global Wilson algebra]
\label{thm:newS7V33-algebra}
On the lumped Hilbert space,
\begin{equation}
 (K_a^{\rm FE})^*=-K_a^{\rm FE},\qquad
 L_a^{\rm FE}=(L_a^{\rm FE})^*\ge0,                    \label{eq:newS7V33-adjoints}
\end{equation}
and
\begin{equation}
 K_a^{\rm FE}\Gamma=-\Gamma K_a^{\rm FE},\qquad
 L_a^{\rm FE}\Gamma=\Gamma L_a^{\rm FE},\qquad
 (D_{W,a}^{\rm FE})^*=\Gamma D_{W,a}^{\rm FE}\Gamma. \label{eq:newS7V33-chiral-algebra}
\end{equation}
The operators are exactly spin--gauge covariant and have matrix blocks only
between vertices belonging to a common small simplex.
\end{theorem}

\begin{proof}
The global interpolants are continuous and piecewise \(H^1\), so weak
integration by parts on the closed manifold has no uncancelled interior-face
term.  Skew-adjointness of \(\mathcal D_{g,A}\) gives
\(\mathfrak d_a(\eta,\psi)=-\overline{\mathfrak d_a(\psi,\eta)}\).
The stiffness form is Hermitian and nonnegative.  Riesz representation in
the positive lumped inner product proves \eqref{eq:newS7V33-adjoints}.

The interpolant commutes with chirality, the geometric Dirac anticommutes
with it, and the connection Laplacian form commutes with it.  This proves
\eqref{eq:newS7V33-chiral-algebra}.  Intrinsic covariance follows from
\Cref{lem:newS7V33-interpolant}.  Two nodal basis functions have disjoint
supports unless their vertices share a simplex, proving locality.
\end{proof}

\begin{theorem}[Uniform cutoff scaling]
\label{thm:newS7V33-scaling}
There is \(C\), uniform on bounded smooth metric--gauge families, such that
\begin{equation}
 \|K_a^{\rm FE}\|\le{C\over a},\qquad
 \|L_a^{\rm FE}\|\le{C\over a^2},\qquad
 \|D_{W,a}^{\rm FE}\|\le{C\over a}.                   \label{eq:newS7V33-scaling}
\end{equation}
The matrix degree is uniformly bounded.
\end{theorem}

\begin{proof}
Shape-regular finite-element inverse estimates give
\(\|\nabla I_a\psi\|_{L^2}\le Ca^{-1}\|I_a\psi\|_{L^2}\).
Mass-lumping equivalence gives
\(c\|\psi\|_{a,w}\le\|I_a\psi\|_{L^2}\le C\|\psi\|_{a,w}\).
The Dirac form contains one derivative and bounded connection coefficients;
the stiffness form contains two first derivatives.  Cauchy--Schwarz and
Riesz representation yield the first two estimates, and the third follows
from \eqref{eq:newS7V33-FE-Wilson}.  Bounded degree is
\Cref{lem:newS7V33-mesh} plus locality from
\Cref{thm:newS7V33-algebra}.
\end{proof}

\subsection{Moment consistency}
\label{sec:newS7V33-consistency}

Let \(S_a\psi=(\psi(x))_{x\in X_a}\) denote sampling of a smooth spinor.

\begin{lemma}[Interpolation consistency]
\label{lem:newS7V33-interpolation-error}
For every \(C^2\) spinor \(\psi\),
\begin{equation}
 \|I_aS_a\psi-\psi\|_{L^2}
 +a\|\nabla(I_aS_a\psi-\psi)\|_{L^2}
 \le Ca^2\|\psi\|_{H^2}.                               \label{eq:newS7V33-interpolation-error}
\end{equation}
The estimate is uniform through two smooth metric parameter derivatives if
the corresponding spin connection jets are bounded.
\end{lemma}

\begin{proof}
Trivialize by parallel transport in each convex normal simplex.  The
interpolant is the ordinary affine nodal interpolant plus an \(O(a^2)\)
curvature/connection remainder.  The Bramble--Hilbert affine interpolation
estimate on the finite reference-shape family gives the displayed \(L^2\)
and \(H^1\) bounds.  Differentiating the parallel transports uses the V26
Volterra contact formulas and preserves the estimate on a bounded family.
\end{proof}

\begin{theorem}[Weak Dirac and Wilson moment convergence]
\label{thm:newS7V33-moments}
For smooth spinors \(\eta,\psi\),
\begin{align}
 \left|\langle S_a\eta,K_a^{\rm FE}S_a\psi\rangle_{a,w}
 -\langle\eta,\mathcal D_{g,A}\psi\rangle_{L^2}\right|
 &\le Ca\|\eta\|_{H^1}\|\psi\|_{H^2},                \label{eq:newS7V33-Dirac-consistency}\
 \left|\langle S_a\eta,L_a^{\rm FE}S_a\psi\rangle_{a,w}
 -\langle\nabla\eta,\nabla\psi\rangle_{L^2}\right|
 &\le Ca\|\eta\|_{H^2}\|\psi\|_{H^2}.               \label{eq:newS7V33-Laplace-consistency}
\end{align}
The same statements hold after one or two metric variations with the
corresponding differentiated continuum forms.
\end{theorem}

\begin{proof}
By definition the two discrete pairings are the continuum forms evaluated
on \(I_aS_a\eta,I_aS_a\psi\).  Subtract the continuum forms at
\(\eta,\psi\), use \Cref{lem:newS7V33-interpolation-error}, and apply
Cauchy--Schwarz.  The Dirac form needs one \(H^1\) error and the stiffness
form two.  Metric differentiation acts on the volume, Clifford map,
connection, and parallel transports; the weight contacts are
\Cref{lem:newS7V33-quadrature} and the transport contacts are V26.  The
differentiated interpolation estimates give the final assertion.
\end{proof}

\begin{corollary}[The Wilson term vanishes on smooth modes]
\label{cor:newS7V33-Wilson-vanishes}
For sampled smooth spinors, the Wilson contribution satisfies weakly
\begin{equation}
 {ra\over2}\langle S_a\eta,L_a^{\rm FE}S_a\psi\rangle_{a,w}
 =O(a)                                                     \label{eq:newS7V33-Wilson-smooth}
\end{equation}
and therefore does not change the continuum Dirac principal or
subprincipal action on smooth modes.
\end{corollary}

\begin{proof}
Apply \eqref{eq:newS7V33-Laplace-consistency} and multiply by \(ra/2\).
\end{proof}

\subsection{The remaining stratified frozen gap problem}
\label{sec:newS7V33-frozen-gap}

The finite one-star catalogue in
\Cref{lem:newS7V33-finite-stars} is enough for locality, scaling, and moment
bookkeeping.  It is not, by itself, a spectral-gap certificate on a collar
whose physical radius is fixed while its number of vertices tends to
infinity.

\begin{proposition}[One-star gaps do not imply a collar gap]
\label{prop:newS7V33-star-not-enough}
Uniform invertibility of finitely many isolated one-star matrices does not
imply a uniform gap for their coupled union.  In particular, replacing the
SLGR local closing in V31 by independent singular-value tests of its vertex
stars is invalid.
\end{proposition}

\begin{proof}
The two-block example in
\Cref{prop:newS7V30-gap-no-go,thm:newS7V32-affine-no-go} already has
invertible isolated diagonal blocks and a singular coupled matrix.  A
physical collar contains order \((R/a)^4\) coupled stars, so a bound on each
isolated star omits precisely these off-diagonal couplings.  V31 requires a
gap for the whole closing operator acting on \(\chi_i\psi\), not for its
individual block rows.
\end{proof}

The original finite triangulation has finitely many open face strata.  Near
the interior of a stratum of dimension \(j\), the rescaled edgewise mesh is
translation-periodic in \(j\) tangential directions and has a fixed
polyhedral normal pattern.  At an original vertex there is no tangential
direction and the rescaled model is a finite collection of conical sectors.
The finite combinatorics in \Cref{lem:newS7V33-finite-stars} therefore lead
to finitely many \emph{stratified tangent operators}, not merely finitely
many matrices.

\begin{definition}[Stratified frozen gap card, SFGC]
\label{def:newS7V33-SFGC}
For every incidence type \(\tau\) of an open face stratum of
\(\mathcal K_0\), freeze the coframe \(E\) and the connection at a point of
that stratum and rescale the V33 Dirac and stiffness forms by \(a\).  Denote
the resulting tangent sign operator by
\begin{equation}
 \mathbb H_\tau(E,k_\parallel),                          \label{eq:newS7V33-tangent-operator}
\end{equation}
after Fourier transformation in the tangential directions; it acts on the
remaining normal discrete fibres.  SFGC consists of:
\begin{enumerate}[label=\textup{(\roman*)},leftmargin=3.5em]
\item the uniform tangent gap
      \begin{equation}
       \inf_\tau\inf_{E\in\mathscr E}
       \inf_{k_\parallel}inf\sigma
       |\mathbb H_\tau(E,k_\parallel)|
       =:\delta_{\rm FE}>0;                            \label{eq:newS7V33-SFGC}
      \end{equation}
\item absence of a zero-energy normal bound state for every positive-
      codimension face, edge, and vertex tangent model;
\item finite-section stability: collars cut from a tangent model and closed
      by the declared FSGC boundary have gap at least
      \(\delta_{\rm FE}/2\) once their rescaled radius is large; and
\item uniformity of the preceding rows through the compact coframe set and
      the finite set of original incidence types.
\end{enumerate}
Here \(\inf\sigma|T|\) denotes the distance of the spectrum of the
self-adjoint operator \(T\) from zero.
\end{definition}

\begin{lemma}[Why SFGC is a finite family of spectral problems]
\label{lem:newS7V33-SFGC-finite}
The index set of tangent incidence types in SFGC is finite.  For each type,
the only continuous parameters are the compact coframe variable \(E\) and a
finite-dimensional tangential Brillouin variable \(k_\parallel\).
\end{lemma}

\begin{proof}
The original complex has finitely many simplices and hence finitely many
face-incidence links.  The edgewise subdivision is periodic in barycentric
integer coordinates parallel to any fixed open face.  Its normal link is
one of the finite combinatorial patterns from
\Cref{lem:newS7V33-finite-stars}.  Fourier transformation replaces the
tangential translations by \(k_\parallel\) on a compact torus.  Uniform
ellipticity confines \(E\) to the declared compact set.
\end{proof}

\begin{theorem}[SFGC is sufficient for the V31 seam gap]
\label{thm:newS7V33-SFGC-sufficient}
Assume SFGC and bounded smooth coframe--connection jets.  Then the V33
edgewise finite-element kernel satisfies SLGR for all sufficiently small
\(a\), with \(g_0=\delta_{\rm FE}/4\).  Consequently
\begin{equation}
 \|H_{G,a}^{\rm FE}\psi\|
 \ge{\delta_{\rm FE}\over8a}\|\psi\|                  \label{eq:newS7V33-global-gap}
\end{equation}
after decreasing \(a\) once more to absorb the V31 commutator constant.
\end{theorem}

\begin{proof}
Choose a fixed physical collar radius \(R\) inside the good cover.  Suppose
the asserted local gap fails along \(a_n\downarrow0\).  Rescale a violating
collar by \(a_n^{-1}\), choose a point where its normalized mass does not
vanish, and pass to a subsequence with one fixed original stratum incidence
type; this is possible by
\Cref{lem:newS7V33-SFGC-finite}.  Smoothness freezes the coframe and removes
the connection at leading rescaled order.  The local finite-element blocks
converge on every fixed rescaled ball to the corresponding tangent operator.

If the mass remains in bounded normal distance, the limiting vector gives a
zero generalized mode or normal bound state, contradicting SFGC(i)--(ii).
If it escapes to the closing boundary, SFGC(iii) gives the same
contradiction.  Uniformity in SFGC(iv) prevents loss along the compact
coframe or tangential-momentum parameters.  Thus the local closing gap is at
least \(\delta_{\rm FE}/(4a)\) for small \(a\), which is SLGR.  Apply the
V31 IMS theorem and choose \(a\) so that
\(C_\chi\le\delta_{\rm FE}/(8a)\).
\end{proof}

\begin{assessment}[Acquired rows and live gap]
\label{ass:newS7V33-status}
V33 constructs one global cofinal simplicial spin mesh and one intrinsic
finite-element Wilson kernel.  It proves GSMC's mesh, quadrature,
spin-cocycle, adjoint, positivity, scaling, and weak first-moment rows,
including two metric contacts.  The remaining sign question is SFGC: a
finite family of bulk and stratified tangent-operator gaps, normal bound-
state exclusions, and finite-section stability estimates.  Weak consistency
and isolated one-star singular values do not prove those spectral rows.
\end{assessment}

\begin{programme}[V34: bulk symbol and codimension-one tangent model]
\label{prog:newS7V33-V34}
V34 will first compute the translation-periodic bulk symbol of the
mass-lumped edgewise finite-element forms and test whether its positive
stiffness Wilson term removes every nonphysical zero.  It will then derive
the codimension-one face tangent operator and separate its tangential
Fourier momentum from its normal recurrence.  A possible zero-energy
surface mode will be decided by the normal transfer matrix before higher-
codimension strata are considered.
\end{programme}

\begin{firewall}[Claim boundary after seventh-series V33]
\label{fire:newS7V33-boundary}
V33 fixes a face-compatible cofinal edgewise subdivision of one finite
global triangulation, proves bounded degree and finite local combinatorics,
constructs intrinsic nodal spin interpolation, positive mass-lumped
quadrature, and the global Galerkin Dirac and connection-stiffness Wilson
forms, proves their exact covariance and gamma-five algebra, cutoff scaling,
and weak moment consistency through two metric variations.  It also proves
that isolated one-star gaps are insufficient, formulates the correct finite
family of stratified tangent spectral problems, and proves SFGC sufficient
for SLGR and the global overlap gap.

V33 does not prove SFGC, the global overlap gap for this finite-element
kernel, strong local heat-symbol consistency, Gaussian heat locality, the
common Einstein stress, pure gravitational owners, rough metric or AOP
control, measure concentration, the determinant-line phase, interacting
cutoff removal, or the Standard Model--Einstein continuum.
\end{firewall}

\begin{theorem}[Seventh-series V33 result]
\label{thm:newS7V33-result}
Every compact smooth spin four-manifold with a fixed finite triangulation
admits the explicit cofinal edgewise spin meshes and finite-element Wilson
operators of V33.  They have exact global spin--gauge covariance,
gamma-five Hermiticity, a positive Wilson form, bounded degree, natural
\(a^{-1}\)/\(a^{-2}\) scaling, positive second-order volume quadrature, and
the correct weak Dirac and Laplace moments through two metric variations.
Their remaining overlap-sign obstruction is the stratified tangent
certificate \eqref{eq:newS7V33-SFGC}, including its normal bound-state and
finite-section rows.
\end{theorem}

\begin{proof}
Use \Cref{lem:newS7V33-mesh,
lem:newS7V33-finite-stars,
lem:newS7V33-interpolant,
lem:newS7V33-quadrature,
thm:newS7V33-algebra,
thm:newS7V33-scaling,
lem:newS7V33-interpolation-error,
thm:newS7V33-moments,
cor:newS7V33-Wilson-vanishes,
prop:newS7V33-star-not-enough,
lem:newS7V33-SFGC-finite,
thm:newS7V33-SFGC-sufficient}.
\end{proof}


\paragraph{Parent-version record.}
V33 was formed from
\texttt{Renewal\_SM\_Einstein\_Continuum\_Research\_Notes\_V32.tex}.
The V32 parent had \(13\,100\) logical source lines, \(526\,117\) bytes,
and SHA--256
\[
 \texttt{5509DD0AE6F3FB30997C68FABBF0AF429F2E775A356420BD67D73F75932C2B5F}.
\]
The next compulsory companion audit is V35.

% =============================================================================
% END APPENDED SEVENTH-SERIES V33 MATERIAL
% =============================================================================


% =============================================================================
% BEGIN APPENDED SEVENTH-SERIES V34 MATERIAL
% =============================================================================

\clearpage
\section{Seventh-series V34: the regular simplicial bulk gap and the
codimension-one Evans gate}
\label{sec:v7v34}

V33 replaces the chart-seam question by stratified tangent operators.  V34
closes the translation-periodic bulk row.  In the interior of an original
simplex, edgewise subdivision is a one-vertex Bravais lattice.  The
mass-lumped Dirac symbol is therefore a Clifford vector and the stiffness
symbol is a positive scalar.  This makes their contributions to the Wilson
sign square commute and gives an exact bulk gap once one common Wilson
parameter separates the nonphysical Dirac zeros.  For a codimension-one
original face, V34 derives the normal recurrence and proves that a surface
zero exists exactly when a finite Evans determinant vanishes.

\subsection{The regular edgewise bulk symbols}
\label{sec:newS7V34-bulk-symbol}

Fix one translation-periodic interior type
\(\tau\in\mathfrak T_{\rm bulk}\subset\mathfrak T_{\rm ed}\), freeze the
coframe \(E\), set the connection to zero, and use barycentric integer
coordinates on its rank-four Bravais lattice.  After half-density
conjugation and rescaling, write
\begin{equation}
 \mathbb K_\tau(k;E):=aK_{\tau,a}^{\rm FE}(k),\qquad
 \mathbb L_\tau(k;E):=a^2L_{\tau,a}^{\rm FE}(k),         \label{eq:newS7V34-rescaled-symbols}
\end{equation}
for \(k\) in the Brillouin torus \(\mathbb T^4\).

\begin{lemma}[Clifford-vector Dirac symbol]
\label{lem:newS7V34-Dirac-symbol}
There are real trigonometric polynomials
\(s_{\tau,a}(k;E)\), \(a=1,\ldots,4\), such that
\begin{equation}
 \mathbb K_\tau(k;E)
 =i\sum_{a=1}^4\gamma^a s_{\tau,a}(k;E),qquad
 s_\tau(-k;E)=-s_\tau(k;E).                             \label{eq:newS7V34-K-symbol}
\end{equation}
Moreover,
\begin{equation}
 s_\tau(k;E)=Ek+O(|k|^3)                                \label{eq:newS7V34-K-consistency}
\end{equation}
in a fixed linear coordinate identification at \(k=0\), uniformly for \(E\)
in the compact elliptic set.
\end{lemma}

\begin{proof}
The frozen nodal space has one spinor fibre per Bravais vertex.  Translation
invariance makes every hopping coefficient depend only on the displacement
vector \(v\).  Skew-adjointness pairs \(v\) and \(-v\) with opposite
adjoints, and chirality excludes a scalar spin component.  The flat frozen
Galerkin Dirac form is linear in Clifford multiplication, hence its symbol
has the form \eqref{eq:newS7V34-K-symbol}; reversal makes the coefficients
real and odd.  Exactness of affine interpolation on linear functions gives
the derivative \(Ds_\tau(0;E)=E\).  Pairing opposite displacements removes
the quadratic Taylor term, giving the stated cubic remainder.
\end{proof}

\begin{lemma}[Positive scalar stiffness symbol]
\label{lem:newS7V34-stiffness-symbol}
There is an even real trigonometric polynomial
\(\ell_\tau(k;E)\) such that
\begin{equation}
 \mathbb L_\tau(k;E)=\ell_\tau(k;E)I,qquad
 \ell_\tau(k;E)\ge0,                                   \label{eq:newS7V34-L-symbol}
\end{equation}
and
\begin{equation}
 \ell_\tau(k;E)=0\quad\Longleftrightarrow\quad k=0,qquad
 \ell_\tau(k;E)=k^TG(E)k+O(|k|^4),                    \label{eq:newS7V34-L-properties}
\end{equation}
where \(G(E)\) is uniformly positive definite.
\end{lemma}

\begin{proof}
The frozen connection-stiffness form acts identically on spin components,
so its one-site Bloch symbol is scalar.  Positivity of the form gives
\(\ell_\tau\ge0\).  Equality means that the piecewise-affine Bloch
interpolant has zero gradient on every simplex and is therefore constant on
the connected periodic mesh.  A constant Bloch section has \(k=0\), and the
converse is immediate.  Affine exactness gives the continuum stiffness
quadratic form at the origin; shape regularity and ellipticity make \(G(E)\)
uniformly positive.  Evenness removes odd Taylor terms.
\end{proof}

\subsection{Removal of every regular bulk doubler}
\label{sec:newS7V34-bulk-gap}

Let
\begin{equation}
 Z_\tau(E):=\{k\in\mathbb T^4:s_\tau(k;E)=0\}.           \label{eq:newS7V34-zero-set}
\end{equation}
By \eqref{eq:newS7V34-K-consistency}, \(k=0\) is isolated uniformly on a
small neighbourhood.  Define, with the convention that the infimum of the
empty set is \(+\infty\),
\begin{equation}
 \ell_*:=inf_{\tau\in\mathfrak T_{\rm bulk}}
 \inf_{E\in\mathscr E}
 \inf_{k\in Z_\tau(E)\setminus\{0\}}\ell_\tau(k;E).     \label{eq:newS7V34-ell-star}
\end{equation}

\begin{lemma}[Uniform separation of nonphysical zeros]
\label{lem:newS7V34-zero-separation}
The number \(\ell_*\) in \eqref{eq:newS7V34-ell-star} is strictly positive
provided the zero sets vary closedly with \(E\); this closedness holds for
the finite trigonometric symbols above.
\end{lemma}

\begin{proof}
Uniform invertibility of \(Ds_\tau(0;E)=E\) and compactness of
\(\mathscr E\) give one neighbourhood of the origin containing no other
zero.  The remaining joint zero set in
\(\mathbb T^4\times\mathscr E\) is compact because it is the inverse image
of zero under a continuous finite family.  By
\eqref{eq:newS7V34-L-properties}, \(\ell_\tau\) is positive there.  Its
minimum over the finite compact union is positive.
\end{proof}

Choose one Wilson parameter satisfying
\begin{equation}
 r>{2\rho\over\ell_*}                                   \label{eq:newS7V34-r-choice}
\end{equation}
when nonphysical zeros occur; otherwise take any \(r>0\).

\begin{theorem}[Uniform regular-bulk Wilson sign gap]
\label{thm:newS7V34-bulk-gap}
For the common choice \eqref{eq:newS7V34-r-choice}, there is
\(\delta_{\rm bulk}>0\) such that every regular interior type obeys
\begin{equation}
 \left\|\Gamma\left\{
 \mathbb K_\tau(k;E)+{r\over2}\mathbb L_\tau(k;E)-\rho
 \right\}v\right\|
 \ge\delta_{\rm bulk}\|v\|                             \label{eq:newS7V34-bulk-gap}
\end{equation}
uniformly in \(k,E,\tau\).  The corresponding dimensional gap is
\(\delta_{\rm bulk}/a\).
\end{theorem}

\begin{proof}
The stiffness symbol is scalar and therefore commutes with the
Clifford-vector Dirac symbol.  Gamma-five Hermiticity gives the exact square
\begin{equation}
 \left|s_\tau(k;E)\right|^2
 +\left\{{r\over2}\ell_\tau(k;E)-\rho\right\}^2.       \label{eq:newS7V34-bulk-square}
\end{equation}
A zero would require \(k\in Z_\tau(E)\).  At the physical origin the second
term is \(\rho^2>0\).  At a nonphysical zero,
\((r/2)\ell_\tau\ge(r/2)\ell_*>\rho\), so the second term is again
nonzero.  The continuous expression is positive on the finite compact
parameter space; its minimum has a positive square root
\(\delta_{\rm bulk}\).
\end{proof}

\begin{corollary}[The bulk part of SFGC]
\label{cor:newS7V34-bulk-SFGC}
All codimension-zero tangent operators in V33 satisfy SFGC with the one
common Wilson parameter \eqref{eq:newS7V34-r-choice}.  No enumeration of
their individual doubler momenta is required.
\end{corollary}

\begin{proof}
Every codimension-zero tangent operator is a regular translation-periodic
type from the finite set \(\mathfrak T_{\rm bulk}\).  Apply
\Cref{thm:newS7V34-bulk-gap}.
\end{proof}

\subsection{The codimension-one face operator}
\label{sec:newS7V34-face}

Fix an open original three-face \(F\) shared by two four-simplices.  The
edgewise subdivision is translation-periodic along \(F\).  Fourier transform
the three tangential integer coordinates and retain a normal layer index
\(m\in\mathbb Z\).  The frozen dimensionless sign operator has finite normal
range \(J\):
\begin{equation}
 (\mathbb H_F(E,k_\parallel)u)_m
 =\sum_{j=-J}^J A_j(m;E,k_\parallel)u_{m+j},             \label{eq:newS7V34-face-recurrence}
\end{equation}
where \(A_j(m)\) equals one constant Laurent coefficient family for
\(m\ge m_+\), another for \(m\le m_-\), and differs only on finitely many
interface layers.

\begin{lemma}[No unit-circle characteristic root]
\label{lem:newS7V34-no-unit-root}
For either constant half-space recurrence, the characteristic pencil
\begin{equation}
 \mathcal P_\pm(z;E,k_\parallel)
 :=\sum_{j=-J}^JA_j^\pm(E,k_\parallel)z^j              \label{eq:newS7V34-pencil}
\end{equation}
has no nonzero root with \(|z|=1\).
\end{lemma}

\begin{proof}
Writing \(z=e^{ik_\perp}\) turns the pencil into the bulk Bloch sign symbol
of the adjacent original simplex.  A unit-circle root would give a zero of
that bulk symbol, contrary to \Cref{thm:newS7V34-bulk-gap}.
\end{proof}

After the standard companion linearization of the finite-range recurrence,
let \(\mathcal S_+(E,k_\parallel)\) be the space of Cauchy data which decay as
\(m\to+\infty\), and \(\mathcal U_-(E,k_\parallel)\) the data which decay as
\(m\to-\infty\).  Lemma \ref{lem:newS7V34-no-unit-root} makes these
exponential dichotomy spaces well-defined.  Propagate both to one common
interface section using \eqref{eq:newS7V34-face-recurrence}.  Let
\(B_+(E,k_\parallel)\) and \(B_-(E,k_\parallel)\) be matrices whose columns
are bases of the propagated spaces.

\begin{definition}[Face Evans determinant]
\label{def:newS7V34-Evans}
When the stable and unstable dimensions sum to the Cauchy-data dimension,
define
\begin{equation}
 \mathcal E_F(E,k_\parallel)
 :=\det\bigl(B_-(E,k_\parallel)\  B_+(E,k_\parallel)\bigr).             \label{eq:newS7V34-Evans}
\end{equation}
If the dimensions do not sum correctly, set the face gap card to fail.
Changing either basis multiplies \(\mathcal E_F\) by a nonzero determinant,
so its zero set is intrinsic.
\end{definition}

\begin{theorem}[Exact surface-zero criterion]
\label{thm:newS7V34-surface-criterion}
For fixed \(E,k_\parallel\), the face tangent operator has a nonzero
square-summable zero mode if and only if
\begin{equation}
 \mathcal E_F(E,k_\parallel)=0.                          \label{eq:newS7V34-Evans-zero}
\end{equation}
\end{theorem}

\begin{proof}
A square-summable solution of the constant recurrence on the positive
half-line has interface Cauchy data in \(\mathcal S_+\); on the negative
half-line its data lie in \(\mathcal U_-\).  The finitely many interface
equations propagate these data to the same section.  A global decaying
solution exists exactly when the two propagated subspaces intersect
nontrivially.  Under the complementary-dimension condition, this happens
exactly when their concatenated basis matrix is singular, which is
\eqref{eq:newS7V34-Evans-zero}.
\end{proof}

\begin{definition}[Codimension-one face gap card, C1FGC]
\label{def:newS7V34-C1FGC}
Choose orthonormal bases of the propagated dichotomy spaces in the common
Cauchy-data inner product.  C1FGC is the uniform nonvanishing statement
\begin{equation}
 \inf_F\inf_{E\in\mathscr E}\inf_{k_\parallel}
 |\mathcal E_F(E,k_\parallel)|>0,                       \label{eq:newS7V34-C1FGC}
\end{equation}
together with a uniform lower bound on the bulk dichotomy exponents.  A
change between orthonormal bases is unitary within each subspace and changes
the determinant only by a scalar of modulus one.  Thus its modulus is
intrinsic; equivalently it is the product of the sines of the principal
angles between the two propagated subspaces.
\end{definition}

\begin{proposition}[C1FGC closes the codimension-one SFGC row]
\label{prop:newS7V34-C1FGC-sufficient}
The bulk theorem and C1FGC imply a uniform zero-energy gap and finite-section
stability for every open-face tangent model.
\end{proposition}

\begin{proof}
The bulk gap gives uniform exponential dichotomies by compactness.  C1FGC
gives a uniform angle between the propagated stable and unstable spaces, so
the interface matching operator has a uniformly bounded inverse at zero.
The same dichotomy estimates show that cutting the two tails at distance
\(N\) changes the inverse by \(O(e^{-cN})\).  Hence sufficiently long finite
sections preserve a fixed fraction of the infinite-model gap.
\end{proof}

\begin{assessment}[Bulk closed; face determinant open]
\label{ass:newS7V34-status}
V34 proves the regular bulk portion of SFGC for one common sufficiently
large Wilson parameter.  It also replaces the vague possibility of a
codimension-one surface mode by the exact finite Evans determinant
\eqref{eq:newS7V34-Evans}.  C1FGC has not been proved: the determinant must
still be evaluated or bounded uniformly over the finite original faces,
compact coframes, and tangential Brillouin torus.
\end{assessment}

\begin{programme}[V35: compulsory audit and the face determinant]
\label{prog:newS7V34-V35}
V35 will perform the compulsory YM/NS audit, then compute the interface
coefficient matrices \(A_j\) from the V33 mass-lumped Dirac and stiffness
forms for one oriented reference face.  It will test whether reflection of
the two adjacent simplex orderings gives a unitary conjugacy which makes the
Evans determinant manifestly nonzero.  If not, the first possible
zero-energy surface mode will be isolated as a finite generalized
eigenvalue problem before codimension-two strata are attempted.
\end{programme}

\begin{firewall}[Claim boundary after seventh-series V34]
\label{fire:newS7V34-boundary}
V34 derives the regular edgewise Dirac and stiffness Bloch symbols, proves
uniform separation of nonphysical Dirac zeros by positive stiffness, chooses
one common Wilson parameter, and proves the full codimension-zero tangent
gap.  It derives the codimension-one normal recurrence, proves absence of
unit-circle bulk roots, defines the face Evans determinant, proves its exact
surface-zero criterion, and proves C1FGC sufficient for the face SFGC row.

V34 does not prove C1FGC, the higher-codimension tangent gaps, full SFGC, the
global finite-element overlap gap, strong local heat consistency, Gaussian
heat locality, the common Einstein stress, pure gravitational owners, rough
metric or AOP control, measure concentration, the determinant-line phase,
interacting cutoff removal, or the Standard Model--Einstein continuum.
\end{firewall}

\begin{theorem}[Seventh-series V34 result]
\label{thm:newS7V34-result}
For every regular interior type of the V33 edgewise finite-element mesh,
one common Wilson parameter satisfying \eqref{eq:newS7V34-r-choice} removes
all nonphysical Dirac zeros and gives a uniform sign-kernel gap
\(\delta_{\rm bulk}/a\).  For every codimension-one original face, a
square-summable zero-energy surface mode exists exactly at a zero of the
finite Evans determinant \eqref{eq:newS7V34-Evans}.  The next gap is the
uniform nonvanishing card C1FGC.
\end{theorem}

\begin{proof}
Use \Cref{lem:newS7V34-Dirac-symbol,
lem:newS7V34-stiffness-symbol,
lem:newS7V34-zero-separation,
thm:newS7V34-bulk-gap,
cor:newS7V34-bulk-SFGC,
lem:newS7V34-no-unit-root,
thm:newS7V34-surface-criterion,
prop:newS7V34-C1FGC-sufficient}.
\end{proof}

\paragraph{Parent-version record.}
V34 was formed from
\texttt{Renewal\_SM\_Einstein\_Continuum\_Research\_Notes\_V33.tex}.
The V33 parent had \(13\,592\) logical source lines, \(548\,612\) bytes,
and SHA--256
\[
 \texttt{2120D65F6EDCA90CBEB18D28AA329EED912C2528940A422B98F473B33C522252}.
\]
The next compulsory companion audit is V35.

% =============================================================================
% END APPENDED SEVENTH-SERIES V34 MATERIAL
% =============================================================================

% =============================================================================
% BEGIN APPENDED SEVENTH-SERIES V35 MATERIAL
% =============================================================================

\section{Seventh-series V35: compulsory audit and the face-interface acquisition}
\label{sec:v7v35}
\label{sec:newS7V35}

\subsection{Purpose and companion audit}
\label{sec:newS7V35-audit}

The active obstruction after V34 is no longer a formal continuum
coefficient.  It is the possibility that a frozen codimension-one mesh
interface carries a square-summable zero mode even though both adjacent
bulk Wilson sign kernels are uniformly gapped.  This note first performs
the compulsory fifth-version comparison with the active Yang--Mills and
Navier--Stokes programmes.  It then identifies a class of faces for which
the obstruction vanishes exactly and a larger perturbative class for which
it vanishes quantitatively.

\begin{audit}[Companion snapshots used at V35]
\label{aud:newS7V35-snapshots}
The companion files inspected before beginning the new argument were:
\begin{center}
\begin{tabular}{p{0.20\linewidth}p{0.52\linewidth}r}
programme & source & bytes\\ \hline
Yang--Mills &
\texttt{Renewal\_Yang\_Mills\_Mass\_Gap\_Research\_Notes\_V11.tex}
& (176,232)\\
Navier--Stokes &
\texttt{Renewal\_NS\_Stability\_Research\_Notes\_V46.tex}
& (462,114)
\end{tabular}
\end{center}
The Yang--Mills source had (5,027) newline characters and SHA--256
\[
\texttt{223BED98D69854FA8AE2AA8351EDEE32CB7B6CCA14AC8D51A82C8EF9E8F52FCB}.
\]
The Navier--Stokes source had (9,154) newline characters and SHA--256
\[
\texttt{06BC83CB7707FBB6731FC528B5D63C741F9BFA894E0A3F582BA57C2DE7F2075F}.
\]
These hashes fix the exact documents used in the audit; later edits of a
companion programme cannot silently change the present inference.
\end{audit}

The Yang--Mills note proves an exact matrix--forest determinant identity,
identifies the graded coefficient (e_g(\lambda_2,\ldots,\lambda_N)), and
uses the bound (\Sigma_x^g/g!) to close its diffuse pair-mass RFAC row.
Its concentrated component-Laplacian card remains open.  The useful lesson
here is structural: a local decomposition does not permit one to discard
the determinant which retains the global matching correlations.

The Navier--Stokes note rewrites localized pressure work as an
energy--Hessian pairing, obtains a far-pressure contribution of order
(O(\varepsilon)), and resolves each cell into multipoles with an explicit
sign geometry.  Its sign problem remains global.  The useful lesson here is
again representational: the variables should be adapted to the interface
or localization geometry before estimating them.

\begin{decision}[Effect of the V35 audit]
\label{dec:newS7V35-audit}
Neither companion source supplies an SM metric coefficient or a missing
Wilson gap estimate.  The Yang--Mills determinant argument supports
retaining the exact Evans determinant rather than replacing it by
independent entrywise estimates.  The Navier--Stokes localization argument
supports the tangential Fourier/normal recurrence representation already
introduced in V34.  Accordingly, the SM direction remains the exact face
matching problem.
\end{decision}

\subsection{Operator setting for one frozen face}
\label{sec:newS7V35-setting}

Fix a face type (F), a compact coframe parameter (E\in\mathscr E), and
tangential momentum (k_\parallel\in\mathbb T^3).  Write
\[
 \mathscr H_\perp:=\ell^2(\mathbb Z;\mathbb C^q)
\]
for the normal-layer Hilbert space, after enlarging (q) if a finite-range
recurrence has been put into a nearest-layer block form.  Let
(H_F(E,k_\parallel)) be the dimensionless self-adjoint frozen Wilson sign
kernel.  All operator norms in this section are norms on
(\mathscr H_\perp).

\begin{definition}[Transparent reference face]
\label{def:newS7V35-transparent}
A face family (H_{\rm tr}(E,k_\parallel)) is \emph{transparent} with bulk
reference (H_{\rm b}(E,k_\parallel)) if there is a unitary layer operator
(U(E,k_\parallel)) such that
\begin{equation}
 H_{\rm tr}=U H_{\rm b}U^*,                            \label{eq:newS7V35-transparent}
\end{equation}
where (H_{\rm b}) is a full-line translation-periodic bulk tangent
operator satisfying
\begin{equation}
 \|H_{\rm b}u\|\ge\delta_{\rm b}\|u\|                \label{eq:newS7V35-bulk-lower}
\end{equation}
with one (delta_{\rm b}>0) uniform over its parameter family.  We allow
(U) to be layer dependent, but require it to have uniformly bounded
finite propagation when locality is invoked.
\end{definition}

The definition includes a change of spin frame on one side of a face.  In
that case (U_m) is the corresponding spin lift on each layer and
\eqref{eq:newS7V35-transparent} is exactly the cocycle covariance proved in
V30.  A change of simplex incidence stencil is not automatically of this
form.

\begin{theorem}[Transparent faces inherit the bulk gap]
\label{thm:newS7V35-transparent-gap}
Every transparent face satisfies
\begin{equation}
 \|H_{\rm tr}u\|\ge\delta_{\rm b}\|u\|.               \label{eq:newS7V35-transparent-gap}
\end{equation}
It has no square-summable zero mode, and its V34 Evans determinant is
nonzero for every parameter.
\end{theorem}

\begin{proof}
Unitary invariance and \eqref{eq:newS7V35-bulk-lower} give
\[
 \|H_{\rm tr}u\|
 =\|H_{\rm b}U^*u\|
 \ge\delta_{\rm b}\|U^*u\|
 =\delta_{\rm b}\|u\|.
\]
Thus the kernel is trivial.  The exact surface-zero criterion of
\Cref{thm:newS7V34-surface-criterion} then gives nonvanishing of the Evans
determinant.
\end{proof}

\begin{corollary}[Pure spin-frame seams are harmless]
\label{cor:newS7V35-spin-seam}
If the two descriptions of a frozen face differ only by the nodal spin
cocycle of V30, while their incidence, mass-lumping, and stiffness
coefficients agree after relabelling, then the face is transparent and
obeys \eqref{eq:newS7V35-transparent-gap}.
\end{corollary}

\begin{proof}
The nodal spin transformation is unitary in the half-density Hilbert space.
The V30 covariance identities conjugate every transport, Clifford, Wilson,
and sign-kernel entry.  Coefficient agreement after relabelling therefore
gives \eqref{eq:newS7V35-transparent}; apply
\Cref{thm:newS7V35-transparent-gap}.
\end{proof}

\subsection{A quantitative acquisition neighbourhood}
\label{sec:newS7V35-acquisition}

Let a general face be written relative to a transparent reference as
\begin{equation}
 H_F=H_{\rm tr}+W_F.                                  \label{eq:newS7V35-split}
\end{equation}
For the frozen finite-element construction of V33, (W_F) has finite
normal support whenever the two operators have the same two asymptotic
half-space stencils.  The following estimate does not require finite
support.

\begin{theorem}[Relative-norm face gap]
\label{thm:newS7V35-relative-gap}
Suppose (H_{\rm tr}) is transparent with gap (delta_{\rm b}) and
\begin{equation}
 \|W_F\|\le\eta<\delta_{\rm b}.                       \label{eq:newS7V35-smallness}
\end{equation}
Then
\begin{equation}
 \|H_Fu\|\ge(\delta_{\rm b}-\eta)\|u\|               \label{eq:newS7V35-face-gap}
\end{equation}
for every (u\in\mathscr H_\perp).  In particular, (H_F) is invertible,
\begin{equation}
 \|H_F^{-1}\|\le {1\over\delta_{\rm b}-\eta},         \label{eq:newS7V35-inverse}
\end{equation}
and it has no surface zero mode.
\end{theorem}

\begin{proof}
The reverse triangle inequality and
\Cref{thm:newS7V35-transparent-gap} give
\[
 \|H_Fu\|
 \ge\|H_{\rm tr}u\|-\|W_Fu\|
 \ge(\delta_{\rm b}-\eta)\|u\|.
\]
Because (H_F) is self-adjoint, the same positive lower bound applies to
the distance from zero to its spectrum.  Hence it is onto as well as
one-to-one and \eqref{eq:newS7V35-inverse} follows from the spectral
theorem.
\end{proof}

\begin{definition}[Interface relative-norm card, IRNC]
\label{def:newS7V35-IRNC}
For every original face (F), choose a transparent reference with bulk gap
(delta_{{\rm b},F}).  IRNC is the existence of one (c_{\rm I}>0) such
that
\begin{equation}
 \sup_{F,E,k_\parallel}
 \|H_F-H_{{\rm tr},F}\|
 \le
 \inf_F\delta_{{\rm b},F}-c_{\rm I}.                 \label{eq:newS7V35-IRNC}
\end{equation}
The card is finite and explicit: every difference has only finitely many
nonzero block rows, and every block is determined by the two adjacent
simplex coframes and their face incidence data.
\end{definition}

\begin{proposition}[IRNC closes C1FGC]
\label{prop:newS7V35-IRNC-C1FGC}
If IRNC holds, then all face tangent kernels have gap at least (c_{\rm I}),
and the codimension-one face gap card C1FGC of V34 holds.
\end{proposition}

\begin{proof}
The first assertion is \Cref{thm:newS7V35-relative-gap}, uniformly in the
finite face set and compact parameter set.  The bulk gap also gives uniform
exponential dichotomies.  If the normalized Evans determinants had a
sequence tending to zero, compactness would give a limiting parameter at
which the propagated stable and unstable spaces intersect.  The V34 exact
surface-zero criterion would then produce a zero mode, contradicting the
uniform operator gap.  Hence the normalized determinants have a positive
minimum, which is C1FGC.
\end{proof}

\begin{lemma}[A computable row-sum sufficient test]
\label{lem:newS7V35-row-sum}
Write the self-adjoint defect in normal blocks as
\[
 (W_Fu)_m=\sum_n W_{mn}u_n
\]
and suppose only finitely many (W_{mn}) are nonzero.  Then
\begin{equation}
 \|W_F\|
 \le
 \left(
   \sup_m\sum_n\|W_{mn}\|
   \sup_n\sum_m\|W_{mn}\|
 \right)^{1/2}.                                      \label{eq:newS7V35-Schur}
\end{equation}
For a self-adjoint block matrix, the two suprema agree.  Thus the maximum
block row sum being smaller than (delta_{\rm b}) is sufficient for
\eqref{eq:newS7V35-smallness}.
\end{lemma}

\begin{proof}
For nonnegative scalar sequences (a_m=\|u_m\|), the triangle inequality
gives
\(|(W_Fu)_m\|\le\sum_n\|W_{mn}\|a_n\).
The scalar Schur test yields \eqref{eq:newS7V35-Schur}.  Self-adjointness
gives (|W_{mn}\|=|W_{nm}\|), so the row and column suprema coincide.
\end{proof}

\subsection{What the acquisition theorem does and does not cover}
\label{sec:newS7V35-scope}

The V33 face construction has three logically distinct sources of an
interface defect:
\begin{enumerate}
\item a spin-frame change, which is removed exactly by
\Cref{cor:newS7V35-spin-seam};
\item a relabelling or reflection which may be removed by a unitary
permutation if it preserves the mass and stiffness entries; and
\item a genuine change in the affine simplex coframe or incidence pattern,
which survives as (W_F).
\end{enumerate}
Only the third source enters IRNC.  Smoothness of the continuum coframe
does not by itself make its frozen jump across an original face small: the
two affine simplex charts can have an order-one relative shear even when
the underlying metric is smooth.  Therefore no IRNC estimate is asserted
without evaluating the actual V33 blocks.

\begin{proposition}[Why chart smoothness alone cannot close IRNC]
\label{prop:newS7V35-no-smoothness-shortcut}
There is no estimate of the form
\[
 \|W_F\|\le C a\|\nabla E\|_{L^\infty}
\]
which follows solely from smoothness of the continuum coframe and shape
regularity of two independently ordered affine simplices.
\end{proposition}

\begin{proof}
Take a constant Euclidean coframe, so the right side is zero.  Two adjacent
simplices may use different vertex orderings.  Their pulled-back edge
generators and local mass/stiffness matrices are then related only after
the corresponding incidence permutation.  If that permutation is omitted
from the comparison, the block matrices differ by an order-one amount.
Thus the claimed estimate is false before the exact combinatorial unitary
has been extracted.  After it is extracted, any remaining defect must be
estimated from the explicit affine data rather than from coframe
smoothness alone.
\end{proof}

\begin{assessment}[V35 acquisition]
\label{ass:newS7V35-status}
V35 closes every transparent face and every face lying in the explicit
IRNC neighbourhood of a transparent reference.  The general face problem
is not closed because the required norm smallness need not hold.  The
correct nonperturbative next object is a finite-dimensional determinant
built from the transparent resolvent and the finitely supported defect.
\end{assessment}

\begin{programme}[V36: Birman--Schwinger face determinant]
\label{prog:newS7V35-V36}
V36 will factor the finite-layer defect through its interface space and
prove that a face zero mode exists exactly when a finite
Birman--Schwinger matrix has eigenvalue (-1).  It will relate that
determinant to the V34 Evans zero set and isolate the precise finite matrix
whose nonvanishing must be checked when IRNC fails.
\end{programme}

\begin{firewall}[Claim boundary after seventh-series V35]
\label{fire:newS7V35-boundary}
V35 records the compulsory YM/NS audit, proves exact gap inheritance for
transparent faces, proves a quantitative gap for defects smaller than the
bulk gap, supplies a computable block Schur test, and proves that the
resulting IRNC implies C1FGC.  It also proves why continuum coframe
smoothness cannot replace the missing combinatorial comparison.

V35 does not prove IRNC for every V33 face, evaluate the face Evans
determinants, exclude nonperturbative face modes, prove the higher
codimension tangent gaps, full SFGC, the global finite-element overlap gap,
strong local heat consistency, Gaussian heat locality, the common Einstein
stress, pure gravitational owners, rough metric or AOP control, measure
concentration, the determinant-line phase, interacting cutoff removal, or
the Standard Model--Einstein continuum.
\end{firewall}

\begin{theorem}[Seventh-series V35 result]
\label{thm:newS7V35-result}
Pure spin-frame seams are transparent and inherit the V34 bulk gap.
More generally, every face whose finite-layer defect satisfies IRNC has a
uniform gap and satisfies C1FGC.  A genuine order-one incidence or affine
defect remains a finite-dimensional nonperturbative matching problem.
\end{theorem}

\begin{proof}
Combine \Cref{thm:newS7V35-transparent-gap,
cor:newS7V35-spin-seam,
thm:newS7V35-relative-gap,
prop:newS7V35-IRNC-C1FGC,
lem:newS7V35-row-sum,
prop:newS7V35-no-smoothness-shortcut}.
\end{proof}

\paragraph{Parent-version record.}
V35 was formed from
\texttt{Renewal\_SM\_Einstein\_Continuum\_Research\_Notes\_V34.tex}.
The V34 parent had (13,940) logical source lines, (563,881) bytes,
and SHA--256
\[
 \texttt{488C5745C8042EA4C3BD6A85CAD1A4A0A843A95E49B1F68469BB11C34CF6A1EA}.
\]
The next compulsory companion audit is V40.

% =============================================================================
% END APPENDED SEVENTH-SERIES V35 MATERIAL
% =============================================================================

% =============================================================================
% BEGIN APPENDED SEVENTH-SERIES V36 MATERIAL
% =============================================================================

\section{Seventh-series V36: finite Birman--Schwinger reduction of a face}
\label{sec:v7v36}
\label{sec:newS7V36}

\subsection{The reference-interface hypothesis}
\label{sec:newS7V36-reference}

V35 settles transparent and relatively small face defects.  This note does
not assume smallness.  Instead it eliminates the two infinite normal tails
exactly.  The only additional input is an invertible comparison interface
with the same frozen tails.

Fix (F,E,k_\parallel) as in V35.  Let (H_F) denote the dimensionless
self-adjoint face sign kernel on
(mathscr H_\perp=\ell^2(\mathbb Z;\mathbb C^q)).

\begin{definition}[Admissible face reference]
\label{def:newS7V36-admissible-reference}
An operator (H_0=H_{0,F}(E,k_\parallel)) is an admissible reference for
(H_F) if:
\begin{enumerate}
\item (H_0) is bounded and self-adjoint;
\item (H_0) agrees with (H_F) outside a fixed finite set of normal
layers;
\item there is a parameter-uniform (delta_0>0) such that
\begin{equation}
 |H_0u|\ge\delta_0|u|.                            \label{eq:newS7V36-H0gap}
\end{equation}
\end{enumerate}
A transparent face from V35 is an admissible reference whenever it has the
same two asymptotic half-space stencils as (H_F).
\end{definition}

The second condition makes (W:=H_F-H_0) finite rank.  Choose a finite set
of layers (I_F\subset\mathbb Z) containing every block row and block
column of (W), let
\begin{equation}
 P_F:\mathscr H_\perp\longrightarrow
 \mathscr I_F:=\bigoplus_{m\in I_F}\mathbb C^q          \label{eq:newS7V36-P}
\end{equation}
be coordinate restriction, and write uniquely
\begin{equation}
 W=P_F^*V_FP_F,qquad V_F=V_F^*.                       \label{eq:newS7V36-factor}
\end{equation}
Here (P_FP_F^*=I_{\mathscr I_F}) and (|P_F|=1).

\begin{definition}[Compressed reference Green matrix]
\label{def:newS7V36-Green}
Define the finite matrix
\begin{equation}
 G_F:=P_FH_0^{-1}P_F^*                                \label{eq:newS7V36-G}
\end{equation}
and the face Birman--Schwinger matrix
\begin{equation}
 M_F:=I_{\mathscr I_F}+G_FV_F.                        \label{eq:newS7V36-M}
\end{equation}
Both retain the dependence on (E,k_\parallel).
\end{definition}

\subsection{Exact zero-mode equivalence}
\label{sec:newS7V36-equivalence}

\begin{theorem}[Finite face Birman--Schwinger principle]
\label{thm:newS7V36-BS}
Under \Cref{def:newS7V36-admissible-reference}, the maps
\begin{align}
 ker H_F&\longrightarrow\ker M_F,&
 \psi&\longmapsto P_F\psi,                            \label{eq:newS7V36-map-forward}\\
 ker M_F&\longrightarrow\ker H_F,&
 u&\longmapsto-H_0^{-1}P_F^*V_Fu                      \label{eq:newS7V36-map-back}
\end{align}
are mutually inverse.  Consequently
\begin{equation}
 0\in\sigma(H_F)quad\Longleftrightarrow\quad
 \det M_F=0,                                          \label{eq:newS7V36-det-criterion}
\end{equation}
and the two kernels have the same dimension.
\end{theorem}

\begin{proof}
If (H_F\psi=0), then
\[
 \psi=-H_0^{-1}P_F^*V_FP_F\psi.
\]
Applying (P_F) gives (M_F(P_F\psi)=0).  If
(P_F\psi=0), the displayed identity gives (psi=0), so the forward map
is injective.

Conversely, let (M_Fu=0) and set
(psi=-H_0^{-1}P_F^*V_Fu).  Then
\[
 P_F\psi=-G_FV_Fu=u.
\]
It follows that
\[
 H_F\psi
 =H_0\psi+P_F^*V_FP_F\psi
 =-P_F^*V_Fu+P_F^*V_Fu=0.
\]
The two constructions preserve (u=P_F\psi), so they are inverse.  The
finite-dimensional determinant criterion and equality of nullities follow.
\end{proof}

\begin{corollary}[Independence of interface-window padding]
\label{cor:newS7V36-padding}
Enlarging (I_F) by layers on which (W) vanishes multiplies the
Birman--Schwinger determinant by (1).  In particular, the determinant in
\eqref{eq:newS7V36-det-criterion} is unchanged.
\end{corollary}

\begin{proof}
In the orthogonal decomposition into the old interface space and the added
coordinates, the enlarged defect matrix has the form
\(widetilde V_F=\left(\begin{smallmatrix}V_F&0\\0&0\end{smallmatrix}\right)\).
Thus
\[
 I+\widetilde G_F\widetilde V_F
 =\begin{pmatrix}I+G_FV_F&0\\ *&I\end{pmatrix},
\]
whose determinant is (det M_F).
\end{proof}

\begin{corollary}[Reference invariance of the zero set]
\label{cor:newS7V36-reference-invariance}
If (H_0) and (widetilde H_0) are two admissible references, their finite
Birman--Schwinger determinants vanish at exactly the same parameters and
with the same kernel dimension.
\end{corollary}

\begin{proof}
Apply \Cref{thm:newS7V36-BS} to each reference.  Both finite kernels are
canonically isomorphic to (ker H_F).
\end{proof}

\subsection{Relation to the Evans determinant}
\label{sec:newS7V36-Evans}

\begin{theorem}[Evans and Birman--Schwinger zero sets coincide]
\label{thm:newS7V36-Evans-BS}
Assume the V34 exponential dichotomies and an admissible reference.  Then
\begin{equation}
 \mathcal E_F(E,k_\parallel)=0
 \quad\Longleftrightarrow\quad
 \det M_F(E,k_\parallel)=0.                           \label{eq:newS7V36-zero-set}
\end{equation}
The nullity on either side is the dimension of the square-summable surface
zero-mode space.
\end{theorem}

\begin{proof}
By the exact surface-zero criterion of
\Cref{thm:newS7V34-surface-criterion}, the Evans determinant vanishes
exactly when (ker H_F\ne0).  By
\Cref{thm:newS7V36-BS}, this is exactly when (det M_F=0), and both
matching constructions identify their nullspace with (ker H_F).
\end{proof}

This theorem permits two complementary calculations.  The Evans form uses
only the finite-range recurrence and decaying bulk subspaces.  The
Birman--Schwinger form uses the reference Green matrix on finitely many
layers.  Neither replaces the determinant by unrelated local estimates.

\subsection{Resolvent formula and a quantitative card}
\label{sec:newS7V36-resolvent}

\begin{theorem}[Finite-interface resolvent formula]
\label{thm:newS7V36-Woodbury}
The face operator (H_F) is invertible if and only if (M_F) is
invertible.  In that event
\begin{equation}
 H_F^{-1}
 =H_0^{-1}
 -H_0^{-1}P_F^*V_FM_F^{-1}P_FH_0^{-1}.               \label{eq:newS7V36-Woodbury}
\end{equation}
If
(sigma_{min}(M_F)\ge\mu_F>0), then
\begin{equation}
 |H_F^{-1}|
 \le {1\over\delta_0}
 +{\|V_F\|\over\delta_0^2\mu_F}.                    \label{eq:newS7V36-inverse-bound}
\end{equation}
\end{theorem}

\begin{proof}
Direct multiplication of the right side of
\eqref{eq:newS7V36-Woodbury} by
(H_0+P_F^*V_FP_F), using
(M_F=I+P_FH_0^{-1}P_F^*V_F), gives the identity.  Conversely, singularity
of (M_F) is equivalent to a nonzero kernel of (H_F) by
\Cref{thm:newS7V36-BS}; self-adjoint Fredholmness then precludes
invertibility.  Finally,
(|H_0^{-1}|\le\delta_0^{-1}), (|P_F|=1), and
(|M_F^{-1}|=\sigma_{\min}(M_F)^{-1}) give
\eqref{eq:newS7V36-inverse-bound}.
\end{proof}

\begin{definition}[Face Birman--Schwinger card, FBSC]
\label{def:newS7V36-FBSC}
For a finite family of admissible references with common gap (delta_0),
FBSC is the pair of uniform bounds
\begin{equation}
 \inf_{F,E,k_\parallel}\sigma_{\min}(M_F)\ge\mu_*>0,
 qquad
 \sup_{F,E,k_\parallel}\|V_F\|\le v_*<\infty.        \label{eq:newS7V36-FBSC}
\end{equation}
\end{definition}

\begin{corollary}[FBSC closes the face gap]
\label{cor:newS7V36-FBSC-gap}
FBSC implies
\begin{equation}
 \|H_Fu\|\ge
 \left({1\over\delta_0}+{v_*\over\delta_0^2\mu_*}\right)^{-1}
 \|u\|                                                  \label{eq:newS7V36-FBSC-gap}
\end{equation}
uniformly in all face parameters.  Hence FBSC implies C1FGC and the
codimension-one row of SFGC.
\end{corollary}

\begin{proof}
Use \Cref{thm:newS7V36-Woodbury} and the spectral characterization of the
gap of an invertible self-adjoint operator.  The implication to C1FGC is
the compactness argument in \Cref{prop:newS7V35-IRNC-C1FGC}.
\end{proof}

\begin{proposition}[IRNC is the Neumann subregion of FBSC]
\label{prop:newS7V36-IRNC-subregion}
If (|V_F|<delta_0), then
\begin{equation}
 |G_FV_F|\le{|V_F|\overdelta_0}<1,qquad
 sigma_{\min}(M_F)\ge1-{|V_F|\overdelta_0}.        \label{eq:newS7V36-Neumann}
\end{equation}
Thus the V35 relative-norm criterion is exactly a readily checked
Neumann-series subregion of the present finite determinant problem.
\end{proposition}

\begin{proof}
Compression cannot increase norm, so
(|G_F|\le|H_0^{-1}|\ledelta_0^{-1}).  The reverse triangle inequality
for (M_Fu=u+G_FV_Fu) gives \eqref{eq:newS7V36-Neumann}.
\end{proof}

\subsection{A self-adjoint finite matrix when the defect has fixed rank}
\label{sec:newS7V36-selfadjoint}

Although (M_F=I+G_FV_F) need not be self-adjoint, the zero problem has a
self-adjoint form.  On a parameter region where the rank of (V_F) is
constant, discard its zero eigenspace and factor
\begin{equation}
 P_F^*V_FP_F=Q_F^*J_FQ_F,                             \label{eq:newS7V36-sign-factor}
\end{equation}
where (J_F=J_F^*=J_F^{-1}) is the signature matrix and (Q_F) maps onto
the rank space.

\begin{proposition}[Self-adjoint interface Hessian]
\label{prop:newS7V36-interface-Hessian}
With
\begin{equation}
 B_F:=J_F+Q_FH_0^{-1}Q_F^*,                           \label{eq:newS7V36-B}
\end{equation}
the matrix (B_F) is self-adjoint and
\begin{equation}
 ker H_F\cong\ker B_F.                               \label{eq:newS7V36-B-kernel}
\end{equation}
In particular, the signed eigenvalues of one finite Hermitian matrix can
be followed continuously until the defect rank changes.
\end{proposition}

\begin{proof}
Self-adjointness is immediate.  If (H_F\psi=0), put
(y=J_FQ_F\psi).  Then
(psi=-H_0^{-1}Q_F^*y) and
\[
 Q_FH_0^{-1}Q_F^*y=-Q_F\psi=-J_Fy,
\]
so (B_Fy=0).  Conversely, if (B_Fy=0), set
(psi=-H_0^{-1}Q_F^*y).  Then
(Q_F\psi=J_Fy), whence
((H_0+Q_F^*J_FQ_F)\psi=0).  Neither map kills a nonzero vector, and they
are inverse after the stated identifications.
\end{proof}

The Hermitian form is useful because a possible crossing is now an actual
eigenvalue crossing through zero.  It also makes clear why entrywise
absolute-value estimates may lose the sign correlation which determines
the interface mode.

\subsection{Finite acquisition target}
\label{sec:newS7V36-target}

For each fixed original face, all matrices in
\eqref{eq:newS7V36-M} are finite.  The remaining analytical variable is
(k_\parallel\in\mathbb T^3), and the coframe lies in the compact class
(mathscr E).  A proof of nonvanishing can therefore proceed by one of
three rigorous routes:
\begin{enumerate}
\item exact conjugacy to a transparent reference;
\item the V35/V36 norm bound (|G_FV_F|<1);
\item a direct lower bound for (sigma_{\min}(M_F)), or equivalently for
the absolute eigenvalues of (B_F) on constant-rank regions.
\end{enumerate}
These routes exhaust the face problem once an admissible reference has
been constructed.  Constructing such a reference with the correct two
tails is a separate acquisition row when the adjacent bulk stencils are
not unitarily conjugate.

\begin{definition}[Two-tail reference card, TTRC]
\label{def:newS7V36-TTRC}
TTRC is the construction, for every adjacent pair of V33 bulk stencils, of
one self-adjoint finite-range full-line interface (H_0) which has those
two stencils as its tails and satisfies the uniform gap
\eqref{eq:newS7V36-H0gap}.
\end{definition}

\begin{assessment}[Exact reduction, two finite acquisitions]
\label{ass:newS7V36-status}
V36 replaces the infinite surface-mode equation by the exact finite matrix
(M_F).  The remaining codimension-one work is now separated cleanly into
TTRC, which supplies a gapped comparison interface, and FBSC, which proves
that the actual finite interface defect does not close its gap.  Neither
card is assumed proved for arbitrary adjacent affine simplex types.
\end{assessment}

\begin{programme}[V37: acquire the two-tail reference]
\label{prog:newS7V36-V37}
V37 will compare the gapped bulk symbols on the two sides.  It will prove
that a uniformly gapped self-adjoint homotopy of bulk symbols produces an
admissible slowly varying two-tail reference, and will identify the
topological obstruction when no such homotopy exists.  This goes upstream
of FBSC: without TTRC, the finite determinant cannot yet cover every face.
\end{programme}

\begin{firewall}[Claim boundary after seventh-series V36]
\label{fire:newS7V36-boundary}
V36 proves an exact finite Birman--Schwinger principle for every face with
an admissible gapped reference, proves interface-window and zero-set
invariance, relates the finite determinant to the V34 Evans determinant,
derives a resolvent formula and quantitative gap, and separates the
remaining work into TTRC and FBSC.

V36 does not prove TTRC or FBSC for all V33 face types, exclude every
nonperturbative surface mode, prove higher-codimension tangent gaps, full
SFGC, the global finite-element overlap gap, strong local heat consistency,
Gaussian heat locality, the common Einstein stress, pure gravitational
owners, rough metric or AOP control, measure concentration, the
determinant-line phase, interacting cutoff removal, or the Standard
Model--Einstein continuum.
\end{firewall}

\begin{theorem}[Seventh-series V36 result]
\label{thm:newS7V36-result}
For every frozen face admitting a uniformly gapped reference with the same
tails, the existence and multiplicity of a square-summable zero mode equal
the nullity of the explicit finite matrix
(M_F=I+P_FH_0^{-1}P_F^*V_F).  A uniform lower singular-value bound for
this matrix gives the quantitative face gap
\eqref{eq:newS7V36-FBSC-gap}.
\end{theorem}

\begin{proof}
Use \Cref{thm:newS7V36-BS,
thm:newS7V36-Evans-BS,
thm:newS7V36-Woodbury,
cor:newS7V36-FBSC-gap,
prop:newS7V36-interface-Hessian}.
\end{proof}

\paragraph{Parent-version record.}
V36 was formed from
\texttt{Renewal\_SM\_Einstein\_Continuum\_Research\_Notes\_V35.tex}.
The V35 parent had (14,299) logical source lines, (578,958) bytes,
and SHA--256
\[
 \texttt{E4AFADA548B340A7A05A0D1D862AD4F47D45FA10D1D38CD15E46DEE221D83F72}.
\]
The next compulsory companion audit is V40.

% =============================================================================
% END APPENDED SEVENTH-SERIES V36 MATERIAL
% =============================================================================

% =============================================================================
% BEGIN APPENDED SEVENTH-SERIES V37 MATERIAL
% =============================================================================

\section{Seventh-series V37: adiabatic acquisition of a two-tail reference}
\label{sec:v7v37}
\label{sec:newS7V37}

\subsection{Why the construction is upstream of FBSC}
\label{sec:newS7V37-purpose}

The Birman--Schwinger reduction of V36 is exact once a gapped comparison
operator with the correct two tails is known.  When the left and right
bulk stencils are not exactly unitarily conjugate, a sharp transparent seam
need not exist.  The present note proves that a uniformly gapped path of
bulk symbols is enough: spreading that path across a sufficiently wide but
finite normal collar produces the required two-tail reference.

Fix the tangential parameter
\(\zeta=(E,k_\parallel)\) in a compact set \(\mathscr Z\).  Let
\(s\in[0,1]\) parameterize a path of one-dimensional finite-range Laurent
operators
\begin{equation}
 (A_s(\zeta)u)_m
 =\sum_{j=-J}^J h_j(s,\zeta)u_{m+j},\qquad
 h_{-j}(s,\zeta)=h_j(s,\zeta)^*.                     \label{eq:newS7V37-Laurent}
\end{equation}
The fibre dimension and range \(J\) are fixed.

\begin{hypothesis}[Uniform gapped bulk path, UGBP]
\label{hyp:newS7V37-UGBP}
There are constants \(\delta>0\), \(K_0<\infty\), and \(K_1<\infty\) such
that, for all \(s,\zeta\),
\begin{align}
 \|A_s(\zeta)u\|&\ge\delta\|u\|,                       \label{eq:newS7V37-path-gap}\\
 \sum_{j=-J}^J\|h_j(s,\zeta)\|&\le K_0,               \label{eq:newS7V37-path-size}\\
 \sum_{j=-J}^J\|h_j(s,\zeta)-h_j(t,\zeta)\|
 &\le K_1|s-t|.                                       \label{eq:newS7V37-path-Lipschitz}
\end{align}
All constants are independent of \(\zeta\).
\end{hypothesis}

Choose a Lipschitz function
\(\theta:\mathbb R\to[0,1]\) satisfying
\begin{equation}
 \theta(x)=0\ (x\le-1),\qquad
 \theta(x)=1\ (x\ge1),\qquad
 \operatorname{Lip}(\theta)\le C_\theta.              \label{eq:newS7V37-theta}
\end{equation}
For an integer collar width \(L\ge1\), set
\(s_x=\theta(x/L)\) for half-integers as well as integers and define the
midpoint-quantized interface
\begin{equation}
 (H_L(\zeta)u)_m
 :=\sum_{j=-J}^J h_j(s_{m+j/2},\zeta)u_{m+j}.          \label{eq:newS7V37-HL}
\end{equation}

\begin{lemma}[Self-adjointness and exact tails]
\label{lem:newS7V37-selfadjoint}
The operator \(H_L(\zeta)\) is bounded and self-adjoint.  It equals the
\(s=0\) Laurent operator on all matrix entries whose two endpoints lie to
the left of \(-L\), and equals the \(s=1\) Laurent operator on all entries
whose endpoints lie to the right of \(L\).  Hence it differs from its two
tails only in finitely many normal layers.
\end{lemma}

\begin{proof}
The matrix entry from \(m\) to \(m+j\) is
\(h_j(s_{m+j/2},\zeta)\).  The reverse entry is
\[
 h_{-j}(s_{m+j-j/2},\zeta)
 =h_j(s_{m+j/2},\zeta)^*,
\]
so the matrix is self-adjoint.  Its norm is bounded by the finite row sum
in \eqref{eq:newS7V37-path-size}.  The tail statements follow from
\eqref{eq:newS7V37-theta} because the midpoint of every indicated edge is
in the constant region.
\end{proof}

\subsection{A discrete localization estimate}
\label{sec:newS7V37-localization}

For an integer scale \(\ell\ge2J+1\), choose a quadratic partition of unity
\(\{\chi_p\}_{p\in\mathbb Z}\) on \(\mathbb Z\) with
\begin{equation}
 \sum_p\chi_p(m)^2=1,\qquad
 \operatorname{diam}(\operatorname{supp}\chi_p)\le C_0\ell,\qquad
 \sum_p|\chi_p(m)-\chi_p(n)|^2
 \le {C_1|m-n|^2\over\ell^2}.                         \label{eq:newS7V37-partition}
\end{equation}
Such a partition is obtained by translating a piecewise-linear tent and
normalizing the locally finite square sum; the constants are numerical.

\begin{lemma}[Square-summed commutator bound]
\label{lem:newS7V37-commutator}
There is \(C_{\rm com}=C_{\rm com}(J,K_0,C_1)\) such that
\begin{equation}
 \left(
  \sum_p\|[H_L,\chi_p]u\|^2
 \right)^{1/2}
 \le {C_{\rm com}\over\ell}\|u\|                    \label{eq:newS7V37-commutator}
\end{equation}
uniformly in \(L,\zeta\).
\end{lemma}

\begin{proof}
The \(m\)-th component of the commutator is
\[
 \sum_{|j|\le J}h_j(s_{m+j/2},\zeta)
 \{\chi_p(m+j)-\chi_p(m)\}u_{m+j}.
\]
Cauchy--Schwarz over the finite \(j\)-sum, followed by the last bound in
\eqref{eq:newS7V37-partition}, gives a factor at most
\(C J/\ell\).  Summing in \(m,p\), translating \(m+j\) back to \(m\), and
using \eqref{eq:newS7V37-path-size} proves the stated bound.  All overlap
multiplicities are controlled by \(C_0,C_1\) and are independent of the
scales.
\end{proof}

Choose one integer \(m_p\) in each nonempty support of \(\chi_p\) and put
\(s_p=s_{m_p}\).

\begin{lemma}[Square-summed freezing error]
\label{lem:newS7V37-freezing}
There is \(C_{\rm fr}=C_{\rm fr}(J,K_1,C_0,C_\theta)\) such that
\begin{equation}
 \left(
  \sum_p\|(H_L-A_{s_p})\chi_pu\|^2
 \right)^{1/2}
 \le C_{\rm fr}{\ell\over L}\|u\|.                  \label{eq:newS7V37-freezing}
\end{equation}
\end{lemma}

\begin{proof}
If either endpoint of a matrix entry meets
\(\operatorname{supp}\chi_p\), its midpoint is within
\(C_0\ell+J\) of \(m_p\).  Equations
\eqref{eq:newS7V37-path-Lipschitz} and
\eqref{eq:newS7V37-theta} therefore bound the corresponding row sum of
\(H_L-A_{s_p}\) by \(C K_1C_\theta\ell/L\).  The block Schur test and the
uniformly bounded overlap of the supports give
\eqref{eq:newS7V37-freezing} after summing in \(p\).
\end{proof}

\subsection{The adiabatic two-tail gap}
\label{sec:newS7V37-gap}

\begin{theorem}[Discrete adiabatic gluing theorem]
\label{thm:newS7V37-adiabatic}
Under UGBP,
\begin{equation}
 \|H_L(\zeta)u\|
 \ge
 \left(
  \delta-{C_{\rm com}\over\ell}
  -C_{\rm fr}{\ell\over L}
 \right)\|u\|                                          \label{eq:newS7V37-two-scale-gap}
\end{equation}
for every \(\ell\ge2J+1\).  In particular, there are constants \(C_*>0\)
and \(L_0<\infty\) such that
\begin{equation}
 \|H_L(\zeta)u\|
 \ge(\delta-C_*L^{-1/2})\|u\|                        \label{eq:newS7V37-adiabatic-gap}
\end{equation}
for \(L\ge L_0\), uniformly in \(\zeta\).  For sufficiently large \(L\),
\(H_L\) is therefore an admissible two-tail reference with gap at least
\(\delta/2\).
\end{theorem}

\begin{proof}
For each \(p\), the frozen gap gives
\[
 \delta\|\chi_pu\|
 \le\|A_{s_p}\chi_pu\|
 \le\|\chi_pH_Lu\|
    +\|[H_L,\chi_p]u\|
    +\|(A_{s_p}-H_L)\chi_pu\|.
\]
Take the \(\ell^2\)-norm of this inequality in \(p\) and use Minkowski's
inequality.  The partition identity gives
\[
 \sum_p\|\chi_pu\|^2=\|u\|^2,\qquad
 \sum_p\|\chi_pH_Lu\|^2=\|H_Lu\|^2.
\]
Apply \Cref{lem:newS7V37-commutator,lem:newS7V37-freezing} and rearrange to
obtain \eqref{eq:newS7V37-two-scale-gap}.  Taking
\(\ell\) to be an integer comparable with \(L^{1/2}\) gives
\eqref{eq:newS7V37-adiabatic-gap}.  Enlarge \(L_0\) so that the error is at
most \(\delta/2\), and use \Cref{lem:newS7V37-selfadjoint} for the exact-tail
condition.
\end{proof}

\begin{corollary}[UGBP implies TTRC]
\label{cor:newS7V37-UGBP-TTRC}
Every pair of V33 bulk stencils joined by a UGBP path satisfies the V36
two-tail reference card TTRC.
\end{corollary}

\begin{proof}
Use \(H_L\) with \(L\) large enough in
\Cref{thm:newS7V37-adiabatic}.  It is self-adjoint, uniformly gapped, and
has the prescribed two tails.
\end{proof}

\subsection{Bulk paths already supplied by V34}
\label{sec:newS7V37-existing-paths}

For a fixed regular interior combinatorial type \(\tau\), V34 proves a
uniform bulk gap over the compact coframe class \(\mathscr E\), provided the
common Wilson parameter satisfies the V34 large-\(r\) choice.  This becomes
a path statement on each path-connected part of \(\mathscr E\).

\begin{proposition}[Same-type adjacent bulks satisfy TTRC]
\label{prop:newS7V37-same-type}
Suppose the two sides of a face have, after the exact spin lift and an
incidence-preserving unitary relabelling, the same regular interior type
\(\tau\).  Suppose their coframes \(E_-\) and \(E_+\) are joined by a
Lipschitz path \(E_s\in\mathscr E\).  Then their Laurent coefficients form
a UGBP path, and the pair satisfies TTRC.
\end{proposition}

\begin{proof}
Use \(h_j(s,k_\parallel)=h_{j,\tau}(E_s,k_\parallel)\).  The V34 bulk
theorem supplies one positive gap along the compact path.  The explicit
finite Laurent formula of V34 is smooth in the coframe and momentum, so
compactness supplies the uniform size and Lipschitz constants in UGBP.
Apply \Cref{cor:newS7V37-UGBP-TTRC}.  Unitary spin and incidence changes do
not alter the gap.
\end{proof}

This proposition covers coframe variation within a common stencil type.
It does not prove that the two Freudenthal orderings induced on every
original face are incidence-preservingly equivalent.  That is a finite
combinatorial question and remains to be checked from the V33 assembly.

\subsection{The obstruction to a bulk path}
\label{sec:newS7V37-obstruction}

The adiabatic theorem cannot connect two bulk phases through a path which
must close its gap.  This obstruction can be stated without importing a
classification theorem.

Let the full four-momentum Bloch symbol of a bulk endpoint be
\(\widehat A_s(k)\), \(k\in\mathbb T^4\), and define its negative spectral
projection
\begin{equation}
 \Pi_s(k):=\mathbf 1_{(-\infty,0)}(\widehat A_s(k)).   \label{eq:newS7V37-negative-projection}
\end{equation}

\begin{proposition}[Endpoint negative bundles are a necessary invariant]
\label{prop:newS7V37-bundle-obstruction}
Along a continuous uniformly gapped bulk path, the ranges of
\(\Pi_0\) and \(\Pi_1\) are isomorphic complex vector bundles over
\(\mathbb T^4\).  Consequently every Chern number formed from these bundles
is the same at the two endpoints.  If any such invariant differs, UGBP is
impossible.
\end{proposition}

\begin{proof}
Uniform separation from zero permits the Riesz formula
\[
 \Pi_s(k)={1\over2\pi i}\int_\Gamma
   (z-\widehat A_s(k))^{-1}\,dz,
\]
where one contour \(\Gamma\) surrounds the negative spectrum for every
\(s,k\).  Hence \(\Pi_s(k)\) is continuous in \((s,k)\), and smooth in \(k\)
when the coefficients are smooth.  Its ranges form a vector bundle over
\([0,1]\times\mathbb T^4\).  Restriction to the two ends of a product with
an interval gives isomorphic bundles.  Chern numbers are invariant under
bundle isomorphism, proving the obstruction statement.
\end{proof}

\begin{definition}[Bulk path card, BPC]
\label{def:newS7V37-BPC}
BPC asserts that every pair of adjacent V33 regular bulk types has equal
negative-bundle invariants and admits an explicit UGBP path after the exact
spin and incidence identifications.
\end{definition}

\begin{assessment}[TTRC reduced to BPC]
\label{ass:newS7V37-status}
The analytic part of the two-tail reference problem is closed: UGBP
produces a reference with a quantitative gap.  For same-type adjacent
bulks connected within the V34 coframe class, TTRC follows.  For genuinely
different Freudenthal types, BPC remains a finite-symbol acquisition.  No
equality of their negative bundles or existence of a gapped path is
assumed here.
\end{assessment}

\begin{programme}[V38: compare the finite bulk types]
\label{prog:newS7V37-V38}
V38 will use the V34 scalar-stiffness form of the Wilson symbol to construct
a candidate path between different regular Freudenthal types.  The first
route will hold a large common positive Wilson stiffness while interpolating
the Dirac and stiffness coefficients.  The proof must prevent the physical
mode at \(k=0\) and all nonphysical zeros from crossing zero throughout the
path.  If convex interpolation is unsafe, the note will isolate the exact
joint zero-separation card rather than assume it.
\end{programme}

\begin{firewall}[Claim boundary after seventh-series V37]
\label{fire:newS7V37-boundary}
V37 proves a discrete adiabatic gluing theorem with an
\(O(L^{-1/2})\) gap loss, proves UGBP sufficient for the two-tail reference
card, closes TTRC for same-type coframe-connected bulks, and identifies the
negative spectral bundle as a necessary invariant of any remaining path.

V37 does not prove BPC for distinct regular types, TTRC for every V33 face,
FBSC, exclusion of every surface mode, higher-codimension tangent gaps,
full SFGC, the global finite-element overlap gap, strong local heat
consistency, Gaussian heat locality, the common Einstein stress, pure
gravitational owners, rough metric or AOP control, measure concentration,
the determinant-line phase, interacting cutoff removal, or the Standard
Model--Einstein continuum.
\end{firewall}

\begin{theorem}[Seventh-series V37 result]
\label{thm:newS7V37-result}
A uniformly gapped Lipschitz path between two finite-range bulk Laurent
symbols produces a self-adjoint finite-collar interface with the prescribed
tails and gap at least \(\delta-C_*L^{-1/2}\).  Thus BPC implies TTRC, and
TTRC is already proved for adjacent bulks of the same regular type whose
coframes are path-connected inside the V34 compact class.
\end{theorem}

\begin{proof}
Combine \Cref{lem:newS7V37-selfadjoint,
lem:newS7V37-commutator,
lem:newS7V37-freezing,
thm:newS7V37-adiabatic,
cor:newS7V37-UGBP-TTRC,
prop:newS7V37-same-type,
prop:newS7V37-bundle-obstruction}.
\end{proof}

\paragraph{Parent-version record.}
V37 was formed from
\texttt{Renewal\_SM\_Einstein\_Continuum\_Research\_Notes\_V36.tex}.
The V36 parent had \(14\,684\) logical source lines, \(593\,205\) bytes,
and SHA--256
\[
 \texttt{EE42AA4EAA7A21B4D5860AAEF80A25E2E257829E28557A0D7AD4D000241978F9}.
\]
The next compulsory companion audit is V40.

% =============================================================================
% END APPENDED SEVENTH-SERIES V37 MATERIAL
% =============================================================================

% =============================================================================
% BEGIN APPENDED SEVENTH-SERIES V38 MATERIAL
% =============================================================================

\section{Seventh-series V38: a universal gapped bulk path}
\label{sec:v7v38}
\label{sec:newS7V38}

\subsection{The joint infrared--stiffness mechanism}
\label{sec:newS7V38-mechanism}

V37 reduced the two-tail reference card to a gapped path between the finite
set of regular bulk symbols.  Direct convex interpolation of two Hermitian
symbols is unsafe.  The special scalar-stiffness square from V34 supplies a
different mechanism: the Wilson mass can vanish only on a small stiffness
level set, and a sufficiently large common Wilson parameter moves that
level set into the region where the physical Dirac slope is uniformly
injective.

Let \(k\in\mathbb T^4=(-\pi,\pi]^4\), and put
\begin{equation}
 q(k):=2\sum_{\alpha=1}^4(1-\cos k_\alpha).            \label{eq:newS7V38-q}
\end{equation}
Consider a compact family, indexed by \(t\in[0,1]\) and auxiliary
parameters \(\xi\in\Xi\), of Hermitian Wilson symbols whose exact square is
\begin{equation}
 \mathcal H_{t,\xi}(k)^2
 =
 |s_{t,\xi}(k)|^2
 +
 \left\{{r\over2}\ell_{t,\xi}(k)-\rho\right\}^2.       \label{eq:newS7V38-square}
\end{equation}

\begin{hypothesis}[Joint infrared--stiffness data, JISD]
\label{hyp:newS7V38-JISD}
There are \(c_\ell,b_0,C_2>0\) such that, uniformly in \(t,\xi,k\),
\begin{align}
 \ell_{t,\xi}(k)&\ge c_\ell q(k),                      \label{eq:newS7V38-stiff-coercive}\\
 s_{t,\xi}(0)&=0,                                     \label{eq:newS7V38-s-zero}\\
 s_{t,\xi}(k)&=B_{t,\xi}k+R_{t,\xi}(k),               \label{eq:newS7V38-expansion}\\
 \sigma_{\min}(B_{t,\xi})&\ge b_0,\qquad
 |R_{t,\xi}(k)|\le C_2|k|^2
 \quad\hbox{for }|k|\le1.                             \label{eq:newS7V38-remainder}
\end{align}
We also require \(\ell_{t,\xi}(0)=0\).
\end{hypothesis}

\begin{lemma}[Uniform physical-cone injectivity]
\label{lem:newS7V38-cone}
Under JISD, if
\begin{equation}
 \varepsilon:=\min\left\{1,{b_0\over2C_2}\right\},     \label{eq:newS7V38-epsilon}
\end{equation}
then
\begin{equation}
 |s_{t,\xi}(k)|\ge {b_0\over2}|k|                    \label{eq:newS7V38-cone-bound}
\end{equation}
for \(0<|k|\le\varepsilon\).
\end{lemma}

\begin{proof}
Equations \eqref{eq:newS7V38-expansion}--\eqref{eq:newS7V38-remainder}
give
\[
 |s_{t,\xi}(k)|
 \ge b_0|k|-C_2|k|^2
 \ge {b_0\over2}|k|.
\]
\end{proof}

Define
\begin{equation}
 q_\varepsilon:=\min_{\varepsilon\le|k|\le2\pi}q(k)>0,              \label{eq:newS7V38-qeps}
\end{equation}
where the upper radius merely contains the chosen torus representative.

\begin{theorem}[Large-\(r\) universal path gap]
\label{thm:newS7V38-large-r-path}
Assume JISD and choose
\begin{equation}
 r>{2\rho\over c_\ell q_\varepsilon}.                 \label{eq:newS7V38-r-threshold}
\end{equation}
Then \(\mathcal H_{t,\xi}(k)\) has no zero for any \(t,\xi,k\).
Consequently compactness gives
\begin{equation}
 \|\mathcal H_{t,\xi}(k)v\|
 \ge\delta_{\rm path}\|v\|                            \label{eq:newS7V38-path-gap}
\end{equation}
with one \(\delta_{\rm path}>0\).
\end{theorem}

\begin{proof}
A zero of \eqref{eq:newS7V38-square} would require
\begin{equation}
 s_{t,\xi}(k)=0,\qquad
 \ell_{t,\xi}(k)={2\rho\over r}.                      \label{eq:newS7V38-zero-system}
\end{equation}
The stiffness inequality and \eqref{eq:newS7V38-r-threshold} imply
\[
 q(k)\le {2\rho\over r c_\ell}<q_\varepsilon,
\]
so \(|k|<\varepsilon\) on the chosen torus representative.  The point
\(k=0\) cannot satisfy the second equation in
\eqref{eq:newS7V38-zero-system}, while every nonzero point in this ball is
excluded by \Cref{lem:newS7V38-cone}.  Thus the square is everywhere
positive.  It is continuous on the compact parameter space, so its square
root has a positive minimum.
\end{proof}

\subsection{Coercivity of every regular stiffness symbol}
\label{sec:newS7V38-stiffness}

For a regular V33 bulk type \(\tau\), the mass-lumped connection stiffness
has the scalar Bloch form
\begin{equation}
 \ell_{\tau,E}(k)
 =\sum_{d\in\mathcal D_\tau}w_{\tau,E}(d)
   |e^{ik\cdot d}-1|^2,                               \label{eq:newS7V38-edge-stiffness}
\end{equation}
where the finite edge set \(\mathcal D_\tau\subset\mathbb Z^4\) generates
the Bravais lattice and the weights are strictly positive.  Shape
regularity and compactness of \(\mathscr E\) bound every required weight
away from zero and infinity.

\begin{lemma}[Uniform stiffness comparison]
\label{lem:newS7V38-stiffness-comparison}
For the finite type set and compact coframe class, there is \(c_*>0\) such
that
\begin{equation}
 \ell_{\tau,E}(k)\ge c_*q(k)                          \label{eq:newS7V38-uniform-stiffness}
\end{equation}
for every \(\tau,E,k\).
\end{lemma}

\begin{proof}
If \(\ell_{\tau,E}(k)=0\), every summand in
\eqref{eq:newS7V38-edge-stiffness} vanishes.  Hence
\(k\cdot d\in2\pi\mathbb Z\) for every lattice generator \(d\), which means
\(k=0\) in \(\mathbb T^4\).  Near zero,
\[
 \ell_{\tau,E}(k)
 =\sum_dw_{\tau,E}(d)(k\cdot d)^2+O(|k|^4).
\]
The quadratic form is positive definite because the edge directions span
\(\mathbb R^4\), uniformly so over the finite type set and compact
shape-regular parameter class.  Thus \(\ell/q\) has a positive lower limit
at zero.  Away from zero it is a positive continuous function on a compact
set.  Taking the minimum over the finite types proves the claim.
\end{proof}

\subsection{Canonical deformation of a Dirac--stiffness pair}
\label{sec:newS7V38-canonical}

The V33 weak Dirac consistency, applied to a translation-periodic frozen
bulk, gives
\begin{equation}
 s_{\tau,E}(k)=B_{\tau,E}k+O(|k|^2),                  \label{eq:newS7V38-principal}
\end{equation}
where \(B_{\tau,E}\) is the invertible physical coframe map.  Shape
regularity gives a uniform lower singular-value bound and the finite
trigonometric formulas give uniform second-derivative bounds.

Let
\begin{equation}
 s^{\rm can}_{B}(k):=B(\sin k_1,\ldots,\sin k_4)^T.    \label{eq:newS7V38-canonical-Dirac}
\end{equation}

\begin{definition}[Type-to-canonical path]
\label{def:newS7V38-type-path}
For \(u\in[0,1]\), set
\begin{align}
 s_u(k)&=(1-u)s_{\tau,E}(k)+u s^{\rm can}_{B_{\tau,E}}(k),            \label{eq:newS7V38-s-path}\\
 \ell_u(k)&=(1-u)\ell_{\tau,E}(k)+u q(k).             \label{eq:newS7V38-l-path}
\end{align}
\end{definition}

\begin{lemma}[The type-to-canonical path satisfies JISD]
\label{lem:newS7V38-type-JISD}
The family \eqref{eq:newS7V38-s-path}--\eqref{eq:newS7V38-l-path}
satisfies JISD uniformly over all regular types and allowed coframes.
\end{lemma}

\begin{proof}
The stiffness assertion follows from
\Cref{lem:newS7V38-stiffness-comparison}:
\[
 \ell_u(k)\ge\min\{c_*,1\}q(k).
\]
Both Dirac symbols vanish at zero and have the same derivative
\(B_{\tau,E}\) there.  Their convex combination therefore has that same
derivative.  The uniform lower singular-value and second-remainder bounds
follow from shape regularity, compactness, and finiteness of the
trigonometric stencils.
\end{proof}

\subsection{Removal of the coframe map}
\label{sec:newS7V38-coframe}

Use the polar decomposition
\begin{equation}
 B_{\tau,E}=R_{\tau,E}S_{\tau,E},\qquad
 R_{\tau,E}\in SO(4),\quad S_{\tau,E}>0.              \label{eq:newS7V38-polar}
\end{equation}
The orientation is fixed by the spin structure.  The orthogonal factor is
an exact Clifford-frame change: its spin lift conjugates
\(\gamma\cdot B_{\tau,E}v\) to
\(\gamma\cdot S_{\tau,E}v\).  This unitary identification does not change
the bulk gap or its negative spectral bundle.

\begin{definition}[Canonical-coframe path]
\label{def:newS7V38-coframe-path}
After removing \(R_{\tau,E}\), put
\begin{equation}
 S_u=(1-u)S_{\tau,E}+uI,\qquad
 s_u(k)=S_u(\sin k_1,\ldots,\sin k_4)^T,\qquad
 \ell_u(k)=q(k).                                      \label{eq:newS7V38-S-path}
\end{equation}
\end{definition}

\begin{lemma}[The canonical-coframe path satisfies JISD]
\label{lem:newS7V38-coframe-JISD}
The path \eqref{eq:newS7V38-S-path} satisfies JISD uniformly over the
compact coframe class.
\end{lemma}

\begin{proof}
Convex combinations of positive-definite matrices are positive definite,
and
\[
 \sigma_{\min}(S_u)
 \ge(1-u)\sigma_{\min}(S_{\tau,E})+u.
\]
Compact shape regularity supplies one positive lower bound.  The sine
remainder is uniformly quadratic near zero, and the stiffness equality
gives \eqref{eq:newS7V38-stiff-coercive} with \(c_\ell=1\).
\end{proof}

At \(u=1\), every regular type reaches the same universal pair
\begin{equation}
 s_{\rm univ}(k)=(\sin k_1,\ldots,\sin k_4)^T,\qquad
 \ell_{\rm univ}(k)=q(k).                             \label{eq:newS7V38-universal}
\end{equation}

\subsection{Closure of the bulk path and two-tail cards}
\label{sec:newS7V38-closure}

\begin{theorem}[Universal bulk path theorem]
\label{thm:newS7V38-universal-path}
After exact spin-frame and lattice-incidence identifications, any two
regular V33 bulk symbols in the same oriented spin sector are connected by
a finite-range Lipschitz UGBP path.  One common Wilson parameter \(r\),
chosen larger than the maximum of the V34 and
\eqref{eq:newS7V38-r-threshold} thresholds, works for the finite type set
and compact coframe class.
\end{theorem}

\begin{proof}
Deform the first endpoint to its canonical Dirac symbol using
\Cref{def:newS7V38-type-path}, remove its orthogonal coframe factor by the
exact spin unitary, and deform its positive polar factor to the identity
using \Cref{def:newS7V38-coframe-path}.  Perform the reverse construction
for the second endpoint and concatenate the paths through
\eqref{eq:newS7V38-universal}.  Lemmas
\ref{lem:newS7V38-type-JISD} and
\ref{lem:newS7V38-coframe-JISD} supply uniform JISD data on every segment.
The finite concatenation also has uniform JISD data after a harmless
Lipschitz reparameterization.  Apply
\Cref{thm:newS7V38-large-r-path}.  Each symbol remains a finite
trigonometric polynomial, so the corresponding Laurent range is finite.
\end{proof}

\begin{corollary}[BPC and TTRC for all regular faces]
\label{cor:newS7V38-BPC-TTRC}
The V37 bulk path card BPC holds for the regular V33 bulk family.
Consequently every codimension-one face whose two tails are regular bulk
types satisfies the V36 two-tail reference card TTRC.
\end{corollary}

\begin{proof}
The explicit path in \Cref{thm:newS7V38-universal-path} proves BPC.  The
discrete adiabatic gluing theorem
\Cref{thm:newS7V37-adiabatic} converts that path into an admissible
two-tail reference.
\end{proof}

\begin{corollary}[No bulk topological mismatch among the regular types]
\label{cor:newS7V38-no-bulk-mismatch}
All regular V33 bulk types have isomorphic negative spectral bundles after
the allowed spin and incidence identifications.  Thus the V37 necessary
bundle obstruction does not occur inside this family.
\end{corollary}

\begin{proof}
Use the UGBP path and
\Cref{prop:newS7V37-bundle-obstruction}.
\end{proof}

\subsection{Remaining codimension-one acquisition}
\label{sec:newS7V38-remaining}

The reference produced by V37 may use a wide collar, but its difference
from the actual sharp V33 face still has finite normal support.  Therefore
the exact V36 Birman--Schwinger matrix exists for every regular face.  The
only remaining face row is its uniform nonvanishing.

\begin{assessment}[TTRC closed; FBSC remains]
\label{ass:newS7V38-status}
V38 closes BPC and TTRC for the regular V33 family.  It does not prove that
the actual sharp face equals the adiabatic reference or that the finite
Birman--Schwinger matrix stays invertible.  Accidental, non-topologically
protected surface zero modes remain possible and are exactly the FBSC
problem.
\end{assessment}

\begin{programme}[V39: algebraic certificate for FBSC]
\label{prog:newS7V38-V39}
V39 will expose the parameter dependence of the finite matrix
\(M_F(E,k_\parallel)\).  It will clear the gapped reference denominators,
convert the zero condition into a finite Laurent-polynomial determinant,
and formulate a checkable positive certificate on the tangential torus and
compact coframe cells.  This is the correct finite calculation before any
higher-codimension tangent model is attempted.
\end{programme}

\begin{firewall}[Claim boundary after seventh-series V38]
\label{fire:newS7V38-boundary}
V38 proves a general large-\(r\) path-gap theorem, uniform stiffness
coercivity for the regular edgewise stencils, a type-to-canonical
deformation, a positive-polar coframe deformation, and hence BPC and TTRC
for all regular V33 bulk pairs.

V38 does not prove FBSC, exclude accidental sharp-face zero modes, prove
higher-codimension tangent gaps, full SFGC, the global finite-element
overlap gap, strong local heat consistency, Gaussian heat locality, the
common Einstein stress, pure gravitational owners, rough metric or AOP
control, measure concentration, the determinant-line phase, interacting
cutoff removal, or the Standard Model--Einstein continuum.
\end{firewall}

\begin{theorem}[Seventh-series V38 result]
\label{thm:newS7V38-result}
All regular V33 bulk types in the oriented spin sector connect to the
universal pair \eqref{eq:newS7V38-universal} through one uniformly gapped
finite-range Wilson path after strengthening the common large-\(r\)
choice.  The adiabatic V37 construction therefore supplies a uniformly
gapped two-tail reference for every regular codimension-one face.
\end{theorem}

\begin{proof}
Use \Cref{lem:newS7V38-cone,
thm:newS7V38-large-r-path,
lem:newS7V38-stiffness-comparison,
lem:newS7V38-type-JISD,
lem:newS7V38-coframe-JISD,
thm:newS7V38-universal-path,
cor:newS7V38-BPC-TTRC}.
\end{proof}

\paragraph{Parent-version record.}
V38 was formed from
\texttt{Renewal\_SM\_Einstein\_Continuum\_Research\_Notes\_V37.tex}.
The V37 parent had \(15\,033\) logical source lines, \(606\,947\) bytes,
and SHA--256
\[
 \texttt{2A65FCC77AEB026758CA8F8E6BE1EA419FFC787926C2EC15239AA32A1E8210D5}.
\]
The next compulsory companion audit is V40.

% =============================================================================
% END APPENDED SEVENTH-SERIES V38 MATERIAL
% =============================================================================

% =============================================================================
% BEGIN APPENDED SEVENTH-SERIES V39 MATERIAL
% =============================================================================

\section{Seventh-series V39: certified finite-range approximation of FBSC}
\label{sec:v7v39}
\label{sec:newS7V39}

\subsection{Correction of the naive denominator-clearing route}
\label{sec:newS7V39-correction}

V38 proposed clearing the denominators of the compressed reference inverse
in order to obtain a finite Laurent determinant.  That proposal requires a
correction.  Even when a two-tail reference is finite range, eliminating
its two infinite tails generally produces matrix Weyl functions determined
by characteristic roots.  These functions need not be rational Laurent
functions of \(e^{ik_\parallel}\).  Direct denominator clearing is therefore
not justified.

The reference gap supplies a better route.  Its inverse can be approximated
in operator norm by an explicit odd polynomial.  Because the reference is
finite range, every polynomial approximant is finite range and its
compressed matrix is a finite Laurent expression.  The approximation error
is known exactly and decays geometrically.

\subsection{An explicit polynomial inverse}
\label{sec:newS7V39-polynomial}

Let \(H_0=H_{0,F}(E,k_\parallel)\) be the admissible two-tail reference
constructed through V37--V38.  In addition to its uniform gap, fix a
uniform row-sum bound
\begin{equation}
 \delta_0\|u\|\le\|H_0u\|,\qquad
 \|H_0\|\le M_0,                                      \label{eq:newS7V39-two-sided}
\end{equation}
where \(0<\delta_0\le M_0<\infty\).  Define
\begin{equation}
 \alpha:={2\over M_0^2+\delta_0^2},\qquad
 \kappa:={M_0^2-\delta_0^2\over M_0^2+\delta_0^2}<1.  \label{eq:newS7V39-alpha-kappa}
\end{equation}
For \(n\ge0\), put
\begin{equation}
 p_n(x):=\alpha x\sum_{j=0}^n(1-\alpha x^2)^j
 ={1-(1-\alpha x^2)^{n+1}\over x},                   \label{eq:newS7V39-pn}
\end{equation}
where the first expression defines the polynomial also at \(x=0\).

\begin{lemma}[Uniform inverse approximation]
\label{lem:newS7V39-inverse-approximation}
For every admissible reference satisfying
\eqref{eq:newS7V39-two-sided},
\begin{equation}
 \|H_0^{-1}-p_n(H_0)\|
 \le {\kappa^{n+1}\over\delta_0}.                     \label{eq:newS7V39-inverse-error}
\end{equation}
\end{lemma}

\begin{proof}
The spectrum of the self-adjoint \(H_0\) lies in
\[
 [-M_0,-\delta_0]\cup[\delta_0,M_0].
\]
For every spectral value \(x\) in this set,
\[
 {1\over x}-p_n(x)={(1-\alpha x^2)^{n+1}\over x}.
\]
The choice \eqref{eq:newS7V39-alpha-kappa} gives
\[
 |1-\alpha x^2|\le\kappa,\qquad |x|^{-1}\le\delta_0^{-1}.
\]
The continuous functional calculus proves
\eqref{eq:newS7V39-inverse-error}.
\end{proof}

\begin{lemma}[Finite propagation of the approximant]
\label{lem:newS7V39-propagation}
If \(H_0\) has graph propagation at most \(J_0\), then
\(p_n(H_0)\) has propagation at most \((2n+1)J_0\).  After tangential
Fourier transform, every entry of
\begin{equation}
 G_{F,n}:=P_Fp_n(H_0)P_F^*                            \label{eq:newS7V39-Gn}
\end{equation}
is a finite Laurent expression in
\(z_\alpha=e^{ik_{\parallel,\alpha}}\).
\end{lemma}

\begin{proof}
The polynomial \(p_n\) has degree \(2n+1\).  A product of \(d\)
range-\(J_0\) matrices has propagation at most \(dJ_0\).  Tangential
Fourier transformation turns each tangential shift into a Laurent
monomial; finite sums and products remain finite Laurent expressions.
\end{proof}

Any explicit local denominators already present in the V33 mass-lumped
weights may be retained for interval evaluation or cleared after their
strict positivity has been recorded.  No denominator from an infinite
tail resolvent is cleared.

\subsection{A certified surrogate for the face matrix}
\label{sec:newS7V39-surrogate}

Retain the V36 factorization
\[
 H_F=H_0+P_F^*V_FP_F
\]
and define
\begin{equation}
 M_F=I+P_FH_0^{-1}P_F^*V_F,\qquad
 M_{F,n}:=I+G_{F,n}V_F.                               \label{eq:newS7V39-Mn}
\end{equation}

\begin{theorem}[Certified face-matrix approximation]
\label{thm:newS7V39-M-error}
If \(\|V_F\|\le v_*\), then
\begin{equation}
 \|M_F-M_{F,n}\|
 \le \varepsilon_n,\qquad
 \varepsilon_n:={v_*\over\delta_0}\kappa^{n+1}.       \label{eq:newS7V39-epsilon}
\end{equation}
Consequently
\begin{equation}
 \sigma_{\min}(M_F)
 \ge\sigma_{\min}(M_{F,n})-\varepsilon_n.             \label{eq:newS7V39-singular-lower}
\end{equation}
\end{theorem}

\begin{proof}
Compression by \(P_F\) has norm one.  Thus
\[
 \|M_F-M_{F,n}\|
 \le\|H_0^{-1}-p_n(H_0)\|\,\|V_F\|,
\]
and \Cref{lem:newS7V39-inverse-approximation} gives
\eqref{eq:newS7V39-epsilon}.  The smallest singular value is
one-Lipschitz in the operator norm, which proves
\eqref{eq:newS7V39-singular-lower}.
\end{proof}

\begin{corollary}[Finite-range sufficient certificate for FBSC]
\label{cor:newS7V39-finite-certificate}
If, for some \(n\) and \(\beta_n>\varepsilon_n\),
\begin{equation}
 \inf_{F,E,k_\parallel}
 \sigma_{\min}(M_{F,n}(E,k_\parallel))
 \ge\beta_n,                                          \label{eq:newS7V39-beta}
\end{equation}
then FBSC holds with
\begin{equation}
 \mu_*=\beta_n-\varepsilon_n>0.                       \label{eq:newS7V39-mu}
\end{equation}
The V36 quantitative face gap follows with this value of \(\mu_*\).
\end{corollary}

\begin{proof}
Apply \eqref{eq:newS7V39-singular-lower} uniformly and then
\Cref{cor:newS7V36-FBSC-gap}.
\end{proof}

\begin{corollary}[Detection side of the certificate]
\label{cor:newS7V39-detection}
If the exact face matrix is singular at a parameter \(x_0\), then
\begin{equation}
 \sigma_{\min}(M_{F,n}(x_0))\le\varepsilon_n           \label{eq:newS7V39-detection}
\end{equation}
for every \(n\).  Thus a sequence of surrogate minima separated uniformly
from zero excludes an exact surface mode, while minima converging at the
certified error scale identify where an exact calculation must be refined.
\end{corollary}

\begin{proof}
Use the reverse singular-value perturbation inequality with
\(\sigma_{\min}(M_F(x_0))=0\).
\end{proof}

\subsection{A finite interval certificate on parameter cells}
\label{sec:newS7V39-cells}

Let \(\mathscr D_F\) be a compact coordinate domain for the allowed coframe
cell and tangential torus, using finitely many angular charts for the
torus.  Since \(M_{F,n}\) is a finite expression, its derivative has a
computable bound
\begin{equation}
 \|M_{F,n}(x)-M_{F,n}(y)\|\le L_{F,n}|x-y|.            \label{eq:newS7V39-Lipschitz}
\end{equation}

\begin{proposition}[Cellwise lower bound]
\label{prop:newS7V39-cell-bound}
Let \(C\subset\mathscr D_F\) have centre \(x_C\) and radius \(h_C\).
If a rigorous pointwise computation gives
\begin{equation}
 \sigma_{\min}(M_{F,n}(x_C))\ge b_C,                  \label{eq:newS7V39-center}
\end{equation}
then
\begin{equation}
 \inf_{x\in C}\sigma_{\min}(M_{F,n}(x))
 \ge b_C-L_{F,n}h_C.                                 \label{eq:newS7V39-cell}
\end{equation}
\end{proposition}

\begin{proof}
For \(x\in C\), the singular-value perturbation inequality and
\eqref{eq:newS7V39-Lipschitz} give
\[
 \sigma_{\min}(M_{F,n}(x))
 \ge\sigma_{\min}(M_{F,n}(x_C))
   -\|M_{F,n}(x)-M_{F,n}(x_C)\|
 \ge b_C-L_{F,n}h_C.
\]
\end{proof}

\begin{definition}[Polynomial face certificate, PFC]
\label{def:newS7V39-PFC}
PFC consists of:
\begin{enumerate}
\item a polynomial degree \(n\);
\item a finite cover by parameter cells \(C\);
\item outward-rounded rational or interval bounds \(b_C\) and
\(L_{F,n}\);
\item the strict inequalities
\begin{equation}
 b_C-L_{F,n}h_C>\varepsilon_n                         \label{eq:newS7V39-PFC}
\end{equation}
for every face and cell.
\end{enumerate}
\end{definition}

\begin{theorem}[PFC implies FBSC]
\label{thm:newS7V39-PFC-FBSC}
PFC implies FBSC, C1FGC, and the codimension-one SFGC row.
\end{theorem}

\begin{proof}
The finite minimum
\[
 \beta_n:=\min_{F,C}(b_C-L_{F,n}h_C)
\]
is strictly larger than \(\varepsilon_n\).  Apply
\Cref{prop:newS7V39-cell-bound,cor:newS7V39-finite-certificate}, followed
by the V36 implication from FBSC to C1FGC.
\end{proof}

\subsection{Completeness of refinement when a gap exists}
\label{sec:newS7V39-completeness}

\begin{theorem}[Eventual termination of certified refinement]
\label{thm:newS7V39-termination}
Assume FBSC is true with exact lower bound \(\mu_*>0\), the parameter
dependence is uniformly continuous, and the interval arithmetic used for
the finite expressions is convergent under cell subdivision.  Then there
exist a finite degree \(n\) and a finite cell refinement which satisfy PFC.
\end{theorem}

\begin{proof}
Choose \(n\) so large that
\(\varepsilon_n<\mu_*/4\).  The perturbation inequality then gives
\[
 \sigma_{\min}(M_{F,n})\ge3\mu_*/4.
\]
Uniform continuity and convergence of the interval enclosures permit a
finite cell cover on which the loss
\(L_{F,n}h_C\) together with the outward-rounding loss is smaller than
\(\mu_*/4\).  The resulting certified lower bound is greater than
\(\mu_*/2>\varepsilon_n\).  Compactness makes the refinement finite.
\end{proof}

This is a completeness statement conditional on the mathematical truth of
FBSC.  It does not decide which alternative holds.  If an exact zero is
present, \Cref{cor:newS7V39-detection} forces every sufficiently accurate
surrogate to display a minimum at the same geometric error scale.

\subsection{Data still needed from the V33 assembly}
\label{sec:newS7V39-data}

PFC requires no infinite matrix.  Its complete input ledger is:
\begin{enumerate}
\item the finite face set and its normal-layer interface windows;
\item the explicit blocks of the actual sharp face \(H_F\);
\item the finite-collar reference blocks \(H_0\) from the V37 path;
\item certified constants \(\delta_0,M_0,v_*\);
\item the compact coframe cells and their denominator-positivity bounds.
\end{enumerate}
Every item is determined by the V33 mass-lumped Dirac and stiffness forms
and the V38 bulk path.  The next work is therefore coefficient extraction,
not another abstract spectral reduction.

\begin{assessment}[An exact finite certification route]
\label{ass:newS7V39-status}
V39 gives an exponentially convergent finite-range certificate which is
sufficient for FBSC and which must eventually succeed whenever FBSC is
true.  It corrects the unjustified claim that the exact infinite-tail Green
matrix is rational.  The actual V33 face blocks have not yet been inserted,
so FBSC itself remains open.
\end{assessment}

\begin{programme}[V40: compulsory audit and reference-face block ledger]
\label{prog:newS7V39-V40}
V40 will perform the compulsory YM/NS audit, then begin the explicit PFC
input ledger for one oriented reference face.  It will derive the interface
window, the sharp-face defect blocks, and row-sum constants directly from
the V33 bilinear forms.  Any missing normalization will be pushed upstream
before a numerical or interval determinant is asserted.
\end{programme}

\begin{firewall}[Claim boundary after seventh-series V39]
\label{fire:newS7V39-boundary}
V39 proves a uniform polynomial approximation to the gapped reference
inverse, finite propagation of the approximant, a certified singular-value
error for the face matrix, a cellwise interval certificate, and eventual
termination of refinement whenever FBSC is true.

V39 does not prove PFC or FBSC for the explicit V33 face data, exclude every
surface mode, prove higher-codimension tangent gaps, full SFGC, the global
finite-element overlap gap, strong local heat consistency, Gaussian heat
locality, the common Einstein stress, pure gravitational owners, rough
metric or AOP control, measure concentration, the determinant-line phase,
interacting cutoff removal, or the Standard Model--Einstein continuum.
\end{firewall}

\begin{theorem}[Seventh-series V39 result]
\label{thm:newS7V39-result}
For every regular face, the exact V36 Birman--Schwinger matrix admits a
finite Laurent surrogate \(M_{F,n}\) with certified uniform error
\((v_*/\delta_0)\kappa^{n+1}\).  A finite cell cover satisfying
\eqref{eq:newS7V39-PFC} proves FBSC and the codimension-one SFGC row; if
FBSC is true, sufficiently high degree and sufficiently fine validated
subdivision must produce such a certificate.
\end{theorem}

\begin{proof}
Use \Cref{lem:newS7V39-inverse-approximation,
lem:newS7V39-propagation,
thm:newS7V39-M-error,
cor:newS7V39-finite-certificate,
prop:newS7V39-cell-bound,
thm:newS7V39-PFC-FBSC,
thm:newS7V39-termination}.
\end{proof}

\paragraph{Parent-version record.}
V39 was formed from
\texttt{Renewal\_SM\_Einstein\_Continuum\_Research\_Notes\_V38.tex}.
The V38 parent had \(15\,405\) logical source lines, \(621\,188\) bytes,
and SHA--256
\[
 \texttt{DC075129C4F4E2BB8075871DF79617CE4E718B4A4D19334571731B35F99DCCC1}.
\]
The next compulsory companion audit is V40.

% =============================================================================
% END APPENDED SEVENTH-SERIES V39 MATERIAL
% =============================================================================

% =============================================================================
% BEGIN APPENDED SEVENTH-SERIES V40 MATERIAL
% =============================================================================

\section{Seventh-series V40: compulsory audit and the face-block ledger}
\label{sec:v7v40}
\label{sec:newS7V40}

\subsection{Compulsory companion audit}
\label{sec:newS7V40-audit}

\begin{audit}[Companion snapshots used at V40]
\label{aud:newS7V40-snapshots}
The exact companion snapshots inspected before the V40 argument were:
\begin{center}
\begin{tabular}{p{0.20\linewidth}p{0.52\linewidth}r}
programme & source & bytes\\ \hline
Yang--Mills &
\texttt{Renewal\_Yang\_Mills\_Mass\_Gap\_Research\_Notes\_V15.tex}
& \(238\,512\)\\
Navier--Stokes &
\texttt{Renewal\_NS\_Stability\_Research\_Notes\_V62.tex}
& \(621\,237\)
\end{tabular}
\end{center}
The Yang--Mills source had \(6\,743\) newline characters and SHA--256
\[
 \texttt{BD4243291F737FAF4E835607F11742554C44B4AF1A781716073D021FA926531B}.
\]
The Navier--Stokes source had \(12\,185\) newline characters and SHA--256
\[
 \texttt{C0ACA4D330DE68579FEDD179F3CC286DA5124997B58267D2B48F810BBD59B491}.
\]
The active NS file advanced while the audit began; the second hash records
the complete V62 bytes actually read and makes the comparison reproducible.
\end{audit}

The Yang--Mills note proves a three-sign counterexample in which every
one- and two-coordinate marginal is independent of a parameter while the
three-body moment equals that parameter.  It therefore corrects a
pair-kernel argument and replaces it by a Cayley-weighted hypergraph
activity card.  Its stratified form also records normal-interface
resonances and finite-collar stability, explicitly borrowing the SM
interface lesson.

The Navier--Stokes note proves that a radial cutoff integrates the Hessian
of a localized Newtonian potential to its mass contribution, while every
higher multipole cancels in the spherical mean.  It then bounds the full
cross term on a nearly constant anchor and derives conditional defect and
measure consequences.  No global regularity result is claimed.

\begin{decision}[Effect of the V40 audit]
\label{dec:newS7V40-audit}
The YM counterexample rules out replacing the full face determinant by
pairwise or entrywise surrogate conditions: PFC must retain the complete
matrix \(M_{F,n}\).  The NS cancellation shows the value of performing an
exact symmetry projection before taking absolute bounds.  Accordingly,
V40 records the full assembled face blocks first and permits
block-diagonalization only by exact automorphisms of the oriented face,
coframe, and spin data.  Neither companion imports a missing SM gap.
\end{decision}

\subsection{Normalized affine-simplex data}
\label{sec:newS7V40-simplex-data}

Work in one frozen codimension-one tangent model.  Smooth spin and gauge
connections are lower order after multiplication by \(a\), so the leading
tangent operator uses the frozen trivial transport.  For every small
four-simplex \(\sigma\), define
\begin{equation}
 \nu_\sigma:={\operatorname{Vol}_g(\sigma)\over a^4},\qquad
 b_{\sigma x}:=a\,d\phi_x|_\sigma
 \quad(x\in\operatorname{Vert}\sigma).                \label{eq:newS7V40-normalized-data}
\end{equation}
Both quantities are independent of \(a\) in the frozen affine model.
Write \(c_E(b)\) for Clifford multiplication by the covector \(b\) in the
frozen coframe convention of V33--V34.

\begin{lemma}[Normalized lumped masses]
\label{lem:newS7V40-lumped-mass}
For a vertex \(x\),
\begin{equation}
 w_x=a^4\omega_x,\qquad
 \omega_x={1\over5}\sum_{\sigma\ni x}\nu_\sigma>0.     \label{eq:newS7V40-omega}
\end{equation}
\end{lemma}

\begin{proof}
On an affine four-simplex,
\(\int_\sigma\phi_x\,d\mu_g=\operatorname{Vol}_g(\sigma)/5\).
Sum this identity over the incident simplices and use the V33 definition
of \(w_x\).
\end{proof}

For vertices \(x,y\), define the assembled dimensionless form blocks
\begin{align}
 \mathsf D_{xy}
 &:={1\over5}\sum_{\sigma\supset\{x,y\}}
       \nu_\sigma c_E(b_{\sigma y}),                  \label{eq:newS7V40-D-form}\\
 \mathsf Q_{xy}
 &:=\sum_{\sigma\supset\{x,y\}}
       \nu_\sigma
       \langle b_{\sigma x},b_{\sigma y}\rangle_{g^{-1}}\,I_{\rm spin}.
                                                               \label{eq:newS7V40-Q-form}
\end{align}
An empty incidence sum is zero.

\begin{lemma}[Exact affine element ledger]
\label{lem:newS7V40-element-ledger}
In the frozen trivialization, the leading tangent contributions of one
simplex to the V33 Dirac and stiffness forms are precisely the summands in
\eqref{eq:newS7V40-D-form} and
\eqref{eq:newS7V40-Q-form}.  After assembly and half-density conjugation,
\begin{align}
 \mathbb K_{xy}
 &=\omega_x^{-1/2}\mathsf D_{xy}\omega_y^{-1/2},       \label{eq:newS7V40-K-block}\\
 \mathbb L_{xy}
 &=\omega_x^{-1/2}\mathsf Q_{xy}\omega_y^{-1/2}.       \label{eq:newS7V40-L-block}
\end{align}
\end{lemma}

\begin{proof}
On \(\sigma\),
\[
 I_a\eta=\sum_x\phi_x\eta_x,\qquad
 \mathcal D_EI_a\psi
 =\sum_y a^{-1}c_E(b_{\sigma y})\psi_y.
\]
Hence the \(xy\) Dirac form entry is
\[
 {a^4\nu_\sigma\over5}\,a^{-1}c_E(b_{\sigma y})
 =a^3{\nu_\sigma\over5}c_E(b_{\sigma y}).
\]
Similarly, the stiffness form entry is
\[
 a^4\nu_\sigma\,a^{-2}
 \langle b_{\sigma x},b_{\sigma y}\rangle_{g^{-1}}
 =a^2\nu_\sigma
 \langle b_{\sigma x},b_{\sigma y}\rangle_{g^{-1}}.
\]
The Riesz operators in the weighted nodal space have a factor \(w_x^{-1}\)
on the output row.  Conjugating to counting \(\ell^2\) by
\(\operatorname{diag}(w_x^{1/2})\), then multiplying the Dirac operator by
\(a\) and the stiffness operator by \(a^2\), gives
\eqref{eq:newS7V40-K-block}--\eqref{eq:newS7V40-L-block}.
\end{proof}

\begin{corollary}[Dimensionless sharp-face sign blocks]
\label{cor:newS7V40-H-block}
The actual sharp tangent sign kernel has blocks
\begin{equation}
 \mathbb H^{\rm sh}_{xy}
 =\Gamma\left\{
   \mathbb K_{xy}+{r\over2}\mathbb L_{xy}
   -\rho\,\delta_{xy}I
 \right\}.                                             \label{eq:newS7V40-H-block}
\end{equation}
These blocks are self-adjoint after full assembly, have finite incidence
range, and contain all spin correlations required by PFC.
\end{corollary}

\begin{proof}
Insert \eqref{eq:newS7V40-K-block} and
\eqref{eq:newS7V40-L-block} into the V33 Wilson definition, apply the
half-density conjugation, and use the V33 exact Wilson algebra.
\end{proof}

\subsection{Tangential Fourier and normal-layer blocks}
\label{sec:newS7V40-normal-blocks}

Choose a primitive tangential cell of the face and label its vertex orbits
by \(a=1,\ldots,q_F\).  Label normal layers by \(m\in\mathbb Z\).  If
\(d\in\mathbb Z^3\) is a tangential translation, define
\begin{equation}
 \mathcal H^{\rm sh}_{j,d;ab}
 :=\mathbb H^{\rm sh}_{(0,a),(j,b+d)}.                 \label{eq:newS7V40-real-block}
\end{equation}
Only finitely many \(j,d\) occur.

\begin{definition}[Sharp face coefficient ledger]
\label{def:newS7V40-sharp-ledger}
The normal recurrence blocks are
\begin{equation}
 A_j^{\rm sh}(E,k_\parallel)_{ab}
 :=\sum_d e^{ik_\parallel\cdot d}
       \mathcal H^{\rm sh}_{j,d;ab}.                  \label{eq:newS7V40-Aj}
\end{equation}
Every real-space entry in \(\mathcal H^{\rm sh}_{j,d;ab}\) is obtained by
the finite incidence sums
\eqref{eq:newS7V40-omega}--\eqref{eq:newS7V40-H-block}.
\end{definition}

\begin{lemma}[Exact interface range]
\label{lem:newS7V40-interface-range}
Let
\begin{equation}
 J_F:=\max\{|m(y)-m(x)|:
 x,y\hbox{ belong to one small simplex meeting }F\}.   \label{eq:newS7V40-JF}
\end{equation}
Then \(A_j^{\rm sh}=0\) for \(|j|>J_F\).  A sharp block row differs from
its left or right regular bulk formula only if the closed star of its
vertex orbit meets the original face.
\end{lemma}

\begin{proof}
V33 locality says that a nonzero block joins vertices of a common small
simplex, which gives the first assertion by definition of \(J_F\).
If a closed star lies strictly in one original four-simplex, every incident
small simplex belongs to its regular edgewise bulk subdivision.  The
incidence sums in \eqref{eq:newS7V40-D-form} and
\eqref{eq:newS7V40-Q-form} are then exactly the corresponding bulk sums.
\end{proof}

Let \(H_{0,L}\) be the V37 adiabatic reference with collar width \(L\).
Its two tails agree with the same regular bulks.

\begin{corollary}[Finite defect window]
\label{cor:newS7V40-defect-window}
After translating the face to normal layer zero, the defect
\begin{equation}
 W_F=H_F^{\rm sh}-H_{0,L}                             \label{eq:newS7V40-W}
\end{equation}
has no nonzero block row or column outside
\begin{equation}
 I_F=\{-L-J_F-C_{\rm st},\ldots,L+J_F+C_{\rm st}\},   \label{eq:newS7V40-window}
\end{equation}
where \(C_{\rm st}\) is the finite maximum normal radius of a V33 vertex
star.  Thus the V36 interface dimension is
\begin{equation}
 \dim\mathscr I_F=q_F\,|I_F|\,\dim S_{\rm SM},          \label{eq:newS7V40-interface-dimension}
\end{equation}
before exact symmetry reduction.
\end{corollary}

\begin{proof}
Outside the displayed window, the reference interpolation has reached its
constant tail and the vertex star of the sharp mesh lies wholly in the same
regular bulk.  The block formulas therefore agree.  Coordinate restriction
to all remaining layer and orbit fibres gives the stated dimension.
\end{proof}

\subsection{Uniform row-sum constants for PFC}
\label{sec:newS7V40-constants}

The finite type set and compact shape-regular coframe class supply constants
\begin{equation}
 0<\omega_-\le\omega_x\le\omega_+,\quad
 0<\nu_\sigma\le\nu_+,\quad
 \|c_E(b_{\sigma x})\|\le c_+b_+,\quad
 |\langle b_{\sigma x},b_{\sigma y}\rangle_{g^{-1}}|
 \le g_+b_+^2.                                        \label{eq:newS7V40-geometric-bounds}
\end{equation}
Let \(N_{\rm pair}\) bound the number of small simplices containing a
fixed vertex pair and \(d_{\rm row}\) the number of coupled vertices in a
row.

\begin{lemma}[Explicit block and operator bounds]
\label{lem:newS7V40-row-bounds}
Set
\begin{align}
 C_K&:={N_{\rm pair}\nu_+c_+b_+\over5\omega_-},        \label{eq:newS7V40-CK}\\
 C_L&:={N_{\rm pair}\nu_+g_+b_+^2\over\omega_-},       \label{eq:newS7V40-CL}\\
 M_{\rm sh}&:=\rho+d_{\rm row}
              \left(C_K+{r\over2}C_L\right).          \label{eq:newS7V40-Msh}
\end{align}
Then
\begin{equation}
 \|\mathbb K_{xy}\|\le C_K,\qquad
 \|\mathbb L_{xy}\|\le C_L,\qquad
 \|H_F^{\rm sh}\|\le M_{\rm sh}.                      \label{eq:newS7V40-row-result}
\end{equation}
If the adiabatic reference blocks obey the analogous row bound \(M_0\),
then the V36 defect satisfies
\begin{equation}
 \|V_F\|=\|W_F\|\le v_*:=M_{\rm sh}+M_0.              \label{eq:newS7V40-vstar}
\end{equation}
\end{lemma}

\begin{proof}
Each incidence sum has at most \(N_{\rm pair}\) terms.  Apply
\eqref{eq:newS7V40-geometric-bounds} and
\(\omega_x^{-1/2}\omega_y^{-1/2}\le\omega_-^{-1}\) to
\eqref{eq:newS7V40-K-block}--\eqref{eq:newS7V40-L-block}.  The block Schur
row test gives \(M_{\rm sh}\).  Finally,
\(\|H_F^{\rm sh}-H_{0,L}\|\le\|H_F^{\rm sh}\|+\|H_{0,L}\|\).
\end{proof}

\subsection{Exact symmetry before absolute bounds}
\label{sec:newS7V40-symmetry}

\begin{definition}[Face automorphism group]
\label{def:newS7V40-face-group}
Let \(\mathcal G_F(E)\) be the finite group of oriented incidence
automorphisms which preserve the face, the frozen coframe data, the lumped
weights, the Clifford action through its spin lift, and both the sharp and
reference block families.
\end{definition}

\begin{proposition}[Exact sector reduction]
\label{prop:newS7V40-sector}
If the interface restriction is chosen \(\mathcal G_F(E)\)-invariant, then
\(V_F\), \(G_{F,n}\), and \(M_{F,n}\) commute with the finite unitary
representation of \(\mathcal G_F(E)\).  Hence
\begin{equation}
 M_{F,n}\cong\bigoplus_{\lambda\in\widehat{\mathcal G_F(E)}}
 M_{F,n}^{(\lambda)},\qquad
 \sigma_{\min}(M_{F,n})
 =\min_\lambda\sigma_{\min}(M_{F,n}^{(\lambda)}).      \label{eq:newS7V40-sector}
\end{equation}
\end{proposition}

\begin{proof}
Every ingredient is assembled from invariant incidence sums and the exact
spin lift, so it intertwines the group action.  Polynomial functional
calculus preserves commutants, giving the assertion for \(G_{F,n}\).
Finite-group unitary representation theory supplies the orthogonal
isotypic decomposition, and singular values of a block direct sum are the
union of the block singular values.
\end{proof}

For a generic coframe the exact automorphism group may be trivial.  No
spherical, reflection, or pairwise cancellation is inserted unless it is
an automorphism of the complete oriented coefficient ledger.  This is the
precise SM use of the two companion lessons.

\begin{assessment}[The abstract PFC inputs are now assembled formulas]
\label{ass:newS7V40-status}
V40 derives every sharp-face block, the normal recurrence, a finite defect
window, and uniform norm constants directly from the V33 bilinear forms.
What remains is finite enumeration of the small simplices and evaluation
of these formulas for the chosen Freudenthal reference face.  The note has
not supplied those incidence tables, so PFC and FBSC remain open.
\end{assessment}

\begin{programme}[V41: finite incidence extraction and validation rows]
\label{prog:newS7V40-V41}
V41 will specify an orientation-independent incidence table for a
Freudenthal four-simplex and its shared three-face, prove that the table
reconstructs \(\omega_x,\mathsf D_{xy},\mathsf Q_{xy}\) without duplicate
or omitted simplex contributions, and state exact validation identities:
partition of unity, affine exactness, Dirac skew-adjointness, stiffness
row-sum zero, and positivity.  These identities must be checked before the
PFC matrices are trusted.
\end{programme}

\begin{firewall}[Claim boundary after seventh-series V40]
\label{fire:newS7V40-boundary}
V40 performs the compulsory YM/NS audit, derives the normalized lumped
mass, Dirac, stiffness, and sign-kernel blocks, derives tangential Fourier
and normal-layer coefficients, bounds the exact interface window and
operator norms, and proves exact finite-group sector reduction.

V40 does not enumerate the reference-face incidence table, prove PFC or
FBSC, exclude every surface mode, prove higher-codimension tangent gaps,
full SFGC, the global finite-element overlap gap, strong local heat
consistency, Gaussian heat locality, the common Einstein stress, pure
gravitational owners, rough metric or AOP control, measure concentration,
the determinant-line phase, interacting cutoff removal, or the Standard
Model--Einstein continuum.
\end{firewall}

\begin{theorem}[Seventh-series V40 result]
\label{thm:newS7V40-result}
The PFC input for every regular sharp face is determined by the finite
incidence formulas
\eqref{eq:newS7V40-omega}--\eqref{eq:newS7V40-H-block}.  Its exact defect
window is finite as in \eqref{eq:newS7V40-window}, its uniform defect norm
is bounded by \eqref{eq:newS7V40-vstar}, and every exact face automorphism
reduces the complete surrogate matrix by
\eqref{eq:newS7V40-sector} without losing determinant correlations.
\end{theorem}

\begin{proof}
Use \Cref{lem:newS7V40-lumped-mass,
lem:newS7V40-element-ledger,
cor:newS7V40-H-block,
lem:newS7V40-interface-range,
cor:newS7V40-defect-window,
lem:newS7V40-row-bounds,
prop:newS7V40-sector}.
\end{proof}

\paragraph{Parent-version record.}
V40 was formed from
\texttt{Renewal\_SM\_Einstein\_Continuum\_Research\_Notes\_V39.tex}.
The V39 parent had \(15\,757\) logical source lines, \(634\,205\) bytes,
and SHA--256
\[
 \texttt{EF17E5F9B5C6F2877CAFA777D01132055AE161B612C198CAE2427D8A6471CF40}.
\]
The next compulsory companion audit is V45.

% =============================================================================
% END APPENDED SEVENTH-SERIES V40 MATERIAL
% =============================================================================

% =============================================================================
% BEGIN APPENDED SEVENTH-SERIES V41 MATERIAL
% =============================================================================

\section{Seventh-series V41: incidence extraction and assembly validation}
\label{sec:v7v41}
\label{sec:newS7V41}

\subsection{The table must precede the determinant}
\label{sec:newS7V41-purpose}

V40 expresses every sharp-face coefficient as a finite incidence sum.
Those formulas are rigorous only if the small-simplex incidence list is
itself auditable.  A missing simplex changes the lumped mass, while a
duplicate simplex changes both the Dirac and stiffness blocks; either error
can create or remove a numerical surface zero.  This note defines an
orientation-independent table, proves how it reconstructs the affine
blocks, and gives identities which must hold before PFC is attempted.
\begin{correction}[The normal row suppressed in the V40 ledger]
\label{corr:newS7V41-normal-row}
The V40 notation
\(\mathcal H^{\rm sh}_{j,d;ab}\) and
\(A_j^{\rm sh}(E,k_\parallel)\) displayed the row based at layer zero.
The actual face recurrence is not translation invariant in the normal
direction.  Its complete coefficient ledger is
\begin{align}
 \mathcal H^{\rm sh}_{j,d;ab}(m)
 &:=\mathbb H^{\rm sh}_{(m,a),(m+j,b+d)},              \label{eq:newS7V41-real-block-m}\\
 A_j^{\rm sh}(m,E,k_\parallel)_{ab}
 &:=\sum_de^{ik_\parallel\cdot d}
       \mathcal H^{\rm sh}_{j,d;ab}(m).                \label{eq:newS7V41-Aj-m}
\end{align}
The V40 formulas are the \(m=0\) row.  Every claim about finite range,
tail stabilization, and incidence reconstruction applies row by row to
\eqref{eq:newS7V41-Aj-m}.  No normal translation invariance is assumed at
the face.
\end{correction}

\subsection{Global combinatorial records}
\label{sec:newS7V41-records}

Let \(\mathcal V\) be the set of global barycentric-lattice vertex
identifiers in one finite normal/tangential fundamental window.  A vertex
identifier records:
\begin{equation}
 x=(\Sigma;\alpha_0,\ldots,\alpha_4),\qquad
 \alpha_i\in\mathbb N_0,\quad\sum_i\alpha_i=n,         \label{eq:newS7V41-vertex-id}
\end{equation}
modulo the declared identifications on common original faces and modulo
tangential lattice translation.  Here \(\Sigma\) is an original
four-simplex label.  On a common face, the global identifier is the
equivalence class, not either local copy.

\begin{definition}[Small-simplex record]
\label{def:newS7V41-simplex-record}
A small-simplex record is
\begin{equation}
 \mathfrak s=(\{x_0,\ldots,x_4\},[v_0,\ldots,v_4],
              \epsilon_{\mathfrak s},\mathbf t_{\mathfrak s}),        \label{eq:newS7V41-simplex-record}
\end{equation}
where:
\begin{enumerate}
\item \(\{x_0,\ldots,x_4\}\) is the unordered set of five distinct global
vertex identifiers;
\item \([v_0,\ldots,v_4]\) is one local ordering used only to compute an
edge matrix;
\item \(\epsilon_{\mathfrak s}\in\{-1,1\}\) records that ordering relative
to the oriented original simplex; and
\item \(\mathbf t_{\mathfrak s}\in\mathbb Z^3\) is the tangential
translation bringing the record to the chosen fundamental window.
\end{enumerate}
Two records represent the same simplex exactly when their unordered global
vertex sets and translation classes agree.
\end{definition}

The standard edgewise/Freudenthal rule can be stored without reference to
the arbitrary order \([v_0,\ldots,v_4]\): for each permitted edgewise cell,
store its base lattice point and its ordered increment word; convert the
five resulting lattice points to global identifiers; then use the
unordered set as the primary key.  Near an original face only the increment
words allowed by the nonnegative barycentric coordinates are retained.

\begin{definition}[Oriented three-face record]
\label{def:newS7V41-face-record}
For each unordered four-vertex set
\(f=\{y_0,\ldots,y_3\}\) occurring as a face, store the incidence list
\begin{equation}
 \mathcal I(f)=
 \{(\mathfrak s,i,\varepsilon_{\mathfrak s,f}) :
 f=\mathfrak s\setminus\{v_i\}\},                     \label{eq:newS7V41-face-incidence}
\end{equation}
where \(\varepsilon_{\mathfrak s,f}\) is the sign of the boundary
orientation induced by \(\mathfrak s\), expressed in one fixed ordering of
the global identifiers of \(f\).
\end{definition}

\begin{lemma}[Orientation independence of the primary table]
\label{lem:newS7V41-orientation-independence}
Permuting the local order \([v_0,\ldots,v_4]\) changes the stored simplex
orientation and the corresponding face signs, but does not change the
primary simplex key, the global face keys, or any assembled V40 block.
\end{lemma}

\begin{proof}
The primary keys are unordered global vertex sets.  A local permutation
changes the determinant sign and every induced boundary sign by the same
standard alternating rule.  Barycentric basis functions are tied to their
global vertex identifiers, so the geometric covector attached to a vertex
is merely relabelled.  The V40 sums run over global \(x,y\) and are
therefore invariant under that relabelling.
\end{proof}

\subsection{Geometry reconstructed from one record}
\label{sec:newS7V41-geometry}

Choose frozen affine coordinates in the containing original simplex and
write
\begin{equation}
 r_j:={X(v_j)-X(v_0)\over a},\qquad
 R_{\mathfrak s}:=(r_1\ r_2\ r_3\ r_4).               \label{eq:newS7V41-edge-matrix}
\end{equation}
Let \(e^j\) be the coordinate covectors in \(\mathbb R^4\).

\begin{lemma}[Affine reconstruction formulas]
\label{lem:newS7V41-affine-reconstruction}
If \(\det R_{\mathfrak s}\ne0\), then
\begin{align}
 b_{\mathfrak s,v_j}&=R_{\mathfrak s}^{-T}e^j
 &&(j=1,\ldots,4),                                    \label{eq:newS7V41-bj}\\
 b_{\mathfrak s,v_0}&=-\sum_{j=1}^4b_{\mathfrak s,v_j},              \label{eq:newS7V41-b0}\\
 \nu_{\mathfrak s}
 &= {|\det R_{\mathfrak s}|\sqrt{\det g_{\mathfrak s}}\over4!}.
                                                               \label{eq:newS7V41-volume}
\end{align}
They are invariant under a change of local vertex order after the
vertex labels are carried along.
\end{lemma}

\begin{proof}
The affine function \(\phi_{v_j}\), \(j\ge1\), is characterized by
\(\phi_{v_j}(v_0)=0\) and
\(d\phi_{v_j}(X(v_i)-X(v_0))=\delta_{ij}\).  Multiplication by \(a\) gives
\eqref{eq:newS7V41-bj}.  Partition of unity gives
\eqref{eq:newS7V41-b0}.  The Euclidean coordinate volume of the normalized
simplex is \(|\det R_{\mathfrak s}|/4!\), and the frozen metric density
multiplies it by \(\sqrt{\det g_{\mathfrak s}}\).  All three quantities are
geometric data attached to labelled vertices, so a local reordering only
changes their presentation.
\end{proof}

\begin{corollary}[Element identities]
\label{cor:newS7V41-element-identities}
Every valid simplex record obeys
\begin{align}
 \sum_{x\in\mathfrak s}b_{\mathfrak s,x}&=0,           \label{eq:newS7V41-partition-gradient}\\
 \sum_{x\in\mathfrak s}
 {X(x)-X(v_0)\over a}\otimes b_{\mathfrak s,x}
 &=I_{\mathbb R^4},                                   \label{eq:newS7V41-affine-exact}\\
 \sum_{x\in\mathfrak s}{\nu_{\mathfrak s}\over5}
 &=\nu_{\mathfrak s}.                                 \label{eq:newS7V41-mass-sum}
\end{align}
\end{corollary}

\begin{proof}
The first and third are immediate.  In the second identity the \(v_0\)
term vanishes and \(\sum_{j=1}^4r_j\otimes R^{-T}e^j=I\).
\end{proof}

\subsection{Incidence validation rows}
\label{sec:newS7V41-validation}

\begin{definition}[Face incidence validation card, FIVC]
\label{def:newS7V41-FIVC}
FIVC consists of the following finite checks on the fundamental-window
table:
\begin{enumerate}
\item[\textup{(F1)}] every simplex key occurs exactly once;
\item[\textup{(F2)}] every simplex has five distinct vertices and
\(\det R_{\mathfrak s}\ne0\);
\item[\textup{(F3)}] every internal small three-face has exactly two
incidences and their induced signs sum to zero;
\item[\textup{(F4)}] every small face on an original shared three-face has
the same global vertex key from both original simplices;
\item[\textup{(F5)}] tangential translation of a boundary record re-enters
the fundamental table with the declared phase displacement;
\item[\textup{(F6)}] the element identities
\eqref{eq:newS7V41-partition-gradient}--
\eqref{eq:newS7V41-mass-sum} hold; and
\item[\textup{(F7)}] every vertex star obtained from the table belongs to
the finite V33 star catalogue and contains the declared number of incident
simplex records.
\end{enumerate}
\end{definition}

\begin{lemma}[FIVC excludes duplicate and omitted element contributions]
\label{lem:newS7V41-no-duplicate}
If FIVC holds, then every small four-simplex contribution enters each V40
mass, Dirac, and stiffness incidence sum exactly once, including at an
original shared face.
\end{lemma}

\begin{proof}
F1 excludes duplicate four-simplex records.  An omitted simplex would leave
one of its internal faces with a single incidence or would break the
declared face subdivision or star count, contradicting F3, F4, or F7.
F4 identifies the two local descriptions of a shared-face vertex, so the
assembly uses one global nodal degree of freedom rather than two copies.
Thus the sums indexed by \(\sigma\supset\{x,y\}\) are exact.
\end{proof}

\subsection{Operator validation identities}
\label{sec:newS7V41-operator-identities}

Assemble \(\omega_x,\mathsf D,\mathsf Q\) from the validated table by the
V40 formulas.

\begin{lemma}[Stiffness conservation and positivity]
\label{lem:newS7V41-stiffness-validation}
The assembled stiffness form satisfies
\begin{align}
 \sum_y\mathsf Q_{xy}&=0,                              \label{eq:newS7V41-Q-row-sum}\\
 \sum_{x,y}\langle z_x,\mathsf Q_{xy}z_y\rangle
 &=\sum_{\mathfrak s}\nu_{\mathfrak s}
   \left\|\sum_{x\in\mathfrak s}
        b_{\mathfrak s,x}z_x\right\|_{g^{-1}}^2
 \ge0.                                                 \label{eq:newS7V41-Q-positive}
\end{align}
\end{lemma}

\begin{proof}
For a fixed element contribution, summing over \(y\) inserts
\(\sum_yb_{\mathfrak s,y}=0\), proving the row sum.  Expanding the square
on the right gives exactly the assembled quadratic form.  The equality
uses F1 to count each element once.
\end{proof}

\begin{lemma}[Dirac constant-mode and affine rows]
\label{lem:newS7V41-Dirac-validation}
In the frozen zero-connection tangent model,
\begin{align}
 \sum_y\mathsf D_{xy}&=0,                              \label{eq:newS7V41-D-constant}\\
 \sum_y\mathsf D_{xy}\,(\lambda\cdot X(y))
 &=a\,\omega_x\,c_E(\lambda)
   +\mathsf B_{x,\lambda},                             \label{eq:newS7V41-D-affine}
\end{align}
where \(\mathsf B_{x,\lambda}\) consists only of boundary faces of the
finite extraction window.  On the full periodic/tangent assembly, opposite
boundary translates cancel and \(\mathsf B_{x,\lambda}=0\) in the weak
periodic sense.
\end{lemma}

\begin{proof}
The constant row follows from
\(\sum_yc_E(b_{\mathfrak s,y})=0\) on each element.  For an affine nodal
function, \Cref{cor:newS7V41-element-identities} reconstructs its exact
constant gradient.  Integrating the test hat function gives
\(a\omega_x c_E(\lambda)\).  If the computation is performed on a finite
window before periodic identification, discrete integration by parts
leaves only its boundary faces.  F3 cancels every interior face and F5
cancels the translated boundary pairs.
\end{proof}

\begin{lemma}[Assembled Dirac adjoint check]
\label{lem:newS7V41-Dirac-adjoint}
On the full tangentially periodic face tangent,
\begin{equation}
 \mathsf D+\mathsf D^*=0.                              \label{eq:newS7V41-D-adjoint}
\end{equation}
Equivalently, after Fourier transform,
\begin{equation}
 A_j^K(m,k_\parallel)^*
 =-A_{-j}^K(m+j,k_\parallel),                          \label{eq:newS7V41-K-block-adjoint}
\end{equation}
where \(A_j^K\) denotes the Dirac part of the normal block.
\end{lemma}

\begin{proof}
Elementwise integration by parts expresses
\(\mathsf D+\mathsf D^*\) as a sum of oriented small-face traces.  F3 pairs
each internal face with the opposite induced sign, F4 uses the same trace
degrees of freedom across an original face, and F5 pairs the boundary of a
tangential fundamental window.  All traces cancel.  Fourier transformation
turns the matrix adjoint identity into
\eqref{eq:newS7V41-K-block-adjoint}.
\end{proof}

\begin{theorem}[Validated assembly theorem]
\label{thm:newS7V41-validated-assembly}
If FIVC and the affine reconstruction formulas hold, then the extracted
matrices reproduce the V33 frozen Galerkin forms and satisfy
\begin{equation}
 \mathbb K^*=-\mathbb K,\qquad
 \mathbb L^*=\mathbb L\ge0,\qquad
 \mathbb K\Gamma=-\Gamma\mathbb K,\qquad
 [\mathbb L,\Gamma]=0.                                \label{eq:newS7V41-full-algebra}
\end{equation}
The resulting sign blocks
\(\Gamma(\mathbb K+r\mathbb L/2-\rho)\) are therefore self-adjoint.
\end{theorem}

\begin{proof}
\Cref{lem:newS7V41-no-duplicate} identifies the extracted incidence sums
with the defining V33 element sums.  Lemmas
\ref{lem:newS7V41-stiffness-validation} and
\ref{lem:newS7V41-Dirac-adjoint} give the adjoint and positivity rows.
Clifford multiplication anticommutes with chirality, while the scalar
stiffness blocks commute with it.  The final self-adjointness is the V33
Wilson algebra.
\end{proof}

\subsection{A machine-independent certificate manifest}
\label{sec:newS7V41-manifest}

\begin{definition}[Incidence certificate manifest, ICM]
\label{def:newS7V41-ICM}
For each face type, ICM contains:
\begin{enumerate}
\item the sorted global vertex identifiers and their tangential
translations;
\item the sorted unordered simplex keys;
\item the oriented three-face incidence lists;
\item the exact or outward-rounded edge matrices, normalized volumes, and
barycentric covectors;
\item residuals for F1--F7 and
\eqref{eq:newS7V41-Q-row-sum}--
\eqref{eq:newS7V41-K-block-adjoint}; and
\item the hashes of the table and of the code or hand calculation which
formed the V40 blocks.
\end{enumerate}
\end{definition}

\begin{proposition}[ICM makes PFC reproducible]
\label{prop:newS7V41-ICM-PFC}
An ICM satisfying FIVC and the operator identities determines
\(A_j^{\rm sh}\), the defect window, \(V_F\), and every polynomial surrogate
\(M_{F,n}\) without an unrecorded combinatorial choice.
\end{proposition}

\begin{proof}
The vertex and simplex tables determine
\eqref{eq:newS7V41-edge-matrix}--\eqref{eq:newS7V41-volume}.  The V40
incidence formulas then determine the half-density blocks.  Tangential
translations determine their Fourier monomials and normal labels determine
\(A_j^{\rm sh}\).  The declared V37 reference path determines \(H_0\);
subtraction and polynomial functional calculus determine \(V_F\) and
\(M_{F,n}\).  FIVC removes ambiguity from duplicate local descriptions.
\end{proof}

\subsection{The remaining concrete input}
\label{sec:newS7V41-remaining}

The general edgewise rule fixes the local increment words, but the attached
manuscripts do not specify a machine-readable ordered triangulation
\(\mathcal K_0\), its original-simplex affine vertex coordinates, or a
chosen compact coframe cell decomposition.  Those data affect the list of
face types and the numerical interval constants.  They cannot be invented
without changing the discretization.

\begin{definition}[Mesh incidence data card, MIDC]
\label{def:newS7V41-MIDC}
MIDC is the provision, for the chosen finite ordered base triangulation, of
the original vertex identifications, affine reference maps, original-face
incidences, and compact coframe cells needed to instantiate ICM.
\end{definition}

\begin{assessment}[Validation layer closed; instantiation remains]
\label{ass:newS7V41-status}
V41 closes the logical gap between an edgewise simplex list and the V40
operator blocks: FIVC plus the reconstruction identities is a complete
assembly validation layer.  MIDC is not present in the manuscripts, so no
numerical face determinant is asserted.  This is an upstream data
bottleneck rather than a proof that FBSC fails.
\end{assessment}

\begin{programme}[V42: remove dependence on one base triangulation]
\label{prog:newS7V41-V42}
Rather than wait for a numerical \(\mathcal K_0\), V42 will seek a
shape-regular uniform face theorem.  It will bound the sharp interface
against the V38 universal reference using only the geometric constants in
\eqref{eq:newS7V40-geometric-bounds}.  If those bounds cannot reach the
Neumann region, it will formulate a compact universal ICM parameter space
whose PFC certificate covers every admissible base triangulation at once.
\end{programme}

\begin{firewall}[Claim boundary after seventh-series V41]
\label{fire:newS7V41-boundary}
V41 defines orientation-independent simplex and face records, proves the
affine reconstruction formulas, formulates FIVC, proves exact counting,
stiffness conservation and positivity, Dirac constant and adjoint rows,
the validated assembly theorem, and the reproducible ICM manifest.

V41 does not supply MIDC for a concrete base triangulation, prove PFC or
FBSC, exclude every surface mode, prove higher-codimension tangent gaps,
full SFGC, the global finite-element overlap gap, strong local heat
consistency, Gaussian heat locality, the common Einstein stress, pure
gravitational owners, rough metric or AOP control, measure concentration,
the determinant-line phase, interacting cutoff removal, or the Standard
Model--Einstein continuum.
\end{firewall}

\begin{theorem}[Seventh-series V41 result]
\label{thm:newS7V41-result}
Every V40 face block can be reconstructed without orientation ambiguity
from the global simplex and face records.  FIVC and identities
\eqref{eq:newS7V41-partition-gradient}--
\eqref{eq:newS7V41-K-block-adjoint} certify that the extraction counts every
element once and preserves the exact Dirac--Wilson adjoint, chirality, and
positivity algebra required by PFC.
\end{theorem}

\begin{proof}
Use \Cref{lem:newS7V41-orientation-independence,
lem:newS7V41-affine-reconstruction,
cor:newS7V41-element-identities,
lem:newS7V41-no-duplicate,
lem:newS7V41-stiffness-validation,
lem:newS7V41-Dirac-validation,
lem:newS7V41-Dirac-adjoint,
thm:newS7V41-validated-assembly,
prop:newS7V41-ICM-PFC}.
\end{proof}

\paragraph{Parent-version record.}
V41 was formed from
\texttt{Renewal\_SM\_Einstein\_Continuum\_Research\_Notes\_V40.tex}.
The V40 parent had \(16\,153\) logical source lines, \(650\,056\) bytes,
and SHA--256
\[
 \texttt{D7F10502AE670168B210EA307AC1A00C6706C767E1359E3B90038CFCBE5A898A}.
\]
The next compulsory companion audit is V45.

% =============================================================================
% END APPENDED SEVENTH-SERIES V41 MATERIAL
% =============================================================================

% =============================================================================
% BEGIN APPENDED SEVENTH-SERIES V42 MATERIAL
% =============================================================================

\section{Seventh-series V42: the universal shape-regular face space}
\label{sec:v7v42}
\label{sec:newS7V42}

\subsection{Replacing missing mesh data by a uniform theorem}
\label{sec:newS7V42-purpose}

V41 identifies MIDC because the manuscripts do not fix one
machine-readable base triangulation.  The continuum programme should not
depend on a favorable hidden choice of that triangulation.  This note
therefore enlarges the acquisition: it constructs a compact parameter
space containing every normalized face star allowed by declared
shape-regularity and ellipticity bounds.  A certificate on this universal
space covers every concrete base mesh satisfying those bounds.

\subsection{Normalized shape class}
\label{sec:newS7V42-shape-class}

Fix constants \(C_{\rm sh}\ge1\), \(N_{\rm inc}\in\mathbb N\), and
\(0<\lambda_-\le\lambda_+<\infty\).  A normalized small four-simplex is
represented by four edge columns \(R=(r_1,\ldots,r_4)\) after translation
of one vertex to zero and normalization of its largest edge to one.  Impose
\begin{equation}
 \|R\|\le C_{\rm sh},\qquad
 |\det R|\ge C_{\rm sh}^{-1},\qquad
 \|R^{-1}\|\le C_{\rm sh}.                             \label{eq:newS7V42-shape-bounds}
\end{equation}
The coframe/metric parameter obeys
\begin{equation}
 \lambda_-I\le g\le\lambda_+I.                        \label{eq:newS7V42-metric-bounds}
\end{equation}
All face identifications are required to be affine isometries on the common
geometric three-face with the prescribed orientation reversal.

\begin{definition}[Universal face-star class]
\label{def:newS7V42-universal-class}
\(\mathfrak U_{\rm face}\) is the set of all V41 ICM records based at one
open original three-face which:
\begin{enumerate}
\item have at most \(N_{\rm inc}\) incident small simplices at every vertex;
\item satisfy FIVC;
\item satisfy \eqref{eq:newS7V42-shape-bounds} and
\eqref{eq:newS7V42-metric-bounds};
\item use the fixed edgewise/Freudenthal subdivision rule on both sides;
and
\item are identified modulo global translation, tangential period choice,
and orientation-preserving relabelling.
\end{enumerate}
\end{definition}

\begin{lemma}[A shape bound supplies an incidence bound]
\label{lem:newS7V42-incidence-bound}
In fixed dimension four, a uniform lower bound on simplex shape quality
supplies a finite upper bound \(N_{\rm inc}=N_{\rm inc}(C_{\rm sh})\) on
the number of shape-regular simplices meeting a vertex in a conforming
triangulation.
\end{lemma}

\begin{proof}
Intersect every incident simplex with a sufficiently small sphere about
the vertex.  Shape regularity gives a positive lower bound, depending only
on \(C_{\rm sh}\), for the spherical three-volume of its solid-angle
region.  The interiors of these regions are disjoint and their total
volume is at most the volume of \(S^3\).  Dividing the latter by the former
gives the required finite bound.
\end{proof}

\begin{theorem}[Finite combinatorics and compact geometry]
\label{thm:newS7V42-compactness}
The universal face-star class is a finite disjoint union of compact
geometric parameter spaces, one for each abstract incidence type.
\end{theorem}

\begin{proof}
The incidence bound and the fixed dimension bound the number of vertices,
edges, faces, and four-simplices in a closed star of the face fundamental
cell.  Only finitely many finite abstract complexes exist on a bounded
number of labelled vertices, and FIVC selects a subset of them.

For a fixed abstract type, translate one vertex to zero and use the edge
normalization.  Equation \eqref{eq:newS7V42-shape-bounds} bounds every
coordinate and keeps every required determinant away from zero.
Face-matching equations and the non-strict versions of the declared
bounds are closed.  Thus the coordinate set is closed and bounded in a
finite-dimensional Euclidean space.  The metric interval
\eqref{eq:newS7V42-metric-bounds} is compact.  Taking the quotient by the
finite relabelling group preserves compactness.  The finite disjoint union
is compact.
\end{proof}

\begin{corollary}[Every admissible base mesh embeds]
\label{cor:newS7V42-embedding}
Let \(\mathcal K_0\) be any finite ordered base triangulation whose
edgewise subdivisions obey the declared constants.  Every normalized open
face tangent ICM of \(\mathcal K_0\) is a point of
\(\mathfrak U_{\rm face}\).
\end{corollary}

\begin{proof}
V33 supplies face compatibility and the finite star structure; the
declared shape, incidence, and ellipticity bounds supply the remaining
conditions in \Cref{def:newS7V42-universal-class}.
\end{proof}

\subsection{Continuity of the complete face certificate}
\label{sec:newS7V42-continuity}

\begin{lemma}[Continuity of the assembled blocks]
\label{lem:newS7V42-block-continuity}
On each compact incidence component of \(\mathfrak U_{\rm face}\), the
quantities
\[
 \nu_{\mathfrak s},\quad b_{\mathfrak s,x},\quad
 \omega_x,\quad \mathbb K_{xy},\quad\mathbb L_{xy},\quad
 A_j^{\rm sh}(m,E,k_\parallel)
\]
are continuous in the geometric and coframe parameters and finite Laurent
polynomials in \(e^{ik_\parallel}\).
\end{lemma}

\begin{proof}
The V41 reconstruction uses determinants and inverse edge matrices.
Their denominators stay away from zero by
\eqref{eq:newS7V42-shape-bounds}.  The lumped weights stay positive by the
volume and incidence bounds.  The V40 assembly uses finite sums, products,
and positive square roots of these quantities.  Tangential translation
adds only finitely many Laurent monomials.
\end{proof}

\begin{lemma}[Continuous universal references]
\label{lem:newS7V42-reference-continuity}
The V38 type-to-canonical and positive-polar paths can be chosen on a
finite compact chart cover so that the V37 adiabatic reference blocks,
their gap constants, and their polynomial inverse surrogates are continuous
on each chart.
\end{lemma}

\begin{proof}
The type-to-canonical path is algebraic in the endpoint coefficients.
The positive polar factor
\((B^*B)^{1/2}\) is continuous and uniformly positive on the compact
invertible coframe set.  Orthogonal Clifford conjugation is well-defined
independently of the sign of its spin lift; a finite spin-chart cover
chooses continuous lifts locally.  The V37 midpoint construction uses
finite sums of these coefficients.  The gap is uniform by V38, and
polynomial functional calculus is continuous in the operator
coefficients.
\end{proof}

\begin{corollary}[Compact universal PFC family]
\label{cor:newS7V42-universal-PFC-family}
For every fixed polynomial degree \(n\), the complete matrices
\[
 M_{F,n}(x),\qquad x\in\mathfrak U_{\rm face}\times\mathbb T^3,
\]
form a finite collection of continuous finite-dimensional matrix families
on compact parameter spaces.
\end{corollary}

\begin{proof}
Combine \Cref{thm:newS7V42-compactness,
lem:newS7V42-block-continuity,
lem:newS7V42-reference-continuity} with the V39 construction.
\end{proof}

\subsection{The universal face gap card}
\label{sec:newS7V42-UFGC}

\begin{definition}[Universal face gap card, UFGC]
\label{def:newS7V42-UFGC}
UFGC is the strict inequality
\begin{equation}
 \inf_{x\in\mathfrak U_{\rm face}}\inf_{k_\parallel\in\mathbb T^3}
 \sigma_{\min}M_F(x,k_\parallel)>0.                   \label{eq:newS7V42-UFGC}
\end{equation}
\end{definition}

\begin{theorem}[Universal certification theorem]
\label{thm:newS7V42-universal-certificate}
If UFGC is true, then a finite V39 polynomial face certificate exists on
\(\mathfrak U_{\rm face}\).  That single certificate proves FBSC and
C1FGC for every base triangulation satisfying the declared shape and
ellipticity bounds.
\end{theorem}

\begin{proof}
The universal parameter space is a finite union of compact components.
UFGC supplies one positive exact lower bound.  The V39 approximation error
is uniform because \(\delta_0,M_0,v_*\) have finite universal bounds by
compactness.  Apply the eventual-termination theorem
\Cref{thm:newS7V39-termination} on each component and take the finite union
of the resulting cell covers.  Every concrete mesh embeds by
\Cref{cor:newS7V42-embedding}.
\end{proof}

\begin{corollary}[A universal counterexample is decisive]
\label{cor:newS7V42-counterexample}
If one point of \(\mathfrak U_{\rm face}\) has \(\det M_F=0\), then UFGC is
false.  Such a point need not invalidate a particular \(\mathcal K_0\), but
it proves that shape regularity alone cannot close its face gap.
\end{corollary}

\begin{proof}
The infimum in \eqref{eq:newS7V42-UFGC} then contains a zero.
\end{proof}

\subsection{Why the crude Neumann bound cannot settle UFGC}
\label{sec:newS7V42-neumann}

The V40 universal defect estimate is
\(v_*=M_{\rm sh}+M_0\).  It deliberately ignores cancellation between the
sharp and reference blocks.

\begin{proposition}[Absolute row bounds lie outside the Neumann region]
\label{prop:newS7V42-row-no-go}
For the universal Wilson family,
\begin{equation}
 v_*\ge2\rho,\qquad \delta_0\le\rho.                  \label{eq:newS7V42-row-no-go}
\end{equation}
Consequently the crude estimate \(v_*<\delta_0\) required by the V35
relative-norm card can never follow from these separate row bounds.
\end{proposition}

\begin{proof}
Both row bounds \(M_{\rm sh}\) and \(M_0\) contain the diagonal Wilson mass
contribution \(\rho\), so their declared sum is at least \(2\rho\).
The essential spectrum of a two-tail reference contains the spectrum of
each limiting bulk.  At full bulk momentum \(k=0\), the Dirac and stiffness
symbols vanish and the sign-kernel magnitude is \(\rho\).  Hence the
distance of the reference spectrum from zero is at most \(\rho\), so
\(\delta_0\le\rho\).
\end{proof}

This no-go concerns the separate absolute upper bounds, not the true defect
norm.  Exact spin, incidence, and moment cancellations can make
\(H_F^{\rm sh}-H_0\) much smaller.  They must be retained in the full
matrix, in agreement with the V40 audit.

\subsection{A finite alternative even when UFGC fails}
\label{sec:newS7V42-strata}

\begin{definition}[Gap-safe face subspace]
\label{def:newS7V42-safe-subspace}
For \(\mu>0\), define
\begin{equation}
 \mathfrak U_{\rm face}^{(\mu)}
 :=\{x\in\mathfrak U_{\rm face}:
 \inf_{k_\parallel}\sigma_{\min}M_F(x,k_\parallel)\ge\mu\}.            \label{eq:newS7V42-safe-subspace}
\end{equation}
\end{definition}

\begin{proposition}[Safe subspaces are closed and certifiable]
\label{prop:newS7V42-safe-closed}
Every \(\mathfrak U_{\rm face}^{(\mu)}\) is compact.  A sufficiently refined
PFC certifies any compact subset lying in its interior.
\end{proposition}

\begin{proof}
The minimum singular value is continuous in all finite parameters after
the exact V36 reduction and the continuous-reference chart construction.
The defining superlevel set is therefore closed in a compact space.
The V39 approximation and cell argument apply with the fixed margin
\(\mu\).
\end{proof}

Thus, if shape regularity is too broad, the programme can declare a
verifiable open mesh-admissibility condition rather than choose an
unrecorded favorable triangulation.

\begin{definition}[Certified mesh-admissibility card, CMAC]
\label{def:newS7V42-CMAC}
CMAC is membership of every face-star parameter of the chosen base mesh in
a finite PFC-certified compact subset of
\(\mathfrak U_{\rm face}^{(\mu)}\) for one \(\mu>0\).
\end{definition}

\begin{corollary}[CMAC is sufficient for the face row]
\label{cor:newS7V42-CMAC}
CMAC implies FBSC, C1FGC, and the codimension-one SFGC row for the chosen
base mesh, with constants uniform along its cofinal subdivisions.
\end{corollary}

\begin{proof}
There are finitely many original faces and their normalized edgewise
parameters remain in the certified sets.  Apply the PFC implication of V39
and the V36 face-gap theorem.
\end{proof}

\begin{assessment}[The missing data have become a universal alternative]
\label{ass:newS7V42-status}
MIDC is no longer logically necessary for formulating a complete
certificate: UFGC is a compact universal problem, and CMAC is a
mesh-specific open alternative if UFGC fails.  Neither has been proved.
The absolute row-sum route is rigorously insufficient, so the next progress
must retain exact interface cancellations or modify the face discretization
within the same continuum class.
\end{assessment}

\begin{programme}[V43: adiabatic face replacement]
\label{prog:newS7V42-V43}
V43 will pursue the second route.  It will replace the sharp stencil only
inside a collar of \(L(a)\) lattice layers by the already-gapped V37
reference, with
\[
 L(a)\longrightarrow\infty,\qquad aL(a)\longrightarrow0.
\]
It will prove the open-face tangent gap by construction and establish the
precise moment and norm conditions under which this collar replacement has
the same smooth continuum Dirac--Wilson limit.  Intersections of face
collars at edges and vertices will remain separate strata.
\end{programme}

\begin{firewall}[Claim boundary after seventh-series V42]
\label{fire:newS7V42-boundary}
V42 constructs the compact universal shape-regular face space, proves
continuity of all PFC data, proves that UFGC admits a finite certificate,
shows every admissible base mesh embeds, proves the crude row-bound
Neumann no-go, and formulates the certifiable mesh-admissibility
alternative CMAC.

V42 does not prove UFGC or CMAC, exclude all sharp-face surface modes,
justify a modified face stencil, prove higher-codimension tangent gaps,
full SFGC, the global finite-element overlap gap, strong local heat
consistency, Gaussian heat locality, the common Einstein stress, pure
gravitational owners, rough metric or AOP control, measure concentration,
the determinant-line phase, interacting cutoff removal, or the Standard
Model--Einstein continuum.
\end{firewall}

\begin{theorem}[Seventh-series V42 result]
\label{thm:newS7V42-result}
All shape-regular V33 open-face tangents with fixed geometric bounds form a
finite union of compact parameter families.  If their exact face matrices
are uniformly invertible, the V39 refinement theorem produces one finite
certificate covering every admissible base triangulation.  Separate
absolute row bounds cannot prove that invertibility because their Neumann
ratio is at least two.
\end{theorem}

\begin{proof}
Use \Cref{lem:newS7V42-incidence-bound,
thm:newS7V42-compactness,
cor:newS7V42-embedding,
lem:newS7V42-block-continuity,
lem:newS7V42-reference-continuity,
thm:newS7V42-universal-certificate,
prop:newS7V42-row-no-go,
cor:newS7V42-CMAC}.
\end{proof}

\paragraph{Parent-version record.}
V42 was formed from
\texttt{Renewal\_SM\_Einstein\_Continuum\_Research\_Notes\_V41.tex}.
The V41 parent had \(16\,589\) logical source lines, \(668\,767\) bytes,
and SHA--256
\[
 \texttt{C54883E08FDC32E403EBC2DD79785694C4844BA55EE129D7473F50F852313E25}.
\]
The next compulsory companion audit is V45.

% =============================================================================
% END APPENDED SEVENTH-SERIES V42 MATERIAL
% =============================================================================

% =============================================================================
% BEGIN APPENDED SEVENTH-SERIES V43 MATERIAL
% =============================================================================

\section{Seventh-series V43: adiabatic face replacement and its moment gate}
\label{sec:v7v43}
\label{sec:newS7V43}

\subsection{The proposed replacement}
\label{sec:newS7V43-replacement}

Let \(F^\circ\) be the interior of an original three-face, away from its
two- and lower-dimensional boundary strata.  Choose integers \(L(a)\) such
that
\begin{equation}
 L(a)\longrightarrow\infty,\qquad
 w(a):=aL(a)\longrightarrow0.                         \label{eq:newS7V43-scales}
\end{equation}
In a normal collar of \(F^\circ\) with \(L(a)\) lattice layers, replace the
sharp-face dimensionless sign kernel \(H_a^{\rm sh}\) by the V37 adiabatic
two-tail reference \(H_{a,L}^{\rm ad}\).  Outside a slightly larger collar,
retain \(H_a^{\rm sh}\).  Tangential cutoff zones near
\(\partial F^\circ\) belong to the edge and vertex strata and are not
included in the open-face assertion.

\begin{definition}[Open-face adiabatic kernel]
\label{def:newS7V43-open-face-kernel}
On a face chart whose tangential support stays a fixed positive physical
distance from \(\partial F^\circ\), define
\begin{equation}
 \widetilde H_a=H_{a,L}^{\rm ad}.                     \label{eq:newS7V43-Htilde-core}
\end{equation}
On the two bulk sides beyond the collar,
\(\widetilde H_a=H_a^{\rm sh}\).  A global definition will require
additional edge and vertex collar rules.
\end{definition}

\begin{theorem}[Open-face tangent gap by construction]
\label{thm:newS7V43-open-face-gap}
For every compact tangential subset of \(F^\circ\), the frozen tangent
operators of \(\widetilde H_a\) have a uniform gap
\begin{equation}
 \|\widetilde H_a^{\rm tan}u\|
 \ge\left(\delta_{\rm path}-C L(a)^{-1/2}\right)\|u\|. \label{eq:newS7V43-tangent-gap}
\end{equation}
For sufficiently small \(a\), this is at least
\(\delta_{\rm path}/2\).
\end{theorem}

\begin{proof}
At a frozen tangential point, \(\widetilde H_a^{\rm tan}\) is precisely the
midpoint-quantized V37 interface formed from the uniformly gapped V38 bulk
path.  Apply \Cref{thm:newS7V37-adiabatic}.  Equation
\eqref{eq:newS7V43-scales} makes \(L(a)^{-1/2}\) tend to zero.
\end{proof}

\begin{proposition}[Tangentially varying open-face gap]
\label{prop:newS7V43-varying-gap}
Suppose the smooth geometric coefficients vary by \(O(a)\) between
neighbouring tangential cells.  Localize on boxes of \(\ell(a)\) lattice
units, where
\begin{equation}
 \ell(a)\longrightarrow\infty,\qquad
 {\ell(a)\over L(a)}+a\ell(a)\longrightarrow0.         \label{eq:newS7V43-three-scales}
\end{equation}
Then, away from the tangential cutoff zones,
\begin{equation}
 \|\widetilde H_a u\|
 \ge\left(
 \delta_{\rm path}
 -C\{\ell(a)^{-1}+\ell(a)/L(a)+a\ell(a)\}
 \right)\|u\|.                                        \label{eq:newS7V43-varying-gap}
\end{equation}
\end{proposition}

\begin{proof}
Use a quadratic partition of unity in the normal and tangential lattice
coordinates.  The finite-range commutator costs \(C/\ell\), exactly as in
V37.  Freezing the normal path over a box costs \(C\ell/L\), while freezing
the smooth tangential geometry costs \(Ca\ell\).  The frozen box operator
has gap \(\delta_{\rm path}\).  The direct-sum Minkowski argument from
\Cref{thm:newS7V37-adiabatic} gives
\eqref{eq:newS7V43-varying-gap}.
\end{proof}

For example, \(L(a)=\lfloor a^{-1/2}\rfloor\) and
\(\ell(a)=\lfloor a^{-1/4}\rfloor\) satisfy both
\eqref{eq:newS7V43-scales} and
\eqref{eq:newS7V43-three-scales}.

\subsection{The smooth-continuum comparison}
\label{sec:newS7V43-continuum}

The physical sign-kernel scale is \(a^{-1}\).  Put
\begin{equation}
 B_a:=\widetilde H_a-H_a^{\rm sh}.                    \label{eq:newS7V43-B}
\end{equation}
The common Wilson mass \(-\rho\Gamma\) cancels in this difference.  Write
its finite stencil in a smooth local trivialization as
\begin{equation}
 (B_a\psi)_x=\sum_{|y-x|\le Ca}B_{xy}^{(a)}\psi_y,
 \qquad \sup_x\sum_y\|B_{xy}^{(a)}\|\le C_B.           \label{eq:newS7V43-stencil}
\end{equation}
Define the zeroth and first discrete moments
\begin{align}
 Z_a(x)&:=\sum_yB_{xy}^{(a)},                          \label{eq:newS7V43-Z}\\
 \mathcal M_a(x)&:=a^{-1}\sum_y
 B_{xy}^{(a)}\otimes\exp_x^{-1}(y).                   \label{eq:newS7V43-M}
\end{align}

\begin{lemma}[Pointwise moment expansion]
\label{lem:newS7V43-moment-expansion}
For every \(C^2\) spinor in the local trivialization,
\begin{equation}
 a^{-1}(B_a\psi)_x
 =a^{-1}Z_a(x)\psi(x)
  +\mathcal M_a(x)\nabla\psi(x)
  +O(a)\|\psi\|_{C^2},                                \label{eq:newS7V43-expansion}
\end{equation}
uniformly in the open-face collar.
\end{lemma}

\begin{proof}
Parallel-transport \(\psi_y\) to \(x\) and use its covariant Taylor
expansion
\[
 \psi_y=\psi(x)+\nabla\psi(x)\exp_x^{-1}(y)
       +O(a^2)\|\psi\|_{C^2}.
\]
Insert this in \eqref{eq:newS7V43-stencil}, use the row bound, and divide
by \(a\).
\end{proof}

Let \(\mathcal C_a(F)\) denote the open-face collar.  Its physical
four-volume is bounded by \(C_Fw(a)=C_FaL(a)\).

\begin{theorem}[Strong smooth-core comparison]
\label{thm:newS7V43-strong-comparison}
If
\begin{equation}
 Z_a(x)=0,\qquad
 \sup_{a,x}\|\mathcal M_a(x)\|<\infty,                \label{eq:newS7V43-null-moment}
\end{equation}
then, for every fixed smooth spinor \(\psi\),
\begin{equation}
 \left\|a^{-1}B_a\psi\right\|_{L^2(\mathcal C_a(F))}
 \le C_\psi\,\{aL(a)\}^{1/2}\longrightarrow0.          \label{eq:newS7V43-strong-comparison}
\end{equation}
The same conclusion holds after one or two metric variations if the
corresponding moment and row bounds are uniform.
\end{theorem}

\begin{proof}
Under \eqref{eq:newS7V43-null-moment},
\Cref{lem:newS7V43-moment-expansion} bounds the physical difference
pointwise by a constant depending on \(\psi\), uniformly in the collar.
Multiplication by the square root of the collar volume gives
\eqref{eq:newS7V43-strong-comparison}.  Metric differentiation repeats the
same finite Taylor argument with the differentiated blocks.
\end{proof}

\begin{definition}[Collar null-moment card, CNMC]
\label{def:newS7V43-CNMC}
CNMC is the existence of a self-adjoint, spin--gauge covariant open-face
adiabatic replacement which:
\begin{enumerate}
\item has the V37--V38 tangent gap;
\item agrees exactly with the two sharp bulk tails;
\item satisfies the zeroth-moment and bounded first-moment conditions
\eqref{eq:newS7V43-null-moment}; and
\item satisfies those conditions through the two metric contacts required
by the determinant modulus.
\end{enumerate}
\end{definition}

\begin{corollary}[Conditional continuum safety of the replacement]
\label{cor:newS7V43-CNMC-safety}
Under CNMC and \eqref{eq:newS7V43-scales}, the open-face replacement has the
same action as the sharp V33 sign kernel on every fixed smooth core, up to
an \(L^2\) error tending to zero, while its open-face tangent gap is
uniformly positive.
\end{corollary}

\begin{proof}
Combine \Cref{thm:newS7V43-open-face-gap,
prop:newS7V43-varying-gap,
thm:newS7V43-strong-comparison}.
\end{proof}

\subsection{Why adiabatic variation alone does not prove CNMC}
\label{sec:newS7V43-no-go}

For every frozen path value \(s\), the derivative and stiffness parts of
the translation-invariant symbol annihilate a constant.  Midpoint
quantization samples different \(s\)-values in different hopping
directions, so the cancellations in a single frozen row need not remain
exact.

\begin{lemma}[Generic size of the midpoint zeroth moment]
\label{lem:newS7V43-generic-Z}
For a Lipschitz coefficient path with normal scale \(L\), midpoint
quantization alone gives only
\begin{equation}
 \|Z_a(x)\|\le {C\over L}.                             \label{eq:newS7V43-Z-generic}
\end{equation}
There is no formal reason for the coefficient of \(L^{-1}\) to vanish.
\end{lemma}

\begin{proof}
At a fixed path value, the row sum of the nonmass stencil is zero.
Replace every midpoint value \(s_{m+j/2}\) by \(s_m\); the resulting row
sum vanishes.  Lipschitz continuity changes each of the finitely many
coefficients by at most \(C|j|/L\), proving the estimate.  A path with a
nonzero derivative of one hopping coefficient and no compensating
derivative in its paired row realizes a nonzero first-order term.
\end{proof}

\begin{proposition}[Slow variation is insufficient for a shrinking collar]
\label{prop:newS7V43-slow-insufficient}
The generic bound \eqref{eq:newS7V43-Z-generic} does not imply either
strong or weak disappearance of the physical perturbation when
\(aL(a)\to0\).  Indeed, its zeroth-order contribution has the bounds
\begin{align}
 \left\|a^{-1}Z_a\psi\right\|_{L^2(\mathcal C_a(F))}
 &\lesssim {1\over\sqrt{aL(a)}}\|\psi\|_{L^\infty},    \label{eq:newS7V43-strong-bad}\\
 \left|\left\langle\eta,a^{-1}Z_a\psi\right\rangle_
 {L^2(\mathcal C_a(F))}\right|
 &\lesssim 1                                           \label{eq:newS7V43-weak-bad}
\end{align}
for fixed smooth \(\eta,\psi\).  The first bound diverges and the second
need not tend to zero.
\end{proposition}

\begin{proof}
The collar volume is \(O(aL)\).  Multiply the pointwise magnitude
\(O((aL)^{-1})\) by the square root of that volume to obtain
\eqref{eq:newS7V43-strong-bad}, and by the volume itself to obtain
\eqref{eq:newS7V43-weak-bad}.
\end{proof}

This is the precise obstruction to declaring the adiabatic reference to be
the physical discretization merely because its collar shrinks.

\subsection{Exact algebra which a CNMC construction must preserve}
\label{sec:newS7V43-algebra}

Let \(\mathbf 1\otimes v\) denote a locally constant nodal spinor.  Since
the common mass cancels in \(B_a\), the zeroth-moment condition is
equivalent to
\begin{equation}
 B_a(\mathbf 1\otimes v)=0\quad\hbox{for every spinor }v.              \label{eq:newS7V43-constant-kernel}
\end{equation}
It must coexist with
\begin{equation}
 \widetilde H_a^*=\widetilde H_a,\qquad
 \widetilde H_a\ \hbox{spin--gauge covariant},\qquad
 \|\widetilde H_a u\|\ge c\|u\|                       \label{eq:newS7V43-three-requirements}
\end{equation}
on the open-face tangent.  Enforcing
\eqref{eq:newS7V43-constant-kernel} by an arbitrary diagonal row correction
can destroy self-adjointness; symmetrizing an arbitrary row correction can
restore a zeroth moment.  The three requirements must therefore be
constructed together.

\begin{assessment}[Spectral acquisition gained; continuum gate exposed]
\label{ass:newS7V43-status}
The V37 reference gives a uniform open-face gap even with slowly varying
tangential geometry.  Replacing the sharp stencil by that reference is
continuum-safe under CNMC.  Midpoint adiabaticity alone gives only an
\(O(L^{-1})\) zeroth moment, which is too large on a collar whose physical
width tends to zero.  CNMC is therefore the first open row of the modified
discretization route.
\end{assessment}

\begin{programme}[V44: conservative self-adjoint collar forms]
\label{prog:newS7V43-V44}
V44 will construct the collar at the level of sesquilinear edge and
simplex forms rather than by rowwise symbol interpolation.  It will test
whether a discrete conservative form can preserve self-adjointness,
spin--gauge covariance, and the exact constant-mode moment while remaining
within \(o(1)\) of the gapped adiabatic reference.  If the required
correction is not norm-small, the programme will return to UFGC rather
than assume the gap survives.
\end{programme}

\begin{firewall}[Claim boundary after seventh-series V43]
\label{fire:newS7V43-boundary}
V43 proves the open-face adiabatic tangent gap, its slowly tangentially
varying version, the exact smooth-core moment expansion, continuum safety
under CNMC, the generic \(O(L^{-1})\) zeroth-moment estimate, and the
strong/weak no-go showing that this estimate alone is insufficient.

V43 does not prove CNMC, justify the adiabatic collar as a final physical
discretization, handle collar intersections at edges or vertices, prove
UFGC or CMAC for the sharp stencil, full SFGC, the global overlap gap,
strong local heat consistency, Gaussian heat locality, the common Einstein
stress, pure gravitational owners, rough metric or AOP control, measure
concentration, the determinant-line phase, interacting cutoff removal, or
the Standard Model--Einstein continuum.
\end{firewall}

\begin{theorem}[Seventh-series V43 result]
\label{thm:newS7V43-result}
The adiabatic replacement has a uniform open-face tangent gap, but it shares
the sharp V33 smooth continuum action on a shrinking collar only after an
exact zeroth-moment cancellation and bounded first moment.  The generic
midpoint construction misses that cancellation by \(O(L^{-1})\), a size
which can survive weakly and diverge strongly after the physical
\(a^{-1}\) scaling.
\end{theorem}

\begin{proof}
Use \Cref{thm:newS7V43-open-face-gap,
prop:newS7V43-varying-gap,
lem:newS7V43-moment-expansion,
thm:newS7V43-strong-comparison,
lem:newS7V43-generic-Z,
prop:newS7V43-slow-insufficient}.
\end{proof}

\paragraph{Parent-version record.}
V43 was formed from
\texttt{Renewal\_SM\_Einstein\_Continuum\_Research\_Notes\_V42.tex}.
The V42 parent had \(16\,957\) logical source lines, \(683\,875\) bytes,
and SHA--256
\[
 \texttt{6272340736BFABE608EDDA44644A5F49D434BA01C8538C3212A127989D876B06}.
\]
The next compulsory companion audit is V45.

% =============================================================================
% END APPENDED SEVENTH-SERIES V43 MATERIAL
% =============================================================================

% =============================================================================
% BEGIN APPENDED SEVENTH-SERIES V44 MATERIAL
% =============================================================================

\section{Seventh-series V44: a conservative self-adjoint collar completion}
\label{sec:v7v44}
\label{sec:newS7V44}

\subsection{Finite normal fibre at zero tangential momentum}
\label{sec:newS7V44-fibre}

Fix a frozen open-face parameter and let
\[
 B_L(k_\parallel):=H_L^{\rm ad}(k_\parallel)
                   -H^{\rm sh}(k_\parallel)
\]
be the dimensionless difference between the V37 adiabatic interface and
the sharp V33 face.  It is self-adjoint and supported in a finite set
\(S_L\) of normal layers.  Enlarge \(S_L\), if necessary, so that it
contains every nonzero block row and column.  Write
\[
 N_L:=|S_L|,\qquad c_1L\le N_L\le c_2L.
\]
The common Wilson mass cancels in \(B_L\).

Let \(q\) be the spin--internal fibre dimension and define the normalized
constant-layer isometry
\begin{equation}
 Q_L:\mathbb C^q\longrightarrow
 \ell^2(S_L;\mathbb C^q),\qquad
 (Q_Lv)_m=N_L^{-1/2}v.                                \label{eq:newS7V44-Q}
\end{equation}
At zero tangential momentum, set
\begin{equation}
 Y_L:=-B_L(0)Q_L.                                     \label{eq:newS7V44-Y}
\end{equation}

\begin{lemma}[Size of the constant-mode defect]
\label{lem:newS7V44-Y-bound}
The midpoint estimate of V43 implies
\begin{equation}
 \|Y_L\|\le {C_Y\over L}.                             \label{eq:newS7V44-Y-bound}
\end{equation}
The same estimate holds after the two declared metric variations.
\end{lemma}

\begin{proof}
Let
\[
 z_m:=\sum_nB_L(0)_{mn}.
\]
By \Cref{lem:newS7V43-generic-Z}, \(\|z_m\|\le C/L\), and only
\(N_L=O(L)\) rows occur.  For \(\|v\|=1\),
\[
 \|B_L(0)Q_Lv\|^2
 ={1\over N_L}\sum_{m\in S_L}\|z_mv\|^2
 \le {1\over N_L}N_L{C^2\over L^2}.
\]
Metric differentiation preserves the Lipschitz path and finite-row
bounds, so the same calculation applies to the differentiated residuals.
\end{proof}

\subsection{The self-adjoint completion}
\label{sec:newS7V44-completion}

\begin{definition}[Conservative collar correction]
\label{def:newS7V44-R}
Define on the normal fibre, and extend by zero outside \(S_L\),
\begin{equation}
 R_L:=Y_LQ_L^*+Q_LY_L^*
      -Q_L(Q_L^*Y_L)Q_L^*.                            \label{eq:newS7V44-R}
\end{equation}
Use the same normal-layer matrix at every tangential lattice translate, so
\(R_L\) has zero tangential displacement and is independent of
\(k_\parallel\).
\end{definition}

\begin{lemma}[Hermitian compatibility]
\label{lem:newS7V44-compatibility}
The finite matrix \(Q_L^*Y_L\) is self-adjoint.
\end{lemma}

\begin{proof}
By \eqref{eq:newS7V44-Y},
\[
 Q_L^*Y_L=-Q_L^*B_L(0)Q_L.
\]
The zero-momentum fibre \(B_L(0)\) is self-adjoint, so its compression is
self-adjoint.
\end{proof}

\begin{theorem}[Exact self-adjoint moment completion]
\label{thm:newS7V44-completion}
The correction \(R_L\) is self-adjoint, is supported in \(S_L\times S_L\),
and satisfies
\begin{align}
 (B_L(0)+R_L)Q_L&=0,                                  \label{eq:newS7V44-moment-zero}\\
 \|R_L\|&\le3\|Y_L\|\le {3C_Y\over L}.                \label{eq:newS7V44-R-bound}
\end{align}
\end{theorem}

\begin{proof}
Self-adjointness follows from \Cref{lem:newS7V44-compatibility}.  Moreover,
\[
 R_LQ_L
 =Y_L+Q_LY_L^*Q_L-Q_LQ_L^*Y_L.
\]
Since \(Y_L^*Q_L=Q_L^*Y_L\), the last two terms cancel, so
\(R_LQ_L=Y_L=-B_L(0)Q_L\).  This proves
\eqref{eq:newS7V44-moment-zero}.  Finally,
\[
 \|R_L\|
 \le\|Y_LQ_L^*\|+\|Q_LY_L^*\|
    +\|Q_L(Q_L^*Y_L)Q_L^*\|
 \le3\|Y_L\|.
\]
Use \Cref{lem:newS7V44-Y-bound}.
\end{proof}

\begin{corollary}[The completion is finite rank in the normal fibre]
\label{cor:newS7V44-rank}
For each \(k_\parallel\),
\begin{equation}
 \operatorname{rank}R_L\le2q.                         \label{eq:newS7V44-rank}
\end{equation}
\end{corollary}

\begin{proof}
The range of \eqref{eq:newS7V44-R} lies in
\(\operatorname{Ran}Y_L+\operatorname{Ran}Q_L\), whose dimension is at
most \(2q\).
\end{proof}

\subsection{Gap preservation}
\label{sec:newS7V44-gap}

Define the completed adiabatic sign kernel
\begin{equation}
 H_L^{\rm con}(k_\parallel)
 :=H_L^{\rm ad}(k_\parallel)+R_L.                     \label{eq:newS7V44-Hcon}
\end{equation}

\begin{theorem}[Conservative open-face gap]
\label{thm:newS7V44-gap}
Uniformly in \(k_\parallel\) and the compact frozen geometric parameters,
\begin{equation}
 \|H_L^{\rm con}(k_\parallel)u\|
 \ge
 \left(
 \delta_{\rm path}-C L^{-1/2}-{3C_Y\over L}
 \right)\|u\|.                                        \label{eq:newS7V44-gap}
\end{equation}
Thus the completed kernel has gap at least
\(\delta_{\rm path}/2\) for all sufficiently large \(L\).
\end{theorem}

\begin{proof}
The V37 estimate gives
\[
 \|H_L^{\rm ad}u\|
 \ge(\delta_{\rm path}-CL^{-1/2})\|u\|.
\]
Apply the reverse triangle inequality and
\eqref{eq:newS7V44-R-bound}.
\end{proof}

\begin{corollary}[The overlap sign changes by \(O(L^{-1})\)]
\label{cor:newS7V44-sign-change}
Once both kernels have gap at least \(\delta_{\rm path}/2\),
\begin{equation}
 \|\operatorname{sgn}(H_L^{\rm con})
   -\operatorname{sgn}(H_L^{\rm ad})\|
 \le {C\over\delta_{\rm path}}\|R_L\|
 \le {C'\over L}.                                     \label{eq:newS7V44-sign-change}
\end{equation}
\end{corollary}

\begin{proof}
Use the resolvent integral for the sign function on two self-adjoint
operators whose spectra avoid
\((-\delta_{\rm path}/2,\delta_{\rm path}/2)\), insert the resolvent
identity, and integrate the product of the two resolvent bounds.  Then use
\eqref{eq:newS7V44-R-bound}.
\end{proof}

\subsection{First-moment control despite the long normal range}
\label{sec:newS7V44-first-moment}

Although \(R_L\) connects all layers in \(S_L\), its blocks are small.
Writing \(Q_{L,m}=N_L^{-1/2}I\) and
\(Y_{L,m}=-N_L^{-1/2}z_m\), one has
\begin{equation}
 \|Q_{L,m}\|\le C L^{-1/2},\qquad
 \|Y_{L,m}\|\le C L^{-3/2}.                           \label{eq:newS7V44-block-scales}
\end{equation}

\begin{lemma}[Weighted row bound for the completion]
\label{lem:newS7V44-weighted-row}
There is a constant independent of \(L\) such that
\begin{equation}
 \sup_{m\in S_L}
 \sum_{n\in S_L}(1+|m-n|)\|(R_L)_{mn}\|\le C_R.       \label{eq:newS7V44-weighted-row}
\end{equation}
The unweighted row sum is \(O(L^{-1})\).
\end{lemma}

\begin{proof}
Each block in the first two terms of \eqref{eq:newS7V44-R} is bounded by
\(CL^{-2}\) using \eqref{eq:newS7V44-block-scales}.  The compressed matrix
\(Q_L^*Y_L\) is \(O(L^{-1})\), so each block of the third term is also
\(O(L^{-2})\).  There are \(O(L)\) terms in an unweighted row and
\(\sum_{n\in S_L}|m-n|=O(L^2)\), proving both assertions.
\end{proof}

\begin{corollary}[Bounded physical first moment]
\label{cor:newS7V44-first-moment}
The physical first moment of \(R_L\),
\begin{equation}
 a^{-1}\sum_n(R_L)_{mn}\otimes\{a(n-m)\},
\end{equation}
is uniformly bounded.  Its physical interaction range is
\(O(aL)=O(w(a))\), which tends to zero under
\eqref{eq:newS7V43-scales}.
\end{corollary}

\begin{proof}
Cancel the factors \(a^{-1}\) and \(a\), then apply the weighted row bound.
\end{proof}

\subsection{Covariant curved lift}
\label{sec:newS7V44-curved}

On a smooth face collar, identify the normal-layer fibres with the central
layer by the same short-geodesic spin--gauge transports used in V33.
Define \(Q_L\) by transporting one central fibre vector to every layer
before multiplying by \(N_L^{-1/2}\).  Form \(Y_L\) and \(R_L\) in that
transported frame and transport the resulting blocks back to their
endpoints.

\begin{lemma}[Spin--gauge covariance]
\label{lem:newS7V44-covariance}
The curved construction of \(Q_L,Y_L,R_L\), and \(H_L^{\rm con}\) is
covariant under changes of spin and Standard Model gauge frame.
\end{lemma}

\begin{proof}
Endpoint frame changes conjugate every short-geodesic transport.  Therefore
\(Q_L\) intertwines the central and nodal frame actions.  The defect
\(B_L\) is covariant, hence so is \(Y_L=-B_LQ_L\).  Every product in
\eqref{eq:newS7V44-R} has the required endpoint conjugations.
\end{proof}

\begin{lemma}[Curved zeroth-moment tolerance]
\label{lem:newS7V44-curved-moment}
After applying the frozen completion pointwise along the face, smooth
tangential variation leaves
\begin{equation}
 \|Z_a^{\rm con}(x)\|\le C a,                         \label{eq:newS7V44-curved-Z}
\end{equation}
while the first moment stays uniformly bounded.  The statement holds
through two metric contacts on a bounded smooth family.
\end{lemma}

\begin{proof}
At a frozen tangential point,
\eqref{eq:newS7V44-moment-zero} cancels the leading zeroth moment exactly.
Neighbouring tangential coefficient and transport data differ by \(O(a)\);
only finitely many tangential shifts occur, so their uncancelled row sum is
\(O(a)\).  The V43 stencil bound and
\Cref{lem:newS7V44-weighted-row} give the first-moment bound.  Differentiated
coframes, transports, and polar paths remain uniformly smooth, so the same
argument applies to the two metric contacts.
\end{proof}

\begin{correction}[Curved tolerance in the V43 comparison]
\label{corr:newS7V44-curved-tolerance}
The exact condition \(Z_a=0\) in
\eqref{eq:newS7V43-null-moment} is sufficient but stronger than necessary
for a curved collar.  The estimate \(\|Z_a\|\le Ca\) is also sufficient:
the first term in \eqref{eq:newS7V43-expansion} is then uniformly bounded,
and its \(L^2\) norm on the collar is \(O((aL)^{1/2})\).
\end{correction}

\begin{theorem}[Open-face CNMC acquisition]
\label{thm:newS7V44-CNMC}
The conservative completion satisfies CNMC in every frozen open-face
tangent and its curved \(O(a)\)-tolerance through two metric contacts.
For any scales satisfying \eqref{eq:newS7V43-scales} and
\eqref{eq:newS7V43-three-scales}, it has:
\begin{enumerate}
\item a uniform open-face sign-kernel gap;
\item the same smooth-core continuum action as the sharp V33 kernel, with
error \(O((aL)^{1/2})\);
\item exact spin--gauge covariance and gamma-five Hermiticity; and
\item physical interaction range tending to zero.
\end{enumerate}
\end{theorem}

\begin{proof}
The gap is \Cref{thm:newS7V44-gap} combined with the tangential freezing
argument of \Cref{prop:newS7V43-varying-gap}.  The frozen zeroth moment is
\eqref{eq:newS7V44-moment-zero}; its curved tolerance is
\Cref{lem:newS7V44-curved-moment}.  Insert that tolerance and the bounded
first moment into the V43 Taylor proof, using
\Cref{corr:newS7V44-curved-tolerance}.  Covariance is
\Cref{lem:newS7V44-covariance}.  Self-adjointness of
\(H_L^{\rm con}\) means that the associated Wilson kernel retains
gamma-five Hermiticity.  The range statement is
\Cref{cor:newS7V44-first-moment}.
\end{proof}

\subsection{What has moved to the lower strata}
\label{sec:newS7V44-strata}

The completed collar is defined on compact subsets of an open original
face.  Near an original edge, several face collars meet and no single
normal coordinate or constant-layer isometry covers all of them.  The
construction therefore closes the codimension-one tangent row for the
modified discretization but does not yet define a global operator at the
codimension-two skeleton.

\begin{assessment}[Open faces acquired by a continuum-safe modification]
\label{ass:newS7V44-status}
V44 proves CNMC and the open-face tangent gap for the conservative
adiabatic discretization.  This bypasses the unproved sharp-face UFGC while
retaining the same smooth-core continuum action.  Its low-rank completion
has increasing lattice range, but vanishing physical range and vanishing
operator norm relative to the adiabatic kernel.  Uniform Gaussian kernel
locality has not been proved for the completed family.
\end{assessment}

\begin{programme}[V45: compulsory audit and codimension-two collar model]
\label{prog:newS7V44-V45}
V45 will perform the compulsory YM/NS audit, then formulate the tangent
model where several completed face collars meet at an original edge.  It
will test a two-normal adiabatic interpolation on a sector fan, identify
its joint moment-completion space, and separate its gap card from the
open-face theorem proved here.
\end{programme}

\begin{firewall}[Claim boundary after seventh-series V44]
\label{fire:newS7V44-boundary}
V44 constructs an explicit self-adjoint rank-\(2q\) normal-fibre
completion, proves its \(O(L^{-1})\) norm, exact frozen zeroth moment,
bounded first moment, gap preservation, sign-function stability,
spin--gauge covariance, curved \(O(a)\) tolerance, and open-face CNMC.

V44 does not handle intersections of face collars at edges or vertices,
prove Gaussian locality for the growing lattice-range completion, prove
the sharp-stencil UFGC, full SFGC, the global overlap gap, strong local heat
consistency, the common Einstein stress, pure gravitational owners, rough
metric or AOP control, measure concentration, the determinant-line phase,
interacting cutoff removal, or the Standard Model--Einstein continuum.
\end{firewall}

\begin{theorem}[Seventh-series V44 result]
\label{thm:newS7V44-result}
The \(O(L^{-1})\) midpoint row defect admits the explicit self-adjoint
completion \eqref{eq:newS7V44-R}.  The completion cancels the frozen
zeroth moment exactly, preserves a uniform open-face gap, has bounded
physical first moment and vanishing physical range, and yields the same
smooth-core continuum action through two metric contacts.
\end{theorem}

\begin{proof}
Use \Cref{lem:newS7V44-Y-bound,
lem:newS7V44-compatibility,
thm:newS7V44-completion,
thm:newS7V44-gap,
cor:newS7V44-sign-change,
lem:newS7V44-weighted-row,
lem:newS7V44-covariance,
lem:newS7V44-curved-moment,
thm:newS7V44-CNMC}.
\end{proof}

\paragraph{Parent-version record.}
V44 was formed from
\texttt{Renewal\_SM\_Einstein\_Continuum\_Research\_Notes\_V43.tex}.
The V43 parent had \(17\,291\) logical source lines, \(697\,455\) bytes,
and SHA--256
\[
 \texttt{B63AC090512C073588AA0CFBD3F2DA1B3C0B740C114E122B7D7345DD3A9E7606}.
\]
The next compulsory companion audit is V45.

% =============================================================================
% END APPENDED SEVENTH-SERIES V44 MATERIAL
% =============================================================================

% =============================================================================
% BEGIN APPENDED SEVENTH-SERIES V45 MATERIAL
% =============================================================================

\section{Seventh-series V45: compulsory audit and the codimension-two fan}
\label{sec:v7v45}
\label{sec:newS7V45}

\subsection{Compulsory companion audit}
\label{sec:newS7V45-audit}

\begin{audit}[Companion snapshots used at V45]
\label{aud:newS7V45-snapshots}
The exact active companion snapshots read before this note were:
\begin{center}
\begin{tabular}{p{0.20\linewidth}p{0.52\linewidth}r}
programme & source & bytes\\ \hline
Yang--Mills &
\texttt{Renewal\_Yang\_Mills\_Mass\_Gap\_Research\_Notes\_V21.tex}
& \(323\,789\)\\
Navier--Stokes &
\texttt{Renewal\_NS\_Stability\_Research\_Notes\_V75.tex}
& \(764\,980\)
\end{tabular}
\end{center}
The Yang--Mills source had \(9\,142\) newline characters and SHA--256
\[
 \texttt{4393EA3951991D7E1BCAA67210893033D54B9B272B3260A38D1F90A3989BF150}.
\]
The Navier--Stokes source had \(14\,908\) newline characters and SHA--256
\[
 \texttt{654135F3D7C2203CB486A38C910DFBDB51F2E5EFBEC9ADB0E38CC19A797E739B}.
\]
\end{audit}

The Yang--Mills note derives its exact one-link conditional law from the
complete incident interaction star.  It proves that response matching and
one-site marginals do not control mixed conditional influence, and that
even perfect mixing does not supply the missing static Hessian sign.  Its
positive procedural result is exact: all response-matched writers sharing
two variables must be summed pointwise before oscillation.

The Navier--Stokes note derives a mean-vorticity identity and closes two
extremal estimates through a nonnegative \(2\times2\) matrix loop.  A
spectral radius below one would force both extremal observables to vanish;
the class forbids that, yielding a constants-only lower constraint.  It
does not prove global regularity.

\begin{decision}[Effect of the V45 audit]
\label{dec:newS7V45-audit}
At a codimension-two edge, separate estimates for the incident face
collars can miss a joint localized mode, just as separate Yang--Mills
writers can miss their pointwise cancellation or reinforcement.  The
complete cyclic sector fan must be assembled before any norm or
determinant is taken.  The NS matrix-loop method remains available as a
sufficient coupled-sector estimate, but it is unnecessary for the first
construction below: the V38 universal bulk symbol gives a common gapped
hub and makes the entire sector loop contractible inside a known gapped
family.
\end{decision}

\subsection{The frozen edge sector fan}
\label{sec:newS7V45-fan}

Fix an open original two-simplex \(e\).  Fourier transform the two lattice
directions tangent to \(e\), with momentum \(k_e\in\mathbb T^2\), and retain
the two transverse lattice coordinates
\[
 p=(p_1,p_2)\in\mathbb Z^2.
\]
The original four-simplices incident to \(e\) occur in a cyclic order
\(\Sigma_1,\ldots,\Sigma_m\).  Their regular bulk symbols are
\(H_1,\ldots,H_m\); the ray between sectors \(i\) and \(i+1\) is the
original face \(F_i\), with indices modulo \(m\).

\begin{definition}[Gapped path tree]
\label{def:newS7V45-tree}
Let \(H_{\rm u}\) be the universal symbol of V38.  For each sector choose
the V38 UGBP path
\begin{equation}
 P_i:[0,1]\longrightarrow\mathfrak G,\qquad
 P_i(0)=H_{\rm u},\quad P_i(1)=H_i,                   \label{eq:newS7V45-Pi}
\end{equation}
where \(\mathfrak G\) is the space of self-adjoint finite-range symbols
with gap at least \(\delta_{\rm path}\).  Pad all hopping ranges to one
common finite range.  The abstract union of the intervals \(P_i\), joined
only at their \(H_{\rm u}\) endpoints, is the finite metric tree
\(\mathcal T_e\).
\end{definition}

The coefficient map from the abstract tree to symbols may identify
additional points; that does not affect the following contraction.

\begin{lemma}[The cyclic face path is null through gapped symbols]
\label{lem:newS7V45-null-loop}
Choose the face path from sector \(i\) to sector \(i+1\) to be
\begin{equation}
 P_i^{-1}*P_{i+1}.                                    \label{eq:newS7V45-face-path}
\end{equation}
The concatenation of these \(m\) face paths is a loop in
\(\mathcal T_e\) and is null-homotopic there.  Every symbol in the
null-homotopy has gap at least \(\delta_{\rm path}\).
\end{lemma}

\begin{proof}
The concatenated word is
\[
 P_1^{-1}P_2P_2^{-1}P_3\cdots P_m^{-1}P_1.
\]
Successive path/backtrack pairs contract, or equivalently every loop in a
tree contracts along the unique geodesic to its root \(H_{\rm u}\).
The coefficient map sends the tree into the union of the V38 UGBP paths,
where the declared uniform gap holds.
\end{proof}

\subsection{A Lipschitz filling of the transverse fan}
\label{sec:newS7V45-filling}

Scale the transverse plane by \(L\).  Outside the unit disk, partition it
into the \(m\) sector interiors and strips of fixed scaled width about the
separating rays.  Map each sector interior to its endpoint \(H_i\) and
map a cross-section of the \(i\)-th strip along
\eqref{eq:newS7V45-face-path}.  On the boundary of the unit disk this gives
the cyclic tree loop of \Cref{lem:newS7V45-null-loop}.

\begin{lemma}[Gapped fan filling]
\label{lem:newS7V45-filling}
There is a piecewise-Lipschitz map
\begin{equation}
 \Phi_e:\mathbb R^2\longrightarrow\mathcal T_e         \label{eq:newS7V45-Phi}
\end{equation}
which has the prescribed sector and face-strip values outside the unit
disk, maps the central disk into \(\mathcal T_e\), and has a Lipschitz
constant depending only on the finite fan geometry and V38 path constants.
\end{lemma}

\begin{proof}
Triangulate the boundary circle at every endpoint and midpoint of a
face-strip interval.  Map its edges piecewise along the corresponding
geodesics of the finite tree.  Cone this triangulation to the centre of the
disk and send the centre to the root \(H_{\rm u}\).  On each radial
triangle, map a point first to its boundary value and then a proportional
distance along the unique tree geodesic from the root.  A finite tree has
piecewise-linear geodesics, so after a finite subdivision the map is
Lipschitz on every triangle with a uniform maximum constant.  Attach the
already prescribed constant-sector and strip maps outside the disk.
\end{proof}

Let \(h_d(\tau;E,k_e)\), \(|d|\le J\), denote the common-range transverse
hopping coefficients of the symbol represented by
\(\tau\in\mathcal T_e\), including its finite orbit fibre.  They obey
\(h_{-d}(\tau)=h_d(\tau)^*\).

\begin{definition}[Two-normal midpoint edge operator]
\label{def:newS7V45-edge-operator}
Set \(\tau_x=\Phi_e(x/L)\) for half-lattice points and define
\begin{equation}
 (H_{e,L}^{\rm ad}u)_p
 :=\sum_{|d|\le J}
 h_d(\tau_{p+d/2};E,k_e)u_{p+d}.                      \label{eq:newS7V45-edge-operator}
\end{equation}
\end{definition}

\begin{lemma}[Self-adjointness and exact sector tails]
\label{lem:newS7V45-edge-algebra}
\(H_{e,L}^{\rm ad}\) is self-adjoint.  Beyond an \(O(L)\) transverse
neighbourhood of the edge and the \(O(L)\)-wide face strips, it equals the
appropriate regular bulk operator; inside each face strip far from the
edge it equals the V37 face interpolation, up to the declared
piecewise-Lipschitz reparameterization.
\end{lemma}

\begin{proof}
The coefficient on the unoriented pair \(\{p,p+d\}\) is evaluated at its
midpoint.  Reversing the pair replaces \(d\) by \(-d\) and takes the
adjoint, proving self-adjointness.  The tail assertions follow from the
definition of \(\Phi_e\) outside the central disk.
\end{proof}

\subsection{Two-normal adiabatic localization}
\label{sec:newS7V45-adiabatic}

Choose a quadratic partition of unity
\(\{\chi_\alpha\}\) on \(\mathbb Z^2\), supported on boxes of side
\(O(\ell)\), such that
\begin{equation}
 \sum_\alpha\chi_\alpha(p)^2=1,\qquad
 \sum_\alpha|\chi_\alpha(p)-\chi_\alpha(q)|^2
 \le {C|p-q|^2\over\ell^2}.                           \label{eq:newS7V45-partition}
\end{equation}

\begin{lemma}[Two-normal commutator and freezing errors]
\label{lem:newS7V45-errors}
There are constants independent of \(L,\ell,E,k_e\) such that
\begin{align}
 \left(\sum_\alpha
 \|[H_{e,L}^{\rm ad},\chi_\alpha]u\|^2\right)^{1/2}
 &\le {C_{\rm com}\over\ell}\|u\|,                    \label{eq:newS7V45-commutator}\\
 \left(\sum_\alpha
 \|(H_{e,L}^{\rm ad}-H_{\tau_\alpha})
       \chi_\alpha u\|^2\right)^{1/2}
 &\le C_{\rm fr}{\ell\over L}\|u\|,                  \label{eq:newS7V45-freezing}
\end{align}
where \(H_{\tau_\alpha}\) is the translation-invariant symbol frozen at
one point in the support of \(\chi_\alpha\).
\end{lemma}

\begin{proof}
The commutator calculation is the V37 finite-range proof with
\(p,d\in\mathbb Z^2\); the partition estimate supplies \(1/\ell\).
The map \(\Phi_e(p/L)\) varies by at most \(C\ell/L\) on a box, and the
tree coefficient map is Lipschitz on its finite path segments.  The block
Schur test and bounded partition overlap give the second estimate.
\end{proof}

\begin{theorem}[Codimension-two adiabatic gap]
\label{thm:newS7V45-edge-gap}
The two-normal edge operator satisfies
\begin{equation}
 \|H_{e,L}^{\rm ad}u\|
 \ge
 \left(
 \delta_{\rm path}-{C_{\rm com}\over\ell}
 -C_{\rm fr}{\ell\over L}
 \right)\|u\|.                                        \label{eq:newS7V45-edge-gap}
\end{equation}
Taking \(\ell\asymp L^{1/2}\) gives
\begin{equation}
 \|H_{e,L}^{\rm ad}u\|
 \ge(\delta_{\rm path}-C_eL^{-1/2})\|u\|.             \label{eq:newS7V45-edge-gap-root}
\end{equation}
\end{theorem}

\begin{proof}
On every localization box, the frozen symbol belongs to
\(\mathcal T_e\) and has gap \(\delta_{\rm path}\).  Apply that gap to
\(\chi_\alpha u\), sum in \(\ell^2(\alpha)\), use Minkowski's inequality,
and insert \Cref{lem:newS7V45-errors}, exactly as in the V37 proof.
\end{proof}

\subsection{The joint moment defect}
\label{sec:newS7V45-moment}

The construction assembled the complete cyclic fan before estimating it.
It has not yet imposed the physical zeroth moment.

\begin{lemma}[Size and support of the edge moment defect]
\label{lem:newS7V45-edge-moment}
Let \(B_{e,L}=H_{e,L}^{\rm ad}-H_e^{\rm sh}\), after using the completed
open-face stencil in the remote face strips.  Its nonbulk part is supported
on \(O(L^2)\) transverse sites.  At zero edge momentum and on a constant
transverse spinor, its row residual obeys
\begin{equation}
 \left\|\sum_q(B_{e,L})_{pq}\right\|\le {C\over L}.    \label{eq:newS7V45-edge-row}
\end{equation}
\end{lemma}

\begin{proof}
The central region has radius \(O(L)\), hence \(O(L^2)\) sites.  At a
frozen value of \(\Phi_e\), the nonmass translation-invariant symbol has
zero row sum.  Replacing the midpoint values in one row by the central
value changes finitely many coefficients by \(O(L^{-1})\).  The already
completed remote face strips have zero frozen row defect by V44.
\end{proof}

Pairwise application of the V44 face corrections does not cancel
\eqref{eq:newS7V45-edge-row} in the central region, because every row sees
the complete sector fan.  A joint two-normal completion is required.

\begin{assessment}[The edge gap is acquired before its continuum moment]
\label{ass:newS7V45-status}
V45 constructs a joint codimension-two adiabatic operator and proves its
uniform frozen edge gap.  The universal hub removes the possible cyclic
bulk-path obstruction.  Its \(O(L^{-1})\) zeroth moment on \(O(L^2)\)
sites remains.  As in V43, that moment must be removed before the
shrinking edge tube can be declared continuum-safe.
\end{assessment}

\begin{programme}[V46: joint two-normal conservative completion]
\label{prog:newS7V45-V46}
V46 will extend the V44 normalized-constant-subspace construction from
\(O(L)\) normal layers to the \(O(L^2)\) transverse edge disk.  It will
prove that the complete fan residual still yields an \(O(L^{-1})\)
self-adjoint correction, establish a bounded transverse first moment, and
lift it covariantly along the edge.  This is a genuinely joint completion,
not a sum of facewise corrections.
\end{programme}

\begin{firewall}[Claim boundary after seventh-series V45]
\label{fire:newS7V45-boundary}
V45 performs the compulsory YM/NS audit, constructs the complete cyclic
edge sector fan, proves its path loop contracts through the universal
gapped hub, builds a Lipschitz two-normal filling, and proves the
codimension-two adiabatic gap with \(O(L^{-1/2})\) loss.

V45 does not cancel the edge zeroth moment, prove continuum safety of the
edge replacement, handle vertex strata, prove Gaussian locality for the
growing-range completions, full SFGC, the global overlap gap, strong local
heat consistency, the common Einstein stress, pure gravitational owners,
rough metric or AOP control, measure concentration, the determinant-line
phase, interacting cutoff removal, or the Standard Model--Einstein
continuum.
\end{firewall}

\begin{theorem}[Seventh-series V45 result]
\label{thm:newS7V45-result}
The face paths around every original edge form a null loop in the V38
universal gapped path tree.  A Lipschitz filling of that loop gives a
self-adjoint two-normal finite-range tangent operator with gap
\(\delta_{\rm path}-O(L^{-1/2})\).  Its sole new continuum obstruction is
the joint \(O(L^{-1})\) zeroth moment on the \(O(L^2)\)-site central disk.
\end{theorem}

\begin{proof}
Use \Cref{lem:newS7V45-null-loop,
lem:newS7V45-filling,
lem:newS7V45-edge-algebra,
lem:newS7V45-errors,
thm:newS7V45-edge-gap,
lem:newS7V45-edge-moment}.
\end{proof}

\paragraph{Parent-version record.}
V45 was formed from
\texttt{Renewal\_SM\_Einstein\_Continuum\_Research\_Notes\_V44.tex}.
The V44 parent had \(17\,681\) logical source lines, \(711\,656\) bytes,
and SHA--256
\[
 \texttt{92C8B334E6A04B879B4B6DB25AF983C095F68453C0A6007D2A21FF589F7D36D3}.
\]
The next compulsory companion audit is V50.

% =============================================================================
% END APPENDED SEVENTH-SERIES V45 MATERIAL
% =============================================================================

% =============================================================================
% BEGIN APPENDED SEVENTH-SERIES V46 MATERIAL
% =============================================================================

\section{Seventh-series V46: conservative completion of the codimension-two fan}
\label{sec:v7v46}
\label{sec:newS7V46}

\begin{correction}[The V45 codimension-two stratum is a triangle]
\label{corr:newS7V46-terminology}
V45 denotes its stratum by \(e\) but defines it as an open original
two-simplex.  In a four-dimensional simplicial complex this is a triangle,
also called a codimension-two hinge, not a one-simplex edge.  All V45
formulas use two tangential and two transverse coordinates and therefore
apply to the triangle/hinge.  Every use of the word edge in the V45
codimension-two discussion is to be read as triangle or hinge.
The genuine one-simplex edge is codimension three and remains downstream.
\end{correction}

\subsection{The joint transverse constant subspace}
\label{sec:newS7V46-subspace}

Fix a frozen open triangle \(T\), its tangential momentum
\(k_T\in\mathbb T^2\), and the complete V45 sector-fan operator
\(H_{T,L}^{\rm ad}\).  Let \(H_T^{\rm ref}\) denote the physical stencil
which equals the sharp bulk in sector interiors and the V44 completed face
stencil in remote face strips.  The difference
\begin{equation}
 B_{T,L}(k_T):=H_{T,L}^{\rm ad}(k_T)-H_T^{\rm ref}(k_T)              \label{eq:newS7V46-B}
\end{equation}
is self-adjoint and has all nonzero transverse block rows and columns in a
disk-like set \(S_{T,L}\subset\mathbb Z^2\) satisfying
\begin{equation}
 c_1L^2\le N_{T,L}:=|S_{T,L}|\le c_2L^2.              \label{eq:newS7V46-N}
\end{equation}

Define the normalized constant-transverse isometry
\begin{equation}
 Q_{T,L}:\mathbb C^q\longrightarrow
 \ell^2(S_{T,L};\mathbb C^q),\qquad
 (Q_{T,L}v)_p=N_{T,L}^{-1/2}v.                        \label{eq:newS7V46-Q}
\end{equation}
At zero tangential momentum put
\begin{equation}
 Y_{T,L}:=-B_{T,L}(0)Q_{T,L}.                         \label{eq:newS7V46-Y}
\end{equation}

\begin{lemma}[The two-normal residual still has size \(L^{-1}\)]
\label{lem:newS7V46-Y-bound}
There is a uniform constant \(C_T\) such that
\begin{equation}
 \|Y_{T,L}\|\le {C_T\over L}.                         \label{eq:newS7V46-Y-bound}
\end{equation}
The estimate is uniform over the finite hinge fans, compact coframe class,
and two metric variations.
\end{lemma}

\begin{proof}
Let
\[
 z_p:=\sum_qB_{T,L}(0)_{pq}.
\]
V45 gives \(\|z_p\|\le C/L\) on at most \(N_{T,L}=O(L^2)\) rows.  Thus,
for \(\|v\|=1\),
\[
 \|B_{T,L}(0)Q_{T,L}v\|^2
 ={1\over N_{T,L}}\sum_{p\in S_{T,L}}\|z_pv\|^2
 \le {1\over N_{T,L}}N_{T,L}{C^2\over L^2}.
\]
Finiteness of the fan types and smooth compact dependence of the V38 paths
make the constant uniform.  Differentiating the finite coefficient paths
and midpoint rows preserves the same estimate.
\end{proof}

\subsection{Joint completion and gap}
\label{sec:newS7V46-completion}

\begin{definition}[Hinge conservative correction]
\label{def:newS7V46-R}
Define
\begin{equation}
 R_{T,L}:=
 Y_{T,L}Q_{T,L}^*+Q_{T,L}Y_{T,L}^*
 -Q_{T,L}(Q_{T,L}^*Y_{T,L})Q_{T,L}^*,                \label{eq:newS7V46-R}
\end{equation}
extend it by zero outside \(S_{T,L}\), and use the same transverse matrix
at every translate along the hinge.
\end{definition}

\begin{theorem}[Joint self-adjoint moment completion]
\label{thm:newS7V46-completion}
The correction \(R_{T,L}\) is self-adjoint, has rank at most \(2q\) in
each \(k_T\)-fibre, and obeys
\begin{align}
 (B_{T,L}(0)+R_{T,L})Q_{T,L}&=0,                      \label{eq:newS7V46-zero}\\
 \|R_{T,L}\|&\le {3C_T\over L}.                       \label{eq:newS7V46-norm}
\end{align}
\end{theorem}

\begin{proof}
As in V44,
\[
 Q_{T,L}^*Y_{T,L}
 =-Q_{T,L}^*B_{T,L}(0)Q_{T,L}
\]
is self-adjoint.  Therefore \eqref{eq:newS7V46-R} is self-adjoint and
\[
 R_{T,L}Q_{T,L}
 =Y_{T,L}+Q_{T,L}
   \{Y_{T,L}^*Q_{T,L}-Q_{T,L}^*Y_{T,L}\}
 =Y_{T,L}.
\]
This proves \eqref{eq:newS7V46-zero}.  The range lies in
\(\operatorname{Ran}Y_{T,L}+\operatorname{Ran}Q_{T,L}\), and the same
three-term norm estimate as V44 gives
\(\|R_{T,L}\|\le3\|Y_{T,L}\|\).  Apply
\Cref{lem:newS7V46-Y-bound}.
\end{proof}

Define
\begin{equation}
 H_{T,L}^{\rm con}:=H_{T,L}^{\rm ad}+R_{T,L}.          \label{eq:newS7V46-Hcon}
\end{equation}

\begin{theorem}[Uniform codimension-two gap]
\label{thm:newS7V46-gap}
For every hinge fan and all \(k_T\),
\begin{equation}
 \|H_{T,L}^{\rm con}(k_T)u\|
 \ge
 \left(
 \delta_{\rm path}-C_eL^{-1/2}-{3C_T\over L}
 \right)\|u\|.                                        \label{eq:newS7V46-gap}
\end{equation}
For sufficiently large \(L\), the gap is at least
\(\delta_{\rm path}/2\).  Moreover,
\begin{equation}
 \|\operatorname{sgn}(H_{T,L}^{\rm con})
 -\operatorname{sgn}(H_{T,L}^{\rm ad})\|
 \le {C\over L}.                                      \label{eq:newS7V46-sign}
\end{equation}
\end{theorem}

\begin{proof}
Combine the V45 gap
\eqref{eq:newS7V45-edge-gap-root} with
\eqref{eq:newS7V46-norm} and the reverse triangle inequality.  Once both
gaps exceed a fixed fraction of \(\delta_{\rm path}\), the sign-function
resolvent estimate used in V44 gives \eqref{eq:newS7V46-sign}.
\end{proof}

\subsection{Transverse first moment}
\label{sec:newS7V46-first-moment}

The block sizes satisfy
\begin{equation}
 \|(Q_{T,L})_p\|\le CL^{-1},\qquad
 \|(Y_{T,L})_p\|\le CL^{-2}.                          \label{eq:newS7V46-block-scales}
\end{equation}

\begin{lemma}[Two-normal weighted row estimate]
\label{lem:newS7V46-weighted-row}
There is \(C_R\), independent of \(L\), such that
\begin{align}
 \sup_{p\in S_{T,L}}\sum_{q\in S_{T,L}}
 \|(R_{T,L})_{pq}\|&\le {C_R\over L},                 \label{eq:newS7V46-row}\\
 \sup_{p\in S_{T,L}}\sum_{q\in S_{T,L}}
 |p-q|\|(R_{T,L})_{pq}\|&\le C_R.                    \label{eq:newS7V46-first}
\end{align}
\end{lemma}

\begin{proof}
Every block of the first two terms of \eqref{eq:newS7V46-R} is
\(O(L^{-3})\).  Since
\(\|Q_{T,L}^*Y_{T,L}\|=O(L^{-1})\), the third term has the same block
size.  A row contains \(O(L^2)\) sites, giving \(O(L^{-1})\).
The sum of their distances from a fixed site is \(O(L^3)\), giving the
second bound.
\end{proof}

\begin{corollary}[Bounded physical moment and shrinking range]
\label{cor:newS7V46-physical-moment}
After restoring lattice spacing \(a\), the physical first transverse
moment of \(R_{T,L}\) is uniformly bounded and its physical range is
\(O(aL)\).  Under \eqref{eq:newS7V43-scales}, that range tends to zero.
\end{corollary}

\begin{proof}
In the first moment, the factor \(a\) from displacement cancels the
physical \(a^{-1}\) sign-kernel scale.  Apply
\eqref{eq:newS7V46-first}.  The support diameter is \(O(L)\) lattice
units.
\end{proof}

\subsection{Covariant lift along a curved hinge}
\label{sec:newS7V46-curved}

For each tangential hinge vertex, choose the central fibre and identify the
fibres in its transverse disk by V33 short-geodesic spin--gauge transport.
Use those transports in \(Q_{T,L}\), form \eqref{eq:newS7V46-Y} and
\eqref{eq:newS7V46-R}, and transport every output block back to its endpoint.

\begin{lemma}[Covariance and curved residual]
\label{lem:newS7V46-curved}
The lifted \(R_{T,L}\) is spin--gauge covariant and self-adjoint.  Smooth
variation along the hinge leaves a zeroth moment bounded by
\begin{equation}
 \|Z_{T,a}^{\rm con}(x)\|\le Ca,                      \label{eq:newS7V46-curved-Z}
\end{equation}
and a uniformly bounded first moment, through two metric variations.
\end{lemma}

\begin{proof}
The covariance proof is identical to V44: \(Q_{T,L}\) intertwines the
central and endpoint frames, and every factor in the completion has the
corresponding endpoint conjugation.  Self-adjointness is invariant under
unitary transport.  The frozen transverse moment vanishes by
\eqref{eq:newS7V46-zero}.  Neighbouring tangential coefficients and
transports differ by \(O(a)\), leaving
\eqref{eq:newS7V46-curved-Z}.  The finite-range fan part and
\Cref{lem:newS7V46-weighted-row} bound the first moment.  Smooth compact
dependence gives the differentiated claims.
\end{proof}

The physical tube around a two-dimensional hinge has four-volume
\begin{equation}
 \operatorname{Vol}(\mathcal C_a(T))\le C_T(aL)^2.     \label{eq:newS7V46-tube-volume}
\end{equation}

\begin{theorem}[Continuum safety at codimension two]
\label{thm:newS7V46-continuum}
Let \(L(a)\to\infty\) and \(aL(a)\to0\), and choose the localization scale
as in \eqref{eq:newS7V43-three-scales}.  For every fixed smooth spinor
\(\psi\),
\begin{equation}
 \left\|a^{-1}
 (H_{T,L}^{\rm con}-H_T^{\rm ref})\psi
 \right\|_{L^2(\mathcal C_a(T))}
 \le C_\psi\,aL(a)\longrightarrow0.                  \label{eq:newS7V46-continuum}
\end{equation}
The estimate holds through two metric contacts.
\end{theorem}

\begin{proof}
Use the covariant Taylor expansion from V43.  The curved zeroth moment is
\(O(a)\), so after multiplication by \(a^{-1}\) it is bounded.  The first
moment is bounded by \Cref{lem:newS7V46-weighted-row}, and the second
remainder is smaller.  The pointwise physical difference on a smooth
spinor is therefore \(O(1)\).  Multiply by the square root of the tube
volume \eqref{eq:newS7V46-tube-volume}.  The metric-contact proof is the
same.
\end{proof}

\subsection{The codimension-two acquisition}
\label{sec:newS7V46-acquisition}

\begin{theorem}[Hinge collar acquisition]
\label{thm:newS7V46-acquisition}
The V45 sector fan plus the joint completion
\eqref{eq:newS7V46-R} defines a codimension-two collar kernel which:
\begin{enumerate}
\item matches the completed face and bulk stencils outside its transverse
central disk;
\item is self-adjoint and spin--gauge covariant;
\item has a uniform frozen and slowly tangentially varying gap;
\item has bounded first moment and curved zeroth moment \(O(a)\) through
two metric contacts; and
\item has the same smooth-core continuum action as the previously
constructed bulk and face kernels.
\end{enumerate}
\end{theorem}

\begin{proof}
Tail matching is the V45 filling construction.  Self-adjointness, gap, and
sign stability are \Cref{thm:newS7V46-completion,thm:newS7V46-gap}.
Covariance and moments are \Cref{lem:newS7V46-curved}; continuum safety is
\Cref{thm:newS7V46-continuum}.  Slow tangential variation is controlled by
the V43/V45 localization estimate.
\end{proof}

\begin{assessment}[Codimension two closed for the modified discretization]
\label{ass:newS7V46-status}
The modified conservative adiabatic discretization now has acquired bulk,
open-face, and open-triangle tangent gaps with compatible smooth continuum
moments.  The next stratum is a genuine original one-simplex edge, with
one tangential and three transverse lattice directions.  Gaussian locality
for the growing-range low-rank completions remains unproved.
\end{assessment}

\begin{programme}[V47: codimension-three sector complex]
\label{prog:newS7V46-V47}
V47 will replace the planar fan by the three-dimensional link of an
original edge.  It will prove that the already chosen universal-hub paths
and triangle fillings extend over that link, or isolate the first
\(\pi_2\)-type obstruction in the gapped symbol space.  On a successful
extension, the \(d=3\) adiabatic localization and normalized-constant
completion will be derived before the vertex stratum is attempted.
\end{programme}

\begin{firewall}[Claim boundary after seventh-series V46]
\label{fire:newS7V46-boundary}
V46 corrects the V45 stratum terminology, constructs the joint
two-normal rank-\(2q\) completion, proves its \(O(L^{-1})\) norm, exact
frozen moment, uniform hinge gap, sign stability, bounded transverse first
moment, covariant curved lift, and continuum safety through two metric
contacts.

V46 does not handle genuine one-simplex edges or original vertices, prove
Gaussian locality for the growing-range completions, full SFGC, the global
overlap gap, strong local heat consistency, the common Einstein stress,
pure gravitational owners, rough metric or AOP control, measure
concentration, the determinant-line phase, interacting cutoff removal, or
the Standard Model--Einstein continuum.
\end{firewall}

\begin{theorem}[Seventh-series V46 result]
\label{thm:newS7V46-result}
The complete codimension-two sector fan admits the explicit joint
self-adjoint completion \eqref{eq:newS7V46-R}.  It preserves a gap
\(\delta_{\rm path}-O(L^{-1/2})\), cancels the frozen zeroth moment,
has bounded physical first moment and shrinking range, and is
continuum-equivalent on smooth cores through two metric contacts.
\end{theorem}

\begin{proof}
Use \Cref{lem:newS7V46-Y-bound,
thm:newS7V46-completion,
thm:newS7V46-gap,
lem:newS7V46-weighted-row,
lem:newS7V46-curved,
thm:newS7V46-continuum,
thm:newS7V46-acquisition}.
\end{proof}

\paragraph{Parent-version record.}
V46 was formed from
\texttt{Renewal\_SM\_Einstein\_Continuum\_Research\_Notes\_V45.tex}.
The V45 parent had \(18\,017\) logical source lines, \(725\,650\) bytes,
and SHA--256
\[
 \texttt{5552283B0D2AB6F97F486C89143C35A84A2EF50B4187ACE0D775E83B5283C121}.
\]
The next compulsory companion audit is V50.

% =============================================================================
% END APPENDED SEVENTH-SERIES V46 MATERIAL
% =============================================================================

% =============================================================================
% BEGIN APPENDED SEVENTH-SERIES V47 MATERIAL
% =============================================================================

\section{Seventh-series V47: the codimension-three edge completion}
\label{sec:v7v47}
\label{sec:newS7V47}

\subsection{The link of an original edge}
\label{sec:newS7V47-link}

Let \(E^\circ\) be an open original one-simplex.  Fourier transform its
one tangential lattice direction, with momentum \(k_E\in\mathbb T\), and
retain three transverse lattice coordinates \(p\in\mathbb Z^3\).

\begin{lemma}[Combinatorial transverse link]
\label{lem:newS7V47-link-sphere}
Because the original triangulation is a closed combinatorial
four-manifold, the link \(\operatorname{Lk}(E)\) is a finite triangulated
two-sphere.  Its two-cells correspond to original four-simplices incident
to \(E\), its one-cells correspond to original three-faces containing
\(E\), and its vertices correspond to original two-simplices containing
\(E\).
\end{lemma}

\begin{proof}
The link of an interior \(j\)-simplex in a combinatorial \(d\)-manifold is
a sphere of dimension \(d-j-1\).  Here \(d=4\) and \(j=1\), giving
\(S^2\).  Removing \(E\) from an incident simplex lowers its dimension by
two and gives the stated incidence correspondence.
\end{proof}

Let \(H_1,\ldots,H_m\) be the regular bulk symbols associated with the
two-cells of this link.  Join each to the V38 universal symbol \(H_{\rm u}\)
by its fixed UGBP path \(P_i\).

\begin{definition}[Edge universal path tree]
\label{def:newS7V47-tree}
\(\mathcal T_E\) is the finite rooted metric tree obtained by taking one
abstract copy of every \(P_i\) and identifying their universal endpoints
with the root \(H_{\rm u}\).
\end{definition}

The leading face interpolation on a link edge is
\(P_i^{-1}*P_j\).  The leading hinge filling at a link vertex is the V45
tree-valued disk filling for the cyclic fan of the bulk cells meeting
there.  Thus the already constructed leading coefficient map on the
stratified link takes values in \(\mathcal T_E\).

\begin{lemma}[No \(\pi_2\) obstruction]
\label{lem:newS7V47-no-pi2}
The complete leading coefficient map
\begin{equation}
 f_E:\operatorname{Lk}(E)\simeq S^2\longrightarrow\mathcal T_E          \label{eq:newS7V47-link-map}
\end{equation}
extends continuously, and after finite subdivision Lipschitz-continuously,
to the three-ball.
\end{lemma}

\begin{proof}
A tree is contractible by motion along the unique geodesic to its root.
Compose this contraction with \(f_E\) to obtain a cone extension over
\(B^3\).  The link map is piecewise Lipschitz on a finite triangulation and
the tree has finitely many path edges.  Subdivide the cone whenever an
image crosses a tree vertex; on each resulting simplex the geodesic
parameter is piecewise linear.  The finite maximum of the simplex
Lipschitz constants is finite.
\end{proof}

\subsection{A three-normal coefficient field}
\label{sec:newS7V47-field}

In the scaled transverse space \(\mathbb R^3\), use conical regions over
the cells of \(\operatorname{Lk}(E)\).  Away from the origin, prescribe:
\begin{enumerate}
\item the constant bulk endpoint in the interior of each three-dimensional
sector;
\item the V37 path in an \(O(1)\)-wide scaled sheet about each sector
interface; and
\item the V45 disk filling in an \(O(1)\)-radius scaled tube about each
hinge ray.
\end{enumerate}
On the unit central ball use the extension in
\Cref{lem:newS7V47-no-pi2}.

\begin{lemma}[Lipschitz edge filling]
\label{lem:newS7V47-filling}
These prescriptions can be joined into a piecewise-Lipschitz map
\begin{equation}
 \Phi_E:\mathbb R^3\longrightarrow\mathcal T_E         \label{eq:newS7V47-Phi}
\end{equation}
whose coefficient image has uniform gap
\(\delta_{\rm path}\), with a Lipschitz constant uniform over the finite
original edge types.
\end{lemma}

\begin{proof}
Choose nested conical sheets and tubes whose boundary cell structure is a
finite subdivision of the link triangulation.  The lower-dimensional
prescriptions agree on common boundaries because they were all formed
from the same paths \(P_i\) and root \(H_{\rm u}\).  Apply the cone
extension on the central ball and a finite piecewise-linear interpolation
of tree geodesic parameters in the transition annulus.  Finiteness of the
base triangulation and compactness of the coefficient paths give uniform
constants.  Every image point lies on a V38 gapped path.
\end{proof}

After padding to one common transverse hopping range \(J\), let
\(h_d(\tau;E,k_E)\), \(d\in\mathbb Z^3\), denote the symbol blocks.  Put
\(\tau_x=\Phi_E(x/L)\) at half-lattice points and define
\begin{equation}
 (H_{E,L}^{\rm lead}u)_p
 :=\sum_{|d|\le J}h_d(\tau_{p+d/2};E,k_E)u_{p+d}.      \label{eq:newS7V47-leading}
\end{equation}
Midpoint evaluation makes this operator self-adjoint.

\subsection{Matching the completed lower strata}
\label{sec:newS7V47-boundary}

The physical face and hinge collars include the V44 and V46 low-rank
completions, while \eqref{eq:newS7V47-leading} contains only their
tree-valued leading fields.  Let \(S_{\partial,L}\) be the sum of those
completed lower-stratum corrections, multiplied on both sides by
self-adjoint cutoffs which equal one outside the central edge ball and
zero in a smaller central ball.

\begin{lemma}[Boundary-completion debit]
\label{lem:newS7V47-boundary-debit}
The cutoffs can be chosen on scale \(L\) so that
\begin{align}
 S_{\partial,L}^*&=S_{\partial,L},                    \label{eq:newS7V47-boundary-self}\\
 \|S_{\partial,L}\|&\le {C_\partial\over L},           \label{eq:newS7V47-boundary-norm}\\
 \left\|\sum_q(S_{\partial,L})_{pq}\right\|
 &\le {C_\partial\over L}.                            \label{eq:newS7V47-boundary-row}
\end{align}
Its transverse weighted first row moment is uniformly bounded.
\end{lemma}

\begin{proof}
There are finitely many incident faces and hinges.  Each V44 or V46
completion has norm \(O(L^{-1})\), so sandwiching it by a contraction and
summing finitely many terms gives \eqref{eq:newS7V47-boundary-norm} and
self-adjointness.  Before cutoff, its frozen row moment is zero and its
weighted first row moment is bounded.  A cutoff changes by \(O(L^{-1})\)
over every hopping displacement.  Commuting the cutoff through the
completion therefore costs its weighted first moment times \(L^{-1}\),
which proves \eqref{eq:newS7V47-boundary-row}.  The same calculation
preserves the weighted moment bound.
\end{proof}

Define the boundary-compatible adiabatic operator
\begin{equation}
 H_{E,L}^{\rm ad}:=H_{E,L}^{\rm lead}+S_{\partial,L}. \label{eq:newS7V47-Had}
\end{equation}
It agrees with the completed bulk/face/hinge hierarchy outside the central
edge ball.

\subsection{Three-normal adiabatic gap}
\label{sec:newS7V47-gap}

\begin{theorem}[Uncompleted codimension-three gap]
\label{thm:newS7V47-adiabatic-gap}
Uniformly in the finite edge types, compact geometric parameters, and
\(k_E\),
\begin{equation}
 \|H_{E,L}^{\rm ad}u\|
 \ge
 \left(
 \delta_{\rm path}-C L^{-1/2}-{C_\partial\over L}
 \right)\|u\|.                                        \label{eq:newS7V47-adiabatic-gap}
\end{equation}
\end{theorem}

\begin{proof}
Use a quadratic partition of unity on \(\mathbb Z^3\) with box scale
\(\ell\).  The finite-range commutator bound is \(C/\ell\); freezing
\(\Phi_E(p/L)\) costs \(C\ell/L\).  Every frozen leading symbol has gap
\(\delta_{\rm path}\).  The V37 direct-sum localization argument gives
\[
 \|H_{E,L}^{\rm lead}u\|
 \ge(\delta_{\rm path}-C/\ell-C\ell/L)\|u\|.
\]
Take \(\ell\asymp L^{1/2}\), then use
\eqref{eq:newS7V47-boundary-norm} and the reverse triangle inequality.
\end{proof}

\subsection{Joint three-normal conservative completion}
\label{sec:newS7V47-completion}

Let \(H_E^{\rm ref}\) be the completed lower-stratum physical stencil and
put
\[
 B_{E,L}:=H_{E,L}^{\rm ad}-H_E^{\rm ref}.
\]
Choose a transverse support set \(S_{E,L}\) containing all nonzero rows and
columns of this difference.  It has
\begin{equation}
 c_1L^3\le N_{E,L}:=|S_{E,L}|\le c_2L^3.              \label{eq:newS7V47-N}
\end{equation}
At \(k_E=0\), define
\begin{align}
 (Q_{E,L}v)_p&:=N_{E,L}^{-1/2}v,                      \label{eq:newS7V47-Q}\\
 Y_{E,L}&:=-B_{E,L}(0)Q_{E,L},                        \label{eq:newS7V47-Y}\\
 R_{E,L}&:=Y_{E,L}Q_{E,L}^*+Q_{E,L}Y_{E,L}^*
 -Q_{E,L}(Q_{E,L}^*Y_{E,L})Q_{E,L}^*.                \label{eq:newS7V47-R}
\end{align}

\begin{lemma}[Dimension-independent completion norm]
\label{lem:newS7V47-completion-norm}
The leading midpoint residual and the boundary debit satisfy
\[
 \left\|\sum_q(B_{E,L}(0))_{pq}\right\|\le {C\over L}.
\]
Consequently
\begin{equation}
 \|Y_{E,L}\|\le {C\over L},\qquad
 \|R_{E,L}\|\le {3C\over L}.                          \label{eq:newS7V47-R-norm}
\end{equation}
\end{lemma}

\begin{proof}
Freezing the tree parameter in one row makes its nonmass row sum zero;
midpoint variation costs \(O(L^{-1})\).  Add
\eqref{eq:newS7V47-boundary-row}.  There are \(N_{E,L}=O(L^3)\) residual
rows.  Normalization by \(N_{E,L}^{-1/2}\) gives
\[
 \|B_{E,L}(0)Q_{E,L}\|^2
 \le {1\over N_{E,L}}N_{E,L}{C^2\over L^2}.
\]
The V44 algebra gives the three-term bound for \(R_{E,L}\).
\end{proof}

\begin{theorem}[Codimension-three moment and gap completion]
\label{thm:newS7V47-completion}
\(R_{E,L}\) is self-adjoint, has rank at most \(2q\) in each tangential
fibre, and
\begin{align}
 (B_{E,L}(0)+R_{E,L})Q_{E,L}&=0,                      \label{eq:newS7V47-zero}\\
 \|(H_{E,L}^{\rm ad}+R_{E,L})u\|
 &\ge(\delta_{\rm path}-CL^{-1/2})\|u\|.              \label{eq:newS7V47-final-gap}
\end{align}
\end{theorem}

\begin{proof}
The compression \(Q_{E,L}^*Y_{E,L}\) is self-adjoint because
\(B_{E,L}(0)\) is.  The completion proof of V44/V46 gives
self-adjointness, rank, and \eqref{eq:newS7V47-zero}.  Combine
\Cref{thm:newS7V47-adiabatic-gap} with
\eqref{eq:newS7V47-R-norm}.
\end{proof}

\subsection{First moment and continuum safety}
\label{sec:newS7V47-continuum}

The normalized blocks satisfy
\[
 \|(Q_{E,L})_p\|=O(L^{-3/2}),\qquad
 \|(Y_{E,L})_p\|=O(L^{-5/2}),
\]
so every block of \(R_{E,L}\) is \(O(L^{-4})\).

\begin{lemma}[Three-normal weighted row bound]
\label{lem:newS7V47-weighted-row}
\begin{align}
 \sup_p\sum_q\|(R_{E,L})_{pq}\|&=O(L^{-1}),            \label{eq:newS7V47-row}\\
 \sup_p\sum_q|p-q|\|(R_{E,L})_{pq}\|&=O(1).           \label{eq:newS7V47-first}
\end{align}
\end{lemma}

\begin{proof}
A row has \(O(L^3)\) blocks of size \(O(L^{-4})\), proving the first
bound.  The sum of transverse distances from a fixed site over a radius
\(O(L)\) three-ball is \(O(L^4)\), proving the second.
\end{proof}

Lift the construction along the curved edge using short-geodesic
spin--gauge transports from a central fibre.  Smooth variation along the
edge leaves a zeroth moment \(O(a)\) after the frozen cancellation, and
\Cref{lem:newS7V47-weighted-row} bounds the first moment through two metric
contacts.

\begin{theorem}[Continuum safety at a codimension-three edge]
\label{thm:newS7V47-continuum}
Let \(L(a)\to\infty\) and \(aL(a)\to0\).  On the physical tube
\(\mathcal C_a(E)\), whose volume is \(O((aL)^3)\), every fixed smooth
spinor satisfies
\begin{equation}
 \left\|a^{-1}
 \{H_{E,L}^{\rm ad}+R_{E,L}-H_E^{\rm ref}\}\psi
 \right\|_{L^2(\mathcal C_a(E))}
 \le C_\psi\{aL(a)\}^{3/2}\longrightarrow0.           \label{eq:newS7V47-continuum}
\end{equation}
The result holds through two metric variations, and the physical range of
the completion is \(O(aL)\).
\end{theorem}

\begin{proof}
The covariant Taylor expansion of V43 applies.  The physical zeroth and
first-moment contributions are uniformly bounded, and the remainder is
smaller.  Multiply the pointwise bound by the square root of the tube
volume.  Metric differentiation preserves the uniform coefficient and
moment estimates.  The support diameter is \(O(L)\) lattice units.
\end{proof}

\begin{assessment}[Codimension three acquired]
\label{ass:newS7V47-status}
The conservative adiabatic hierarchy now has compatible bulk, face,
triangle, and one-simplex-edge tangent gaps and smooth continuum moments.
No \(\pi_2\) obstruction occurs because the entire leading link map lands
in the universal path tree.  The final simplicial stratum is an original
vertex, whose transverse link is \(S^3\).
\end{assessment}

\begin{programme}[V48: the original-vertex four-ball]
\label{prog:newS7V47-V48}
V48 will extend the tree-valued \(S^3\) link data over a transverse
four-ball, prove the four-normal adiabatic gap, attach the finite lower
strata with an \(O(L^{-1})\) debit, and perform the normalized-constant
completion on \(O(L^4)\) sites.  It will then state the complete
stratified tangent gap for the modified discretization.
\end{programme}

\begin{firewall}[Claim boundary after seventh-series V47]
\label{fire:newS7V47-boundary}
V47 identifies the edge link as \(S^2\), proves its complete leading map
extends through the universal path tree, constructs the three-normal
operator, controls the completed lower-stratum boundary debit, and proves
the codimension-three gap, moment completion, covariance, bounded first
moment, and continuum safety.

V47 does not handle original vertices, prove one global stratified gap,
prove Gaussian locality for the growing-range completions, the global
overlap heat calculus, strong local heat consistency, the common Einstein
stress, pure gravitational owners, rough metric or AOP control, measure
concentration, the determinant-line phase, interacting cutoff removal, or
the Standard Model--Einstein continuum.
\end{firewall}

\begin{theorem}[Seventh-series V47 result]
\label{thm:newS7V47-result}
The completed lower-stratum data on the \(S^2\) link of every original
one-simplex edge admit a joint three-normal extension.  Its self-adjoint
sign kernel has gap \(\delta_{\rm path}-O(L^{-1/2})\), exact frozen zeroth
moment after a rank-\(2q\) correction, bounded physical first moment,
shrinking range, and the same smooth-core continuum action through two
metric contacts.
\end{theorem}

\begin{proof}
Use \Cref{lem:newS7V47-link-sphere,
lem:newS7V47-no-pi2,
lem:newS7V47-filling,
lem:newS7V47-boundary-debit,
thm:newS7V47-adiabatic-gap,
lem:newS7V47-completion-norm,
thm:newS7V47-completion,
lem:newS7V47-weighted-row,
thm:newS7V47-continuum}.
\end{proof}

\paragraph{Parent-version record.}
V47 was formed from
\texttt{Renewal\_SM\_Einstein\_Continuum\_Research\_Notes\_V46.tex}.
The V46 parent had \(18\,365\) logical source lines, \(738\,961\) bytes,
and SHA--256
\[
 \texttt{93DFC743BBC32A3301821412F85EEF89960BA7297FC2B92B755AD53176AD4A97}.
\]
The next compulsory companion audit is V50.

% =============================================================================
% END APPENDED SEVENTH-SERIES V47 MATERIAL
% =============================================================================

% =============================================================================
% BEGIN APPENDED SEVENTH-SERIES V48 MATERIAL
% =============================================================================

\section{Seventh-series V48: the original-vertex four-normal completion}
\label{sec:v7v48}
\label{sec:newS7V48}

\subsection{The final simplicial link}
\label{sec:newS7V48-link}

Let \(v\) be an original vertex of the closed four-dimensional
triangulation.  There is no tangential lattice direction; all four lattice
directions are transverse.

\begin{lemma}[The vertex link]
\label{lem:newS7V48-link}
The link \(\operatorname{Lk}(v)\) is a finite triangulated three-sphere.
Its strata encode, in descending cell dimension, the original
four-simplices, three-faces, two-simplex hinges, and one-simplex edges
incident to \(v\).
\end{lemma}

\begin{proof}
The link of an interior \(j\)-simplex in a combinatorial \(d\)-manifold is
\(S^{d-j-1}\).  Taking \(d=4,j=0\) gives \(S^3\).  The incidence
correspondence follows by deleting \(v\) from each incident simplex.
\end{proof}

Join every incident regular bulk endpoint to the V38 universal symbol
\(H_{\rm u}\) by its fixed path \(P_i\), and let \(\mathcal T_v\) be the
resulting finite rooted path tree.  By construction:
\begin{enumerate}
\item every leading face sheet uses paths in \(\mathcal T_v\);
\item every leading hinge disk is the V45 tree-valued filling; and
\item every leading edge three-ball is the V47 tree-valued filling.
\end{enumerate}
Consequently the complete leading coefficient data on
\(\operatorname{Lk}(v)\) define a map
\begin{equation}
 f_v:S^3\longrightarrow\mathcal T_v.                  \label{eq:newS7V48-link-map}
\end{equation}

\begin{lemma}[No \(\pi_3\) obstruction]
\label{lem:newS7V48-no-pi3}
The link map \eqref{eq:newS7V48-link-map} extends piecewise-Lipschitz
continuously to \(B^4\), and every coefficient symbol in the extension has
gap at least \(\delta_{\rm path}\).
\end{lemma}

\begin{proof}
Contract the tree along unique geodesics to its root and cone \(f_v\) with
that contraction.  A finite subdivision makes the map piecewise linear on
tree edges and hence Lipschitz on each four-simplex of the cone.  Its image
lies in the V38 gapped path tree.
\end{proof}

\subsection{Four-normal adiabatic field}
\label{sec:newS7V48-field}

Use the stratified conical extension of the already constructed bulk,
face, hinge, and edge leading fields outside the unit ball in
\(\mathbb R^4\), and the extension of
\Cref{lem:newS7V48-no-pi3} inside.  This gives a piecewise-Lipschitz map
\begin{equation}
 \Phi_v:\mathbb R^4\longrightarrow\mathcal T_v.        \label{eq:newS7V48-Phi}
\end{equation}
After padding the finite hopping ranges, let \(h_d(\tau)\),
\(d\in\mathbb Z^4\), be the corresponding self-adjoint coefficient
family.  Define
\begin{equation}
 (H_{v,L}^{\rm lead}u)_p
 :=\sum_{|d|\le J}h_d(\Phi_v((p+d/2)/L))u_{p+d}.       \label{eq:newS7V48-leading}
\end{equation}

\begin{lemma}[Algebra and leading gap]
\label{lem:newS7V48-leading-gap}
\(H_{v,L}^{\rm lead}\) is self-adjoint and
\begin{equation}
 \|H_{v,L}^{\rm lead}u\|
 \ge(\delta_{\rm path}-C_vL^{-1/2})\|u\|.             \label{eq:newS7V48-leading-gap}
\end{equation}
\end{lemma}

\begin{proof}
Midpoint evaluation and \(h_{-d}=h_d^*\) give self-adjointness.  Localize
\(\mathbb Z^4\) on boxes of side \(\ell\).  Finite hopping range gives a
commutator cost \(C/\ell\), while the Lipschitz variation of
\(\Phi_v(p/L)\) gives a freezing cost \(C\ell/L\).  Every frozen symbol has
gap \(\delta_{\rm path}\).  The V37 direct-sum argument and
\(\ell\asymp L^{1/2}\) prove the estimate.
\end{proof}

\subsection{Attachment of completed boundary strata}
\label{sec:newS7V48-boundary}

Let \(S_{\partial v,L}\) be the finite sum of the V44, V46, and V47
completed lower-stratum corrections, sandwiched by self-adjoint cutoffs
which equal one in their remote conical collars and vanish in a smaller
central vertex ball.

\begin{lemma}[Vertex boundary debit]
\label{lem:newS7V48-boundary}
The cutoffs may be chosen on scale \(L\) so that
\begin{align}
 S_{\partial v,L}^*&=S_{\partial v,L},                 \label{eq:newS7V48-boundary-self}\\
 \|S_{\partial v,L}\|&\le {C_\partial\over L},         \label{eq:newS7V48-boundary-norm}\\
 \left\|\sum_q(S_{\partial v,L})_{pq}\right\|
 &\le {C_\partial\over L},                            \label{eq:newS7V48-boundary-row}
\end{align}
and its four-normal weighted first row moment is uniformly bounded.
\end{lemma}

\begin{proof}
Only finitely many lower strata meet a fixed original vertex.  Every
completion has norm \(O(L^{-1})\), zero frozen row moment, and bounded
weighted first moment.  Sandwiching by cutoffs preserves self-adjointness
and norm.  The cutoff commutator has size \(L^{-1}\) times the weighted
moment, exactly as in \Cref{lem:newS7V47-boundary-debit}.
\end{proof}

Put
\begin{equation}
 H_{v,L}^{\rm ad}:=H_{v,L}^{\rm lead}+S_{\partial v,L}.               \label{eq:newS7V48-Had}
\end{equation}
It matches the complete lower-stratum hierarchy outside the central
vertex ball and obeys
\begin{equation}
 \|H_{v,L}^{\rm ad}u\|
 \ge(\delta_{\rm path}-CL^{-1/2})\|u\|.               \label{eq:newS7V48-Had-gap}
\end{equation}

\subsection{Four-normal conservative completion}
\label{sec:newS7V48-completion}

Let \(H_v^{\rm ref}\) denote the already completed lower-stratum physical
stencil and set
\[
 B_{v,L}:=H_{v,L}^{\rm ad}-H_v^{\rm ref}.
\]
Choose a four-dimensional support set \(S_{v,L}\) containing all nonzero
rows and columns, with
\begin{equation}
 c_1L^4\le N_{v,L}:=|S_{v,L}|\le c_2L^4.              \label{eq:newS7V48-N}
\end{equation}
Define
\begin{align}
 (Q_{v,L}z)_p&:=N_{v,L}^{-1/2}z,                      \label{eq:newS7V48-Q}\\
 Y_{v,L}&:=-B_{v,L}Q_{v,L},                           \label{eq:newS7V48-Y}\\
 R_{v,L}&:=Y_{v,L}Q_{v,L}^*+Q_{v,L}Y_{v,L}^*
 -Q_{v,L}(Q_{v,L}^*Y_{v,L})Q_{v,L}^*.                \label{eq:newS7V48-R}
\end{align}

\begin{lemma}[Four-normal completion estimates]
\label{lem:newS7V48-completion-estimates}
The midpoint and boundary row residuals give
\begin{equation}
 \|Y_{v,L}\|\le {C\over L},\qquad
 \|R_{v,L}\|\le {3C\over L}.                          \label{eq:newS7V48-R-norm}
\end{equation}
Moreover, \(R_{v,L}\) is self-adjoint, has rank at most \(2q\), and
\begin{equation}
 (B_{v,L}+R_{v,L})Q_{v,L}=0.                          \label{eq:newS7V48-zero}
\end{equation}
\end{lemma}

\begin{proof}
The leading midpoint row residual is \(O(L^{-1})\), and
\eqref{eq:newS7V48-boundary-row} has the same size.  Normalizing its
\(N_{v,L}=O(L^4)\) rows by \(N_{v,L}^{-1/2}\) gives
\(\|Y_{v,L}\|\le C/L\), independently of dimension.  The compression
\(Q_{v,L}^*Y_{v,L}=-Q_{v,L}^*B_{v,L}Q_{v,L}\) is self-adjoint.  The
algebra of the V44 completion now gives self-adjointness, rank, the exact
moment identity, and the three-term norm bound.
\end{proof}

\begin{theorem}[Uniform original-vertex gap]
\label{thm:newS7V48-vertex-gap}
The completed vertex kernel
\begin{equation}
 H_{v,L}^{\rm con}:=H_{v,L}^{\rm ad}+R_{v,L}           \label{eq:newS7V48-Hcon}
\end{equation}
satisfies
\begin{equation}
 \|H_{v,L}^{\rm con}u\|
 \ge(\delta_{\rm path}-CL^{-1/2})\|u\|.               \label{eq:newS7V48-vertex-gap}
\end{equation}
It therefore has gap at least \(\delta_{\rm path}/2\) for sufficiently
large \(L\), and its sign differs from that of \(H_{v,L}^{\rm ad}\) by
\(O(L^{-1})\) in operator norm.
\end{theorem}

\begin{proof}
Combine \eqref{eq:newS7V48-Had-gap} and
\eqref{eq:newS7V48-R-norm}.  Apply the gapped sign-function resolvent
estimate for the final assertion.
\end{proof}

\subsection{Moment scaling and continuum safety}
\label{sec:newS7V48-continuum}

Here
\[
 \|(Q_{v,L})_p\|=O(L^{-2}),\qquad
 \|(Y_{v,L})_p\|=O(L^{-3}),
\]
so each completion block is \(O(L^{-5})\).

\begin{lemma}[Four-normal weighted row bound]
\label{lem:newS7V48-weighted-row}
\begin{align}
 \sup_p\sum_q\|(R_{v,L})_{pq}\|&=O(L^{-1}),            \label{eq:newS7V48-row}\\
 \sup_p\sum_q|p-q|\|(R_{v,L})_{pq}\|&=O(1).           \label{eq:newS7V48-first}
\end{align}
\end{lemma}

\begin{proof}
A row has \(O(L^4)\) blocks of size \(O(L^{-5})\).  The sum of distances
from a fixed point over a four-ball of radius \(O(L)\) is \(O(L^5)\).
\end{proof}

Lift all fibres to the vertex by V33 short-geodesic spin--gauge transport.
Smooth variation of the background over the physical ball of radius
\(w=aL\) changes the finite coefficient rows by \(O(w)\) in norm.  This is
a vanishing perturbation of the frozen gap.  Neighbour-to-neighbour
variation is \(O(a)\), so after the exact frozen completion the curved
zeroth moment is \(O(a)\); the first moment remains bounded through two
metric variations.

\begin{theorem}[Continuum safety at an original vertex]
\label{thm:newS7V48-continuum}
If \(L(a)\to\infty\) and \(aL(a)\to0\), then the curved vertex kernel has
gap
\begin{equation}
 \delta_{\rm path}-C\{L(a)^{-1/2}+aL(a)\},             \label{eq:newS7V48-curved-gap}
\end{equation}
and, for every fixed smooth spinor,
\begin{equation}
 \left\|a^{-1}
 (H_{v,L}^{\rm con}-H_v^{\rm ref})\psi
 \right\|_{L^2(B_{CaL}(v))}
 \le C_\psi\{aL(a)\}^2\longrightarrow0.               \label{eq:newS7V48-continuum}
\end{equation}
Both statements hold through two metric contacts.
\end{theorem}

\begin{proof}
The row-sum norm of the smooth geometric perturbation over the ball is
\(O(aL)\), giving \eqref{eq:newS7V48-curved-gap} by the reverse triangle
inequality.  The covariant Taylor argument from V43 has bounded physical
zeroth and first-moment terms.  The vertex ball has volume \(O((aL)^4)\);
its square root gives \eqref{eq:newS7V48-continuum}.  Smooth differentiated
coefficient bounds give the metric-contact statement.
\end{proof}

\subsection{All-stratum local acquisition}
\label{sec:newS7V48-all-strata}

\begin{definition}[Multiscale stratified local gap card, MSLGC]
\label{def:newS7V48-MSLGC}
MSLGC is the existence of one choice \(L(a)\to\infty\),
\(aL(a)\to0\), and one \(c_{\rm loc}>0\) such that every frozen bulk,
open-face, open-hinge, open-edge, and vertex model of the completed
adiabatic hierarchy has dimensionless sign-kernel gap at least
\(c_{\rm loc}\), with curved freezing and two metric contacts controlled
on their corresponding physical collars.
\end{definition}

\begin{theorem}[MSLGC is proved for the completed hierarchy]
\label{thm:newS7V48-MSLGC}
Take, for example,
\begin{equation}
 L(a)=\lfloor a^{-1/2}\rfloor.                         \label{eq:newS7V48-scale-choice}
\end{equation}
Then the bulk theorem V38 and the completed collar theorems V44, V46, V47,
and V48 prove MSLGC with \(c_{\rm loc}=\delta_{\rm path}/4\) for all
sufficiently small \(a\).
\end{theorem}

\begin{proof}
For the chosen scale,
\[
 L^{-1/2}=O(a^{1/4}),\qquad aL=O(a^{1/2}).
\]
Every local gap estimate is
\(\delta_{\rm path}-O(L^{-1/2})-O(aL)\), apart from the regular bulk gap,
which is already uniform.  There are finitely many original strata and
compact parameter families, so one sufficiently small \(a\) makes every
debit at most \(3\delta_{\rm path}/4\).  The two-contact bounds were proved
at each preceding stratum.
\end{proof}

\begin{assessment}[Every simplicial stratum is locally acquired]
\label{ass:newS7V48-status}
The modified conservative adiabatic discretization now has a proved local
sign-kernel gap at every simplicial stratum and matching smooth continuum
moments.  The remaining gap task is global: glue the overlapping collar
regions without losing the lower bound.  Because completion ranges grow
like \(L(a)\), the original fixed-range IMS estimate must be replaced by a
multiscale commutator calculation.
\end{assessment}

\begin{programme}[V49: global multiscale IMS gap]
\label{prog:newS7V48-V49}
V49 will construct one global stratified collar operator in descending
codimension order, verify that the local definitions agree in their
overlap zones, and choose an IMS scale \(R(a)\) satisfying
\[
 L(a)\ll R(a)\ll a^{-1}.
\]
It will estimate commutators of both the fixed-range leading stencil and
the \(L(a)\)-range norm-\(O(L(a)^{-1})\) completions and attempt to derive
one global gap of order \(a^{-1}\) for the physical sign kernel.
\end{programme}

\begin{firewall}[Claim boundary after seventh-series V48]
\label{fire:newS7V48-boundary}
V48 identifies the vertex link as \(S^3\), proves its tree-valued extension,
constructs the four-normal leading and boundary-compatible operators,
performs the rank-\(2q\) completion on \(O(L^4)\) sites, proves the vertex
gap, moment and continuum estimates, and proves the all-stratum local card
MSLGC.

V48 does not prove one global sign-kernel gap, Gaussian locality for the
growing-range completions, the global overlap heat calculus, strong local
heat consistency, the common Einstein stress, pure gravitational owners,
rough metric or AOP control, measure concentration, the determinant-line
phase, interacting cutoff removal, or the Standard Model--Einstein
continuum.
\end{firewall}

\begin{theorem}[Seventh-series V48 result]
\label{thm:newS7V48-result}
The complete \(S^3\) link data at every original vertex extend through the
universal gapped tree to a four-normal collar.  Its conservative completion
has a uniform gap, exact frozen zeroth moment, bounded first moment,
shrinking physical range, and continuum-equivalent smooth action.  Together
with V38, V44, V46, and V47, this proves MSLGC at every simplicial stratum.
\end{theorem}

\begin{proof}
Use \Cref{lem:newS7V48-link,
lem:newS7V48-no-pi3,
lem:newS7V48-leading-gap,
lem:newS7V48-boundary,
lem:newS7V48-completion-estimates,
thm:newS7V48-vertex-gap,
lem:newS7V48-weighted-row,
thm:newS7V48-continuum,
thm:newS7V48-MSLGC}.
\end{proof}

\paragraph{Parent-version record.}
V48 was formed from
\texttt{Renewal\_SM\_Einstein\_Continuum\_Research\_Notes\_V47.tex}.
The V47 parent had \(18\,739\) logical source lines, \(753\,815\) bytes,
and SHA--256
\[
 \texttt{5852CCFEB55A3675EBB86D87BC66243C3B03D19EB81BC346232572A494694B80}.
\]
The next compulsory companion audit is V50.

% =============================================================================
% END APPENDED SEVENTH-SERIES V48 MATERIAL
% =============================================================================

% =============================================================================
% BEGIN APPENDED SEVENTH-SERIES V49 MATERIAL
% =============================================================================

\section{Seventh-series V49: the global multiscale stratified gap}
\label{sec:v7v49}
\label{sec:newS7V49}

\subsection{One global collar hierarchy}
\label{sec:newS7V49-hierarchy}

Let \(\mathcal K_0^{(c)}\) denote the union of original simplices of
codimension at least \(c\), \(c=1,\ldots,4\).  Choose nested tubular cutoff
profiles on physical width \(w(a)=aL(a)\), with fixed separated profile
radii.  The open-face construction is used off the hinge tubes, the hinge
construction overrides it near \(\mathcal K_0^{(2)}\), the one-simplex
edge construction overrides both near \(\mathcal K_0^{(3)}\), and the
vertex construction is used near \(\mathcal K_0^{(4)}\).

\begin{definition}[Global conservative adiabatic sign kernel]
\label{def:newS7V49-global-kernel}
Construct \(H_a^{\rm strat}\) inductively in increasing codimension.  At
stage \(c\), for every open codimension-\(c\) stratum \(S\), let
\(D_{S,a}\) be the difference between its completed local collar model and
the hierarchy already constructed on its boundary strata.  With a
spin--gauge scalar cutoff \(\chi_{S,a}\), add
\begin{equation}
 \chi_{S,a}D_{S,a}\chi_{S,a}.                         \label{eq:newS7V49-splice}
\end{equation}
The cutoff is one in the core of the \(S\)-collar and zero before another
stratum of higher codimension is reached.  At the next stage the higher
stratum model is defined relative to the hierarchy just constructed, so
the induction telescopes in every core.
\end{definition}

\begin{lemma}[Global algebra and matching]
\label{lem:newS7V49-global-algebra}
\(H_a^{\rm strat}\) is self-adjoint and spin--gauge covariant.  It equals
the regular bulk kernel outside the skeleton collars and equals the
corresponding completed face, hinge, edge, or vertex model in the core of
each stratified collar.
\end{lemma}

\begin{proof}
Every \(D_{S,a}\) is self-adjoint because it is a difference of
self-adjoint kernels, and
\eqref{eq:newS7V49-splice} is self-adjoint.  The cutoffs are scalar in
spin--gauge fibres and all local models are covariant, proving covariance.
The lower-stratum models were constructed to match their boundary
hierarchy outside their central regions.  Where \(\chi_{S,a}=1\), the
difference in \eqref{eq:newS7V49-splice} replaces the preceding hierarchy
by the \(S\)-model; where it is zero, nothing changes.  The fixed nested
profiles make the statements compatible on every transition annulus.
\end{proof}

\begin{corollary}[Global gamma-five Hermiticity]
\label{cor:newS7V49-gamma5}
Define the associated Wilson kernel by
\begin{equation}
 D_{W,a}^{\rm strat}-{\rho\over a}
 :={1\over a}\Gamma H_a^{\rm strat}.                  \label{eq:newS7V49-DW}
\end{equation}
Then
\begin{equation}
 (D_{W,a}^{\rm strat}-\rho/a)^*
 =\Gamma(D_{W,a}^{\rm strat}-\rho/a)\Gamma.           \label{eq:newS7V49-gamma5}
\end{equation}
\end{corollary}

\begin{proof}
Use self-adjointness of \(H_a^{\rm strat}\) and \(\Gamma^2=I\).
\end{proof}

\subsection{Uniform weighted row ledger}
\label{sec:newS7V49-row-ledger}

Split the global dimensionless kernel as
\begin{equation}
 H_a^{\rm strat}=H_a^{\rm lead}+\mathcal R_a,          \label{eq:newS7V49-split}
\end{equation}
where \(H_a^{\rm lead}\) contains the finite-range midpoint coefficient
fields and \(\mathcal R_a\) contains every conservative low-rank
completion and its hierarchical cutoff.

\begin{lemma}[Global completion bounds]
\label{lem:newS7V49-global-completion}
There is a constant independent of \(a,L\) such that
\begin{align}
 \|\mathcal R_a\|&\le {C\over L},                     \label{eq:newS7V49-R-norm}\\
 \sup_x\sum_y\|\mathcal R_a(x,y)\|&\le {C\over L},    \label{eq:newS7V49-R-row}\\
 \sup_x\sum_y d_a(x,y)\|\mathcal R_a(x,y)\|&\le C,    \label{eq:newS7V49-R-moment}
\end{align}
and the analogous column estimates hold.  Here \(d_a\) is lattice distance
in units of \(a\).
\end{lemma}

\begin{proof}
At any vertex of the cofinal mesh, only a uniformly bounded number of
original-stratum collars overlap.  V44, V46, V47, and V48 give the three
bounds for each completion.  Multiplication by cutoffs does not increase
the norm or unweighted Schur row bound.  Its commutation debit is controlled
by the already bounded weighted row moment, and the cutoffs change only on
scale \(L\).  Summing the bounded overlap proves the result.  Column bounds
follow from self-adjointness.
\end{proof}

\begin{lemma}[Leading finite-range bound]
\label{lem:newS7V49-leading-bound}
The leading kernel has one fixed lattice range \(J\) and uniform row and
column sums:
\begin{equation}
 \sup_x\sum_y\|H_a^{\rm lead}(x,y)\|
 +\sup_y\sum_x\|H_a^{\rm lead}(x,y)\|\le C_{\rm lead}.               \label{eq:newS7V49-leading-bound}
\end{equation}
\end{lemma}

\begin{proof}
All V38 tree paths were padded to one finite hopping range.  Their compact
coefficient image has a finite row bound.  The base triangulation has
bounded degree, and only boundedly many hierarchical profiles meet a row.
\end{proof}

\subsection{A partition at the intermediate scale}
\label{sec:newS7V49-partition}

Choose integers \(R(a)\) satisfying
\begin{equation}
 L(a)\ll R(a)\ll a^{-1}.                              \label{eq:newS7V49-scales}
\end{equation}
There is a quadratic partition of unity
\(\{\chi_\alpha\}\) on the global vertex set, supported in lattice balls
of radius \(C R\), with bounded overlap and
\begin{equation}
 \sum_\alpha|\chi_\alpha(x)-\chi_\alpha(y)|^2
 \le C\min\left\{1,{d_a(x,y)^2\over R^2}\right\}.      \label{eq:newS7V49-partition}
\end{equation}

\begin{lemma}[Multiscale commutator estimate]
\label{lem:newS7V49-commutator}
The complete global kernel obeys
\begin{equation}
 \left(\sum_\alpha
 \|[H_a^{\rm strat},\chi_\alpha]u\|^2\right)^{1/2}
 \le {C_{\rm IMS}\over R}\|u\|.                       \label{eq:newS7V49-commutator}
\end{equation}
\end{lemma}

\begin{proof}
For the fixed-range leading kernel, use
\eqref{eq:newS7V49-partition} and the row/column Schur bounds to obtain
\(C/R\).  For \(\mathcal R_a\),
\[
 [\mathcal R_a,\chi_\alpha](x,y)
 =\mathcal R_a(x,y)\{\chi_\alpha(y)-\chi_\alpha(x)\}.
\]
The square-summed partition difference is bounded by
\(Cd_a(x,y)^2/R^2\).  Apply the vector-valued Schur test with the weighted
row and column bounds \eqref{eq:newS7V49-R-moment}; this again gives
\(C/R\).  Minkowski's inequality combines the two parts.
\end{proof}

The growing lattice range \(L\) therefore causes no factor \(L/R\):
its \(O(L^{-1})\) block mass and bounded first moment have already paid for
that range.

\subsection{Comparison with one acquired local model}
\label{sec:newS7V49-local-comparison}

For every partition centre \(x_\alpha\), choose the deepest original
stratum whose \(CR\)-ball meets its core.  Let \(H_\alpha^{\rm mod}\) be
the corresponding completed frozen bulk, face, hinge, edge, or vertex
model from MSLGC, translated to that centre.

\begin{lemma}[One-model coverage]
\label{lem:newS7V49-one-model}
For sufficiently small \(a\), each partition ball meets at most one
original vertex and only the finite incidence fan of one deepest stratum.
The model \(H_\alpha^{\rm mod}\) contains every lower-stratum collar which
can occur in that ball.
\end{lemma}

\begin{proof}
The original triangulation is finite, so distinct disjoint original
strata have a positive separation after removing fixed neighbourhoods of
their common boundary.  The physical partition radius is \(aR\to0\), while
the complete collar width \(aL\) satisfies \(aL\ll aR\).  Thus a small
partition ball sees only the local link of its deepest incident stratum.
The recursive constructions V44--V48 included the entire lower-stratum
fan in that link.
\end{proof}

\begin{lemma}[Curved-to-frozen comparison]
\label{lem:newS7V49-freezing}
On the support of \(\chi_\alpha\),
\begin{equation}
 \|(H_a^{\rm strat}-H_\alpha^{\rm mod})\chi_\alpha\|
 \le C_{\rm fr}\,aR.                                  \label{eq:newS7V49-freezing}
\end{equation}
The constant is uniform through two metric variations.
\end{lemma}

\begin{proof}
Over a physical ball of radius \(O(aR)\), the smooth coframe, metric,
spin--gauge transports, and finite tree-path coefficients differ from
their central frozen values by \(O(aR)\).  The leading row sum is uniformly
bounded.  The completion formula is a finite algebraic expression in its
defect and normalized constant isometry; V44--V48 give uniformly bounded
parameter derivatives and an overall completion norm \(O(L^{-1})\).
Cutoff profiles are the same functions of the normalized collar
coordinates in the model and global operator.  The block Schur test gives
\eqref{eq:newS7V49-freezing}.  Differentiating the smooth data twice gives
the metric-contact version.
\end{proof}

\subsection{The global dimensionless gap}
\label{sec:newS7V49-gap}

\begin{theorem}[Global multiscale IMS gap]
\label{thm:newS7V49-global-gap}
Let \(c_{\rm loc}>0\) be the MSLGC gap from V48.  Then
\begin{equation}
 \|H_a^{\rm strat}u\|
 \ge
 \left(c_{\rm loc}-{C_{\rm IMS}\over R}-C_{\rm fr}aR\right)\|u\|.
                                                               \label{eq:newS7V49-global-gap}
\end{equation}
\end{theorem}

\begin{proof}
For every \(\alpha\), the local model gap gives
\[
 c_{\rm loc}\|\chi_\alpha u\|
 \le\|H_\alpha^{\rm mod}\chi_\alpha u\|
 \le\|\chi_\alpha H_a^{\rm strat}u\|
   +\|[H_a^{\rm strat},\chi_\alpha]u\|
   +\|(H_\alpha^{\rm mod}-H_a^{\rm strat})\chi_\alpha u\|.
\]
Take the \(\ell^2\)-norm in \(\alpha\).  The quadratic partition gives
\[
 \sum_\alpha\|\chi_\alpha u\|^2=\|u\|^2,\qquad
 \sum_\alpha\|\chi_\alpha H_a^{\rm strat}u\|^2
 =\|H_a^{\rm strat}u\|^2.
\]
Use \Cref{lem:newS7V49-commutator,lem:newS7V49-freezing} and Minkowski's
inequality, then rearrange.
\end{proof}

\begin{corollary}[Explicit scale choice and physical gap]
\label{cor:newS7V49-scale-gap}
With
\begin{equation}
 L(a)=\lfloor a^{-1/2}\rfloor,\qquad
 R(a)=\lfloor a^{-3/4}\rfloor,                        \label{eq:newS7V49-scale-choice}
\end{equation}
one has \(L/R=O(a^{1/4})\), \(R^{-1}=O(a^{3/4})\), and
\(aR=O(a^{1/4})\).  For sufficiently small \(a\),
\begin{equation}
 \|H_a^{\rm strat}u\|\ge {c_{\rm loc}\over2}\|u\|.    \label{eq:newS7V49-dimensionless-gap}
\end{equation}
The physical Wilson sign kernel \(a^{-1}H_a^{\rm strat}\) therefore has
gap at least \(c_{\rm loc}/(2a)\).
\end{corollary}

\begin{proof}
The scale relations are immediate.  Insert them in
\eqref{eq:newS7V49-global-gap}.
\end{proof}

\subsection{Global smooth-core equivalence}
\label{sec:newS7V49-continuum}

\begin{theorem}[Global continuum consistency of the stratified modification]
\label{thm:newS7V49-continuum}
For every fixed smooth spinor \(\psi\),
\begin{equation}
 \left\|a^{-1}(H_a^{\rm strat}-H_a^{\rm sh})\psi\right\|_{L^2(M)}
 \longrightarrow0,                                   \label{eq:newS7V49-continuum}
\end{equation}
and the convergence holds through the two metric contacts used in the
determinant modulus.
\end{theorem}

\begin{proof}
Outside the skeleton collars the two kernels agree.  On open-face collars,
V44 gives an \(O((aL)^{1/2})\) norm; on hinge tubes V46 gives
\(O(aL)\); on edge tubes V47 gives \(O((aL)^{3/2})\); and on vertex balls
V48 gives \(O((aL)^2)\).  The original triangulation has finitely many
strata.  Assign every overlap to its deepest stratum; the resulting regions
are disjoint and the squared estimates add.  Since \(aL\to0\), the sum
tends to zero.  Each stratum theorem includes two metric contacts.
\end{proof}

\begin{assessment}[The global gap is closed for the modified kernel]
\label{ass:newS7V49-status}
V49 constructs one global self-adjoint covariant stratified kernel and
proves its uniform dimensionless gap, hence its \(a^{-1}\) physical gap.
It also proves smooth-core equivalence with the sharp V33 kernel through
two metric contacts.  This closes the former SLGR/SFGC gap obstruction for
the modified discretization.

The next obstruction is locality at the strength needed for heat-kernel
coefficients.  The leading operator has fixed range, but each conservative
completion has lattice range \(L(a)\) and norm \(O(L(a)^{-1})\).
Operator-norm convergence of sign functions does not by itself control
local diagonal traces after the number of lattice modes grows.
\end{assessment}

\begin{programme}[V50: compulsory audit and completion-locality ledger]
\label{prog:newS7V49-V50}
V50 will perform the compulsory YM/NS audit.  It will then decompose the
overlap sign and heat functionals into the fixed-range leading kernel plus
the conservative completion, use resolvent identities to derive weighted
off-diagonal bounds, and identify the exact Schatten/local-trace rate
needed for the completion to disappear from the continuum heat
coefficients.
\end{programme}

\begin{firewall}[Claim boundary after seventh-series V49]
\label{fire:newS7V49-boundary}
V49 defines the global hierarchical collar kernel, proves its algebra and
matching, establishes global weighted row and commutator estimates,
constructs the intermediate-scale partition, proves the curved-to-frozen
comparison, obtains a uniform global dimensionless and physical gap, and
proves global smooth-core continuum equivalence through two metric
contacts.

V49 does not prove Gaussian or local-trace locality for the growing-range
completions, the overlap heat coefficients for the modified kernel, strong
local heat consistency, the common Einstein stress, pure gravitational
owners, rough metric or AOP control, measure concentration, the
determinant-line phase, interacting cutoff removal, or the Standard
Model--Einstein continuum.
\end{firewall}

\begin{theorem}[Seventh-series V49 result]
\label{thm:newS7V49-result}
For \(L=a^{-1/2}\) and \(R=a^{-3/4}\), the global conservative adiabatic
kernel satisfies
\[
 \|H_a^{\rm strat}u\|\ge(c_{\rm loc}/2)\|u\|,
\]
so its physical sign kernel has gap \(c_{\rm loc}/(2a)\).  It is
self-adjoint, spin--gauge covariant, gamma-five Hermitian in Wilson form,
and agrees with the sharp V33 kernel on every fixed smooth core in the
continuum limit through two metric contacts.
\end{theorem}

\begin{proof}
Use \Cref{lem:newS7V49-global-algebra,
cor:newS7V49-gamma5,
lem:newS7V49-global-completion,
lem:newS7V49-commutator,
lem:newS7V49-one-model,
lem:newS7V49-freezing,
thm:newS7V49-global-gap,
cor:newS7V49-scale-gap,
thm:newS7V49-continuum}.
\end{proof}

\paragraph{Parent-version record.}
V49 was formed from
\texttt{Renewal\_SM\_Einstein\_Continuum\_Research\_Notes\_V48.tex}.
The V48 parent had \(19\,105\) logical source lines, \(768\,084\) bytes,
and SHA--256
\[
 \texttt{6DD672DD1826488399383DE86191740D6661DE66E3BA36BC7DBDF74ACFF7049C}.
\]
The next compulsory companion audit is V50.

% =============================================================================
% END APPENDED SEVENTH-SERIES V49 MATERIAL
% =============================================================================

% =============================================================================
% BEGIN APPENDED SEVENTH-SERIES V50 MATERIAL
% =============================================================================

\section{Seventh-series V50: compulsory audit and the completion-locality ledger}
\label{sec:v7v50}
\label{sec:newS7V50}

\subsection{Compulsory companion audit}
\label{sec:newS7V50-audit}

\begin{audit}[Companion snapshots used at V50]
\label{aud:newS7V50-snapshots}
The exact coherent companion snapshots read before this note were:
\begin{center}
\begin{tabular}{p{0.20\linewidth}p{0.52\linewidth}r}
programme & source & bytes\\ \hline
Yang--Mills &
\texttt{Renewal\_Yang\_Mills\_Mass\_Gap\_Research\_Notes\_V25.tex}
& \(392\,293\)\\
Navier--Stokes &
\texttt{Renewal\_NS\_Stability\_Research\_Notes\_V83.tex}
& \(864\,809\)
\end{tabular}
\end{center}
The Yang--Mills source had \(10\,996\) newline characters and SHA--256
\[
 \texttt{8D2598532555A9239FEA347533E57AE802F43D4DEC214F065BB8F91539D8FB08}.
\]
The Navier--Stokes source had \(16\,820\) newline characters and SHA--256
\[
 \texttt{28DD8FF205A2D211A76496E833BC4DEE4DF49ED944C70865AE238B56FA098ACB}.
\]
The NS programme advanced from the initially listed V82 to V83 before the
atomic read; the V83 digest is the snapshot actually audited.
\end{audit}

The Yang--Mills note proves an exact sector-weight/within-sector
total-variation decomposition and a disagreement-path theorem with
absorption criterion
\((\Delta-1)(a+p)<1\).  It explicitly warns that changing sector weights
would change the theory.  The Navier--Stokes note computes an Oseen kernel
constant, distinguishes a certified bound from quadrature, and propagates
the certified value into explicit regularity floors without claiming
global regularity.

\begin{decision}[Effect of the V50 audit]
\label{dec:newS7V50-audit}
For the SM completion, a small operator norm is the analogue of
within-sector contraction: it does not count how many ultraviolet modes
see the perturbation.  Schatten rank and heat smoothing must be tracked
separately.  In addition, a computed diagonal decay cannot replace a
certified weighted resolvent estimate.  V50 therefore records rigorous
weighted and trace-class bounds and states explicitly where they stop
short of the heat coefficient.
\end{decision}

\subsection{Leading and completed global kernels}
\label{sec:newS7V50-kernels}

Use the V49 decomposition
\[
 H_a:=H_a^{\rm strat}=H_{0,a}+\mathcal R_a,
 \qquad H_{0,a}:=H_a^{\rm lead}.
\]
Both are dimensionless.  V49 gives
\(\|\mathcal R_a\|\le C/L\) and
\(\|H_a u\|\ge c_{\rm loc}\|u\|/2\).

\begin{lemma}[The leading kernel is also uniformly gapped]
\label{lem:newS7V50-leading-gap}
For sufficiently small \(a\),
\begin{equation}
 \|H_{0,a}u\|\ge {c_{\rm loc}\over3}\|u\|.             \label{eq:newS7V50-leading-gap}
\end{equation}
\end{lemma}

\begin{proof}
By the reverse triangle inequality,
\[
 \|H_{0,a}u\|
 \ge\|H_au\|-\|\mathcal R_au\|
 \ge(c_{\rm loc}/2-C/L)\|u\|.
\]
Take \(L\) large enough that \(C/L\le c_{\rm loc}/6\).
\end{proof}

For a block matrix \(A\) on the lattice, define the two-sided weighted
Schur norm
\begin{equation}
 \|A\|_{\mu,{\rm Sch}}
 :=\max\left\{
 \sup_x\sum_y e^{\mu d_a(x,y)}\|A(x,y)\|,
 \sup_y\sum_x e^{\mu d_a(x,y)}\|A(x,y)\|
 \right\}.                                             \label{eq:newS7V50-Schur}
\end{equation}

\begin{lemma}[Weighted size of the completion]
\label{lem:newS7V50-weighted-R}
There is \(\gamma_0>0\) such that, for
\(\mu_L=\gamma_0/L\),
\begin{equation}
 \|\mathcal R_a\|_{\mu_L,{\rm Sch}}\le {C_R\over L}.   \label{eq:newS7V50-weighted-R}
\end{equation}
The estimate holds through two metric contacts.
\end{lemma}

\begin{proof}
Every completion block has lattice range at most \(C_0L\).  Therefore
\(e^{\mu_Ld_a(x,y)}\le e^{C_0\gamma_0}\) on its support.  Apply the
unweighted row and column estimates of
\Cref{lem:newS7V49-global-completion}.  The preceding stratum theorems
give the same estimates for two differentiated completions.
\end{proof}

\subsection{Certified resolvent and sign locality}
\label{sec:newS7V50-resolvent}

\begin{lemma}[Leading Combes--Thomas bound]
\label{lem:newS7V50-CT-leading}
Let \(\Gamma_\pm\) be fixed contours enclosing the positive and negative
spectra of every \(H_{0,a}\), respectively, while staying a fixed distance
from those spectra and from zero.  For sufficiently small fixed
\(\mu_0>0\),
\begin{equation}
 \sup_{z\in\Gamma_+\cup\Gamma_-}
 \|(H_{0,a}-z)^{-1}\|_{\mu_0,{\rm Sch}}\le C_{\rm CT}. \label{eq:newS7V50-CT-leading}
\end{equation}
\end{lemma}

\begin{proof}
\(H_{0,a}\) has fixed lattice range, a uniform Schur bound, and the gap
\eqref{eq:newS7V50-leading-gap}.  Conjugating by
\(e^{\mu d_a(\,\cdot\,,x_0)}\) changes the operator by \(O(\mu)\) in norm.
Choose \(\mu_0\) so that this is smaller than half the contour distance,
and use the resolvent identity.  The resulting pointwise exponential bound
has an exponent independent of \(a\); summing it over the bounded-geometry
four-dimensional lattice at a slightly smaller exponent gives the
two-sided Schur estimate.
\end{proof}

\begin{theorem}[Completed resolvent locality]
\label{thm:newS7V50-resolvent}
For \(\mu_L=\gamma_0/L\) with \(\gamma_0\) sufficiently small,
\begin{equation}
 \sup_{z\in\Gamma_+\cup\Gamma_-}
 \|(H_a-z)^{-1}\|_{\mu_L,{\rm Sch}}\le C              \label{eq:newS7V50-resolvent}
\end{equation}
uniformly in \(a\).  Moreover,
\begin{equation}
 \|(H_a-z)^{-1}-(H_{0,a}-z)^{-1}\|_{\mu_L,{\rm Sch}}
 \le {C\over L}.                                      \label{eq:newS7V50-resolvent-difference}
\end{equation}
\end{theorem}

\begin{proof}
The leading resolvent bound remains valid at the smaller weight
\(\mu_L\).  The weighted Schur matrices form a Banach algebra.  By
\Cref{lem:newS7V50-weighted-R},
\[
 \|(H_{0,a}-z)^{-1}\mathcal R_a\|_{\mu_L,{\rm Sch}}
 \le C/L<1/2
\]
for large \(L\).  The Neumann resolvent formula gives
\eqref{eq:newS7V50-resolvent}.  The resolvent identity with one factor
\(\mathcal R_a\) gives
\eqref{eq:newS7V50-resolvent-difference}.
\end{proof}

\begin{theorem}[Weighted overlap-sign locality]
\label{thm:newS7V50-sign-locality}
Let \(S_a=\operatorname{sgn}(H_a)\) and
\(S_{0,a}=\operatorname{sgn}(H_{0,a})\).  Then
\begin{align}
 \|S_a\|_{\mu_L,{\rm Sch}}&\le C,                     \label{eq:newS7V50-sign-weighted}\\
 \|S_a-S_{0,a}\|_{\mu_L,{\rm Sch}}&\le {C\over L}.    \label{eq:newS7V50-sign-difference}
\end{align}
Thus the completion-created sign tail has lattice decay length \(O(L)\),
physical decay length \(O(aL)\), and total weighted row mass \(O(L^{-1})\).
The assertions hold through two metric contacts.
\end{theorem}

\begin{proof}
Represent the positive and negative spectral projections by the two fixed
contours and use \Cref{thm:newS7V50-resolvent}.  Their difference uses
\eqref{eq:newS7V50-resolvent-difference}.  Metric differentiation inserts
the bounded contacts and additional resolvents; the V49/V50 weighted
bounds are stable under the finite products.
\end{proof}

\subsection{Global rank and trace norm}
\label{sec:newS7V50-trace}

\begin{lemma}[Rank ledger by stratum]
\label{lem:newS7V50-rank}
On a fixed compact four-volume, the global completion satisfies
\begin{equation}
 \operatorname{rank}\mathcal R_a
 \le C\sum_{c=1}^4a^{-(4-c)}
 \le C'a^{-3}.                                        \label{eq:newS7V50-rank}
\end{equation}
\end{lemma}

\begin{proof}
For a codimension-\(c\) stratum, each tangential fundamental fibre carries
a completion of rank at most a fixed multiple of the spin--internal
dimension.  A quasi-uniform discretization of its compact
\((4-c)\)-dimensional volume has \(O(a^{-(4-c)})\) tangential fibres.
Cutoff multiplication cannot increase rank.  There are finitely many
original strata; the face contribution \(c=1\) dominates.
\end{proof}

\begin{corollary}[Trace-class completion and sign difference]
\label{cor:newS7V50-trace-class}
\begin{align}
 \|\mathcal R_a\|_1&\le {C a^{-3}\over L},             \label{eq:newS7V50-R-trace}\\
 \|S_a-S_{0,a}\|_1&\le {C a^{-3}\over L}.             \label{eq:newS7V50-sign-trace}
\end{align}
The same scale holds through two metric contacts.
\end{corollary}

\begin{proof}
For a finite-rank operator,
\(\|A\|_1\le\operatorname{rank}(A)\|A\|\); use V49 and
\Cref{lem:newS7V50-rank}.  The contour resolvent identity gives
\[
 S_a-S_{0,a}
 ={1\over2\pi i}\sum_{\pm}(\pm1)
 \int_{\Gamma_\pm}
 (H_a-z)^{-1}\mathcal R_a(H_{0,a}-z)^{-1}\,dz.
\]
Bound the outer factors in operator norm and the middle factor in trace
norm.  Differentiated contour formulas contain only a fixed number of
such products.
\end{proof}

\begin{corollary}[Bounded spectral densities do not see the completion]
\label{cor:newS7V50-bounded-density}
For every uniformly bounded insertion \(B_a\),
\begin{equation}
 a^4\left|
 \operatorname{Tr}\{B_a(S_a-S_{0,a})\}
 \right|
 \le C\|B_a\|\,{a\over L}\longrightarrow0.            \label{eq:newS7V50-bounded-density}
\end{equation}
\end{corollary}

\begin{proof}
Use trace duality and \eqref{eq:newS7V50-sign-trace}.
\end{proof}

\subsection{Why this does not yet prove heat-coefficient locality}
\label{sec:newS7V50-heat-barrier}

Define the dimensionless overlap blocks
\begin{equation}
 Q_a:=\rho(I+\Gamma S_a),\qquad
 Q_{0,a}:=\rho(I+\Gamma S_{0,a}),                     \label{eq:newS7V50-Q}
\end{equation}
and \(A_a=Q_a^*Q_a\), \(A_{0,a}=Q_{0,a}^*Q_{0,a}\).
Then
\[
 \|A_a-A_{0,a}\|_1\le Ca^{-3}/L.
\]
For dimensionless heat time \(T\), Duhamel's formula gives only
\begin{equation}
 \left|\operatorname{Tr}
 (e^{-TA_a}-e^{-TA_{0,a}})\right|
 \le C T\,{a^{-3}\over L}.                            \label{eq:newS7V50-crude-heat}
\end{equation}
At physical heat time \(t\), \(T=t/a^2\), so after volume normalization
the right side is
\begin{equation}
 a^4\times C{t\over a^2}{a^{-3}\over L}
 ={Ct\over aL}.                                       \label{eq:newS7V50-crude-bad}
\end{equation}
This bound is useless when \(aL\to0\).  It is an insufficiency of the
estimate, not evidence that the heat traces actually diverge.

\begin{proposition}[Operator-norm smooth-core convergence does not fix \(a_4\)]
\label{prop:newS7V50-core-not-heat}
A family of differential or lattice operators may differ from a reference
only in a collar of width \(w\to0\) and converge on every fixed smooth core,
while derivatives of its principal coefficients grow like \(w^{-1}\).
Therefore smooth-core convergence and a uniform gap do not by themselves
imply convergence of the local \(a_4\) heat coefficient.
\end{proposition}

\begin{proof}
Let a principal coefficient make a fixed-amplitude smooth transition
\(g_w(x)=g_-+\chi(x/w)(g_+-g_-)\) across a hypersurface.  Applied to a fixed
smooth function, the operator difference is bounded and supported in
volume \(O(w)\), so its \(L^2\) norm is \(O(w^{1/2})\).  But
\(\partial g_w=O(w^{-1})\) and
\(\partial^2g_w=O(w^{-2})\).  The local heat polynomial of a
second-order operator with variable principal symbol contains curvature
and connection expressions formed from such derivatives.  Their
integrals are not controlled by the smooth-core \(L^2\) estimate.  Hence
the latter cannot imply \(a_4\) convergence.
\end{proof}

In the current construction, the V38 universal path may change the
physical first Dirac moment inside the shrinking collar, and the V44--V48
completions cancel constants but have only a bounded, not vanishing, first
moment.  These facts are harmless for V49's smooth-core estimate but are
not harmless enough for a local heat theorem.

\subsection{The heat-flat jet card}
\label{sec:newS7V50-HFJC}

\begin{definition}[Heat-flat jet card, HFJC]
\label{def:newS7V50-HFJC}
HFJC consists of:
\begin{enumerate}
\item a replacement of every V38 path by a gapped path whose invariant
physical Dirac first jet is independent of the path parameter and whose
Wilson second jet is uniformly controlled;
\item a self-adjoint covariant completion which matches constants and
all physical affine nodal modes, not constants alone;
\item through two metric contacts, vanishing bounds for the spatial
derivatives of every residual principal-jet coefficient at the orders
entering \(a_4\); and
\item a discrete Gaussian/Schatten smoothing estimate which improves
\eqref{eq:newS7V50-crude-heat} to a collar-volume bound on the MATC
mesoscopic interval.
\end{enumerate}
\end{definition}

\begin{definition}[Completion heat smoothing card, CHSC]
\label{def:newS7V50-CHSC}
For \(w=aL\), CHSC is a scale \(\varepsilon(a)\downarrow0\), with
\(a^2/\varepsilon(a)\to0\), such that the inserted relative heat
difference \(\Delta H_a^{\rm rel}(t)\) between the completed and leading
kernels has a local expansion whose coefficient differences tend to zero
and whose remainder satisfies the MATC logarithmic integrability row.
A sufficient quantitative leading bound is
\begin{equation}
 |\partial_g^j\Delta H_a^{\rm rel}(t)|
 \le C_j w\,t^{-2-j/2},
 \qquad j=0,1,2,\quad
 \varepsilon(a)\le t\le t_0,                          \label{eq:newS7V50-CHSC}
\end{equation}
after the declared lower heat coefficients are subtracted, together with
the corresponding integrable remainder improvement.
\end{definition}

The displayed rate is a target, not a result.  With \(w=a^{1/2}\), a
mesoscopic choice \(\varepsilon(a)=a^\beta\) with sufficiently small
\(\beta>0\) can make each fixed power \(w\varepsilon(a)^{-N}\) vanish,
while still satisfying \(a^2/\varepsilon(a)\to0\).  The exact exponent
needed must come from the heat-word estimate, not be chosen retroactively.

\begin{assessment}[Certified locality obtained; heat-flatness remains]
\label{ass:newS7V50-status}
V50 proves weighted resolvent and sign locality on physical range \(aL\),
and proves trace-class invisibility for bounded normalized spectral
densities.  It also proves that the available trace-norm Duhamel bound is
too weak at physical heat time and identifies a principal-jet defect not
seen by smooth-core convergence.  HFJC and CHSC are the first open heat
rows for the modified discretization.
\end{assessment}

\begin{programme}[V51: affine-moment completion]
\label{prog:newS7V50-V51}
V51 will first replace the V38 hub paths by moment-preserving paths in the
affine space of symbols with one common physical Dirac first jet.  It will
then generalize the V44 completion from the normalized constant subspace
to normalized constant and affine transverse modes, and add the finite
tangential Fourier-jet corrections required to match all four physical
first moments.  The gap survives only if the enlarged correction remains
\(O(L^{-1})\); that estimate will be proved rather than assumed.
\end{programme}

\begin{firewall}[Claim boundary after seventh-series V50]
\label{fire:newS7V50-boundary}
V50 performs the compulsory YM/NS audit, proves a gap for the leading
kernel, weighted completion, resolvent, and sign locality, derives the
global rank and trace-class ledger, and proves invisibility for bounded
normalized spectral densities.  It derives the exact crude heat bound and
proves why it is insufficient, then formulates HFJC and CHSC.

V50 does not prove affine-jet flatness, CHSC, MATC for the completed
kernel, convergence of its \(a_4\) coefficients, the common Einstein
stress, pure gravitational owners, rough metric or AOP control, measure
concentration, the determinant-line phase, interacting cutoff removal, or
the Standard Model--Einstein continuum.
\end{firewall}

\begin{theorem}[Seventh-series V50 result]
\label{thm:newS7V50-result}
The conservative completion changes the overlap sign by
\(O(L^{-1})\) in an exponential Schur norm of physical range \(aL\), and
by \(O(a^{-3}/L)\) in trace norm on a fixed four-volume.  These estimates
prove bounded-density locality but do not control the mesoscopic heat
trace, because the direct Duhamel bound scales as \(t/(aL)\) after volume
normalization and the current collar does not preserve all affine
principal jets.
\end{theorem}

\begin{proof}
Use \Cref{lem:newS7V50-leading-gap,
lem:newS7V50-weighted-R,
thm:newS7V50-resolvent,
thm:newS7V50-sign-locality,
lem:newS7V50-rank,
cor:newS7V50-trace-class,
cor:newS7V50-bounded-density,
prop:newS7V50-core-not-heat}.
\end{proof}

\paragraph{Parent-version record.}
V50 was formed from
\texttt{Renewal\_SM\_Einstein\_Continuum\_Research\_Notes\_V49.tex}.
The V49 parent had \(19\,478\) logical source lines, \(783\,145\) bytes,
and SHA--256
\[
 \texttt{32CDDCF742BB53F6AA195C9E57D5571B861EE7B370883DBD02CE5A5486DC92F3}.
\]
The next compulsory companion audit is V55.

% =============================================================================
% END APPENDED SEVENTH-SERIES V50 MATERIAL
% =============================================================================

% =============================================================================
% BEGIN APPENDED SEVENTH-SERIES V51 MATERIAL
% =============================================================================

\section{Seventh-series V51: invariant affine-jet completion}
\label{sec:v7v51}
\label{sec:newS7V51}

\subsection{Why the hub must retain the physical coframe}
\label{sec:newS7V51-hub}

V50 isolated an order-one principal-jet excursion in the stratified
collars.  The type-to-canonical part of V38 already preserves the derivative
of the Dirac symbol, but the subsequent polar deformation to the identity
changes that derivative.  The latter deformation was useful for comparing
arbitrary coframes.  It is unnecessary for the geometric seams used here:
after short-geodesic spin transport, all simplex types meeting at one
physical point see the same continuum coframe.  This note therefore uses a
coframe-dependent hub and never deforms the physical Dirac jet.

Because the regular type set is finite, pass once to a common periodic
supercell.  Every frozen type is then a matrix-valued Laurent operator on
one lattice \(\Lambda\), possibly with a finite cell fibre.  Let
\(\kappa\in\mathbb T^4\) be the corresponding Bloch variable and let
\[
 J_E:T_x^*M\longrightarrow\mathbb R^4
\]
denote the oriented physical coframe map in the Clifford convention of
V33--V34.  Physical covectors are represented on the common cell by the
nodal restrictions of affine functions.

\begin{lemma}[Invariant endpoint affine jet]
\label{lem:newS7V51-endpoint-jet}
After the exact incidence and spin identifications, every regular frozen
V33 type \(\tau\) at the same physical coframe \(E\) has
\begin{align}
 s_{\tau,E}(0)&=0,                                    \label{eq:newS7V51-s0}\\
 D_\kappa s_{\tau,E}(0)\,\kappa&=J_E\kappa,            \label{eq:newS7V51-s1}\\
 \ell_{\tau,E}(0)&=0,\qquad
 D_\kappa\ell_{\tau,E}(0)=0.                           \label{eq:newS7V51-l01}
\end{align}
Here \(\kappa\) is read as the physical covector represented by the
supercell affine mode.  The identities, including their first two metric
variations, are independent of \(\tau\).
\end{lemma}

\begin{proof}
The element identity
\eqref{eq:newS7V41-affine-exact} says that the nodal interpolant reproduces
every physical affine function exactly.  Applying the assembled Dirac form
and using \eqref{eq:newS7V41-D-affine} on the periodic cell removes the
finite-window boundary term and gives Clifford multiplication by the same
physical covector.  Constants give \eqref{eq:newS7V51-s0}.  The stiffness
row sum vanishes, and its symmetric edge form kills the first derivative at
zero, giving \eqref{eq:newS7V51-l01}.  These are algebraic finite-element
identities, so differentiating the common coframe and mass normalization
preserves them through the two declared metric contacts.
\end{proof}

Choose four generating supercell displacements and their reciprocal
coordinates.  Define
\begin{equation}
 s_E^{\rm hub}(\kappa)
 :=J_E(\sin\kappa_1,\ldots,\sin\kappa_4)^T,\qquad
 \ell^{\rm hub}(\kappa):=
 2\sum_{\mu=1}^4(1-\cos\kappa_\mu).                   \label{eq:newS7V51-hub}
\end{equation}
If a larger cell fibre is present, use the same expression on its
translation-invariant affine sector and add a fixed positive scalar
stiffness on the orthogonal internal cell modes.  The latter is chosen once
on the finite common cell and does not alter
\eqref{eq:newS7V51-s0}--\eqref{eq:newS7V51-l01}.

\begin{definition}[Physical-jet-preserving bulk path]
\label{def:newS7V51-PJP}
For \(u\in[0,1]\), join a type \(\tau\) to the hub at the same \(E\) by
\begin{align}
 s_{\tau,E,u}&=(1-u)s_{\tau,E}+u s_E^{\rm hub},        \label{eq:newS7V51-spath}\\
 \ell_{\tau,E,u}&=(1-u)\ell_{\tau,E}
                    +u\ell^{\rm hub}.                 \label{eq:newS7V51-lpath}
\end{align}
Two endpoint types are joined by their two such branches through the
single hub \(H_E^{\rm hub}\).
\end{definition}

\begin{theorem}[Gapped coframe-dependent hub]
\label{thm:newS7V51-hub-gap}
The paths in \Cref{def:newS7V51-PJP} satisfy JISD with uniform constants
over the finite regular type set and the compact coframe class.  For the
large Wilson parameter already allowed in V38, they have one gap
\(\delta_{\rm jet}>0\).  Along every path,
\begin{equation}
 D_\kappa s_{\tau,E,u}(0)=J_E                    \label{eq:newS7V51-fixed-jet}
\end{equation}
exactly, while the second \(\kappa\)-derivatives of \(s\) and \(\ell\)
are uniformly bounded through two metric variations.
\end{theorem}

\begin{proof}
Both endpoint and hub stiffnesses dominate the common
\(q(\kappa)\) by \Cref{lem:newS7V38-stiffness-comparison}; their convex
combination does as well.  By
\Cref{lem:newS7V51-endpoint-jet}, both Dirac symbols vanish at zero and
have the same derivative \(J_E\).  Compact ellipticity of \(E\) supplies
the JISD lower singular-value bound.  Finiteness of the Laurent stencils,
compactness, and convexity supply the uniform second-remainder bounds.
\Cref{thm:newS7V38-large-r-path} gives the gap.  Equation
\eqref{eq:newS7V51-fixed-jet} follows by differentiating
\eqref{eq:newS7V51-spath}.  All second derivatives are finite sums of
bounded trigonometric coefficients, and the same argument applies after
the two metric differentiations.
\end{proof}

\begin{corollary}[The stratified path tree retains the physical jet]
\label{cor:newS7V51-tree}
At a fixed stratum point, every incident type reaches the same
\(E\)-dependent hub.  Hence the V45--V48 link maps still land in a
contractible rooted tree and their fillings remain available with
\(\delta_{\rm path}\) replaced by \(\delta_{\rm jet}\).  Every point of
the tree has the common invariant first Dirac jet \(J_E\).
\end{corollary}

\begin{proof}
Take one abstract copy of each branch in
\Cref{def:newS7V51-PJP} and identify its hub endpoint.  This is a finite
tree.  Contractibility and the link-filling arguments are unchanged.
The last statement is \eqref{eq:newS7V51-fixed-jet}.
\end{proof}

\subsection{The affine transverse subspace}
\label{sec:newS7V51-transverse}

Treat all strata at once.  Let \(\Sigma\) have codimension
\(d\in\{1,2,3,4\}\), put \(t=4-d\), and freeze its tangential geometry.
The \(d\)-normal collar fibre is a regular set
\(S_{\Sigma,L}\subset\mathbb Z^d\) with
\begin{equation}
 c_1L^d\le N_{\Sigma,L}:=|S_{\Sigma,L}|\le c_2L^d.     \label{eq:newS7V51-N}
\end{equation}
The rescaled covariance matrix of its centred coordinates has eigenvalues
in a fixed compact subset of \((0,\infty)\).  This holds for the interval,
disk, ball, and four-ball fibres used in V44--V48 after changing their
boundary by only \(O(1)\) layers.

Let \(\mathcal A_{\Sigma,L}\) be the scalar span on
\(S_{\Sigma,L}\) of
\[
 1,\quad {n_1-\bar n_1\over L},\ldots,
 {n_d-\bar n_d\over L}.
\]
Choose its orthonormal basis with the normalized constant first, and
tensor it with the spin--internal fibre \(\mathbb C^q\).  Denote the
resulting isometry by
\begin{equation}
 P_{\Sigma,L}:\mathbb C^{(d+1)q}\longrightarrow
 \ell^2(S_{\Sigma,L};\mathbb C^q),                    \label{eq:newS7V51-P}
\end{equation}
and its constant summand by
\(Q_{\Sigma,L}:\mathbb C^q\to\operatorname{Ran}P_{\Sigma,L}\).
Uniform nondegeneracy of the covariance matrix gives
\begin{equation}
 \|(P_{\Sigma,L})_n\|+\|(Q_{\Sigma,L})_n\|
 \le C N_{\Sigma,L}^{-1/2}.                           \label{eq:newS7V51-P-block}
\end{equation}

Let
\[
 B_{\Sigma,L}(k_\parallel)
 :=H_{\Sigma,L}^{\rm ad}(k_\parallel)
   -H_\Sigma^{\rm ref}(k_\parallel)
\]
be the leading adiabatic-minus-reference difference formed with the
jet-preserving tree.  At lower-dimensional boundaries,
\(H_\Sigma^{\rm ref}\) uses the already completed lower-codimension model;
the construction below is therefore read inductively in increasing
codimension.

\begin{lemma}[Affine midpoint defect]
\label{lem:newS7V51-affine-defect}
Uniformly for \(d=1,2,3,4\),
\begin{align}
 \|B_{\Sigma,L}(0)P_{\Sigma,L}\|
 &\le {C\over L},                                     \label{eq:newS7V51-BP}\\
 \|(\partial_{k_\alpha}B_{\Sigma,L})(0)Q_{\Sigma,L}\|
 &\le {C\over L},\qquad 1\le\alpha\le t.              \label{eq:newS7V51-dBQ}
\end{align}
The estimates hold through two metric variations.
\end{lemma}

\begin{proof}
Apply a frozen path stencil at a normal site \(n\) to a constant or
physical affine nodal mode.  By
\Cref{lem:newS7V51-endpoint-jet,thm:newS7V51-hub-gap}, its value is
independent of the path parameter and agrees with every sector reference
stencil.  Thus the pointwise frozen difference has zero constant and
affine moments.

The actual self-adjoint collar places an edge coefficient at its midpoint.
For a displacement \(e\) from the fixed finite stencil, its path parameter
differs from the row parameter by \(O(|e|/L)\).  The path is uniformly
Lipschitz, so the remaining row action on an affine mode is \(O(L^{-1})\)
times its normalized value.  By
\eqref{eq:newS7V51-P-block}, each of the \(O(L^d)\) nonzero row blocks is
therefore \(O(N_{\Sigma,L}^{-1/2}L^{-1})\).  Squaring and summing proves
\eqref{eq:newS7V51-BP}.  Differentiation at zero tangential momentum
inserts one of finitely many tangential displacements.  The invariant
first moment again cancels its frozen part, and the same midpoint estimate
proves \eqref{eq:newS7V51-dBQ}.

For \(d>1\), angular interpolation parameters have bounded derivatives in
the rescaled normal variables after using the sector charts from
V45--V48; the central chart is chosen Lipschitz, so the estimate remains
\(O(L^{-1})\).  The completed boundary models satisfy the same affine
identities by the induction below and contribute only the already bounded
\(O(L^{-1})\) boundary debit.  Differentiating the finite coefficient
identities and the common \(J_E\) gives the two metric-contact statements.
\end{proof}

\subsection{Hermitian completion of the transverse affine modes}
\label{sec:newS7V51-perp-completion}

Set
\begin{equation}
 Y_{\Sigma,L}:=-B_{\Sigma,L}(0)P_{\Sigma,L}            \label{eq:newS7V51-Y}
\end{equation}
and define
\begin{equation}
 R_{\Sigma,L}^{\perp}:=
 Y_{\Sigma,L}P_{\Sigma,L}^*
 +P_{\Sigma,L}Y_{\Sigma,L}^*
 -P_{\Sigma,L}(P_{\Sigma,L}^*Y_{\Sigma,L})
  P_{\Sigma,L}^*.                                     \label{eq:newS7V51-Rperp}
\end{equation}

\begin{theorem}[Transverse affine completion]
\label{thm:newS7V51-transverse-completion}
The operator \(R_{\Sigma,L}^{\perp}\) is self-adjoint, has rank at most
\(2(d+1)q\), and satisfies
\begin{align}
 (B_{\Sigma,L}(0)+R_{\Sigma,L}^{\perp})P_{\Sigma,L}
 &=0,                                                  \label{eq:newS7V51-transverse-zero}\\
 \|R_{\Sigma,L}^{\perp}\|&\le {C\over L}.              \label{eq:newS7V51-Rperp-norm}
\end{align}
\end{theorem}

\begin{proof}
Because \(B_{\Sigma,L}(0)\) is self-adjoint,
\[
 P_{\Sigma,L}^*Y_{\Sigma,L}
 =-P_{\Sigma,L}^*B_{\Sigma,L}(0)P_{\Sigma,L}
\]
is self-adjoint.  Consequently
\eqref{eq:newS7V51-Rperp} is self-adjoint and multiplication on the right
by \(P_{\Sigma,L}\) gives \(R_{\Sigma,L}^{\perp}P_{\Sigma,L}
=Y_{\Sigma,L}\).  This proves
\eqref{eq:newS7V51-transverse-zero}.  Its range lies in
\(\operatorname{Ran}Y_{\Sigma,L}+\operatorname{Ran}P_{\Sigma,L}\), and
\[
 \|R_{\Sigma,L}^{\perp}\|\le3\|Y_{\Sigma,L}\|
 \le {C\over L}
\]
by \Cref{lem:newS7V51-affine-defect}.
\end{proof}

\subsection{Tangential Fourier-jet completion}
\label{sec:newS7V51-tangent-completion}

The transverse correction is independent of \(k_\parallel\).  Put
\begin{equation}
 D_{\Sigma,L,\alpha}
 :=(\partial_{k_\alpha}B_{\Sigma,L})(0),\qquad
 Y_{\Sigma,L,\alpha}:=-D_{\Sigma,L,\alpha}Q_{\Sigma,L}.              \label{eq:newS7V51-Yalpha}
\end{equation}
Since \(B_{\Sigma,L}(k_\parallel)\) is self-adjoint for real
\(k_\parallel\), every \(D_{\Sigma,L,\alpha}\) is self-adjoint.  Define
\begin{equation}
 K_{\Sigma,L,\alpha}:=
 Y_{\Sigma,L,\alpha}Q_{\Sigma,L}^*
 +Q_{\Sigma,L}Y_{\Sigma,L,\alpha}^*
 -Q_{\Sigma,L}
  (Q_{\Sigma,L}^*Y_{\Sigma,L,\alpha})
  Q_{\Sigma,L}^*.                                     \label{eq:newS7V51-Kalpha}
\end{equation}
Finally set
\begin{equation}
 R_{\Sigma,L}^{\rm aff}(k_\parallel)
 :=R_{\Sigma,L}^{\perp}
   +\sum_{\alpha=1}^{t}\sin(k_\alpha)
      K_{\Sigma,L,\alpha}.                            \label{eq:newS7V51-Raff}
\end{equation}

\begin{theorem}[Full physical affine-jet completion]
\label{thm:newS7V51-full-completion}
For real \(k_\parallel\),
\(R_{\Sigma,L}^{\rm aff}(k_\parallel)\) is self-adjoint and
\begin{align}
 (B_{\Sigma,L}(0)+R_{\Sigma,L}^{\rm aff}(0))
 P_{\Sigma,L}&=0,                                     \label{eq:newS7V51-full-zero}\\
 \partial_{k_\alpha}
 \{B_{\Sigma,L}+R_{\Sigma,L}^{\rm aff}\}(0)
 Q_{\Sigma,L}&=0,\quad 1\le\alpha\le t,               \label{eq:newS7V51-full-first}\\
 \sup_{k_\parallel}\|R_{\Sigma,L}^{\rm aff}(k_\parallel)\|
 &\le {C\over L}.                                     \label{eq:newS7V51-Raff-norm}
\end{align}
Its rank in each tangential fibre is at most
\[
 2\{(d+1)+t\}q=10q.
\]
Equations \eqref{eq:newS7V51-full-zero} and
\eqref{eq:newS7V51-full-first} say exactly that the completed difference
annihilates all four physical affine nodal modes.
\end{theorem}

\begin{proof}
As in the preceding proof,
\(Q_{\Sigma,L}^*Y_{\Sigma,L,\alpha}
=-Q_{\Sigma,L}^*D_{\Sigma,L,\alpha}Q_{\Sigma,L}\)
is self-adjoint.  Hence \(K_{\Sigma,L,\alpha}\) is self-adjoint,
\[
 K_{\Sigma,L,\alpha}Q_{\Sigma,L}
 =Y_{\Sigma,L,\alpha},
\qquad
 \|K_{\Sigma,L,\alpha}\|
 \le3\|Y_{\Sigma,L,\alpha}\|\le {C\over L}.            \label{eq:newS7V51-K-bound}
\]
The real scalar \(\sin k_\alpha\) preserves self-adjointness.
At \(k_\parallel=0\), all sine terms vanish, so
\eqref{eq:newS7V51-full-zero} is
\eqref{eq:newS7V51-transverse-zero}.  Differentiating
\eqref{eq:newS7V51-Raff}, using
\(\partial_{k_\alpha}\sin k_\beta|_0=\delta_{\alpha\beta}\), and applying
\eqref{eq:newS7V51-K-bound} gives
\[
 (D_{\Sigma,L,\alpha}+K_{\Sigma,L,\alpha})
 Q_{\Sigma,L}=0,
\]
which is \eqref{eq:newS7V51-full-first}.  The norm bound follows from the
triangle inequality and \(t\le3\).  Each Hermitian completion has range in
the sum of two spaces of the stated dimensions; summing the rank bounds
gives \(2(d+1)q+2tq=10q\).

A normal affine mode lies in \(\operatorname{Ran}P_{\Sigma,L}\).
A tangential affine mode is obtained by differentiating a Bloch constant
mode at \(k_\parallel=0\).  Thus
\eqref{eq:newS7V51-full-zero}--\eqref{eq:newS7V51-full-first} exhaust the
constant and four affine physical modes.
\end{proof}

\paragraph{Real-space interpretation.}
The Fourier coefficient of
\(\sin k_\alpha K_{\Sigma,L,\alpha}\) at the positive tangential shift is
\((2i)^{-1}K_{\Sigma,L,\alpha}\), and that at the negative shift is its
adjoint.  Thus \eqref{eq:newS7V51-Raff} is a genuine self-adjoint
real-space operator of one-step tangential range.  Under reversal of the
tangential coordinate, both \(\sin k_\alpha\) and the derivative defining
\(K_{\Sigma,L,\alpha}\) change sign, so the construction is independent
of the chosen orientation.

\subsection{Row moments, covariance, and the gap}
\label{sec:newS7V51-control}

\begin{lemma}[Affine completion row ledger]
\label{lem:newS7V51-row-ledger}
For the real-space kernel of \(R_{\Sigma,L}^{\rm aff}\),
\begin{align}
 \sup_x\sum_y\|R_{\Sigma,L}^{\rm aff}(x,y)\|
 &\le {C\over L},                                     \label{eq:newS7V51-row0}\\
 \sup_x\sum_y(1+|x-y|)
       \|R_{\Sigma,L}^{\rm aff}(x,y)\|
 &\le C,                                               \label{eq:newS7V51-row1}\\
 \sup_x\sum_y e^{\gamma|x-y|/L}
       \|R_{\Sigma,L}^{\rm aff}(x,y)\|
 &\le {C_\gamma\over L}                               \label{eq:newS7V51-rowexp}
\end{align}
for every fixed sufficiently small \(\gamma>0\).  The bounds hold through
two metric variations.
\end{lemma}

\begin{proof}
The proof of \Cref{lem:newS7V51-affine-defect} gives the block estimates
\[
 \|(P_{\Sigma,L})_x\|+\|(Q_{\Sigma,L})_x\|
 \le CN_{\Sigma,L}^{-1/2},\qquad
 \|Y_x\|+\sum_\alpha\|Y_{\alpha,x}\|
 \le {C\over L N_{\Sigma,L}^{1/2}}.
\]
Every block in
\eqref{eq:newS7V51-Rperp} and
\eqref{eq:newS7V51-Kalpha} is therefore
\(O((LN_{\Sigma,L})^{-1})\).  Summing \(N_{\Sigma,L}\)
blocks proves \eqref{eq:newS7V51-row0}.  Their normal distance is
\(O(L)\), so inserting one distance proves
\eqref{eq:newS7V51-row1}; the tangential sine terms add only a unit
displacement.  The diameter of the transverse fibre is \(O(L)\), so the
exponential is uniformly bounded by a constant depending on \(\gamma\),
which proves \eqref{eq:newS7V51-rowexp}.  Metric derivatives preserve all
displayed block estimates.
\end{proof}

\begin{lemma}[Covariant curved lift]
\label{lem:newS7V51-covariance}
Define \(P_{\Sigma,L}\) and \(Q_{\Sigma,L}\) after identifying the
transverse fibres with a central fibre by the V33 short-geodesic
spin--gauge transports.  Then
\(R_{\Sigma,L}^{\rm aff}\) is spin--gauge covariant, self-adjoint, and
gamma-five Hermitian in the associated Wilson operator.  Its frozen
constant and affine difference moments vanish exactly.  Smooth curvature
leaves an \(O(a)\) transport tolerance, through two metric contacts.
\end{lemma}

\begin{proof}
The transported \(P\) and \(Q\) intertwine central and endpoint frame
actions.  Therefore \(B,Y,Y_\alpha,R^\perp,K_\alpha\), and \(R^{\rm aff}\)
all acquire exactly the endpoint conjugations.  Unitary transport
preserves self-adjointness.  If \(H\) is self-adjoint and
\(D=\Gamma H\), then \(D^*=\Gamma D\Gamma\), giving gamma-five
Hermiticity.  The frozen moment statement is
\Cref{thm:newS7V51-full-completion}.  On a curved mesh, neighbouring
transports and coefficients differ from their frozen values by \(O(a)\)
over the fixed tangential shifts; the transverse transports are already
included in \(P\) and \(Q\).  Finite smooth dependence on the metric gives
the same tolerance after two variations.
\end{proof}

Define the new hierarchical stratified kernel by replacing, at every
stratum, the V44--V48 constant completion with
\eqref{eq:newS7V51-Raff}, in increasing codimension.

\begin{theorem}[Uniform gap after affine completion]
\label{thm:newS7V51-gap}
There are constants \(C,L_0\), uniform over every stratum type and the
compact geometric family, such that
\begin{equation}
 \|H_{\Sigma,L}^{\rm aff}(k_\parallel)u\|
 \ge\{\delta_{\rm jet}-CL^{-1/2}-CL^{-1}\}\|u\|.       \label{eq:newS7V51-gap}
\end{equation}
For \(L\ge L_0\), all local tangent models have gap at least
\(\delta_{\rm jet}/2\).  The global IMS argument of V49 then gives
\begin{equation}
 \|H_a^{\rm aff}u\|
 \ge\left\{c_{\rm aff}-{C\over R}-CaR\right\}\|u\|,
 \qquad L\ll R\ll a^{-1},                             \label{eq:newS7V51-global-gap}
\end{equation}
with \(c_{\rm aff}>0\).  Hence the completed global dimensionless kernel
is uniformly gapped for any cofinal scales satisfying these inequalities.
\end{theorem}

\begin{proof}
The jet-preserving tree has gap \(\delta_{\rm jet}\) by
\Cref{thm:newS7V51-hub-gap}.  The V37 multi-normal freezing estimates used
in V45--V48 lose \(CL^{-1/2}\).  The new completion loses at most
\(CL^{-1}\) by \eqref{eq:newS7V51-Raff-norm} and the reverse triangle
inequality.  Finiteness of all stratum types makes the constants uniform.
For the global statement, the IMS commutator of the growing-range part is
\(O(R^{-1})\) by \eqref{eq:newS7V51-row1}; freezing smooth geometric data
on an \(R\)-cell is \(O(aR)\).  This is exactly the V49 argument with the
new local gap.
\end{proof}

\begin{corollary}[Frozen HFJC affine rows acquired]
\label{cor:newS7V51-HFJC}
Items one and the frozen part of item two in HFJC hold for the modified
stratified kernel:
\begin{enumerate}
\item the gapped interpolation keeps the invariant physical Dirac first
jet exactly and has uniformly bounded Wilson second jet; and
\item the covariant Hermitian completion matches constants and all four
physical affine modes exactly, with norm \(O(L^{-1})\), rank at most
\(10q\) per stratum fibre, and physical range \(O(aL)\).
\end{enumerate}
\end{corollary}

\begin{proof}
Use \Cref{thm:newS7V51-hub-gap,thm:newS7V51-full-completion,
lem:newS7V51-row-ledger,lem:newS7V51-covariance}.
\end{proof}

\subsection{The remaining heat-flat derivative problem}
\label{sec:newS7V51-remaining}

The principal Dirac jet no longer makes an order-one transition across a
collar.  This removes the counterexample mechanism in
\Cref{prop:newS7V50-core-not-heat} for that jet.  It does not yet prove
local heat convergence.  The Wilson second jet is bounded but may vary by
order one across physical width \(w=aL\).  Because the Wilson term carries
one explicit power of \(a\) in the continuum expansion, its first spatial
derivative is expected at scale \(a/w=L^{-1}\), whereas successive
derivatives cost further powers of \(w^{-1}\).  The exact heat words must
be listed before choosing \(L(a)\); otherwise a scale choice could hide a
divergent differentiated irrelevant coefficient.

\begin{assessment}[Affine heat-flatness acquired; CHSC remains]
\label{ass:newS7V51-status}
V51 removes the physical first-Dirac-jet excursion from every stratum
path and constructs one explicit self-adjoint correction that annihilates
all frozen affine modes.  The correction remains \(O(L^{-1})\), so the
local and global gaps survive.  The remaining HFJC row is the
spatial-derivative ledger for the Wilson and curvature residuals.  CHSC
and mesoscopic heat convergence are not inferred from the affine
identities alone.
\end{assessment}

\begin{programme}[V52: differentiated irrelevant-jet ledger]
\label{prog:newS7V51-V52}
V52 will square the locally completed Wilson sign kernel, identify every
coefficient containing a spatial derivative of the collar profile through
the order that can enter the four-dimensional \(a_4\) density, and attach
its exact power of \(a\), \(L\), and \(w=aL\).  It will then choose, or
prove impossible, one scale \(L=a^{-\alpha}\) that makes all such
coefficients vanish while retaining the V51 gap and IMS inequalities.
Only after that ledger is closed will CHSC be attacked.
\end{programme}

\begin{firewall}[Claim boundary after seventh-series V51]
\label{fire:newS7V51-boundary}
V51 constructs a coframe-dependent universal hub, proves its uniform gap
and exact invariant first Dirac jet, extends the conservative completion
to the constant and transverse-affine subspace, adds explicit tangential
Fourier-jet corrections, and proves self-adjointness, covariance, rank,
row-moment, weighted-locality, and gap bounds.  It closes the frozen
constant and affine rows of HFJC.

V51 does not prove the differentiated irrelevant-jet row of HFJC, CHSC,
MATC for the completed kernel, convergence of its \(a_4\) coefficients,
the common Einstein stress, pure gravitational owners, rough metric or
AOP control, measure concentration, the determinant-line phase,
interacting cutoff removal, or the Standard Model--Einstein continuum.
\end{firewall}

\begin{theorem}[Seventh-series V51 result]
\label{thm:newS7V51-result}
The regular V33 types meeting at a physical point admit a uniformly gapped
common-hub interpolation whose invariant physical Dirac first jet is
constant along the path.  At every stratum of codimension one through
four, the resulting \(O(L^{-1})\) affine defect admits the explicit
self-adjoint correction \eqref{eq:newS7V51-Raff}.  The corrected difference
annihilates constants and all physical affine modes, has fibre rank at
most \(10q\), exponential Schur norm \(O(L^{-1})\) on range \(O(L)\), and
preserves both the local tangent gaps and the global IMS gap.
\end{theorem}

\begin{proof}
Combine \Cref{lem:newS7V51-endpoint-jet,
thm:newS7V51-hub-gap,
cor:newS7V51-tree,
lem:newS7V51-affine-defect,
thm:newS7V51-transverse-completion,
thm:newS7V51-full-completion,
lem:newS7V51-row-ledger,
lem:newS7V51-covariance,
thm:newS7V51-gap,
cor:newS7V51-HFJC}.
\end{proof}

\paragraph{Parent-version record.}
V51 was formed from
\texttt{Renewal\_SM\_Einstein\_Continuum\_Research\_Notes\_V50.tex}.
The V50 parent had \(19\,902\) logical source lines, \(799\,888\) bytes,
and SHA--256
\[
 \texttt{5931C1A01125C21735AC326608EF856F01124AC9D2439155A5C4CD04871495F2}.
\]
The next compulsory companion audit is V55.

% =============================================================================
% END APPENDED SEVENTH-SERIES V51 MATERIAL
% =============================================================================

% =============================================================================
% BEGIN APPENDED SEVENTH-SERIES V52 MATERIAL
% =============================================================================

\section{Seventh-series V52: local jet-balanced quantization and the heat scale}
\label{sec:v7v52}
\label{sec:newS7V52}

\subsection{Going upstream from the growing-range completion}
\label{sec:newS7V52-upstream}

V51 proves that an affine-moment completion can be made self-adjoint with
norm \(O(L^{-1})\).  Its transverse range is nevertheless \(O(L)\).
Although its zeroth and first moments vanish, its second lattice moment can
be \(O(L)\); after physical rescaling this is an \(O(aL)\) second-order
coefficient varying across a collar of the same width \(aL\).  Repeated
spatial derivatives of that artificial coefficient are not controlled by
operator norm.  This is precisely the kind of ultraviolet multiplicity
identified in V50.

The fixed physical first jet from V51 permits a better construction.
Midpoint quantization of a symbol whose zeroth and first momentum jets are
independent of position has an \(O(L^{-2})\), rather than
\(O(L^{-1})\), constant-mode defect.  A scale with \(L^{-2}=o(a)\)
then removes the need for any nonlocal conservative completion.

\begin{correction}[The completion is auxiliary, not part of the final
candidate]
\label{corr:newS7V52-no-completion}
The operators \(R_{\Sigma,L}^{\rm aff}\) constructed in V51 and all of
their stated algebraic estimates remain valid.  Beginning with this note,
they are used only as a comparison proving that affine conservation is not
algebraically obstructed.  The SM--Einstein continuum candidate is instead
the fixed-range jet-balanced midpoint kernel defined below.
\end{correction}

\subsection{Smooth fixed-jet stratum symbols}
\label{sec:newS7V52-smooth-symbol}

Let \(\mathscr G_E^{(1)}\) be the affine space of finite Laurent symbols
whose mass, value at zero momentum, and physical first Dirac jet equal
those of the V33 kernel at coframe \(E\), and whose stiffness value and
first jet vanish at zero.  Every point of the V51 path tree lies in
\(\mathscr G_E^{(1)}\).

\begin{lemma}[Smooth fixed-jet filling]
\label{lem:newS7V52-smooth-filling}
For every stratum link used in V45--V48, its map into the V51 path tree
extends over the transverse unit ball to a \(C^\infty\) map
\begin{equation}
 z\longmapsto h_E(z,\kappa)\in\mathscr G_E^{(1)}       \label{eq:newS7V52-h}
\end{equation}
with the prescribed sector values near the boundary.  All \(z\)-derivatives
are bounded uniformly over the finite stratum types and through two metric
variations.  The corresponding frozen symbols have gap at least
\(3\delta_{\rm jet}/4\).
\end{lemma}

\begin{proof}
The tree filling of
\Cref{cor:newS7V51-tree} is Lipschitz and uniformly gapped by
\(\delta_{\rm jet}\).  Smooth it away from the prescribed boundary in a
finite-dimensional coefficient space.  Near a tree vertex, first retract
the map into a coefficient ball of radius
\(\delta_{\rm jet}/8\) about the common hub and smooth there by ordinary
convolution; elsewhere smooth separately along a single branch.  A
partition of unity is used only inside such hub balls.  The smoothed
symbol differs in operator norm by less than
\(\delta_{\rm jet}/4\), so the reverse triangle inequality preserves the
stated gap.

The constraints defining \(\mathscr G_E^{(1)}\) are affine linear.
Convolution and convex averaging therefore preserve them exactly.
There are finitely many coefficient charts and stratum types, and the
coframe class is compact, so the smoothing radii and all required
derivative bounds may be chosen uniformly.  The construction is smooth in
the coframe parameters, giving the differentiated assertions.
\end{proof}

Write the finite Laurent expansion as
\begin{equation}
 h_E(z,\kappa)=\sum_{e\in\mathcal D}h_{E,e}(z)e^{i\kappa\cdot e},
 \qquad
 h_{E,-e}(z)=h_{E,e}(z)^*,                            \label{eq:newS7V52-Laurent}
\end{equation}
after enlarging one fixed displacement set \(\mathcal D\) to contain the
finite family.  Cell-fibre indices are included in the coefficient
matrices.

\begin{definition}[Jet-balanced midpoint kernel, JBMK]
\label{def:newS7V52-JBMK}
In a \(d\)-normal stratum chart, define
\begin{equation}
 (H_{\Sigma,L}^{\rm jb}u)_n
 :=\sum_{e\in\mathcal D}
 h_{E(a n),e}\left({n_\perp+e_\perp/2\over L}\right)
 U_{n,n+e}u_{n+e},                                    \label{eq:newS7V52-JBMK}
\end{equation}
where \(U_{n,n+e}\) is the V33 short-geodesic spin--gauge transport.
The diagonal mass block is included in \(e=0\).  On overlaps, use the
hierarchical smooth fixed-jet filling of
\Cref{lem:newS7V52-smooth-filling}; outside the collars use the sharp V33
kernel.
\end{definition}

\begin{lemma}[Exact structure and fixed range]
\label{lem:newS7V52-structure}
The JBMK is self-adjoint and spin--gauge covariant, has the same fixed
lattice range \(\mathcal D\) for every \(a,L\), and gives a gamma-five
Hermitian Wilson operator.  These properties hold through two metric
contacts.
\end{lemma}

\begin{proof}
The coefficient at the reversed edge is the adjoint coefficient evaluated
at the same midpoint by \eqref{eq:newS7V52-Laurent}; the endpoint
transports are mutual adjoints.  This proves self-adjointness.
Frame changes conjugate the endpoint transports and coefficient blocks,
which proves covariance.  The displacement set is fixed by definition.
As in \Cref{lem:newS7V51-covariance}, a self-adjoint sign kernel
\(H\) gives a gamma-five Hermitian Wilson operator \(D=\Gamma H\).
Metric differentiation preserves these algebraic identities.
\end{proof}

\subsection{The cancelled midpoint term}
\label{sec:newS7V52-midpoint}

The following elementary estimate is the reason for changing direction.
For clarity, first suppress the slowly varying physical coframe and
transports and work in one frozen chart.

\begin{lemma}[Second-order constant defect]
\label{lem:newS7V52-constant-defect}
Let \(h_e(z)\) be \(C^2\), finite range, and suppose
\begin{equation}
 \sum_e h_e(z)=M,\qquad
 \sum_e e_\mu h_e(z)=J_\mu                         \label{eq:newS7V52-fixed-moments}
\end{equation}
are independent of \(z\).  For the midpoint operator
\[
 (H_Lu)_n=\sum_e h_e((n+e/2)/L)u_{n+e},
\]
one has
\begin{equation}
 \left\|\sum_eh_e((n+e/2)/L)-M\right\|
 \le {C\over L^2}.                                    \label{eq:newS7V52-row-defect}
\end{equation}
If \(f(n)=v_0+\sum_\mu n_\mu v_\mu\) is affine and
\(|n|\le CL\), then
\begin{equation}
 \|(H_L-H_{z_n}^{\rm fr})f(n)\|
 \le {C\over L}\left(\|v_0\|+\sum_\mu L\|v_\mu\|\right),              \label{eq:newS7V52-affine-defect}
\end{equation}
where \(H_{z_n}^{\rm fr}\) freezes every coefficient at \(z_n=n/L\).
\end{lemma}

\begin{proof}
Taylor's formula with integral remainder gives
\[
 h_e(z_n+e/2L)
 =h_e(z_n)+{e_j\over2L}\partial_jh_e(z_n)
 +O(L^{-2}|e|^2).
\]
After summing in \(e\), the zeroth term is \(M\), while the coefficient of
\(L^{-1}\) is
\[
 {1\over2L}\partial_j\sum_e e_jh_e(z_n)=0
\]
by \eqref{eq:newS7V52-fixed-moments}.  Finite range and the uniform
second-derivative bound prove
\eqref{eq:newS7V52-row-defect}.

For the affine mode,
\[
 f(n+e)=f(n)+e_\mu v_\mu.
\]
The part multiplying \(f(n)\) is bounded by
\eqref{eq:newS7V52-row-defect}.  In the remaining part, freezing gives
\(\sum_e e_\mu h_e(z_n)v_\mu\).  The difference between midpoint and
frozen first moments is \(O(L^{-1})\sum_\mu\|v_\mu\|\).
Since \(|f(n)|\le C(\|v_0\|+\sum_\mu L\|v_\mu\|)\) in the collar, the
displayed bound follows.
\end{proof}

\begin{corollary}[Physical consistency rate]
\label{cor:newS7V52-physical-rate}
For a fixed smooth spinor \(\psi\), the physical kernel
\(a^{-1}H_{\Sigma,L}^{\rm jb}\) differs from its frozen fixed-jet action
by
\begin{equation}
 O\left({1\over aL^2}+{1\over L}+a\right)              \label{eq:newS7V52-consistency-rate}
\end{equation}
locally, through two metric contacts.  Thus this error tends to zero if
\begin{equation}
 L\longrightarrow\infty,\qquad aL\longrightarrow0,
 \qquad aL^2\longrightarrow\infty.                    \label{eq:newS7V52-three-scale}
\end{equation}
\end{corollary}

\begin{proof}
The constant Taylor term gives \(a^{-1}L^{-2}\) by
\Cref{lem:newS7V52-constant-defect}.  The physical affine increment has
coefficient \(a\nabla\psi\), so
\eqref{eq:newS7V52-affine-defect}, after multiplication by \(a^{-1}\),
gives \(L^{-1}\).  The quadratic Taylor remainder of \(\psi\) gives
\(O(a)\) for a fixed stencil.  Smooth coframe and parallel-transport
variation is the ordinary V33 covariant consistency term and obeys the
same bound.  The three limits are precisely those making all terms vanish.
\end{proof}

\subsection{Gap without a conservative completion}
\label{sec:newS7V52-gap}

\begin{theorem}[Local and global JBMK gaps]
\label{thm:newS7V52-gap}
For every stratum tangent,
\begin{equation}
 \|H_{\Sigma,L}^{\rm jb}u\|
 \ge\{3\delta_{\rm jet}/4-CL^{-1/2}\}\|u\|.            \label{eq:newS7V52-local-gap}
\end{equation}
Consequently the global hierarchical JBMK satisfies
\begin{equation}
 \|H_a^{\rm jb}u\|
 \ge\left\{c_{\rm jb}-{C\over R}-CaR\right\}\|u\|,
 \qquad L\ll R\ll a^{-1},                             \label{eq:newS7V52-global-gap}
\end{equation}
for one \(c_{\rm jb}>0\) and all sufficiently small \(a\).
\end{theorem}

\begin{proof}
The frozen symbol gap is
\Cref{lem:newS7V52-smooth-filling}.  The coefficients of the midpoint
operator change by \(O(L^{-1})\) per lattice step and have fixed range, so
the localization/freezing proof of
\Cref{thm:newS7V37-adiabatic} gives the
\(CL^{-1/2}\) debit.  No completion debit occurs.

Use an IMS partition on lattice scale \(R\).  Fixed range makes the
commutator \(O(R^{-1})\), and freezing the smooth geometric data costs
\(O(aR)\).  Applying the local tangent estimate in each IMS cell and
summing proves \eqref{eq:newS7V52-global-gap}.
\end{proof}

\begin{corollary}[Fixed-range sign locality]
\label{cor:newS7V52-sign-locality}
The resolvent and sign of \(H_a^{\rm jb}\) have exponential lattice
off-diagonal decay with an \(a\)-independent decay exponent.  Their first
two metric variations obey the corresponding differentiated exponential
Schur bounds.
\end{corollary}

\begin{proof}
The hopping range and row norm are uniformly bounded by
\Cref{lem:newS7V52-structure}, and the gap is uniform by
\Cref{thm:newS7V52-gap}.  The standard conjugated-resolvent argument used
in V50 now permits a fixed exponential weight, rather than the
\(\gamma/L\) weight required by the completion.  Insert the sign
resolvent formula.  Two metric derivatives insert finitely many uniformly
bounded covariant coefficient derivatives between such resolvents, so
the same convolution estimate applies.
\end{proof}

\subsection{Differentiated irrelevant-jet ledger}
\label{sec:newS7V52-jet-ledger}

Put \(w=aL\).  On low physical momenta, fixed-jet equality means that every
collar-dependent term beyond the continuum Dirac operator carries a
positive power of \(a\).  More explicitly, analytic functional calculus
for the uniformly gapped fixed-range kernel gives, on smooth cores,
\begin{equation}
 \mathscr D_a^{\rm jb}
 =\mathscr D_E+\sum_{m\ge1}a^m
   P_{m+1}\bigl(z=x_\perp/w,\nabla\bigr),              \label{eq:newS7V52-low-expansion}
\end{equation}
as an asymptotic covariant differential expansion.  The operator
\(P_{m+1}\) has order at most \(m+1\); its coefficient functions have
uniformly bounded \(z\)-derivatives.  The Wilson second jet begins at
\(m=1\).

\begin{lemma}[Profile-derivative power counting]
\label{lem:newS7V52-power-counting}
A term from order \(a^mP_{m+1}\) containing \(j\) physical derivatives of
a collar coefficient is bounded by
\begin{equation}
 C_{m,j}a^mw^{-j}.                                    \label{eq:newS7V52-power}
\end{equation}
In the local four-dimensional heat coefficient of weight four, the
order-\(a^m\) contribution differentiates its collar coefficients at most
\(m+4\) times.  Hence all collar-dependent local \(a_4\)-jet terms are
bounded by finite sums of
\begin{equation}
 C_m a^m w^{-(m+4)}.                                  \label{eq:newS7V52-worst}
\end{equation}
The same bounds hold through two metric contacts.
\end{lemma}

\begin{proof}
Every physical normal derivative of a coefficient
\(F(x_\perp/w)\) contributes \(w^{-1}\), proving
\eqref{eq:newS7V52-power}.  For the second statement, use the local
resolvent parametrix recursion for the squared Dirac operator.  The
coefficient of heat weight four is obtained after four units of covariant
symbol differentiation beyond the order of the inserted perturbation.
An insertion \(a^mP_{m+1}\) has differential excess \(m\) relative to the
first-order Dirac operator, so no monomial at that order contains more
than \(m+4\) coefficient derivatives.  This is also the usual dimensional
weight count: \(a^m\) has length weight \(m\), and a local density of heat
weight four can compensate it with at most \(m+4\) derivatives.
Applying \eqref{eq:newS7V52-power} with \(j\le m+4\) gives
\eqref{eq:newS7V52-worst}.  Metric contacts differentiate only the bounded
coefficient functions and transports, not the scale \(w\), so they do not
change the powers.
\end{proof}

\begin{definition}[Heat-compatible collar scales]
\label{def:newS7V52-scales}
For the remainder of the local heat programme, set
\begin{equation}
 L(a):=a^{-7/8},\qquad
 w(a):=aL=a^{1/8},\qquad
 R(a):=a^{-15/16}.                                    \label{eq:newS7V52-scales}
\end{equation}
\end{definition}

\begin{theorem}[Closure of the differentiated local jet ledger]
\label{thm:newS7V52-jet-closure}
The scales \eqref{eq:newS7V52-scales} satisfy every JBMK consistency and
gap inequality.  Moreover, for every \(m\ge1\),
\begin{equation}
 a^mw^{-(m+4)}
 =a^{(7m-4)/8}\longrightarrow0.                       \label{eq:newS7V52-jet-vanish}
\end{equation}
The slowest local heat-jet rate is \(a^{3/8}\), at \(m=1\).
\end{theorem}

\begin{proof}
Directly,
\[
 L^{-1/2}=a^{7/16},\quad
 {1\over aL^2}=a^{3/4},\quad
 L^{-1}=a^{7/8},\quad
 R^{-1}=a^{15/16},\quad
 aR=a^{1/16}.
\]
All tend to zero, while
\[
 L/R=a^{1/16}\to0,\qquad aR=a^{1/16}\to0.
\]
Thus
\eqref{eq:newS7V52-three-scale},
\eqref{eq:newS7V52-local-gap}, and
\eqref{eq:newS7V52-global-gap} hold.  Since \(w=a^{1/8}\),
\[
 a^mw^{-(m+4)}
 =a^{m-(m+4)/8}=a^{(7m-4)/8}.
\]
This exponent is increasing in \(m\) and equals \(3/8\) at \(m=1\).
\end{proof}

\begin{corollary}[HFJC local rows acquired for the JBMK]
\label{cor:newS7V52-HFJC}
For the fixed-range JBMK, the first three rows of HFJC hold:
\begin{enumerate}
\item the physical Dirac first jet is independent of the collar profile;
\item the midpoint constant defect is \(o(a)\), and the physical affine
defect tends to zero without a growing-range completion; and
\item every differentiated collar coefficient capable of entering the
formal local \(a_4\) density tends to its sharp value through two metric
contacts.
\end{enumerate}
\end{corollary}

\begin{proof}
The first row is \Cref{thm:newS7V51-hub-gap}.  The second is
\Cref{cor:newS7V52-physical-rate,thm:newS7V52-jet-closure}.  The third is
\Cref{lem:newS7V52-power-counting,thm:newS7V52-jet-closure}.
\end{proof}

\subsection{What is still missing from the actual heat trace}
\label{sec:newS7V52-remaining}

The preceding result is a local jet statement.  It does not justify
interchanging the lattice limit, the small-time heat expansion, and the
metric variations.  In particular, the asymptotic series
\eqref{eq:newS7V52-low-expansion} is controlled on smooth low-momentum
cores, whereas the heat trace at time \(t\downarrow0\) samples momenta of
size \(t^{-1/2}\).  A mesoscopic interval
\[
 a^2\ll\varepsilon(a)\le t\le t_0
\]
must be chosen, and the high-momentum tail must be bounded uniformly
before the \(a_4\) coefficient can be identified with the continuum one.
The fixed-range exponential locality in
\Cref{cor:newS7V52-sign-locality} removes the growing-range obstruction
from V50 but does not by itself provide that uniform heat remainder.

\begin{assessment}[The local heat jets close; the mesoscopic estimate is
next]
\label{ass:newS7V52-status}
V52 replaces the long-range completion by a fixed-range midpoint
quantization.  The fixed first moment cancels its nominal \(L^{-1}\)
constant defect, and the declared scales make the remaining defect
continuum-negligible.  The global gap, fixed exponential sign locality,
and all local \(a_4\)-jet powers are now compatible.  The remaining heat
bottleneck is a uniform mesoscopic lattice parametrix and its integrable
remainder, namely the fixed-range version of CHSC.
\end{assessment}

\begin{programme}[V53: mesoscopic heat parametrix]
\label{prog:newS7V52-V53}
V53 will use the fixed exponential sign locality and the scale
\eqref{eq:newS7V52-scales} to split the heat trace into momenta below and
above \(a/\sqrt t\).  It will construct the covariant low-momentum
parametrix through heat weight four, estimate the Brillouin-zone tail, and
determine an explicit \(\varepsilon(a)=a^\beta\).  The target is the CHSC
bound with two metric contacts and a remainder integrable in
\(dt/t\); failure of any exponent will be recorded rather than hidden.
\end{programme}

\begin{firewall}[Claim boundary after seventh-series V52]
\label{fire:newS7V52-boundary}
V52 constructs the smooth fixed-jet stratum filling and its self-adjoint
fixed-range midpoint quantization, proves the \(O(L^{-2})\) constant
defect, physical consistency, local and global gaps, fixed exponential
sign locality, and the complete formal differentiated \(a_4\)-jet power
ledger for \(L=a^{-7/8}\).  It removes the V51 growing-range completion
from the continuum candidate.

V52 does not prove the uniform mesoscopic heat parametrix, CHSC, MATC for
the JBMK, convergence of the actual lattice \(a_4\) coefficient, the
common Einstein stress, pure gravitational owners, rough metric or AOP
control, measure concentration, the determinant-line phase, interacting
cutoff removal, or the Standard Model--Einstein continuum.
\end{firewall}

\begin{theorem}[Seventh-series V52 result]
\label{thm:newS7V52-result}
The V51 fixed-jet tree admits a self-adjoint covariant fixed-range midpoint
quantization with uniform local and global gaps.  Its constant defect is
\(O(L^{-2})\) because the \(O(L^{-1})\) midpoint term is the derivative of
the fixed physical first moment and vanishes.  With
\(L=a^{-7/8}\) and \(R=a^{-15/16}\), its physical consistency errors,
IMS errors, and every collar-dependent local heat-weight-four jet through
two metric contacts tend to zero.  No growing-range conservative
completion is required.
\end{theorem}

\begin{proof}
Combine \Cref{lem:newS7V52-smooth-filling,
lem:newS7V52-structure,
lem:newS7V52-constant-defect,
cor:newS7V52-physical-rate,
thm:newS7V52-gap,
cor:newS7V52-sign-locality,
lem:newS7V52-power-counting,
thm:newS7V52-jet-closure,
cor:newS7V52-HFJC}.
\end{proof}

\paragraph{Parent-version record.}
V52 was formed from
\texttt{Renewal\_SM\_Einstein\_Continuum\_Research\_Notes\_V51.tex}.
The V51 parent had \(20\,470\) logical source lines, \(824\,281\) bytes,
and SHA--256
\[
 \texttt{699C607E81D2A1837079D255E56E5D7C0F76E4BEE93975197D5BCFC792A01D05}.
\]
The next compulsory companion audit is V55.

% =============================================================================
% END APPENDED SEVENTH-SERIES V52 MATERIAL
% =============================================================================

% =============================================================================
% BEGIN APPENDED SEVENTH-SERIES V53 MATERIAL
% =============================================================================

\section{Seventh-series V53: the stratified mesoscopic heat transfer}
\label{sec:v7v53}
\label{sec:newS7V53}

\subsection{Uniform overlap ellipticity on the fixed-jet family}
\label{sec:newS7V53-ellipticity}

V52 supplies a fixed-range, uniformly gapped Wilson sign kernel.  Heat
transfer also requires the massless overlap square to have exactly one
low-energy branch.  This is stronger than invertibility of the Wilson
kernel and must be recorded separately.

For a frozen path symbol write
\[
 m(z,\kappa):={r\over2}\ell(z,\kappa)-\rho,\qquad
 \omega(z,\kappa):=\{\,|s(z,\kappa)|^2+m(z,\kappa)^2\,\}^{1/2}.
\]
On the translation-invariant physical cell sector, the dimensionless
overlap square has scalar eigenvalue
\begin{equation}
 \lambda_{\rm ov}(z,\kappa)
 =2\rho^2\left(1+{m(z,\kappa)\over\omega(z,\kappa)}\right).           \label{eq:newS7V53-overlap-eigenvalue}
\end{equation}
The additional common-cell modes in
\Cref{lem:newS7V52-smooth-filling} are uniformly massive.

\begin{lemma}[Uniform unique overlap cone]
\label{lem:newS7V53-overlap-cone}
There are constants \(c_0,C_0,\kappa_0>0\), uniform in the stratum,
coframe, path, and two metric variations, such that
\begin{align}
 c_0|\kappa|^2\le\lambda_{\rm ov}(z,\kappa)
 &\le C_0|\kappa|^2,\qquad |\kappa|\le\kappa_0,         \label{eq:newS7V53-low-cone}\\
 \lambda_{\rm ov}(z,\kappa)&\ge c_0,
 \qquad |\kappa|\ge\kappa_0.                           \label{eq:newS7V53-high-gap}
\end{align}
Thus \(\kappa=0\) is the only overlap zero on the complete fixed-jet
family.
\end{lemma}

\begin{proof}
Near zero,
\[
 s(z,\kappa)=J_E\kappa+O(|\kappa|^2),\qquad
 m(z,\kappa)=-\rho+O(|\kappa|^2)
\]
uniformly by \Cref{thm:newS7V51-hub-gap}.  Expanding
\eqref{eq:newS7V53-overlap-eigenvalue} gives
\[
 \lambda_{\rm ov}(z,\kappa)
 =|J_E\kappa|^2+O(|\kappa|^3).
\]
Compact ellipticity of \(E\) proves
\eqref{eq:newS7V53-low-cone} after decreasing \(\kappa_0\).

If \(\lambda_{\rm ov}=0\), then \(s=0\) and \(m<0\).  The large-\(r\)
argument in \Cref{thm:newS7V38-large-r-path} says that every zero of \(s\)
outside the physical cone has
\((r/2)\ell>\rho\), hence \(m>0\), a contradiction.  On the compact
complement of the cone, continuity gives
\eqref{eq:newS7V53-high-gap}.  The common-cell complement was assigned a
fixed positive stiffness and obeys the same compactness argument.
Differentiated uniformity follows from the compact smooth coefficient
family and the Wilson gap.
\end{proof}

\begin{corollary}[Physical high-branch energy]
\label{cor:newS7V53-high-energy}
For the physical overlap square
\(A_a^{\rm jb}:=(\mathscr D_a^{\rm jb})^*
\mathscr D_a^{\rm jb}\), every frozen mode with
\(|\kappa|\ge\kappa_0\) has energy at least \(c_0/a^2\).  The low branch
has the continuum quadratic cone with uniform ellipticity.
\end{corollary}

\begin{proof}
Physical overlap eigenvalues are \(a^{-2}\lambda_{\rm ov}\).  Apply
\Cref{lem:newS7V53-overlap-cone}.
\end{proof}

\subsection{The phase-space heat split}
\label{sec:newS7V53-split}

Let \(\Pi_{\rm lo}\) and \(\Pi_{\rm hi}\) be a smooth periodic momentum
partition, with \(\Pi_{\rm lo}\) supported in
\(|\kappa|<2\kappa_0\) and equal to one on
\(|\kappa|\le\kappa_0\).  Quantize it in the same midpoint calculus as the
JBMK.

\begin{lemma}[Variable-coefficient high-branch suppression]
\label{lem:newS7V53-high-tail}
For every \(N\) and \(0<t\le t_0\), the normalized local or integrated
relative heat contribution containing \(\Pi_{\rm hi}\) obeys
\begin{equation}
 |\mathcal H_{a,{\rm hi}}^{\rm rel}(t)|
 \le C_N\left\{
 \left({a^2\over t}\right)^N+
 L^{-N}\right\}
 +C\left({t\over a^2}\right)^2e^{-ct/a^2}.             \label{eq:newS7V53-high-tail}
\end{equation}
The same estimate holds after two metric variations.
\end{lemma}

\begin{proof}
On the support of the high cutoff,
\Cref{cor:newS7V53-high-energy} gives the frozen lower bound
\(c/a^2\).  Commuting the cutoff through the midpoint operator produces
one factor \(L^{-1}\) from a profile derivative; iterating the
almost-invariant phase-space projection \(N\) times leaves \(O(L^{-N})\).
On the invariant high block, the spectral theorem gives
\(\exp(-ct/a^2)\).  There are \(O(a^{-4})\) lattice states per fixed
four-volume; normalized relative heat power counting converts this to the
polynomial prefactor displayed, exactly as in
\Cref{thm:newS7V8-heat-comparison}.

On the transition annulus, repeated resolvent integration by parts in the
periodic symbol calculus gives \((a^2/t)^N\).  All symbol derivatives are
uniform because the sign kernel has the fixed exponential strip from
\Cref{cor:newS7V52-sign-locality}; an \(x\)-derivative is charged to the
explicit \(L^{-1}\) remainder above.  Metric differentiation inserts at
most two uniformly bounded resolvent words and preserves the energy
denominator, proving the contact estimates.
\end{proof}

\subsection{Low-branch parametrix with a shrinking profile}
\label{sec:newS7V53-low}

Let \(C_{0,a}^{\rm jb},C_{2,a}^{\rm jb},C_{4,a}^{\rm jb}\) denote the
owned local relative coefficient functionals obtained by the covariant
low-symbol recursion.  Here \emph{owned} means the matter, gauge, Higgs,
and mixed scalar-curvature terms already computed in V23 and V29; the pure
gravitational terms are not silently included.

Define the single error scale
\begin{equation}
 \eta_a:=
 {1\over aL^2}+{1\over L}+a w^{-5}.                   \label{eq:newS7V53-eta}
\end{equation}
For the V52 scales,
\begin{equation}
 \eta_a\le C a^{3/8}.                                  \label{eq:newS7V53-eta-rate}
\end{equation}

\begin{theorem}[Stratified low-branch heat parametrix]
\label{thm:newS7V53-low-parametrix}
On every bounded smooth gauge--Higgs--metric family, the relative JBMK
heat trace has, for \(a^2\le t\le t_0\),
\begin{equation}
 H_{a,{\rm lo}}^{\rm rel,jb}(t)
 =(4\pi t)^{-2}
 \{C_{0,a}^{\rm jb}+tC_{2,a}^{\rm jb}
   +t^2C_{4,a}^{\rm jb}+r_{a,{\rm lo}}(t)\},           \label{eq:newS7V53-low-expansion}
\end{equation}
where, for every \(N\),
\begin{equation}
 {|\partial_g^j r_{a,{\rm lo}}(t)|\over t^2}
 \le C_j\left\{
 {a\over\sqrt t}+\eta_a+\sqrt t
 \right\}
 +C_{j,N}\left({a^2\over t}\right)^N,
 \qquad j=0,1,2.                                      \label{eq:newS7V53-low-remainder}
\end{equation}
The constants are independent of the shrinking collars.
\end{theorem}

\begin{proof}
On the low branch put \(\kappa=ap\).  The uniform Wilson gap permits the
sign resolvent to be expanded by the Cauchy formula; the fixed exponential
strip makes all momentum moments finite.  Retaining covariant symbol
weight at most four produces
\(C_{0,a}^{\rm jb},C_{2,a}^{\rm jb},C_{4,a}^{\rm jb}\).
The factorial-free heat divided-difference estimate of V15 sums the
remaining insertions, exactly as in
\Cref{thm:newS7V16-local-parametrix}.

There are four error classes.  First, replacing the lattice dispersion and
vertices by their continuum values on heat momentum
\(|p|=O(t^{-1/2})\) costs \(Ca/\sqrt t\); this is the V16 lattice term.
Second, the midpoint constant and affine defects cost
\[
 C\{(aL^2)^{-1}+L^{-1}\}
\]
by \Cref{cor:newS7V52-physical-rate}.  Third, after the weight-four
coefficient has been retained, the worst differentiated collar insertion
has one explicit \(a\) and five profile derivatives.  Its pointwise size
is \(Ca w^{-5}\) by
\Cref{lem:newS7V52-power-counting}; higher lattice orders are smaller by
\Cref{thm:newS7V52-jet-closure}.  These three contributions give
\(\eta_a\).

Fourth, the first omitted continuum heat weight gives the ordinary
\(C\sqrt t\) relative-parametrix remainder on a bounded smooth family.
Terms with one additional shrinking-profile derivative have global
coefficient at most \(Ca w^{-5}\): their pointwise bound
\(Ca w^{-6}\) is supported in a set of volume \(O(w)\).
Thus they are already covered by \(\eta_a\sqrt t\le C\sqrt t\).
Higher codimension collars have volume \(O(w^d)\) and are smaller.
Taylor expansion of the low symbol to arbitrary order leaves
\(C_N(a^2/t)^N\).

Every metric contact differentiates a finite resolvent or transport word.
V52 gives the same jet powers through two contacts, and the factorial-free
majorant is unchanged.  This proves
\eqref{eq:newS7V53-low-remainder}.
\end{proof}

\begin{lemma}[Owned coefficient convergence through the strata]
\label{lem:newS7V53-C4}
Let \(C_4^{\rm owned}\) be the smooth continuum coefficient assembled in
V23 and V29.  Then
\begin{equation}
 |\partial_g^j(C_{4,a}^{\rm jb}-C_4^{\rm owned})|
 \le C_j\eta_a,\qquad j=0,1,2.                        \label{eq:newS7V53-C4}
\end{equation}
The same statement holds for the owned parts of \(C_0\) and \(C_2\) after
their declared local subtractions.
\end{lemma}

\begin{proof}
Away from the original skeleton, this is the smooth local lattice transfer
proved in V23 and V29.  In a stratum collar, the physical first Dirac jet
is the same as in the bulk by
\Cref{thm:newS7V51-hub-gap}; all collar-dependent coefficient monomials of
heat weight four are bounded by the differentiated ledger of
\Cref{thm:newS7V52-jet-closure}.  The worst rate is \(a^{3/8}\), and the
midpoint defects contribute the first two terms of
\eqref{eq:newS7V53-eta}.  A finite collar cover and the codimension
volume bounds sum these estimates.  Metric differentiation is already
included in both cited results.
\end{proof}

\subsection{An explicit mesoscopic endpoint}
\label{sec:newS7V53-window}

\begin{definition}[JBMK mesoscopic scale]
\label{def:newS7V53-epsilon}
Set
\begin{equation}
 \varepsilon_{\rm jb}(a):=a^{1/2}.                    \label{eq:newS7V53-epsilon}
\end{equation}
\end{definition}

Then \(\varepsilon_{\rm jb}\downarrow0\) and
\(a^2/\varepsilon_{\rm jb}=a^{3/2}\to0\).

\begin{theorem}[Logarithmic integrability of the stratified remainder]
\label{thm:newS7V53-integrability}
Let \(r_a^{\rm jb}\) be the sum of the low remainder in
\eqref{eq:newS7V53-low-expansion} and the high contribution of
\Cref{lem:newS7V53-high-tail}, written in the normalization of MATC.
Then, through two metric contacts,
\begin{equation}
 \lim_{\tau\downarrow0}\limsup_{a\downarrow0}
 \int_{\varepsilon_{\rm jb}(a)}^\tau
 {|\partial_g^j r_a^{\rm jb}(t)|\over t^3}\,dt=0,
 \qquad j=0,1,2.                                      \label{eq:newS7V53-integrability}
\end{equation}
\end{theorem}

\begin{proof}
Divide \eqref{eq:newS7V53-low-remainder} by \(t\) and integrate.  The
lattice term is
\[
 a\int_{\varepsilon_{\rm jb}}^\tau t^{-3/2}\,dt
 \le {2a\over\sqrt{\varepsilon_{\rm jb}}}
 =2a^{3/4}.
\]
The collar term is bounded by
\[
 \eta_a\log{\tau\over\varepsilon_{\rm jb}}
 \le C a^{3/8}|\log a|,
\]
and therefore tends to zero.  The continuum remainder gives
\[
 \int_{\varepsilon_{\rm jb}}^\tau t^{-1/2}\,dt
 \le2\sqrt\tau.
\]
Finally,
\[
 \int_{\varepsilon_{\rm jb}}^\tau
 \left({a^2\over t}\right)^N{dt\over t}
 \le {1\over N}
 \left({a^2\over\varepsilon_{\rm jb}}\right)^N
 ={a^{3N/2}\over N}.
\]
Choose the arbitrary \(N\) in
\Cref{lem:newS7V53-high-tail} so that
\(L^{-N}|\log\varepsilon_{\rm jb}|\to0\).
The exponentially small term is handled after
\(u=t/a^2\), as in V8.  Taking \(a\downarrow0\), then
\(\tau\downarrow0\), proves the claim.
\end{proof}

\begin{corollary}[Stratified owned rows of MATC]
\label{cor:newS7V53-MATC}
For the smooth JBMK family, rows (i)--(iii) of MATC hold for every
previously owned matter and mixed metric coefficient, with
\(\varepsilon=\varepsilon_{\rm jb}\), through two metric contacts.
Replacing the smooth-locus overlap kernel by the complete stratified JBMK
does not change those continuum \(C_4\) owners.
\end{corollary}

\begin{proof}
The expansion is
\Cref{thm:newS7V53-low-parametrix} plus
\Cref{lem:newS7V53-high-tail}.  Coefficient convergence is
\Cref{lem:newS7V53-C4}.  The MATC remainder row is
\Cref{thm:newS7V53-integrability}.
\end{proof}

\begin{corollary}[Fixed-range CHSC is acquired for the seam modification]
\label{cor:newS7V53-CHSC}
The fixed-range replacement of CHSC holds: relative to the smooth owned
continuum parametrix, the JBMK seam modification has convergent local
coefficients and a logarithmically integrable mesoscopic remainder through
two metric contacts.  No growing-range completion heat estimate is needed.
\end{corollary}

\begin{proof}
Use \Cref{corr:newS7V52-no-completion,
lem:newS7V53-C4,thm:newS7V53-integrability}.
\end{proof}

\subsection{The remaining coefficient ownership problem}
\label{sec:newS7V53-remaining}

This note transfers only coefficients already computed and assigned in
the smooth programme.  It does not manufacture the uncomputed purely
gravitational \(R^2\), \(|{\rm Ric}|^2\),
\(|{\rm Riem}|^2\), and total-divergence owners of the complete
overlap--Yukawa square.  Nor does it prove MATC row (iv) for those owners.
Those terms are required before the two metric contacts can be assembled
into one common Einstein stress and before the microscopic interval can be
declared exhausted.

\begin{assessment}[The seam heat bottleneck is closed]
\label{ass:newS7V53-status}
The fixed-jet, fixed-range stratified kernel has a uniform mesoscopic heat
parametrix.  Its high Brillouin branch is exponentially suppressed, its
low branch has an explicit integrable remainder, and its previously owned
\(C_4\) coefficients converge through two metric contacts.  The former
completion-locality obstruction is no longer active.  The next bottleneck
lies upstream in the smooth pure-gravity coefficient ledger.
\end{assessment}

\begin{programme}[V54: pure gravitational owners]
\label{prog:newS7V53-V54}
V54 will return to the smooth overlap square and compute the complete
parity-even pure gravitational heat-weight-four coefficient with all
spin, internal multiplicity, determinant-modulus, and subtraction factors
visible.  It will separate the Euler density and total divergence from the
metric-active Weyl and scalar-curvature combinations, differentiate the
owned functional twice, and state the exact residual needed for a common
Einstein stress.  The V53 transfer theorem will then carry any acquired
owner through the stratified mesh.
\end{programme}

\begin{firewall}[Claim boundary after seventh-series V53]
\label{fire:newS7V53-boundary}
V53 proves uniform overlap ellipticity on the fixed-jet family,
high-branch suppression, the stratified low-branch heat parametrix, an
explicit \(a^{1/2}\) mesoscopic endpoint, logarithmic remainder
integrability, and convergence through two metric contacts for every
matter and mixed metric \(C_4\) owner already established in V23 and V29.
It acquires the fixed-range seam version of CHSC.

V53 does not compute the pure gravitational heat coefficient, prove its
microscopic counterterm exhaustion, assemble the full common Einstein
stress, control rough metric or AOP fields, prove measure concentration,
construct the determinant-line phase, remove the interacting cutoff, or
construct the Standard Model--Einstein continuum.
\end{firewall}

\begin{theorem}[Seventh-series V53 result]
\label{thm:newS7V53-result}
For the complete fixed-range stratified JBMK and
\(\varepsilon(a)=a^{1/2}\), the high overlap branch is uniformly
negligible and the low branch has a covariant heat expansion through
weight four with
\[
 {|\partial_g^j r_a(t)|\over t^2}
 \le C_j\{a/\sqrt t+a^{3/8}+\sqrt t\}
 +C_{j,N}(a^2/t)^N,
 \qquad j=0,1,2.
\]
Its proper-time remainder is logarithmically integrable, and every
previously owned matter or mixed metric \(C_4\) coefficient converges to
the smooth continuum value through two metric contacts.
\end{theorem}

\begin{proof}
Combine \Cref{lem:newS7V53-overlap-cone,
lem:newS7V53-high-tail,
thm:newS7V53-low-parametrix,
lem:newS7V53-C4,
thm:newS7V53-integrability,
cor:newS7V53-MATC,
cor:newS7V53-CHSC}.
\end{proof}

\paragraph{Parent-version record.}
V53 was formed from
\texttt{Renewal\_SM\_Einstein\_Continuum\_Research\_Notes\_V52.tex}.
The V52 parent had \(20\,943\) logical source lines, \(843\,759\) bytes,
and SHA--256
\[
 \texttt{75F815D20C23682DF9FD2EE243D878359207D96574409D61A64266DC1F3B12D4}.
\]
The next compulsory companion audit is V55.

% =============================================================================
% END APPENDED SEVENTH-SERIES V53 MATERIAL
% =============================================================================

% =============================================================================
% BEGIN APPENDED SEVENTH-SERIES V54 MATERIAL
% =============================================================================

\section{Seventh-series V54: pure gravitational heat owners}
\label{sec:v7v54}
\label{sec:newS7V54}

\subsection{The complete local coefficient in the established convention}
\label{sec:newS7V54-general}

The V7 and V17 calculations retained only terms involving gauge or Yukawa
fields.  This note restores the identity, spin-curvature, and total
derivative terms in the same convention
\[
 L=\nabla^*\nabla+E,\qquad
 \Tr e^{-tL}\sim(4\pi t)^{-2}
 \sum_{j\ge0}t^jA_{2j}(L).
\]
Write \(\Delta=\nabla^\mu\nabla_\mu\) on scalar functions.

\begin{lemma}[Full local \(A_4\) formula]
\label{lem:newS7V54-full-A4}
For a Laplace-type operator \(L=\nabla^*\nabla+E\), the local coefficient
in the convention above is
\begin{align}
 a_4(x,L)={1\over360}\tr\Big[
 &(5R^2-2|{\rm Ric}|^2+2|{\rm Riem}|^2+12\Delta R)I
 -60RE+180E^2                                      \notag\\
 &\quad+30\Omega_{\mu\nu}\Omega^{\mu\nu}
 -60\Delta E\Big].                                    \label{eq:newS7V54-full-A4}
\end{align}
On a closed manifold the two Laplacian terms integrate to zero.
\end{lemma}

\begin{proof}
Write the standard Laplace operator as
\(-(\nabla^2+\mathcal E)\).  The diagonal heat recursion at weight four is
\[
 {1\over360}\tr\big[
 (5R^2-2|{\rm Ric}|^2+2|{\rm Riem}|^2+12\Delta R)I
 +60R\mathcal E+180\mathcal E^2
 +30\Omega^2+60\Delta\mathcal E\big].
\]
The convention of this manuscript has \(\mathcal E=-E\), as recorded in
\Cref{lem:newS7V17-general-A4}.  Substitution gives
\eqref{eq:newS7V54-full-A4}.  Stokes' theorem removes the scalar
Laplacians after integration on a closed manifold.
\end{proof}

\subsection{Spin curvature and the Dirac coefficient}
\label{sec:newS7V54-spin}

For the untwisted spin connection,
\begin{equation}
 \Omega_{\mu\nu}^S={1\over4}R_{\mu\nu ab}\gamma^{ab}.   \label{eq:newS7V54-spin-curvature}
\end{equation}

\begin{lemma}[Spin-curvature trace]
\label{lem:newS7V54-spin-trace}
With the Clifford conventions of V7,
\begin{equation}
 \tr_S(\Omega_{\mu\nu}^S\Omega^{S,\mu\nu})
 =-{1\over2}|{\rm Riem}|^2.                            \label{eq:newS7V54-spin-trace}
\end{equation}
\end{lemma}

\begin{proof}
Insert \eqref{eq:newS7V54-spin-curvature} and use
\eqref{eq:newS7V7-two-gamma}:
\begin{align*}
 \tr_S(\Omega_{\mu\nu}^S\Omega^{S,\mu\nu})
 &={1\over16}R_{\mu\nu ab}R^{\mu\nu}{}_{cd}
 4(g^{ad}g^{bc}-g^{ac}g^{bd})\\
 &=-{1\over2}R_{\mu\nu ab}R^{\mu\nu ab}.
\end{align*}
The last equality uses antisymmetry in \(a,b\).
\end{proof}

\begin{theorem}[Pure-gravity Dirac heat coefficients]
\label{thm:newS7V54-Dirac}
For one complex four-component Dirac square tensored with an internal
space of complex dimension \(n\),
\begin{align}
 a_0^{D}&=4n,                                           \label{eq:newS7V54-Dirac-a0}\\
 a_2^{D}&=-{n\over3}R,                                  \label{eq:newS7V54-Dirac-a2}\\
 a_4^{D}\big|_{\rm grav}
 &={n\over360}\left(
 5R^2-8|{\rm Ric}|^2-7|{\rm Riem}|^2-12\Delta R
 \right).                                               \label{eq:newS7V54-Dirac-a4}
\end{align}
The \(a_0\) and \(a_2\) equations are local densities.
\end{theorem}

\begin{proof}
The Lichnerowicz endomorphism is \(E=(R/4)I_S\), so
\[
 a_2=\tr_{S\otimes\mathbb C^n}(RI/6-E)
 =n(4R/6-R)=-{n\over3}R.
\]
For \(a_4\), the identity term in
\eqref{eq:newS7V54-full-A4} contributes
\[
 n(20R^2-8|{\rm Ric}|^2+8|{\rm Riem}|^2+48\Delta R).
\]
The \(-60RE+180E^2\) terms contribute
\[
 n(-60R^2+45R^2),
\]
the spin curvature contributes
\(-15n|{\rm Riem}|^2\) by
\Cref{lem:newS7V54-spin-trace}, and
\(-60\Delta E\) contributes \(-60n\Delta R\).
The numerator is therefore
\[
 n\{5R^2-8|{\rm Ric}|^2-7|{\rm Riem}|^2-12\Delta R\},
\]
which proves \eqref{eq:newS7V54-Dirac-a4}.  The spin rank gives
\eqref{eq:newS7V54-Dirac-a0}.
\end{proof}

\begin{corollary}[Parity-even Weyl coefficient]
\label{cor:newS7V54-Weyl}
For one physical complex Weyl block in an internal representation of
dimension \(n\), the parity-even pure gravitational coefficients are one
half of \eqref{eq:newS7V54-Dirac-a0}--\eqref{eq:newS7V54-Dirac-a4}:
\begin{align}
 a_0^{W}&=2n,                                           \label{eq:newS7V54-Weyl-a0}\\
 a_2^{W}&=-{n\over6}R,                                  \label{eq:newS7V54-Weyl-a2}\\
 a_4^{W}\big|_{\rm grav,even}
 &={n\over720}\left(
 5R^2-8|{\rm Ric}|^2-7|{\rm Riem}|^2-12\Delta R
 \right).                                               \label{eq:newS7V54-Weyl-a4}
\end{align}
The chiral spin trace also has a gravitational Pontryagin density.  It is
parity odd and belongs to the determinant-line phase ledger, not to the
positive modulus computed here.
\end{corollary}

\begin{proof}
Insert \(P_\chi=(I+\chi\gamma_5)/2\).  Every parity-even Clifford trace is
one half of the full spin trace.  The \(\gamma_5\) part is proportional to
the completely antisymmetric contraction of two Riemann tensors and is
therefore the separate Pontryagin density.
\end{proof}

\subsection{Euler--Weyl decomposition and active metric owner}
\label{sec:newS7V54-decomposition}

In four dimensions set
\begin{align}
 \mathcal E_4&:=|{\rm Riem}|^2-4|{\rm Ric}|^2+R^2,      \label{eq:newS7V54-Euler}\\
 |C|^2&:=|{\rm Riem}|^2-2|{\rm Ric}|^2+{1\over3}R^2.    \label{eq:newS7V54-Weyl-square}
\end{align}

\begin{lemma}[No independent \(R^2\) owner]
\label{lem:newS7V54-decomposition}
The Dirac numerator obeys the exact identity
\begin{equation}
 5R^2-8|{\rm Ric}|^2-7|{\rm Riem}|^2
 =11\mathcal E_4-18|C|^2.                              \label{eq:newS7V54-decomposition}
\end{equation}
Thus the massless Dirac square has no independent pure \(R^2\) coefficient.
\end{lemma}

\begin{proof}
Substitute
\eqref{eq:newS7V54-Euler}--\eqref{eq:newS7V54-Weyl-square}.
The Riemann, Ricci, and scalar coefficients on the right are respectively
\(11-18=-7\), \(-44+36=-8\), and \(11-6=5\).
\end{proof}

On a closed oriented four-manifold,
\(\int\mathcal E_4\,d\mu_g=32\pi^2\chi(M)\) and
\(\int\Delta R\,d\mu_g=0\).  Hence only the Weyl-square term contributes
to a local metric variation.

\begin{definition}[Bach tensor convention]
\label{def:newS7V54-Bach}
Define \(B_{\mu\nu}\) by the variational identity
\begin{equation}
 \delta\int_M|C|^2\,d\mu_g
 =-2\int_M B_{\mu\nu}\,\delta g^{\mu\nu}\,d\mu_g.       \label{eq:newS7V54-Bach-variation}
\end{equation}
Equivalently,
\begin{equation}
 B_{\mu\nu}
 =\left(\nabla^\rho\nabla^\sigma
       +{1\over2}R^{\rho\sigma}\right)
 C_{\mu\rho\nu\sigma}                                  \label{eq:newS7V54-Bach}
\end{equation}
in this curvature-sign convention.
\end{definition}

\begin{theorem}[Pure-gravity heat stress and Hessian]
\label{thm:newS7V54-Bach-owner}
For one physical Weyl block of internal dimension \(n\), the integrated
parity-even \(A_4\) functional has first variation
\begin{equation}
 \delta A_4^W\big|_{\rm grav,even}
 ={n\over20}\int_M
 B_{\mu\nu}\,\delta g^{\mu\nu}\,d\mu_g.                \label{eq:newS7V54-Weyl-variation}
\end{equation}
For the determinant modulus, this is multiplied by \(1/2\).
Its second variation is the bilinear form obtained by linearizing the
Bach tensor:
\begin{equation}
 \delta_k\delta_h A_4^W\big|_{\rm grav,even}
 ={n\over20}\,
 \delta_k\int_M B_{\mu\nu}h^{\mu\nu}\,d\mu_g.           \label{eq:newS7V54-Weyl-Hessian}
\end{equation}
On a bounded \(C^6\) metric family it is a bounded, symmetric,
fourth-order local bilinear form in \(h,k\).
\end{theorem}

\begin{proof}
By \Cref{cor:newS7V54-Weyl,lem:newS7V54-decomposition}, the integrated
metric-active term is
\[
 -{n\over40}\int_M|C|^2\,d\mu_g.
\]
The Euler term is topological and the Laplacian term integrates to zero.
Apply \eqref{eq:newS7V54-Bach-variation} to obtain
\eqref{eq:newS7V54-Weyl-variation}.  Differentiate once more to obtain
\eqref{eq:newS7V54-Weyl-Hessian}.  Formula
\eqref{eq:newS7V54-Bach} contains two derivatives of curvature and hence
four derivatives of the metric.  Smooth differentiation under the
integral is valid on a bounded \(C^6\) family.  Symmetry follows because
\eqref{eq:newS7V54-Weyl-Hessian} is the Hessian of a scalar functional.
\end{proof}

\subsection{The Standard Model multiplicity ledger}
\label{sec:newS7V54-SM}

Let \(N_W^\bullet\) be the total complex internal dimension of the
all-left-handed physical Standard Model multiplets, including generations
and colour:
\begin{equation}
 N_W^{\min}
 =3(6+3+3+2+1)=45,\qquad
 N_W^\nu=N_W^{\min}+3=48.                              \label{eq:newS7V54-NW}
\end{equation}
The summands are \(Q_L,u_R^c,d_R^c,L_L,e_R^c\); the optional branch adds
one \(\nu_R^c\) per generation.

\begin{theorem}[Physical and determinant pure-gravity ledgers]
\label{thm:newS7V54-SM-ledger}
Before the common heat prefactor \((4\pi)^{-2}\), the physical chiral-block
Standard Model coefficients are
\begin{align}
 A_{0,{\rm SM}}^{\rm phys}&=2N_W^\bullet\int_Md\mu_g,   \label{eq:newS7V54-SM-a0}\\
 A_{2,{\rm SM}}^{\rm phys}\big|_{\rm grav}
 &=-{N_W^\bullet\over6}\int_MR\,d\mu_g,                \label{eq:newS7V54-SM-a2}\\
 A_{4,{\rm SM}}^{\rm phys}\big|_{\rm grav,even}
 &={N_W^\bullet\over720}\int_M
 \{11\mathcal E_4-18|C|^2-12\Delta R\}\,d\mu_g.        \label{eq:newS7V54-SM-a4}
\end{align}
For
\(\log|\det\mathcal B|=(1/2)\log\det(\mathcal B^*\mathcal B)\),
each right side is multiplied by \(1/2\).  In particular, the
metric-active determinant-modulus \(A_4\) owner is
\begin{equation}
 -{N_W^\bullet\over80}\int_M|C|^2\,d\mu_g,             \label{eq:newS7V54-SM-C2-det}
\end{equation}
and its heat-owner metric gradient is
\begin{equation}
 {N_W^\bullet\over40}B_{\mu\nu}.                       \label{eq:newS7V54-SM-Bach}
\end{equation}
The fermionic functional-integral sign and the sign of the proper-time
representation remain separate from this heat/multiplicity ledger.
\end{theorem}

\begin{proof}
Sum \Cref{cor:newS7V54-Weyl} over the representation dimensions in
\eqref{eq:newS7V54-NW}.  The determinant-modulus factor is the convention
fixed in V20--V24.  Since
\(-18/1440=-1/80\), \eqref{eq:newS7V54-SM-C2-det} follows.
Its variation follows from
\eqref{eq:newS7V54-Bach-variation}.
\end{proof}

\begin{corollary}[Complete parity-even physical smooth \(A_4\) ledger]
\label{cor:newS7V54-complete-smooth}
The physical smooth Standard Model parity-even heat coefficient through
weight four is
\begin{align}
 {1\over(4\pi)^2}\int_M\Big[
 &{N_W^\bullet\over720}
   \{11\mathcal E_4-18|C|^2-12\Delta R\}
 +{1\over3}\mathfrak G_{\rm SM}(F)                    \notag\\
 &+2\mathfrak a_\bullet|DH|^2
  +2\mathfrak b_\bullet|H|^4
  +{\mathfrak a_\bullet\over3}R|H|^2
 \Big]\,d\mu_g.                                        \label{eq:newS7V54-complete-smooth}
\end{align}
The determinant modulus is one half of this physical heat ledger.
\end{corollary}

\begin{proof}
Add the pure-gravity term
\eqref{eq:newS7V54-SM-a4} to the mutually independent owners in
\eqref{eq:newS7V29-smooth-total}.  The Lichnerowicz trace calculation in
V17 proves that no omitted parity-even gauge--Higgs--curvature cross term
exists at this weight.
\end{proof}

\subsection{Transfer through the stratified JBMK}
\label{sec:newS7V54-transfer}

\begin{theorem}[Stratified pure-gravity coefficient transfer]
\label{thm:newS7V54-transfer}
For the V52 JBMK and the V53 endpoint
\(\varepsilon(a)=a^{1/2}\), the lattice pure-gravity coefficient satisfies
\begin{equation}
 \left|\partial_g^j
 \left(C_{4,a}^{\rm grav,jb}
 -{N_W^\bullet\over720}
  \int_M\{11\mathcal E_4-18|C|^2-12\Delta R\}\,d\mu_g
 \right)\right|
 \le C_j a^{3/8},\quad j=0,1,2.                        \label{eq:newS7V54-transfer}
\end{equation}
Its mesoscopic remainder obeys the logarithmic integrability row of MATC.
\end{theorem}

\begin{proof}
The low-symbol recursion in
\Cref{thm:newS7V53-low-parametrix} retains every covariant invariant of
heat weight at most four; the word estimates did not depend on whether the
retained curvature was gauge or Riemannian.  The spin trace identifying
the pure Riemannian invariant is
\Cref{thm:newS7V54-Dirac}, and the physical multiplicity is
\Cref{thm:newS7V54-SM-ledger}.  The V52 differentiated collar ledger gives
the \(a^{3/8}\) bound through two contacts.  High-branch suppression and
remainder integrability are
\Cref{lem:newS7V53-high-tail,thm:newS7V53-integrability}.
\end{proof}

\begin{corollary}[All smooth parity-even \(C_4\) owners cross the seams]
\label{cor:newS7V54-all-C4}
Every parity-even term in
\eqref{eq:newS7V54-complete-smooth} transfers to the complete stratified
JBMK, with two metric contacts and logarithmically integrable mesoscopic
error.
\end{corollary}

\begin{proof}
Use \Cref{cor:newS7V53-MATC} for the matter and mixed terms and
\Cref{thm:newS7V54-transfer} for the pure-gravity terms.
\end{proof}

\subsection{The microscopic gravitational subtraction remains}
\label{sec:newS7V54-remaining}

The calculation closes the continuum coefficient and its mesoscopic
stratified transfer.  MATC row (iv) is stronger: on
\(0<t<\varepsilon(a)\), it requires an exact finite-cutoff local
subtraction whose metric gradient and Hessian agree with the declared
cosmological, Einstein--Hilbert, and curvature-square owners and leave no
nonlocal ultraviolet remainder.  V8 controls the free rank, and V29
controls the smooth \(A_2\) response, but the complete nonlinear metric
contact subtraction has not yet been assembled.  Therefore
\eqref{eq:newS7V54-SM-Bach} is the acquired curvature-square component of
the common heat stress, not yet the full cutoff-renormalized Einstein
stress.

\begin{assessment}[The pure-gravity coefficient is acquired]
\label{ass:newS7V54-status}
The parity-even gravitational \(A_4\) owner is now explicit, including
spin trace, Weyl multiplicity, determinant factor, Euler and divergence
terms, Bach variation, and the second metric contact.  V53 transfers it
through every stratified seam.  The next obstruction is the exact
microscopic nonlinear metric subtraction, followed by rough-field
control.
\end{assessment}

\begin{programme}[V55: compulsory audit and microscopic metric
subtraction]
\label{prog:newS7V54-V55}
V55 will first perform the compulsory YM/NS audit.  It will then define
the exact finite-cutoff metric Taylor owner through second contact at each
JBMK cell, test its diffeomorphism and frame Ward identities, and determine
whether the microscopic interval reduces to the local
\(A_0,A_2,A_4\) counterterms.  If the cellwise Hessian fails to glue
covariantly, the construction will go upstream to a relative determinant
with a common metric base rather than assert MATC row (iv).
\end{programme}

\begin{firewall}[Claim boundary after seventh-series V54]
\label{fire:newS7V54-boundary}
V54 proves the full pure-gravity Dirac and Weyl heat coefficients, their
Euler--Weyl decomposition, the Bach first variation and symmetric second
contact, the minimal and optional-neutrino Standard Model multiplicity
ledgers, and the stratified mesoscopic transfer of every parity-even
\(C_4\) owner.

V54 does not prove the microscopic nonlinear metric subtraction, the full
cutoff-renormalized common Einstein stress, rough metric or AOP control,
measure concentration, the determinant-line phase, interacting cutoff
removal, or the Standard Model--Einstein continuum.
\end{firewall}

\begin{theorem}[Seventh-series V54 result]
\label{thm:newS7V54-result}
For \(N_W^\bullet=45\) in the minimal Standard Model and \(48\) in the
optional Dirac-neutrino branch, the physical chiral-block pure
gravitational heat coefficient is
\[
 {N_W^\bullet\over720(4\pi)^2}
 \int_M\{11\mathcal E_4-18|C|^2-12\Delta R\}\,d\mu_g.
\]
The determinant modulus carries one half.  Its only metric-active
curvature-square term is
\(-N_W^\bullet\int|C|^2/80\) before the common heat prefactor, with
gradient \(N_W^\bullet B_{\mu\nu}/40\).  This coefficient and its first
two metric contacts transfer to the complete stratified JBMK with error
\(O(a^{3/8})\) and a logarithmically integrable mesoscopic remainder.
\end{theorem}

\begin{proof}
Combine \Cref{lem:newS7V54-full-A4,
lem:newS7V54-spin-trace,
thm:newS7V54-Dirac,
cor:newS7V54-Weyl,
lem:newS7V54-decomposition,
thm:newS7V54-Bach-owner,
thm:newS7V54-SM-ledger,
cor:newS7V54-complete-smooth,
thm:newS7V54-transfer,
cor:newS7V54-all-C4}.
\end{proof}

\paragraph{Parent-version record.}
V54 was formed from
\texttt{Renewal\_SM\_Einstein\_Continuum\_Research\_Notes\_V53.tex}.
The V53 parent had \(21\,361\) logical source lines, \(860\,179\) bytes,
and SHA--256
\[
 \texttt{F99B1260D184D4D4F7AA257CA398F8816EF4F2C6DDFCDCAF4393C2A73F051C4B}.
\]
The next compulsory companion audit is V55.

% =============================================================================
% END APPENDED SEVENTH-SERIES V54 MATERIAL
% =============================================================================

% =============================================================================
% BEGIN APPENDED SEVENTH-SERIES V55 MATERIAL
% =============================================================================

\section{Seventh-series V55: compulsory audit and exact microscopic metric contacts}
\label{sec:v7v55}
\label{sec:newS7V55}

\subsection{Compulsory companion audit}
\label{sec:newS7V55-audit}

\begin{audit}[Companion snapshots used at V55]
\label{aud:newS7V55-snapshots}
The newest coherent companion sources read before this note were:
\begin{center}
\begin{tabular}{p{0.18\linewidth}p{0.55\linewidth}r}
programme & source & bytes\\ \hline
Yang--Mills &
\texttt{Renewal\_Yang\_Mills\_Mass\_Gap\_Research\_Notes\_V29.tex}
& \(462\,845\)\\
Navier--Stokes &
\texttt{Renewal\_NS\_Stability\_Research\_Notes\_V5.tex}
& \(83\,663\)
\end{tabular}
\end{center}
The Yang--Mills source had \(13\,056\) newline characters and SHA--256
\[
 \texttt{32FEF34DA4B92B56CF90BFC0D14C28164FE92F6D3F7253D8DD2D5C9294F52E3D}.
\]
The file named
\texttt{Renewal\_Yang\_Mills\_Mass\_Gap\_Research\_Notes\_V30.tex}
had exactly the same byte count, newline count, and SHA--256 digest.
It is therefore a byte-identical alias, not a new mathematical append;
V29 is the coherent mathematical snapshot audited here.

The Navier--Stokes source had \(1\,594\) newline characters and SHA--256
\[
 \texttt{7A079B7C8EB8D22621FE980640B110828CFD9A6BDF2708C67A53626495ED7404}.
\]
Both sources ended in exactly one \(\verb|\end{document}|\).
\end{audit}

The Yang--Mills note proves a conditional multi-defect estimate from a
vacuum-fillability radius.  For its canonical collar, the incidence
constants are \(A=4\), \(B=24\), and the strict geometric threshold is
\(\rho<\delta/\sqrt{24}\).  It also proves that the required radius is an
attained finite compact optimization.  Its firewall is essential: the
occurrence of a sufficiently fillable local vacuum, and the matched
good-gate transmission estimate, remain open.  In particular, a periodic
or unconditional suppression estimate cannot replace the conditional
partition-function denominator.

The Navier--Stokes V5 note records the certified second and third Oseen
constants and the resulting type-I floor \(0.3438\).  Its next block
separates a near kernel part from an \(L^3\) far tail and uses
\(L^6\times L^2\) control rather than a global sup norm.  It explicitly
makes no global regularity claim.

\begin{decision}[V55 direction after the audit]
\label{dec:newS7V55-audit}
No numerical or probabilistic assertion is imported from either companion.
Two proof-order lessons are used:
\begin{enumerate}
\item retain the actual common normalization when comparing local metric
sectors, rather than subtracting independently normalized absolute
determinants; and
\item split microscopic locality by the heat radius and estimate its far
tail in a trace ideal, rather than controlling all metric contacts by one
pointwise supremum.
\end{enumerate}
Accordingly, the microscopic metric object below is a common-base relative
heat trace with both Duhamel contacts assembled before localization.
\end{decision}

\subsection{The common-base relative heat trace}
\label{sec:newS7V55-relative}

Use the half-density identification from V26 so that all metrics act on
one finite lattice Hilbert space \(\mathcal H_a\).  Let
\[
 A_a(g):=\mathscr D_a^{\rm jb}(g)^*
          \mathscr D_a^{\rm jb}(g)+\mu^2
\]
with a fixed \(\mu>0\) used only as an infrared regulator.  It does not
alter ultraviolet coefficients.  For two metrics in one positive
bounded-geometry chart define
\begin{equation}
 \Theta_a(t;g,g_0)
 :=\Tr_{\mathcal H_a}
 \{e^{-tA_a(g)}-e^{-tA_a(g_0)}\}.                      \label{eq:newS7V55-Theta}
\end{equation}

\begin{lemma}[Exact rank cancellation at the microscopic endpoint]
\label{lem:newS7V55-rank}
For every fixed \(a\),
\begin{equation}
 \Theta_a(0;g,g_0)=0,\qquad
 |\Theta_a(t;g,g_0)|\le
 t\,\|A_a(g)-A_a(g_0)\|_1.                             \label{eq:newS7V55-rank}
\end{equation}
Consequently \(\Theta_a(t;g,g_0)/t\) is integrable at \(t=0\).
\end{lemma}

\begin{proof}
The two identity operators have the same finite rank.  Duhamel's formula
gives
\[
 e^{-tA}-e^{-tA_0}
 =-\int_0^t e^{-(t-s)A}(A-A_0)e^{-sA_0}\,ds.
\]
Both heat factors are contractions.  Take the trace norm and integrate.
\end{proof}

\begin{definition}[Exact relative microscopic functional]
\label{def:newS7V55-mic}
With \(\varepsilon_{\rm jb}=a^{1/2}\), set
\begin{equation}
 \mathcal M_a(g,g_0)
 :={1\over2}\int_0^{\varepsilon_{\rm jb}(a)}
 {\Theta_a(t;g,g_0)\over t}\,dt.                       \label{eq:newS7V55-mic}
\end{equation}
The factor \(1/2\) is the determinant-modulus factor.  The fermionic
effective-action sign remains a separate ledger entry.
\end{definition}

\begin{corollary}[Finite-cutoff existence]
\label{cor:newS7V55-mic-exists}
\(\mathcal M_a(g,g_0)\) is finite for every \(a>0\) and vanishes when
\(g=g_0\).
\end{corollary}

\begin{proof}
Use \Cref{lem:newS7V55-rank} and
\eqref{eq:newS7V55-mic}.
\end{proof}

\subsection{Both metric contacts before localization}
\label{sec:newS7V55-contacts}

Let \(g_s\) be a \(C^2\) metric path and put
\(A_s=A_a(g_s)\), \(\dot A_s=\partial_sA_s\).

\begin{theorem}[Exact first and second heat contacts]
\label{thm:newS7V55-contacts}
At every fixed \(a,t>0\),
\begin{align}
 {d\over ds}\Tr e^{-tA_s}
 &=-t\Tr(\dot A_s e^{-tA_s}),                          \label{eq:newS7V55-first}\\
 {d^2\over ds^2}\Tr e^{-tA_s}
 &=-t\Tr(\ddot A_s e^{-tA_s})                         \notag\\
 &\quad+t^2\int_0^1
 \Tr\{\dot A_s e^{-(1-\theta)tA_s}
       \dot A_s e^{-\theta tA_s}\}\,d\theta.           \label{eq:newS7V55-second}
\end{align}
The separated and contact terms in
\eqref{eq:newS7V55-second} must be retained together.
\end{theorem}

\begin{proof}
Differentiate the exponential by Duhamel's formula,
\[
 {d\over ds}e^{-tA_s}
 =-t\int_0^1e^{-(1-\theta)tA_s}
 \dot A_s e^{-\theta tA_s}\,d\theta.
\]
Cyclicity of the finite-dimensional trace gives
\eqref{eq:newS7V55-first}.  Differentiate once more.  The derivative of
\(\dot A_s\) gives the first term of
\eqref{eq:newS7V55-second}; differentiating either heat factor and using
the two simplex parameters combines, by cyclicity, into the displayed
single \(\theta\)-integral.
\end{proof}

\begin{corollary}[Contacts of the microscopic functional]
\label{cor:newS7V55-mic-contacts}
On a bounded \(C^2\) metric chart, the first two path derivatives of
\(\mathcal M_a(g_s,g_0)\) are obtained by inserting
\eqref{eq:newS7V55-first}--\eqref{eq:newS7V55-second} under the integral
in \eqref{eq:newS7V55-mic}.
\end{corollary}

\begin{proof}
At fixed \(a\), all matrices and their first two metric derivatives are
bounded on the compact path.  The extra factor \(t\) in
\eqref{eq:newS7V55-first}--\eqref{eq:newS7V55-second} cancels the
proper-time denominator at zero, so dominated differentiation applies.
\end{proof}

\subsection{Exact frame and diffeomorphism Ward cancellation}
\label{sec:newS7V55-Ward}

\begin{theorem}[Unitary-orbit Ward identity]
\label{thm:newS7V55-Ward}
Suppose a frame change, gauge change, or simplicial diffeomorphism lift acts
by a \(C^2\) unitary family \(U_s\) on the common half-density space and
\[
 A_s=U_sA_0U_s^*.
\]
Then
\begin{equation}
 {d\over ds}\Tr e^{-tA_s}
 ={d^2\over ds^2}\Tr e^{-tA_s}=0                       \label{eq:newS7V55-Ward}
\end{equation}
for every \(t\), and the corresponding first two contacts of
\(\mathcal M_a\) vanish exactly.
\end{theorem}

\begin{proof}
Functional calculus gives
\(e^{-tA_s}=U_se^{-tA_0}U_s^*\), whose trace is constant.
Differentiating twice proves \eqref{eq:newS7V55-Ward}.  Equivalently, at
first order \(\dot A=[K,A]\), and cyclicity kills
\(\Tr([K,A]e^{-tA})\); at second order the contact and separated terms in
\eqref{eq:newS7V55-second} cancel.  Integrate in \(t\).
\end{proof}

\begin{warning}[Why the separated Hessian is not a stress]
\label{warn:newS7V55-separated}
The positive-looking second term in
\eqref{eq:newS7V55-second} is not separately invariant.  On a unitary
orbit it is cancelled by the \(\ddot A_s\) contact.  Dropping that contact
would create a spurious metric or gauge mass and violate the exact Ward
identity.
\end{warning}

\subsection{Microscopic heat-radius locality}
\label{sec:newS7V55-locality}

\begin{lemma}[Two-contact off-diagonal bound]
\label{lem:newS7V55-offdiag}
Let a metric variation be supported in a physical set \(K\), and let
\(\chi_{K_r^c}\) be supported at distance at least \(r\) from \(K\).
For \(j=1,2\) and \(0<t\le\varepsilon_{\rm jb}(a)\),
\begin{equation}
 \left\|\chi_{K_r^c}\,
 \partial_g^j e^{-tA_a(g)}\right\|_1
 \le C_j a^{-4}(1+t/a^2)^{m_j}
 \left\{e^{-c r^2/t}+e^{-cr/a}\right\}.                \label{eq:newS7V55-offdiag}
\end{equation}
\end{lemma}

\begin{proof}
The JBMK sign and both metric contacts have fixed exponential lattice
decay by \Cref{cor:newS7V52-sign-locality}.  Insert the first or second
Duhamel formula.  Davies conjugation by a truncated distance function
gives the Gaussian factor for chains of heat-scale length
\(\sqrt t\); a single long exponentially local hop gives
\(e^{-cr/a}\).  Each differentiated heat word has a fixed polynomial in
\(t/a^2\), as in \Cref{lem:newS7V29-vertex-bounds}.  The finite-volume
trace rank is \(O(a^{-4})\).  Hölder's inequality for trace ideals and
simplex integration give the stated bound.
\end{proof}

\begin{theorem}[Shrinking physical locality of the microscopic contacts]
\label{thm:newS7V55-mic-local}
Let \(r_a=a^{1/8}\).  The contribution to either metric contact of
\(\mathcal M_a\) from points at distance at least \(r_a\) from the
variation support tends to zero faster than every power of \(a\).
Thus the exact microscopic relative functional is local on a physical
radius tending to zero.
\end{theorem}

\begin{proof}
For \(t\le a^{1/2}\),
\[
 {r_a^2\over t}\ge a^{-1/4},\qquad
 {r_a\over a}=a^{-7/8}.
\]
Insert \eqref{eq:newS7V55-offdiag} in the contact formulas and integrate
with \(dt/t\).  Both exponentials dominate the fixed powers of
\(a^{-1}\) and \(t/a^2\), uniformly on the interval.
\end{proof}

\subsection{The exact remaining microscopic card}
\label{sec:newS7V55-card}

\begin{definition}[Microscopic metric subtraction card, MMSC]
\label{def:newS7V55-MMSC}
MMSC consists of:
\begin{enumerate}
\item a covariant finite-\(a\) local decomposition of
\(\mathcal M_a(g,g_0)\), through two metric contacts, into the relative
cosmological, Einstein--Hilbert, and complete V54 \(A_4\) owners plus a
remainder;
\item convergence of the \(A_4\) coefficient to
\eqref{eq:newS7V54-complete-smooth} and absorption of the power terms by
the declared \(A_0,A_2\) counterterms;
\item vanishing of the remainder, with the exact unitary-orbit Ward
identities of \Cref{thm:newS7V55-Ward}; and
\item compatibility on overlapping stratum charts with the same common
metric base \(g_0\).
\end{enumerate}
\end{definition}

\begin{theorem}[Rows of MMSC acquired at V55]
\label{thm:newS7V55-MMSC-partial}
The common-base construction acquires:
\begin{enumerate}
\item finite-cutoff existence and exact rank cancellation;
\item the complete first and second Duhamel contact;
\item exact frame, gauge, and lifted-diffeomorphism Ward cancellation; and
\item shrinking-radius trace locality through two metric contacts.
\end{enumerate}
The unproved row is the finite invariant power-counting decomposition into
only \(A_0,A_2,A_4\) owners.
\end{theorem}

\begin{proof}
Use \Cref{lem:newS7V55-rank,
cor:newS7V55-mic-exists,
thm:newS7V55-contacts,
cor:newS7V55-mic-contacts,
thm:newS7V55-Ward,
thm:newS7V55-mic-local}.
\end{proof}

\begin{assessment}[The audit selects a common-base subtraction]
\label{ass:newS7V55-status}
The microscopic functional is now an exact, finite, common-base relative
object.  Its two contacts obey the Ward identities and are local on a
shrinking physical radius.  Locality alone does not prove that its
cutoff-divergent part contains only the finite V54 invariant list.
That power-counting statement is the remaining MMSC row and cannot be
replaced by an independently normalized cell determinant.
\end{assessment}

\begin{programme}[V56: microscopic invariant power counting]
\label{prog:newS7V55-V56}
V56 will rescale \(t=a^2u\), expand the common-base local symbol in
covariant metric jets, and use hypercubic/supercell symmetry plus the Ward
identity to classify all scalar contact terms of engineering dimension at
most four.  It will prove that dimensions zero and two are precisely the
cosmological and Einstein--Hilbert owners and that dimension four is the
V54 Euler--Weyl ledger, or else isolate an additional lattice anisotropy
owner.  The \(u\)-tail will be matched to V53 at
\(\varepsilon_{\rm jb}/a^2=a^{-3/2}\).
\end{programme}

\begin{firewall}[Claim boundary after seventh-series V55]
\label{fire:newS7V55-boundary}
V55 performs the compulsory YM/NS audit, defines the exact common-base
microscopic relative heat functional, proves rank cancellation, its first
and second Duhamel contacts, exact unitary-orbit Ward identities, and
shrinking-radius trace locality.  It acquires all structural rows of MMSC
except invariant power counting.

V55 does not prove the microscopic \(A_0,A_2,A_4\) decomposition, complete
MATC row (iv), the full cutoff-renormalized common Einstein stress, rough
metric or AOP control, measure concentration, the determinant-line phase,
interacting cutoff removal, or the Standard Model--Einstein continuum.
\end{firewall}

\begin{theorem}[Seventh-series V55 result]
\label{thm:newS7V55-result}
At every finite cutoff, the common-base microscopic determinant-modulus
functional \eqref{eq:newS7V55-mic} is finite and has exact first and second
metric contacts.  The assembled contacts vanish on frame, gauge, and
lifted-diffeomorphism unitary orbits and are trace-local on physical radius
\(a^{1/8}\to0\).  Hence no nonlocal microscopic metric response survives
outside a shrinking neighbourhood; the remaining task is to prove that
the local response is exhausted by the V54 \(A_0,A_2,A_4\) invariant
owners.
\end{theorem}

\begin{proof}
Combine \Cref{dec:newS7V55-audit,
lem:newS7V55-rank,
thm:newS7V55-contacts,
thm:newS7V55-Ward,
lem:newS7V55-offdiag,
thm:newS7V55-mic-local,
thm:newS7V55-MMSC-partial}.
\end{proof}

\paragraph{Parent-version record.}
V55 was formed from
\texttt{Renewal\_SM\_Einstein\_Continuum\_Research\_Notes\_V54.tex}.
The V54 parent had \(21\,806\) logical source lines, \(877\,110\) bytes,
and SHA--256
\[
 \texttt{C9077B6FC65193E2F9396E1162C226A0E3260EA32A7EAFFBA968420CC9B27EA9}.
\]
The next compulsory companion audit is V60.

% =============================================================================
% END APPENDED SEVENTH-SERIES V55 MATERIAL
% =============================================================================

% =============================================================================
% BEGIN APPENDED SEVENTH-SERIES V56 MATERIAL
% =============================================================================

\section{Seventh-series V56: the microscopic mesh-Ward reduction}
\label{sec:v7v56}
\label{sec:newS7V56}

\begin{correction}[Scope of the V55 diffeomorphism statement]
\label{corr:newS7V56-Ward-scope}
\Cref{thm:newS7V55-Ward} is exact for frame and gauge changes and for
those simplicial relabellings which act by a unitary on the fixed lattice
space.  A general continuum diffeomorphism moves the embedded mesh and
need not be such a relabelling.  Therefore V55 does not prove the full
continuum diffeomorphism Ward identity at finite cutoff.
\end{correction}

\subsection{Microscopic metric polarization}
\label{sec:newS7V56-polarization}

Freeze one regular supercell at a constant metric \(g_0\).  For compactly
supported metric waves \(h_{\mu\nu}e^{iq\cdot x}\) and
\(k_{\rho\sigma}e^{-iq\cdot x}\), let
\begin{equation}
 \Pi_a^{\mu\nu,\rho\sigma}(q)
 :=\partial_h\partial_k\mathcal M_a(g_0+h+k,g_0)\big|_{h=k=0}.        \label{eq:newS7V56-Pi}
\end{equation}
The definition uses the complete contact-plus-separated expression
\eqref{eq:newS7V55-second}.

\begin{lemma}[Even finite-symbol expansion]
\label{lem:newS7V56-expansion}
For \(|aq|\) sufficiently small, the polarization is an even analytic
function of \(aq\) and has
\begin{equation}
 \Pi_a(q)
 =a^{-4}\mathsf P_0
  +a^{-2}\mathsf P_2(q)
  +\mathsf P_4(q)
  +O(a^2|q|^6),                                        \label{eq:newS7V56-expansion}
\end{equation}
where \(\mathsf P_{2j}\) is homogeneous of degree \(2j\) in \(q\).
Each tensor coefficient is a finite Brillouin-zone integral of the frozen
supercell overlap symbol and its first two metric contacts.
\end{lemma}

\begin{proof}
Set \(t=a^2u\) in \eqref{eq:newS7V55-mic}.  On a constant supercell,
Bloch transform makes every heat factor a finite matrix function of
\(\kappa\), while a metric wave shifts momentum by \(aq\).
The Wilson gap makes the sign and both metric derivatives analytic on a
fixed complex neighbourhood.  Taylor expansion in \(aq\) may therefore be
performed under the compact Brillouin integral and the finite
\(\theta\)-simplex integral in \eqref{eq:newS7V55-second}.
Reflection pairs \(\kappa\) with \(-\kappa\), so odd powers vanish.
Four-dimensional lattice normalization supplies \(a^{-4}\); each pair of
external momenta supplies \(a^2\), giving
\eqref{eq:newS7V56-expansion}.  The coefficient formulas contain only
finite matrix products, sign resolvents, and the \(u,\theta,\kappa\)
integrals, hence are finite-symbol quantities.
\end{proof}

\begin{remark}[The rank cancellation does not determine
\(\mathsf P_0\)]
\label{rem:newS7V56-rank-not-volume}
The identity \(\Theta_a(0;g,g_0)=0\) cancels the common finite Hilbert-space
rank.  It does not prove that the metric response of the microscopic
cosmological owner has the continuum volume form.  A fixed combinatorial
mesh has a metric-independent number of sites; its volume counterterm must
be fixed by the metric polarization, not inferred from rank.
\end{remark}

\subsection{The exact Ward defect}
\label{sec:newS7V56-defect}

For a continuum infinitesimal diffeomorphism
\(h_{\mu\nu}=q_\mu\xi_\nu+q_\nu\xi_\mu\), define
\begin{equation}
 \mathcal W_a^{\nu,\rho\sigma}(q)
 :=q_\mu\Pi_a^{\mu\nu,\rho\sigma}(q)
   -\mathcal C_a^{\nu,\rho\sigma}(q),                  \label{eq:newS7V56-W}
\end{equation}
where \(\mathcal C_a\) is the metric contact dictated by differentiating
the one-point response.  This is the lattice analogue of assembling the
seagull/contact term before applying the Ward identity.

\begin{lemma}[What exact unitary covariance proves]
\label{lem:newS7V56-liftable}
If the metric wave is generated by a liftable simplicial automorphism,
\(\mathcal W_a(q)=0\) exactly.  For a general continuum metric wave,
V55 proves only that \(\mathcal W_a\) is supported on a shrinking physical
neighbourhood; it gives no small bound on its local coefficient.
\end{lemma}

\begin{proof}
The first assertion is \Cref{thm:newS7V55-Ward}, with the complete contact
included.  The second follows from
\Cref{thm:newS7V55-mic-local}: off-diagonal contributions vanish, but the
diagonal trace may remain order one or power divergent.  Locality alone
does not impose rotational or diffeomorphism tensor identities on that
coefficient.
\end{proof}

\begin{definition}[Discrete microscopic Ward card, DMWC]
\label{def:newS7V56-DMWC}
For every regular bulk and stratum supercell, DMWC requires:
\begin{align}
 q_\mu\mathsf P_0^{\mu\nu,\rho\sigma}
 &=\mathsf C_0^{\nu,\rho\sigma}(q),                    \label{eq:newS7V56-DMWC0}\\
 q_\mu\mathsf P_2^{\mu\nu,\rho\sigma}(q)
 &=\mathsf C_2^{\nu,\rho\sigma}(q),                    \label{eq:newS7V56-DMWC2}\\
 q_\mu\mathsf P_4^{\mu\nu,\rho\sigma}(q)
 &=\mathsf C_4^{\nu,\rho\sigma}(q),                    \label{eq:newS7V56-DMWC4}
\end{align}
after subtracting respectively the cosmological, Einstein--Hilbert, and
V54 curvature-square contacts, with one uniform remainder bound through
two metric variations.  It also requires agreement of these tensors under
the declared face identifications.
\end{definition}

\begin{theorem}[DMWC is a finite certified acquisition]
\label{thm:newS7V56-finite}
Once the mesh incidence data card MIDC of V41 is instantiated, every
identity in DMWC reduces to equality or a strict interval enclosure for
finitely many integrals of analytic matrix functions on compact
Brillouin tori.  A successful enclosure proves that no independent
mesh-anisotropic scalar of engineering dimension at most four survives
the microscopic metric Hessian.
\end{theorem}

\begin{proof}
MIDC fixes the finite simplex records, supercell translations, coframes,
and face identifications.  V41 reconstructs every Dirac and stiffness
matrix from those records.  Differentiating the V26 half-density,
transport, and sign-resolvent formulas twice gives the finite matrices in
\Cref{lem:newS7V56-expansion}.  Compactness of the Brillouin torus and the
uniform Wilson gap make all derivatives bounded and permit convergent
interval quadrature.  There are finitely many type and contact indices.

If \eqref{eq:newS7V56-DMWC0}--\eqref{eq:newS7V56-DMWC4} hold, the
polarization annihilates every infinitesimal diffeomorphism after the
standard invariant contacts are removed.  The parity-even local scalar
classification in four dimensions then leaves only
\[
 1,\quad R,\quad R^2,\quad |{\rm Ric}|^2,\quad
 |{\rm Riem}|^2,\quad\Delta R,
\]
whose weight-four combination is fixed by V54.  Thus no independent
supercell tensor remains.
\end{proof}

\begin{proposition}[Why DMWC cannot be certified from the present sources]
\label{prop:newS7V56-missing-data}
The attached manuscripts do not provide the machine-readable ordered base
triangulation, affine vertex coordinates, or complete supercell incidence
tables required by MIDC.  Therefore the finite integrals in DMWC are
well-defined algorithms but cannot be evaluated or interval-certified
without adding geometric data not present in the programme.
\end{proposition}

\begin{proof}
This is the same data boundary proved in V41: the general edgewise
subdivision rule does not determine the ordered finite base complex or its
affine reference maps.  Those entries change the frozen matrices and hence
the integrands in \eqref{eq:newS7V56-Pi}.  Choosing them here would define
a new regulator rather than prove a statement about the attached one.
\end{proof}

\subsection{Conditional closure of the microscopic row}
\label{sec:newS7V56-conditional}

\begin{theorem}[DMWC implies MMSC]
\label{thm:newS7V56-implication}
Assume MIDC and DMWC with uniform certified margins.  Then the microscopic
metric subtraction card MMSC holds through two metric contacts.  Combined
with V53--V54, rows (i)--(iv) of MATC hold for the complete parity-even
smooth Standard Model block on the stratified JBMK.
\end{theorem}

\begin{proof}
By \Cref{thm:newS7V55-mic-local}, the microscopic response is local on a
shrinking physical neighbourhood.  Rescaling and
\Cref{lem:newS7V56-expansion} show that only engineering dimensions
zero, two, and four can fail to vanish.  DMWC identifies those terms with
the declared invariant contacts and excludes a mesh-spurion remainder.
Terms of dimension at least six carry \(a^2\) and vanish after the local
Taylor remainder estimate; the \(u\)-tail at
\(\varepsilon_{\rm jb}/a^2=a^{-3/2}\) matches the V53 high/low parametrix.
This proves MMSC.  V53 supplies MATC rows (i)--(iii), and MMSC is row (iv).
\end{proof}

\begin{assessment}[A finite data bottleneck, not a hidden theorem]
\label{ass:newS7V56-status}
The nonlinear microscopic metric problem has been reduced to the finite
supercell Ward card DMWC.  Exact gauge/frame covariance and liftable mesh
symmetries are insufficient to prove it for arbitrary continuum
diffeomorphisms.  The missing input is concrete MIDC geometry followed by
certified finite-symbol integration.  No mesh-anisotropy cancellation is
assumed.
\end{assessment}

\begin{programme}[V57: regulator-independent alternative]
\label{prog:newS7V56-V57}
Because MIDC is absent, V57 will test an upstream alternative: define the
microscopic metric counterterm by covariant continuum matching at the
common physical first jet and prove that the difference between any two
regular supercells is a finite local scheme change.  The objective is to
show that possible DMWC defects can be absorbed into explicitly declared
local mesh counterterms without changing the V54 logarithmic owner or its
Bach stress.  If such a universality theorem fails, MIDC will remain an
irreducible input to the continuum construction.
\end{programme}

\begin{firewall}[Claim boundary after seventh-series V56]
\label{fire:newS7V56-boundary}
V56 corrects the scope of the finite-cutoff diffeomorphism Ward identity,
defines the complete microscopic polarization and its contact Ward defect,
proves its even finite-symbol expansion, formulates DMWC, proves its finite
certifiability from MIDC, and proves DMWC sufficient for MMSC and MATC
row (iv).

V56 does not instantiate MIDC, certify DMWC, exclude mesh-anisotropic
microscopic counterterms, prove the full common Einstein stress, control
rough metric or AOP fields, prove measure concentration, construct the
determinant-line phase, remove the interacting cutoff, or construct the
Standard Model--Einstein continuum.
\end{firewall}

\begin{theorem}[Seventh-series V56 result]
\label{thm:newS7V56-result}
The missing microscopic metric subtraction is equivalent, after V55
locality and contact assembly, to a finite supercell Ward test through
external momentum degree four.  With instantiated MIDC, this test is a
finite family of certifiable compact Brillouin integrals; if it holds, the
only ultraviolet metric owners are the cosmological,
Einstein--Hilbert, and V54 Euler--Weyl terms and MATC row (iv) follows.
The current sources do not contain enough base-mesh data to certify that
test.
\end{theorem}

\begin{proof}
Use \Cref{corr:newS7V56-Ward-scope,
lem:newS7V56-expansion,
lem:newS7V56-liftable,
thm:newS7V56-finite,
prop:newS7V56-missing-data,
thm:newS7V56-implication}.
\end{proof}

\paragraph{Parent-version record.}
V56 was formed from
\texttt{Renewal\_SM\_Einstein\_Continuum\_Research\_Notes\_V55.tex}.
The V55 parent had \(22\,194\) logical source lines, \(891\,965\) bytes,
and SHA--256
\[
 \texttt{60634700B3F7F53EFCCFD013F91AA100572506F74C0316AD2768C97F3354A19F}.
\]
The next compulsory companion audit is V60.

% =============================================================================
% END APPENDED SEVENTH-SERIES V56 MATERIAL
% =============================================================================

% =============================================================================
% BEGIN APPENDED SEVENTH-SERIES V57 MATERIAL
% =============================================================================

\section{Seventh-series V57: local scheme equivalence without MIDC}
\label{sec:v7v57}
\label{sec:newS7V57}

\subsection{What can be universal without evaluating a supercell}
\label{sec:newS7V57-purpose}

V56 shows that a fixed simplicial regulator can possess local
mesh-anisotropic metric contacts.  It does not follow that the continuum
construction requires those contacts to vanish.  Renormalization may
subtract regulator-dependent local terms, provided three facts are proved:
the list is finite through engineering dimension four, its coefficients
are defined from the regulator rather than guessed, and the subtraction
does not alter the universal logarithmic owner.  This note proves those
facts for the fixed-range JBMK class.

\begin{definition}[Admissible JBMK regulator]
\label{def:newS7V57-admissible}
An admissible regulator \(\mathcal R\) is a finite-supercell JBMK family
with:
\begin{enumerate}
\item the common physical mass and first Dirac jet of V51;
\item a fixed exponential symbol strip, uniform Wilson gap, and unique
overlap cone as in V52--V53;
\item exact spin-frame and Standard Model gauge covariance, reflection
pairing, and two bounded metric contacts; and
\item a finite set of normalized cell shapes in a compact
shape-regular class.
\end{enumerate}
No continuum diffeomorphism Ward identity is imposed at finite \(a\).
\end{definition}

\subsection{A logarithmic matching scale}
\label{sec:newS7V57-scale}

For \(0<a<e^{-1}\), put
\begin{equation}
 K_a:=|\log a|^4,\qquad
 \tau_a:=a^2K_a.                                       \label{eq:newS7V57-scale}
\end{equation}
Then
\begin{equation}
 K_a\to\infty,\qquad \tau_a\to0,\qquad
 {a\over\sqrt{\tau_a}}=|\log a|^{-2}\to0,\qquad
 a\sqrt{K_a}=a|\log a|^2\to0.                          \label{eq:newS7V57-scale-limits}
\end{equation}

\begin{definition}[Core and annular microscopic pieces]
\label{def:newS7V57-split}
For the common-base trace \eqref{eq:newS7V55-Theta}, define
\begin{align}
 \mathcal M_a^{\rm core,\mathcal R}(g,g_0)
 &:={1\over2}\int_0^{\tau_a}
 {\Theta_a^\mathcal R(t;g,g_0)\over t}\,dt,             \label{eq:newS7V57-core}\\
 \mathcal M_a^{\rm ann,\mathcal R}(g,g_0)
 &:={1\over2}\int_{\tau_a}^{\varepsilon_{\rm jb}(a)}
 {\Theta_a^\mathcal R(t;g,g_0)\over t}\,dt.             \label{eq:newS7V57-annulus}
\end{align}
\end{definition}

\begin{lemma}[The annulus is regulator independent]
\label{lem:newS7V57-annulus}
Let \(\mathcal R_1,\mathcal R_2\) be two admissible regulators with the
same physical fields and first Dirac jet.  After the common continuum
\(A_0,A_2,A_4\) owners are inserted, their annular difference satisfies,
through two metric contacts,
\begin{equation}
 \left|\partial_g^j
 \{\mathcal M_a^{\rm ann,\mathcal R_1}
  -\mathcal M_a^{\rm ann,\mathcal R_2}\}\right|
 \le {C_j\over\sqrt{K_a}}+o(1)
 =O(|\log a|^{-2})+o(1).                               \label{eq:newS7V57-annulus-bound}
\end{equation}
\end{lemma}

\begin{proof}
On \(t\ge\tau_a\), each regulator lies in the low-momentum overlap window
because \(t/a^2\ge K_a\to\infty\).  Apply the V53 parametrix to each one.
Their retained continuum coefficients agree: the first jet fixes the
Dirac principal operator, V23 and V29 fix the matter owners, and V54 fixes
the pure-gravity owner.  The leading lattice remainder integrates as
\[
 \int_{\tau_a}^{\varepsilon_{\rm jb}}
 {a\over\sqrt t}{dt\over t}
 \le {2a\over\sqrt{\tau_a}}={2\over\sqrt{K_a}}.
\]
The differentiated collar terms tend to zero by V52; the ordinary
\(\sqrt t\) remainder is \(o(1)\) as the upper endpoint tends to zero.
High Brillouin branches are exponentially small by V53.  Subtract the two
estimates and use \eqref{eq:newS7V57-scale-limits}.
\end{proof}

\subsection{Local expansion of the microscopic core}
\label{sec:newS7V57-core}

At \(t=a^2u\), the core has \(0<u<K_a\).  Its heat radius is
\(O(\sqrt u)\) lattice steps and hence at most
\(O(a\sqrt{K_a})\) in physical units.  Fix a normal coordinate chart at
\(x\), parallel transport every fibre to \(x\), and Taylor expand the
metric, gauge, and Higgs data there.

\begin{definition}[Mesh-spurion jet space]
\label{def:newS7V57-spurions}
For a normalized supercell type \(\tau\), let
\(\mathfrak J_{\tau,\le4}\) be the finite-dimensional vector space of
parity-even, spin-frame and gauge-invariant scalar contractions formed
from:
\begin{enumerate}
\item physical background jets of total engineering dimension at most
four; and
\item the finitely many normalized edge and cell moment tensors of
\(\tau\).
\end{enumerate}
The covariant subspace contains
\[
 1,\ R,\ |F|^2,\ |DH|^2,\ |H|^4,\ R|H|^2,\
 R^2,\ |{\rm Ric}|^2,\ |{\rm Riem}|^2,\ \Delta R.
\]
The complementary vectors are called mesh spurions.
\end{definition}

\begin{lemma}[Finite-dimensionality]
\label{lem:newS7V57-finite-jets}
\(\mathfrak J_{\tau,\le4}\) is finite dimensional, uniformly over the
finite type set.  Reflection symmetry removes every odd engineering
dimension.
\end{lemma}

\begin{proof}
Only derivatives of total order at most four occur.  At each order there
are finitely many tensor indices, finite spin--internal rank, and finitely
many edge directions.  Tensor contraction therefore spans a finite vector
space.  The reflected supercell changes every odd displacement moment by
its negative while leaving the trace invariant, so odd terms cancel.
\end{proof}

\begin{theorem}[Microscopic core locality and power counting]
\label{thm:newS7V57-core-expansion}
For every admissible regulator and bounded \(C^6\) background family there
are regulator-defined coefficient functions
\(c_{\tau,J}(a)\), \(J\in\mathfrak J_{\tau,\le4}\), such that
\begin{equation}
 \mathcal M_a^{\rm core,\mathcal R}(g,g_0)
 =\sum_{\tau}\sum_{J\in\mathfrak J_{\tau,\le4}}
 c_{\tau,J}(a)
 \int_M\{J(g,A,H;\tau)-J(g_0,A,H;\tau)\}\,d\mu
 +\mathcal E_a^\mathcal R.                             \label{eq:newS7V57-core-expansion}
\end{equation}
The coefficients are exact convergent finite-symbol heat integrals on
\(0<u<K_a\), and
\begin{equation}
 |\partial_g^j\mathcal E_a^\mathcal R|
 \le C_j a^2(1+K_a)^{N_j}\longrightarrow0,
 \qquad j=0,1,2.                                       \label{eq:newS7V57-core-error}
\end{equation}
for fixed integers \(N_j\).
\end{theorem}

\begin{proof}
Use the exact first and second contact formula of V55 before expanding.
After \(t=a^2u\), Bloch transformation on each frozen supercell expresses
the local trace as an integral of finite matrix heat words.  The Wilson
gap and exponential strip make every background derivative integrable
uniformly for \(0<u<K_a\).

Within a heat word, all vertices lie with exponentially small exception
inside \(C\sqrt u\) lattice steps of \(x\).  Taylor expand the transported
background data through physical derivative order four.  Every retained
coefficient is a compact Brillouin integral followed by a convergent
finite \(u\)- and Duhamel-simplex integral.  Spin-frame and gauge
covariance and reflection put the resulting scalar in
\(\mathfrak J_{\tau,\le4}\).

The first omitted parity-even term has engineering dimension six and
therefore carries \(a^2\).  Gaussian moments and differentiated Duhamel
words bound its \(u\)-dependence by a fixed polynomial
\((1+u)^{N_j}\).  Integration to \(K_a\) proves
\eqref{eq:newS7V57-core-error}.  Since every power of
\(|\log a|\) is \(o(a^{-\epsilon})\), its right side tends to zero.
The same proof applies after two metric contacts.
\end{proof}

\begin{remark}[No numerical mesh data are needed for the definition]
\label{rem:newS7V57-no-MIDC}
MIDC is required to evaluate the numbers \(c_{\tau,J}(a)\).  It is not
required to define them: they are the Taylor coefficients of the exact
finite-\(a\) core functional for whichever admissible base mesh is chosen.
Thus they may serve as nonperturbative renormalization prescriptions.
\end{remark}

\subsection{Declared mesh counterterms}
\label{sec:newS7V57-counterterms}

\begin{definition}[Minimal mesh counterterm]
\label{def:newS7V57-counterterm}
Let
\begin{equation}
 \mathcal C_a^{\rm mesh,\mathcal R}(g,g_0)
 :=\sum_\tau\sum_{J\in\mathfrak J_{\tau,\le4}}
 c_{\tau,J}(a)
 \int_M\{J(g,A,H;\tau)-J(g_0,A,H;\tau)\}\,d\mu          \label{eq:newS7V57-counterterm}
\end{equation}
with the following renormalization conditions:
\begin{enumerate}
\item on the covariant subspace, its finite parts are fixed by the chosen
cosmological, Einstein--Hilbert, gauge, Higgs, and V54 curvature-square
normalizations; and
\item on the mesh-spurion complement, the renormalized coefficient is set
to zero.
\end{enumerate}
\end{definition}

\begin{corollary}[Core removal]
\label{cor:newS7V57-core-removal}
For every admissible regulator,
\begin{equation}
 \partial_g^j
 \{\mathcal M_a^{\rm core,\mathcal R}
   -\mathcal C_a^{\rm mesh,\mathcal R}\}
 \longrightarrow0,\qquad j=0,1,2.                     \label{eq:newS7V57-core-removal}
\end{equation}
\end{corollary}

\begin{proof}
This is \eqref{eq:newS7V57-core-error}.
\end{proof}

\begin{theorem}[Regulator scheme-equivalence theorem]
\label{thm:newS7V57-scheme}
Let \(\mathcal R_1,\mathcal R_2\) be any two admissible JBMK regulators
with the same continuum fields and physical first jet.  After subtracting
their minimal mesh counterterms,
\begin{equation}
 \partial_g^j\Big[
 \mathcal M_a^{\mathcal R_1}
 -\mathcal C_a^{\rm mesh,\mathcal R_1}
 -\mathcal M_a^{\mathcal R_2}
 +\mathcal C_a^{\rm mesh,\mathcal R_2}
 \Big]\longrightarrow0,\qquad j=0,1,2.                \label{eq:newS7V57-scheme}
\end{equation}
Thus the microscopic difference between regular supercell regulators is
a finite local scheme change through engineering dimension four.
\end{theorem}

\begin{proof}
Split each microscopic functional by
\Cref{def:newS7V57-split}.  The subtracted core difference tends to zero
by \Cref{cor:newS7V57-core-removal}.  The annular difference tends to zero
by \Cref{lem:newS7V57-annulus}.  Add the estimates.
\end{proof}

\subsection{Why the V54 logarithmic owner is unchanged}
\label{sec:newS7V57-log}

\begin{theorem}[Universality of the logarithmic coefficient]
\label{thm:newS7V57-log}
The coefficient of \(|\log a|\) in the renormalized parity-even
determinant modulus of every admissible regulator is the covariant V54
coefficient.  In particular its pure-gravity metric-active part is
\begin{equation}
 -{N_W^\bullet\over80(4\pi)^2}
 \int_M|C|^2\,d\mu_g,                                  \label{eq:newS7V57-log-C2}
\end{equation}
with heat-owner gradient
\(N_W^\bullet B_{\mu\nu}/\{40(4\pi)^2\}\).
Mesh-spurion counterterms change only regulator-dependent power
coefficients and finite or \(o(|\log a|)\) scheme terms.
\end{theorem}

\begin{proof}
On the matching annulus
\(\tau_a\le t\le t_0\), V53--V54 identify the heat coefficient with the
continuum one.  Its proper-time integral contains
\[
 \log{t_0\over\tau_a}
 =2|\log a|-4\log|\log a|+O(1).
\]
The regulator difference on this annulus is
\(O(|\log a|^{-2})+o(1)\) by
\Cref{lem:newS7V57-annulus}, so it cannot change the coefficient of
\(|\log a|\).

The rescaled core has \(0<u<K_a\).  Its dimension-four coefficients may
grow at most logarithmically in \(K_a\), hence as
\(O(\log|\log a|)\), by the heat-word power count.  Lower-dimensional
terms are power counterterms.  Neither can change the
\(|\log a|\) coefficient after the subtraction
\eqref{eq:newS7V57-counterterm}.  Insert the determinant-modulus V54
coefficient and its Bach variation.
\end{proof}

\begin{corollary}[MIDC is not a logical existence prerequisite]
\label{cor:newS7V57-MIDC}
DMWC remains a useful test for whether a chosen mesh needs no
noncovariant counterterms.  It is not necessary for regulator-independent
existence of the renormalized smooth determinant modulus: the exact
coefficients \eqref{eq:newS7V57-counterterm} subtract any local DMWC
defect, and \Cref{thm:newS7V57-scheme} identifies the common limit.
\end{corollary}

\begin{proof}
If DMWC holds, every mesh-spurion coefficient vanishes after the standard
contacts.  If it fails, the complementary terms are still members of the
finite space \(\mathfrak J_{\tau,\le4}\) and are removed by the second
renormalization condition.  Scheme equivalence then applies.
\end{proof}

\subsection{Closure of the microscopic parity-even smooth card}
\label{sec:newS7V57-closure}

\begin{theorem}[Renormalized MMSC and MATC row (iv)]
\label{thm:newS7V57-MMSC}
With the declared minimal mesh counterterm, MMSC holds for every
admissible JBMK regulator through two metric contacts.  Consequently MATC
row (iv) holds for the complete parity-even smooth Standard Model
determinant modulus, and rows (i)--(iv) have a regulator-independent
continuum limit.
\end{theorem}

\begin{proof}
The core is exhausted by
\Cref{thm:newS7V57-core-expansion,cor:newS7V57-core-removal}.  The
annulus matches the continuum by
\Cref{lem:newS7V57-annulus}, and its logarithmic owner is
\Cref{thm:newS7V57-log}.  Exact contact assembly and microscopic locality
are V55.  These are precisely the MMSC rows and MATC row (iv).  MATC
rows (i)--(iii) are V53--V54.
\end{proof}

\begin{assessment}[The mesh-data obstruction becomes a scheme choice]
\label{ass:newS7V57-status}
The absence of MIDC prevents numerical evaluation of a particular
regulator's microscopic coefficients, but no longer blocks the existence
theorem for the smooth parity-even determinant modulus.  Every supercell
difference is a declared finite local scheme change, and the V54
logarithmic coefficient and Bach stress are universal.  The next
mathematical obstruction is no longer smooth heat transfer; it is control
of rough metric, gauge, and Higgs configurations under the interacting
cutoff measure.
\end{assessment}

\begin{programme}[V58: the common renormalized Einstein heat stress]
\label{prog:newS7V57-V58}
V58 will assemble the cosmological, Einstein--Hilbert, gauge, Higgs,
nonminimal, and Bach metric gradients in the determinant-modulus
convention, state their common weak stress functional, and prove
regulator-independent convergence on bounded smooth families through the
second contact.  It will then identify the precise coercive/no-escape
estimate needed to extend that stress from smooth families to the rough
AOP support of the interacting measures.
\end{programme}

\begin{firewall}[Claim boundary after seventh-series V57]
\label{fire:newS7V57-boundary}
V57 proves a logarithmic core/annulus split, a finite local microscopic
jet expansion through dimension four, the exact minimal mesh counterterm,
scheme equivalence of all admissible JBMK regulators, universality of the
V54 logarithmic owner and Bach stress, and the renormalized smooth MMSC
and MATC row (iv).

V57 does not numerically evaluate a mesh counterterm, prove that a chosen
mesh satisfies DMWC without spurions, construct the full common
renormalized stress including its lower-owner normalization conditions,
control rough metric or AOP fields, prove measure concentration, construct
the determinant-line phase, remove the interacting cutoff, or construct
the Standard Model--Einstein continuum.
\end{firewall}

\begin{theorem}[Seventh-series V57 result]
\label{thm:newS7V57-result}
For any two admissible fixed-range stratified JBMK regulators with the
same physical first Dirac jet, their microscopic determinant-modulus
difference is, through two metric contacts, a finite local scheme change
spanned by background and mesh jets of engineering dimension at most
four.  Subtracting the exact regulator-defined mesh coefficients makes
their continuum limits agree.  The coefficient of \(|\log a|\) is
regulator independent and equals the complete V54 parity-even owner;
therefore smooth MMSC and MATC row (iv) hold without instantiating MIDC.
\end{theorem}

\begin{proof}
Use \Cref{lem:newS7V57-annulus,
lem:newS7V57-finite-jets,
thm:newS7V57-core-expansion,
cor:newS7V57-core-removal,
thm:newS7V57-scheme,
thm:newS7V57-log,
cor:newS7V57-MIDC,
thm:newS7V57-MMSC}.
\end{proof}

\paragraph{Parent-version record.}
V57 was formed from
\texttt{Renewal\_SM\_Einstein\_Continuum\_Research\_Notes\_V56.tex}.
The V56 parent had \(22\,457\) logical source lines, \(903\,787\) bytes,
and SHA--256
\[
 \texttt{7747EEE8D58E9CD49FCF58AD1AC6BCC5D83CD71755C75BAA554DD1BB22716C64}.
\]
The next compulsory companion audit is V60.

% =============================================================================
% END APPENDED SEVENTH-SERIES V57 MATERIAL
% =============================================================================

% =============================================================================
% BEGIN APPENDED SEVENTH-SERIES V58 MATERIAL
% =============================================================================

\section{Seventh-series V58: the common renormalized Einstein heat stress}
\label{sec:v7v58}
\label{sec:newS7V58}

\subsection{Metric-gradient convention}
\label{sec:newS7V58-convention}

For an integrated Euclidean functional \(\mathcal F(g,A,H)\), define its
inverse-metric gradient \(\mathcal E_{\mu\nu}\) and stress
\(\mathcal T_{\mu\nu}\) by
\begin{equation}
 \delta_g\mathcal F[h]
 =\int_M\mathcal E_{\mu\nu}h^{\mu\nu}\,d\mu_g,\qquad
 \mathcal T_{\mu\nu}:=-2\mathcal E_{\mu\nu}.            \label{eq:newS7V58-gradient}
\end{equation}
This fixes every sign below.  The fermionic functional-integral sign is
applied to the assembled effective action, not hidden in the heat owner.

Define the Standard Model gauge contractions
\begin{align}
 \mathfrak G_{\rm SM}(F)
 &:=\sum_f\tr_{R_f}(F_{\alpha\beta}^*F^{\alpha\beta}),  \label{eq:newS7V58-Gscalar}\\
 \mathfrak G_{\mu\nu}(F)
 &:=\sum_f\operatorname{Re}\tr_{R_f}
   (F_{\mu\rho}^*F_\nu{}^\rho),                        \label{eq:newS7V58-Gtensor}
\end{align}
with the same physical Weyl multiplicities as V23.  Put
\begin{equation}
 \mathfrak K_{\mu\nu}(H)
 :=\operatorname{Re}\langle D_\mu H,D_\nu H\rangle.    \label{eq:newS7V58-Ktensor}
\end{equation}

\begin{lemma}[Elementary metric variations]
\label{lem:newS7V58-variations}
On a closed four-manifold,
\begin{align}
 \delta\int d\mu_g
 &=-{1\over2}\int g_{\mu\nu}h^{\mu\nu}\,d\mu_g,         \label{eq:newS7V58-var-volume}\\
 \delta\int R\,d\mu_g
 &=\int G_{\mu\nu}h^{\mu\nu}\,d\mu_g,                  \label{eq:newS7V58-var-R}\\
 \delta\int\mathfrak G_{\rm SM}(F)\,d\mu_g
 &=\int\{2\mathfrak G_{\mu\nu}
 -\tfrac12g_{\mu\nu}\mathfrak G_{\rm SM}\}
 h^{\mu\nu}\,d\mu_g,                                  \label{eq:newS7V58-var-F}\\
 \delta\int|DH|^2\,d\mu_g
 &=\int\{\mathfrak K_{\mu\nu}
 -\tfrac12g_{\mu\nu}|DH|^2\}h^{\mu\nu}\,d\mu_g,        \label{eq:newS7V58-var-DH}\\
 \delta\int|H|^4\,d\mu_g
 &=-{1\over2}\int g_{\mu\nu}|H|^4h^{\mu\nu}\,d\mu_g,   \label{eq:newS7V58-var-H4}\\
 \delta\int R|H|^2\,d\mu_g
 &=\int\big[
 |H|^2G_{\mu\nu}
 +(g_{\mu\nu}\Delta-\nabla_\mu\nabla_\nu)|H|^2
 \big]h^{\mu\nu}\,d\mu_g.                              \label{eq:newS7V58-var-RH}
\end{align}
\end{lemma}

\begin{proof}
The first two formulas are the determinant variation of the volume form
and the Palatini identity, with the divergence discarded by Stokes'
theorem.  In \eqref{eq:newS7V58-var-F}, each of the two inverse metrics
contracting \(F_{\alpha\beta}^*F_{\rho\sigma}\) contributes one copy of
\(\mathfrak G_{\mu\nu}\), while the volume contributes the trace term.
The same calculation with one inverse metric proves
\eqref{eq:newS7V58-var-DH}; the Higgs field and connection are held fixed
during the metric contact.  Equation \eqref{eq:newS7V58-var-H4} comes only
from the volume.  Multiplying the Palatini identity by \(|H|^2\) and
integrating its two derivatives by parts proves
\eqref{eq:newS7V58-var-RH}.
\end{proof}

\subsection{The complete logarithmic heat gradient}
\label{sec:newS7V58-log}

By V54, the parity-even determinant-modulus weight-four functional,
with the common \((4\pi)^{-2}\) factor included and with topological and
integrated-divergence terms omitted, is
\begin{align}
 \mathcal F_{\log}^{\det}
 ={1\over(4\pi)^2}\int_M\Big[
 &-{N_W^\bullet\over80}|C|^2
 +{1\over6}\mathfrak G_{\rm SM}(F)
 +\mathfrak a_\bullet|DH|^2                            \notag\\
 &+\mathfrak b_\bullet|H|^4
 +{\mathfrak a_\bullet\over6}R|H|^2
 \Big]\,d\mu_g.                                        \label{eq:newS7V58-Flog}
\end{align}

\begin{theorem}[Common logarithmic metric gradient]
\label{thm:newS7V58-log-gradient}
The inverse-metric gradient of \(\mathcal F_{\log}^{\det}\) is
\begin{align}
 \mathcal E_{\mu\nu}^{\log}
 ={1\over(4\pi)^2}\Big[
 &{N_W^\bullet\over40}B_{\mu\nu}
 +{1\over6}\{2\mathfrak G_{\mu\nu}
       -\tfrac12g_{\mu\nu}\mathfrak G_{\rm SM}\}       \notag\\
 &+\mathfrak a_\bullet
   \{\mathfrak K_{\mu\nu}
       -\tfrac12g_{\mu\nu}|DH|^2\}
 -{\mathfrak b_\bullet\over2}g_{\mu\nu}|H|^4           \notag\\
 &+{\mathfrak a_\bullet\over6}
 \big\{|H|^2G_{\mu\nu}
 +(g_{\mu\nu}\Delta-\nabla_\mu\nabla_\nu)|H|^2\big\}
 \Big].                                                \label{eq:newS7V58-Elog}
\end{align}
Its stress is \(-2\mathcal E_{\mu\nu}^{\log}\) by
\eqref{eq:newS7V58-gradient}.
\end{theorem}

\begin{proof}
The Weyl-square variation is
\eqref{eq:newS7V54-Bach-variation}; multiplying by
\(-N_W^\bullet/\{80(4\pi)^2\}\) gives the first term.
Apply \Cref{lem:newS7V58-variations} to the remaining terms in
\eqref{eq:newS7V58-Flog}.  The Euler density and \(\Delta R\) integrate
to metric-independent or zero functionals on the closed manifold and
therefore contribute no gradient.
\end{proof}

\subsection{Lower owners and renormalization conditions}
\label{sec:newS7V58-lower}

Let \(\Lambda_R,\kappa_R,m_R^2\) be the chosen finite renormalized
coefficients of the volume, scalar-curvature, and Higgs-mass owners after
the exact lattice power subtractions.  Define
\begin{equation}
 \mathcal F_{\le2}^R
 :=\int_M\{\Lambda_R+\kappa_R R+m_R^2|H|^2\}\,d\mu_g.   \label{eq:newS7V58-Flower}
\end{equation}

\begin{lemma}[Lower-owner gradient]
\label{lem:newS7V58-lower-gradient}
\begin{equation}
 \mathcal E_{\mu\nu}^{\le2}
 =-\tfrac12\Lambda_Rg_{\mu\nu}
 +\kappa_RG_{\mu\nu}
 -\tfrac12m_R^2|H|^2g_{\mu\nu}.                        \label{eq:newS7V58-Elower}
\end{equation}
\end{lemma}

\begin{proof}
Use \eqref{eq:newS7V58-var-volume}--\eqref{eq:newS7V58-var-R}; the
pointwise Higgs norm is metric independent in the fixed internal Hermitian
structure.
\end{proof}

Finite renormalizations of the gauge, Higgs kinetic, quartic,
nonminimal, and Weyl-square owners have the same tensor forms as
\eqref{eq:newS7V58-Elog}.  Their coefficients are fixed by physical
normalization conditions.  Mesh-spurion finite parts are fixed to zero by
\Cref{def:newS7V57-counterterm}.

\begin{definition}[Common renormalized heat gradient]
\label{def:newS7V58-common-gradient}
\(\mathcal E_{\mu\nu}^{\rm heat,R}\) is the sum of
\eqref{eq:newS7V58-Elower}, the logarithmic gradient
\eqref{eq:newS7V58-Elog} multiplied by the chosen renormalization-scale
logarithm, and the finite covariant owner gradients fixed by the
renormalization conditions.
\end{definition}

\begin{theorem}[Regulator-independent smooth stress convergence]
\label{thm:newS7V58-smooth-convergence}
Let \(\Gamma_a^{\rm ren,\mathcal R}\) be the parity-even
determinant-modulus functional of any admissible stratified JBMK regulator,
with its V57 minimal mesh counterterm and the common finite normalization
conditions.  On every bounded \(C^6\) metric, \(C^3\) gauge, and \(C^3\)
Higgs family,
\begin{align}
 \partial_g\Gamma_a^{\rm ren,\mathcal R}[h]
 &\longrightarrow
 \int_M\mathcal E_{\mu\nu}^{\rm heat,R}h^{\mu\nu}\,d\mu_g,            \label{eq:newS7V58-first-limit}\\
 \partial_g^2\Gamma_a^{\rm ren,\mathcal R}[h,k]
 &\longrightarrow
 \partial_k\int_M
 \mathcal E_{\mu\nu}^{\rm heat,R}h^{\mu\nu}\,d\mu_g                 \label{eq:newS7V58-second-limit}
\end{align}
uniformly for \(h,k\) in bounded \(C^6\) sets.  The limits are independent
of the supercell regulator and of the stratified collar choices.
\end{theorem}

\begin{proof}
V53--V54 give the mesoscopic coefficient and remainder convergence through
two metric contacts.  V57 supplies the microscopic subtraction, proves
scheme equivalence, and fixes all mesh-spurion finite parts to zero.
V5 controls the large-time relative tail on the protected gap chart.
The surviving covariant local gradients are exactly
\Cref{thm:newS7V58-log-gradient,lem:newS7V58-lower-gradient} and their
declared finite renormalizations.  The second limit follows from the
assembled second Duhamel contact and uniform differentiated remainders,
not from separately differentiating divergent terms.
\end{proof}

\subsection{The off-shell conservation identity}
\label{sec:newS7V58-conservation}

Let \(\mathcal J^\mu\) and \(\mathcal H\) be the gauge and Higgs
Euler--Lagrange gradients of the same renormalized local functional.

\begin{theorem}[Diffeomorphism Noether identity]
\label{thm:newS7V58-Noether}
The common continuum gradients obey, distributionally,
\begin{equation}
 2\nabla^\mu\mathcal E_{\mu\nu}^{\rm heat,R}
 +\operatorname{Re}\langle\mathcal H,D_\nu H\rangle
 +\langle\mathcal J^\mu,F_{\nu\mu}\rangle
 =0,                                                   \label{eq:newS7V58-Noether}
\end{equation}
up to the conventional gauge transformation included in the Lie
derivative of \(A\).  In particular the heat stress is covariantly
conserved when the gauge and Higgs equations hold.
\end{theorem}

\begin{proof}
Apply an infinitesimal compactly supported diffeomorphism generated by
\(X\) to the scalar functional
\(\mathcal F_{\le2}^R+\mathcal F_{\log}^{\det}\) and its finite covariant
owners.  Its metric variation is
\(h^{\mu\nu}=-(\nabla^\mu X^\nu+\nabla^\nu X^\mu)\).
The gauge Lie derivative is
\(\iota_XF+D(\iota_XA)\); gauge invariance removes the last term.
The Higgs variation is \(X^\nu D_\nu H\), up to the same gauge rotation.
Integrating the metric term by parts and using arbitrary \(X\) gives
\eqref{eq:newS7V58-Noether}.  The Bach tensor is divergence free because
\(\int|C|^2\) is diffeomorphism invariant.
\end{proof}

\subsection{The natural finite-energy extension}
\label{sec:newS7V58-energy-space}

The smooth stress formula has a canonical weak interpretation.  Define
\(\mathcal X\) to consist of fields with:
\begin{equation}
\begin{split}
 &\lambda_-g_*\le g\le\lambda_+g_*,
 \qquad g\in W^{2,2}\cap W^{1,4},\\
 &F_A\in L^2,\qquad
 H\in W^{1,2}\cap L^4.                                 \label{eq:newS7V58-X}
\end{split}
\end{equation}
All spaces are local in a fixed finite atlas, with the usual compatibility
on overlaps.

\begin{lemma}[Finite-energy weak stress]
\label{lem:newS7V58-weak-stress}
For every field in \(\mathcal X\), the functional
\[
 \int_M\{|C|^2+\mathfrak G_{\rm SM}(F)+|DH|^2+|H|^4
          +R|H|^2\}\,d\mu_g
\]
is finite, and its first metric variation defines a distribution on
smooth symmetric test tensors.  On smooth fields that distribution equals
\eqref{eq:newS7V58-Elog}.
\end{lemma}

\begin{proof}
Uniform ellipticity and
\(g\in W^{2,2}\cap W^{1,4}\) give Riemann curvature in \(L^2\).
The gauge and Higgs terms are integrable by
\eqref{eq:newS7V58-X}; \(R|H|^2\in L^1\) by
\(R\in L^2\) and \(H\in L^4\).  Define the variation by differentiating
the integrated functional against a smooth \(h\); integrations by parts
place all derivatives of curvature and \(|H|^2\) on \(h\).
Every resulting pairing is \(L^1\).  Density of smooth fields and
\Cref{lem:newS7V58-variations,thm:newS7V54-Bach-owner} identify the smooth
formula.
\end{proof}

\begin{lemma}[Strong finite-energy continuity]
\label{lem:newS7V58-strong}
If \(g_n\to g\) strongly in
\(W^{2,2}\cap W^{1,4}\) with common ellipticity,
\(F_{A_n}\to F_A\) strongly in \(L^2\), and
\(H_n\to H\) strongly in \(W^{1,2}\cap L^4\), then the weak heat stresses
converge on every smooth test tensor.
\end{lemma}

\begin{proof}
The assumptions give strong \(L^2\) convergence of curvature and strong
convergence of every quadratic contraction after multiplication by the
uniformly controlled inverse metrics.  Hölder's inequality treats
\(R|H|^2\), the gauge square, Higgs kinetic term, and quartic term.
The weak variational definition in
\Cref{lem:newS7V58-weak-stress} moves derivatives onto the fixed smooth
test tensor.  Each pairing therefore converges.
\end{proof}

\subsection{The exact no-escape card}
\label{sec:newS7V58-noescape}

Let \(\boldsymbol\mu_a\) denote the proposed interacting cutoff measures
after gauge fixing or on the corresponding quotient.

\begin{definition}[Renormalized stress no-escape card, RSNC]
\label{def:newS7V58-RSNC}
RSNC consists of:
\begin{enumerate}
\item tightness of \(\boldsymbol\mu_a\) in the weak finite-energy topology
of \(\mathcal X\), with uniform metric ellipticity;
\item for every \(\epsilon>0\), a compact
\(\mathcal K_\epsilon\subset\mathcal X\) such that
\[
 \sup_a\boldsymbol\mu_a(\mathcal K_\epsilon^c)<\epsilon;
\]
\item uniform integrability of the lattice renormalized stress:
\[
 \lim_{M\to\infty}\sup_a
 \int_{\{\|\mathcal T_a^{\rm ren}\|_{(W^{2,\infty})^*}>M\}}
 \|\mathcal T_a^{\rm ren}\|_{(W^{2,\infty})^*}\,
 d\boldsymbol\mu_a=0;                                  \label{eq:newS7V58-UI}
\]
\item absence of curvature, gauge-energy, or Higgs-energy defect measures,
equivalently strong convergence in the energy quantities on every
Skorohod realization of a convergent subsequence; and
\item negligibility of the complement of the protected overlap-gap chart
for the determinant and both metric contacts.
\end{enumerate}
\end{definition}

\begin{theorem}[RSNC transfers the common stress to the measure limit]
\label{thm:newS7V58-RSNC}
Assume cutoff measures converge weakly along a subsequence and satisfy
RSNC.  Then for every smooth test tensor \(h\),
\begin{equation}
 \lim_{a\downarrow0}
 \int\langle\mathcal T_a^{\rm ren},h\rangle\,
 d\boldsymbol\mu_a
 =
 \int\langle\mathcal T^{\rm heat,R},h\rangle\,
 d\boldsymbol\mu,                                      \label{eq:newS7V58-stress-measure}
\end{equation}
where the right side is the finite-energy weak stress of
\Cref{lem:newS7V58-weak-stress}.
\end{theorem}

\begin{proof}
Use tightness and Skorohod representation on a subsequence.  The
no-defect row upgrades the energy variables to the strong convergences in
\Cref{lem:newS7V58-strong}, while the protected-chart row permits use of
\Cref{thm:newS7V58-smooth-convergence} after smooth approximation on each
compact set.  Hence the paired stresses converge in probability.
Uniform integrability \eqref{eq:newS7V58-UI} upgrades this to convergence
of expectations.  Uniqueness of the limit removes the Skorohod
representation.
\end{proof}

\begin{proposition}[Fixed AOP admissibility does not imply RSNC]
\label{prop:newS7V58-AOP-no}
A cutoff-independent plaquette bound
\(\|I-U_p\|\le\delta\) and uniform metric ellipticity do not imply any of
the energy-tightness, no-defect, or stress-uniform-integrability rows of
RSNC.
\end{proposition}

\begin{proof}
Choose admissible plaquette holonomies whose distance from the identity is
a fixed nonzero number below \(\delta\), alternating at lattice scale.
Their physical curvature is of order \(a^{-2}\), so the Riemann-sum
\(L^2\) gauge energy diverges.  Likewise, uniformly elliptic edge metrics
may alternate by a fixed small amplitude; their discrete second
differences are \(O(a^{-2})\), so curvature energy and the paired
curvature-square stress need not be tight.  These examples remain in the
fixed admissible set.  Thus a measure/action estimate, rather than the AOP
inequality alone, is necessary.
\end{proof}

\begin{assessment}[The smooth common stress is closed]
\label{ass:newS7V58-status}
The parity-even determinant-modulus metric gradient is now one explicit
regulator-independent tensor through two contacts, and it extends as a
weak distribution to the natural finite-energy space.  Passing it through
an interacting measure is exactly RSNC.  The fixed AOP bound does not
supply RSNC; the next step must use the complete bosonic action and
conditional measure, including the Euclidean conformal-direction issue.
\end{assessment}

\begin{programme}[V59: coercivity and the conformal direction]
\label{prog:newS7V58-V59}
V59 will decompose the metric into a unit-determinant part and conformal
factor, compute the signs of the renormalized Einstein--Hilbert and
Weyl-square owners on that decomposition, and test whether the declared
renewal/grand-tensor gravitational sector supplies a positive nonlinear
control strong enough for RSNC.  If the real Euclidean conformal mode
remains unbounded below, the note will isolate the smallest additional
positive gravitational owner or contour input needed rather than assume
measure tightness.
\end{programme}

\begin{firewall}[Claim boundary after seventh-series V58]
\label{fire:newS7V58-boundary}
V58 computes the complete common determinant-modulus heat gradient,
including lower owners, gauge, Higgs, nonminimal, and Bach terms; proves
regulator-independent convergence through two metric contacts on smooth
families; proves its Noether identity and finite-energy weak extension;
and formulates RSNC with a conditional transfer theorem.  It proves fixed
AOP admissibility insufficient for RSNC.

V58 does not prove RSNC for the interacting measures, control the
Euclidean conformal direction, prove measure concentration, construct the
determinant-line phase, remove the interacting cutoff, or construct the
Standard Model--Einstein continuum.
\end{firewall}

\begin{theorem}[Seventh-series V58 result]
\label{thm:newS7V58-result}
After the V57 mesh subtraction and common finite renormalization
conditions, every admissible stratified JBMK regulator has the same smooth
parity-even metric gradient \(\mathcal E_{\mu\nu}^{\rm heat,R}\).  Its
logarithmic part is the explicit tensor
\eqref{eq:newS7V58-Elog}; the first two lattice metric contacts converge
to its gradient and Hessian.  The tensor obeys the off-shell
diffeomorphism Noether identity, extends weakly to the finite-energy space
\eqref{eq:newS7V58-X}, and passes to an interacting measure limit if RSNC
holds.  Fixed AOP admissibility alone does not imply that card.
\end{theorem}

\begin{proof}
Use \Cref{lem:newS7V58-variations,
thm:newS7V58-log-gradient,
lem:newS7V58-lower-gradient,
thm:newS7V58-smooth-convergence,
thm:newS7V58-Noether,
lem:newS7V58-weak-stress,
lem:newS7V58-strong,
thm:newS7V58-RSNC,
prop:newS7V58-AOP-no}.
\end{proof}

\paragraph{Parent-version record.}
V58 was formed from
\texttt{Renewal\_SM\_Einstein\_Continuum\_Research\_Notes\_V57.tex}.
The V57 parent had \(22\,866\) logical source lines, \(920\,426\) bytes,
and SHA--256
\[
 \texttt{ED4A995EB7788DD46323CCBB32EF994FB6348664514CBD64F1A421435C3B64AF}.
\]
The next compulsory companion audit is V60.

% =============================================================================
% END APPENDED SEVENTH-SERIES V58 MATERIAL
% =============================================================================

% =============================================================================
% APPENDED SEVENTH-SERIES V59 MATERIAL
% =============================================================================

\section{The Euclidean conformal obstruction and a coercive gravity card}
\label{sec:newS7V59}
\label{sec:v7v59}

The preceding note reduced passage of the common renormalized heat stress
through the interacting limit to RSNC.  This note attacks the first genuinely
gravitational obstruction in that card.  It separates an arbitrary metric
into a volume-changing conformal factor and a unit-volume-form metric,
computes the action on that factor, and records exactly what a positive
finite-cutoff gravitational construction would still have to prove.  The
calculation is deliberately prior to any claimed continuum measure: it
prevents positivity or tightness from being imported from the formal
Einstein action.

\subsection{Conformal decomposition and exact signs}
\label{sec:newS7V59-conformal}

Let \(M\) be a closed oriented four-manifold.  Fix a smooth positive volume
form \(d\mu_0\).  Every smooth Riemannian metric can be written uniquely as

\begin{equation}
 g=e^{2\phi}\widehat g,
 \qquad
 d\mu_{\widehat g}=d\mu_0,
 \qquad
 \phi={1\over4}\log {d\mu_g\over d\mu_0}.
 \label{eq:newS7V59-split}
\end{equation}

Thus \(\widehat g\) is the unit-determinant part relative to \(d\mu_0\),
while \(\phi\) is the conformal factor.

\begin{lemma}[Four-dimensional conformal formulae]
\label{lem:newS7V59-formulae}
With the convention
\(\widehat\Delta=\widehat\nabla^\mu\widehat\nabla_\mu\),

\begin{align}
 d\mu_g&=e^{4\phi}d\mu_{\widehat g},
 \label{eq:newS7V59-volume}\\
 R_g&=e^{-2\phi}
 \bigl(\widehat R-6\widehat\Delta\phi
                 -6|\widehat\nabla\phi|^2\bigr),
 \label{eq:newS7V59-scalar}\\
 \int_M R_g\,d\mu_g
 &=\int_M e^{2\phi}
 \bigl(\widehat R+6|\widehat\nabla\phi|^2\bigr)
 \,d\mu_{\widehat g},
 \label{eq:newS7V59-EH}\\
 \int_M |C_g|_g^2\,d\mu_g
 &=\int_M |C_{\widehat g}|_{\widehat g}^2\,d\mu_{\widehat g}.
 \label{eq:newS7V59-Weyl}
\end{align}
\end{lemma}

\begin{proof}
The determinant in dimension four gives
\eqref{eq:newS7V59-volume}.  The difference tensor of the two
Levi--Civita connections is

\[
 \Gamma(g)^\rho_{\mu\nu}-\Gamma(\widehat g)^\rho_{\mu\nu}
 =\delta^\rho_\mu\phi_\nu+\delta^\rho_\nu\phi_\mu
  -\widehat g_{\mu\nu}\widehat g^{\rho\sigma}\phi_\sigma.
\]

Contracting the curvature difference gives
\eqref{eq:newS7V59-scalar}.  On the closed manifold,

\[
 -6\int e^{2\phi}\widehat\Delta\phi\,d\mu_{\widehat g}
 =12\int e^{2\phi}|\widehat\nabla\phi|^2\,d\mu_{\widehat g},
\]

which combines with the last term of
\eqref{eq:newS7V59-scalar} to prove \eqref{eq:newS7V59-EH}.
The Weyl tensor with one index raised is conformally invariant; lowering
and contracting its indices contributes \(e^{4\phi}e^{-8\phi}\), which is
cancelled by the volume factor.  This proves
\eqref{eq:newS7V59-Weyl}.
\end{proof}

\begin{theorem}[Uniformly elliptic conformal instability]
\label{thm:newS7V59-instability}
Let \(\kappa>0\).  On a flat four-torus the Euclidean action

\begin{equation}
 S_E(g)=-\kappa\int_M R_g\,d\mu_g
       +\Lambda\int_Md\mu_g
       +\beta\int_M|C_g|^2\,d\mu_g
 \label{eq:newS7V59-SE}
\end{equation}

is unbounded below, for every real \(\Lambda,\beta\), even on a family
having one common ellipticity constant and bounded volume.  The same is
true after adding any zeroth-order local density which stays uniformly
bounded on bounded conformal factors.
\end{theorem}

\begin{proof}
In periodic flat coordinates set
\(\phi_N=\varepsilon\sin(Nx^1)\), where
\(0<\varepsilon<1\) is fixed and \(N\) is a positive integer, and put
\(g_N=e^{2\phi_N}\delta\).  All eigenvalues of \(g_N\) lie between
\(e^{-2\varepsilon}\) and \(e^{2\varepsilon}\).  Its volume is independent
of \(N\), and \(C_{g_N}=0\).  By \eqref{eq:newS7V59-EH},

\[
 \int_M R_{g_N}\,d\mu_{g_N}
 =6\varepsilon^2N^2
   \int_M e^{2\varepsilon\sin(Nx^1)}\cos^2(Nx^1)\,d\mu_\delta
 =c_\varepsilon N^2
\]

for a constant \(c_\varepsilon>0\) independent of \(N\).  The Einstein
term therefore tends to \(-\infty\), while the cosmological, Weyl, and
stated zeroth-order terms remain bounded.
\end{proof}

\begin{corollary}[The determinant Weyl owner cannot cure the conformal mode]
\label{cor:newS7V59-Weyl-no}
No choice of coefficient multiplying the \(C^2\) owner computed in V54--V58
controls the real Euclidean conformal direction.  In particular, the
universal determinant-modulus logarithm does not imply gravitational
measure tightness.
\end{corollary}

\begin{proof}
Equation \eqref{eq:newS7V59-Weyl} makes the owner independent of \(\phi\),
and it vanishes identically on the sequence in
\Cref{thm:newS7V59-instability}.
\end{proof}

\subsection{The first local coercive owner}
\label{sec:newS7V59-coercive}

The preceding sequence also identifies the derivative order which is
missing.  Bounded functions of the conformal factor are \(O(1)\), the
Einstein term is \(-O(N^2)\), and

\begin{equation}
 \int_M R_{g_N}^2\,d\mu_{g_N}
 =\int_M
 \bigl(-6\Delta\phi_N-6|\nabla\phi_N|^2\bigr)^2\,d\mu_\delta
 =O(N^4)
 \label{eq:newS7V59-R2scale}
\end{equation}

with a strictly positive \(N^4\) leading coefficient.  Hence a positive
\(R^2\) owner is the smallest scalar curvature-square addition that blocks
this high-frequency conformal escape.  Control of the nonconformal
curvature additionally requires a positive Weyl-square coefficient (or an
equivalent traceless-Ricci owner).

\begin{theorem}[Deterministic curvature--matter coercivity]
\label{thm:newS7V59-coercivity}
Let \(\alpha,\beta>0\), let \(\kappa,\xi,q,\Lambda\in\mathbb R\), and let
\(\lambda>0\).  Consider the local part

\begin{align}
 S_{\rm loc}(g,A,H)=\int_M\bigl\{
 &\alpha R^2+\beta|C|^2-\kappa R+\Lambda
  +c_F|F_A|^2+c_H|D_AH|^2 \notag\\
 &+q|H|^2+\lambda|H|^4+\xi R|H|^2
 \bigr\}\,d\mu_g,
 \label{eq:newS7V59-local-action}
\end{align}

where \(c_F,c_H>0\).  Choose
\(\epsilon_0,\epsilon_H,\eta>0\) so that

\begin{equation}
 c_R:=\alpha-\epsilon_0-\epsilon_H>0,
 \qquad
 c_4:=\lambda-{\xi^2\over4\epsilon_H}-\eta>0.
 \label{eq:newS7V59-coeff-condition}
\end{equation}

Writing \(q_-:=\max\{0,-q\}\), one has

\begin{align}
 S_{\rm loc}\geq{}&
 c_R\|R\|_{L^2(g)}^2+\beta\|C\|_{L^2(g)}^2
 +c_F\|F_A\|_{L^2(g)}^2+c_H\|D_AH\|_{L^2(g)}^2
 +c_4\|H\|_{L^4(g)}^4 \notag\\
 &-\left(|\Lambda|+{\kappa^2\over4\epsilon_0}
                  +{q_-^2\over4\eta}\right)\operatorname{Vol}_g(M).
 \label{eq:newS7V59-coercive-bound}
\end{align}
\end{theorem}

\begin{proof}
Pointwise Young inequalities give

\[
 -\kappa R\geq-\epsilon_0R^2-{\kappa^2\over4\epsilon_0},
 \qquad
 \xi R|H|^2\geq-\epsilon_HR^2
                    -{\xi^2\over4\epsilon_H}|H|^4,
\]

and
\(q|H|^2\geq-\eta|H|^4-q_-^2/(4\eta)\).
Also \(\Lambda\geq-|\Lambda|\).  Integration and
\eqref{eq:newS7V59-coeff-condition} prove the claim.
\end{proof}

\begin{lemma}[Full curvature control modulo topology]
\label{lem:newS7V59-full-curvature}
For every smooth metric on a closed oriented four-manifold,

\begin{equation}
 \int_M|\operatorname{Rm}|^2\,d\mu_g
 =2\int_M|C|^2\,d\mu_g
  +{1\over3}\int_MR^2\,d\mu_g-32\pi^2\chi(M).
 \label{eq:newS7V59-GB-control}
\end{equation}

Consequently, the first two positive terms in
\eqref{eq:newS7V59-coercive-bound} control the complete \(L^2\) curvature,
up to a fixed additive topological constant.
\end{lemma}

\begin{proof}
In four dimensions, with
\(\operatorname{Ric}_0=\operatorname{Ric}-{1\over4}Rg\),

\[
 |\operatorname{Rm}|^2=|C|^2+2|\operatorname{Ric}_0|^2+{1\over6}R^2,
 \quad
 32\pi^2\chi(M)=\int_M
 \left(|C|^2-2|\operatorname{Ric}_0|^2+{1\over6}R^2\right)d\mu_g.
\]

Eliminating \(\|\operatorname{Ric}_0\|_2^2\) gives
\eqref{eq:newS7V59-GB-control}.
\end{proof}

The theorem is a deterministic energy estimate, not a construction of a
measure.  In particular, a curvature-square action has higher derivatives;
the displayed lower bound neither supplies a reflection-positive lattice
realization nor proves that unwanted higher-derivative modes decouple from
the Lorentzian theory.

\subsection{Audit of the supplied gravitational sources}
\label{sec:newS7V59-source-audit}

The local source audit used the following exact files:

\begin{itemize}
\item
\(\texttt{Renewal\_Grand\_Tensor\_Monograph\_V067.tex}\),
\(2\,804\,854\) bytes, SHA--256

\[
\texttt{FACD058EC8BB210AD8B296C9D7F4D78A229D08D7A2AE378808700655CB8BC0FE};
\]

\item the locally newest continuation of the physical-source note,
\(\texttt{Renewal\_Einstein\_Standard\_Model\_Physical\_Source\_Metric\_}\)
\(\texttt{Duhamel\_Ancestry\_and\_Quantum\_Continuum\_Programme\_V35.tex}\),
\(1\,410\,598\) bytes, SHA--256

\[
\texttt{AE1FBC0A8125C37465E38E7649E2992725F4579A86F0BB17483B9AE829CFB17C}.
\]
\end{itemize}

The Grand Tensor manuscript explicitly leaves the real Euclidean conformal
direction, the nonlinear global Einstein measure, and Lorentzian
reconstruction open.  Its positive master-state branches are conditional
on supplied actions; its structural positive-form machinery does not select
a local gravitational action or conformal contour.  The physical-source
note develops metric Duhamel response and reflection-positivity firewalls,
but supplies no coercive conformal action or contour either.  Thus neither
source discharges \Cref{thm:newS7V59-instability}; importing such a discharge
would repeat a conditional statement as a proof.

\subsection{A finite-cutoff conformal gravity measure card}
\label{sec:newS7V59-card}

\begin{definition}[Conformal gravity measure card, CGMC]
\label{def:newS7V59-CGMC}
A proposed lattice gravitational sector satisfies CGMC if, after the stated
gauge fixing or quotient, it supplies all of the following with constants
independent of the mesh \(a\):

\begin{enumerate}
\item a finite, normalized, positive cutoff probability measure
\(d\boldsymbol\mu_a=Z_a^{-1}e^{-S_a}d\rho_a\), together with the required
reflection-positive involution;
\item a nonnegative reconstructed energy \(\mathcal E_a\) containing
\(\|R_a\|_2^2+\|C_a\|_2^2+\|F_a\|_2^2+
\|D_aH_a\|_2^2+\|H_a\|_4^4\), volume and ellipticity control, and an action
bound \(S_a\geq c\mathcal E_a-C\);
\item an entropy estimate strong enough to upgrade that action bound to
an exponential probability moment

\[
 \sup_a\int e^{\theta\mathcal E_a}\,d\boldsymbol\mu_a<\infty
 \quad\hbox{for some }\theta>0;
 \label{eq:newS7V59-exp-moment}
\]

\item compactness of bounded \(\mathcal E_a\)-sublevels in the weak
finite-energy topology and reconstruction with no curvature, gauge, or
Higgs energy defect;
\item a polynomial stress bound

\[
 \|\mathcal T_a^{\rm ren}\|_{(W^{2,\infty})^*}
 \leq C(1+\mathcal E_a^p)
 \label{eq:newS7V59-stress-bound}
\]

for fixed \(p\), and vanishing probability of the complement of the
protected overlap-gap chart; and
\item a declared low-energy normalization in which the Einstein coefficient
has its physical sign and any curvature-square modes either remain part of
the continuum theory or are proved to decouple, without losing the preceding
rows.
\end{enumerate}
\end{definition}

\begin{proposition}[CGMC implies the analytic rows of RSNC]
\label{prop:newS7V59-CGMC-RSNC}
CGMC implies tightness, compact capture, stress uniform integrability, the
no-defect row, and protected-chart negligibility in
\Cref{def:newS7V58-RSNC}.  Hence, together with the regulator convergence of
V57--V58, it implies the expectation convergence
\eqref{eq:newS7V58-stress-measure} along every weakly convergent subsequence.
\end{proposition}

\begin{proof}
Markov's inequality applied to \eqref{eq:newS7V59-exp-moment} gives

\[
 \sup_a\boldsymbol\mu_a\{\mathcal E_a>L\}
 \leq e^{-\theta L}
 \sup_a\int e^{\theta\mathcal E_a}\,d\boldsymbol\mu_a.
\]

Compactness of the sublevels therefore gives tightness and compact capture.
For every \(p<\infty\), the same exponential moment makes
\(\mathcal E_a^p\) uniformly integrable; the stress bound gives
\eqref{eq:newS7V58-UI}.  CGMC row 4 is precisely strong reconstructed energy
convergence and hence excludes the three defect measures, while the last
part of row 5 is the protected-chart row.  The conclusion now follows from
\Cref{thm:newS7V58-RSNC}.
\end{proof}

\begin{assessment}[What has and has not moved]
\label{ass:newS7V59-status}
The real Einstein--Weyl functional cannot be the missing finite-cutoff
measure input: its conformal instability persists under uniform ellipticity.
A positive \(R^2\) coefficient blocks that particular escape, and positive
\(R^2+C^2\) coefficients give full deterministic \(L^2\)-curvature control,
subject to the explicit Higgs mixing inequality
\eqref{eq:newS7V59-coeff-condition}.  The remaining problem is constructive:
one must prove CGMC for a reflection-positive positive action, or specify a
complex conformal contour and prove that it restores the required reality,
positivity, projectivity, and stress bounds.  None of the audited manuscripts
currently supplies either proof.
\end{assessment}

\begin{programme}[V60: companion audit and normalization attack]
\label{prog:newS7V59-V60}
V60 will first perform the compulsory five-version audit of the newest
coherent Yang--Mills and Navier--Stokes notes.  It will then test whether
their conditional-kernel and no-escape machinery yields a noncircular route
from a local coercive action to the exponential moment
\eqref{eq:newS7V59-exp-moment}, separating single-cell normalization from
global gravitational entropy.
\end{programme}

\begin{firewall}[Claim boundary after seventh-series V59]
\label{fire:newS7V59-boundary}
V59 proves the conformal transformation formulae, an explicit uniformly
elliptic conformal instability of the Euclidean Einstein--Weyl action, a
curvature--matter coercive bound after positive \(R^2+C^2\) completion, and
the conditional implication CGMC to RSNC.  It also verifies that the two
audited gravitational sources leave the needed conformal construction open.

V59 does not prove CGMC, reflection positivity of a curvature-square
discretization, an entropy bound, a conformal contour theorem, RSNC for the
interacting measures, removal of the cutoff, or construction of the
Standard Model--Einstein continuum.
\end{firewall}

\begin{theorem}[Seventh-series V59 result]
\label{thm:newS7V59-result}
The Euclidean Einstein action with physical sign is unbounded below on a
uniformly elliptic, fixed-amplitude conformal sequence, and a Weyl-square
owner is exactly blind to that sequence.  A positive scalar-curvature-square
owner is therefore the first local curvature owner that blocks this escape;
together with a positive Weyl-square owner it controls full \(L^2\)
curvature modulo topology.  Under
\eqref{eq:newS7V59-coeff-condition} the coupled curvature--gauge--Higgs
action obeys \eqref{eq:newS7V59-coercive-bound}.  Promoting this deterministic
bound to RSNC requires precisely the finite-cutoff normalization, entropy,
compactness, positivity, and low-energy rows of CGMC.
\end{theorem}

\begin{proof}
Combine \Cref{lem:newS7V59-formulae,
thm:newS7V59-instability,
cor:newS7V59-Weyl-no,
thm:newS7V59-coercivity,
lem:newS7V59-full-curvature,
prop:newS7V59-CGMC-RSNC}.
\end{proof}

\paragraph{Parent-version record.}
V59 was formed from
\texttt{Renewal\_SM\_Einstein\_Continuum\_Research\_Notes\_V58.tex}.
The V58 parent had \(23\,313\) logical source lines, \(938\,816\) bytes,
and SHA--256
\[
 \texttt{EAB9A2107808E75E863046422C5680B9ACFB9D0C046079668DA6A6DC73489112}.
\]
The next compulsory companion audit is V60.

% =============================================================================
% END APPENDED SEVENTH-SERIES V59 MATERIAL
% =============================================================================

% =============================================================================
% APPENDED SEVENTH-SERIES V60 MATERIAL
% =============================================================================

\section{Companion audit, normalization entropy, and the quantum no-escape fork}
\label{sec:newS7V60}
\label{sec:v7v60}

This is the compulsory five-version audit.  It is followed by the promised
attack on the passage from a deterministic action lower bound to a Gibbs
probability estimate.  The main outcome is a correction of logical level:
V59's exponential raw-energy moment is a valid sufficient condition for a
finite-energy branch, but it is not a plausible generic property of a
four-dimensional quantum field as the number of cutoff modes diverges.

\subsection{Compulsory V60 companion audit}
\label{sec:newS7V60-audit}

The newest complete companion checkpoints inspected read-only were

\begin{align}
 &\texttt{Renewal\_Yang\_Mills\_Mass\_Gap\_Research\_Notes\_V34.tex},
 \notag\\
 &\qquad 542\,585\ {\rm bytes},\quad 15\,389\ {\rm logical\ lines},\quad
 \texttt{9E3CD3090F7B839A7F62F545E3EBF78BD5BF9D1A5E1638B5F056375453510098},
 \label{eq:newS7V60-YM-hash}\\
 &\texttt{Renewal\_NS\_Stability\_Research\_Notes\_V20.tex},
 \notag\\
 &\qquad 212\,147\ {\rm bytes},\quad 4\,118\ {\rm logical\ lines},\quad
 \texttt{C57F860D202D2E8CA4DE280197A5F36974DAB6DCDBA3882D2395087F469A56BD}.
 \label{eq:newS7V60-NS-hash}
\end{align}

Each checkpoint has one live document terminator.

\paragraph{Yang--Mills V34.}
The note identifies the response-matched plaquette coefficient
\(b^{\rm marg}=-n/d\), proves bounded local derivative stocks when the
response slope has a uniform floor and its numerator is bounded, and keeps
root regularity, generated-action stocks, and the topological short as
separate unproved inputs.  Its relevant lesson here is exact: finite
incidence and a bounded response coefficient control differentiated local
writers, while their undifferentiated oscillation remains extensive in the
number of cells.

\paragraph{Navier--Stokes V20.}
The note consolidates a local gate and a nearly constant bound, but records
that neither a sign nor a quantitative defect-rate floor follows.  Its
conditional measure lower bound is useful as a logical analogue: a strong
bound on a selected configuration does not show that the measure assigns
enough mass to that configuration class.

\begin{theorem}[V60 companion-audit decision]
\label{thm:newS7V60-audit}
The YM finite-incidence stock argument may be reused only for the local
stress-domination row after a gravitational response coefficient is proved
bounded.  It gives no uniform partition-function or entropy estimate.  The
NS gate gives no direct gravitational estimate, but confirms that local
smallness cannot replace a probability-mass lower bound.  Consequently
neither companion note proves CGMC, RSNC, or their quantum replacement
below.
\end{theorem}

\begin{proof}
YM V34 bounds each differentiated coordinate row because it meets only
finitely many plaquettes, while its own zeroth stock is \(O(M^4)\).  A
partition function sums over all coordinates and is sensitive to that
extensive stock, so the derivative estimate cannot bound its normalization.
NS V20 explicitly leaves the lower defect rate conditional.  The asserted
limitations are therefore stated in the audited results themselves.
\end{proof}

\subsection{What coercivity does to a normalized Gibbs measure}
\label{sec:newS7V60-normalization}

\begin{lemma}[Exact coercivity--normalization quotient]
\label{lem:newS7V60-quotient}
Let \((\Omega_a,\rho_a)\) be a finite measure space,
\(\mathcal E_a\geq0\), and

\[
 d\mu_a=Z_a^{-1}e^{-S_a}\,d\rho_a,
 \qquad
 Z_a=\int_{\Omega_a}e^{-S_a}\,d\rho_a.
\]

If \(S_a\geq c\mathcal E_a-C\), then for every \(0<\theta<c\),

\begin{equation}
 \int e^{\theta\mathcal E_a}\,d\mu_a
 \leq {e^C\over Z_a}
 \int e^{-(c-\theta)\mathcal E_a}\,d\rho_a.
 \label{eq:newS7V60-partition-ratio}
\end{equation}

If, in addition, there is a measurable \(B_a\) with
\(\rho_a(B_a)\geq m_a>0\) and \(S_a\leq K_a\) on \(B_a\), then

\begin{equation}
 \int e^{\theta\mathcal E_a}\,d\mu_a
 \leq {e^{C+K_a}\over m_a}
 \int e^{-(c-\theta)\mathcal E_a}\,d\rho_a.
 \label{eq:newS7V60-smallball}
\end{equation}
\end{lemma}

\begin{proof}
The action lower bound gives
\[
 e^{-S_a+\theta\mathcal E_a}
 \leq e^C e^{-(c-\theta)\mathcal E_a}.
\]
Integrating and dividing by \(Z_a\) proves
\eqref{eq:newS7V60-partition-ratio}.  On \(B_a\),
\(e^{-S_a}\geq e^{-K_a}\), whence
\(Z_a\geq m_ae^{-K_a}\), which proves
\eqref{eq:newS7V60-smallball}.
\end{proof}

\begin{corollary}[The missing datum is global entropy]
\label{cor:newS7V60-entropy}
The deterministic coercive estimate
\eqref{eq:newS7V59-coercive-bound} proves the exponential moment
\eqref{eq:newS7V59-exp-moment} only after a cutoff-uniform bound for the
partition ratio on the right of
\eqref{eq:newS7V60-partition-ratio}.  A one-cell action estimate or a
finite-incidence derivative bound does not provide that ratio.  A small-ball
proof must control the dimension-dependent quantity \(e^{K_a}/m_a\).
\end{corollary}

\begin{proof}
This is the content of \Cref{lem:newS7V60-quotient}.  Under tensorization,
small-ball masses multiply, so \(m_a\) generally decays exponentially in
the number of cells even when the one-cell mass is fixed.
\end{proof}

\subsection{The extensive-mode obstruction}
\label{sec:newS7V60-extensive}

\begin{proposition}[A coercive Gaussian fails the raw uniform moment]
\label{prop:newS7V60-Gaussian}
Let \(D_a\to\infty\), let \(\mu_a\) be standard Gaussian probability on
\(\mathbb R^{D_a}\), and set
\(\mathcal E_a(x)=|x|^2/2\).  Although
\(S_a=\mathcal E_a\) is strictly coercive and every cutoff measure is
normalized, for \(0<\theta<1\),

\begin{equation}
 \int_{\mathbb R^{D_a}}e^{\theta\mathcal E_a}\,d\mu_a
 =(1-\theta)^{-D_a/2}\longrightarrow\infty,
 \qquad
 \int\mathcal E_a\,d\mu_a={D_a\over2}.
 \label{eq:newS7V60-Gaussian-moment}
\end{equation}
\end{proposition}

\begin{proof}
The Gaussian integral factorizes into \(D_a\) copies of
\[
 (2\pi)^{-1/2}\int_{\mathbb R}
 e^{-(1-\theta)x^2/2}\,dx=(1-\theta)^{-1/2}.
\]
Differentiating at \(\theta=0\) gives the mean energy.
\end{proof}

\begin{assessment}[Correction of the V59 sufficient card]
\label{ass:newS7V60-CGMC-correction}
CGMC remains a correct sufficient card for a measure supported on the
classical finite-energy space, but
\Cref{prop:newS7V60-Gaussian} shows that its cutoff-uniform raw total-energy
moment is already false for the simplest growing Gaussian system.  A
four-dimensional quantum continuum normally requires subtraction and
renormalization of composite energy densities and a distributional state
space.  Therefore failure of CGMC does not by itself disprove a quantum
continuum; it closes only the direct finite-energy route used in V58.
\end{assessment}

The same point affects the no-defect row of RSNC.  Strong \(L^2\) convergence
of raw cutoff curvatures and gauge fields is a classical compactness
mechanism.  Quantum composite fields can converge after local subtraction
even while their unrenormalized energies diverge.  The correct alternative
must make those subtractions part of the object being shown tight.

\subsection{Renormalized composite no-escape}
\label{sec:newS7V60-RCNC}

\begin{definition}[Renormalized composite no-escape card, RCNC]
\label{def:newS7V60-RCNC}
Let \(\mathscr Y\) be a declared nuclear distribution space for the
gauge-fixed or quotient fundamental fields, and let
\(\mathscr D'(S^2T^*M)\) be the distribution space for symmetric
two-tensors.  RCNC consists of:
\begin{enumerate}
\item normalized cutoff probabilities satisfying the declared gauge and
Osterwalder--Schrader positivity conditions, whose cylinder laws are tight
on \(\mathscr Y\);
\item a finite renormalized composite family
\[
 [R^2]_a,\ [C^2]_a,\ [|F|^2]_a,\ [|DH|^2]_a,\ [|H|^4]_a,
 \quad\hbox{and}\quad \mathcal T_a^{\rm ren},
\]
with all local subtractions fixed by the common V57 renormalization
conditions rather than chosen subsequence by subsequence;
\item joint convergence in law of every finite set of smeared fundamental
and composite fields, including
\(\langle\mathcal T_a^{\rm ren},h\rangle\) for every smooth test tensor
\(h\);
\item for each such \(h\), numbers \(\delta_h>0\) and \(C_h<\infty\) with
\[
 \sup_a\int
 |\langle\mathcal T_a^{\rm ren},h\rangle|^{1+\delta_h}\,d\mu_a
 \leq C_h;
 \label{eq:newS7V60-RCNC-UI}
\]
\item cutoff Ward/BRST and diffeomorphism identities whose breaking terms
vanish after smearing, plus vanishing probability of the complement of the
protected overlap-gap chart; and
\item agreement, when the metric and bosonic fields are restricted to a
smooth external-background chart, with the heat-kernel stress and its two
contacts computed in V57--V58.
\end{enumerate}
\end{definition}

\begin{theorem}[RCNC transfers the renormalized stress without raw energy
compactness]
\label{thm:newS7V60-RCNC}
Assume RCNC and write \(\mathcal T^{\rm ren}\) for the limiting random
distribution.  Then for every smooth test tensor \(h\),

\begin{equation}
 \lim_{a\downarrow0}
 \int\langle\mathcal T_a^{\rm ren},h\rangle\,d\mu_a
 =
 \int\langle\mathcal T^{\rm ren},h\rangle\,d\mu.
 \label{eq:newS7V60-RCNC-limit}
\end{equation}

The limiting stress obeys the distributional identities obtained as limits
of the row-5 Ward identities.  On smooth external backgrounds its
parity-even local representative is the common V58 heat stress.
\end{theorem}

\begin{proof}
RCNC row 3 gives convergence in law of each paired stress.  The
\(L^{1+\delta_h}\) bound in row 4 is uniform integrability, so Vitali's
theorem gives convergence of expectations.  Smearing each cutoff identity
against a fixed smooth test field and using joint convergence plus uniform
integrability passes the identity to the limit.  Row 6 identifies the
smooth-background representative by the uniqueness theorem of V57--V58.
\end{proof}

\begin{proposition}[Neither no-escape card subsumes the other]
\label{prop:newS7V60-fork}
RSNC plus the V58 finite-energy topology identifies nonlinear stresses by
strong classical energy convergence.  RCNC identifies them as independently
renormalized composite random distributions.  Without an additional theorem
equating the composite renormalization with a finite-energy product, neither
card implies the other.
\end{proposition}

\begin{proof}
\Cref{prop:newS7V60-Gaussian} gives systems with well-defined Gaussian
distributional limits but divergent raw energy moments, excluding the
RSNC mechanism.  Conversely, strong finite-energy convergence fixes
classical products but supplies neither a chosen composite subtraction nor
Osterwalder--Schrader convergence for quantum cylinder laws.  The missing
implications are therefore substantive.
\end{proof}

\begin{assessment}[The revised upstream bottleneck]
\label{ass:newS7V60-status}
There are now two honest routes.  The classical route must prove the very
strong CGMC and then RSNC.  The quantum route must construct renormalized
composite fields and prove RCNC.  Action coercivity helps finite-cutoff
normalization but does not overcome the extensive-mode entropy or define
the limiting stress.  The YM stock compiler can contribute to RCNC's local
stress bound only after its response inputs are populated; the NS gate
provides no missing probability floor.
\end{assessment}

\begin{programme}[V61: reflection positivity of the conformal quadratic
sector]
\label{prog:newS7V60-V61}
V61 will linearize the conformal sector about a flat background and test
the candidate positive \(R^2\) completion against
Osterwalder--Schrader reflection positivity.  The inverse polynomial in
momentum will be factorized and its residues tested.  If the stabilizing
fourth-order covariance has a negative spectral weight, the note will
isolate the exact alternative required: a second-order positive auxiliary
field representation with acceptable signs, a nonlocal positive spectral
form factor, or a declared conformal contour.  It will not infer quantum
positivity from the V59 lower bound.
\end{programme}

\begin{firewall}[Claim boundary after seventh-series V60]
\label{fire:newS7V60-boundary}
V60 performs the compulsory companion audit; proves the exact
coercivity-to-partition-ratio inequality; proves by a growing Gaussian model
that a uniform raw total-energy exponential moment is generally too strong;
and formulates RCNC with a conditional stress-transfer theorem that does not
assume raw strong-energy convergence.

V60 does not construct the cutoff gravitational probability, prove an
entropy ratio, prove RCNC, prove reflection positivity of the conformal
sector, construct a renormalized quantum stress on random metrics, remove
the interacting cutoff, or construct the Standard Model--Einstein
continuum.
\end{firewall}

\begin{theorem}[Seventh-series V60 result]
\label{thm:newS7V60-result}
A deterministic coercive action bound yields a cutoff-uniform energy moment
only together with the global partition ratio
\eqref{eq:newS7V60-partition-ratio}.  Even strict Gaussian coercivity fails
the uniform raw total-energy moment when the number of modes diverges.
Therefore the finite-energy CGMC--RSNC route and the quantum
composite-field RCNC route must be kept distinct.  RCNC suffices to pass
the smeared renormalized stress and its Ward identities to a distributional
limit, while agreeing with the V58 common heat stress on smooth external
backgrounds.
\end{theorem}

\begin{proof}
Use \Cref{thm:newS7V60-audit,
lem:newS7V60-quotient,
cor:newS7V60-entropy,
prop:newS7V60-Gaussian,
thm:newS7V60-RCNC,
prop:newS7V60-fork}.
\end{proof}

\paragraph{Parent-version record.}
V60 was formed from
\texttt{Renewal\_SM\_Einstein\_Continuum\_Research\_Notes\_V59.tex}.
The V59 parent had \(23\,743\) logical source lines, \(954\,504\) bytes,
and SHA--256
\[
 \texttt{C0892B99F05245F80AE108987C1BF69676278C8E4DCE67D9519986A8BF535FA4}.
\]
The next compulsory companion audit is V65.

% =============================================================================
% END APPENDED SEVENTH-SERIES V60 MATERIAL
% =============================================================================

% =============================================================================
% APPENDED SEVENTH-SERIES V61 MATERIAL
% =============================================================================

\section{Quadratic conformal stabilization versus reflection positivity}
\label{sec:newS7V61}
\label{sec:v7v61}

V59 showed that a positive scalar-curvature-square owner blocks the explicit
high-frequency conformal escape of the physical-sign Euclidean Einstein
action.  V60 showed that this deterministic fact is not yet a quantum
measure estimate.  The next upstream test is sharper: even if the quadratic
form is made positive, can its Gaussian covariance satisfy
Osterwalder--Schrader reflection positivity?  The answer is negative when
the conformal factor itself is required to be a positive-metric scalar
field with a local fourth-order inverse propagator.

\subsection{The exact flat-background Hessian}
\label{sec:newS7V61-Hessian}

Work first on a flat four-torus, write \(g=e^{2\phi}\delta\), and suppose
the cosmological tadpole has been cancelled by the chosen background
equation or a volume constraint.  For

\[
 S_{\rm conf}(\phi)
 =\alpha\int R_g^2\,d\mu_g
  -\kappa\int R_g\,d\mu_g
  +\Lambda\int d\mu_g,
 \qquad \alpha,\kappa>0,
\]

the identities of V59 give

\begin{align}
 \int R_g^2\,d\mu_g
 &=36\int(\Delta\phi+|\nabla\phi|^2)^2\,dx,
 \label{eq:newS7V61-R2-exact}\\
 -\kappa\int R_g\,d\mu_g
 &=-6\kappa\int e^{2\phi}|\nabla\phi|^2\,dx,
 \label{eq:newS7V61-EH-exact}\\
 \Lambda\int d\mu_g
 &=\Lambda\int e^{4\phi}\,dx.
 \label{eq:newS7V61-Lambda-exact}
\end{align}

\begin{lemma}[Conformal quadratic polynomial]
\label{lem:newS7V61-polynomial}
After removal of the constant and linear terms, the quadratic action is

\begin{equation}
 S_{\rm conf}^{(2)}(\phi)
 =\int\left(36\alpha(\Delta\phi)^2
            -6\kappa|\nabla\phi|^2+8\Lambda\phi^2\right)dx.
 \label{eq:newS7V61-quadratic}
\end{equation}

Its Fourier multiplier is

\begin{equation}
 Q(z)=36\alpha z^2-6\kappa z+8\Lambda,
 \qquad z=|p|^2.
 \label{eq:newS7V61-Q}
\end{equation}

On a continuum of momenta \(Q(z)>0\) for all \(z\geq0\) exactly when
\[
 \Lambda>{\kappa^2\over32\alpha}.
 \label{eq:newS7V61-positive-condition}
\]
Equality gives a zero mode at \(z=\kappa/(12\alpha)\).
\end{lemma}

\begin{proof}
Taking the quadratic parts of
\eqref{eq:newS7V61-R2-exact}--\eqref{eq:newS7V61-Lambda-exact} gives
\eqref{eq:newS7V61-quadratic}; Fourier transformation gives
\eqref{eq:newS7V61-Q}.  The upward-opening polynomial has its minimum at
\(z=\kappa/(12\alpha)\), where
\[
 Q_{\min}=8\Lambda-{\kappa^2\over4\alpha}.
\]
This proves the condition and the equality statement.
\end{proof}

\begin{corollary}[What the positive \(R^2\) term actually fixes]
\label{cor:newS7V61-fix}
The positive \(R^2\) owner controls the ultraviolet conformal modes, but with
the physical Einstein sign it does not by itself make the quadratic action
positive: a positive zeroth-order Hessian, a finite-volume spectral gap, or
another infrared mechanism is also required.
\end{corollary}

\begin{proof}
The \(z^2\) coefficient in \eqref{eq:newS7V61-Q} is positive and the
\(z\) coefficient is negative.  The minimum calculation in
\Cref{lem:newS7V61-polynomial} proves the assertion.
\end{proof}

\subsection{A spectral test for Gaussian reflection positivity}
\label{sec:newS7V61-spectral}

\begin{lemma}[Positive scalar spectral representation]
\label{lem:newS7V61-KL}
Let a nonzero, translation- and Euclidean-invariant scalar Gaussian
covariance on \(\mathbb R^4\) be tempered and
Osterwalder--Schrader reflection positive.  Up to contact terms invisible
between strictly positive and strictly negative times, its momentum
covariance has a positive Källén--Lehmann representation

\begin{equation}
 C(z)=\int_{[0,\infty)}{d\rho(s)\over z+s},
 \qquad z>0,
 \qquad \rho\geq0.
 \label{eq:newS7V61-KL}
\end{equation}

Consequently a noncontact covariance cannot obey \(C(z)=O(z^{-2})\) as
\(z\to\infty\).
\end{lemma}

\begin{proof}
Reflection positivity reconstructs a positive Hilbert space, a
nonnegative Hamiltonian, and a positive translation spectral measure.
Euclidean invariance disintegrates that measure by mass squared and gives
\eqref{eq:newS7V61-KL}; equivalently, for each spatial momentum the
positive-time kernel is a Laplace transform of a positive measure.
For the final assertion, multiply \eqref{eq:newS7V61-KL} by \(z\):
\[
 zC(z)=\int {z\over z+s}\,d\rho(s).
\]
The integrand increases pointwise to one.  Monotone convergence gives
\(\lim_{z\to\infty}zC(z)=\rho([0,\infty])\), allowing \(+\infty\).
If \(C(z)=O(z^{-2})\), this limit is zero, so \(\rho=0\), contradicting
the nonzero covariance.  Polynomial contact terms cannot equal a decaying
rational covariance and do not alter this conclusion.
\end{proof}

\begin{theorem}[Fourth-order conformal Gaussian no-go]
\label{thm:newS7V61-no-go}
Let \(Q\) be a real polynomial of degree two with positive leading
coefficient and \(Q(z)>0\) for every \(z\geq0\).  The nonzero scalar
Gaussian covariance

\begin{equation}
 C(p)={1\over Q(|p|^2)}
 \label{eq:newS7V61-covariance}
\end{equation}

is not Osterwalder--Schrader reflection positive.  In particular, whenever
\eqref{eq:newS7V61-positive-condition} makes the local conformal Hessian
positive, its elementary \(\phi\)-Gaussian still fails scalar reflection
positivity.
\end{theorem}

\begin{proof}
Since \(Q(z)\sim q_2z^2\) with \(q_2>0\),
\(C(z)=O(z^{-2})\).  If the Gaussian were reflection positive,
\Cref{lem:newS7V61-KL} would force its positive spectral measure to vanish,
and hence \(C=0\), a contradiction.
\end{proof}

\begin{proposition}[The negative residue in the factorized case]
\label{prop:newS7V61-residue}
Suppose, after any sign choice or infrared completion, the quadratic
polynomial factorizes as
\[
 Q(z)=\alpha_0(z+m_1^2)(z+m_2^2),
 \qquad
 \alpha_0>0,\quad0\leq m_1<m_2.
\]
Then
\begin{equation}
 {1\over Q(z)}
 ={1\over\alpha_0(m_2^2-m_1^2)}
 \left({1\over z+m_1^2}-{1\over z+m_2^2}\right).
 \label{eq:newS7V61-residues}
\end{equation}
The second pole has negative spectral weight.  There is a smooth test
function supported at positive Euclidean times whose reflected Gaussian
quadratic form is strictly negative.
\end{proposition}

\begin{proof}
The partial fraction identity is immediate.  Fix a spatial momentum and
write \(\omega_j=(|\mathbf p|^2+m_j^2)^{1/2}\).  The reflected time form,
apart from one positive common factor, is
\[
 {1\over2\omega_1}
 \left|\int_0^\infty f(t)e^{-\omega_1t}\,dt\right|^2
 -
 {1\over2\omega_2}
 \left|\int_0^\infty f(t)e^{-\omega_2t}\,dt\right|^2.
\]
Choose two distinct positive times and a linear combination of narrow
smooth bumps around them whose \(\omega_1\)-Laplace transform vanishes.
The two exponential evaluation vectors are linearly independent because
\(\omega_1\ne\omega_2\); hence the combination can be chosen with nonzero
\(\omega_2\)-Laplace transform.  The displayed form is then negative.
The bump widths can be taken small enough to preserve strict negativity.
\end{proof}

\begin{corollary}[Degenerate and complex roots do not evade the test]
\label{cor:newS7V61-roots}
A double pole cannot be the Stieltjes transform of a positive measure, and
a nonreal pole cannot lie in the support form of
\eqref{eq:newS7V61-KL}.  Thus coincident or complex roots of \(Q\) do not
restore reflection positivity.
\end{corollary}

\begin{proof}
A rational Stieltjes transform of a positive measure has only simple poles
on the nonpositive real axis, with nonnegative residues: each pole is an
atom of the measure in \eqref{eq:newS7V61-KL}.  A double or nonreal pole
has neither form.  The general conclusion also follows directly from
\Cref{thm:newS7V61-no-go}.
\end{proof}

\subsection{The precise scope of the no-go}
\label{sec:newS7V61-scope}

The theorem applies when \(\phi\) is itself an elementary scalar in the
positive Osterwalder--Schrader algebra.  In gravity, that assumption must be
declared rather than hidden.  A gauge-fixed metric component can live in an
indefinite auxiliary space provided a separate theorem proves positivity on
the diffeomorphism-invariant observable algebra.  Likewise, a complex
conformal contour is outside the real Gaussian hypothesis, but it must
eventually produce real, positive physical Schwinger functions.

\begin{definition}[Conformal positivity resolution card, CPRC]
\label{def:newS7V61-CPRC}
A resolution of the combined V59 and V61 obstructions must choose and prove
one of the following branches.
\begin{enumerate}
\item \emph{Physical-algebra branch:} construct the stabilized gauge-fixed
measure, identify a diffeomorphism-invariant Osterwalder--Schrader
subalgebra, and prove positivity there despite the conformal covariance.
\item \emph{Second-order branch:} replace the fourth-order elementary
conformal covariance by a bounded, second-order positive field system and
prove that integrating or constraining its additional fields yields the
required Einstein low-energy response without a negative spectral weight.
\item \emph{Positive spectral branch:} specify a possibly nonlocal
covariance with a positive measure in \eqref{eq:newS7V61-KL}, prove
nonlinear stability of the real conformal direction, and recover the
Einstein response at low momentum.
\item \emph{Contour branch:} specify the integration cycle, prove
convergence and cutoff compatibility, and reconstruct a real
reflection-positive physical observable algebra.
\end{enumerate}
Every branch must also preserve the anomaly coefficients and common smooth
stress fixed in V54--V58 and must populate RCNC rather than merely give a
formal propagator.
\end{definition}

\begin{proposition}[Local auxiliary fields require an additional theorem]
\label{prop:newS7V61-auxiliary}
No positive local second-order auxiliary-field representation can reproduce
\eqref{eq:newS7V61-covariance} as the covariance of an elementary
positive-metric scalar \(\phi\).  Any proposed auxiliary resolution must
therefore change the physical observable assignment, impose a constraint,
or change the \(\phi\) covariance.
\end{proposition}

\begin{proof}
If such a representation existed with \(\phi\) in its positive physical
algebra, marginalizing the other Gaussian fields would preserve reflection
positivity of the \(\phi\) subalgebra.  Its marginal covariance would be
\eqref{eq:newS7V61-covariance}, contradicting
\Cref{thm:newS7V61-no-go}.
\end{proof}

\begin{assessment}[The conformal trilemma]
\label{ass:newS7V61-status}
The physical-sign Einstein action on a real conformal contour is
ultraviolet-unbounded below without a higher-derivative stabilizer.  A
positive \(R^2\) stabilizer gives a fourth-order elementary covariance, and
that covariance fails scalar reflection positivity.  Reversing the
Einstein sign does not cure the fourth-order spectral obstruction.  Thus
the programme now needs a CPRC branch in which positivity is proved on the
physical gravitational algebra, or the conformal construction is changed;
local action positivity alone cannot close RCNC.
\end{assessment}

\begin{programme}[V62: the physical observable algebra]
\label{prog:newS7V61-V62}
V62 will go upstream to the linearized diffeomorphism quotient.  It will
decompose the metric Hessian into spin projectors, identify which scalar
combination is constrained or gauge invariant, and test whether restricting
Osterwalder--Schrader positivity to conserved gauge-invariant sources
removes the negative conformal spectral contribution.  This is the first
CPRC branch and must be settled before attempting a nonlinear entropy
estimate.
\end{programme}

\begin{firewall}[Claim boundary after seventh-series V61]
\label{fire:newS7V61-boundary}
V61 computes the flat conformal Hessian of the \(R^2\)-completed physical
Einstein action; gives its exact quadratic positivity condition; proves
that every nonzero degree-two inverse polynomial covariance fails scalar
Osterwalder--Schrader positivity; and formulates the four admissible CPRC
branches.

V61 does not prove that the gauge-invariant gravitational observable
algebra is nonpositive, rule out a contour or nonlocal construction, choose
a CPRC branch, prove RCNC, construct an interacting gravitational measure,
remove the cutoff, or construct the Standard Model--Einstein continuum.
\end{firewall}

\begin{theorem}[Seventh-series V61 result]
\label{thm:newS7V61-result}
The positive \(R^2\) owner solves only the ultraviolet lower-bound part of
the Euclidean conformal problem.  With the physical Einstein sign, the
quadratic action additionally needs
\(\Lambda>\kappa^2/(32\alpha)\) on continuous flat momenta.  Whenever the
resulting elementary conformal covariance is the reciprocal of a nonzero
quadratic polynomial in \(p^2\), its \(p^{-4}\) decay contradicts a positive
Källén--Lehmann representation, so scalar Osterwalder--Schrader positivity
fails.  A valid continuation must populate one branch of CPRC and establish
positivity on the actual physical observable algebra.
\end{theorem}

\begin{proof}
Use \Cref{lem:newS7V61-polynomial,
cor:newS7V61-fix,
lem:newS7V61-KL,
thm:newS7V61-no-go,
prop:newS7V61-residue,
cor:newS7V61-roots,
prop:newS7V61-auxiliary}.
\end{proof}

\paragraph{Parent-version record.}
V61 was formed from
\texttt{Renewal\_SM\_Einstein\_Continuum\_Research\_Notes\_V60.tex}.
The V60 parent had \(24\,088\) logical source lines, \(968\,713\) bytes,
and SHA--256
\[
 \texttt{694915230CF53282E3C89F5BE64A7B65F1CF99C7244614BB3543EFFBB8D56740}.
\]
The next compulsory companion audit is V65.

% =============================================================================
% END APPENDED SEVENTH-SERIES V61 MATERIAL
% =============================================================================

% =============================================================================
% APPENDED SEVENTH-SERIES V62 MATERIAL
% =============================================================================

\section{The linearized diffeomorphism quotient does not remove the
curvature-square ghost}
\label{sec:newS7V62}
\label{sec:v7v62}

V61 left open whether scalar reflection positivity was simply being imposed
on the wrong, gauge-fixed variable.  This note tests the first CPRC branch
on the linearized physical algebra.  The conformal trace is not itself a
diffeomorphism-invariant field, but its linearized scalar curvature is.
More decisively, the positive Weyl-square owner required by the V59
full-curvature coercive estimate creates a negative massive spin-two
spectral weight which is visible to the gauge-invariant linearized Weyl
tensor.

\subsection{Gauge-invariant linearized curvatures}
\label{sec:newS7V62-curvatures}

Let \(g_{\mu\nu}=\delta_{\mu\nu}+h_{\mu\nu}\) on a flat chart.  The
linearized diffeomorphism action is

\begin{equation}
 \delta_\xi h_{\mu\nu}=\partial_\mu\xi_\nu+\partial_\nu\xi_\mu.
 \label{eq:newS7V62-gauge}
\end{equation}

The linearized scalar and Riemann curvatures are

\begin{align}
 R^{(1)}(h)
 &=\partial_\mu\partial_\nu h_{\mu\nu}-\Delta h,
 \label{eq:newS7V62-R1}\\
 R^{(1)}_{\mu\nu\rho\sigma}(h)
 &={1\over2}\bigl(
 \partial_\rho\partial_\nu h_{\mu\sigma}
 +\partial_\sigma\partial_\mu h_{\nu\rho}
 -\partial_\sigma\partial_\nu h_{\mu\rho}
 -\partial_\rho\partial_\mu h_{\nu\sigma}\bigr).
 \label{eq:newS7V62-Riem1}
\end{align}

\begin{lemma}[The curvature observables descend to the quotient]
\label{lem:newS7V62-invariance}
The tensors \(R^{(1)}\), \(R_{\mu\nu\rho\sigma}^{(1)}\), and
\(C_{\mu\nu\rho\sigma}^{(1)}\) are invariant under
\eqref{eq:newS7V62-gauge}.  For a pure conformal perturbation
\(h_{\mu\nu}=2\phi\delta_{\mu\nu}\),

\begin{equation}
 R^{(1)}(2\phi\delta)=-6\Delta\phi.
 \label{eq:newS7V62-conformal-R}
\end{equation}

Thus the high-frequency conformal direction is not erased merely by passing
to the linearized diffeomorphism quotient.
\end{lemma}

\begin{proof}
Insert \eqref{eq:newS7V62-gauge} into
\eqref{eq:newS7V62-Riem1}.  Every third derivative occurs twice with
opposite sign, since partial derivatives commute.  All contractions and the
trace-free Weyl combination are therefore invariant.  Substituting
\(h_{\mu\nu}=2\phi\delta_{\mu\nu}\) into
\eqref{eq:newS7V62-R1} gives
\(2\Delta\phi-8\Delta\phi=-6\Delta\phi\).
\end{proof}

A conserved source also defines a quotient observable: if
\(\partial_\mu J_{\mu\nu}=0\), integration by parts gives
\[
 \delta_\xi\int h_{\mu\nu}J_{\mu\nu}\,dx
 =-2\int\xi_\nu\partial_\mu J_{\mu\nu}\,dx=0.
\]
Traceless conserved sources isolate the spin-two sector.

\subsection{The transverse-traceless Hessian}
\label{sec:newS7V62-TT}

Let \(h^{\rm TT}\) obey
\(\partial_\mu h_{\mu\nu}^{\rm TT}=0\) and
\(h_{\mu\mu}^{\rm TT}=0\).  On this slice,

\[
 R^{(1)}(h^{\rm TT})=0,
 \qquad
 \operatorname{Ric}^{(1)}(h^{\rm TT})
 =-{1\over2}\Delta h^{\rm TT}.
\]

\begin{lemma}[Einstein--Weyl TT quadratic action]
\label{lem:newS7V62-TT-action}
For
\[
 S_{\rm EW}(g)=-\kappa\int R_g\,d\mu_g
                +\beta\int|C_g|^2\,d\mu_g,
 \qquad \kappa,\beta>0,
\]
the flat-background transverse-traceless quadratic part is

\begin{equation}
 S_{\rm EW}^{(2)}(h^{\rm TT})
 ={1\over2}\int h^{\rm TT}
 \left(\kappa_2(-\Delta)+\beta_2(-\Delta)^2\right)
 h^{\rm TT}\,dx,
 \label{eq:newS7V62-TT-action}
\end{equation}

where \(\kappa_2,\beta_2>0\).  With the normalization
\(g=\delta+h\), one may take
\(\kappa_2=\kappa/2\) and \(\beta_2=\beta\).
\end{lemma}

\begin{proof}
The second variation of the negative-sign Euclidean Einstein functional at
a Ricci-flat metric, restricted to TT variations, is
\[
 {\kappa\over4}\int|\partial h^{\rm TT}|^2\,dx.
\]
For the Weyl term, the quadratic Gauss--Bonnet identity and
\(R^{(1)}=0\) give
\[
 \int|C^{(1)}|^2\,dx
 =2\int|\operatorname{Ric}^{(1)}|^2\,dx
 ={1\over2}\int|\Delta h^{\rm TT}|^2\,dx.
\]
Writing their sum in the one-half convention proves
\eqref{eq:newS7V62-TT-action} and the stated coefficients.
\end{proof}

\begin{theorem}[The Weyl-square pole is present on the physical spin-two
algebra]
\label{thm:newS7V62-spin2-ghost}
The spin-two covariance coefficient determined by
\eqref{eq:newS7V62-TT-action} is

\begin{equation}
 D_2(z)={1\over z(\kappa_2+\beta_2z)}
 ={1\over\kappa_2}
 \left({1\over z}-{1\over z+M_2^2}\right),
 \qquad
 M_2^2={\kappa_2\over\beta_2}.
 \label{eq:newS7V62-D2}
\end{equation}

Its massive spin-two pole has negative spectral weight.  This violates
Osterwalder--Schrader positivity not only for a gauge-fixed metric
component, but also on the gauge-invariant algebra generated by smeared
linearized Weyl tensors.
\end{theorem}

\begin{proof}
The partial fraction identity proves the negative residue.  Conserved
traceless sources project onto the TT coefficient, so the positive spectral
condition applies separately to \(D_2\).  Equivalently, apply the
linearized Weyl operator to both entries of the metric covariance.  On a
spin-two polarization this operator is nonzero and has two derivatives.
The noncontact scalar coefficient of the resulting curvature covariance
contains
\[
 {z\over\kappa_2+\beta_2z}
 ={1\over\beta_2}
 -{\kappa_2\over\beta_2^2}{1\over z+M_2^2}.
\]
The first term is a contact term, while the massive pole has strictly
negative coefficient.  A positive-time test can annihilate the massless
or contact contribution while retaining the massive Laplace mode, exactly
as in \Cref{prop:newS7V61-residue}.  Since the linearized Weyl tensor is
gauge invariant by \Cref{lem:newS7V62-invariance}, the negative norm occurs
inside the physical curvature algebra.
\end{proof}

\begin{corollary}[The first CPRC branch does not rescue finite positive
\(C^2\)]
\label{cor:newS7V62-branch1}
For a finite coefficient \(\beta>0\), restricting positivity from the
gauge-fixed metric algebra to the linearized diffeomorphism-invariant
curvature algebra does not remove the Weyl-square obstruction.
\end{corollary}

\begin{proof}
This is the last assertion of
\Cref{thm:newS7V62-spin2-ghost}.
\end{proof}

\subsection{The scalar constraint is not an automatic cure}
\label{sec:newS7V62-scalar}

\begin{proposition}[Conformal curvature requires a declared constraint]
\label{prop:newS7V62-scalar}
Passing to diffeomorphism equivalence classes does not set
\(R^{(1)}\) to zero.  Removing the conformal scalar from the physical
spectrum therefore requires an equation-of-motion constraint, a reduced
phase-space construction, or a BRST cohomology theorem; it does not follow
from gauge invariance alone.
\end{proposition}

\begin{proof}
\Cref{lem:newS7V62-invariance} shows simultaneously that \(R^{(1)}\) is
gauge invariant and that it equals \(-6\Delta\phi\) on a conformal
perturbation.  It is nonzero for every nonharmonic \(\phi\).  Hence no
linearized diffeomorphism can eliminate that curvature observable.
\end{proof}

For pure Einstein gravity, lapse and shift constraints in a canonical
formulation leave the two transverse spin-two helicities rather than a
propagating conformal scalar.  That statement does not define a positive
covariant Euclidean measure over all real metrics.  A renewal construction
which uses it must exhibit the constraint before integration and prove that
the induced physical Schwinger functions satisfy the required positivity.

\subsection{A decoupling limit and its price}
\label{sec:newS7V62-decoupling}

\begin{lemma}[The negative spin-two pole decouples distributionally if
\(\beta_2\downarrow0\)]
\label{lem:newS7V62-decoupling}
For fixed \(\kappa_2>0\),
\[
 D_{2,\beta}(p)-{1\over\kappa_2|p|^2}
 =-{1\over\kappa_2}\,
 {1\over |p|^2+\kappa_2/\beta_2}.
\]
For every Schwartz test function whose massless pairing is defined, the
right side tends to zero as \(\beta_2\downarrow0\).
\end{lemma}

\begin{proof}
The identity is \eqref{eq:newS7V62-D2}.  The absolute value of the
right side is bounded by
\(\beta_2/\kappa_2^2\).  Multiplication by a Schwartz Fourier test
density and dominated convergence prove distributional convergence.
\end{proof}

\begin{assessment}[Counterterm versus resummed kinetic owner]
\label{ass:newS7V62-counterterm}
The lemma permits a logically distinct use of \(C^2\): its coefficient may
serve as a cutoff counterterm or perturbative local contact owner while its
extra pole is sent to infinite mass and never promoted to a finite
fundamental propagator.  This is compatible with the local heat-kernel
anomaly bookkeeping of V54--V58.  It forfeits the finite positive
Weyl-square coercivity used in V59, so another nonperturbative mechanism
must then control traceless curvature and populate RCNC.
\end{assessment}

\begin{definition}[Spin-two decoupling card, S2DC]
\label{def:newS7V62-S2DC}
A cutoff family using a positive Weyl-square regulator satisfies S2DC if:
\begin{enumerate}
\item \(\beta_2(a)>0\) at finite cutoff and
\(M_2(a)^2=\kappa_2(a)/\beta_2(a)\to\infty\);
\item every fixed collection of gauge-invariant separated-point
observables has a limit in which the negative-pole contribution vanishes,
uniformly enough to pass the Osterwalder--Schrader inequalities;
\item the local \(C^2\) anomaly and stress contacts converge with the
renormalization fixed in V54--V58; and
\item tightness and nonlinear curvature control are proved by estimates
which do not require a positive lower bound for \(\beta_2(a)\).
\end{enumerate}
\end{definition}

\begin{proposition}[S2DC is necessary for the local-regulator branch]
\label{prop:newS7V62-S2DC}
Within a local polynomial Einstein--Weyl cutoff action, if the limiting
physical algebra is Osterwalder--Schrader positive and contains the
linearized Weyl observable, then either S2DC removes the finite massive
pole or another sector cancels it through a proved physical constraint.
\end{proposition}

\begin{proof}
If a finite uncancelled pole remains, its negative physical spectral weight
contradicts \Cref{thm:newS7V62-spin2-ghost}.  The two stated alternatives
are the ways in which the hypothesis of that theorem can fail: the pole
leaves every fixed spectral window, or it is absent from the constrained
physical algebra.
\end{proof}

\begin{assessment}[Updated conformal/gravitational fork]
\label{ass:newS7V62-status}
The physical-algebra branch by itself does not reconcile the V59 coercive
completion with reflection positivity.  The scalar curvature detects the
conformal direction on the quotient, and the Weyl curvature detects the
negative massive spin-two residue.  The least radical remaining local
option is S2DC: use curvature-square terms as finite-cutoff regulators and
renormalized contact owners, send their unwanted poles beyond every fixed
physical scale, and obtain nonlinear no-escape from a different estimate.
No such estimate is yet present in the attached programme.
\end{assessment}

\begin{programme}[V63: can a positive spectral form factor control
curvature?]
\label{prog:newS7V62-V63}
V63 will attack CPRC's positive-spectral branch.  Starting from a
Källén--Lehmann covariance, it will characterize the ultraviolet growth
allowed for its inverse and test whether a nonlocal complete-Bernstein form
factor can be both reflection positive and strong enough to suppress the
V59 conformal sequence.  If positive spectral weight forbids the required
fourth-order coercivity for an elementary metric mode, that no-go will be
stated and the programme will return to constraints or contour
reconstruction.
\end{programme}

\begin{firewall}[Claim boundary after seventh-series V62]
\label{fire:newS7V62-boundary}
V62 proves that linearized scalar and Weyl curvatures descend to the
diffeomorphism quotient; computes the positive Einstein--Weyl TT Hessian;
exhibits its negative massive spin-two residue; proves that the residue is
visible in the gauge-invariant Weyl algebra; and gives a distributional
decoupling lemma and the exact S2DC requirements.

V62 does not prove nonlinear failure of every constrained physical algebra,
construct S2DC, supply alternative curvature tightness, prove RCNC,
construct an interacting measure, remove the cutoff, or construct the
Standard Model--Einstein continuum.
\end{firewall}

\begin{theorem}[Seventh-series V62 result]
\label{thm:newS7V62-result}
The linearized diffeomorphism quotient does not erase either obstruction
found in V59--V61.  The conformal mode is detected by the gauge-invariant
scalar curvature, while a finite positive Weyl-square kinetic coefficient
produces a negative massive spin-two pole detected by the gauge-invariant
linearized Weyl tensor.  The pole can disappear from fixed scales when its
coefficient tends to zero, but then positive Weyl-square coercivity cannot
supply the uniform nonlinear no-escape estimate.  A valid local-regulator
continuation must therefore prove S2DC and obtain that estimate elsewhere,
or construct a physical constraint which cancels the pole.
\end{theorem}

\begin{proof}
Use \Cref{lem:newS7V62-invariance,
lem:newS7V62-TT-action,
thm:newS7V62-spin2-ghost,
cor:newS7V62-branch1,
prop:newS7V62-scalar,
lem:newS7V62-decoupling,
prop:newS7V62-S2DC}.
\end{proof}

\paragraph{Parent-version record.}
V62 was formed from
\texttt{Renewal\_SM\_Einstein\_Continuum\_Research\_Notes\_V61.tex}.
The V61 parent had \(24\,433\) logical source lines, \(982\,445\) bytes,
and SHA--256
\[
 \texttt{103E39C5B66373DB7D0EFEEBE6D478B403AF9B1B634C50F426E618F5263A854A}.
\]
The next compulsory companion audit is V65.

% =============================================================================
% END APPENDED SEVENTH-SERIES V62 MATERIAL
% =============================================================================

% =============================================================================
% APPENDED SEVENTH-SERIES V63 MATERIAL
% =============================================================================

\section{Positive spectral form factors cannot retain the physical
conformal Einstein slope}
\label{sec:newS7V63}
\label{sec:v7v63}

The third CPRC branch proposed a nonlocal covariance with positive spectral
measure.  This note tests that possibility at the quadratic level without
assuming locality or rationality.  Positivity itself forces two facts: the
inverse covariance is increasing in \(p^2\), and it grows at most linearly
along the positive momentum axis.  The first contradicts the physical-sign
Einstein conformal slope; the second forbids fourth-order ultraviolet
coercivity for an elementary positive-metric conformal scalar.

\subsection{Monotonicity and the ultraviolet ceiling}
\label{sec:newS7V63-ceiling}

\begin{theorem}[Positive spectral inverse theorem]
\label{thm:newS7V63-inverse}
Let
\begin{equation}
 C(z)=\int_{[0,\infty)}{d\rho(s)\over z+s},
 \qquad z>0,
 \label{eq:newS7V63-Stieltjes}
\end{equation}
where \(\rho\) is a nonzero positive measure and \(C(z)<\infty\) for
every \(z>0\).  Put \(F(z)=C(z)^{-1}\).  Then:
\begin{enumerate}
\item \(C\) is strictly decreasing and \(F\) is strictly increasing;
\item there exist \(S<\infty\) and \(m>0\) such that
\begin{equation}
 C(z)\geq {m\over z+S},
 \qquad
 F(z)\leq {z+S\over m}
 \quad(z>0);
 \label{eq:newS7V63-linear-ceiling}
\end{equation}
\item in particular, \(F(z)\) cannot have a positive \(z^2\) leading
growth and cannot have a negative derivative on any interval.
\end{enumerate}
\end{theorem}

\begin{proof}
For \(z_2>z_1>0\), the integrand
\((z+s)^{-1}\) is strictly smaller at \(z_2\) for every \(s\).
Integration against a nonzero positive measure gives
\(C(z_2)<C(z_1)\), and reciprocation gives
\(F(z_2)>F(z_1)\).  Because \(\rho\ne0\), monotone exhaustion of
\([0,\infty)\) by \([0,S]\) gives an \(S<\infty\) for which
\(m:=\rho([0,S])>0\).  Then
\[
 C(z)\geq\int_{[0,S]}{d\rho(s)\over z+s}
 \geq {m\over z+S}.
\]
Reciprocation proves \eqref{eq:newS7V63-linear-ceiling}.  The last claims
follow immediately.
\end{proof}

\begin{corollary}[No positive-spectral fourth-order elementary mode]
\label{cor:newS7V63-no-fourth}
An elementary scalar Gaussian satisfying Euclidean invariance and
Osterwalder--Schrader positivity cannot have an inverse covariance which is
coercive like \(c|p|^4\) at high momentum.
\end{corollary}

\begin{proof}
The Källén--Lehmann representation of
\Cref{lem:newS7V61-KL} puts its covariance in the form
\eqref{eq:newS7V63-Stieltjes}.  The inverse is \(O(|p|^2)\) by
\eqref{eq:newS7V63-linear-ceiling}, rather than bounded below by a positive
multiple of \(|p|^4\).
\end{proof}

This is stronger than the polynomial no-go of V61: adding infinitely many
positive massive poles or a positive continuous spectral density does not
remove the ultraviolet ceiling.

\subsection{The low-momentum Einstein sign}
\label{sec:newS7V63-low}

\begin{theorem}[Physical conformal slope obstruction]
\label{thm:newS7V63-slope}
Suppose the conformal factor is an elementary positive-metric scalar and
its stable quadratic inverse covariance has a finite expansion
\[
 F(z)=F(0)+f_1z+o(z)
 \qquad(z\downarrow0).
\]
Reflection positivity requires \(f_1\geq0\).  The physical-sign Euclidean
Einstein term contributes \(f_1^{\rm EH}=-6\kappa<0\).  Hence a
reflection-positive total covariance cannot retain that Einstein
coefficient as its complete low-energy conformal kinetic slope.
\end{theorem}

\begin{proof}
By \Cref{thm:newS7V63-inverse}, \(F\) is increasing.  Therefore every
right derivative at zero, when finite, is nonnegative.  The conformal
formula \eqref{eq:newS7V61-quadratic} gives the coefficient
\(-6\kappa\) of \(z\).  If all other terms are \(o(z)\), the total slope is
negative, a contradiction.  If another sector changes the total slope to a
nonnegative number, the complete low-energy conformal kinetic coefficient
is no longer the physical Einstein one; that is a different resolution
which must be declared.
\end{proof}

\begin{corollary}[A mass term does not hide the sign]
\label{cor:newS7V63-mass}
Choosing a sufficiently positive zeroth-order Hessian so that
\(F(z)=36\alpha z^2-6\kappa z+8\Lambda\) stays positive, as in V61, does
not restore reflection positivity: \(F\) still decreases near \(z=0\).
\end{corollary}

\begin{proof}
\(F'(0)=-6\kappa<0\), contrary to
\Cref{thm:newS7V63-inverse}.
\end{proof}

\subsection{Why ultraviolet linear growth is not the V59 coercive estimate}
\label{sec:newS7V63-sequence}

\begin{proposition}[Spectral quadratic forms have only second-order
control on the V59 sequence]
\label{prop:newS7V63-sequence}
Let \(F\) satisfy \eqref{eq:newS7V63-linear-ceiling}, and let
\(\phi_N=\varepsilon\sin(Nx^1)\) on the flat torus.  Its translation-
invariant quadratic action obeys
\begin{equation}
 \langle\phi_N,F(-\Delta)\phi_N\rangle
 \leq C\varepsilon^2(N^2+S)
 \label{eq:newS7V63-sequence-bound}
\end{equation}
for a constant independent of \(N\).  It cannot yield the positive
\(N^4\) control supplied by \(R^2\) in
\eqref{eq:newS7V59-R2scale}.
\end{proposition}

\begin{proof}
\(\phi_N\) is supported on the two Fourier modes with
\(|p|^2=N^2\).  Hence the quadratic form is a fixed normalization factor
times \(\varepsilon^2F(N^2)\).  Apply
\eqref{eq:newS7V63-linear-ceiling}.
\end{proof}

Second-order control could dominate a negative second-order Einstein term
only by changing its net ultraviolet coefficient.  The slope theorem shows
that, if the conformal factor is physical at low momentum, positive
spectral weight also changes the forbidden negative infrared slope.  There
is no interpolation through a positive elementary covariance which keeps
that slope and becomes fourth-order in the ultraviolet.

\subsection{Contact terms and subtractions}
\label{sec:newS7V63-contact}

\begin{lemma}[Local contact terms do not repair spectral positivity]
\label{lem:newS7V63-contact}
Adding a polynomial \(P(z)\) to a covariance changes only distributions
supported at coincident points.  It cannot change the sign of an isolated
noncontact pole, and it cannot turn the fourth-order covariance
\(1/Q(z)\) of V61 into a positive Källén--Lehmann covariance at separated
times.
\end{lemma}

\begin{proof}
Fourier transformation sends a polynomial in momentum to a finite linear
combination of derivatives of the delta distribution.  Test functions
whose reflected supports are separated from the reflection plane pair to
zero with those terms.  The negative-pole test in
\Cref{prop:newS7V61-residue} uses such separated supports, so its sign is
unchanged.
\end{proof}

\begin{assessment}[Anomaly owners are not propagator cures]
\label{ass:newS7V63-anomaly}
The local \(C^2\), \(R^2\), and Euler densities required to renormalize
curved-background composite observables may remain as contact owners in
the effective action and stress.  Their existence does not justify
resumming them into a positive fundamental metric covariance.  V54--V58
fix those local owners; V61--V63 show that a finite fourth-order resummation
creates a different spectral problem.
\end{assessment}

\subsection{Decision on the positive-spectral CPRC branch}
\label{sec:newS7V63-decision}

\begin{theorem}[Closure of the elementary positive-spectral branch]
\label{thm:newS7V63-closure}
There is no nonzero Euclidean-invariant, reflection-positive elementary
conformal Gaussian which simultaneously has:
\begin{enumerate}
\item the physical Einstein conformal slope \(-6\kappa z\) at low
momentum with \(\kappa>0\);
\item a stable positive inverse covariance; and
\item fourth-order ultraviolet coercivity sufficient to reproduce the
\(R^2\) control of the V59 sequence.
\end{enumerate}
\end{theorem}

\begin{proof}
Items 1 and 2 contradict the monotonicity conclusion of
\Cref{thm:newS7V63-slope}.  Independently, items 2 and 3 contradict the
linear ceiling of \Cref{cor:newS7V63-no-fourth} and the explicit sequence
bound \eqref{eq:newS7V63-sequence-bound}.
\end{proof}

\begin{corollary}[Remaining CPRC routes]
\label{cor:newS7V63-routes}
At quadratic level, the unresolved routes are reduced to:
\begin{enumerate}
\item impose and quantize gravitational constraints so the conformal factor
is not an elementary member of the physical positive algebra;
\item use a complex conformal contour and reconstruct positivity only after
the contour integral; or
\item change the low-energy theory or observable assignment explicitly.
\end{enumerate}
An unconstrained positive-spectral form factor for the elementary conformal
factor is not a remaining route.
\end{corollary}

\begin{proof}
Combine \Cref{thm:newS7V63-closure} with the scope statement of V61 and the
physical Weyl-algebra result of V62.
\end{proof}

\begin{definition}[Linear constraint-branch card, LCBC]
\label{def:newS7V63-LCBC}
Before a nonlinear constraint construction can be used in RCNC, its
linearized limit must:
\begin{enumerate}
\item give a symplectic or Euclidean quotient with exactly the two
transverse spin-two polarizations per nonzero momentum in vacuum;
\item impose the Hamiltonian and momentum constraints before assigning the
positive physical inner product;
\item yield a positive massless spin-two spectral measure on
gauge-invariant observables;
\item show how matter sources determine the nondynamical scalar and vector
potentials without reintroducing a negative physical norm; and
\item reproduce the V58 smooth metric response and Ward identity.
\end{enumerate}
\end{definition}

\begin{programme}[V64: solve the free linear constraint branch]
\label{prog:newS7V63-V64}
V64 will carry out the York decomposition for every nonzero flat-torus
momentum, impose the linear Hamiltonian and momentum constraints, quotient
linearized diffeomorphisms, and compute the reduced Euclidean quadratic
form.  It will determine exactly which rows of LCBC hold in free vacuum and
which fail or become nonlocal after matter sources are admitted.
\end{programme}

\begin{firewall}[Claim boundary after seventh-series V63]
\label{fire:newS7V63-boundary}
V63 proves monotonicity and a linear ultraviolet ceiling for the inverse of
every nonzero positive Källén--Lehmann covariance; proves that the physical
Euclidean conformal Einstein slope violates that monotonicity; applies the
ceiling to the explicit high-frequency V59 sequence; and closes the
unconstrained elementary positive-spectral CPRC branch.

V63 does not close the gravitational constraint or contour branches, prove
LCBC, construct RCNC, give nonlinear no-escape, remove the cutoff, or
construct the Standard Model--Einstein continuum.
\end{firewall}

\begin{theorem}[Seventh-series V63 result]
\label{thm:newS7V63-result}
Positive spectral weight cannot reconcile the physical conformal Einstein
sign with ultraviolet curvature-square stabilization when the conformal
factor is an elementary physical scalar.  Its inverse covariance must be
increasing and at most linear in \(p^2\), whereas the Einstein conformal
slope is negative and the V59 stabilizer grows quadratically in \(p^2\).
The programme must therefore impose the gravitational constraints before
positivity, use a reconstructed contour, or explicitly change the theory.
\end{theorem}

\begin{proof}
Use \Cref{thm:newS7V63-inverse,
cor:newS7V63-no-fourth,
thm:newS7V63-slope,
prop:newS7V63-sequence,
lem:newS7V63-contact,
thm:newS7V63-closure}.
\end{proof}

\paragraph{Parent-version record.}
V63 was formed from
\texttt{Renewal\_SM\_Einstein\_Continuum\_Research\_Notes\_V62.tex}.
The V62 parent had \(24\,788\) logical source lines, \(996\,182\) bytes,
and SHA--256
\[
 \texttt{C32830E8A27B50FF7AC935164D69A1969F70DB9A0D5E7E6555D94F2B98DB637E}.
\]
The next compulsory companion audit is V65.

% =============================================================================
% END APPENDED SEVENTH-SERIES V63 MATERIAL
% =============================================================================

% =============================================================================
% APPENDED SEVENTH-SERIES V64 MATERIAL
% =============================================================================

\section{The free constraint branch and a reflection-positive spatial
regulator}
\label{sec:newS7V64}
\label{sec:v7v64}

V63 reduced the remaining local choices to a constraint-first
quantization or a reconstructed contour.  This note solves the linear
vacuum part of the constraint branch.  On every nonzero spatial momentum,
the Hamiltonian and momentum constraints together with their gauge orbits
remove the scalar and vector canonical pairs before Wick rotation.  The
reduced Euclidean field has two transverse-traceless polarizations and a
positive second-order time covariance.  This also identifies a viable
finite-cutoff regulator pattern: higher \emph{spatial} derivatives can
improve ultraviolet control while preserving time-reflection positivity.

\subsection{York decomposition on a flat spatial torus}
\label{sec:newS7V64-York}

Work on \(\mathbb R_t\times\mathbb T^3\) and linearize the ADM spatial
metric and momentum as
\[
 \gamma_{ij}=\delta_{ij}+q_{ij},
 \qquad
 \pi^{ij}=p^{ij}.
\]
For each nonzero spatial Fourier mode the York decomposition is

\begin{equation}
 q_{ij}=q_{ij}^{\rm TT}
 +2\partial_{(i}v^{T}_{j)}
 +\left(\partial_i\partial_j-{1\over3}\delta_{ij}\Delta\right)\sigma
 +{1\over3}\delta_{ij}\tau,
 \label{eq:newS7V64-York-q}
\end{equation}

where
\[
 \partial_iq_{ij}^{\rm TT}=0,\quad q_{ii}^{\rm TT}=0,\quad
 \partial_iv_i^T=0,\quad \tau=q_{ii}.
\]
There is an analogous decomposition

\begin{equation}
 p_{ij}=p_{ij}^{\rm TT}
 +2\partial_{(i}w^{T}_{j)}
 +\left(\partial_i\partial_j-{1\over3}\delta_{ij}\Delta\right)\rho
 +{1\over3}\delta_{ij}p.
 \label{eq:newS7V64-York-p}
\end{equation}

\begin{lemma}[Uniqueness at nonzero momentum]
\label{lem:newS7V64-York}
For every nonzero Fourier momentum \(k\), decompositions
\eqref{eq:newS7V64-York-q} and
\eqref{eq:newS7V64-York-p} are unique and orthogonal in the Euclidean
tensor inner product.
\end{lemma}

\begin{proof}
Choose an orthonormal basis whose third vector is \(k/|k|\).
Symmetric tensors split orthogonally into the two-dimensional TT block,
two mixed transverse-longitudinal entries, one traceless scalar, and one
trace scalar.  Their dimensions are \(2+2+1+1=6\).  The coefficients in
\eqref{eq:newS7V64-York-q} are nonzero on \(k\ne0\), so the decomposition is
unique.  The same argument applies to \(p_{ij}\).
\end{proof}

\subsection{Constraints and gauge orbits}
\label{sec:newS7V64-constraints}

Up to a positive normalization, the linear vacuum ADM constraints are

\begin{equation}
 \mathcal H^{(1)}
 =\partial_i\partial_jq_{ij}-\Delta q_{ii}=0,
 \qquad
 \mathcal H_i^{(1)}=-2\partial_jp_{ij}=0.
 \label{eq:newS7V64-constraints}
\end{equation}

\begin{lemma}[The scalar and vector constraint equations]
\label{lem:newS7V64-constraint-components}
On the York variables at nonzero momentum,

\begin{align}
 \mathcal H^{(1)}
 &={2\over3}\Delta(\Delta\sigma-\tau),
 \label{eq:newS7V64-H-scalar}\\
 \partial_jp_{ij}
 &=\Delta w_i^T
 +{1\over3}\partial_i(2\Delta\rho+p).
 \label{eq:newS7V64-M-components}
\end{align}

Hence the vacuum constraints imply
\[
 \tau=\Delta\sigma,\qquad w_i^T=0,\qquad p=-2\Delta\rho.
\]
\end{lemma}

\begin{proof}
Double divergence of the traceless scalar term in
\eqref{eq:newS7V64-York-q} is
\((2/3)\Delta^2\sigma\), while double divergence of the trace term is
\((1/3)\Delta\tau\).  Subtracting \(\Delta\tau\) gives
\eqref{eq:newS7V64-H-scalar}.  Taking one divergence of
\eqref{eq:newS7V64-York-p} gives
\eqref{eq:newS7V64-M-components}.  Transverse and longitudinal parts are
orthogonal and \(\Delta\) is invertible at \(k\ne0\), proving the last
claims.
\end{proof}

Write a spatial gauge parameter as
\(\xi_i=\xi_i^T+\partial_i\chi\).  Its action is
\[
 \delta q_{ij}=2\partial_{(i}\xi_{j)},\qquad
 \delta v_i^T=\xi_i^T,\qquad
 \delta\sigma=2\chi,\qquad
 \delta\tau=2\Delta\chi.
\]
Thus \(\tau-\Delta\sigma\) is gauge invariant and is exactly the scalar
set to zero by the Hamiltonian constraint.

\begin{theorem}[Linear vacuum reduced phase space]
\label{thm:newS7V64-reduction}
For each nonzero spatial momentum, imposing the four linear ADM
constraints and quotienting their four gauge flows leaves precisely the
canonical pair
\[
 (q_{ij}^{\rm TT},p_{ij}^{\rm TT}),
\]
with two independent polarizations.  No scalar conformal canonical pair is
present in the reduced vacuum phase space.
\end{theorem}

\begin{proof}
The three momentum constraints remove the two transverse components
\(w^T\) and one scalar momentum combination; their spatial-diffeomorphism
orbits remove \(v^T\) and the corresponding scalar coordinate.
The Hamiltonian constraint sets the gauge-invariant scalar
\(\tau-\Delta\sigma\) to zero, and its normal-deformation gauge flow removes
the conjugate scalar momentum.  Equivalently, four first-class constraints
remove eight dimensions from the twelve-dimensional symmetric-tensor
phase space, leaving four dimensions, exactly the two TT canonical pairs.
The component equations in
\Cref{lem:newS7V64-constraint-components} identify those pairs explicitly.
\end{proof}

\begin{remark}[Zero modes]
\label{rem:newS7V64-zero}
The theorem excludes \(k=0\), where \(\Delta\) is not invertible.  Global
volume, flat-torus moduli, and their conjugate momenta require a separate
finite-dimensional constraint and normalization.  They cannot be discarded
as if covered by the nonzero-mode argument.
\end{remark}

\subsection{Positive reduced vacuum covariance}
\label{sec:newS7V64-positive}

\begin{lemma}[Reduced Euclidean action]
\label{lem:newS7V64-reduced-action}
After solving the linear vacuum constraints and Wick rotating the reduced
time variable, the quadratic action has the form

\begin{equation}
 S_{E,{\rm red}}^{(2)}
 ={K\over2}\sum_{\lambda=1}^2
 \int_{\mathbb R\times\mathbb T^3}
 \left(|\partial_\tau q_\lambda^{\rm TT}|^2
       |\nabla q_\lambda^{\rm TT}|^2\right)\,d\tau\,d^3x,
 \qquad K>0.
 \label{eq:newS7V64-red-action}
\end{equation}
\end{lemma}

\begin{proof}
The ADM kinetic form restricted to TT momenta is positive, and the
physical-sign Euclidean Einstein Hessian restricted to TT metric
perturbations is positive by \Cref{lem:newS7V62-TT-action} with
\(\beta=0\).  Eliminating \(p^{\rm TT}\) from the first-order reduced
action and Wick rotating gives \eqref{eq:newS7V64-red-action}; the positive
constant \(K\) depends only on the normalization of \(\kappa\) and \(q\).
\end{proof}

\begin{theorem}[Free reduced reflection positivity]
\label{thm:newS7V64-free-OS}
The Gaussian measure determined by
\eqref{eq:newS7V64-red-action}, with its spatial zero mode treated
separately, is reflection positive in Euclidean time on the algebra
generated by the two TT modes.
\end{theorem}

\begin{proof}
For fixed nonzero spatial momentum \(k\) and either polarization, the
covariance is
\[
 C_k(\tau)={e^{-|k||\tau|}\over2K|k|}.
\]
If \(f(\tau,k)\) is supported at \(\tau>0\), its reflected quadratic form is
\[
 \int {d^3k\over2K|k|}
 \left|\int_0^\infty e^{-|k|\tau}f(\tau,k)\,d\tau\right|^2\geq0.
\]
Summing the two polarizations proves Gaussian reflection positivity.
\end{proof}

\begin{assessment}[Free-vacuum LCBC rows]
\label{ass:newS7V64-LCBC}
The nonzero-mode free vacuum satisfies LCBC rows 1--3: constraints precede
the inner product, exactly two TT polarizations remain, and their massless
spectral weight is positive.  The zero modes, matter-sourced constraints,
and agreement with the nonlinear V58 response remain open, so LCBC is not
closed.
\end{assessment}

\subsection{Matter makes the reduced action nonlocal}
\label{sec:newS7V64-matter}

With matter, the constraints take the schematic linear form
\[
 \mathcal H^{(1)}=c_E T_{00},
 \qquad
 \mathcal H_i^{(1)}=c_P T_{0i}.
\]

\begin{proposition}[Sourced scalar and vector potentials]
\label{prop:newS7V64-sources}
At every nonzero spatial momentum, the matter energy and momentum densities
determine the gauge-invariant scalar and vector constraint variables through
inverse powers of \(\Delta\).  Substitution into the reduced action therefore
produces instantaneous spatially nonlocal matter interactions.  Free TT
reflection positivity alone does not determine the sign or
Osterwalder--Schrader positivity of the coupled reduced measure.
\end{proposition}

\begin{proof}
Equation \eqref{eq:newS7V64-H-scalar} gives
\[
 \Delta\sigma-\tau={3c_E\over2}\Delta^{-1}T_{00}.
\]
The transverse and longitudinal parts of
\eqref{eq:newS7V64-M-components} similarly solve for constraint momenta by
\(\Delta^{-1}T_{0i}\).  Eliminating these variables inserts the Green
operator \(\Delta^{-1}\) into the source action.  The proof of
\Cref{thm:newS7V64-free-OS} contains no estimate for that nonlocal
interaction or for the determinant and Jacobian of the constraint map, so
it cannot imply positivity of the coupled measure.
\end{proof}

\subsection{Higher spatial derivatives preserve time positivity}
\label{sec:newS7V64-spatial-regulator}

The four-derivative obstruction in V61--V62 came from a polynomial in the
full Euclidean \(p_0^2+|k|^2\), which introduced a second time pole.  A
transfer-matrix cutoff can instead keep two time derivatives and add
higher spatial derivatives.

\begin{theorem}[Reflection-positive spatial fourth-order regulator]
\label{thm:newS7V64-spatial-OS}
Let
\[
 \omega_a(k)^2=c_2(a)|k|^2+c_4(a)|k|^4+m_a^2,
 \qquad
 c_2(a)>0,\quad c_4(a)\geq0,\quad m_a^2\geq0.
\]
The quadratic action

\begin{equation}
 S_{a,{\rm TT}}
 ={1\over2}\sum_{\lambda=1}^2\int
 \left(|\partial_\tau q_\lambda|^2
 q_\lambda\,\omega_a(-i\nabla)^2q_\lambda\right)
 d\tau\,d^3x
 \label{eq:newS7V64-anisotropic-action}
\end{equation}

is time-reflection positive for every cutoff \(a\).  Its spatial
high-frequency action grows like \(c_4(a)|k|^4\) when \(c_4(a)>0\), without
introducing the negative second time pole of
\eqref{eq:newS7V62-D2}.
\end{theorem}

\begin{proof}
At fixed \(k\), the covariance is
\[
 C_{a,k}(\tau)
 ={e^{-\omega_a(k)|\tau|}\over2\omega_a(k)}
\]
when \(\omega_a(k)>0\), with the massless zero mode separated.  For a
positive-time test function, the reflected form is
\[
 \int {d^3k\over2\omega_a(k)}
 \left|\int_0^\infty
 e^{-\omega_a(k)\tau}f(\tau,k)\,d\tau\right|^2\geq0.
\]
There is one positive time pole for each \(k\).  The last assertion follows
from the definition of \(\omega_a(k)^2\).
\end{proof}

\begin{corollary}[A nonempty regulator route]
\label{cor:newS7V64-route}
The V62 spin-two no-go does not rule out a constraint-first transfer-matrix
regularization with higher spatial curvature and only second-order time
evolution.  If \(c_4(a)\to0\) at fixed physical momentum, its free
dispersion converges to the relativistic massless TT dispersion while
remaining reflection positive at every cutoff.
\end{corollary}

\begin{proof}
\Cref{thm:newS7V64-spatial-OS} proves cutoff positivity.  Pointwise
\(\omega_a(k)^2\to c_2|k|^2\) under the stated limit, so the fixed-momentum
covariances converge to those of \Cref{thm:newS7V64-free-OS}.
\end{proof}

\begin{definition}[Reduced transfer-matrix gravity card, RTGC]
\label{def:newS7V64-RTGC}
The constraint-first branch satisfies RTGC if it provides:
\begin{enumerate}
\item a finite lattice phase space, constraint surface, gauge quotient or
BRST complex, and a positive one-step transfer operator on the physical
space;
\item a higher-spatial-derivative or other ultraviolet estimate uniform
enough for physical cylinder-law tightness, without a higher-time-derivative
pole;
\item a treatment of global zero modes and a proof that matter-sourced
constraint kernels preserve normalization and physical positivity;
\item restoration of spacetime Euclidean invariance and S2DC at fixed
physical momenta;
\item convergence of the common renormalized determinant composites and
stress contacts of V57--V58; and
\item RCNC for the resulting physical observable algebra.
\end{enumerate}
\end{definition}

\begin{assessment}[What the constraint branch has gained]
\label{ass:newS7V64-status}
The linear vacuum obstruction is resolved on reduced phase space: the
conformal scalar is removed by a constraint and gauge flow, the two TT
modes have positive spectral weight, and a spatial fourth-order regulator
can improve cutoff control without a negative time pole.  The price is an
anisotropic, constraint-first construction.  Its nonlinear constraint
Jacobian, zero modes, matter-induced nonlocality, restoration of spacetime
symmetry, and renormalized composite stress are exactly the unresolved
RTGC rows.
\end{assessment}

\begin{programme}[V65: companion audit and the finite transfer kernel]
\label{prog:newS7V64-V65}
V65 will perform the compulsory companion audit and then write the
finite-dimensional transfer-kernel criterion which makes a constrained
one-step kernel a positive contraction.  It will test whether the current
renewal conditional kernels supply the required symmetry, normalization,
constraint projection, and spatial-coercivity rows of RTGC, recording the
first absent row rather than assuming a transfer matrix.
\end{programme}

\begin{firewall}[Claim boundary after seventh-series V64]
\label{fire:newS7V64-boundary}
V64 proves the nonzero-mode linear York decomposition; solves the vacuum
linear constraints and quotient; obtains exactly two TT canonical pairs;
proves reflection positivity of their reduced Euclidean Gaussian; and
proves time-reflection positivity of a higher-spatial-derivative regulator.

V64 does not handle global zero modes, prove positivity after matter
constraints, construct a nonlinear constraint measure or transfer
operator, restore spacetime Euclidean invariance, prove RTGC or RCNC,
remove the cutoff, or construct the Standard Model--Einstein continuum.
\end{firewall}

\begin{theorem}[Seventh-series V64 result]
\label{thm:newS7V64-result}
Constraint-first quantization resolves the free nonzero-mode conformal
problem at linear order: the ADM constraints and their gauge flows remove
all scalar and vector canonical pairs, leaving two positive massless TT
polarizations.  Higher spatial derivatives can regulate these modes while
preserving a single positive time pole and exact time-reflection
positivity.  Extending this nonempty route to the interacting
Standard Model--Einstein system is precisely RTGC, whose matter,
zero-mode, nonlinear, symmetry-restoration, and composite-stress rows
remain open.
\end{theorem}

\begin{proof}
Use \Cref{lem:newS7V64-York,
lem:newS7V64-constraint-components,
thm:newS7V64-reduction,
lem:newS7V64-reduced-action,
thm:newS7V64-free-OS,
prop:newS7V64-sources,
thm:newS7V64-spatial-OS,
cor:newS7V64-route}.
\end{proof}

\paragraph{Parent-version record.}
V64 was formed from
\texttt{Renewal\_SM\_Einstein\_Continuum\_Research\_Notes\_V63.tex}.
The V63 parent had \(25\,087\) logical source lines, \(1\,008\,227\) bytes,
and SHA--256
\[
 \texttt{E68FED8420C2661B55CC071C87B0939FAA313F9038CF76BD1B6FEE793500F8C2}.
\]
The next compulsory companion audit is V65.

% =============================================================================
% END APPENDED SEVENTH-SERIES V64 MATERIAL
% =============================================================================

% =============================================================================
% APPENDED SEVENTH-SERIES V65 MATERIAL
% =============================================================================

\section{Companion audit and the constrained one-step kernel criterion}
\label{sec:newS7V65}
\label{sec:v7v65}

This is the compulsory V65 companion audit.  Its Yang--Mills input meets
the present bottleneck unusually closely: the flat inverse and local
writer are controlled by a uniform gap, but the nonlinear covariance still
requires a connected expansion in which vacuum volume cancels.  The audit
is followed by an exact finite-cutoff transfer theorem and a comparison
with the conditional-kernel claims in the attached gravitational sources.

\subsection{Compulsory V65 companion audit}
\label{sec:newS7V65-audit}

The newest complete companion checkpoints inspected read-only were

\begin{align}
 &\texttt{Renewal\_Yang\_Mills\_Mass\_Gap\_Research\_Notes\_V37.tex},
 \notag\\
 &\qquad 590\,146\ {\rm bytes},\quad 16\,816\ {\rm logical\ lines},\quad
 \texttt{6B2A98BEF53BCAD8B6ECA6D27E2131FA11A3056AF7B80261512D298AAD9E93F6},
 \label{eq:newS7V65-YM-hash}\\
 &\texttt{Renewal\_NS\_Stability\_Research\_Notes\_V26.tex},
 \notag\\
 &\qquad 278\,283\ {\rm bytes},\quad 5\,319\ {\rm logical\ lines},\quad
 \texttt{82C887F8C2C69E4F28FAD7C3DC8DE5B9099CB5A4BF431EBA1FFF3E91935978AD}.
 \label{eq:newS7V65-NS-hash}
\end{align}

Each has one live document terminator.

\paragraph{Yang--Mills V37.}
From a finite-range vertical Hessian with the uniform spectral interval
\([1/60,8]\), the note proves exponential decay of the inverse by a
Neumann-series range argument.  It then proves an exact covariance Duhamel
identity and identifies the nonlinear target as a connected three-point
cumulant with a spatially summable tail.  Its firewall is essential here:
an extensive absolute Taylor remainder is not a connected bound, and a
small bad-link density does not control coherent bad polymers.

\paragraph{Navier--Stokes V26.}
The note computes exact rank-one kernel norms and obtains a modest
improvement in one second-derivative constant, while explicitly preserving
the open global-regularity boundary.  This changes no gravitational
constraint or transfer-kernel row.

\begin{theorem}[V65 companion-audit decision]
\label{thm:newS7V65-audit}
The finite-range inverse proof and connected Duhamel identity of YM V37 are
valid templates for a reduced gravitational detail field only after RTGC
provides a uniformly positive reduced Hessian and a common integration
chart.  None of its numerical gap, Wilson compactness, or polymer
activities transfers to the gravitational variables.  NS V26 supplies no
input to RTGC.
\end{theorem}

\begin{proof}
The Neumann proof uses only finite range and a two-sided spectral bound, so
its algebra is reusable.  Its constants and nonlinear estimates use the
compact Wilson fibre and its action, absent in the gravitational problem.
The NS result concerns Oseen-kernel rank-one norms, which do not occur in
the ADM transfer kernel.
\end{proof}

\subsection{Reversibility is weaker than transfer positivity}
\label{sec:newS7V65-reversibility}

\begin{proposition}[A reversible Markov kernel may have a negative transfer
eigenvalue]
\label{prop:newS7V65-reversible-no}
On the two-point space with uniform stationary measure, the kernel
\[
 P=\begin{pmatrix}0&1\\1&0\end{pmatrix}
\]
is nonnegative, normalized, stationary, and reversible, but its transfer
operator is not positive semidefinite.
\end{proposition}

\begin{proof}
The matrix is symmetric and each row sums to one.  Its eigenvectors
\((1,1)\) and \((1,-1)\) have eigenvalues \(1\) and \(-1\), respectively.
\end{proof}

Thus pointwise positivity, normalization, and detailed balance do not by
themselves give link-reflection positivity.  The missing statement is
operator positivity.

\begin{theorem}[Finite one-step Osterwalder--Schrader criterion]
\label{thm:newS7V65-one-step}
Let \((X,\pi)\) be a finite probability space and let \(T\) be the
operator of a stationary Markov kernel.  Assume
\[
 T1=1,\qquad T=T^*,\qquad T\succeq0
 \quad\hbox{on }L^2(X,\pi).
 \label{eq:newS7V65-T-card}
\]
Then \(0\preceq T\preceq I\).  Every finite stationary path measure built
from \(T\) is reflection positive across a link, and on the support of
\(T\)
\[
 H_a=-a^{-1}\log T
\]
is a nonnegative self-adjoint transfer Hamiltonian.
\end{theorem}

\begin{proof}
A Markov operator is an \(L^2(\pi)\) contraction by Jensen's inequality
and stationarity.  Self-adjointness puts its spectrum in \([-1,1]\), while
\(T\succeq0\) improves this to \([0,1]\), proving the first assertion and
the spectral definition of \(H_a\).

For a positive-time cylinder function \(F\), integrate all variables
strictly inside each reflected half-path.  The Markov property produces
one boundary function \(f_F\in L^2(X,\pi)\).  Reflection across the central
link gives
\[
 \mathbb E[\overline{\Theta F}F]
 =\langle f_F,Tf_F\rangle_{L^2(\pi)}\geq0.
\]
Linear combinations and density give reflection positivity for the finite
cylinder algebra.
\end{proof}

\begin{corollary}[The reflection Gram is the finite boundary form]
\label{cor:newS7V65-Gram}
In a basis of boundary functions, link-reflection positivity is equivalent
to positive semidefiniteness of the matrix representing \(T\).  Pointwise
positive matrix entries are insufficient.
\end{corollary}

\begin{proof}
The reflected form in the proof of
\Cref{thm:newS7V65-one-step} is the quadratic form of \(T\).  The
equivalence follows by testing every coefficient vector; the insufficiency
is \Cref{prop:newS7V65-reversible-no}.
\end{proof}

\subsection{A constructive positive factorization}
\label{sec:newS7V65-factorization}

\begin{theorem}[Sandwiched heat-kernel construction]
\label{thm:newS7V65-sandwich}
Let \(G_a=e^{-aK_a}\) be a positive self-adjoint contraction on a
finite-dimensional Hilbert space, let \(W_a=e^{-aV_a/2}\) be a bounded
positive multiplication operator, and let \(P_a=P_a^*=P_a^2\) be an
orthogonal constraint projector.  Then on
\(\mathcal H_{a,\rm phys}=\operatorname{Ran}P_a\),

\begin{equation}
 T_a=P_aW_aG_aW_aP_a
 \label{eq:newS7V65-sandwich}
\end{equation}

is self-adjoint and positive semidefinite.  After division by its largest
nonzero eigenvalue it is a positive contraction.  If its coordinate kernel
is positivity improving, a positive top eigenvector gives a reversible
normalized Markov kernel by the ground-state transform.
\end{theorem}

\begin{proof}
For \(\psi\in\operatorname{Ran}P_a\),
\[
 \langle\psi,T_a\psi\rangle
 =\langle W_aP_a\psi,G_aW_aP_a\psi\rangle\geq0.
\]
The product is manifestly self-adjoint.  Spectral normalization gives a
positive contraction.  If \(T_a\psi_0=\lambda_0\psi_0\) with
\(\psi_0(x)>0\), define
\[
 {\cal P}_a(x,y)
 ={T_a(x,y)\psi_0(y)\over\lambda_0\psi_0(x)}.
\]
The eigenvalue equation makes its row sum one.  Symmetry of \(T_a\) gives
detailed balance with stationary weight proportional to
\(\psi_0(x)^2\), proving the last assertion.
\end{proof}

\begin{remark}[Where the construction can fail]
\label{rem:newS7V65-fail}
For compact spatial gauge transformations, group averaging gives an
orthogonal projector.  The gravitational Hamiltonian constraint is not a
compact group average, and a formal delta function of the constraint is
not automatically a bounded orthogonal projector.  Moreover, BRST
projections can use an indefinite auxiliary inner product.  RTGC therefore
requires an actual physical Hilbert space and projector or an equivalent
reduced-coordinate construction before
\eqref{eq:newS7V65-sandwich} can be invoked.
\end{remark}

\subsection{What conditional detail transport preserves}
\label{sec:newS7V65-conditional}

\begin{theorem}[Positive seed transport through an exact conditional
extension]
\label{thm:newS7V65-extension}
Let \(\mu(dx)\) be a probability, let \(K(dy\mid x)\) be a normalized
conditional kernel, and put
\(\widetilde\mu(dx,dy)=\mu(dx)K(dy\mid x)\).  Define the isometry
\[
 (If)(x,y)=f(x)
\]
from \(L^2(\mu)\) to \(L^2(\widetilde\mu)\), and let
\(E=I^*\) be conditional expectation over \(y\).  If \(T\) is a
self-adjoint positive Markov contraction on \(L^2(\mu)\), then

\begin{equation}
 \widetilde T=ITE
 \label{eq:newS7V65-extension}
\end{equation}

is a self-adjoint positive Markov contraction on
\(L^2(\widetilde\mu)\).  Its nonzero spectrum is the spectrum of \(T\), and
it intertwines old observables:
\(\widetilde T I=IT\).
\end{theorem}

\begin{proof}
\(E=I^*\) and \(EI=1\).  Therefore
\[
 \langle F,\widetilde TF\rangle_{\widetilde\mu}
 =\langle EF,T\,EF\rangle_\mu\geq0,
\]
and \(\widetilde T^*=IT^*E=\widetilde T\).  Conditional expectation and
inclusion are contractions, so is \(\widetilde T\); normalization gives
\(\widetilde T1=1\).  The intertwining identity follows from \(EI=1\).
The range orthogonal to \(I L^2(\mu)\) is killed by \(E\), while the
restriction to that range is unitarily equivalent to \(T\), proving the
spectrum statement.
\end{proof}

\begin{corollary}[Conditional transport preserves but does not create the
seed]
\label{cor:newS7V65-seed}
An exactly normalized conditional detail law preserves reflection-positive
one-step dynamics if the old one-step operator is already positive.  It
does not construct the initial constrained gravitational transfer
operator.
\end{corollary}

\begin{proof}
Preservation is \Cref{thm:newS7V65-extension} followed by
\Cref{thm:newS7V65-one-step}.  If \(T\) has a negative eigenvector \(f\),
then
\[
 \langle If,\widetilde TIf\rangle
 =\langle f,Tf\rangle<0,
\]
so a normalized conditional extension cannot repair a nonpositive seed.
\end{proof}

\subsection{Audit against the attached gravity sources}
\label{sec:newS7V65-source}

The Grand Tensor source with SHA--256 recorded in V59 contains an exact
normalized conditional refinement and a reversible one-step extension.  It
also contains the correct finite reflection-Gram criterion and explicitly
warns that pointwise positive weights need not give reflection positivity.
The physical-source V35 note, also hashed in V59, states that a positive
partition-of-unity density does not supply a reflection-positive control
section.

\begin{theorem}[First absent RTGC interface in the supplied sources]
\label{thm:newS7V65-missing}
The audited gravitational sources supply abstract preservation and finite
testing machinery, but do not supply the seed data required by
\Cref{thm:newS7V65-sandwich}: a nonlinear gravitational constraint surface
or bounded physical projector \(P_a\), a positive reduced kinetic
semigroup \(G_a\), a spatially coercive physical potential \(V_a\), and a
proof that their composition has a nonzero positive top state compatible
with the Standard Model determinant.  Hence RTGC row 1 is not populated by
the conditional-kernel theorem.
\end{theorem}

\begin{proof}
The conditional refinement starts with an old reversible one-step kernel;
\Cref{cor:newS7V65-seed} shows why this is a preservation hypothesis rather
than a seed construction.  The reflection-Gram theorem is an equivalence
and diagnostic, not a positivity proof for a proposed gravitational
matrix.  The physical-source firewall explicitly leaves reflection
positivity separate from positive density.  None of the audited passages
constructs the four displayed seed objects.
\end{proof}

\subsection{The nonlinear connected target after a seed is supplied}
\label{sec:newS7V65-connected}

Suppose a finite-cutoff reduced seed has a quadratic reference
\(Q_a\), nonlinear local remainder \(\mathcal R_a\), and smeared physical
writers \(F_a,G_a\).  The exact interpolation
\[
 d\mu_{a,t}=Z_{a,t}^{-1}e^{-Q_a-t\mathcal R_a}\,d\rho_a
\]
obeys

\begin{equation}
 \operatorname{Cov}_{a,1}(F_a,G_a)
 -\operatorname{Cov}_{a,0}(F_a,G_a)
 =-\int_0^1
 \operatorname{Cum}_{a,t}(F_a,G_a,\mathcal R_a)\,dt.
 \label{eq:newS7V65-Duhamel}
\end{equation}

\begin{proof}
Differentiate the normalized expectation of \(F_aG_a\), \(F_a\), and
\(G_a\), and combine the five centered terms.  This is the
finite-dimensional covariance Duhamel identity independently proved in the
audited YM V37 note.
\end{proof}

\begin{definition}[Reduced transfer connected card, RTCC]
\label{def:newS7V65-RTCC}
RTCC requires, after the seed of
\Cref{thm:newS7V65-sandwich} is constructed:
\begin{enumerate}
\item a finite-range reduced quadratic Hessian with a uniform gap on detail
modes and an explicitly separated massless TT sector;
\item exponentially summable inverse blocks uniformly in spatial volume;
\item connected bounds for \eqref{eq:newS7V65-Duhamel}, score remainders,
constraint Jacobians, and bad-chart polymers, with vacuum clusters
cancelled against \(Z_{a,t}\) before absolute values are taken; and
\item enough decay and cutoff powers to prove RCNC stress uniform
integrability and restoration of spatially nonlocal matter constraints.
\end{enumerate}
\end{definition}

\begin{assessment}[The present first missing object]
\label{ass:newS7V65-status}
The exact preservation algebra is now closed: a positive one-step seed
survives conditional refinement, and a sandwiched heat kernel with a true
orthogonal constraint projector is positive.  The attached construction
does not furnish that gravitational seed.  The first task is therefore not
a polymer estimate but the finite nonlinear reduced configuration space
and its physical one-step operator.  RTCC becomes relevant only after that
operator exists.
\end{assessment}

\begin{programme}[V66: construct the finite reduced TT seed]
\label{prog:newS7V65-V66}
V66 will build the seed in the only sector currently controlled: a finite
spatial lattice of reduced TT variables with the anisotropic dispersion of
V64.  It will give the explicit heat-kernel factorization, ground-state
normalization, and spatial-energy tail bound, then identify why coupling
the scalar and vector Standard Model constraints is not a bounded
perturbation of that seed.
\end{programme}

\begin{firewall}[Claim boundary after seventh-series V65]
\label{fire:newS7V65-boundary}
V65 performs the compulsory companion audit; proves that reversible
pointwise-positive kernels need not have positive transfer spectrum; proves
the exact one-step reflection-positive criterion, sandwiched heat-kernel
construction, and positive conditional-extension theorem; audits the
gravitational sources against those criteria; and isolates the seed and
connected RTCC interfaces.

V65 does not construct the nonlinear gravitational constraint projector,
the coupled physical transfer kernel, RTCC, RTGC, RCNC, cutoff removal, or
the Standard Model--Einstein continuum.
\end{firewall}

\begin{theorem}[Seventh-series V65 result]
\label{thm:newS7V65-result}
A constrained finite-cutoff gravitational transfer matrix is obtained if a
true physical projector \(P_a\), positive kinetic semigroup \(G_a\), and
positive potential multiplier \(W_a\) make the sandwiched operator
\(P_aW_aG_aW_aP_a\) nonzero.  Such an operator yields link-reflection
positivity and transports exactly through normalized conditional detail
kernels.  Reversibility, pointwise-positive weights, and conditional
normalization do not construct the seed.  The supplied sources leave that
seed open; after it is constructed, the nonlinear no-escape target is the
connected card RTCC rather than an extensive Taylor bound.
\end{theorem}

\begin{proof}
Use \Cref{thm:newS7V65-audit,
prop:newS7V65-reversible-no,
thm:newS7V65-one-step,
thm:newS7V65-sandwich,
thm:newS7V65-extension,
cor:newS7V65-seed,
thm:newS7V65-missing}.
\end{proof}

\paragraph{Parent-version record.}
V65 was formed from
\texttt{Renewal\_SM\_Einstein\_Continuum\_Research\_Notes\_V64.tex}.
The V64 parent had \(25\,498\) logical source lines, \(1\,023\,510\) bytes,
and SHA--256
\[
 \texttt{C4546270A2AA0C09495A03CB96437BD7D568DCB7B82D50C0BF6577E85F5C06F6}.
\]
The next compulsory companion audit is V70.

% =============================================================================
% END APPENDED SEVENTH-SERIES V65 MATERIAL
% =============================================================================

% =============================================================================
% APPENDED SEVENTH-SERIES V66 MATERIAL
% =============================================================================

\section{The exact finite reduced TT seed and the sourced-constraint
obstruction}
\label{sec:newS7V66}
\label{sec:v7v66}

V65 identified the first missing object as a positive constrained
one-step seed.  It can be constructed completely in the finite linear
vacuum sector isolated by V64.  The construction below gives the exact
Mehler kernel, its ground-state normalization, reflection positivity, and
its spatial-energy moment.  It also shows why this seed is not yet stable
under the Standard Model constraints: solving those constraints introduces
an inverse Laplacian whose norm diverges with volume and an unbounded
matter-dependent interaction.

\subsection{The finite lattice TT space}
\label{sec:newS7V66-TT-space}

Let \(\Lambda_{a,L}\) be a periodic three-dimensional spatial lattice with
the zero Fourier mode removed.  Let \(\nabla_i\) and \(\nabla_i^*\) be
opposite difference operators and put
\[
 -\Delta_a=\sum_i\nabla_i^*\nabla_i>0
\]
on the zero-mean subspace.  Define the transverse vector projector

\begin{equation}
 P_{ij}=\delta_{ij}
 -\nabla_i(-\Delta_a)^{-1}\nabla_j^*,
 \label{eq:newS7V66-vector-projector}
\end{equation}

and the symmetric transverse-traceless projector

\begin{equation}
 (\Pi_a^{\rm TT})_{ij,kl}
 ={1\over2}(P_{ik}P_{jl}+P_{il}P_{jk})
 -{1\over2}P_{ij}P_{kl}.
 \label{eq:newS7V66-TT-projector}
\end{equation}

\begin{lemma}[Discrete TT orthogonal projection]
\label{lem:newS7V66-TT-projector}
On every nonzero lattice momentum,
\(\Pi_a^{\rm TT}\) is self-adjoint and idempotent.  Its range consists of
symmetric tensors obeying
\(\nabla_i^*q_{ij}=0\) and \(q_{ii}=0\), and it has two polarization
dimensions per momentum.
\end{lemma}

\begin{proof}
Fourier transformation sends \(P\) to
\[
 P_{ij}(k)=\delta_{ij}
 -{\widehat k_i\overline{\widehat k_j}\over|\widehat k|^2},
\]
the orthogonal projection onto the two-dimensional plane perpendicular to
\(\widehat k\).  On symmetric tensors over that plane,
\eqref{eq:newS7V66-TT-projector} subtracts one half of the two-dimensional
trace.  It is therefore the orthogonal projection onto the
two-dimensional symmetric traceless subspace.  Fourier inversion proves
the assertions.
\end{proof}

Let \(\mathcal H_{a,L}^{\rm TT}=\operatorname{Ran}\Pi_a^{\rm TT}\), a
finite-dimensional real Euclidean space.  On it define

\begin{equation}
 A_{a,L}
 =c_2(a)(-\Delta_a)+c_4(a)(-\Delta_a)^2+m_a^2I,
 \quad
 c_2(a)>0,\quad c_4(a)\geq0,\quad m_a^2\geq0.
 \label{eq:newS7V66-A}
\end{equation}

The removed zero mode and \(c_2(a)>0\) make \(A_{a,L}\) strictly positive.
Write \(\Omega_{a,L}=A_{a,L}^{1/2}\).

\subsection{Mehler seed and ground-state transform}
\label{sec:newS7V66-Mehler}

\begin{theorem}[Exact finite reduced TT transfer kernel]
\label{thm:newS7V66-Mehler}
On \(L^2(\mathcal H_{a,L}^{\rm TT},dq)\), let

\begin{equation}
 \mathsf H_{a,L}
 ={1\over2}\left(-\Delta_q+q\cdot A_{a,L}q\right).
 \label{eq:newS7V66-Hamiltonian}
\end{equation}

For every time step \(\epsilon>0\),
\(\mathsf T_{a,L}=e^{-\epsilon\mathsf H_{a,L}}\) is positive,
self-adjoint, trace class, and has the strictly positive kernel

\begin{align}
 K_\epsilon(q,q')={}&
 \det\left({\Omega_{a,L}\over
 2\pi\sinh(\epsilon\Omega_{a,L})}\right)^{1/2}
 \notag\\
 &\times\exp\left\{-{1\over2}\left[
 q\cdot\Omega\coth(\epsilon\Omega)q
 +q'\cdot\Omega\coth(\epsilon\Omega)q'
 -2q\cdot\Omega\operatorname{csch}(\epsilon\Omega)q'
 \right]\right\}.
 \label{eq:newS7V66-Mehler}
\end{align}

Here every matrix function is evaluated at
\(\Omega=\Omega_{a,L}\).  The normalized operator

\begin{equation}
 \widehat{\mathsf T}_{a,L}
 =e^{\epsilon E_{0,a,L}}\mathsf T_{a,L},
 \qquad
 E_{0,a,L}={1\over2}\operatorname{Tr}\Omega_{a,L},
 \label{eq:newS7V66-normalized-T}
\end{equation}

is a positive contraction with a simple top eigenvalue one.
\end{theorem}

\begin{proof}
Diagonalize the positive matrix \(A_{a,L}\).  The Hamiltonian becomes a
finite sum of independent one-dimensional harmonic oscillators with
frequencies \(\omega_j>0\).  The one-dimensional heat kernel is the
corresponding factor in \eqref{eq:newS7V66-Mehler}; multiplying the factors
gives the determinant and matrix expression.  Spectral calculus makes
\(e^{-\epsilon\mathsf H}\) positive and self-adjoint, while the oscillator
spectrum makes it trace class.  Its lowest energy is
\((1/2)\sum_j\omega_j\), is simple, and all other eigenvalues become
\(e^{-\epsilon\sum_jn_j\omega_j}<1\) after the normalization
\eqref{eq:newS7V66-normalized-T}.
\end{proof}

The normalized ground state is

\begin{equation}
 \psi_{0,a,L}(q)
 =\det\left({\Omega_{a,L}\over\pi}\right)^{1/4}
 \exp\left(-{1\over2}q\cdot\Omega_{a,L}q\right).
 \label{eq:newS7V66-ground}
\end{equation}

\begin{corollary}[Ornstein--Uhlenbeck physical Markov kernel]
\label{cor:newS7V66-OU}
The ground-state transform of
\(\widehat{\mathsf T}_{a,L}\) is the Gaussian transition

\begin{equation}
 q'\mid q\sim
 \mathcal N\left(
 e^{-\epsilon\Omega}q,\,
 {1\over2}\Omega^{-1}(I-e^{-2\epsilon\Omega})
 \right),
 \label{eq:newS7V66-OU}
\end{equation}

reversible with respect to

\begin{equation}
 d\pi_{a,L}(q)
 =|\psi_{0,a,L}(q)|^2\,dq
 =\det\left({\Omega\over\pi}\right)^{1/2}
 e^{-q\cdot\Omega q}\,dq.
 \label{eq:newS7V66-invariant}
\end{equation}

Its stationary path measure is time-reflection positive.
\end{corollary}

\begin{proof}
Apply the ground-state transform in
\Cref{thm:newS7V65-sandwich} to the Mehler kernel and complete the square
in \(q'\).  This gives \eqref{eq:newS7V66-OU}.  The covariance identity
\[
 {1\over2}\Omega^{-1}
 =e^{-\epsilon\Omega}
 \left({1\over2}\Omega^{-1}\right)e^{-\epsilon\Omega}
 +{1\over2}\Omega^{-1}(I-e^{-2\epsilon\Omega})
\]
proves stationarity; symmetry of the Mehler kernel proves detailed balance.
Reflection positivity follows from
\Cref{thm:newS7V65-one-step}.
\end{proof}

\begin{theorem}[The free finite TT seed populates RTGC row 1]
\label{thm:newS7V66-seed}
With the global zero modes excluded as declared, the finite reduced TT
system has an explicit nonzero positive one-step operator, physical
probability, and positive transfer Hamiltonian.  Thus RTGC row 1 is
populated for the linear vacuum TT sector.
\end{theorem}

\begin{proof}
Use \Cref{lem:newS7V66-TT-projector,
thm:newS7V66-Mehler,cor:newS7V66-OU}.  The reduction itself supplies the
physical projector; no formal gravitational delta constraint is used.
\end{proof}

\subsection{Exact finite-cutoff spatial-energy entropy}
\label{sec:newS7V66-energy}

Define
\[
 \mathcal E_{a,L}^{\rm sp}(q)
 ={1\over2}q\cdot A_{a,L}q.
\]
Let \(\omega_{\max}=\|\Omega_{a,L}\|\) and let
\(\omega_1,\ldots,\omega_D\) be its eigenvalues.

\begin{lemma}[Spatial-energy moment and tail]
\label{lem:newS7V66-energy}
For \(0<\theta<2/\omega_{\max}\),

\begin{equation}
 \int e^{\theta\mathcal E_{a,L}^{\rm sp}}\,d\pi_{a,L}
 =\prod_{j=1}^D
 \left(1-{\theta\omega_j\over2}\right)^{-1/2}.
 \label{eq:newS7V66-energy-mgf}
\end{equation}

Consequently,

\begin{equation}
 \pi_{a,L}\{\mathcal E_{a,L}^{\rm sp}\geq R\}
 \leq e^{-\theta R}
 \prod_{j=1}^D
 \left(1-{\theta\omega_j\over2}\right)^{-1/2}.
 \label{eq:newS7V66-energy-tail}
\end{equation}
\end{lemma}

\begin{proof}
In an eigenbasis, the invariant density is the product of
\((\omega_j/\pi)^{1/2}e^{-\omega_jq_j^2}dq_j\), while the spatial energy is
\((1/2)\sum_j\omega_j^2q_j^2\).  Each Gaussian integral equals
\((1-\theta\omega_j/2)^{-1/2}\).  Their product proves
\eqref{eq:newS7V66-energy-mgf}; exponential Markov inequality proves
\eqref{eq:newS7V66-energy-tail}.
\end{proof}

\begin{assessment}[The entropy remains extensive]
\label{ass:newS7V66-entropy}
The seed is exactly normalized, but the logarithm of
\eqref{eq:newS7V66-energy-mgf} is a sum over all TT modes and grows with
spatial volume.  This is the TT instance of V60's extensive-mode
obstruction.  Uniform continuum control must use local cylinder fields,
energy density, and connected estimates; the finite formula cannot be
promoted to a uniform raw-energy moment.
\end{assessment}

\subsection{Fixed-momentum regulator removal}
\label{sec:newS7V66-removal}

\begin{proposition}[Free fixed-momentum restoration]
\label{prop:newS7V66-restoration}
Suppose the lattice Laplacian converges to \(|k|^2\) at each fixed
nonzero physical momentum, \(c_2(a)\to c^2>0\),
\(m_a^2\to0\), and \(c_4(a)|k|^4\to0\).  Then the oscillator frequency,
Mehler covariance, and all finite collections of fixed-momentum TT
cylinder laws converge to those of the relativistic massless reduced
graviton with frequency \(c|k|\).
\end{proposition}

\begin{proof}
Functional calculus in the finite set of selected Fourier modes gives
\[
 \omega_a(k)^2
 =c_2(a)\widehat k_a^2+c_4(a)\widehat k_a^4+m_a^2
 \longrightarrow c^2|k|^2.
\]
The covariance and transition matrices in
\eqref{eq:newS7V66-OU} are continuous functions of every positive selected
frequency.  Gaussian characteristic functions therefore converge.
\end{proof}

This is only fixed-cylinder restoration.  It supplies no volume-uniform
tightness and no nonlinear restoration of spacetime diffeomorphism
invariance.

\subsection{Why the Standard Model constraint is not a bounded
perturbation}
\label{sec:newS7V66-SM}

Let \(k_{\min}(L)\) denote the smallest nonzero spatial momentum.  On a
torus of physical side length \(L\),
\[
 \|(-\Delta_a)^{-1}\|={1\over\widehat k_{\min}(L)^2}\asymp L^2.
\]

\begin{theorem}[Sourced-constraint volume obstruction]
\label{thm:newS7V66-source-obstruction}
The linear solution of the Hamiltonian and momentum constraints with
Standard Model sources is not a volume-uniform bounded perturbation of the
free TT seed.  More precisely:
\begin{enumerate}
\item the solution operator contains
\((-\Delta_a)^{-1}T_{00}\) and
\((-\Delta_a)^{-1}T_{0i}\), whose linear operator norm grows like
\(L^2\);
\item the energy and momentum densities are unbounded composite functions
of the gauge, Higgs, and fermion variables; and
\item substituting the solution into the reduced action produces
instantaneous nonlocal interactions and a constraint Jacobian not present
in the Mehler measure.
\end{enumerate}
Therefore no cutoff- and volume-uniform bound
\(\|V_{\rm source}\|_\infty<C\) can justify the multiplier hypothesis of
\Cref{thm:newS7V65-sandwich}.
\end{theorem}

\begin{proof}
The inverse-Laplacian formulas follow from
\Cref{prop:newS7V64-sources}; the displayed norm is the reciprocal of the
smallest nonzero eigenvalue.  The Standard Model stress tensor contains
quadratic gauge and covariant-derivative terms and a quartic Higgs term, so
it is unbounded on its finite-dimensional field space.  Eliminating a
linearly constrained variable inserts its Green operator between two
source factors and contributes the determinant of the constraint map.
These three facts exclude a uniform sup-norm perturbation estimate.
\end{proof}

\begin{corollary}[The free seed cannot simply be reweighted]
\label{cor:newS7V66-no-reweight}
Multiplying the free TT Mehler kernel by the exponential of the
matter-sourced reduced interaction does not prove that the resulting
operator is finite, nonzero, positive semidefinite, or volume-uniform.
Each property requires a separate stability or factorization estimate.
\end{corollary}

\begin{proof}
Unbounded multiplication falls outside
\Cref{thm:newS7V65-sandwich}.  Pointwise positivity would still not prove
operator positivity by \Cref{prop:newS7V65-reversible-no}.  The
inverse-Laplacian growth prevents a volume-uniform conclusion from a
fixed-volume calculation.
\end{proof}

\begin{definition}[Sourced constraint stability card, SCSC]
\label{def:newS7V66-SCSC}
The finite TT seed extends to the linearized Standard Model only if the
coupled cutoff construction proves:
\begin{enumerate}
\item zero-total-source conditions or a separate global-mode law making
the inverse Laplacian well-defined;
\item stability of the induced nonlocal source quadratic form relative to
the positive gauge--Higgs--fermion Hamiltonian, uniformly in volume;
\item a positive transfer factorization of the combined constraint kernel,
not merely a positive pointwise density;
\item uniform control of the constraint determinant, gauge Jacobian, and
BRST or reduced measure;
\item connected spatial decay strong enough to replace the \(L^2\)
operator norm growth by local neutral-cluster estimates; and
\item compatibility with the V57--V58 determinant stress and RCNC.
\end{enumerate}
\end{definition}

\begin{assessment}[Exact position after the seed construction]
\label{ass:newS7V66-status}
The linear vacuum seed is no longer hypothetical: its finite physical
space, kernel, top state, Markov law, and reflection positivity are
explicit.  Its raw energy entropy is extensive but known exactly, and its
fixed-momentum regulator limit is elementary.  The first coupling
obstruction is SCSC, especially the volume-growing inverse Laplacian and
the sign/factorization of the induced Standard Model source interaction.
\end{assessment}

\begin{programme}[V67: neutral clusters for the sourced constraint]
\label{prog:newS7V66-V67}
V67 will decompose a zero-mean lattice source into neutral spatial blocks
and prove the sharp inverse-Laplacian quadratic identity.  It will test
whether local charge neutrality plus exponential connected decay can
replace the \(L^2\) norm \(O(L^2)\) by a volume-uniform energy-density
bound.  If Standard Model energy density has no signed neutrality, the note
will record that failure and move upstream to the global Hamiltonian
constraint rather than assume screening.
\end{programme}

\begin{firewall}[Claim boundary after seventh-series V66]
\label{fire:newS7V66-boundary}
V66 constructs the exact finite reduced TT Mehler seed, ground-state
Markov kernel, reflection-positive path law, spatial-energy moment and
tail, and fixed-momentum regulator limit.  It proves that the
matter-sourced ADM constraints are not a volume-uniform bounded
perturbation and formulates SCSC.

V66 does not solve SCSC, handle global modes, prove a coupled Standard
Model transfer factorization, establish volume-uniform connected bounds,
prove RTGC or RCNC, remove the interacting cutoff, or construct the
Standard Model--Einstein continuum.
\end{firewall}

\begin{theorem}[Seventh-series V66 result]
\label{thm:newS7V66-result}
The free nonzero-mode reduced TT sector has an explicit positive
finite-cutoff transfer seed with two physical polarizations, exact
ground-state normalization, and reflection-positive Ornstein--Uhlenbeck
path measure.  Its higher-spatial-derivative regulator converges on every
fixed momentum when its coefficient is removed.  Coupling the Standard
Model is not a bounded reweighting: the sourced constraints introduce a
volume-growing inverse Laplacian, unbounded composite sources, nonlocal
interactions, and a new determinant.  Extending the seed therefore requires
SCSC rather than formal conditional normalization.
\end{theorem}

\begin{proof}
Use \Cref{lem:newS7V66-TT-projector,
thm:newS7V66-Mehler,
cor:newS7V66-OU,
thm:newS7V66-seed,
lem:newS7V66-energy,
prop:newS7V66-restoration,
thm:newS7V66-source-obstruction,
cor:newS7V66-no-reweight}.
\end{proof}

\paragraph{Parent-version record.}
V66 was formed from
\texttt{Renewal\_SM\_Einstein\_Continuum\_Research\_Notes\_V65.tex}.
The V65 parent had \(25\,899\) logical source lines, \(1\,039\,617\) bytes,
and SHA--256
\[
 \texttt{98CCDE56C34D6A902BE3CEACE2AF93DE63B3A048DF8C5A58EA8C90DBAEDD6208}.
\]
The next compulsory companion audit is V70.

% =============================================================================
% END APPENDED SEVENTH-SERIES V66 MATERIAL
% =============================================================================

% =============================================================================
% APPENDED SEVENTH-SERIES V67 MATERIAL
% =============================================================================

\section{Neutral constraint sources, connected Coulomb bounds, and the
global Hamiltonian mode}
\label{sec:newS7V67}
\label{sec:v7v67}

V66 estimated the sourced constraints by the operator norm of the inverse
spatial Laplacian and obtained an \(O(L^2)\) volume debit.  That estimate is
sharp for the lowest Fourier mode but unnecessarily pessimistic for locally
neutral or mixing sources.  This note proves two replacements: a
deterministic divergence representation and a probabilistic connected
covariance bound.  It then tests the Standard Model energy source against
their hypotheses.  The nonzero modes can be controlled conditionally; the
positive mean energy cannot.  It belongs to the global Hamiltonian
constraint and forces the background geometry to move.

\subsection{The inverse-Laplacian quadratic identity}
\label{sec:newS7V67-identity}

Let \(\Lambda\) be a finite periodic spatial lattice, let
\(\ell^2_0(\Lambda)\) be the zero-mean scalar subspace, and set
\[
 L=-\Delta=\nabla^*\nabla>0
\]
there.  For \(\rho\in\ell^2_0(\Lambda)\), define

\begin{equation}
 \mathcal C(\rho)
 =\langle\rho,L^{-1}\rho\rangle.
 \label{eq:newS7V67-Coulomb}
\end{equation}

\begin{lemma}[Longitudinal-flow identity]
\label{lem:newS7V67-flow}
If \(\rho=\nabla^*j\) for an edge field \(j\), then

\begin{equation}
 \mathcal C(\rho)
 =\|\nabla L^{-1}\nabla^*j\|_2^2
 =\|\mathsf P_{\rm long}j\|_2^2
 \leq\|j\|_2^2,
 \label{eq:newS7V67-flow}
\end{equation}

where
\(\mathsf P_{\rm long}=\nabla L^{-1}\nabla^*\) is the orthogonal
projection onto gradient edge fields.
\end{lemma}

\begin{proof}
\(\mathsf P_{\rm long}\) is self-adjoint and
\[
 \mathsf P_{\rm long}^2
 =\nabla L^{-1}(\nabla^*\nabla)L^{-1}\nabla^*
 =\mathsf P_{\rm long}.
\]
Moreover,
\[
 \langle\nabla^*j,L^{-1}\nabla^*j\rangle
 =\langle j,\nabla L^{-1}\nabla^*j\rangle
 =\langle j,\mathsf P_{\rm long}j\rangle
 =\|\mathsf P_{\rm long}j\|_2^2.
\]
An orthogonal projection is a contraction.
\end{proof}

This removes the \(L^2\) loss whenever a source has a local flow
representation with volume-uniform norm.

\subsection{Block neutrality}
\label{sec:newS7V67-block}

\begin{theorem}[Deterministic block-neutral Coulomb bound]
\label{thm:newS7V67-block}
Partition the periodic lattice into disjoint connected blocks \(B\) of
diameter at most \(r\).  Suppose
\[
 \sum_{x\in B}\rho(x)=0
 \quad\hbox{for every }B.
 \label{eq:newS7V67-block-neutral}
\]
Then there is an edge field \(j\), supported within the blocks, such that
\(\rho=\nabla^*j\) and

\begin{equation}
 \mathcal C(\rho)
 \leq\|j\|_2^2
 \leq C_d r^2\|\rho\|_2^2,
 \label{eq:newS7V67-block-bound}
\end{equation}

where \(C_d\) depends only on the lattice dimension and block shape
regularity, not on the number of blocks or the torus side length.
\end{theorem}

\begin{proof}
On each block solve the zero-mean Neumann problem
\[
 L_Bu_B=\rho|_B,\qquad \sum_{x\in B}u_B(x)=0,
\]
and set \(j_B=\nabla u_B\) on internal block edges.  The compatibility
condition is \eqref{eq:newS7V67-block-neutral}, and
\(\nabla^*j_B=\rho|_B\).  The first nonzero Neumann eigenvalue obeys
\(\lambda_1(L_B)\geq c_dr^{-2}\).  Hence
\[
 \|j_B\|_2^2
 =\langle u_B,\rho\rangle
 \leq\|u_B\|_2\|\rho\|_{2,B}
 \leq\lambda_1(L_B)^{-1/2}
      \|j_B\|_2\|\rho\|_{2,B},
\]
so
\(\|j_B\|_2\leq c_d^{-1/2}r\|\rho\|_{2,B}\).
Sum the squared bounds over disjoint blocks and apply
\Cref{lem:newS7V67-flow}.
\end{proof}

\begin{corollary}[Volume-uniform density bound]
\label{cor:newS7V67-density}
If the block diameter is fixed in physical units and
\(|\Lambda|^{-1}\mathbb E\|\rho\|_2^2\leq C\), then
\[
 {1\over|\Lambda|}\mathbb E\mathcal C(\rho)
 \leq C_d r^2C
\]
uniformly in the spatial volume.
\end{corollary}

\begin{proof}
Take expectations in \eqref{eq:newS7V67-block-bound} and divide by
\(|\Lambda|\).
\end{proof}

\subsection{Connected decay without samplewise block neutrality}
\label{sec:newS7V67-connected}

Let \(G_L(x,y)\) be the zero-mean Green kernel of \(L^{-1}\).
On three-dimensional tori it obeys a volume-uniform local bound of the form
\[
 |G_L(x,y)|\leq C\left(1+{1\over d(x,y)}\right)
\]
away from the removed constant mode, with the diagonal interpreted by the
lattice cutoff.

\begin{theorem}[Connected Coulomb energy density]
\label{thm:newS7V67-connected}
Let \(\rho\) be a stationary random zero-mean lattice field with
\(\mathbb E\rho(x)=0\) and connected covariance
\[
 |\mathbb E[\rho(x)\rho(y)]|
 \leq C_0e^{-c_0d(x,y)}
 \label{eq:newS7V67-cov-decay}
\]
uniformly in the torus volume.  Then

\begin{equation}
 {1\over|\Lambda|}
 \mathbb E\mathcal C(\rho)
 \leq
 C_0\sup_x\sum_y|G_L(x,y)|e^{-c_0d(x,y)}
 \leq C_1,
 \label{eq:newS7V67-connected-bound}
\end{equation}

where \(C_1\) is independent of the volume.
\end{theorem}

\begin{proof}
Expand the normalized quadratic form:
\[
 {1\over|\Lambda|}\mathbb E\mathcal C(\rho)
 ={1\over|\Lambda|}\sum_{x,y}
 G_L(x,y)\mathbb E[\rho(x)\rho(y)].
\]
Take absolute values and use stationarity or a supremum over \(x\).
The number of sites at distance \(n\) is \(O((1+n)^2)\), the
three-dimensional Green kernel is \(O(1+n^{-1})\), and the exponential
factor makes the resulting series summable uniformly in the torus size.
\end{proof}

\begin{remark}[Why this is the YM V37 type of estimate]
\label{rem:newS7V67-YM}
The volume disappears only after retaining the connected covariance
between the two marked sources.  Bounding the full nonlinear remainder or
the absolute number of bad sites before forming that connected quantity
would restore a volume factor.  This is the same logical order isolated by
the audited YM V37 covariance Duhamel argument.
\end{remark}

\begin{proposition}[A mean source is not screened by connected decay]
\label{prop:newS7V67-mean}
If \(\mathbb E\rho(x)=m\ne0\), then the constant component
\(m\mathbf1\) is outside the domain of the zero-mean inverse \(L^{-1}\).
Subtracting it by hand replaces \(\rho\) with
\(\rho-m\) in the nonzero-mode estimate but leaves a separate global
constraint for \(m\).
\end{proposition}

\begin{proof}
\(L\mathbf1=0\), so \(L^{-1}\) is defined only on the orthogonal complement
of \(\mathbf1\).  The orthogonal decomposition
\(\rho=(\rho-m)+m\mathbf1\) proves the statement.
\end{proof}

\subsection{The integrated Hamiltonian constraint}
\label{sec:newS7V67-global}

On a closed spatial slice \(\Sigma\), the ADM Hamiltonian constraint has
the form

\begin{equation}
 R^{(3)}+K^2-K_{ij}K^{ij}-2\Lambda
 =16\pi G\,T_{nn}.
 \label{eq:newS7V67-Hamiltonian}
\end{equation}

\begin{theorem}[No nontrivial positive-energy matter on the flat
time-symmetric zero-mode background]
\label{thm:newS7V67-flat-no-matter}
Let \(\Sigma=\mathbb T^3\) carry a flat spatial metric, let
\(K_{ij}=0\), and let \(\Lambda=0\).  If the matter source obeys
\(T_{nn}\geq0\), then the Hamiltonian constraint implies
\(T_{nn}=0\) almost everywhere.  Thus the V66 flat vacuum seed cannot be
coupled to a nontrivial positive-energy Standard Model state while its
global geometric mode is kept fixed.
\end{theorem}

\begin{proof}
Under the stated assumptions the left side of
\eqref{eq:newS7V67-Hamiltonian} vanishes pointwise, so
\(T_{nn}=0\) pointwise.  More generally, integrating the constraint shows
that the positive mean matter energy must be balanced by spatial
curvature, extrinsic curvature, or the cosmological term.
\end{proof}

\begin{corollary}[The energy mean belongs to the background equation]
\label{cor:newS7V67-background}
Writing
\[
 T_{nn}=\overline T_{nn}+
 (T_{nn}-\overline T_{nn})
\]
places only the fluctuation in the nonzero-mode inverse-Laplacian problem.
The mean must be coupled to a volume, mean-curvature, curvature, or
cosmological variable satisfying the integrated Hamiltonian constraint.
\end{corollary}

\begin{proof}
Apply \Cref{prop:newS7V67-mean} and integrate
\eqref{eq:newS7V67-Hamiltonian}.
\end{proof}

\subsection{Test against Standard Model sources}
\label{sec:newS7V67-SM}

\begin{proposition}[Local neutrality is not an identity for Standard Model
energy]
\label{prop:newS7V67-SM-neutrality}
Neither diffeomorphism conservation nor Standard Model gauge neutrality
implies
\eqref{eq:newS7V67-block-neutral} for the energy density \(T_{nn}\).
For a nonzero classical gauge or Higgs wave packet,
\(T_{nn}\) can be nonnegative with strictly positive integral on one
block.
\end{proposition}

\begin{proof}
Stress conservation is a spacetime equation
\(\nabla_\mu T^{\mu\nu}=0\); it does not state that the energy on each
spatial block vanishes.  Internal gauge charges may cancel while the gauge
and Higgs kinetic energies remain positive.  A smooth wave packet supported
inside one block gives the announced example.
\end{proof}

\begin{proposition}[Conditional usefulness of connected decay]
\label{prop:newS7V67-SM-connected}
If a constructed cutoff Standard Model state proves exponential connected
decay for the centered, renormalized source
\[
 [T_{nn}]_a-\mathbb E[T_{nn}]_a
\]
with a bound uniform in volume and cutoff after smearing, then
\Cref{thm:newS7V67-connected} controls its expected nonzero-mode
constraint energy per unit volume.  This conclusion does not control the
global mean, higher cumulants, or the sign of the induced gravitational
interaction.
\end{proposition}

\begin{proof}
The centered source meets the mean-zero hypothesis, so the theorem applies.
The three excluded quantities do not enter its two-point calculation.
\end{proof}

At a massless or critical Standard Model point, exponential clustering may
itself fail; then the decay rate must be replaced by a power for which
\(\sum_y|G(x,y)||\operatorname{Cov}(x,y)|\) remains finite.  No clustering
estimate of either kind is presently supplied for the coupled measure.

\subsection{Separated local and global constraint cards}
\label{sec:newS7V67-cards}

\begin{definition}[Connected nonzero-mode constraint card, CNMCC]
\label{def:newS7V67-CNMCC}
CNMCC consists of:
\begin{enumerate}
\item a centered renormalized source composite with a cutoff-consistent
zero-mode subtraction;
\item either a block-flow representation satisfying
\eqref{eq:newS7V67-block-bound} or connected decay satisfying
\eqref{eq:newS7V67-connected-bound};
\item corresponding higher connected cumulant bounds for the exponential
of the induced source interaction;
\item stability of that interaction relative to the Standard Model
Hamiltonian and positivity of the resulting physical transfer operator;
and
\item compatibility with RCNC stress and Ward identities.
\end{enumerate}
\end{definition}

\begin{definition}[Global Hamiltonian-mode card, GHMC]
\label{def:newS7V67-GHMC}
GHMC consists of:
\begin{enumerate}
\item a finite-cutoff global volume or scale variable, its conjugate
mean-curvature variable, and the global lapse constraint;
\item a normalized law or physical projector imposing the integrated
version of \eqref{eq:newS7V67-Hamiltonian} with the mean Standard Model
energy and cosmological term;
\item lower-bounded physical evolution or a specified contour for this
global mode;
\item compatibility with the local TT transfer kernel and with refinement;
and
\item tightness and restoration of the intended Einstein low-energy
normalization.
\end{enumerate}
\end{definition}

\begin{theorem}[SCSC splits into a local connected problem and a global
background problem]
\label{thm:newS7V67-split}
The inverse-Laplacian part of SCSC can be volume-uniformly controlled by
CNMCC, but CNMCC cannot impose the mean Hamiltonian constraint.  Coupling a
nontrivial positive-energy Standard Model state to the V66 seed therefore
requires both CNMCC and GHMC.
\end{theorem}

\begin{proof}
\Cref{thm:newS7V67-block,thm:newS7V67-connected} give the two
volume-uniform nonzero-mode mechanisms.  \Cref{prop:newS7V67-mean} shows
that both omit the constant mode, and
\Cref{thm:newS7V67-flat-no-matter} shows that nonzero mean matter energy
cannot be added while retaining the flat time-symmetric background.
\end{proof}

\begin{assessment}[The correct upstream move]
\label{ass:newS7V67-status}
The \(O(L^2)\) inverse-Laplacian debit is not an unavoidable local
obstruction: exact block neutrality or exponential connected decay removes
it at energy-density level.  Standard Model energy has no block-neutrality
identity, so its plausible route is a renormalized connected estimate.
Before that estimate can couple to gravity, however, the mean energy must
be absorbed by a genuine solution of the global Hamiltonian constraint.
The next step must therefore move upstream from the flat TT seed to a
background with a dynamical global scale or mean curvature.
\end{assessment}

\begin{programme}[V68: homogeneous global constraint seed]
\label{prog:newS7V67-V68}
V68 will introduce the homogeneous scale factor and its canonical momentum
on a compact constant-curvature slice, derive the finite-dimensional
Hamiltonian constraint with cosmological and mean Standard Model energy
terms, and determine whether deparametrization by a monotone clock yields a
positive physical Hamiltonian.  It will keep the conformal-sign and contour
choice explicit and will not attach the TT Mehler seed until GHMC rows
1--3 are proved.
\end{programme}

\begin{firewall}[Claim boundary after seventh-series V67]
\label{fire:newS7V67-boundary}
V67 proves the exact longitudinal-flow identity, a volume-uniform
block-neutral inverse-Laplacian bound, and a volume-uniform expected
Coulomb-energy-density bound under exponential connected decay.  It proves
that Standard Model energy is not block neutral and that nontrivial
positive mean energy is incompatible with the fixed flat
time-symmetric vacuum background.  It splits SCSC into CNMCC and GHMC.

V67 does not prove Standard Model clustering, higher connected source
bounds, transfer positivity after sourced constraints, GHMC, RTGC, RCNC,
cutoff removal, or the Standard Model--Einstein continuum.
\end{firewall}

\begin{theorem}[Seventh-series V67 result]
\label{thm:newS7V67-result}
The nonzero-mode gravitational constraint need not pay the worst-case
\(O(L^2)\) inverse-Laplacian norm: local block neutrality gives the
deterministic bound \eqref{eq:newS7V67-block-bound}, and connected source
decay gives the expected density bound
\eqref{eq:newS7V67-connected-bound}.  Standard Model energy satisfies
neither samplewise block neutrality nor zero mean.  Its centered
fluctuation therefore requires CNMCC, while its positive mean requires a
new global geometric seed satisfying GHMC.  The flat TT seed alone cannot
host nontrivial positive-energy matter.
\end{theorem}

\begin{proof}
Use \Cref{lem:newS7V67-flow,
thm:newS7V67-block,
thm:newS7V67-connected,
prop:newS7V67-mean,
thm:newS7V67-flat-no-matter,
prop:newS7V67-SM-neutrality,
thm:newS7V67-split}.
\end{proof}

\paragraph{Parent-version record.}
V67 was formed from
\texttt{Renewal\_SM\_Einstein\_Continuum\_Research\_Notes\_V66.tex}.
The V66 parent had \(26\,333\) logical source lines, \(1\,055\,461\) bytes,
and SHA--256
\[
 \texttt{FE6523301385AFCE2E05D7BF8C30D2F918B74D816A9AC4966D649D3550C61957}.
\]
The next compulsory companion audit is V70.

% =============================================================================
% END APPENDED SEVENTH-SERIES V67 MATERIAL
% =============================================================================

% =============================================================================
% APPENDED SEVENTH-SERIES V68 MATERIAL
% =============================================================================

\section{The homogeneous Hamiltonian constraint and intrinsic scale time}
\label{sec:newS7V68}
\label{sec:v7v68}

V67 proved that positive mean Standard Model energy cannot be added to the
flat time-symmetric seed while the global geometric mode is frozen.  This
note restores that mode in the homogeneous sector.  The ADM scale momentum
has the negative gravitational kinetic sign, but that sign can be used to
solve the Hamiltonian constraint for the scale momentum.  On a monotone
branch the scale factor becomes an intrinsic clock and the remaining
Hamiltonian is a positive square root, provided an explicit operator under
the root is positive.  This gives a local solution of the first three GHMC
rows; turning points and reflection positivity of a nonstationary
scale-time kernel remain open.

\subsection{Exact homogeneous constraint}
\label{sec:newS7V68-constraint}

Let \((\Sigma,\gamma^{(k)})\) be a compact homogeneous three-manifold of
reference volume \(V_0\) and scalar curvature \(6k\).  Use

\[
 ds^2=-N(t)^2dt^2+a(t)^2\gamma^{(k)}_{ij}dx^idx^j,
 \qquad a>0.
\]

After the usual boundary term is removed, the gravitational Lagrangian is

\begin{equation}
 L_g
 =-{3V_0a\dot a^2\over8\pi GN}
 +{3kV_0Na\over8\pi G}
 -{\Lambda V_0Na^3\over8\pi G}.
 \label{eq:newS7V68-Lagrangian}
\end{equation}

\begin{lemma}[Scale momentum and constraint]
\label{lem:newS7V68-constraint}
The momentum and lapse constraint are

\begin{align}
 p_a&=-{3V_0a\dot a\over4\pi GN},
 \label{eq:newS7V68-pa}\\
 \mathcal C(a,p_a;z)
 &=-A(a)p_a^2+B(a;z)=0,
 \label{eq:newS7V68-C}\\
 A(a)&={2\pi G\over3V_0a}>0,
 \label{eq:newS7V68-A}\\
 B(a;z)&=
 -{3kV_0a\over8\pi G}
 +{\Lambda V_0a^3\over8\pi G}
 +H_{\rm SM}(a;z).
 \label{eq:newS7V68-B}
\end{align}

Here \(z\) denotes the finite-cutoff Standard Model canonical variables and
\(H_{\rm SM}\) is their homogeneous-slice Hamiltonian.
\end{lemma}

\begin{proof}
Differentiate \eqref{eq:newS7V68-Lagrangian} with respect to \(\dot a\) to
obtain \eqref{eq:newS7V68-pa}.  Solving for \(\dot a\) and forming
\(p_a\dot a-L_g+NH_{\rm SM}\) gives
\[
 H=N\left(
 -{2\pi G\over3V_0a}p_a^2
 -{3kV_0a\over8\pi G}
 +{\Lambda V_0a^3\over8\pi G}
 +H_{\rm SM}\right).
\]
Variation of \(N\) proves the constraint and the definitions.
\end{proof}

\begin{corollary}[Extrinsic curvature balances positive mean energy]
\label{cor:newS7V68-balance}
For \(k=0\), \(\Lambda\geq0\), and \(H_{\rm SM}\geq0\), the constraint has
the two real branches

\begin{equation}
 p_a=\sigma\sqrt{B(a;z)\over A(a)},
 \qquad \sigma\in\{-1,+1\}.
 \label{eq:newS7V68-branches}
\end{equation}

Nonzero mean matter energy is balanced by nonzero scale momentum and hence
nonzero mean extrinsic curvature.  The time-symmetric choice \(p_a=0\)
forces \(B=0\), recovering the obstruction of V67.
\end{corollary}

\begin{proof}
\eqref{eq:newS7V68-C} is equivalent to
\(A(a)p_a^2=B(a;z)\).  Since \(A>0\), the stated sign assumptions give
\eqref{eq:newS7V68-branches}.  Equation \eqref{eq:newS7V68-pa} identifies
nonzero \(p_a\) with nonzero \(\dot a/N\).
\end{proof}

\subsection{Intrinsic scale deparametrization}
\label{sec:newS7V68-scale-time}

\begin{theorem}[Local positive scale-time Hamiltonian]
\label{thm:newS7V68-scale-H}
Let \(\mathcal U\) be a branch of the finite constrained phase space on
which \(B(a;z)>0\) and \(p_a\) has a fixed sign.  Then \(a\) is a local
monotone clock.  Gauge fixing \(a=s\) and solving the constraint gives a
reduced Hamiltonian

\begin{equation}
 H_{\rm phys}^{(a)}(s;z)
 =\sqrt{B(s;z)\over A(s)}
 \label{eq:newS7V68-Hphys}
\end{equation}

after choosing the orientation for which the reduced Hamiltonian is
positive.
\end{theorem}

\begin{proof}
Hamilton's equation gives
\[
 {da\over dt}
 =N{\partial\mathcal C\over\partial p_a}
 =-2NA(a)p_a.
\]
It is nonzero on \(\mathcal U\), so \(a\) is a valid local gauge clock.
The canonical action contains
\[
 \int(p_a\dot a+p_z\dot z-N\mathcal C)\,dt.
\]
Set \(a=s\), insert one branch of
\eqref{eq:newS7V68-branches}, and orient \(s\) so that the term
\(p_a\,ds\) becomes \(-H_{\rm phys}^{(a)}ds\).  This yields
\eqref{eq:newS7V68-Hphys}.
\end{proof}

\begin{corollary}[Finite-cutoff operator square root]
\label{cor:newS7V68-operator}
At a fixed \(s>0\), suppose the renormalized finite-cutoff Standard Model
Hamiltonian is self-adjoint and the operator
\[
 B_s=
 -{3kV_0s\over8\pi G}
 +{\Lambda V_0s^3\over8\pi G}
 +H_{\rm SM}(s)
\]
is bounded below by a positive number.  Then

\begin{equation}
 \widehat H_{\rm phys}^{(a)}(s)
 =A(s)^{-1/2}B_s^{1/2}
 \label{eq:newS7V68-operator-H}
\end{equation}

is a positive self-adjoint operator and
\(e^{-\epsilon\widehat H_{\rm phys}^{(a)}(s)}\) is a positive
self-adjoint contraction.
\end{corollary}

\begin{proof}
\(A(s)\) is a positive scalar.  The spectral theorem gives the unique
positive square root \(B_s^{1/2}\); multiplication by \(A(s)^{-1/2}\)
preserves self-adjointness and positivity.  Functional calculus gives the
last assertion.
\end{proof}

\begin{remark}[Vacuum-energy renormalization is part of \(B_s\)]
\label{rem:newS7V68-vacuum}
A lower shift of \(H_{\rm SM}\) cannot be discarded independently of
\(\Lambda\): it changes the same homogeneous constraint.  The positivity
hypothesis in \Cref{cor:newS7V68-operator} is therefore a renormalized
cosmological-plus-matter condition, not a free normalization convention.
\end{remark}

\subsection{Clock alternatives do not automatically improve positivity}
\label{sec:newS7V68-clocks}

\begin{proposition}[A positive-energy scalar clock gives a restricted
radicand]
\label{prop:newS7V68-scalar-clock}
Add a homogeneous massless scalar clock \(T\) with momentum \(p_T\) and
Hamiltonian \(p_T^2/(2V_0a^3)\).  Solving the total constraint for \(p_T\)
gives

\begin{equation}
 p_T^2
 =2V_0a^3\bigl(A(a)p_a^2-B(a;z)\bigr).
 \label{eq:newS7V68-scalar-radicand}
\end{equation}

Hence scalar-clock evolution is defined only on the region where the
right side is nonnegative.  Increasing positive Standard Model energy
decreases this radicand.
\end{proposition}

\begin{proof}
Add \(p_T^2/(2V_0a^3)\) to
\eqref{eq:newS7V68-C}, set the sum to zero, and rearrange.
\end{proof}

\begin{proposition}[A linear clock retains the conformal lower-bound
problem]
\label{prop:newS7V68-linear-clock}
If an auxiliary clock produces a constraint
\(p_T+\mathcal C=0\), then its reduced Hamiltonian is
\(\mathcal C\).  Because \(\mathcal C\) contains
\(-A(a)p_a^2\), it is unbounded below on the unreduced real scale momentum
unless a further constraint, contour, or lower-bound mechanism is supplied.
\end{proposition}

\begin{proof}
Solving gives \(p_T=-\mathcal C\), so the Hamiltonian generating \(T\)
translations is \(\mathcal C\) up to orientation.  At fixed other
variables, \(-A(a)p_a^2\to-\infty\) as \(|p_a|\to\infty\).
\end{proof}

Thus the scale clock is locally better adapted to the negative
gravitational kinetic sign: it solves for that momentum rather than leaving
it inside the physical Hamiltonian.

\subsection{Nonstationary Euclidean steps}
\label{sec:newS7V68-nonstationary}

For an interval \(s_0<s_1<\cdots<s_n\), a natural scale-time Euclidean
evolution is a product

\begin{equation}
 \mathsf U_{n,0}
 =e^{-\Delta s_nH(s_{n-1/2})}\cdots
  e^{-\Delta s_1H(s_{1/2})}.
 \label{eq:newS7V68-product}
\end{equation}

\begin{lemma}[Products of positive steps need not be positive
self-adjoint]
\label{lem:newS7V68-product}
If \(A\) and \(B\) are positive self-adjoint matrices which do not commute,
then \(AB\) need not be self-adjoint and therefore need not define a
positive one-step reflection form.
\end{lemma}

\begin{proof}
\((AB)^*=BA\), which differs from \(AB\) when \([A,B]\ne0\).  A positive
operator must be self-adjoint.
\end{proof}

\begin{corollary}[Positive instantaneous Hamiltonians do not prove a
scale-time OS state]
\label{cor:newS7V68-no-OS}
The positivity of every
\(\widehat H_{\rm phys}^{(a)}(s)\) proves contractivity of each
instantaneous step but does not by itself prove reflection positivity of
the nonstationary product \eqref{eq:newS7V68-product}.  A reflected
palindromic factorization or a direct reflection-Gram theorem is still
required.
\end{corollary}

\begin{proof}
Apply \Cref{lem:newS7V68-product}; the Hamiltonians at different scale
values need not commute because \(H_{\rm SM}(s)\) changes with \(s\).
\end{proof}

\begin{theorem}[Palindromic finite slab is positive]
\label{thm:newS7V68-palindrome}
Let \(X_1,\ldots,X_n\) be bounded operators and define
\[
 \mathsf K_{\rm slab}
 =X_1^*\cdots X_n^*X_n\cdots X_1.
\]
Then \(\mathsf K_{\rm slab}\succeq0\).  In particular, a scale-time slab
whose negative half is the adjoint reflection of its positive half has a
positive boundary reflection kernel.
\end{theorem}

\begin{proof}
Put \(Y=X_n\cdots X_1\).  Then
\(\mathsf K_{\rm slab}=Y^*Y\), so
\(\langle\psi,\mathsf K_{\rm slab}\psi\rangle=\|Y\psi\|^2\geq0\).
\end{proof}

This gives finite reflected slabs but not a stationary infinite-volume
state or a translation generator in scale time.

\subsection{Homogeneous global-mode card}
\label{sec:newS7V68-GHMC}

\begin{definition}[Refined global Hamiltonian-mode card, GHMC\(^{+}\)]
\label{def:newS7V68-GHMCplus}
The homogeneous scale branch satisfies GHMC\(^{+}\) if:
\begin{enumerate}
\item the finite-cutoff constraint is
\eqref{eq:newS7V68-C}, with the Standard Model vacuum shift and
cosmological term renormalized together;
\item a branch and interval have \(B_s\succeq b_*I>0\), so \(a\) is
monotone and \eqref{eq:newS7V68-operator-H} exists;
\item finite slabs have a declared palindromic or direct positive
reflection factorization;
\item turning points \(B=0\), the \(a\downarrow0\) boundary, and expanding
versus contracting branches have a self-adjoint matching rule;
\item the scale law couples consistently to the TT Mehler seed and the
centered constraint card CNMCC; and
\item refinement, physical-clock independence, Einstein normalization,
and RCNC are proved.
\end{enumerate}
\end{definition}

\begin{theorem}[Local homogeneous progress]
\label{thm:newS7V68-progress}
On a flat or suitably curved branch for which \(B_s\succeq b_*I>0\),
GHMC\(^{+}\) rows 1--2 hold, and row 3 holds for every explicitly
palindromic finite slab.  Rows 4--6 remain independent requirements.
\end{theorem}

\begin{proof}
\Cref{lem:newS7V68-constraint} gives row 1.
\Cref{thm:newS7V68-scale-H,cor:newS7V68-operator} give row 2.
\Cref{thm:newS7V68-palindrome} gives the stated finite-slab part of row 3.
None of those arguments treats a turning point, couples TT and local source
fluctuations, or takes a refinement limit.
\end{proof}

\begin{assessment}[What the global mode resolves]
\label{ass:newS7V68-status}
Allowing nonzero scale momentum removes the contradiction between a flat
spatial slice and positive mean matter energy: the mean extrinsic curvature
balances the homogeneous source.  On a fixed-sign branch, intrinsic scale
time converts the negative scale kinetic term into a positive square-root
Hamiltonian.  This is local in clock space and conditional on
\(B_s>0\).  It does not yet supply a stationary reflection-positive state,
cross a turning point, or couple the nonlinear local constraint fields.
\end{assessment}

\begin{programme}[V69: turning points and branch matching]
\label{prog:newS7V68-V69}
V69 will analyze the zero set \(B_s=0\).  It will determine when the
square-root Hamiltonian is essentially self-adjoint on a common domain,
whether a contracting branch can be joined to an expanding branch by a
unitary or reflection-positive rule, and whether the simplest Airy-type
turning model loses scale-clock evolution.  If no canonical matching
exists, GHMC\(^{+}\) will be restricted to monotone sectors and the
superselection input will be stated explicitly.
\end{programme}

\begin{firewall}[Claim boundary after seventh-series V68]
\label{fire:newS7V68-boundary}
V68 derives the exact homogeneous Einstein--Standard-Model Hamiltonian
constraint; proves local deparametrization by the scale factor and the
positive square-root Hamiltonian on \(B_s>0\); proves that scalar and linear
clock alternatives retain distinct positivity restrictions; and proves a
palindromic finite-slab reflection factorization.

V68 does not prove positivity of \(B_s\) in the renormalized interacting
theory, cross scale turning points, construct a stationary physical state,
couple the TT and centered constraint sectors, prove GHMC\(^{+}\), RTGC, or
RCNC, remove the cutoff, or construct the Standard Model--Einstein
continuum.
\end{firewall}

\begin{theorem}[Seventh-series V68 result]
\label{thm:newS7V68-result}
The global Hamiltonian mode supplies a local mechanism for positive mean
Standard Model energy which the flat time-symmetric TT seed lacked.
On any branch with \(B_s>0\), the scale factor is monotone and the
constraint yields the positive physical Hamiltonian
\(\sqrt{B_s/A(s)}\).  Positive instantaneous steps are insufficient for a
nonstationary Osterwalder--Schrader state, but every reflected palindromic
finite slab has a positive boundary kernel.  Turning points, branch
matching, and coupling to the local constraint fluctuations are the
remaining GHMC\(^{+}\) interfaces.
\end{theorem}

\begin{proof}
Use \Cref{lem:newS7V68-constraint,
cor:newS7V68-balance,
thm:newS7V68-scale-H,
cor:newS7V68-operator,
prop:newS7V68-scalar-clock,
prop:newS7V68-linear-clock,
lem:newS7V68-product,
thm:newS7V68-palindrome,
thm:newS7V68-progress}.
\end{proof}

\paragraph{Parent-version record.}
V68 was formed from
\texttt{Renewal\_SM\_Einstein\_Continuum\_Research\_Notes\_V67.tex}.
The V67 parent had \(26\,764\) logical source lines, \(1\,070\,915\) bytes,
and SHA--256
\[
 \texttt{42AEE27D812B8CC7F5C3C694F519FD1BAC24246008AD1E9CFDCC1CEEBBC5E8B2}.
\]
The next compulsory companion audit is V70.

% =============================================================================
% END APPENDED SEVENTH-SERIES V68 MATERIAL
% =============================================================================

% =============================================================================
% APPENDED SEVENTH-SERIES V69 MATERIAL
% =============================================================================

\section{Scale-clock turning points and the missing branch-matching datum}
\label{sec:newS7V69}
\label{sec:v7v69}

V68 constructed a positive intrinsic-scale Hamiltonian on a branch where
\(B(a;z)>0\) and \(p_a\) has fixed sign.  This note studies the boundary
\(B=0\).  A simple zero is a regular turning point of the full
Hamiltonian-constraint trajectory but a singularity of the scale clock:
the expanding and contracting descriptions meet there, while the
first-order square-root evolution does not determine their relative
matching.  The second-order constraint has an Airy normal form and crosses
the point only after its domain or forbidden-side boundary condition is
specified.

\subsection{Classical normal form}
\label{sec:newS7V69-classical}

Retain the homogeneous constraint
\[
 \mathcal C(a,p_a;z)=-A(a)p_a^2+B(a;z)=0,
 \qquad A(a)>0.
\]
First freeze the remaining canonical variables at their value at a turning
point; their evolution contributes higher-order coupled terms but does not
change the transverse simple-zero calculation.

\begin{lemma}[Simple turning-point expansion]
\label{lem:newS7V69-simple}
Suppose
\[
 B(a_*;z_*)=0,
 \qquad
 b_*:=\partial_aB(a_*;z_*)\ne0.
\]
Then on the allowed side \(b_*(a-a_*)>0\),

\begin{equation}
 p_a
 =\sigma\left({b_*\over A_*}(a-a_*)\right)^{1/2}
 +O(|a-a_*|),
 \qquad
 A_*:=A(a_*),
 \label{eq:newS7V69-p-expansion}
\end{equation}

where \(\sigma=\pm1\).
\end{lemma}

\begin{proof}
Taylor expansion gives
\[
 {B(a;z_*)\over A(a)}
 ={b_*\over A_*}(a-a_*)+O((a-a_*)^2).
\]
Taking either real square root proves
\eqref{eq:newS7V69-p-expansion}.
\end{proof}

\begin{theorem}[The full trajectory reflects while scale time stops]
\label{thm:newS7V69-reflection}
Choose a smooth nonzero lapse \(N\) near the simple turning point.  Along a
full Hamiltonian trajectory,

\begin{align}
 \dot a&=-2NA(a)p_a,
 \label{eq:newS7V69-adot}\\
 \dot p_a&=N A'(a)p_a^2-N\partial_aB(a;z).
 \label{eq:newS7V69-pdot}
\end{align}

At the turning point,

\begin{equation}
 \dot a_*=0,\qquad
 \dot p_{a,*}=-N_*b_*,
 \qquad
 \ddot a_*=2N_*^2A_*b_*.
 \label{eq:newS7V69-turn-data}
\end{equation}

Thus, for \(b_*\ne0\),
\[
 a(t)-a_*
 =N_*^2A_*b_*(t-t_*)^2+o((t-t_*)^2).
\]
The parameter \(t\) crosses the point and \(p_a\) changes sign, but \(a\)
has an extremum and cannot be a coordinate through the crossing.
\end{theorem}

\begin{proof}
Hamilton's equations for \(N\mathcal C\) give
\eqref{eq:newS7V69-adot}--\eqref{eq:newS7V69-pdot}.  At
\(B=p_a=0\), the first two identities in
\eqref{eq:newS7V69-turn-data} follow.  Differentiating
\eqref{eq:newS7V69-adot}, every term containing \(p_a\) or \(\dot a\)
vanishes at the turn, leaving
\[
 \ddot a_*=-2N_*A_*\dot p_{a,*}
 =2N_*^2A_*b_*.
\]
Taylor's theorem gives the trajectory expansion.  Since
\(\dot p_{a,*}\ne0\), \(p_a\) changes sign.
\end{proof}

\begin{corollary}[No single scale-clock branch crosses a simple turn]
\label{cor:newS7V69-no-clock-crossing}
The gauge condition \(a=s\) has vanishing Faddeev--Popov derivative
\[
 \{a,\mathcal C\}=-2A(a)p_a
\]
at a simple turning point.  Hence the scale-clock reduction of V68 is not a
valid gauge chart there.
\end{corollary}

\begin{proof}
The displayed bracket is \(\partial\mathcal C/\partial p_a\), and it
vanishes at \(p_a=0\).  Gauge transversality therefore fails.
\end{proof}

\subsection{Finite-cutoff square roots at a spectral zero}
\label{sec:newS7V69-square-root}

\begin{lemma}[A nonnegative spectral zero is not by itself an operator
singularity]
\label{lem:newS7V69-continuity}
Let \(B(s)\) be a continuous finite-dimensional family of positive
semidefinite Hermitian matrices.  Then \(H(s)=B(s)^{1/2}\) is continuous
and obeys

\begin{equation}
 \|B(s)^{1/2}-B(t)^{1/2}\|
 \leq\|B(s)-B(t)\|^{1/2}.
 \label{eq:newS7V69-holder}
\end{equation}

The nonautonomous equation
\[
 i\partial_sU(s,s_0)=H(s)U(s,s_0),\qquad U(s_0,s_0)=I,
\]
therefore has a unique unitary propagator on every interval on which
\(B(s)\succeq0\).
\end{lemma}

\begin{proof}
The square-root inequality for positive matrices follows from functional
calculus and the operator monotonicity of the square root.  Continuity of
\(H\) follows.  The time-ordered Dyson series converges uniformly on a
compact interval because \(H(s)\) is bounded and continuous; differentiating
\(U^*U\) gives zero, so the propagator is unitary.
\end{proof}

\begin{proposition}[A sign-changing eigenvalue ends self-adjoint
scale-time evolution]
\label{prop:newS7V69-negative}
If an eigenvalue of \(B(s)\) crosses transversely through zero, then on one
side \(B(s)\) has negative spectrum.  There is no self-adjoint operator
\(H(s)\) satisfying \(H(s)^2=B(s)\) on that side.
\end{proposition}

\begin{proof}
The square of a self-adjoint operator is positive semidefinite.  A matrix
with a negative eigenvalue cannot equal such a square.
\end{proof}

Thus a zero which merely touches and leaves \(B\geq0\) permits formal
finite-cutoff square-root propagation, although the classical clock still
loses transversality.  A generic simple sign change ends the positive
square-root sector.

\subsection{The second-order constraint and Airy matching}
\label{sec:newS7V69-Airy}

Choose the symmetric one-dimensional factor ordering

\begin{equation}
 \widehat{\mathcal C}\Psi
 =\partial_a\!\left(A(a)\partial_a\Psi\right)
 +B(a)\Psi.
 \label{eq:newS7V69-WDW}
\end{equation}

\begin{lemma}[Conserved constraint current]
\label{lem:newS7V69-current}
For real \(A,B\), any two solutions \(\Psi_1,\Psi_2\) of
\eqref{eq:newS7V69-WDW} have constant Wronskian current

\begin{equation}
 J(\Psi_1,\Psi_2)
 =iA(a)\left(
 \overline{\Psi_1}\,\partial_a\Psi_2
 -(\partial_a\overline{\Psi_1})\Psi_2
 \right).
 \label{eq:newS7V69-current}
\end{equation}
\end{lemma}

\begin{proof}
Differentiate \eqref{eq:newS7V69-current}.  Substitution of
\(\partial_a(A\partial_a\Psi_j)=-B\Psi_j\) cancels the two real potential
terms.
\end{proof}

\begin{theorem}[Airy normal form at a simple turn]
\label{thm:newS7V69-Airy}
At a simple scalar turning point, freezing \(A=A_*\) and retaining the
leading term \(B=b_*(a-a_*)\) reduces
\eqref{eq:newS7V69-WDW} to

\begin{equation}
 \Psi''+\gamma(a-a_*)\Psi=0,
 \qquad
 \gamma={b_*\over A_*}.
 \label{eq:newS7V69-Airy-eq}
\end{equation}

For \(\gamma>0\), its solution space is

\begin{equation}
 \Psi(a)
 =c_1\operatorname{Ai}\!\left(-\gamma^{1/3}(a-a_*)\right)
 +c_2\operatorname{Bi}\!\left(-\gamma^{1/3}(a-a_*)\right).
 \label{eq:newS7V69-Airy-solutions}
\end{equation}

Both \(\Psi(a_*)\) and \(\Psi'(a_*)\) are required to fix a solution across
the turn.
\end{theorem}

\begin{proof}
Taylor expansion and division by \(A_*\) give
\eqref{eq:newS7V69-Airy-eq}.  With
\(x=-\gamma^{1/3}(a-a_*)\), it becomes
\(d^2\Psi/dx^2-x\Psi=0\), whose independent solutions are
\(\operatorname{Ai}(x)\) and \(\operatorname{Bi}(x)\).  A second-order
regular ordinary differential equation is uniquely determined by two
Cauchy data.
\end{proof}

\begin{corollary}[The first-order branch omits one matching datum]
\label{cor:newS7V69-missing-datum}
The first-order equation
\[
 i\partial_a\psi
 =\sigma\sqrt{B(a)/A(a)}\,\psi
\]
on the allowed side selects one local WKB current orientation.  It supplies
only one complex boundary value at \(a_*\) and does not determine the two
coefficients in \eqref{eq:newS7V69-Airy-solutions}.  A branch-matching or
forbidden-side boundary condition is additional data.
\end{corollary}

\begin{proof}
The first-order equation fixes \(\partial_a\psi\) from \(\psi\) and one
choice of \(\sigma\).  The second-order Airy equation permits arbitrary
\((\Psi(a_*),\Psi'(a_*))\).  Hence the first-order datum determines only a
one-dimensional subspace unless a second relation is imposed.
\end{proof}

\begin{proposition}[Decay in the forbidden Airy region is a genuine
boundary choice]
\label{prop:newS7V69-decay}
For \(\gamma>0\), the side \(a<a_*\) is exponentially forbidden in the
Airy approximation.  Requiring decay as
\(a_*-a\to+\infty\) sets \(c_2=0\) and fixes a unit-flux reflection phase
on the oscillatory side.  Without that requirement, arbitrary
\(\operatorname{Bi}\) admixture changes the matching.
\end{proposition}

\begin{proof}
For positive argument \(x\to+\infty\),
\(\operatorname{Ai}(x)\) decays and
\(\operatorname{Bi}(x)\) grows.  Thus decay forces \(c_2=0\).
For negative argument the Airy asymptotics express
\(\operatorname{Ai}(-r)\) as equal-amplitude incoming and outgoing
oscillations with a fixed relative phase, so the conserved current has
unit-modulus reflection.  A nonzero \(c_2\) changes those coefficients.
\end{proof}

The decay condition may be physically reasonable in a specific forbidden
region, but it is not implied by the local constraint or by conditional
kernel normalization.

\subsection{Self-adjoint domains and the \(a=0\) endpoint}
\label{sec:newS7V69-domain}

\begin{proposition}[A turning point and a singular endpoint are different]
\label{prop:newS7V69-endpoint}
A simple interior zero of \(B\) is a regular point of the differential
expression \eqref{eq:newS7V69-WDW} when \(A_*>0\); solutions cross it
uniquely from two Cauchy data.  The boundary \(a=0\), where
\(A(a)\propto a^{-1}\) and the matter and curvature potentials can be
singular, requires a separate limit-point/limit-circle or boundary-form
analysis.
\end{proposition}

\begin{proof}
At \(a_*>0\), \(A\) is nonzero and all coefficients are smooth, so ordinary
differential-equation existence and uniqueness apply.  At \(a=0\) the
coefficient \(A\) and terms in \(B\) are not regular in general, so that
argument does not apply.  Self-adjointness is determined by the endpoint
boundary form obtained from \eqref{eq:newS7V69-current}.
\end{proof}

\begin{definition}[Turning-point and branch card, TPBC]
\label{def:newS7V69-TPBC}
A global scale-sector construction satisfies TPBC if it specifies:
\begin{enumerate}
\item one symmetric factor ordering and a self-adjoint domain for the full
constraint, including \(a=0\) and large \(a\);
\item the physical inner product or current sector selected from
\eqref{eq:newS7V69-current};
\item a matching condition at every zero of the operator \(B_a\), including
mode conversion when its spectral projections vary with \(a\);
\item the relation between the second-order constraint solutions and the
expanding/contracting first-order scale-clock branches;
\item a reflection-positive finite-slab realization of that matching; and
\item compatibility with the TT seed, CNMCC, refinement, and RCNC.
\end{enumerate}
\end{definition}

\begin{theorem}[Two honest global choices]
\label{thm:newS7V69-choices}
The present scale construction has exactly two established continuations:
\begin{enumerate}
\item restrict to monotone sectors with
\(B_a\succeq b_*I>0\), treating the branch sign as a superselection input;
or
\item retain the second-order constraint and supply TPBC.
\end{enumerate}
Identifying the two branches at \(B=0\) without either restriction is not
determined by V68.
\end{theorem}

\begin{proof}
The first choice is the domain of
\Cref{thm:newS7V68-scale-H}.  For a crossing zero,
\Cref{prop:newS7V69-negative} ends the positive square-root evolution and
\Cref{cor:newS7V69-missing-datum} proves that additional matching data are
required.  TPBC lists those data.
\end{proof}

\begin{assessment}[Turning-point decision]
\label{ass:newS7V69-status}
A simple turn is not a curvature singularity of the full constraint; it is
a coordinate singularity of intrinsic scale time.  The full Airy-type
constraint can cross it, but only after two Cauchy data and a global domain
or forbidden-side condition are fixed.  No attached result selects that
condition.  Until TPBC is constructed, the rigorous positive Hamiltonian
of V68 applies only to monotone superselection sectors separated from
\(B=0\).
\end{assessment}

\begin{programme}[V70: companion audit and monotone-sector projectivity]
\label{prog:newS7V69-V70}
V70 will perform the compulsory companion audit.  It will then test whether
exact conditional renewal preserves a fixed monotone scale sector:
refinement must keep a uniform spectral floor for \(B_a\), preserve the
branch sign, and intertwine the scale-time transfer kernels.  If detail
energy can close the gap or move a turning surface into the sector, the
first missing uniform estimate will be isolated.
\end{programme}

\begin{firewall}[Claim boundary after seventh-series V69]
\label{fire:newS7V69-boundary}
V69 derives the classical simple-turn normal form and proves failure of the
scale gauge there; proves finite-cutoff square-root continuity while
\(B\geq0\) and impossibility beyond negative spectrum; derives the
second-order Airy model, conserved current, and missing matching datum; and
formulates TPBC and the monotone-sector alternative.

V69 does not choose a turning-point boundary condition, prove a global
self-adjoint constraint, cross \(a=0\), prove monotone-sector stability
under refinement, close GHMC\(^{+}\), RTGC, or RCNC, remove the cutoff, or
construct the Standard Model--Einstein continuum.
\end{firewall}

\begin{theorem}[Seventh-series V69 result]
\label{thm:newS7V69-result}
At a simple zero of \(B\), the full constrained trajectory reverses its
scale momentum and crosses in an external parameter, while the scale factor
has an extremum and ceases to be a valid clock.  The positive square-root
Hamiltonian exists continuously only where \(B\succeq0\) and has no
self-adjoint continuation into negative spectrum.  The full second-order
constraint has an Airy normal form which crosses the turn from two Cauchy
data, but a single scale-clock branch omits one matching datum.  Rigorous
continuation therefore requires TPBC or restriction to a uniformly
positive monotone sector.
\end{theorem}

\begin{proof}
Use \Cref{lem:newS7V69-simple,
thm:newS7V69-reflection,
cor:newS7V69-no-clock-crossing,
lem:newS7V69-continuity,
prop:newS7V69-negative,
lem:newS7V69-current,
thm:newS7V69-Airy,
cor:newS7V69-missing-datum,
prop:newS7V69-decay,
thm:newS7V69-choices}.
\end{proof}

\paragraph{Parent-version record.}
V69 was formed from
\texttt{Renewal\_SM\_Einstein\_Continuum\_Research\_Notes\_V68.tex}.
The V68 parent had \(27\,163\) logical source lines, \(1\,085\,117\) bytes,
and SHA--256
\[
 \texttt{494AEA51000721D8CC0BE002B96268C192528B91C9E66F624063421D409C7BB7}.
\]
The next compulsory companion audit is V70.

% =============================================================================
% END APPENDED SEVENTH-SERIES V69 MATERIAL
% =============================================================================

% =============================================================================
% APPENDED SEVENTH-SERIES V70 MATERIAL
% =============================================================================

\section{Companion audit and monotone-sector refinement}
\label{sec:newS7V70}
\label{sec:v7v70}

This is the compulsory V70 companion audit.  The mathematical continuation
tests whether the monotone scale sectors left by V69 survive exact
conditional renewal.  Compression of the old constraint into the fine
space is not enough: detail and cross blocks can close the positive
spectral floor, and the operator square root does not commute with a mere
compression.  The correct interface is a uniform Schur margin together
with projectivity of the transfer operator itself.

\subsection{Compulsory V70 companion audit}
\label{sec:newS7V70-audit}

At the audit snapshot, the newest distinct complete Yang--Mills checkpoint
was

\begin{align}
 &\texttt{Renewal\_Yang\_Mills\_Mass\_Gap\_Research\_Notes\_V41.tex},
 \notag\\
 &\qquad 642\,499\ {\rm bytes},\quad 18\,431\ {\rm logical\ lines},\quad
 \texttt{845DCFE971A22DB495F669224E046DED32ECF5D9D41FFB796282494D9443094D}.
 \label{eq:newS7V70-YM-hash}
\end{align}

The contemporaneous V42 file was byte-identical to V41 and therefore was
not treated as a mathematical append.  The newest Navier--Stokes checkpoint
remained

\begin{align}
 &\texttt{Renewal\_NS\_Stability\_Research\_Notes\_V26.tex},
 \notag\\
 &\qquad 278\,283\ {\rm bytes},\quad 5\,319\ {\rm logical\ lines},\quad
 \texttt{82C887F8C2C69E4F28FAD7C3DC8DE5B9099CB5A4BF431EBA1FFF3E91935978AD}.
 \label{eq:newS7V70-NS-hash}
\end{align}

Each distinct checkpoint has one live document terminator.

\paragraph{Yang--Mills V38--V41.}
The new block advances from local alignment to a strict recentered
conditional influence row and then to an explicit multi-star
misalignment-cluster tail for the periodic pure-Wilson law.  Its final
firewall is the relevant point here: unconditional contour suppression is
not yet boundary-stable conditional suppression and does not prove the
marked connected expansion.

\paragraph{Navier--Stokes V26.}
This file is unchanged from the V65 audit.  Its rank-one Oseen-kernel
constant refinement does not act on the gravitational scale constraint.

\begin{theorem}[V70 companion-audit decision]
\label{thm:newS7V70-audit}
The YM good-star/bad-cluster split supplies a useful architecture for
monotone-sector stability: prove a strict local floor on aligned details
and a conditional connected activity for the complement.  Its Wilson
constants and unconditional chessboard bounds do not transfer to the
gravitational constraint.  NS V26 supplies no new input.
\end{theorem}

\begin{proof}
The YM local response calculation and contour estimate use compact
\(\mathrm{SU}(3)\), the Wilson plaquette identity, and reflection
chessboards.  The scale constraint has noncompact metric variables and the
operator \(B_a\), so none of those estimates applies numerically.
The distinction between unconditional and conditional activities is
explicit in the YM V41 firewall.
\end{proof}

\subsection{Compression does not preserve a fine spectral floor}
\label{sec:newS7V70-compression}

Let \(I_n:\mathcal H_n\to\mathcal H_{n+1}\) be the isometry induced by an
exact conditional extension, let \(E_n=I_n^*\), and decompose

\[
 \mathcal H_{n+1}=I_n\mathcal H_n\oplus\mathcal D_n.
\]

Suppose the fine radicand has block form

\begin{equation}
 B_{n+1}
 =\begin{pmatrix}
 B_n&C_n^*\\
 C_n&D_n
 \end{pmatrix},
 \qquad
 E_nB_{n+1}I_n=B_n.
 \label{eq:newS7V70-block}
\end{equation}

\begin{proposition}[Exact old compression is insufficient]
\label{prop:newS7V70-compression-no}
The identity \(E_nB_{n+1}I_n=B_n\succeq bI\) does not imply
\(B_{n+1}\succeq0\), nor any positive fine floor.
\end{proposition}

\begin{proof}
Take one-dimensional old and detail spaces and
\[
 B_{n+1}=\begin{pmatrix}b&c\\c&d\end{pmatrix}
\]
with \(b>0\).  Its old compression is \(b\), but its determinant is
\(bd-c^2\).  Choosing \(c^2>bd\) gives one negative eigenvalue.
\end{proof}

\begin{theorem}[Uniform block Schur margin]
\label{thm:newS7V70-Schur}
Assume
\[
 B_n\succeq bI,\qquad D_n\succeq dI,\qquad\|C_n\|\leq c,
 \qquad b,d>0.
 \label{eq:newS7V70-block-bounds}
\]
Then

\begin{equation}
 B_{n+1}\succeq
 \lambda_*(b,d,c)I,
 \qquad
 \lambda_*={b+d-\sqrt{(b-d)^2+4c^2}\over2},
 \label{eq:newS7V70-floor}
\end{equation}

provided \(c^2<bd\).  The constant is optimal given only the three bounds.
\end{theorem}

\begin{proof}
For \(x\in\mathcal H_n\) and \(y\in\mathcal D_n\),
\[
 \left\langle
 \binom{x}{y},B_{n+1}\binom{x}{y}\right\rangle
 \geq b\|x\|^2+d\|y\|^2-2c\|x\|\|y\|.
\]
The smallest eigenvalue of the scalar matrix
\(\left(\begin{smallmatrix}b&-c\\-c&d\end{smallmatrix}\right)\)
is \(\lambda_*\).  It is positive exactly when \(c^2<bd\), proving the
bound.  Equality is attained by scalar blocks
\(B_n=b,D_n=d,C_n=-c\), so no larger universal constant follows from
these data.
\end{proof}

\begin{corollary}[Detail energy can create a turning surface]
\label{cor:newS7V70-turn}
Even when the old sector obeys \(B_n\succeq bI>0\), the fine theory can
acquire a zero or negative mode if the detail floor vanishes or the
old--detail coupling reaches \(c^2\geq bd\).  Exact marginal projectivity
does not prevent this event.
\end{corollary}

\begin{proof}
Use \Cref{prop:newS7V70-compression-no,thm:newS7V70-Schur}.
\end{proof}

\subsection{Square roots do not commute with compression}
\label{sec:newS7V70-square-root}

\begin{proposition}[Nonlinear functional-calculus obstruction]
\label{prop:newS7V70-sqrt}
Even when both \(B_n\) and \(B_{n+1}\) are strictly positive and
\(E_nB_{n+1}I_n=B_n\), one can have

\begin{equation}
 E_nB_{n+1}^{1/2}I_n\ne B_n^{1/2}.
 \label{eq:newS7V70-sqrt-no}
\end{equation}
\end{proposition}

\begin{proof}
Let
\[
 B_n=(1),\qquad
 B_{n+1}=\begin{pmatrix}1&c\\c&1\end{pmatrix},
 \qquad0<|c|<1.
\]
The old compression is one.  Diagonalizing the fine matrix gives
\[
 (B_{n+1}^{1/2})_{11}
 ={\,\sqrt{1+c}+\sqrt{1-c}\over2}<1=B_n^{1/2},
\]
where strict concavity of the square root proves the inequality.
\end{proof}

\begin{lemma}[A reducing embedding does transport the square root]
\label{lem:newS7V70-reducing}
If \(B_{n+1}I_n=I_nB_n\), then for every bounded Borel function \(f\),
\[
 f(B_{n+1})I_n=I_nf(B_n).
\]
In particular, positive square roots and their exponential transfer steps
intertwine exactly.
\end{lemma}

\begin{proof}
The intertwining holds for polynomials by induction.  Uniform approximation
gives continuous functions on the joint compact spectral interval, and
the spectral theorem extends it to bounded Borel functions.
\end{proof}

The reducing identity is much stronger than exact compression.  Generic
old--detail cross terms violate it.

\subsection{Projectivity belongs at transfer-operator level}
\label{sec:newS7V70-transfer}

At scale \(n\), put

\[
 H_n(s)=A(s)^{-1/2}B_n(s)^{1/2},
 \qquad
 T_n(s;\epsilon)=e^{-\epsilon H_n(s)}.
\]

\begin{theorem}[Direct transfer projectivity is sufficient]
\label{thm:newS7V70-transfer}
Suppose, on a common interval \(J\):
\begin{enumerate}
\item \(B_n(s)\succeq b_*I\) for every \(n\) and \(s\in J\);
\item each \(T_n(s;\epsilon)\) is a positive self-adjoint contraction; and
\item
\begin{equation}
 E_nT_{n+1}(s;\epsilon)I_n=T_n(s;\epsilon)
 \label{eq:newS7V70-transfer-projectivity}
\end{equation}
for every allowed step and scale.
\end{enumerate}
Then every finite palindromic scale-time slab has projective
reflection-positive boundary forms on old observables, and the branch sign
may be chosen consistently throughout the tower.
\end{theorem}

\begin{proof}
The uniform floor keeps the scale clock transverse and the positive square
root defined.  A palindromic slab is \(Y_n^*Y_n\) by
\Cref{thm:newS7V68-palindrome}.  Repeated use of
\eqref{eq:newS7V70-transfer-projectivity} shows that its quadratic form on
\(I_n f\) equals the old slab form on \(f\).  Positivity is preserved, and
the absence of a zero mode prevents a change of the fixed classical branch
sign.
\end{proof}

\begin{proposition}[Constraint projectivity does not imply transfer
projectivity]
\label{prop:newS7V70-constraint-not-transfer}
The compression identity in \eqref{eq:newS7V70-block} does not imply
\eqref{eq:newS7V70-transfer-projectivity}.
\end{proposition}

\begin{proof}
The example in \Cref{prop:newS7V70-sqrt} already violates the infinitesimal
identity
\[
 E_nH_{n+1}I_n=H_n.
\]
Expanding both exponentials at \(\epsilon=0\) shows that exact transfer
projectivity would imply that identity, so it fails in the example.
\end{proof}

\begin{corollary}[The shell normalizer must be chosen after the square root]
\label{cor:newS7V70-shell}
When fine constraint cross terms are present, a shell counterterm which
only makes the constraint or static measure projective need not make the
scale-time transfer operator projective.  Exact renewal must either define
the conditional law directly from the physical transfer kernel or add the
normalization induced by its square-root functional calculus.
\end{corollary}

\begin{proof}
This is \Cref{prop:newS7V70-constraint-not-transfer}; the logarithm of the
conditional normalization is the operator-level shell correction required
to restore \eqref{eq:newS7V70-transfer-projectivity}.
\end{proof}

\subsection{Good details and bad turning clusters}
\label{sec:newS7V70-good-bad}

\begin{definition}[Monotone-sector refinement card, MSRC]
\label{def:newS7V70-MSRC}
On a scale-time interval \(J\), MSRC consists of:
\begin{enumerate}
\item an old floor \(B_n(s)\succeq b_*I\) uniform in \(n,s\);
\item on every good detail block,
\[
 D_n(s)\succeq d_*I,\qquad
 \|C_n(s)\|\leq c_*,
 \qquad c_*^2<b_*d_*;
\]
\item a boundary-stable conditional activity for connected clusters on
which the good block bounds fail, summable after two marked physical
writers are inserted;
\item exact transfer projectivity
\eqref{eq:newS7V70-transfer-projectivity}, including its shell normalizer;
\item preservation of the chosen branch orientation and zero-mode law; and
\item compatibility with CNMCC, the TT transfer seed, and RCNC.
\end{enumerate}
\end{definition}

\begin{theorem}[Deterministic good-sector preservation]
\label{thm:newS7V70-good}
MSRC rows 1--2 give the fine good-sector floor
\[
 B_{n+1}(s)\succeq
 \lambda_*(b_*,d_*,c_*)I>0
\]
uniformly on \(J\).  With row 4, all good-sector palindromic slabs are
projective and reflection positive.  Removing the bad-sector restriction
requires the independent connected activity in row 3.
\end{theorem}

\begin{proof}
The floor is \Cref{thm:newS7V70-Schur}; projective reflected positivity is
\Cref{thm:newS7V70-transfer}.  Neither proof integrates configurations on
which the block hypotheses fail, so row 3 is necessary for the full law.
\end{proof}

\begin{assessment}[First unpopulated refinement estimate]
\label{ass:newS7V70-status}
Conditional renewal can preserve a monotone scale sector, but static
projectivity is not enough.  The first absent gravitational estimate is a
uniform fine-detail Schur margin for the physical radicand \(B\).  Even
after that margin is supplied, transfer projectivity must be imposed after
the square root, and exceptional turning clusters need a boundary-stable
marked connected activity.  The audited YM work supplies the architecture
but not these gravitational bounds.
\end{assessment}

\begin{programme}[V71: the scale-detail Schur blocks]
\label{prog:newS7V70-V71}
V71 will compute the second variation of the homogeneous radicand
\(B(a;z)\) with respect to one shell of TT, gauge, and Higgs detail
variables around the free product seed.  It will identify the diagonal
detail floors and old--detail cross operators, determine which blocks obey
the Schur criterion, and isolate the fermion determinant and global energy
blocks that cannot be bounded by the free Hessian.
\end{programme}

\begin{firewall}[Claim boundary after seventh-series V70]
\label{fire:newS7V70-boundary}
V70 performs the compulsory companion audit; proves the optimal block
Schur floor; proves that exact constraint compression neither preserves
fine positivity nor commutes with the square root; gives sufficient direct
transfer-projectivity conditions; and formulates MSRC with its good-detail
and bad-cluster rows.

V70 does not prove the gravitational Schur bounds, boundary-stable turning
cluster activity, transfer shell normalizer, MSRC, TPBC, GHMC\(^{+}\),
RTGC, or RCNC, remove the cutoff, or construct the Standard Model--Einstein
continuum.
\end{firewall}

\begin{theorem}[Seventh-series V70 result]
\label{thm:newS7V70-result}
A monotone scale sector survives one refinement step if the old and detail
radicands have positive floors and their cross norm satisfies
\(c^2<bd\), giving the optimal fine floor
\eqref{eq:newS7V70-floor}.  Exact compression of the constraint is
insufficient because the square root does not commute with compression.
Projective reflection-positive scale-time slabs require direct
intertwining of the physical transfer operators, plus a conditional
connected bound for exceptional turning clusters.  These are precisely the
unpopulated rows of MSRC.
\end{theorem}

\begin{proof}
Use \Cref{thm:newS7V70-audit,
prop:newS7V70-compression-no,
thm:newS7V70-Schur,
prop:newS7V70-sqrt,
lem:newS7V70-reducing,
thm:newS7V70-transfer,
prop:newS7V70-constraint-not-transfer,
thm:newS7V70-good}.
\end{proof}

\paragraph{Parent-version record.}
V70 was formed from
\texttt{Renewal\_SM\_Einstein\_Continuum\_Research\_Notes\_V69.tex}.
The V69 parent had \(27\,583\) logical source lines, \(1\,099\,886\) bytes,
and SHA--256
\[
 \texttt{1341F38CF6F3096BA8B357B3AC00DC0CC5BE7ADC1E50B2A2D8792670439828D7}.
\]
The next compulsory companion audit is V75.

% =============================================================================
% END APPENDED SEVENTH-SERIES V70 MATERIAL
% =============================================================================

% =============================================================================
% APPENDED SEVENTH-SERIES V71 MATERIAL
% =============================================================================

\section{Free Standard Model shell floors and the fermionic Schur
obstruction}
\label{sec:newS7V71}
\label{sec:v7v71}

V70 reduced preservation of a monotone scale sector to a block Schur
estimate for the physical radicand \(B_s\).  This note computes that
estimate at the free product seed.  Every nonzero-momentum bosonic detail
shell has a positive quadratic floor after the physical gauge quotient and
the electroweak vacuum are chosen; orthogonal Fourier shells have zero
old--detail cross block.  Small bosonic interactions preserve the floor by
an explicit inequality.  The fermion determinant is different: its second
variation has an indefinite term and cannot be assigned a positive detail
floor from an overlap gap alone.

\subsection{The extended radicand}
\label{sec:newS7V71-radicand}

On a fixed monotone scale value \(s>0\), extend the homogeneous expression
of V68 by including the physical inhomogeneous modes:

\begin{equation}
 B_s
 =U_{\rm grav}(s)
 +H_{\rm TT}(s)+H_{\rm YM}(s)+H_{\rm H}(s)
 +H_{\rm F}(s)+H_{\rm int}(s),
 \label{eq:newS7V71-B}
\end{equation}

where
\[
 U_{\rm grav}(s)
 =-{3kV_0s\over8\pi G}
  +{\Lambda_RV_0s^3\over8\pi G}.
\]
The vacuum-energy subtraction in \(H_{\rm F}\) and \(H_{\rm H}\) is part
of \(\Lambda_R\), as required in V68.

Use one orthogonal momentum-shell decomposition
\[
 \mathcal H_{\rm phys}
 =\mathcal H_{\rm old}\oplus\mathcal H_{\rm det}
\]
after the linear constraints and gauge quotient.  The shell excludes all
global zero modes and has lattice momenta
\[
 0<\varkappa_-\leq|\widehat k|\leq\varkappa_+.
 \label{eq:newS7V71-shell}
\]

\subsection{Bosonic quadratic blocks}
\label{sec:newS7V71-bosons}

\begin{lemma}[TT detail floor]
\label{lem:newS7V71-TT}
In canonically normalized TT coordinates, the free shell Hamiltonian is

\begin{equation}
 H_{\rm TT,det}^{(2)}
 ={1\over2}\sum_{k,\lambda}
 \left(|p_{\lambda,k}|^2+
 \omega_{\rm TT}(s,k)^2|q_{\lambda,k}|^2\right),
 \label{eq:newS7V71-TT}
\end{equation}

with
\[
 \omega_{\rm TT}(s,k)^2
 =c_2(s)|\widehat k|^2+c_4(s)|\widehat k|^4+m_{\rm TT}(s)^2,
\quad c_2(s)>0,\quad c_4(s),m_{\rm TT}(s)^2\geq0.
\]
Its phase-space Hessian has floor

\begin{equation}
 d_{\rm TT}(s)
 =\min\{1,c_2(s)\varkappa_-^2+
 c_4(s)\varkappa_-^4+m_{\rm TT}(s)^2\}>0.
 \label{eq:newS7V71-TT-floor}
\end{equation}
\end{lemma}

\begin{proof}
The modes diagonalize the quadratic Hamiltonian.  Each momentum block has
Hessian \(\operatorname{diag}(1,\omega_{\rm TT}^2)\), whose least
eigenvalue is the displayed minimum.  The shell lower momentum bound makes
it positive.
\end{proof}

\begin{lemma}[Physical gauge detail floor]
\label{lem:newS7V71-gauge}
After imposing the linear Gauss law and quotienting longitudinal gauge
modes, each transverse gauge detail mode has a Hamiltonian

\begin{equation}
 H_{\rm YM,det}^{(2)}
 ={1\over2}\sum_{k,\alpha}
 \left(g_\alpha(s)^2|E_{\alpha,k}^{T}|^2+
 g_\alpha(s)^{-2}\omega_\alpha(s,k)^2
 |A_{\alpha,k}^{T}|^2\right),
 \label{eq:newS7V71-gauge}
\end{equation}

where \(g_\alpha(s)>0\) and
\(\omega_\alpha(s,k)^2\geq c_\alpha(s)\varkappa_-^2+
m_\alpha(s)^2>0\).  Hence its Hessian floor is

\begin{equation}
 d_{\rm YM}(s)
 =\min_\alpha\{g_\alpha(s)^2,\,
 g_\alpha(s)^{-2}
 [c_\alpha(s)\varkappa_-^2+m_\alpha(s)^2]\}>0.
 \label{eq:newS7V71-gauge-floor}
\end{equation}
\end{lemma}

\begin{proof}
The Gauss constraint removes the longitudinal canonical pair.  The
remaining transverse electric and magnetic quadratic forms are positive;
on the shell their magnetic symbol is bounded below by the stated
nonzero-momentum value.  Electroweak vector masses add a nonnegative term,
while the photon and gluon floors still come from \(\varkappa_->0\).
Diagonalization gives \eqref{eq:newS7V71-gauge-floor}.
\end{proof}

\begin{lemma}[Higgs detail floor about the physical vacuum]
\label{lem:newS7V71-Higgs}
Let the Higgs potential be
\[
 V(H)=\lambda\left(|H|^2-{v^2\over2}\right)^2,
 \qquad\lambda>0,
\]
and expand about one electroweak vacuum.  After quotienting the Goldstone
directions together with the broken gauge directions, the physical radial
detail mode \(h\) has

\begin{equation}
 H_{\rm H,det}^{(2)}
 ={1\over2}\sum_k
 \left(s^{-3}|\pi_{h,k}|^2+
 [s|\widehat k|^2+s^3m_h^2]|h_k|^2\right),
 \qquad m_h^2=2\lambda v^2>0.
 \label{eq:newS7V71-Higgs}
\end{equation}

Its Hessian floor is

\begin{equation}
 d_{\rm H}(s)
 =\min\{s^{-3},s\varkappa_-^2+s^3m_h^2\}>0.
 \label{eq:newS7V71-Higgs-floor}
\end{equation}
\end{lemma}

\begin{proof}
The Hessian of \(V\) at a minimum has eigenvalue \(2\lambda v^2\) in the
radial direction and zero along the gauge orbit.  The physical quotient
removes the latter, while the FRW volume factors give the displayed
canonical kinetic, spatial-gradient, and mass coefficients.  The shell
minimum proves the floor.
\end{proof}

\begin{remark}[Expansion at \(H=0\) would be the wrong seed]
\label{rem:newS7V71-Higgs-zero}
For the broken-phase potential, the Hessian at \(H=0\) has a negative
direction.  A positive Schur estimate therefore requires the physical
vacuum chart or a symmetric-phase parameter regime; it cannot be obtained
by expanding blindly at the origin.
\end{remark}

\begin{theorem}[Free bosonic product shell]
\label{thm:newS7V71-free-boson}
At the translation-invariant free product seed, orthogonal old and detail
Fourier shells have
\[
 C_{\rm bos}^{(0)}=0
\]
in the block decomposition \eqref{eq:newS7V70-block}, and the bosonic
detail block obeys

\begin{equation}
 D_{\rm bos}^{(0)}(s)\succeq
 d_0(s)I,\qquad
 d_0(s):=\min\{d_{\rm TT}(s),d_{\rm YM}(s),d_{\rm H}(s)\}>0.
 \label{eq:newS7V71-d0}
\end{equation}

Consequently the free bosonic refinement preserves every old radicand
floor \(b(s)>0\), with fine floor \(\min\{b(s),d_0(s)\}\).
\end{theorem}

\begin{proof}
Translation invariance diagonalizes every free quadratic operator in
momentum.  Disjoint Fourier bands are orthogonal, so their mixed Hessian
vanishes.  The direct sum floor is the minimum of
\Cref{lem:newS7V71-TT,lem:newS7V71-gauge,lem:newS7V71-Higgs}.
Apply \Cref{thm:newS7V70-Schur} with \(c=0\).
\end{proof}

\subsection{Small bosonic interactions}
\label{sec:newS7V71-perturbation}

Write the bosonic Hessian on a protected field chart as

\begin{equation}
 \operatorname{Hess}B_s
 =\operatorname{Hess}B_s^{(0)}+\mathcal R_s.
 \label{eq:newS7V71-R}
\end{equation}

\begin{theorem}[Quantitative interacting Schur criterion]
\label{thm:newS7V71-interacting}
Suppose on a field chart

\begin{align}
 \|\mathcal R_{DD}(s)\|&\leq r_D(s)<d_0(s),
 \label{eq:newS7V71-rD}\\
 \|\mathcal R_{DO}(s)\|&\leq r_C(s),
 \label{eq:newS7V71-rC}\\
 B_{\rm old}(s)&\succeq b(s)I>0.
 \label{eq:newS7V71-b}
\end{align}

Then the fine bosonic radicand has the uniform lower bound

\begin{equation}
 B_{\rm fine,bos}(s)\succeq
 \lambda_*\bigl(b(s),d_0(s)-r_D(s),r_C(s)\bigr)I>0
 \label{eq:newS7V71-bos-floor}
\end{equation}

whenever

\begin{equation}
 r_C(s)^2<b(s)[d_0(s)-r_D(s)].
 \label{eq:newS7V71-bos-Schur}
\end{equation}
\end{theorem}

\begin{proof}
Weyl's eigenvalue inequality gives
\[
 D_{\rm bos}(s)\succeq[d_0(s)-r_D(s)]I.
\]
The old--detail block has norm at most \(r_C(s)\).  Apply the optimal
Schur theorem \Cref{thm:newS7V70-Schur}.
\end{proof}

\begin{lemma}[Fixed-cutoff small-chart bounds]
\label{lem:newS7V71-small-chart}
At a fixed finite lattice cutoff, the local bosonic Standard Model and
uniformly elliptic metric Hamiltonians are smooth functions of their
coordinates on every compact protected chart.  Since the free
old--detail cross Hessian vanishes, there are chart constants
\(K_D,K_C<\infty\) such that on a radius-\(\varepsilon\) ball,

\begin{equation}
 r_D\leq K_D\varepsilon,\qquad
 r_C\leq K_C\varepsilon.
 \label{eq:newS7V71-small-chart}
\end{equation}
\end{lemma}

\begin{proof}
The lattice field space is finite dimensional.  Plaquette, covariant
difference, Higgs polynomial, and uniformly elliptic inverse-metric
expressions have bounded third derivatives on a compact chart.  The mean
value theorem bounds the Hessian difference from its vacuum value by the
third-derivative supremum times the chart radius.  The cross block is zero
at the vacuum by \Cref{thm:newS7V71-free-boson}.
\end{proof}

\begin{firewall}[Fixed-cutoff smoothness is not a uniform shell theorem]
\label{fire:newS7V71-uniform}
The constants \(K_D,K_C\) in
\eqref{eq:newS7V71-small-chart} may grow with the number and physical scale
of lattice variables.  A cutoff-uniform result requires scaled local
incidence norms and a boundary-stable bad-chart activity, exactly the
unpopulated MSRC rows.  Finite-dimensional continuity alone does not prove
them.
\end{firewall}

\subsection{Fermions do not form a positive Hessian block}
\label{sec:newS7V71-fermions}

There are two inequivalent finite-cutoff organizations.  A canonical
fermionic Fock Hamiltonian may be bounded below after normal ordering, but
Grassmann variables do not give an ordinary positive bosonic Hessian.  If
fermions are integrated out, the determinant modulus contributes an
indefinite second variation.

Let \(Q(z)=D(z)^*D(z)\) on the protected overlap-gap chart and define

\[
 \Phi_{\rm F}(z)=-{1\over2}\operatorname{Tr}\log Q(z).
\]

\begin{lemma}[Exact determinant-modulus Hessian]
\label{lem:newS7V71-det-Hessian}
For a bosonic tangent \(u\),

\begin{equation}
 \delta_u^2\Phi_{\rm F}
 ={1\over2}\operatorname{Tr}
 \left(Q^{-1}\delta_uQ\,Q^{-1}\delta_uQ\right)
 -{1\over2}\operatorname{Tr}
 \left(Q^{-1}\delta_u^2Q\right).
 \label{eq:newS7V71-det-Hessian}
\end{equation}

The first term is nonnegative, while the second has no fixed sign.
\end{lemma}

\begin{proof}
Differentiate
\[
 \delta\operatorname{Tr}\log Q
 =\operatorname{Tr}(Q^{-1}\delta Q)
\]
and use
\(\delta Q^{-1}=-Q^{-1}(\delta Q)Q^{-1}\).
For Hermitian \(\delta Q\), the first trace is the squared Hilbert--Schmidt
norm of
\(Q^{-1/2}(\delta Q)Q^{-1/2}\).  The sign of
\(\operatorname{Tr}(Q^{-1}\delta^2Q)\) is unrestricted.
\end{proof}

\begin{proposition}[An overlap gap is not a fermionic Schur floor]
\label{prop:newS7V71-gap-no-floor}
The spectral bound \(Q\succeq\gamma^2I\) controls the norm of
\(Q^{-1}\) but does not imply
\(\delta^2\Phi_{\rm F}\succeq0\), nor a cutoff-uniform lower Hessian bound.
\end{proposition}

\begin{proof}
The gap gives \(\|Q^{-1}\|\leq\gamma^{-2}\), but
\eqref{eq:newS7V71-det-Hessian} contains the indefinite second term.
For the scalar example \(Q(t)=1+t^2\), the gap is one and
\[
 \Phi_{\rm F}(t)=-{1\over2}\log(1+t^2),
 \qquad
 \Phi_{\rm F}''(0)=-1.
\]
Thus a protected gap coexists with a negative determinant Hessian.  Trace
norm estimates additionally carry the number of fermion modes unless
local connected cancellation is proved.
\end{proof}

\begin{corollary}[The V57 smooth matching theorem has a different role]
\label{cor:newS7V71-V57}
The V57 regulator theorem fixes the local renormalized determinant owners
and their smooth contacts.  It does not supply the nonlinear positive
detail Hessian needed in \eqref{eq:newS7V71-bos-Schur}.
\end{corollary}

\begin{proof}
Smooth regulator equivalence identifies the finite local subtraction and
continuum variation.  The counterexample in
\Cref{prop:newS7V71-gap-no-floor} shows that neither a gap nor a defined
variation implies positivity.
\end{proof}

\subsection{The exact shell card}
\label{sec:newS7V71-card}

\begin{definition}[Scale-detail floor card, SDFC]
\label{def:newS7V71-SDFC}
On a monotone interval \(J\), SDFC requires:
\begin{enumerate}
\item the physical TT, gauge, and Higgs detail coordinates and their
cutoff-uniform free floor \(d_0>0\);
\item scaled local bounds \(r_D,r_C\) satisfying
\eqref{eq:newS7V71-bos-Schur} uniformly in \(s\in J\) and refinement;
\item a canonical positive fermion transfer factorization, or a direct
lower bound for the determinant contribution which compensates the
indefinite term in \eqref{eq:newS7V71-det-Hessian};
\item a boundary-stable connected activity for the complement of the
protected field and overlap-gap charts;
\item the direct transfer projectivity row of MSRC; and
\item compatibility with anomaly contacts, CNMCC, and RCNC.
\end{enumerate}
\end{definition}

\begin{theorem}[What is populated at the product seed]
\label{thm:newS7V71-populated}
At every fixed cutoff and scale value, the free bosonic product seed
populates SDFC row 1 and has zero cross block.  A sufficiently small fixed
chart populates the bosonic part of row 2.  No result through V71
populates row 3 or makes the chart constants uniform in the cutoff.
\end{theorem}

\begin{proof}
Row 1 and the zero cross block are
\Cref{thm:newS7V71-free-boson}.  The fixed-chart statement follows from
\Cref{thm:newS7V71-interacting,lem:newS7V71-small-chart}.  The fermionic
failure is \Cref{prop:newS7V71-gap-no-floor}, and
\Cref{fire:newS7V71-uniform} records the missing uniformity.
\end{proof}

\begin{assessment}[First unresolved shell block]
\label{ass:newS7V71-status}
The TT, physical gauge, and radial Higgs shells are not the first
obstruction: at the free product seed they have explicit positive floors
and no old--detail cross term.  Their interacting Schur condition is
quantitative and nonempty.  The first unresolved block is fermionic.
Integrating it out produces an indefinite determinant Hessian, while
retaining it canonically requires a positive chiral fermion transfer
operator compatible with the determinant line, anomaly, and overlap
refinement.
\end{assessment}

\begin{programme}[V72: canonical fermion transfer seed]
\label{prog:newS7V71-V72}
V72 will retain the finite lattice fermions in Fock space rather than
insert their determinant into a bosonic Hessian.  It will state and prove
the positive-transfer criterion for a quadratic number-conserving fermion
Hamiltonian, include normal-ordering and vacuum-energy transport, and test
which chiral Standard Model features obstruct a gauge-covariant,
reflection-positive one-step factorization.
\end{programme}

\begin{firewall}[Claim boundary after seventh-series V71]
\label{fire:newS7V71-boundary}
V71 computes explicit positive free-shell floors for TT, physical gauge,
and Higgs detail modes; proves zero old--detail cross block at the free
product seed; gives the exact interacting Schur condition and fixed-cutoff
small-chart bounds; and proves by the determinant Hessian formula and a
counterexample that an overlap gap is not a fermionic positive floor.

V71 does not prove cutoff-uniform bosonic chart bounds, construct a
positive chiral fermion transfer seed, control bad charts, prove SDFC,
MSRC, GHMC\(^{+}\), RTGC, or RCNC, remove the cutoff, or construct the
Standard Model--Einstein continuum.
\end{firewall}

\begin{theorem}[Seventh-series V71 result]
\label{thm:newS7V71-result}
Every nonzero-momentum free bosonic detail shell of the reduced
TT--Standard-Model product seed has a positive Hessian floor after the
physical gauge quotient and expansion about the electroweak vacuum.
Orthogonal shells have zero quadratic cross block, and small interactions
preserve the monotone radicand whenever
\(r_C^2<b(d_0-r_D)\).  The integrated fermion determinant does not share
this property: its exact Hessian contains an indefinite second-variation
term, and a spectral gap alone supplies no lower floor.  The next required
seed is therefore a canonical positive fermion transfer construction.
\end{theorem}

\begin{proof}
Use \Cref{lem:newS7V71-TT,
lem:newS7V71-gauge,
lem:newS7V71-Higgs,
thm:newS7V71-free-boson,
thm:newS7V71-interacting,
lem:newS7V71-small-chart,
lem:newS7V71-det-Hessian,
prop:newS7V71-gap-no-floor,
thm:newS7V71-populated}.
\end{proof}

\paragraph{Parent-version record.}
V71 was formed from
\texttt{Renewal\_SM\_Einstein\_Continuum\_Research\_Notes\_V70.tex}.
The V70 parent had \(27\,977\) logical source lines, \(1\,113\,876\) bytes,
and SHA--256
\[
 \texttt{EBDB6ED0FA0D9CEC79EDD48A7236ECFBD98F8766E918FAE6FF34B31400617577}.
\]
The next compulsory companion audit is V75.

% =============================================================================
% END APPENDED SEVENTH-SERIES V71 MATERIAL
% =============================================================================

% =============================================================================
% APPENDED SEVENTH-SERIES V72 MATERIAL
% =============================================================================

\section{Canonical fermion transfer positivity and the chiral carrier gap}
\label{sec:newS7V72}
\label{sec:v7v72}

V71 showed that the integrated determinant is not a positive Schur block.
This note retains fermions on their finite Fock space.  Every Hermitian
finite one-particle Hamiltonian gives a positive transfer contraction after
its filled-sea energy is subtracted, and an operator-valued fermion
multiplier can be inserted into the bosonic heat-kernel sandwich without
losing transfer positivity.  This solves the algebraic positivity problem
for a supplied common carrier.  It does not construct the required chiral
Standard Model carrier, its history transport, or its determinant-line
phase; the attached physical-source programme explicitly leaves that same
input unpopulated.

\subsection{Finite number-conserving Fock seed}
\label{sec:newS7V72-Fock}

Let \(\mathfrak h\) be a finite-dimensional one-particle Hilbert space and
\[
 \mathcal F_-(\mathfrak h)
 =\bigoplus_{r=0}^{\dim\mathfrak h}\wedge^r\mathfrak h
\]
its fermionic Fock space.  For a Hermitian operator \(k\) on
\(\mathfrak h\), let \(d\Gamma(k)\) be its number-conserving second
quantization.

\begin{theorem}[Filled-sea positive transfer]
\label{thm:newS7V72-Fock}
Let \(\varepsilon_1,\ldots,\varepsilon_N\) be the eigenvalues of \(k\) and
put

\begin{equation}
 E_0(k)=\sum_{\varepsilon_j<0}\varepsilon_j,
 \qquad
 H_{\rm F}(k)=d\Gamma(k)-E_0(k)I.
 \label{eq:newS7V72-normal-order}
\end{equation}

Then \(H_{\rm F}(k)\succeq0\).  For every \(\epsilon>0\),

\begin{equation}
 T_{\rm F}(\epsilon;k)=e^{-\epsilon H_{\rm F}(k)}
 \label{eq:newS7V72-TF}
\end{equation}

is a positive self-adjoint contraction with top eigenvalue one.  Moreover,

\begin{equation}
 \operatorname{Tr}_{\mathcal F_-}
 e^{-\beta H_{\rm F}(k)}
 =e^{\beta E_0(k)}
 \det_{\mathfrak h}(I+e^{-\beta k})>0.
 \label{eq:newS7V72-trace}
\end{equation}
\end{theorem}

\begin{proof}
Diagonalize \(k\) and write
\[
 d\Gamma(k)=\sum_{j=1}^N\varepsilon_j n_j,
 \qquad n_j\in\{0,1\}.
\]
Its minimum is attained by filling exactly the negative-energy levels and
equals \(E_0(k)\).  Thus \eqref{eq:newS7V72-normal-order} is nonnegative,
and functional calculus proves the transfer assertions.  The trace
factorizes:
\[
 \operatorname{Tr}e^{-\beta d\Gamma(k)}
 =\prod_j(1+e^{-\beta\varepsilon_j})
 =\det(I+e^{-\beta k}).
\]
Multiplication by \(e^{\beta E_0(k)}\) gives
\eqref{eq:newS7V72-trace}; every factor is positive.
\end{proof}

\begin{corollary}[Grassmann kernels need no pointwise positivity]
\label{cor:newS7V72-Grassmann}
In fermionic coherent states,
\[
 \langle\overline\eta'|
 e^{-\epsilon d\Gamma(k)}|\eta\rangle
 =\exp\{\overline\eta'\,e^{-\epsilon k}\eta\}.
\]
Reflection positivity of the canonical seed follows from the positive Fock
operator \eqref{eq:newS7V72-TF}; it is not a statement that Grassmann
coefficients form a pointwise positive density.
\end{corollary}

\begin{proof}
The coherent-state identity follows by diagonalizing \(k\) and checking
one occupation mode.  The Fock inner product represents the reflected
boundary pairing, so positivity is the operator inequality
\(T_{\rm F}\succeq0\).
\end{proof}

\subsection{Gauge covariance and the physical projector}
\label{sec:newS7V72-gauge}

\begin{lemma}[Second-quantized gauge covariance]
\label{lem:newS7V72-gauge}
Suppose a compact lattice gauge group acts on \(\mathfrak h\) by unitaries
\(u(g)\) and
\[
 k(g\cdot z)=u(g)k(z)u(g)^*.
 \label{eq:newS7V72-one-gauge}
\]
Then
\[
 H_{\rm F}(k(g\cdot z))
 =\Gamma(u(g))H_{\rm F}(k(z))\Gamma(u(g))^*.
 \label{eq:newS7V72-Fock-gauge}
\]
The group average

\begin{equation}
 P_{\rm G}=\int_G\Gamma(u(g))\,dg
 \label{eq:newS7V72-gauge-projector}
\end{equation}

is the orthogonal projector onto gauge-invariant Fock states.
\end{lemma}

\begin{proof}
Second quantization commutes with unitary conjugation, and \(E_0(k)\)
depends only on the spectrum, which is invariant under
\eqref{eq:newS7V72-one-gauge}.  Haar invariance gives
\(P_{\rm G}^*=P_{\rm G}\) and
\[
 P_{\rm G}^2
 =\int_{G\times G}\Gamma(u(gh))\,dg\,dh=P_{\rm G}.
\]
Its range is the invariant subspace by the standard averaging argument.
\end{proof}

\begin{corollary}[Gauge-projected positive step]
\label{cor:newS7V72-gauge-step}
For every positive fermion transfer \(T_{\rm F}\),
\[
 P_{\rm G}T_{\rm F}P_{\rm G}\succeq0.
\]
If \(T_{\rm F}\) commutes with the gauge action, this is its restriction to
the physical gauge-invariant Fock space.
\end{corollary}

\begin{proof}
For every \(\Psi\),
\[
 \langle\Psi,P_{\rm G}T_{\rm F}P_{\rm G}\Psi\rangle
 =\langle P_{\rm G}\Psi,T_{\rm F}P_{\rm G}\Psi\rangle\geq0.
\]
Commutation makes the invariant range reducing.
\end{proof}

\subsection{Joint boson--fermion transfer sandwich}
\label{sec:newS7V72-joint}

Let \(\mathcal Q_a\) be a finite bosonic configuration space with Hilbert
space \(L^2(\mathcal Q_a,dq)\), and let \(G_{\rm B}\succeq0\) be its kinetic
heat contraction.  Suppose \(q\mapsto H_{\rm F}(q)\) is a measurable
nonnegative Fock Hamiltonian and \(V_{\rm B}(q)\) is a bosonic potential
such that

\[
 \mathcal V(q)=V_{\rm B}(q)I+H_{\rm F}(q)\succeq0
\]
after one declared constant shift.

\begin{theorem}[Operator-valued fermion sandwich]
\label{thm:newS7V72-joint}
On
\(\mathcal H=L^2(\mathcal Q_a,dq)\otimes\mathcal F_-(\mathfrak h)\), define
\[
 W=e^{-\epsilon\mathcal V/2},
 \qquad
 \mathsf T=P_{\rm phys}\,W(G_{\rm B}\otimes I)W\,P_{\rm phys},
 \label{eq:newS7V72-joint-T}
\]
where \(P_{\rm phys}\) is an orthogonal gauge and reduced-gravity
projector.  Then \(\mathsf T\) is a positive self-adjoint contraction.
If it is nonzero, spectral normalization gives a positive physical
one-step seed of the form required in V65.
\end{theorem}

\begin{proof}
\(0\preceq W\preceq I\) by functional calculus.  For every \(\Psi\),
\[
 \langle\Psi,\mathsf T\Psi\rangle
 =\langle WP_{\rm phys}\Psi,
 (G_{\rm B}\otimes I)WP_{\rm phys}\Psi\rangle\geq0.
\]
The expression is self-adjoint.  Its norm is at most one because all three
factors are contractions.  A nonzero positive operator may be divided by
its top eigenvalue at finite cutoff.
\end{proof}

\begin{lemma}[Finite-cutoff Higgs domination of the sea energy]
\label{lem:newS7V72-Higgs-bound}
Assume compact gauge links, a uniformly elliptic finite lattice metric, and
a one-particle Hamiltonian
\[
 k(q)=k_{\rm hop}(U,g)+k_{\rm Yuk}(H)
\]
with fixed finite internal multiplicity and Yukawa term linear in the
site Higgs field.  Then

\begin{equation}
 E_0(k(q))
 \geq-C_0|\Lambda_a|-C_1\sum_x|H_x|.
 \label{eq:newS7V72-sea-bound}
\end{equation}

If the bosonic potential contains
\(\lambda\sum_x|H_x|^4-C_2\sum_x|H_x|^2\) with \(\lambda>0\), the combined
finite-cutoff boson--fermion Hamiltonian is bounded below by
\(-C_3|\Lambda_a|\).
\end{lemma}

\begin{proof}
The filled-sea energy satisfies
\[
 E_0(k)\geq-\operatorname{Tr}|k|.
\]
Finite range, compact links, and metric ellipticity bound the trace norm of
the hopping part by \(C_0|\Lambda_a|\); fixed internal multiplicity and
linearity give
\(\operatorname{Tr}|k_{\rm Yuk}|\leq C_1\sum_x|H_x|\).
For every \(\delta>0\),
\[
 C_1|H|+C_2|H|^2\leq\delta|H|^4+C_\delta.
\]
Choose \(\delta<\lambda\), sum over sites, and absorb the remaining constant
into \(C_3|\Lambda_a|\).
\end{proof}

\begin{corollary}[Finite canonical fermions do not destroy lower
boundedness]
\label{cor:newS7V72-lower}
At every finite cutoff, the canonical Fock organization and positive Higgs
quartic permit the constant shift required in
\Cref{thm:newS7V72-joint}.  The shift is extensive and must be transported
together with the renormalized cosmological term.
\end{corollary}

\begin{proof}
Apply \Cref{lem:newS7V72-Higgs-bound}; adding \(C_3|\Lambda_a|\) makes the
combined operator nonnegative.  A volume-proportional constant is a vacuum
energy and hence belongs to the cosmological owner.
\end{proof}

\subsection{Exact refinement at a block-diagonal free seed}
\label{sec:newS7V72-refinement}

\begin{proposition}[Fock factorization of orthogonal shells]
\label{prop:newS7V72-factor}
If
\[
 \mathfrak h_{n+1}=\mathfrak h_n\oplus\mathfrak d_n,
 \qquad
 k_{n+1}=k_n\oplus k_{\rm d},
\]
then, up to the canonical graded tensor identification,

\begin{align}
 \mathcal F_-(\mathfrak h_{n+1})
 &\cong\mathcal F_-(\mathfrak h_n)
 \widehat\otimes\mathcal F_-(\mathfrak d_n),
 \label{eq:newS7V72-Fock-factor}\\
 H_{\rm F}(k_{n+1})
 &=H_{\rm F}(k_n)\otimes I+
 I\otimes H_{\rm F}(k_{\rm d}),
 \label{eq:newS7V72-H-factor}\\
 T_{\rm F}(k_{n+1})
 &=T_{\rm F}(k_n)\otimes T_{\rm F}(k_{\rm d}).
 \label{eq:newS7V72-T-factor}
\end{align}

Taking the detail ground-state conditional expectation exactly returns the
old normalized transfer.
\end{proposition}

\begin{proof}
Exterior algebras factor over direct sums.  Second quantization is additive
and the filled-sea energy is additive under a direct sum, proving
\eqref{eq:newS7V72-H-factor}; exponentiation gives
\eqref{eq:newS7V72-T-factor}.  The normalized detail ground state has
transfer eigenvalue one, so its matrix element leaves the old factor.
\end{proof}

Interactions which mix old and detail one-particle modes destroy this exact
factorization.  Their treatment belongs at the transfer-operator level,
not in a claimed positive determinant Hessian.

\subsection{Why this is not yet the chiral Standard Model transfer}
\label{sec:newS7V72-chiral}

\begin{lemma}[Finite exact chirality has zero square index]
\label{lem:newS7V72-index}
Let a finite-dimensional one-particle space split as
\(\mathfrak h_+\oplus\mathfrak h_-\) with equal dimensions, and let an
exactly chirality-odd operator have form
\[
 D=\begin{pmatrix}0&D_-\\D_+&0\end{pmatrix},
 \qquad D_-=D_+^*.
\]
Then
\[
 \dim\ker D_+-\dim\ker D_-=0.
\]
\end{lemma}

\begin{proof}
\(D_-=D_+^*\).  In finite equal dimensions,
\[
 \dim\ker D_+
 =\dim\operatorname{coker}D_+
 =\dim\ker D_+^*
 =\dim\ker D_-.
\]
\end{proof}

Thus nonzero topological index on a finite lattice requires a modified
chirality relation, unequal boundary degrees, or an additional dimension.
The overlap/Ginsparg--Wilson construction chooses the first option; it does
not automatically provide a local-in-time positive Hamiltonian transfer.

\begin{proposition}[A positive thermal trace does not determine the chiral
determinant line]
\label{prop:newS7V72-phase}
For fixed Hermitian \(k\), the trace
\eqref{eq:newS7V72-trace} is a positive real number.  Therefore it cannot
by itself encode a bosonic-history-dependent complex phase of a chiral
Euclidean determinant.  Such a phase must occur through history transport,
a Berry/line holonomy, boundary degrees, or an equivalent extra structure.
\end{proposition}

\begin{proof}
The right side of \eqref{eq:newS7V72-trace} is strictly positive.
A nonconstant complex phase is information not contained in a positive real
scalar function.
\end{proof}

The attached physical-source V35 manuscript already supplies abstract
transfer-log, conditional reflected-kernel, and record-incidence compilers.
Its own integrated V8 theorem states that no result there identifies the
active chiral common-action transfer, physical connection, or line
incidence of the original Einstein--Standard-Model regulator.  The present
Fock theorem supplies a generic seed criterion; it does not repeat or
silently populate that missing active carrier.

\begin{definition}[Canonical chiral transfer card, CCTC]
\label{def:newS7V72-CCTC}
CCTC consists of:
\begin{enumerate}
\item one common finite physical one-particle carrier for every reflected
bosonic history in a shell, with a Hermitian gauge-covariant Hamiltonian;
\item locality in the time direction, the declared
Ginsparg--Wilson/domain-wall chiral relation, correct species content, and
the Standard Model anomaly cancellation on that same carrier;
\item the positive Fock step and physical gauge projection of
\Cref{thm:newS7V72-joint};
\item a reflection-covariant Hilbert-bundle connection whose holonomy gives
the determinant-line phase and whose real structure gives the required
Pfaffian orientation;
\item exact shell transfer projectivity, including filled-sea vacuum energy
and the cosmological counterterm; and
\item uniform locality, protected-gap probability, and RCNC stress
contacts.
\end{enumerate}
\end{definition}

\begin{theorem}[What the canonical construction closes]
\label{thm:newS7V72-closure}
If CCTC rows 1--2 provide the active one-particle Hamiltonian, then the
finite Fock construction proves the algebraic positivity and lower-bound
parts of row 3 and the free block-diagonal part of row 5.  It does not prove
the chiral carrier, line phase, interacting projectivity, or uniform
continuum rows.
\end{theorem}

\begin{proof}
Use \Cref{thm:newS7V72-Fock,
lem:newS7V72-gauge,
thm:newS7V72-joint,
lem:newS7V72-Higgs-bound,
prop:newS7V72-factor}.  The omitted claims are precisely CCTC rows not used
in those proofs.
\end{proof}

\begin{assessment}[The fermionic bottleneck after canonicalization]
\label{ass:newS7V72-status}
Finite canonical fermions are compatible with a positive physical transfer
operator and with the Higgs lower bound.  The determinant-Hessian sign
problem of V71 was caused by integrating them out before the Schur
argument.  The remaining obstruction is specifically chiral and
history-dependent: construct one active local-in-time
Ginsparg--Wilson or boundary carrier, its positive reflected Fock step, and
its determinant-line connection on the same histories.
\end{assessment}

\begin{programme}[V73: temporal locality of the overlap carrier]
\label{prog:newS7V72-V73}
V73 will split the overlap/Ginsparg--Wilson operator into temporal blocks
and state the necessary nearest-slab factorization for a positive transfer
step.  It will test whether an overlap sign function built from the full
four-dimensional Wilson kernel is local enough in time for that
factorization, and isolate the domain-wall enlargement if direct
nearest-time transfer fails.
\end{programme}

\begin{firewall}[Claim boundary after seventh-series V72]
\label{fire:newS7V72-boundary}
V72 proves the finite filled-sea Fock positivity theorem, thermal
determinant identity, gauge projection, joint operator-valued transfer
sandwich, Higgs domination of the sea energy, and free shell
factorization.  It proves that a positive thermal trace does not encode the
chiral determinant phase and formulates CCTC.

V72 does not construct the active chiral one-particle carrier, prove
overlap/domain-wall time locality and reflection positivity, trivialize
the determinant line, prove interacting transfer projectivity, CCTC,
SDFC, MSRC, RTGC, or RCNC, remove the cutoff, or construct the Standard
Model--Einstein continuum.
\end{firewall}

\begin{theorem}[Seventh-series V72 result]
\label{thm:newS7V72-result}
Any supplied finite Hermitian one-particle Standard Model Hamiltonian has a
nonnegative filled-sea Fock Hamiltonian and a positive transfer
contraction.  Its operator-valued Boltzmann factor can be sandwiched with
the bosonic heat kernel and physical projector while preserving transfer
positivity, and the Higgs quartic controls its finite-cutoff sea-energy
lower bound.  What remains is not generic fermion positivity but the
construction of the common chiral, local-in-time, gauge-covariant carrier
and determinant-line transport required by CCTC.
\end{theorem}

\begin{proof}
Use \Cref{thm:newS7V72-Fock,
cor:newS7V72-Grassmann,
lem:newS7V72-gauge,
thm:newS7V72-joint,
lem:newS7V72-Higgs-bound,
prop:newS7V72-factor,
lem:newS7V72-index,
prop:newS7V72-phase,
thm:newS7V72-closure}.
\end{proof}

\paragraph{Parent-version record.}
V72 was formed from
\texttt{Renewal\_SM\_Einstein\_Continuum\_Research\_Notes\_V71.tex}.
The V71 parent had \(28\,454\) logical source lines, \(1\,130\,191\) bytes,
and SHA--256
\[
 \texttt{ABA4750FB0C1F72E004B2D601C8169D9360E96F89F4BED009F1AF639AF3472C9}.
\]
The next compulsory companion audit is V75.

% =============================================================================
% END APPENDED SEVENTH-SERIES V72 MATERIAL
% =============================================================================

% =============================================================================
% APPENDED SEVENTH-SERIES V73 MATERIAL
% =============================================================================

\section{Gauge-active overlap boundary locality and the conditional collar
card}
\label{sec:newS7V73}
\label{sec:v7v73}

The attached physical-source V13 already proves that the free overlap
crossing is an integral of reflected one-sided squares and gives a
fixed-cutoff gauge-near-free Hilbert--Schmidt window.  Repeating that
factorization would not advance the programme.  This note attacks its
declared gauge-active stop instead.  A uniform Wilson gap and finite range
give an exponentially local Riesz projection.  If two gauge histories
differ in a set \(X\), their overlap signs differ by a kernel exponentially
small in the distances of both endpoints from \(X\).  This converts
near-free positivity from a volume-dependent Hilbert--Schmidt estimate to
a local Schur estimate and identifies the missing probabilistic input as a
boundary-stable connected activity for bad gauge collars.

\subsection{Why the four-dimensional overlap is not a nearest-time
operator}
\label{sec:newS7V73-nonlocal}

For the free Wilson kernel with mass parameter \(m\), the Hermitian symbol
satisfies

\begin{equation}
 H_W(p)^2
 =\sum_{\mu=0}^3\sin^2p_\mu+
 \left(\sum_{\mu=0}^3(1-\cos p_\mu)-m\right)^2.
 \label{eq:newS7V73-HW2}
\end{equation}

\begin{proposition}[Generic temporal non-ultralocality of the overlap sign]
\label{prop:newS7V73-nonultralocal}
For a generic fixed spatial momentum and admissible Wilson mass, the
Fourier symbol
\[
 \operatorname{sgn}H_W(p_0,\mathbf p)
 ={H_W(p_0,\mathbf p)\over
 \sqrt{H_W(p_0,\mathbf p)^2}}
\]
is not a finite Laurent polynomial in \(e^{ip_0}\).  Hence its temporal
kernel has nonzero entries at arbitrarily large time separations.
\end{proposition}

\begin{proof}
At fixed \(\mathbf p\), the denominator squared has form
\[
 \alpha(\mathbf p)-\beta(\mathbf p)\cos p_0
\]
after expanding \eqref{eq:newS7V73-HW2}, with
\(\beta(\mathbf p)\ne0\) for generic \(\mathbf p,m\).  Its inverse square
root has branch points in the complex \(e^{ip_0}\)-plane and is not a
Laurent polynomial.  A translation-invariant kernel has finite temporal
range exactly when its Fourier symbol is a finite Laurent polynomial.
Therefore the overlap sign is not temporally finite range.
\end{proof}

\begin{assessment}[Nearest-slab Fock transfer is not the direct overlap
representation]
\label{ass:newS7V73-nearest}
The generic four-dimensional overlap operator cannot be inserted directly
as the nearest-time one-particle Hamiltonian assumed in the simplest V72
Fock step.  This does not contradict reflection positivity: the supplied
V13 free theorem factorizes the complete nonlocal cross-plane kernel.
For the active gauge theory, one must either prove positivity of that full
boundary kernel or use a local domain-wall/spatial-Hamiltonian enlargement.
\end{assessment}

\subsection{Exponential locality from a uniform Wilson gap}
\label{sec:newS7V73-locality}

Let \(H\) be a self-adjoint matrix on a bounded-degree lattice graph, with
range at most \(R\), row norm bounded by \(M\), and

\begin{equation}
 \operatorname{spec}H\subset[-M,-\gamma]\cup[\gamma,M],
 \qquad\gamma>0.
 \label{eq:newS7V73-gap}
\end{equation}

\begin{lemma}[Finite-range Combes--Thomas bound]
\label{lem:newS7V73-CT}
Let \(z\) have distance at least \(\delta>0\) from
\(\operatorname{spec}H\).  There are \(C_{\delta},c_{\delta}>0\), depending
only on \(\delta,M,R\) and the graph degree, such that

\begin{equation}
 \|(z-H)^{-1}_{xy}\|
 \leq C_\delta e^{-c_\delta d(x,y)}.
 \label{eq:newS7V73-CT}
\end{equation}
\end{lemma}

\begin{proof}
Fix a vertex \(x\) and let \(W_\eta\) multiply the component at \(y\) by
\(e^{\eta d(x,y)}\).  Finite range and the inequality
\(|d(x,u)-d(x,v)|\leq R\) on a nonzero matrix block give
\[
 \|W_\eta HW_\eta^{-1}-H\|
 \leq C_H(e^{\eta R}-1).
\]
Choose \(\eta>0\) so the right side is below \(\delta/2\).  The resolvent
identity and \(\|(z-H)^{-1}\|\leq\delta^{-1}\) then give
\[
 \|W_\eta(z-H)^{-1}W_\eta^{-1}\|\leq2\delta^{-1}.
\]
Taking the \(xy\) block yields \eqref{eq:newS7V73-CT} with
\(c_\delta=\eta\).
\end{proof}

\begin{theorem}[Exponential locality of the Wilson sign]
\label{thm:newS7V73-sign-local}
Under \eqref{eq:newS7V73-gap}, there are \(C,c>0\), uniform in the lattice
volume, such that

\begin{equation}
 \|(\operatorname{sgn}H)_{xy}\|
 \leq Ce^{-cd(x,y)}
 \quad(x\ne y).
 \label{eq:newS7V73-sign-local}
\end{equation}
\end{theorem}

\begin{proof}
Let \(\Gamma_+\) be a fixed contour enclosing
\([\gamma,M]\) and staying a positive distance from both spectral
intervals.  The positive spectral projector is
\[
 P_+(H)={1\over2\pi i}\int_{\Gamma_+}(z-H)^{-1}\,dz,
 \qquad
 \operatorname{sgn}H=2P_+(H)-I.
\]
Apply \Cref{lem:newS7V73-CT} uniformly on the compact contour and integrate.
The identity term is diagonal, proving the off-diagonal estimate.
\end{proof}

\subsection{A localized gauge-history difference}
\label{sec:newS7V73-difference}

Let \(H\) and \(\widetilde H=H+E\) obey the same range, norm, and gap
bounds.  Suppose the nonzero blocks of \(E\) meet a set \(X\) and write
\[
 \|E\|_{\rm row}
 =\sup_u\sum_v\|E_{uv}\|.
\]

\begin{theorem}[Two-sided localized sign perturbation]
\label{thm:newS7V73-local-difference}
There are constants \(C,c>0\), depending only on the common local Wilson
bounds and gap, such that

\begin{equation}
 \|[\operatorname{sgn}\widetilde H-\operatorname{sgn}H]_{xy}\|
 \leq C\|E\|_{\rm row}
 \sum_{u\in X}e^{-c[d(x,u)+d(u,y)]}.
 \label{eq:newS7V73-local-difference}
\end{equation}

For a uniformly locally finite \(X\), the associated operator between two
boundary half-spaces has Schur norm bounded by
\(C'\|E\|_{\rm row}\), independently of total volume.
\end{theorem}

\begin{proof}
Use the common contour and the resolvent identity:
\[
 P_+(\widetilde H)-P_+(H)
 ={1\over2\pi i}\int_{\Gamma_+}
 (z-\widetilde H)^{-1}E(z-H)^{-1}\,dz.
\]
The two Combes--Thomas factors bound the \(xy\) block by a sum over
nonzero \(E_{uv}\).  Finite range lets \(v\) be replaced by a uniformly
bounded neighbourhood of \(u\), which changes only \(C,c\); summing its
row norm gives \eqref{eq:newS7V73-local-difference}.  The graph convolution
of two exponentials has uniformly summable rows and columns.  The Schur
test proves the final statement.
\end{proof}

\begin{corollary}[Remote gauge changes have exponentially small crossing
effect]
\label{cor:newS7V73-remote}
If \(X\) is at graph distance \(r\) from the reflection boundary, then the
positive-to-negative crossing perturbation obeys

\begin{equation}
 \|T_+[\operatorname{sgn}\widetilde H-\operatorname{sgn}H]T_-\|
 \leq C\|E\|_{\rm row}e^{-2cr}.
 \label{eq:newS7V73-remote}
\end{equation}
\end{corollary}

\begin{proof}
Every path from a positive-time endpoint to \(X\) and then to a
negative-time endpoint has total distance at least \(2r\).  Sum
\eqref{eq:newS7V73-local-difference} and apply the Schur test.
\end{proof}

\subsection{A volume-uniform good-collar positivity theorem}
\label{sec:newS7V73-good}

Let \(K_{\times,0}=R_0^*R_0\succeq0\) be the free represented crossing
Gram already proved in the supplied V13 source.  Fix a finite local
positive-time screen \(P\) with

\begin{equation}
 PK_{\times,0}P\succeq\kappa_0P,
 \qquad\kappa_0>0.
 \label{eq:newS7V73-free-floor}
\end{equation}

\begin{theorem}[Good gauge collar preserves the overlap crossing]
\label{thm:newS7V73-good}
Assume:
\begin{enumerate}
\item the Wilson kernel along the straight history path from the free
field to \(U\) has the uniform local bounds and gap
\eqref{eq:newS7V73-gap};
\item in a fixed collar of the reflection plane,
\[
 \|H_W(U)-H_W(I)\|_{\rm row}\leq\delta;
 \]
\item outside that collar, the perturbation starts at distance \(r\); and
\item the Hodge, chiral, reflection, and one-letter synthesis factors have
been transported to one common frame with local perturbation norm at most
\(\delta_{\rm frame}\).
\end{enumerate}
Then the represented gauge crossing satisfies

\begin{equation}
 PK_\times(U)P\succeq
 [\kappa_0-C_1\delta-C_2e^{-2cr}
 -C_3\delta_{\rm frame}]P.
 \label{eq:newS7V73-gauge-floor}
\end{equation}

In particular, the crossing remains positive whenever the bracket is
positive, with constants independent of the total spatial volume.
\end{theorem}

\begin{proof}
The active crossing is the transported block
\[
 K_\times(U)
 =-\varepsilon_\Theta a^{-1}
 T_+\star_0P_+\operatorname{sgn}H_W(U)T_-R_\Theta
\]
from the supplied V13 reduction.  Split the Wilson perturbation into its
collar and remote parts.  Apply
\Cref{thm:newS7V73-local-difference} to the first and
\Cref{cor:newS7V73-remote} to the second.  The product rule bounds the
transported frame changes.  Compression to \(P\) cannot increase operator
norm.  Weyl's lower-eigenvalue inequality applied to
\eqref{eq:newS7V73-free-floor} gives
\eqref{eq:newS7V73-gauge-floor}.
\end{proof}

\begin{assessment}[Improvement over the fixed-cutoff Hilbert--Schmidt
window]
\label{ass:newS7V73-improvement}
The supplied V13 bound uses a global Hilbert--Schmidt norm and can grow
with the number of lattice modes.  The local resolvent proof above uses
row sums and exponential distance from the reflection collar, so its good
screen margin is volume independent.  It still assumes a uniform Wilson
gap and a locally small gauge collar; it is not an arbitrary-gauge theorem.
\end{assessment}

\subsection{The conditional bad-collar problem}
\label{sec:newS7V73-bad}

\begin{definition}[Gauge-active overlap collar card, GAOCC]
\label{def:newS7V73-GAOCC}
GAOCC consists of:
\begin{enumerate}
\item a cutoff-uniform Wilson gap and finite local row bounds on the
protected history set;
\item a local free crossing screen with floor \(\kappa_0>0\);
\item good-collar parameters satisfying the positive bracket in
\eqref{eq:newS7V73-gauge-floor};
\item a boundary-stable conditional multi-defect estimate for connected
bad collar clusters, after conditioning on every explored exterior;
\item cancellation of unmarked bad clusters against the common
normalizer, leaving summable two-writer connected activities;
\item determinant-line phase, reflection frame, and adjacent-shell
transport on the same gauge histories; and
\item compatibility with the canonical Fock seed, SDFC, MSRC, and RCNC.
\end{enumerate}
\end{definition}

\begin{proposition}[The audited YM contour is not GAOCC row 4]
\label{prop:newS7V73-YM-no}
The periodic pure-Wilson multi-star bound audited at V70 does not imply
GAOCC row 4 for the coupled overlap--Higgs--gravity conditional law.
\end{proposition}

\begin{proof}
That estimate is unconditional, uses the bare Wilson action and periodic
reflection chessboards, and its own firewall leaves exterior-stable
conditioning and response-generated interactions open.  GAOCC requires
the estimate after arbitrary allowed exterior conditioning and on the
same coupled histories that define \(H_W\) and its determinant line.
\end{proof}

\begin{theorem}[Conditional completion compiler]
\label{thm:newS7V73-compiler}
If GAOCC rows 1--6 hold, then every fixed represented boundary screen has a
reflection-positive gauge-active overlap crossing after summing good and
bad collars, and its two marked crossing writers have a volume-uniform
connected bound.  This populates the reflection-crossing part of CCTC;
canonical time evolution and continuum projectivity remain separate.
\end{theorem}

\begin{proof}
On the good event, \Cref{thm:newS7V73-good} supplies a strict positive
Gram.  Rows 4--5 give an absolutely summable connected polymer correction
to each two-writer matrix entry, uniformly in the exterior.  If its
operator norm is smaller than the good margin, Weyl's inequality preserves
positivity.  Row 6 places every term in one line and reflection frame.
These statements populate the boundary crossing but do not construct the
canonical generator or refinement maps.
\end{proof}

\begin{assessment}[The active chiral stopping point]
\label{ass:newS7V73-status}
Temporal non-ultralocality is not the decisive obstruction: the free
overlap boundary kernel is already reflection positive, and a uniform
Wilson gap makes gauge changes exponentially local.  The first active
missing estimate is probabilistic and conditional: bad boundary collars
must have a coupled, exterior-stable connected activity small enough to
preserve the local crossing floor.  The current Yang--Mills contour is a
promising bare-law precursor but does not prove that row.
\end{assessment}

\begin{programme}[V74: determinant reweighting of a boundary collar]
\label{prog:newS7V73-V74}
V74 will compare the coupled gauge conditional law to the bare Wilson
collar law.  It will derive the exact Radon--Nikodym oscillation contributed
by the Higgs action and fermion determinant modulus on a finite collar,
then determine which local determinant-resolvent norm would preserve the
YM multi-defect activity under exterior conditioning.  The determinant
phase will remain a separate line row.
\end{programme}

\begin{firewall}[Claim boundary after seventh-series V73]
\label{fire:newS7V73-boundary}
V73 proves generic temporal non-ultralocality of the four-dimensional
overlap sign; proves uniform exponential locality and a two-sided localized
gauge-history perturbation bound from a Wilson gap; derives a
volume-independent good-collar crossing margin; and formulates GAOCC and
its conditional completion compiler without repeating the supplied free
overlap factorization.

V73 does not prove a uniform Wilson gap under the coupled measure,
conditional bad-collar activity, determinant-line transport, GAOCC, CCTC,
SDFC, MSRC, RTGC, or RCNC, remove the cutoff, or construct the Standard
Model--Einstein continuum.
\end{firewall}

\begin{theorem}[Seventh-series V73 result]
\label{thm:newS7V73-result}
Although the four-dimensional overlap sign is not nearest-time local, a
uniform finite-range Wilson gap makes its kernel exponentially local.
A gauge perturbation supported in \(X\) changes the sign kernel by the
two-sided bound \eqref{eq:newS7V73-local-difference}; consequently a
locally near-free reflection collar preserves every fixed free positive
crossing screen with the volume-independent margin
\eqref{eq:newS7V73-gauge-floor}.  Extending this to the active coupled law
requires the exterior-stable bad-collar and line-transport rows of GAOCC.
\end{theorem}

\begin{proof}
Use \Cref{prop:newS7V73-nonultralocal,
lem:newS7V73-CT,
thm:newS7V73-sign-local,
thm:newS7V73-local-difference,
cor:newS7V73-remote,
thm:newS7V73-good,
prop:newS7V73-YM-no,
thm:newS7V73-compiler}.
\end{proof}

\paragraph{Parent-version record.}
V73 was formed from
\texttt{Renewal\_SM\_Einstein\_Continuum\_Research\_Notes\_V72.tex}.
The V72 parent had \(28\,921\) logical source lines, \(1\,146\,496\) bytes,
and SHA--256
\[
 \texttt{43E761121A335CE077DCD86E7054BAAF9BFCBE0C7E8D3857F7CCC47945496FD5}.
\]
The next compulsory companion audit is V75.

% =============================================================================
% END APPENDED SEVENTH-SERIES V73 MATERIAL
% =============================================================================

% =============================================================================
% APPENDED SEVENTH-SERIES V74 MATERIAL
% =============================================================================

\section{Local collar reweighting by Higgs and determinant modulus}
\label{sec:newS7V74}
\label{sec:v7v74}

V73 reduced gauge-active overlap positivity to a conditional bad-collar
activity.  The available Yang--Mills contour estimate concerns the bare
Wilson law, whereas the active collar law is reweighted by Higgs, metric,
and fermion factors.  This note proves the exact comparison theorem needed
to transport a conditional multi-defect bound.  A local logarithmic
oscillation \(r\) changes the activity \(p\) to at most \(e^rp\).  On a
protected fermion singular-value chart, the determinant modulus has such a
local oscillation bound.  A Wilson-kernel gap alone does not provide it:
massless Dirac divisors make the logarithmic oscillation infinite.

\subsection{Conditional likelihood-ratio comparison}
\label{sec:newS7V74-comparison}

Let \(S\) be a finite collar region, \(\eta\) an exterior configuration,
and \(\nu_S^\eta\) a reference conditional probability.  Let

\begin{equation}
 d\mu_S^\eta
 ={e^{-W_S^\eta}\over
 \nu_S^\eta(e^{-W_S^\eta})}\,d\nu_S^\eta.
 \label{eq:newS7V74-RN}
\end{equation}

\begin{lemma}[Exact oscillation comparison]
\label{lem:newS7V74-osc}
For every event \(E\) in the collar,

\begin{equation}
 \mu_S^\eta(E)
 \leq e^{\operatorname{osc}_S W_S^\eta}\nu_S^\eta(E),
 \qquad
 \operatorname{osc}_S W
 :=\sup_SW-\inf_SW.
 \label{eq:newS7V74-osc}
\end{equation}
\end{lemma}

\begin{proof}
The numerator in \eqref{eq:newS7V74-RN} is at most
\(e^{-\inf W}\nu_S^\eta(E)\), while its normalizing denominator is at
least \(e^{-\sup W}\).  Their ratio proves
\eqref{eq:newS7V74-osc}.
\end{proof}

Paying the oscillation of the entire collar would reintroduce its volume.
The comparison must be applied only in the neighbourhood of the marked
defects.

\begin{theorem}[Local reweighting of a conditional multi-defect activity]
\label{thm:newS7V74-local-activity}
Let \(J\) be any finite set of marked boundary cells and let \(N(J)\) be a
fixed-radius neighbourhood satisfying
\[
 |N(J)|\leq \nu_0|J|.
 \]
Assume, uniformly in every exterior condition:
\begin{enumerate}
\item the reference law obeys
\begin{equation}
 \nu^\eta\!\left(\bigcap_{j\in J}\mathcal B_j\right)
 \leq p^{|J|};
 \label{eq:newS7V74-base-activity}
\end{equation}
\item changing only the variables in \(N(J)\) changes the coupled
logarithmic weight by at most
\begin{equation}
 \operatorname{osc}_{N(J)}W^\eta
 \leq r|J|.
 \label{eq:newS7V74-local-osc}
\end{equation}
\end{enumerate}
Then the reweighted conditional law satisfies

\begin{equation}
 \mu^\eta\!\left(\bigcap_{j\in J}\mathcal B_j\right)
 \leq (e^rp)^{|J|}.
 \label{eq:newS7V74-reweighted-activity}
\end{equation}
\end{theorem}

\begin{proof}
Condition first on all variables outside \(N(J)\).  The event depends only
on the marked neighbourhood and the assumed reference estimate remains
valid under that exterior.  Apply
\Cref{lem:newS7V74-osc} in \(N(J)\) and use
\eqref{eq:newS7V74-local-osc}.  Integrating the resulting bound over the
outer variables gives \eqref{eq:newS7V74-reweighted-activity}.
\end{proof}

\begin{corollary}[Subcritical cluster condition after reweighting]
\label{cor:newS7V74-subcritical}
On a dependency graph of maximum forward branching \(D-1\), bad-cluster
paths have an exponentially summable tail whenever

\begin{equation}
 (D-1)e^rp<1.
 \label{eq:newS7V74-subcritical}
\end{equation}
\end{corollary}

\begin{proof}
There are at most \(D(D-1)^{n-1}\) self-avoiding paths of length \(n\)
from a root.  Apply \eqref{eq:newS7V74-reweighted-activity} to each path
and sum the geometric series.
\end{proof}

\subsection{Bosonic collar oscillation on a protected chart}
\label{sec:newS7V74-bosonic}

\begin{lemma}[Higgs--metric local oscillation]
\label{lem:newS7V74-Higgs}
Fix a collar chart with:
\[
 \|H_x\|\leq R_H,\qquad
 \lambda_-I\preceq g_x\preceq\lambda_+I,
\]
compact gauge links, and fixed bounds on the local metric differences
entering the discretized action.  For a finite-range gauge--Higgs--metric
action \(S_{\rm B}\), changing variables only in a set \(X\) gives

\begin{equation}
 \operatorname{osc}_X S_{\rm B}
 \leq L_{\rm B}(R_H,\lambda_\pm)|N_R(X)|
 \leq L_{\rm B}\nu_R|X|,
 \label{eq:newS7V74-bos-osc}
\end{equation}

where the constants are independent of total volume.
\end{lemma}

\begin{proof}
Only action cells within the interaction range \(R\) of \(X\) can change.
Every such local cell is a continuous function on the compact protected
coordinate set, so its oscillation is bounded by one common
\(L_{\rm B}\).  Bounded graph degree gives
\(|N_R(X)|\leq\nu_R|X|\).  Summing the affected cell oscillations proves
\eqref{eq:newS7V74-bos-osc}.
\end{proof}

\begin{firewall}[Noncompact fields require a bad-chart activity]
\label{fire:newS7V74-bos-chart}
There is no uniform \(L_{\rm B}\) on the full Higgs or metric field space.
The probability that a marked collar leaves the amplitude and ellipticity
chart must be included in the bad-cluster activity; a finite-cutoff
pointwise potential bound is not a substitute for that probability.
\end{firewall}

\subsection{Determinant-modulus variation on a singular-value chart}
\label{sec:newS7V74-det}

Let \(D(u)\) be the active finite chiral coefficient along a differentiable
collar interpolation and
\[
 Q(u)=D(u)^*D(u),\qquad
 \Phi_{\rm det}(u)=-{1\over2}\operatorname{Tr}\log Q(u).
\]

\begin{lemma}[Trace-norm determinant variation]
\label{lem:newS7V74-det-var}
If
\[
 Q(u)\succeq\delta_D^2I
 \quad(0\leq u\leq1),
 \label{eq:newS7V74-Dirac-gap}
\]
then

\begin{equation}
 |\Phi_{\rm det}(1)-\Phi_{\rm det}(0)|
 \leq {1\over2\delta_D^2}
 \int_0^1\|\dot Q(u)\|_1\,du.
 \label{eq:newS7V74-det-var}
\end{equation}
\end{lemma}

\begin{proof}
Differentiate:
\[
 {d\Phi_{\rm det}\over du}
 =-{1\over2}\operatorname{Tr}(Q^{-1}\dot Q).
\]
Trace duality and \(\|Q^{-1}\|\leq\delta_D^{-2}\) give
\[
 \left|{d\Phi_{\rm det}\over du}\right|
 \leq {1\over2\delta_D^2}\|\dot Q\|_1.
\]
Integrate over \(u\).
\end{proof}

\begin{lemma}[Localized overlap variation is nuclear per changed cell]
\label{lem:newS7V74-nuclear}
Assume the uniform Wilson gap and localized sign bound of V73.  If a gauge
interpolation changes links in a set \(X\), and all local Hodge and Yukawa
coefficients remain in a protected bounded chart, then

\begin{equation}
 \int_0^1\|\dot Q(u)\|_1\,du
 \leq L_{\rm F}|N_R(X)|
 \leq L_{\rm F}\nu_R|X|,
 \label{eq:newS7V74-nuclear}
\end{equation}

with constants independent of total volume.
\end{lemma}

\begin{proof}
The derivative of each finite-range Wilson block is supported near one
changed link.  Differentiating the Riesz formula used in V73 places that
local derivative between two exponentially decaying resolvents.  Hence its
matrix blocks are bounded by
\[
 Ce^{-c[d(x,X)+d(y,X)]}.
\]
The sum of their trace norms over \(x,y\) is \(O(|X|)\) on a bounded-degree
graph.  The product rule for \(Q=D^*D\), together with bounded protected
Hodge, overlap, and Yukawa factors, preserves this nuclear bound.  Summing
over the fixed-radius action neighbourhood proves
\eqref{eq:newS7V74-nuclear}.
\end{proof}

\begin{theorem}[Protected determinant collar oscillation]
\label{thm:newS7V74-det-osc}
Under \eqref{eq:newS7V74-Dirac-gap} and
\eqref{eq:newS7V74-nuclear}, changing a collar set \(X\) changes the
determinant logarithmic modulus by at most

\begin{equation}
 \operatorname{osc}_X\Phi_{\rm det}
 \leq {L_{\rm F}\nu_R\over2\delta_D^2}|X|.
 \label{eq:newS7V74-det-osc}
\end{equation}
\end{theorem}

\begin{proof}
Join any two protected configurations by an allowed interpolation, apply
\Cref{lem:newS7V74-det-var,lem:newS7V74-nuclear}, and take the supremum.
\end{proof}

\begin{proposition}[A Wilson gap is not a Dirac singular-value gap]
\label{prop:newS7V74-two-gaps}
The lower spectral gap of the Hermitian Wilson kernel used to define
\(\operatorname{sgn}H_W\) does not imply
\eqref{eq:newS7V74-Dirac-gap} for the chiral Dirac coefficient.
\end{proposition}

\begin{proof}
The Wilson gap prevents a discontinuity of its sign and controls overlap
locality.  The overlap Dirac operator can nevertheless have physical or
topological zero modes; indeed its index construction requires such modes
in nontrivial sectors.  Algebraically, a gapped involution
\(\varepsilon_W\) may have a vector on which
\(I+\gamma_5\varepsilon_W\) vanishes.  Thus the two spectral statements are
independent.
\end{proof}

\begin{countertheorem}[Divisors destroy logarithmic oscillation]
\label{cth:newS7V74-divisor}
Let \(D(t)=\operatorname{diag}(t,1,\ldots,1)\).  Then the auxiliary Wilson
kernel may be kept fixed and gapped, while
\[
 -\log|\det D(t)|=-\log|t|
\]
has infinite oscillation on every chart containing \(t=0\).
\end{countertheorem}

\begin{proof}
As \(t\to0\), \(-\log|t|\to+\infty\).  The auxiliary kernel does not occur
in this scalar example and may be chosen independently.
\end{proof}

\subsection{Transport of a bare conditional activity}
\label{sec:newS7V74-transport}

Let \(p_{\rm W}\) be a boundary-stable conditional bare-Wilson defect
activity, if such an estimate is proved.  On a protected collar, combine
\Cref{lem:newS7V74-Higgs,thm:newS7V74-det-osc} and any bounded local
counterterm oscillation into

\begin{equation}
 r_{\rm cpl}
 =\nu_0\left(
 L_{\rm B}\nu_R+
 {L_{\rm F}\nu_R\over2\delta_D^2}
 +L_{\rm ct}\right).
 \label{eq:newS7V74-r}
\end{equation}

\begin{theorem}[Coupled collar activity compiler]
\label{thm:newS7V74-compiler}
If the bare reference bound
\eqref{eq:newS7V74-base-activity} holds conditionally on every exterior
with \(p=p_{\rm W}\), and the interpolation remains in the protected
charts, then the coupled determinant-modulus--Higgs--metric law obeys

\begin{equation}
 \mu^\eta\!\left(\bigcap_{j\in J}\mathcal B_j\right)
 \leq p_{\rm cpl}^{|J|},
 \qquad
 p_{\rm cpl}:=e^{r_{\rm cpl}}p_{\rm W}.
 \label{eq:newS7V74-pcpl}
\end{equation}

It has an exponential bad-collar path tail whenever

\begin{equation}
 (D-1)e^{r_{\rm cpl}}p_{\rm W}<1.
 \label{eq:newS7V74-coupled-subcritical}
\end{equation}
\end{theorem}

\begin{proof}
The three logarithmic weight oscillations add.  Their bound on the
fixed-radius neighbourhood \(N(J)\) is \(r_{\rm cpl}|J|\).  Apply
\Cref{thm:newS7V74-local-activity,cor:newS7V74-subcritical}.
\end{proof}

\begin{firewall}[The compiler has two absent hypotheses]
\label{fire:newS7V74-absent}
The current notes do not prove the exterior-uniform bare activity
\(p_{\rm W}\), and massless/topological chiral sectors do not have a
uniform \(\delta_D>0\).  Therefore
\eqref{eq:newS7V74-pcpl} is a conditional compiler, not a bound for the
active law.
\end{firewall}

\subsection{Divisor sectors and phase separation}
\label{sec:newS7V74-divisor-card}

\begin{definition}[Collar reweighting and divisor card, CRDC]
\label{def:newS7V74-CRDC}
CRDC requires:
\begin{enumerate}
\item a boundary-stable bare conditional multi-defect activity;
\item protected Higgs amplitude and metric ellipticity charts with
summable complement clusters;
\item on every nondivisor stratum, a Dirac singular-value floor and the
localized nuclear derivative bound
\eqref{eq:newS7V74-nuclear};
\item on divisor strata, a reduced determinant or zero-mode-saturated
minor with its own locally bounded logarithmic weight;
\item the subcritical inequality
\eqref{eq:newS7V74-coupled-subcritical} on each stratum;
\item a determinant-line phase and Pfaffian orientation compatible across
stratum boundaries; and
\item the GAOCC marked connected and RCNC stress rows.
\end{enumerate}
\end{definition}

\begin{proposition}[Modulus transport does not solve the phase row]
\label{prop:newS7V74-phase}
All estimates in V74 concern the Radon--Nikodym probability built from the
determinant modulus.  They neither construct nor bound the complex chiral
determinant-line phase.
\end{proposition}

\begin{proof}
\(\Phi_{\rm det}=-(1/2)\operatorname{Tr}\log(D^*D)\) is invariant under
multiplication of \(\det D\) by a unit complex number.  Hence every modulus
estimate is identical for all phases and cannot select one.
\end{proof}

\begin{assessment}[Exact coupled-collar frontier]
\label{ass:newS7V74-status}
The cost of passing from a bare conditional collar law to the coupled
modulus law is now explicit: each marked bad cell pays the finite factor
\(e^{r_{\rm cpl}}\) on a protected nondivisor chart.  The two upstream
inputs are equally sharp.  The Yang--Mills activity must first be made
exterior stable, and physical zero modes require divisor-stratified reduced
determinants.  A Wilson sign gap supplies locality but cannot replace the
Dirac singular-value or divisor row.
\end{assessment}

\begin{programme}[V75: companion audit and reduced determinants on a
divisor]
\label{prog:newS7V74-V75}
V75 will perform the compulsory companion audit.  It will then derive the
finite-dimensional determinant factorization near a constant-rank divisor,
separate the vanishing singular values from the reduced determinant, and
test whether zero-mode insertions give a locally bounded collar weight
without choosing a determinant-line phase.
\end{programme}

\begin{firewall}[Claim boundary after seventh-series V74]
\label{fire:newS7V74-boundary}
V74 proves the exact conditional oscillation comparison, the local
multi-defect reweighting theorem, protected bosonic and determinant-modulus
collar oscillation bounds, and the coupled subcritical activity compiler.
It proves that a Wilson gap is distinct from a Dirac singular-value gap and
that divisor charts require reduced determinants.

V74 does not prove the bare exterior-stable activity, control massless or
topological divisors, construct the determinant phase, prove CRDC, GAOCC,
CCTC, SDFC, MSRC, RTGC, or RCNC, remove the cutoff, or construct the
Standard Model--Einstein continuum.
\end{firewall}

\begin{theorem}[Seventh-series V74 result]
\label{thm:newS7V74-result}
A boundary-stable bare multi-defect activity \(p_{\rm W}\) survives
coupled Higgs--metric--determinant-modulus reweighting with activity
\(p_{\rm cpl}=e^{r_{\rm cpl}}p_{\rm W}\), where
\eqref{eq:newS7V74-r} is volume independent on a protected nondivisor
chart.  The proof uses the localized overlap nuclear derivative and a
Dirac singular-value floor.  The Wilson gap alone does not give that floor,
and determinant divisors make the logarithmic oscillation infinite.
Divisor-stratified reduced determinants and an independent line-phase row
are therefore necessary.
\end{theorem}

\begin{proof}
Use \Cref{lem:newS7V74-osc,
thm:newS7V74-local-activity,
cor:newS7V74-subcritical,
lem:newS7V74-Higgs,
lem:newS7V74-det-var,
lem:newS7V74-nuclear,
thm:newS7V74-det-osc,
prop:newS7V74-two-gaps,
cth:newS7V74-divisor,
thm:newS7V74-compiler,
prop:newS7V74-phase}.
\end{proof}

\paragraph{Parent-version record.}
V74 was formed from
\texttt{Renewal\_SM\_Einstein\_Continuum\_Research\_Notes\_V73.tex}.
The V73 parent had \(29\,319\) logical source lines, \(1\,161\,738\) bytes,
and SHA--256
\[
 \texttt{564826715D07D2211EF192D8408A8534D1888C00F25279FC833190809C57FFE3}.
\]
The next compulsory companion audit is V75.

% =============================================================================
% END APPENDED SEVENTH-SERIES V74 MATERIAL
% =============================================================================

% =============================================================================
% APPENDED SEVENTH-SERIES V75 MATERIAL
% =============================================================================

\section{Companion audit and reduced determinants on constant-rank
divisors}
\label{sec:newS7V75}
\label{sec:v7v75}

This is the compulsory V75 companion audit.  It is followed by the divisor
construction required by V74.  Near a constant-rank chiral divisor, the
ordinary determinant separates into a vanishing small Schur determinant
and a nonvanishing reduced determinant.  Equivalently, the product of
nonzero singular values has a bounded logarithmic derivative under a
nonzero singular-value gap.  Fermionic zero-mode insertions saturate the
missing Grassmann directions and return this reduced determinant.  This
does not license replacing a vanishing vacuum determinant by a reduced one:
the latter is the coefficient of a specified saturated observable and
still carries a determinant-line orientation.

\subsection{Compulsory V75 companion audit}
\label{sec:newS7V75-audit}

At the audit snapshot, the newest distinct complete Yang--Mills checkpoint
was

\begin{align}
 &\texttt{Renewal\_Yang\_Mills\_Mass\_Gap\_Research\_Notes\_V44.tex},
 \notag\\
 &\qquad 680\,399\ {\rm bytes},\quad 19\,544\ {\rm logical\ lines},\quad
 \texttt{A58EF4B86D8E2BA521FE2FC5CA357E39D9302C5E2A26A7AA6F278F0D003EF9B7}.
 \label{eq:newS7V75-YM-hash}
\end{align}

The contemporaneous V45 file was byte-identical and was not counted as a
new mathematical append.  The newest Navier--Stokes checkpoint remained

\begin{align}
 &\texttt{Renewal\_NS\_Stability\_Research\_Notes\_V26.tex},
 \notag\\
 &\qquad 278\,283\ {\rm bytes},\quad 5\,319\ {\rm logical\ lines},\quad
 \texttt{82C887F8C2C69E4F28FAD7C3DC8DE5B9099CB5A4BF431EBA1FFF3E91935978AD}.
 \label{eq:newS7V75-NS-hash}
\end{align}

Each distinct checkpoint has one live document terminator.

\paragraph{Yang--Mills V42--V44.}
The new block turns aligned one-link response into a disagreement coupling
and covariance-decay criterion, then introduces three generated
conditional stocks: a negative Hessian debit, a mixed quotient-response
derivative, and a collar oscillation.  Its response-matched plaquette
marginal is exactly an effective Wilson-coupling shift and adds no separate
leading quotient debit.  The actual generated pushforward has not yet been
shown to populate those stocks.

\paragraph{Navier--Stokes V26.}
The file is unchanged from V65 and V70; it adds no divisor or fermion-line
input.

\begin{theorem}[V75 companion-audit decision]
\label{thm:newS7V75-audit}
The YM collar-oscillation stock independently confirms the local
Radon--Nikodym architecture used in V74.  Its bounded finite-range
compiler can be reused only after the active determinant and metric
pushforward are written as a uniformly summable local interaction.  It
does not control chiral divisors or determinant phases.  NS V26 supplies no
new input.
\end{theorem}

\begin{proof}
YM V44 assumes a compact one-link fibre and an explicit generated
interaction with bounded local derivatives.  At a chiral divisor,
\(-\log|\det D|\) has infinite ordinary oscillation by V74, so it does not
meet that hypothesis until a reduced saturated sector is constructed.
\end{proof}

\subsection{Finite-dimensional divisor factorization}
\label{sec:newS7V75-factor}

Let \(D_0:E\to F\) be a square complex linear map of rank \(N-r\).  Put
\[
 K=\ker D_0,\qquad C=\ker D_0^*,
\qquad \dim K=\dim C=r,
\]
and choose the orthogonal decompositions
\[
 E=K\oplus K^\perp,\qquad F=C\oplus C^\perp.
\]
For a nearby \(D\), write

\begin{equation}
 D=\begin{pmatrix}
 A&B\\ C_1&D_\perp
 \end{pmatrix}:
 K\oplus K^\perp\longrightarrow C\oplus C^\perp.
 \label{eq:newS7V75-block}
\end{equation}

At \(D=D_0\), one has \(A=B=C_1=0\) and
\(D_{\perp,0}:K^\perp\to C^\perp\) invertible.

\begin{theorem}[Exact divisor Schur factorization]
\label{thm:newS7V75-Schur}
If \(D_\perp\) is invertible, then

\begin{equation}
 \det D
 =\det D_\perp\,
 \det S_K(D),
 \qquad
 S_K(D):=A-BD_\perp^{-1}C_1.
 \label{eq:newS7V75-Schur}
\end{equation}

The factor \(\det S_K(D)\) contains every determinant zero near \(D_0\);
\(\det D_\perp\) is nonzero on a sufficiently small neighbourhood.
\end{theorem}

\begin{proof}
Multiply \eqref{eq:newS7V75-block} on the left by the determinant-one block
matrix
\[
 \begin{pmatrix}I&-BD_\perp^{-1}\\0&I\end{pmatrix}.
\]
The result is block lower triangular with diagonal blocks
\(S_K(D)\) and \(D_\perp\), proving \eqref{eq:newS7V75-Schur}.  Invertibility
of \(D_\perp\) is open.
\end{proof}

\begin{definition}[Reduced determinant modulus]
\label{def:newS7V75-reduced}
On a rank-\(N-r\) stratum, let
\(\sigma_1(D),\ldots,\sigma_{N-r}(D)>0\) be the nonzero singular values.
Define

\begin{equation}
 |\det{}'D|
 :=\prod_{j=1}^{N-r}\sigma_j(D)
 =\det{}'(D^*D)^{1/2}.
 \label{eq:newS7V75-reduced}
\end{equation}
\end{definition}

This modulus is basis independent.  A complex reduced determinant is a
section of the determinant line built from the kernel and cokernel; its
phase is not fixed by \eqref{eq:newS7V75-reduced}.

\subsection{Constant-rank logarithmic variation}
\label{sec:newS7V75-variation}

\begin{lemma}[Pseudodeterminant derivative]
\label{lem:newS7V75-pseudodet}
Let \(Q(u)\geq0\) be a differentiable Hermitian family of constant rank,
and suppose every nonzero eigenvalue is at least \(\delta^2>0\).  Then

\begin{equation}
 {d\over du}\log\det{}'Q(u)
 =\operatorname{Tr}\bigl(Q(u)^\dagger\dot Q(u)\bigr),
 \label{eq:newS7V75-pseudodet}
\end{equation}

where \(Q^\dagger\) is the Moore--Penrose inverse.
\end{lemma}

\begin{proof}
Choose a differentiable local spectral frame for the positive spectral
subspace.  If \(\lambda_j(u)>0\) are the positive eigenvalues, then
\[
 {d\over du}\log\det{}'Q
 =\sum_j{\dot\lambda_j\over\lambda_j}.
\]
The diagonal matrix elements of \(\dot Q\) in the spectral frame are
\(\dot\lambda_j\); off-diagonal frame derivatives have zero trace.
Since \(Q^\dagger\) has eigenvalues \(\lambda_j^{-1}\) on the positive
subspace and zero on the kernel, the sum equals the trace in
\eqref{eq:newS7V75-pseudodet}.  Degenerate eigenvalues are handled by
their spectral projections and the same trace identity.
\end{proof}

\begin{theorem}[Reduced-determinant collar oscillation]
\label{thm:newS7V75-reduced-osc}
Let \(D(u)\) have constant rank \(N-r\), put \(Q(u)=D(u)^*D(u)\), and
assume its nonzero singular values are at least \(\delta>0\).  Then

\begin{equation}
 \left|
 \log|\det{}'D(1)|-\log|\det{}'D(0)|
 \right|
 \leq {1\over2\delta^2}
 \int_0^1\|\dot Q(u)\|_1\,du.
 \label{eq:newS7V75-reduced-osc}
\end{equation}

If the localized nuclear bound of V74 is \(L_{\rm F}\nu_R|X|\), the
right side is at most
\(L_{\rm F}\nu_R|X|/(2\delta^2)\), even though the ordinary determinant
vanishes identically on the stratum.
\end{theorem}

\begin{proof}
By \eqref{eq:newS7V75-reduced},
\[
 \log|\det{}'D|={1\over2}\log\det{}'Q.
\]
Use \Cref{lem:newS7V75-pseudodet}, trace duality, and
\(\|Q^\dagger\|\leq\delta^{-2}\), then integrate.
\end{proof}

\begin{corollary}[The divisor singularity has been isolated]
\label{cor:newS7V75-isolated}
Near \(D_0\), the divergent ordinary logarithm splits as
\[
 -\log|\det D|
 =-\log|\det D_\perp|-\log|\det S_K(D)|.
\]
The first term has a bounded local oscillation under a nonzero singular
gap; all divergence transverse to the divisor lies in the finite
\(r\times r\) Schur matrix \(S_K(D)\).
\end{corollary}

\begin{proof}
Take absolute values and logarithms in
\eqref{eq:newS7V75-Schur}.  The reduced factor is controlled by
\Cref{thm:newS7V75-reduced-osc}.
\end{proof}

\subsection{Grassmann zero-mode saturation}
\label{sec:newS7V75-saturation}

\begin{theorem}[Exact saturated Gaussian integral]
\label{thm:newS7V75-saturation}
Choose singular-vector coordinates in which
\[
 D_0=\operatorname{diag}(0,\ldots,0,
 \sigma_{r+1},\ldots,\sigma_N),
 \qquad \sigma_j>0.
\]
Let \((\xi_i,\overline\xi_i)_{i=1}^r\) be the Grassmann coordinates of the
zero modes.  With one fixed Berezin ordering,

\begin{equation}
 \int d\overline\psi\,d\psi\,
 e^{-\overline\psi D_0\psi}
 \prod_{i=1}^r\overline\xi_i\xi_i
 =\varepsilon_{\rm B}
 \prod_{j=r+1}^N\sigma_j
 =\varepsilon_{\rm B}|\det{}'D_0|,
 \label{eq:newS7V75-saturation}
\end{equation}

where \(\varepsilon_{\rm B}=\pm1\) depends only on the declared Berezin
ordering.
\end{theorem}

\begin{proof}
The exponential contains no zero-mode variables.  The inserted product
saturates each zero-mode Berezin integral exactly once.  Every nonzero
mode contributes
\[
 \int d\overline\eta_j\,d\eta_j\,
 e^{-\sigma_j\overline\eta_j\eta_j}
\]
equal to \(\sigma_j\) up to the common ordering sign.  Multiplication gives
\eqref{eq:newS7V75-saturation}.
\end{proof}

\begin{corollary}[Unsaturated zero modes make the vacuum coefficient
vanish]
\label{cor:newS7V75-unsaturated}
With \(r>0\) and no zero-mode insertions,
\[
 \int d\overline\psi\,d\psi\,
 e^{-\overline\psi D_0\psi}=0.
\]
\end{corollary}

\begin{proof}
Each zero-mode Berezin integral contains the integral of one with respect
to an unsaturated Grassmann variable and is zero.
\end{proof}

\begin{firewall}[The reduced determinant is an observable coefficient]
\label{fire:newS7V75-not-vacuum}
Replacing the vanishing vacuum determinant on a divisor by
\(|\det{}'D|\) changes the theory.  The reduced determinant occurs only
after the zero-mode insertions appropriate to the chosen correlation
function or after an independently declared mass/source saturation.  Its
normalization cannot be used as a vacuum probability without that record.
\end{firewall}

\subsection{Kernel bundles and the determinant line}
\label{sec:newS7V75-line}

\begin{proposition}[Change of zero-mode frame]
\label{prop:newS7V75-frame}
Let the kernel and cokernel bases change by
\[
 \xi'=U_K\xi,\qquad\overline\xi'= \overline\xi\,U_C^*,
\qquad U_K,U_C\in U(r).
\]
Then the zero-mode insertion wedge and the complex reduced determinant
transform by reciprocal determinant characters.  Their product in the
saturated integral is invariant, while neither phase is individually
canonical.
\end{proposition}

\begin{proof}
The ordered kernel wedge gains \(\det U_K\), and the cokernel wedge gains
\(\overline{\det U_C}\).  The determinant of the complementary map changes
by the inverse factors because the determinant of the full coordinate
transformation is fixed.  The factors cancel in the product.  This is
exactly the transition law of the determinant line
\(\det C\otimes(\det K)^*\).
\end{proof}

\begin{definition}[Divisor saturation and line card, DSLC]
\label{def:newS7V75-DSLC}
On every constant-index stratum, DSLC requires:
\begin{enumerate}
\item smooth kernel and cokernel spectral projectors separated from the
nonzero singular spectrum by \(\delta>0\);
\item a declared zero-mode insertion or mass/source saturation matching
the physical observable;
\item the reduced-modulus oscillation
\eqref{eq:newS7V75-reduced-osc} with a local nuclear stock;
\item transition functions for the kernel, cokernel, determinant, and
Pfaffian lines, including the Berezin orientation;
\item compatibility on overlaps of divisor charts and under adjacent
renewal shells; and
\item the GAOCC, CCTC, and RCNC crossing and stress rows on the same
saturated histories.
\end{enumerate}
\end{definition}

\begin{theorem}[Divisor contribution to the collar compiler]
\label{thm:newS7V75-divisor-compiler}
On a fixed constant-rank stratum satisfying DSLC rows 1--3, a specified
zero-mode-saturated modulus has the same local reweighting form as the
nondivisor determinant in V74, with \(\delta_D\) replaced by the smallest
nonzero singular value \(\delta\).  Therefore the V74 activity compiler
applies to that saturated observable coefficient.  It does not apply to
the unsaturated vacuum coefficient and does not populate the phase row.
\end{theorem}

\begin{proof}
\Cref{thm:newS7V75-saturation} identifies the modulus as
\(|\det{}'D|\), and
\Cref{thm:newS7V75-reduced-osc} gives its local logarithmic oscillation.
Insert that bound into
\Cref{thm:newS7V74-local-activity}.  The two exclusions are
\Cref{cor:newS7V75-unsaturated,prop:newS7V75-frame}.
\end{proof}

\begin{assessment}[What the divisor construction resolves]
\label{ass:newS7V75-status}
The logarithmic singularity of a determinant divisor is now separated
from its regular nonzero-mode factor.  On a constant-rank stratum, the
zero-mode-saturated coefficient has a volume-local reduced-determinant
oscillation and can enter the V74 collar compiler.  What remains is
physical rather than algebraic: specify which Standard Model observables
saturate each index sector, transport their determinant-line phase, and
control transitions where the rank changes.
\end{assessment}

\begin{programme}[V76: rank-changing crossings]
\label{prog:newS7V75-V76}
V76 will analyze a transverse simple singular-value crossing.  It will
derive the signed small Schur eigenvalue, show how the determinant-line
orientation changes across the divisor, and test whether a mass or source
regularization has a limit independent of the side of approach.  The
result will determine the exact additional datum needed to glue the
constant-rank DSLC charts.
\end{programme}

\begin{firewall}[Claim boundary after seventh-series V75]
\label{fire:newS7V75-boundary}
V75 performs the compulsory companion audit; proves the finite divisor
Schur factorization, constant-rank pseudodeterminant derivative and local
oscillation bound, exact Grassmann zero-mode saturation identity, and
determinant-line frame transformation; and gives the conditional divisor
collar compiler.

V75 does not specify the physical saturation in every Standard Model index
sector, cross a rank-changing divisor, construct the line phase, prove
DSLC, CRDC, GAOCC, CCTC, MSRC, RTGC, or RCNC, remove the cutoff, or
construct the Standard Model--Einstein continuum.
\end{firewall}

\begin{theorem}[Seventh-series V75 result]
\label{thm:newS7V75-result}
Near a rank-\(N-r\) chiral divisor, the full determinant is the product of
a regular reduced determinant and the \(r\times r\) small Schur
determinant.  Along a constant-rank stratum with nonzero singular-value
floor \(\delta\), the reduced modulus has the local logarithmic bound
\eqref{eq:newS7V75-reduced-osc}.  Exactly saturated zero-mode insertions
produce that reduced determinant, while the unsaturated vacuum coefficient
remains zero.  The saturated modulus may therefore enter the V74
reweighting compiler, but only with its observable record and independent
determinant-line phase.
\end{theorem}

\begin{proof}
Use \Cref{thm:newS7V75-audit,
thm:newS7V75-Schur,
lem:newS7V75-pseudodet,
thm:newS7V75-reduced-osc,
cor:newS7V75-isolated,
thm:newS7V75-saturation,
cor:newS7V75-unsaturated,
prop:newS7V75-frame,
thm:newS7V75-divisor-compiler}.
\end{proof}

\paragraph{Parent-version record.}
V75 was formed from
\texttt{Renewal\_SM\_Einstein\_Continuum\_Research\_Notes\_V74.tex}.
The V74 parent had \(29\,764\) logical source lines, \(1\,177\,192\) bytes,
and SHA--256
\[
 \texttt{5D679D35BF95DE1DE9193B233E8E696AB581C7E4DF0E61077864F21134D2B0D1}.
\]
The next compulsory companion audit is V80.

% =============================================================================
% END APPENDED SEVENTH-SERIES V75 MATERIAL
% =============================================================================

% =============================================================================
% APPENDED SEVENTH-SERIES V76 MATERIAL
% =============================================================================

\section{Transverse rank crossings and determinant-line parity}
\label{sec:newS7V76}
\label{sec:v7v76}

V75 controlled the nonzero-mode determinant on a constant-rank divisor.
That result leaves a codimension-one gluing problem: when one mode crosses
zero, the ordinary determinant vanishes linearly, its modulus is continuous,
and its phase changes between the two adjacent invertible chambers.  This
section computes that change from the small Schur complement.  It also
shows that a positive mass regulator does not choose a limit independent
of the order of limits.  Thus a regulator cannot replace the missing
determinant-line orientation.

\subsection{The scalar Schur coordinate at a simple crossing}
\label{sec:newS7V76-Schur}

Let \(D(t):E\to F\), \(|t|<t_0\), be a \(C^2\) family of complex linear
maps between Hermitian spaces of the same finite dimension \(N\).  Suppose
that \(D(0)\) has rank \(N-1\).  Choose unit vectors
\[
 e_0\in\ker D(0),\qquad f_0\in\ker D(0)^*,
\]
and use the fixed orthogonal decompositions
\[
 E=\mathbb C e_0\oplus e_0^\perp,
 \qquad F=\mathbb C f_0\oplus f_0^\perp.
\]
Write
\begin{equation}
 D(t)=
 \begin{pmatrix}a(t)&b(t)\\ c(t)&G(t)\end{pmatrix},
 \qquad
 s(t):=a(t)-b(t)G(t)^{-1}c(t).
 \label{eq:newS7V76-block}
\end{equation}
Here \(a\) and \(s\) are scalars, \(b\) is a row, \(c\) is a column, and
\(G(0):e_0^\perp\to f_0^\perp\) is invertible.

\begin{lemma}[First variation of the small Schur coordinate]
\label{lem:newS7V76-first}
For sufficiently small \(|t|\), \(G(t)\) is invertible and
\begin{equation}
 s(t)=\alpha t+O(t^2),
 \qquad
 \alpha:=\langle f_0,\dot D(0)e_0\rangle.
 \label{eq:newS7V76-alpha}
\end{equation}
In particular, \(\alpha\ne0\) is the intrinsic transversality condition
for this framed divisor chart.
\end{lemma}

\begin{proof}
At \(t=0\), the kernel and cokernel equations give
\(a(0)=b(0)=c(0)=0\).  Openness of invertibility gives the first claim.
Differentiate \(s=a-bG^{-1}c\) at zero.  Every term obtained from the
second summand contains \(b(0)\) or \(c(0)\), hence vanishes.  Therefore
\(\dot s(0)=\dot a(0)=\langle f_0,\dot D(0)e_0\rangle\).  Taylor's theorem
proves \eqref{eq:newS7V76-alpha}.
\end{proof}

\begin{theorem}[Linear determinant zero and chamber phase jump]
\label{thm:newS7V76-phase-jump}
Assume \(\alpha\ne0\), and put \(g_0:=\det G(0)\ne0\).  Then
\begin{equation}
 \det D(t)=g_0\alpha t+O(t^2).
 \label{eq:newS7V76-det-linear}
\end{equation}
Consequently \(D(t)\) is invertible for all sufficiently small
\(t\ne0\), and
\begin{equation}
 \lim_{t\downarrow0}{\det D(t)\over|\det D(t)|}
 =-{\lim}_{t\uparrow0}{\det D(t)\over|\det D(t)|}
 ={g_0\alpha\over|g_0\alpha|}.
 \label{eq:newS7V76-phase-limits}
\end{equation}
Thus a transverse simple crossing has determinant-line transition factor
\(-1\) between the two real parameter chambers.
\end{theorem}

\begin{proof}
The V75 Schur identity gives \(\det D(t)=\det G(t)s(t)\).  Since
\(\det G(t)=g_0+O(t)\), \Cref{lem:newS7V76-first} yields
\eqref{eq:newS7V76-det-linear}.  Its leading coefficient is nonzero, so
the determinant has no other zero after the neighbourhood is reduced.
Division by the modulus and one-sided passage to the limit gives
\eqref{eq:newS7V76-phase-limits}.
\end{proof}

\begin{corollary}[Crossing multiplicity]
\label{cor:newS7V76-multiplicity}
If an isolated crossing has a small \(r\times r\) Schur matrix satisfying
\[
 \det S(t)=\beta t^m+O(t^{m+1}),\qquad \beta\ne0,
\]
then the chamber transition factor is \((-1)^m\).  Only the parity of the
vanishing order is visible to this two-chamber comparison.
\end{corollary}

\begin{proof}
The complementary determinant is continuous and nonzero.  Replacing
\(t\) by \(-t\) multiplies the leading small determinant by \((-1)^m\).
\end{proof}

\subsection{The small singular value}
\label{sec:newS7V76-singular}

\begin{theorem}[Sharp first order of the crossing singular value]
\label{thm:newS7V76-singular}
Under the hypotheses of \Cref{thm:newS7V76-phase-jump}, the unique
singular value of \(D(t)\) tending to zero obeys
\begin{equation}
 \sigma_{\min}(D(t))=|\alpha|\,|t|+O(t^2).
 \label{eq:newS7V76-singular}
\end{equation}
Every other singular value stays above a positive constant.
\end{theorem}

\begin{proof}
Block Gaussian elimination gives the exact factorization
\begin{equation}
 D(t)=
 \begin{pmatrix}1&bG^{-1}\\0&I\end{pmatrix}
 \begin{pmatrix}s&0\\0&G\end{pmatrix}
 \begin{pmatrix}1&0\\G^{-1}c&I\end{pmatrix}.
 \label{eq:newS7V76-LDU}
\end{equation}
The two triangular factors and their inverses have norm \(1+O(|t|)\),
because \(b(t),c(t)=O(t)\).  Hence the least singular value of \(D(t)\)
differs from that of \(\operatorname{diag}(s(t),G(t))\) by a multiplicative
factor \(1+O(|t|)\).  The singular values of \(G(t)\) remain uniformly
positive, while \(|s(t)|=|\alpha||t|+O(t^2)\).  This proves
\eqref{eq:newS7V76-singular}.  The same factorization and the min--max
principle keep the remaining singular values away from zero.
\end{proof}

\begin{corollary}[Optimal logarithmic blow-up]
\label{cor:newS7V76-log}
Along either invertible chamber,
\begin{equation}
 {d\over dt}\log|\det D(t)|={1\over t}+O(1),
 \qquad t\to0.
 \label{eq:newS7V76-log}
\end{equation}
Thus the ordinary determinant oscillation across any collar containing the
divisor is infinite, whereas the complementary reduced determinant stays
regular.
\end{corollary}

\begin{proof}
Factor \(\det D(t)=t(g_0\alpha+O(t))\) and differentiate its logarithmic
modulus.  The regular complementary factor is the V75 reduced factor.
\end{proof}

\subsection{Mass and source regulators do not remove the gluing datum}
\label{sec:newS7V76-regulator}

The local issue is already present in the scalar Schur coordinate.  For
clarity, suppose a real structure has been chosen and \(s(t)\) and
\(\alpha\) are real with \(\alpha>0\).  A positive scalar mass replaces
\(s(t)\) by \(s(t)+m\).

\begin{theorem}[Noncommuting mass and crossing limits]
\label{thm:newS7V76-noncommuting}
Let \(s(t)=\alpha t+O(t^2)\), \(\alpha>0\), and define
\[
 q(t,m):={s(t)+m\over|s(t)+m|}
\]
where the denominator is nonzero.  Then
\begin{align}
 \lim_{m\downarrow0}\lim_{t\to0}q(t,m)&=+1,
 \label{eq:newS7V76-limit-a}\\
 \lim_{t\downarrow0}\lim_{m\downarrow0}q(t,m)&=+1,
 \qquad
 \lim_{t\uparrow0}\lim_{m\downarrow0}q(t,m)=-1.
 \label{eq:newS7V76-limit-b}
\end{align}
More generally, along paths \(m=\kappa t+o(t)\) the limiting sign is the
sign of \(\alpha+\kappa\) on the positive side whenever
\(\alpha+\kappa\ne0\).  Hence the regulator limit is not independent of
the order or path of approach.
\end{theorem}

\begin{proof}
For fixed \(m>0\), continuity gives \(s(t)+m>0\) near zero, proving
\eqref{eq:newS7V76-limit-a}.  For fixed nonzero \(t\) sufficiently close
to zero, taking \(m\downarrow0\) returns \(\operatorname{sgn}s(t)\), which
has the two limits in \eqref{eq:newS7V76-limit-b}.  Substitution of
\(m=\kappa t+o(t)\) gives
\(s(t)+m=(\alpha+\kappa)t+o(t)\).
\end{proof}

\begin{proposition}[Complex regulator ray dependence]
\label{prop:newS7V76-complex-ray}
For the complex small factor \(s(t)=\alpha t+O(t^2)\), a regulator
\(m=\rho e^{i\theta}\) selects phase \(e^{i\theta}\) if \(t\to0\) first
and \(\rho\downarrow0\) second.  Taking \(\rho\downarrow0\) first returns
the two phases in \eqref{eq:newS7V76-phase-limits}.  Therefore changing
the regulator ray changes the extension across the zero.
\end{proposition}

\begin{proof}
At fixed \(\rho>0\), \(s(t)+\rho e^{i\theta}\to\rho e^{i\theta}\).
The reversed order removes the regulator before the chamber limit and is
then \Cref{thm:newS7V76-phase-jump}.
\end{proof}

\begin{firewall}[A mass prescription is extra structure]
\label{fire:newS7V76-mass}
A mass or source may define a physically intended deformation, but its
phase, sign, and order of removal must be recorded.  The existence of a
regulated nonzero determinant does not prove that the unregulated chiral
measure has a canonical chamber-independent phase.
\end{firewall}

\subsection{Saturated value, normal derivative, and line gluing}
\label{sec:newS7V76-saturated}

\begin{theorem}[Value and normal derivative carry different data]
\label{thm:newS7V76-value-derivative}
Fix the kernel and cokernel frames used above.  At the crossing, the
zero-mode-saturated Gaussian coefficient is, up to the declared Berezin
sign,
\begin{equation}
 Z_{\rm sat}(0)=\varepsilon_{\rm B}\det G(0).
 \label{eq:newS7V76-sat-value}
\end{equation}
The ordinary determinant instead satisfies
\begin{equation}
 {d\over dt}\det D(t)\big|_{t=0}=\det G(0)\,
 \langle f_0,\dot D(0)e_0\rangle.
 \label{eq:newS7V76-normal-derivative}
\end{equation}
Thus saturation supplies the divisor value of a specified observable,
while the normal crossing map supplies the additional gluing factor.
Neither determines the other.
\end{theorem}

\begin{proof}
Equation \eqref{eq:newS7V76-sat-value} is the V75 saturated Gaussian
identity expressed in the chosen complementary bases.  Differentiate the
Schur factorization at zero and use \(s(0)=0\) and
\(\dot s(0)=\langle f_0,\dot D(0)e_0\rangle\).  This yields
\eqref{eq:newS7V76-normal-derivative}.
\end{proof}

\begin{definition}[Rank-crossing line card, RCLC]
\label{def:newS7V76-RCLC}
For each rank-changing chart, RCLC records:
\begin{enumerate}
\item fixed kernel and cokernel projectors at the crossing and their
oriented or unitary frames on overlaps;
\item the small Schur map \(S\), its vanishing order, and for a simple
crossing the normal map \(\alpha=\langle f_0,\dot D e_0\rangle\);
\item the chamber transition factor \((-1)^m\), together with any complex
phase holonomy not visible to its parity;
\item the physical zero-mode insertions and Berezin ordering on the
divisor;
\item every mass or source regulator, its phase, and its prescribed order
of removal; and
\item compatibility with DSLC on the divisor and CRDC on both adjacent
invertible collars.
\end{enumerate}
\end{definition}

\begin{theorem}[Finite crossing-gluing criterion]
\label{thm:newS7V76-gluing}
Consider finitely many fermion factors crossing the same real
codimension-one wall, with determinant vanishing orders \(m_j\).  The
tensor-product determinant phase has chamber transition
\begin{equation}
 \tau_{\rm wall}=(-1)^{\sum_j m_j}.
 \label{eq:newS7V76-wall-parity}
\end{equation}
A real nonvanishing scalar trivialization can agree on both chambers only
if \(\sum_jm_j\) is even.  If it is even, the parity obstruction on this
single overlap vanishes; complex holonomy and compatibility on multiple
overlaps remain separate conditions.
\end{theorem}

\begin{proof}
Tensor-product transition functions multiply.  Apply
\Cref{cor:newS7V76-multiplicity} to every factor to obtain
\eqref{eq:newS7V76-wall-parity}.  A scalar trivialization agreeing across
the overlap requires transition \(+1\), giving necessity.  When the sum
is even, this particular transition is \(+1\), which proves only the
stated local sufficiency.  Nothing in the parity computation controls a
phase accumulated around a loop or a higher overlap cocycle.
\end{proof}

\begin{assessment}[What has and has not been glued]
\label{ass:newS7V76-status}
The rank-changing singularity is now reduced to finite data: the normal
small Schur map, its parity, the saturated divisor observable, and the
regulator prescription.  The local sign obstruction cancels for an even
total crossing multiplicity.  To use this in the Standard Model one must
next compute the tensor-product anomaly line for its actual chiral
representations.  Local perturbative anomaly cancellation and the global
\(SU(2)\) parity are necessary tests; they are not consequences of the
abstract crossing algebra.
\end{assessment}

\begin{programme}[V77: Standard Model anomaly-line arithmetic]
\label{prog:newS7V76-V77}
V77 will use left-handed one-generation Standard Model multiplets to
compute the \(SU(3)^3\), \(SU(3)^2U(1)\), \(SU(2)^2U(1)\), \(U(1)^3\),
and mixed gravitational--\(U(1)\) coefficients, and the mod-two
\(SU(2)\) doublet count.  It will state exactly what these cancellations
prove about the local determinant line and what still requires a
regulator-level common-action construction and global holonomy control.
\end{programme}

\begin{firewall}[Claim boundary after seventh-series V76]
\label{fire:newS7V76-boundary}
V76 proves the simple-crossing Schur expansion, determinant phase jump,
first order of the vanishing singular value, logarithmic blow-up,
noncommuting regulator limits, distinction between saturated value and
normal derivative, and the finite wall-parity criterion.

V76 does not prove Standard Model anomaly cancellation inside the renewal
regulator, trivialize the full complex determinant or Pfaffian line,
control nontransverse or accumulated crossings, prove RCLC, DSLC, CRDC,
GAOCC, CCTC, MSRC, RTGC, or RCNC, remove the cutoff, or construct the
Standard Model--Einstein continuum.
\end{firewall}

\begin{theorem}[Seventh-series V76 result]
\label{thm:newS7V76-result}
At a transverse one-mode rank crossing,
\(\det D(t)=\det G(0)\alpha t+O(t^2)\) and
\(\sigma_{\min}(D(t))=|\alpha||t|+O(t^2)\).  The determinant phases in the
two chambers differ by \(-1\); mass or source regularization has a
side-, path-, or ray-dependent removal limit.  The saturated divisor
coefficient supplies \(\det G(0)\), while the normal map supplies
\(\alpha\).  For a product of crossings, the local real gluing obstruction
is the parity \((-1)^{\sum_jm_j}\).
\end{theorem}

\begin{proof}
Combine \Cref{lem:newS7V76-first,
thm:newS7V76-phase-jump,
thm:newS7V76-singular,
thm:newS7V76-noncommuting,
prop:newS7V76-complex-ray,
thm:newS7V76-value-derivative,
thm:newS7V76-gluing}.
\end{proof}

\paragraph{Parent-version record.}
V76 was formed from
\texttt{Renewal\_SM\_Einstein\_Continuum\_Research\_Notes\_V75.tex}.
The V75 parent had \(30\,192\) logical source lines, \(1\,192\,622\) bytes,
and SHA--256
\[
 \texttt{792860B564969268DF126EC2DAED20D5871E892E13820943B67285E928BEB60C}.
\]
The next compulsory companion audit remains V80.

% =============================================================================
% END APPENDED SEVENTH-SERIES V76 MATERIAL
% =============================================================================
% =============================================================================
% APPENDED SEVENTH-SERIES V77 MATERIAL
% =============================================================================

\section{Standard Model anomaly-line arithmetic at one generation}
\label{sec:newS7V77}
\label{sec:v7v77}

V76 reduced a simple rank-changing wall to a determinant-line transition.
A tensor product of fermion factors can glue only if its phase
obstructions cancel.  The first necessary test is representation
arithmetic.  This section performs that test for one Standard Model
generation using left-handed Weyl fields.  All local gauge and mixed
hypercharge anomaly coefficients vanish, and the number of weak doublets
is even.  These are exact algebraic cancellations.  They identify the
correct trivial determinant-line tensor product, but they do not by
themselves construct a local chiral lattice measure on arbitrary gauge
histories.

\subsection{Conventions and the left-handed multiplet}
\label{sec:newS7V77-conventions}

Use the gauge Lie algebra
\[
 \mathfrak g=\mathfrak{su}(3)\oplus\mathfrak{su}(2)\oplus\mathfrak u(1)_Y
\]
and write every fermion as a left-handed Weyl field.  One generation is
\begin{equation}
 \begin{array}{c|ccccc}
 \text{field}&Q&u^c&d^c&L&e^c\\ \hline
 SU(3)&\mathbf3&\overline{\mathbf3}&\overline{\mathbf3}&\mathbf1&\mathbf1\\
 SU(2)&\mathbf2&\mathbf1&\mathbf1&\mathbf2&\mathbf1\\
 Y&\frac16&-\frac23&\frac13&-\frac12&1
 \end{array}.
 \label{eq:newS7V77-spectrum}
\end{equation}
A sterile \(\nu^c\sim(\mathbf1,\mathbf1)_0\), if included, contributes
zero to every coefficient below.

Normalize the quadratic indices by
\[
 T_3(\mathbf3)=T_3(\overline{\mathbf3})={1\over2},
 \qquad T_2(\mathbf2)={1\over2},
\]
and the cubic color index by
\[
 A_3(\mathbf3)=1,
 \qquad A_3(\overline{\mathbf3})=-1.
\]
Dimensions of spectator representations supply multiplicities.

\begin{lemma}[Anomaly trace decomposition]
\label{lem:newS7V77-traces}
For a Weyl multiplet \((R_3,R_2)_y\), its contributions, up to common
normalization constants, are
\begin{align}
 \mathcal A_{333}&=A_3(R_3)\dim R_2,
 \label{eq:newS7V77-a333}\\
 \mathcal A_{33Y}&=yT_3(R_3)\dim R_2,
 \label{eq:newS7V77-a33y}\\
 \mathcal A_{22Y}&=yT_2(R_2)\dim R_3,
 \label{eq:newS7V77-a22y}\\
 \mathcal A_{YYY}&=y^3\dim R_3\dim R_2,
 \label{eq:newS7V77-ayyy}\\
 \mathcal A_{{\rm grav}{\rm grav}Y}&=y\dim R_3\dim R_2.
 \label{eq:newS7V77-aggy}
\end{align}
The perturbative \(SU(2)^3\) coefficient vanishes for every multiplet.
Coefficients with one simple-group generator and otherwise only
hypercharge generators vanish representation by representation.
\end{lemma}

\begin{proof}
The consistent four-dimensional gauge anomaly is proportional to the
symmetrized trace of three representation generators.  Tracing spectator
identity matrices gives the displayed dimensions.  For two simple-group
generators and one hypercharge generator, hypercharge is scalar and the
quadratic trace gives \(yT(R)\).  Three hypercharge insertions give
\(y^3\dim R\), while the mixed gravitational term has a single
hypercharge trace.  The fundamental of \(SU(2)\) is pseudoreal, so its
symmetric cubic invariant is zero; the same holds for every \(SU(2)\)
representation.  Finally, a term with one generator of a simple group is
zero because that generator is traceless.
\end{proof}

\subsection{Exact cancellation of the local coefficients}
\label{sec:newS7V77-local}

\begin{theorem}[One-generation perturbative anomaly cancellation]
\label{thm:newS7V77-local-cancel}
For the spectrum \eqref{eq:newS7V77-spectrum},
\begin{equation}
 \mathcal A_{333}=\mathcal A_{33Y}=\mathcal A_{22Y}
 =\mathcal A_{YYY}=\mathcal A_{{\rm grav}{\rm grav}Y}=0.
 \label{eq:newS7V77-all-zero}
\end{equation}
Together with the identities in \Cref{lem:newS7V77-traces}, every local
cubic gauge anomaly and the mixed gravitational--hypercharge anomaly
vanishes.
\end{theorem}

\begin{proof}
For the pure color coefficient, the weak doublet multiplicity of \(Q\)
gives
\begin{equation}
 \mathcal A_{333}=2A_3(\mathbf3)
 +A_3(\overline{\mathbf3})+A_3(\overline{\mathbf3})
 =2-1-1=0.
 \label{eq:newS7V77-333-calc}
\end{equation}
For the color--hypercharge coefficient,
\begin{equation}
 \mathcal A_{33Y}
 =2\left({1\over6}\right){1\over2}
 +\left(-{2\over3}\right){1\over2}
 +\left({1\over3}\right){1\over2}
 ={1\over6}-{1\over3}+{1\over6}=0.
 \label{eq:newS7V77-33y-calc}
\end{equation}
For the weak--hypercharge coefficient, color supplies three copies of the
quark doublet:
\begin{equation}
 \mathcal A_{22Y}
 =3\left({1\over6}\right){1\over2}
 +\left(-{1\over2}\right){1\over2}
 ={1\over4}-{1\over4}=0.
 \label{eq:newS7V77-22y-calc}
\end{equation}
The cubic hypercharge trace is
\begin{align}
 \mathcal A_{YYY}
 &=6\left({1\over6}\right)^3
 +3\left(-{2\over3}\right)^3
 +3\left({1\over3}\right)^3
 +2\left(-{1\over2}\right)^3+1
 \notag\\
 &={1-32+4-9+36\over36}=0.
 \label{eq:newS7V77-yyy-calc}
\end{align}
The linear hypercharge trace is
\begin{equation}
 \mathcal A_{{\rm grav}{\rm grav}Y}
 =6\left({1\over6}\right)+3\left(-{2\over3}\right)
 +3\left({1\over3}\right)+2\left(-{1\over2}\right)+1
 =1-2+1-1+1=0.
 \label{eq:newS7V77-ggy-calc}
\end{equation}
The remaining local coefficients vanish by
\Cref{lem:newS7V77-traces}.
\end{proof}

\begin{corollary}[Generation multiplicity and sterile neutrinos]
\label{cor:newS7V77-generations}
Any number \(n_g\) of complete generations has the same cancellations.
Adding any number of sterile \((\mathbf1,\mathbf1)_0\) left-handed fields
does not change them.
\end{corollary}

\begin{proof}
Anomaly traces add under direct sums.  Each generation contributes zero,
and a sterile neutral singlet has zero generators.
\end{proof}

\subsection{The global weak mod-two test}
\label{sec:newS7V77-global}

\begin{theorem}[Even weak-doublet parity]
\label{thm:newS7V77-Witten}
One Standard Model generation contains four left-handed fundamental
\(SU(2)\) doublets when color multiplicity is counted: three copies of
\(Q\) and one copy of \(L\).  Hence its mod-two large-gauge-transformation
phase is
\begin{equation}
 (-1)^{3+1}=+1.
 \label{eq:newS7V77-Witten}
\end{equation}
For \(n_g\) generations the number is \(4n_g\), still even.
\end{theorem}

\begin{proof}
A single left-handed fundamental doublet has the mod-two sign associated
with the nontrivial class of \(\pi_4(SU(2))\).  Independent fermion
Pfaffians multiply their signs.  The quark doublet carries a three
-dimensional color spectator space and therefore occurs three times; the
lepton doublet occurs once.  Their product is \((-1)^4=+1\).  Repetition
for \(n_g\) generations preserves even parity.
\end{proof}

\begin{remark}[Scope of the global test]
\label{rem:newS7V77-global-scope}
The calculation proves cancellation of the traditional fundamental
\(SU(2)\) mod-two obstruction.  It is not a bordism classification for
every global form of the Standard Model gauge group, every spin or
spin-charge structure, or every gravitational background.  Those choices
must be fixed before a full global anomaly theorem is stated.
\end{remark}

\subsection{Compatibility with gauge-invariant mass generation}
\label{sec:newS7V77-Yukawa}

Let the Higgs transform as \(H\sim(\mathbf1,\mathbf2)_{1/2}\).

\begin{proposition}[Hypercharge neutrality of the Yukawa vertices]
\label{prop:newS7V77-Yukawa}
In left-handed notation the hypercharge sums of the up, down, and charged
lepton Yukawa monomials are respectively
\begin{align}
 Y(Q)+Y(H)+Y(u^c)&={1\over6}+{1\over2}-{2\over3}=0,
 \label{eq:newS7V77-up}\\
 Y(Q)-Y(H)+Y(d^c)&={1\over6}-{1\over2}+{1\over3}=0,
 \label{eq:newS7V77-down}\\
 Y(L)-Y(H)+Y(e^c)&=-{1\over2}-{1\over2}+1=0.
 \label{eq:newS7V77-lepton}
\end{align}
The corresponding color and weak indices can be contracted to singlets.
Thus electroweak mass generation respects the same gauge representation
bookkeeping used in the anomaly sum.
\end{proposition}

\begin{proof}
The displayed sums prove \(U(1)_Y\) invariance.  The up monomial contracts
the two weak fundamentals with the antisymmetric tensor; the down and
charged-lepton monomials use the conjugate Higgs doublet.  Fundamental and
antifundamental color indices contract by the identity.
\end{proof}

\subsection{Consequences for the determinant line}
\label{sec:newS7V77-line}

\begin{theorem}[Continuum local determinant-line curvature cancels]
\label{thm:newS7V77-curvature}
For a smooth continuum family of four-dimensional chiral Dirac operators
with the representation content \eqref{eq:newS7V77-spectrum}, the sum of
the standard local gauge-anomaly curvature polynomials of the individual
determinant lines is zero.  The mixed local curvature term containing one
hypercharge field strength and two Riemann curvatures is also zero.
\end{theorem}

\begin{proof}
The family-index curvature is linear in the fermion representation and
its gauge part is built from the symmetrized cubic traces listed in
\Cref{lem:newS7V77-traces}; its mixed gravitational--hypercharge part is
built from the linear hypercharge trace.  Every coefficient is zero by
\Cref{thm:newS7V77-local-cancel}.  Additivity of first Chern forms under
tensor products completes the proof.
\end{proof}

\begin{corollary}[Flatness is not a chosen trivialization]
\label{cor:newS7V77-flat-not-trivial}
The tensor-product determinant line has zero local anomaly curvature on
the smooth continuum gauge-family charts covered by
\Cref{thm:newS7V77-curvature}.  This does not select a global phase:
a flat line may retain nontrivial holonomy, and a finite regulator still
must supply compatible transition functions.
\end{corollary}

\begin{proof}
Vanishing curvature is local flatness.  Holonomy of a flat line is a
representation of the fundamental group and need not be trivial.  A
choice of nonzero section is additional data even when curvature
vanishes.
\end{proof}

\begin{definition}[Standard Model anomaly-line card, SMALC]
\label{def:newS7V77-SMALC}
At each cutoff and renewal shell, SMALC requires:
\begin{enumerate}
\item the exact left-handed multiplet list, hypercharge normalization, and
global gauge-group form;
\item regulator-level representatives of all local anomaly traces and a
proof that their tensor sum vanishes before the cutoff is removed;
\item the mod-two weak-doublet transition on every admissible large gauge
loop;
\item transition functions and holonomies of the full determinant and
Pfaffian tensor product, including divisor RCLC charts;
\item compatibility of the chosen phase with Yukawa couplings, zero-mode
saturations, reflection or transfer structure, and gauge averaging; and
\item stability of all preceding rows under refinement and the continuum
limit.
\end{enumerate}
\end{definition}

\begin{firewall}[Representation arithmetic is necessary, not a regulator]
\label{fire:newS7V77-regulator}
The cancellations in \Cref{thm:newS7V77-local-cancel,thm:newS7V77-Witten}
identify the anomaly-free Standard Model tensor product.  They do not
prove that a chosen overlap, domain-wall, renewal, or other finite-cutoff
construction realizes a local gauge-covariant phase on arbitrary gauge
histories.  That construction must populate SMALC rows 2--6.
\end{firewall}

\begin{assessment}[Upstream obstruction after the arithmetic]
\label{ass:newS7V77-status}
No uncancelled conventional Standard Model representation coefficient
forces the V76 local crossing sign or local curvature to survive in the
full fermion product.  The bottleneck has therefore moved upstream to a
common-action statement: construct one finite regulator for all chiral
multiplets, compute its Berry curvature and transition cocycle, and show
that the algebraic cancellations occur with uniform locality on the same
admissible gauge set used by the bosonic collar compiler.
\end{assessment}

\begin{programme}[V78: finite determinant-projector geometry]
\label{prog:newS7V77-V78}
V78 will derive the Berry connection and curvature of a finite-dimensional
occupied-state determinant line directly from its spectral projector.  It
will prove additivity under tensor products and identify the precise trace
identity that a common Standard Model regulator must satisfy.  This goes
upstream of arbitrary-history overlap factorization and converts SMALC
row 2 into a finite matrix estimate.
\end{programme}

\begin{firewall}[Claim boundary after seventh-series V77]
\label{fire:newS7V77-boundary}
V77 proves all conventional one-generation local Standard Model anomaly
coefficient cancellations, the even fundamental \(SU(2)\) doublet parity,
Yukawa hypercharge compatibility, and cancellation of the corresponding
smooth continuum local determinant-line curvature polynomial.

V77 does not classify all global anomalies for a chosen quotient gauge
group, construct a common finite regulator or its phase, prove SMALC,
RCLC, DSLC, CRDC, GAOCC, CCTC, MSRC, RTGC, or RCNC, remove the cutoff, or
construct the Standard Model--Einstein continuum.
\end{firewall}

\begin{theorem}[Seventh-series V77 result]
\label{thm:newS7V77-result}
For each complete Standard Model generation written as left-handed Weyl
multiplets, the \(SU(3)^3\), \(SU(3)^2U(1)\), \(SU(2)^2U(1)\), \(U(1)^3\),
and mixed gravitational--\(U(1)\) anomaly coefficients vanish exactly,
and the weak fundamental-doublet count is four.  Hence the conventional
local anomaly curvature and the fundamental \(SU(2)\) mod-two sign cancel
in the full generation.  A regulator-level phase and global
trivialization remain to be constructed.
\end{theorem}

\begin{proof}
Use \Cref{lem:newS7V77-traces,
thm:newS7V77-local-cancel,
thm:newS7V77-Witten,
prop:newS7V77-Yukawa,
thm:newS7V77-curvature,
cor:newS7V77-flat-not-trivial}.
\end{proof}

\paragraph{Parent-version record.}
V77 was formed from
\texttt{Renewal\_SM\_Einstein\_Continuum\_Research\_Notes\_V76.tex}.
The V76 parent had \(30\,556\) logical source lines, \(1\,207\,033\) bytes,
and SHA--256
\[
 \texttt{0F56EF495F87F4B9800EEB722EAF5F28B041E3F45EEC2209B5BFED4E9AE3E751}.
\]
The next compulsory companion audit remains V80.

% =============================================================================
% END APPENDED SEVENTH-SERIES V77 MATERIAL
% =============================================================================
% =============================================================================
% APPENDED SEVENTH-SERIES V78 MATERIAL
% =============================================================================

\section{Finite occupied-state determinant lines and projector curvature}
\label{sec:newS7V78}
\label{sec:v7v78}

V77 showed that the continuum anomaly polynomial cancels in a complete
Standard Model generation.  A finite chiral regulator, however, presents
a family of occupied subspaces rather than an anomaly polynomial.  Its
phase is the determinant line of those subspaces.  This section derives
its connection and curvature directly from the orthogonal spectral
projector, proves frame covariance and tensor-product additivity, and
gives gap-dependent bounds.  The result identifies a concrete finite
matrix row that a common Standard Model regulator must populate.

\subsection{Connection and curvature from a projector}
\label{sec:newS7V78-projector}

Let \(M\) be a smooth finite-dimensional parameter manifold and let
\(P(x)\) be a smooth family of Hermitian projections on \(\mathbb C^n\)
of constant rank \(r\).  On a coordinate patch choose an orthonormal
frame \(W(x):\mathbb C^r\to\mathbb C^n\), so that
\begin{equation}
 W^*W=I_r,
 \qquad WW^*=P.
 \label{eq:newS7V78-frame}
\end{equation}
The wedge \(w_1\wedge\cdots\wedge w_r\) is a local unit section of
\(\det\operatorname{Ran}P\).

\begin{theorem}[Finite determinant-line connection]
\label{thm:newS7V78-connection}
Define the scalar one-form and two-form
\begin{equation}
 \mathcal A_W:=\operatorname{Tr}(W^*dW),
 \qquad
 \mathcal F_P:=\operatorname{Tr}(P\,dP\wedge dP).
 \label{eq:newS7V78-AF}
\end{equation}
Then \(\mathcal A_W\) is purely imaginary and
\begin{equation}
 d\mathcal A_W=\mathcal F_P.
 \label{eq:newS7V78-curvature}
\end{equation}
If \(W'=Wg\) with \(g:M\to U(r)\), then
\begin{equation}
 \mathcal A_{W'}=\mathcal A_W+operatorname{Tr}(g^{-1}dg),
 \qquad \mathcal F_{P}'=\mathcal F_P.
 \label{eq:newS7V78-gauge}
\end{equation}
Thus \(\mathcal F_P\) is the frame-independent curvature of the occupied
state determinant line.
\end{theorem}

\begin{proof}
Differentiating \(W^*W=I\) gives
\((W^*dW)^*=-W^*dW\), so its trace is imaginary.  Also
\[
 d\mathcal A_W=\operatorname{Tr}(dW^*\wedge dW).
\]
Insert \(I=P+(I-P)\) between the two differentials.  The \(P\) part is
minus \(\operatorname{Tr}((W^*dW)\wedge(W^*dW))\), which is zero because
the trace of a matrix-valued one-form wedged with itself is the trace of a
commutator.  Hence
\[
 d\mathcal A_W
 =\operatorname{Tr}(dW^*(I-P)\wedge dW).
\]
From \(P=WW^*\), one has
\((I-P)dP\,W=(I-P)dW\) and
\(W^*dP(I-P)=dW^*(I-P)\).  Cyclicity of the trace now gives
\(d\mathcal A_W=\operatorname{Tr}(P,dP\wedge dP)\).
For \(W'=Wg\), direct differentiation yields
\[
 (Wg)^*d(Wg)=g^{-1}(W^*dW)g+g^{-1}dg.
\]
Taking the trace proves the first transformation law; the curvature is
unchanged because \(P\) is unchanged.
\end{proof}

\begin{corollary}[Closedness and integral periods]
\label{cor:newS7V78-periods}
The form \(\mathcal F_P\) is closed.  For every closed oriented surface
\(\Sigma\subset M\),
\begin{equation}
 {1\over2\pi i}\int_\Sigma\mathcal F_P\in\mathbb Z.
 \label{eq:newS7V78-Chern}
\end{equation}
A global nowhere-zero unit section of the determinant line can exist only
if all these integers vanish.
\end{corollary}

\begin{proof}
Locally \(\mathcal F_P=d\mathcal A_W\), so it is closed.  On overlaps the
connections differ by \(d\log\det g\).  Applying Stokes' theorem to a
patch decomposition of \(\Sigma\), the interior terms cancel and the
remaining boundary windings of \(\det g:\partial U\to U(1)\) are integers.
A global unit section gives a global connection potential and therefore
zero periods.
\end{proof}

\subsection{Tensor products, direct sums, and complements}
\label{sec:newS7V78-additivity}

\begin{theorem}[Additivity of finite anomaly curvature]
\label{thm:newS7V78-additivity}
For projector families \(P_1\) and \(P_2\), let
\(P=P_1\oplus P_2\).  Then
\begin{equation}
 \det\operatorname{Ran}P
 \simeq
 \det\operatorname{Ran}P_1\otimes
 \det\operatorname{Ran}P_2,
 \qquad
 \mathcal F_P=\mathcal F_{P_1}+\mathcal F_{P_2}.
 \label{eq:newS7V78-additivity}
\end{equation}
For \(P^\perp=I-P\),
\begin{equation}
 \mathcal F_{P^\perp}=-\mathcal F_P.
 \label{eq:newS7V78-complement}
\end{equation}
\end{theorem}

\begin{proof}
Concatenate local frames for the direct sum.  The determinant wedge is the
ordered product of the two wedges and the block-diagonal connection trace
is the sum, proving \eqref{eq:newS7V78-additivity}.  Since
\(dP^\perp=-dP\),
\[
 \mathcal F_{P^\perp}=\operatorname{Tr}((I-P)dP\wedge dP).
\]
But \(\operatorname{Tr}(dP\wedge dP)=0\), again by the trace-commutator
identity.  Subtraction gives \eqref{eq:newS7V78-complement}.
\end{proof}

\begin{corollary}[Vectorlike cancellation]
\label{cor:newS7V78-vectorlike}
If two regulated chiral sectors use complementary occupied projectors on
the same finite one-particle space, their determinant-line curvature
cancels exactly at finite cutoff.
\end{corollary}

\begin{proof}
Apply \eqref{eq:newS7V78-complement} and tensor-product additivity.
\end{proof}

\begin{firewall}[Standard Model cancellation is not complementary pairing]
\label{fire:newS7V78-not-vectorlike}
The Standard Model spectrum is chiral.  Its V77 cancellation occurs among
different representations and hypercharges, not by pairing every
projector with its complement.  Therefore
\Cref{cor:newS7V78-vectorlike} does not prove finite-cutoff Standard Model
curvature cancellation.
\end{firewall}

\subsection{Spectral projectors and gap bounds}
\label{sec:newS7V78-gap}

Let \(H(x)=H(x)^*\) be a smooth finite Hermitian matrix family with no
zero eigenvalue, and set
\begin{equation}
 P_-(x)=\mathbf1_{(-\infty,0)}(H(x)).
 \label{eq:newS7V78-negative}
\end{equation}

\begin{lemma}[Resolvent derivative of the occupied projector]
\label{lem:newS7V78-resolvent}
Let \(\Gamma\) be a positively oriented contour enclosing the negative
spectrum and excluding the nonnegative spectrum.  For every tangent
variation \(X\),
\begin{equation}
 d_XP_-={1\over2\pi i}\oint_\Gamma
 (z-H)^{-1}(d_XH)(z-H)^{-1}\,dz.
 \label{eq:newS7V78-resolvent}
\end{equation}
If \(\operatorname{dist}(\Gamma,\operatorname{spec}H)\geq\eta>0\) and
\(\Gamma\) has length \(L_\Gamma\), then
\begin{equation}
 \|d_XP_-\|
 \leq {L_\Gamma\over2\pi\eta^2}\|d_XH\|.
 \label{eq:newS7V78-dP-bound}
\end{equation}
\end{lemma}

\begin{proof}
The Riesz formula is
\(P_-=(2\pi i)^{-1}\oint_\Gamma(z-H)^{-1}dz\).  Differentiate it and use
\(d(z-H)^{-1}=(z-H)^{-1}(dH)(z-H)^{-1}\).  The resolvent norm is at most
\(\eta^{-1}\); integrating its product bound proves
\eqref{eq:newS7V78-dP-bound}.
\end{proof}

\begin{theorem}[Finite curvature bound]
\label{thm:newS7V78-curvature-bound}
For tangent vectors \(X,Y\),
\begin{equation}
 |\mathcal F_{P_-}(X,Y)|
 \leq 2r\,\|d_XP_-\|\,\|d_YP_-\|
 \leq {rL_\Gamma^2\over2\pi^2\eta^4}
 \|d_XH\|\,\|d_YH\|,
 \label{eq:newS7V78-curvature-bound}
\end{equation}
where \(r=\operatorname{rank}P_-\).
\end{theorem}

\begin{proof}
Evaluation of the wedge gives
\[
 \mathcal F_{P_-}(X,Y)
 =\operatorname{Tr}P_-\bigl[(d_XP_-)(d_YP_-)
 -(d_YP_-)(d_XP_-)\bigr].
\]
The absolute trace of either term is at most
\(r\|d_XP_-\|\|d_YP_-\|\).  Sum the two bounds and use
\eqref{eq:newS7V78-dP-bound}.
\end{proof}

\begin{corollary}[Why the Wilson-kernel gap matters]
\label{cor:newS7V78-Wilson-gap}
On every admissible gauge set with a uniform Hermitian Wilson-kernel gap,
the overlap occupied projector and its Berry curvature are smooth and
finite.  Approaching a kernel-gap closing can make the resolvent estimate
and the curvature bound diverge.
\end{corollary}

\begin{proof}
A uniform gap permits a contour with uniform positive \(\eta\).
\Cref{lem:newS7V78-resolvent,thm:newS7V78-curvature-bound} then apply.
No such contour persists at a gap closing.
\end{proof}

\subsection{Gauge tangents and the common-regulator trace row}
\label{sec:newS7V78-gauge}

Suppose a gauge transformation acts through a unitary representation
\(U(g)\) and
\begin{equation}
 P_-(A^g)=U(g)P_-(A)U(g)^*.
 \label{eq:newS7V78-covariance}
\end{equation}
For an infinitesimal generator \(T_X\), the vertical tangent is
\(\delta_XP_-=[T_X,P_-]\).

\begin{proposition}[Vertical projector cocycle]
\label{prop:newS7V78-vertical}
On two infinitesimal gauge tangents \(X,Y\), the determinant-line
curvature is the finite trace
\begin{equation}
 \mathcal F_{P_-}(X,Y)
 =\operatorname{Tr}P_-
 \left([T_X,P_-][T_Y,P_-]-[T_Y,P_-][T_X,P_-]\right).
 \label{eq:newS7V78-vertical}
\end{equation}
It is additive over all left-handed Standard Model multiplets.
\end{proposition}

\begin{proof}
Differentiate \eqref{eq:newS7V78-covariance} along the gauge orbit and
insert the resulting commutators into \eqref{eq:newS7V78-AF}.  Additivity
is \Cref{thm:newS7V78-additivity}.
\end{proof}

\begin{definition}[Projector-curvature compiler, PCC]
\label{def:newS7V78-PCC}
A common finite Standard Model chiral regulator satisfies PCC on an
admissible gauge chart if:
\begin{enumerate}
\item all multiplet kernels are defined on the same lattice, gauge field,
metric or frame data, boundary convention, and admissible set;
\item every occupied projector has a uniform spectral gap and a localized
resolvent derivative;
\item the summed finite curvature
\begin{equation}
 \mathcal F_{\rm SM}:=\sum_f
 \operatorname{Tr}(P_f\,dP_f\wedge dP_f)
 \label{eq:newS7V78-FSM}
\end{equation}
is either zero or the exterior derivative of an explicitly constructed
local measure-current counterterm;
\item the counterterm is gauge compatible and has uniform collar
oscillation and refinement bounds;
\item all two-cycle periods and residual flat holonomies are computed;
and
\item the resulting local phase agrees with RCLC at rank divisors, CCTC
for transfer structure, and the fermionic stress row used by RCNC.
\end{enumerate}
\end{definition}

\begin{theorem}[Exact finite target supplied by PCC]
\label{thm:newS7V78-target}
If PCC rows 1--5 hold on a gauge-family region and the residual flat
holonomies are trivial, then the tensor-product occupied-state determinant
line admits a global unit section on that region.  Its connection differs
from the sum of the frame Berry connections by the declared local
measure-current counterterm.
\end{theorem}

\begin{proof}
Row 3 cancels the curvature after subtracting the counterterm.  Row 5
removes its integral Chern periods and the holonomy of the remaining flat
line.  Parallel transport from a base point is then path independent and
defines a global unit section.  The transformation law
\eqref{eq:newS7V78-gauge} shows that this section has the stated corrected
connection.
\end{proof}

\begin{assessment}[The finite obstruction is now explicit]
\label{ass:newS7V78-status}
V77 cancels the continuum representation traces.  V78 shows what must
cancel before a continuum limit is taken: the finite projector two-form
\eqref{eq:newS7V78-FSM}, possibly after an explicit local measure current,
together with its periods and flat holonomy.  A kernel gap supplies
smoothness and quantitative bounds, but it does not force the Standard
Model sum to vanish.
\end{assessment}

\begin{programme}[V79: local primitive on a good gauge chart]
\label{prog:newS7V78-V79}
V79 will go upstream to the local cohomological step.  On a star-shaped
good gauge chart it will construct a primitive of a closed projector
curvature by the radial homotopy formula, prove quantitative bounds, and
track gauge covariance.  This will separate the locally solvable phase
problem from overlap periods and global holonomy before the V80 companion
audit.
\end{programme}

\begin{firewall}[Claim boundary after seventh-series V78]
\label{fire:newS7V78-boundary}
V78 proves the finite determinant-line connection and projector-curvature
formula, frame covariance, Chern-period condition, direct-sum and
complement additivity, resolvent derivative, spectral-gap curvature bound,
and vertical finite trace cocycle.  It defines the exact PCC target.

V78 does not construct the Standard Model measure current, prove finite
curvature cancellation or locality on arbitrary gauge histories, compute
all periods or holonomies, prove PCC, SMALC, RCLC, DSLC, CRDC, GAOCC,
CCTC, MSRC, RTGC, or RCNC, remove the cutoff, or construct the Standard
Model--Einstein continuum.
\end{firewall}

\begin{theorem}[Seventh-series V78 result]
\label{thm:newS7V78-result}
For every smooth constant-rank finite occupied projector \(P\), the
curvature of its determinant line is
\(\mathcal F_P=\operatorname{Tr}(P\,dP\wedge dP)\).  It is frame
independent and additive under direct sums.  For a gapped spectral
projector it has the explicit resolvent and norm bounds
\eqref{eq:newS7V78-resolvent}--\eqref{eq:newS7V78-curvature-bound}.  Thus a
common Standard Model regulator must cancel or locally trivialize the
finite sum \eqref{eq:newS7V78-FSM}, then separately remove its periods and
flat holonomy.
\end{theorem}

\begin{proof}
Combine \Cref{thm:newS7V78-connection,
cor:newS7V78-periods,
thm:newS7V78-additivity,
lem:newS7V78-resolvent,
thm:newS7V78-curvature-bound,
prop:newS7V78-vertical,
thm:newS7V78-target}.
\end{proof}

\paragraph{Parent-version record.}
V78 was formed from
\texttt{Renewal\_SM\_Einstein\_Continuum\_Research\_Notes\_V77.tex}.
The V77 parent had \(30\,914\) logical source lines, \(1\,221\,264\) bytes,
and SHA--256
\[
 \texttt{B9D560A03BA0E1E3EB6B7CB2071B4066F927DBE8CF5B1959C9DC4C4933AEA6AD}.
\]
The next compulsory companion audit remains V80.

% =============================================================================
% END APPENDED SEVENTH-SERIES V78 MATERIAL
% =============================================================================
% =============================================================================
% APPENDED SEVENTH-SERIES V79 MATERIAL
% =============================================================================

\section{Radial local primitives on good gauge charts}
\label{sec:newS7V79}
\label{sec:v7v79}

V78 expressed the finite chiral phase obstruction as the closed projector
two-form \(\mathcal F_{\rm SM}\).  On a contractible chart, closedness has
an explicit solution: integrate the curvature along radial segments from
a base configuration.  This section proves the homotopy formula, derives
volume-local bounds from exponentially summable curvature components, and
isolates the remaining overlap and holonomy problem.  The construction is
conditional on a single star-shaped admissible chart; it does not assert
that the full gauge orbit space admits such a chart.

\subsection{The finite-dimensional radial homotopy}
\label{sec:newS7V79-homotopy}

Let \(V\) be a finite-dimensional real normed space and let
\(\mathcal U\subset V\) be open and star-shaped about \(x_0\).  Put
\(q=x-x_0\) and \(H_t(x)=x_0+tq\).

\begin{definition}[Radial homotopy operator]
\label{def:newS7V79-K}
For a smooth \(k\)-form \(\omega\), \(k\geq1\), define
\begin{equation}
 (K\omega)_x(v_1,\ldots,v_{k-1})
 :=\int_0^1 t^{k-1}
 \omega_{H_t(x)}(q,v_1,\ldots,v_{k-1})\,dt.
 \label{eq:newS7V79-K}
\end{equation}
\end{definition}

\begin{theorem}[Radial Poincare formula]
\label{thm:newS7V79-Poincare}
For every smooth form of positive degree,
\begin{equation}
 dK\omega+Kd\omega=\omega.
 \label{eq:newS7V79-homotopy}
\end{equation}
Consequently, if a smooth two-form \(\mathcal F\) is closed, then
\begin{equation}
 \mathcal J:=K\mathcal F,
 \qquad
 \mathcal J_x(v)=\int_0^1t\,
 \mathcal F_{x_0+t(x-x_0)}(x-x_0,v)\,dt
 \label{eq:newS7V79-current}
\end{equation}
satisfies \(d\mathcal J=\mathcal F\).
\end{theorem}

\begin{proof}
Consider the map \(H:[0,1]\times\mathcal U\to\mathcal U\),
\(H(t,x)=H_t(x)\).  Cartan's formula on the product gives
\[
 {d\over dt}H_t^*\omega
 =d\,H_t^*(\iota_{\partial_tH}\omega)
 +H_t^*(\iota_{\partial_tH}d\omega).
\]
Integrate from zero to one.  The pullback by \(H_1\) is the identity, and
the pullback of a positive-degree form by the constant map \(H_0\) is
zero.  Evaluating the contraction shows that the first argument contributes
\(q\), while each of the remaining \(k-1\) tangent vectors contributes a
factor \(t\).  The integrated contraction is therefore exactly
\eqref{eq:newS7V79-K}, proving \eqref{eq:newS7V79-homotopy}.  If
\(d\mathcal F=0\), the second term vanishes and
\(dK\mathcal F=\mathcal F\).
\end{proof}

\begin{corollary}[A flat corrected connection]
\label{cor:newS7V79-flat}
Let \(\mathcal A_W\) be a local determinant-line Berry connection with
\(d\mathcal A_W=\mathcal F\).  Then
\begin{equation}
 \mathcal A_{\rm flat}:=\mathcal A_W-\mathcal J
 \label{eq:newS7V79-flat}
\end{equation}
is closed on \(\mathcal U\).  Since \(\mathcal U\) is star-shaped, there
is a scalar \(\phi\), unique after fixing \(\phi(x_0)\), with
\(\mathcal A_{\rm flat}=d\phi\).  Multiplication of the local determinant
section by \(e^{-\phi}\) produces a section whose connection is
\(\mathcal J\).
\end{corollary}

\begin{proof}
The curvature of \eqref{eq:newS7V79-flat} is
\(\mathcal F-d\mathcal J=0\).  Apply
\Cref{thm:newS7V79-Poincare} to the closed one-form.  The usual connection
transformation under scalar multiplication gives the last assertion.
\end{proof}

\subsection{Quantitative chart bounds}
\label{sec:newS7V79-bounds}

\begin{proposition}[Uniform norm bound for the radial current]
\label{prop:newS7V79-norm}
Suppose \(\|\mathcal F_x\|\leq M\) throughout the radial hull of
\(\mathcal U\).  Then
\begin{equation}
 |\mathcal J_x(v)|
 \leq {M\over2}\|x-x_0\|\,\|v\|.
 \label{eq:newS7V79-norm}
\end{equation}
If the projector resolvent data satisfy the V78 bound with common
constants, one may take \(M\) to be the corresponding summed curvature
bound.
\end{proposition}

\begin{proof}
Use \eqref{eq:newS7V79-current}, the two-form norm, and
\(\int_0^1t\,dt=1/2\).  The second assertion is
\Cref{thm:newS7V78-curvature-bound} summed over multiplets.
\end{proof}

A global Euclidean norm makes \eqref{eq:newS7V79-norm} grow with volume.
The relevant estimate instead resolves the curvature into cell
components.  Let \(\Lambda\) be a finite cell set with distance \(d\), and
write \(x=(x_z)_{z\in\Lambda}\).  Assume tangent vectors have the same
block decomposition.

\begin{theorem}[Volume-local radial-current bound]
\label{thm:newS7V79-local}
Assume that for every point on the radial hull,
\begin{equation}
 |\mathcal F_x(u_y,v_z)|
 \leq C_{\mathcal F}e^{-\mu d(y,z)}
 \|u_y\|\,\|v_z\|
 \label{eq:newS7V79-kernel}
\end{equation}
for cell-supported tangents, and that
\begin{equation}
 R_0:=\sup_{x\in\mathcal U}\sup_{y\in\Lambda}
 \|x_y-(x_0)_y\|<\infty,
 \qquad
 S_\mu:=\sup_{z\in\Lambda}\sum_{y\in\Lambda}
 e^{-\mu d(y,z)}<\infty.
 \label{eq:newS7V79-row-sum}
\end{equation}
If \(v\) is supported in \(X\subset\Lambda\), then
\begin{equation}
 |\mathcal J_x(v)|
 \leq {1\over2}C_{\mathcal F}R_0S_\mu
 |X|\,\|v\|_{\infty}.
 \label{eq:newS7V79-local-current}
\end{equation}
The constant is independent of \(|\Lambda|\) whenever
\(C_{\mathcal F},R_0,S_\mu\) are.
\end{theorem}

\begin{proof}
By bilinearity, expand the two-form in
\eqref{eq:newS7V79-current} over the cell components of
\(q=x-x_0\) and \(v\).  Since \(v_z=0\) off \(X\),
\begin{align*}
 |\mathcal J_x(v)|
 &\leq\int_0^1t\,dt
 \sum_{z\in X}\sum_{y\in\Lambda}
 C_{\mathcal F}e^{-\mu d(y,z)}
 \|q_y\|\,\|v_z\|\\
 &\leq {1\over2}C_{\mathcal F}R_0S_\mu
 |X|\,\|v\|_\infty.
\end{align*}
This is \eqref{eq:newS7V79-local-current}.
\end{proof}

\begin{corollary}[Local phase response]
\label{cor:newS7V79-phase-response}
Under \Cref{thm:newS7V79-local}, the infinitesimal phase response to a
variation supported on \(X\) is bounded by a constant times \(|X|\), not
the total volume.  Thus the radial primitive is compatible with a collar
oscillation compiler at the level of first variations.
\end{corollary}

\begin{proof}
The measure-current response is \(\mathcal J_x(v)\); apply
\eqref{eq:newS7V79-local-current}.
\end{proof}

\begin{firewall}[First-variation locality is not full interaction locality]
\label{fire:newS7V79-first-only}
The bound \eqref{eq:newS7V79-local-current} controls the response one-form.
It does not yet express the integrated phase as a uniformly summable
finite-range interaction.  That stronger conclusion needs corresponding
bounds on derivatives of the curvature kernel and on overlap transition
functions.
\end{firewall}

\subsection{Equivariance of the homotopy}
\label{sec:newS7V79-equivariance}

\begin{theorem}[Equivariant radial primitive]
\label{thm:newS7V79-equivariant}
Let a group \(G\) act linearly on \(V\), preserve \(\mathcal U\), and fix
\(x_0\).  If \(g^*\mathcal F=\mathcal F\) for every \(g\in G\), then
\begin{equation}
 g^*\mathcal J=\mathcal J.
 \label{eq:newS7V79-equivariant}
\end{equation}
If the curvature is only covariant after adding an exact two-form, the
radial primitive acquires the corresponding exact one-form and must be
corrected on overlaps.
\end{theorem}

\begin{proof}
Linearity and \(gx_0=x_0\) imply
\(gH_t(x)=H_t(gx)\).  Pullback therefore commutes with the radial homotopy
operator: \(g^*K=Kg^*\).  Applying this identity to the invariant
\(\mathcal F\) gives \eqref{eq:newS7V79-equivariant}.  The final statement
follows by applying the homotopy identity to the exact covariance defect.
\end{proof}

\begin{remark}[Gauge charts]
\label{rem:newS7V79-gauge-chart}
A generic local gauge action on link variables is not a linear action
fixing an arbitrary chart center.  \Cref{thm:newS7V79-equivariant} applies
directly only after an equivariant coordinate choice with an invariant
center, or on a gauge-fixed slice with a separately proved reconstruction
law.  This is a genuine PCC row, not a harmless coordinate convention.
\end{remark}

\subsection{Overlap cocycles and the surviving global problem}
\label{sec:newS7V79-overlaps}

\begin{theorem}[Difference of chart primitives]
\label{thm:newS7V79-overlap}
Let \(\mathcal U_a,\mathcal U_b\) be star-shaped charts carrying radial
primitives \(\mathcal J_a,\mathcal J_b\) of the same closed curvature
\(\mathcal F\).  On their overlap,
\begin{equation}
 d(\mathcal J_a-\mathcal J_b)=0.
 \label{eq:newS7V79-overlap-closed}
\end{equation}
On every simply connected overlap component there is a phase function
\(\chi_{ab}\) with
\begin{equation}
 \mathcal J_a-\mathcal J_b=d\chi_{ab}.
 \label{eq:newS7V79-overlap-phase}
\end{equation}
On triple overlaps the sum
\(\chi_{ab}+\chi_{bc}+\chi_{ca}\) is locally constant.  These constants
and the periods on nonsimply connected overlaps are precisely the gluing
data not fixed by the local radial construction.
\end{theorem}

\begin{proof}
Both primitives differentiate to \(\mathcal F\), proving
\eqref{eq:newS7V79-overlap-closed}.  The Poincare lemma for one-forms on a
simply connected component gives \eqref{eq:newS7V79-overlap-phase}.
Taking the differential of the triple sum gives zero, so it is locally
constant.  A closed one-form on a nonsimply connected overlap is exact
exactly when all of its loop periods vanish.
\end{proof}

\begin{definition}[Good-chart measure-current card, GCMCC]
\label{def:newS7V79-GCMCC}
A finite chiral gauge-family cover satisfies GCMCC if:
\begin{enumerate}
\item every chart is radially contractible inside the common admissible
kernel-gap region;
\item the summed Standard Model projector curvature is closed and obeys
the uniform exponential component bound
\eqref{eq:newS7V79-kernel};
\item each radial current obeys the local estimate
\eqref{eq:newS7V79-local-current} and the required derivative analogues;
\item gauge covariance is proved either by an equivariant contraction or
by explicit slice reconstruction;
\item all overlap phases, triple constants, loop periods, and divisor
RCLC transitions are compatible; and
\item the resulting phase current has uniform renewal-refinement,
transfer, stress, and collar-oscillation bounds.
\end{enumerate}
\end{definition}

\begin{assessment}[Local cohomology is no longer the bottleneck]
\label{ass:newS7V79-status}
On a star-shaped kernel-gap chart, the closed projector curvature always
has the explicit primitive \eqref{eq:newS7V79-current}.  Exponential
component decay converts it into the volume-local response bound
\eqref{eq:newS7V79-local-current}.  The unresolved work is now sharply
located in uniform curvature locality, equivariant chart construction,
and compatibility of overlap cocycles and rank-divisor transitions.
\end{assessment}

\begin{programme}[V80: companion audit and coupled chart compiler]
\label{prog:newS7V79-V80}
V80 will perform the compulsory YM and NS audit before adding new
mathematics.  It will compare the newest YM generated-quotient and collar
stocks with GCMCC, then formulate the combined local activity bound for a
bosonic modulus and a chiral chart phase.  Any missing uniform curvature
kernel estimate will remain an explicit debit rather than being inferred
from continuum anomaly arithmetic.
\end{programme}

\begin{firewall}[Claim boundary after seventh-series V79]
\label{fire:newS7V79-boundary}
V79 proves the radial Poincare formula, constructs a local primitive of the
finite projector curvature, proves norm and exponentially localized
response bounds, establishes equivariance under a fixed-center linear
action, and derives the overlap cocycle.

V79 does not prove a uniform Standard Model curvature-kernel estimate,
construct a globally equivariant gauge cover, eliminate overlap periods
or holonomy, prove GCMCC, PCC, SMALC, RCLC, DSLC, CRDC, GAOCC, CCTC,
MSRC, RTGC, or RCNC, remove the cutoff, or construct the Standard
Model--Einstein continuum.
\end{firewall}

\begin{theorem}[Seventh-series V79 result]
\label{thm:newS7V79-result}
On every star-shaped finite gauge chart, the closed Standard Model
projector curvature has the explicit radial primitive
\eqref{eq:newS7V79-current}.  If its cell components decay exponentially
with a uniform row sum, the primitive responds to an \(X\)-supported
variation by at most
\(C_{\mathcal F}R_0S_\mu|X|\|v\|_\infty/2\).  Local construction therefore
costs no total-volume factor.  Global phase existence still requires
compatible equivariant overlaps, vanishing periods, and divisor gluing.
\end{theorem}

\begin{proof}
Combine \Cref{thm:newS7V79-Poincare,
cor:newS7V79-flat,
prop:newS7V79-norm,
thm:newS7V79-local,
thm:newS7V79-equivariant,
thm:newS7V79-overlap}.
\end{proof}

\paragraph{Parent-version record.}
V79 was formed from
\texttt{Renewal\_SM\_Einstein\_Continuum\_Research\_Notes\_V78.tex}.
The V78 parent had \(31\,293\) logical source lines, \(1\,235\,062\) bytes,
and SHA--256
\[
 \texttt{443A6A66A9612A039BE3396F13A0497ADBCAADED1F1784AAF7FA1052CCCF4ECF}.
\]
The next compulsory companion audit is V80.

% =============================================================================
% END APPENDED SEVENTH-SERIES V79 MATERIAL
% =============================================================================
% =============================================================================
% APPENDED SEVENTH-SERIES V80 MATERIAL
% =============================================================================

\section{Companion audit and a coupled modulus--phase collar compiler}
\label{sec:newS7V80}
\label{sec:v7v80}

This is the compulsory V80 companion audit.  The newest Yang--Mills note
has advanced: it proves an abstract weighted interaction compiler and,
more significantly here, differentiates an eliminated-shell logarithmic
marginal exactly.  Its connected covariance term can destroy convexity.
The analogous warning for chiral fermions is that a local determinant
modulus and a local Berry current do not make the normalized complex
measure positive or even keep its partition function away from zero.  This
section proves the exact conditional phase margin needed to transfer the
existing bad-cluster estimates to complex amplitudes and records the
retained-carrier alternative.

\subsection{Compulsory V80 companion audit}
\label{sec:newS7V80-audit}

At the audit snapshot, the newest distinct complete Yang--Mills checkpoint
was
\begin{align}
 &\texttt{Renewal\_Yang\_Mills\_Mass\_Gap\_Research\_Notes\_V45.tex},
 \notag\\
 &\qquad 704\,056\ {\rm bytes},\quad 20\,173\ {\rm logical\ lines},\quad
 \texttt{793583EFCD7EBDD7E96D77D5A9E7C6CE78CECDF693BA6196A7C60D5B3AF1E6AC}.
 \label{eq:newS7V80-YM-hash}
\end{align}
It has one live document terminator.  The newest Navier--Stokes checkpoint
remained
\begin{align}
 &\texttt{Renewal\_NS\_Stability\_Research\_Notes\_V26.tex},
 \notag\\
 &\qquad 278\,283\ {\rm bytes},\quad 5\,319\ {\rm logical\ lines},\quad
 \texttt{82C887F8C2C69E4F28FAD7C3DC8DE5B9099CB5A4BF431EBA1FFF3E91935978AD}.
 \label{eq:newS7V80-NS-hash}
\end{align}
It also has one live document terminator.

\paragraph{Yang--Mills V45.}
The note proves that an exponentially weighted support-local interaction
norm controls its negative Hessian, mixed quotient response, and collar
oscillation stocks, with exponentially small truncation tails.  It then
differentiates the exact eliminated-detail marginal and obtains a direct
second derivative minus a connected conditional covariance.  It offers
two honest routes: prove the required weighted interaction norm, or prove
a same-law conditional detail-mixing estimate for the score covariances.
If either route becomes circular, it directs the programme upstream to a
retained-detail positive transfer dilation.  The physical Yang--Mills
pushforward has not yet populated either route.

\paragraph{Navier--Stokes V26.}
The note is unchanged.  Its rank-one norm improvement applies only after
an actual invariant rank-one input is proved; it gives no chiral phase or
complex normalization estimate.

\begin{theorem}[V80 audit decision]
\label{thm:newS7V80-audit}
The YM V45 differentiation identity applies structurally to the
fermionic logarithmic marginal: connected fluctuations enter after the
carrier is eliminated and have no positivity sign in a complex measure.
Accordingly, GCMCC first-variation locality will not be promoted to a
probability theorem.  The active programme must either prove a uniform
conditional phase margin for the normalized complex marginal or retain a
positive fermionic transfer carrier and establish projectivity before
marginalization.  NS V26 supplies no additional estimate.
\end{theorem}

\begin{proof}
The Yang--Mills identity is an algebraic differentiation of a normalized
integral, so the same calculation holds with a complex action whenever
the partition function is nonzero.  Positivity was essential only for the
sign of a real covariance and for probability comparison.  A chiral phase
removes that sign.  The two stated routes are precisely the alternatives
that preserve, respectively, a controlled normalization or an upstream
positive carrier.  The NS conclusion follows from its unchanged scope.
\end{proof}

\subsection{Absolute law and phase margin}
\label{sec:newS7V80-margin}

Let \(\nu\) be a probability measure representing the already controlled
bosonic action times the fermion determinant modulus in one finite
conditional cylinder.  Let \(\Theta\) be a measurable real phase and set
\begin{equation}
 z_\nu:=\mathbb E_\nu e^{i\Theta},
 \qquad
 \rho_\nu:=|z_\nu|,
 \qquad
 \mathbb E_{\mathbb C}F
 :={\mathbb E_\nu(Fe^{i\Theta})\over z_\nu}
 \label{eq:newS7V80-complex-expectation}
\end{equation}
when \(z_\nu\ne0\).  The number \(\rho_\nu\in[0,1]\) is the phase margin.

\begin{theorem}[Exact phase-margin comparison]
\label{thm:newS7V80-margin}
For every bounded measurable \(F\),
\begin{equation}
 |\mathbb E_{\mathbb C}F|
 \leq {\mathbb E_\nu|F|\over\rho_\nu}.
 \label{eq:newS7V80-margin-bound}
\end{equation}
In particular, for every event \(A\),
\begin{equation}
 \left|{\mathbb E_\nu(\mathbf1_Ae^{i\Theta})\over z_\nu}\right|
 \leq {\nu(A)\over\rho_\nu}.
 \label{eq:newS7V80-event}
\end{equation}
Both constants are sharp.
\end{theorem}

\begin{proof}
The triangle inequality gives
\(|\mathbb E_\nu(Fe^{i\Theta})|\leq\mathbb E_\nu|F|\), because the phase
has unit modulus.  Divide by \(|z_\nu|\).  Taking \(F\) constant with
phase aligned to \(z_\nu\) gives equality in
\eqref{eq:newS7V80-margin-bound}; indicator examples supported on a
constant-phase set show that the event coefficient cannot be improved
without more information.
\end{proof}

\begin{countertheorem}[Local phase regularity does not prevent a zero]
\label{cth:newS7V80-zero}
A bounded phase with finite oscillation can have zero phase margin.  Hence
no finite GCMCC response bound alone guarantees that
\eqref{eq:newS7V80-complex-expectation} exists.
\end{countertheorem}

\begin{proof}
Let \(\nu\) give probability \(1/2\) to each of two points and take
\(\Theta=0\) on one point and \(\Theta=\pi\) on the other.  Then
\(\mathbb E_\nu e^{i\Theta}=(1-1)/2=0\), while \(\Theta\) is bounded and
has finite oscillation \(\pi\).
\end{proof}

\begin{lemma}[Arc sufficient condition]
\label{lem:newS7V80-arc}
If there is a real \(\theta_0\) and \(0\leq\eta<\pi/2\) such that
\begin{equation}
 |\Theta-\theta_0|\leq\eta
 \quad\nu\text{-almost surely},
 \label{eq:newS7V80-arc}
\end{equation}
then
\begin{equation}
 \rho_\nu\geq\cos\eta>0.
 \label{eq:newS7V80-arc-margin}
\end{equation}
\end{lemma}

\begin{proof}
The real part of
\(e^{-i\theta_0}\mathbb E_\nu e^{i\Theta}\) is
\(\mathbb E_\nu\cos(\Theta-\theta_0)\geq\cos\eta\).  The modulus is at
least this real part.
\end{proof}

\begin{lemma}[Variance sufficient condition]
\label{lem:newS7V80-variance}
If \(\Theta\in L^2(\nu)\), then with
\(\overline\Theta=\mathbb E_\nu\Theta\),
\begin{equation}
 \rho_\nu
 \geq 1-{1\over2}\operatorname{Var}_\nu(\Theta).
 \label{eq:newS7V80-variance}
\end{equation}
Thus variance strictly below two is a sufficient, though generally
nonuniform, phase-margin condition.
\end{lemma}

\begin{proof}
Rotate by \(e^{-i\overline\Theta}\).  Since
\(\cos u\geq1-u^2/2\) for real \(u\),
\[
 \rho_\nu\geq
 \operatorname{Re}\mathbb E_\nu e^{i(\Theta-\overline\Theta)}
 \geq1-{1\over2}\mathbb E_\nu
 (\Theta-\overline\Theta)^2.
\]
\end{proof}

\begin{firewall}[An extensive phase defeats the elementary criteria]
\label{fire:newS7V80-extensive}
If \(\Theta\) is a sum of order-volume local terms, its total oscillation
and variance may grow with volume even when every local derivative is
bounded.  Neither \Cref{lem:newS7V80-arc} nor
\Cref{lem:newS7V80-variance} then supplies a uniform phase margin.  A
cluster expansion, symmetry, exact pairing, or retained-carrier argument
is needed.
\end{firewall}

\subsection{Complex defect-path amplitudes}
\label{sec:newS7V80-defects}

\begin{theorem}[Transfer of an absolute defect-path bound]
\label{thm:newS7V80-path}
Suppose that, uniformly in admissible boundary conditions, the absolute
conditional law obeys
\begin{equation}
 \nu(\Gamma\text{ is bad}\mid\mathcal F_\partial)
 \leq p^{|\Gamma|}
 \label{eq:newS7V80-absolute-path}
\end{equation}
for every selected path \(\Gamma\), and its conditional phase margin obeys
\begin{equation}
 \rho_\nu(\mathcal F_\partial)\geq\rho_0>0.
 \label{eq:newS7V80-uniform-margin}
\end{equation}
Then the normalized complex conditional amplitude satisfies
\begin{equation}
 \left|
 {\mathbb E_\nu[\mathbf1_{\{\Gamma\ {
m bad}\}}e^{i\Theta}
 \mid\mathcal F_\partial]
 \over
 \mathbb E_\nu[e^{i\Theta}\mid\mathcal F_\partial]}
 \right|
 \leq \rho_0^{-1}p^{|\Gamma|}.
 \label{eq:newS7V80-complex-path}
\end{equation}
If at most \(D^n\) selected paths of length \(n\) meet a fixed cell and
\(Dp<1\), their total complex amplitude has an exponentially decaying
tail.
\end{theorem}

\begin{proof}
Apply \eqref{eq:newS7V80-event} to the conditional probability space and
use \eqref{eq:newS7V80-absolute-path}--\eqref{eq:newS7V80-uniform-margin}.
Summing over paths of length \(n\) gives at most
\(\rho_0^{-1}(Dp)^n\), whose tail is geometric.
\end{proof}

\begin{corollary}[Coupling to the V74 modulus compiler]
\label{cor:newS7V80-coupled}
Let a local Radon--Nikodym perturbation of the absolute law have logarithmic
oscillation at most \(r|X|\), and let the unperturbed conditional defect
activity be \(p\).  Then the V74 comparison gives absolute activity
\(p_{\rm abs}=e^rp\).  Under the phase-margin row
\eqref{eq:newS7V80-uniform-margin}, the normalized complex path amplitude
is bounded by
\begin{equation}
 \rho_0^{-1}(e^rp)^{|\Gamma|}.
 \label{eq:newS7V80-coupled-activity}
\end{equation}
The subcritical condition remains \(D e^rp<1\); the phase margin changes
the prefactor but not the exponential rate.
\end{corollary}

\begin{proof}
Use the local activity theorem of V74 to obtain the absolute path bound,
then apply \Cref{thm:newS7V80-path}.
\end{proof}

\begin{remark}[Why conditional margins are required]
\label{rem:newS7V80-conditional}
Cluster exploration repeatedly conditions on exposed collars and boundary
values.  A nonzero phase margin for the full-volume partition function
does not imply a margin for those conditional cylinders.  The uniform row
\eqref{eq:newS7V80-uniform-margin} must therefore cover the actual family
of exploration conditionals.
\end{remark}

\subsection{Exact complex marginal differentiation}
\label{sec:newS7V80-differentiation}

Let
\begin{equation}
 Z(q)=\int e^{-S(q,y)+i\Theta(q,y)}\,d\lambda(y),
 \qquad \Phi(q)=-\log Z(q),
 \label{eq:newS7V80-Z}
\end{equation}
where \(Z\ne0\) on the chart and derivatives are dominated.  Put
\(K=S-i\Theta\) and use the normalized complex expectation associated to
\(e^{-K}\).

\begin{theorem}[Complex logarithmic-marginal identities]
\label{thm:newS7V80-complex-derivatives}
For tangent directions \(u,v\),
\begin{align}
 D\Phi[u]&=\mathbb E_{\mathbb C}[DK[u]],
 \label{eq:newS7V80-first}\\
 D^2\Phi[u,v]
 &=\mathbb E_{\mathbb C}[D^2K[u,v]]
 -\operatorname{Cov}_{\mathbb C}(DK[u],DK[v]).
 \label{eq:newS7V80-second}
\end{align}
The complex covariance in \eqref{eq:newS7V80-second} has no positivity
order.  Moreover, by \Cref{thm:newS7V80-margin},
\begin{equation}
 |\mathbb E_{\mathbb C}F|
 \leq \rho_\nu^{-1}\mathbb E_\nu|F|,
 \label{eq:newS7V80-moment-margin}
\end{equation}
so every derivative loader through the marginal pays an inverse phase
margin.
\end{theorem}

\begin{proof}
Differentiate \(Z\):
\(DZ[u]=-Z\mathbb E_{\mathbb C}DK[u]\).  This proves the first identity.
For every observable \(F\), direct quotient differentiation gives
\[
 D\mathbb E_{\mathbb C}F
 =\mathbb E_{\mathbb C}DF
 -\operatorname{Cov}_{\mathbb C}(F,DK).
\]
Apply it to \(F=DK[u]\) in direction \(v\).  Complex covariances are
complex scalars or forms and therefore have no positive-semidefinite
ordering.  The last bound is \eqref{eq:newS7V80-margin-bound}.
\end{proof}

\begin{corollary}[Small phase margin amplifies response]
\label{cor:newS7V80-amplification}
Even if all unnormalized score moments are volume-local, the normalized
marginal derivatives can diverge when \(\rho_\nu\to0\).  Thus the GCMCC
current bound and the modulus singular-value gap do not control the
complex logarithmic marginal without an independent phase-margin row.
\end{corollary}

\begin{proof}
Apply \eqref{eq:newS7V80-moment-margin} to every term in
\eqref{eq:newS7V80-first}--\eqref{eq:newS7V80-second}.  The preceding
hypotheses do not bound \(\rho_\nu^{-1}\).
\end{proof}

\subsection{The retained-carrier alternative}
\label{sec:newS7V80-carrier}

\begin{proposition}[Marginalization is the source of the phase denominator]
\label{prop:newS7V80-carrier}
Before division by \(Z\), the unnormalized functional
\begin{equation}
 \mathcal L(F):=\int F e^{-S+i\Theta}\,d\lambda
 \label{eq:newS7V80-unnormalized}
\end{equation}
obeys \(|\mathcal L(F)|\leq\int|F|e^{-S}d\lambda\) with no inverse phase
margin.  The factor \(\rho_\nu^{-1}\) enters exactly when one normalizes by
the eliminated-carrier partition function.
\end{proposition}

\begin{proof}
The first statement is the triangle inequality.  Since
\(Z=Z_{\rm abs}z_\nu\), division by \(Z\) introduces
\(|z_\nu|^{-1}=\rho_\nu^{-1}\).
\end{proof}

\begin{assessment}[Upstream route]
\label{ass:newS7V80-carrier}
V72 constructed a positive finite filled-sea Fock transfer while retaining
the fermionic carrier.  YM V45 independently identifies retained-detail
transfer dilation as the upstream response when logarithmic marginal
covariances are circular.  The common lesson is exact: prove locality,
positivity, gauge projection, and renewal projectivity for the joint
carrier transfer first; take a scalar determinant or logarithmic marginal
only after a uniform phase margin or an exact cancellation identity has
been established.
\end{assessment}

\begin{definition}[Coupled complex phase card, CCPC]
\label{def:newS7V80-CCPC}
For every finite regulator and exploration conditional, CCPC requires one
of the following two complete routes.

\emph{Marginal route:}
\begin{enumerate}
\item CRDC controls the positive modulus law and its defect activity;
\item PCC and GCMCC construct the local chiral phase current and compatible
RCLC divisor transitions;
\item a volume- and boundary-uniform conditional phase margin
\(\rho_\nu\geq\rho_0>0\) is proved on every cylinder used by exploration;
\item normalized first and second response stocks include the explicit
factor \(\rho_0^{-1}\); and
\item the complex path condition \(De^rp<1\) and all overlap holonomies are
verified.
\end{enumerate}

\emph{Carrier route:}
\begin{enumerate}
\item a common positive joint boson--fermion transfer acts on a declared
finite Hilbert space before determinant marginalization;
\item gauge projection, chiral multiplet content, reflection or transfer
positivity, and divisor sectors are compatible on that same carrier;
\item renewal embeddings intertwine the transfers and observables;
\item volume-local propagation and stress bounds are uniform; and
\item any later scalar marginal is taken only with a proved nonvanishing
normalization or an observable-specific zero-mode saturation.
\end{enumerate}
\end{definition}

\begin{theorem}[Conditional coupled collar compiler]
\label{thm:newS7V80-compiler}
Under the CCPC marginal route, the absolute V74 collar compiler and the
complex phase compiler combine without changing the defect exponential
rate: normalized bad-path amplitudes obey
\eqref{eq:newS7V80-coupled-activity}.  Under the carrier route, no scalar
complex probability is formed at this stage, so the conclusion is instead
an operator-level transfer bound and must not be reported as a Gibbs
probability estimate.
\end{theorem}

\begin{proof}
The marginal-route assertion is
\Cref{cor:newS7V80-coupled}.  The carrier route retains a positive transfer
operator and its matrix elements; by definition it has not divided by a
complex determinant partition function.  Therefore a complex probability
statement is neither needed nor licensed.
\end{proof}

\begin{assessment}[New exact frontier]
\label{ass:newS7V80-status}
The modulus, divisor, crossing, anomaly arithmetic, projector curvature,
and local chart primitive now feed one conditional compiler.  Its first
unpopulated scalar input is a uniform conditional phase margin.  Since
local phase response alone cannot supply that margin, the next attack must
seek a same-law fluctuation criterion for phase increments and determine
whether it can be uniform.  Failure through extensive variance is a
constructive signal to adopt the retained-carrier route, not to assume the
margin.
\end{assessment}

\begin{programme}[V81: relative phase increments and conditional margins]
\label{prog:newS7V80-V81}
V81 will distinguish the extensive absolute phase from the local phase
increment created by changing one collar.  It will derive exact comparison
identities for the ratio of two conditional complex partition functions,
prove arc and variance criteria for the increment under one reference
conditional law, and test whether a local GCMCC current plus mixing can
control the ratio without a volume-wide phase margin.  If the common
background denominator remains, the note will record the failure and move
upstream to the carrier transfer.
\end{programme}

\begin{firewall}[Claim boundary after seventh-series V80]
\label{fire:newS7V80-boundary}
V80 performs the compulsory YM/NS audit; proves the exact phase-margin
comparison, arc and variance sufficient criteria, conditional complex
bad-path bound, coupling to the V74 modulus activity, complex
logarithmic-marginal derivative identities, and the location of the
inverse phase-margin loss.  It defines CCPC and its two honest routes.

V80 does not prove a uniform conditional phase margin, a phase cluster
expansion, a common projective carrier transfer, or CCPC, GCMCC, PCC,
SMALC, RCLC, DSLC, CRDC, GAOCC, CCTC, MSRC, RTGC, or RCNC, remove the
cutoff, or construct the Standard Model--Einstein continuum.
\end{firewall}

\begin{theorem}[Seventh-series V80 result]
\label{thm:newS7V80-result}
A controlled determinant modulus and a local chiral phase current yield a
normalized complex collar bound only after the conditional phase margin is
bounded below.  With margin \(\rho_0\) and absolute activity \(e^rp\), a
bad path \(\Gamma\) has normalized complex amplitude at most
\(\rho_0^{-1}(e^rp)^{|\Gamma|}\).  Exact marginal differentiation shows
that the same inverse margin amplifies response covariances.  The only
honest alternative is to retain the positive fermionic carrier and prove
transfer projectivity before scalar marginalization.
\end{theorem}

\begin{proof}
Combine \Cref{thm:newS7V80-audit,
thm:newS7V80-margin,
cth:newS7V80-zero,
lem:newS7V80-arc,
lem:newS7V80-variance,
thm:newS7V80-path,
cor:newS7V80-coupled,
thm:newS7V80-complex-derivatives,
prop:newS7V80-carrier,
thm:newS7V80-compiler}.
\end{proof}

\paragraph{Parent-version record.}
V80 was formed from
\texttt{Renewal\_SM\_Einstein\_Continuum\_Research\_Notes\_V79.tex}.
The V79 parent had \(31\,636\) logical source lines, \(1\,248\,237\) bytes,
and SHA--256
\[
 \texttt{5A1C183656DCED8210C439FF5A549BBE57E86C8EAE69156906C91CB358E42EFB}.
\]
The next compulsory companion audit is V85.

% =============================================================================
% END APPENDED SEVENTH-SERIES V80 MATERIAL
% =============================================================================
% =============================================================================
% APPENDED SEVENTH-SERIES V81 MATERIAL
% =============================================================================

\section{Local phase lifts and sequential complex conditionals}
\label{sec:newS7V81}
\label{sec:v7v81}

V80 required a lower phase margin for every complex conditional used in a
cluster exploration.  A volume-wide margin is generally too strong: even
independent local phases can make it exponentially small.  The correct
local question is whether the phase restricted to the variables updated
at one exploration step lies in an arc shorter than a semicircle.  Any
phase depending only on the fixed exterior cancels from that conditional.
This section proves the cancellation, converts a local phase oscillation
into a conditional margin, and derives the resulting per-step activity.
It also states the exact factorization hypothesis needed to iterate
complex conditional kernels.

\subsection{Cancellation of the fixed exterior phase}
\label{sec:newS7V81-exterior}

Let \(x\) denote variables in a finite update set \(X\), let \(y\) denote
the fixed exterior data, and let \(\nu_y\) be a positive conditional
probability on \(x\).  Suppose a real lift of the phase on this conditional
chart has the form
\begin{equation}
 \Theta(x,y)=\Theta_{\rm ext}(y)+\theta_X(x;y).
 \label{eq:newS7V81-split}
\end{equation}

\begin{theorem}[Exterior-phase cancellation]
\label{thm:newS7V81-exterior}
Whenever the denominator is nonzero, the normalized complex conditional
of a bounded observable \(F(x)\) is
\begin{equation}
 {\int F(x)e^{i\Theta(x,y)}\,d\nu_y(x)
  \over
  \int e^{i\Theta(x,y)}\,d\nu_y(x)}
 =
 {\int F(x)e^{i\theta_X(x;y)}\,d\nu_y(x)
  \over
  \int e^{i\theta_X(x;y)}\,d\nu_y(x)}.
 \label{eq:newS7V81-exterior-cancel}
\end{equation}
Thus only the phase variation over the updated variables enters the
conditional margin.
\end{theorem}

\begin{proof}
The factor \(e^{i\Theta_{\rm ext}(y)}\) is independent of \(x\) and
multiplies both numerator and denominator.  Cancel it.
\end{proof}

\begin{remark}[A lift is part of the chart data]
\label{rem:newS7V81-lift}
The phase is intrinsically circle valued.  Equation
\eqref{eq:newS7V81-split} uses a real lift on one nonvanishing determinant
chart.  RCLC transition functions must compare different lifts, and a
loop with nontrivial determinant-line holonomy may prevent a global one.
\end{remark}

\subsection{Oscillation gives a local conditional margin}
\label{sec:newS7V81-oscillation}

For a real function \(f\), write
\(\operatorname{osc}f=\sup f-\inf f\).

\begin{theorem}[Oscillation-to-margin theorem]
\label{thm:newS7V81-osc-margin}
If
\begin{equation}
 w_X(y):=\operatorname{osc}_{x}\theta_X(x;y)<\pi,
 \label{eq:newS7V81-width}
\end{equation}
then for every positive conditional probability \(\nu_y\),
\begin{equation}
 \left|\int e^{i\theta_X(x;y)}\,d\nu_y(x)\right|
 \geq \cos\left({w_X(y)\over2}\right)>0.
 \label{eq:newS7V81-local-margin}
\end{equation}
The bound is unchanged after any positive reweighting of \(\nu_y\).
\end{theorem}

\begin{proof}
Let
\(c=(\sup_x\theta_X+\inf_x\theta_X)/2\).  Then
\(|\theta_X-c|\leq w_X/2<\pi/2\).  Rotate the integral by \(e^{-ic}\) and
take its real part:
\[
 \operatorname{Re}\int e^{i(\theta_X-c)}d\nu_y
 =\int\cos(\theta_X-c)d\nu_y
 \geq\cos(w_X/2).
\]
The modulus is at least this real part.  A positive reweighting produces
another probability supported on the same phase range, so the same proof
applies.
\end{proof}

\begin{corollary}[Local complex event bound]
\label{cor:newS7V81-event}
Under \eqref{eq:newS7V81-width}, for every event \(A\) in the update
variables,
\begin{equation}
 \left|
 {\int_Ae^{i\theta_X}\,d\nu_y
  \over\int e^{i\theta_X}\,d\nu_y}
 \right|
 \leq {\nu_y(A)\over\cos(w_X(y)/2)}.
 \label{eq:newS7V81-event}
\end{equation}
\end{corollary}

\begin{proof}
The numerator has modulus at most \(\nu_y(A)\), and the denominator is
bounded by \eqref{eq:newS7V81-local-margin}.
\end{proof}

\subsection{Loading the oscillation from the Berry current}
\label{sec:newS7V81-current}

Suppose a GCMCC chart supplies a real phase lift \(\theta\) whose
differential is the declared real one-form \(-i\mathcal J\).  Let any two
conditional configurations differing only in \(X\) be joined inside the
chart by a path \(\gamma\) of length at most \(\ell_X\).

\begin{proposition}[Current-to-phase-width loader]
\label{prop:newS7V81-current-width}
If
\begin{equation}
 \sup_{z\in\gamma}\sup_{\substack{v\ {
m supported\ in}\ X\\
 \|v\|=1}}
 |\mathcal J_z(v)|\leq L_X,
 \label{eq:newS7V81-current-bound}
\end{equation}
then
\begin{equation}
 w_X(y)\leq L_X\ell_X.
 \label{eq:newS7V81-width-bound}
\end{equation}
Consequently \(L_X\ell_X<\pi\) is a sufficient local phase-margin row.
\end{proposition}

\begin{proof}
For endpoints \(x_0,x_1\) of the path,
\[
 |\theta(x_1,y)-\theta(x_0,y)|
 =\left|\int_\gamma -i\mathcal J\right|
 \leq\int_\gamma L_X\,ds\leq L_X\ell_X.
\]
Take the supremum over endpoint pairs.
\end{proof}

\begin{corollary}[Fixed-size collar criterion]
\label{cor:newS7V81-collar}
If the V79 cell-kernel bound gives a current response at most
\(L_0|X|\|v\|_\infty\), and every update path has
\(\ell_X\leq\ell_0|X|\) in the matching norm, then
\begin{equation}
 w_X\leq L_0\ell_0|X|^2.
 \label{eq:newS7V81-collar-width}
\end{equation}
For a uniformly bounded star or collar size \(q\), the sufficient row is
\(L_0\ell_0q^2<\pi\), independent of total volume.
\end{corollary}

\begin{proof}
Insert the two declared bounds into
\eqref{eq:newS7V81-width-bound} and use \(|X|\leq q\).
\end{proof}

\begin{remark}[Norm matching can improve the square]
\label{rem:newS7V81-norm}
The factor \(|X|^2\) reflects the deliberately crude pairing of an
\(\ell^\infty\) response bound with an extensive path length.  A dual
\(\ell^1\)--\(\ell^\infty\) formulation or a coordinate-by-coordinate
path may reduce it.  Such an improvement must be proved from the actual
local current kernel and path geometry.
\end{remark}

\subsection{Sequential complex kernels}
\label{sec:newS7V81-sequential}

A complex conditional is not a probability kernel.  Iteration therefore
requires an operator norm statement rather than an appeal to probabilistic
tower properties.  Let \(K_j(\omega,dz)\) be normalized complex kernels
for exploration step \(j\), and let \(B_j\) be the bad subset at that
step.  Define
\begin{equation}
 (T_jf)(\omega):=\int_{B_j}f(\omega,z)K_j(\omega,dz).
 \label{eq:newS7V81-Tj}
\end{equation}

\begin{theorem}[Sequential bad-step product bound]
\label{thm:newS7V81-sequential}
Assume each \(K_j\) is obtained from a positive reference conditional
\(\nu_{j,\omega}\) by a phase of width at most \(w<\pi\), and
\begin{equation}
 \nu_{j,\omega}(B_j)\leq p_{\rm abs}
 \label{eq:newS7V81-step-prob}
\end{equation}
uniformly in \(j\) and the exploration history \(\omega\).  Then
\begin{equation}
 \|T_j\|_{\infty\to\infty}
 \leq q_{\mathbb C}:={p_{\rm abs}\over\cos(w/2)}.
 \label{eq:newS7V81-step-norm}
\end{equation}
For a declared factorization in which the amplitude of a fixed
\(n\)-step bad history is
\((T_1\cdots T_n\mathbf1)(\omega_0)\),
\begin{equation}
 |(T_1\cdots T_n\mathbf1)(\omega_0)|
 \leq q_{\mathbb C}^{,n}.
 \label{eq:newS7V81-product}
\end{equation}
\end{theorem}

\begin{proof}
For \(\|f\|_\infty\leq1\), the numerator defining \(T_jf\) has modulus at
most \(\nu_{j,\omega}(B_j)\), while
\Cref{thm:newS7V81-osc-margin} bounds the normalization by
\(\cos(w/2)\).  This proves \eqref{eq:newS7V81-step-norm}.  Submultiplicity
of the operator norm gives \eqref{eq:newS7V81-product}.
\end{proof}

\begin{corollary}[Local complex exploration criterion]
\label{cor:newS7V81-exploration}
If a modulus perturbation of oscillation \(r\) per step turns a base bad
activity \(p\) into \(p_{\rm abs}=e^rp\), and at most \(D^n\) histories
of length \(n\) meet a cell, then the complex exploration tail decays
exponentially provided
\begin{equation}
 \boxed{
 D\,{e^rp\over\cos(w/2)}<1,
 \qquad w<\pi.}
 \label{eq:newS7V81-complex-subcritical}
\end{equation}
This criterion uses local conditional margins and does not contain a
global phase-margin prefactor.
\end{corollary}

\begin{proof}
Combine the V74 modulus reweighting with
\Cref{thm:newS7V81-sequential}, multiply by the history count \(D^n\),
and sum the resulting geometric tail.
\end{proof}

\begin{firewall}[The declared factorization is essential]
\label{fire:newS7V81-factorization}
Equation \eqref{eq:newS7V81-product} is valid only after the complex
amplitude has been represented by the ordered kernels \(K_j\) on the same
exploration filtration.  Pointwise bounds on unrelated one-cell
conditionals do not imply that representation.  Nonlocal determinant
phases, zeros of intermediate normalizations, and overlap holonomy can
all obstruct it.
\end{firewall}

\subsection{Why the global margin can still vanish}
\label{sec:newS7V81-global}

\begin{countertheorem}[Uniform local margins need not give a global margin]
\label{cth:newS7V81-global}
For every \(0<\eta<\pi/2\), there is a product system with one-site phase
width \(2\eta<\pi\) and uniform one-site margin \(\cos\eta\), while its
volume-\(N\) phase margin is \((\cos\eta)^N\) and tends to zero.
\end{countertheorem}

\begin{proof}
Let \(N\) independent spins take values \(\pm1\) with equal probability
and set \(\Theta=\eta\sum_{j=1}^N\sigma_j\).  Conditional on every other
spin, the two one-site phases differ by \(2\eta\), and their normalized
average has modulus \(\cos\eta\).  Independence gives
\[
 \mathbb E e^{i\Theta}
 =\prod_{j=1}^N\mathbb E e^{i\eta\sigma_j}
 =(\cos\eta)^N.
\]
\end{proof}

\begin{corollary}[Local DLR data and global normalization are different
rows]
\label{cor:newS7V81-DLR}
A projectively consistent family of nonzero local complex kernels may have
no volume-uniform absolute phase margin.  Local correlation control must
therefore be proved from kernel factorization such as
\Cref{thm:newS7V81-sequential}, while global partition normalization and
free-energy phase require a separate thermodynamic argument.
\end{corollary}

\begin{proof}
The product example has perfectly consistent local kernels but an
exponentially small global margin.  Hence neither row implies the other.
\end{proof}

\subsection{Local conditional phase card}
\label{sec:newS7V81-card}

\begin{definition}[Local conditional phase-lift card, LCPLC]
\label{def:newS7V81-LCPLC}
For each exploration step and admissible boundary history, LCPLC requires:
\begin{enumerate}
\item a nonvanishing determinant chart and a real phase lift split as in
\eqref{eq:newS7V81-split};
\item a path inside the common kernel-gap chart connecting every pair of
updated configurations;
\item a GCMCC current bound giving a uniform phase width \(w<\pi\);
\item a normalized complex kernel whose denominator obeys
\eqref{eq:newS7V81-local-margin};
\item an exact ordered-kernel factorization of the exploration amplitude;
\item the subcritical row \eqref{eq:newS7V81-complex-subcritical}; and
\item compatibility of phase lifts with RCLC divisor charts, gauge
transformations, chart overlaps, renewal embeddings, and the retained
fermionic observables.
\end{enumerate}
\end{definition}

\begin{theorem}[Conditional replacement for the V80 global-margin row]
\label{thm:newS7V81-replacement}
For bad-cluster and covariance exploration, LCPLC rows 1--6 replace the
V80 volume-uniform global phase-margin hypothesis and give an exponential
complex history bound.  They do not replace the global normalization row
needed to define a finite-volume partition function or its thermodynamic
pressure.
\end{theorem}

\begin{proof}
Rows 1--4 give the one-step bound; row 5 licenses its iteration; row 6
makes the history sum geometric.  This is
\Cref{cor:newS7V81-exploration}.  The counterexample
\Cref{cth:newS7V81-global} proves the final limitation.
\end{proof}

\begin{assessment}[What the relative-phase attack achieved]
\label{ass:newS7V81-status}
The common background phase is harmless inside a fixed-boundary local
conditional because it cancels exactly.  A local phase response can yield
a uniform step margin if its lifted oscillation is below \(\pi\), and the
only cost is the explicit activity factor \(1/\cos(w/2)\).  The remaining
nonformal inputs are an actual local phase lift for the chiral determinant,
an exact exploration-kernel factorization, and a separate global
normalization or carrier-transfer construction.
\end{assessment}

\begin{programme}[V82: one-fibre phase transport and holonomy]
\label{prog:newS7V81-V82}
V82 will restrict the projector connection to a single compact gauge-link
fibre with its boundary fixed.  It will prove the path-dependence formula
for two phase transports, express it as the curvature flux through a
spanning surface, and give a small-flux condition for a single-valued local
lift.  This directly tests LCPLC rows 1--3 and identifies whether the
one-link fibre topology forces an additional holonomy writer.
\end{programme}

\begin{firewall}[Claim boundary after seventh-series V81]
\label{fire:newS7V81-boundary}
V81 proves cancellation of exterior phases in local conditionals, the
phase-oscillation margin, loading of that oscillation from a Berry current,
the sequential complex-kernel product bound, the local subcritical
criterion, and the separation between local kernel margins and global
partition margins.

V81 does not construct the required chiral phase lifts or ordered kernel
factorization, prove a global normalization, LCPLC, CCPC, GCMCC, PCC,
SMALC, RCLC, DSLC, CRDC, GAOCC, CCTC, MSRC, RTGC, or RCNC, remove the
cutoff, or construct the Standard Model--Einstein continuum.
\end{firewall}

\begin{theorem}[Seventh-series V81 result]
\label{thm:newS7V81-result}
On a fixed-boundary update, every exterior determinant phase cancels from
the normalized conditional.  If the remaining lifted phase has width
\(w<\pi\), its conditional margin is at least \(\cos(w/2)\).  Under an
exact sequential-kernel factorization, a modulus activity \(e^rp\) becomes
the complex activity \(e^rp/\cos(w/2)\), yielding exponential exploration
when \eqref{eq:newS7V81-complex-subcritical} holds.  This local result does
not imply a volume-uniform global phase margin.
\end{theorem}

\begin{proof}
Combine \Cref{thm:newS7V81-exterior,
thm:newS7V81-osc-margin,
prop:newS7V81-current-width,
thm:newS7V81-sequential,
cor:newS7V81-exploration,
cth:newS7V81-global,
thm:newS7V81-replacement}.
\end{proof}

\paragraph{Parent-version record.}
V81 was formed from
\texttt{Renewal\_SM\_Einstein\_Continuum\_Research\_Notes\_V80.tex}.
The V80 parent had \(32\,124\) logical source lines, \(1\,267\,475\) bytes,
and SHA--256
\[
 \texttt{E5BC3A22E30FBA23633A50A9E7EBA402B4530E0E59EC2532663B402C9ED7FFE1}.
\]
The next compulsory companion audit is V85.

% =============================================================================
% END APPENDED SEVENTH-SERIES V81 MATERIAL
% =============================================================================
% =============================================================================
% APPENDED SEVENTH-SERIES V82 MATERIAL
% =============================================================================

\section{Radial gauge, one-fibre phase transport, and holonomy}
\label{sec:newS7V82}
\label{sec:v7v82}

V81 converted a local phase width into a conditional complex margin.  A
geometric distinction must now be made explicit.  The projector measure
current is a connection one-form whose exterior derivative is the
projector curvature; unless that curvature vanishes, it cannot itself be
the differential of a scalar phase.  On a star-shaped chart the radial
homotopy instead puts the connection into radial gauge: the transformed
connection equals the radial curvature primitive \(Kf\), while the frame
change has differential \(a-Kf\).  This section proves that identity,
quantifies path dependence by curvature flux, and corrects the precise
interpretation of the V81 phase-width loader.

\subsection{Real connection conventions}
\label{sec:newS7V82-conventions}

Let \(\mathcal A\) and \(\mathcal F=d\mathcal A\) be the imaginary-valued
determinant-line connection and curvature of V78.  Define the real forms
\begin{equation}
 a:=-i\mathcal A,
 \qquad
 f:=-i\mathcal F=da.
 \label{eq:newS7V82-real}
\end{equation}
On a star-shaped chart \(\mathcal U\) about \(x_0\), let \(K\) be the
radial homotopy of V79 and put
\begin{equation}
 \vartheta:=Ka,
 \qquad
 j:=Kf.
 \label{eq:newS7V82-theta-j}
\end{equation}
Here \(\vartheta\) is a scalar and \(j\) is a real one-form.

\begin{theorem}[Exact radial-gauge identity]
\label{thm:newS7V82-radial-gauge}
The gauge-transformed connection
\begin{equation}
 a_{\rm rad}:=a-d\vartheta
 \label{eq:newS7V82-arad-def}
\end{equation}
satisfies
\begin{equation}
 \boxed{a_{\rm rad}=j=Kf,\qquad da_{\rm rad}=f.}
 \label{eq:newS7V82-radial-gauge}
\end{equation}
It also obeys the radial condition
\begin{equation}
 (a_{\rm rad})_x(x-x_0)=0.
 \label{eq:newS7V82-radial-condition}
\end{equation}
The scalar frame-change phase has
\begin{equation}
 d\vartheta=a-j,
 \label{eq:newS7V82-frame-phase}
\end{equation}
not \(d\vartheta=j\) unless special additional conditions hold.
\end{theorem}

\begin{proof}
Apply the V79 homotopy identity to the one-form \(a\):
\[
 dKa+Kda=a.
\]
Since \(Ka=\vartheta\), \(da=f\), and \(Kf=j\), rearrangement gives
\(a-d\vartheta=j\), proving
\eqref{eq:newS7V82-radial-gauge} and
\eqref{eq:newS7V82-frame-phase}.  Differentiation gives
\(da_{\rm rad}=dj=f\).  Finally, by the two-form formula for \(K\),
\[
 j_x(x-x_0)=\int_0^1t\,
 f_{x_0+t(x-x_0)}(x-x_0,x-x_0)\,dt=0
\]
because a two-form vanishes on two equal vectors.
\end{proof}

\begin{corollary}[Radial transport fixes the local frame]
\label{cor:newS7V82-radial-frame}
In radial gauge, parallel transport from \(x_0\) to \(x\) along the radial
segment has phase one.  Thus a base vector in the determinant line defines
a unique radial frame on the star-shaped chart.  Its transverse connection
response is exactly \(j=Kf\).
\end{corollary}

\begin{proof}
The tangent to the radial segment is proportional to \(x-x_0\).
Equation \eqref{eq:newS7V82-radial-condition} makes the connection integral
along it zero.  The remaining assertion is
\eqref{eq:newS7V82-radial-gauge}.
\end{proof}

\begin{correction}[Precise use of the V81 current loader]
\label{corr:newS7V82-V81}
The hypothesis preceding V81
\Cref{prop:newS7V81-current-width} stated that a scalar phase lift could
have differential \(-i\mathcal J\).  Taken literally for the projector
connection, that is possible only when the pulled-back curvature vanishes,
because a scalar differential is closed whereas
\(d(-i\mathcal J)=-i\mathcal F\).

The valid general statement is the following.  In the radial frame,
\(-i\mathcal J=j\) is the connection response along a declared path, so
its line integral controls parallel-transport phase width along that path.
The scalar frame-change phase is instead \(\vartheta\) with
\(d\vartheta=a-j\).  Any use of V81 for an actual scalar determinant
coefficient must specify the radial frame and include the remaining
coefficient response in that frame.  V81's oscillation-to-margin theorem
itself is unchanged.
\end{correction}

\subsection{Path dependence equals curvature flux}
\label{sec:newS7V82-paths}

For a piecewise smooth path \(\gamma\), define its abelian connection
transport phase by
\begin{equation}
 U_a(\gamma):=\exp\left(i\int_\gamma a\right).
 \label{eq:newS7V82-transport}
\end{equation}

\begin{theorem}[Two-path curvature-flux formula]
\label{thm:newS7V82-flux}
Let \(\gamma_0\) and \(\gamma_1\) have the same endpoints and lie in one
connection chart.  If the closed loop
\(\gamma_1\cdot\overline{\gamma_0}\) bounds an oriented piecewise smooth
surface \(\Sigma\) in that chart, then
\begin{equation}
 {U_a(\gamma_1)\over U_a(\gamma_0)}
 =\exp\left(i\oint_{\gamma_1\cdot\overline{\gamma_0}}a\right)
 =\exp\left(i\int_\Sigma f\right).
 \label{eq:newS7V82-flux}
\end{equation}
The same formula holds in radial gauge with \(a\) replaced by \(j=Kf\).
\end{theorem}

\begin{proof}
The quotient of exponentials is the exponential of the oriented closed
line integral.  Stokes' theorem and \(da=f\) give the surface integral.
A gauge change adds an exact one-form to the connection, whose closed-loop
integral is an integer multiple of \(2\pi\) for a single-valued unitary
transition; its exponential is one.  Since radial gauge is such a local
gauge, the formula is unchanged.
\end{proof}

\begin{corollary}[Quantitative path ambiguity]
\label{cor:newS7V82-path-bound}
Suppose
\begin{equation}
 |f_x(u,v)|\leq M_f\|u\|\|v\|
 \label{eq:newS7V82-f-bound}
\end{equation}
on \(\Sigma\), and let \(\operatorname{Area}(\Sigma)\) be its induced
parametric area.  Then the circle distance between the two transport
phases is at most
\begin{equation}
 \operatorname{dist}_{S^1}
 \bigl(U_a(\gamma_1),U_a(\gamma_0)\bigr)
 \leq
 \min\{\pi,M_f\operatorname{Area}(\Sigma)\}.
 \label{eq:newS7V82-path-bound}
\end{equation}
In particular, flux below \(\pi\) selects a unique principal phase
difference, but it does not make the difference zero.
\end{corollary}

\begin{proof}
The absolute flux is at most
\(M_f\operatorname{Area}(\Sigma)\).  The geodesic circle distance of
\(e^{i\alpha}\) from one is at most \(\min\{\pi,|\alpha|\}\).  Apply
\eqref{eq:newS7V82-flux}.
\end{proof}

\begin{theorem}[Path independence criterion]
\label{thm:newS7V82-path-independent}
Transport between endpoints is path independent on a connected region if
and only if
\begin{equation}
 \exp\left(i\oint_\gamma a\right)=1
 \label{eq:newS7V82-trivial-holonomy}
\end{equation}
for every closed loop \(\gamma\) in the region.  On a simply connected
region this holds if \(f=0\).  Conversely, path independence implies
\(f=0\).  Thus a merely small nonzero curvature cannot yield exact path
independence.
\end{theorem}

\begin{proof}
If all closed-loop holonomies are one, the quotient of transports along
any two endpoint paths is one.  The converse follows by comparing any
loop with the constant path.  When the region is simply connected and
\(f=0\), Stokes' theorem or the Poincare lemma gives trivial loop
holonomy.  If transport is path independent, every sufficiently small
contractible loop has zero holonomy.  Dividing its flux by its area and
shrinking it to a point shows every component of \(f\) is zero.
\end{proof}

\begin{corollary}[Flat does not mean globally trivial]
\label{cor:newS7V82-flat-holonomy}
On a nonsimply connected fibre, \(f=0\) does not by itself imply
\eqref{eq:newS7V82-trivial-holonomy}.  A flat determinant line may carry a
nontrivial character of the fibre fundamental group.
\end{corollary}

\begin{proof}
Flat unitary line connections are classified, up to the relevant gauge
equivalence, by their holonomy representations
\(\pi_1\to U(1)\).  Nontrivial representations can exist when
\(\pi_1\ne0\).
\end{proof}

\subsection{A compact one-link gauge fibre}
\label{sec:newS7V82-fibre}

Fix all variables except one lattice link.  Its unfixed fibre is a compact
gauge group \(G\), or a quotient or slice thereof.  A kernel-gap
admissibility condition may remove a closed subset and change its topology.

\begin{proposition}[Local chart versus full-fibre topology]
\label{prop:newS7V82-local-global}
On every sufficiently small coordinate ball contained in the kernel-gap
set, radial gauge exists and the identities
\eqref{eq:newS7V82-radial-gauge}--\eqref{eq:newS7V82-flux} hold.  Extending
the phase around the full one-link fibre additionally requires:
\begin{enumerate}
\item trivial holonomy on noncontractible loops;
\item integral curvature periods compatible with the determinant-line
Chern class; and
\item compatible RCLC transitions around every removed determinant or
kernel-gap divisor.
\end{enumerate}
These conditions do not follow from a pointwise kernel gap or a small
local curvature norm.
\end{proposition}

\begin{proof}
A small coordinate ball is star-shaped, so
\Cref{thm:newS7V82-radial-gauge} applies.  A global extension compares
local frames on overlaps.  Closed loops test their holonomy, closed
surfaces test the integral first Chern class as in V78, and loops around
removed divisors test the transition functions found in V76.  Local norm
bounds control small fillings but do not determine any of these discrete
or global data.
\end{proof}

\begin{remark}[Standard Model fibre warning]
\label{rem:newS7V82-SM-fibre}
The electroweak factor contains a \(U(1)\) direction, and the physical
Standard Model gauge group may be taken with a finite central quotient.
The fundamental group and allowed large loops therefore depend on the
declared global group, exactly as required by SMALC.  Lie-algebra anomaly
arithmetic alone cannot decide their flat holonomies.
\end{remark}

\subsection{A corrected local phase-width compiler}
\label{sec:newS7V82-compiler}

Let the radial-gauge connection response \(j\) satisfy a local bound along
a declared conditional path \(\gamma\), and let
\(r_{\rm coeff}\) denote the real phase-response one-form of the remaining
scalar fermionic coefficient in the radial frame.  Define the total
conditional phase response
\begin{equation}
 q_{\rm ph}:=j+r_{\rm coeff}.
 \label{eq:newS7V82-total-response}
\end{equation}

\begin{theorem}[Corrected phase-width loader]
\label{thm:newS7V82-loader}
If every pair of conditional configurations is joined by a declared path
of length at most \(\ell_X\) and
\begin{equation}
 \sup_{z\in\gamma}\sup_{\|v\|=1}
 |q_{{\rm ph},z}(v)|\leq L_{\rm ph},
 \label{eq:newS7V82-q-bound}
\end{equation}
then the scalar phase coefficient transported in the radial frame has
width
\begin{equation}
 w_X\leq L_{\rm ph}\ell_X.
 \label{eq:newS7V82-width}
\end{equation}
If two allowed path conventions are used, their widths differ by at most
the curvature-flux bound \eqref{eq:newS7V82-path-bound} plus the
corresponding coefficient-response loop integral.
\end{theorem}

\begin{proof}
Integrate the complete response \(q_{\rm ph}\) along the declared path and
use the length bound, exactly as in the fundamental theorem for line
integrals.  For two paths, close them into a loop.  The connection part is
\Cref{cor:newS7V82-path-bound}; the scalar coefficient part is its ordinary
loop integral.
\end{proof}

\begin{definition}[One-fibre holonomy card, OFHC]
\label{def:newS7V82-OFHC}
For every link fibre used by LCPLC, OFHC requires:
\begin{enumerate}
\item a cover by kernel-gap radial charts and the explicit projectors,
connections, curvatures, and radial currents on them;
\item a declared radial frame and the remaining scalar coefficient response
\(r_{\rm coeff}\);
\item a uniform bound on the total response \eqref{eq:newS7V82-total-response};
\item curvature fluxes for competing conditional paths and transition
functions on chart overlaps;
\item all flat holonomies on noncontractible fibre loops;
\item RCLC data for removed rank divisors and SMALC data for the chosen
global gauge group; and
\item the final width \(w_X<\pi\) inserted into the V81 complex activity.
\end{enumerate}
\end{definition}

\begin{assessment}[The one-fibre bottleneck]
\label{ass:newS7V82-status}
The local projector curvature can always be put into radial gauge, and
path ambiguity is exactly its flux.  This repairs the phase-width compiler:
one must bound the full radial-frame scalar response, not identify the
curvature primitive with an exact scalar differential.  The next analytic
input is a localized bound on the projector curvature kernel and hence on
\(j=Kf\); the remaining scalar coefficient and global fibre holonomies
stay independent rows.
\end{assessment}

\begin{programme}[V83: localized projector-curvature kernel]
\label{prog:newS7V82-V83}
V83 will combine the Wilson-kernel spectral gap, the resolvent derivative
of V78, and the Combes--Thomas mechanism of V73.  It will prove exponential
decay of a projector derivative caused by a local link variation and then
of the bilocal curvature components.  This will populate the analytic
part of GCMCC and OFHC while leaving the common-regulator anomaly sum and
global holonomy explicit.
\end{programme}

\begin{firewall}[Claim boundary after seventh-series V82]
\label{fire:newS7V82-boundary}
V82 proves the exact radial-gauge identity, corrects the interpretation of
the V81 current loader, proves the two-path curvature-flux formula and its
quantitative bound, characterizes path independence and flat holonomy,
and gives the corrected total-response phase-width loader.

V82 does not bound the remaining scalar coefficient response, compute the
full one-link fibre holonomies, prove OFHC, LCPLC, CCPC, GCMCC, PCC, SMALC,
RCLC, DSLC, CRDC, GAOCC, CCTC, MSRC, RTGC, or RCNC, remove the cutoff, or
construct the Standard Model--Einstein continuum.
\end{firewall}

\begin{theorem}[Seventh-series V82 result]
\label{thm:newS7V82-result}
On a star-shaped determinant-line chart, radial frame change puts the real
connection into the exact form \(a_{\rm rad}=Kf\), while the frame-change
phase obeys \(d\vartheta=a-Kf\).  Parallel transports along two endpoint
paths differ by the exponential of the curvature flux between them.
Therefore a local conditional phase width must be loaded from the complete
radial-frame response and declared path geometry; small curvature controls
but does not eliminate holonomy.
\end{theorem}

\begin{proof}
Combine \Cref{thm:newS7V82-radial-gauge,
cor:newS7V82-radial-frame,
thm:newS7V82-flux,
cor:newS7V82-path-bound,
thm:newS7V82-path-independent,
prop:newS7V82-local-global,
thm:newS7V82-loader}.
\end{proof}

\paragraph{Parent-version record.}
V82 was formed from
\texttt{Renewal\_SM\_Einstein\_Continuum\_Research\_Notes\_V81.tex}.
The V81 parent had \(32\,526\) logical source lines, \(1\,282\,517\) bytes,
and SHA--256
\[
 \texttt{23DDC59B57E8F7792B4B0743F59D816A5B25BA965EB77873143619DF656DC45C}.
\]
The next compulsory companion audit is V85.

% =============================================================================
% END APPENDED SEVENTH-SERIES V82 MATERIAL
% =============================================================================
% =============================================================================
% APPENDED SEVENTH-SERIES V83 MATERIAL
% =============================================================================

\section{Exponential locality of projector derivatives and curvature}
\label{sec:newS7V83}
\label{sec:v7v83}

V82 reduced the one-fibre phase-width problem to bounds on the radial-gauge
connection \(j=Kf\), the remaining scalar coefficient response, and global
holonomy.  The analytic input for \(j\) is locality of the finite
projector curvature.  This section derives it from a uniform spectral gap,
a Combes--Thomas resolvent bound, and local dependence of the Hermitian
kernel on each link or cell variable.  The constants are independent of
volume on graphs with uniform exponential row sums.

\subsection{Local finite-range kernel hypotheses}
\label{sec:newS7V83-hypotheses}

Let \((\Lambda,d)\) be a finite metric graph with fibre dimension at most
\(n_0\) per cell.  Let \(H=H^*\) act on
\(\bigoplus_{x\in\Lambda}\mathbb C^{n_x}\).  Assume:
\begin{enumerate}
\item \(H_{xy}=0\) when \(d(x,y)>R_H\), and its row norm is uniformly
bounded;
\item \(\operatorname{spec}H\cap[-\delta,\delta]=\varnothing\) for one
\(\delta>0\);
\item a Riesz contour \(\Gamma\) around the negative spectrum has length
\(L_\Gamma\), stays a distance at least \(\eta>0\) from the spectrum, and
its resolvent obeys
\begin{equation}
 \|(z-H)^{-1}_{xy}\|
 \leq C_R e^{-\mu d(x,y)},
 \qquad z\in\Gamma;
 \label{eq:newS7V83-CT}
\end{equation}
\item a unit tangent variation at cell \(p\) gives
\(V_p=d_pH\), supported in
\(N_{R_V}(p)\times N_{R_V}(p)\), with
\begin{equation}
 \sup_p\sum_{u,v}\|(V_p)_{uv}\|\leq C_V;
 \label{eq:newS7V83-V}
\end{equation}
and
\item for every \(\alpha>0\) used below,
\begin{equation}
 S_\alpha:=\sup_{p\in\Lambda}
 \sum_{x\in\Lambda}e^{-\alpha d(x,p)}<\infty
 \label{eq:newS7V83-row}
\end{equation}
uniformly over the regulator family.
\end{enumerate}
The V73 Combes--Thomas theorem supplies item 3 from items 1--2 with
appropriate graph-growth constants.  It is retained here as an explicit
hypothesis so every constant remains visible.

Let
\begin{equation}
 P:=\mathbf1_{(-\infty,0)}(H).
 \label{eq:newS7V83-P}
\end{equation}

\subsection{Localized first derivative of the projector}
\label{sec:newS7V83-dP}

\begin{theorem}[Two-sided localization of a projector derivative]
\label{thm:newS7V83-dP}
Under the hypotheses above,
\begin{equation}
 \|(d_pP)_{xy}\|
 \leq C_P
 e^{-\mu d(x,p)}e^{-\mu d(y,p)},
 \label{eq:newS7V83-dP}
\end{equation}
where one may take
\begin{equation}
 C_P:={L_\Gamma\over2\pi}C_R^2C_Ve^{2\mu R_V}.
 \label{eq:newS7V83-CP}
\end{equation}
In particular,
\begin{equation}
 \|d_pP\|_1\leq n_0 C_PS_\mu^2,
 \qquad
 \|d_pP\|_2\leq n_0^{1/2}C_PS_{2\mu},
 \label{eq:newS7V83-Schatten}
\end{equation}
up to the harmless choice of block Schatten normalization, uniformly in
volume.
\end{theorem}

\begin{proof}
The V78 resolvent derivative formula gives
\[
 (d_pP)_{xy}={1\over2\pi i}\oint_\Gamma
 \sum_{u,v}(z-H)^{-1}_{xu}(V_p)_{uv}
 (z-H)^{-1}_{vy}\,dz.
\]
Only \(u,v\in N_{R_V}(p)\) occur.  For such indices,
\[
 e^{-\mu d(x,u)}\leq e^{\mu R_V}e^{-\mu d(x,p)},
 \qquad
 e^{-\mu d(v,y)}\leq e^{\mu R_V}e^{-\mu d(p,y)}.
\]
Use \eqref{eq:newS7V83-CT}, integrate the contour length, and sum
\eqref{eq:newS7V83-V}; this proves
\eqref{eq:newS7V83-dP}--\eqref{eq:newS7V83-CP}.
Summing the block trace norms over \(x,y\) gives the first Schatten bound.
Squaring and summing the block Hilbert--Schmidt norms gives the second,
with the displayed fibre factor.  Uniform row sums remove any volume
factor.
\end{proof}

\begin{corollary}[Remote link response]
\label{cor:newS7V83-remote}
If both matrix indices lie a distance at least \(L\) from the varied cell,
then
\begin{equation}
 \|(d_pP)_{xy}\|\leq C_Pe^{-2\mu L}.
 \label{eq:newS7V83-remote}
\end{equation}
Thus a local link variation changes the occupied subspace only through
exponentially small remote matrix elements.
\end{corollary}

\begin{proof}
Insert \(d(x,p),d(y,p)\geq L\) into
\eqref{eq:newS7V83-dP}.
\end{proof}

\subsection{Bilocal decay of curvature components}
\label{sec:newS7V83-curvature}

Let \(u_p\) and \(v_q\) be unit tangent variations supported at cells
\(p\) and \(q\).  The real curvature component is, up to the fixed factor
\(-i\),
\begin{equation}
 f_{pq}(u_p,v_q)
 =-i\operatorname{Tr}P
 \left[(d_pP)(d_qP)-(d_qP)(d_pP)\right].
 \label{eq:newS7V83-fpq}
\end{equation}

\begin{lemma}[Exponential overlap sum]
\label{lem:newS7V83-overlap}
For every \(p,q\),
\begin{equation}
 \sum_{z\in\Lambda}
 e^{-\mu d(z,p)}e^{-\mu d(z,q)}
 \leq S_{\mu/2}e^{-(\mu/2)d(p,q)}.
 \label{eq:newS7V83-overlap}
\end{equation}
\end{lemma}

\begin{proof}
The triangle inequality gives
\(d(z,p)+d(z,q)\geq d(p,q)\).  Hence
\[
 e^{-\mu[d(z,p)+d(z,q)]}
 \leq e^{-(\mu/2)d(p,q)}e^{-(\mu/2)d(z,p)}.
\]
Sum over \(z\) and use \eqref{eq:newS7V83-row}.
\end{proof}

\begin{theorem}[Exponential curvature kernel]
\label{thm:newS7V83-curvature}
There is a volume-independent constant
\begin{equation}
 C_f:=2n_0^2C_P^2S_\mu^2S_{\mu/2}
 \label{eq:newS7V83-Cf}
\end{equation}
such that
\begin{equation}
 |f_{pq}(u_p,v_q)|
 \leq C_fe^{-(\mu/2)d(p,q)}.
 \label{eq:newS7V83-curvature}
\end{equation}
Consequently the curvature row sum satisfies
\begin{equation}
 \sup_q\sum_p|f_{pq}(u_p,v_q)|
 \leq C_fS_{\mu/2}.
 \label{eq:newS7V83-curvature-row}
\end{equation}
\end{theorem}

\begin{proof}
Since \(\|P\|\leq1\),
\[
 |\operatorname{Tr}P(d_pP)(d_qP)|
 \leq\|(d_pP)(d_qP)\|_1.
\]
Expand the product into cell blocks and bound its nuclear norm by the sum
of the block nuclear norms.  Using \eqref{eq:newS7V83-dP} twice gives
\begin{align*}
 \|(d_pP)(d_qP)\|_1
 &\leq n_0^2C_P^2
 \left(\sum_xe^{-\mu d(x,p)}\right)
 \left(\sum_ye^{-\mu d(y,q)}\right)\\
 &\qquad\times
 \sum_ze^{-\mu d(z,p)}e^{-\mu d(z,q)}\\
 &\leq n_0^2C_P^2S_\mu^2S_{\mu/2}
 e^{-(\mu/2)d(p,q)}
\end{align*}
by \Cref{lem:newS7V83-overlap}.  The reversed product has the same bound;
the factor two gives \eqref{eq:newS7V83-Cf}.  Sum
\eqref{eq:newS7V83-curvature} over \(p\) to get
\eqref{eq:newS7V83-curvature-row}.
\end{proof}

\begin{corollary}[Finite Standard Model sum remains local]
\label{cor:newS7V83-SM-sum}
For finitely many Standard Model multiplets whose Hermitian regulator
kernels satisfy the preceding hypotheses with common lower gap and common
locality constants, the total projector curvature
\begin{equation}
 f_{\rm SM}=\sum_f f_f
 \label{eq:newS7V83-fSM}
\end{equation}
has the same exponential form, with \(C_f\) replaced by the sum of the
species constants.  This conclusion uses no anomaly cancellation.
\end{corollary}

\begin{proof}
Apply \Cref{thm:newS7V83-curvature} to each species and use the triangle
inequality.  The number of multiplets is fixed independently of volume.
\end{proof}

\begin{firewall}[Locality is not cancellation]
\label{fire:newS7V83-not-cancel}
V77 cancels continuum anomaly coefficients.  V83 proves that every finite
projector curvature, and therefore their sum, is exponentially local on a
common gapped admissible set.  It does not prove that the finite sum is
zero, exact with a gauge-compatible local counterterm, or small in the
cutoff.  Those are separate PCC rows.
\end{firewall}

\subsection{Loading the radial current}
\label{sec:newS7V83-radial}

Assume the star-shaped chart coordinates obey the cellwise radial bound
\begin{equation}
 \sup_{x\in\mathcal U}\sup_p
 \|x_p-(x_0)_p\|\leq R_0.
 \label{eq:newS7V83-R0}
\end{equation}
Let \(j=Kf_{\rm SM}\) be the V82 radial-gauge connection.

\begin{theorem}[Uniform local radial-connection response]
\label{thm:newS7V83-j}
If a tangent \(v\) is supported in \(X\), then
\begin{equation}
 |j_x(v)|
 \leq {1\over2}R_0C_f^{\rm SM}S_{\mu/2}
 |X|\,\|v\|_\infty,
 \label{eq:newS7V83-j}
\end{equation}
where \(C_f^{\rm SM}\) is the summed species constant.  The coefficient is
uniform in volume.
\end{theorem}

\begin{proof}
Insert the curvature kernel
\eqref{eq:newS7V83-curvature} into the V79 radial-current formula.  The
radial coordinate supplies at most \(R_0\) per cell, the sum over its cell
index supplies \(S_{\mu/2}\), the tangent sum contains only
\(p\in X\), and \(\int_0^1t\,dt=1/2\).
\end{proof}

\begin{corollary}[Analytic GCMCC and OFHC rows]
\label{cor:newS7V83-cards}
Under the common gap, finite-range, local-variation, and graph-row
hypotheses, the curvature-locality row of GCMCC and the radial-current
analytic row of OFHC are populated by
\eqref{eq:newS7V83-curvature} and \eqref{eq:newS7V83-j}.  The LCPLC phase
width still needs the scalar coefficient response and a numerical
small-width inequality.
\end{corollary}

\begin{proof}
These are exactly GCMCC row 2 and the connection part of OFHC rows 1--3.
V82 \eqref{eq:newS7V82-total-response} contains the additional coefficient
term.
\end{proof}

\subsection{Metric and frame dependence}
\label{sec:newS7V83-metric}

\begin{proposition}[Uniform elliptic-collar extension]
\label{prop:newS7V83-metric}
Suppose the Hermitian kernel also depends on discrete metric, vierbein, or
spin-connection variables.  If those variables range in a compact
uniformly elliptic collar on which:
\begin{enumerate}
\item the hopping range and row norm of \(H\) are uniform;
\item the spectral gap is at least \(\delta\); and
\item every unit local geometric variation \(d_pH\) obeys
\eqref{eq:newS7V83-V},
\end{enumerate}
then all conclusions of
\Cref{thm:newS7V83-dP,thm:newS7V83-curvature,thm:newS7V83-j} hold jointly
for gauge and geometric tangents with common constants.
\end{proposition}

\begin{proof}
The proofs use only the listed range, resolvent, gap, variation-support,
and row-sum constants.  Compact uniform ellipticity makes their suprema
finite by hypothesis, so the same estimates apply to either tangent type
and to mixed curvature components.
\end{proof}

\begin{firewall}[A geometric collar has not been constructed]
\label{fire:newS7V83-metric-collar}
\Cref{prop:newS7V83-metric} is a uniform conditional theorem.  The active
programme has not yet shown that the Einstein renewal measure remains in
such a compact elliptic collar, nor that the chiral kernel gap survives on
it.  Those are RCNC and GAOCC inputs.
\end{firewall}

\begin{definition}[Localized projector-curvature card, LPCC]
\label{def:newS7V83-LPCC}
LPCC records, for each regulator and species:
\begin{enumerate}
\item finite hopping range, uniform row norm, and graph-growth row sums;
\item a common admissible spectral gap and explicit resolvent contour;
\item support and nuclear bounds for every local gauge and geometric
kernel derivative;
\item the constants \(C_P,C_f\) in
\eqref{eq:newS7V83-CP} and \eqref{eq:newS7V83-Cf};
\item the summed Standard Model curvature row and radial-current bound;
\item the finite anomaly-cancellation or counterterm identity, separately
from locality; and
\item refinement stability and compatibility with PCC, GCMCC, OFHC,
LCPLC, RCLC, and the Einstein stress response.
\end{enumerate}
\end{definition}

\begin{assessment}[Analytic progress and remaining debit]
\label{ass:newS7V83-status}
The kernel-gap hypothesis now yields a volume-uniform exponentially local
projector curvature and radial connection.  This supplies the analytic
locality mechanism anticipated in V73 and V78.  The next debit is not
decay but cutoff scaling: after summing the chiral Standard Model
representations, determine whether the finite curvature equals a local
counterterm derivative plus a remainder that vanishes uniformly as the
cutoff is removed.
\end{assessment}

\begin{programme}[V84: anomaly-cancelled curvature remainder]
\label{prog:newS7V83-V84}
V84 will formulate a common-regulator local expansion of the summed finite
curvature into representation trace tensors and a controlled remainder.
It will use the exact V77 arithmetic to cancel every universal cubic and
mixed trace term, prove an abstract remainder theorem with explicit
cutoff scaling, and leave the regulator-specific expansion as the single
numerical writer.  This avoids inferring finite cancellation directly from
continuum anomaly cancellation.
\end{programme}

\begin{firewall}[Claim boundary after seventh-series V83]
\label{fire:newS7V83-boundary}
V83 proves two-sided exponential localization and uniform Schatten bounds
for local derivatives of a gapped finite-range spectral projector,
exponential bilocal decay and row summability of its curvature, finite
Standard Model sum locality, and a volume-local radial-connection bound.

V83 does not prove finite Standard Model curvature cancellation, its cutoff
scaling, the scalar coefficient response, a common elliptic gauge--metric
collar, or LPCC, OFHC, LCPLC, CCPC, GCMCC, PCC, SMALC, RCLC, DSLC, CRDC,
GAOCC, CCTC, MSRC, RTGC, or RCNC, remove the cutoff, or construct the
Standard Model--Einstein continuum.
\end{firewall}

\begin{theorem}[Seventh-series V83 result]
\label{thm:newS7V83-result}
For a uniformly gapped finite-range Hermitian regulator kernel with local
cell derivatives, a variation at \(p\) changes the occupied projector by
at most
\(C_Pe^{-\mu d(x,p)}e^{-\mu d(y,p)}\) in the \((x,y)\) block.  Curvature
components at \(p,q\) are at most
\(C_fe^{-(\mu/2)d(p,q)}\), uniformly in volume.  The finite Standard Model
sum remains exponentially local, and its radial-gauge connection has the
local response bound \eqref{eq:newS7V83-j}.
\end{theorem}

\begin{proof}
Combine \Cref{thm:newS7V83-dP,
lem:newS7V83-overlap,
thm:newS7V83-curvature,
cor:newS7V83-SM-sum,
thm:newS7V83-j,
prop:newS7V83-metric}.
\end{proof}

\paragraph{Parent-version record.}
V83 was formed from
\texttt{Renewal\_SM\_Einstein\_Continuum\_Research\_Notes\_V82.tex}.
The V82 parent had \(32\,912\) logical source lines, \(1\,297\,677\) bytes,
and SHA--256
\[
 \texttt{1E8669A4C2D652D0ED92C810F113ADD9F546E758C68BA22CA7123A7DCF489B08}.
\]
The next compulsory companion audit is V85.

% =============================================================================
% END APPENDED SEVENTH-SERIES V83 MATERIAL
% =============================================================================
% =============================================================================
% APPENDED SEVENTH-SERIES V84 MATERIAL
% =============================================================================

\section{Common-regulator anomaly cancellation and curvature remainders}
\label{sec:newS7V84}
\label{sec:v7v84}

V83 proves exponential locality of each finite occupied-projector
curvature.  V77 proves cancellation of the continuum Standard Model
anomaly traces.  These statements combine only if every chiral multiplet
uses a common regulator expansion with the same coefficient kernels.
This section makes that alignment hypothesis explicit and proves the
resulting local remainder theorem.  It also shows how small curvature
remainders force integer Chern periods to vanish on fixed two-cycles, while
leaving flat holonomy open.

\subsection{An exponentially weighted curvature norm}
\label{sec:newS7V84-norm}

For a bilocal two-form kernel \(\omega=(\omega_{pq})\) on a finite graph,
define
\begin{equation}
 \|\omega\|_\mu
 :=\sup_q\sum_p e^{\mu d(p,q)}
 \sup_{\|u_p\|=\|v_q\|=1}|\omega_{pq}(u_p,v_q)|.
 \label{eq:newS7V84-norm}
\end{equation}
The V83 estimate gives a finite norm at every exponent strictly smaller
than its decay exponent.

\begin{lemma}[Basic weighted-norm stability]
\label{lem:newS7V84-norm}
For finite families of kernels and scalars \(c_j\),
\begin{equation}
 \left\|\sum_jc_j\omega_j\right\|_\mu
 \leq\sum_j|c_j|\|\omega_j\|_\mu.
 \label{eq:newS7V84-triangle}
\end{equation}
If \(\|\omega^{(\varepsilon)}\|_\mu\leq C\varepsilon^\kappa\),
then every unweighted row sum and every fixed pair component is
\(O(\varepsilon^\kappa)\), uniformly in volume.
\end{lemma}

\begin{proof}
The first assertion is the triangle inequality before the sums and
suprema.  The weight is at least one, so dropping it bounds the unweighted
row sum; a fixed component is one summand of that row.
\end{proof}

\subsection{The common universal expansion}
\label{sec:newS7V84-expansion}

Let \(\varepsilon>0\) be the cutoff scale.  For every left-handed
multiplet \(s\) in one Standard Model generation, let
\(f_s^{(\varepsilon)}\) be its real finite projector curvature.  Denote by
\begin{equation}
 \mathfrak a_I(s),
 \qquad
 I\in\{333,33Y,22Y,YYY,{\rm gg}Y\},
 \label{eq:newS7V84-anomaly-coeff}
\end{equation}
the five representation coefficients computed in V77.

\begin{definition}[Common-regulator curvature expansion, CRCE]
\label{def:newS7V84-CRCE}
CRCE at exponent \(\mu\) and order \(\kappa>0\) consists of:
\begin{enumerate}
\item species-independent local two-form kernels
\(\Omega_I^{(\varepsilon)}\) for the five anomaly types;
\item local one-form counterterms \(\ell_s^{(\varepsilon)}\);
\item local curvature remainders \(r_s^{(\varepsilon)}\); and
\item the exact identity
\begin{equation}
 f_s^{(\varepsilon)}
 =\sum_I\mathfrak a_I(s)\Omega_I^{(\varepsilon)}
 +d\ell_s^{(\varepsilon)}+r_s^{(\varepsilon)},
 \label{eq:newS7V84-expansion}
\end{equation}
with
\begin{equation}
 \|r_s^{(\varepsilon)}\|_\mu
 \leq C_s\varepsilon^\kappa
 \label{eq:newS7V84-remainder}
\end{equation}
uniformly in volume, admissible boundary data, and every gauge--metric
collar used by the renewal construction.
\end{enumerate}
Pure \(SU(2)^3\) and traces with one simple-group generator are absent
because they vanish representation by representation.
\end{definition}

\begin{theorem}[Anomaly-cancelled finite curvature remainder]
\label{thm:newS7V84-cancel}
Assume CRCE for all multiplets in a complete Standard Model generation.
Put
\begin{equation}
 f_{\rm SM}^{(\varepsilon)}:=\sum_sf_s^{(\varepsilon)},
 \qquad
 \ell_{\rm SM}^{(\varepsilon)}:=\sum_s\ell_s^{(\varepsilon)},
 \qquad
 r_{\rm SM}^{(\varepsilon)}:=\sum_sr_s^{(\varepsilon)}.
 \label{eq:newS7V84-sums}
\end{equation}
Then
\begin{equation}
 \boxed{
 f_{\rm SM}^{(\varepsilon)}
 -d\ell_{\rm SM}^{(\varepsilon)}
 =r_{\rm SM}^{(\varepsilon)},
 \qquad
 \|r_{\rm SM}^{(\varepsilon)}\|_\mu
 \leq C_{\rm SM}\varepsilon^\kappa,}
 \label{eq:newS7V84-cancel}
\end{equation}
where \(C_{\rm SM}=\sum_sC_s\).  The same conclusion holds for any
number of complete generations and neutral singlets.
\end{theorem}

\begin{proof}
Sum \eqref{eq:newS7V84-expansion} over the multiplets.  For every anomaly
type \(I\), V77 proves
\(\sum_s\mathfrak a_I(s)=0\).  Because the kernel
\(\Omega_I^{(\varepsilon)}\) is the same for every species, its total
coefficient is zero.  The counterterms and remainders add as in
\eqref{eq:newS7V84-sums}.  Apply
\Cref{lem:newS7V84-norm} and
\eqref{eq:newS7V84-remainder}.  Complete generations repeat zero, and
neutral singlets have every \(\mathfrak a_I=0\).
\end{proof}

\begin{corollary}[Corrected connection curvature]
\label{cor:newS7V84-corrected}
If \(a_{\rm SM}^{(\varepsilon)}\) is the summed real determinant-line
connection, then the locally corrected connection
\begin{equation}
 \widetilde a_{\rm SM}^{(\varepsilon)}
 :=a_{\rm SM}^{(\varepsilon)}-\ell_{\rm SM}^{(\varepsilon)}
 \label{eq:newS7V84-corrected-connection}
\end{equation}
has curvature \(r_{\rm SM}^{(\varepsilon)}\), which tends to zero in the
weighted local norm at rate \(\varepsilon^\kappa\).
\end{corollary}

\begin{proof}
Differentiate \eqref{eq:newS7V84-corrected-connection} and use
\eqref{eq:newS7V84-cancel}.
\end{proof}

\subsection{Why regulator alignment is indispensable}
\label{sec:newS7V84-alignment}

\begin{countertheorem}[Arithmetic cancellation fails for mismatched
coefficient kernels]
\label{cth:newS7V84-mismatch}
Coefficients \(c_1+c_2=0\) do not imply
\(c_1\Omega_1+c_2\Omega_2=0\) when the two species use different regulator
kernels.  Hence V77 arithmetic alone cannot cancel a collection of
species-dependent finite curvature coefficients.
\end{countertheorem}

\begin{proof}
Take \(c_1=1,c_2=-1\) and any two distinct two-forms
\(\Omega_1\ne\Omega_2\).  Then
\(c_1\Omega_1+c_2\Omega_2=\Omega_1-\Omega_2\ne0\), although
\(c_1+c_2=0\).
\end{proof}

\begin{corollary}[Common action means common coefficient writers]
\label{cor:newS7V84-common-action}
A common finite Standard Model regulator must align Wilson masses,
overlap sign conventions, boundary conditions, admissibility domains,
metric couplings, and subtraction prescriptions sufficiently to prove the
species-independent kernels in \eqref{eq:newS7V84-expansion}.  Merely
placing separate regulators on the same lattice does not establish CRCE.
\end{corollary}

\begin{proof}
Any species dependence in an anomaly coefficient kernel produces exactly
the uncancelled difference in
\Cref{cth:newS7V84-mismatch}.  The listed choices are inputs on which that
kernel can depend, so their alignment or controlled comparison is a
necessary part of the identity.
\end{proof}

\subsection{Radial current of the anomaly-cancelled remainder}
\label{sec:newS7V84-radial}

Let \(K\) be the V79 radial homotopy on a cellwise bounded chart with
\(\|x_p-(x_0)_p\|\leq R_0\).  Define
\begin{equation}
 j_{\rm rem}^{(\varepsilon)}:=Kr_{\rm SM}^{(\varepsilon)}.
 \label{eq:newS7V84-jrem}
\end{equation}

\begin{theorem}[Vanishing local radial response]
\label{thm:newS7V84-jrem}
Under CRCE, if \(v\) is supported in \(X\), then
\begin{equation}
 |j_{\rm rem}^{(\varepsilon)}(v)|
 \leq {1\over2}R_0C_{\rm SM}\varepsilon^\kappa
 |X|\,\|v\|_\infty.
 \label{eq:newS7V84-jrem-bound}
\end{equation}
Here the weighted norm convention absorbs its exponential row sum.  Thus
the anomaly-cancelled part of the radial phase response vanishes locally
at the same cutoff rate.
\end{theorem}

\begin{proof}
Use the radial formula
\[
 (Kr)_x(v)=\int_0^1t\,
 r_{x_0+t(x-x_0)}(x-x_0,v)\,dt.
\]
Expand in cell components.  For each cell supporting \(v\), the weighted
row norm bounds the sum over the radial coordinate cells by
\(C_{\rm SM}\varepsilon^\kappa R_0\).  Sum over \(X\) and integrate
\(t\) to obtain the factor \(1/2\).
\end{proof}

\begin{corollary}[Asymptotic LCPLC width contribution]
\label{cor:newS7V84-width}
For a fixed-size update path of length at most \(\ell_X\), the phase width
contributed by \(j_{\rm rem}^{(\varepsilon)}\) is at most
\begin{equation}
 {1\over2}R_0C_{\rm SM}\varepsilon^\kappa
 |X|\ell_X.
 \label{eq:newS7V84-width}
\end{equation}
It is eventually below any prescribed positive margin.  Counterterm,
scalar-coefficient, divisor, and flat-holonomy contributions remain.
\end{corollary}

\begin{proof}
Integrate \eqref{eq:newS7V84-jrem-bound} along the update path and use the
V82 corrected phase-width loader.
\end{proof}

\subsection{Integer periods from a small remainder}
\label{sec:newS7V84-periods}

\begin{theorem}[Small-curvature quantization forces a fixed Chern period
to vanish]
\label{thm:newS7V84-period}
Let \(\Sigma\) be a fixed closed oriented two-cycle in the common
admissible parameter region.  Suppose the corrected determinant line is
defined on \(\Sigma\), and
\begin{equation}
 {1\over2\pi}\int_\Sigma r_{\rm SM}^{(\varepsilon)}\in\mathbb Z.
 \label{eq:newS7V84-integer}
\end{equation}
If
\begin{equation}
 \left|\int_\Sigma r_{\rm SM}^{(\varepsilon)}\right|<2\pi,
 \label{eq:newS7V84-less-2pi}
\end{equation}
then the period is exactly zero.  In particular, a uniform pointwise
\(O(\varepsilon^\kappa)\) bound on a fixed finite-area cycle implies zero
period for all sufficiently small \(\varepsilon\).
\end{theorem}

\begin{proof}
By \eqref{eq:newS7V84-integer}, the integral is an integer multiple of
\(2\pi\).  The only such multiple with absolute value strictly below
\(2\pi\) is zero.  The final assertion follows when the curvature bound
times the fixed area tends to zero.
\end{proof}

\begin{firewall}[Fixed cycles do not control a growing configuration space]
\label{fire:newS7V84-growing-cycles}
The preceding quantization theorem applies to each fixed two-cycle with a
uniform area and a line defined over it.  Regulator refinement and volume
growth can introduce new cycles or cycles whose geometric area grows.
Their periods require a uniform topological inventory; they cannot be
removed by pointwise smallness alone.
\end{firewall}

\begin{corollary}[Zero Chern periods do not remove flat holonomy]
\label{cor:newS7V84-flat}
Even when every relevant curvature period vanishes, the corrected line may
retain a nontrivial flat holonomy on a noncontractible loop.  CRCE therefore
populates the local curvature row of PCC but not its final holonomy row.
\end{corollary}

\begin{proof}
This is V82 \Cref{cor:newS7V82-flat-holonomy}: curvature and flat holonomy
are independent data.
\end{proof}

\subsection{Gauge--metric uniformity and the actual missing writer}
\label{sec:newS7V84-metric}

\begin{proposition}[Differentiated remainder row]
\label{prop:newS7V84-differentiated}
Suppose CRCE is differentiable in a local metric or vierbein direction
\(h_p\), and
\begin{equation}
 \|d_{h_p}r_s^{(\varepsilon)}\|_\mu
 \leq C_s^{(1)}\varepsilon^\kappa
 \label{eq:newS7V84-stress-remainder}
\end{equation}
uniformly.  Then the anomaly-cancelled curvature stress response satisfies
\begin{equation}
 \left\|d_{h_p}
 \left(f_{\rm SM}^{(\varepsilon)}
 -d\ell_{\rm SM}^{(\varepsilon)}\right)\right\|_\mu
 \leq C_{\rm SM}^{(1)}\varepsilon^\kappa,
 \label{eq:newS7V84-stress-sum}
\end{equation}
where \(C_{\rm SM}^{(1)}=\sum_sC_s^{(1)}\).
\end{proposition}

\begin{proof}
Differentiate the exact identity
\eqref{eq:newS7V84-cancel}, commute parameter differentiation with the
finite exterior derivative, and use the weighted-norm triangle inequality.
\end{proof}

\begin{firewall}[Undifferentiated convergence is insufficient for Einstein
coupling]
\label{fire:newS7V84-stress}
A curvature remainder that is small for each fixed metric need not have a
small metric derivative.  The Einstein source and RCNC require the
separate differentiated row \eqref{eq:newS7V84-stress-remainder}; it is not
implied by \eqref{eq:newS7V84-remainder}.
\end{firewall}

\begin{definition}[Common-regulator anomaly remainder card, CRARC]
\label{def:newS7V84-CRARC}
CRARC requires:
\begin{enumerate}
\item the exact species kernels and a proof of the common universal
expansion \eqref{eq:newS7V84-expansion};
\item a declared local counterterm compatible with gauge symmetry and the
chosen Standard Model global group;
\item the uniform weighted remainder bound
\eqref{eq:newS7V84-remainder};
\item its gauge, metric, boundary, and refinement derivative analogues;
\item a uniform inventory of Chern periods, overlap cocycles, and flat
holonomies;
\item compatibility with RCLC zero-mode sectors and OFHC one-fibre charts;
and
\item insertion of the final phase width into LCPLC and of the stress
response into RCNC.
\end{enumerate}
\end{definition}

\begin{assessment}[What is proved and what must be computed]
\label{ass:newS7V84-status}
If a common regulator supplies CRCE, the exact V77 arithmetic cancels all
universal anomaly kernels and leaves a uniformly local
\(O(\varepsilon^\kappa)\) curvature remainder after a declared
counterterm.  This is a complete abstract compiler.  The missing writer is
now singular and concrete: derive \eqref{eq:newS7V84-expansion} with common
kernels and its differentiated local remainder bounds for the actual
renewal overlap or carrier regulator.
\end{assessment}

\begin{programme}[V85: audit and source-level common-regulator search]
\label{prog:newS7V84-V85}
V85 will perform the compulsory YM and NS audit, then inspect the attached
physical-source and Grand Tensor materials for an actual common chiral
regulator expansion, local counterterm, or differentiated remainder row.
If absent, it will not add another abstract synonym: it will move upstream
to a retained finite Fock carrier and formulate the exact refinement
intertwiner needed before determinant marginalization.
\end{programme}

\begin{firewall}[Claim boundary after seventh-series V84]
\label{fire:newS7V84-boundary}
V84 defines the common universal finite-curvature expansion and proves
that exact Standard Model anomaly arithmetic cancels its universal terms,
leaving a weighted local cutoff remainder; it proves the radial remainder
bound, fixed-cycle integer-period vanishing criterion, and differentiated
stress remainder compiler.

V84 does not derive CRCE for the physical regulator, align its species
kernels, construct its local counterterm, inventory all cycles or flat
holonomies, or prove CRARC, LPCC, OFHC, LCPLC, CCPC, GCMCC, PCC, SMALC,
RCLC, DSLC, CRDC, GAOCC, CCTC, MSRC, RTGC, or RCNC, remove the cutoff, or
construct the Standard Model--Einstein continuum.
\end{firewall}

\begin{theorem}[Seventh-series V84 result]
\label{thm:newS7V84-result}
For a common finite chiral regulator admitting
\eqref{eq:newS7V84-expansion}, the complete Standard Model curvature after
the declared local counterterm is exactly the summed remainder and obeys
\(\|r_{\rm SM}^{(\varepsilon)}\|_\mu\leq
C_{\rm SM}\varepsilon^\kappa\).  Its radial local response vanishes at the
same rate.  Fixed integral Chern periods are eventually zero when their
flux becomes smaller than \(2\pi\), while flat holonomy remains a separate
obstruction.
\end{theorem}

\begin{proof}
Combine \Cref{thm:newS7V84-cancel,
cor:newS7V84-corrected,
cth:newS7V84-mismatch,
thm:newS7V84-jrem,
thm:newS7V84-period,
cor:newS7V84-flat,
prop:newS7V84-differentiated}.
\end{proof}

\paragraph{Parent-version record.}
V84 was formed from
\texttt{Renewal\_SM\_Einstein\_Continuum\_Research\_Notes\_V83.tex}.
The V83 parent had \(33\,317\) logical source lines, \(1\,311\,794\) bytes,
and SHA--256
\[
 \texttt{57287BAADF3679F956F5F7E425DF8660AE0678346AB2700B2A1E65B4A5D08484}.
\]
The next compulsory companion audit is V85.

% =============================================================================
% END APPENDED SEVENTH-SERIES V84 MATERIAL
% =============================================================================
% =============================================================================
% APPENDED SEVENTH-SERIES V85 MATERIAL
% =============================================================================

\section{Companion and source audit; retained-Fock refinement transport}
\label{sec:newS7V85}
\label{sec:v7v85}

This is the compulsory V85 companion audit and the source-level search
promised in V84.  The search finds strong abstract line, holonomy,
counterterm, and projective-limit compilers, but no actual common-regulator
curvature expansion with a differentiated vanishing remainder.  The
physical-source manuscript itself records that its numerical phase and
Berry packet is unpopulated.  Accordingly this note moves upstream to the
retained finite Fock carrier.  V72 already proved positivity and free-shell
factorization at one cutoff; the new result here is an exact refinement
intertwiner and a quantitative trace-norm defect when old and detail
carriers fail to reduce the fine Hamiltonian exactly.

\subsection{Compulsory V85 companion audit}
\label{sec:newS7V85-audit}

At the audit snapshot, the newest distinct complete Yang--Mills checkpoint
remained
\begin{align}
 &\texttt{Renewal\_Yang\_Mills\_Mass\_Gap\_Research\_Notes\_V45.tex},
 \notag\\
 &\qquad 704\,056\ {\rm bytes},\quad 20\,173\ {\rm logical\ lines},\quad
 \texttt{793583EFCD7EBDD7E96D77D5A9E7C6CE78CECDF693BA6196A7C60D5B3AF1E6AC}.
 \label{eq:newS7V85-YM-hash}
\end{align}
The newest Navier--Stokes checkpoint remained
\begin{align}
 &\texttt{Renewal\_NS\_Stability\_Research\_Notes\_V26.tex},
 \notag\\
 &\qquad 278\,283\ {\rm bytes},\quad 5\,319\ {\rm logical\ lines},\quad
 \texttt{82C887F8C2C69E4F28FAD7C3DC8DE5B9099CB5A4BF431EBA1FFF3E91935978AD}.
 \label{eq:newS7V85-NS-hash}
\end{align}
Both have one live document terminator.  Their mathematical lessons are
unchanged from V80: YM V45 exposes conditional covariance created by
marginalization and recommends retained-detail transfer when its mixing
estimate is circular; NS V26 forbids a rank-restricted norm without an
actual invariant rank proof.

\begin{theorem}[V85 companion-audit decision]
\label{thm:newS7V85-audit}
The companion files add no new numerical writer after V80.  YM V45 still
supports the upstream retained-carrier route, while NS V26 supplies no
fermionic refinement or phase estimate.
\end{theorem}

\begin{proof}
The file identities, hashes, and mathematical appendices are the same
distinct checkpoints read at V80.  The stated conclusion is their already
recorded scope.
\end{proof}

\subsection{Attached-source search}
\label{sec:newS7V85-source}

The newest attached physical-source file inspected was
\begin{align}
 &\texttt{Renewal\_Einstein\_Standard\_Model\_Physical\_Source\_Metric\_}
 \notag\\
 &\qquad
 \texttt{Duhamel\_Ancestry\_and\_Quantum\_Continuum\_Programme\_V35.tex},
 \notag\\
 &\qquad 1\,410\,598\ {\rm bytes},\quad 36\,310\ {\rm logical\ lines},
 \notag\\
 &\qquad
 \texttt{AE1FBC0A8125C37465E38E7649E2992725F4579A86F0BB17483B9AE829CFB17C}.
 \label{eq:newS7V85-physical-source}
\end{align}
The Grand Tensor source was
\begin{align}
 &\texttt{Renewal\_Grand\_Tensor\_Monograph\_V067.tex},
 \notag\\
 &\qquad 2\,804\,854\ {\rm bytes},\quad 68\,938\ {\rm logical\ lines},
 \notag\\
 &\qquad
 \texttt{FACD058EC8BB210AD8B296C9D7F4D78A229D08D7A2AE378808700655CB8BC0FE}.
 \label{eq:newS7V85-GT-source}
\end{align}
The line counts here are the direct audit reader's counts; the hashes fix
the exact byte packets.

\begin{proposition}[Exact outcome of the source search]
\label{prop:newS7V85-source}
The physical-source V35 manuscript supplies a protected compact
phase-quenched determinant cylinder, abstract phase visibility and Berry
incidence formulas, and a conditional projective line-landing theorem.  It
states explicitly that no numerical phase visibility, phase Gram, Berry
curvature, or gauge-history crossing is supplied by its corpus.

Grand Tensor V067 supplies finite zero-mode saturation, determinant and
Pfaffian line factorization, spectral-flow and holonomy alternatives, an
exact absolute Wilsonian counterterm cocycle with a conditional quasilocal
tail bound, and a reflection-positive projective-limit theorem under
summable defects and further hypotheses.  Its own status ledger leaves the
regulator section and orientation phase, zero-mode and global-holonomy
data, running and matching, counterterm and Ward convergence, and the
Einstein--matter continuum conditional.

Neither file supplies CRCE \eqref{eq:newS7V84-expansion}, a
species-aligned finite projector-curvature counterterm, or the uniform
cutoff and metric-derivative remainders
\eqref{eq:newS7V84-remainder} and
\eqref{eq:newS7V84-stress-remainder} for the physical chiral regulator.
\end{proposition}

\begin{proof}
This is a formula and claim-boundary audit of the exact files frozen in
\eqref{eq:newS7V85-physical-source}--\eqref{eq:newS7V85-GT-source}.  The
physical-source phase acquisition card terminates with the quoted missing
numerical phase and crossing inputs.  The Grand Tensor quantum section
states its phase transport, counterterm, and continuum conclusions as
conditional theorems and later lists the unpopulated physical data.  A
search for a species-common curvature expansion and a locally weighted
\(O(\varepsilon^\kappa)\) differentiated remainder finds no such writer.
Therefore those abstract results cannot populate CRARC.
\end{proof}

\begin{firewall}[Existing compilers are retained, not reproved]
\label{fire:newS7V85-no-repeat}
This note uses the Grand Tensor projective-limit and counterterm results as
upstream architecture and does not restate them as new Standard Model
population.  The new work below is the concrete finite-Fock
old/detail intertwiner and its explicit Hamiltonian-coupling defect.
\end{firewall}

\subsection{Fermionic Fock factorization across one refinement}
\label{sec:newS7V85-Fock}

Let \(\mathfrak h_n\) be a finite one-particle space at cutoff \(n\), and
let \(i_n:\mathfrak h_n\to\mathfrak h_{n+1}\) be an isometry.  Choose the
orthogonal detail space
\begin{equation}
 \mathfrak d_n:=\mathfrak h_{n+1}\ominus i_n\mathfrak h_n,
 \qquad
 \mathfrak h_{n+1}=i_n\mathfrak h_n\oplus\mathfrak d_n.
 \label{eq:newS7V85-one-particle-split}
\end{equation}
Write
\(\mathcal F(\mathfrak h)=\bigoplus_{k=0}^{\dim\mathfrak h}
\Lambda^k\mathfrak h\) for finite fermionic Fock space.

\begin{lemma}[Canonical graded Fock refinement unitary]
\label{lem:newS7V85-Fock-unitary}
There is a canonical unitary, after fixing the old-before-detail wedge
ordering,
\begin{equation}
 U_n:\mathcal F(\mathfrak h_{n+1})
 \longrightarrow
 \mathcal F(\mathfrak h_n)\widehat\otimes
 \mathcal F(\mathfrak d_n),
 \label{eq:newS7V85-Fock-unitary}
\end{equation}
which sends
\begin{equation}
 i_nu_1\wedge\cdots\wedge i_nu_k
 \wedge v_1\wedge\cdots\wedge v_\ell
 \longmapsto
 (u_1\wedge\cdots\wedge u_k)
 \widehat\otimes(v_1\wedge\cdots\wedge v_\ell).
 \label{eq:newS7V85-wedge-map}
\end{equation}
Even number-preserving operators have the ordinary tensor action under
this identification.
\end{lemma}

\begin{proof}
Exterior powers of an orthogonal direct sum decompose orthogonally as
\[
 \Lambda^m(\mathfrak h_n\oplus\mathfrak d_n)
 \simeq\bigoplus_{k+\ell=m}
 \Lambda^k\mathfrak h_n\otimes\Lambda^\ell\mathfrak d_n.
\]
The fixed wedge ordering determines all Koszul signs and makes
\eqref{eq:newS7V85-wedge-map} an isometry on an orthonormal wedge basis.
Taking the direct sum gives the unitary.  Even operators commute with the
grading sign and therefore act by the ordinary tensor rule.
\end{proof}

Let \(k_n=k_n^*\) be a one-particle Hamiltonian and define its many-body
ground-energy shift by
\begin{equation}
 E_0(k_n):=\sum_{\lambda_j(k_n)<0}\lambda_j(k_n),
 \qquad
 K_n:=d\Gamma(k_n)-E_0(k_n)I\geq0.
 \label{eq:newS7V85-filled-sea}
\end{equation}
The positivity in \eqref{eq:newS7V85-filled-sea} was proved at one cutoff
in V72 and is not reproved here.

\begin{theorem}[Exact reducing refinement factorization]
\label{thm:newS7V85-exact-factorization}
Suppose the fine one-particle Hamiltonian reduces the split
\eqref{eq:newS7V85-one-particle-split} and
\begin{equation}
 k_{n+1}=i_nk_ni_n^*\oplus k_n^{\rm d}.
 \label{eq:newS7V85-reducing}
\end{equation}
Then
\begin{align}
 U_nK_{n+1}U_n^*
 &=K_n\otimes I+I\otimes K_n^{\rm d},
 \label{eq:newS7V85-K-factor}\\
 U_ne^{-\beta K_{n+1}}U_n^*
 &=e^{-\beta K_n}\otimes e^{-\beta K_n^{\rm d}}.
 \label{eq:newS7V85-T-factor}
\end{align}
Consequently the normalized Gibbs carrier state factorizes and is exactly
projective:
\begin{equation}
 U_n\rho_{n+1}U_n^*=\rho_n\otimes\rho_n^{\rm d},
 \qquad
 \operatorname{Tr}_{\mathcal F(\mathfrak d_n)}
 (U_n\rho_{n+1}U_n^*)=\rho_n.
 \label{eq:newS7V85-state-projective}
\end{equation}
Every old observable \(A\) has exactly preserved expectation under the
embedding \(A\mapsto A\otimes I\).
\end{theorem}

\begin{proof}
Second quantization is additive on orthogonal direct sums:
\[
 U_nd\Gamma(k_{n+1})U_n^*
 =d\Gamma(k_n)\otimes I+I\otimes d\Gamma(k_n^{\rm d}).
\]
The negative eigenvalues of a direct sum are the union of the negative
eigenvalues of its blocks, so \(E_0(k_{n+1})=E_0(k_n)+E_0(k_n^{\rm d})\).
This proves \eqref{eq:newS7V85-K-factor}.  The two even tensor summands
commute, giving \eqref{eq:newS7V85-T-factor}.  Traces multiply, so
normalization gives the tensor product density.  Partial trace and the
expectation identity are immediate.
\end{proof}

\subsection{Quantitative defect for a nonreducing fine carrier}
\label{sec:newS7V85-defect}

We first prove a finite Gibbs perturbation bound which does not contain the
Hilbert-space dimension.

\begin{lemma}[Normalized Gibbs trace-distance bound]
\label{lem:newS7V85-Gibbs}
Let \(A=A^*\), \(B=B^*\) on a finite Hilbert space, let
\(\|A-B\|\leq\epsilon\), and define
\[
 \rho_A={e^{-\beta A}\over\operatorname{Tr}e^{-\beta A}},
 \qquad
 \rho_B={e^{-\beta B}\over\operatorname{Tr}e^{-\beta B}}.
\]
Then
\begin{equation}
 \|\rho_A-\rho_B\|_1
 \leq 2\beta\epsilon e^{2\beta\epsilon}.
 \label{eq:newS7V85-Gibbs}
\end{equation}
\end{lemma}

\begin{proof}
Weyl's eigenvalue bound gives
\begin{equation}
 e^{-\beta\epsilon}Z_B\leq Z_A\leq e^{\beta\epsilon}Z_B,
 \qquad Z_A=\operatorname{Tr}e^{-\beta A}.
 \label{eq:newS7V85-Z-ratio}
\end{equation}
Duhamel's formula is
\[
 e^{-\beta A}-e^{-\beta B}
 =-\int_0^\beta e^{-(\beta-s)A}(A-B)e^{-sB}\,ds.
\]
Schatten Holder with exponents
\(p=\beta/(\beta-s)\) and \(q=\beta/s\) gives
\[
 \|e^{-(\beta-s)A}(A-B)e^{-sB}\|_1
 \leq\epsilon Z_A^{1-s/\beta}Z_B^{s/\beta}
 \leq\epsilon\max\{Z_A,Z_B\}.
\]
Thus the unnormalized trace distance is at most
\(\beta\epsilon\max\{Z_A,Z_B\}\).  Split the normalized difference into
an unnormalized difference at one denominator and the partition-function
normalization difference.  The trace bounds and
\eqref{eq:newS7V85-Z-ratio} bound each contribution by
\(\beta\epsilon e^{2\beta\epsilon}\), proving
\eqref{eq:newS7V85-Gibbs}.
\end{proof}

\begin{theorem}[Approximate retained-carrier refinement]
\label{thm:newS7V85-approx}
Let
\begin{equation}
 \widetilde K_{n+1}:=
 U_nK_{n+1}U_n^*,
 \qquad
 K_{n+1}^{0}:=K_n\otimes I+I\otimes K_n^{\rm d},
 \label{eq:newS7V85-K0}
\end{equation}
and suppose
\begin{equation}
 \|\widetilde K_{n+1}-K_{n+1}^{0}\|\leq\epsilon_n.
 \label{eq:newS7V85-epsilon}
\end{equation}
Then the old marginal of the fine normalized carrier state obeys
\begin{equation}
 \left\|
 \operatorname{Tr}_{\mathcal F(\mathfrak d_n)}
 (U_n\rho_{n+1}U_n^*)-\rho_n
 \right\|_1
 \leq 2\beta_n\epsilon_ne^{2\beta_n\epsilon_n}.
 \label{eq:newS7V85-projective-defect}
\end{equation}
The same right side bounds the error in every old observable of operator
norm at most one.
\end{theorem}

\begin{proof}
The Gibbs state of \(K_{n+1}^0\) is
\(\rho_n\otimes\rho_n^{\rm d}\) by exact tensor factorization.  Apply
\Cref{lem:newS7V85-Gibbs} to
\(\widetilde K_{n+1}\) and \(K_{n+1}^0\).  Partial trace is a
trace-norm contraction on Hermitian trace-class differences, which gives
\eqref{eq:newS7V85-projective-defect}.  Trace duality gives the observable
bound.
\end{proof}

\begin{corollary}[One-particle coupling loader]
\label{cor:newS7V85-one-particle}
Suppose before the ground-energy shift the fine one-particle Hamiltonian is
\begin{equation}
 U_n^{(1)}k_{n+1}(U_n^{(1)})^*
 =k_n\oplus k_n^{\rm d}+v_n,
 \label{eq:newS7V85-one-particle-coupling}
\end{equation}
where \(N_{n+1}=\dim\mathfrak h_{n+1}\).  After matching the scalar
ground-energy shift as a declared counterterm, the many-body coupling
obeys
\begin{equation}
 \|d\Gamma(v_n)\|\leq N_{n+1}\|v_n\|.
 \label{eq:newS7V85-dGamma-bound}
\end{equation}
Hence \(\epsilon_n\leq N_{n+1}\|v_n\|\) is a sufficient defect writer.
\end{corollary}

\begin{proof}
On the \(k\)-particle sector,
\(d\Gamma(v_n)\) is the sum of \(k\) copies of \(v_n\), so its norm is at
most \(k\|v_n\|\).  Since \(0\leq k\leq N_{n+1}\), taking the maximum over
Fock sectors gives \eqref{eq:newS7V85-dGamma-bound}.  A scalar energy shift
must be compared separately because its Fock operator norm is its absolute
value.
\end{proof}

\begin{firewall}[The dimension factor is a real ultraviolet debit]
\label{fire:newS7V85-dimension}
The crude but rigorous factor \(N_{n+1}\) in
\eqref{eq:newS7V85-dGamma-bound} can diverge under refinement.  A useful
continuum route therefore needs exact reduction, a particle-number or
energy-weighted sector estimate, locality plus a fixed-observable norm, or
an off-diagonal coupling small enough to beat this factor.  It cannot be
dropped because the one-particle coupling is small.
\end{firewall}

\subsection{Cofinal retained-carrier landing}
\label{sec:newS7V85-landing}

Define the adjacent retained-carrier defect
\begin{equation}
 \delta_n^{\rm F}:=
 2\beta_n\epsilon_ne^{2\beta_n\epsilon_n}.
 \label{eq:newS7V85-deltaF}
\end{equation}

\begin{theorem}[Summable Fock refinement defects]
\label{thm:newS7V85-summable}
If the coarse maps are the declared partial traces and
\begin{equation}
 \sum_{n=1}^\infty\delta_n^{\rm F}<\infty,
 \label{eq:newS7V85-summable}
\end{equation}
then for every fixed cutoff \(m\), the marginals at level \(m\) of the
finer carrier states form a trace-norm Cauchy sequence.  Their limits are
exactly compatible under the partial traces.  Every fixed bounded old
observable has a unique limiting expectation, and its difference from the
level-\(m\) expectation is at most
\begin{equation}
 \sum_{n\geq m}\delta_n^{\rm F}.
 \label{eq:newS7V85-tail}
\end{equation}
\end{theorem}

\begin{proof}
Completely positive trace-preserving partial traces contract trace
distance.  Push the adjacent estimate
\eqref{eq:newS7V85-projective-defect} to level \(m\) and telescope between
any two finer levels.  The tail of
\eqref{eq:newS7V85-summable} makes the sequence Cauchy and gives
\eqref{eq:newS7V85-tail}.  Continuity of partial trace gives exact
compatibility of the limits, and trace duality gives the observable
statement.  This is the retained-Fock specialization of the Grand Tensor
summable projective-defect mechanism.
\end{proof}

\begin{corollary}[Exact branch]
\label{cor:newS7V85-exact-branch}
If every refinement is reducing as in
\eqref{eq:newS7V85-reducing}, then \(\epsilon_n=\delta_n^{\rm F}=0\) and
the retained carrier is exactly projective at every cutoff.
\end{corollary}

\begin{proof}
Use \Cref{thm:newS7V85-exact-factorization}.
\end{proof}

\subsection{Gauge, metric, and interaction compatibility}
\label{sec:newS7V85-compatibility}

\begin{proposition}[Gauge projection commutes with coarse marginal under a
factorized action]
\label{prop:newS7V85-gauge}
Suppose a compact gauge group acts on the refinement split by
\begin{equation}
 \mathcal U_{n+1}(g)
 =\mathcal U_n(g)\widehat\otimes\mathcal U_n^{\rm d}(g),
 \label{eq:newS7V85-gauge-factor}
\end{equation}
and the fine density is gauge invariant.  Then its old partial trace is
gauge invariant.  Haar averaging before or after taking the old marginal
gives the same old state.
\end{proposition}

\begin{proof}
Use covariance of partial trace under product unitaries:
\[
 \operatorname{Tr}_{\rm d}
 \bigl[(\mathcal U_n\otimes\mathcal U_n^{\rm d})
 \rho(\mathcal U_n\otimes\mathcal U_n^{\rm d})^*\bigr]
 =\mathcal U_n(\operatorname{Tr}_{\rm d}\rho)\mathcal U_n^*.
\]
If \(\rho\) is invariant, so is its marginal.  Integrating this identity
against Haar measure proves commutation with gauge averaging.
\end{proof}

\begin{firewall}[Quadratic carrier theorem]
\label{fire:newS7V85-quadratic}
The exact factorization theorem applies to number-preserving quadratic
Fock Hamiltonians with a reducing one-particle split.  A physical
Standard Model carrier also contains gauge interaction, Yukawa mixing,
possibly Majorana or Nambu pairing, and metric dependence.  Those terms
must either be included as even fine-carrier interactions in the defect
\(\epsilon_n\), or treated by a corresponding CAR/quasifree
Bogoliubov refinement theorem.  They are not covered by silently calling
the free shell factorization interacting.
\end{firewall}

\begin{definition}[Retained Fock refinement card, RFRC]
\label{def:newS7V85-RFRC}
For every adjacent cutoff pair, RFRC requires:
\begin{enumerate}
\item the one-particle isometry, old/detail orthogonal split, wedge ordering,
and Fock unitary \eqref{eq:newS7V85-Fock-unitary};
\item the complete gauge--Higgs--Yukawa--metric carrier Hamiltonian and its
ground-energy or vacuum counterterm;
\item an exact reducing factorization or the operator defect
\(\epsilon_n\) in \eqref{eq:newS7V85-epsilon};
\item the summability row \eqref{eq:newS7V85-summable}, uniformly on the
bosonic histories used by the joint transfer;
\item factorized gauge action, Gauss projection, reflection action, and
old-observable embedding;
\item compatible zero-mode saturation, determinant-line and Pfaffian data
when a later marginal is taken;
\item uniform local propagation and metric-stress derivatives of the
intertwining defect; and
\item coupling to the bosonic renewal embeddings and RCNC.
\end{enumerate}
\end{definition}

\begin{assessment}[The upstream writer is now concrete]
\label{ass:newS7V85-status}
The retained-carrier route no longer asks vaguely for projectivity.  It asks
for the explicit operator difference
\[
 U_nK_{n+1}U_n^*-
 (K_n\otimes I+I\otimes K_n^{\rm d})
\]
and the summability of its Gibbs defect.  Exact reduction gives zero debit;
a nonreducing coupling pays
\(2\beta_n\epsilon_ne^{2\beta_n\epsilon_n}\).  The first physical task is
to write the actual overlap/Yukawa/gauge--metric carrier in the same
old/detail Fock coordinates and determine whether its off-diagonal term is
summable or whether a different refinement embedding is required.
\end{assessment}

\begin{programme}[V86: quasifree pairing and covariance transport]
\label{prog:newS7V85-V86}
V86 will extend RFRC from number-preserving \(d\Gamma(k)\) carriers to a
finite Nambu--Bogoliubov quadratic Hamiltonian, needed for chiral pairing
and possible Majorana sectors.  It will express the thermal quasifree state
by its one-particle covariance, prove the exact old-block marginal rule,
and derive an adjacent covariance defect from the matrix functional
calculus.  This is the next upstream layer before adding interacting
Yukawa carrier terms.
\end{programme}

\begin{firewall}[Claim boundary after seventh-series V85]
\label{fire:newS7V85-boundary}
V85 performs the compulsory companion and attached-source audits; proves
that the searched corpus does not populate CRCE; constructs the canonical
old/detail Fock refinement unitary; proves exact reducing transfer and
state projectivity; derives the dimension-free normalized Gibbs
perturbation estimate, the adjacent retained-carrier defect, its
one-particle loader, and summable cofinal landing; and proves gauge-average
compatibility under a factorized action.

V85 does not populate the interacting Standard Model carrier defect, prove
RFRC, CRARC, LPCC, OFHC, LCPLC, CCPC, GCMCC, PCC, SMALC, RCLC, DSLC, CRDC,
GAOCC, CCTC, MSRC, RTGC, or RCNC, remove the cutoff, or construct the
Standard Model--Einstein continuum.
\end{firewall}

\begin{theorem}[Seventh-series V85 result]
\label{thm:newS7V85-result}
A reducing old/detail split of a finite one-particle Hamiltonian induces an
exact tensor factorization of the filled-sea Fock transfer and an exactly
projective old carrier state.  If the complete fine many-body Hamiltonian
differs from the factorized one by operator norm \(\epsilon_n\), its old
marginal differs from the intended old Gibbs carrier by at most
\(2\beta_n\epsilon_ne^{2\beta_n\epsilon_n}\).  Summability of these
defects produces an exactly compatible cofinal retained-carrier family and
unique limits for all fixed bounded old observables.
\end{theorem}

\begin{proof}
Combine \Cref{lem:newS7V85-Fock-unitary,
thm:newS7V85-exact-factorization,
lem:newS7V85-Gibbs,
thm:newS7V85-approx,
cor:newS7V85-one-particle,
thm:newS7V85-summable,
prop:newS7V85-gauge}.
\end{proof}

\paragraph{Parent-version record.}
V85 was formed from
\texttt{Renewal\_SM\_Einstein\_Continuum\_Research\_Notes\_V84.tex}.
The V84 parent had \(33\,736\) logical source lines, \(1\,327\,550\) bytes,
and SHA--256
\[
 \texttt{9FD3065DE8F65202721EF938B849B800C01B9C5894A22D2672DBD7232BBE45E5}.
\]
The next compulsory companion audit is V90.

% =============================================================================
% END APPENDED SEVENTH-SERIES V85 MATERIAL
% =============================================================================
% =============================================================================
% APPENDED SEVENTH-SERIES V86 MATERIAL
% =============================================================================

\section{Nambu quasifree refinement and local covariance transport}
\label{sec:newS7V86}
\label{sec:v7v86}

V85 treated number-preserving filled-sea Fock Hamiltonians.  Chiral
pairing, Majorana sectors, and Bogoliubov descriptions require a self-dual
Nambu carrier.  At finite cutoff its thermal state is still quasifree and
is determined by a Fermi--Dirac covariance.  This section proves that
restriction to the old CAR algebra is exact compression of that covariance,
derives a resolvent bound for nonreducing old/detail couplings, and turns it
into dimension-independent bounds for every fixed-degree local Majorana
correlator.

\subsection{Finite self-dual Nambu data}
\label{sec:newS7V86-Nambu}

Let \(\mathcal K\) be a finite-dimensional Nambu space with an antiunitary
involution \(C\).  A BdG one-particle Hamiltonian is a Hermitian matrix
\(\mathcal H\) satisfying
\begin{equation}
 C\mathcal HC^{-1}=-\mathcal H.
 \label{eq:newS7V86-PH}
\end{equation}
At inverse temperature \(\beta>0\), define
\begin{equation}
 f_\beta(t):={1\over1+e^{\beta t}},
 \qquad
 S_\beta(\mathcal H):=f_\beta(\mathcal H).
 \label{eq:newS7V86-FD}
\end{equation}
Since \(f_\beta(-t)=1-f_\beta(t)\),
\begin{equation}
 CS_\beta(\mathcal H)C^{-1}=I-S_\beta(\mathcal H).
 \label{eq:newS7V86-selfdual-cov}
\end{equation}
This is the self-dual covariance condition.

\begin{lemma}[Thermal quadratic state is determined by the covariance]
\label{lem:newS7V86-quasifree}
The normalized Gibbs state of the finite even quadratic BdG Hamiltonian is
an even quasifree CAR state with covariance
\(S_\beta(\mathcal H)\).  Its odd Majorana moments vanish, and every
\(2r\)-point Majorana moment is the Pfaffian of its two-point contraction
matrix.
\end{lemma}

\begin{proof}
A finite Bogoliubov transformation brings \(\mathcal H\) to particle--hole
pairs with energies \(\pm E_j\).  The many-body Hamiltonian becomes a sum
of commuting terms \(E_j(a_j^*a_j-1/2)\).  Its Gibbs density is a tensor
product of two-level thermal densities with occupation
\((1+e^{\beta E_j})^{-1}\).  Direct expansion gives vanishing odd moments
and Wick factorization in that basis.  Bogoliubov covariance transports
the two-point matrix to \(f_\beta(\mathcal H)\), and Wick Pfaffians are
basis covariant.
\end{proof}

\subsection{Exact old-algebra restriction}
\label{sec:newS7V86-restriction}

Let \(P:\mathcal K_{n+1}\to\mathcal K_n\) be the old Nambu projection.  We
assume \(P\) commutes with the particle--hole involution, so both the old
space and its orthogonal detail are self-dual.

\begin{theorem}[Quasifree marginal compression]
\label{thm:newS7V86-compression}
Let \(\omega_S\) be an even quasifree state on the fine CAR algebra with
covariance \(S\).  Its restriction to the old CAR algebra is quasifree
with covariance
\begin{equation}
 S_{\rm old}=PSP\big|_{\mathcal K_n}.
 \label{eq:newS7V86-compression}
\end{equation}
If
\begin{equation}
 \mathcal H_{n+1}=\mathcal H_n\oplus\mathcal H_n^{\rm d},
 \label{eq:newS7V86-reducing}
\end{equation}
then
\begin{equation}
 PS_\beta(\mathcal H_{n+1})P
 =S_\beta(\mathcal H_n),
 \label{eq:newS7V86-exact-covariance}
\end{equation}
and the thermal quasifree carrier is exactly projective on every old CAR
observable.
\end{theorem}

\begin{proof}
Every old Majorana monomial has all of its test vectors in
\(\operatorname{Ran}P\).  Its fine-state Wick Pfaffian therefore uses only
the compressed two-point contractions \(PSP\).  These are precisely the
Wick moments of the quasifree state with covariance
\eqref{eq:newS7V86-compression}; monomials span the finite CAR algebra, so
the restricted states agree.  Under \eqref{eq:newS7V86-reducing},
functional calculus is block diagonal, giving
\eqref{eq:newS7V86-exact-covariance}.
\end{proof}

\begin{corollary}[Pairing does not obstruct exact projectivity when the
split reduces]
\label{cor:newS7V86-pairing}
Anomalous particle--hole pairing inside the old block and inside the detail
block is allowed in \eqref{eq:newS7V86-reducing}.  Exact projectivity fails
only through terms coupling the two blocks or through incompatible
particle--hole embeddings.
\end{corollary}

\begin{proof}
The proof of \Cref{thm:newS7V86-compression} uses block reduction, not
number conservation within a block.
\end{proof}

\subsection{Analytic functional-calculus defect}
\label{sec:newS7V86-functional}

Let
\begin{equation}
 \mathcal H_{n+1}=\mathcal H_0+V,
 \qquad
 \mathcal H_0:=\mathcal H_n\oplus\mathcal H_n^{\rm d},
 \label{eq:newS7V86-perturbation}
\end{equation}
with \(V=V^*\) preserving the particle--hole symmetry.  Assume
\(\|\mathcal H_{n+1}\|,\|\mathcal H_0\|\leq M\).

Choose a contour \(\Gamma_{\beta,M}\) enclosing \([-M,M]\) in a domain
where \(f_\beta\) is holomorphic.  Let its length be \(L\), its distance
from both spectra be at least \(\eta>0\), and put
\begin{equation}
 C_{\beta,M,\Gamma}
 :={L\over2\pi\eta^2}
 \sup_{z\in\Gamma_{\beta,M}}|f_\beta(z)|.
 \label{eq:newS7V86-C}
\end{equation}
Such a contour exists because the nearest poles of \(f_\beta\) lie at
\((2k+1)i\pi/\beta\).

\begin{theorem}[Fermi--Dirac covariance Lipschitz bound]
\label{thm:newS7V86-Lipschitz}
Under the preceding assumptions,
\begin{equation}
 \|S_\beta(\mathcal H_{n+1})-S_\beta(\mathcal H_0)\|
 \leq C_{\beta,M,\Gamma}\|V\|.
 \label{eq:newS7V86-Lipschitz}
\end{equation}
Consequently the old covariance defect satisfies
\begin{equation}
 \left\|
 PS_\beta(\mathcal H_{n+1})P-S_\beta(\mathcal H_n)
 \right\|
 \leq C_{\beta,M,\Gamma}\|V\|.
 \label{eq:newS7V86-old-defect}
\end{equation}
\end{theorem}

\begin{proof}
Holomorphic functional calculus and the resolvent identity give
\begin{align*}
 f_\beta(\mathcal H_{n+1})-f_\beta(\mathcal H_0)
 ={1\over2\pi i}\oint_{\Gamma_{\beta,M}}
 &f_\beta(z)(z-\mathcal H_{n+1})^{-1}V\\
 &\times(z-\mathcal H_0)^{-1}\,dz.
\end{align*}
Each resolvent has norm at most \(\eta^{-1}\).  Taking norms and integrating
proves \eqref{eq:newS7V86-Lipschitz}.  Compression is contractive, and the
old block of \(f_\beta(\mathcal H_0)\) is
\(f_\beta(\mathcal H_n)\), proving
\eqref{eq:newS7V86-old-defect}.
\end{proof}

\begin{remark}[Temperature dependence]
\label{rem:newS7V86-temperature}
The constant deteriorates when a contour approaches the Fermi--Dirac poles,
and its useful scaling must be checked with the physical \(\beta_n\).  The
abstract fact that \(0\leq S_\beta\leq I\) does not make the functional
calculus uniformly Lipschitz in a low-temperature limit.
\end{remark}

\subsection{Fixed-degree correlation bounds without a Fock dimension}
\label{sec:newS7V86-correlations}

Let \(c(u)\) denote normalized Majorana fields with
\(\|u\|=1\).  Let \(\omega\) and \(\widetilde\omega\) be even quasifree
states whose two-point contractions differ in operator norm by at most
\(\delta\).

\begin{lemma}[Pfaffian product-difference bound]
\label{lem:newS7V86-Pfaffian}
For unit vectors \(u_1,\ldots,u_{2r}\),
\begin{equation}
 \left|
 \omega(c(u_1)\cdots c(u_{2r}))
 -\widetilde\omega(c(u_1)\cdots c(u_{2r}))
 \right|
 \leq r(2r-1)!!\,\delta.
 \label{eq:newS7V86-Pfaffian}
\end{equation}
The constant depends on the observable degree \(2r\), not on the total
number of fine modes.
\end{lemma}

\begin{proof}
Wick's theorem expands each moment into \((2r-1)!!\) signed pairings.  Every
pairing is a product of \(r\) two-point contractions, each of modulus at
most one.  For products \(\prod_{j=1}^ra_j\) and
\(\prod_{j=1}^r\widetilde a_j\) with all factors of modulus at most one,
telescoping one factor at a time gives an absolute difference at most
\(r\max_j|a_j-\widetilde a_j|\).  Each contraction difference is at most
\(\delta\).  Sum over pairings.
\end{proof}

\begin{theorem}[Local quasifree refinement defect]
\label{thm:newS7V86-local-defect}
For the fine state of \(\mathcal H_{n+1}=\mathcal H_0+V\) and the intended
old state of \(\mathcal H_n\), every old normalized Majorana monomial of
degree \(2r\) obeys
\begin{equation}
 |\omega_{n+1}(O_{2r})-\omega_n(O_{2r})|
 \leq r(2r-1)!!\,
 C_{\beta,M,\Gamma}\|V\|.
 \label{eq:newS7V86-local-defect}
\end{equation}
Odd moments vanish in both states.  Finite linear combinations inherit the
corresponding coefficient-weighted bound.
\end{theorem}

\begin{proof}
The old covariance difference is bounded by
\eqref{eq:newS7V86-old-defect}.  Apply
\Cref{lem:newS7V86-Pfaffian}.  Even quasifree states have zero odd moments,
and linearity handles finite combinations.
\end{proof}

\begin{corollary}[Improvement over full-state trace distance]
\label{cor:newS7V86-improvement}
For any fixed local polynomial observable, the V86 defect contains no
factor \(\dim\mathfrak h_{n+1}\).  It can therefore be summable even when
the full Fock-state trace-distance loader in V85 is too crude because its
many-body coupling norm grows with the number of modes.
\end{corollary}

\begin{proof}
The right side of \eqref{eq:newS7V86-local-defect} depends only on the fixed
degree, the functional-calculus constant, and \(\|V\|\).  Compare with the
V85 one-particle loader \(N_{n+1}\|v_n\|\).
\end{proof}

\subsection{Cofinal local CAR projectivity}
\label{sec:newS7V86-cofinal}

For each refinement let
\begin{equation}
 \epsilon_n^{\rm cov}
 :=C_{\beta_n,M_n,\Gamma_n}\|V_n\|.
 \label{eq:newS7V86-eps-cov}
\end{equation}

\begin{theorem}[Summable covariance-defect landing]
\label{thm:newS7V86-landing}
If
\begin{equation}
 \sum_n\epsilon_n^{\rm cov}<\infty,
 \label{eq:newS7V86-cov-sum}
\end{equation}
then every fixed finite-degree old CAR polynomial has a unique cofinal
expectation along the declared embeddings.  For a degree-\(2r\) normalized
Majorana monomial, the tail from level \(m\) is at most
\begin{equation}
 r(2r-1)!!\sum_{n\geq m}\epsilon_n^{\rm cov}.
 \label{eq:newS7V86-cov-tail}
\end{equation}
The limiting moments obey the CAR relations and Wick consistency on the
quasilocal union of the finite old algebras.
\end{theorem}

\begin{proof}
Apply \eqref{eq:newS7V86-local-defect} at each adjacent refinement and
telescope.  The summable tail makes every fixed moment Cauchy and gives
\eqref{eq:newS7V86-cov-tail}.  Each finite CAR relation and Wick identity
is a polynomial equality among finitely many moments; pass it to the
limit.
\end{proof}

\begin{firewall}[Moment projectivity is not full-state projectivity]
\label{fire:newS7V86-moments}
The theorem controls the quasilocal algebra one fixed polynomial at a time.
It does not give trace-norm convergence of the full old density matrices
when the old algebra dimension grows, nor does it prove tightness of
unbounded fermionic composites.  Those need uniform norm, energy, or
noncommutative \(L^p\) estimates.
\end{firewall}

\subsection{Compatibility rows and the next interaction}
\label{sec:newS7V86-compatibility}

\begin{proposition}[Symmetry-covariant compression]
\label{prop:newS7V86-symmetry}
Suppose a gauge or reflection symmetry is implemented on Nambu space by an
operator commuting with the old projection \(P\), and both fine and old
BdG Hamiltonians transform covariantly.  Then the compressed covariance
and every limiting old Wick moment transform covariantly under that
symmetry.
\end{proposition}

\begin{proof}
Functional calculus commutes with unitary conjugation.  Commutation with
\(P\) passes the conjugation through the compression.  Wick Pfaffians of
the transformed two-point matrix give the transformed moments; cofinal
limits preserve the finite identities.
\end{proof}

\begin{definition}[Nambu covariance refinement card, NCRC]
\label{def:newS7V86-NCRC}
NCRC requires:
\begin{enumerate}
\item self-dual old/detail embeddings commuting with particle--hole,
gauge, and reflection structures;
\item the complete finite BdG matrix and the norm bounds \(M_n\);
\item the coupling block \(V_n\), a valid holomorphic contour, and the
constant \(C_{\beta_n,M_n,\Gamma_n}\);
\item either exact reduction or the summability row
\eqref{eq:newS7V86-cov-sum};
\item a declared local CAR observable core and any stronger full-state or
energy topology needed later;
\item compatible zero-mode, Pfaffian orientation, RCLC, and anomaly-line
phase data; and
\item uniform dependence on the gauge--Higgs--metric histories and their
stress variations.
\end{enumerate}
\end{definition}

\begin{assessment}[What Nambu transport resolves]
\label{ass:newS7V86-status}
Pairing and Majorana doubling do not obstruct refinement when the Nambu
old/detail split reduces the BdG Hamiltonian.  A nonreducing quadratic
coupling produces the explicit covariance defect
\(C_{\beta,M,\Gamma}\|V\|\), and fixed local correlations pay only a
degree constant.  The next missing layer is the nonquadratic carrier:
Yukawa and gauge-dependent interactions can change the state away from
quasifree Wick form and require a local Duhamel or conditional-mixing
estimate rather than a covariance matrix alone.
\end{assessment}

\begin{programme}[V87: interacting retained-carrier Duhamel response]
\label{prog:newS7V86-V87}
V87 will compare a factorized old/detail carrier with an even local
interaction by an imaginary-time Duhamel expansion.  It will derive an
exact derivative formula for an old observable, express the refinement
defect as a connected thermal correlation with the coupling, and prove a
summable local bound under a declared imaginary-time clustering kernel.
This is the retained-carrier analogue of the YM V45 upstream covariance
route, without taking a determinant logarithm.
\end{programme}

\begin{firewall}[Claim boundary after seventh-series V86]
\label{fire:newS7V86-boundary}
V86 proves exact self-dual quasifree marginal compression, exact reducing
BdG projectivity, a holomorphic Fermi--Dirac covariance perturbation bound,
fixed-degree Pfaffian correlation estimates without a total-mode factor,
summable local CAR landing, and symmetry-covariant compression.

V86 does not control nonquadratic Yukawa or gauge carrier interactions,
full-state trace distance in a growing algebra, unbounded composites, or
prove NCRC, RFRC, CRARC, LPCC, OFHC, LCPLC, CCPC, GCMCC, PCC, SMALC, RCLC,
DSLC, CRDC, GAOCC, CCTC, MSRC, RTGC, or RCNC, remove the cutoff, or
construct the Standard Model--Einstein continuum.
\end{firewall}

\begin{theorem}[Seventh-series V86 result]
\label{thm:newS7V86-result}
A finite thermal BdG carrier restricts to the old CAR algebra by compressing
its Fermi--Dirac covariance.  A reducing Nambu refinement is exactly
projective.  If the fine BdG matrix differs from the old/detail direct sum
by \(V_n\), the old covariance defect is at most
\(C_{\beta_n,M_n,\Gamma_n}\|V_n\|\), and every fixed degree-\(2r\)
Majorana moment differs by at most
\(r(2r-1)!!C_{\beta_n,M_n,\Gamma_n}\|V_n\|\).  Summability of these
covariance defects gives a unique cofinal quasilocal CAR moment family.
\end{theorem}

\begin{proof}
Combine \Cref{lem:newS7V86-quasifree,
thm:newS7V86-compression,
thm:newS7V86-Lipschitz,
lem:newS7V86-Pfaffian,
thm:newS7V86-local-defect,
thm:newS7V86-landing,
prop:newS7V86-symmetry}.
\end{proof}

\paragraph{Parent-version record.}
V86 was formed from
\texttt{Renewal\_SM\_Einstein\_Continuum\_Research\_Notes\_V85.tex}.
The V85 parent had \(34\,293\) logical source lines, \(1\,349\,056\) bytes,
and SHA--256
\[
 \texttt{F82D3621B5120FFE4FAB8F8ABE2DEE56A2539E1791C76B9FF35BBA3C23DE3C0A}.
\]
The next compulsory companion audit is V90.

% =============================================================================
% END APPENDED SEVENTH-SERIES V86 MATERIAL
% =============================================================================
% =============================================================================
% APPENDED SEVENTH-SERIES V87 MATERIAL
% =============================================================================

\section{Interacting retained carriers and Duhamel refinement response}
\label{sec:newS7V87}
\label{sec:v7v87}

V86 controls quadratic old/detail coupling through quasifree covariances.
A Yukawa interaction or a genuinely interacting gauge-dependent carrier
need not remain quasifree.  The correct positive-carrier response is the
Kubo--Mori, or imaginary-time Duhamel, covariance.  This section derives
the exact response identity, applies it to one refinement, and proves a
volume-local projectivity bound under a same-carrier imaginary-time
clustering kernel.  No determinant logarithm and no inverse complex phase
margin occur.

\subsection{Duhamel covariance}
\label{sec:newS7V87-Duhamel}

Let \(K(\lambda)=K_0+\lambda V\) be a differentiable finite-dimensional
self-adjoint Hamiltonian family, \(0\leq\lambda\leq1\), and let
\begin{equation}
 \rho_\lambda={e^{-\beta K(\lambda)}\over Z_\lambda},
 \qquad
 Z_\lambda=\operatorname{Tr}e^{-\beta K(\lambda)},
 \qquad
 \omega_\lambda(A)=\operatorname{Tr}(\rho_\lambda A).
 \label{eq:newS7V87-Gibbs}
\end{equation}
For operators \(A,B\), define
\begin{align}
 \langle A;B\rangle_{\rho}
 &:=\int_0^1
 \operatorname{Tr}(\rho^{1-s}A\rho^sB)\,ds,
 \label{eq:newS7V87-Duhamel-product}\\
 \operatorname{Cov}_{\rho}^{\rm D}(A,B)
 &:=\langle A;B\rangle_\rho-omega_\rho(A)\omega_\rho(B).
 \label{eq:newS7V87-Duhamel-cov}
\end{align}

\begin{lemma}[Duhamel derivative of the exponential]
\label{lem:newS7V87-exp-derivative}
For a differentiable self-adjoint matrix family,
\begin{equation}
 {d\over d\lambda}e^{-\beta K(\lambda)}
 =-\beta\int_0^1
 e^{-\beta(1-s)K(\lambda)}\dot K(\lambda)
 e^{-\beta sK(\lambda)}\,ds.
 \label{eq:newS7V87-exp-derivative}
\end{equation}
\end{lemma}

\begin{proof}
Apply the fundamental theorem of calculus to
\(e^{-(1-s)\beta K(\lambda+h)}e^{-s\beta K(\lambda)}\), divide by \(h\),
and pass to the finite-dimensional norm limit.  Equivalently, expand both
exponentials in convergent power series and collect the position of the
single differentiated factor.  The integral over \(s\) gives
\eqref{eq:newS7V87-exp-derivative}.
\end{proof}

\begin{theorem}[Exact interacting Gibbs response]
\label{thm:newS7V87-response}
For a differentiable observable \(A(\lambda)\),
\begin{equation}
 {d\over d\lambda}\omega_\lambda(A)
 =\omega_\lambda(\dot A)
 -\beta\operatorname{Cov}^{\rm D}_{\rho_\lambda}(A,V).
 \label{eq:newS7V87-response}
\end{equation}
In particular, for a fixed \(A\),
\begin{equation}
 \omega_1(A)-\omega_0(A)
 =-\beta\int_0^1
 \operatorname{Cov}^{\rm D}_{\rho_\lambda}(A,V)\,d\lambda.
 \label{eq:newS7V87-integrated-response}
\end{equation}
\end{theorem}

\begin{proof}
Differentiate the numerator and denominator in
\eqref{eq:newS7V87-Gibbs} using
\Cref{lem:newS7V87-exp-derivative}.  Cyclicity of the trace turns the
numerator derivative into
\(-\beta Z_\lambda\langle A;V\rangle_{\rho_\lambda}\), plus
\(Z_\lambda\omega_\lambda(\dot A)\).  The partition derivative is
\(-\beta Z_\lambda\omega_\lambda(V)\).  The quotient rule gives
\eqref{eq:newS7V87-response}; integrate it for fixed \(A\).
\end{proof}

\begin{proposition}[Positivity and a dimension-free fallback]
\label{prop:newS7V87-positive}
For every operator \(A\),
\begin{equation}
 \langle A^*;A\rangle_\rho\geq0.
 \label{eq:newS7V87-positive}
\end{equation}
Moreover,
\begin{equation}
 |\operatorname{Cov}^{\rm D}_\rho(A,B)|
 \leq2\|A\|\|B\|.
 \label{eq:newS7V87-trivial}
\end{equation}
Consequently
\begin{equation}
 |\omega_1(A)-\omega_0(A)|
 \leq2\beta\|A\|\|V\|.
 \label{eq:newS7V87-global-bound}
\end{equation}
\end{proposition}

\begin{proof}
Diagonalize \(\rho\).  The Duhamel product of \(A^*\) and \(A\) is a sum
of \(|A_{jk}|^2\) times the nonnegative logarithmic mean of the two
eigenvalues of \(\rho\), proving positivity.  Schatten Holder gives
\[
 |\operatorname{Tr}(\rho^{1-s}A\rho^sB)|
 \leq\|A\|\|B\|
\]
for every \(s\); the product of the two expectations has the same bound.
This proves \eqref{eq:newS7V87-trivial}.  Insert it in
\eqref{eq:newS7V87-integrated-response}.
\end{proof}

\begin{firewall}[Positivity of the Duhamel form is not spatial decay]
\label{fire:newS7V87-no-decay}
Equation \eqref{eq:newS7V87-positive} gives a positive thermal response
metric.  It does not imply that the covariance of distant observables is
small.  The global bound \eqref{eq:newS7V87-global-bound} sees the full
coupling norm and is generally insufficient under refinement.
\end{firewall}

\subsection{One interacting old/detail refinement}
\label{sec:newS7V87-refinement}

Use the V85 Fock refinement unitary and put
\begin{equation}
 K_0=K_n\otimes I+I\otimes K_n^{\rm d},
 \qquad
 K(\lambda)=K_0+\lambda V_n.
 \label{eq:newS7V87-refinement}
\end{equation}
The reference Gibbs state factorizes as
\(\rho_0=\rho_n\otimes\rho_n^{\rm d}\).  Let \(A\) be an old observable,
represented as \(A\otimes I\).

\begin{theorem}[Exact retained-carrier refinement defect]
\label{thm:newS7V87-refinement}
The difference between the interacting fine expectation and the intended
old expectation is
\begin{equation}
 \boxed{
 \omega_{n+1}(A\otimes I)-\omega_n(A)
 =-\beta_n\int_0^1
 \operatorname{Cov}^{\rm D}_{\rho_{n,\lambda}}
 (A\otimes I,V_n)\,d\lambda.}
 \label{eq:newS7V87-refinement-defect}
\end{equation}
At \(\lambda=0\), every purely detail term has zero connected covariance
with \(A\otimes I\).  Thus only old terms and genuine old/detail couplings
can initiate the refinement defect.
\end{theorem}

\begin{proof}
The first identity is
\Cref{thm:newS7V87-response} applied to
\eqref{eq:newS7V87-refinement}.  At \(\lambda=0\), factorization gives
\[
 \langle A\otimes I;I\otimes B\rangle_{\rho_0}
 =\omega_n(A)\omega_n^{\rm d}(B),
\]
so the connected Duhamel covariance is zero.
\end{proof}

\begin{corollary}[No determinant phase denominator]
\label{cor:newS7V87-no-phase}
The right side of \eqref{eq:newS7V87-refinement-defect} is evaluated in
positive density matrices \(\rho_{n,\lambda}\).  It contains no factor
\(\rho_\nu^{-1}\) from the V80 complex phase margin.  Determinant-line
phase and anomaly data remain relevant to gauge scalarization, but they do
not condition this retained-carrier response formula.
\end{corollary}

\begin{proof}
Every \(K(\lambda)\) is self-adjoint, so its exponential has positive
trace and defines the normalized positive state
\eqref{eq:newS7V87-Gibbs}.  No complex partition quotient is taken.
\end{proof}

\subsection{Local imaginary-time clustering loader}
\label{sec:newS7V87-clustering}

Let the carrier algebra be indexed by a metric graph.  Write the even
interaction as an absolutely convergent finite-volume sum
\begin{equation}
 V_n=\sum_{X\Subset\Lambda_{n+1}}V_{n,X}.
 \label{eq:newS7V87-interaction}
\end{equation}
For a local observable \(A\), let \(B=\operatorname{supp}A\) and
\(d(B,X)=\min_{b\in B,x\in X}d(b,x)\).

\begin{definition}[Imaginary-time carrier mixing bound]
\label{def:newS7V87-ITCM}
ITCM with constants \(C_{\rm D},\mu>0\) means that, uniformly in
\(n,\lambda\in[0,1]\), admissible bosonic histories, and boundary data,
\begin{equation}
 \left|
 \operatorname{Cov}^{\rm D}_{\rho_{n,\lambda}}(A,V_{n,X})
 \right|
 \leq C_{\rm D}\|A\|\|V_{n,X}\|
 |B|\,|X|e^{-\mu d(B,X)}
 \label{eq:newS7V87-ITCM}
\end{equation}
for the declared local old-observable core.
\end{definition}

\begin{theorem}[Local interacting refinement bound]
\label{thm:newS7V87-local}
Under ITCM and absolute convergence of
\eqref{eq:newS7V87-interaction},
\begin{equation}
 |\omega_{n+1}(A)-\omega_n(A)|
 \leq \beta_nC_{\rm D}\|A\||B|\,
 \mathcal V_n(B),
 \label{eq:newS7V87-local}
\end{equation}
where
\begin{equation}
 \mathcal V_n(B):=
 \sum_X|X|\|V_{n,X}\|e^{-\mu d(B,X)}.
 \label{eq:newS7V87-local-stock}
\end{equation}
The bound is independent of total volume whenever the stock is.
\end{theorem}

\begin{proof}
Insert the interaction sum in
\eqref{eq:newS7V87-refinement-defect}.  Absolute convergence permits
termwise covariance and integration.  Apply
\eqref{eq:newS7V87-ITCM}, sum over \(X\), and integrate
\(\lambda\in[0,1]\).
\end{proof}

\begin{corollary}[Summable local interacting projectivity]
\label{cor:newS7V87-summable}
If for every fixed old support \(B\),
\begin{equation}
 \sum_{n\geq n(B)}\beta_n\mathcal V_n(B)<\infty,
 \label{eq:newS7V87-summable}
\end{equation}
then every fixed bounded observable in the declared local carrier core has
a unique cofinal expectation.  Its tail is bounded by
\begin{equation}
 C_{\rm D}\|A\||B|
 \sum_{n\geq m}\beta_n\mathcal V_n(B).
 \label{eq:newS7V87-tail}
\end{equation}
\end{corollary}

\begin{proof}
Apply \Cref{thm:newS7V87-local} at adjacent refinements and telescope.
The summability hypothesis makes the expectations Cauchy and yields the
tail bound.
\end{proof}

\begin{remark}[Old-block renormalization]
\label{rem:newS7V87-old-terms}
If \(V_n\) contains a retained old-only counterterm of nonvanishing size,
its distance from \(B\) is not large and it need not be summable.  Such a
term must be absorbed into the definition of the renormalized old
Hamiltonian \(K_n\).  The stock \(\mathcal V_n(B)\) is reserved for the
residual mismatch after this matching.
\end{remark}

\subsection{Metric response and the Einstein source}
\label{sec:newS7V87-metric}

Let \(q\) be a local metric or frame parameter and allow both the
Hamiltonian and observable to depend on it.

\begin{theorem}[Retained-carrier stress response]
\label{thm:newS7V87-stress}
At every finite cutoff,
\begin{equation}
 D_q\omega(A)
 =\omega(D_qA)
 -\beta\operatorname{Cov}^{\rm D}_\rho(A,D_qK).
 \label{eq:newS7V87-stress}
\end{equation}
If \(D_qK=\sum_XW_{q,X}\) satisfies ITCM with the same uniform constants,
then for \(q\) supported at \(p\),
\begin{equation}
 |D_q\omega(A)-\omega(D_qA)|
 \leq\beta C_{\rm D}\|A\||B|
 \sum_X|X|\|W_{q,X}\|e^{-\mu d(B,X)}.
 \label{eq:newS7V87-stress-local}
\end{equation}
\end{theorem}

\begin{proof}
Apply \Cref{thm:newS7V87-response} with \(q\) as the parameter and
\(V=D_qK\).  The local bound is the same summation as in
\Cref{thm:newS7V87-local}.
\end{proof}

\begin{corollary}[Carrier source is not a determinant Hessian]
\label{cor:newS7V87-not-det}
The fermionic Einstein response on the retained carrier is a positive-state
Duhamel covariance with \(D_qK\).  It need not have the sign of the Hessian
of \(-\log|\det D|\), which V71 showed can be indefinite.  Marginal and
carrier source rows must therefore remain distinct.
\end{corollary}

\begin{proof}
Equation \eqref{eq:newS7V87-stress} follows from a positive Gibbs state
before marginalization.  The determinant Hessian follows after eliminating
the carrier and contains a different covariance subtraction.  Equality
would require an additional exact marginal identity.
\end{proof}

\begin{definition}[Interacting retained-carrier mixing card, IRCMC]
\label{def:newS7V87-IRCMC}
IRCMC requires:
\begin{enumerate}
\item the RFRC old/detail carrier and a renormalized old Hamiltonian with
all persistent old terms matched;
\item a support-local decomposition of the residual even interaction;
\item ITCM \eqref{eq:newS7V87-ITCM} uniformly along the interpolation and
on the common gauge--Higgs--metric collar;
\item the summable local stocks \eqref{eq:newS7V87-summable} for a dense
bounded old-observable core;
\item the differentiated stress, current, gauge, and boundary analogues;
\item compatibility with Gauss projection, reflection, CAR grading,
zero-mode and Pfaffian sectors; and
\item a proof that the limiting local carrier states and bosonic renewal
marginals live on one projective history system.
\end{enumerate}
\end{definition}

\begin{assessment}[The interacting bottleneck is a same-carrier covariance]
\label{ass:newS7V87-status}
The retained-carrier refinement defect is now an exact connected
imaginary-time covariance, not an unspecified interaction error.  A uniform
ITCM kernel loads a summable local projective limit and its metric response
without any complex normalization.  The first missing estimate is therefore
the decay of these same-carrier Duhamel correlations for the actual
interpolating Standard Model Hamiltonian.  Proving it from the desired
continuum mass gap would be circular; the next note must seek a finite
high-temperature, finite-time, or local-convexity entrance.
\end{assessment}

\begin{programme}[V88: finite-time Duhamel locality without a mass-gap
assumption]
\label{prog:newS7V87-V88}
V88 will use a finite-range Hamiltonian and a Lieb--Robinson bound to
control imaginary-time response over a bounded \(\beta\) interval.  It
will separate the exponential spatial tail from the growth in imaginary
time and derive an explicit ITCM candidate.  If analytic continuation to
imaginary time loses volume-uniform control, the exact failure will be
recorded and the route will move further upstream to a Trotterized positive
slab with a classical dependency graph.
\end{programme}

\begin{firewall}[Claim boundary after seventh-series V87]
\label{fire:newS7V87-boundary}
V87 proves the exact Duhamel Gibbs response, its positivity and global
bound, the interacting retained-carrier refinement identity, a local
projectivity theorem under ITCM, summable cofinal landing, and the matching
retained-carrier metric-stress response.

V87 does not prove ITCM for the physical Standard Model carrier, control
unbounded bosonic histories, or prove IRCMC, NCRC, RFRC, CRARC, LPCC, OFHC,
LCPLC, CCPC, GCMCC, PCC, SMALC, RCLC, DSLC, CRDC, GAOCC, CCTC, MSRC,
RTGC, or RCNC, remove the cutoff, or construct the Standard Model--Einstein
continuum.
\end{firewall}

\begin{theorem}[Seventh-series V87 result]
\label{thm:newS7V87-result}
For an interacting retained carrier interpolating from the factorized
old/detail Gibbs state, the exact old-observable refinement defect is the
integral \eqref{eq:newS7V87-refinement-defect} of a same-carrier Duhamel
covariance.  If these covariances decay as
\eqref{eq:newS7V87-ITCM}, the defect is bounded by the local stock
\eqref{eq:newS7V87-local-stock}; summability gives a unique cofinal local
state.  Metric response obeys the same Duhamel formula and local compiler.
\end{theorem}

\begin{proof}
Combine \Cref{lem:newS7V87-exp-derivative,
thm:newS7V87-response,
prop:newS7V87-positive,
thm:newS7V87-refinement,
thm:newS7V87-local,
cor:newS7V87-summable,
thm:newS7V87-stress}.
\end{proof}

\paragraph{Parent-version record.}
V87 was formed from
\texttt{Renewal\_SM\_Einstein\_Continuum\_Research\_Notes\_V86.tex}.
The V86 parent had \(34\,701\) logical source lines, \(1\,364\,739\) bytes,
and SHA--256
\[
 \texttt{0C38606545B25907342A989280033E81FFDFB1B53A800CB0758ED42C4BA49831}.
\]
The next compulsory companion audit is V90.

% =============================================================================
% END APPENDED SEVENTH-SERIES V87 MATERIAL
% =============================================================================
% =============================================================================
% APPENDED SEVENTH-SERIES V88 MATERIAL
% =============================================================================

\section{Finite imaginary-time locality and the missing thermal-mixing row}
\label{sec:newS7V88}
\label{sec:v7v88}

V87 reduces interacting carrier projectivity to decay of Duhamel
covariances.  A finite-range Hamiltonian controls how far an observable can
spread in imaginary time, but this kinematic statement does not by itself
make the Gibbs state spatially mixing.  This section proves an elementary
nested-commutator quasi-locality theorem in an explicit small-complex-time
window, shows how equal-time Gibbs mixing would then load ITCM, and gives a
one-dimensional low-temperature counterexample to any uniform inference
from finite range and a Hamiltonian gap alone.

\subsection{Interaction and graph hypotheses}
\label{sec:newS7V88-hypotheses}

Let \(\Lambda\) be a finite graph and
\begin{equation}
 K=\sum_{Z\Subset\Lambda}\Phi_Z,
 \qquad \Phi_Z=\Phi_Z^*,
 \label{eq:newS7V88-H}
\end{equation}
where
\begin{equation}
 \operatorname{diam}Z\leq R,
 \qquad |Z|\leq q,
 \qquad
 J:=\sup_{x\in\Lambda}\sum_{Z\ni x}\|\Phi_Z\|<\infty.
 \label{eq:newS7V88-local-interaction}
\end{equation}
Assume \(q\geq2\); one-site terms can be included separately without
support growth.  For a local observable \(A\) with support \(B\), put
\(b=|B|\) and
\begin{equation}
 \tau_s(A):=e^{sK}Ae^{-sK},
 \qquad s\geq0.
 \label{eq:newS7V88-imaginary}
\end{equation}

\begin{lemma}[Nested-commutator chain bound]
\label{lem:newS7V88-chain}
The \(m\)-fold commutator obeys
\begin{equation}
 \|\operatorname{ad}_K^m(A)\|
 \leq (2J)^m\|A\|
 \prod_{j=0}^{m-1}\bigl(b+j(q-1)\bigr).
 \label{eq:newS7V88-chain}
\end{equation}
Every nonzero term in its interaction expansion is supported within graph
distance \(mR\) of \(B\).
\end{lemma}

\begin{proof}
In the first commutator only interactions meeting \(B\) contribute, and
the sum of their norms is at most \(Jb\).  The commutator norm costs a
factor two.  After \(j\) connected interactions have been appended, their
union with \(B\) has at most \(b+j(q-1)\) cells.  The sum of norms of the
next interactions meeting that union is at most
\(J[b+j(q-1)]\).  Iteration gives
\eqref{eq:newS7V88-chain}.  Each appended interaction has diameter at most
\(R\) and meets the current support, so the support radius grows by at most
\(R\) per step.
\end{proof}

\subsection{A small-complex-time quasi-locality theorem}
\label{sec:newS7V88-quasilocality}

Put
\begin{equation}
 a_B:={b\over q-1},
 \qquad
 x_s:=2J(q-1)s,
 \label{eq:newS7V88-x}
\end{equation}
and for \(0\leq x<1\) define the negative-binomial tail
\begin{equation}
 T_a(m_0,x):=
 \sum_{m=m_0}^\infty{(a)_m\over m!}x^m,
 \label{eq:newS7V88-tail}
\end{equation}
where \((a)_m=a(a+1)\cdots(a+m-1)\).

\begin{theorem}[Imaginary-time quasi-local expansion]
\label{thm:newS7V88-quasilocal}
If \(x_s<1\), the commutator series converges absolutely and
\begin{equation}
 \tau_s(A)=\sum_{m=0}^\infty{s^m\over m!}
 \operatorname{ad}_K^m(A),
 \qquad
 \|\tau_s(A)\|
 \leq\|A\|(1-x_s)^{-a_B}.
 \label{eq:newS7V88-series}
\end{equation}
Let \(\tau_s^{(<m_0)}(A)\) be the sum of the interaction-chain terms of
length below \(m_0\).  It is supported in the
\((m_0-1)R\)-neighbourhood of \(B\) and
\begin{equation}
 \|\tau_s(A)-\tau_s^{(<m_0)}(A)\|
 \leq\|A\|T_{a_B}(m_0,x_s).
 \label{eq:newS7V88-tail-bound}
\end{equation}
For every \(x<x'<1\) there is a finite
\(C(a,x,x')\) such that
\begin{equation}
 T_a(m_0,x)\leq C(a,x,x')(x')^{m_0}.
 \label{eq:newS7V88-exponential-tail}
\end{equation}
\end{theorem}

\begin{proof}
By \Cref{lem:newS7V88-chain}, the norm of the order-\(m\) term is at most
\[
 \|A\|{(a_B)_m\over m!}x_s^m,
\]
because
\(\prod_{j=0}^{m-1}[b+j(q-1)]=(q-1)^m(a_B)_m\).
The binomial series sums to \((1-x_s)^{-a_B}\), proving convergence and
the norm bound.  Omitting every chain of length at least \(m_0\) gives
\eqref{eq:newS7V88-tail-bound}, and the support statement is
\Cref{lem:newS7V88-chain}.  Finally,
\((a)_m/m!\) grows only polynomially in \(m\).  Therefore
\(((a)_m/m!)(x/x')^m\) is bounded, and summing the resulting geometric
tail proves \eqref{eq:newS7V88-exponential-tail}.
\end{proof}

\begin{corollary}[Uniform Duhamel-time window]
\label{cor:newS7V88-window}
The entire imaginary-time interval appearing in the Duhamel product is
covered whenever
\begin{equation}
 x_\beta:=2J(q-1)\beta<1.
 \label{eq:newS7V88-window}
\end{equation}
In that case the bounds above hold uniformly for
\(0\leq s\leq\beta\).
\end{corollary}

\begin{proof}
The Duhamel product uses
\(e^{tK}Ae^{-tK}\) with \(0\leq t\leq\beta\).  The quantity \(x_t\) is
increasing in \(t\), so \eqref{eq:newS7V88-window} controls the whole
interval.
\end{proof}

\begin{firewall}[This is a sufficient window, not a sharp temperature]
\label{fire:newS7V88-window}
The threshold \eqref{eq:newS7V88-window} comes from an absolute
nested-commutator count and is not claimed sharp.  Failure of the inequality
means this proof gives no uniform complex-time bound; it does not prove
that imaginary-time evolution is nonlocal.
\end{firewall}

\subsection{From equal-time mixing to Duhamel mixing}
\label{sec:newS7V88-ETCM}

Let \(\omega(A)=\operatorname{Tr}(e^{-\beta K}A)/Z\).  Cyclicity gives the
exact representation
\begin{equation}
 \langle A;D\rangle_\rho
 =\int_0^1\omega\bigl(\tau_{s\beta}(A)D\bigr)\,ds.
 \label{eq:newS7V88-Duhamel-imaginary}
\end{equation}

\begin{definition}[Equal-time carrier mixing, ETCM]
\label{def:newS7V88-ETCM}
ETCM with constants \(C_0,\mu_0>0\) means
\begin{equation}
 |\omega(CD)-\omega(C)\omega(D)|
 \leq C_0\|C\|\|D\||\operatorname{supp}C|
 |\operatorname{supp}D|e^{-\mu_0d(\operatorname{supp}C,
 \operatorname{supp}D)}
 \label{eq:newS7V88-ETCM}
\end{equation}
for the declared local even carrier observables, uniformly in volume,
interpolation parameter, bosonic histories, and boundary data.
\end{definition}

Assume the graph balls obey
\begin{equation}
 |N_L(B)|\leq C_g|B|e^{\kappa_gL}.
 \label{eq:newS7V88-growth}
\end{equation}

\begin{theorem}[ETCM plus imaginary-time locality loads ITCM]
\label{thm:newS7V88-loader}
Assume ETCM, \eqref{eq:newS7V88-growth}, and the uniform window
\eqref{eq:newS7V88-window}.  Restrict the declared exploration core to
\(|B|\leq b_0\), and fix \(x'\) with \(x_\beta<x'<1\).  Then there
are finite constants \(C_{\rm D}\) and \(\mu_{\rm D}>0\), depending only on
the displayed uniform data, such that for local \(A,D\),
\begin{equation}
 |\operatorname{Cov}^{\rm D}_\rho(A,D)|
 \leq C_{\rm D}\|A\|\|D\||B||X|
 e^{-\mu_{\rm D}d(B,X)},
 \label{eq:newS7V88-ITCM}
\end{equation}
provided \(\mu_0>\kappa_g\).  One may choose any
\begin{equation}
 \mu_{\rm D}<
 \min\left\{{\mu_0-\kappa_g\over2},
 { -\log x'\over2R}\right\}.
 \label{eq:newS7V88-muD}
\end{equation}
Thus the V87 ITCM row follows from a separate equal-time mixing theorem and
the explicit high-temperature complex-time window.
\end{theorem}

\begin{proof}
Let \(L=d(B,X)\) and choose
\(m_0=\max\{1,\lfloor L/(2R)\rfloor\}\).  Decompose
\(\tau_{s\beta}(A)=A_{<}+A_{\geq}\) as in
\Cref{thm:newS7V88-quasilocal}.  The first term is supported in
\(N_{L/2}(B)\), has norm at most
\(\|A\|(1-x_\beta)^{-a_B}\), and remains at distance at least \(L/2\)
from \(X\).  ETCM and \eqref{eq:newS7V88-growth} therefore bound its
connected correlation with \(D\) by a constant times
\[
 \|A\|\|D\||B||X|
 e^{-(\mu_0-\kappa_g)L/2}.
\]
The covariance involving \(A_{\geq}\) is at most
\(2\|A_{\geq}\|\|D\|\).  Equations
\eqref{eq:newS7V88-tail-bound}--\eqref{eq:newS7V88-exponential-tail}
bound it by a constant times
\(\|A\|\|D\|(x')^{L/(2R)}\).  Enlarge the constant for bounded \(L\),
combine the two exponents, and integrate
\eqref{eq:newS7V88-Duhamel-imaginary} over \(s\).  Any exponent strictly
below \eqref{eq:newS7V88-muD} absorbs the floor functions and polynomial
support factors.
\end{proof}

\begin{corollary}[A noncircular finite entrance]
\label{cor:newS7V88-entrance}
If ETCM is proved at finite cutoff by a high-temperature dependency,
cluster, or positive-slab argument that does not use the desired continuum
mass gap, then \Cref{thm:newS7V88-loader} supplies the same-carrier ITCM
needed by V87 and hence the interacting local refinement bound.
\end{corollary}

\begin{proof}
Apply \Cref{thm:newS7V88-loader} inside the V87 local compiler.
\end{proof}

\subsection{Why locality and an energy gap do not give uniform mixing}
\label{sec:newS7V88-counterexample}

\begin{countertheorem}[Low-temperature finite-range carrier with no uniform
clustering]
\label{cth:newS7V88-Ising}
Let
\begin{equation}
 K_N=-J\sum_{j=1}^{N-1}Z_jZ_{j+1},
 \qquad J>0,
 \label{eq:newS7V88-Ising}
\end{equation}
be the open classical Ising chain on \(N\) qubits.  The interaction range,
row norm, and excitation-energy scale are uniform in \(N\), and its real-
or imaginary-time evolution preserves every \(Z_j\) exactly.  Nevertheless
there is a sequence \(\beta_N\to\infty\) such that
\begin{equation}
 \operatorname{Cov}_{\beta_N}(Z_1,Z_N)
 =\bigl(\tanh(\beta_NJ)\bigr)^{N-1}longrightarrow1.
 \label{eq:newS7V88-Ising-cov}
\end{equation}
Hence finite range, trivial propagation of the chosen observables, and a
nonzero local excitation scale do not imply a volume- and
cutoff-temperature-uniform ETCM or ITCM bound.
\end{countertheorem}

\begin{proof}
The Hamiltonian is diagonal and invariant under the global spin flip, so
\(\omega(Z_j)=0\).  Introduce bond variables
\(b_j=Z_jZ_{j+1}\).  For open boundary conditions their Gibbs weights
factorize, and
\(\omega(b_j)=\tanh(\beta J)\).  Since
\(Z_1Z_N=\prod_{j=1}^{N-1}b_j\), this proves the equality in
\eqref{eq:newS7V88-Ising-cov}.  Choose \(\beta_N\) so that
\(\tanh(\beta_NJ)\geq1-N^{-2}\).  Then the right side is at least
\((1-N^{-2})^{N-1}\to1\).  Every \(Z_j\) commutes with \(K_N\), so its
evolution is exactly local.  Flipping one interior spin costs at most a
fixed multiple of \(J\), independent of \(N\).  The claimed uniform
mixing conclusion therefore fails despite the listed kinematic bounds.
\end{proof}

\begin{firewall}[Hamiltonian spectral information and Gibbs mixing are
different]
\label{fire:newS7V88-gap}
A many-body energy gap, especially with a degenerate low-energy sector,
does not by itself control thermal correlations uniformly when
\(\beta_n\) changes with the cutoff.  The physical clock and temperature
scaling must be entered into ETCM explicitly.
\end{firewall}

\subsection{Finite-time carrier mixing card}
\label{sec:newS7V88-card}

\begin{definition}[Finite-time carrier mixing card, FTCMC]
\label{def:newS7V88-FTCMC}
FTCMC requires:
\begin{enumerate}
\item a support-local bounded interaction with uniform \(R,q,J\) and graph
growth constants;
\item the explicit complex-time window
\(2J(q-1)\beta_n<1\), or a stronger replacement theorem beyond that
window;
\item an independently proved uniform ETCM estimate along the complete
old/detail interpolation;
\item the ITCM constants produced by
\Cref{thm:newS7V88-loader};
\item the V87 summable local coupling stocks and their gauge, metric,
reflection, and boundary derivatives;
\item a declared physical-clock scaling explaining the role of
\(\beta_n\); and
\item compatibility with RFRC, NCRC, IRCMC, and the bosonic renewal
history.
\end{enumerate}
\end{definition}

\begin{assessment}[The upstream bottleneck has split cleanly]
\label{ass:newS7V88-status}
Finite-range carrier dynamics supplies exponential imaginary-time
quasi-locality in the explicit window
\(2J(q-1)\beta<1\).  It upgrades equal-time mixing to the Duhamel mixing
needed for interacting refinement.  It does not supply equal-time mixing.
The first remaining writer is therefore a positive finite-cutoff ETCM
proof for the actual interpolating carrier, with physical \(\beta_n\)
scaling.  If no direct quantum proof is available, the next upstream object
is a Trotterized slab whose dependency graph can be tested by a Dobrushin
row without invoking the continuum gap.
\end{assessment}

\begin{programme}[V89: Trotter-slab dependency compiler]
\label{prog:newS7V88-V89}
V89 will formulate a finite checkerboard or Trotter slab with explicit
local kernels, prove a Dobrushin influence bound for strictly positive slab
weights, and show how its equal-time boundary marginal inherits ETCM.  It
will also isolate the fermionic sign or nonstoquastic obstruction: if the
slab kernels are not positive, the classical dependency theorem cannot be
applied and the retained operator route must keep a separate matrix-valued
mixing norm.
\end{programme}

\begin{firewall}[Claim boundary after seventh-series V88]
\label{fire:newS7V88-boundary}
V88 proves a nested-commutator imaginary-time quasi-locality theorem with
explicit convergence and exponential-tail bounds, proves that ETCM plus
that window loads ITCM, and gives an exact low-temperature finite-range
counterexample to automatic uniform mixing.

V88 does not prove ETCM or the complex-time window for the physical
Standard Model scaling, construct a positive Trotter slab, or prove FTCMC,
IRCMC, NCRC, RFRC, CRARC, LPCC, OFHC, LCPLC, CCPC, GCMCC, PCC, SMALC,
RCLC, DSLC, CRDC, GAOCC, CCTC, MSRC, RTGC, or RCNC, remove the cutoff, or
construct the Standard Model--Einstein continuum.
\end{firewall}

\begin{theorem}[Seventh-series V88 result]
\label{thm:newS7V88-result}
For a finite-range bounded carrier interaction, imaginary-time evolution of
a local observable has the exponential chain tail
\eqref{eq:newS7V88-tail-bound} whenever
\(2J(q-1)\beta<1\).  If the Gibbs state separately satisfies ETCM, this
quasi-locality yields the Duhamel mixing bound
\eqref{eq:newS7V88-ITCM} and closes the V87 local refinement row.  Finite
range and a fixed local energy scale alone cannot provide the uniform
mixing, as shown by \Cref{cth:newS7V88-Ising}.
\end{theorem}

\begin{proof}
Combine \Cref{lem:newS7V88-chain,
thm:newS7V88-quasilocal,
cor:newS7V88-window,
thm:newS7V88-loader,
cth:newS7V88-Ising}.
\end{proof}

\paragraph{Parent-version record.}
V88 was formed from
\texttt{Renewal\_SM\_Einstein\_Continuum\_Research\_Notes\_V87.tex}.
The V87 parent had \(35\,113\) logical source lines, \(1\,379\,780\) bytes,
and SHA--256
\[
 \texttt{9B328FA64CE8CC77310C22F73C3F6571BF4AED355CE337377C323B815C1A0B1B}.
\]
The next compulsory companion audit is V90.

% =============================================================================
% END APPENDED SEVENTH-SERIES V88 MATERIAL
% =============================================================================
% =============================================================================
% APPENDED SEVENTH-SERIES V89 MATERIAL
% =============================================================================

\section{Positive Trotter slabs and a dependency-to-ETCM compiler}
\label{sec:newS7V89}
\label{sec:v7v89}

V88 shows that imaginary-time quasi-locality upgrades equal-time mixing to
Duhamel mixing, but does not prove the equal-time row.  A possible
noncircular entrance is a finite Trotter or checkerboard slab with strictly
positive local weights.  On that branch the problem is classical and its
single-site influences can be computed from mixed oscillations of the slab
factors.  This section proves that loader, derives exponential boundary
influence and equal-slice covariance under a strict Dobrushin row, and
states why a fermionic sign, complex gate, or nonuniform time-slice limit
blocks the argument.

\subsection{Strictly positive slab factorization}
\label{sec:newS7V89-slab}

Let \(I\) be a finite space--time cell graph and let each cell variable
\(z_i\) take values in a compact measured space \((E_i,m_i)\).  Suppose the
finite slab weight is
\begin{equation}
 W(z)=\prod_{A\Subset I}w_A(z_A),
 \qquad
 0<w_A<\infty,
 \qquad
 \phi_A:=-\log w_A,
 \label{eq:newS7V89-weight}
\end{equation}
with finite-range factors and a finite partition function.  Its probability
law is \(d\mu=Z^{-1}W\prod_i dm_i\).

For two exterior configurations \(\eta,\eta'\) that differ only at cell
\(j\), define the conditional energy difference seen at \(i\) by
\begin{equation}
 \Delta_{ij}(x_i;\eta,\eta')
 :=\sum_{A\ni i,j}
 \left[\phi_A(x_i,\eta_{A\setminus\{i\}})
 -\phi_A(x_i,\eta'_{A\setminus\{i\}})\right]
 \label{eq:newS7V89-Delta}
\end{equation}
and the mixed oscillation stock
\begin{equation}
 \kappa_{ij}:=
 \sup_{\eta\sim_j\eta'}
 \operatorname{osc}_{x_i}\Delta_{ij}(x_i;\eta,\eta').
 \label{eq:newS7V89-kappa}
\end{equation}
The sum is zero if no factor contains both \(i\) and \(j\).

Let \(\mu_i(\cdot\mid\eta)\) be the one-cell conditional law and use the
total-variation convention
\(\|\nu-\nu'\|_{\rm TV}=\sup_B|\nu(B)-\nu'(B)|\).  Define
\begin{equation}
 c_{ij}:=
 \sup_{\eta\sim_j\eta'}
 \|\mu_i(\cdot\mid\eta)-\mu_i(\cdot\mid\eta')\|_{\rm TV}.
 \label{eq:newS7V89-cij}
\end{equation}

\begin{lemma}[Log-oscillation comparison for normalized laws]
\label{lem:newS7V89-TV}
Let \(\nu\) and \(\nu'\) be equivalent probability measures and suppose
\begin{equation}
 \operatorname{osc}\log{d\nu\over d\nu'}\leq\kappa.
 \label{eq:newS7V89-log-osc}
\end{equation}
Then
\begin{equation}
 \|\nu-\nu'\|_{\rm TV}
 \leq \min\left\{1,{e^\kappa-1\over2}\right\}.
 \label{eq:newS7V89-TV}
\end{equation}
\end{lemma}

\begin{proof}
Let \(r=d\nu/d\nu'\).  Since \(\int r,d\nu'=1\), either \(r=1\)
almost surely or its essential range meets both sides of one.  Thus zero
lies between the essential infimum and supremum of \(\log r\).  The
oscillation hypothesis gives \(|\log r|\leq\kappa\), hence
\(|r-1|\leq e^\kappa-1\).  Therefore
\[
 \|\nu-\nu'\|_{\rm TV}
 ={1\over2}\int|r-1|d\nu'
 \leq{e^\kappa-1\over2}.
\]
Total variation is at most one.
\end{proof}

\begin{theorem}[Slab factor-to-influence loader]
\label{thm:newS7V89-loader}
For the positive factorization \eqref{eq:newS7V89-weight},
\begin{equation}
 \boxed{
 c_{ij}\leq
 \widehat c_{ij}:=
 \min\left\{1,{e^{\kappa_{ij}}-1\over2}\right\}.}
 \label{eq:newS7V89-loader}
\end{equation}
Consequently the explicitly checkable row
\begin{equation}
 \widehat\alpha:=\sup_i\sum_{j\ne i}\widehat c_{ij}<1
 \label{eq:newS7V89-Dobrushin-row}
\end{equation}
implies the true Dobrushin row
\(\alpha:=\sup_i\sum_{j\ne i}c_{ij}<1\).
\end{theorem}

\begin{proof}
All factors not containing \(i\) cancel from the one-cell conditional.
When the exterior changes only at \(j\), the logarithm of the ratio of the
two unnormalized conditional densities is
\(-\Delta_{ij}(x_i;\eta,\eta')\).  Normalization adds a constant in
\(x_i\), so the log-ratio oscillation is exactly bounded by
\(\kappa_{ij}\).  Apply \Cref{lem:newS7V89-TV}.  Summation gives the final
claim.
\end{proof}

\subsection{Influence propagation}
\label{sec:newS7V89-propagation}

Let \(C=(c_{ij})\) and, when \(\alpha<1\), define the nonnegative
resolvent
\begin{equation}
 D:=(I-C)^{-1}=\sum_{n=0}^\infty C^n.
 \label{eq:newS7V89-D}
\end{equation}
For a bounded function \(F\), let
\begin{equation}
 \delta_i(F):=sup_{z\sim_i z'}|F(z)-F(z')|.
 \label{eq:newS7V89-osc}
\end{equation}

\begin{theorem}[Boundary influence under the strict row]
\label{thm:newS7V89-boundary}
Let \(\Lambda\subset I\) be a finite conditional region and let exterior
conditions \(\eta,\eta'\) differ on \(Y\subset I\setminus\Lambda\).  There
is a coupling of the two conditional slab laws such that the disagreement
probabilities \(p_i\), \(i\in\Lambda\), obey
\begin{equation}
 p_i\leq\sum_{j\in Y}D_{ij}.
 \label{eq:newS7V89-disagreement}
\end{equation}
Consequently
\begin{equation}
 |\mathbb E_\eta F-\mathbb E_{\eta'}F|
 \leq\sum_{i\in\Lambda}\delta_i(F)
 \sum_{j\in Y}D_{ij}.
 \label{eq:newS7V89-boundary-influence}
\end{equation}
\end{theorem}

\begin{proof}
Couple each one-cell heat-bath update maximally.  If the current
configurations disagree at a set with indicator vector \(v\), the
probability that an update at \(i\) disagrees is at most
\(\sum_jc_{ij}v_j\), plus the direct exterior discrepancies.  Successive
sweeps and monotone iteration from the all-disagreement vector give the
minimal nonnegative inequality
\(p\leq Cp+b\), where
\(b_i\leq\sum_{j\in Y}c_{ij}\) together with the boundary-to-interior
incidence.  Extending \(C\) to include the frozen boundary sites makes this
notation simply \(p\leq C p+\mathbf1_Y\).  Since
\(\|C\|_{\infty\to\infty}\leq\alpha<1\), iteration gives
\(p\leq\sum_{n\geq0}C^n\mathbf1_Y=D\mathbf1_Y\), which is
\eqref{eq:newS7V89-disagreement}.  Under this coupling,
\(|F(z)-F(z')|\leq\sum_i\delta_i(F)\mathbf1_{\{z_i\ne z_i'\}}\).
Take expectations to obtain
\eqref{eq:newS7V89-boundary-influence}.
\end{proof}

\begin{remark}[Boundary matrix convention]
\label{rem:newS7V89-boundary-matrix}
Equivalently one may separate an interior influence matrix from a
boundary-incidence matrix.  The proof then reads
\(p\leq(I-C_{\Lambda\Lambda})^{-1}C_{\Lambda Y}\mathbf1\).
The extended notation in \eqref{eq:newS7V89-disagreement} only shortens
this standard expression.
\end{remark}

Assume \(c_{ij}=0\) when \(d(i,j)>R_c\).

\begin{lemma}[Exponential decay of the influence resolvent]
\label{lem:newS7V89-D-decay}
If \(0<\alpha<1\), then for \(i\ne j\),
\begin{equation}
 D_{ij}\leq
 {\alpha^{\lceil d(i,j)/R_c\rceil}\over1-\alpha}
 \leq {\alpha^{-1}\over1-\alpha}
 \exp\left(-{|\log\alpha|\over R_c}d(i,j)\right).
 \label{eq:newS7V89-D-decay}
\end{equation}
\end{lemma}

\begin{proof}
A product of fewer than \(\lceil d(i,j)/R_c\rceil\) finite-range influence
steps cannot connect \(i\) to \(j\), so the corresponding entries of
\(C^n\) vanish.  Every row sum of \(C^n\) is at most \(\alpha^n\), hence
\((C^n)_{ij}\leq\alpha^n\).  Sum the geometric tail and use the ceiling
bound.
\end{proof}

\subsection{Equal-slice covariance}
\label{sec:newS7V89-covariance}

\begin{theorem}[Positive-slab Dobrushin row implies ETCM]
\label{thm:newS7V89-ETCM}
Let \(F,G\) be bounded real functions supported in disjoint cell sets
\(B,X\) of the positive slab.  Under
\eqref{eq:newS7V89-Dobrushin-row},
\begin{equation}
 |\operatorname{Cov}_\mu(F,G)|
 \leq {1\over4}
 \sum_{i\in B}\sum_{j\in X}
 \delta_i(F)D_{ij}\delta_j(G).
 \label{eq:newS7V89-covariance}
\end{equation}
In particular, for fixed-size supports and finite-range influences, the
right side decays exponentially with \(d(B,X)\) by
\Cref{lem:newS7V89-D-decay}.  Restricting \(B,X\) to one equal-time
boundary slice gives the ETCM row needed in V88.
\end{theorem}

\begin{proof}
Reveal the variables in \(X\) one at a time and write the Doob martingale
for \(\mathbb E(F\mid z_X)\).  Changing the revealed value at \(j\)
changes this conditional expectation by at most
\(\sum_{i\in B}\delta_i(F)D_{ij}\), by
\Cref{thm:newS7V89-boundary}.  The standard covariance telescoping identity
pairs that martingale increment with the corresponding change of \(G\).
For two bounded one-variable functions, the absolute covariance is at most
one quarter of the product of their oscillations.  Summing the increments
over \(j\in X\) yields
\eqref{eq:newS7V89-covariance}.  The decay statement is
\Cref{lem:newS7V89-D-decay}; an equal-time slice is a subset of the slab,
so the same bound restricts to it.
\end{proof}

\begin{corollary}[Positive-slab route to interacting carrier projectivity]
\label{cor:newS7V89-route}
If the Trotter-slab representation of every interpolating carrier is
strictly positive, has a finite-range factor graph, and satisfies
\eqref{eq:newS7V89-Dobrushin-row} uniformly in volume, boundary data,
interpolation parameter, and Trotter depth, then its equal-time marginal
satisfies a uniform ETCM estimate.  If the V88 complex-time window also
holds, ITCM and the V87 local projectivity compiler follow.
\end{corollary}

\begin{proof}
Apply \Cref{thm:newS7V89-ETCM} uniformly and then
\Cref{thm:newS7V88-loader,thm:newS7V87-local}.
\end{proof}

\subsection{Trotter origin and its obstructions}
\label{sec:newS7V89-Trotter}

Let a finite carrier Hamiltonian be decomposed into checkerboard layers
\(K=K_1+\cdots+K_m\) whose terms within each layer commute.  At Trotter
depth \(M\), insert a fixed product basis between the factors in
\begin{equation}
 \left(e^{-\beta K_1/M}\cdots e^{-\beta K_m/M}\right)^M.
 \label{eq:newS7V89-Trotter}
\end{equation}
The resulting matrix-element product is a slab factorization.

\begin{proposition}[Positivity gate]
\label{prop:newS7V89-positivity}
The classical probability construction
\eqref{eq:newS7V89-weight} applies to
\eqref{eq:newS7V89-Trotter} only if all matrix-element products can be made
real and nonnegative by one compatible basis and gauge convention, and the
retained support has strictly positive local conditionals.  A negative or
complex fermionic matrix element cannot be inserted into
\eqref{eq:newS7V89-cij}; total-variation Dobrushin comparison is then
unlicensed.
\end{proposition}

\begin{proof}
A probability density must be nonnegative, and the conditional laws in
\eqref{eq:newS7V89-cij} require positive normalizations.  Negative or
complex factor products define a signed or complex functional rather than
a probability.  Total variation between its normalized one-site
``conditionals'' is therefore not the probability metric used in
\Cref{lem:newS7V89-TV,thm:newS7V89-boundary}.
\end{proof}

\begin{countertheorem}[Small Trotter time does not automatically give a
uniform influence row]
\label{cth:newS7V89-small-time}
Even when \(e^{-\Delta\tau K}\to I\) as
\(\Delta\tau\downarrow0\), off-diagonal matrix elements in a fixed basis
may vanish like a power of \(\Delta\tau\).  Their negative logarithms then
have oscillations of order \(|\log\Delta\tau|\).  Consequently the mixed
stocks \(\kappa_{ij}\) and the row
\eqref{eq:newS7V89-Dobrushin-row} need not improve as Trotter depth grows.
\end{countertheorem}

\begin{proof}
Take a two-level hopping Hamiltonian with a nonzero off-diagonal entry.
The off-diagonal matrix element of
\(e^{-\Delta\tau K}=I-\Delta\tau K+O(\Delta\tau^2)\) is
\(-\Delta\tau K_{12}+O(\Delta\tau^2)\).  Its absolute logarithmic weight is
\(-\log\Delta\tau+O(1)\), while a diagonal matrix element has logarithm
\(O(\Delta\tau)\).  A conditional comparison involving the two transition
types can therefore have logarithmic oscillation growing like
\(|\log\Delta\tau|\).  Small operator time alone does not bound
\eqref{eq:newS7V89-kappa}.
\end{proof}

\begin{firewall}[Finite-depth mixing must pass to the quantum limit]
\label{fire:newS7V89-Trotter-limit}
A strict row at each fixed Trotter depth is insufficient if
\(\widehat\alpha_M\uparrow1\), the decay length diverges, or local Trotter
errors fail to vanish uniformly.  The ETCM constants and equal-slice
marginals must have a controlled \(M\to\infty\) limit before they are used
in V88.
\end{firewall}

\subsection{Slab dependency card}
\label{sec:newS7V89-card}

\begin{definition}[Positive slab dependency card, PSDC]
\label{def:newS7V89-PSDC}
For every cutoff, interpolation parameter, and bosonic history, PSDC
requires:
\begin{enumerate}
\item a declared checkerboard decomposition, product basis, Trotter depth,
and exact local factor graph;
\item real strict positivity of every retained slab conditional, including
fermionic and gauge signs;
\item the mixed oscillations \(\kappa_{ij}\), loaded influence matrix
\(\widehat C\), and a uniform row margin
\(\sup_i\sum_j\widehat c_{ij}\leq\alpha_*<1\);
\item uniform finite influence range and the resolvent decay
\eqref{eq:newS7V89-D-decay};
\item convergence of equal-time marginals and their ETCM constants as
Trotter depth tends to infinity;
\item compatibility with gauge projection, reflection, CAR grading,
metric variation, and the old/detail refinement; and
\item insertion of the resulting ETCM constants into FTCMC, ITCM, and the
V87 summable local stock.
\end{enumerate}
\end{definition}

\begin{assessment}[The finite positive entrance is exact but conditional]
\label{ass:newS7V89-status}
A strictly positive Trotter slab with a uniform mixed-oscillation row below
one supplies the missing equal-time mixing theorem and therefore the
interacting retained-carrier refinement route.  The active physical
debits are now sharply testable: positivity of the full chiral/gauge slab,
uniformity of the Dobrushin margin in Trotter depth and cutoff, and
compatibility with the physical clock.  A failed positivity gate sends the
programme back to a matrix-valued operator mixing norm rather than to a
classical probability claim.
\end{assessment}

\begin{programme}[V90: companion audit and sign-resolved slab alternative]
\label{prog:newS7V89-V90}
V90 will perform the compulsory YM and NS audit.  It will then split the
actual carrier slab into a positive modulus kernel and a unitary or signed
local factor, derive the exact phase-margin loss for one slab conditional,
and determine whether the V81 local complex-kernel criterion can coexist
with the V89 Dobrushin row.  If the local sign width is not below \(\pi\),
the note will retain the operator carrier and formulate a completely
bounded map influence instead.
\end{programme}

\begin{firewall}[Claim boundary after seventh-series V89]
\label{fire:newS7V89-boundary}
V89 proves the slab mixed-oscillation influence loader, strict-row boundary
influence, exponential influence-resolvent decay, positive-slab ETCM
compiler, and the exact positivity and Trotter-uniformity gates.  It also
proves that small Trotter time need not improve logarithmic influences.

V89 does not construct a positive Standard Model slab, prove a uniform
Dobrushin row or Trotter limit, or prove PSDC, FTCMC, IRCMC, NCRC, RFRC,
CRARC, LPCC, OFHC, LCPLC, CCPC, GCMCC, PCC, SMALC, RCLC, DSLC, CRDC,
GAOCC, CCTC, MSRC, RTGC, or RCNC, remove the cutoff, or construct the
Standard Model--Einstein continuum.
\end{firewall}

\begin{theorem}[Seventh-series V89 result]
\label{thm:newS7V89-result}
For a strictly positive finite-range slab factorization, the single-site
influence is bounded explicitly by
\(\widehat c_{ij}=\min\{1,(e^{\kappa_{ij}}-1)/2\}\).  A uniform row sum
below one gives exponential boundary influence and equal-slice covariance
decay, thereby loading ETCM and, inside the V88 complex-time window, the
V87 interacting projectivity compiler.  Negative or complex slab weights
and a nonuniform Trotter-depth row invalidate this probability route.
\end{theorem}

\begin{proof}
Combine \Cref{lem:newS7V89-TV,
thm:newS7V89-loader,
thm:newS7V89-boundary,
lem:newS7V89-D-decay,
thm:newS7V89-ETCM,
prop:newS7V89-positivity,
cth:newS7V89-small-time}.
\end{proof}

\paragraph{Parent-version record.}
V89 was formed from
\texttt{Renewal\_SM\_Einstein\_Continuum\_Research\_Notes\_V88.tex}.
The V88 parent had \(35\,506\) logical source lines, \(1\,394\,386\) bytes,
and SHA--256
\[
 \texttt{5BA51DE1E5843959F3338FCC69F0B66D3E12B599F9DAC18922AA29980670BDFE}.
\]
The next compulsory companion audit is V90.

% =============================================================================
% END APPENDED SEVENTH-SERIES V89 MATERIAL
% =============================================================================
% =============================================================================
% APPENDED SEVENTH-SERIES V90 MATERIAL
% =============================================================================

\section{Companion audit and sign-resolved slab influences}
\label{sec:newS7V90}
\label{sec:v7v90}

This is the compulsory V90 companion audit.  No companion file has advanced
since V85.  The mathematical continuation tests whether the positive-slab
compiler of V89 survives a local sign or phase.  A fixed exterior phase
cancels, as in V81, but boundary changes affect both the positive one-site
law and its local phase.  This section derives an explicit complex
single-site influence bound.  Its row can support an algebraic complex
Dobrushin contraction only when every local phase width is below \(\pi\)
and the accumulated modulus and phase influences remain strictly
subcritical.

\subsection{Compulsory V90 companion audit}
\label{sec:newS7V90-audit}

At the audit snapshot, the newest distinct complete Yang--Mills checkpoint
was still
\begin{align}
 &\texttt{Renewal\_Yang\_Mills\_Mass\_Gap\_Research\_Notes\_V45.tex},
 \notag\\
 &\qquad 704\,056\ {\rm bytes},\quad 20\,173\ {\rm logical\ lines},\quad
 \texttt{793583EFCD7EBDD7E96D77D5A9E7C6CE78CECDF693BA6196A7C60D5B3AF1E6AC}.
 \label{eq:newS7V90-YM-hash}
\end{align}
The newest Navier--Stokes checkpoint was still
\begin{align}
 &\texttt{Renewal\_NS\_Stability\_Research\_Notes\_V26.tex},
 \notag\\
 &\qquad 278\,283\ {\rm bytes},\quad 5\,319\ {\rm logical\ lines},\quad
 \texttt{82C887F8C2C69E4F28FAD7C3DC8DE5B9099CB5A4BF431EBA1FFF3E91935978AD}.
 \label{eq:newS7V90-NS-hash}
\end{align}
Both have one live document terminator.

\begin{theorem}[V90 audit decision]
\label{thm:newS7V90-audit}
The companion audit imports no new writer.  YM V45 continues to support a
same-law covariance or retained-carrier route; NS V26 continues to forbid
unproved rank reductions.  The V89 positive-slab row and the present phase
row must therefore be populated within the Standard Model carrier itself.
\end{theorem}

\begin{proof}
The newest distinct filenames, sizes, and hashes agree with the V85 audit,
and inspection reveals no later append.  The conclusion is their unchanged
claim boundary.
\end{proof}

\subsection{Local modulus--phase conditional}
\label{sec:newS7V90-conditional}

Let the slab factors be nonzero complex functions
\begin{equation}
 g_A(z_A)=|g_A(z_A)|e^{i\theta_A(z_A)}
 \label{eq:newS7V90-factor}
\end{equation}
on one declared phase-lift chart.  Their moduli define the positive slab
law of V89.  For a cell \(i\) and fixed exterior \(\eta\), let
\(\nu_i^\eta\) be its positive one-cell conditional and put
\begin{equation}
 \theta_i(x;\eta):=sum_{A\ni i}\theta_A(x,\eta_{A\setminus\{i\}}),
 \qquad
 z_i(\eta):=\int e^{i\theta_i(x;\eta)}d\nu_i^\eta(x).
 \label{eq:newS7V90-local-phase}
\end{equation}
All factors not meeting \(i\) contribute a phase depending only on
\(\eta\) and cancel from the normalized conditional.  When
\(z_i(\eta)\ne0\), define the mass-one complex kernel
\begin{equation}
 Q_i^\eta(dx):={e^{i\theta_i(x;\eta)}\over z_i(\eta)}
 \nu_i^\eta(dx).
 \label{eq:newS7V90-Q}
\end{equation}

Define the local phase width and margin
\begin{equation}
 w_i:=\sup_\eta\operatorname{osc}_x\theta_i(x;\eta),
 \qquad
 \rho_i:=\cos(w_i/2).
 \label{eq:newS7V90-width}
\end{equation}
When \(w_i<\pi\), V81 gives \(|z_i(\eta)|\geq\rho_i>0\).

\begin{lemma}[Factor phase-width loader]
\label{lem:newS7V90-width}
On a common real lift,
\begin{equation}
 w_i\leq\sum_{A\ni i}
 \sup_\eta\operatorname{osc}_x
 \theta_A(x,\eta_{A\setminus\{i\}}).
 \label{eq:newS7V90-width-loader}
\end{equation}
Thus a uniformly summable local phase interaction gives a
volume-independent width.  It gives a useful conditional margin only if
the right side is strictly below \(\pi\).
\end{lemma}

\begin{proof}
Oscillation is subadditive under sums.  Apply this to
\eqref{eq:newS7V90-local-phase} and take the exterior supremum.
\end{proof}

\subsection{Complex influence bound}
\label{sec:newS7V90-influence}

For two exteriors \(\eta\sim_j\eta'\), define the phase-boundary stock
\begin{equation}
 \chi_{ij}:=sup_{\eta\sim_j\eta'}\sup_x
 \left|e^{i\theta_i(x;\eta)}
 -e^{i\theta_i(x;\eta')}\right|.
 \label{eq:newS7V90-chi}
\end{equation}
It vanishes unless a phase factor meets both \(i\) and \(j\), and it is at
most
\begin{equation}
 \min\left\{2,
 \sup_{\eta\sim_j\eta'}\sup_x
 |\theta_i(x;\eta)-\theta_i(x;\eta')|
 \right\}
 \label{eq:newS7V90-chi-lift}
\end{equation}
for the chosen lift.

For a complex measure \(\sigma\), write \(\|\sigma\|_{\rm var}\) for its
full total variation.  Define
\begin{equation}
 c_{ij}^{\mathbb C}:={1\over2}
 \sup_{\eta\sim_j\eta'}
 \|Q_i^\eta-Q_i^{\eta'}\|_{\rm var}.
 \label{eq:newS7V90-complex-influence}
\end{equation}

\begin{theorem}[Modulus--phase influence compiler]
\label{thm:newS7V90-influence}
Assume \(w_i<\pi\).  Let \(c_{ij}\) be the positive modulus influence from
V89.  Then
\begin{equation}
 \boxed{
 c_{ij}^{\mathbb C}
 \leq
 \left(c_{ij}+{\chi_{ij}\over2}\right)
 \left({1\over\rho_i}+{1\over\rho_i^2}\right).}
 \label{eq:newS7V90-influence}
\end{equation}
Therefore the explicit row
\begin{equation}
 \alpha_{\mathbb C}:=sup_i\sum_{j\ne i}
 \left(c_{ij}+{\chi_{ij}\over2}\right)
 \left({1\over\rho_i}+{1\over\rho_i^2}\right)<1
 \label{eq:newS7V90-complex-row}
\end{equation}
is a sufficient complex single-site contraction condition.
\end{theorem}

\begin{proof}
Put
\(q=e^{i\theta_i(\cdot;\eta)}\nu_i^\eta\) and define \(q'\) similarly.
Since the positive total-variation convention in V89 is half the full
variation,
\begin{equation}
 \|q-q'\|_{\rm var}
 \leq2c_{ij}+\chi_{ij}.
 \label{eq:newS7V90-q-difference}
\end{equation}
Indeed, add and subtract
\(e^{i\theta_i(\cdot;\eta)}\nu_i^{\eta'}\); the first difference is the
positive-law variation and the second is bounded by the phase supremum.
The normalizations are \(z=q(1)\) and \(z'=q'(1)\), so
\(|z-z'|\leq\|q-q'\|_{\rm var}\), while
\(|z|,|z'|\geq\rho_i\).  Hence
\begin{align*}
 \left\|{q\over z}-{q'\over z'}\right\|_{\rm var}
 &\leq {\|q-q'\|_{\rm var}\over|z|}
 +{|z-z'|\over|z||z'|}\|q'\|_{\rm var}\\
 &\leq(2c_{ij}+\chi_{ij})
 \left({1\over\rho_i}+{1\over\rho_i^2}\right),
\end{align*}
because \(\|q'\|_{\rm var}=1\).  Divide by two to obtain
\eqref{eq:newS7V90-influence}.  Summation gives
\eqref{eq:newS7V90-complex-row}.
\end{proof}

\begin{corollary}[The phase loss is per influence row]
\label{cor:newS7V90-loss}
Even if the modulus row is far below one, the complex row can fail because
\(w_i\uparrow\pi\), so \(\rho_i\downarrow0\), or because the boundary
phase response \(\sum_j\chi_{ij}\) is large.  A local nonzero determinant
and finite phase oscillation do not prevent either failure.
\end{corollary}

\begin{proof}
Both losses are explicit on the right side of
\eqref{eq:newS7V90-complex-row}.
\end{proof}

\subsection{Algebraic complex Dobrushin contraction}
\label{sec:newS7V90-complex-Dobrushin}

A complex mass-one kernel acts on bounded functions by
\((Q_iF)(\eta)=\int F(x_i,\eta)Q_i^\eta(dx_i)\).  Positivity is unavailable,
but total-variation operator norms still telescope.

\begin{theorem}[Complex-kernel boundary resolvent]
\label{thm:newS7V90-resolvent}
Assume the kernels \(Q_i\) form the one-site conditionals of one declared
finite complex slab functional, every local normalization is nonzero, and
\(\alpha_{\mathbb C}<1\).  Let
\begin{equation}
 D_{\mathbb C}:=(I-C_{\mathbb C})^{-1}
 =\sum_{n=0}^\infty C_{\mathbb C}^n,
 \qquad C_{\mathbb C}=(c_{ij}^{\mathbb C}).
 \label{eq:newS7V90-DC}
\end{equation}
Then the oscillation of any bounded local function under a change of
boundary data is bounded by the same resolvent propagation as in V89, with
\(D\) replaced by \(D_{\mathbb C}\).  If the complex influences have
range \(R_c\),
\begin{equation}
 (D_{\mathbb C})_{ij}
 \leq {\alpha_{\mathbb C}^{\lceil d(i,j)/R_c\rceil}
 \over1-\alpha_{\mathbb C}}
 \label{eq:newS7V90-DC-decay}
\end{equation}
for \(i\ne j\).
\end{theorem}

\begin{proof}
For a mass-one complex kernel, constants are fixed.  If two exteriors differ
at \(j\), trace duality between bounded functions and complex variation
shows that the change of a one-cell update is at most
\(2c_{ij}^{\mathbb C}\) times the relevant cell oscillation.  Successive
single-cell replacements therefore give the same nonnegative oscillation
recursion as the V89 disagreement proof, with coefficient matrix
\(C_{\mathbb C}\).  Its Neumann series converges because its row norm is
below one.  Finite range makes the first possible connecting power at
least \(\lceil d/R_c\rceil\), and the geometric row bound proves
\eqref{eq:newS7V90-DC-decay}.
\end{proof}

\begin{firewall}[Complex contraction is not a probability coupling]
\label{fire:newS7V90-not-probability}
The preceding theorem controls algebraic conditional response of a declared
complex functional.  It does not create a positive Gibbs measure, a
probabilistic coupling, reflection positivity, or a positive Osterwalder--
Schrader form.  Those properties require the retained positive operator
carrier or an independently proved scalar positive-shadow branch.
\end{firewall}

\begin{corollary}[Compatibility with the V81 bad-step criterion]
\label{cor:newS7V90-V81}
If the modulus bad activity at one slab step is \(p_{\rm abs}\), then its
normalized complex bad amplitude is at most
\(p_{\rm abs}/\rho_i\).  Thus the V81 path criterion and the stronger
complex influence row can be tested from the same stocks
\((c_{ij},\chi_{ij},w_i)\), but neither implies positivity of the slab.
\end{corollary}

\begin{proof}
The event bound is V81
\Cref{cor:newS7V81-event}; the shared stocks are the definitions above.
\end{proof}

\subsection{The fermionic sign wall}
\label{sec:newS7V90-sign-wall}

\begin{countertheorem}[A binary sign can saturate the local phase wall]
\label{cth:newS7V90-sign}
Let one cell have two values of equal positive conditional probability and
phases \(0\) and \(\pi\).  Then \(w_i=\pi\), \(z_i=0\), and the complex
one-cell conditional \eqref{eq:newS7V90-Q} is undefined.  No modulus
Dobrushin margin repairs this failure.
\end{countertheorem}

\begin{proof}
The phase average is \((1+e^{i\pi})/2=0\).  The width and conclusion follow
directly.
\end{proof}

\begin{corollary}[Permutation signs require an independent treatment]
\label{cor:newS7V90-permutation}
A Trotter representation in which a local fermionic permutation or loop
changes the factor sign can realize the phase pair in
\Cref{cth:newS7V90-sign}.  Before using the scalar complex row, one must
prove that the actual conditional weights avoid such balanced sign
sectors.  Otherwise the fermionic degrees of freedom must remain in the
operator or Grassmann carrier.
\end{corollary}

\begin{proof}
A sign change is a phase change by \(\pi\).  If the positive conditional
weights of the two sectors balance, the local normalization vanishes; near
balance its inverse is arbitrarily large.
\end{proof}

\subsection{Completely bounded carrier alternative}
\label{sec:newS7V90-cb}

Let \(\mathcal T_i^\eta\) be a local completely positive unital update map
on the retained observable algebra, obtained from the positive operator
carrier rather than a scalar sign expansion.  Define
\begin{equation}
 c_{ij}^{\rm cb}:={1\over2}
 \sup_{\eta\sim_j\eta'}
 \|\mathcal T_i^\eta-\mathcal T_i^{\eta'}\|_{\rm cb}.
 \label{eq:newS7V90-cb-influence}
\end{equation}

\begin{proposition}[Completely bounded influence resolvent]
\label{prop:newS7V90-cb}
Suppose the declared local update scheme has an oscillation seminorm vector
\(v(F)\) satisfying, after one update,
\begin{equation}
 v_i(\mathcal TF)\leq\sum_jc_{ij}^{\rm cb}v_j(F),
 \label{eq:newS7V90-cb-recursion}
\end{equation}
and
\begin{equation}
 \alpha_{\rm cb}:=\sup_i\sum_jc_{ij}^{\rm cb}<1.
 \label{eq:newS7V90-cb-row}
\end{equation}
Then boundary-response oscillations are bounded by
\begin{equation}
 D_{\rm cb}:=(I-C_{\rm cb})^{-1},
 \label{eq:newS7V90-Dcb}
\end{equation}
and decay exponentially when \(C_{\rm cb}\) has finite range.  This
statement contains no scalar phase margin.
\end{proposition}

\begin{proof}
Iterate \eqref{eq:newS7V90-cb-recursion}.  The row norm below one makes the
Neumann series converge and sums all influence paths.  Finite-range decay
is the same path-length argument as
\Cref{lem:newS7V89-D-decay}.  Complete boundedness keeps the estimate stable
under adjoining untouched fermionic ancillas.
\end{proof}

\begin{firewall}[The recursion must be proved for the physical update]
\label{fire:newS7V90-cb-recursion}
The cb row is not automatic from complete positivity.  The actual
old/detail carrier update, its conditioning or boundary convention, and the
seminorm vector must be written, and
\eqref{eq:newS7V90-cb-recursion} must be proved.  A small spectral gap or a
finite local Hilbert space alone does not supply it.
\end{firewall}

\begin{definition}[Sign-resolved slab card, SRSC]
\label{def:newS7V90-SRSC}
SRSC has two disjoint branches.

\emph{Scalar complex branch:}
\begin{enumerate}
\item a common local phase lift and factorization
\eqref{eq:newS7V90-factor};
\item widths \(w_i<\pi\), margins \(\rho_i\), and boundary phase stocks
\(\chi_{ij}\);
\item the positive modulus influences \(c_{ij}\) and the strict complex row
\eqref{eq:newS7V90-complex-row};
\item nonzero global and intermediate normalizations and an exact complex
conditional factorization;
\item compatible V81 bad-step, overlap, RCLC, gauge, reflection, and Trotter
limits.
\end{enumerate}

\emph{Operator carrier branch:}
\begin{enumerate}
\item the positive retained carrier and local completely positive update
maps;
\item a declared stable cb oscillation seminorm;
\item the proved recursion \eqref{eq:newS7V90-cb-recursion} and strict row
\eqref{eq:newS7V90-cb-row};
\item uniform Trotter, cutoff, gauge--metric, and boundary limits; and
\item insertion of the resulting decay into ETCM, ITCM, IRCMC, and RFRC.
\end{enumerate}
\end{definition}

\begin{assessment}[Decision after the sign split]
\label{ass:newS7V90-status}
The scalar sign-resolved route is viable only with a strict local phase arc
and a much stronger row than the positive modulus Dobrushin condition.  A
balanced fermionic sign reaches the exact wall \(w=\pi\) and destroys the
conditional normalization.  The operator carrier route avoids that
normalization but now owes a concrete completely bounded influence
recursion.  The next upstream attack should construct this recursion for a
finite local Gibbs or transfer update rather than assume a scalar sign
cancellation.
\end{assessment}

\begin{programme}[V91: local carrier channel and cb contraction]
\label{prog:newS7V90-V91}
V91 will define a finite local heat-bath or Petz-type carrier replacement
channel relative to a faithful local Gibbs density, prove complete
positivity and preservation of constants, and estimate its boundary
variation by a Duhamel derivative.  It will identify a checkable cb row in
terms of local Hamiltonian boundary derivatives and a local density floor.
If the density floor collapses with cutoff, that collapse will be retained
as the exact upstream obstruction.
\end{programme}

\begin{firewall}[Claim boundary after seventh-series V90]
\label{fire:newS7V90-boundary}
V90 performs the compulsory companion audit; proves the local phase-width
loader, the modulus--phase complex influence compiler, its algebraic
Dobrushin resolvent, the binary sign wall, and a conditional completely
bounded carrier alternative.

V90 does not populate either SRSC branch for the physical Standard Model
carrier, prove a cb update recursion, a local density floor, or prove SRSC,
PSDC, FTCMC, IRCMC, NCRC, RFRC, CRARC, LPCC, OFHC, LCPLC, CCPC, GCMCC,
PCC, SMALC, RCLC, DSLC, CRDC, GAOCC, CCTC, MSRC, RTGC, or RCNC, remove
the cutoff, or construct the Standard Model--Einstein continuum.
\end{firewall}

\begin{theorem}[Seventh-series V90 result]
\label{thm:newS7V90-result}
A sign-resolved one-cell slab conditional with phase width \(w_i<\pi\)
has margin \(\rho_i=\cos(w_i/2)\), and its complex boundary influence is
bounded by \eqref{eq:newS7V90-influence}.  A strict row
\eqref{eq:newS7V90-complex-row} yields exponential algebraic conditional
response, but no positivity.  A balanced binary sign has zero conditional
normalization.  The non-scalar alternative is a positive retained-carrier
update with a separately proved completely bounded influence row.
\end{theorem}

\begin{proof}
Combine \Cref{thm:newS7V90-audit,
lem:newS7V90-width,
thm:newS7V90-influence,
thm:newS7V90-resolvent,
cth:newS7V90-sign,
prop:newS7V90-cb}.
\end{proof}

\paragraph{Parent-version record.}
V90 was formed from
\texttt{Renewal\_SM\_Einstein\_Continuum\_Research\_Notes\_V89.tex}.
The V89 parent had \(35\,924\) logical source lines, \(1\,410\,841\) bytes,
and SHA--256
\[
 \texttt{4F4213FC0B9009024F7A84C99B9EAEE43A07667825AC3819B4E675C28D1366A9}.
\]
The next compulsory companion audit is V95.

% =============================================================================
% END APPENDED SEVENTH-SERIES V90 MATERIAL
% =============================================================================
% =============================================================================
% APPENDED SEVENTH-SERIES V91 MATERIAL
% =============================================================================

\section{Faithful local carrier channels and completely bounded influence}
\label{sec:newS7V91}
\label{sec:v7v91}

V90 leaves an operator-carrier alternative when scalar slab signs destroy
local normalization.  This section constructs two finite local update
channels.  A product replacement channel has an exact diamond-norm
influence equal to the trace distance between its local Gibbs densities.
A Petz recovery channel retains a correlated cell--collar Gibbs block; on a
uniformly faithful density set its channel varies Lipschitz-continuously.
Combining these facts with the V85 Gibbs perturbation bound gives a
checkable completely bounded influence row.  The Petz constant exposes a
new ultraviolet debit: a collapsing local density floor.

\subsection{Product replacement channel}
\label{sec:newS7V91-replacement}

Let \(\mathcal H_i\) be a finite cell Hilbert space and let \(\sigma_i\) be
a density matrix.  Define the Schrodinger replacement channel
\begin{equation}
 \mathcal R_{\sigma_i}(X):=\operatorname{Tr}(X)\sigma_i.
 \label{eq:newS7V91-replacement}
\end{equation}
Its Heisenberg dual is
\(\mathcal R_{\sigma_i}^*(A)=\operatorname{Tr}(\sigma_iA)I\).

\begin{theorem}[Exact replacement-channel influence]
\label{thm:newS7V91-replacement}
The map \(\mathcal R_{\sigma_i}\) is completely positive and trace
preserving, and for any two density matrices
\begin{equation}
 \boxed{
 \|\mathcal R_{\sigma_i}-\mathcal R_{\tau_i}\|_\diamond
 =\|\sigma_i-\tau_i\|_1.}
 \label{eq:newS7V91-replacement-diamond}
\end{equation}
Thus a boundary-dependent product conditional
\(\sigma_i(\eta)\) has completely bounded influence
\begin{equation}
 c_{ij}^{\rm rep}={1\over2}
 \sup_{\eta\sim_j\eta'}
 \|\sigma_i(\eta)-\sigma_i(\eta')\|_1.
 \label{eq:newS7V91-replacement-influence}
\end{equation}
\end{theorem}

\begin{proof}
A Kraus representation follows by diagonalizing \(\sigma_i\) and mapping
an input basis to its eigenvectors with the square roots of the
eigenvalues; trace preservation is immediate from
\(\operatorname{Tr}\mathcal R_{\sigma_i}(X)=\operatorname{Tr}X\).
For an arbitrary ancilla and trace-class \(Y\),
\[
 [ (\mathcal R_{\sigma_i}-\mathcal R_{\tau_i})\otimes
 \operatorname{id}](Y)
 =(\sigma_i-\tau_i)\otimes\operatorname{Tr}_{i}Y.
\]
The trace norm is at most
\(\|\sigma_i-\tau_i\|_1\|Y\|_1\), because partial trace followed by this
replacement has the stated induced norm.  Equality is attained on a
product density input.  This proves
\eqref{eq:newS7V91-replacement-diamond} and then
\eqref{eq:newS7V91-replacement-influence}.
\end{proof}

\begin{corollary}[Local Gibbs Hamiltonian loader]
\label{cor:newS7V91-Gibbs-loader}
Let
\begin{equation}
 \sigma_i(\eta)={e^{-\beta h_i(\eta)}
 \over\operatorname{Tr}e^{-\beta h_i(\eta)}}
 \label{eq:newS7V91-local-Gibbs}
\end{equation}
and suppose
\begin{equation}
 \epsilon_{ij}:=sup_{\eta\sim_j\eta'}
 \|h_i(\eta)-h_i(\eta')\|<\infty.
 \label{eq:newS7V91-epsilon}
\end{equation}
Then
\begin{equation}
 c_{ij}^{\rm rep}
 \leq\beta\epsilon_{ij}e^{2\beta\epsilon_{ij}}.
 \label{eq:newS7V91-Gibbs-loader}
\end{equation}
\end{corollary}

\begin{proof}
Apply the normalized Gibbs trace-distance lemma of V85 to the two local
Hamiltonians and divide its bound by two.
\end{proof}

\begin{firewall}[Replacement discards quantum cell--collar correlations]
\label{fire:newS7V91-replacement}
The product channel is exact for a conditionally decoupled or classical
cell specification.  Applied to an entangled carrier it destroys the input
cell--collar correlations and need not preserve the intended interacting
Gibbs state.  Its good influence row is therefore not a certificate for
the physical carrier unless this preservation is separately proved.
\end{firewall}

\subsection{Petz recovery on a correlated collar}
\label{sec:newS7V91-Petz}

Let \(A\) be the updated cell and \(B\) a finite collar.  Let
\(\sigma_{AB}>0\) be a faithful density matrix and
\(\sigma_B=\operatorname{Tr}_A\sigma_{AB}\).  Define
\begin{equation}
 \mathcal P_{\sigma}:\mathcal T_1(\mathcal H_B)
 \longrightarrow\mathcal T_1(\mathcal H_A\otimes\mathcal H_B)
 \label{eq:newS7V91-Petz-domain}
\end{equation}
by
\begin{equation}
 \mathcal P_{\sigma}(X)
 :=\sigma_{AB}^{1/2}
 \left[I_A\otimes\sigma_B^{-1/2}X\sigma_B^{-1/2}\right]
 \sigma_{AB}^{1/2}.
 \label{eq:newS7V91-Petz}
\end{equation}
The associated cell--collar update is
\begin{equation}
 \mathcal H_A^{\sigma}:=
 \mathcal P_{\sigma}\circ\operatorname{Tr}_A.
 \label{eq:newS7V91-heatbath}
\end{equation}

\begin{theorem}[Petz heat-bath channel]
\label{thm:newS7V91-Petz}
The maps \(\mathcal P_{\sigma}\) and
\(\mathcal H_A^{\sigma}\) are completely positive and trace preserving.
They satisfy
\begin{equation}
 \mathcal P_{\sigma}(\sigma_B)=\sigma_{AB},
 \qquad
 \mathcal H_A^{\sigma}(\sigma_{AB})=\sigma_{AB}.
 \label{eq:newS7V91-Petz-recovery}
\end{equation}
Thus the correlated local Gibbs block is a fixed point of the update.
\end{theorem}

\begin{proof}
Equation \eqref{eq:newS7V91-Petz} is a composition of conjugations by fixed
operators and the embedding \(X\mapsto I_A\otimes X\), hence is completely
positive.  For trace-class \(X\), cyclicity and the definition of the
marginal give
\begin{align*}
 \operatorname{Tr}_{AB}\mathcal P_\sigma(X)
 &=\operatorname{Tr}_{AB}
 \left[I_A\otimes\sigma_B^{-1/2}X\sigma_B^{-1/2}\right]\sigma_{AB}\\
 &=\operatorname{Tr}_{B}
 \sigma_B^{-1/2}X\sigma_B^{-1/2}\sigma_B
 =\operatorname{Tr}_BX.
\end{align*}
So \(\mathcal P_\sigma\) is trace preserving; partial trace is CPTP, hence
so is their composition.  Substituting \(X=\sigma_B\) proves recovery, and
then
\(\mathcal H_A^\sigma(\sigma_{AB})=
\mathcal P_\sigma(\sigma_B)=\sigma_{AB}\).
\end{proof}

\begin{remark}[Local preservation is not global Gibbs preservation]
\label{rem:newS7V91-global}
When the channel is embedded in a larger carrier, preservation of one
collar density does not automatically preserve the global Gibbs state.
That requires a quantum Markov compatibility, commuting-square identity,
or an explicitly proved invariant global update scheme.
\end{remark}

\subsection{Faithful-density Lipschitz bound}
\label{sec:newS7V91-Lipschitz}

Let \(d_A=\dim\mathcal H_A\), \(d_B=\dim\mathcal H_B\), and define the
compact faithful set
\begin{equation}
 \mathfrak D_\lambda
 :=\{\sigma_{AB}\geq\lambda I_{AB}:
 \operatorname{Tr}\sigma_{AB}=1\},
 \qquad \lambda>0.
 \label{eq:newS7V91-Dlambda}
\end{equation}
For \(\sigma\in\mathfrak D_\lambda\), its marginal obeys
\begin{equation}
 \sigma_B\geq d_A\lambda I_B.
 \label{eq:newS7V91-marginal-floor}
\end{equation}

\begin{lemma}[Finite Petz Lipschitz constant]
\label{lem:newS7V91-Petz-Lipschitz}
There is a finite constant
\(L_{\rm P}(d_A,d_B,\lambda)\) such that for all
\(\sigma,\tau\in\mathfrak D_\lambda\),
\begin{equation}
 \|\mathcal H_A^\sigma-\mathcal H_A^\tau\|_\diamond
 \leq L_{\rm P}(d_A,d_B,\lambda)
 \|\sigma-\tau\|_1.
 \label{eq:newS7V91-Petz-Lipschitz}
\end{equation}
The constant is the supremum of the derivative norm of the explicit map
\(\sigma\mapsto\mathcal H_A^\sigma\) on
\(\mathfrak D_\lambda\) and is therefore a finite, directly computable
matrix quantity.  It diverges at worst by a finite inverse power of
\(\lambda\) as the floor tends to zero.
\end{lemma}

\begin{proof}
The maps \(X\mapsto X^{1/2}\) and
\(X\mapsto X^{-1/2}\) are continuously Frechet differentiable on matrices
whose spectra lie in \([\lambda,1]\) and
\([d_A\lambda,1]\), respectively.  Their derivatives have finite operator
norms there by the resolvent integral formulas for the two matrix
functions.  Partial trace is linear, and multiplication and tensor
embedding are bilinear bounded maps in finite dimension.  Therefore the
formula \eqref{eq:newS7V91-Petz} defines a continuously differentiable map
from \(\mathfrak D_\lambda\) into the finite-dimensional space of
superoperators equipped with diamond norm.  The derivative norm attains a
finite supremum on this compact set; call it \(L_{\rm P}\).  The segment
between \(\sigma\) and \(\tau\) stays in
\(\mathfrak D_\lambda\), so the fundamental theorem of calculus proves
\eqref{eq:newS7V91-Petz-Lipschitz}.  The resolvent derivatives of the
inverse square root contain only finite powers of the inverse spectral
floor, giving the final statement.
\end{proof}

\begin{theorem}[Petz boundary-influence loader]
\label{thm:newS7V91-Petz-loader}
Let a boundary-dependent faithful local Gibbs block be
\begin{equation}
 \sigma_{AB}(\eta)
 ={e^{-\beta h_{AB}(\eta)}
 \over\operatorname{Tr}e^{-\beta h_{AB}(\eta)}}
 \in\mathfrak D_\lambda
 \label{eq:newS7V91-block-Gibbs}
\end{equation}
uniformly.  If
\begin{equation}
 \epsilon_{ij}^{AB}:=sup_{\eta\sim_j\eta'}
 \|h_{AB}(\eta)-h_{AB}(\eta')\|,
 \label{eq:newS7V91-block-epsilon}
\end{equation}
then the channel influence
\begin{equation}
 c_{ij}^{\rm P}:={1\over2}
 \sup_{\eta\sim_j\eta'}
 \|\mathcal H_A^{\sigma(\eta)}-
 \mathcal H_A^{\sigma(\eta')}\|_\diamond
 \label{eq:newS7V91-P-influence}
\end{equation}
obeys
\begin{equation}
 \boxed{
 c_{ij}^{\rm P}
 \leq L_{\rm P}(d_A,d_B,\lambda)
 \beta\epsilon_{ij}^{AB}
 e^{2\beta\epsilon_{ij}^{AB}}.}
 \label{eq:newS7V91-Petz-loader}
\end{equation}
\end{theorem}

\begin{proof}
The V85 Gibbs perturbation lemma gives
\[
 \|\sigma_{AB}(\eta)-\sigma_{AB}(\eta')\|_1
 \leq2\beta\epsilon_{ij}^{AB}
 e^{2\beta\epsilon_{ij}^{AB}}.
\]
Insert this in \Cref{lem:newS7V91-Petz-Lipschitz} and divide by two.
\end{proof}

\subsection{An explicit local density floor}
\label{sec:newS7V91-floor}

\begin{lemma}[Gibbs floor from a local Hamiltonian norm]
\label{lem:newS7V91-floor}
If \(h=h^*\) acts on a \(d\)-dimensional collar and \(\|h\|\leq M\), then
its normalized Gibbs density satisfies
\begin{equation}
 {e^{-\beta h}\over\operatorname{Tr}e^{-\beta h}}
 \geq {e^{-2\beta M}\over d}I.
 \label{eq:newS7V91-floor}
\end{equation}
\end{lemma}

\begin{proof}
Every eigenvalue of \(e^{-\beta h}\) is at least \(e^{-\beta M}\), while
its trace is at most \(de^{\beta M}\).  Divide the two bounds.
\end{proof}

\begin{corollary}[The Petz ultraviolet debit]
\label{cor:newS7V91-floor-collapse}
For a fixed-size collar with uniformly bounded \(\beta M\),
\eqref{eq:newS7V91-floor} supplies a uniform \(\lambda>0\).  If the collar
dimension grows or \(\beta M\to\infty\), this elementary floor can collapse
and the Petz Lipschitz constant can diverge.  A sharper sector, energy, or
modular estimate is then required.
\end{corollary}

\begin{proof}
Substitute the bound \(\lambda=e^{-2\beta M}/d\) into
\Cref{lem:newS7V91-Petz-Lipschitz}.
\end{proof}

\subsection{A checkable completely bounded row}
\label{sec:newS7V91-row}

\begin{theorem}[Finite collar cb contraction]
\label{thm:newS7V91-row}
Suppose a local update scheme uses either the exact replacement branch or
the Petz branch, preserves the intended global carrier state, and satisfies
the V90 oscillation recursion.  Define the appropriate explicit upper
bounds
\begin{equation}
 \widehat c_{ij}^{\rm cb}:=
 \begin{cases}
 \beta\epsilon_{ij}e^{2\beta\epsilon_{ij}},
 &\text{replacement branch},\\
 L_{\rm P}(d_A,d_B,\lambda)
 \beta\epsilon_{ij}^{AB}e^{2\beta\epsilon_{ij}^{AB}},
 &\text{Petz branch}.
 \end{cases}
 \label{eq:newS7V91-cb-loader}
\end{equation}
If
\begin{equation}
 \widehat\alpha_{\rm cb}:=sup_i\sum_{j\ne i}
 \widehat c_{ij}^{\rm cb}<1,
 \label{eq:newS7V91-cb-row}
\end{equation}
then the completely bounded influence resolvent of V90 converges.  With
finite influence range it has exponential spatial decay, uniformly in
ancillary fermionic degrees of freedom.
\end{theorem}

\begin{proof}
The replacement bounds are
\Cref{thm:newS7V91-replacement,cor:newS7V91-Gibbs-loader}; the Petz bounds
are \Cref{thm:newS7V91-Petz-loader}.  They dominate the actual diamond
influences.  The strict row and the declared recursion invoke V90
\Cref{prop:newS7V90-cb}.  Diamond or cb norms are stable under tensoring
with an identity ancilla, giving the last assertion.
\end{proof}

\begin{firewall}[State preservation and the recursion remain independent]
\label{fire:newS7V91-preservation}
Complete positivity, a faithful local fixed point, and a strict numerical
row do not prove that the embedded updates preserve the full interacting
carrier or satisfy the V90 oscillation recursion.  Quantum conditional
expectations can fail to commute, and overlapping Petz maps can generate
additional memory.  Those two identities must be checked for the actual
update schedule.
\end{firewall}

\begin{definition}[Local carrier channel card, LCCC]
\label{def:newS7V91-LCCC}
LCCC requires:
\begin{enumerate}
\item fixed-size cell and collar algebras, the local Gibbs blocks, and their
boundary Hamiltonians;
\item a declared replacement, Petz, or other CPTP update with an exact
fixed-state theorem;
\item a uniform local density floor or a sharper modular substitute;
\item the boundary Hamiltonian differences and loaded diamond influences;
\item a proved invariant global carrier and the cb oscillation recursion;
\item the strict row \eqref{eq:newS7V91-cb-row} and its finite-range
resolvent decay;
\item uniform cutoff, Trotter, gauge, reflection, metric, and boundary
control; and
\item insertion into ETCM, ITCM, IRCMC, RFRC, and the Einstein stress row.
\end{enumerate}
\end{definition}

\begin{assessment}[What the carrier channel construction resolves]
\label{ass:newS7V91-status}
The operator alternative now has explicit local channels and a checkable
boundary sensitivity.  Product replacement is stable but generally loses
entanglement.  Petz recovery preserves the correlated local Gibbs block,
but its influence pays a faithfulness constant that may collapse in the
ultraviolet.  The next task is the overlap problem: determine when a family
of local Petz updates for overlapping collars shares the global Gibbs state
and satisfies a contraction recursion.
\end{assessment}

\begin{programme}[V92: commuting-square and approximate Markov conditions]
\label{prog:newS7V91-V92}
V92 will prove exact global preservation when overlapping local carrier
Hamiltonians form a commuting Gibbs specification.  It will then quantify
the failure for noncommuting collars through a conditional mutual
information or recovery defect, and insert that defect as an additive term
in the cb recursion.  This is the precise noncommutative replacement for
assuming that local heat baths define one global Markov field.
\end{programme}

\begin{firewall}[Claim boundary after seventh-series V91]
\label{fire:newS7V91-boundary}
V91 constructs and analyzes exact product replacement and correlated Petz
local channels; proves complete positivity, trace preservation, local
Gibbs recovery, faithful-set Lipschitz continuity, local Hamiltonian and
density-floor loaders, and a conditional strict cb row.

V91 does not prove global Gibbs preservation or the cb oscillation
recursion for overlapping physical collars, or prove LCCC, SRSC, PSDC,
FTCMC, IRCMC, NCRC, RFRC, CRARC, LPCC, OFHC, LCPLC, CCPC, GCMCC, PCC,
SMALC, RCLC, DSLC, CRDC, GAOCC, CCTC, MSRC, RTGC, or RCNC, remove the
cutoff, or construct the Standard Model--Einstein continuum.
\end{firewall}

\begin{theorem}[Seventh-series V91 result]
\label{thm:newS7V91-result}
A product replacement channel has diamond influence exactly equal to the
trace distance of its local target densities.  A correlated Petz update is
CPTP and fixes its faithful cell--collar Gibbs block; on a density floor
\(\lambda\) its channel variation is bounded by the finite constant
\(L_{\rm P}(d_A,d_B,\lambda)\).  Local Hamiltonian boundary changes load
both influences through the V85 Gibbs estimate, yielding the explicit cb
row \eqref{eq:newS7V91-cb-row}.  The remaining physical conditions are
global state preservation and contraction for overlapping updates.
\end{theorem}

\begin{proof}
Combine \Cref{thm:newS7V91-replacement,
cor:newS7V91-Gibbs-loader,
thm:newS7V91-Petz,
lem:newS7V91-Petz-Lipschitz,
thm:newS7V91-Petz-loader,
lem:newS7V91-floor,
thm:newS7V91-row}.
\end{proof}

\paragraph{Parent-version record.}
V91 was formed from
\texttt{Renewal\_SM\_Einstein\_Continuum\_Research\_Notes\_V90.tex}.
The V90 parent had \(36\,376\) logical source lines, \(1\,428\,044\) bytes,
and SHA--256
\[
 \texttt{AAD05866695C80843AB1B1D6EBD719BFE077849194CFD8C65CC465026AF556B8}.
\]
The next compulsory companion audit is V95.

% =============================================================================
% END APPENDED SEVENTH-SERIES V91 MATERIAL
% =============================================================================
% =============================================================================
% APPENDED SEVENTH-SERIES V92 MATERIAL
% =============================================================================

\section{Petz overlap consistency and approximate Markov defects}
\label{sec:newS7V92}
\label{sec:v7v92}

V91 constructs a faithful local Petz update but correctly leaves global
Gibbs preservation open.  This section identifies the exact state-level
condition.  On a tripartite collar \(A\!:\!B\!:\!C\), a commuting Markov
logarithm identity makes the Petz recovery of \(A\) from \(B\) reconstruct
the full state from its \(BC\) marginal.  Away from that branch, the trace
norm of the recovery error is the exact stationarity defect.  Relative
entropy bounds it, and a contracting update sweep converts the sum of local
defects into a controlled distance from its invariant state.

\subsection{Tripartite Petz recovery}
\label{sec:newS7V92-tripartite}

Let \(\sigma_{ABC}>0\) be a faithful finite density matrix, with faithful
marginals \(\sigma_{AB},\sigma_{BC},\sigma_B\).  Define the Petz recovery
map associated with \(\sigma_{AB}\) by
\begin{equation}
 \mathcal P_{B\to AB}^{\sigma}(X_{BC})
 :=\sigma_{AB}^{1/2}\sigma_B^{-1/2}
 X_{BC}\sigma_B^{-1/2}\sigma_{AB}^{1/2},
 \label{eq:newS7V92-Petz}
\end{equation}
where identities on untouched tensor factors are implicit.  This is the
V91 map tensored with the identity on \(C\).  Put
\begin{equation}
 \widehat\sigma_{ABC}:=
 \mathcal P_{B\to AB}^{\sigma}(\sigma_{BC}),
 \qquad
 \varepsilon_A^{\rm rec}:=
 \|\sigma_{ABC}-\widehat\sigma_{ABC}\|_1.
 \label{eq:newS7V92-recovery-defect}
\end{equation}

\begin{theorem}[Commuting Markov logarithm identity]
\label{thm:newS7V92-commuting}
Assume the four density matrices, embedded on \(ABC\), commute and obey
\begin{equation}
 \log\sigma_{ABC}+\log\sigma_B
 =\log\sigma_{AB}+\log\sigma_{BC}.
 \label{eq:newS7V92-log-Markov}
\end{equation}
Then
\begin{equation}
 \widehat\sigma_{ABC}=\sigma_{ABC},
 \qquad
 \varepsilon_A^{\rm rec}=0.
 \label{eq:newS7V92-exact-recovery}
\end{equation}
The conditional mutual information also vanishes:
\begin{equation}
 I_\sigma(A:C\mid B)
 :=S(AB)+S(BC)-S(B)-S(ABC)=0.
 \label{eq:newS7V92-CMI-zero}
\end{equation}
\end{theorem}

\begin{proof}
Commutativity permits ordinary functional calculus and reordering in
\eqref{eq:newS7V92-Petz}.  Therefore
\[
 \widehat\sigma_{ABC}
 =\exp(\log\sigma_{AB}+\log\sigma_{BC}-\log\sigma_B)
 =\sigma_{ABC}
\]
by \eqref{eq:newS7V92-log-Markov}.  For the entropy identity, use
\(S(\rho)=-\operatorname{Tr}\rho\log\rho\) and the fact that the trace of
an embedded marginal logarithm against \(\sigma_{ABC}\) equals the trace
against that marginal.  Thus
\[
 I_\sigma(A:C\mid B)
 =\operatorname{Tr}\sigma_{ABC}
 (\log\sigma_{ABC}+\log\sigma_B-log\sigma_{AB}-\log\sigma_{BC})=0.
\]
\end{proof}

\begin{remark}[Exact quantum Markov states are more general]
\label{rem:newS7V92-general-Markov}
The commuting logarithm identity is a sufficient branch chosen for a
direct proof.  General zero-conditional-mutual-information quantum states
also admit exact recovery, but their structure theorem is not needed here
and is not inferred for the physical carrier.
\end{remark}

\subsection{The recovery defect is the stationarity defect}
\label{sec:newS7V92-stationarity}

Define the globalized local heat-bath channel
\begin{equation}
 \mathcal H_A^\sigma
 :=\mathcal P_{B\to AB}^{\sigma}\circ\operatorname{Tr}_A
 \label{eq:newS7V92-heatbath}
\end{equation}
on the \(ABC\) block.

\begin{theorem}[Exact stationarity-defect identity]
\label{thm:newS7V92-stationarity}
For the target state \(\sigma_{ABC}\),
\begin{equation}
 \boxed{
 \|\mathcal H_A^\sigma(\sigma_{ABC})-sigma_{ABC}\|_1
 =\varepsilon_A^{\rm rec}.}
 \label{eq:newS7V92-stationarity}
\end{equation}
Hence the commuting Markov branch is fixed exactly, while every
noncommuting failure is measured without approximation by the recovery
defect.
\end{theorem}

\begin{proof}
Partial trace gives
\(\operatorname{Tr}_A\sigma_{ABC}=\sigma_{BC}\).  Substitution in
\eqref{eq:newS7V92-heatbath} yields
\(\mathcal H_A^\sigma(\sigma_{ABC})=widehat\sigma_{ABC}\), so the two
sides of \eqref{eq:newS7V92-stationarity} are identical by definition.
\end{proof}

\begin{corollary}[Overlapping exact updates share the target state]
\label{cor:newS7V92-overlap}
If every cell update in an overlapping family has zero recovery defect for
the same global state \(\sigma\), then every ordered composition of those
updates preserves \(\sigma\), whether or not the channels commute.
\end{corollary}

\begin{proof}
Each channel maps \(\sigma\) to itself.  Induction over the ordered
composition proves the claim.
\end{proof}

\subsection{Relative-entropy loader}
\label{sec:newS7V92-relative-entropy}

When \(\operatorname{supp}\sigma\subseteq
\operatorname{supp}\widehat\sigma\), define
\begin{equation}
 \Delta_A^{\rm P}:=
 D(\sigma_{ABC}\|\widehat\sigma_{ABC})
 =\operatorname{Tr}\sigma_{ABC}
 (\log\sigma_{ABC}-\log\widehat\sigma_{ABC}).
 \label{eq:newS7V92-Petz-entropy}
\end{equation}

\begin{lemma}[Finite quantum Pinsker bound]
\label{lem:newS7V92-Pinsker}
For finite density matrices \(\rho,\tau\),
\begin{equation}
 \|\rho-\tau\|_1\leq\sqrt{2D(\rho\|\tau)}.
 \label{eq:newS7V92-Pinsker}
\end{equation}
\end{lemma}

\begin{proof}
Measure the two states with the Helstrom two-outcome projection onto the
positive and negative parts of \(\rho-\tau\).  The resulting classical
total-variation distance is
\(\frac12\|\rho-\tau\|_1\).  Monotonicity of relative entropy under this
measurement and the classical Pinsker inequality give
\[
 D(\rho\|\tau)
 \geq2\left({1\over2}\|\rho-\tau\|_1\right)^2.
\]
Rearrange.
\end{proof}

\begin{corollary}[Entropy-to-stationarity loader]
\label{cor:newS7V92-entropy-loader}
The Petz update defect obeys
\begin{equation}
 \varepsilon_A^{\rm rec}
 \leq\sqrt{2\Delta_A^{\rm P}}.
 \label{eq:newS7V92-entropy-loader}
\end{equation}
Thus a local relative-entropy recovery estimate supplies the missing
stationarity row of V91.
\end{corollary}

\begin{proof}
Apply \Cref{lem:newS7V92-Pinsker} to
\(\sigma_{ABC}\) and \(\widehat\sigma_{ABC}\), then use
\Cref{thm:newS7V92-stationarity}.
\end{proof}

\begin{firewall}[Conditional mutual information is not substituted for the
Petz defect]
\label{fire:newS7V92-CMI}
For a general noncommuting state, this note does not claim
\(\Delta_A^{\rm P}=I(A:C\mid B)\) or that the ordinary Petz map is the
optimal approximate recovery.  The directly computed quantities are
\(\varepsilon_A^{\rm rec}\) and \(\Delta_A^{\rm P}\).  Conditional mutual
information is an exact zero diagnostic on the commuting branch only.
\end{firewall}

\subsection{Accumulation through an update sweep}
\label{sec:newS7V92-sweep}

Let \(\mathcal H_1,\ldots,\mathcal H_m\) be CPTP local updates and define
the ordered sweep
\begin{equation}
 \mathcal T:=\mathcal H_m\circ\cdots\circ\mathcal H_1.
 \label{eq:newS7V92-sweep}
\end{equation}
Suppose their defects on a target state \(\sigma\) are
\begin{equation}
 \varepsilon_i:=\|\mathcal H_i(\sigma)-\sigma\|_1.
 \label{eq:newS7V92-local-defects}
\end{equation}

\begin{lemma}[Sweep stationarity defect]
\label{lem:newS7V92-sweep-defect}
The full sweep obeys
\begin{equation}
 \|\mathcal T(\sigma)-\sigma\|_1
 \leq\sum_{i=1}^m\varepsilon_i.
 \label{eq:newS7V92-sweep-defect}
\end{equation}
\end{lemma}

\begin{proof}
Insert and subtract the partially applied sweeps with \(\sigma\) at the
next input.  Each later CPTP composition contracts trace distance on
Hermitian differences.  The triangle inequality then bounds the telescoping
sum by \(\sum_i\|\mathcal H_i(\sigma)-\sigma\|_1\).
\end{proof}

\begin{theorem}[Contracting sweep with approximate target]
\label{thm:newS7V92-contracting}
Assume \(\mathcal T\) has an invariant density \(\tau\) and contracts
traceless Hermitian differences by
\begin{equation}
 \|\mathcal T(X)\|_1\leq\alpha\|X\|_1,
 \qquad \operatorname{Tr}X=0,
 \qquad 0\leq\alpha<1.
 \label{eq:newS7V92-contraction}
\end{equation}
Then the invariant state is unique and
\begin{equation}
 \boxed{
 \|\tau-\sigma\|_1
 \leq{\sum_i\varepsilon_i\over1-\alpha}.}
 \label{eq:newS7V92-fixed-point-distance}
\end{equation}
If each update is a Petz channel, one may replace
\(\varepsilon_i\) by either its direct recovery defect or
\(\sqrt{2\Delta_i^{\rm P}}\).
\end{theorem}

\begin{proof}
If \(\tau_1,\tau_2\) are invariant, their difference is traceless and
\(\|\tau_1-\tau_2\|_1\leq\alpha\|\tau_1-\tau_2\|_1\), so they agree.  For
the target comparison,
\begin{align*}
 \|\tau-\sigma\|_1
 &\leq\|\mathcal T(\tau)-\mathcal T(\sigma)\|_1
 +\|\mathcal T(\sigma)-\sigma\|_1\\
 &\leq\alpha\|\tau-\sigma\|_1+\sum_i\varepsilon_i
\end{align*}
by \Cref{lem:newS7V92-sweep-defect}.  Rearrangement proves
\eqref{eq:newS7V92-fixed-point-distance}.  Use
\Cref{cor:newS7V92-entropy-loader} for the final option.
\end{proof}

\begin{corollary}[Observable bias]
\label{cor:newS7V92-observable}
For every observable \(O\),
\begin{equation}
 |\operatorname{Tr}(\tau-\sigma)O|
 \leq\|O\|{\sum_i\varepsilon_i\over1-\alpha}.
 \label{eq:newS7V92-observable}
\end{equation}
\end{corollary}

\begin{proof}
Apply trace duality to
\eqref{eq:newS7V92-fixed-point-distance}.
\end{proof}

\subsection{Local rather than extensive recovery sums}
\label{sec:newS7V92-local}

The sum over every cell in \eqref{eq:newS7V92-fixed-point-distance} is
extensive.  For local continuum work one needs a spatially weighted version.
Let \(O\) have support \(B\) and suppose the response of its Heisenberg
iterate to update \(i\) is bounded by a kernel \(G(B,i)\).

\begin{proposition}[Localized recovery-bias compiler]
\label{prop:newS7V92-local}
Assume an ordered-update telescoping in which
\begin{equation}
 \|\mathcal H_{<i}^*(O)-
 \mathcal H_{<i}^{*,\sigma}(O)\|
 \leq\|O\|G(B,i)
 \label{eq:newS7V92-response-kernel}
\end{equation}
controls how the local observable reaches update \(i\).  Then the one-sweep
bias obeys
\begin{equation}
 |\operatorname{Tr}[\mathcal T(\sigma)-\sigma]O|
 \leq\|O\|\sum_iG(B,i)\varepsilon_i.
 \label{eq:newS7V92-local-bias}
\end{equation}
If \(G(B,i)\leq C_G|B|e^{-\mu d(B,i)}\) and the recovery defects have a
uniform weighted row sum, the right side is volume local.
\end{proposition}

\begin{proof}
Telescope the sweep one update at a time in the Schrodinger picture and
move all preceding or following channels to the observable by duality.
At the insertion of update \(i\), the state difference has trace norm
\(\varepsilon_i\), while the reached observable has the response bound
\eqref{eq:newS7V92-response-kernel}.  Trace duality and summation give
\eqref{eq:newS7V92-local-bias}.  The final assertion is the exponential
row-sum estimate.
\end{proof}

\begin{firewall}[The response kernel is an additional mixing row]
\label{fire:newS7V92-response}
The localized bound avoids an extensive defect only after the Heisenberg
response kernel \(G\) is proved.  A strict local channel influence row is a
candidate loader, but the exact update ordering and overlap geometry must
show that it yields \eqref{eq:newS7V92-response-kernel}.
\end{firewall}

\subsection{Approximate quantum Markov card}
\label{sec:newS7V92-card}

\begin{definition}[Approximate carrier Markov card, ACMC]
\label{def:newS7V92-ACMC}
For every overlapping carrier collar, ACMC requires:
\begin{enumerate}
\item faithful \(ABC,AB,BC,B\) density matrices on one common global
carrier state;
\item either the exact commuting logarithm identity
\eqref{eq:newS7V92-log-Markov} or the direct Petz recovery defect
\eqref{eq:newS7V92-recovery-defect};
\item an optional relative-entropy certificate
\(\Delta_A^{\rm P}\) and the Pinsker loader;
\item the local density floors and boundary channel influences of LCCC;
\item a proved invariant or approximately invariant ordered update sweep;
\item a global contraction or the localized response kernel
\eqref{eq:newS7V92-response-kernel};
\item volume-local summability of recovery defects and boundary influences;
and
\item uniform gauge, reflection, cutoff, metric, and refinement
compatibility.
\end{enumerate}
\end{definition}

\begin{assessment}[The overlap obstruction is now measurable]
\label{ass:newS7V92-status}
Overlapping Petz updates preserve one global carrier exactly on the
commuting Markov branch.  In general, the direct recovery distance is the
exact local stationarity debit.  A contracting sweep turns its sum into
the explicit factor \((1-\alpha)^{-1}\), while a local observable needs a
spatial response kernel to avoid an extensive sum.  The next analytic task
is therefore to derive that response kernel from the cb influence matrix
and insert the recovery defects into an ETCM bound.
\end{assessment}

\begin{programme}[V93: local channel resolvent with recovery sources]
\label{prog:newS7V92-V93}
V93 will solve the inhomogeneous cb recursion
\(v\leq C_{\rm cb}v+b^{\rm rec}\), derive the explicit local resolvent
\((I-C_{\rm cb})^{-1}b^{\rm rec}\), and prove an equal-time covariance and
refinement-bias bound for local observables.  This will combine LCCC and
ACMC into a candidate ETCM loader while keeping all recovery sources
visible.
\end{programme}

\begin{firewall}[Claim boundary after seventh-series V92]
\label{fire:newS7V92-boundary}
V92 proves exact Petz recovery under a commuting Markov logarithm identity,
the stationarity-defect identity, relative-entropy recovery loader, sweep
defect accumulation, contracting fixed-point comparison, and a conditional
localized recovery-bias compiler.

V92 does not prove the physical Standard Model carrier is Markov, bound its
recovery defects or response kernel, or prove ACMC, LCCC, SRSC, PSDC,
FTCMC, IRCMC, NCRC, RFRC, CRARC, LPCC, OFHC, LCPLC, CCPC, GCMCC, PCC,
SMALC, RCLC, DSLC, CRDC, GAOCC, CCTC, MSRC, RTGC, or RCNC, remove the
cutoff, or construct the Standard Model--Einstein continuum.
\end{firewall}

\begin{theorem}[Seventh-series V92 result]
\label{thm:newS7V92-result}
A commuting Markov logarithm identity makes the local Petz channel recover
and preserve the full tripartite carrier exactly.  Away from that branch,
its trace-norm recovery error equals the local stationarity defect and is
at most \(\sqrt{2\Delta_A^{\rm P}}\).  An ordered contracting sweep has a
unique invariant state within
\((1-\alpha)^{-1}\sum_i\varepsilon_i\) of the target, and a declared local
response kernel replaces this extensive sum by the weighted local bound
\eqref{eq:newS7V92-local-bias}.
\end{theorem}

\begin{proof}
Combine \Cref{thm:newS7V92-commuting,
thm:newS7V92-stationarity,
lem:newS7V92-Pinsker,
cor:newS7V92-entropy-loader,
lem:newS7V92-sweep-defect,
thm:newS7V92-contracting,
prop:newS7V92-local}.
\end{proof}

\paragraph{Parent-version record.}
V92 was formed from
\texttt{Renewal\_SM\_Einstein\_Continuum\_Research\_Notes\_V91.tex}.
The V91 parent had \(36\,821\) logical source lines, \(1\,444\,746\) bytes,
and SHA--256
\[
 \texttt{B052023F3B32BCAEE354CB75968025B370A05489A67C7E675313BFEE811619DE}.
\]
The next compulsory companion audit is V95.

% =============================================================================
% END APPENDED SEVENTH-SERIES V92 MATERIAL
% =============================================================================
% =============================================================================
% APPENDED SEVENTH-SERIES V93 MATERIAL
% =============================================================================

\section{Inhomogeneous channel resolvents and local Duhamel response}
\label{sec:newS7V93}
\label{sec:v7v93}

V92 leaves two local sources in an overlapping update scheme: boundary
variation of the channels and their recovery defect on the intended Gibbs
state.  Both enter an inhomogeneous contraction recursion.  This section
solves that recursion, proves exponential propagation for finite-range
influences, and differentiates a Gibbs-preserving channel fixed-point
equation.  The derivative is the same Duhamel response found in V87, so a
local channel source bound becomes a direct ITCM loader without separately
assuming ETCM.

\subsection{Abstract local seminorm packet}
\label{sec:newS7V93-seminorms}

Let \(\mathfrak X\) be the real vector space of traceless Hermitian
trace-class operators on a finite carrier.  Fix nonnegative seminorms
\(v_i:\mathfrak X\to[0,\infty)\), indexed by cells, with the following
declared dual property: for every local observable \(A\) supported in
\(B\),
\begin{equation}
 |\operatorname{Tr}(XA)|
 \leq\|A\|\sum_{i\in B}v_i(X).
 \label{eq:newS7V93-dual}
\end{equation}
The packet may be built from local conditional trace norms, quotient
oscillations, or another cb-stable construction, but
\eqref{eq:newS7V93-dual} must be proved for the chosen version.

Let \(\mathcal T\) be a CPTP update sweep and suppose its action on
traceless differences satisfies
\begin{equation}
 v_i(\mathcal T X)\leq\sum_jC_{ij}v_j(X),
 \qquad C_{ij}\geq0,
 \qquad
 \alpha:=\sup_i\sum_jC_{ij}<1.
 \label{eq:newS7V93-contraction}
\end{equation}
Define
\begin{equation}
 D:=(I-C)^{-1}=\sum_{n=0}^\infty C^n.
 \label{eq:newS7V93-D}
\end{equation}

\begin{lemma}[Inhomogeneous resolvent]
\label{lem:newS7V93-resolvent}
If nonnegative vectors \(u,b\) obey
\begin{equation}
 u\leq Cu+b
 \label{eq:newS7V93-inhomogeneous}
\end{equation}
componentwise, and \(u\) is bounded, then
\begin{equation}
 \boxed{u\leq Db.}
 \label{eq:newS7V93-resolvent}
\end{equation}
\end{lemma}

\begin{proof}
Iterate \eqref{eq:newS7V93-inhomogeneous}:
\[
 u\leq C^Nu+\sum_{n=0}^{N-1}C^nb.
\]
The row norm of \(C^N\) is at most \(\alpha^N\), so the first term tends
to zero for bounded \(u\).  The finite sum converges monotonically to
\(Db\).
\end{proof}

\subsection{Localized decay of an inhomogeneous source}
\label{sec:newS7V93-decay}

Assume \(C_{ij}=0\) for \(d(i,j)>R_c\).  Then, as in V89,
\begin{equation}
 D_{ij}\leq C_D e^{-\gamma d(i,j)},
 \qquad
 C_D={\alpha^{-1}\over1-\alpha},
 \qquad
 \gamma={|\log\alpha|\over R_c}
 \label{eq:newS7V93-D-decay}
\end{equation}
for \(i\ne j\).

\begin{theorem}[Recovery-source localization]
\label{thm:newS7V93-source}
Let a source localized near \(X\) obey
\begin{equation}
 b_j\leq b_0\sum_{x\in X}e^{-\nu d(j,x)}.
 \label{eq:newS7V93-source}
\end{equation}
Suppose the graph has uniform exponential convolution sums: for every
\(0<\mu<\min\{\gamma,\nu\}\),
\begin{equation}
 S_{\gamma,\nu,\mu}:=\sup_{i,x}
 e^{\mu d(i,x)}\sum_j
 e^{-\gamma d(i,j)}e^{-\nu d(j,x)}<\infty.
 \label{eq:newS7V93-convolution}
\end{equation}
Then every bounded solution of \eqref{eq:newS7V93-inhomogeneous} satisfies
\begin{equation}
 u_i\leq C_Db_0S_{\gamma,\nu,\mu}
 |X|e^{-\mu d(i,X)}
 \label{eq:newS7V93-source-decay}
\end{equation}
away from the harmless diagonal contribution, which is absorbed by
enlarging the constant.
\end{theorem}

\begin{proof}
Use \Cref{lem:newS7V93-resolvent} and expand
\((Db)_i=\sum_jD_{ij}b_j\).  Insert
\eqref{eq:newS7V93-D-decay} and
\eqref{eq:newS7V93-source}; the convolution hypothesis gives
\eqref{eq:newS7V93-source-decay}.  The term \(D_{ii}\) is bounded by
\((1-\alpha)^{-1}\) and satisfies the same estimate after the constant is
enlarged.
\end{proof}

\begin{corollary}[Local observable response]
\label{cor:newS7V93-observable}
For \(A\) supported in \(B\),
\begin{equation}
 |\operatorname{Tr}(XA)|
 \leq C_*\|A\||B||X|e^{-\mu d(B,X)},
 \label{eq:newS7V93-observable}
\end{equation}
whenever \(u_i=v_i(X)\) has the source bound of
\Cref{thm:newS7V93-source}.  The constant \(C_*\) is independent of total
volume.
\end{corollary}

\begin{proof}
Insert \eqref{eq:newS7V93-source-decay} into
\eqref{eq:newS7V93-dual}, sum over \(i\in B\), and use
\(d(i,X)\geq d(B,X)\).
\end{proof}

\subsection{Approximate recovery as an inhomogeneous source}
\label{sec:newS7V93-recovery}

Let \(\sigma\) be the intended carrier state and let \(\tau\) be an
invariant state of \(\mathcal T\).  Define the local sweep recovery source
\begin{equation}
 b_i^{\rm rec}:=v_i(\mathcal T\sigma-\sigma).
 \label{eq:newS7V93-recovery-source}
\end{equation}

\begin{theorem}[Local invariant-state bias]
\label{thm:newS7V93-bias}
Under \eqref{eq:newS7V93-contraction},
\begin{equation}
 v_i(\tau-\sigma)
 \leq(D b^{\rm rec})_i.
 \label{eq:newS7V93-bias}
\end{equation}
If the local Petz recovery defects and the update ordering load
\eqref{eq:newS7V93-source}, then the bias of every local observable obeys
\eqref{eq:newS7V93-observable}.
\end{theorem}

\begin{proof}
Since \(\tau=\mathcal T\tau\),
\[
 \tau-\sigma=\mathcal T(\tau-\sigma)
 +(\mathcal T\sigma-\sigma).
\]
Apply the seminorms, use \eqref{eq:newS7V93-contraction}, and obtain
\(u\leq Cu+b^{\rm rec}\).  Invoke
\Cref{lem:newS7V93-resolvent}.  The final statement is
\Cref{thm:newS7V93-source,cor:newS7V93-observable}.
\end{proof}

\begin{remark}[Loading the recovery source]
\label{rem:newS7V93-loading}
V92 bounds the global trace defect of each Petz update by
\(\varepsilon_i^{\rm rec}\) or
\(\sqrt{2\Delta_i^{\rm P}}\).  To load
\eqref{eq:newS7V93-recovery-source} locally, one must retain the support of
each update and propagate only through the preceding update channels.  A
sum of all scalar defects discards this geometry and returns the extensive
V92 bound.
\end{remark}

\subsection{Differentiated fixed-point equation}
\label{sec:newS7V93-fixed-point}

Let \(\lambda\mapsto K_\lambda=K+\lambda B_X\) be a local Hamiltonian
perturbation and let
\begin{equation}
 \rho_\lambda={e^{-\beta K_\lambda}
 \over\operatorname{Tr}e^{-\beta K_\lambda}}.
 \label{eq:newS7V93-rho-lambda}
\end{equation}
Suppose a differentiable channel sweep \(\mathcal T_\lambda\) preserves
this Gibbs state exactly:
\begin{equation}
 \mathcal T_\lambda(\rho_\lambda)=\rho_\lambda.
 \label{eq:newS7V93-fixed-point}
\end{equation}
Assume the same contraction matrix \(C\) works uniformly in \(\lambda\).

Define the differentiated channel source at \(\lambda=0\) by
\begin{equation}
 b_i^{B_X}:=
 v_i\bigl(\dot{\mathcal T}_0(\rho_0)\bigr).
 \label{eq:newS7V93-channel-source}
\end{equation}

\begin{theorem}[Channel source controls Gibbs linear response]
\label{thm:newS7V93-linear-response}
Let \(u_i=v_i(\dot\rho_0)\).  Then
\begin{equation}
 u\leq Cu+b^{B_X},
 \qquad
 u\leq D b^{B_X}.
 \label{eq:newS7V93-linear-response}
\end{equation}
If \(b^{B_X}\) satisfies the localized source bound
\eqref{eq:newS7V93-source}, then for every local \(A_B\),
\begin{equation}
 \left|{d\over d\lambda}
 \operatorname{Tr}(\rho_\lambda A_B)\bigg|_{\lambda=0}\right|
 \leq C_*\|A_B\||B||X|e^{-\mu d(B,X)}.
 \label{eq:newS7V93-response-decay}
\end{equation}
\end{theorem}

\begin{proof}
Differentiate \eqref{eq:newS7V93-fixed-point}:
\[
 \dot\rho_0=\mathcal T_0(\dot\rho_0)
 +\dot{\mathcal T}_0(\rho_0).
\]
Apply the seminorm contraction to obtain the first inequality in
\eqref{eq:newS7V93-linear-response}; the inhomogeneous resolvent gives the
second.  The local response is
\(\operatorname{Tr}(\dot\rho_0A_B)\), so
\Cref{thm:newS7V93-source,cor:newS7V93-observable} prove
\eqref{eq:newS7V93-response-decay}.
\end{proof}

\begin{corollary}[Direct ITCM loader]
\label{cor:newS7V93-ITCM}
Under the hypotheses of \Cref{thm:newS7V93-linear-response},
\begin{equation}
 \beta\left|
 \operatorname{Cov}^{\rm D}_{\rho_0}(A_B,B_X)
 \right|
 \leq C_*\|A_B\||B||X|e^{-\mu d(B,X)}.
 \label{eq:newS7V93-ITCM}
\end{equation}
Thus localized differentiated channel sources load the V87 ITCM estimate
directly.  No separate equal-time mixing theorem or scalar phase margin is
needed on this operator branch.
\end{corollary}

\begin{proof}
The V87 Gibbs response identity for the perturbation
\(K+\lambda B_X\) gives
\[
 {d\over d\lambda}\omega_\lambda(A_B)\big|_0
 =-\beta\operatorname{Cov}^{\rm D}_{\rho_0}(A_B,B_X).
\]
Combine it with \eqref{eq:newS7V93-response-decay}.
\end{proof}

\begin{firewall}[The channel source is not automatically local]
\label{fire:newS7V93-source}
A local Hamiltonian perturbation \(B_X\) can change normalization and
modular inverses in every Petz update.  The locality of
\(b_i^{B_X}=v_i(\dot{\mathcal T}\rho)\) must be proved from the actual
faithful collar maps, density floors, and update schedule.  It is not a
formal consequence of the support of \(B_X\).
\end{firewall}

\subsection{Boundary perturbations and covariance}
\label{sec:newS7V93-boundary}

\begin{proposition}[Fixed-point boundary influence]
\label{prop:newS7V93-boundary}
Let \(\mathcal T^\eta\) and \(\mathcal T^{\eta'}\) have invariant states
\(\rho^\eta,\rho^{\eta'}\), share the contraction matrix \(C\), and let
\begin{equation}
 b_i^{\eta,\eta'}:=
 v_i\bigl((\mathcal T^\eta-\mathcal T^{\eta'})
ho^{\eta'}\bigr).
 \label{eq:newS7V93-boundary-source}
\end{equation}
Then
\begin{equation}
 v(\rho^\eta-\rho^{\eta'})
 \leq D b^{\eta,\eta'}.
 \label{eq:newS7V93-boundary}
\end{equation}
The V91 diamond channel differences give a direct upper bound on the source
components after the chosen seminorm normalization.
\end{proposition}

\begin{proof}
Subtract the two fixed-point identities and write
\[
 \rho^\eta-\rho^{\eta'}
 =\mathcal T^\eta(\rho^\eta-\rho^{\eta'})
 +(\mathcal T^\eta-\mathcal T^{\eta'})\rho^{\eta'}.
\]
Apply the contraction and
\Cref{lem:newS7V93-resolvent}.  Diamond norm controls a channel difference
on every density input, including ancillary collar systems.
\end{proof}

\begin{definition}[Inhomogeneous carrier contraction card, ICCC]
\label{def:newS7V93-ICCC}
ICCC requires:
\begin{enumerate}
\item a cb-stable local seminorm packet with the dual bound
\eqref{eq:newS7V93-dual};
\item a physical Gibbs-preserving or recovery-controlled channel sweep;
\item the uniform finite-range contraction matrix \(C\) with row margin
\(\alpha<1\);
\item localized recovery, boundary, Hamiltonian, current, and metric
channel sources;
\item uniform exponential convolution sums and the resolvent response
bounds;
\item differentiated fixed-point identities and the direct ITCM loader;
\item summable refinement tails for the V87 old/detail couplings; and
\item compatibility with LCCC, ACMC, RFRC, gauge projection, reflection,
zero modes, and RCNC.
\end{enumerate}
\end{definition}

\begin{assessment}[The operator mixing target is now finite and local]
\label{ass:newS7V93-status}
The missing completely bounded route is no longer a request for an
unspecified mixing theorem.  It requires a proved seminorm contraction
matrix and localized inhomogeneous sources.  Their Neumann resolvent gives
both recovery bias and Duhamel response decay.  The first physical
calculation is the differentiated Petz or heat-bath channel source created
by a local Standard Model Hamiltonian or metric insertion.
\end{assessment}

\begin{programme}[V94: metric-differentiated channel source]
\label{prog:newS7V93-V94}
V94 will differentiate the finite Petz channel with respect to a local
metric or vierbein parameter.  Using resolvent formulas for square roots
and inverse square roots, it will bound the channel source in terms of the
local Gibbs derivative, the density floor, and distance from the metric
insertion.  This will connect ICCC directly to the retained-carrier
Einstein stress row before the V95 companion audit.
\end{programme}

\begin{firewall}[Claim boundary after seventh-series V93]
\label{fire:newS7V93-boundary}
V93 solves the inhomogeneous local channel recursion, proves exponential
localization of recovery and boundary sources, derives local invariant-state
bias, differentiates a Gibbs-preserving fixed-point equation, and turns a
localized channel source directly into the V87 ITCM bound.

V93 does not construct the physical seminorm contraction matrix or bound
its Petz channel sources, or prove ICCC, ACMC, LCCC, SRSC, PSDC, FTCMC,
IRCMC, NCRC, RFRC, CRARC, LPCC, OFHC, LCPLC, CCPC, GCMCC, PCC, SMALC,
RCLC, DSLC, CRDC, GAOCC, CCTC, MSRC, RTGC, or RCNC, remove the cutoff, or
construct the Standard Model--Einstein continuum.
\end{firewall}

\begin{theorem}[Seventh-series V93 result]
\label{thm:newS7V93-result}
For a finite-range local contraction matrix with row sum below one, every
recovery, boundary, or differentiated Hamiltonian source propagates through
\((I-C)^{-1}\) and remains exponentially localized.  If a differentiable
channel sweep preserves the Gibbs family, its fixed-point derivative obeys
this inhomogeneous recursion; the exact Gibbs response identity then gives
the Duhamel mixing bound \eqref{eq:newS7V93-ITCM}.  This supplies a direct
operator-channel route to interacting retained-carrier projectivity.
\end{theorem}

\begin{proof}
Combine \Cref{lem:newS7V93-resolvent,
thm:newS7V93-source,
cor:newS7V93-observable,
thm:newS7V93-bias,
thm:newS7V93-linear-response,
cor:newS7V93-ITCM,
prop:newS7V93-boundary}.
\end{proof}

\paragraph{Parent-version record.}
V93 was formed from
\texttt{Renewal\_SM\_Einstein\_Continuum\_Research\_Notes\_V92.tex}.
The V92 parent had \(37\,244\) logical source lines, \(1\,460\,015\) bytes,
and SHA--256
\[
 \texttt{804DB4C9D5C3D7D2E169F0BF382A3632B7FD0DD39823473EF5FB8C546BB2D6A7}.
\]
The next compulsory companion audit is V95.

% =============================================================================
% END APPENDED SEVENTH-SERIES V93 MATERIAL
% =============================================================================
% =============================================================================
% APPENDED SEVENTH-SERIES V94 MATERIAL
% =============================================================================

\section{Metric-differentiated Petz channels and local fermionic stress}
\label{sec:newS7V94}
\label{sec:v7v94}

V93 identifies the differentiated channel source
\(b_i^{B_X}=v_i(\dot{\mathcal T}\rho)\) as the missing input for direct
Duhamel mixing.  This section computes a finite local bound for a Petz
update whose cell--collar Gibbs density depends on a metric, vierbein, or
spin-connection parameter.  The derivative of the local Gibbs density is
a Duhamel integral; the faithful-set Lipschitz constant of V91 converts it
into a diamond-norm derivative of the Petz channel.  Finite-range or
exponentially local metric dependence then loads the V93 source kernel and
the retained-carrier Einstein stress row.

\subsection{Derivative of a normalized local Gibbs density}
\label{sec:newS7V94-Gibbs}

Let \(q\mapsto h(q)=h(q)^*\) be a differentiable Hamiltonian on a fixed
finite collar and define
\begin{equation}
 \sigma(q)={e^{-\beta h(q)}\over Z(q)},
 \qquad
 Z(q)=\operatorname{Tr}e^{-\beta h(q)}.
 \label{eq:newS7V94-sigma}
\end{equation}

\begin{theorem}[Exact local Gibbs derivative]
\label{thm:newS7V94-Gibbs-derivative}
The density derivative is
\begin{equation}
 \dot\sigma
 =-\beta\int_0^1
 \sigma^{1-s}
 \bigl(\dot h-\operatorname{Tr}(\sigma\dot h)I\bigr)
 \sigma^s\,ds.
 \label{eq:newS7V94-Gibbs-derivative}
\end{equation}
It has trace zero and satisfies the dimension-free trace-norm bound
\begin{equation}
 \boxed{
 \|\dot\sigma\|_1\leq2\beta\|\dot h\|.}
 \label{eq:newS7V94-Gibbs-bound}
\end{equation}
Adding a scalar function of \(q\) to \(h(q)\) leaves
\(\dot\sigma\) unchanged.
\end{theorem}

\begin{proof}
Differentiate the exponential by the V87 Duhamel formula and use
\(\dot Z=-\beta Z\operatorname{Tr}(\sigma\dot h)\).  Combining the
numerator and normalization derivatives gives
\eqref{eq:newS7V94-Gibbs-derivative}; its trace is the derivative of
\(\operatorname{Tr}\sigma=1\).  Generalized Schatten Holder gives
\[
 \|\sigma^{1-s}X\sigma^s\|_1\leq\|X\|
\]
for \(0\leq s\leq1\), while
\(\|\dot h-\operatorname{Tr}(\sigma\dot h)I\|leq2\|\dot h\|\).  Integrate
to obtain \eqref{eq:newS7V94-Gibbs-bound}.  A scalar shift cancels inside
the centered derivative.
\end{proof}

\begin{corollary}[Local metric insertion]
\label{cor:newS7V94-local-Gibbs}
If a unit metric tangent at cell \(x\) gives
\begin{equation}
 \|d_xh_i\|\leq C_he^{-\nu d(i,x)},
 \label{eq:newS7V94-h-local}
\end{equation}
for the local Gibbs block used by update \(i\), then
\begin{equation}
 \|d_x\sigma_i\|_1
 \leq2\beta C_he^{-\nu d(i,x)}.
 \label{eq:newS7V94-sigma-local}
\end{equation}
For a strictly finite-range dependence, the right side is zero outside
that range.
\end{corollary}

\begin{proof}
Apply \eqref{eq:newS7V94-Gibbs-bound} to
\eqref{eq:newS7V94-h-local}.
\end{proof}

\subsection{Derivative of the Petz update}
\label{sec:newS7V94-Petz}

Let \(\mathcal H_A^{\sigma(q)}\) be the V91 Petz heat-bath channel on a
fixed cell--collar algebra.  Assume
\begin{equation}
 \sigma(q)\geq\lambda I
 \label{eq:newS7V94-floor}
\end{equation}
throughout a parameter interval, with fixed dimensions.

\begin{theorem}[Petz channel derivative bound]
\label{thm:newS7V94-Petz-derivative}
The channel is differentiable in diamond norm and
\begin{equation}
 \boxed{
 \left\|{d\over dq}\mathcal H_A^{\sigma(q)}\right\|_\diamond
 \leq L_{\rm P}(d_A,d_B,\lambda)\|\dot\sigma(q)\|_1
 \leq2\beta L_{\rm P}(d_A,d_B,\lambda)\|\dot h(q)\|.}
 \label{eq:newS7V94-Petz-derivative}
\end{equation}
Here \(L_{\rm P}\) is the faithful-set derivative supremum defined in
V91.
\end{theorem}

\begin{proof}
By definition,
\(L_{\rm P}\) bounds the derivative of the finite smooth map
\(\sigma\mapsto\mathcal H_A^\sigma\) from trace norm to diamond norm on
the faithful set.  The chain rule gives the first inequality.  Use
\Cref{thm:newS7V94-Gibbs-derivative} for the second.
\end{proof}

\begin{corollary}[Exponential metric-channel source]
\label{cor:newS7V94-channel-local}
Under \eqref{eq:newS7V94-h-local},
\begin{equation}
 \|d_x\mathcal H_i\|_\diamond
 \leq C_{\mathcal H}e^{-\nu d(i,x)},
 \qquad
 C_{\mathcal H}:=2\beta L_{\rm P}C_h.
 \label{eq:newS7V94-channel-local}
\end{equation}
For every density matrix \(\rho\),
\begin{equation}
 \|(d_x\mathcal H_i)(\rho)\|_1
 \leq C_{\mathcal H}e^{-\nu d(i,x)}.
 \label{eq:newS7V94-channel-state}
\end{equation}
\end{corollary}

\begin{proof}
Insert \eqref{eq:newS7V94-h-local} into
\eqref{eq:newS7V94-Petz-derivative}.  The diamond norm dominates the
trace norm on every unancilled density input.
\end{proof}

\begin{firewall}[The density floor is differentiated only through a uniform
set]
\label{fire:newS7V94-floor}
The proof requires one fixed \(\lambda>0\) on the whole parameter path.  A
pointwise faithful density whose least eigenvalue approaches zero can have
an unbounded inverse-square-root derivative.  Differentiability at each
cutoff does not supply a uniform Einstein stress bound.
\end{firewall}

\subsection{Derivative of an ordered update sweep}
\label{sec:newS7V94-sweep}

Let
\begin{equation}
 \mathcal T(q)=\mathcal H_m(q)\circ\cdots\circ\mathcal H_1(q).
 \label{eq:newS7V94-sweep}
\end{equation}

\begin{lemma}[Sweep derivative telescoping]
\label{lem:newS7V94-sweep}
For every density \(\rho\),
\begin{equation}
 \dot{\mathcal T}(\rho)
 =\sum_{k=1}^m
 \mathcal H_m\circ\cdots\circ\mathcal H_{k+1}
 \circ\dot{\mathcal H}_k\circ
 \mathcal H_{k-1}\circ\cdots\circ\mathcal H_1(\rho).
 \label{eq:newS7V94-sweep-derivative}
\end{equation}
The global trace norm is bounded by
\begin{equation}
 \|\dot{\mathcal T}(\rho)\|_1
 \leq\sum_{k=1}^m\|\dot{\mathcal H}_k\|_\diamond.
 \label{eq:newS7V94-sweep-global}
\end{equation}
\end{lemma}

\begin{proof}
Differentiate the finite composition by the product rule.  Every undifferentiated
channel is CPTP and contracts trace norm on Hermitian differences; the input
to each differentiated channel is a density.  Apply the triangle inequality.
\end{proof}

The global sum is extensive.  Let \(v_i\) be the V93 local seminorm packet,
and suppose the undifferentiated channel schedule propagates an insertion
at update \(k\) to cell \(i\) by a nonnegative kernel \(G_{ik}\):
\begin{equation}
 v_i\bigl(
 \mathcal H_m\circ\cdots\circ\mathcal H_{k+1}(X_k)
 \bigr)
 \leq G_{ik}\|X_k\|_1.
 \label{eq:newS7V94-propagation}
\end{equation}

\begin{theorem}[Localized differentiated sweep source]
\label{thm:newS7V94-sweep-source}
If
\begin{equation}
 G_{ik}\leq C_Ge^{-\gamma d(i,k)}
 \label{eq:newS7V94-G}
\end{equation}
and the local channel derivatives obey
\eqref{eq:newS7V94-channel-local}, then
\begin{equation}
 b_i^{x}:=v_i(\dot{\mathcal T}(\rho))
 \leq C_GC_{\mathcal H}
 \sum_k e^{-\gamma d(i,k)}e^{-\nu d(k,x)}.
 \label{eq:newS7V94-source-convolution}
\end{equation}
Under the V93 exponential convolution row, for every
\(0<\mu<\min\{\gamma,\nu\}\),
\begin{equation}
 b_i^{x}\leq C_be^{-\mu d(i,x)}
 \label{eq:newS7V94-source-local}
\end{equation}
with \(C_b\) independent of total volume.
\end{theorem}

\begin{proof}
Apply \eqref{eq:newS7V94-propagation} to each term of
\eqref{eq:newS7V94-sweep-derivative}.  The trace norm of the differentiated
channel output is bounded by
\eqref{eq:newS7V94-channel-state}; this gives
\eqref{eq:newS7V94-source-convolution}.  The convolution row is exactly
V93 \eqref{eq:newS7V93-convolution} for a singleton source.
\end{proof}

\begin{corollary}[Metric source propagated by the cb resolvent]
\label{cor:newS7V94-resolvent}
If the sweep preserves the differentiable Gibbs family and has the V93
finite-range contraction matrix \(C\) with row sum below one, then
\begin{equation}
 v(\dot\rho)\leq(I-C)^{-1}b^x,
 \label{eq:newS7V94-resolvent}
\end{equation}
and every local observable \(A_B\) obeys
\begin{equation}
 |D_x\omega(A_B)-\omega(D_xA_B)|
 \leq C_{\rm met}\|A_B\||B|e^{-\mu' d(B,x)}
 \label{eq:newS7V94-metric-response}
\end{equation}
for some \(\mu'>0\) fixed by the source and influence exponents.
\end{corollary}

\begin{proof}
Use \Cref{thm:newS7V93-linear-response} with the source
\eqref{eq:newS7V94-source-local}, then apply the V93 dual local observable
bound.  The direct derivative \(D_xA_B\) is separated on the left.
\end{proof}

\subsection{Fermionic Einstein stress row}
\label{sec:newS7V94-stress}

Let \(D_xK\) be the retained-carrier metric insertion at cell \(x\).

\begin{theorem}[Localized retained-carrier stress covariance]
\label{thm:newS7V94-stress}
Under the hypotheses of
\Cref{cor:newS7V94-resolvent},
\begin{equation}
 \beta\left|
 \operatorname{Cov}^{\rm D}_\rho(A_B,D_xK)
 \right|
 \leq C_{\rm met}\|A_B\||B|e^{-\mu' d(B,x)}.
 \label{eq:newS7V94-stress}
\end{equation}
If \(A_B\) itself is metric independent, this is the complete derivative
of its retained-carrier expectation.  For a metric-dependent observable,
\(\omega(D_xA_B)\) is a separate direct contact term.
\end{theorem}

\begin{proof}
The V87 stress identity is
\[
 D_x\omega(A_B)-\omega(D_xA_B)
 =-\beta\operatorname{Cov}^{\rm D}_\rho(A_B,D_xK).
\]
Apply \eqref{eq:newS7V94-metric-response}.
\end{proof}

\begin{corollary}[Einstein source locality]
\label{cor:newS7V94-Einstein}
The retained fermionic contribution to a local metric-response observable
is exponentially insensitive to distant bounded carrier insertions under
the complete channel hypotheses.  This populates the analytic locality row
of the fermionic RCNC stress packet.  It does not prove the Einstein
constraint or the cutoff convergence of the stress itself.
\end{corollary}

\begin{proof}
Equation \eqref{eq:newS7V94-stress} is the stated exponential locality.
The Einstein constraint also contains the gravitational variation,
bosonic stresses, counterterms, and a continuum limit, none of which is an
algebraic consequence of the covariance bound.
\end{proof}

\subsection{Unitary covariance}
\label{sec:newS7V94-covariance}

\begin{proposition}[Naturality of the Petz channel]
\label{prop:newS7V94-natural}
Let \(U_{AB}=U_A\otimes U_B\), let
\(\sigma'_{AB}=U_{AB}\sigma_{AB}U_{AB}^*\), and let
\(X'_B=U_BX_BU_B^*\).  Then
\begin{equation}
 \mathcal P_{\sigma'}(X'_B)
 =U_{AB}\mathcal P_\sigma(X_B)U_{AB}^*.
 \label{eq:newS7V94-natural}
\end{equation}
Consequently product-implemented gauge and local frame rotations preserve
the Petz update and all unitarily invariant influence bounds.
\end{proposition}

\begin{proof}
The marginal transforms as
\(\sigma'_B=U_B\sigma_BU_B^*\), and functional calculus transports its
square root and inverse square root by the same conjugation.  Substitute
these identities in the V91 Petz formula and collect the unitaries.
\end{proof}

\begin{firewall}[Transformations must preserve the tensor split]
\label{fire:newS7V94-tensor-split}
The naturality theorem assumes a product unitary on the declared cell and
collar factors.  A refinement, gauge fixing, or metric frame change that
mixes old, detail, cell, and collar subspaces requires a transported tensor
factorization and a separate intertwining residual.
\end{firewall}

\begin{definition}[Metric-differentiated carrier channel card, MDCCC]
\label{def:newS7V94-MDCCC}
MDCCC requires:
\begin{enumerate}
\item a fixed faithful cell--collar algebra and uniform density floor;
\item local Gibbs Hamiltonians with finite-range or exponentially local
metric, vierbein, spin-connection, gauge, and Higgs derivatives;
\item the Petz derivative bound \eqref{eq:newS7V94-Petz-derivative};
\item a proved update propagation kernel \(G\) and localized sweep source;
\item a uniform cb contraction resolvent and exact or recovery-controlled
Gibbs preservation;
\item the localized Duhamel stress bound and all direct contact terms;
\item unitary covariance plus any tensor-split intertwining residual; and
\item cutoff summability and compatibility with ICCC, ACMC, LCCC, RFRC,
and RCNC.
\end{enumerate}
\end{definition}

\begin{assessment}[Metric response has reached the local channel level]
\label{ass:newS7V94-status}
A local metric derivative of a faithful Gibbs collar changes its Petz
channel by at most
\(2\beta L_{\rm P}\|D_xh_i\|\) in diamond norm.  An exponentially local
update propagation kernel and the cb resolvent then give an exponentially
local fermionic stress response.  The remaining physical debits are the
uniform density floor, actual update propagation and contraction, global
Gibbs recovery, and cutoff summability on the dynamical Einstein collar.
\end{assessment}

\begin{programme}[V95: audit and cofinal stress-projectivity card]
\label{prog:newS7V94-V95}
V95 will perform the compulsory YM and NS audit, then combine the V85
adjacent carrier defect with the V94 metric derivative to prove a
summable cofinal stress-response theorem.  It will separate convergence of
expectations from convergence of metric derivatives and require an
explicit differentiated adjacent defect, preventing an undifferentiated
projective limit from being reported as an Einstein source limit.
\end{programme}

\begin{firewall}[Claim boundary after seventh-series V94]
\label{fire:newS7V94-boundary}
V94 proves the exact normalized Gibbs density derivative and trace-norm
bound, Petz channel diamond derivative, localized differentiated sweep
source, cb-resolvent metric response, retained-carrier Duhamel stress
locality, and unitary naturality.

V94 does not populate the physical density floor, propagation or
contraction matrices, global recovery, or cutoff stress convergence, or
prove MDCCC, ICCC, ACMC, LCCC, SRSC, PSDC, FTCMC, IRCMC, NCRC, RFRC,
CRARC, LPCC, OFHC, LCPLC, CCPC, GCMCC, PCC, SMALC, RCLC, DSLC, CRDC,
GAOCC, CCTC, MSRC, RTGC, or RCNC, remove the cutoff, or construct the
Standard Model--Einstein continuum.
\end{firewall}

\begin{theorem}[Seventh-series V94 result]
\label{thm:newS7V94-result}
For a faithful finite cell--collar Gibbs block,
\(\|D_x\sigma_i\|_1\leq2\beta\|D_xh_i\|\) and its Petz update changes by
at most \(2\beta L_{\rm P}\|D_xh_i\|\) in diamond norm.  If those local
sources and the update propagation decay exponentially and the cb row is
strictly subcritical, the retained fermionic metric response obeys the
localized stress estimate \eqref{eq:newS7V94-stress}.  Scalar shifts cancel
from the density derivative, while direct metric contact terms remain
separate.
\end{theorem}

\begin{proof}
Combine \Cref{thm:newS7V94-Gibbs-derivative,
thm:newS7V94-Petz-derivative,
lem:newS7V94-sweep,
thm:newS7V94-sweep-source,
cor:newS7V94-resolvent,
thm:newS7V94-stress,
prop:newS7V94-natural}.
\end{proof}

\paragraph{Parent-version record.}
V94 was formed from
\texttt{Renewal\_SM\_Einstein\_Continuum\_Research\_Notes\_V93.tex}.
The V93 parent had \(37\,646\) logical source lines, \(1\,474\,002\) bytes,
and SHA--256
\[
 \texttt{C18D3856067E30FD99FFE23A940476D3814C54A95F4C6A95C885AFF10E7CD2B6}.
\]
The next compulsory companion audit is V95.

% =============================================================================
% END APPENDED SEVENTH-SERIES V94 MATERIAL
% =============================================================================
% =============================================================================
% APPENDED SEVENTH-SERIES V95 MATERIAL
% =============================================================================

\section{Companion audit and cofinal stress-response projectivity}
\label{sec:newS7V95}
\label{sec:v7v95}

This is the compulsory V95 companion audit.  It is followed by the
cross-cutoff theorem required by V94.  Projectivity of carrier expectations
at each fixed metric does not imply convergence of their metric derivatives.
The latter is the fermionic Einstein source.  This section proves a
Banach-chart \(C^1\) landing theorem from one summable base-value row and
one summable adjacent derivative row, then applies it to retained-carrier
Duhamel stress and expectation-level Einstein residuals.

\subsection{Compulsory V95 companion audit}
\label{sec:newS7V95-audit}

At the audit snapshot, the newest distinct complete Yang--Mills checkpoint
remained
\begin{align}
 &\texttt{Renewal\_Yang\_Mills\_Mass\_Gap\_Research\_Notes\_V45.tex},
 \notag\\
 &\qquad 704\,056\ {\rm bytes},\quad 20\,173\ {\rm logical\ lines},\quad
 \texttt{793583EFCD7EBDD7E96D77D5A9E7C6CE78CECDF693BA6196A7C60D5B3AF1E6AC}.
 \label{eq:newS7V95-YM-hash}
\end{align}
The newest Navier--Stokes checkpoint remained
\begin{align}
 &\texttt{Renewal\_NS\_Stability\_Research\_Notes\_V26.tex},
 \notag\\
 &\qquad 278\,283\ {\rm bytes},\quad 5\,319\ {\rm logical\ lines},\quad
 \texttt{82C887F8C2C69E4F28FAD7C3DC8DE5B9099CB5A4BF431EBA1FFF3E91935978AD}.
 \label{eq:newS7V95-NS-hash}
\end{align}
Both have one live document terminator.

\begin{theorem}[V95 audit decision]
\label{thm:newS7V95-audit}
No new companion writer is available.  YM V45's distinction between a
static marginal and its response covariance reinforces the need for a
separate differentiated projectivity row.  NS V26 supplies no carrier
stress transport and licenses no rank shortcut.
\end{theorem}

\begin{proof}
The newest distinct file identities and hashes agree with the V90 audit,
and no later active append is present.  The stated lessons are their
unchanged mathematical scope.
\end{proof}

\subsection{A Banach-chart \(C^1\) landing theorem}
\label{sec:newS7V95-C1}

Let \(Q\) be a convex subset of a real Banach space \(E\), let
\(q_0\in Q\), and suppose
\begin{equation}
 R_Q:=\sup_{q\in Q}\|q-q_0\|_E<\infty.
 \label{eq:newS7V95-radius}
\end{equation}
Let \(F_n:Q\to\mathbb R\) be continuously Frechet differentiable.

\begin{theorem}[Summable adjacent derivatives give a \(C^1\) limit]
\label{thm:newS7V95-C1}
Assume
\begin{align}
 \sum_{n=1}^\infty
 |F_{n+1}(q_0)-F_n(q_0)|&<\infty,
 \label{eq:newS7V95-base-sum}\\
 \sum_{n=1}^\infty
 \sup_{q\in Q}\|DF_{n+1}(q)-DF_n(q)\|_{E^*}&<\infty.
 \label{eq:newS7V95-derivative-sum}
\end{align}
Then \(F_n\) converges uniformly on \(Q\) to a differentiable function
\(F\), the derivatives converge uniformly to \(DF\), and
\begin{align}
 \sup_{q\in Q}|F(q)-F_m(q)|
 &\leq
 \sum_{n\geq m}|F_{n+1}(q_0)-F_n(q_0)|
 \notag\\
 &\quad+R_Q\sum_{n\geq m}
 \sup_{q\in Q}\|DF_{n+1}(q)-DF_n(q)\|,
 \label{eq:newS7V95-function-tail}\\
 \sup_{q\in Q}\|DF(q)-DF_m(q)\|
 &\leq\sum_{n\geq m}
 \sup_{q\in Q}\|DF_{n+1}(q)-DF_n(q)\|.
 \label{eq:newS7V95-derivative-tail}
\end{align}
\end{theorem}

\begin{proof}
The derivative series is uniformly Cauchy by
\eqref{eq:newS7V95-derivative-sum}; let its continuous limit be \(G(q)\).
For \(q\in Q\), convexity and the fundamental theorem of calculus give
\begin{align*}
 F_{n+1}(q)-F_n(q)
 &=F_{n+1}(q_0)-F_n(q_0)\\
 &\quad+\int_0^1
 [DF_{n+1}-DF_n]_{q_0+t(q-q_0)}(q-q_0)\,dt.
\end{align*}
The two summable rows and \eqref{eq:newS7V95-radius} make \(F_n\) uniformly
Cauchy and prove \eqref{eq:newS7V95-function-tail}.  Pass to the limit in
\[
 F_n(q+h)-F_n(q)=\int_0^1DF_n(q+th)h\,dt
\]
on every segment contained in \(Q\).  Uniform convergence of \(DF_n\)
gives
\(F(q+h)-F(q)=\int_0^1G(q+th)h\,dt\), so \(DF=G\).  The derivative tail
is the tail of the uniformly convergent derivative series.
\end{proof}

\begin{countertheorem}[Uniform function convergence does not transport
stress]
\label{cth:newS7V95-no-derivative}
There are smooth functions \(F_n:[0,1]\to\mathbb R\) converging uniformly
to zero whose derivatives do not converge to the derivative of the limit at the
origin.
\end{countertheorem}

\begin{proof}
Take \(F_n(q)=n^{-1}\sin(nq)\).  Then
\(\|F_n\|_\infty\leq n^{-1}\to0\), while
\(F_n'(0)=1\) for every \(n\), whereas the derivative of the zero limit is
zero.
\end{proof}

\begin{firewall}[Undifferentiated projectivity is not an Einstein source]
\label{fire:newS7V95-C1}
Trace-norm or expectation convergence of finite carrier states at each
fixed metric cannot be differentiated without a row such as
\eqref{eq:newS7V95-derivative-sum}.  The counterexample remains valid even
with uniform convergence.
\end{firewall}

\subsection{Application to retained-carrier observables}
\label{sec:newS7V95-carrier}

Let \(\rho_n(q)\) be the finite retained carrier state on cutoff \(n\), and
let \(A_n(q)\) be a coherently embedded bounded local observable.  Put
\begin{equation}
 F_n(q):=\operatorname{Tr}(\rho_n(q)A_n(q)).
 \label{eq:newS7V95-Fn}
\end{equation}
For a tangent \(v\in E\), the V87 identity gives
\begin{equation}
 DF_n(q)[v]
 =\omega_{n,q}(DA_n(q)[v])
 -\beta_n\operatorname{Cov}^{\rm D}_{n,q}
 \bigl(A_n(q),DK_n(q)[v]\bigr).
 \label{eq:newS7V95-stress-formula}
\end{equation}

\begin{definition}[Adjacent stress defect]
\label{def:newS7V95-stress-defect}
After transporting levels \(n+1\) and \(n\) to the same old observable and
metric tangent spaces, define
\begin{equation}
 \delta_n^{(1)}(A;Q):=
 \sup_{q\in Q}\sup_{\|v\|_E\leq1}
 |DF_{n+1}(q)[v]-DF_n(q)[v]|.
 \label{eq:newS7V95-stress-defect}
\end{equation}
Its direct and covariance parts are
\begin{align}
 \delta_{n,{\rm dir}}^{(1)}
 &:=\sup_{q,v}
 \left|\omega_{n+1,q}(DA_{n+1}[v])
 -\omega_{n,q}(DA_n[v])\right|,
 \label{eq:newS7V95-direct-defect}\\
 \delta_{n,{\rm cov}}^{(1)}
 &:=\sup_{q,v}
 \left|\beta_{n+1}\operatorname{Cov}^{\rm D}_{n+1,q}
 (A_{n+1},DK_{n+1}[v])
 -\beta_n\operatorname{Cov}^{\rm D}_{n,q}
 (A_n,DK_n[v])\right|.
 \label{eq:newS7V95-cov-defect}
\end{align}
Then
\begin{equation}
 \delta_n^{(1)}\leq
 \delta_{n,{\rm dir}}^{(1)}+delta_{n,{\rm cov}}^{(1)}.
 \label{eq:newS7V95-defect-split}
\end{equation}
\end{definition}

\begin{theorem}[Cofinal retained-carrier stress landing]
\label{thm:newS7V95-stress-landing}
Suppose for each fixed protected local observable \(A\):
\begin{enumerate}
\item its transported basepoint expectation defects are summable;
\item \(\sum_n\delta_{n,{\rm dir}}^{(1)}<\infty\); and
\item \(\sum_n\delta_{n,{\rm cov}}^{(1)}<\infty\),
\end{enumerate}
uniformly on the bounded convex metric chart \(Q\).  Then the cofinal
expectation \(F(q)\) exists in \(C^1(Q)\), and its metric derivative is the
uniform limit of the finite direct contact and retained-carrier Duhamel
stress terms in \eqref{eq:newS7V95-stress-formula}.
\end{theorem}

\begin{proof}
The basepoint row is \eqref{eq:newS7V95-base-sum}.  The two derivative rows
and \eqref{eq:newS7V95-defect-split} imply
\eqref{eq:newS7V95-derivative-sum}.  Apply
\Cref{thm:newS7V95-C1} to \eqref{eq:newS7V95-Fn}; uniform convergence in
\eqref{eq:newS7V95-stress-formula} identifies the limiting derivative.
\end{proof}

\begin{corollary}[Local channel loaders are not yet adjacent loaders]
\label{cor:newS7V95-local-versus-adjacent}
V94 bounds each finite stress response locally in space.  To populate
\(\delta_{n,{\rm cov}}^{(1)}\), one must additionally compare the
Duhamel covariances, metric insertions, temperatures, embeddings, and
channel sources at adjacent cutoffs.  Uniform boundedness of each term does
not imply summability of their differences.
\end{corollary}

\begin{proof}
A bounded sequence need not be Cauchy.  The defect
\eqref{eq:newS7V95-cov-defect} is a difference of the complete adjacent
terms, whereas V94 bounds their individual magnitudes.
\end{proof}

\subsection{Expectation-level Einstein residual passage}
\label{sec:newS7V95-Einstein}

Let \(v\) be a compactly supported metric tangent.  On cutoff \(n\), write
the regulated expectation-level residual as
\begin{equation}
 \mathcal R_n(q)[v]
 :=2D\mathcal S_{g,n}(q)[v]
 -\mathcal T_{{\rm cl},n}(q)[v]
 -\mathcal T_{{\rm car},n}(q)[v]
 -\mathcal T_{{\rm ct},n}(q)[v],
 \label{eq:newS7V95-residual}
\end{equation}
where the carrier term is the finite derivative supplied by
\eqref{eq:newS7V95-stress-formula}, with every direct contact term included.

\begin{theorem}[Weak Einstein residual landing]
\label{thm:newS7V95-Einstein}
Assume on a common tangent core:
\begin{enumerate}
\item \(D\mathcal S_{g,n}\), \(\mathcal T_{{\rm cl},n}\),
\(\mathcal T_{{\rm car},n}\), and \(\mathcal T_{{\rm ct},n}\) converge
uniformly on bounded metric charts and unit tangent sets;
\item the carrier convergence is supplied by
\Cref{thm:newS7V95-stress-landing}; and
\item
\begin{equation}
 \sup_{q\in Q}\sup_{\|v\|\leq1}|\mathcal R_n(q)[v]|
 \longrightarrow0.
 \label{eq:newS7V95-residual-zero}
\end{equation}
\end{enumerate}
Then the limiting functionals satisfy
\begin{equation}
 2D\mathcal S_g
 =\mathcal T_{\rm cl}+\mathcal T_{\rm car}+
 \mathcal T_{\rm ct}
 \label{eq:newS7V95-Einstein-limit}
\end{equation}
on that tangent core.
\end{theorem}

\begin{proof}
Uniform convergence permits termwise passage to the limit in
\eqref{eq:newS7V95-residual}.  The residual tends to zero by
\eqref{eq:newS7V95-residual-zero}, so the limiting difference is zero on
every declared tangent.
\end{proof}

\begin{firewall}[Expectation identity is not the full continuum]
\label{fire:newS7V95-expectation}
Equation \eqref{eq:newS7V95-Einstein-limit} is a weak expectation-level
identity on a declared tangent and observable core.  It does not prove a
pointwise Einstein equation, construct a continuum metric measure, establish
Lorentzian dynamics, or control all composite-operator counterterms.
\end{firewall}

\subsection{Cofinal stress-projectivity card}
\label{sec:newS7V95-card}

\begin{definition}[Cofinal stress-projectivity card, CSPC]
\label{def:newS7V95-CSPC}
For each protected observable and bounded convex metric chart, CSPC
requires:
\begin{enumerate}
\item coherent observable, carrier, metric-tangent, and counterterm
embeddings across cutoffs;
\item summable basepoint expectation defects;
\item the direct and Duhamel covariance defects
\eqref{eq:newS7V95-direct-defect}--\eqref{eq:newS7V95-cov-defect};
\item summability of both derivative rows, uniformly on the chart;
\item V94 spatial locality and uniform-integrability bounds for every
stress insertion;
\item convergence of gravitational, classical, carrier, and counterterm
functionals on one tangent core;
\item vanishing regulated residual \eqref{eq:newS7V95-residual-zero}; and
\item compatibility with MDCCC, ICCC, RFRC, CRARC, RCNC, and the physical
clock.
\end{enumerate}
\end{definition}

\begin{assessment}[The Einstein handoff now has the correct topology]
\label{ass:newS7V95-status}
The finite carrier programme has reached an honest cofinal stress criterion:
base expectation projectivity and adjacent derivative projectivity are
separate summable rows.  Together they give a \(C^1\) metric-chart limit and
permit an expectation-level Einstein residual to pass.  The first missing
calculation is the adjacent Duhamel covariance defect
\eqref{eq:newS7V95-cov-defect} for the actual old/detail refinement.
\end{assessment}

\begin{programme}[V96: differentiated adjacent Gibbs defect]
\label{prog:newS7V95-V96}
V96 will differentiate the V85 comparison between the interacting fine
Gibbs carrier and the factorized old/detail carrier.  It will bound the
adjacent derivative by the Hamiltonian mismatch, its metric derivative,
and Duhamel covariance continuity, producing an explicit candidate for
\(\delta_{n,{\rm cov}}^{(1)}\).  This avoids treating the V94 within-cutoff
stress bound as a cross-cutoff theorem.
\end{programme}

\begin{firewall}[Claim boundary after seventh-series V95]
\label{fire:newS7V95-boundary}
V95 performs the compulsory companion audit; proves the Banach-chart
\(C^1\) landing theorem and counterexample; defines and splits the adjacent
carrier stress defect; proves conditional cofinal stress landing; and
passes a uniformly vanishing expectation-level Einstein residual to the
limit.

V95 does not populate the adjacent stress defects or continuum residual,
or prove CSPC, MDCCC, ICCC, ACMC, LCCC, SRSC, PSDC, FTCMC, IRCMC, NCRC,
RFRC, CRARC, LPCC, OFHC, LCPLC, CCPC, GCMCC, PCC, SMALC, RCLC, DSLC,
CRDC, GAOCC, CCTC, MSRC, RTGC, or RCNC, remove the cutoff, or construct
the Standard Model--Einstein continuum.
\end{firewall}

\begin{theorem}[Seventh-series V95 result]
\label{thm:newS7V95-result}
A summable basepoint expectation defect together with a summable uniform
adjacent metric-derivative defect yields a \(C^1\) cofinal carrier
expectation on every bounded convex metric chart.  For Gibbs carriers the
derivative splits exactly into a direct contact term and a Duhamel stress
covariance.  If the gravitational, classical, carrier, and counterterm
functionals converge and the regulated expectation residual vanishes
uniformly, their limits satisfy the weak identity
\eqref{eq:newS7V95-Einstein-limit} on the declared tangent core.
\end{theorem}

\begin{proof}
Combine \Cref{thm:newS7V95-audit,
thm:newS7V95-C1,
cth:newS7V95-no-derivative,
thm:newS7V95-stress-landing,
thm:newS7V95-Einstein}.
\end{proof}

\paragraph{Parent-version record.}
V95 was formed from
\texttt{Renewal\_SM\_Einstein\_Continuum\_Research\_Notes\_V94.tex}.
The V94 parent had \(38\,068\) logical source lines, \(1\,489\,176\) bytes,
and SHA--256
\[
 \texttt{914F091028EBD3DFAEB562F36F4C28DB5EBC3217BE4A2D893E45EB10CE9B7502}.
\]
The next compulsory companion audit is V100.

% =============================================================================
% END APPENDED SEVENTH-SERIES V95 MATERIAL
% =============================================================================
% =============================================================================
% BEGIN APPENDED SEVENTH-SERIES V96 MATERIAL
% =============================================================================

\section{Differentiated adjacent Gibbs comparison}
\label{sec:v7v96}

V95 isolated the adjacent Duhamel-covariance row which is needed for a
cofinal stress tensor.  This appendix computes that row for two finite
Gibbs carriers on one transported local algebra.  The estimate is deliberately
finite-volume: its density-floor constant must subsequently be replaced by a
local channel estimate, or shown to be summable after it is combined with the
old/detail mismatch.

Let \(\mathfrak A=M_d(\mathbb C)\), let \(\mathfrak A_{\rm sa}\) denote its
self-adjoint part, and put
\begin{equation}
 \rho_H:=\frac{e^{-\beta H}}{\Tr(e^{-\beta H})},\qquad
 \omega_H(A):=\Tr(\rho_HA),\qquad \beta>0.
 \label{eq:newS7V96-Gibbs}
\end{equation}
For a faithful density matrix \(\rho\), define the centred Duhamel form
\begin{equation}
 \mathcal B_\rho(A,X)
 :=\int_0^1\Tr\!\left(
 \rho^s(A-\Tr(\rho A)\mathbf 1)
 \rho^{1-s}(X-\Tr(\rho X)\mathbf 1)\right)\,ds .
 \label{eq:newS7V96-B}
\end{equation}
By expanding the centring terms and using cyclicity,
\begin{equation}
 \mathcal B_\rho(A,X)
 =\int_0^1\Tr(\rho^sA\rho^{1-s}X)\,ds
   -\Tr(\rho A)\Tr(\rho X).
 \label{eq:newS7V96-B-uncentred}
\end{equation}

\begin{lemma}[Uniform elementary bounds]
\label{lem:newS7V96-elementary}
For every density matrix \(\rho\) and \(A,X\in\mathfrak A\),
\begin{equation}
 |\mathcal B_\rho(A,X)|\leq2\|A\|\,\|X\|.
 \label{eq:newS7V96-B-bound}
\end{equation}
Consequently, for fixed \(\rho\),
\begin{align}
 |\mathcal B_\rho(A,X)-\mathcal B_\rho(A',X)|
 &\leq2\|A-A'\|\,\|X\|,\label{eq:newS7V96-A-Lip}\\
 |\mathcal B_\rho(A,X)-\mathcal B_\rho(A,X')|
 &\leq2\|A\|\,\|X-X'\|.\label{eq:newS7V96-X-Lip}
\end{align}
\end{lemma}

\begin{proof}
For \(0<s<1\), Schatten H\"older gives
\[
 |\Tr(\rho^sA\rho^{1-s}X)|
 \leq \|\rho^s\|_{1/s}\|A\|
       \|\rho^{1-s}\|_{1/(1-s)}\|X\|
 =\|A\|\|X\|.
\]
The same bound holds at the endpoints by continuity.  The product of the
two expectations in \eqref{eq:newS7V96-B-uncentred} is also bounded by
\(\|A\|\|X\|\), proving \eqref{eq:newS7V96-B-bound}.  Bilinearity of
\(\mathcal B_\rho\) and the same estimate give
\eqref{eq:newS7V96-A-Lip}--\eqref{eq:newS7V96-X-Lip}.
\end{proof}

\begin{definition}[Density-floor Duhamel constant]
\label{def:newS7V96-LD}
For \(0<\lambda\leq d^{-1}\), let
\begin{equation}
 \mathcal K_{d,\lambda}
 :=\{\rho\in\mathfrak A_{\rm sa}:\rho\geq\lambda\mathbf1,
                      \ \Tr\rho=1\}.
 \label{eq:newS7V96-K}
\end{equation}
On the affine trace-one space set
\begin{equation}
 L_D(d,\lambda):=
 \sup_{\substack{\rho\in\mathcal K_{d,\lambda},\ \
                  Z=Z^*,\ \Tr Z=0,\ \|Z\|_1\leq1\\
                  \|A\|\leq1,\ \|X\|\leq1}}
 \big|D_\rho\mathcal B_\rho(A,X)[Z]\big|.
 \label{eq:newS7V96-LD}
\end{equation}
\end{definition}

\begin{lemma}[Continuity in the state]
\label{lem:newS7V96-state-Lip}
The number \(L_D(d,\lambda)\) is finite.  If
\(\rho,\sigma\in\mathcal K_{d,\lambda}\), then
\begin{equation}
 |\mathcal B_\rho(A,X)-\mathcal B_\sigma(A,X)|
 \leq L_D(d,\lambda)\|A\|\|X\|\,\|\rho-\sigma\|_1.
 \label{eq:newS7V96-state-Lip}
\end{equation}
\end{lemma}

\begin{proof}
On the open positive cone, \((\rho,s)\mapsto\rho^s\) is continuously
Fr\'echet differentiable for \(s\in[0,1]\); this follows, for example,
from the holomorphic functional calculus on a contour separated from the
nonpositive real axis.  The contour and resolvent bounds can be chosen
uniformly when \(\rho\geq\lambda\mathbf1\).  Hence the integrand in
\eqref{eq:newS7V96-B-uncentred}, and its derivative in \(\rho\), are
continuous on the compact product of \(\mathcal K_{d,\lambda}\), the two
operator-norm unit balls, the trace-norm unit ball in the trace-zero tangent
space, and \([0,1]\).  Differentiation under the \(s\)-integral is therefore
valid and the supremum \eqref{eq:newS7V96-LD} is finite.

The segment \(\sigma+t(\rho-\sigma)\) remains in
\(\mathcal K_{d,\lambda}\).  The fundamental theorem of calculus along
that segment, followed by homogeneity in \(A\) and \(X\), gives
\eqref{eq:newS7V96-state-Lip}.
\end{proof}

\begin{lemma}[Trace-norm motion of a Gibbs density]
\label{lem:newS7V96-Gibbs-motion}
If \(H_0,H_1\in\mathfrak A_{\rm sa}\), then
\begin{equation}
 \|\rho_{H_1}-\rho_{H_0}\|_1
 \leq2\beta\|H_1-H_0\|.
 \label{eq:newS7V96-trace-motion}
\end{equation}
If \(\|H_j\|\leq M\), then
\begin{equation}
 \rho_{H_j}\geq\lambda_{d,M}\mathbf1,
 \qquad \lambda_{d,M}:=d^{-1}e^{-2\beta M}.
 \label{eq:newS7V96-floor}
\end{equation}
\end{lemma}

\begin{proof}
Put \(H_t=H_0+t(H_1-H_0)\) and \(V=H_1-H_0\).  Duhamel's formula,
including the derivative of the partition function, gives
\begin{equation}
 \dot\rho_t=-\beta\int_0^1
 \rho_t^s\big(V-\omega_{H_t}(V)\mathbf1\big)
 \rho_t^{1-s}\,ds .
 \label{eq:newS7V96-rho-dot}
\end{equation}
Noncommutative H\"older yields
\[
 \|\rho_t^sY\rho_t^{1-s}\|_1\leq\|Y\|
 \quad(0\leq s\leq1),
\]
and \(\|V-\omega_{H_t}(V)\mathbf1\|\leq2\|V\|\).  Integrating
\eqref{eq:newS7V96-rho-dot} first in \(s\) and then in \(t\) proves
\eqref{eq:newS7V96-trace-motion}.  If \(\|H_j\|\leq M\), the smallest
eigenvalue of \(e^{-\beta H_j}\) is at least \(e^{-\beta M}\), while its
trace is at most \(d e^{\beta M}\).  Their quotient is
\eqref{eq:newS7V96-floor}.
\end{proof}

\subsection{The adjacent metric-derivative estimate}
\label{sec:newS7V96-adjacent}

Let \(Q\) be a convex open subset of a real Banach space.  For
\(j\in\{0,1\}\), let
\[
 H_j:Q\to\mathfrak A_{\rm sa},\qquad
 O_j:Q\to\mathfrak A
\]
be \(C^1\), and define \(F_j(q)=\omega_{H_j(q)}(O_j(q))\).  For a tangent
\(v\), abbreviate
\begin{equation}
 X_j(q;v):=DH_j(q)[v],\qquad
 Y_j(q;v):=DO_j(q)[v].
 \label{eq:newS7V96-XY}
\end{equation}
The exact finite Gibbs derivative is
\begin{equation}
 DF_j(q)[v]
 =\omega_{H_j(q)}(Y_j(q;v))
  -\beta\mathcal B_{\rho_{H_j(q)}}(O_j(q),X_j(q;v)).
 \label{eq:newS7V96-derivative}
\end{equation}

\begin{theorem}[Differentiated adjacent Gibbs bound]
\label{thm:newS7V96-adjacent}
Fix \(q\in Q\) and \(v\).  Suppose
\begin{align}
 \|H_j(q)\|&\leq M,&
 \|H_1(q)-H_0(q)\|&\leq\varepsilon_0,\label{eq:newS7V96-ass-H}\\
 \|O_j(q)\|&\leq M_O,&
 \|O_1(q)-O_0(q)\|&\leq a_0,\label{eq:newS7V96-ass-O}\\
 \|X_j(q;v)\|&\leq M_X,&
 \|X_1(q;v)-X_0(q;v)\|&\leq\varepsilon_1,\label{eq:newS7V96-ass-X}\\
 \|Y_0(q;v)\|&\leq M_Y,&
 \|Y_1(q;v)-Y_0(q;v)\|&\leq a_1.
 \label{eq:newS7V96-ass-Y}
\end{align}
Then, with \(\lambda_{d,M}\) from \eqref{eq:newS7V96-floor},
\begin{align}
 |DF_1(q)[v]-DF_0(q)[v]|
 \leq{}&a_1+2\beta M_Y\varepsilon_0
       +2\beta a_0M_X+2\beta M_O\varepsilon_1
 \notag\\
 &\quad+2\beta^2 L_D(d,\lambda_{d,M})
             M_OM_X\varepsilon_0.
 \label{eq:newS7V96-main-bound}
\end{align}
\end{theorem}

\begin{proof}
Write \(\rho_j=\rho_{H_j(q)}\).  The direct term obeys
\begin{align*}
 |\Tr(\rho_1Y_1)-\Tr(\rho_0Y_0)|
 &\leq\|Y_1-Y_0\|+\|\rho_1-\rho_0\|_1\|Y_0\|\\
 &\leq a_1+2\beta M_Y\varepsilon_0
\end{align*}
by \Cref{lem:newS7V96-Gibbs-motion}.  Insert and subtract successively
\(\mathcal B_{\rho_1}(O_0,X_1)\) and
\(\mathcal B_{\rho_1}(O_0,X_0)\).  Then
\Cref{lem:newS7V96-elementary,lem:newS7V96-state-Lip} give
\begin{align*}
 &|\mathcal B_{\rho_1}(O_1,X_1)
       -\mathcal B_{\rho_0}(O_0,X_0)|\\
 &\qquad\leq2a_0M_X+2M_O\varepsilon_1
 +L_D(d,\lambda_{d,M})M_OM_X\|\rho_1-\rho_0\|_1\\
 &\qquad\leq2a_0M_X+2M_O\varepsilon_1
 +2\beta L_D(d,\lambda_{d,M})M_OM_X\varepsilon_0.
\end{align*}
Multiply this covariance estimate by \(\beta\), combine it with the
direct estimate, and use \eqref{eq:newS7V96-derivative}.
\end{proof}

\begin{corollary}[Scalar-energy gauge]
\label{cor:newS7V96-scalar}
Suppose a \(C^1\) scalar field \(c:Q\to\mathbb R\) and a remainder
\(W:Q\to\mathfrak A_{\rm sa}\) satisfy
\begin{equation}
 H_1-H_0=c\mathbf1+W.
 \label{eq:newS7V96-scalar-split}
\end{equation}
Then \Cref{thm:newS7V96-adjacent} remains valid with
\(\varepsilon_0=\|W(q)\|\) and
\(\varepsilon_1=\|DW(q)[v]\|\), after replacing \(H_1\) by
\(H_1-c\mathbf1\) in the harmless common norm bound \(M\).
\end{corollary}

\begin{proof}
Adding a scalar to a Hamiltonian leaves its normalized Gibbs density
unchanged.  Moreover
\(\mathcal B_\rho(A,z\mathbf1)=0\).  Thus both \(c\mathbf1\) and
\(Dc[v]\mathbf1\) disappear from \eqref{eq:newS7V96-derivative}; apply the
theorem to \(H_1-c\mathbf1=H_0+W\).
\end{proof}

\subsection{Candidate adjacent row for the old/detail refinement}
\label{sec:newS7V96-candidate}

Transport the V85 fine carrier and its factorized old/detail comparison to
the same finite collar algebra.  At scale \(n\), let
\(\varepsilon_{0,n}\), \(\varepsilon_{1,n}\), \(a_{0,n}\), and
\(a_{1,n}\) be the four mismatch bounds in
\eqref{eq:newS7V96-ass-H}--\eqref{eq:newS7V96-ass-Y}, uniformly on the
chosen metric chart and its unit tangent set.  Define
\begin{align}
 \widehat\delta^{(1)}_{n,\mathrm{Gibbs}}
 :={}&a_{1,n}+2\beta M_Y\varepsilon_{0,n}
 +2\beta M_Xa_{0,n}+2\beta M_O\varepsilon_{1,n}
 \notag\\
 &+2\beta^2L_D(d_n,\lambda_{d_n,M_n})
      M_OM_X\varepsilon_{0,n}.
 \label{eq:newS7V96-delta}
\end{align}

\begin{corollary}[Conditional population of the V95 covariance row]
\label{cor:newS7V96-populate}
If the transported V85 pair satisfies the hypotheses of
\Cref{thm:newS7V96-adjacent} at every scale and
\begin{equation}
 \sum_{n\geq n_0}\widehat\delta^{(1)}_{n,\mathrm{Gibbs}}<\infty,
 \label{eq:newS7V96-sum}
\end{equation}
then the corresponding adjacent carrier derivatives are uniformly
summable.  Together with the V95 basepoint row, they produce the \(C^1\)
cofinal carrier expectation supplied by V95.
\end{corollary}

\begin{proof}
Theorem \ref{thm:newS7V96-adjacent} bounds the norm of every adjacent
derivative functional by \eqref{eq:newS7V96-delta}.  Sum over \(n\) and
apply the Weierstrass test, followed by the V95 Banach-chart \(C^1\)
landing theorem.
\end{proof}

\begin{firewall}[The density-floor obstruction]
\label{fire:newS7V96-floor}
The constant \(L_D(d_n,\lambda_{d_n,M_n})\) may grow with the collar
dimension and with the inverse spectral floor.  Therefore small
\(\varepsilon_{0,n}\) alone does not establish
\eqref{eq:newS7V96-sum}.  The estimate is a correct finite-scale loader,
not a uniform thermodynamic theorem.  V97 must go upstream to the local
recovery-channel recursion, where mixed metric/refinement sources can be
propagated in local seminorms without a global density floor.
\end{firewall}

\begin{assessment}[Progress at V96]
\label{ass:newS7V96-status}
The missing V95 comparison has been differentiated.  Four independently
checkable defects now control the adjacent stress response, and scalar
vacuum-energy shifts have been quotiented out exactly.  The remaining
bottleneck is visible in one coefficient: global finite-volume smoothness
introduces \(L_D(d_n,\lambda_{d_n,M_n})\), which is not known to be
summable along the cofinal sequence.
\end{assessment}

\begin{programme}[V97: mixed-source local resolvent]
\label{prog:newS7V96-V97}
V97 will differentiate the V93 inhomogeneous channel fixed-point equation
simultaneously in a metric direction and in the old/detail interpolation.
The resulting mixed source will be propagated by the same spatial
resolvent.  This is the upstream replacement for the global density-floor
constant in \eqref{eq:newS7V96-delta}.
\end{programme}

\begin{firewall}[Claim boundary after seventh-series V96]
\label{fire:newS7V96-boundary}
V96 proves finite-dimensional state continuity of the Duhamel form, the
trace-norm Gibbs perturbation bound, the differentiated adjacent Gibbs
estimate, and its scalar-energy reduction.  It supplies the explicit
candidate \eqref{eq:newS7V96-delta} for the V95 adjacent derivative row.

V96 does not prove that this candidate is summable, control its density-floor
constant uniformly, populate the actual metric/refinement mismatch rows,
remove the cutoff, or construct the Standard Model--Einstein continuum.
\end{firewall}

\begin{theorem}[Seventh-series V96 result]
\label{thm:newS7V96-result}
On a finite transported collar algebra, the difference of two Gibbs
expectation metric derivatives is bounded by
\eqref{eq:newS7V96-main-bound}.  Thus Hamiltonian, observable, metric
insertion, and direct-contact mismatches form a sufficient adjacent stress
loader.  Scalar Hamiltonian shifts make no contribution.  Uniform cofinal
landing follows if and only if the programme separately proves a summable
majorant such as \eqref{eq:newS7V96-sum}; no such majorant is inferred from
finite dimensionality.
\end{theorem}

\begin{proof}
Combine \Cref{thm:newS7V96-adjacent,cor:newS7V96-scalar,
cor:newS7V96-populate,fire:newS7V96-floor}.
\end{proof}

\paragraph{Parent-version record.}
V96 was formed from
\texttt{Renewal\_SM\_Einstein\_Continuum\_Research\_Notes\_V95.tex}.
The V95 parent had \(38\,440\) logical source lines, \(1\,503\,216\) bytes,
and SHA--256
\[
 \texttt{34C8806F873E3D961488FDD856DC354435F1449194FD48E3DD6D7203CB744518}.
\]
The next compulsory companion audit is V100.

% =============================================================================
% END APPENDED SEVENTH-SERIES V96 MATERIAL
% =============================================================================

% =============================================================================
% BEGIN APPENDED SEVENTH-SERIES V97 MATERIAL
% =============================================================================

\section{Mixed metric/refinement sources in a local channel resolvent}
\label{sec:v7v97}

V96 gave a finite-collar derivative comparison, but its state-continuity
constant can deteriorate with the global spectral floor.  We now replace
that step by differentiating the local fixed-point equation in two
directions.  The metric parameter produces stress response; the refinement
parameter interpolates between the factorized old/detail carrier and the
fine carrier.  Their mixed derivative is exactly the adjacent stress
quantity required by V95.

\begin{correction}[Logical strength of the last sentence of V96]
\label{cor:newS7V97-v96-correction}
In \Cref{thm:newS7V96-result}, the words ``if and only if'' are to be read
as ``whenever''.  Summability of the displayed V96 majorant is sufficient
for cofinal landing; it was not proved necessary, because cancellations
could make the true adjacent differences summable even when that particular
majorant is not.
\end{correction}

\subsection{Two-parameter differentiated fixed points}
\label{sec:newS7V97-fixed}

Let \(I,J\subset\mathbb R\) be intervals, let
\((s,q)\mapsto\mathcal T_{s,q}\) be a \(C^2\) family of CPTP maps on a
finite carrier algebra, and let \((s,q)\mapsto\rho_{s,q}\) be a \(C^2\)
family of density matrices satisfying
\begin{equation}
 \mathcal T_{s,q}(\rho_{s,q})=\rho_{s,q}.
 \label{eq:newS7V97-fixed}
\end{equation}
The parameter \(s\) is the refinement interpolation and \(q\) is one
metric-chart coordinate.  Directional derivatives on a Banach metric chart
are obtained by restricting to a line in the chosen tangent direction.
Write
\[
 \rho_s:=\partial_s\rho,
 \quad \rho_q:=\partial_q\rho,
 \quad \rho_{sq}:=\partial_s\partial_q\rho,
\]
and use the same subscripts for channel derivatives.  All three density
derivatives are traceless Hermitian.

Retain the cell seminorms \(v_i\) and dual estimate
\eqref{eq:newS7V93-dual}.  Suppose, uniformly in \((s,q)\),
\begin{align}
 v_i(\mathcal T X)&\leq\sum_jC_{ij}v_j(X),
 \label{eq:newS7V97-C}\\
 v_i(\mathcal T_s X)&\leq\sum_jA^{s}_{ij}v_j(X),
 \qquad
 v_i(\mathcal T_q X)\leq\sum_jA^{q}_{ij}v_j(X)
 \label{eq:newS7V97-A}
\end{align}
for traceless Hermitian \(X\), where all matrices on the right are
nonnegative and \(\sup_i\sum_jC_{ij}<1\).  Define the point sources
\begin{equation}
 b_i^s:=v_i(\mathcal T_s\rho),\qquad
 b_i^q:=v_i(\mathcal T_q\rho),\qquad
 b_i^{sq}:=v_i(\mathcal T_{sq}\rho).
 \label{eq:newS7V97-b}
\end{equation}

\begin{theorem}[Mixed-source resolvent]
\label{thm:newS7V97-mixed}
Let \(D=(I-C)^{-1}\) as in V93 and set
\(u_i^a=v_i(\rho_a)\) for
\(a\in\{s,q,sq\}\).  Then, componentwise,
\begin{align}
 u^s&\leq Db^s, & u^q&\leq Db^q,
 \label{eq:newS7V97-first}\\
 u^{sq}&\leq D\bigl(b^{sq}+A^sDb^q+A^qDb^s\bigr).
 \label{eq:newS7V97-mixed-resolvent}
\end{align}
\end{theorem}

\begin{proof}
Differentiate \eqref{eq:newS7V97-fixed} once:
\begin{equation}
 \rho_s=\mathcal T\rho_s+\mathcal T_s\rho,
 \qquad
 \rho_q=\mathcal T\rho_q+\mathcal T_q\rho.
 \label{eq:newS7V97-first-equations}
\end{equation}
The seminorm bounds and the V93 inhomogeneous resolvent give
\eqref{eq:newS7V97-first}.  Differentiating the second identity in
\eqref{eq:newS7V97-first-equations} with respect to \(s\) gives the exact
mixed equation
\begin{equation}
 \rho_{sq}=\mathcal T\rho_{sq}
 +\mathcal T_s\rho_q+\mathcal T_q\rho_s
 +\mathcal T_{sq}\rho.
 \label{eq:newS7V97-mixed-equation}
\end{equation}
Thus
\[
 u^{sq}\leq Cu^{sq}+A^su^q+A^qu^s+b^{sq}.
\]
Apply the inhomogeneous resolvent and then
\eqref{eq:newS7V97-first}; nonnegativity of all matrices preserves the
inequalities and proves \eqref{eq:newS7V97-mixed-resolvent}.
\end{proof}

\begin{remark}[The cross terms cannot be discarded]
\label{rem:newS7V97-cross}
The terms \(A^sDb^q\) and \(A^qDb^s\) record, respectively, refinement of
the metric response and metric variation of the refinement response.
A bound only on \(\mathcal T_{sq}\rho\) misses both.  This is the precise
reason that separately bounded refinement and stress rows do not by
themselves control the differentiated adjacent defect.
\end{remark}

\subsection{Volume-independent weighted localization}
\label{sec:newS7V97-weighted}

For a nonempty cell set \(Z\) and \(\mu>0\), put
\begin{equation}
 \|z\|_{\mu,Z}:=\sup_i e^{\mu d(i,Z)}|z_i|,
 \qquad
 \|K\|_\mu:=\sup_i\sum_j e^{\mu d(i,j)}|K_{ij}|.
 \label{eq:newS7V97-weighted-norms}
\end{equation}
The triangle inequality implies
\begin{equation}
 \|Kz\|_{\mu,Z}\leq\|K\|_\mu\|z\|_{\mu,Z}.
 \label{eq:newS7V97-weighted-product}
\end{equation}

\begin{theorem}[Localized mixed response]
\label{thm:newS7V97-local}
Assume for some \(\mu>0\)
\begin{equation}
 c_\mu:=\|C\|_\mu<1,qquad
 a_s:=\|A^s\|_\mu<\infty,qquad
 a_q:=\|A^q\|_\mu<\infty,
 \label{eq:newS7V97-weighted-hyp}
\end{equation}
and let \(R_\mu=(1-c_\mu)^{-1}\).  If the three sources in
\eqref{eq:newS7V97-b} have finite \(\|\cdot\|_{\mu,Z}\)-norm, then
\begin{align}
 \|u^s\|_{\mu,Z}&\leq R_\mu\|b^s\|_{\mu,Z},
 &
 \|u^q\|_{\mu,Z}&\leq R_\mu\|b^q\|_{\mu,Z},
 \label{eq:newS7V97-first-local}\\
 \|u^{sq}\|_{\mu,Z}
 &\leq R_\mu\left(
 \|b^{sq}\|_{\mu,Z}
 +a_sR_\mu\|b^q\|_{\mu,Z}
 +a_qR_\mu\|b^s\|_{\mu,Z}
 \right).
 \label{eq:newS7V97-mixed-local}
\end{align}
All constants are independent of the total number of cells.
\end{theorem}

\begin{proof}
Equation \eqref{eq:newS7V97-weighted-product} gives
\(\|C^n\|_\mu\leq c_\mu^n\), hence
\begin{equation}
 \|D\|_\mu\leq\sum_{n\geq0}c_\mu^n=R_\mu.
 \label{eq:newS7V97-D-weighted}
\end{equation}
Apply this estimate to \eqref{eq:newS7V97-first} and
\eqref{eq:newS7V97-mixed-resolvent}, then use
\eqref{eq:newS7V97-weighted-product} for \(A^s\) and \(A^q\).
\end{proof}

\begin{corollary}[Observable mixed response]
\label{cor:newS7V97-observable}
Let \(A_{s,q}\) be a \(C^2\) observable family supported in a fixed finite
cell set \(B\).  Then
\begin{align}
 \left|\partial_s\partial_q\Tr(\rho_{s,q}A_{s,q})\right|
 \leq{}&\|A\|\sum_{i\in B}u_i^{sq}
 +\|A_s\|\sum_{i\in B}u_i^q
 +\|A_q\|\sum_{i\in B}u_i^s
 +\|A_{sq}\|.
 \label{eq:newS7V97-observable-mixed}
\end{align}
In particular, if the hypotheses of \Cref{thm:newS7V97-local} hold, each
of the first three sums is bounded by \(|B|e^{-\mu d(B,Z)}\) times the
corresponding right side of
\eqref{eq:newS7V97-first-local} or
\eqref{eq:newS7V97-mixed-local}.
\end{corollary}

\begin{proof}
Twice differentiating the expectation gives
\begin{align*}
 \partial_s\partial_q\Tr(\rho A)
 ={}&\Tr(\rho_{sq}A)+\Tr(\rho_qA_s)
    +\Tr(\rho_sA_q)+\Tr(\rho A_{sq}).
\end{align*}
Apply the V93 dual seminorm bound to the first three terms and
\(|\Tr(\rho A_{sq})|\leq\|A_{sq}\|\) to the last.  For (i\in B),
(d(i,Z)\geq d(B,Z)); insert the weighted estimates.
\end{proof}

\subsection{Recovery defects and the adjacent stress row}
\label{sec:newS7V97-recovery}

Exact preservation in \eqref{eq:newS7V97-fixed} is a strong hypothesis.
For a declared target state define the signed recovery defect
\begin{equation}
 e_{s,q}:=\mathcal T_{s,q}(\rho_{s,q})-\rho_{s,q}.
 \label{eq:newS7V97-e}
\end{equation}

\begin{proposition}[Differentiated recovery-defect loader]
\label{prop:newS7V97-recovery}
If \(e\) is \(C^2\), the conclusions of
\Cref{thm:newS7V97-mixed} remain true after replacing
\begin{align}
 b_i^s&\ 
 \text{by }\widetilde b_i^s:=v_i(\mathcal T_s\rho-e_s),
 &
 b_i^q&\ 
 \text{by }\widetilde b_i^q:=v_i(\mathcal T_q\rho-e_q),
 \label{eq:newS7V97-recovery-first}\\
 b_i^{sq}&\ 
 \text{by }\widetilde b_i^{sq}:=v_i(\mathcal T_{sq}\rho-e_{sq}).
 \label{eq:newS7V97-recovery-mixed}
\end{align}
The weighted theorem also remains true if these modified sources obey its
hypotheses.
\end{proposition}

\begin{proof}
Equation \eqref{eq:newS7V97-e} is equivalent to
\(\rho=\mathcal T\rho-e\).  Its first and mixed derivatives are
\begin{align*}
 \rho_s&=\mathcal T\rho_s+\mathcal T_s\rho-e_s,\\
 \rho_q&=\mathcal T\rho_q+\mathcal T_q\rho-e_q,\\
 \rho_{sq}&=\mathcal T\rho_{sq}+\mathcal T_s\rho_q
 +\mathcal T_q\rho_s+\mathcal T_{sq}\rho-e_{sq}.
\end{align*}
Repeat the proof of \Cref{thm:newS7V97-mixed}; seminorms remove the signs.
The proof of \Cref{thm:newS7V97-local} is unchanged.
\end{proof}

For scale \(n\), let \(s\in[0,1]\) interpolate from the transported
factorized old/detail carrier at \(s=0\) to the fine carrier at \(s=1\).
For a protected local observable, write
\begin{equation}
 F_n(s,q):=\Tr(\rho_{n,s,q}A_{n,s,q}).
 \label{eq:newS7V97-F}
\end{equation}
The adjacent derivative defect is the exact integral
\begin{equation}
 \partial_qF_n(1,q)-\partial_qF_n(0,q)
 =\int_0^1\partial_s\partial_qF_n(s,q)\,ds.
 \label{eq:newS7V97-FTC}
\end{equation}

\begin{theorem}[Local sufficient row for adjacent stress projectivity]
\label{thm:newS7V97-adjacent}
Assume the transported interpolation satisfies, uniformly in
\(s\in[0,1]\), the hypotheses of
\Cref{thm:newS7V97-local,prop:newS7V97-recovery}.  Let
\(m_n(s,q;v)\) denote the right side of
\eqref{eq:newS7V97-observable-mixed}, with the three response vectors
replaced by the explicit bounds
\eqref{eq:newS7V97-first-local}--\eqref{eq:newS7V97-mixed-local}.  If
\begin{equation}
 \sum_{n\geq n_0}
 \sup_{q\in Q,\ \|v\|\leq1}
 \int_0^1m_n(s,q;v)\,ds<\infty,
 \label{eq:newS7V97-adjacent-sum}
\end{equation}
then the V95 adjacent carrier-derivative row is summable, without using the
global density-floor constant \(L_D(d_n,\lambda_{d_n,M_n})\).
\end{theorem}

\begin{proof}
Apply \Cref{cor:newS7V97-observable} pointwise in \(s,q,v\), integrate
with \eqref{eq:newS7V97-FTC}, and take the stated suprema.  The Weierstrass
test and \eqref{eq:newS7V97-adjacent-sum} give the V95 row.
\end{proof}

\begin{firewall}[What has moved from global to local]
\label{fire:newS7V97-local}
The global density floor has disappeared from the conclusion, but only
because it has been replaced by concrete local obligations:
\(c_\mu<1\), finite weighted derivative kernels \(A^s,A^q\), localized
first and mixed channel sources, differentiated recovery defects, and the
summation \eqref{eq:newS7V97-adjacent-sum}.  None of those numerical rows
is supplied by abstract complete positivity.
\end{firewall}

\begin{programme}[V98: cofinal mixed-source compiler]
\label{prog:newS7V97-V98}
V98 will turn the scale-wise theorem into a reusable cofinal criterion.  It
will identify a scalar summability ledger for channel contraction margins,
refinement shells, metric sources, recovery defects, and observable contact
terms, and prove that this ledger implies both the V95 base and derivative
rows on a countable local core.
\end{programme}

\begin{firewall}[Claim boundary after seventh-series V97]
\label{fire:newS7V97-boundary}
V97 proves the exact two-parameter fixed-point derivative equation, its
mixed-source local resolvent, a volume-independent exponentially weighted
bound, the observable mixed-response formula, the differentiated recovery
defect extension, and a sufficient local adjacent-stress row.

V97 does not prove the physical channel contraction margin or derivative
kernels, load the actual old/detail and recovery sources, establish their
cofinal summability, remove the cutoff, or construct the Standard
Model--Einstein continuum.
\end{firewall}

\begin{theorem}[Seventh-series V97 result]
\label{thm:newS7V97-result}
For a two-parameter locally contractive channel family, the mixed
metric/refinement response obeys
\eqref{eq:newS7V97-mixed-resolvent}; under the weighted contraction
hypotheses it obeys the volume-independent estimate
\eqref{eq:newS7V97-mixed-local}.  Integrated along the refinement
interpolation, this mixed response bounds the adjacent metric derivative
and supplies the V95 stress-projectivity row whenever
\eqref{eq:newS7V97-adjacent-sum} is finite.
\end{theorem}

\begin{proof}
Combine \Cref{thm:newS7V97-mixed,thm:newS7V97-local,
cor:newS7V97-observable,prop:newS7V97-recovery,
thm:newS7V97-adjacent}.
\end{proof}

\paragraph{Parent-version record.}
V97 was formed from
\texttt{Renewal\_SM\_Einstein\_Continuum\_Research\_Notes\_V96.tex}.
The V96 parent had \(38\,802\) logical source lines, \(1\,516\,545\) bytes,
and SHA--256
\[
 \texttt{1174D33054C9B48FF16BFCC2C6C672945450C15234174A85791CF94CDE83D447}.
\]
The next compulsory companion audit is V100.

% =============================================================================
% END APPENDED SEVENTH-SERIES V97 MATERIAL
% =============================================================================
% =============================================================================
% BEGIN APPENDED SEVENTH-SERIES V98 MATERIAL
% =============================================================================

\section{A cofinal mixed-source compiler on a countable local core}
\label{sec:v7v98}

V97 turns a two-parameter channel interpolation into local first and mixed
response bounds.  The remaining bookkeeping problem is cofinal: the same
bounds must be summable for every protected observable and every declared
metric tangent.  This appendix packages the required quantities into one
nonnegative scalar ledger and proves what its summability actually yields.

\subsection{Scale-wise scalarization}
\label{sec:newS7V98-scalar}

Let \(\mathfrak A_{\rm loc}^{0}\) be a countable unital local
\(*\)-algebra with enumeration \((A^k)_{k\geq1}\).  At adjacent scale
\(n\), suppose the V97 interpolation
\(s\mapsto(\rho_{n,s,q},\mathcal T_{n,s,q})\), \(0\leq s\leq1\), joins
the transported scale-\(n\) carrier to the scale-\(n+1\) carrier on a
common collar algebra.  Write \(B_k\) for the cell support of \(A^k\).
All bounds below are understood after the declared adjacent embeddings.

Fix an exponential weight \(\mu>0\).  For each \((n,s,q)\), define
\begin{align}
 R_n&:=(1-c_{\mu,n})^{-1},
 \label{eq:newS7V98-R}\\
 U_n^s&:=R_n b_n^s,qquad
 U_n^q:=R_n b_n^q,
 \label{eq:newS7V98-first-U}\\
 U_n^{sq}&:=R_n\bigl(
 b_n^{sq}+a_n^sR_nb_n^q+a_n^qR_nb_n^s\bigr),
 \label{eq:newS7V98-mixed-U}
\end{align}
where \(c_{\mu,n}\), \(a_n^s\), and \(a_n^q\) are the weighted kernel
norms of \(C_n,A_n^s,A_n^q\), while
\(b_n^s,b_n^q,b_n^{sq}\) are the corresponding weighted norms of the
modified sources from \Cref{prop:newS7V97-recovery}.  Thus V97 gives
\begin{equation}
 \|u_n^s\|_{\mu,Z_n}\leq U_n^s,qquad
 \|u_n^q\|_{\mu,Z_n}\leq U_n^q,qquad
 \|u_n^{sq}\|_{\mu,Z_n}\leq U_n^{sq}.
 \label{eq:newS7V98-U-meaning}
\end{equation}
The quantities may depend on \((s,q)\), and on a metric tangent in the
entries carrying a \(q\) superscript.

Set
\begin{equation}
 \chi_{n,k}:=|B_k|e^{-\mu d(B_k,Z_n)}.
 \label{eq:newS7V98-chi}
\end{equation}
For the transported observable family \(A^k_{n,s,q}\), use the pointwise
abbreviations
\begin{align}
 o^k_n&:=\|A^k_{n,s,q}\|,&
 o^{k,s}_n&:=\|\partial_sA^k_{n,s,q}\|,
 \label{eq:newS7V98-o-first}\\
 o^{k,q}_n&:=\|\partial_qA^k_{n,s,q}\|,&
 o^{k,sq}_n&:=\|\partial_s\partial_qA^k_{n,s,q}\|.
 \label{eq:newS7V98-o-mixed}
\end{align}
Define the nonnegative base and derivative loaders
\begin{align}
 B_{n,k}(q)
 &:={\int_0^1}\left[
 \chi_{n,k}o^k_nU_n^s+o^{k,s}_n\right]ds,
 \label{eq:newS7V98-B}\\
 G_{n,k}(q;v)
 &:={\int_0^1}\left[
 \chi_{n,k}\bigl(
 o^k_nU_n^{sq}+o^{k,s}_nU_n^q+o^{k,q}_nU_n^s\bigr)
 +o^{k,sq}_n\right]ds.
 \label{eq:newS7V98-G}
\end{align}
Here the \(q\)-derivatives in \eqref{eq:newS7V98-G} are directional
derivatives along \(v\).

\begin{lemma}[Adjacent expectation and stress bounds]
\label{lem:newS7V98-adjacent}
Let \(F_n^k(q)\) be the scale-\(n\) expectation of \(A^k\).  Under the
V97 hypotheses,
\begin{align}
 |F_{n+1}^k(q)-F_n^k(q)|&\leq B_{n,k}(q),
 \label{eq:newS7V98-base-bound}\\
 |DF_{n+1}^k(q)[v]-DF_n^k(q)[v]|&\leq G_{n,k}(q;v).
 \label{eq:newS7V98-derivative-bound}
\end{align}
\end{lemma}

\begin{proof}
Once differentiating the interpolated expectation gives
\[
 \partial_s\Tr(\rho A^k)
 =\Tr(\rho_sA^k)+\Tr(\rho A_s^k).
\]
The V93 dual seminorm estimate, \eqref{eq:newS7V98-U-meaning}, and the
support factor \eqref{eq:newS7V98-chi} bound its absolute value by the
integrand in \eqref{eq:newS7V98-B}.  Integrate from \(s=0\) to \(s=1\).
The twice differentiated identity is
\Cref{cor:newS7V97-observable}; applying the same bounds gives the
integrand in \eqref{eq:newS7V98-G}, and V97's fundamental-theorem identity
\eqref{eq:newS7V97-FTC} proves \eqref{eq:newS7V98-derivative-bound}.
\end{proof}

\subsection{One ledger for countably many rows}
\label{sec:newS7V98-ledger}

Let \((v^\ell)_{\ell\geq1}\) be a countable collection of metric tangents.
For a bounded convex chart \(Q_m\), define
\begin{align}
 \overline B_{n,k}^{(m)}
 &:=\sup_{q\in Q_m}B_{n,k}(q),
 \label{eq:newS7V98-Bbar}\\
 \overline G_{n,k,\ell}^{(m)}
 &:=\sup_{q\in Q_m}G_{n,k}(q;v^\ell).
 \label{eq:newS7V98-Gbar}
\end{align}
Choose an increasing exhaustion \(Q_1\subset Q_2\subset\cdots\) of the
declared charts and put
\begin{equation}
 \Lambda_n:=
 \max\left\{
 \max_{\substack{1\leq m,k\leq n}}2^{-(m+k)}
       \overline B_{n,k}^{(m)},
 \quad
 \max_{\substack{1\leq m,k,\ell\leq n}}2^{-(m+k+\ell)}
       \overline G_{n,k,\ell}^{(m)}
 \right\}.
 \label{eq:newS7V98-Lambda}
\end{equation}

\begin{theorem}[Master-ledger implication]
\label{thm:newS7V98-ledger}
If
\begin{equation}
 \sum_{n\geq n_0}\Lambda_n<\infty,
 \label{eq:newS7V98-Lambda-sum}
\end{equation}
then, for every fixed \((m,k,\ell)\),
\begin{equation}
 \sum_n\overline B_{n,k}^{(m)}<\infty,
 \qquad
 \sum_n\overline G_{n,k,\ell}^{(m)}<\infty.
 \label{eq:newS7V98-all-rows}
\end{equation}
Consequently \(F_n^k\) converges uniformly on \(Q_m\), and its directional
derivatives along \(v^\ell\) converge uniformly there.
\end{theorem}

\begin{proof}
For \(n\geq\max\{m,k\}\), definition
\eqref{eq:newS7V98-Lambda} gives
\[
 \overline B_{n,k}^{(m)}\leq2^{m+k}\Lambda_n.
\]
For \(n\geq\max\{m,k,\ell\}\), it likewise gives
\[
 \overline G_{n,k,\ell}^{(m)}
 \leq2^{m+k+\ell}\Lambda_n.
\]
The omitted initial portions are finite sums of finite numbers, so
\eqref{eq:newS7V98-Lambda-sum} proves \eqref{eq:newS7V98-all-rows}.
Now apply \Cref{lem:newS7V98-adjacent} and the Weierstrass test.
\end{proof}

\begin{remark}[Why the weights do not weaken a fixed row]
\label{rem:newS7V98-weights}
The factors \(2^{-(m+k)}\) and \(2^{-(m+k+\ell)}\) merely permit one
scalar ledger to enumerate countably many obligations.  Once a row is
fixed, its reciprocal weight is a fixed finite constant.  No uniform claim
over all observables or all tangents follows from
\eqref{eq:newS7V98-Lambda-sum}.
\end{remark}

\subsection{Landing as a state and completion of the tangent core}
\label{sec:newS7V98-landing}

Let \(\mathfrak A_{\rm ql}\) be the norm closure of
\(\mathfrak A_{\rm loc}^{0}\), assuming the adjacent embeddings make this
a coherent inductive local algebra.

\begin{theorem}[Algebraic state landing]
\label{thm:newS7V98-state}
Suppose \eqref{eq:newS7V98-Lambda-sum} holds and each \(F_n(\,cdot\,;q)\)
is a state on the represented local algebra.  Then, for every \(q\) in a
declared chart, the limits on \(\mathfrak A_{\rm loc}^{0}\) extend
uniquely to a state \(F(\,\cdot\,;q)\) on
\(\mathfrak A_{\rm ql}\).
\end{theorem}

\begin{proof}
Theorem \ref{thm:newS7V98-ledger} gives a limit for every enumerated local
observable.  Linearity follows by passing to the limit in the finite-scale
linear identities, using the coherent embeddings.  If \(A\geq0\) belongs
to the local algebra, then \(F_n(A;q)\geq0\), hence \(F(A;q)\geq0\).
Also \(F(\mathbf1;q)=1\), and
\(|F(A;q)|\leq\|A\|\).  Thus the local functional is bounded with norm
one and extends uniquely by continuity to the norm closure; positivity is
preserved under norm approximation.
\end{proof}

\begin{lemma}[Dense tangent-core completion]
\label{lem:newS7V98-tangent}
Let \(E\) be a separable Banach tangent space and let
\((v^\ell)\) be dense in \(E\).  Suppose \(L_n\in E^*\),
\begin{equation}
 \sup_n\|L_n\|_{E^*}\leq M<\infty,
 \label{eq:newS7V98-uniform-L}
\end{equation}
and \((L_n(v^\ell))_n\) is Cauchy for every \(\ell\).  Then \(L_n(v)\)
converges for every \(v\in E\), defining an element \(L\in E^*\) with
\(\|L\|\leq M\).
\end{lemma}

\begin{proof}
For \(v\in E\), choose \(v^\ell\) with \(\|v-v^\ell\|<\varepsilon\).
Then
\[
 |L_n(v)-L_r(v)|
 \leq2M\varepsilon+|L_n(v^\ell)-L_r(v^\ell)|.
\]
First take \(n,r\) large and then let \(\varepsilon\downarrow0\).  The
limit is linear and its norm is at most \(M\).
\end{proof}

\begin{corollary}[Stress functional on the completed tangent core]
\label{cor:newS7V98-stress}
Fix \(A^k\) and \(q\).  If the finite carrier derivatives
\(DF_n^k(q)\) are uniformly bounded in \(E^*\), then the derivative limits
supplied on \((v^\ell)\) by \Cref{thm:newS7V98-ledger} extend uniquely to a
bounded stress functional on all of \(E\).
\end{corollary}

\begin{proof}
Apply \Cref{lem:newS7V98-tangent} to
\(L_n=DF_n^k(q)\).
\end{proof}

\begin{firewall}[Topology of the conclusion]
\label{fire:newS7V98-topology}
The master ledger gives uniform convergence on each enumerated chart for
each enumerated observable and tangent.  The state extension is a
\(C^*\)-algebraic state, and the tangent extension is pointwise weak
convergence under the separate uniform dual-norm bound
\eqref{eq:newS7V98-uniform-L}.  These statements do not give operator-norm
convergence of stress tensors, a continuum spacetime measure, or a
Lorentzian field equation.
\end{firewall}

\begin{assessment}[What remains to load]
\label{ass:newS7V98-status}
The abstract cofinal logic is closed: one explicit scalar sequence
\(\Lambda_n\) dominates every fixed protected base and stress row.
The physical work is now concentrated in proving
\(c_{\mu,n}<1\), estimating the first and mixed source norms in
\eqref{eq:newS7V98-first-U}--\eqref{eq:newS7V98-mixed-U}, controlling
observable contact derivatives, and showing
\(\sum_n\Lambda_n<\infty\) for the actual SM--Einstein regulator.
\end{assessment}

\begin{programme}[V99: residual and counterterm coherence]
\label{prog:newS7V98-V99}
V99 will add the gravitational, classical, carrier, and counterterm rows to
the same countable-core compiler.  It will prove a cocycle criterion under
which adjacent residual defects telescope to a regulator-independent weak
Einstein residual, while keeping finite counterterm freedom explicit.
\end{programme}

\begin{firewall}[Claim boundary after seventh-series V98]
\label{fire:newS7V98-boundary}
V98 proves the scalarization of V97, the master-ledger implication for
countably many charts, observables, and tangents, algebraic state landing,
and completion of a dense tangent core under a uniform dual-norm bound.

V98 does not populate or sum the physical ledger, prove uniqueness beyond
the declared inductive algebra, construct a continuum metric, establish
Lorentzian dynamics, or construct the Standard Model--Einstein continuum.
\end{firewall}

\begin{theorem}[Seventh-series V98 result]
\label{thm:newS7V98-result}
The scalar quantities \eqref{eq:newS7V98-B} and
\eqref{eq:newS7V98-G} dominate, respectively, every adjacent carrier
expectation and metric-derivative defect.  Summability of the single master
ledger \eqref{eq:newS7V98-Lambda} therefore yields all fixed protected
rows, a state on the norm-closed inductive local algebra, and bounded
stress functionals on the completed tangent core whenever
\eqref{eq:newS7V98-uniform-L} also holds.
\end{theorem}

\begin{proof}
Combine \Cref{lem:newS7V98-adjacent,thm:newS7V98-ledger,
thm:newS7V98-state,lem:newS7V98-tangent,cor:newS7V98-stress}.
\end{proof}

\paragraph{Parent-version record.}
V98 was formed from
\texttt{Renewal\_SM\_Einstein\_Continuum\_Research\_Notes\_V97.tex}.
The V97 parent had \(39\,147\) logical source lines, \(1\,529\,158\) bytes,
and SHA--256
\[
 \texttt{8D2BA4D562E661E6BF85EF147602C66F08646F081151631500F7DDD2489C39A9}.
\]
The next compulsory companion audit is V100.

% =============================================================================
% END APPENDED SEVENTH-SERIES V98 MATERIAL
% =============================================================================

% =============================================================================
% BEGIN APPENDED SEVENTH-SERIES V99 MATERIAL
% =============================================================================

\section{Counterterm cocycles and regulator-independent weak residuals}
\label{sec:v7v99}

The V98 ledger lands the carrier state and its metric derivative on a
countable local core.  An Einstein statement also requires coherent limits
of the gravitational, classical matter, and counterterm variations.  This
appendix formulates that requirement in a Banach space of weak residuals.
It separates two questions which must not be conflated: convergence after
a counterterm lift, and independence of the regulator modulo the allowed
finite counterterms.

\subsection{Residual Banach space and finite counterterm sector}
\label{sec:newS7V99-space}

Fix a bounded metric chart \(Q_m\) with tangent Banach space \(E_m\).  Let
\begin{equation}
 \mathcal X_m:=C_b(Q_m;E_m^*)
 \label{eq:newS7V99-X}
\end{equation}
with the uniform dual norm.  If only a protected observable sector is being
considered, \(\mathcal X_m\) is replaced by the corresponding finite or
countable product, equipped with a declared complete weighted norm.  Let
\(\mathcal C_m\subset\mathcal X_m\) be the finite-dimensional space of
allowed finite local counterterm variations.  It is closed.  Denote the
quotient map by
\begin{equation}
 \pi_m:\mathcal X_m\longrightarrow
 \mathcal X_m/\mathcal C_m.
 \label{eq:newS7V99-pi}
\end{equation}

At regulator \(r\), define the raw and renormalized residuals
\begin{align}
 U_r&:=2D\mathcal S_{g,r}
 -\mathcal T_{{\rm cl},r}-\mathcal T_{{\rm car},r},
 \label{eq:newS7V99-U}\\
 R_r&:=U_r-\mathcal T_{{\rm ct},r},
 \qquad \mathcal T_{{\rm ct},r}\in\mathcal C_m.
 \label{eq:newS7V99-R}
\end{align}
Every equality in this section is an equality in \(\mathcal X_m\), hence
is uniform on \(Q_m\) and its unit tangent set.

\begin{lemma}[Absolute counterterm lift]
\label{lem:newS7V99-lift}
Let \((U_n)\subset\mathcal X_m\).  Suppose its adjacent differences admit
\begin{equation}
 U_{n+1}-U_n=c_n+\eta_n,
 \qquad c_n\in\mathcal C_m,
 \qquad
 \sum_n\|\eta_n\|_{\mathcal X_m}<\infty.
 \label{eq:newS7V99-split}
\end{equation}
Choose any \(P_{n_0}\in\mathcal C_m\) and define
\begin{equation}
 P_{n+1}:=P_n+c_n.
 \label{eq:newS7V99-P-recursion}
\end{equation}
Then \(P_n\in\mathcal C_m\), and the lifted residuals
\(R_n:=U_n-P_n\) converge in \(\mathcal X_m\).  More precisely,
\begin{equation}
 R_{n+1}-R_n=\eta_n,
 \qquad
 \|R_\infty-R_N\|\leq\sum_{n\geq N}\|\eta_n\|.
 \label{eq:newS7V99-tail}
\end{equation}
\end{lemma}

\begin{proof}
Finite partial sums of the \(c_n\) remain in the vector space
\(\mathcal C_m\), so \eqref{eq:newS7V99-P-recursion} is admissible.
Subtract it from \eqref{eq:newS7V99-split}.  The resulting increment is
\(\eta_n\); absolute summability makes \((R_n)\) Cauchy in the Banach
space \(\mathcal X_m\) and gives the tail estimate.
\end{proof}

\begin{proposition}[Finite-scheme ambiguity]
\label{prop:newS7V99-scheme}
Suppose two counterterm choices \(P_n,P_n'\in\mathcal C_m\) produce limits
\(R_\infty=\lim(U_n-P_n)\) and
\(R_\infty'=\lim(U_n-P_n')\) in \(\mathcal X_m\).  Then
\begin{equation}
 R_\infty'-R_\infty\in\mathcal C_m,
 \qquad
 \pi_m(R_\infty')=\pi_m(R_\infty).
 \label{eq:newS7V99-scheme}
\end{equation}
\end{proposition}

\begin{proof}
For every \(n\),
\((U_n-P_n')-(U_n-P_n)=P_n-P_n'\in\mathcal C_m\).  The left side
converges to \(R_\infty'-R_\infty\).  Closedness of the finite-dimensional
subspace \(\mathcal C_m\) places the limit in that subspace, and the
quotient identity follows.
\end{proof}

\subsection{The exact quotient cocycle}
\label{sec:newS7V99-cocycle}

Let \(\mathfrak R\) be a directed family of regulators whose residuals are
transported into the same \(\mathcal X_m\).  For \(r,t\in\mathfrak R\)
define
\begin{equation}
 \delta_m(r,t):=\pi_m(U_t-U_r).
 \label{eq:newS7V99-delta}
\end{equation}

\begin{lemma}[Residual cocycle]
\label{lem:newS7V99-cocycle}
For all \(r,s,t\in\mathfrak R\),
\begin{equation}
 \delta_m(r,t)=\delta_m(r,s)+\delta_m(s,t),
 \qquad
 \delta_m(r,r)=0.
 \label{eq:newS7V99-cocycle}
\end{equation}
The cocycle is unaffected by any regulator-dependent choice of
\(\mathcal C_m\)-valued counterterm.
\end{lemma}

\begin{proof}
Both identities are the telescoping identity for the differences of the
single-valued map \(r\mapsto\pi_m(U_r)\).  Since
\(\pi_m(P_r)=0\) for \(P_r\in\mathcal C_m\), replacing \(U_r\) by
\(U_r-P_r\) leaves \(\delta_m\) unchanged.
\end{proof}

\begin{theorem}[Cofinal regulator independence in the quotient]
\label{thm:newS7V99-cofinal}
Suppose there is a cofinal sequence \((r_N)\) in \(\mathfrak R\) such that
\begin{equation}
 \epsilon_N^{(m)}:=
 \sup_{r,t\succeq r_N}\|\delta_m(r,t)\|_{
       \mathcal X_m/\mathcal C_m}
 \longrightarrow0.
 \label{eq:newS7V99-tail-diameter}
\end{equation}
Then \(\pi_m(U_r)\) converges as a net.  Every cofinal regulator sequence
has the same limit in \(\mathcal X_m/\mathcal C_m\).
\end{theorem}

\begin{proof}
For \(r,t\succeq r_N\), the quotient distance between
\(\pi_m(U_r)\) and \(\pi_m(U_t)\) is exactly
\(\|\delta_m(r,t)\|\), which is at most
\(\epsilon_N^{(m)}\).  Thus the net is Cauchy.  The quotient of a Banach
space by a closed subspace is Banach, so the net converges.  A cofinal
subnet of a convergent net has the same limit; every cofinal sequence is
such a subnet after the usual directed reindexing.
\end{proof}

\begin{corollary}[Common-refinement bridge]
\label{cor:newS7V99-bridge}
Let \((r_n)\) and \((s_n)\) be two regulator sequences for which
\(\pi_m(U_{r_n})\to u\) and \(\pi_m(U_{s_n})\to v\).  If there are common
refinements \(t_n\succeq r_n,s_n\) satisfying
\begin{equation}
 \|\delta_m(r_n,t_n)\|+
 \|\delta_m(s_n,t_n)\|\longrightarrow0,
 \label{eq:newS7V99-bridge}
\end{equation}
then \(u=v\).
\end{corollary}

\begin{proof}
The cocycle and the triangle inequality give
\[
 \|\pi_m(U_{r_n})-\pi_m(U_{s_n})\|
 \leq\|\delta_m(r_n,t_n)\|+
       \|\delta_m(s_n,t_n)\|.
\]
Pass to the limit.
\end{proof}

\begin{firewall}[One chain is not regulator independence]
\label{fire:newS7V99-chain}
Absolute summability of adjacent defects along one chosen sequence proves
convergence along that sequence.  It does not imply
\eqref{eq:newS7V99-tail-diameter} for a directed regulator family and does
not compare a second discretization.  A common-refinement bridge or an
equivalent tail-diameter estimate is a separate required row.
\end{firewall}

\subsection{A combined residual ledger}
\label{sec:newS7V99-ledger}

At adjacent scale \(n\), suppose nonnegative numbers
\(g_n^{(m)},c_n^{(m)},k_n^{(m)},h_n^{(m)}\) satisfy
\begin{align}
 2\|D\mathcal S_{g,n+1}-D\mathcal S_{g,n}\|_{\mathcal X_m}
 &\leq g_n^{(m)},
 \label{eq:newS7V99-g}\\
 \|\mathcal T_{{\rm cl},n+1}-\mathcal T_{{\rm cl},n}\|_{\mathcal X_m}
 &\leq c_n^{(m)},
 \label{eq:newS7V99-c}\\
 \|\mathcal T_{{\rm car},n+1}-\mathcal T_{{\rm car},n}\|_{\mathcal X_m}
 &\leq k_n^{(m)}.
 \label{eq:newS7V99-k}
\end{align}
The carrier number \(k_n^{(m)}\) may be supplied by the V98 mixed-source
ledger only after the dense tangent-core limit has the uniform dual-norm
bound required there.  Assume further that the raw adjacent residual can
be decomposed as
\begin{equation}
 U_{n+1}-U_n=c_n^{\rm ct}+\eta_n,qquad
 c_n^{\rm ct}\in\mathcal C_m,qquad
 \|\eta_n\|_{\mathcal X_m}
 \leq g_n^{(m)}+c_n^{(m)}+k_n^{(m)}+h_n^{(m)},
 \label{eq:newS7V99-residual-loader}
\end{equation}
where \(h_n^{(m)}\) contains embedding, contact, and tensor-split
remainders not already present in the three displayed rows.  Define
\begin{equation}
 \Omega_n:=\max_{1\leq m\leq n}2^{-m}
 \bigl(g_n^{(m)}+c_n^{(m)}+k_n^{(m)}+h_n^{(m)}\bigr).
 \label{eq:newS7V99-Omega}
\end{equation}

\begin{theorem}[Residual master-ledger landing]
\label{thm:newS7V99-ledger}
If
\begin{equation}
 \sum_n\Omega_n<\infty,
 \label{eq:newS7V99-Omega-sum}
\end{equation}
then, on every fixed chart \(Q_m\), the recursive counterterm lift
\begin{equation}
 \mathcal T_{{\rm ct},n+1}-\mathcal T_{{\rm ct},n}
 =c_n^{\rm ct}
 \label{eq:newS7V99-ct-recursion}
\end{equation}
produces a limit \(R_\infty^{(m)}\in\mathcal X_m\).  Its tail is bounded
by
\begin{equation}
 \|R_\infty^{(m)}-R_N\|_{\mathcal X_m}
 \leq2^m\sum_{n\geq N}\Omega_n
 \qquad(N\geq m).
 \label{eq:newS7V99-residual-tail}
\end{equation}
\end{theorem}

\begin{proof}
For \(n\geq m\), definition \eqref{eq:newS7V99-Omega} and
\eqref{eq:newS7V99-residual-loader} imply
\(\|\eta_n\|\leq2^m\Omega_n\).  The finite initial segment is harmless.
Apply \Cref{lem:newS7V99-lift}; its tail estimate gives
\eqref{eq:newS7V99-residual-tail}.
\end{proof}

\begin{theorem}[When the weak Einstein identity can be normalized]
\label{thm:newS7V99-Einstein}
Assume \Cref{thm:newS7V99-ledger} and suppose
\begin{equation}
 \pi_m(R_\infty^{(m)})=0.
 \label{eq:newS7V99-zero-class}
\end{equation}
Then \(R_\infty^{(m)}\in\mathcal C_m\).  The finite counterterm change
\begin{equation}
 \mathcal T_{{\rm ct},n}'
 :=\mathcal T_{{\rm ct},n}+R_\infty^{(m)}
 \label{eq:newS7V99-finite-shift}
\end{equation}
produces \(R_n'\to0\) in \(\mathcal X_m\), hence the limiting weak
identity
\begin{equation}
 2D\mathcal S_g
 =\mathcal T_{\rm cl}+\mathcal T_{\rm car}
  +\mathcal T_{\rm ct}'
 \label{eq:newS7V99-Einstein}
\end{equation}
uniformly on the declared chart and unit tangent set.
\end{theorem}

\begin{proof}
The kernel of \(\pi_m\) is \(\mathcal C_m\), so
\eqref{eq:newS7V99-zero-class} places \(R_\infty^{(m)}\) in the allowed
finite counterterm sector.  Subtracting the shift
\eqref{eq:newS7V99-finite-shift} from the residual gives
\(R_n'=R_n-R_\infty^{(m)}\to0\).  Expanding the definition of the
residual yields \eqref{eq:newS7V99-Einstein}.
\end{proof}

\begin{firewall}[A vanishing class is an additional physical equation]
\label{fire:newS7V99-zero}
The summable ledger proves existence of a residual limit; it does not prove
\eqref{eq:newS7V99-zero-class}.  The latter asserts that every surviving
component is an allowed finite local counterterm.  Components outside
\(\mathcal C_m\), including an anomalous or nonlocal remainder, cannot be
removed by the finite shift \eqref{eq:newS7V99-finite-shift}.
\end{firewall}

\begin{assessment}[Status at the V100 threshold]
\label{ass:newS7V99-status}
Carrier projectivity and stress projectivity now feed a complete abstract
residual compiler.  It yields a convergent weak residual under the scalar
ledger \(\Omega_n\), identifies regulator independence with a quotient
cocycle tail condition, and states the exact extra condition needed to
normalize the limiting residual to zero.  None of the physical source rows
or bridge estimates has yet been numerically populated.
\end{assessment}

\begin{programme}[V100: companion audit and seventh-series ledger]
\label{prog:newS7V99-V100}
V100 will perform the compulsory YM and NS audit, compare any new companion
loaders with \(\Lambda_n\), \(\Omega_n\), and the quotient bridge, and
close the seventh series with a precise theorem/debit ledger.  The file will
then be frozen as the next \texttt{\_f7} archive, and a compact new V1 will
restart with context, the proved-results summary, and the next upstream
attack.
\end{programme}

\begin{firewall}[Claim boundary after seventh-series V99]
\label{fire:newS7V99-boundary}
V99 proves absolute counterterm lifting, finite-scheme ambiguity, the exact
quotient residual cocycle, a cofinal regulator-independence criterion, a
combined residual master ledger, and the precise finite-renormalization
criterion for a vanishing weak Einstein residual.

V99 does not populate the gravitational, classical, carrier, counterterm,
bridge, or zero-class rows for the physical regulator; prove regulator
independence; remove the cutoff; or construct the Standard Model--Einstein
continuum.
\end{firewall}

\begin{theorem}[Seventh-series V99 result]
\label{thm:newS7V99-result}
An adjacent raw residual split into an allowed local-counterterm increment
and an absolutely summable remainder admits a convergent renormalized lift.
Its quotient class is independent of finite counterterm scheme.  A
regulator-independent quotient limit follows from
\eqref{eq:newS7V99-tail-diameter}; if that limit has zero class, one allowed
finite counterterm shift makes the limiting weak Einstein residual vanish.
\end{theorem}

\begin{proof}
Combine \Cref{lem:newS7V99-lift,prop:newS7V99-scheme,
lem:newS7V99-cocycle,thm:newS7V99-cofinal,
thm:newS7V99-ledger,thm:newS7V99-Einstein}.
\end{proof}

\paragraph{Parent-version record.}
V99 was formed from
\texttt{Renewal\_SM\_Einstein\_Continuum\_Research\_Notes\_V98.tex}.
The V98 parent had \(39\,458\) logical source lines, \(1\,540\,698\) bytes,
and SHA--256
\[
 \texttt{06B94A29F47BC30B704DA79716A33C31C77F39B4BCE945E6F923904CEB6E1CF4}.
\]
The next compulsory companion audit is V100.

% =============================================================================
% END APPENDED SEVENTH-SERIES V99 MATERIAL
% =============================================================================
% =============================================================================
% BEGIN APPENDED SEVENTH-SERIES V100 MATERIAL
% =============================================================================

\section{V100 companion audit and the physical-clock critical row}
\label{sec:v7v100}

This is the compulsory V100 audit and the closing note of the seventh
series.  The companion audit changes the analytic direction.  The newest
Yang--Mills note proves that a Ward-critical quadratic owner cannot have a
uniform strict absolute one-site influence margin.  Therefore the fixed
margin used in the V97 local resolvent can only be a finite-regulator or
block hypothesis.  The cofinal branch must allow the margin to close at the
physical-clock rate and must scale each channel source by the same clock.

\subsection{Compulsory companion audit}
\label{sec:newS7V100-audit}

The newest coherent active companion sources in the shared folder were
read without modification:
\begin{align}
 &\texttt{Renewal\_Yang\_Mills\_Mass\_Gap\_Research\_Notes\_V49.tex},
 \notag\\
 &\qquad 22\,057\ \text{logical lines},\quad769\,656\ \text{bytes},
 \notag\\
 &\qquad
 \texttt{FBB6EF594D5AC09F3D02906C53288E09589001489650CD6834E6435C9B0C8C6A},
 \label{eq:newS7V100-YM-hash}\\
 &\texttt{Renewal\_NS\_Stability\_Research\_Notes\_V26.tex},
 \notag\\
 &\qquad 5\,319\ \text{logical lines},\quad278\,283\ \text{bytes},
 \notag\\
 &\qquad
 \texttt{82C887F8C2C69E4F28FAD7C3DC8DE5B9099CB5A4BF431EBA1FFF3E91935978AD}.
 \label{eq:newS7V100-NS-hash}
\end{align}
Each has one live document terminator and ends with
\verb|\end{document}|.  Older active-looking and frozen companion files
were left untouched.

YM V49 proves three facts relevant here.  First, one scalar plaquette
coefficient cannot reproduce the momentum-dependent flat Schur Hessian.
Second, the complete flat Schur owner is exponentially quasilocal.  Third,
its Ward zero forces the absolute Gaussian one-site influence row to be at
least one.  Hence a family converging to that owner cannot simultaneously
have convergent one-site linear response and a uniform Dobrushin margin
strictly below one.  Its proposed alternatives are a scale-dependent
near-critical row, a physical-duration block, or a cancellation norm not
dominated by absolute coordinate influences.

NS V26 computes rank-one norms for the first two derivatives of the Oseen
kernel.  Its first-derivative refinement gives no improvement over the
symmetric-traceless norm, while its second-derivative ratio is smaller on a
bounded radial range.  Those constants do not populate any SM carrier,
anomaly, counterterm, or Einstein-residual row.

\begin{audit}[V100 companion decision]
\label{aud:newS7V100-decision}
The YM Ward obstruction is imported as a negative compatibility test: the
V97--V98 fixed weighted contraction margin must not be asserted uniformly
on a cofinal one-site gauge carrier approaching the Ward-critical owner.
The programme switches to the physical-clock scaling proved below.  No NS
constant is imported.
\end{audit}

\begin{proof}
The YM implication is exactly its zero-row-sum theorem and
fixed-contraction versus Ward-critical alternative.  The NS quantities are
bounds for an Oseen-kernel stability ledger and have no identified map to
the SM local channel sources.  Importing them would therefore be an
unsupported identification.
\end{proof}

\subsection{Nonautonomous decay at physical time}
\label{sec:newS7V100-clock}

Use the V97 weighted vector norm and induced matrix norm.  Let
\(\tau_j>0\) be microscopic clock steps and set
\begin{equation}
 t_0=0,\qquad t_N:=\sum_{j=0}^{N-1}\tau_j.
 \label{eq:newS7V100-time}
\end{equation}

\begin{lemma}[Critical one-step rows produce physical-time decay]
\label{lem:newS7V100-product}
Suppose nonnegative propagation matrices \(C_j\) obey
\begin{equation}
 \|C_j\|_\mu\leq1-m\tau_j+K\tau_j^2,
 \qquad
 0<\tau_j\leq\tau_*,
 \qquad
 \eta:=m-K\tau_*>0.
 \label{eq:newS7V100-critical-row}
\end{equation}
Then
\begin{equation}
 \|C_{N-1}\cdots C_k\|_\mu
 \leq e^{-\eta(t_N-t_k)}.
 \label{eq:newS7V100-product}
\end{equation}
If a nonnegative recursion satisfies
\begin{equation}
 u_{j+1}\leq C_ju_j+\tau_jb_j,
 \qquad \|b_j\|_{\mu,Z}\leq B,
 \label{eq:newS7V100-recursion}
\end{equation}
then, with \(\tau_{\max,N}:=\max_{j<N}\tau_j\),
\begin{equation}
 \|u_N\|_{\mu,Z}
 \leq e^{-\eta t_N}\|u_0\|_{\mu,Z}
 +{B e^{\eta\tau_{\max,N}}\over\eta}.
 \label{eq:newS7V100-forced}
\end{equation}
\end{lemma}

\begin{proof}
The first inequality in \eqref{eq:newS7V100-critical-row} is at most
\(1-\eta\tau_j\), hence at most \(e^{-\eta\tau_j}\).  Submultiplicativity
gives \eqref{eq:newS7V100-product}.  Iterating
\eqref{eq:newS7V100-recursion} and applying that product bound gives
\begin{align*}
 \|u_N\|_{\mu,Z}
 &\leq e^{-\eta t_N}\|u_0\|_{\mu,Z}
 +B\sum_{k=0}^{N-1}\tau_k
 e^{-\eta(t_N-t_{k+1})}.
\end{align*}
For \(s\in[t_k,t_{k+1}]\),
\[
 e^{-\eta(t_N-t_{k+1})}
 \leq e^{\eta\tau_{\max,N}}e^{-\eta(t_N-s)}.
\]
Integrate on each clock interval and sum.  The resulting integral is at
most \(1/\eta\), proving \eqref{eq:newS7V100-forced}.
\end{proof}

\begin{remark}[What closes and what remains positive]
\label{rem:newS7V100-critical}
The microscopic margin in \eqref{eq:newS7V100-critical-row} tends to zero
when \(\tau_j\to0\); the cumulative decay remains
\(e^{-\eta t}\) in physical time.  This is compatible with the YM V49
obstruction because no uniform one-step margin is claimed.
\end{remark}

\subsection{Cancellation of the critical resolvent by source scaling}
\label{sec:newS7V100-resolvent}

At a fixed regulator scale \(n\), let \(\tau_n\) be the physical duration
of one channel sweep.  Suppose the V97 weighted contraction coefficient
satisfies
\begin{equation}
 c_{\mu,n}\leq1-m\tau_n+K\tau_n^2,
 \qquad 0<\tau_n\leq\tau_*,
 \qquad \eta=m-K\tau_*>0.
 \label{eq:newS7V100-scale-row}
\end{equation}
Then its resolvent factor obeys
\begin{equation}
 R_n=(1-c_{\mu,n})^{-1}\leq{1\over\eta\tau_n}.
 \label{eq:newS7V100-R}
\end{equation}
Let \(\varepsilon_n\geq0\) measure the old/detail refinement amplitude.
Assume the modified V97 sources and channel-derivative kernels satisfy
\begin{align}
 b_n^s&\leq\tau_n\varepsilon_n\beta_s,&
 b_n^q&\leq\tau_n\beta_q,&
 b_n^{sq}&\leq\tau_n\varepsilon_n\beta_{sq},
 \label{eq:newS7V100-source-scale}\\
 a_n^s&\leq\tau_n\varepsilon_n\alpha_s,&
 a_n^q&\leq\tau_n\alpha_q,
 \label{eq:newS7V100-kernel-scale}
\end{align}
with constants independent of \(n\).  The recovery-defect derivatives are
included in the three \(b\)-terms, as required by V97.

\begin{theorem}[Physical-clock mixed-source cancellation]
\label{thm:newS7V100-cancellation}
Under \eqref{eq:newS7V100-scale-row}--
\eqref{eq:newS7V100-kernel-scale}, the V98 response scalars obey
\begin{align}
 U_n^s&\leq{\beta_s\over\eta}\varepsilon_n,
 &
 U_n^q&\leq{\beta_q\over\eta},
 \label{eq:newS7V100-first-U}\\
 U_n^{sq}&\leq{\varepsilon_n\over\eta}
 \left(
 \beta_{sq}+{\alpha_s\beta_q+\alpha_q\beta_s\over\eta}
 \right).
 \label{eq:newS7V100-mixed-U}
\end{align}
In particular, the factor \((\eta\tau_n)^{-1}\) from the closing
contraction margin cancels against the clock factor in every first or mixed
source.  No inverse power of \(\tau_n\) remains.
\end{theorem}

\begin{proof}
Insert \eqref{eq:newS7V100-R} and
\eqref{eq:newS7V100-source-scale} into
\(U_n^s=R_nb_n^s\) and \(U_n^q=R_nb_n^q\).  For the mixed term use the
exact V98 formula
\[
 U_n^{sq}=R_n\bigl(
 b_n^{sq}+a_n^sR_nb_n^q+a_n^qR_nb_n^s\bigr).
\]
The three terms in parentheses are bounded, respectively, by
\[
 \tau_n\varepsilon_n\beta_{sq},\qquad
 \tau_n\varepsilon_n{\alpha_s\beta_q\over\eta},\qquad
 \tau_n\varepsilon_n{\alpha_q\beta_s\over\eta}.
\]
Multiplication by \((\eta\tau_n)^{-1}\) proves
\eqref{eq:newS7V100-mixed-U}.
\end{proof}

\begin{corollary}[Summable refinement amplitude loads V98]
\label{cor:newS7V100-V98}
Fix a chart, protected observable, and tangent.  Suppose
\(\chi_{n,k}\), \(o_n^k\), and \(o_n^{k,q}\) are uniformly bounded and
that the refinement contact derivatives satisfy
\begin{equation}
 o_n^{k,s}\leq\varepsilon_n\gamma_s^k,qquad
 o_n^{k,sq}\leq\varepsilon_n\gamma_{sq}^k.
 \label{eq:newS7V100-contact-scale}
\end{equation}
Then constants \(C_{B,k}\) and \(C_{G,k,v}\), independent of \(n\),
satisfy
\begin{equation}
 \overline B_{n,k}^{(m)}\leq C_{B,k}\varepsilon_n,
 \qquad
 \overline G_{n,k,v}^{(m)}\leq C_{G,k,v}\varepsilon_n.
 \label{eq:newS7V100-BG}
\end{equation}
Thus \(\sum_n\varepsilon_n<\infty\) gives every fixed V98 base and stress
row.  If the weighted maxima of these constants in
\eqref{eq:newS7V98-Lambda} are uniformly bounded, then
\(\Lambda_n\leq C\varepsilon_n\) and the single V98 master ledger is
summable.
\end{corollary}

\begin{proof}
In the base loader \eqref{eq:newS7V98-B}, both \(U_n^s\) and
\(o_n^{k,s}\) are \(O(\varepsilon_n)\) by
\eqref{eq:newS7V100-first-U} and
\eqref{eq:newS7V100-contact-scale}.  In the derivative loader
\eqref{eq:newS7V98-G}, the four products are, in order,
\(O(\varepsilon_n)\),
\(O(\varepsilon_n)O(1)\),
\(O(1)O(\varepsilon_n)\), and
\(O(\varepsilon_n)\), using
\eqref{eq:newS7V100-first-U}--\eqref{eq:newS7V100-mixed-U}.
Integrating over the unit refinement interval proves
\eqref{eq:newS7V100-BG}.  The two summability conclusions follow by
comparison with \(\sum_n\varepsilon_n\).
\end{proof}

\begin{consequence}[Corrected cofinal channel alternatives]
\label{cons:newS7V100-alternatives}
The local channel branch can feed the SM--Einstein continuum compiler by
any of the following separately proved mechanisms:
\begin{enumerate}
\item a fixed weighted contraction margin on a non-Ward or finite-volume
carrier;
\item the critical physical-clock packet
\eqref{eq:newS7V100-scale-row}--\eqref{eq:newS7V100-kernel-scale};
\item a block channel of nonvanishing physical duration; or
\item a cancellation seminorm for which a strict cofinal row is proved and
which still has the V93 dual-observable property.
\end{enumerate}
For the Ward-critical one-site absolute-influence branch, item 1 is ruled
out by the audited YM V49 alternative.
\end{consequence}

\begin{proof}
Items 1, 3, and 4 are applications of the V93--V98 local resolvent once
their respective contraction and duality hypotheses are proved.  Item 2 is
\Cref{thm:newS7V100-cancellation,cor:newS7V100-V98}.  The final exclusion
is \Cref{aud:newS7V100-decision}.
\end{proof}

\subsection{Seventh-series result and debit ledger}
\label{sec:newS7V100-ledger}

\begin{theorem}[Major results proved in V75--V100]
\label{thm:newS7V100-ledger}
The seventh-series block V75--V100 establishes the following finite and
conditional continuum statements.
\begin{enumerate}
\item It factors determinant crossings into transverse and Schur sectors,
tracks a simple transverse rank crossing and determinant-line phase, and
separates zero-mode saturation from the reduced determinant.
\item It verifies the complete one-generation local SM anomaly arithmetic
and the even number of left-handed \(SU(2)\) doublets.  This removes those
finite representation-theoretic obstructions; it does not construct the
regulated chiral measure.
\item It constructs the finite determinant-line Berry curvature, proves
local radial primitives and holonomy formulae, and derives exponential
projector and curvature locality under a uniform finite-range spectral gap.
\item It proves conditional common-regulator anomaly cancellation and
retains the exact finite Fock, quasifree, Nambu, and Gibbs comparison
identities needed by the carrier sector.
\item It derives exact Gibbs Duhamel metric response, a high-temperature
commutator locality branch, Petz/replacement recovery channels, local
channel fixed-point resolvents, and localized stress response under the
stated contraction, recovery, and differentiability packets.
\item It proves the base-plus-derivative \(C^1\) cofinal landing theorem,
the differentiated adjacent Gibbs bound, the mixed metric/refinement
channel equation, the countable-core master ledger, algebraic state
landing, and tangent-core completion.
\item It proves the counterterm quotient cocycle, distinguishes
single-chain convergence from regulator independence, and gives the exact
zero-class condition for a finite counterterm normalization of the weak
Einstein residual.
\item After the V100 YM audit, it replaces an inadmissible uniform
Ward-critical one-site margin by the physical-clock theorem
\eqref{eq:newS7V100-forced} and the mixed-source cancellation
\eqref{eq:newS7V100-mixed-U}.
\end{enumerate}
\end{theorem}

\begin{proof}
The items are, respectively, the result theorems of V75--V76, V77,
V78--V83, V84--V86, V87--V94, V95--V98, V99, and the present
\Cref{lem:newS7V100-product,thm:newS7V100-cancellation}.  Each item retains
the hypotheses and firewalls of its cited source result.
\end{proof}

\begin{firewall}[Open debit ledger at the end of the seventh series]
\label{fire:newS7V100-debits}
The programme has not yet constructed the full Standard Model--Einstein
continuum.  The principal unpaid rows are:
\begin{enumerate}
\item one common regulator with gauge-covariant chiral measure, zero-mode
patching, determinant-line trivialization, and reflection/positivity data;
\item physical old/detail channel and recovery maps satisfying the
critical source scalings
\eqref{eq:newS7V100-source-scale}--\eqref{eq:newS7V100-kernel-scale};
\item a proved summable refinement amplitude \(\varepsilon_n\), uniform
integrability, and completion from the protected algebra/tangent core;
\item gravitational and classical adjacent rows, all composite-operator
contact counterterms, and a compatible finite counterterm space;
\item common-refinement bridge estimates for regulator independence and
the zero quotient class \eqref{eq:newS7V99-zero-class}; and
\item construction of a continuum metric probability law, Lorentzian or
reconstruction theorem, physical clock, and nonperturbative continuum
field equations.
\end{enumerate}
No theorem in this series supplies these data by naming them.
\end{firewall}

\begin{assessment}[Direction after the restart]
\label{ass:newS7V100-next}
The first new-series target is the second debit above.  The programme should
construct a local channel as a short physical-time step
\(\mathcal T_{n}=I+\tau_n\mathcal L_n+O(\tau_n^2)\), derive the metric and
refinement derivatives of its generator, and prove that the recovery
residuals carry the same \(\tau_n\) factor.  If that construction fails in
the one-site norm, the next upstream branch is a fixed-physical-duration
block or a dual cancellation seminorm, not a return to a uniform
Dobrushin margin forbidden by the Ward limit.
\end{assessment}

\begin{firewall}[Claim boundary after seventh-series V100]
\label{fire:newS7V100-boundary}
V100 performs the compulsory companion audit, imports the YM Ward
obstruction as a negative compatibility test, proves nonautonomous
physical-time contraction, proves cancellation of a closing critical
resolvent by clock-scaled first and mixed sources, loads the V98 rows from a
summable refinement amplitude, and records the seventh-series theorem and
debit ledger.

V100 does not construct or populate the physical generator, prove its
source scalings, establish summability, regulator independence, the zero
residual class, cutoff removal, or the Standard Model--Einstein continuum.
\end{firewall}

\begin{theorem}[Seventh-series V100 result]
\label{thm:newS7V100-result}
A weighted channel row of the form
\(1-m\tau+O(\tau^2)\) yields exponential decay in accumulated physical
time.  At a fixed scale its resolvent grows like \(\tau^{-1}\), but this
growth cancels exactly when the metric, refinement, mixed, and recovery
sources are derivatives of a short-time channel and hence are
\(O(\tau)\).  If refinement-tagged sources carry an additional summable
amplitude \(\varepsilon_n\), the adjacent expectation and stress defects
are \(O(\varepsilon_n)\) and feed the V98 cofinal compiler without a
uniform microscopic contraction margin.
\end{theorem}

\begin{proof}
Combine \Cref{lem:newS7V100-product,
thm:newS7V100-cancellation,cor:newS7V100-V98,
cons:newS7V100-alternatives}.
\end{proof}

\paragraph{Parent-version record.}
V100 was formed from
\texttt{Renewal\_SM\_Einstein\_Continuum\_Research\_Notes\_V99.tex}.
The V99 parent had \(39\,816\) logical source lines, \(1\,553\,794\) bytes,
and SHA--256
\[
 \texttt{2EC68514815F62ECEA1EA478676E89870599C6D334DF1D628C31AB2395DAE014}.
\]
After validation, this large seventh-series file is to be frozen as
\texttt{Renewal\_SM\_Einstein\_Continuum\_Research\_Notes\_V100\_f7.tex}.
The next active file restarts at V1 with a compact context and result ledger.

% =============================================================================
% END APPENDED SEVENTH-SERIES V100 MATERIAL
% =============================================================================


\end{document}
